% Volumen documental compañero. Fuente independiente, sin compilación automática.
\documentclass[a4paper,10pt]{article}
\usepackage[top=23mm,bottom=23mm,left=20mm,right=20mm,headheight=15pt]{geometry}
\usepackage{fontspec,amsmath,unicode-math}
\usepackage[spanish,es-nodecimaldot,es-noquoting]{babel}
\usepackage{microtype,booktabs,longtable,array,needspace,fancyhdr}
\usepackage[unicode,hidelinks,bookmarksopen=true,bookmarksopenlevel=1,bookmarksnumbered=false]{hyperref}
\setmainfont{STIXTwoText-Regular.otf}[BoldFont=STIXTwoText-Bold.otf,ItalicFont=STIXTwoText-Italic.otf,BoldItalicFont=STIXTwoText-BoldItalic.otf]
\setmathfont{STIXTwoMath-Regular.otf}
\renewcommand{\normalsize}{\fontsize{10.5}{13.4}\selectfont}\normalsize
\setlength{\parindent}{1em}\setlength{\parskip}{3pt}
\setlength{\emergencystretch}{2em}
\clubpenalty=10000\widowpenalty=10000\raggedbottom
\clubpenalties 3 10000 10000 0
\widowpenalties 3 10000 10000 0
\AddToHook{cmd/section/before}{\clearpage}
\AddToHook{cmd/subsection/before}{\Needspace{7\baselineskip}}
\AddToHook{cmd/subsubsection/before}{\Needspace{6\baselineskip}}
\pagestyle{fancy}\fancyhf{}
\fancyhead[L]{Catálogo estructural y metrológico de partículas}
\fancyfoot[C]{\thepage}
\hypersetup{pdftitle={Catálogo estructural y metrológico de partículas},pdfauthor={Oumar Haidara Fall; Rubén Ramos Balsa},pdfsubject={Estados, rutas, masas y anchuras}}
\newcommand{\CatalogRouteReference}{Las secciones~\ref{sec:vi-324-rutas} y~\ref{sec:vi-inventario-metrologico} de este catálogo expresan, respectivamente, las 324 rutas orientadas y los 855 registros metrológicos.}
% Apertura independiente del índice: corrección quirúrgica 20260928.
\let\HMTIndiceOriginal\tableofcontents
\renewcommand{\tableofcontents}{\clearpage\begingroup\setlength{\parskip}{0pt}\fontsize{9.2}{10.5}\selectfont\HMTIndiceOriginal\endgroup\clearpage}

\usepackage[table]{xcolor}
\usepackage{tabularx}
\definecolor{routeblue}{HTML}{234B70}
\definecolor{routepale}{HTML}{F0F5F9}
\definecolor{routegold}{HTML}{B69B62}
\setsansfont{TeX Gyre Heros}[Scale=0.92]
\newcommand{\RouteGroup}[1]{\multicolumn{2}{@{}l@{}}{\rule{0pt}{17pt}\textcolor{routeblue}{\sffamily\bfseries #1}}\\}
\newcommand{\RouteRow}[2]{\textcolor{routeblue}{\sffamily\bfseries #1}& #2\\}
\newcommand{\RouteTitle}[2]{\par\noindent{\color{routeblue}\sffamily\bfseries\fontsize{15}{18}\selectfont #1}\hfill{\sffamily\fontsize{10}{12}\selectfont #2}\par\vspace{4pt}{\color{routegold}\hrule height .6pt}\vspace{6pt}}
\newcommand{\CatalogTableLeading}{11.3}
\AtBeginEnvironment{longtable}{\fontsize{9.4}{\CatalogTableLeading}\selectfont}

\begin{document}
\thispagestyle{plain}
\begin{center}
{\fontsize{10}{12}\selectfont Manuscrito de recopilación para transferencia académica\par}
\vspace{12mm}
{\fontsize{17}{20.5}\selectfont\bfseries Catálogo estructural y metrológico de partículas\par}
\vspace{6pt}
{\fontsize{11}{13.5}\selectfont Estados, rutas, masas y anchuras\par}
\vspace{11pt}
{\fontsize{11.5}{14}\selectfont Oumar Haidara Fall \textperiodcentered\ Rubén Ramos Balsa\par}
\vspace{4pt}
{\fontsize{8}{10}\selectfont Coautoría sin prelación de contribución\par}
\vspace{3pt}
{\fontsize{7.5}{9.5}\selectfont Integración y recopilación documental generadas por ChatGPT -- Astra.\par}
\end{center}
\vspace{7pt}
\begin{center}\bfseries Presentación\end{center}
\begingroup
\fontsize{10.5}{12.8}\selectfont
\setlength{\parskip}{2pt plus 1pt minus .5pt}
\noindent
\begin{minipage}{\textwidth}
Este catálogo reúne el inventario estructural de \(324\) rutas orientadas
y \(855\) registros metrológicos de masas y anchuras. Su organización
permite consultar la identidad y los campos de cada ruta y localizar las
observaciones metrológicas, sin equiparar el número de registros documentales con el
de predicciones individualizadas. Acompaña a \emph{Partículas,
persistencia y derivación del espectro de masas}, donde se desarrollan
las definiciones y demostraciones que construyen la partícula, el atlas
y el operador de masa. Aquí se expresa íntegro el dominio de datos
estructurales y de contraste utilizado por ese desarrollo.

La primera sección distingue \(13\) multisecciones, \(56\) colecciones de
ruta y \(23\) clases del inventario. La
sección~\ref{sec:vi-324-rutas} presenta las \(324\) rutas, conservando
todos los campos de sus dos registros concordantes. La
sección~\ref{sec:vi-inventario-metrologico} reúne \(471\) masas y
\(384\) anchuras de la edición metrológica fijada, con sus identificadores,
estados, unidades y fuentes. Valores centrales, intervalos, límites y
anchuras conservan su significado propio. La presencia de una
observación en el inventario permite localizarla para el contraste;
su identificación física requiere el lector y la coordenatización especificados
en el artículo.

Cada ruta se presenta en una ficha completa, con campos alineados y secuencias íntegras. El índice enlazado y los marcadores permiten localizar cada celda y orientación. El apéndice técnico conserva en este mismo volumen las huellas de procedencia, los ordinales de ambas fuentes y sus códigos originales. Las tablas metrológicas agrupan la identidad, las especies, el valor, la incertidumbre y el estado documental.\par
\end{minipage}
\par\endgroup
\clearpage
\tableofcontents
\clearpage
% Delta focal 2026-09-26. Los cuerpos 11/12 del catalogo permanecen inalterados.
\section{Longitud estructural, acci\'on y respuesta gravitatoria}
\label{sec:vi-catalogo-dependencia-l-g}

El cat\'alogo re\'une 324 rutas y 855 registros metrol\'ogicos. Su
interpretaci\'on sigue la construcci\'on del mismo estado HMT--MD, en el orden
causal
\[
 \mathrm{APP}\longrightarrow\mathrm{TRIT}\longrightarrow\mathrm{TPK}
 \longrightarrow X_{\mathrm{enr}}
 \longrightarrow\mathcal C^{\mathrm{disc}}_{\mathrm{joint}}
 \longrightarrow L_{\alpha}
 \longrightarrow(H\Lambda,R\Lambda)
 \longrightarrow G .
 \tag{A.1}\label{eq:vi-cat-cadena-lg}
\]
Las coordenadas \(\pi\), \(\varphi\), \(e\) y
\(\alpha_{\mathrm{HMT}}\) son salidas correlacionadas de la imagen non\'adica
tipada anterior; no se introducen como valores convencionales para seleccionar
esta realizaci\'on. Para \(N=12\cdot4\cdot120=5760\), la longitud marcada es
\[
 L_{\alpha}=\alpha_{\mathrm{HMT}}^{16}
 \frac{N-1}{4N}L_*
 =\frac{5759}{23040}\alpha_{\mathrm{HMT}}^{16}L_* .
 \tag{A.2}\label{eq:vi-cat-longitud-alpha}
\]

Abreviamos en esta secci\'on \(\hbar=\hbar_{\mathrm{ret}}\) y
\(c=c_{\mathrm{int}}\). Sobre la fibra positiva no nula, la energ\'ia
\(H\), la longitud de acci\'on reducida \(\Lambda\) y la lectura radial
conjugada \(R\) satisfacen las dos relaciones operatorias siguientes.
Definimos adem\'as \(M:=H/c^2\) y declaramos la identificaci\'on f\'isica
radial \(R:=GM/c^2\). Al compararla, sobre la fibra invertible, con la
relaci\'on operatoria \(R=L_{\alpha}^{2}H/(\hbar c)\), se obtiene
\[
 H\Lambda=\hbar c\,I,\qquad
 R\Lambda=L_{\alpha}^{2}I,\qquad
 R=\frac{L_{\alpha}^{2}}{\hbar c}H,\qquad
 \boxed{G=\frac{c^{3}L_{\alpha}^{2}}{\hbar}} .
 \tag{A.3}\label{eq:vi-cat-realizacion-lg}
\]
El reconocimiento metrol\'ogico de \(G\) es posterior y no interviene en la
elecci\'on del estado, del lector ni de sus coeficientes.

\subsection{Registro M10 y conservaci\'on de la ley de masas}
\label{sec:vi-catalogo-m10-activo}

El c\'odigo M10 del repertorio documental remite a la
realizaci\'on tipada de \eqref{eq:vi-cat-realizacion-lg}. M10 no postula una
part\'icula gravitatoria ni convierte en salida HMT el registro externo que el
inventario conserva bajo la clase
\texttt{GRAVITATIONAL\_OR\_HYPOTHETICAL}: dicho registro permanece como dato
metrol\'ogico externo, separado de la construcci\'on
\(L_{\alpha}\mapsto G\).

La escala energ\'etica inducida por la misma longitud es
\[
 Q_L:=\frac{\hbar c}{L_{\alpha}}.
 \tag{A.4}\label{eq:vi-cat-escala-q}
\]
Si una coordenatización equivalente sustituye \(Q_L\) por \(Q_L'\), el coeficiente
electr\'onico cotransforma seg\'un
\[
 b_e'=b_e\frac{Q_L}{Q_L'},\qquad Q_L'b_e'=Q_Lb_e .
 \tag{A.5}\label{eq:vi-cat-cotransformacion-be}
\]
Las definiciones de la energ\'ia del reloj y del factor electr\'onico dan
\[
 E_{\mathrm{clk}}=Q_L,\qquad
 m_{\mathrm{clk}}=\frac{E_{\mathrm{clk}}}{c^2},\qquad
 b_e=\frac{E_e^{(0)}}{E_{\mathrm{clk}}},\qquad
 m_{\mathrm{clk}}b_e=\frac{E_e^{(0)}}{c^2}.
\]
La cancelaci\'on exacta de \(Q_L\) demuestra que la realizaci\'on de \(G\) no
deriva de nuevo las masas, no selecciona rutas y no modifica las tablas ya
construidas. La especializaci\'on \(\Xi_\Gamma=1\) del teorema de
compatibilidad equivale a
\(\bar\lambda_\Gamma=R_{108}=L_{\alpha}\): identifica una coincidencia de
longitudes y no afirma ni crea la existencia de una especie de masa de Planck
en el cat\'alogo.

La d\'ecada \(10^{-34}\) pertenece exclusivamente a la rama absoluta
acci\'on--reloj--masa,
\(\hbar=\eta_{\mathrm{ret}}10^{-34}\mathcal U_S\); no se reinserta en el
exponente relativo, en el car\'acter de ruta ni como correcci\'on por especie.

\subsection{Canales constitutivos y momento orientado de Catal\'an}
\label{sec:vi-catalogo-propietario-constitutivo}

Sean \(0<q_+,q_-<1\) los canales etiquetados de transporte y \(P_+,P_-\)
sus proyectores ortogonales de hoja. Para \(s\in(0,1)\), definimos
\[
 \mathcal C_2(s)=\sum_{k\geq0}
 \frac{(-1)^k s^{2k+1}}{(2k+1)^2},\qquad
 \mathcal D=s\frac{\mathrm d}{\mathrm ds},\qquad
 f(s)=\frac{1-s^3}{1-s^4}.
 \tag{A.6}\label{eq:vi-cat-constitutivas}
\]
La convergencia local uniforme de la serie y de sus derivadas permite aplicar
dos veces \(\mathcal D\). La serie geom\'etrica resultante da
\[
 \mathcal D^2\mathcal C_2(s)=\frac{s}{1+s^2},\qquad
 f(s)=\frac{1+s+s^2}{s(1+s)}\mathcal D^2\mathcal C_2(s).
\]
Las identidades \(90=3\cdot30\) y \(120=4\cdot30\) determinan
\[
 s_{\pm}=q_{\pm}^{30},\qquad
 r_{\pm}=f(s_{\pm})
 =\frac{1-q_{\pm}^{90}}{1-q_{\pm}^{120}}.
\]
Por tanto, \(V=r_+P_++r_-P_-\) re\'une las respuestas conservando sus
etiquetas. La serie de \(\mathcal C_2\) converge absolutamente en \(s=1\);
su valor es \(G_{\mathrm{Cat}}\), distinguido del coeficiente gravitatorio
\(G\).

La funci\'on completa se recupera de \(f\) mediante
\[
 \mathcal C_2(s)=\int_0^s\log\!\left(\frac{s}{t}\right)
 \frac{1+t}{1+t+t^2}f(t)\,\mathrm dt.
\]
En efecto, \((1+t)f(t)/(1+t+t^2)=1/(1+t^2)\); dos aplicaciones de
\(\mathcal D\) restituyen \(s/(1+s^2)\). Las soluciones de la ecuaci\'on
homog\'enea son \(a\log s+b\); la condici\'on
\(\mathcal C_2(s)\to0\) cuando \(s\downarrow0\) fija ambos coeficientes a
cero. Quedan expl\'icitos la respuesta, su restituci\'on regular y el
significado de la coordenada de Catal\'an usada en las fichas.

% Generado por technical/generar_apendices_catalogo.py. No editar a mano.
\clearpage
\section{Catálogos de multisecciones, colecciones y clases}
\label{sec:vi-catalogos-estructurales}

Este apéndice organiza tres niveles distintos del inventario interno. Las
13 multisecciones reúnen celdas del atlas; las 56 colecciones clasifican rutas
internas; las 23 clases documentan la articulación con sectores y
representaciones físicas. Estos cardinales no cuentan el mismo objeto y
no se suman para obtener un número de partículas. Las tablas acompañan las
definiciones y demostraciones del artículo: su función es permitir localizar
cada registro, sus incidencias y sus lectores, no sustituir una prueba por un censo.

Se conserva el orden documental de las filas, sus identificadores y todas
las repeticiones que figuren en las fuentes. Los identificadores literales
permiten cotejar las tablas con los archivos JSON. Una raya indica una celda
vacía de la fuente; no representa un cero. Los archivos conservan además todos
los campos y cifras, los estados originales y la procedencia verificable.
\CatalogRouteReference
% La remisión depende del punto de entrada: artículo o catálogo.

\subsection{Las 13 multisecciones del atlas}

Para una multisección \(M\), sea \(\mathcal C(M)\) su conjunto de celdas y
sea \(\gamma_{B_1}(c)\) la ruta primaria que el atlas conserva en la celda
\(c\). En este apéndice estructural se define la fibra
\[
 B_1(M):=
 \operatorname{span}\bigl\{e_{\gamma_{B_1}(c)}:c\in\mathcal C(M)\bigr\}
 \subseteq \ell^2(\mathcal R_{\mathrm{adm}}).
\]
Las rutas primarias de celdas distintas determinan vectores básicos distintos;
por consiguiente,
\(\operatorname{rank}B_1(M)=\dim B_1(M)=\#\mathcal C(M)\).
Esta definición fija el objeto al que se refiere el rango de la tabla; la
columna de colecciones cuenta las
colecciones de ruta presentes, no especies físicas. Se conservan separadamente
los espectros enteros de los lectores \(K_D\) y \(K_\Omega\).
En la notación literal \texttt{x} indica multiplicidad. Las 13 filas tienen
el estado documental de operador diagonal exacto sobre la fibra \(B_1\).
Su expresión como escalar físico requiere el transporte declarado, salvo
la selección de rango uno prevista en el propio registro. Este enunciado
reproduce el alcance registrado y no convierte todos los espectros en masas únicas.

\begingroup
\setlength{\tabcolsep}{4pt}
\setlength{\extrarowheight}{4pt}
\renewcommand{\arraystretch}{1.13}
\begin{longtable}{@{}>{\raggedright\arraybackslash}p{\dimexpr 0.19871795\linewidth-2\tabcolsep\relax}>{\raggedright\arraybackslash}p{\dimexpr 0.19871795\linewidth-2\tabcolsep\relax}>{\raggedright\arraybackslash}p{\dimexpr 0.30769231\linewidth-2\tabcolsep\relax}>{\raggedright\arraybackslash}p{\dimexpr 0.29487179\linewidth-2\tabcolsep\relax}@{}}
\caption{Las 13 multisecciones, sin selección de filas.}\label{tab:vi-multisecciones}\\
\toprule
Multisección & Rango e incidencia & Clasificación y profundidad & Lectores enteros \\
\midrule
\endfirsthead
\multicolumn{4}{l}{\textit{Continuación}}\\
\toprule
Multisección & Rango e incidencia & Clasificación y profundidad & Lectores enteros \\
\midrule
\endhead
\midrule
\multicolumn{4}{r}{Continúa en la página siguiente}\\
\endfoot
\bottomrule
\endlastfoot
% VI-CATALOG-ROW multisecciones 1
charged\_\allowbreak{}lepton\_\allowbreak{}G0 & \textit{Rango}: 1\par \textit{Celdas (identificadores)}: 41\par \textit{Colecciones}: 1 & \textit{Colección}: charged\_\allowbreak{}lepton\par \textit{Profundidad}: G0\_\allowbreak{}center\par \textit{Color}: none & \textit{K D}: (80,\allowbreak{}54,\allowbreak{}6)\allowbreak{}x1\par \textit{K Ω}: (80,\allowbreak{}54,\allowbreak{}6)\allowbreak{}x1 \\
% VI-CATALOG-ROW multisecciones 2
charged\_\allowbreak{}lepton\_\allowbreak{}G2 & \textit{Rango}: 8\par \textit{Celdas (identificadores)}: 21,\allowbreak{}23,\allowbreak{}25,\allowbreak{}39,\allowbreak{}43,\allowbreak{}57,\allowbreak{}59,\allowbreak{}61\par \textit{Colecciones}: 8 & \textit{Colección}: charged\_\allowbreak{}lepton\par \textit{Profundidad}: G2\par \textit{Color}: none & \textit{K D}: (0,\allowbreak{}24,\allowbreak{}0)\allowbreak{}x1; (40,\allowbreak{}78,\allowbreak{}27)\allowbreak{}x2; (80,\allowbreak{}54,\allowbreak{}7)\allowbreak{}x1; (80,\allowbreak{}66,\allowbreak{}2)\allowbreak{}x1; (80,\allowbreak{}66,\allowbreak{}3)\allowbreak{}x1; (100,\allowbreak{}66,\allowbreak{}0)\allowbreak{}x1; (100,\allowbreak{}66,\allowbreak{}2)\allowbreak{}x1\par \textit{K Ω}: (80,\allowbreak{}54,\allowbreak{}6)\allowbreak{}x8 \\
% VI-CATALOG-ROW multisecciones 3
charged\_\allowbreak{}lepton\_\allowbreak{}G4 & \textit{Rango}: 16\par \textit{Celdas (identificadores)}: 1,\allowbreak{}3,\allowbreak{}5,\allowbreak{}7,\allowbreak{}9,\allowbreak{}19,\allowbreak{}27,\allowbreak{}37,\allowbreak{}45,\allowbreak{}55,\allowbreak{}63,\allowbreak{}73,\allowbreak{}75,\allowbreak{}77,\allowbreak{}79,\allowbreak{}81\par \textit{Colecciones}: 16 & \textit{Colección}: charged\_\allowbreak{}lepton\par \textit{Profundidad}: G4\par \textit{Color}: none & \textit{K D}: (0,\allowbreak{}24,\allowbreak{}0)\allowbreak{}x1; (40,\allowbreak{}24,\allowbreak{}2)\allowbreak{}x2; (40,\allowbreak{}54,\allowbreak{}6)\allowbreak{}x2; (40,\allowbreak{}84,\allowbreak{}30)\allowbreak{}x2; (60,\allowbreak{}78,\allowbreak{}25)\allowbreak{}x3; (80,\allowbreak{}66,\allowbreak{}2)\allowbreak{}x2; (80,\allowbreak{}66,\allowbreak{}3)\allowbreak{}x3; (80,\allowbreak{}66,\allowbreak{}4)\allowbreak{}x1\par \textit{K Ω}: (24,\allowbreak{}54,\allowbreak{}2)\allowbreak{}x1; (56,\allowbreak{}54,\allowbreak{}5)\allowbreak{}x4; (80,\allowbreak{}54,\allowbreak{}6)\allowbreak{}x11 \\
% VI-CATALOG-ROW multisecciones 4
neutrino\_\allowbreak{}G1 & \textit{Rango}: 2\par \textit{Celdas (identificadores)}: 40,\allowbreak{}42\par \textit{Colecciones}: 2 & \textit{Colección}: neutrino\par \textit{Profundidad}: G1\par \textit{Color}: none & \textit{K D}: (40,\allowbreak{}24,\allowbreak{}2)\allowbreak{}x1; (40,\allowbreak{}78,\allowbreak{}27)\allowbreak{}x1\par \textit{K Ω}: (40,\allowbreak{}54,\allowbreak{}4)\allowbreak{}x1; (80,\allowbreak{}54,\allowbreak{}6)\allowbreak{}x1 \\
% VI-CATALOG-ROW multisecciones 5
neutrino\_\allowbreak{}G2 & \textit{Rango}: 4\par \textit{Celdas (identificadores)}: 22,\allowbreak{}24,\allowbreak{}58,\allowbreak{}60\par \textit{Colecciones}: 4 & \textit{Colección}: neutrino\par \textit{Profundidad}: G2\par \textit{Color}: none & \textit{K D}: (0,\allowbreak{}24,\allowbreak{}0)\allowbreak{}x1; (40,\allowbreak{}84,\allowbreak{}30)\allowbreak{}x1; (60,\allowbreak{}78,\allowbreak{}25)\allowbreak{}x1; (80,\allowbreak{}66,\allowbreak{}3)\allowbreak{}x1\par \textit{K Ω}: (4,\allowbreak{}54,\allowbreak{}2)\allowbreak{}x1; (80,\allowbreak{}54,\allowbreak{}6)\allowbreak{}x3 \\
% VI-CATALOG-ROW multisecciones 6
neutrino\_\allowbreak{}G3 & \textit{Rango}: 6\par \textit{Celdas (identificadores)}: 20,\allowbreak{}26,\allowbreak{}38,\allowbreak{}44,\allowbreak{}56,\allowbreak{}62\par \textit{Colecciones}: 6 & \textit{Colección}: neutrino\par \textit{Profundidad}: G3\par \textit{Color}: none & \textit{K D}: (40,\allowbreak{}24,\allowbreak{}2)\allowbreak{}x2; (40,\allowbreak{}78,\allowbreak{}27)\allowbreak{}x1; (60,\allowbreak{}84,\allowbreak{}28)\allowbreak{}x1; (80,\allowbreak{}54,\allowbreak{}6)\allowbreak{}x1; (80,\allowbreak{}66,\allowbreak{}3)\allowbreak{}x1\par \textit{K Ω}: (36,\allowbreak{}54,\allowbreak{}4)\allowbreak{}x1; (56,\allowbreak{}54,\allowbreak{}5)\allowbreak{}x1; (80,\allowbreak{}54,\allowbreak{}6)\allowbreak{}x4 \\
% VI-CATALOG-ROW multisecciones 7
neutrino\_\allowbreak{}G4 & \textit{Rango}: 8\par \textit{Celdas (identificadores)}: 2,\allowbreak{}4,\allowbreak{}6,\allowbreak{}8,\allowbreak{}74,\allowbreak{}76,\allowbreak{}78,\allowbreak{}80\par \textit{Colecciones}: 8 & \textit{Colección}: neutrino\par \textit{Profundidad}: G4\par \textit{Color}: none & \textit{K D}: (0,\allowbreak{}24,\allowbreak{}0)\allowbreak{}x1; (40,\allowbreak{}24,\allowbreak{}2)\allowbreak{}x1; (40,\allowbreak{}78,\allowbreak{}27)\allowbreak{}x2; (40,\allowbreak{}84,\allowbreak{}30)\allowbreak{}x1; (80,\allowbreak{}54,\allowbreak{}7)\allowbreak{}x1; (80,\allowbreak{}66,\allowbreak{}2)\allowbreak{}x1; (100,\allowbreak{}66,\allowbreak{}0)\allowbreak{}x1\par \textit{K Ω}: (36,\allowbreak{}54,\allowbreak{}4)\allowbreak{}x1; (56,\allowbreak{}54,\allowbreak{}5)\allowbreak{}x1; (80,\allowbreak{}54,\allowbreak{}6)\allowbreak{}x6 \\
% VI-CATALOG-ROW multisecciones 8
quark\_\allowbreak{}down\_\allowbreak{}G1 & \textit{Rango}: 2\par \textit{Celdas (identificadores)}: 32,\allowbreak{}50\par \textit{Colecciones}: 2 & \textit{Colección}: quark\_\allowbreak{}down\par \textit{Profundidad}: G1\par \textit{Color}: B | G & \textit{K D}: (40,\allowbreak{}24,\allowbreak{}2)\allowbreak{}x1; (40,\allowbreak{}78,\allowbreak{}27)\allowbreak{}x1\par \textit{K Ω}: (40,\allowbreak{}54,\allowbreak{}4)\allowbreak{}x1; (80,\allowbreak{}54,\allowbreak{}6)\allowbreak{}x1 \\
% VI-CATALOG-ROW multisecciones 9
quark\_\allowbreak{}down\_\allowbreak{}G2 & \textit{Rango}: 4\par \textit{Celdas (identificadores)}: 30,\allowbreak{}34,\allowbreak{}48,\allowbreak{}52\par \textit{Colecciones}: 4 & \textit{Colección}: quark\_\allowbreak{}down\par \textit{Profundidad}: G2\par \textit{Color}: B | G | R & \textit{K D}: (0,\allowbreak{}24,\allowbreak{}0)\allowbreak{}x1; (40,\allowbreak{}84,\allowbreak{}30)\allowbreak{}x1; (60,\allowbreak{}78,\allowbreak{}25)\allowbreak{}x1; (80,\allowbreak{}66,\allowbreak{}2)\allowbreak{}x1\par \textit{K Ω}: (4,\allowbreak{}54,\allowbreak{}2)\allowbreak{}x1; (80,\allowbreak{}54,\allowbreak{}6)\allowbreak{}x3 \\
% VI-CATALOG-ROW multisecciones 10
quark\_\allowbreak{}down\_\allowbreak{}G3 & \textit{Rango}: 6\par \textit{Celdas (identificadores)}: 12,\allowbreak{}14,\allowbreak{}16,\allowbreak{}66,\allowbreak{}68,\allowbreak{}70\par \textit{Colecciones}: 6 & \textit{Colección}: quark\_\allowbreak{}down\par \textit{Profundidad}: G3\par \textit{Color}: B | G | R & \textit{K D}: (40,\allowbreak{}24,\allowbreak{}2)\allowbreak{}x2; (40,\allowbreak{}78,\allowbreak{}27)\allowbreak{}x1; (60,\allowbreak{}84,\allowbreak{}28)\allowbreak{}x1; (80,\allowbreak{}54,\allowbreak{}6)\allowbreak{}x1; (80,\allowbreak{}66,\allowbreak{}4)\allowbreak{}x1\par \textit{K Ω}: (36,\allowbreak{}54,\allowbreak{}4)\allowbreak{}x1; (80,\allowbreak{}54,\allowbreak{}6)\allowbreak{}x5 \\
% VI-CATALOG-ROW multisecciones 11
quark\_\allowbreak{}down\_\allowbreak{}G4 & \textit{Rango}: 8\par \textit{Celdas (identificadores)}: 10,\allowbreak{}18,\allowbreak{}28,\allowbreak{}36,\allowbreak{}46,\allowbreak{}54,\allowbreak{}64,\allowbreak{}72\par \textit{Colecciones}: 8 & \textit{Colección}: quark\_\allowbreak{}down\par \textit{Profundidad}: G4\par \textit{Color}: B | G | R & \textit{K D}: (0,\allowbreak{}24,\allowbreak{}0)\allowbreak{}x1; (40,\allowbreak{}24,\allowbreak{}2)\allowbreak{}x1; (40,\allowbreak{}78,\allowbreak{}27)\allowbreak{}x2; (40,\allowbreak{}84,\allowbreak{}30)\allowbreak{}x1; (80,\allowbreak{}54,\allowbreak{}7)\allowbreak{}x1; (80,\allowbreak{}66,\allowbreak{}3)\allowbreak{}x1; (100,\allowbreak{}66,\allowbreak{}2)\allowbreak{}x1\par \textit{K Ω}: (36,\allowbreak{}54,\allowbreak{}4)\allowbreak{}x1; (80,\allowbreak{}54,\allowbreak{}6)\allowbreak{}x7 \\
% VI-CATALOG-ROW multisecciones 12
quark\_\allowbreak{}up\_\allowbreak{}G1 & \textit{Rango}: 4\par \textit{Celdas (identificadores)}: 31,\allowbreak{}33,\allowbreak{}49,\allowbreak{}51\par \textit{Colecciones}: 3 & \textit{Colección}: quark\_\allowbreak{}up\par \textit{Profundidad}: G1\par \textit{Color}: B | G | R & \textit{K D}: (40,\allowbreak{}54,\allowbreak{}6)\allowbreak{}x1; (60,\allowbreak{}78,\allowbreak{}25)\allowbreak{}x1; (80,\allowbreak{}66,\allowbreak{}2)\allowbreak{}x1; (80,\allowbreak{}66,\allowbreak{}3)\allowbreak{}x1\par \textit{K Ω}: (80,\allowbreak{}54,\allowbreak{}6)\allowbreak{}x4 \\
% VI-CATALOG-ROW multisecciones 13
quark\_\allowbreak{}up\_\allowbreak{}G3 & \textit{Rango}: 12\par \textit{Celdas (identificadores)}: 11,\allowbreak{}13,\allowbreak{}15,\allowbreak{}17,\allowbreak{}29,\allowbreak{}35,\allowbreak{}47,\allowbreak{}53,\allowbreak{}65,\allowbreak{}67,\allowbreak{}69,\allowbreak{}71\par \textit{Colecciones}: 12 & \textit{Colección}: quark\_\allowbreak{}up\par \textit{Profundidad}: G3\par \textit{Color}: B | G | R & \textit{K D}: (40,\allowbreak{}24,\allowbreak{}2)\allowbreak{}x2; (40,\allowbreak{}78,\allowbreak{}27)\allowbreak{}x2; (60,\allowbreak{}84,\allowbreak{}28)\allowbreak{}x1; (80,\allowbreak{}54,\allowbreak{}6)\allowbreak{}x3; (80,\allowbreak{}66,\allowbreak{}3)\allowbreak{}x1; (80,\allowbreak{}66,\allowbreak{}4)\allowbreak{}x1; (100,\allowbreak{}66,\allowbreak{}0)\allowbreak{}x1; (100,\allowbreak{}66,\allowbreak{}2)\allowbreak{}x1\par \textit{K Ω}: (56,\allowbreak{}54,\allowbreak{}5)\allowbreak{}x3; (80,\allowbreak{}54,\allowbreak{}6)\allowbreak{}x9 \\
\end{longtable}
\endgroup
\subsection{Las 56 colecciones de ruta}

La firma que identifica una colección se conserva literalmente, con sus signos
y separadores. Las celdas y multisecciones hacen visible que una colección de
ruta no equivale por definición a una especie física. Los perfiles conjuntos
cuentan las combinaciones registradas dentro de la colección. Los espectros
decimales del carácter y los espectros completos de firma permanecen íntegros
en \texttt{catalogo\_familias.json}; aquí se expresan la incidencia y los
lectores enteros que permiten orientarse en esa información.

\begingroup
\setlength{\tabcolsep}{4pt}
\setlength{\extrarowheight}{4pt}
\renewcommand{\arraystretch}{1.13}
\begin{longtable}{@{}>{\raggedright\arraybackslash}p{\dimexpr 0.21794872\linewidth-2\tabcolsep\relax}>{\raggedright\arraybackslash}p{\dimexpr 0.34615385\linewidth-2\tabcolsep\relax}>{\raggedright\arraybackslash}p{\dimexpr 0.21794872\linewidth-2\tabcolsep\relax}>{\raggedright\arraybackslash}p{\dimexpr 0.21794872\linewidth-2\tabcolsep\relax}@{}}
\caption{Las 56 colecciones internas de ruta.}\label{tab:vi-familias}\\
\toprule
Colección y firma & Incidencia & \(K_D\) & \(K_\Omega\) \\
\midrule
\endfirsthead
\multicolumn{4}{l}{\textit{Continuación}}\\
\toprule
Colección y firma & Incidencia & \(K_D\) & \(K_\Omega\) \\
\midrule
\endhead
\midrule
\multicolumn{4}{r}{Continúa en la página siguiente}\\
\endfoot
\bottomrule
\endlastfoot
% VI-CATALOG-ROW colecciones 1
1\par 24|\allowbreak{}0|\allowbreak{}12|\allowbreak{}0|\allowbreak{}-12|\allowbreak{}+++-\kern0pt{}-\kern0pt{}-+++-\kern0pt{}-\kern0pt{}- & \textit{Número de celdas}: 5\par \textit{Celdas (identificadores)}: 21,\allowbreak{}31,\allowbreak{}41,\allowbreak{}51,\allowbreak{}81\par \textit{Multisecciones}: charged\_\allowbreak{}lepton\_\allowbreak{}G0 | charged\_\allowbreak{}lepton\_\allowbreak{}G2 | charged\_\allowbreak{}lepton\_\allowbreak{}G4 | quark\_\allowbreak{}up\_\allowbreak{}G1\par \textit{Colecciones}: charged\_\allowbreak{}lepton | quark\_\allowbreak{}up\par \textit{Color}: R | none\par \textit{Perfiles conjuntos}: 5 & (0,\allowbreak{}24,\allowbreak{}0)\allowbreak{}x2; (40,\allowbreak{}54,\allowbreak{}6)\allowbreak{}x1; (60,\allowbreak{}78,\allowbreak{}25)\allowbreak{}x1; (80,\allowbreak{}54,\allowbreak{}6)\allowbreak{}x1 & (24,\allowbreak{}54,\allowbreak{}2)\allowbreak{}x1; (80,\allowbreak{}54,\allowbreak{}6)\allowbreak{}x4 \\
% VI-CATALOG-ROW colecciones 2
2\par 24|\allowbreak{}0|\allowbreak{}4|\allowbreak{}0|\allowbreak{}-12|\allowbreak{}+++-\kern0pt{}-\kern0pt{}-+++-\kern0pt{}-\kern0pt{}- & \textit{Número de celdas}: 2\par \textit{Celdas (identificadores)}: 11,\allowbreak{}61\par \textit{Multisecciones}: charged\_\allowbreak{}lepton\_\allowbreak{}G2 | quark\_\allowbreak{}up\_\allowbreak{}G3\par \textit{Colecciones}: charged\_\allowbreak{}lepton | quark\_\allowbreak{}up\par \textit{Color}: R | none\par \textit{Perfiles conjuntos}: 2 & (60,\allowbreak{}84,\allowbreak{}28)\allowbreak{}x1; (80,\allowbreak{}54,\allowbreak{}7)\allowbreak{}x1 & (80,\allowbreak{}54,\allowbreak{}6)\allowbreak{}x2 \\
% VI-CATALOG-ROW colecciones 3
3\par 24|\allowbreak{}0|\allowbreak{}8|\allowbreak{}0|\allowbreak{}-12|\allowbreak{}+++-\kern0pt{}-\kern0pt{}-+++-\kern0pt{}-\kern0pt{}- & \textit{Número de celdas}: 2\par \textit{Celdas (identificadores)}: 1,\allowbreak{}71\par \textit{Multisecciones}: charged\_\allowbreak{}lepton\_\allowbreak{}G4 | quark\_\allowbreak{}up\_\allowbreak{}G3\par \textit{Colecciones}: charged\_\allowbreak{}lepton | quark\_\allowbreak{}up\par \textit{Color}: R | none\par \textit{Perfiles conjuntos}: 2 & (60,\allowbreak{}78,\allowbreak{}25)\allowbreak{}x1; (80,\allowbreak{}54,\allowbreak{}6)\allowbreak{}x1 & (80,\allowbreak{}54,\allowbreak{}6)\allowbreak{}x2 \\
% VI-CATALOG-ROW colecciones 4
4\par 28|\allowbreak{}-6|\allowbreak{}-2|\allowbreak{}-2|\allowbreak{}-8|\allowbreak{}0++-\kern0pt{}-\kern0pt{}-0++-\kern0pt{}-\kern0pt{}- & \textit{Número de celdas}: 1\par \textit{Celdas (identificadores)}: 53\par \textit{Multisecciones}: quark\_\allowbreak{}up\_\allowbreak{}G3\par \textit{Colecciones}: quark\_\allowbreak{}up\par \textit{Color}: G\par \textit{Perfiles conjuntos}: 1 & (80,\allowbreak{}66,\allowbreak{}3)\allowbreak{}x1 & (56,\allowbreak{}54,\allowbreak{}5)\allowbreak{}x1 \\
% VI-CATALOG-ROW colecciones 5
5\par 28|\allowbreak{}-6|\allowbreak{}-4|\allowbreak{}-2|\allowbreak{}-8|\allowbreak{}0++-\kern0pt{}-\kern0pt{}-0++-\kern0pt{}-\kern0pt{}- & \textit{Número de celdas}: 1\par \textit{Celdas (identificadores)}: 20\par \textit{Multisecciones}: neutrino\_\allowbreak{}G3\par \textit{Colecciones}: neutrino\par \textit{Color}: none\par \textit{Perfiles conjuntos}: 1 & (80,\allowbreak{}66,\allowbreak{}3)\allowbreak{}x1 & (56,\allowbreak{}54,\allowbreak{}5)\allowbreak{}x1 \\
% VI-CATALOG-ROW colecciones 6
6\par 28|\allowbreak{}-6|\allowbreak{}0|\allowbreak{}0|\allowbreak{}-12|\allowbreak{}+++-\kern0pt{}-\kern0pt{}-+++-\kern0pt{}-\kern0pt{}- & \textit{Número de celdas}: 2\par \textit{Celdas (identificadores)}: 30,\allowbreak{}70\par \textit{Multisecciones}: quark\_\allowbreak{}down\_\allowbreak{}G2 | quark\_\allowbreak{}down\_\allowbreak{}G3\par \textit{Colecciones}: quark\_\allowbreak{}down\par \textit{Color}: G\par \textit{Perfiles conjuntos}: 2 & (40,\allowbreak{}24,\allowbreak{}2)\allowbreak{}x1; (80,\allowbreak{}66,\allowbreak{}2)\allowbreak{}x1 & (80,\allowbreak{}54,\allowbreak{}6)\allowbreak{}x2 \\
% VI-CATALOG-ROW colecciones 7
7\par 28|\allowbreak{}-6|\allowbreak{}10|\allowbreak{}0|\allowbreak{}-12|\allowbreak{}+++-\kern0pt{}-\kern0pt{}-+++-\kern0pt{}-\kern0pt{}- & \textit{Número de celdas}: 2\par \textit{Celdas (identificadores)}: 33,\allowbreak{}64\par \textit{Multisecciones}: quark\_\allowbreak{}down\_\allowbreak{}G4 | quark\_\allowbreak{}up\_\allowbreak{}G1\par \textit{Colecciones}: quark\_\allowbreak{}down | quark\_\allowbreak{}up\par \textit{Color}: G\par \textit{Perfiles conjuntos}: 2 & (40,\allowbreak{}24,\allowbreak{}2)\allowbreak{}x1; (80,\allowbreak{}66,\allowbreak{}2)\allowbreak{}x1 & (36,\allowbreak{}54,\allowbreak{}4)\allowbreak{}x1; (80,\allowbreak{}54,\allowbreak{}6)\allowbreak{}x1 \\
% VI-CATALOG-ROW colecciones 8
8\par 28|\allowbreak{}-6|\allowbreak{}12|\allowbreak{}0|\allowbreak{}-12|\allowbreak{}+++-\kern0pt{}-\kern0pt{}-+++-\kern0pt{}-\kern0pt{}- & \textit{Número de celdas}: 1\par \textit{Celdas (identificadores)}: 10\par \textit{Multisecciones}: quark\_\allowbreak{}down\_\allowbreak{}G4\par \textit{Colecciones}: quark\_\allowbreak{}down\par \textit{Color}: G\par \textit{Perfiles conjuntos}: 1 & (100,\allowbreak{}66,\allowbreak{}2)\allowbreak{}x1 & (80,\allowbreak{}54,\allowbreak{}6)\allowbreak{}x1 \\
% VI-CATALOG-ROW colecciones 9
9\par 28|\allowbreak{}-6|\allowbreak{}2|\allowbreak{}0|\allowbreak{}-12|\allowbreak{}+++-\kern0pt{}-\kern0pt{}-+++-\kern0pt{}-\kern0pt{}- & \textit{Número de celdas}: 1\par \textit{Celdas (identificadores)}: 3\par \textit{Multisecciones}: charged\_\allowbreak{}lepton\_\allowbreak{}G4\par \textit{Colecciones}: charged\_\allowbreak{}lepton\par \textit{Color}: none\par \textit{Perfiles conjuntos}: 1 & (40,\allowbreak{}84,\allowbreak{}30)\allowbreak{}x1 & (80,\allowbreak{}54,\allowbreak{}6)\allowbreak{}x1 \\
% VI-CATALOG-ROW colecciones 10
10\par 28|\allowbreak{}-6|\allowbreak{}4|\allowbreak{}-2|\allowbreak{}-8|\allowbreak{}++0-\kern0pt{}-\kern0pt{}-++0-\kern0pt{}-\kern0pt{}- & \textit{Número de celdas}: 1\par \textit{Celdas (identificadores)}: 50\par \textit{Multisecciones}: quark\_\allowbreak{}down\_\allowbreak{}G1\par \textit{Colecciones}: quark\_\allowbreak{}down\par \textit{Color}: G\par \textit{Perfiles conjuntos}: 1 & (40,\allowbreak{}78,\allowbreak{}27)\allowbreak{}x1 & (80,\allowbreak{}54,\allowbreak{}6)\allowbreak{}x1 \\
% VI-CATALOG-ROW colecciones 11
11\par 28|\allowbreak{}-6|\allowbreak{}4|\allowbreak{}0|\allowbreak{}-12|\allowbreak{}+++-\kern0pt{}-\kern0pt{}-+++-\kern0pt{}-\kern0pt{}- & \textit{Número de celdas}: 2\par \textit{Celdas (identificadores)}: 9,\allowbreak{}40\par \textit{Multisecciones}: charged\_\allowbreak{}lepton\_\allowbreak{}G4 | neutrino\_\allowbreak{}G1\par \textit{Colecciones}: charged\_\allowbreak{}lepton | neutrino\par \textit{Color}: none\par \textit{Perfiles conjuntos}: 2 & (40,\allowbreak{}24,\allowbreak{}2)\allowbreak{}x1; (80,\allowbreak{}66,\allowbreak{}2)\allowbreak{}x1 & (40,\allowbreak{}54,\allowbreak{}4)\allowbreak{}x1; (80,\allowbreak{}54,\allowbreak{}6)\allowbreak{}x1 \\
% VI-CATALOG-ROW colecciones 12
12\par 28|\allowbreak{}-6|\allowbreak{}4|\allowbreak{}0|\allowbreak{}-8|\allowbreak{}+0+-\kern0pt{}-\kern0pt{}-+0+-\kern0pt{}-\kern0pt{}- & \textit{Número de celdas}: 1\par \textit{Celdas (identificadores)}: 80\par \textit{Multisecciones}: neutrino\_\allowbreak{}G4\par \textit{Colecciones}: neutrino\par \textit{Color}: none\par \textit{Perfiles conjuntos}: 1 & (40,\allowbreak{}78,\allowbreak{}27)\allowbreak{}x1 & (80,\allowbreak{}54,\allowbreak{}6)\allowbreak{}x1 \\
% VI-CATALOG-ROW colecciones 13
13\par 28|\allowbreak{}-6|\allowbreak{}6|\allowbreak{}-2|\allowbreak{}-8|\allowbreak{}++0-\kern0pt{}-\kern0pt{}-++0-\kern0pt{}-\kern0pt{}- & \textit{Número de celdas}: 1\par \textit{Celdas (identificadores)}: 74\par \textit{Multisecciones}: neutrino\_\allowbreak{}G4\par \textit{Colecciones}: neutrino\par \textit{Color}: none\par \textit{Perfiles conjuntos}: 1 & (40,\allowbreak{}78,\allowbreak{}27)\allowbreak{}x1 & (80,\allowbreak{}54,\allowbreak{}6)\allowbreak{}x1 \\
% VI-CATALOG-ROW colecciones 14
14\par 28|\allowbreak{}-6|\allowbreak{}6|\allowbreak{}0|\allowbreak{}-12|\allowbreak{}+++-\kern0pt{}-\kern0pt{}-+++-\kern0pt{}-\kern0pt{}- & \textit{Número de celdas}: 3\par \textit{Celdas (identificadores)}: 13,\allowbreak{}43,\allowbreak{}63\par \textit{Multisecciones}: charged\_\allowbreak{}lepton\_\allowbreak{}G2 | charged\_\allowbreak{}lepton\_\allowbreak{}G4 | quark\_\allowbreak{}up\_\allowbreak{}G3\par \textit{Colecciones}: charged\_\allowbreak{}lepton | quark\_\allowbreak{}up\par \textit{Color}: G | none\par \textit{Perfiles conjuntos}: 3 & (40,\allowbreak{}24,\allowbreak{}2)\allowbreak{}x1; (80,\allowbreak{}66,\allowbreak{}2)\allowbreak{}x1; (100,\allowbreak{}66,\allowbreak{}2)\allowbreak{}x1 & (80,\allowbreak{}54,\allowbreak{}6)\allowbreak{}x3 \\
% VI-CATALOG-ROW colecciones 15
15\par 28|\allowbreak{}-6|\allowbreak{}6|\allowbreak{}0|\allowbreak{}-8|\allowbreak{}+0+-\kern0pt{}-\kern0pt{}-+0+-\kern0pt{}-\kern0pt{}- & \textit{Número de celdas}: 1\par \textit{Celdas (identificadores)}: 23\par \textit{Multisecciones}: charged\_\allowbreak{}lepton\_\allowbreak{}G2\par \textit{Colecciones}: charged\_\allowbreak{}lepton\par \textit{Color}: none\par \textit{Perfiles conjuntos}: 1 & (40,\allowbreak{}78,\allowbreak{}27)\allowbreak{}x1 & (80,\allowbreak{}54,\allowbreak{}6)\allowbreak{}x1 \\
% VI-CATALOG-ROW colecciones 16
16\par 28|\allowbreak{}-6|\allowbreak{}8|\allowbreak{}0|\allowbreak{}-12|\allowbreak{}+++-\kern0pt{}-\kern0pt{}-+++-\kern0pt{}-\kern0pt{}- & \textit{Número de celdas}: 1\par \textit{Celdas (identificadores)}: 60\par \textit{Multisecciones}: neutrino\_\allowbreak{}G2\par \textit{Colecciones}: neutrino\par \textit{Color}: none\par \textit{Perfiles conjuntos}: 1 & (40,\allowbreak{}84,\allowbreak{}30)\allowbreak{}x1 & (80,\allowbreak{}54,\allowbreak{}6)\allowbreak{}x1 \\
% VI-CATALOG-ROW colecciones 17
17\par 28|\allowbreak{}6|\allowbreak{}-10|\allowbreak{}0|\allowbreak{}-12|\allowbreak{}+++-\kern0pt{}-\kern0pt{}-+++-\kern0pt{}-\kern0pt{}- & \textit{Número de celdas}: 1\par \textit{Celdas (identificadores)}: 19\par \textit{Multisecciones}: charged\_\allowbreak{}lepton\_\allowbreak{}G4\par \textit{Colecciones}: charged\_\allowbreak{}lepton\par \textit{Color}: none\par \textit{Perfiles conjuntos}: 1 & (40,\allowbreak{}84,\allowbreak{}30)\allowbreak{}x1 & (80,\allowbreak{}54,\allowbreak{}6)\allowbreak{}x1 \\
% VI-CATALOG-ROW colecciones 18
18\par 28|\allowbreak{}6|\allowbreak{}-2|\allowbreak{}0|\allowbreak{}-12|\allowbreak{}+++-\kern0pt{}-\kern0pt{}-+++-\kern0pt{}-\kern0pt{}- & \textit{Número de celdas}: 5\par \textit{Celdas (identificadores)}: 18,\allowbreak{}29,\allowbreak{}49,\allowbreak{}59,\allowbreak{}79\par \textit{Multisecciones}: charged\_\allowbreak{}lepton\_\allowbreak{}G2 | charged\_\allowbreak{}lepton\_\allowbreak{}G4 | quark\_\allowbreak{}down\_\allowbreak{}G4 | quark\_\allowbreak{}up\_\allowbreak{}G1 | quark\_\allowbreak{}up\_\allowbreak{}G3\par \textit{Colecciones}: charged\_\allowbreak{}lepton | quark\_\allowbreak{}down | quark\_\allowbreak{}up\par \textit{Color}: B | none\par \textit{Perfiles conjuntos}: 4 & (40,\allowbreak{}24,\allowbreak{}2)\allowbreak{}x1; (40,\allowbreak{}78,\allowbreak{}27)\allowbreak{}x1; (80,\allowbreak{}66,\allowbreak{}3)\allowbreak{}x2; (100,\allowbreak{}66,\allowbreak{}0)\allowbreak{}x1 & (80,\allowbreak{}54,\allowbreak{}6)\allowbreak{}x5 \\
% VI-CATALOG-ROW colecciones 19
19\par 28|\allowbreak{}6|\allowbreak{}-4|\allowbreak{}0|\allowbreak{}-12|\allowbreak{}+++-\kern0pt{}-\kern0pt{}-+++-\kern0pt{}-\kern0pt{}- & \textit{Número de celdas}: 2\par \textit{Celdas (identificadores)}: 2,\allowbreak{}42\par \textit{Multisecciones}: neutrino\_\allowbreak{}G1 | neutrino\_\allowbreak{}G4\par \textit{Colecciones}: neutrino\par \textit{Color}: none\par \textit{Perfiles conjuntos}: 2 & (40,\allowbreak{}78,\allowbreak{}27)\allowbreak{}x1; (100,\allowbreak{}66,\allowbreak{}0)\allowbreak{}x1 & (80,\allowbreak{}54,\allowbreak{}6)\allowbreak{}x2 \\
% VI-CATALOG-ROW colecciones 20
20\par 28|\allowbreak{}6|\allowbreak{}-6|\allowbreak{}0|\allowbreak{}-12|\allowbreak{}+++-\kern0pt{}-\kern0pt{}-+++-\kern0pt{}-\kern0pt{}- & \textit{Número de celdas}: 1\par \textit{Celdas (identificadores)}: 69\par \textit{Multisecciones}: quark\_\allowbreak{}up\_\allowbreak{}G3\par \textit{Colecciones}: quark\_\allowbreak{}up\par \textit{Color}: B\par \textit{Perfiles conjuntos}: 1 & (80,\allowbreak{}66,\allowbreak{}4)\allowbreak{}x1 & (80,\allowbreak{}54,\allowbreak{}6)\allowbreak{}x1 \\
% VI-CATALOG-ROW colecciones 21
21\par 28|\allowbreak{}6|\allowbreak{}-8|\allowbreak{}0|\allowbreak{}-12|\allowbreak{}+++-\kern0pt{}-\kern0pt{}-+++-\kern0pt{}-\kern0pt{}- & \textit{Número de celdas}: 1\par \textit{Celdas (identificadores)}: 52\par \textit{Multisecciones}: quark\_\allowbreak{}down\_\allowbreak{}G2\par \textit{Colecciones}: quark\_\allowbreak{}down\par \textit{Color}: B\par \textit{Perfiles conjuntos}: 1 & (40,\allowbreak{}84,\allowbreak{}30)\allowbreak{}x1 & (80,\allowbreak{}54,\allowbreak{}6)\allowbreak{}x1 \\
% VI-CATALOG-ROW colecciones 22
22\par 28|\allowbreak{}6|\allowbreak{}0|\allowbreak{}0|\allowbreak{}-12|\allowbreak{}+++-\kern0pt{}-\kern0pt{}-+++-\kern0pt{}-\kern0pt{}- & \textit{Número de celdas}: 4\par \textit{Celdas (identificadores)}: 22,\allowbreak{}32,\allowbreak{}72,\allowbreak{}73\par \textit{Multisecciones}: charged\_\allowbreak{}lepton\_\allowbreak{}G4 | neutrino\_\allowbreak{}G2 | quark\_\allowbreak{}down\_\allowbreak{}G1 | quark\_\allowbreak{}down\_\allowbreak{}G4\par \textit{Colecciones}: charged\_\allowbreak{}lepton | neutrino | quark\_\allowbreak{}down\par \textit{Color}: B | none\par \textit{Perfiles conjuntos}: 3 & (40,\allowbreak{}24,\allowbreak{}2)\allowbreak{}x1; (40,\allowbreak{}78,\allowbreak{}27)\allowbreak{}x1; (80,\allowbreak{}66,\allowbreak{}3)\allowbreak{}x2 & (40,\allowbreak{}54,\allowbreak{}4)\allowbreak{}x1; (80,\allowbreak{}54,\allowbreak{}6)\allowbreak{}x3 \\
% VI-CATALOG-ROW colecciones 23
23\par 28|\allowbreak{}6|\allowbreak{}2|\allowbreak{}0|\allowbreak{}-12|\allowbreak{}+++-\kern0pt{}-\kern0pt{}-+++-\kern0pt{}-\kern0pt{}- & \textit{Número de celdas}: 2\par \textit{Celdas (identificadores)}: 8,\allowbreak{}39\par \textit{Multisecciones}: charged\_\allowbreak{}lepton\_\allowbreak{}G2 | neutrino\_\allowbreak{}G4\par \textit{Colecciones}: charged\_\allowbreak{}lepton | neutrino\par \textit{Color}: none\par \textit{Perfiles conjuntos}: 2 & (40,\allowbreak{}24,\allowbreak{}2)\allowbreak{}x1; (40,\allowbreak{}78,\allowbreak{}27)\allowbreak{}x1 & (36,\allowbreak{}54,\allowbreak{}4)\allowbreak{}x1; (80,\allowbreak{}54,\allowbreak{}6)\allowbreak{}x1 \\
% VI-CATALOG-ROW colecciones 24
24\par 28|\allowbreak{}6|\allowbreak{}4|\allowbreak{}0|\allowbreak{}-12|\allowbreak{}+++-\kern0pt{}-\kern0pt{}-+++-\kern0pt{}-\kern0pt{}- & \textit{Número de celdas}: 1\par \textit{Celdas (identificadores)}: 12\par \textit{Multisecciones}: quark\_\allowbreak{}down\_\allowbreak{}G3\par \textit{Colecciones}: quark\_\allowbreak{}down\par \textit{Color}: B\par \textit{Perfiles conjuntos}: 1 & (80,\allowbreak{}66,\allowbreak{}4)\allowbreak{}x1 & (80,\allowbreak{}54,\allowbreak{}6)\allowbreak{}x1 \\
% VI-CATALOG-ROW colecciones 25
25\par 28|\allowbreak{}6|\allowbreak{}8|\allowbreak{}0|\allowbreak{}-12|\allowbreak{}+++-\kern0pt{}-\kern0pt{}-+++-\kern0pt{}-\kern0pt{}- & \textit{Número de celdas}: 1\par \textit{Celdas (identificadores)}: 62\par \textit{Multisecciones}: neutrino\_\allowbreak{}G3\par \textit{Colecciones}: neutrino\par \textit{Color}: none\par \textit{Perfiles conjuntos}: 1 & (40,\allowbreak{}24,\allowbreak{}2)\allowbreak{}x1 & (80,\allowbreak{}54,\allowbreak{}6)\allowbreak{}x1 \\
% VI-CATALOG-ROW colecciones 26
26\par 30|\allowbreak{}-6|\allowbreak{}-4|\allowbreak{}-2|\allowbreak{}-8|\allowbreak{}0++-\kern0pt{}-\kern0pt{}-0++-\kern0pt{}-\kern0pt{}- & \textit{Número de celdas}: 1\par \textit{Celdas (identificadores)}: 4\par \textit{Multisecciones}: neutrino\_\allowbreak{}G4\par \textit{Colecciones}: neutrino\par \textit{Color}: none\par \textit{Perfiles conjuntos}: 1 & (80,\allowbreak{}54,\allowbreak{}7)\allowbreak{}x1 & (56,\allowbreak{}54,\allowbreak{}5)\allowbreak{}x1 \\
% VI-CATALOG-ROW colecciones 27
27\par 30|\allowbreak{}-6|\allowbreak{}-8|\allowbreak{}0|\allowbreak{}-8|\allowbreak{}+0+-\kern0pt{}-\kern0pt{}-+0+-\kern0pt{}-\kern0pt{}- & \textit{Número de celdas}: 1\par \textit{Celdas (identificadores)}: 55\par \textit{Multisecciones}: charged\_\allowbreak{}lepton\_\allowbreak{}G4\par \textit{Colecciones}: charged\_\allowbreak{}lepton\par \textit{Color}: none\par \textit{Perfiles conjuntos}: 1 & (60,\allowbreak{}78,\allowbreak{}25)\allowbreak{}x1 & (80,\allowbreak{}54,\allowbreak{}6)\allowbreak{}x1 \\
% VI-CATALOG-ROW colecciones 28
28\par 30|\allowbreak{}-6|\allowbreak{}0|\allowbreak{}-2|\allowbreak{}-8|\allowbreak{}++0-\kern0pt{}-\kern0pt{}-++0-\kern0pt{}-\kern0pt{}- & \textit{Número de celdas}: 2\par \textit{Celdas (identificadores)}: 24,\allowbreak{}34\par \textit{Multisecciones}: neutrino\_\allowbreak{}G2 | quark\_\allowbreak{}down\_\allowbreak{}G2\par \textit{Colecciones}: neutrino | quark\_\allowbreak{}down\par \textit{Color}: R | none\par \textit{Perfiles conjuntos}: 2 & (0,\allowbreak{}24,\allowbreak{}0)\allowbreak{}x1; (60,\allowbreak{}78,\allowbreak{}25)\allowbreak{}x1 & (4,\allowbreak{}54,\allowbreak{}2)\allowbreak{}x1; (80,\allowbreak{}54,\allowbreak{}6)\allowbreak{}x1 \\
% VI-CATALOG-ROW colecciones 29
29\par 30|\allowbreak{}-6|\allowbreak{}0|\allowbreak{}-2|\allowbreak{}-8|\allowbreak{}0++-\kern0pt{}-\kern0pt{}-0++-\kern0pt{}-\kern0pt{}- & \textit{Número de celdas}: 1\par \textit{Celdas (identificadores)}: 75\par \textit{Multisecciones}: charged\_\allowbreak{}lepton\_\allowbreak{}G4\par \textit{Colecciones}: charged\_\allowbreak{}lepton\par \textit{Color}: none\par \textit{Perfiles conjuntos}: 1 & (40,\allowbreak{}54,\allowbreak{}6)\allowbreak{}x1 & (56,\allowbreak{}54,\allowbreak{}5)\allowbreak{}x1 \\
% VI-CATALOG-ROW colecciones 30
30\par 30|\allowbreak{}-6|\allowbreak{}0|\allowbreak{}0|\allowbreak{}-12|\allowbreak{}+++-\kern0pt{}-\kern0pt{}-+++-\kern0pt{}-\kern0pt{}- & \textit{Número de celdas}: 1\par \textit{Celdas (identificadores)}: 14\par \textit{Multisecciones}: quark\_\allowbreak{}down\_\allowbreak{}G3\par \textit{Colecciones}: quark\_\allowbreak{}down\par \textit{Color}: R\par \textit{Perfiles conjuntos}: 1 & (80,\allowbreak{}54,\allowbreak{}6)\allowbreak{}x1 & (80,\allowbreak{}54,\allowbreak{}6)\allowbreak{}x1 \\
% VI-CATALOG-ROW colecciones 31
31\par 30|\allowbreak{}-6|\allowbreak{}0|\allowbreak{}0|\allowbreak{}-8|\allowbreak{}+0+-\kern0pt{}-\kern0pt{}-+0+-\kern0pt{}-\kern0pt{}- & \textit{Número de celdas}: 1\par \textit{Celdas (identificadores)}: 54\par \textit{Multisecciones}: quark\_\allowbreak{}down\_\allowbreak{}G4\par \textit{Colecciones}: quark\_\allowbreak{}down\par \textit{Color}: R\par \textit{Perfiles conjuntos}: 1 & (0,\allowbreak{}24,\allowbreak{}0)\allowbreak{}x1 & (80,\allowbreak{}54,\allowbreak{}6)\allowbreak{}x1 \\
% VI-CATALOG-ROW colecciones 32
32\par 30|\allowbreak{}-6|\allowbreak{}4|\allowbreak{}0|\allowbreak{}-12|\allowbreak{}+++-\kern0pt{}-\kern0pt{}-+++-\kern0pt{}-\kern0pt{}- & \textit{Número de celdas}: 1\par \textit{Celdas (identificadores)}: 65\par \textit{Multisecciones}: quark\_\allowbreak{}up\_\allowbreak{}G3\par \textit{Colecciones}: quark\_\allowbreak{}up\par \textit{Color}: R\par \textit{Perfiles conjuntos}: 1 & (80,\allowbreak{}54,\allowbreak{}6)\allowbreak{}x1 & (80,\allowbreak{}54,\allowbreak{}6)\allowbreak{}x1 \\
% VI-CATALOG-ROW colecciones 33
33\par 30|\allowbreak{}-6|\allowbreak{}8|\allowbreak{}0|\allowbreak{}-12|\allowbreak{}+++-\kern0pt{}-\kern0pt{}-+++-\kern0pt{}-\kern0pt{}- & \textit{Número de celdas}: 1\par \textit{Celdas (identificadores)}: 44\par \textit{Multisecciones}: neutrino\_\allowbreak{}G3\par \textit{Colecciones}: neutrino\par \textit{Color}: none\par \textit{Perfiles conjuntos}: 1 & (60,\allowbreak{}84,\allowbreak{}28)\allowbreak{}x1 & (80,\allowbreak{}54,\allowbreak{}6)\allowbreak{}x1 \\
% VI-CATALOG-ROW colecciones 34
34\par 30|\allowbreak{}6|\allowbreak{}-4|\allowbreak{}0|\allowbreak{}-12|\allowbreak{}+++-\kern0pt{}-\kern0pt{}-+++-\kern0pt{}-\kern0pt{}- & \textit{Número de celdas}: 1\par \textit{Celdas (identificadores)}: 68\par \textit{Multisecciones}: quark\_\allowbreak{}down\_\allowbreak{}G3\par \textit{Colecciones}: quark\_\allowbreak{}down\par \textit{Color}: R\par \textit{Perfiles conjuntos}: 1 & (60,\allowbreak{}84,\allowbreak{}28)\allowbreak{}x1 & (80,\allowbreak{}54,\allowbreak{}6)\allowbreak{}x1 \\
% VI-CATALOG-ROW colecciones 35
35\par 30|\allowbreak{}6|\allowbreak{}-8|\allowbreak{}0|\allowbreak{}-12|\allowbreak{}+++-\kern0pt{}-\kern0pt{}-+++-\kern0pt{}-\kern0pt{}- & \textit{Número de celdas}: 1\par \textit{Celdas (identificadores)}: 17\par \textit{Multisecciones}: quark\_\allowbreak{}up\_\allowbreak{}G3\par \textit{Colecciones}: quark\_\allowbreak{}up\par \textit{Color}: R\par \textit{Perfiles conjuntos}: 1 & (80,\allowbreak{}54,\allowbreak{}6)\allowbreak{}x1 & (80,\allowbreak{}54,\allowbreak{}6)\allowbreak{}x1 \\
% VI-CATALOG-ROW colecciones 36
36\par 30|\allowbreak{}6|\allowbreak{}0|\allowbreak{}0|\allowbreak{}-12|\allowbreak{}+++-\kern0pt{}-\kern0pt{}-+++-\kern0pt{}-\kern0pt{}- & \textit{Número de celdas}: 2\par \textit{Celdas (identificadores)}: 38,\allowbreak{}58\par \textit{Multisecciones}: neutrino\_\allowbreak{}G2 | neutrino\_\allowbreak{}G3\par \textit{Colecciones}: neutrino\par \textit{Color}: none\par \textit{Perfiles conjuntos}: 2 & (60,\allowbreak{}78,\allowbreak{}25)\allowbreak{}x1; (80,\allowbreak{}54,\allowbreak{}6)\allowbreak{}x1 & (80,\allowbreak{}54,\allowbreak{}6)\allowbreak{}x2 \\
% VI-CATALOG-ROW colecciones 37
37\par 30|\allowbreak{}6|\allowbreak{}0|\allowbreak{}0|\allowbreak{}-8|\allowbreak{}+++-0-+++-0- & \textit{Número de celdas}: 1\par \textit{Celdas (identificadores)}: 78\par \textit{Multisecciones}: neutrino\_\allowbreak{}G4\par \textit{Colecciones}: neutrino\par \textit{Color}: none\par \textit{Perfiles conjuntos}: 1 & (0,\allowbreak{}24,\allowbreak{}0)\allowbreak{}x1 & (80,\allowbreak{}54,\allowbreak{}6)\allowbreak{}x1 \\
% VI-CATALOG-ROW colecciones 38
38\par 30|\allowbreak{}6|\allowbreak{}0|\allowbreak{}2|\allowbreak{}-8|\allowbreak{}+++-\kern0pt{}-0+++-\kern0pt{}-0 & \textit{Número de celdas}: 1\par \textit{Celdas (identificadores)}: 27\par \textit{Multisecciones}: charged\_\allowbreak{}lepton\_\allowbreak{}G4\par \textit{Colecciones}: charged\_\allowbreak{}lepton\par \textit{Color}: none\par \textit{Perfiles conjuntos}: 1 & (40,\allowbreak{}54,\allowbreak{}6)\allowbreak{}x1 & (56,\allowbreak{}54,\allowbreak{}5)\allowbreak{}x1 \\
% VI-CATALOG-ROW colecciones 39
39\par 30|\allowbreak{}6|\allowbreak{}0|\allowbreak{}2|\allowbreak{}-8|\allowbreak{}+++0-\kern0pt{}-+++0-\kern0pt{}- & \textit{Número de celdas}: 1\par \textit{Celdas (identificadores)}: 48\par \textit{Multisecciones}: quark\_\allowbreak{}down\_\allowbreak{}G2\par \textit{Colecciones}: quark\_\allowbreak{}down\par \textit{Color}: R\par \textit{Perfiles conjuntos}: 1 & (0,\allowbreak{}24,\allowbreak{}0)\allowbreak{}x1 & (4,\allowbreak{}54,\allowbreak{}2)\allowbreak{}x1 \\
% VI-CATALOG-ROW colecciones 40
40\par 30|\allowbreak{}6|\allowbreak{}4|\allowbreak{}0|\allowbreak{}-12|\allowbreak{}+++-\kern0pt{}-\kern0pt{}-+++-\kern0pt{}-\kern0pt{}- & \textit{Número de celdas}: 1\par \textit{Celdas (identificadores)}: 7\par \textit{Multisecciones}: charged\_\allowbreak{}lepton\_\allowbreak{}G4\par \textit{Colecciones}: charged\_\allowbreak{}lepton\par \textit{Color}: none\par \textit{Perfiles conjuntos}: 1 & (60,\allowbreak{}78,\allowbreak{}25)\allowbreak{}x1 & (80,\allowbreak{}54,\allowbreak{}6)\allowbreak{}x1 \\
% VI-CATALOG-ROW colecciones 41
41\par 30|\allowbreak{}6|\allowbreak{}8|\allowbreak{}0|\allowbreak{}-12|\allowbreak{}+++-\kern0pt{}-\kern0pt{}-+++-\kern0pt{}-\kern0pt{}- & \textit{Número de celdas}: 1\par \textit{Celdas (identificadores)}: 28\par \textit{Multisecciones}: quark\_\allowbreak{}down\_\allowbreak{}G4\par \textit{Colecciones}: quark\_\allowbreak{}down\par \textit{Color}: R\par \textit{Perfiles conjuntos}: 1 & (80,\allowbreak{}54,\allowbreak{}7)\allowbreak{}x1 & (80,\allowbreak{}54,\allowbreak{}6)\allowbreak{}x1 \\
% VI-CATALOG-ROW colecciones 42
42\par 34|\allowbreak{}0|\allowbreak{}-12|\allowbreak{}-2|\allowbreak{}-8|\allowbreak{}0++-\kern0pt{}-\kern0pt{}-0++-\kern0pt{}-\kern0pt{}- & \textit{Número de celdas}: 1\par \textit{Celdas (identificadores)}: 45\par \textit{Multisecciones}: charged\_\allowbreak{}lepton\_\allowbreak{}G4\par \textit{Colecciones}: charged\_\allowbreak{}lepton\par \textit{Color}: none\par \textit{Perfiles conjuntos}: 1 & (80,\allowbreak{}66,\allowbreak{}4)\allowbreak{}x1 & (56,\allowbreak{}54,\allowbreak{}5)\allowbreak{}x1 \\
% VI-CATALOG-ROW colecciones 43
43\par 34|\allowbreak{}0|\allowbreak{}-12|\allowbreak{}0|\allowbreak{}-12|\allowbreak{}+++-\kern0pt{}-\kern0pt{}-+++-\kern0pt{}-\kern0pt{}- & \textit{Número de celdas}: 2\par \textit{Celdas (identificadores)}: 57,\allowbreak{}76\par \textit{Multisecciones}: charged\_\allowbreak{}lepton\_\allowbreak{}G2 | neutrino\_\allowbreak{}G4\par \textit{Colecciones}: charged\_\allowbreak{}lepton | neutrino\par \textit{Color}: none\par \textit{Perfiles conjuntos}: 2 & (40,\allowbreak{}84,\allowbreak{}30)\allowbreak{}x1; (80,\allowbreak{}66,\allowbreak{}2)\allowbreak{}x1 & (80,\allowbreak{}54,\allowbreak{}6)\allowbreak{}x2 \\
% VI-CATALOG-ROW colecciones 44
44\par 34|\allowbreak{}0|\allowbreak{}-16|\allowbreak{}0|\allowbreak{}-8|\allowbreak{}+++-0-+++-0- & \textit{Número de celdas}: 1\par \textit{Celdas (identificadores)}: 37\par \textit{Multisecciones}: charged\_\allowbreak{}lepton\_\allowbreak{}G4\par \textit{Colecciones}: charged\_\allowbreak{}lepton\par \textit{Color}: none\par \textit{Perfiles conjuntos}: 1 & (40,\allowbreak{}24,\allowbreak{}2)\allowbreak{}x1 & (80,\allowbreak{}54,\allowbreak{}6)\allowbreak{}x1 \\
% VI-CATALOG-ROW colecciones 45
45\par 34|\allowbreak{}0|\allowbreak{}-16|\allowbreak{}0|\allowbreak{}-8|\allowbreak{}+0+-\kern0pt{}-\kern0pt{}-+0+-\kern0pt{}-\kern0pt{}- & \textit{Número de celdas}: 1\par \textit{Celdas (identificadores)}: 15\par \textit{Multisecciones}: quark\_\allowbreak{}up\_\allowbreak{}G3\par \textit{Colecciones}: quark\_\allowbreak{}up\par \textit{Color}: B\par \textit{Perfiles conjuntos}: 1 & (40,\allowbreak{}78,\allowbreak{}27)\allowbreak{}x1 & (80,\allowbreak{}54,\allowbreak{}6)\allowbreak{}x1 \\
% VI-CATALOG-ROW colecciones 46
46\par 34|\allowbreak{}0|\allowbreak{}-20|\allowbreak{}-2|\allowbreak{}-8|\allowbreak{}0++-\kern0pt{}-\kern0pt{}-0++-\kern0pt{}-\kern0pt{}- & \textit{Número de celdas}: 1\par \textit{Celdas (identificadores)}: 35\par \textit{Multisecciones}: quark\_\allowbreak{}up\_\allowbreak{}G3\par \textit{Colecciones}: quark\_\allowbreak{}up\par \textit{Color}: B\par \textit{Perfiles conjuntos}: 1 & (100,\allowbreak{}66,\allowbreak{}0)\allowbreak{}x1 & (56,\allowbreak{}54,\allowbreak{}5)\allowbreak{}x1 \\
% VI-CATALOG-ROW colecciones 47
47\par 34|\allowbreak{}0|\allowbreak{}-4|\allowbreak{}-2|\allowbreak{}-8|\allowbreak{}++0-\kern0pt{}-\kern0pt{}-++0-\kern0pt{}-\kern0pt{}- & \textit{Número de celdas}: 1\par \textit{Celdas (identificadores)}: 56\par \textit{Multisecciones}: neutrino\_\allowbreak{}G3\par \textit{Colecciones}: neutrino\par \textit{Color}: none\par \textit{Perfiles conjuntos}: 1 & (40,\allowbreak{}24,\allowbreak{}2)\allowbreak{}x1 & (36,\allowbreak{}54,\allowbreak{}4)\allowbreak{}x1 \\
% VI-CATALOG-ROW colecciones 48
48\par 34|\allowbreak{}0|\allowbreak{}-4|\allowbreak{}0|\allowbreak{}-12|\allowbreak{}+++-\kern0pt{}-\kern0pt{}-+++-\kern0pt{}-\kern0pt{}- & \textit{Número de celdas}: 3\par \textit{Celdas (identificadores)}: 6,\allowbreak{}25,\allowbreak{}46\par \textit{Multisecciones}: charged\_\allowbreak{}lepton\_\allowbreak{}G2 | neutrino\_\allowbreak{}G4 | quark\_\allowbreak{}down\_\allowbreak{}G4\par \textit{Colecciones}: charged\_\allowbreak{}lepton | neutrino | quark\_\allowbreak{}down\par \textit{Color}: B | none\par \textit{Perfiles conjuntos}: 2 & (80,\allowbreak{}66,\allowbreak{}2)\allowbreak{}x1; (80,\allowbreak{}66,\allowbreak{}3)\allowbreak{}x2 & (80,\allowbreak{}54,\allowbreak{}6)\allowbreak{}x3 \\
% VI-CATALOG-ROW colecciones 49
49\par 34|\allowbreak{}0|\allowbreak{}-4|\allowbreak{}2|\allowbreak{}-8|\allowbreak{}+++-\kern0pt{}-0+++-\kern0pt{}-0 & \textit{Número de celdas}: 1\par \textit{Celdas (identificadores)}: 67\par \textit{Multisecciones}: quark\_\allowbreak{}up\_\allowbreak{}G3\par \textit{Colecciones}: quark\_\allowbreak{}up\par \textit{Color}: G\par \textit{Perfiles conjuntos}: 1 & (100,\allowbreak{}66,\allowbreak{}2)\allowbreak{}x1 & (56,\allowbreak{}54,\allowbreak{}5)\allowbreak{}x1 \\
% VI-CATALOG-ROW colecciones 50
50\par 34|\allowbreak{}0|\allowbreak{}-8|\allowbreak{}-2|\allowbreak{}-8|\allowbreak{}++0-\kern0pt{}-\kern0pt{}-++0-\kern0pt{}-\kern0pt{}- & \textit{Número de celdas}: 1\par \textit{Celdas (identificadores)}: 66\par \textit{Multisecciones}: quark\_\allowbreak{}down\_\allowbreak{}G3\par \textit{Colecciones}: quark\_\allowbreak{}down\par \textit{Color}: B\par \textit{Perfiles conjuntos}: 1 & (40,\allowbreak{}78,\allowbreak{}27)\allowbreak{}x1 & (80,\allowbreak{}54,\allowbreak{}6)\allowbreak{}x1 \\
% VI-CATALOG-ROW colecciones 51
51\par 34|\allowbreak{}0|\allowbreak{}-8|\allowbreak{}-2|\allowbreak{}-8|\allowbreak{}0++-\kern0pt{}-\kern0pt{}-0++-\kern0pt{}-\kern0pt{}- & \textit{Número de celdas}: 1\par \textit{Celdas (identificadores)}: 77\par \textit{Multisecciones}: charged\_\allowbreak{}lepton\_\allowbreak{}G4\par \textit{Colecciones}: charged\_\allowbreak{}lepton\par \textit{Color}: none\par \textit{Perfiles conjuntos}: 1 & (80,\allowbreak{}66,\allowbreak{}3)\allowbreak{}x1 & (56,\allowbreak{}54,\allowbreak{}5)\allowbreak{}x1 \\
% VI-CATALOG-ROW colecciones 52
52\par 34|\allowbreak{}0|\allowbreak{}-8|\allowbreak{}0|\allowbreak{}-12|\allowbreak{}+++-\kern0pt{}-\kern0pt{}-+++-\kern0pt{}-\kern0pt{}- & \textit{Número de celdas}: 1\par \textit{Celdas (identificadores)}: 36\par \textit{Multisecciones}: quark\_\allowbreak{}down\_\allowbreak{}G4\par \textit{Colecciones}: quark\_\allowbreak{}down\par \textit{Color}: G\par \textit{Perfiles conjuntos}: 1 & (40,\allowbreak{}84,\allowbreak{}30)\allowbreak{}x1 & (80,\allowbreak{}54,\allowbreak{}6)\allowbreak{}x1 \\
% VI-CATALOG-ROW colecciones 53
53\par 34|\allowbreak{}0|\allowbreak{}-8|\allowbreak{}0|\allowbreak{}-8|\allowbreak{}+0+-\kern0pt{}-\kern0pt{}-+0+-\kern0pt{}-\kern0pt{}- & \textit{Número de celdas}: 1\par \textit{Celdas (identificadores)}: 5\par \textit{Multisecciones}: charged\_\allowbreak{}lepton\_\allowbreak{}G4\par \textit{Colecciones}: charged\_\allowbreak{}lepton\par \textit{Color}: none\par \textit{Perfiles conjuntos}: 1 & (40,\allowbreak{}24,\allowbreak{}2)\allowbreak{}x1 & (80,\allowbreak{}54,\allowbreak{}6)\allowbreak{}x1 \\
% VI-CATALOG-ROW colecciones 54
54\par 34|\allowbreak{}0|\allowbreak{}0|\allowbreak{}-2|\allowbreak{}-8|\allowbreak{}++0-\kern0pt{}-\kern0pt{}-++0-\kern0pt{}-\kern0pt{}- & \textit{Número de celdas}: 1\par \textit{Celdas (identificadores)}: 26\par \textit{Multisecciones}: neutrino\_\allowbreak{}G3\par \textit{Colecciones}: neutrino\par \textit{Color}: none\par \textit{Perfiles conjuntos}: 1 & (40,\allowbreak{}78,\allowbreak{}27)\allowbreak{}x1 & (80,\allowbreak{}54,\allowbreak{}6)\allowbreak{}x1 \\
% VI-CATALOG-ROW colecciones 55
55\par 34|\allowbreak{}0|\allowbreak{}0|\allowbreak{}0|\allowbreak{}-8|\allowbreak{}+0+-\kern0pt{}-\kern0pt{}-+0+-\kern0pt{}-\kern0pt{}- & \textit{Número de celdas}: 1\par \textit{Celdas (identificadores)}: 47\par \textit{Multisecciones}: quark\_\allowbreak{}up\_\allowbreak{}G3\par \textit{Colecciones}: quark\_\allowbreak{}up\par \textit{Color}: G\par \textit{Perfiles conjuntos}: 1 & (40,\allowbreak{}78,\allowbreak{}27)\allowbreak{}x1 & (80,\allowbreak{}54,\allowbreak{}6)\allowbreak{}x1 \\
% VI-CATALOG-ROW colecciones 56
56\par 34|\allowbreak{}0|\allowbreak{}0|\allowbreak{}2|\allowbreak{}-8|\allowbreak{}+++0-\kern0pt{}-+++0-\kern0pt{}- & \textit{Número de celdas}: 1\par \textit{Celdas (identificadores)}: 16\par \textit{Multisecciones}: quark\_\allowbreak{}down\_\allowbreak{}G3\par \textit{Colecciones}: quark\_\allowbreak{}down\par \textit{Color}: G\par \textit{Perfiles conjuntos}: 1 & (40,\allowbreak{}24,\allowbreak{}2)\allowbreak{}x1 & (36,\allowbreak{}54,\allowbreak{}4)\allowbreak{}x1 \\
\end{longtable}
\endgroup
\subsection{Las 23 clases documentadas}

Estas filas reúnen sectores identificados, residencias y representación. La clave
de estado remite a la leyenda contigua y conserva el estatuto de la fuente.
En particular, un operador interno exacto, una evaluación condicionada y una
realización escalar cerrada se mantienen como afirmaciones distintas. La tabla
no actualiza su estatuto ni transforma un registro condicionado en una predicción.

\begingroup
\setlength{\tabcolsep}{4pt}
\setlength{\extrarowheight}{4pt}
\renewcommand{\arraystretch}{1.13}
\begin{longtable}{@{}>{\raggedright\arraybackslash}p{\dimexpr 0.08860759\linewidth-2\tabcolsep\relax}>{\raggedright\arraybackslash}p{\dimexpr 0.41772152\linewidth-2\tabcolsep\relax}>{\raggedright\arraybackslash}p{\dimexpr 0.49367089\linewidth-2\tabcolsep\relax}@{}}
\caption{Estados originales de las clases.}\label{tab:vi-estados-clases}\\
\toprule
Clave & Estado literal & Lectura documental \\
\midrule
\endfirsthead
\multicolumn{3}{l}{\textit{Continuación}}\\
\toprule
Clave & Estado literal & Lectura documental \\
\midrule
\endhead
\midrule
\multicolumn{3}{r}{Continúa en la página siguiente}\\
\endfoot
\bottomrule
\endlastfoot
C01 & CLOSED\_\allowbreak{}ELECTRON\_\allowbreak{}TWO\_\allowbreak{}WAY & Electrón: cierre bidireccional documentado. \\
C02 & CLOSED\_\allowbreak{}NULL\_\allowbreak{}CHART & Carta nula cerrada en este corte. \\
C03 & NOT\_\allowbreak{}A\_\allowbreak{}UNIVERSAL\_\allowbreak{}MASS\_\allowbreak{}SCALAR & La clase no se expresa como un escalar universal de masa. \\
C04 & OPEN\_\allowbreak{}FLAVOR\_\allowbreak{}COLOR\_\allowbreak{}SCALARIZATION & Escalarización de sabor y color abierta en este corte. \\
C05 & OPEN\_\allowbreak{}MIXING\_\allowbreak{}AND\_\allowbreak{}SCALE & Mezcla y escala abiertas en este corte. \\
C06 & OPEN\_\allowbreak{}NAMED\_\allowbreak{}STATE\_\allowbreak{}SCALARIZATION & Escalarización del estado individualizado abierta en este corte. \\
C07 & OPEN\_\allowbreak{}PHYSICAL\_\allowbreak{}SCALARIZATION & Escalarización física abierta en este corte. \\
C08 & OPEN\_\allowbreak{}ROUTE\_\allowbreak{}AND\_\allowbreak{}SCALARIZATION & Ruta y escalarización abiertas en este corte. \\
\end{longtable}
\endgroup

\begingroup\setlength{\parskip}{0pt}\renewcommand{\CatalogTableLeading}{10.6}
\setlength{\tabcolsep}{4pt}
\setlength{\extrarowheight}{2pt}
\renewcommand{\arraystretch}{1.04}
\begin{longtable}{@{}>{\raggedright\arraybackslash}p{\dimexpr 0.17307692\linewidth-2\tabcolsep\relax}>{\raggedright\arraybackslash}p{\dimexpr 0.38461538\linewidth-2\tabcolsep\relax}>{\raggedright\arraybackslash}p{\dimexpr 0.33333333\linewidth-2\tabcolsep\relax}>{\raggedright\arraybackslash}p{\dimexpr 0.10897436\linewidth-2\tabcolsep\relax}@{}}
\caption{Las 23 clases del catálogo, en su orden documental.}\label{tab:vi-clases}\\
\toprule
Clase & Representantes y residencia & Incidencia y representación & Estado \\
\midrule
\endfirsthead
\multicolumn{4}{l}{\textit{Continuación}}\\
\toprule
Clase & Representantes y residencia & Incidencia y representación & Estado \\
\midrule
\endhead
\midrule
\multicolumn{4}{r}{Continúa en la página siguiente}\\
\endfoot
\bottomrule
\endlastfoot
% VI-CATALOG-ROW clases 1
SM-LC-1\par lepton\_\allowbreak{}cargado & \textit{Representantes}: e− / e+\par \textit{Residencia}: centro (5,\allowbreak{}5) y antisección & \textit{Celdas seleccionadas}: 1\par \textit{Multisecciones}: charged\_\allowbreak{}lepton\_\allowbreak{}G0\par \textit{Carga}: ±1 e\par \textit{Representación}: sin color & C01 \\
% VI-CATALOG-ROW clases 2
SM-LC-2\par lepton\_\allowbreak{}cargado & \textit{Representantes}: μ− / μ+\par \textit{Residencia}: multisección leptón cargado & \textit{Celdas seleccionadas}: 8\par \textit{Multisecciones}: charged\_\allowbreak{}lepton\_\allowbreak{}G2\par \textit{Carga}: ±1 e\par \textit{Representación}: sin color & C07 \\
% VI-CATALOG-ROW clases 3
SM-LC-3\par lepton\_\allowbreak{}cargado & \textit{Representantes}: τ− / τ+\par \textit{Residencia}: multisección leptón cargado & \textit{Celdas seleccionadas}: 16\par \textit{Multisecciones}: charged\_\allowbreak{}lepton\_\allowbreak{}G4\par \textit{Carga}: ±1 e\par \textit{Representación}: sin color & C07 \\
% VI-CATALOG-ROW clases 4
SM-NU-1\par neutrino & \textit{Representantes}: ν\_\allowbreak{}e / anti-ν\_\allowbreak{}e\par \textit{Residencia}: multisecciones neutrino & \textit{Celdas seleccionadas}: 20\par \textit{Multisecciones}: neutrino\_\allowbreak{}G1 | neutrino\_\allowbreak{}G2 | neutrino\_\allowbreak{}G3 | neutrino\_\allowbreak{}G4\par \textit{Carga}: 0\par \textit{Representación}: sin color & C05 \\
% VI-CATALOG-ROW clases 5
SM-NU-2\par neutrino & \textit{Representantes}: ν\_\allowbreak{}μ / anti-ν\_\allowbreak{}μ\par \textit{Residencia}: multisecciones neutrino & \textit{Celdas seleccionadas}: 20\par \textit{Multisecciones}: neutrino\_\allowbreak{}G1 | neutrino\_\allowbreak{}G2 | neutrino\_\allowbreak{}G3 | neutrino\_\allowbreak{}G4\par \textit{Carga}: 0\par \textit{Representación}: sin color & C05 \\
% VI-CATALOG-ROW clases 6
SM-NU-3\par neutrino & \textit{Representantes}: ν\_\allowbreak{}τ / anti-ν\_\allowbreak{}τ\par \textit{Residencia}: multisecciones neutrino & \textit{Celdas seleccionadas}: 20\par \textit{Multisecciones}: neutrino\_\allowbreak{}G1 | neutrino\_\allowbreak{}G2 | neutrino\_\allowbreak{}G3 | neutrino\_\allowbreak{}G4\par \textit{Carga}: 0\par \textit{Representación}: sin color & C05 \\
% VI-CATALOG-ROW clases 7
SM-Q-U1\par quark & \textit{Representantes}: u / anti-u\par \textit{Residencia}: multisección tipo arriba & \textit{Celdas seleccionadas}: 16\par \textit{Multisecciones}: quark\_\allowbreak{}up\_\allowbreak{}G1 | quark\_\allowbreak{}up\_\allowbreak{}G3\par \textit{Carga}: ±2/\allowbreak{}3 e\par \textit{Representación}: R,\allowbreak{}G,\allowbreak{}B = órbita Z/\allowbreak{}3 & C04 \\
% VI-CATALOG-ROW clases 8
SM-Q-D1\par quark & \textit{Representantes}: d / anti-d\par \textit{Residencia}: multisección tipo abajo & \textit{Celdas seleccionadas}: 20\par \textit{Multisecciones}: quark\_\allowbreak{}down\_\allowbreak{}G1 | quark\_\allowbreak{}down\_\allowbreak{}G2 | quark\_\allowbreak{}down\_\allowbreak{}G3 | quark\_\allowbreak{}down\_\allowbreak{}G4\par \textit{Carga}: ±1/\allowbreak{}3 e\par \textit{Representación}: R,\allowbreak{}G,\allowbreak{}B = órbita Z/\allowbreak{}3 & C04 \\
% VI-CATALOG-ROW clases 9
SM-Q-U2\par quark & \textit{Representantes}: c / anti-c\par \textit{Residencia}: multisección tipo arriba & \textit{Celdas seleccionadas}: 16\par \textit{Multisecciones}: quark\_\allowbreak{}up\_\allowbreak{}G1 | quark\_\allowbreak{}up\_\allowbreak{}G3\par \textit{Carga}: ±2/\allowbreak{}3 e\par \textit{Representación}: R,\allowbreak{}G,\allowbreak{}B = órbita Z/\allowbreak{}3 & C04 \\
% VI-CATALOG-ROW clases 10
SM-Q-D2\par quark & \textit{Representantes}: s / anti-s\par \textit{Residencia}: multisección tipo abajo & \textit{Celdas seleccionadas}: 20\par \textit{Multisecciones}: quark\_\allowbreak{}down\_\allowbreak{}G1 | quark\_\allowbreak{}down\_\allowbreak{}G2 | quark\_\allowbreak{}down\_\allowbreak{}G3 | quark\_\allowbreak{}down\_\allowbreak{}G4\par \textit{Carga}: ±1/\allowbreak{}3 e\par \textit{Representación}: R,\allowbreak{}G,\allowbreak{}B = órbita Z/\allowbreak{}3 & C04 \\
% VI-CATALOG-ROW clases 11
SM-Q-U3\par quark & \textit{Representantes}: t / anti-t\par \textit{Residencia}: multisección tipo arriba & \textit{Celdas seleccionadas}: 16\par \textit{Multisecciones}: quark\_\allowbreak{}up\_\allowbreak{}G1 | quark\_\allowbreak{}up\_\allowbreak{}G3\par \textit{Carga}: ±2/\allowbreak{}3 e\par \textit{Representación}: R,\allowbreak{}G,\allowbreak{}B = órbita Z/\allowbreak{}3 & C04 \\
% VI-CATALOG-ROW clases 12
SM-Q-D3\par quark & \textit{Representantes}: b / anti-b\par \textit{Residencia}: multisección tipo abajo & \textit{Celdas seleccionadas}: 20\par \textit{Multisecciones}: quark\_\allowbreak{}down\_\allowbreak{}G1 | quark\_\allowbreak{}down\_\allowbreak{}G2 | quark\_\allowbreak{}down\_\allowbreak{}G3 | quark\_\allowbreak{}down\_\allowbreak{}G4\par \textit{Carga}: ±1/\allowbreak{}3 e\par \textit{Representación}: R,\allowbreak{}G,\allowbreak{}B = órbita Z/\allowbreak{}3 & C04 \\
% VI-CATALOG-ROW clases 13
SM-GA-EM\par boson\_\allowbreak{}gauge & \textit{Representantes}: γ\par \textit{Residencia}: sector gauge nulo & \textit{Celdas seleccionadas}: 0\par \textit{Multisecciones}: \textemdash{}\par \textit{Carga}: 0\par \textit{Representación}: sin color & C02 \\
% VI-CATALOG-ROW clases 14
SM-GA-GL\par boson\_\allowbreak{}gauge & \textit{Representantes}: g\par \textit{Residencia}: sector gauge de color & \textit{Celdas seleccionadas}: 0\par \textit{Multisecciones}: \textemdash{}\par \textit{Carga}: 0\par \textit{Representación}: octete gauge & C02 \\
% VI-CATALOG-ROW clases 15
SM-GA-W\par boson\_\allowbreak{}gauge & \textit{Representantes}: W+ / W−\par \textit{Residencia}: sector gauge masivo & \textit{Celdas seleccionadas}: 0\par \textit{Multisecciones}: \textemdash{}\par \textit{Carga}: ±1 e\par \textit{Representación}: sin color & C08 \\
% VI-CATALOG-ROW clases 16
SM-GA-Z\par boson\_\allowbreak{}gauge & \textit{Representantes}: Z0\par \textit{Residencia}: sector gauge masivo & \textit{Celdas seleccionadas}: 0\par \textit{Multisecciones}: \textemdash{}\par \textit{Carga}: 0\par \textit{Representación}: sin color & C08 \\
% VI-CATALOG-ROW clases 17
SM-H-1\par escalar & \textit{Representantes}: H\par \textit{Residencia}: sector escalar & \textit{Celdas seleccionadas}: 0\par \textit{Multisecciones}: \textemdash{}\par \textit{Carga}: 0\par \textit{Representación}: sin color & C08 \\
% VI-CATALOG-ROW clases 18
COMP-MESON\par meson & \textit{Representantes}: π0, π±, K0, K±, η, ρ, φ, J/\allowbreak{}ψ, Υ…\par \textit{Residencia}: producto de sección y antisección & \textit{Celdas seleccionadas}: 0\par \textit{Multisecciones}: \textemdash{}\par \textit{Carga}: dependiente del estado compuesto\par \textit{Representación}: singlete u otras representaciones & C06 \\
% VI-CATALOG-ROW clases 19
COMP-BARYON\par barion & \textit{Representantes}: p, n y colecciones bariónicas\par \textit{Residencia}: producto de tres secciones & \textit{Celdas seleccionadas}: 0\par \textit{Multisecciones}: \textemdash{}\par \textit{Carga}: dependiente del estado compuesto\par \textit{Representación}: singlete físico & C06 \\
% VI-CATALOG-ROW clases 20
TOPO-FIB\par anyon\_\allowbreak{}no\_\allowbreak{}abeliano & \textit{Representantes}: 1, τ\par \textit{Residencia}: sector topológico & \textit{Celdas seleccionadas}: 0\par \textit{Multisecciones}: \textemdash{}\par \textit{Carga}: definida por el codominio topológico\par \textit{Representación}: no aplica & C03 \\
% VI-CATALOG-ROW clases 21
TOPO-ISING\par sector\_\allowbreak{}majorana & \textit{Representantes}: 1, ψ, σ\par \textit{Residencia}: sector topológico / cero & \textit{Celdas seleccionadas}: 0\par \textit{Multisecciones}: \textemdash{}\par \textit{Carga}: dependiente de la realización\par \textit{Representación}: sin color & C03 \\
% VI-CATALOG-ROW clases 22
TOPO-Z6\par anyones\_\allowbreak{}abelianos & \textit{Representantes}: bosón, fermión y cuatro sectores trenzados\par \textit{Residencia}: sector topológico Z/\allowbreak{}6 & \textit{Celdas seleccionadas}: 0\par \textit{Multisecciones}: \textemdash{}\par \textit{Carga}: Z/\allowbreak{}6 definida por el trenzado\par \textit{Representación}: sin color & C03 \\
% VI-CATALOG-ROW clases 23
VIRT-1\par virtual & \textit{Representantes}: secciones virtuales\par \textit{Residencia}: sistema inverso TPK & \textit{Celdas seleccionadas}: 0\par \textit{Multisecciones}: \textemdash{}\par \textit{Carga}: dependiente del sector de prolongación\par \textit{Representación}: según sector & C03 \\
\end{longtable}
\endgroup

\section{Rutas orientadas: guía de lectura y fichas}\label{sec:vi-324-rutas}
Cada ruta ocupa una página y conserva su identidad dentro de las 81 celdas del atlas. Las cuatro orientaciones de cada celda se consultan por separado. La ficha avanza desde la posición y la clasificación hacia los retornos, la memoria, la firma y los dos lectores; las secuencias completas permiten seguir la evolución registrada.

Los retornos de clase, órbita, celda y estado cinemático son magnitudes distintas. También se distinguen la repetición de fase y el cambio de memoria. Los lectores \(K_D\) y \(K_\Omega\) aparecen juntos para facilitar su comparación, sin identificarlos cuando sus registros difieren. Los números decimales conservan íntegramente las cifras de las fuentes.

Una raya indica un campo vacío, no un cero. Las denominaciones de clasificación se expresan en español; el repertorio técnico conserva sus códigos originales. Las huellas de procedencia y los dos ordinales de origen figuran en el apéndice enlazado de la misma ruta, dentro de este volumen. El emparejamiento de ambos registros se realiza por identidad exacta y conserva todos sus campos.

\textbf{Lectura del estado documental.} Las rutas son internas e independientes de una masa externa objetivo. La unión de registros utiliza la clave exacta, sin consultar masas externas. El estado de torre conservado indica que la memoria de la ruta está disponible y que el transductor \(\eta\) requiere el certificado correspondiente. Estas tres indicaciones se aplican a todas las fichas; la persistencia de Pareto, que puede variar, se consigna individualmente.
\clearpage\subsection*{Índice de rutas}\label{sec:vi-indice-rutas}
\addcontentsline{toc}{subsection}{Índice de las 324 rutas}
\begingroup\setlength{\tabcolsep}{5pt}\renewcommand{\arraystretch}{1.04}\begin{longtable}{@{}lr@{\hspace{15mm}}lr@{\hspace{15mm}}lr@{}}
Ruta & Página & Ruta & Página & Ruta & Página\\\toprule
\endfirsthead
Ruta & Página & Ruta & Página & Ruta & Página\\\toprule
\endhead
\hyperref[nav:vi-ruta-1]{r1c1\_\allowbreak{}h-1v-1} & \pageref{nav:vi-ruta-1} & \hyperref[nav:vi-ruta-2]{r1c1\_\allowbreak{}h-1v+1} & \pageref{nav:vi-ruta-2} & \hyperref[nav:vi-ruta-3]{r1c1\_\allowbreak{}h+1v-1} & \pageref{nav:vi-ruta-3}\\
\hyperref[nav:vi-ruta-4]{r1c1\_\allowbreak{}h+1v+1} & \pageref{nav:vi-ruta-4} & \hyperref[nav:vi-ruta-5]{r1c2\_\allowbreak{}h-1v-1} & \pageref{nav:vi-ruta-5} & \hyperref[nav:vi-ruta-6]{r1c2\_\allowbreak{}h-1v+1} & \pageref{nav:vi-ruta-6}\\
\hyperref[nav:vi-ruta-7]{r1c2\_\allowbreak{}h+1v-1} & \pageref{nav:vi-ruta-7} & \hyperref[nav:vi-ruta-8]{r1c2\_\allowbreak{}h+1v+1} & \pageref{nav:vi-ruta-8} & \hyperref[nav:vi-ruta-9]{r1c3\_\allowbreak{}h-1v-1} & \pageref{nav:vi-ruta-9}\\
\hyperref[nav:vi-ruta-10]{r1c3\_\allowbreak{}h-1v+1} & \pageref{nav:vi-ruta-10} & \hyperref[nav:vi-ruta-11]{r1c3\_\allowbreak{}h+1v-1} & \pageref{nav:vi-ruta-11} & \hyperref[nav:vi-ruta-12]{r1c3\_\allowbreak{}h+1v+1} & \pageref{nav:vi-ruta-12}\\
\hyperref[nav:vi-ruta-13]{r1c4\_\allowbreak{}h-1v-1} & \pageref{nav:vi-ruta-13} & \hyperref[nav:vi-ruta-14]{r1c4\_\allowbreak{}h-1v+1} & \pageref{nav:vi-ruta-14} & \hyperref[nav:vi-ruta-15]{r1c4\_\allowbreak{}h+1v-1} & \pageref{nav:vi-ruta-15}\\
\hyperref[nav:vi-ruta-16]{r1c4\_\allowbreak{}h+1v+1} & \pageref{nav:vi-ruta-16} & \hyperref[nav:vi-ruta-17]{r1c5\_\allowbreak{}h-1v-1} & \pageref{nav:vi-ruta-17} & \hyperref[nav:vi-ruta-18]{r1c5\_\allowbreak{}h-1v+1} & \pageref{nav:vi-ruta-18}\\
\hyperref[nav:vi-ruta-19]{r1c5\_\allowbreak{}h+1v-1} & \pageref{nav:vi-ruta-19} & \hyperref[nav:vi-ruta-20]{r1c5\_\allowbreak{}h+1v+1} & \pageref{nav:vi-ruta-20} & \hyperref[nav:vi-ruta-21]{r1c6\_\allowbreak{}h-1v-1} & \pageref{nav:vi-ruta-21}\\
\hyperref[nav:vi-ruta-22]{r1c6\_\allowbreak{}h-1v+1} & \pageref{nav:vi-ruta-22} & \hyperref[nav:vi-ruta-23]{r1c6\_\allowbreak{}h+1v-1} & \pageref{nav:vi-ruta-23} & \hyperref[nav:vi-ruta-24]{r1c6\_\allowbreak{}h+1v+1} & \pageref{nav:vi-ruta-24}\\
\hyperref[nav:vi-ruta-25]{r1c7\_\allowbreak{}h-1v-1} & \pageref{nav:vi-ruta-25} & \hyperref[nav:vi-ruta-26]{r1c7\_\allowbreak{}h-1v+1} & \pageref{nav:vi-ruta-26} & \hyperref[nav:vi-ruta-27]{r1c7\_\allowbreak{}h+1v-1} & \pageref{nav:vi-ruta-27}\\
\hyperref[nav:vi-ruta-28]{r1c7\_\allowbreak{}h+1v+1} & \pageref{nav:vi-ruta-28} & \hyperref[nav:vi-ruta-29]{r1c8\_\allowbreak{}h-1v-1} & \pageref{nav:vi-ruta-29} & \hyperref[nav:vi-ruta-30]{r1c8\_\allowbreak{}h-1v+1} & \pageref{nav:vi-ruta-30}\\
\hyperref[nav:vi-ruta-31]{r1c8\_\allowbreak{}h+1v-1} & \pageref{nav:vi-ruta-31} & \hyperref[nav:vi-ruta-32]{r1c8\_\allowbreak{}h+1v+1} & \pageref{nav:vi-ruta-32} & \hyperref[nav:vi-ruta-33]{r1c9\_\allowbreak{}h-1v-1} & \pageref{nav:vi-ruta-33}\\
\hyperref[nav:vi-ruta-34]{r1c9\_\allowbreak{}h-1v+1} & \pageref{nav:vi-ruta-34} & \hyperref[nav:vi-ruta-35]{r1c9\_\allowbreak{}h+1v-1} & \pageref{nav:vi-ruta-35} & \hyperref[nav:vi-ruta-36]{r1c9\_\allowbreak{}h+1v+1} & \pageref{nav:vi-ruta-36}\\
\hyperref[nav:vi-ruta-37]{r2c1\_\allowbreak{}h-1v-1} & \pageref{nav:vi-ruta-37} & \hyperref[nav:vi-ruta-38]{r2c1\_\allowbreak{}h-1v+1} & \pageref{nav:vi-ruta-38} & \hyperref[nav:vi-ruta-39]{r2c1\_\allowbreak{}h+1v-1} & \pageref{nav:vi-ruta-39}\\
\hyperref[nav:vi-ruta-40]{r2c1\_\allowbreak{}h+1v+1} & \pageref{nav:vi-ruta-40} & \hyperref[nav:vi-ruta-41]{r2c2\_\allowbreak{}h-1v-1} & \pageref{nav:vi-ruta-41} & \hyperref[nav:vi-ruta-42]{r2c2\_\allowbreak{}h-1v+1} & \pageref{nav:vi-ruta-42}\\
\hyperref[nav:vi-ruta-43]{r2c2\_\allowbreak{}h+1v-1} & \pageref{nav:vi-ruta-43} & \hyperref[nav:vi-ruta-44]{r2c2\_\allowbreak{}h+1v+1} & \pageref{nav:vi-ruta-44} & \hyperref[nav:vi-ruta-45]{r2c3\_\allowbreak{}h-1v-1} & \pageref{nav:vi-ruta-45}\\
\hyperref[nav:vi-ruta-46]{r2c3\_\allowbreak{}h-1v+1} & \pageref{nav:vi-ruta-46} & \hyperref[nav:vi-ruta-47]{r2c3\_\allowbreak{}h+1v-1} & \pageref{nav:vi-ruta-47} & \hyperref[nav:vi-ruta-48]{r2c3\_\allowbreak{}h+1v+1} & \pageref{nav:vi-ruta-48}\\
\hyperref[nav:vi-ruta-49]{r2c4\_\allowbreak{}h-1v-1} & \pageref{nav:vi-ruta-49} & \hyperref[nav:vi-ruta-50]{r2c4\_\allowbreak{}h-1v+1} & \pageref{nav:vi-ruta-50} & \hyperref[nav:vi-ruta-51]{r2c4\_\allowbreak{}h+1v-1} & \pageref{nav:vi-ruta-51}\\
\hyperref[nav:vi-ruta-52]{r2c4\_\allowbreak{}h+1v+1} & \pageref{nav:vi-ruta-52} & \hyperref[nav:vi-ruta-53]{r2c5\_\allowbreak{}h-1v-1} & \pageref{nav:vi-ruta-53} & \hyperref[nav:vi-ruta-54]{r2c5\_\allowbreak{}h-1v+1} & \pageref{nav:vi-ruta-54}\\
\hyperref[nav:vi-ruta-55]{r2c5\_\allowbreak{}h+1v-1} & \pageref{nav:vi-ruta-55} & \hyperref[nav:vi-ruta-56]{r2c5\_\allowbreak{}h+1v+1} & \pageref{nav:vi-ruta-56} & \hyperref[nav:vi-ruta-57]{r2c6\_\allowbreak{}h-1v-1} & \pageref{nav:vi-ruta-57}\\
\hyperref[nav:vi-ruta-58]{r2c6\_\allowbreak{}h-1v+1} & \pageref{nav:vi-ruta-58} & \hyperref[nav:vi-ruta-59]{r2c6\_\allowbreak{}h+1v-1} & \pageref{nav:vi-ruta-59} & \hyperref[nav:vi-ruta-60]{r2c6\_\allowbreak{}h+1v+1} & \pageref{nav:vi-ruta-60}\\
\hyperref[nav:vi-ruta-61]{r2c7\_\allowbreak{}h-1v-1} & \pageref{nav:vi-ruta-61} & \hyperref[nav:vi-ruta-62]{r2c7\_\allowbreak{}h-1v+1} & \pageref{nav:vi-ruta-62} & \hyperref[nav:vi-ruta-63]{r2c7\_\allowbreak{}h+1v-1} & \pageref{nav:vi-ruta-63}\\
\hyperref[nav:vi-ruta-64]{r2c7\_\allowbreak{}h+1v+1} & \pageref{nav:vi-ruta-64} & \hyperref[nav:vi-ruta-65]{r2c8\_\allowbreak{}h-1v-1} & \pageref{nav:vi-ruta-65} & \hyperref[nav:vi-ruta-66]{r2c8\_\allowbreak{}h-1v+1} & \pageref{nav:vi-ruta-66}\\
\hyperref[nav:vi-ruta-67]{r2c8\_\allowbreak{}h+1v-1} & \pageref{nav:vi-ruta-67} & \hyperref[nav:vi-ruta-68]{r2c8\_\allowbreak{}h+1v+1} & \pageref{nav:vi-ruta-68} & \hyperref[nav:vi-ruta-69]{r2c9\_\allowbreak{}h-1v-1} & \pageref{nav:vi-ruta-69}\\
\hyperref[nav:vi-ruta-70]{r2c9\_\allowbreak{}h-1v+1} & \pageref{nav:vi-ruta-70} & \hyperref[nav:vi-ruta-71]{r2c9\_\allowbreak{}h+1v-1} & \pageref{nav:vi-ruta-71} & \hyperref[nav:vi-ruta-72]{r2c9\_\allowbreak{}h+1v+1} & \pageref{nav:vi-ruta-72}\\
\hyperref[nav:vi-ruta-73]{r3c1\_\allowbreak{}h-1v-1} & \pageref{nav:vi-ruta-73} & \hyperref[nav:vi-ruta-74]{r3c1\_\allowbreak{}h-1v+1} & \pageref{nav:vi-ruta-74} & \hyperref[nav:vi-ruta-75]{r3c1\_\allowbreak{}h+1v-1} & \pageref{nav:vi-ruta-75}\\
\hyperref[nav:vi-ruta-76]{r3c1\_\allowbreak{}h+1v+1} & \pageref{nav:vi-ruta-76} & \hyperref[nav:vi-ruta-77]{r3c2\_\allowbreak{}h-1v-1} & \pageref{nav:vi-ruta-77} & \hyperref[nav:vi-ruta-78]{r3c2\_\allowbreak{}h-1v+1} & \pageref{nav:vi-ruta-78}\\
\hyperref[nav:vi-ruta-79]{r3c2\_\allowbreak{}h+1v-1} & \pageref{nav:vi-ruta-79} & \hyperref[nav:vi-ruta-80]{r3c2\_\allowbreak{}h+1v+1} & \pageref{nav:vi-ruta-80} & \hyperref[nav:vi-ruta-81]{r3c3\_\allowbreak{}h-1v-1} & \pageref{nav:vi-ruta-81}\\
\hyperref[nav:vi-ruta-82]{r3c3\_\allowbreak{}h-1v+1} & \pageref{nav:vi-ruta-82} & \hyperref[nav:vi-ruta-83]{r3c3\_\allowbreak{}h+1v-1} & \pageref{nav:vi-ruta-83} & \hyperref[nav:vi-ruta-84]{r3c3\_\allowbreak{}h+1v+1} & \pageref{nav:vi-ruta-84}\\
\hyperref[nav:vi-ruta-85]{r3c4\_\allowbreak{}h-1v-1} & \pageref{nav:vi-ruta-85} & \hyperref[nav:vi-ruta-86]{r3c4\_\allowbreak{}h-1v+1} & \pageref{nav:vi-ruta-86} & \hyperref[nav:vi-ruta-87]{r3c4\_\allowbreak{}h+1v-1} & \pageref{nav:vi-ruta-87}\\
\hyperref[nav:vi-ruta-88]{r3c4\_\allowbreak{}h+1v+1} & \pageref{nav:vi-ruta-88} & \hyperref[nav:vi-ruta-89]{r3c5\_\allowbreak{}h-1v-1} & \pageref{nav:vi-ruta-89} & \hyperref[nav:vi-ruta-90]{r3c5\_\allowbreak{}h-1v+1} & \pageref{nav:vi-ruta-90}\\
\hyperref[nav:vi-ruta-91]{r3c5\_\allowbreak{}h+1v-1} & \pageref{nav:vi-ruta-91} & \hyperref[nav:vi-ruta-92]{r3c5\_\allowbreak{}h+1v+1} & \pageref{nav:vi-ruta-92} & \hyperref[nav:vi-ruta-93]{r3c6\_\allowbreak{}h-1v-1} & \pageref{nav:vi-ruta-93}\\
\hyperref[nav:vi-ruta-94]{r3c6\_\allowbreak{}h-1v+1} & \pageref{nav:vi-ruta-94} & \hyperref[nav:vi-ruta-95]{r3c6\_\allowbreak{}h+1v-1} & \pageref{nav:vi-ruta-95} & \hyperref[nav:vi-ruta-96]{r3c6\_\allowbreak{}h+1v+1} & \pageref{nav:vi-ruta-96}\\
\hyperref[nav:vi-ruta-97]{r3c7\_\allowbreak{}h-1v-1} & \pageref{nav:vi-ruta-97} & \hyperref[nav:vi-ruta-98]{r3c7\_\allowbreak{}h-1v+1} & \pageref{nav:vi-ruta-98} & \hyperref[nav:vi-ruta-99]{r3c7\_\allowbreak{}h+1v-1} & \pageref{nav:vi-ruta-99}\\
\hyperref[nav:vi-ruta-100]{r3c7\_\allowbreak{}h+1v+1} & \pageref{nav:vi-ruta-100} & \hyperref[nav:vi-ruta-101]{r3c8\_\allowbreak{}h-1v-1} & \pageref{nav:vi-ruta-101} & \hyperref[nav:vi-ruta-102]{r3c8\_\allowbreak{}h-1v+1} & \pageref{nav:vi-ruta-102}\\
\hyperref[nav:vi-ruta-103]{r3c8\_\allowbreak{}h+1v-1} & \pageref{nav:vi-ruta-103} & \hyperref[nav:vi-ruta-104]{r3c8\_\allowbreak{}h+1v+1} & \pageref{nav:vi-ruta-104} & \hyperref[nav:vi-ruta-105]{r3c9\_\allowbreak{}h-1v-1} & \pageref{nav:vi-ruta-105}\\
\hyperref[nav:vi-ruta-106]{r3c9\_\allowbreak{}h-1v+1} & \pageref{nav:vi-ruta-106} & \hyperref[nav:vi-ruta-107]{r3c9\_\allowbreak{}h+1v-1} & \pageref{nav:vi-ruta-107} & \hyperref[nav:vi-ruta-108]{r3c9\_\allowbreak{}h+1v+1} & \pageref{nav:vi-ruta-108}\\
\hyperref[nav:vi-ruta-109]{r4c1\_\allowbreak{}h-1v-1} & \pageref{nav:vi-ruta-109} & \hyperref[nav:vi-ruta-110]{r4c1\_\allowbreak{}h-1v+1} & \pageref{nav:vi-ruta-110} & \hyperref[nav:vi-ruta-111]{r4c1\_\allowbreak{}h+1v-1} & \pageref{nav:vi-ruta-111}\\
\hyperref[nav:vi-ruta-112]{r4c1\_\allowbreak{}h+1v+1} & \pageref{nav:vi-ruta-112} & \hyperref[nav:vi-ruta-113]{r4c2\_\allowbreak{}h-1v-1} & \pageref{nav:vi-ruta-113} & \hyperref[nav:vi-ruta-114]{r4c2\_\allowbreak{}h-1v+1} & \pageref{nav:vi-ruta-114}\\
\hyperref[nav:vi-ruta-115]{r4c2\_\allowbreak{}h+1v-1} & \pageref{nav:vi-ruta-115} & \hyperref[nav:vi-ruta-116]{r4c2\_\allowbreak{}h+1v+1} & \pageref{nav:vi-ruta-116} & \hyperref[nav:vi-ruta-117]{r4c3\_\allowbreak{}h-1v-1} & \pageref{nav:vi-ruta-117}\\
\hyperref[nav:vi-ruta-118]{r4c3\_\allowbreak{}h-1v+1} & \pageref{nav:vi-ruta-118} & \hyperref[nav:vi-ruta-119]{r4c3\_\allowbreak{}h+1v-1} & \pageref{nav:vi-ruta-119} & \hyperref[nav:vi-ruta-120]{r4c3\_\allowbreak{}h+1v+1} & \pageref{nav:vi-ruta-120}\\
\hyperref[nav:vi-ruta-121]{r4c4\_\allowbreak{}h-1v-1} & \pageref{nav:vi-ruta-121} & \hyperref[nav:vi-ruta-122]{r4c4\_\allowbreak{}h-1v+1} & \pageref{nav:vi-ruta-122} & \hyperref[nav:vi-ruta-123]{r4c4\_\allowbreak{}h+1v-1} & \pageref{nav:vi-ruta-123}\\
\hyperref[nav:vi-ruta-124]{r4c4\_\allowbreak{}h+1v+1} & \pageref{nav:vi-ruta-124} & \hyperref[nav:vi-ruta-125]{r4c5\_\allowbreak{}h-1v-1} & \pageref{nav:vi-ruta-125} & \hyperref[nav:vi-ruta-126]{r4c5\_\allowbreak{}h-1v+1} & \pageref{nav:vi-ruta-126}\\
\hyperref[nav:vi-ruta-127]{r4c5\_\allowbreak{}h+1v-1} & \pageref{nav:vi-ruta-127} & \hyperref[nav:vi-ruta-128]{r4c5\_\allowbreak{}h+1v+1} & \pageref{nav:vi-ruta-128} & \hyperref[nav:vi-ruta-129]{r4c6\_\allowbreak{}h-1v-1} & \pageref{nav:vi-ruta-129}\\
\hyperref[nav:vi-ruta-130]{r4c6\_\allowbreak{}h-1v+1} & \pageref{nav:vi-ruta-130} & \hyperref[nav:vi-ruta-131]{r4c6\_\allowbreak{}h+1v-1} & \pageref{nav:vi-ruta-131} & \hyperref[nav:vi-ruta-132]{r4c6\_\allowbreak{}h+1v+1} & \pageref{nav:vi-ruta-132}\\
\hyperref[nav:vi-ruta-133]{r4c7\_\allowbreak{}h-1v-1} & \pageref{nav:vi-ruta-133} & \hyperref[nav:vi-ruta-134]{r4c7\_\allowbreak{}h-1v+1} & \pageref{nav:vi-ruta-134} & \hyperref[nav:vi-ruta-135]{r4c7\_\allowbreak{}h+1v-1} & \pageref{nav:vi-ruta-135}\\
\hyperref[nav:vi-ruta-136]{r4c7\_\allowbreak{}h+1v+1} & \pageref{nav:vi-ruta-136} & \hyperref[nav:vi-ruta-137]{r4c8\_\allowbreak{}h-1v-1} & \pageref{nav:vi-ruta-137} & \hyperref[nav:vi-ruta-138]{r4c8\_\allowbreak{}h-1v+1} & \pageref{nav:vi-ruta-138}\\
\hyperref[nav:vi-ruta-139]{r4c8\_\allowbreak{}h+1v-1} & \pageref{nav:vi-ruta-139} & \hyperref[nav:vi-ruta-140]{r4c8\_\allowbreak{}h+1v+1} & \pageref{nav:vi-ruta-140} & \hyperref[nav:vi-ruta-141]{r4c9\_\allowbreak{}h-1v-1} & \pageref{nav:vi-ruta-141}\\
\hyperref[nav:vi-ruta-142]{r4c9\_\allowbreak{}h-1v+1} & \pageref{nav:vi-ruta-142} & \hyperref[nav:vi-ruta-143]{r4c9\_\allowbreak{}h+1v-1} & \pageref{nav:vi-ruta-143} & \hyperref[nav:vi-ruta-144]{r4c9\_\allowbreak{}h+1v+1} & \pageref{nav:vi-ruta-144}\\
\hyperref[nav:vi-ruta-145]{r5c1\_\allowbreak{}h-1v-1} & \pageref{nav:vi-ruta-145} & \hyperref[nav:vi-ruta-146]{r5c1\_\allowbreak{}h-1v+1} & \pageref{nav:vi-ruta-146} & \hyperref[nav:vi-ruta-147]{r5c1\_\allowbreak{}h+1v-1} & \pageref{nav:vi-ruta-147}\\
\hyperref[nav:vi-ruta-148]{r5c1\_\allowbreak{}h+1v+1} & \pageref{nav:vi-ruta-148} & \hyperref[nav:vi-ruta-149]{r5c2\_\allowbreak{}h-1v-1} & \pageref{nav:vi-ruta-149} & \hyperref[nav:vi-ruta-150]{r5c2\_\allowbreak{}h-1v+1} & \pageref{nav:vi-ruta-150}\\
\hyperref[nav:vi-ruta-151]{r5c2\_\allowbreak{}h+1v-1} & \pageref{nav:vi-ruta-151} & \hyperref[nav:vi-ruta-152]{r5c2\_\allowbreak{}h+1v+1} & \pageref{nav:vi-ruta-152} & \hyperref[nav:vi-ruta-153]{r5c3\_\allowbreak{}h-1v-1} & \pageref{nav:vi-ruta-153}\\
\hyperref[nav:vi-ruta-154]{r5c3\_\allowbreak{}h-1v+1} & \pageref{nav:vi-ruta-154} & \hyperref[nav:vi-ruta-155]{r5c3\_\allowbreak{}h+1v-1} & \pageref{nav:vi-ruta-155} & \hyperref[nav:vi-ruta-156]{r5c3\_\allowbreak{}h+1v+1} & \pageref{nav:vi-ruta-156}\\
\hyperref[nav:vi-ruta-157]{r5c4\_\allowbreak{}h-1v-1} & \pageref{nav:vi-ruta-157} & \hyperref[nav:vi-ruta-158]{r5c4\_\allowbreak{}h-1v+1} & \pageref{nav:vi-ruta-158} & \hyperref[nav:vi-ruta-159]{r5c4\_\allowbreak{}h+1v-1} & \pageref{nav:vi-ruta-159}\\
\hyperref[nav:vi-ruta-160]{r5c4\_\allowbreak{}h+1v+1} & \pageref{nav:vi-ruta-160} & \hyperref[nav:vi-ruta-161]{r5c5\_\allowbreak{}h-1v-1} & \pageref{nav:vi-ruta-161} & \hyperref[nav:vi-ruta-162]{r5c5\_\allowbreak{}h-1v+1} & \pageref{nav:vi-ruta-162}\\
\hyperref[nav:vi-ruta-163]{r5c5\_\allowbreak{}h+1v-1} & \pageref{nav:vi-ruta-163} & \hyperref[nav:vi-ruta-164]{r5c5\_\allowbreak{}h+1v+1} & \pageref{nav:vi-ruta-164} & \hyperref[nav:vi-ruta-165]{r5c6\_\allowbreak{}h-1v-1} & \pageref{nav:vi-ruta-165}\\
\hyperref[nav:vi-ruta-166]{r5c6\_\allowbreak{}h-1v+1} & \pageref{nav:vi-ruta-166} & \hyperref[nav:vi-ruta-167]{r5c6\_\allowbreak{}h+1v-1} & \pageref{nav:vi-ruta-167} & \hyperref[nav:vi-ruta-168]{r5c6\_\allowbreak{}h+1v+1} & \pageref{nav:vi-ruta-168}\\
\hyperref[nav:vi-ruta-169]{r5c7\_\allowbreak{}h-1v-1} & \pageref{nav:vi-ruta-169} & \hyperref[nav:vi-ruta-170]{r5c7\_\allowbreak{}h-1v+1} & \pageref{nav:vi-ruta-170} & \hyperref[nav:vi-ruta-171]{r5c7\_\allowbreak{}h+1v-1} & \pageref{nav:vi-ruta-171}\\
\hyperref[nav:vi-ruta-172]{r5c7\_\allowbreak{}h+1v+1} & \pageref{nav:vi-ruta-172} & \hyperref[nav:vi-ruta-173]{r5c8\_\allowbreak{}h-1v-1} & \pageref{nav:vi-ruta-173} & \hyperref[nav:vi-ruta-174]{r5c8\_\allowbreak{}h-1v+1} & \pageref{nav:vi-ruta-174}\\
\hyperref[nav:vi-ruta-175]{r5c8\_\allowbreak{}h+1v-1} & \pageref{nav:vi-ruta-175} & \hyperref[nav:vi-ruta-176]{r5c8\_\allowbreak{}h+1v+1} & \pageref{nav:vi-ruta-176} & \hyperref[nav:vi-ruta-177]{r5c9\_\allowbreak{}h-1v-1} & \pageref{nav:vi-ruta-177}\\
\hyperref[nav:vi-ruta-178]{r5c9\_\allowbreak{}h-1v+1} & \pageref{nav:vi-ruta-178} & \hyperref[nav:vi-ruta-179]{r5c9\_\allowbreak{}h+1v-1} & \pageref{nav:vi-ruta-179} & \hyperref[nav:vi-ruta-180]{r5c9\_\allowbreak{}h+1v+1} & \pageref{nav:vi-ruta-180}\\
\hyperref[nav:vi-ruta-181]{r6c1\_\allowbreak{}h-1v-1} & \pageref{nav:vi-ruta-181} & \hyperref[nav:vi-ruta-182]{r6c1\_\allowbreak{}h-1v+1} & \pageref{nav:vi-ruta-182} & \hyperref[nav:vi-ruta-183]{r6c1\_\allowbreak{}h+1v-1} & \pageref{nav:vi-ruta-183}\\
\hyperref[nav:vi-ruta-184]{r6c1\_\allowbreak{}h+1v+1} & \pageref{nav:vi-ruta-184} & \hyperref[nav:vi-ruta-185]{r6c2\_\allowbreak{}h-1v-1} & \pageref{nav:vi-ruta-185} & \hyperref[nav:vi-ruta-186]{r6c2\_\allowbreak{}h-1v+1} & \pageref{nav:vi-ruta-186}\\
\hyperref[nav:vi-ruta-187]{r6c2\_\allowbreak{}h+1v-1} & \pageref{nav:vi-ruta-187} & \hyperref[nav:vi-ruta-188]{r6c2\_\allowbreak{}h+1v+1} & \pageref{nav:vi-ruta-188} & \hyperref[nav:vi-ruta-189]{r6c3\_\allowbreak{}h-1v-1} & \pageref{nav:vi-ruta-189}\\
\hyperref[nav:vi-ruta-190]{r6c3\_\allowbreak{}h-1v+1} & \pageref{nav:vi-ruta-190} & \hyperref[nav:vi-ruta-191]{r6c3\_\allowbreak{}h+1v-1} & \pageref{nav:vi-ruta-191} & \hyperref[nav:vi-ruta-192]{r6c3\_\allowbreak{}h+1v+1} & \pageref{nav:vi-ruta-192}\\
\hyperref[nav:vi-ruta-193]{r6c4\_\allowbreak{}h-1v-1} & \pageref{nav:vi-ruta-193} & \hyperref[nav:vi-ruta-194]{r6c4\_\allowbreak{}h-1v+1} & \pageref{nav:vi-ruta-194} & \hyperref[nav:vi-ruta-195]{r6c4\_\allowbreak{}h+1v-1} & \pageref{nav:vi-ruta-195}\\
\hyperref[nav:vi-ruta-196]{r6c4\_\allowbreak{}h+1v+1} & \pageref{nav:vi-ruta-196} & \hyperref[nav:vi-ruta-197]{r6c5\_\allowbreak{}h-1v-1} & \pageref{nav:vi-ruta-197} & \hyperref[nav:vi-ruta-198]{r6c5\_\allowbreak{}h-1v+1} & \pageref{nav:vi-ruta-198}\\
\hyperref[nav:vi-ruta-199]{r6c5\_\allowbreak{}h+1v-1} & \pageref{nav:vi-ruta-199} & \hyperref[nav:vi-ruta-200]{r6c5\_\allowbreak{}h+1v+1} & \pageref{nav:vi-ruta-200} & \hyperref[nav:vi-ruta-201]{r6c6\_\allowbreak{}h-1v-1} & \pageref{nav:vi-ruta-201}\\
\hyperref[nav:vi-ruta-202]{r6c6\_\allowbreak{}h-1v+1} & \pageref{nav:vi-ruta-202} & \hyperref[nav:vi-ruta-203]{r6c6\_\allowbreak{}h+1v-1} & \pageref{nav:vi-ruta-203} & \hyperref[nav:vi-ruta-204]{r6c6\_\allowbreak{}h+1v+1} & \pageref{nav:vi-ruta-204}\\
\hyperref[nav:vi-ruta-205]{r6c7\_\allowbreak{}h-1v-1} & \pageref{nav:vi-ruta-205} & \hyperref[nav:vi-ruta-206]{r6c7\_\allowbreak{}h-1v+1} & \pageref{nav:vi-ruta-206} & \hyperref[nav:vi-ruta-207]{r6c7\_\allowbreak{}h+1v-1} & \pageref{nav:vi-ruta-207}\\
\hyperref[nav:vi-ruta-208]{r6c7\_\allowbreak{}h+1v+1} & \pageref{nav:vi-ruta-208} & \hyperref[nav:vi-ruta-209]{r6c8\_\allowbreak{}h-1v-1} & \pageref{nav:vi-ruta-209} & \hyperref[nav:vi-ruta-210]{r6c8\_\allowbreak{}h-1v+1} & \pageref{nav:vi-ruta-210}\\
\hyperref[nav:vi-ruta-211]{r6c8\_\allowbreak{}h+1v-1} & \pageref{nav:vi-ruta-211} & \hyperref[nav:vi-ruta-212]{r6c8\_\allowbreak{}h+1v+1} & \pageref{nav:vi-ruta-212} & \hyperref[nav:vi-ruta-213]{r6c9\_\allowbreak{}h-1v-1} & \pageref{nav:vi-ruta-213}\\
\hyperref[nav:vi-ruta-214]{r6c9\_\allowbreak{}h-1v+1} & \pageref{nav:vi-ruta-214} & \hyperref[nav:vi-ruta-215]{r6c9\_\allowbreak{}h+1v-1} & \pageref{nav:vi-ruta-215} & \hyperref[nav:vi-ruta-216]{r6c9\_\allowbreak{}h+1v+1} & \pageref{nav:vi-ruta-216}\\
\hyperref[nav:vi-ruta-217]{r7c1\_\allowbreak{}h-1v-1} & \pageref{nav:vi-ruta-217} & \hyperref[nav:vi-ruta-218]{r7c1\_\allowbreak{}h-1v+1} & \pageref{nav:vi-ruta-218} & \hyperref[nav:vi-ruta-219]{r7c1\_\allowbreak{}h+1v-1} & \pageref{nav:vi-ruta-219}\\
\hyperref[nav:vi-ruta-220]{r7c1\_\allowbreak{}h+1v+1} & \pageref{nav:vi-ruta-220} & \hyperref[nav:vi-ruta-221]{r7c2\_\allowbreak{}h-1v-1} & \pageref{nav:vi-ruta-221} & \hyperref[nav:vi-ruta-222]{r7c2\_\allowbreak{}h-1v+1} & \pageref{nav:vi-ruta-222}\\
\hyperref[nav:vi-ruta-223]{r7c2\_\allowbreak{}h+1v-1} & \pageref{nav:vi-ruta-223} & \hyperref[nav:vi-ruta-224]{r7c2\_\allowbreak{}h+1v+1} & \pageref{nav:vi-ruta-224} & \hyperref[nav:vi-ruta-225]{r7c3\_\allowbreak{}h-1v-1} & \pageref{nav:vi-ruta-225}\\
\hyperref[nav:vi-ruta-226]{r7c3\_\allowbreak{}h-1v+1} & \pageref{nav:vi-ruta-226} & \hyperref[nav:vi-ruta-227]{r7c3\_\allowbreak{}h+1v-1} & \pageref{nav:vi-ruta-227} & \hyperref[nav:vi-ruta-228]{r7c3\_\allowbreak{}h+1v+1} & \pageref{nav:vi-ruta-228}\\
\hyperref[nav:vi-ruta-229]{r7c4\_\allowbreak{}h-1v-1} & \pageref{nav:vi-ruta-229} & \hyperref[nav:vi-ruta-230]{r7c4\_\allowbreak{}h-1v+1} & \pageref{nav:vi-ruta-230} & \hyperref[nav:vi-ruta-231]{r7c4\_\allowbreak{}h+1v-1} & \pageref{nav:vi-ruta-231}\\
\hyperref[nav:vi-ruta-232]{r7c4\_\allowbreak{}h+1v+1} & \pageref{nav:vi-ruta-232} & \hyperref[nav:vi-ruta-233]{r7c5\_\allowbreak{}h-1v-1} & \pageref{nav:vi-ruta-233} & \hyperref[nav:vi-ruta-234]{r7c5\_\allowbreak{}h-1v+1} & \pageref{nav:vi-ruta-234}\\
\hyperref[nav:vi-ruta-235]{r7c5\_\allowbreak{}h+1v-1} & \pageref{nav:vi-ruta-235} & \hyperref[nav:vi-ruta-236]{r7c5\_\allowbreak{}h+1v+1} & \pageref{nav:vi-ruta-236} & \hyperref[nav:vi-ruta-237]{r7c6\_\allowbreak{}h-1v-1} & \pageref{nav:vi-ruta-237}\\
\hyperref[nav:vi-ruta-238]{r7c6\_\allowbreak{}h-1v+1} & \pageref{nav:vi-ruta-238} & \hyperref[nav:vi-ruta-239]{r7c6\_\allowbreak{}h+1v-1} & \pageref{nav:vi-ruta-239} & \hyperref[nav:vi-ruta-240]{r7c6\_\allowbreak{}h+1v+1} & \pageref{nav:vi-ruta-240}\\
\hyperref[nav:vi-ruta-241]{r7c7\_\allowbreak{}h-1v-1} & \pageref{nav:vi-ruta-241} & \hyperref[nav:vi-ruta-242]{r7c7\_\allowbreak{}h-1v+1} & \pageref{nav:vi-ruta-242} & \hyperref[nav:vi-ruta-243]{r7c7\_\allowbreak{}h+1v-1} & \pageref{nav:vi-ruta-243}\\
\hyperref[nav:vi-ruta-244]{r7c7\_\allowbreak{}h+1v+1} & \pageref{nav:vi-ruta-244} & \hyperref[nav:vi-ruta-245]{r7c8\_\allowbreak{}h-1v-1} & \pageref{nav:vi-ruta-245} & \hyperref[nav:vi-ruta-246]{r7c8\_\allowbreak{}h-1v+1} & \pageref{nav:vi-ruta-246}\\
\hyperref[nav:vi-ruta-247]{r7c8\_\allowbreak{}h+1v-1} & \pageref{nav:vi-ruta-247} & \hyperref[nav:vi-ruta-248]{r7c8\_\allowbreak{}h+1v+1} & \pageref{nav:vi-ruta-248} & \hyperref[nav:vi-ruta-249]{r7c9\_\allowbreak{}h-1v-1} & \pageref{nav:vi-ruta-249}\\
\hyperref[nav:vi-ruta-250]{r7c9\_\allowbreak{}h-1v+1} & \pageref{nav:vi-ruta-250} & \hyperref[nav:vi-ruta-251]{r7c9\_\allowbreak{}h+1v-1} & \pageref{nav:vi-ruta-251} & \hyperref[nav:vi-ruta-252]{r7c9\_\allowbreak{}h+1v+1} & \pageref{nav:vi-ruta-252}\\
\hyperref[nav:vi-ruta-253]{r8c1\_\allowbreak{}h-1v-1} & \pageref{nav:vi-ruta-253} & \hyperref[nav:vi-ruta-254]{r8c1\_\allowbreak{}h-1v+1} & \pageref{nav:vi-ruta-254} & \hyperref[nav:vi-ruta-255]{r8c1\_\allowbreak{}h+1v-1} & \pageref{nav:vi-ruta-255}\\
\hyperref[nav:vi-ruta-256]{r8c1\_\allowbreak{}h+1v+1} & \pageref{nav:vi-ruta-256} & \hyperref[nav:vi-ruta-257]{r8c2\_\allowbreak{}h-1v-1} & \pageref{nav:vi-ruta-257} & \hyperref[nav:vi-ruta-258]{r8c2\_\allowbreak{}h-1v+1} & \pageref{nav:vi-ruta-258}\\
\hyperref[nav:vi-ruta-259]{r8c2\_\allowbreak{}h+1v-1} & \pageref{nav:vi-ruta-259} & \hyperref[nav:vi-ruta-260]{r8c2\_\allowbreak{}h+1v+1} & \pageref{nav:vi-ruta-260} & \hyperref[nav:vi-ruta-261]{r8c3\_\allowbreak{}h-1v-1} & \pageref{nav:vi-ruta-261}\\
\hyperref[nav:vi-ruta-262]{r8c3\_\allowbreak{}h-1v+1} & \pageref{nav:vi-ruta-262} & \hyperref[nav:vi-ruta-263]{r8c3\_\allowbreak{}h+1v-1} & \pageref{nav:vi-ruta-263} & \hyperref[nav:vi-ruta-264]{r8c3\_\allowbreak{}h+1v+1} & \pageref{nav:vi-ruta-264}\\
\hyperref[nav:vi-ruta-265]{r8c4\_\allowbreak{}h-1v-1} & \pageref{nav:vi-ruta-265} & \hyperref[nav:vi-ruta-266]{r8c4\_\allowbreak{}h-1v+1} & \pageref{nav:vi-ruta-266} & \hyperref[nav:vi-ruta-267]{r8c4\_\allowbreak{}h+1v-1} & \pageref{nav:vi-ruta-267}\\
\hyperref[nav:vi-ruta-268]{r8c4\_\allowbreak{}h+1v+1} & \pageref{nav:vi-ruta-268} & \hyperref[nav:vi-ruta-269]{r8c5\_\allowbreak{}h-1v-1} & \pageref{nav:vi-ruta-269} & \hyperref[nav:vi-ruta-270]{r8c5\_\allowbreak{}h-1v+1} & \pageref{nav:vi-ruta-270}\\
\hyperref[nav:vi-ruta-271]{r8c5\_\allowbreak{}h+1v-1} & \pageref{nav:vi-ruta-271} & \hyperref[nav:vi-ruta-272]{r8c5\_\allowbreak{}h+1v+1} & \pageref{nav:vi-ruta-272} & \hyperref[nav:vi-ruta-273]{r8c6\_\allowbreak{}h-1v-1} & \pageref{nav:vi-ruta-273}\\
\hyperref[nav:vi-ruta-274]{r8c6\_\allowbreak{}h-1v+1} & \pageref{nav:vi-ruta-274} & \hyperref[nav:vi-ruta-275]{r8c6\_\allowbreak{}h+1v-1} & \pageref{nav:vi-ruta-275} & \hyperref[nav:vi-ruta-276]{r8c6\_\allowbreak{}h+1v+1} & \pageref{nav:vi-ruta-276}\\
\hyperref[nav:vi-ruta-277]{r8c7\_\allowbreak{}h-1v-1} & \pageref{nav:vi-ruta-277} & \hyperref[nav:vi-ruta-278]{r8c7\_\allowbreak{}h-1v+1} & \pageref{nav:vi-ruta-278} & \hyperref[nav:vi-ruta-279]{r8c7\_\allowbreak{}h+1v-1} & \pageref{nav:vi-ruta-279}\\
\hyperref[nav:vi-ruta-280]{r8c7\_\allowbreak{}h+1v+1} & \pageref{nav:vi-ruta-280} & \hyperref[nav:vi-ruta-281]{r8c8\_\allowbreak{}h-1v-1} & \pageref{nav:vi-ruta-281} & \hyperref[nav:vi-ruta-282]{r8c8\_\allowbreak{}h-1v+1} & \pageref{nav:vi-ruta-282}\\
\hyperref[nav:vi-ruta-283]{r8c8\_\allowbreak{}h+1v-1} & \pageref{nav:vi-ruta-283} & \hyperref[nav:vi-ruta-284]{r8c8\_\allowbreak{}h+1v+1} & \pageref{nav:vi-ruta-284} & \hyperref[nav:vi-ruta-285]{r8c9\_\allowbreak{}h-1v-1} & \pageref{nav:vi-ruta-285}\\
\hyperref[nav:vi-ruta-286]{r8c9\_\allowbreak{}h-1v+1} & \pageref{nav:vi-ruta-286} & \hyperref[nav:vi-ruta-287]{r8c9\_\allowbreak{}h+1v-1} & \pageref{nav:vi-ruta-287} & \hyperref[nav:vi-ruta-288]{r8c9\_\allowbreak{}h+1v+1} & \pageref{nav:vi-ruta-288}\\
\hyperref[nav:vi-ruta-289]{r9c1\_\allowbreak{}h-1v-1} & \pageref{nav:vi-ruta-289} & \hyperref[nav:vi-ruta-290]{r9c1\_\allowbreak{}h-1v+1} & \pageref{nav:vi-ruta-290} & \hyperref[nav:vi-ruta-291]{r9c1\_\allowbreak{}h+1v-1} & \pageref{nav:vi-ruta-291}\\
\hyperref[nav:vi-ruta-292]{r9c1\_\allowbreak{}h+1v+1} & \pageref{nav:vi-ruta-292} & \hyperref[nav:vi-ruta-293]{r9c2\_\allowbreak{}h-1v-1} & \pageref{nav:vi-ruta-293} & \hyperref[nav:vi-ruta-294]{r9c2\_\allowbreak{}h-1v+1} & \pageref{nav:vi-ruta-294}\\
\hyperref[nav:vi-ruta-295]{r9c2\_\allowbreak{}h+1v-1} & \pageref{nav:vi-ruta-295} & \hyperref[nav:vi-ruta-296]{r9c2\_\allowbreak{}h+1v+1} & \pageref{nav:vi-ruta-296} & \hyperref[nav:vi-ruta-297]{r9c3\_\allowbreak{}h-1v-1} & \pageref{nav:vi-ruta-297}\\
\hyperref[nav:vi-ruta-298]{r9c3\_\allowbreak{}h-1v+1} & \pageref{nav:vi-ruta-298} & \hyperref[nav:vi-ruta-299]{r9c3\_\allowbreak{}h+1v-1} & \pageref{nav:vi-ruta-299} & \hyperref[nav:vi-ruta-300]{r9c3\_\allowbreak{}h+1v+1} & \pageref{nav:vi-ruta-300}\\
\hyperref[nav:vi-ruta-301]{r9c4\_\allowbreak{}h-1v-1} & \pageref{nav:vi-ruta-301} & \hyperref[nav:vi-ruta-302]{r9c4\_\allowbreak{}h-1v+1} & \pageref{nav:vi-ruta-302} & \hyperref[nav:vi-ruta-303]{r9c4\_\allowbreak{}h+1v-1} & \pageref{nav:vi-ruta-303}\\
\hyperref[nav:vi-ruta-304]{r9c4\_\allowbreak{}h+1v+1} & \pageref{nav:vi-ruta-304} & \hyperref[nav:vi-ruta-305]{r9c5\_\allowbreak{}h-1v-1} & \pageref{nav:vi-ruta-305} & \hyperref[nav:vi-ruta-306]{r9c5\_\allowbreak{}h-1v+1} & \pageref{nav:vi-ruta-306}\\
\hyperref[nav:vi-ruta-307]{r9c5\_\allowbreak{}h+1v-1} & \pageref{nav:vi-ruta-307} & \hyperref[nav:vi-ruta-308]{r9c5\_\allowbreak{}h+1v+1} & \pageref{nav:vi-ruta-308} & \hyperref[nav:vi-ruta-309]{r9c6\_\allowbreak{}h-1v-1} & \pageref{nav:vi-ruta-309}\\
\hyperref[nav:vi-ruta-310]{r9c6\_\allowbreak{}h-1v+1} & \pageref{nav:vi-ruta-310} & \hyperref[nav:vi-ruta-311]{r9c6\_\allowbreak{}h+1v-1} & \pageref{nav:vi-ruta-311} & \hyperref[nav:vi-ruta-312]{r9c6\_\allowbreak{}h+1v+1} & \pageref{nav:vi-ruta-312}\\
\hyperref[nav:vi-ruta-313]{r9c7\_\allowbreak{}h-1v-1} & \pageref{nav:vi-ruta-313} & \hyperref[nav:vi-ruta-314]{r9c7\_\allowbreak{}h-1v+1} & \pageref{nav:vi-ruta-314} & \hyperref[nav:vi-ruta-315]{r9c7\_\allowbreak{}h+1v-1} & \pageref{nav:vi-ruta-315}\\
\hyperref[nav:vi-ruta-316]{r9c7\_\allowbreak{}h+1v+1} & \pageref{nav:vi-ruta-316} & \hyperref[nav:vi-ruta-317]{r9c8\_\allowbreak{}h-1v-1} & \pageref{nav:vi-ruta-317} & \hyperref[nav:vi-ruta-318]{r9c8\_\allowbreak{}h-1v+1} & \pageref{nav:vi-ruta-318}\\
\hyperref[nav:vi-ruta-319]{r9c8\_\allowbreak{}h+1v-1} & \pageref{nav:vi-ruta-319} & \hyperref[nav:vi-ruta-320]{r9c8\_\allowbreak{}h+1v+1} & \pageref{nav:vi-ruta-320} & \hyperref[nav:vi-ruta-321]{r9c9\_\allowbreak{}h-1v-1} & \pageref{nav:vi-ruta-321}\\
\hyperref[nav:vi-ruta-322]{r9c9\_\allowbreak{}h-1v+1} & \pageref{nav:vi-ruta-322} & \hyperref[nav:vi-ruta-323]{r9c9\_\allowbreak{}h+1v-1} & \pageref{nav:vi-ruta-323} & \hyperref[nav:vi-ruta-324]{r9c9\_\allowbreak{}h+1v+1} & \pageref{nav:vi-ruta-324}\\
\end{longtable}\endgroup
\clearpage\phantomsection\label{nav:vi-ruta-1}
% VI-CATALOG-ROW rutas 1
\pdfbookmark[2]{Celda 1}{cell-1}
\pdfbookmark[3]{r1c1\_h-1v-1}{route-1}
\RouteTitle{Ruta 001}{r1c1\_\allowbreak{}h-1v-1}
\noindent Celda 1\quad (1, 1)\qquad Orientación \(h=-1,\ v=-1\)\hfill\hyperref[sec:vi-indice-rutas]{Índice}\par\vspace{4pt}
\begingroup\fontsize{10.2}{12.5}\selectfont\setlength{\tabcolsep}{5pt}\renewcommand{\arraystretch}{1.1}\rowcolors{1}{routepale}{white}\noindent\begin{tabularx}{\linewidth}{@{}>{\raggedright\arraybackslash}p{.345\linewidth}>{\raggedright\arraybackslash}X@{}}
\RouteGroup{Identidad estructural}
\RouteRow{Clase y multisección}{charged\_\allowbreak{}lepton\_\allowbreak{}G4\quad /\quad charged\_\allowbreak{}lepton\_\allowbreak{}G4}
\RouteRow{Colección y sector}{leptón cargado; leptón}
\RouteRow{Clasificación física}{leptón cargado}
\RouteRow{Generación, profundidad y color}{4; G4; \textemdash{}}
\RouteRow{Ruta primaria y órbita de inversión}{\(B_1\): no\qquad 1,\allowbreak{}1;\allowbreak{}9,\allowbreak{}9}
\RouteRow{Firma de clasificación}{24|\allowbreak{}0|\allowbreak{}8|\allowbreak{}0|\allowbreak{}-12|\allowbreak{}+++-\kern0pt{}-\kern0pt{}-+++-\kern0pt{}-\kern0pt{}-}
\RouteGroup{Retornos, memoria y transporte}
\RouteRow{Retornos con fase}{clase: 27; órbita: 27; celda: 27}
\RouteRow{Retorno cinemático}{en 108: 54; extensión: 216}
\RouteRow{Giros y cambios de hoja}{71; 107}
\RouteRow{Acarreo y diferimiento}{deuda total: 2; final: 1; bloques diferidos: 2}
\RouteRow{Tríadas finales; persistencia de Pareto}{16; no}
\RouteRow{\(K\longrightarrow K+9\)}{repeticiones: 9; cambios de memoria: 9}
\RouteGroup{Firma, fase y diferencias}
\RouteRow{\(n_A,\ s_C,\ \kappa_0,\ \nu_{120},\ \nu_{270}\)}{24, 0, 0, 0, -12}
\RouteRow{\(\tau_{12}\); firma histórica \(\mathbb Z^5\)}{+++-\kern0pt{}-\kern0pt{}-+++-\kern0pt{}-\kern0pt{}-; 24|\allowbreak{}0|\allowbreak{}0|\allowbreak{}0|\allowbreak{}-12}
\RouteRow{\(D_1,\ D_3,\ D_4,\ D_{27},\ D_{36}\)}{54, 56, 70, 48, 32}
\RouteRow{\(\Omega_{90/120},\ P_6\)}{80, 6}
\RouteGroup{Lectores y factores de acción}
\RouteRow{\(K_D\quad /\quad K_\Omega\)}{(80,\allowbreak{}54,\allowbreak{}7)\quad /\quad (80,\allowbreak{}54,\allowbreak{}6)}
\RouteRow{\(\kappa_D\)}{79.60672\allowbreak{}13362175\allowbreak{}19726524\allowbreak{}96}
\RouteRow{\(\kappa_\Omega\)}{79.60661\allowbreak{}01397948\allowbreak{}27586470\allowbreak{}91}
\RouteRow{\(R_{\mathrm{act},D}\)}{1.000000\allowbreak{}05518988\allowbreak{}96223153\allowbreak{}37}
\RouteRow{\(R_{\mathrm{act},\Omega}\)}{1.000000\allowbreak{}05518981\allowbreak{}25318591\allowbreak{}40}
\RouteGroup{Separación de prefijos}
\RouteRow{Área de ambigüedad; área \(\log_2\)}{58; 30.18947\allowbreak{}50100961\allowbreak{}74}
\RouteRow{Primer bloque de separación}{órbita: \textemdash{}; celda: \textemdash{}; ruta: \textemdash{}}
\RouteRow{Ambigüedad final}{3}
\RouteGroup{Secuencias completas}
\RouteRow{\(w_6\)}{486 423 15 486 429 243 657 24 243 486 423 15 486 429 243 657 24 243}
\RouteRow{Tríadas estables por bloque}{0 1 2 3 4 5 6 7 8 9 10 11 12 13 13 15 16 16}
\end{tabularx}\endgroup\par\vspace{5pt}\noindent{\fontsize{8.4}{10}\selectfont\hyperref[trace-1]{Procedencia y códigos originales}\hfill Ficha completa · 001/324}
\clearpage\phantomsection\label{nav:vi-ruta-2}
% VI-CATALOG-ROW rutas 2
\pdfbookmark[3]{r1c1\_h-1v+1}{route-2}
\RouteTitle{Ruta 002}{r1c1\_\allowbreak{}h-1v+1}
\noindent Celda 1\quad (1, 1)\qquad Orientación \(h=-1,\ v=1\)\hfill\hyperref[sec:vi-indice-rutas]{Índice}\par\vspace{4pt}
\begingroup\fontsize{10.2}{12.5}\selectfont\setlength{\tabcolsep}{5pt}\renewcommand{\arraystretch}{1.1}\rowcolors{1}{routepale}{white}\noindent\begin{tabularx}{\linewidth}{@{}>{\raggedright\arraybackslash}p{.345\linewidth}>{\raggedright\arraybackslash}X@{}}
\RouteGroup{Identidad estructural}
\RouteRow{Clase y multisección}{charged\_\allowbreak{}lepton\_\allowbreak{}G4\quad /\quad charged\_\allowbreak{}lepton\_\allowbreak{}G4}
\RouteRow{Colección y sector}{leptón cargado; leptón}
\RouteRow{Clasificación física}{leptón cargado}
\RouteRow{Generación, profundidad y color}{4; G4; \textemdash{}}
\RouteRow{Ruta primaria y órbita de inversión}{\(B_1\): sí\qquad 1,\allowbreak{}1;\allowbreak{}9,\allowbreak{}9}
\RouteRow{Firma de clasificación}{24|\allowbreak{}0|\allowbreak{}8|\allowbreak{}0|\allowbreak{}-12|\allowbreak{}+++-\kern0pt{}-\kern0pt{}-+++-\kern0pt{}-\kern0pt{}-}
\RouteGroup{Retornos, memoria y transporte}
\RouteRow{Retornos con fase}{clase: 27; órbita: 27; celda: 27}
\RouteRow{Retorno cinemático}{en 108: 54; extensión: 216}
\RouteRow{Giros y cambios de hoja}{71; 107}
\RouteRow{Acarreo y diferimiento}{deuda total: 0; final: 0; bloques diferidos: 0}
\RouteRow{Tríadas finales; persistencia de Pareto}{17; no}
\RouteRow{\(K\longrightarrow K+9\)}{repeticiones: 9; cambios de memoria: 9}
\RouteGroup{Firma, fase y diferencias}
\RouteRow{\(n_A,\ s_C,\ \kappa_0,\ \nu_{120},\ \nu_{270}\)}{14, -4, 0, 0, -12}
\RouteRow{\(\tau_{12}\); firma histórica \(\mathbb Z^5\)}{+++-\kern0pt{}-\kern0pt{}-+++-\kern0pt{}-\kern0pt{}-; 14|\allowbreak{}-4|\allowbreak{}0|\allowbreak{}0|\allowbreak{}-12}
\RouteRow{\(D_1,\ D_3,\ D_4,\ D_{27},\ D_{36}\)}{78, 28, 78, 36, 24}
\RouteRow{\(\Omega_{90/120},\ P_6\)}{80, 6}
\RouteGroup{Lectores y factores de acción}
\RouteRow{\(K_D\quad /\quad K_\Omega\)}{(60,\allowbreak{}78,\allowbreak{}25)\quad /\quad (80,\allowbreak{}54,\allowbreak{}6)}
\RouteRow{\(\kappa_D\)}{59.43358\allowbreak{}64101631\allowbreak{}67023563\allowbreak{}04}
\RouteRow{\(\kappa_\Omega\)}{79.60661\allowbreak{}01397948\allowbreak{}27586470\allowbreak{}91}
\RouteRow{\(R_{\mathrm{act},D}\)}{1.000000\allowbreak{}04120422\allowbreak{}24053449\allowbreak{}05}
\RouteRow{\(R_{\mathrm{act},\Omega}\)}{1.000000\allowbreak{}05518981\allowbreak{}25318591\allowbreak{}40}
\RouteGroup{Separación de prefijos}
\RouteRow{Área de ambigüedad; área \(\log_2\)}{54; 28.52932\allowbreak{}50129808}
\RouteRow{Primer bloque de separación}{órbita: \textemdash{}; celda: \textemdash{}; ruta: \textemdash{}}
\RouteRow{Ambigüedad final}{3}
\RouteGroup{Secuencias completas}
\RouteRow{\(w_6\)}{507 504 585 507 510 510 666 672 510 507 504 585 507 510 510 666 672 510}
\RouteRow{Tríadas estables por bloque}{0 1 2 3 4 5 6 7 8 9 10 11 12 13 14 15 16 17}
\end{tabularx}\endgroup\par\vspace{5pt}\noindent{\fontsize{8.4}{10}\selectfont\hyperref[trace-2]{Procedencia y códigos originales}\hfill Ficha completa · 002/324}
\clearpage\phantomsection\label{nav:vi-ruta-3}
% VI-CATALOG-ROW rutas 3
\pdfbookmark[3]{r1c1\_h+1v-1}{route-3}
\RouteTitle{Ruta 003}{r1c1\_\allowbreak{}h+1v-1}
\noindent Celda 1\quad (1, 1)\qquad Orientación \(h=1,\ v=-1\)\hfill\hyperref[sec:vi-indice-rutas]{Índice}\par\vspace{4pt}
\begingroup\fontsize{10.2}{12.5}\selectfont\setlength{\tabcolsep}{5pt}\renewcommand{\arraystretch}{1.1}\rowcolors{1}{routepale}{white}\noindent\begin{tabularx}{\linewidth}{@{}>{\raggedright\arraybackslash}p{.345\linewidth}>{\raggedright\arraybackslash}X@{}}
\RouteGroup{Identidad estructural}
\RouteRow{Clase y multisección}{charged\_\allowbreak{}lepton\_\allowbreak{}G4\quad /\quad charged\_\allowbreak{}lepton\_\allowbreak{}G4}
\RouteRow{Colección y sector}{leptón cargado; leptón}
\RouteRow{Clasificación física}{leptón cargado}
\RouteRow{Generación, profundidad y color}{4; G4; \textemdash{}}
\RouteRow{Ruta primaria y órbita de inversión}{\(B_1\): no\qquad 1,\allowbreak{}1;\allowbreak{}9,\allowbreak{}9}
\RouteRow{Firma de clasificación}{24|\allowbreak{}0|\allowbreak{}8|\allowbreak{}0|\allowbreak{}-12|\allowbreak{}+++-\kern0pt{}-\kern0pt{}-+++-\kern0pt{}-\kern0pt{}-}
\RouteGroup{Retornos, memoria y transporte}
\RouteRow{Retornos con fase}{clase: 27; órbita: 27; celda: 27}
\RouteRow{Retorno cinemático}{en 108: 54; extensión: 216}
\RouteRow{Giros y cambios de hoja}{71; 107}
\RouteRow{Acarreo y diferimiento}{deuda total: 2; final: 0; bloques diferidos: 1}
\RouteRow{Tríadas finales; persistencia de Pareto}{17; no}
\RouteRow{\(K\longrightarrow K+9\)}{repeticiones: 9; cambios de memoria: 9}
\RouteGroup{Firma, fase y diferencias}
\RouteRow{\(n_A,\ s_C,\ \kappa_0,\ \nu_{120},\ \nu_{270}\)}{14, 4, 0, 0, -12}
\RouteRow{\(\tau_{12}\); firma histórica \(\mathbb Z^5\)}{+++-\kern0pt{}-\kern0pt{}-+++-\kern0pt{}-\kern0pt{}-; 14|\allowbreak{}4|\allowbreak{}0|\allowbreak{}0|\allowbreak{}-12}
\RouteRow{\(D_1,\ D_3,\ D_4,\ D_{27},\ D_{36}\)}{78, 28, 78, 36, 24}
\RouteRow{\(\Omega_{90/120},\ P_6\)}{80, 6}
\RouteGroup{Lectores y factores de acción}
\RouteRow{\(K_D\quad /\quad K_\Omega\)}{(60,\allowbreak{}78,\allowbreak{}25)\quad /\quad (80,\allowbreak{}54,\allowbreak{}6)}
\RouteRow{\(\kappa_D\)}{59.43358\allowbreak{}64101631\allowbreak{}67023563\allowbreak{}04}
\RouteRow{\(\kappa_\Omega\)}{79.60661\allowbreak{}01397948\allowbreak{}27586470\allowbreak{}91}
\RouteRow{\(R_{\mathrm{act},D}\)}{1.000000\allowbreak{}04120422\allowbreak{}24053449\allowbreak{}05}
\RouteRow{\(R_{\mathrm{act},\Omega}\)}{1.000000\allowbreak{}05518981\allowbreak{}25318591\allowbreak{}40}
\RouteGroup{Separación de prefijos}
\RouteRow{Área de ambigüedad; área \(\log_2\)}{54; 28.52932\allowbreak{}50129808}
\RouteRow{Primer bloque de separación}{órbita: \textemdash{}; celda: \textemdash{}; ruta: \textemdash{}}
\RouteRow{Ambigüedad final}{3}
\RouteGroup{Secuencias completas}
\RouteRow{\(w_6\)}{666 672 510 666 666 585 507 504 585 666 672 510 666 666 585 507 504 585}
\RouteRow{Tríadas estables por bloque}{0 1 2 3 4 5 6 7 8 9 10 11 12 12 13 15 16 17}
\end{tabularx}\endgroup\par\vspace{5pt}\noindent{\fontsize{8.4}{10}\selectfont\hyperref[trace-3]{Procedencia y códigos originales}\hfill Ficha completa · 003/324}
\clearpage\phantomsection\label{nav:vi-ruta-4}
% VI-CATALOG-ROW rutas 4
\pdfbookmark[3]{r1c1\_h+1v+1}{route-4}
\RouteTitle{Ruta 004}{r1c1\_\allowbreak{}h+1v+1}
\noindent Celda 1\quad (1, 1)\qquad Orientación \(h=1,\ v=1\)\hfill\hyperref[sec:vi-indice-rutas]{Índice}\par\vspace{4pt}
\begingroup\fontsize{10.2}{12.5}\selectfont\setlength{\tabcolsep}{5pt}\renewcommand{\arraystretch}{1.1}\rowcolors{1}{routepale}{white}\noindent\begin{tabularx}{\linewidth}{@{}>{\raggedright\arraybackslash}p{.345\linewidth}>{\raggedright\arraybackslash}X@{}}
\RouteGroup{Identidad estructural}
\RouteRow{Clase y multisección}{charged\_\allowbreak{}lepton\_\allowbreak{}G4\quad /\quad charged\_\allowbreak{}lepton\_\allowbreak{}G4}
\RouteRow{Colección y sector}{leptón cargado; leptón}
\RouteRow{Clasificación física}{leptón cargado}
\RouteRow{Generación, profundidad y color}{4; G4; \textemdash{}}
\RouteRow{Ruta primaria y órbita de inversión}{\(B_1\): no\qquad 1,\allowbreak{}1;\allowbreak{}9,\allowbreak{}9}
\RouteRow{Firma de clasificación}{24|\allowbreak{}0|\allowbreak{}8|\allowbreak{}0|\allowbreak{}-12|\allowbreak{}+++-\kern0pt{}-\kern0pt{}-+++-\kern0pt{}-\kern0pt{}-}
\RouteGroup{Retornos, memoria y transporte}
\RouteRow{Retornos con fase}{clase: 27; órbita: 27; celda: 27}
\RouteRow{Retorno cinemático}{en 108: 54; extensión: 216}
\RouteRow{Giros y cambios de hoja}{71; 107}
\RouteRow{Acarreo y diferimiento}{deuda total: 2; final: 0; bloques diferidos: 1}
\RouteRow{Tríadas finales; persistencia de Pareto}{17; no}
\RouteRow{\(K\longrightarrow K+9\)}{repeticiones: 9; cambios de memoria: 9}
\RouteGroup{Firma, fase y diferencias}
\RouteRow{\(n_A,\ s_C,\ \kappa_0,\ \nu_{120},\ \nu_{270}\)}{24, 0, 0, 0, -12}
\RouteRow{\(\tau_{12}\); firma histórica \(\mathbb Z^5\)}{+++-\kern0pt{}-\kern0pt{}-+++-\kern0pt{}-\kern0pt{}-; 24|\allowbreak{}0|\allowbreak{}0|\allowbreak{}0|\allowbreak{}-12}
\RouteRow{\(D_1,\ D_3,\ D_4,\ D_{27},\ D_{36}\)}{54, 56, 70, 48, 32}
\RouteRow{\(\Omega_{90/120},\ P_6\)}{80, 6}
\RouteGroup{Lectores y factores de acción}
\RouteRow{\(K_D\quad /\quad K_\Omega\)}{(80,\allowbreak{}54,\allowbreak{}7)\quad /\quad (80,\allowbreak{}54,\allowbreak{}6)}
\RouteRow{\(\kappa_D\)}{79.60672\allowbreak{}13362175\allowbreak{}19726524\allowbreak{}96}
\RouteRow{\(\kappa_\Omega\)}{79.60661\allowbreak{}01397948\allowbreak{}27586470\allowbreak{}91}
\RouteRow{\(R_{\mathrm{act},D}\)}{1.000000\allowbreak{}05518988\allowbreak{}96223153\allowbreak{}37}
\RouteRow{\(R_{\mathrm{act},\Omega}\)}{1.000000\allowbreak{}05518981\allowbreak{}25318591\allowbreak{}40}
\RouteGroup{Separación de prefijos}
\RouteRow{Área de ambigüedad; área \(\log_2\)}{58; 30.18947\allowbreak{}50100961\allowbreak{}74}
\RouteRow{Primer bloque de separación}{órbita: \textemdash{}; celda: \textemdash{}; ruta: \textemdash{}}
\RouteRow{Ambigüedad final}{3}
\RouteGroup{Secuencias completas}
\RouteRow{\(w_6\)}{657 24 243 657 18 15 486 423 15 657 24 243 657 18 15 486 423 15}
\RouteRow{Tríadas estables por bloque}{0 1 2 3 4 5 6 7 8 9 10 11 12 13 14 14 15 17}
\end{tabularx}\endgroup\par\vspace{5pt}\noindent{\fontsize{8.4}{10}\selectfont\hyperref[trace-4]{Procedencia y códigos originales}\hfill Ficha completa · 004/324}
\clearpage\phantomsection\label{nav:vi-ruta-5}
% VI-CATALOG-ROW rutas 5
\pdfbookmark[2]{Celda 2}{cell-2}
\pdfbookmark[3]{r1c2\_h-1v-1}{route-5}
\RouteTitle{Ruta 005}{r1c2\_\allowbreak{}h-1v-1}
\noindent Celda 2\quad (1, 2)\qquad Orientación \(h=-1,\ v=-1\)\hfill\hyperref[sec:vi-indice-rutas]{Índice}\par\vspace{4pt}
\begingroup\fontsize{10.2}{12.5}\selectfont\setlength{\tabcolsep}{5pt}\renewcommand{\arraystretch}{1.1}\rowcolors{1}{routepale}{white}\noindent\begin{tabularx}{\linewidth}{@{}>{\raggedright\arraybackslash}p{.345\linewidth}>{\raggedright\arraybackslash}X@{}}
\RouteGroup{Identidad estructural}
\RouteRow{Clase y multisección}{neutrino\_\allowbreak{}G4\quad /\quad neutrino\_\allowbreak{}G4}
\RouteRow{Colección y sector}{neutrino; leptón}
\RouteRow{Clasificación física}{neutrino}
\RouteRow{Generación, profundidad y color}{4; G4; \textemdash{}}
\RouteRow{Ruta primaria y órbita de inversión}{\(B_1\): sí\qquad 1,\allowbreak{}2;\allowbreak{}9,\allowbreak{}8}
\RouteRow{Firma de clasificación}{28|\allowbreak{}6|\allowbreak{}-4|\allowbreak{}0|\allowbreak{}-12|\allowbreak{}+++-\kern0pt{}-\kern0pt{}-+++-\kern0pt{}-\kern0pt{}-}
\RouteGroup{Retornos, memoria y transporte}
\RouteRow{Retornos con fase}{clase: 27; órbita: 27; celda: 27}
\RouteRow{Retorno cinemático}{en 108: 54; extensión: 216}
\RouteRow{Giros y cambios de hoja}{71; 107}
\RouteRow{Acarreo y diferimiento}{deuda total: 1; final: 0; bloques diferidos: 1}
\RouteRow{Tríadas finales; persistencia de Pareto}{17; sí}
\RouteRow{\(K\longrightarrow K+9\)}{repeticiones: 9; cambios de memoria: 9}
\RouteGroup{Firma, fase y diferencias}
\RouteRow{\(n_A,\ s_C,\ \kappa_0,\ \nu_{120},\ \nu_{270}\)}{28, -6, 0, -4, -12}
\RouteRow{\(\tau_{12}\); firma histórica \(\mathbb Z^5\)}{+++-\kern0pt{}-\kern0pt{}-+++-\kern0pt{}-\kern0pt{}-; 28|\allowbreak{}-6|\allowbreak{}0|\allowbreak{}-4|\allowbreak{}-12}
\RouteRow{\(D_1,\ D_3,\ D_4,\ D_{27},\ D_{36}\)}{66, 64, 64, 60, 40}
\RouteRow{\(\Omega_{90/120},\ P_6\)}{80, 6}
\RouteGroup{Lectores y factores de acción}
\RouteRow{\(K_D\quad /\quad K_\Omega\)}{(100,\allowbreak{}66,\allowbreak{}0)\quad /\quad (80,\allowbreak{}54,\allowbreak{}6)}
\RouteRow{\(\kappa_D\)}{99.51837\allowbreak{}47304272\allowbreak{}69134179\allowbreak{}18}
\RouteRow{\(\kappa_\Omega\)}{79.60661\allowbreak{}01397948\allowbreak{}27586470\allowbreak{}91}
\RouteRow{\(R_{\mathrm{act},D}\)}{1.000000\allowbreak{}06899427\allowbreak{}66450218\allowbreak{}95}
\RouteRow{\(R_{\mathrm{act},\Omega}\)}{1.000000\allowbreak{}05518981\allowbreak{}25318591\allowbreak{}40}
\RouteGroup{Separación de prefijos}
\RouteRow{Área de ambigüedad; área \(\log_2\)}{36; 18.0}
\RouteRow{Primer bloque de separación}{órbita: \textemdash{}; celda: \textemdash{}; ruta: \textemdash{}}
\RouteRow{Ambigüedad final}{2}
\RouteGroup{Secuencias completas}
\RouteRow{\(w_6\)}{90 570 264 90 567 657 24 243 657 90 570 264 90 567 657 24 243 657}
\RouteRow{Tríadas estables por bloque}{0 1 2 3 4 5 6 7 8 9 10 11 12 13 13 15 16 17}
\end{tabularx}\endgroup\par\vspace{5pt}\noindent{\fontsize{8.4}{10}\selectfont\hyperref[trace-5]{Procedencia y códigos originales}\hfill Ficha completa · 005/324}
\clearpage\phantomsection\label{nav:vi-ruta-6}
% VI-CATALOG-ROW rutas 6
\pdfbookmark[3]{r1c2\_h-1v+1}{route-6}
\RouteTitle{Ruta 006}{r1c2\_\allowbreak{}h-1v+1}
\noindent Celda 2\quad (1, 2)\qquad Orientación \(h=-1,\ v=1\)\hfill\hyperref[sec:vi-indice-rutas]{Índice}\par\vspace{4pt}
\begingroup\fontsize{10.2}{12.5}\selectfont\setlength{\tabcolsep}{5pt}\renewcommand{\arraystretch}{1.1}\rowcolors{1}{routepale}{white}\noindent\begin{tabularx}{\linewidth}{@{}>{\raggedright\arraybackslash}p{.345\linewidth}>{\raggedright\arraybackslash}X@{}}
\RouteGroup{Identidad estructural}
\RouteRow{Clase y multisección}{neutrino\_\allowbreak{}G4\quad /\quad neutrino\_\allowbreak{}G4}
\RouteRow{Colección y sector}{neutrino; leptón}
\RouteRow{Clasificación física}{neutrino}
\RouteRow{Generación, profundidad y color}{4; G4; \textemdash{}}
\RouteRow{Ruta primaria y órbita de inversión}{\(B_1\): no\qquad 1,\allowbreak{}2;\allowbreak{}9,\allowbreak{}8}
\RouteRow{Firma de clasificación}{28|\allowbreak{}6|\allowbreak{}-4|\allowbreak{}0|\allowbreak{}-12|\allowbreak{}+++-\kern0pt{}-\kern0pt{}-+++-\kern0pt{}-\kern0pt{}-}
\RouteGroup{Retornos, memoria y transporte}
\RouteRow{Retornos con fase}{clase: 27; órbita: 27; celda: 27}
\RouteRow{Retorno cinemático}{en 108: 54; extensión: 216}
\RouteRow{Giros y cambios de hoja}{71; 107}
\RouteRow{Acarreo y diferimiento}{deuda total: 2; final: 0; bloques diferidos: 2}
\RouteRow{Tríadas finales; persistencia de Pareto}{17; sí}
\RouteRow{\(K\longrightarrow K+9\)}{repeticiones: 9; cambios de memoria: 9}
\RouteGroup{Firma, fase y diferencias}
\RouteRow{\(n_A,\ s_C,\ \kappa_0,\ \nu_{120},\ \nu_{270}\)}{50, -6, 2, -4, -8}
\RouteRow{\(\tau_{12}\); firma histórica \(\mathbb Z^5\)}{+0+-\kern0pt{}-\kern0pt{}-+0+-\kern0pt{}-\kern0pt{}-; 50|\allowbreak{}-6|\allowbreak{}2|\allowbreak{}-4|\allowbreak{}-8}
\RouteRow{\(D_1,\ D_3,\ D_4,\ D_{27},\ D_{36}\)}{24, 20, 24, 24, 16}
\RouteRow{\(\Omega_{90/120},\ P_6\)}{48, 5}
\RouteGroup{Lectores y factores de acción}
\RouteRow{\(K_D\quad /\quad K_\Omega\)}{(40,\allowbreak{}24,\allowbreak{}2)\quad /\quad (48,\allowbreak{}54,\allowbreak{}5)}
\RouteRow{\(\kappa_D\)}{39.82508\allowbreak{}59311825\allowbreak{}73056173\allowbreak{}26}
\RouteRow{\(\kappa_\Omega\)}{47.60649\allowbreak{}89433721\allowbreak{}35446416\allowbreak{}86}
\RouteRow{\(R_{\mathrm{act},D}\)}{1.000000\allowbreak{}02761000\allowbreak{}61595142\allowbreak{}15}
\RouteRow{\(R_{\mathrm{act},\Omega}\)}{1.000000\allowbreak{}03300471\allowbreak{}80532430\allowbreak{}84}
\RouteGroup{Separación de prefijos}
\RouteRow{Área de ambigüedad; área \(\log_2\)}{36; 18.0}
\RouteRow{Primer bloque de separación}{órbita: \textemdash{}; celda: \textemdash{}; ruta: \textemdash{}}
\RouteRow{Ambigüedad final}{2}
\RouteGroup{Secuencias completas}
\RouteRow{\(w_6\)}{81 3 0 81 0 3 0 81 3 81 3 0 81 0 3 0 81 3}
\RouteRow{Tríadas estables por bloque}{0 1 2 3 4 5 6 6 8 9 10 11 12 13 14 15 15 17}
\end{tabularx}\endgroup\par\vspace{5pt}\noindent{\fontsize{8.4}{10}\selectfont\hyperref[trace-6]{Procedencia y códigos originales}\hfill Ficha completa · 006/324}
\clearpage\phantomsection\label{nav:vi-ruta-7}
% VI-CATALOG-ROW rutas 7
\pdfbookmark[3]{r1c2\_h+1v-1}{route-7}
\RouteTitle{Ruta 007}{r1c2\_\allowbreak{}h+1v-1}
\noindent Celda 2\quad (1, 2)\qquad Orientación \(h=1,\ v=-1\)\hfill\hyperref[sec:vi-indice-rutas]{Índice}\par\vspace{4pt}
\begingroup\fontsize{10.2}{12.5}\selectfont\setlength{\tabcolsep}{5pt}\renewcommand{\arraystretch}{1.1}\rowcolors{1}{routepale}{white}\noindent\begin{tabularx}{\linewidth}{@{}>{\raggedright\arraybackslash}p{.345\linewidth}>{\raggedright\arraybackslash}X@{}}
\RouteGroup{Identidad estructural}
\RouteRow{Clase y multisección}{neutrino\_\allowbreak{}G4\quad /\quad neutrino\_\allowbreak{}G4}
\RouteRow{Colección y sector}{neutrino; leptón}
\RouteRow{Clasificación física}{neutrino}
\RouteRow{Generación, profundidad y color}{4; G4; \textemdash{}}
\RouteRow{Ruta primaria y órbita de inversión}{\(B_1\): no\qquad 1,\allowbreak{}2;\allowbreak{}9,\allowbreak{}8}
\RouteRow{Firma de clasificación}{28|\allowbreak{}6|\allowbreak{}-4|\allowbreak{}0|\allowbreak{}-12|\allowbreak{}+++-\kern0pt{}-\kern0pt{}-+++-\kern0pt{}-\kern0pt{}-}
\RouteGroup{Retornos, memoria y transporte}
\RouteRow{Retornos con fase}{clase: 27; órbita: 27; celda: 27}
\RouteRow{Retorno cinemático}{en 108: 54; extensión: 216}
\RouteRow{Giros y cambios de hoja}{71; 107}
\RouteRow{Acarreo y diferimiento}{deuda total: 0; final: 0; bloques diferidos: 0}
\RouteRow{Tríadas finales; persistencia de Pareto}{17; sí}
\RouteRow{\(K\longrightarrow K+9\)}{repeticiones: 9; cambios de memoria: 9}
\RouteGroup{Firma, fase y diferencias}
\RouteRow{\(n_A,\ s_C,\ \kappa_0,\ \nu_{120},\ \nu_{270}\)}{50, 6, -2, 4, -8}
\RouteRow{\(\tau_{12}\); firma histórica \(\mathbb Z^5\)}{+++-0-+++-0-; 50|\allowbreak{}6|\allowbreak{}-2|\allowbreak{}4|\allowbreak{}-8}
\RouteRow{\(D_1,\ D_3,\ D_4,\ D_{27},\ D_{36}\)}{24, 20, 24, 24, 16}
\RouteRow{\(\Omega_{90/120},\ P_6\)}{48, 5}
\RouteGroup{Lectores y factores de acción}
\RouteRow{\(K_D\quad /\quad K_\Omega\)}{(40,\allowbreak{}24,\allowbreak{}2)\quad /\quad (48,\allowbreak{}54,\allowbreak{}5)}
\RouteRow{\(\kappa_D\)}{39.82508\allowbreak{}59311825\allowbreak{}73056173\allowbreak{}26}
\RouteRow{\(\kappa_\Omega\)}{47.60649\allowbreak{}89433721\allowbreak{}35446416\allowbreak{}86}
\RouteRow{\(R_{\mathrm{act},D}\)}{1.000000\allowbreak{}02761000\allowbreak{}61595142\allowbreak{}15}
\RouteRow{\(R_{\mathrm{act},\Omega}\)}{1.000000\allowbreak{}03300471\allowbreak{}80532430\allowbreak{}84}
\RouteGroup{Separación de prefijos}
\RouteRow{Área de ambigüedad; área \(\log_2\)}{36; 18.0}
\RouteRow{Primer bloque de separación}{órbita: \textemdash{}; celda: \textemdash{}; ruta: \textemdash{}}
\RouteRow{Ambigüedad final}{2}
\RouteGroup{Secuencias completas}
\RouteRow{\(w_6\)}{0 81 3 0 84 0 81 3 0 0 81 3 0 84 0 81 3 0}
\RouteRow{Tríadas estables por bloque}{0 1 2 3 4 5 6 7 8 9 10 11 12 13 14 15 16 17}
\end{tabularx}\endgroup\par\vspace{5pt}\noindent{\fontsize{8.4}{10}\selectfont\hyperref[trace-7]{Procedencia y códigos originales}\hfill Ficha completa · 007/324}
\clearpage\phantomsection\label{nav:vi-ruta-8}
% VI-CATALOG-ROW rutas 8
\pdfbookmark[3]{r1c2\_h+1v+1}{route-8}
\RouteTitle{Ruta 008}{r1c2\_\allowbreak{}h+1v+1}
\noindent Celda 2\quad (1, 2)\qquad Orientación \(h=1,\ v=1\)\hfill\hyperref[sec:vi-indice-rutas]{Índice}\par\vspace{4pt}
\begingroup\fontsize{10.2}{12.5}\selectfont\setlength{\tabcolsep}{5pt}\renewcommand{\arraystretch}{1.1}\rowcolors{1}{routepale}{white}\noindent\begin{tabularx}{\linewidth}{@{}>{\raggedright\arraybackslash}p{.345\linewidth}>{\raggedright\arraybackslash}X@{}}
\RouteGroup{Identidad estructural}
\RouteRow{Clase y multisección}{neutrino\_\allowbreak{}G4\quad /\quad neutrino\_\allowbreak{}G4}
\RouteRow{Colección y sector}{neutrino; leptón}
\RouteRow{Clasificación física}{neutrino}
\RouteRow{Generación, profundidad y color}{4; G4; \textemdash{}}
\RouteRow{Ruta primaria y órbita de inversión}{\(B_1\): no\qquad 1,\allowbreak{}2;\allowbreak{}9,\allowbreak{}8}
\RouteRow{Firma de clasificación}{28|\allowbreak{}6|\allowbreak{}-4|\allowbreak{}0|\allowbreak{}-12|\allowbreak{}+++-\kern0pt{}-\kern0pt{}-+++-\kern0pt{}-\kern0pt{}-}
\RouteGroup{Retornos, memoria y transporte}
\RouteRow{Retornos con fase}{clase: 27; órbita: 27; celda: 27}
\RouteRow{Retorno cinemático}{en 108: 54; extensión: 216}
\RouteRow{Giros y cambios de hoja}{71; 107}
\RouteRow{Acarreo y diferimiento}{deuda total: 0; final: 0; bloques diferidos: 0}
\RouteRow{Tríadas finales; persistencia de Pareto}{17; sí}
\RouteRow{\(K\longrightarrow K+9\)}{repeticiones: 9; cambios de memoria: 9}
\RouteGroup{Firma, fase y diferencias}
\RouteRow{\(n_A,\ s_C,\ \kappa_0,\ \nu_{120},\ \nu_{270}\)}{28, 6, 0, 4, -12}
\RouteRow{\(\tau_{12}\); firma histórica \(\mathbb Z^5\)}{+++-\kern0pt{}-\kern0pt{}-+++-\kern0pt{}-\kern0pt{}-; 28|\allowbreak{}6|\allowbreak{}0|\allowbreak{}4|\allowbreak{}-12}
\RouteRow{\(D_1,\ D_3,\ D_4,\ D_{27},\ D_{36}\)}{66, 64, 64, 60, 40}
\RouteRow{\(\Omega_{90/120},\ P_6\)}{80, 6}
\RouteGroup{Lectores y factores de acción}
\RouteRow{\(K_D\quad /\quad K_\Omega\)}{(100,\allowbreak{}66,\allowbreak{}0)\quad /\quad (80,\allowbreak{}54,\allowbreak{}6)}
\RouteRow{\(\kappa_D\)}{99.51837\allowbreak{}47304272\allowbreak{}69134179\allowbreak{}18}
\RouteRow{\(\kappa_\Omega\)}{79.60661\allowbreak{}01397948\allowbreak{}27586470\allowbreak{}91}
\RouteRow{\(R_{\mathrm{act},D}\)}{1.000000\allowbreak{}06899427\allowbreak{}66450218\allowbreak{}95}
\RouteRow{\(R_{\mathrm{act},\Omega}\)}{1.000000\allowbreak{}05518981\allowbreak{}25318591\allowbreak{}40}
\RouteGroup{Separación de prefijos}
\RouteRow{Área de ambigüedad; área \(\log_2\)}{40; 20.33985\allowbreak{}00028846\allowbreak{}24}
\RouteRow{Primer bloque de separación}{órbita: \textemdash{}; celda: \textemdash{}; ruta: \textemdash{}}
\RouteRow{Ambigüedad final}{2}
\RouteGroup{Secuencias completas}
\RouteRow{\(w_6\)}{24 243 657 24 246 264 90 570 264 24 243 657 24 246 264 90 570 264}
\RouteRow{Tríadas estables por bloque}{0 1 2 3 4 5 6 7 8 9 10 11 12 13 14 15 16 17}
\end{tabularx}\endgroup\par\vspace{5pt}\noindent{\fontsize{8.4}{10}\selectfont\hyperref[trace-8]{Procedencia y códigos originales}\hfill Ficha completa · 008/324}
\clearpage\phantomsection\label{nav:vi-ruta-9}
% VI-CATALOG-ROW rutas 9
\pdfbookmark[2]{Celda 3}{cell-3}
\pdfbookmark[3]{r1c3\_h-1v-1}{route-9}
\RouteTitle{Ruta 009}{r1c3\_\allowbreak{}h-1v-1}
\noindent Celda 3\quad (1, 3)\qquad Orientación \(h=-1,\ v=-1\)\hfill\hyperref[sec:vi-indice-rutas]{Índice}\par\vspace{4pt}
\begingroup\fontsize{10.2}{12.5}\selectfont\setlength{\tabcolsep}{5pt}\renewcommand{\arraystretch}{1.1}\rowcolors{1}{routepale}{white}\noindent\begin{tabularx}{\linewidth}{@{}>{\raggedright\arraybackslash}p{.345\linewidth}>{\raggedright\arraybackslash}X@{}}
\RouteGroup{Identidad estructural}
\RouteRow{Clase y multisección}{charged\_\allowbreak{}lepton\_\allowbreak{}G4\quad /\quad charged\_\allowbreak{}lepton\_\allowbreak{}G4}
\RouteRow{Colección y sector}{leptón cargado; leptón}
\RouteRow{Clasificación física}{leptón cargado}
\RouteRow{Generación, profundidad y color}{4; G4; \textemdash{}}
\RouteRow{Ruta primaria y órbita de inversión}{\(B_1\): no\qquad 1,\allowbreak{}3;\allowbreak{}9,\allowbreak{}7}
\RouteRow{Firma de clasificación}{28|\allowbreak{}-6|\allowbreak{}2|\allowbreak{}0|\allowbreak{}-12|\allowbreak{}+++-\kern0pt{}-\kern0pt{}-+++-\kern0pt{}-\kern0pt{}-}
\RouteGroup{Retornos, memoria y transporte}
\RouteRow{Retornos con fase}{clase: 9; órbita: 27; celda: 27}
\RouteRow{Retorno cinemático}{en 108: 54; extensión: 216}
\RouteRow{Giros y cambios de hoja}{71; 107}
\RouteRow{Acarreo y diferimiento}{deuda total: 1; final: 1; bloques diferidos: 1}
\RouteRow{Tríadas finales; persistencia de Pareto}{16; no}
\RouteRow{\(K\longrightarrow K+9\)}{repeticiones: 9; cambios de memoria: 9}
\RouteGroup{Firma, fase y diferencias}
\RouteRow{\(n_A,\ s_C,\ \kappa_0,\ \nu_{120},\ \nu_{270}\)}{28, 6, 0, 0, -12}
\RouteRow{\(\tau_{12}\); firma histórica \(\mathbb Z^5\)}{+++-\kern0pt{}-\kern0pt{}-+++-\kern0pt{}-\kern0pt{}-; 28|\allowbreak{}6|\allowbreak{}0|\allowbreak{}0|\allowbreak{}-12}
\RouteRow{\(D_1,\ D_3,\ D_4,\ D_{27},\ D_{36}\)}{66, 60, 64, 48, 32}
\RouteRow{\(\Omega_{90/120},\ P_6\)}{80, 6}
\RouteGroup{Lectores y factores de acción}
\RouteRow{\(K_D\quad /\quad K_\Omega\)}{(80,\allowbreak{}66,\allowbreak{}2)\quad /\quad (80,\allowbreak{}54,\allowbreak{}6)}
\RouteRow{\(\kappa_D\)}{79.51859\allowbreak{}71232726\allowbreak{}53414287\allowbreak{}29}
\RouteRow{\(\kappa_\Omega\)}{79.60661\allowbreak{}01397948\allowbreak{}27586470\allowbreak{}91}
\RouteRow{\(R_{\mathrm{act},D}\)}{1.000000\allowbreak{}05512879\allowbreak{}47092488\allowbreak{}80}
\RouteRow{\(R_{\mathrm{act},\Omega}\)}{1.000000\allowbreak{}05518981\allowbreak{}25318591\allowbreak{}40}
\RouteGroup{Separación de prefijos}
\RouteRow{Área de ambigüedad; área \(\log_2\)}{54; 28.52932\allowbreak{}50129808}
\RouteRow{Primer bloque de separación}{órbita: \textemdash{}; celda: \textemdash{}; ruta: \textemdash{}}
\RouteRow{Ambigüedad final}{3}
\RouteGroup{Secuencias completas}
\RouteRow{\(w_6\)}{423 15 486 423 12 99 327 498 99 423 15 486 423 12 99 327 498 99}
\RouteRow{Tríadas estables por bloque}{0 1 2 3 4 5 6 7 8 9 10 11 12 13 14 15 16 16}
\end{tabularx}\endgroup\par\vspace{5pt}\noindent{\fontsize{8.4}{10}\selectfont\hyperref[trace-9]{Procedencia y códigos originales}\hfill Ficha completa · 009/324}
\clearpage\phantomsection\label{nav:vi-ruta-10}
% VI-CATALOG-ROW rutas 10
\pdfbookmark[3]{r1c3\_h-1v+1}{route-10}
\RouteTitle{Ruta 010}{r1c3\_\allowbreak{}h-1v+1}
\noindent Celda 3\quad (1, 3)\qquad Orientación \(h=-1,\ v=1\)\hfill\hyperref[sec:vi-indice-rutas]{Índice}\par\vspace{4pt}
\begingroup\fontsize{10.2}{12.5}\selectfont\setlength{\tabcolsep}{5pt}\renewcommand{\arraystretch}{1.1}\rowcolors{1}{routepale}{white}\noindent\begin{tabularx}{\linewidth}{@{}>{\raggedright\arraybackslash}p{.345\linewidth}>{\raggedright\arraybackslash}X@{}}
\RouteGroup{Identidad estructural}
\RouteRow{Clase y multisección}{charged\_\allowbreak{}lepton\_\allowbreak{}G4\quad /\quad charged\_\allowbreak{}lepton\_\allowbreak{}G4}
\RouteRow{Colección y sector}{leptón cargado; leptón}
\RouteRow{Clasificación física}{leptón cargado}
\RouteRow{Generación, profundidad y color}{4; G4; \textemdash{}}
\RouteRow{Ruta primaria y órbita de inversión}{\(B_1\): sí\qquad 1,\allowbreak{}3;\allowbreak{}9,\allowbreak{}7}
\RouteRow{Firma de clasificación}{28|\allowbreak{}-6|\allowbreak{}2|\allowbreak{}0|\allowbreak{}-12|\allowbreak{}+++-\kern0pt{}-\kern0pt{}-+++-\kern0pt{}-\kern0pt{}-}
\RouteGroup{Retornos, memoria y transporte}
\RouteRow{Retornos con fase}{clase: 27; órbita: 27; celda: 27}
\RouteRow{Retorno cinemático}{en 108: 54; extensión: 216}
\RouteRow{Giros y cambios de hoja}{71; 107}
\RouteRow{Acarreo y diferimiento}{deuda total: 1; final: 0; bloques diferidos: 1}
\RouteRow{Tríadas finales; persistencia de Pareto}{17; no}
\RouteRow{\(K\longrightarrow K+9\)}{repeticiones: 9; cambios de memoria: 9}
\RouteGroup{Firma, fase y diferencias}
\RouteRow{\(n_A,\ s_C,\ \kappa_0,\ \nu_{120},\ \nu_{270}\)}{8, -2, 0, 0, -12}
\RouteRow{\(\tau_{12}\); firma histórica \(\mathbb Z^5\)}{+++-\kern0pt{}-\kern0pt{}-+++-\kern0pt{}-\kern0pt{}-; 8|\allowbreak{}-2|\allowbreak{}0|\allowbreak{}0|\allowbreak{}-12}
\RouteRow{\(D_1,\ D_3,\ D_4,\ D_{27},\ D_{36}\)}{84, 24, 84, 24, 16}
\RouteRow{\(\Omega_{90/120},\ P_6\)}{80, 6}
\RouteGroup{Lectores y factores de acción}
\RouteRow{\(K_D\quad /\quad K_\Omega\)}{(40,\allowbreak{}84,\allowbreak{}30)\quad /\quad (80,\allowbreak{}54,\allowbreak{}6)}
\RouteRow{\(\kappa_D\)}{39.39035\allowbreak{}82768609\allowbreak{}24917849\allowbreak{}59}
\RouteRow{\(\kappa_\Omega\)}{79.60661\allowbreak{}01397948\allowbreak{}27586470\allowbreak{}91}
\RouteRow{\(R_{\mathrm{act},D}\)}{1.000000\allowbreak{}02730861\allowbreak{}73967087\allowbreak{}05}
\RouteRow{\(R_{\mathrm{act},\Omega}\)}{1.000000\allowbreak{}05518981\allowbreak{}25318591\allowbreak{}40}
\RouteGroup{Separación de prefijos}
\RouteRow{Área de ambigüedad; área \(\log_2\)}{54; 28.52932\allowbreak{}50129808}
\RouteRow{Primer bloque de separación}{órbita: \textemdash{}; celda: \textemdash{}; ruta: \textemdash{}}
\RouteRow{Ambigüedad final}{3}
\RouteGroup{Secuencias completas}
\RouteRow{\(w_6\)}{420 258 414 420 255 252 333 255 252 420 258 414 420 255 252 333 255 252}
\RouteRow{Tríadas estables por bloque}{0 1 2 3 4 5 6 7 8 9 10 11 12 13 14 15 15 17}
\end{tabularx}\endgroup\par\vspace{5pt}\noindent{\fontsize{8.4}{10}\selectfont\hyperref[trace-10]{Procedencia y códigos originales}\hfill Ficha completa · 010/324}
\clearpage\phantomsection\label{nav:vi-ruta-11}
% VI-CATALOG-ROW rutas 11
\pdfbookmark[3]{r1c3\_h+1v-1}{route-11}
\RouteTitle{Ruta 011}{r1c3\_\allowbreak{}h+1v-1}
\noindent Celda 3\quad (1, 3)\qquad Orientación \(h=1,\ v=-1\)\hfill\hyperref[sec:vi-indice-rutas]{Índice}\par\vspace{4pt}
\begingroup\fontsize{10.2}{12.5}\selectfont\setlength{\tabcolsep}{5pt}\renewcommand{\arraystretch}{1.1}\rowcolors{1}{routepale}{white}\noindent\begin{tabularx}{\linewidth}{@{}>{\raggedright\arraybackslash}p{.345\linewidth}>{\raggedright\arraybackslash}X@{}}
\RouteGroup{Identidad estructural}
\RouteRow{Clase y multisección}{charged\_\allowbreak{}lepton\_\allowbreak{}G4\quad /\quad charged\_\allowbreak{}lepton\_\allowbreak{}G4}
\RouteRow{Colección y sector}{leptón cargado; leptón}
\RouteRow{Clasificación física}{leptón cargado}
\RouteRow{Generación, profundidad y color}{4; G4; \textemdash{}}
\RouteRow{Ruta primaria y órbita de inversión}{\(B_1\): no\qquad 1,\allowbreak{}3;\allowbreak{}9,\allowbreak{}7}
\RouteRow{Firma de clasificación}{28|\allowbreak{}-6|\allowbreak{}2|\allowbreak{}0|\allowbreak{}-12|\allowbreak{}+++-\kern0pt{}-\kern0pt{}-+++-\kern0pt{}-\kern0pt{}-}
\RouteGroup{Retornos, memoria y transporte}
\RouteRow{Retornos con fase}{clase: 27; órbita: 27; celda: 27}
\RouteRow{Retorno cinemático}{en 108: 54; extensión: 216}
\RouteRow{Giros y cambios de hoja}{71; 107}
\RouteRow{Acarreo y diferimiento}{deuda total: 0; final: 0; bloques diferidos: 0}
\RouteRow{Tríadas finales; persistencia de Pareto}{17; no}
\RouteRow{\(K\longrightarrow K+9\)}{repeticiones: 9; cambios de memoria: 9}
\RouteGroup{Firma, fase y diferencias}
\RouteRow{\(n_A,\ s_C,\ \kappa_0,\ \nu_{120},\ \nu_{270}\)}{8, 2, 0, 0, -12}
\RouteRow{\(\tau_{12}\); firma histórica \(\mathbb Z^5\)}{+++-\kern0pt{}-\kern0pt{}-+++-\kern0pt{}-\kern0pt{}-; 8|\allowbreak{}2|\allowbreak{}0|\allowbreak{}0|\allowbreak{}-12}
\RouteRow{\(D_1,\ D_3,\ D_4,\ D_{27},\ D_{36}\)}{84, 24, 84, 24, 16}
\RouteRow{\(\Omega_{90/120},\ P_6\)}{80, 6}
\RouteGroup{Lectores y factores de acción}
\RouteRow{\(K_D\quad /\quad K_\Omega\)}{(40,\allowbreak{}84,\allowbreak{}30)\quad /\quad (80,\allowbreak{}54,\allowbreak{}6)}
\RouteRow{\(\kappa_D\)}{39.39035\allowbreak{}82768609\allowbreak{}24917849\allowbreak{}59}
\RouteRow{\(\kappa_\Omega\)}{79.60661\allowbreak{}01397948\allowbreak{}27586470\allowbreak{}91}
\RouteRow{\(R_{\mathrm{act},D}\)}{1.000000\allowbreak{}02730861\allowbreak{}73967087\allowbreak{}05}
\RouteRow{\(R_{\mathrm{act},\Omega}\)}{1.000000\allowbreak{}05518981\allowbreak{}25318591\allowbreak{}40}
\RouteGroup{Separación de prefijos}
\RouteRow{Área de ambigüedad; área \(\log_2\)}{54; 28.52932\allowbreak{}50129808}
\RouteRow{Primer bloque de separación}{órbita: \textemdash{}; celda: \textemdash{}; ruta: \textemdash{}}
\RouteRow{Ambigüedad final}{3}
\RouteGroup{Secuencias completas}
\RouteRow{\(w_6\)}{333 255 252 333 258 414 420 258 414 333 255 252 333 258 414 420 258 414}
\RouteRow{Tríadas estables por bloque}{0 1 2 3 4 5 6 7 8 9 10 11 12 13 14 15 16 17}
\end{tabularx}\endgroup\par\vspace{5pt}\noindent{\fontsize{8.4}{10}\selectfont\hyperref[trace-11]{Procedencia y códigos originales}\hfill Ficha completa · 011/324}
\clearpage\phantomsection\label{nav:vi-ruta-12}
% VI-CATALOG-ROW rutas 12
\pdfbookmark[3]{r1c3\_h+1v+1}{route-12}
\RouteTitle{Ruta 012}{r1c3\_\allowbreak{}h+1v+1}
\noindent Celda 3\quad (1, 3)\qquad Orientación \(h=1,\ v=1\)\hfill\hyperref[sec:vi-indice-rutas]{Índice}\par\vspace{4pt}
\begingroup\fontsize{10.2}{12.5}\selectfont\setlength{\tabcolsep}{5pt}\renewcommand{\arraystretch}{1.1}\rowcolors{1}{routepale}{white}\noindent\begin{tabularx}{\linewidth}{@{}>{\raggedright\arraybackslash}p{.345\linewidth}>{\raggedright\arraybackslash}X@{}}
\RouteGroup{Identidad estructural}
\RouteRow{Clase y multisección}{charged\_\allowbreak{}lepton\_\allowbreak{}G4\quad /\quad charged\_\allowbreak{}lepton\_\allowbreak{}G4}
\RouteRow{Colección y sector}{leptón cargado; leptón}
\RouteRow{Clasificación física}{leptón cargado}
\RouteRow{Generación, profundidad y color}{4; G4; \textemdash{}}
\RouteRow{Ruta primaria y órbita de inversión}{\(B_1\): no\qquad 1,\allowbreak{}3;\allowbreak{}9,\allowbreak{}7}
\RouteRow{Firma de clasificación}{28|\allowbreak{}-6|\allowbreak{}2|\allowbreak{}0|\allowbreak{}-12|\allowbreak{}+++-\kern0pt{}-\kern0pt{}-+++-\kern0pt{}-\kern0pt{}-}
\RouteGroup{Retornos, memoria y transporte}
\RouteRow{Retornos con fase}{clase: 18; órbita: 27; celda: 27}
\RouteRow{Retorno cinemático}{en 108: 54; extensión: 216}
\RouteRow{Giros y cambios de hoja}{71; 107}
\RouteRow{Acarreo y diferimiento}{deuda total: 2; final: 1; bloques diferidos: 2}
\RouteRow{Tríadas finales; persistencia de Pareto}{16; no}
\RouteRow{\(K\longrightarrow K+9\)}{repeticiones: 9; cambios de memoria: 9}
\RouteGroup{Firma, fase y diferencias}
\RouteRow{\(n_A,\ s_C,\ \kappa_0,\ \nu_{120},\ \nu_{270}\)}{28, -6, 0, 0, -12}
\RouteRow{\(\tau_{12}\); firma histórica \(\mathbb Z^5\)}{+++-\kern0pt{}-\kern0pt{}-+++-\kern0pt{}-\kern0pt{}-; 28|\allowbreak{}-6|\allowbreak{}0|\allowbreak{}0|\allowbreak{}-12}
\RouteRow{\(D_1,\ D_3,\ D_4,\ D_{27},\ D_{36}\)}{66, 60, 64, 48, 32}
\RouteRow{\(\Omega_{90/120},\ P_6\)}{80, 6}
\RouteGroup{Lectores y factores de acción}
\RouteRow{\(K_D\quad /\quad K_\Omega\)}{(80,\allowbreak{}66,\allowbreak{}2)\quad /\quad (80,\allowbreak{}54,\allowbreak{}6)}
\RouteRow{\(\kappa_D\)}{79.51859\allowbreak{}71232726\allowbreak{}53414287\allowbreak{}29}
\RouteRow{\(\kappa_\Omega\)}{79.60661\allowbreak{}01397948\allowbreak{}27586470\allowbreak{}91}
\RouteRow{\(R_{\mathrm{act},D}\)}{1.000000\allowbreak{}05512879\allowbreak{}47092488\allowbreak{}80}
\RouteRow{\(R_{\mathrm{act},\Omega}\)}{1.000000\allowbreak{}05518981\allowbreak{}25318591\allowbreak{}40}
\RouteGroup{Separación de prefijos}
\RouteRow{Área de ambigüedad; área \(\log_2\)}{54; 28.52932\allowbreak{}50129808}
\RouteRow{Primer bloque de separación}{órbita: \textemdash{}; celda: \textemdash{}; ruta: \textemdash{}}
\RouteRow{Ambigüedad final}{3}
\RouteGroup{Secuencias completas}
\RouteRow{\(w_6\)}{327 498 99 327 501 486 423 15 486 327 498 99 327 501 486 423 15 486}
\RouteRow{Tríadas estables por bloque}{0 1 2 3 4 5 6 7 7 9 10 11 12 13 14 15 16 16}
\end{tabularx}\endgroup\par\vspace{5pt}\noindent{\fontsize{8.4}{10}\selectfont\hyperref[trace-12]{Procedencia y códigos originales}\hfill Ficha completa · 012/324}
\clearpage\phantomsection\label{nav:vi-ruta-13}
% VI-CATALOG-ROW rutas 13
\pdfbookmark[2]{Celda 4}{cell-4}
\pdfbookmark[3]{r1c4\_h-1v-1}{route-13}
\RouteTitle{Ruta 013}{r1c4\_\allowbreak{}h-1v-1}
\noindent Celda 4\quad (1, 4)\qquad Orientación \(h=-1,\ v=-1\)\hfill\hyperref[sec:vi-indice-rutas]{Índice}\par\vspace{4pt}
\begingroup\fontsize{10.2}{12.5}\selectfont\setlength{\tabcolsep}{5pt}\renewcommand{\arraystretch}{1.1}\rowcolors{1}{routepale}{white}\noindent\begin{tabularx}{\linewidth}{@{}>{\raggedright\arraybackslash}p{.345\linewidth}>{\raggedright\arraybackslash}X@{}}
\RouteGroup{Identidad estructural}
\RouteRow{Clase y multisección}{neutrino\_\allowbreak{}G4\quad /\quad neutrino\_\allowbreak{}G4}
\RouteRow{Colección y sector}{neutrino; leptón}
\RouteRow{Clasificación física}{neutrino}
\RouteRow{Generación, profundidad y color}{4; G4; \textemdash{}}
\RouteRow{Ruta primaria y órbita de inversión}{\(B_1\): no\qquad 1,\allowbreak{}4;\allowbreak{}9,\allowbreak{}6}
\RouteRow{Firma de clasificación}{30|\allowbreak{}-6|\allowbreak{}-4|\allowbreak{}-2|\allowbreak{}-8|\allowbreak{}0++-\kern0pt{}-\kern0pt{}-0++-\kern0pt{}-\kern0pt{}-}
\RouteGroup{Retornos, memoria y transporte}
\RouteRow{Retornos con fase}{clase: 27; órbita: 27; celda: 27}
\RouteRow{Retorno cinemático}{en 108: 54; extensión: 216}
\RouteRow{Giros y cambios de hoja}{71; 107}
\RouteRow{Acarreo y diferimiento}{deuda total: 2; final: 1; bloques diferidos: 2}
\RouteRow{Tríadas finales; persistencia de Pareto}{16; sí}
\RouteRow{\(K\longrightarrow K+9\)}{repeticiones: 9; cambios de memoria: 9}
\RouteGroup{Firma, fase y diferencias}
\RouteRow{\(n_A,\ s_C,\ \kappa_0,\ \nu_{120},\ \nu_{270}\)}{30, 6, -2, 4, -8}
\RouteRow{\(\tau_{12}\); firma histórica \(\mathbb Z^5\)}{+++0-\kern0pt{}-+++0-\kern0pt{}-; 30|\allowbreak{}6|\allowbreak{}-2|\allowbreak{}4|\allowbreak{}-8}
\RouteRow{\(D_1,\ D_3,\ D_4,\ D_{27},\ D_{36}\)}{54, 56, 70, 48, 32}
\RouteRow{\(\Omega_{90/120},\ P_6\)}{56, 5}
\RouteGroup{Lectores y factores de acción}
\RouteRow{\(K_D\quad /\quad K_\Omega\)}{(80,\allowbreak{}54,\allowbreak{}7)\quad /\quad (56,\allowbreak{}54,\allowbreak{}5)}
\RouteRow{\(\kappa_D\)}{79.60672\allowbreak{}13362175\allowbreak{}19726524\allowbreak{}96}
\RouteRow{\(\kappa_\Omega\)}{55.60649\allowbreak{}89433721\allowbreak{}35446416\allowbreak{}86}
\RouteRow{\(R_{\mathrm{act},D}\)}{1.000000\allowbreak{}05518988\allowbreak{}96223153\allowbreak{}37}
\RouteRow{\(R_{\mathrm{act},\Omega}\)}{1.000000\allowbreak{}03855097\allowbreak{}23541416\allowbreak{}45}
\RouteGroup{Separación de prefijos}
\RouteRow{Área de ambigüedad; área \(\log_2\)}{26; 6.339850\allowbreak{}00288462\allowbreak{}4}
\RouteRow{Primer bloque de separación}{órbita: 5; celda: 5; ruta: 5}
\RouteRow{Ambigüedad final}{1}
\RouteGroup{Secuencias completas}
\RouteRow{\(w_6\)}{486 423 15 486 429 243 657 24 243 486 423 15 486 429 243 657 24 243}
\RouteRow{Tríadas estables por bloque}{0 1 2 3 4 5 6 7 8 9 10 11 12 13 13 15 16 16}
\end{tabularx}\endgroup\par\vspace{5pt}\noindent{\fontsize{8.4}{10}\selectfont\hyperref[trace-13]{Procedencia y códigos originales}\hfill Ficha completa · 013/324}
\clearpage\phantomsection\label{nav:vi-ruta-14}
% VI-CATALOG-ROW rutas 14
\pdfbookmark[3]{r1c4\_h-1v+1}{route-14}
\RouteTitle{Ruta 014}{r1c4\_\allowbreak{}h-1v+1}
\noindent Celda 4\quad (1, 4)\qquad Orientación \(h=-1,\ v=1\)\hfill\hyperref[sec:vi-indice-rutas]{Índice}\par\vspace{4pt}
\begingroup\fontsize{10.2}{12.5}\selectfont\setlength{\tabcolsep}{5pt}\renewcommand{\arraystretch}{1.1}\rowcolors{1}{routepale}{white}\noindent\begin{tabularx}{\linewidth}{@{}>{\raggedright\arraybackslash}p{.345\linewidth}>{\raggedright\arraybackslash}X@{}}
\RouteGroup{Identidad estructural}
\RouteRow{Clase y multisección}{neutrino\_\allowbreak{}G4\quad /\quad neutrino\_\allowbreak{}G4}
\RouteRow{Colección y sector}{neutrino; leptón}
\RouteRow{Clasificación física}{neutrino}
\RouteRow{Generación, profundidad y color}{4; G4; \textemdash{}}
\RouteRow{Ruta primaria y órbita de inversión}{\(B_1\): no\qquad 1,\allowbreak{}4;\allowbreak{}9,\allowbreak{}6}
\RouteRow{Firma de clasificación}{30|\allowbreak{}-6|\allowbreak{}-4|\allowbreak{}-2|\allowbreak{}-8|\allowbreak{}0++-\kern0pt{}-\kern0pt{}-0++-\kern0pt{}-\kern0pt{}-}
\RouteGroup{Retornos, memoria y transporte}
\RouteRow{Retornos con fase}{clase: 27; órbita: 27; celda: 27}
\RouteRow{Retorno cinemático}{en 108: 54; extensión: 216}
\RouteRow{Giros y cambios de hoja}{71; 107}
\RouteRow{Acarreo y diferimiento}{deuda total: 0; final: 0; bloques diferidos: 0}
\RouteRow{Tríadas finales; persistencia de Pareto}{17; sí}
\RouteRow{\(K\longrightarrow K+9\)}{repeticiones: 9; cambios de memoria: 9}
\RouteGroup{Firma, fase y diferencias}
\RouteRow{\(n_A,\ s_C,\ \kappa_0,\ \nu_{120},\ \nu_{270}\)}{44, 2, 0, 0, -12}
\RouteRow{\(\tau_{12}\); firma histórica \(\mathbb Z^5\)}{+++-\kern0pt{}-\kern0pt{}-+++-\kern0pt{}-\kern0pt{}-; 44|\allowbreak{}2|\allowbreak{}0|\allowbreak{}0|\allowbreak{}-12}
\RouteRow{\(D_1,\ D_3,\ D_4,\ D_{27},\ D_{36}\)}{78, 28, 78, 36, 24}
\RouteRow{\(\Omega_{90/120},\ P_6\)}{80, 6}
\RouteGroup{Lectores y factores de acción}
\RouteRow{\(K_D\quad /\quad K_\Omega\)}{(60,\allowbreak{}78,\allowbreak{}25)\quad /\quad (80,\allowbreak{}54,\allowbreak{}6)}
\RouteRow{\(\kappa_D\)}{59.43358\allowbreak{}64101631\allowbreak{}67023563\allowbreak{}04}
\RouteRow{\(\kappa_\Omega\)}{79.60661\allowbreak{}01397948\allowbreak{}27586470\allowbreak{}91}
\RouteRow{\(R_{\mathrm{act},D}\)}{1.000000\allowbreak{}04120422\allowbreak{}24053449\allowbreak{}05}
\RouteRow{\(R_{\mathrm{act},\Omega}\)}{1.000000\allowbreak{}05518981\allowbreak{}25318591\allowbreak{}40}
\RouteGroup{Separación de prefijos}
\RouteRow{Área de ambigüedad; área \(\log_2\)}{18; 0.0}
\RouteRow{Primer bloque de separación}{órbita: 1; celda: 1; ruta: 1}
\RouteRow{Ambigüedad final}{1}
\RouteGroup{Secuencias completas}
\RouteRow{\(w_6\)}{507 504 585 507 510 510 666 672 510 507 504 585 507 510 510 666 672 510}
\RouteRow{Tríadas estables por bloque}{0 1 2 3 4 5 6 7 8 9 10 11 12 13 14 15 16 17}
\end{tabularx}\endgroup\par\vspace{5pt}\noindent{\fontsize{8.4}{10}\selectfont\hyperref[trace-14]{Procedencia y códigos originales}\hfill Ficha completa · 014/324}
\clearpage\phantomsection\label{nav:vi-ruta-15}
% VI-CATALOG-ROW rutas 15
\pdfbookmark[3]{r1c4\_h+1v-1}{route-15}
\RouteTitle{Ruta 015}{r1c4\_\allowbreak{}h+1v-1}
\noindent Celda 4\quad (1, 4)\qquad Orientación \(h=1,\ v=-1\)\hfill\hyperref[sec:vi-indice-rutas]{Índice}\par\vspace{4pt}
\begingroup\fontsize{10.2}{12.5}\selectfont\setlength{\tabcolsep}{5pt}\renewcommand{\arraystretch}{1.1}\rowcolors{1}{routepale}{white}\noindent\begin{tabularx}{\linewidth}{@{}>{\raggedright\arraybackslash}p{.345\linewidth}>{\raggedright\arraybackslash}X@{}}
\RouteGroup{Identidad estructural}
\RouteRow{Clase y multisección}{neutrino\_\allowbreak{}G4\quad /\quad neutrino\_\allowbreak{}G4}
\RouteRow{Colección y sector}{neutrino; leptón}
\RouteRow{Clasificación física}{neutrino}
\RouteRow{Generación, profundidad y color}{4; G4; \textemdash{}}
\RouteRow{Ruta primaria y órbita de inversión}{\(B_1\): no\qquad 1,\allowbreak{}4;\allowbreak{}9,\allowbreak{}6}
\RouteRow{Firma de clasificación}{30|\allowbreak{}-6|\allowbreak{}-4|\allowbreak{}-2|\allowbreak{}-8|\allowbreak{}0++-\kern0pt{}-\kern0pt{}-0++-\kern0pt{}-\kern0pt{}-}
\RouteGroup{Retornos, memoria y transporte}
\RouteRow{Retornos con fase}{clase: 27; órbita: 27; celda: 27}
\RouteRow{Retorno cinemático}{en 108: 54; extensión: 216}
\RouteRow{Giros y cambios de hoja}{71; 107}
\RouteRow{Acarreo y diferimiento}{deuda total: 2; final: 0; bloques diferidos: 1}
\RouteRow{Tríadas finales; persistencia de Pareto}{17; sí}
\RouteRow{\(K\longrightarrow K+9\)}{repeticiones: 9; cambios de memoria: 9}
\RouteGroup{Firma, fase y diferencias}
\RouteRow{\(n_A,\ s_C,\ \kappa_0,\ \nu_{120},\ \nu_{270}\)}{44, -2, 0, 0, -12}
\RouteRow{\(\tau_{12}\); firma histórica \(\mathbb Z^5\)}{+++-\kern0pt{}-\kern0pt{}-+++-\kern0pt{}-\kern0pt{}-; 44|\allowbreak{}-2|\allowbreak{}0|\allowbreak{}0|\allowbreak{}-12}
\RouteRow{\(D_1,\ D_3,\ D_4,\ D_{27},\ D_{36}\)}{78, 28, 78, 36, 24}
\RouteRow{\(\Omega_{90/120},\ P_6\)}{80, 6}
\RouteGroup{Lectores y factores de acción}
\RouteRow{\(K_D\quad /\quad K_\Omega\)}{(60,\allowbreak{}78,\allowbreak{}25)\quad /\quad (80,\allowbreak{}54,\allowbreak{}6)}
\RouteRow{\(\kappa_D\)}{59.43358\allowbreak{}64101631\allowbreak{}67023563\allowbreak{}04}
\RouteRow{\(\kappa_\Omega\)}{79.60661\allowbreak{}01397948\allowbreak{}27586470\allowbreak{}91}
\RouteRow{\(R_{\mathrm{act},D}\)}{1.000000\allowbreak{}04120422\allowbreak{}24053449\allowbreak{}05}
\RouteRow{\(R_{\mathrm{act},\Omega}\)}{1.000000\allowbreak{}05518981\allowbreak{}25318591\allowbreak{}40}
\RouteGroup{Separación de prefijos}
\RouteRow{Área de ambigüedad; área \(\log_2\)}{18; 0.0}
\RouteRow{Primer bloque de separación}{órbita: 1; celda: 1; ruta: 1}
\RouteRow{Ambigüedad final}{1}
\RouteGroup{Secuencias completas}
\RouteRow{\(w_6\)}{666 672 510 666 666 585 507 504 585 666 672 510 666 666 585 507 504 585}
\RouteRow{Tríadas estables por bloque}{0 1 2 3 4 5 6 7 8 9 10 11 12 12 13 15 16 17}
\end{tabularx}\endgroup\par\vspace{5pt}\noindent{\fontsize{8.4}{10}\selectfont\hyperref[trace-15]{Procedencia y códigos originales}\hfill Ficha completa · 015/324}
\clearpage\phantomsection\label{nav:vi-ruta-16}
% VI-CATALOG-ROW rutas 16
\pdfbookmark[3]{r1c4\_h+1v+1}{route-16}
\RouteTitle{Ruta 016}{r1c4\_\allowbreak{}h+1v+1}
\noindent Celda 4\quad (1, 4)\qquad Orientación \(h=1,\ v=1\)\hfill\hyperref[sec:vi-indice-rutas]{Índice}\par\vspace{4pt}
\begingroup\fontsize{10.2}{12.5}\selectfont\setlength{\tabcolsep}{5pt}\renewcommand{\arraystretch}{1.1}\rowcolors{1}{routepale}{white}\noindent\begin{tabularx}{\linewidth}{@{}>{\raggedright\arraybackslash}p{.345\linewidth}>{\raggedright\arraybackslash}X@{}}
\RouteGroup{Identidad estructural}
\RouteRow{Clase y multisección}{neutrino\_\allowbreak{}G4\quad /\quad neutrino\_\allowbreak{}G4}
\RouteRow{Colección y sector}{neutrino; leptón}
\RouteRow{Clasificación física}{neutrino}
\RouteRow{Generación, profundidad y color}{4; G4; \textemdash{}}
\RouteRow{Ruta primaria y órbita de inversión}{\(B_1\): sí\qquad 1,\allowbreak{}4;\allowbreak{}9,\allowbreak{}6}
\RouteRow{Firma de clasificación}{30|\allowbreak{}-6|\allowbreak{}-4|\allowbreak{}-2|\allowbreak{}-8|\allowbreak{}0++-\kern0pt{}-\kern0pt{}-0++-\kern0pt{}-\kern0pt{}-}
\RouteGroup{Retornos, memoria y transporte}
\RouteRow{Retornos con fase}{clase: 27; órbita: 27; celda: 27}
\RouteRow{Retorno cinemático}{en 108: 54; extensión: 216}
\RouteRow{Giros y cambios de hoja}{71; 107}
\RouteRow{Acarreo y diferimiento}{deuda total: 2; final: 0; bloques diferidos: 1}
\RouteRow{Tríadas finales; persistencia de Pareto}{17; sí}
\RouteRow{\(K\longrightarrow K+9\)}{repeticiones: 9; cambios de memoria: 9}
\RouteGroup{Firma, fase y diferencias}
\RouteRow{\(n_A,\ s_C,\ \kappa_0,\ \nu_{120},\ \nu_{270}\)}{30, -6, 2, -4, -8}
\RouteRow{\(\tau_{12}\); firma histórica \(\mathbb Z^5\)}{0++-\kern0pt{}-\kern0pt{}-0++-\kern0pt{}-\kern0pt{}-; 30|\allowbreak{}-6|\allowbreak{}2|\allowbreak{}-4|\allowbreak{}-8}
\RouteRow{\(D_1,\ D_3,\ D_4,\ D_{27},\ D_{36}\)}{54, 56, 70, 48, 32}
\RouteRow{\(\Omega_{90/120},\ P_6\)}{56, 5}
\RouteGroup{Lectores y factores de acción}
\RouteRow{\(K_D\quad /\quad K_\Omega\)}{(80,\allowbreak{}54,\allowbreak{}7)\quad /\quad (56,\allowbreak{}54,\allowbreak{}5)}
\RouteRow{\(\kappa_D\)}{79.60672\allowbreak{}13362175\allowbreak{}19726524\allowbreak{}96}
\RouteRow{\(\kappa_\Omega\)}{55.60649\allowbreak{}89433721\allowbreak{}35446416\allowbreak{}86}
\RouteRow{\(R_{\mathrm{act},D}\)}{1.000000\allowbreak{}05518988\allowbreak{}96223153\allowbreak{}37}
\RouteRow{\(R_{\mathrm{act},\Omega}\)}{1.000000\allowbreak{}03855097\allowbreak{}23541416\allowbreak{}45}
\RouteGroup{Separación de prefijos}
\RouteRow{Área de ambigüedad; área \(\log_2\)}{18; 0.0}
\RouteRow{Primer bloque de separación}{órbita: 1; celda: 1; ruta: 1}
\RouteRow{Ambigüedad final}{1}
\RouteGroup{Secuencias completas}
\RouteRow{\(w_6\)}{657 24 243 657 18 15 486 423 15 657 24 243 657 18 15 486 423 15}
\RouteRow{Tríadas estables por bloque}{0 1 2 3 4 5 6 7 8 9 10 11 12 13 14 14 15 17}
\end{tabularx}\endgroup\par\vspace{5pt}\noindent{\fontsize{8.4}{10}\selectfont\hyperref[trace-16]{Procedencia y códigos originales}\hfill Ficha completa · 016/324}
\clearpage\phantomsection\label{nav:vi-ruta-17}
% VI-CATALOG-ROW rutas 17
\pdfbookmark[2]{Celda 5}{cell-5}
\pdfbookmark[3]{r1c5\_h-1v-1}{route-17}
\RouteTitle{Ruta 017}{r1c5\_\allowbreak{}h-1v-1}
\noindent Celda 5\quad (1, 5)\qquad Orientación \(h=-1,\ v=-1\)\hfill\hyperref[sec:vi-indice-rutas]{Índice}\par\vspace{4pt}
\begingroup\fontsize{10.2}{12.5}\selectfont\setlength{\tabcolsep}{5pt}\renewcommand{\arraystretch}{1.1}\rowcolors{1}{routepale}{white}\noindent\begin{tabularx}{\linewidth}{@{}>{\raggedright\arraybackslash}p{.345\linewidth}>{\raggedright\arraybackslash}X@{}}
\RouteGroup{Identidad estructural}
\RouteRow{Clase y multisección}{charged\_\allowbreak{}lepton\_\allowbreak{}G4\quad /\quad charged\_\allowbreak{}lepton\_\allowbreak{}G4}
\RouteRow{Colección y sector}{leptón cargado; leptón}
\RouteRow{Clasificación física}{leptón cargado}
\RouteRow{Generación, profundidad y color}{4; G4; \textemdash{}}
\RouteRow{Ruta primaria y órbita de inversión}{\(B_1\): no\qquad 1,\allowbreak{}5;\allowbreak{}9,\allowbreak{}5}
\RouteRow{Firma de clasificación}{34|\allowbreak{}0|\allowbreak{}-8|\allowbreak{}0|\allowbreak{}-8|\allowbreak{}+0+-\kern0pt{}-\kern0pt{}-+0+-\kern0pt{}-\kern0pt{}-}
\RouteGroup{Retornos, memoria y transporte}
\RouteRow{Retornos con fase}{clase: 27; órbita: 27; celda: 27}
\RouteRow{Retorno cinemático}{en 108: 54; extensión: 216}
\RouteRow{Giros y cambios de hoja}{71; 107}
\RouteRow{Acarreo y diferimiento}{deuda total: 1; final: 0; bloques diferidos: 1}
\RouteRow{Tríadas finales; persistencia de Pareto}{17; sí}
\RouteRow{\(K\longrightarrow K+9\)}{repeticiones: 9; cambios de memoria: 9}
\RouteGroup{Firma, fase y diferencias}
\RouteRow{\(n_A,\ s_C,\ \kappa_0,\ \nu_{120},\ \nu_{270}\)}{34, 0, -2, 0, -8}
\RouteRow{\(\tau_{12}\); firma histórica \(\mathbb Z^5\)}{+++-0-+++-0-; 34|\allowbreak{}0|\allowbreak{}-2|\allowbreak{}0|\allowbreak{}-8}
\RouteRow{\(D_1,\ D_3,\ D_4,\ D_{27},\ D_{36}\)}{66, 64, 64, 60, 40}
\RouteRow{\(\Omega_{90/120},\ P_6\)}{48, 5}
\RouteGroup{Lectores y factores de acción}
\RouteRow{\(K_D\quad /\quad K_\Omega\)}{(100,\allowbreak{}66,\allowbreak{}0)\quad /\quad (48,\allowbreak{}54,\allowbreak{}5)}
\RouteRow{\(\kappa_D\)}{99.51837\allowbreak{}47304272\allowbreak{}69134179\allowbreak{}18}
\RouteRow{\(\kappa_\Omega\)}{47.60649\allowbreak{}89433721\allowbreak{}35446416\allowbreak{}86}
\RouteRow{\(R_{\mathrm{act},D}\)}{1.000000\allowbreak{}06899427\allowbreak{}66450218\allowbreak{}95}
\RouteRow{\(R_{\mathrm{act},\Omega}\)}{1.000000\allowbreak{}03300471\allowbreak{}80532430\allowbreak{}84}
\RouteGroup{Separación de prefijos}
\RouteRow{Área de ambigüedad; área \(\log_2\)}{18; 0.0}
\RouteRow{Primer bloque de separación}{órbita: 1; celda: 1; ruta: 1}
\RouteRow{Ambigüedad final}{1}
\RouteGroup{Secuencias completas}
\RouteRow{\(w_6\)}{90 570 264 90 567 657 24 243 657 90 570 264 90 567 657 24 243 657}
\RouteRow{Tríadas estables por bloque}{0 1 2 3 4 5 6 7 8 9 10 11 12 13 13 15 16 17}
\end{tabularx}\endgroup\par\vspace{5pt}\noindent{\fontsize{8.4}{10}\selectfont\hyperref[trace-17]{Procedencia y códigos originales}\hfill Ficha completa · 017/324}
\clearpage\phantomsection\label{nav:vi-ruta-18}
% VI-CATALOG-ROW rutas 18
\pdfbookmark[3]{r1c5\_h-1v+1}{route-18}
\RouteTitle{Ruta 018}{r1c5\_\allowbreak{}h-1v+1}
\noindent Celda 5\quad (1, 5)\qquad Orientación \(h=-1,\ v=1\)\hfill\hyperref[sec:vi-indice-rutas]{Índice}\par\vspace{4pt}
\begingroup\fontsize{10.2}{12.5}\selectfont\setlength{\tabcolsep}{5pt}\renewcommand{\arraystretch}{1.1}\rowcolors{1}{routepale}{white}\noindent\begin{tabularx}{\linewidth}{@{}>{\raggedright\arraybackslash}p{.345\linewidth}>{\raggedright\arraybackslash}X@{}}
\RouteGroup{Identidad estructural}
\RouteRow{Clase y multisección}{charged\_\allowbreak{}lepton\_\allowbreak{}G4\quad /\quad charged\_\allowbreak{}lepton\_\allowbreak{}G4}
\RouteRow{Colección y sector}{leptón cargado; leptón}
\RouteRow{Clasificación física}{leptón cargado}
\RouteRow{Generación, profundidad y color}{4; G4; \textemdash{}}
\RouteRow{Ruta primaria y órbita de inversión}{\(B_1\): no\qquad 1,\allowbreak{}5;\allowbreak{}9,\allowbreak{}5}
\RouteRow{Firma de clasificación}{34|\allowbreak{}0|\allowbreak{}-8|\allowbreak{}0|\allowbreak{}-8|\allowbreak{}+0+-\kern0pt{}-\kern0pt{}-+0+-\kern0pt{}-\kern0pt{}-}
\RouteGroup{Retornos, memoria y transporte}
\RouteRow{Retornos con fase}{clase: 27; órbita: 27; celda: 27}
\RouteRow{Retorno cinemático}{en 108: 54; extensión: 216}
\RouteRow{Giros y cambios de hoja}{71; 107}
\RouteRow{Acarreo y diferimiento}{deuda total: 2; final: 0; bloques diferidos: 2}
\RouteRow{Tríadas finales; persistencia de Pareto}{17; sí}
\RouteRow{\(K\longrightarrow K+9\)}{repeticiones: 9; cambios de memoria: 9}
\RouteGroup{Firma, fase y diferencias}
\RouteRow{\(n_A,\ s_C,\ \kappa_0,\ \nu_{120},\ \nu_{270}\)}{14, 6, 0, 0, -12}
\RouteRow{\(\tau_{12}\); firma histórica \(\mathbb Z^5\)}{+++-\kern0pt{}-\kern0pt{}-+++-\kern0pt{}-\kern0pt{}-; 14|\allowbreak{}6|\allowbreak{}0|\allowbreak{}0|\allowbreak{}-12}
\RouteRow{\(D_1,\ D_3,\ D_4,\ D_{27},\ D_{36}\)}{24, 20, 24, 24, 16}
\RouteRow{\(\Omega_{90/120},\ P_6\)}{80, 6}
\RouteGroup{Lectores y factores de acción}
\RouteRow{\(K_D\quad /\quad K_\Omega\)}{(40,\allowbreak{}24,\allowbreak{}2)\quad /\quad (80,\allowbreak{}54,\allowbreak{}6)}
\RouteRow{\(\kappa_D\)}{39.82508\allowbreak{}59311825\allowbreak{}73056173\allowbreak{}26}
\RouteRow{\(\kappa_\Omega\)}{79.60661\allowbreak{}01397948\allowbreak{}27586470\allowbreak{}91}
\RouteRow{\(R_{\mathrm{act},D}\)}{1.000000\allowbreak{}02761000\allowbreak{}61595142\allowbreak{}15}
\RouteRow{\(R_{\mathrm{act},\Omega}\)}{1.000000\allowbreak{}05518981\allowbreak{}25318591\allowbreak{}40}
\RouteGroup{Separación de prefijos}
\RouteRow{Área de ambigüedad; área \(\log_2\)}{36; 18.0}
\RouteRow{Primer bloque de separación}{órbita: \textemdash{}; celda: \textemdash{}; ruta: \textemdash{}}
\RouteRow{Ambigüedad final}{2}
\RouteGroup{Secuencias completas}
\RouteRow{\(w_6\)}{81 3 0 81 0 3 0 81 3 81 3 0 81 0 3 0 81 3}
\RouteRow{Tríadas estables por bloque}{0 1 2 3 4 5 6 6 8 9 10 11 12 13 14 15 15 17}
\end{tabularx}\endgroup\par\vspace{5pt}\noindent{\fontsize{8.4}{10}\selectfont\hyperref[trace-18]{Procedencia y códigos originales}\hfill Ficha completa · 018/324}
\clearpage\phantomsection\label{nav:vi-ruta-19}
% VI-CATALOG-ROW rutas 19
\pdfbookmark[3]{r1c5\_h+1v-1}{route-19}
\RouteTitle{Ruta 019}{r1c5\_\allowbreak{}h+1v-1}
\noindent Celda 5\quad (1, 5)\qquad Orientación \(h=1,\ v=-1\)\hfill\hyperref[sec:vi-indice-rutas]{Índice}\par\vspace{4pt}
\begingroup\fontsize{10.2}{12.5}\selectfont\setlength{\tabcolsep}{5pt}\renewcommand{\arraystretch}{1.1}\rowcolors{1}{routepale}{white}\noindent\begin{tabularx}{\linewidth}{@{}>{\raggedright\arraybackslash}p{.345\linewidth}>{\raggedright\arraybackslash}X@{}}
\RouteGroup{Identidad estructural}
\RouteRow{Clase y multisección}{charged\_\allowbreak{}lepton\_\allowbreak{}G4\quad /\quad charged\_\allowbreak{}lepton\_\allowbreak{}G4}
\RouteRow{Colección y sector}{leptón cargado; leptón}
\RouteRow{Clasificación física}{leptón cargado}
\RouteRow{Generación, profundidad y color}{4; G4; \textemdash{}}
\RouteRow{Ruta primaria y órbita de inversión}{\(B_1\): sí\qquad 1,\allowbreak{}5;\allowbreak{}9,\allowbreak{}5}
\RouteRow{Firma de clasificación}{34|\allowbreak{}0|\allowbreak{}-8|\allowbreak{}0|\allowbreak{}-8|\allowbreak{}+0+-\kern0pt{}-\kern0pt{}-+0+-\kern0pt{}-\kern0pt{}-}
\RouteGroup{Retornos, memoria y transporte}
\RouteRow{Retornos con fase}{clase: 27; órbita: 27; celda: 27}
\RouteRow{Retorno cinemático}{en 108: 54; extensión: 216}
\RouteRow{Giros y cambios de hoja}{71; 107}
\RouteRow{Acarreo y diferimiento}{deuda total: 0; final: 0; bloques diferidos: 0}
\RouteRow{Tríadas finales; persistencia de Pareto}{17; sí}
\RouteRow{\(K\longrightarrow K+9\)}{repeticiones: 9; cambios de memoria: 9}
\RouteGroup{Firma, fase y diferencias}
\RouteRow{\(n_A,\ s_C,\ \kappa_0,\ \nu_{120},\ \nu_{270}\)}{14, -6, 0, 0, -12}
\RouteRow{\(\tau_{12}\); firma histórica \(\mathbb Z^5\)}{+++-\kern0pt{}-\kern0pt{}-+++-\kern0pt{}-\kern0pt{}-; 14|\allowbreak{}-6|\allowbreak{}0|\allowbreak{}0|\allowbreak{}-12}
\RouteRow{\(D_1,\ D_3,\ D_4,\ D_{27},\ D_{36}\)}{24, 20, 24, 24, 16}
\RouteRow{\(\Omega_{90/120},\ P_6\)}{80, 6}
\RouteGroup{Lectores y factores de acción}
\RouteRow{\(K_D\quad /\quad K_\Omega\)}{(40,\allowbreak{}24,\allowbreak{}2)\quad /\quad (80,\allowbreak{}54,\allowbreak{}6)}
\RouteRow{\(\kappa_D\)}{39.82508\allowbreak{}59311825\allowbreak{}73056173\allowbreak{}26}
\RouteRow{\(\kappa_\Omega\)}{79.60661\allowbreak{}01397948\allowbreak{}27586470\allowbreak{}91}
\RouteRow{\(R_{\mathrm{act},D}\)}{1.000000\allowbreak{}02761000\allowbreak{}61595142\allowbreak{}15}
\RouteRow{\(R_{\mathrm{act},\Omega}\)}{1.000000\allowbreak{}05518981\allowbreak{}25318591\allowbreak{}40}
\RouteGroup{Separación de prefijos}
\RouteRow{Área de ambigüedad; área \(\log_2\)}{36; 18.0}
\RouteRow{Primer bloque de separación}{órbita: \textemdash{}; celda: \textemdash{}; ruta: \textemdash{}}
\RouteRow{Ambigüedad final}{2}
\RouteGroup{Secuencias completas}
\RouteRow{\(w_6\)}{0 81 3 0 84 0 81 3 0 0 81 3 0 84 0 81 3 0}
\RouteRow{Tríadas estables por bloque}{0 1 2 3 4 5 6 7 8 9 10 11 12 13 14 15 16 17}
\end{tabularx}\endgroup\par\vspace{5pt}\noindent{\fontsize{8.4}{10}\selectfont\hyperref[trace-19]{Procedencia y códigos originales}\hfill Ficha completa · 019/324}
\clearpage\phantomsection\label{nav:vi-ruta-20}
% VI-CATALOG-ROW rutas 20
\pdfbookmark[3]{r1c5\_h+1v+1}{route-20}
\RouteTitle{Ruta 020}{r1c5\_\allowbreak{}h+1v+1}
\noindent Celda 5\quad (1, 5)\qquad Orientación \(h=1,\ v=1\)\hfill\hyperref[sec:vi-indice-rutas]{Índice}\par\vspace{4pt}
\begingroup\fontsize{10.2}{12.5}\selectfont\setlength{\tabcolsep}{5pt}\renewcommand{\arraystretch}{1.1}\rowcolors{1}{routepale}{white}\noindent\begin{tabularx}{\linewidth}{@{}>{\raggedright\arraybackslash}p{.345\linewidth}>{\raggedright\arraybackslash}X@{}}
\RouteGroup{Identidad estructural}
\RouteRow{Clase y multisección}{charged\_\allowbreak{}lepton\_\allowbreak{}G4\quad /\quad charged\_\allowbreak{}lepton\_\allowbreak{}G4}
\RouteRow{Colección y sector}{leptón cargado; leptón}
\RouteRow{Clasificación física}{leptón cargado}
\RouteRow{Generación, profundidad y color}{4; G4; \textemdash{}}
\RouteRow{Ruta primaria y órbita de inversión}{\(B_1\): no\qquad 1,\allowbreak{}5;\allowbreak{}9,\allowbreak{}5}
\RouteRow{Firma de clasificación}{34|\allowbreak{}0|\allowbreak{}-8|\allowbreak{}0|\allowbreak{}-8|\allowbreak{}+0+-\kern0pt{}-\kern0pt{}-+0+-\kern0pt{}-\kern0pt{}-}
\RouteGroup{Retornos, memoria y transporte}
\RouteRow{Retornos con fase}{clase: 27; órbita: 27; celda: 27}
\RouteRow{Retorno cinemático}{en 108: 54; extensión: 216}
\RouteRow{Giros y cambios de hoja}{71; 107}
\RouteRow{Acarreo y diferimiento}{deuda total: 0; final: 0; bloques diferidos: 0}
\RouteRow{Tríadas finales; persistencia de Pareto}{17; sí}
\RouteRow{\(K\longrightarrow K+9\)}{repeticiones: 9; cambios de memoria: 9}
\RouteGroup{Firma, fase y diferencias}
\RouteRow{\(n_A,\ s_C,\ \kappa_0,\ \nu_{120},\ \nu_{270}\)}{34, 0, 2, 0, -8}
\RouteRow{\(\tau_{12}\); firma histórica \(\mathbb Z^5\)}{+0+-\kern0pt{}-\kern0pt{}-+0+-\kern0pt{}-\kern0pt{}-; 34|\allowbreak{}0|\allowbreak{}2|\allowbreak{}0|\allowbreak{}-8}
\RouteRow{\(D_1,\ D_3,\ D_4,\ D_{27},\ D_{36}\)}{66, 64, 64, 60, 40}
\RouteRow{\(\Omega_{90/120},\ P_6\)}{48, 5}
\RouteGroup{Lectores y factores de acción}
\RouteRow{\(K_D\quad /\quad K_\Omega\)}{(100,\allowbreak{}66,\allowbreak{}0)\quad /\quad (48,\allowbreak{}54,\allowbreak{}5)}
\RouteRow{\(\kappa_D\)}{99.51837\allowbreak{}47304272\allowbreak{}69134179\allowbreak{}18}
\RouteRow{\(\kappa_\Omega\)}{47.60649\allowbreak{}89433721\allowbreak{}35446416\allowbreak{}86}
\RouteRow{\(R_{\mathrm{act},D}\)}{1.000000\allowbreak{}06899427\allowbreak{}66450218\allowbreak{}95}
\RouteRow{\(R_{\mathrm{act},\Omega}\)}{1.000000\allowbreak{}03300471\allowbreak{}80532430\allowbreak{}84}
\RouteGroup{Separación de prefijos}
\RouteRow{Área de ambigüedad; área \(\log_2\)}{30; 8.0}
\RouteRow{Primer bloque de separación}{órbita: 5; celda: 5; ruta: 5}
\RouteRow{Ambigüedad final}{1}
\RouteGroup{Secuencias completas}
\RouteRow{\(w_6\)}{24 243 657 24 246 264 90 570 264 24 243 657 24 246 264 90 570 264}
\RouteRow{Tríadas estables por bloque}{0 1 2 3 4 5 6 7 8 9 10 11 12 13 14 15 16 17}
\end{tabularx}\endgroup\par\vspace{5pt}\noindent{\fontsize{8.4}{10}\selectfont\hyperref[trace-20]{Procedencia y códigos originales}\hfill Ficha completa · 020/324}
\clearpage\phantomsection\label{nav:vi-ruta-21}
% VI-CATALOG-ROW rutas 21
\pdfbookmark[2]{Celda 6}{cell-6}
\pdfbookmark[3]{r1c6\_h-1v-1}{route-21}
\RouteTitle{Ruta 021}{r1c6\_\allowbreak{}h-1v-1}
\noindent Celda 6\quad (1, 6)\qquad Orientación \(h=-1,\ v=-1\)\hfill\hyperref[sec:vi-indice-rutas]{Índice}\par\vspace{4pt}
\begingroup\fontsize{10.2}{12.5}\selectfont\setlength{\tabcolsep}{5pt}\renewcommand{\arraystretch}{1.1}\rowcolors{1}{routepale}{white}\noindent\begin{tabularx}{\linewidth}{@{}>{\raggedright\arraybackslash}p{.345\linewidth}>{\raggedright\arraybackslash}X@{}}
\RouteGroup{Identidad estructural}
\RouteRow{Clase y multisección}{neutrino\_\allowbreak{}G4\quad /\quad neutrino\_\allowbreak{}G4}
\RouteRow{Colección y sector}{neutrino; leptón}
\RouteRow{Clasificación física}{neutrino}
\RouteRow{Generación, profundidad y color}{4; G4; \textemdash{}}
\RouteRow{Ruta primaria y órbita de inversión}{\(B_1\): sí\qquad 1,\allowbreak{}6;\allowbreak{}9,\allowbreak{}4}
\RouteRow{Firma de clasificación}{34|\allowbreak{}0|\allowbreak{}-4|\allowbreak{}0|\allowbreak{}-12|\allowbreak{}+++-\kern0pt{}-\kern0pt{}-+++-\kern0pt{}-\kern0pt{}-}
\RouteGroup{Retornos, memoria y transporte}
\RouteRow{Retornos con fase}{clase: 27; órbita: 27; celda: 27}
\RouteRow{Retorno cinemático}{en 108: 54; extensión: 216}
\RouteRow{Giros y cambios de hoja}{71; 107}
\RouteRow{Acarreo y diferimiento}{deuda total: 1; final: 1; bloques diferidos: 1}
\RouteRow{Tríadas finales; persistencia de Pareto}{16; sí}
\RouteRow{\(K\longrightarrow K+9\)}{repeticiones: 9; cambios de memoria: 9}
\RouteGroup{Firma, fase y diferencias}
\RouteRow{\(n_A,\ s_C,\ \kappa_0,\ \nu_{120},\ \nu_{270}\)}{34, 0, 0, 0, -12}
\RouteRow{\(\tau_{12}\); firma histórica \(\mathbb Z^5\)}{+++-\kern0pt{}-\kern0pt{}-+++-\kern0pt{}-\kern0pt{}-; 34|\allowbreak{}0|\allowbreak{}0|\allowbreak{}0|\allowbreak{}-12}
\RouteRow{\(D_1,\ D_3,\ D_4,\ D_{27},\ D_{36}\)}{66, 60, 64, 48, 32}
\RouteRow{\(\Omega_{90/120},\ P_6\)}{80, 6}
\RouteGroup{Lectores y factores de acción}
\RouteRow{\(K_D\quad /\quad K_\Omega\)}{(80,\allowbreak{}66,\allowbreak{}2)\quad /\quad (80,\allowbreak{}54,\allowbreak{}6)}
\RouteRow{\(\kappa_D\)}{79.51859\allowbreak{}71232726\allowbreak{}53414287\allowbreak{}29}
\RouteRow{\(\kappa_\Omega\)}{79.60661\allowbreak{}01397948\allowbreak{}27586470\allowbreak{}91}
\RouteRow{\(R_{\mathrm{act},D}\)}{1.000000\allowbreak{}05512879\allowbreak{}47092488\allowbreak{}80}
\RouteRow{\(R_{\mathrm{act},\Omega}\)}{1.000000\allowbreak{}05518981\allowbreak{}25318591\allowbreak{}40}
\RouteGroup{Separación de prefijos}
\RouteRow{Área de ambigüedad; área \(\log_2\)}{18; 0.0}
\RouteRow{Primer bloque de separación}{órbita: 1; celda: 1; ruta: 1}
\RouteRow{Ambigüedad final}{1}
\RouteGroup{Secuencias completas}
\RouteRow{\(w_6\)}{423 15 486 423 12 99 327 498 99 423 15 486 423 12 99 327 498 99}
\RouteRow{Tríadas estables por bloque}{0 1 2 3 4 5 6 7 8 9 10 11 12 13 14 15 16 16}
\end{tabularx}\endgroup\par\vspace{5pt}\noindent{\fontsize{8.4}{10}\selectfont\hyperref[trace-21]{Procedencia y códigos originales}\hfill Ficha completa · 021/324}
\clearpage\phantomsection\label{nav:vi-ruta-22}
% VI-CATALOG-ROW rutas 22
\pdfbookmark[3]{r1c6\_h-1v+1}{route-22}
\RouteTitle{Ruta 022}{r1c6\_\allowbreak{}h-1v+1}
\noindent Celda 6\quad (1, 6)\qquad Orientación \(h=-1,\ v=1\)\hfill\hyperref[sec:vi-indice-rutas]{Índice}\par\vspace{4pt}
\begingroup\fontsize{10.2}{12.5}\selectfont\setlength{\tabcolsep}{5pt}\renewcommand{\arraystretch}{1.1}\rowcolors{1}{routepale}{white}\noindent\begin{tabularx}{\linewidth}{@{}>{\raggedright\arraybackslash}p{.345\linewidth}>{\raggedright\arraybackslash}X@{}}
\RouteGroup{Identidad estructural}
\RouteRow{Clase y multisección}{neutrino\_\allowbreak{}G4\quad /\quad neutrino\_\allowbreak{}G4}
\RouteRow{Colección y sector}{neutrino; leptón}
\RouteRow{Clasificación física}{neutrino}
\RouteRow{Generación, profundidad y color}{4; G4; \textemdash{}}
\RouteRow{Ruta primaria y órbita de inversión}{\(B_1\): no\qquad 1,\allowbreak{}6;\allowbreak{}9,\allowbreak{}4}
\RouteRow{Firma de clasificación}{34|\allowbreak{}0|\allowbreak{}-4|\allowbreak{}0|\allowbreak{}-12|\allowbreak{}+++-\kern0pt{}-\kern0pt{}-+++-\kern0pt{}-\kern0pt{}-}
\RouteGroup{Retornos, memoria y transporte}
\RouteRow{Retornos con fase}{clase: 27; órbita: 27; celda: 27}
\RouteRow{Retorno cinemático}{en 108: 54; extensión: 216}
\RouteRow{Giros y cambios de hoja}{71; 107}
\RouteRow{Acarreo y diferimiento}{deuda total: 1; final: 0; bloques diferidos: 1}
\RouteRow{Tríadas finales; persistencia de Pareto}{17; sí}
\RouteRow{\(K\longrightarrow K+9\)}{repeticiones: 9; cambios de memoria: 9}
\RouteGroup{Firma, fase y diferencias}
\RouteRow{\(n_A,\ s_C,\ \kappa_0,\ \nu_{120},\ \nu_{270}\)}{50, 4, -2, 2, -8}
\RouteRow{\(\tau_{12}\); firma histórica \(\mathbb Z^5\)}{+++0-\kern0pt{}-+++0-\kern0pt{}-; 50|\allowbreak{}4|\allowbreak{}-2|\allowbreak{}2|\allowbreak{}-8}
\RouteRow{\(D_1,\ D_3,\ D_4,\ D_{27},\ D_{36}\)}{84, 24, 84, 24, 16}
\RouteRow{\(\Omega_{90/120},\ P_6\)}{56, 5}
\RouteGroup{Lectores y factores de acción}
\RouteRow{\(K_D\quad /\quad K_\Omega\)}{(40,\allowbreak{}84,\allowbreak{}30)\quad /\quad (56,\allowbreak{}54,\allowbreak{}5)}
\RouteRow{\(\kappa_D\)}{39.39035\allowbreak{}82768609\allowbreak{}24917849\allowbreak{}59}
\RouteRow{\(\kappa_\Omega\)}{55.60649\allowbreak{}89433721\allowbreak{}35446416\allowbreak{}86}
\RouteRow{\(R_{\mathrm{act},D}\)}{1.000000\allowbreak{}02730861\allowbreak{}73967087\allowbreak{}05}
\RouteRow{\(R_{\mathrm{act},\Omega}\)}{1.000000\allowbreak{}03855097\allowbreak{}23541416\allowbreak{}45}
\RouteGroup{Separación de prefijos}
\RouteRow{Área de ambigüedad; área \(\log_2\)}{18; 0.0}
\RouteRow{Primer bloque de separación}{órbita: 1; celda: 1; ruta: 1}
\RouteRow{Ambigüedad final}{1}
\RouteGroup{Secuencias completas}
\RouteRow{\(w_6\)}{420 258 414 420 255 252 333 255 252 420 258 414 420 255 252 333 255 252}
\RouteRow{Tríadas estables por bloque}{0 1 2 3 4 5 6 7 8 9 10 11 12 13 14 15 15 17}
\end{tabularx}\endgroup\par\vspace{5pt}\noindent{\fontsize{8.4}{10}\selectfont\hyperref[trace-22]{Procedencia y códigos originales}\hfill Ficha completa · 022/324}
\clearpage\phantomsection\label{nav:vi-ruta-23}
% VI-CATALOG-ROW rutas 23
\pdfbookmark[3]{r1c6\_h+1v-1}{route-23}
\RouteTitle{Ruta 023}{r1c6\_\allowbreak{}h+1v-1}
\noindent Celda 6\quad (1, 6)\qquad Orientación \(h=1,\ v=-1\)\hfill\hyperref[sec:vi-indice-rutas]{Índice}\par\vspace{4pt}
\begingroup\fontsize{10.2}{12.5}\selectfont\setlength{\tabcolsep}{5pt}\renewcommand{\arraystretch}{1.1}\rowcolors{1}{routepale}{white}\noindent\begin{tabularx}{\linewidth}{@{}>{\raggedright\arraybackslash}p{.345\linewidth}>{\raggedright\arraybackslash}X@{}}
\RouteGroup{Identidad estructural}
\RouteRow{Clase y multisección}{neutrino\_\allowbreak{}G4\quad /\quad neutrino\_\allowbreak{}G4}
\RouteRow{Colección y sector}{neutrino; leptón}
\RouteRow{Clasificación física}{neutrino}
\RouteRow{Generación, profundidad y color}{4; G4; \textemdash{}}
\RouteRow{Ruta primaria y órbita de inversión}{\(B_1\): no\qquad 1,\allowbreak{}6;\allowbreak{}9,\allowbreak{}4}
\RouteRow{Firma de clasificación}{34|\allowbreak{}0|\allowbreak{}-4|\allowbreak{}0|\allowbreak{}-12|\allowbreak{}+++-\kern0pt{}-\kern0pt{}-+++-\kern0pt{}-\kern0pt{}-}
\RouteGroup{Retornos, memoria y transporte}
\RouteRow{Retornos con fase}{clase: 27; órbita: 27; celda: 27}
\RouteRow{Retorno cinemático}{en 108: 54; extensión: 216}
\RouteRow{Giros y cambios de hoja}{71; 107}
\RouteRow{Acarreo y diferimiento}{deuda total: 0; final: 0; bloques diferidos: 0}
\RouteRow{Tríadas finales; persistencia de Pareto}{17; sí}
\RouteRow{\(K\longrightarrow K+9\)}{repeticiones: 9; cambios de memoria: 9}
\RouteGroup{Firma, fase y diferencias}
\RouteRow{\(n_A,\ s_C,\ \kappa_0,\ \nu_{120},\ \nu_{270}\)}{50, -4, 2, -2, -8}
\RouteRow{\(\tau_{12}\); firma histórica \(\mathbb Z^5\)}{0++-\kern0pt{}-\kern0pt{}-0++-\kern0pt{}-\kern0pt{}-; 50|\allowbreak{}-4|\allowbreak{}2|\allowbreak{}-2|\allowbreak{}-8}
\RouteRow{\(D_1,\ D_3,\ D_4,\ D_{27},\ D_{36}\)}{84, 24, 84, 24, 16}
\RouteRow{\(\Omega_{90/120},\ P_6\)}{56, 5}
\RouteGroup{Lectores y factores de acción}
\RouteRow{\(K_D\quad /\quad K_\Omega\)}{(40,\allowbreak{}84,\allowbreak{}30)\quad /\quad (56,\allowbreak{}54,\allowbreak{}5)}
\RouteRow{\(\kappa_D\)}{39.39035\allowbreak{}82768609\allowbreak{}24917849\allowbreak{}59}
\RouteRow{\(\kappa_\Omega\)}{55.60649\allowbreak{}89433721\allowbreak{}35446416\allowbreak{}86}
\RouteRow{\(R_{\mathrm{act},D}\)}{1.000000\allowbreak{}02730861\allowbreak{}73967087\allowbreak{}05}
\RouteRow{\(R_{\mathrm{act},\Omega}\)}{1.000000\allowbreak{}03855097\allowbreak{}23541416\allowbreak{}45}
\RouteGroup{Separación de prefijos}
\RouteRow{Área de ambigüedad; área \(\log_2\)}{18; 0.0}
\RouteRow{Primer bloque de separación}{órbita: 1; celda: 1; ruta: 1}
\RouteRow{Ambigüedad final}{1}
\RouteGroup{Secuencias completas}
\RouteRow{\(w_6\)}{333 255 252 333 258 414 420 258 414 333 255 252 333 258 414 420 258 414}
\RouteRow{Tríadas estables por bloque}{0 1 2 3 4 5 6 7 8 9 10 11 12 13 14 15 16 17}
\end{tabularx}\endgroup\par\vspace{5pt}\noindent{\fontsize{8.4}{10}\selectfont\hyperref[trace-23]{Procedencia y códigos originales}\hfill Ficha completa · 023/324}
\clearpage\phantomsection\label{nav:vi-ruta-24}
% VI-CATALOG-ROW rutas 24
\pdfbookmark[3]{r1c6\_h+1v+1}{route-24}
\RouteTitle{Ruta 024}{r1c6\_\allowbreak{}h+1v+1}
\noindent Celda 6\quad (1, 6)\qquad Orientación \(h=1,\ v=1\)\hfill\hyperref[sec:vi-indice-rutas]{Índice}\par\vspace{4pt}
\begingroup\fontsize{10.2}{12.5}\selectfont\setlength{\tabcolsep}{5pt}\renewcommand{\arraystretch}{1.1}\rowcolors{1}{routepale}{white}\noindent\begin{tabularx}{\linewidth}{@{}>{\raggedright\arraybackslash}p{.345\linewidth}>{\raggedright\arraybackslash}X@{}}
\RouteGroup{Identidad estructural}
\RouteRow{Clase y multisección}{neutrino\_\allowbreak{}G4\quad /\quad neutrino\_\allowbreak{}G4}
\RouteRow{Colección y sector}{neutrino; leptón}
\RouteRow{Clasificación física}{neutrino}
\RouteRow{Generación, profundidad y color}{4; G4; \textemdash{}}
\RouteRow{Ruta primaria y órbita de inversión}{\(B_1\): no\qquad 1,\allowbreak{}6;\allowbreak{}9,\allowbreak{}4}
\RouteRow{Firma de clasificación}{34|\allowbreak{}0|\allowbreak{}-4|\allowbreak{}0|\allowbreak{}-12|\allowbreak{}+++-\kern0pt{}-\kern0pt{}-+++-\kern0pt{}-\kern0pt{}-}
\RouteGroup{Retornos, memoria y transporte}
\RouteRow{Retornos con fase}{clase: 27; órbita: 27; celda: 27}
\RouteRow{Retorno cinemático}{en 108: 54; extensión: 216}
\RouteRow{Giros y cambios de hoja}{71; 107}
\RouteRow{Acarreo y diferimiento}{deuda total: 2; final: 1; bloques diferidos: 2}
\RouteRow{Tríadas finales; persistencia de Pareto}{16; sí}
\RouteRow{\(K\longrightarrow K+9\)}{repeticiones: 9; cambios de memoria: 9}
\RouteGroup{Firma, fase y diferencias}
\RouteRow{\(n_A,\ s_C,\ \kappa_0,\ \nu_{120},\ \nu_{270}\)}{34, 0, 0, 0, -12}
\RouteRow{\(\tau_{12}\); firma histórica \(\mathbb Z^5\)}{+++-\kern0pt{}-\kern0pt{}-+++-\kern0pt{}-\kern0pt{}-; 34|\allowbreak{}0|\allowbreak{}0|\allowbreak{}0|\allowbreak{}-12}
\RouteRow{\(D_1,\ D_3,\ D_4,\ D_{27},\ D_{36}\)}{66, 60, 64, 48, 32}
\RouteRow{\(\Omega_{90/120},\ P_6\)}{80, 6}
\RouteGroup{Lectores y factores de acción}
\RouteRow{\(K_D\quad /\quad K_\Omega\)}{(80,\allowbreak{}66,\allowbreak{}2)\quad /\quad (80,\allowbreak{}54,\allowbreak{}6)}
\RouteRow{\(\kappa_D\)}{79.51859\allowbreak{}71232726\allowbreak{}53414287\allowbreak{}29}
\RouteRow{\(\kappa_\Omega\)}{79.60661\allowbreak{}01397948\allowbreak{}27586470\allowbreak{}91}
\RouteRow{\(R_{\mathrm{act},D}\)}{1.000000\allowbreak{}05512879\allowbreak{}47092488\allowbreak{}80}
\RouteRow{\(R_{\mathrm{act},\Omega}\)}{1.000000\allowbreak{}05518981\allowbreak{}25318591\allowbreak{}40}
\RouteGroup{Separación de prefijos}
\RouteRow{Área de ambigüedad; área \(\log_2\)}{18; 0.0}
\RouteRow{Primer bloque de separación}{órbita: 1; celda: 1; ruta: 1}
\RouteRow{Ambigüedad final}{1}
\RouteGroup{Secuencias completas}
\RouteRow{\(w_6\)}{327 498 99 327 501 486 423 15 486 327 498 99 327 501 486 423 15 486}
\RouteRow{Tríadas estables por bloque}{0 1 2 3 4 5 6 7 7 9 10 11 12 13 14 15 16 16}
\end{tabularx}\endgroup\par\vspace{5pt}\noindent{\fontsize{8.4}{10}\selectfont\hyperref[trace-24]{Procedencia y códigos originales}\hfill Ficha completa · 024/324}
\clearpage\phantomsection\label{nav:vi-ruta-25}
% VI-CATALOG-ROW rutas 25
\pdfbookmark[2]{Celda 7}{cell-7}
\pdfbookmark[3]{r1c7\_h-1v-1}{route-25}
\RouteTitle{Ruta 025}{r1c7\_\allowbreak{}h-1v-1}
\noindent Celda 7\quad (1, 7)\qquad Orientación \(h=-1,\ v=-1\)\hfill\hyperref[sec:vi-indice-rutas]{Índice}\par\vspace{4pt}
\begingroup\fontsize{10.2}{12.5}\selectfont\setlength{\tabcolsep}{5pt}\renewcommand{\arraystretch}{1.1}\rowcolors{1}{routepale}{white}\noindent\begin{tabularx}{\linewidth}{@{}>{\raggedright\arraybackslash}p{.345\linewidth}>{\raggedright\arraybackslash}X@{}}
\RouteGroup{Identidad estructural}
\RouteRow{Clase y multisección}{charged\_\allowbreak{}lepton\_\allowbreak{}G4\quad /\quad charged\_\allowbreak{}lepton\_\allowbreak{}G4}
\RouteRow{Colección y sector}{leptón cargado; leptón}
\RouteRow{Clasificación física}{leptón cargado}
\RouteRow{Generación, profundidad y color}{4; G4; \textemdash{}}
\RouteRow{Ruta primaria y órbita de inversión}{\(B_1\): no\qquad 1,\allowbreak{}7;\allowbreak{}9,\allowbreak{}3}
\RouteRow{Firma de clasificación}{30|\allowbreak{}6|\allowbreak{}4|\allowbreak{}0|\allowbreak{}-12|\allowbreak{}+++-\kern0pt{}-\kern0pt{}-+++-\kern0pt{}-\kern0pt{}-}
\RouteGroup{Retornos, memoria y transporte}
\RouteRow{Retornos con fase}{clase: 27; órbita: 27; celda: 27}
\RouteRow{Retorno cinemático}{en 108: 54; extensión: 216}
\RouteRow{Giros y cambios de hoja}{71; 107}
\RouteRow{Acarreo y diferimiento}{deuda total: 2; final: 1; bloques diferidos: 2}
\RouteRow{Tríadas finales; persistencia de Pareto}{16; sí}
\RouteRow{\(K\longrightarrow K+9\)}{repeticiones: 9; cambios de memoria: 9}
\RouteGroup{Firma, fase y diferencias}
\RouteRow{\(n_A,\ s_C,\ \kappa_0,\ \nu_{120},\ \nu_{270}\)}{30, -6, 0, 0, -12}
\RouteRow{\(\tau_{12}\); firma histórica \(\mathbb Z^5\)}{+++-\kern0pt{}-\kern0pt{}-+++-\kern0pt{}-\kern0pt{}-; 30|\allowbreak{}-6|\allowbreak{}0|\allowbreak{}0|\allowbreak{}-12}
\RouteRow{\(D_1,\ D_3,\ D_4,\ D_{27},\ D_{36}\)}{54, 56, 70, 48, 32}
\RouteRow{\(\Omega_{90/120},\ P_6\)}{80, 6}
\RouteGroup{Lectores y factores de acción}
\RouteRow{\(K_D\quad /\quad K_\Omega\)}{(80,\allowbreak{}54,\allowbreak{}7)\quad /\quad (80,\allowbreak{}54,\allowbreak{}6)}
\RouteRow{\(\kappa_D\)}{79.60672\allowbreak{}13362175\allowbreak{}19726524\allowbreak{}96}
\RouteRow{\(\kappa_\Omega\)}{79.60661\allowbreak{}01397948\allowbreak{}27586470\allowbreak{}91}
\RouteRow{\(R_{\mathrm{act},D}\)}{1.000000\allowbreak{}05518988\allowbreak{}96223153\allowbreak{}37}
\RouteRow{\(R_{\mathrm{act},\Omega}\)}{1.000000\allowbreak{}05518981\allowbreak{}25318591\allowbreak{}40}
\RouteGroup{Separación de prefijos}
\RouteRow{Área de ambigüedad; área \(\log_2\)}{58; 30.18947\allowbreak{}50100961\allowbreak{}74}
\RouteRow{Primer bloque de separación}{órbita: \textemdash{}; celda: \textemdash{}; ruta: \textemdash{}}
\RouteRow{Ambigüedad final}{3}
\RouteGroup{Secuencias completas}
\RouteRow{\(w_6\)}{486 423 15 486 429 243 657 24 243 486 423 15 486 429 243 657 24 243}
\RouteRow{Tríadas estables por bloque}{0 1 2 3 4 5 6 7 8 9 10 11 12 13 13 15 16 16}
\end{tabularx}\endgroup\par\vspace{5pt}\noindent{\fontsize{8.4}{10}\selectfont\hyperref[trace-25]{Procedencia y códigos originales}\hfill Ficha completa · 025/324}
\clearpage\phantomsection\label{nav:vi-ruta-26}
% VI-CATALOG-ROW rutas 26
\pdfbookmark[3]{r1c7\_h-1v+1}{route-26}
\RouteTitle{Ruta 026}{r1c7\_\allowbreak{}h-1v+1}
\noindent Celda 7\quad (1, 7)\qquad Orientación \(h=-1,\ v=1\)\hfill\hyperref[sec:vi-indice-rutas]{Índice}\par\vspace{4pt}
\begingroup\fontsize{10.2}{12.5}\selectfont\setlength{\tabcolsep}{5pt}\renewcommand{\arraystretch}{1.1}\rowcolors{1}{routepale}{white}\noindent\begin{tabularx}{\linewidth}{@{}>{\raggedright\arraybackslash}p{.345\linewidth}>{\raggedright\arraybackslash}X@{}}
\RouteGroup{Identidad estructural}
\RouteRow{Clase y multisección}{charged\_\allowbreak{}lepton\_\allowbreak{}G4\quad /\quad charged\_\allowbreak{}lepton\_\allowbreak{}G4}
\RouteRow{Colección y sector}{leptón cargado; leptón}
\RouteRow{Clasificación física}{leptón cargado}
\RouteRow{Generación, profundidad y color}{4; G4; \textemdash{}}
\RouteRow{Ruta primaria y órbita de inversión}{\(B_1\): sí\qquad 1,\allowbreak{}7;\allowbreak{}9,\allowbreak{}3}
\RouteRow{Firma de clasificación}{30|\allowbreak{}6|\allowbreak{}4|\allowbreak{}0|\allowbreak{}-12|\allowbreak{}+++-\kern0pt{}-\kern0pt{}-+++-\kern0pt{}-\kern0pt{}-}
\RouteGroup{Retornos, memoria y transporte}
\RouteRow{Retornos con fase}{clase: 18; órbita: 27; celda: 27}
\RouteRow{Retorno cinemático}{en 108: 54; extensión: 216}
\RouteRow{Giros y cambios de hoja}{71; 107}
\RouteRow{Acarreo y diferimiento}{deuda total: 0; final: 0; bloques diferidos: 0}
\RouteRow{Tríadas finales; persistencia de Pareto}{17; sí}
\RouteRow{\(K\longrightarrow K+9\)}{repeticiones: 9; cambios de memoria: 9}
\RouteGroup{Firma, fase y diferencias}
\RouteRow{\(n_A,\ s_C,\ \kappa_0,\ \nu_{120},\ \nu_{270}\)}{14, -4, 0, 0, -12}
\RouteRow{\(\tau_{12}\); firma histórica \(\mathbb Z^5\)}{+++-\kern0pt{}-\kern0pt{}-+++-\kern0pt{}-\kern0pt{}-; 14|\allowbreak{}-4|\allowbreak{}0|\allowbreak{}0|\allowbreak{}-12}
\RouteRow{\(D_1,\ D_3,\ D_4,\ D_{27},\ D_{36}\)}{78, 28, 78, 36, 24}
\RouteRow{\(\Omega_{90/120},\ P_6\)}{80, 6}
\RouteGroup{Lectores y factores de acción}
\RouteRow{\(K_D\quad /\quad K_\Omega\)}{(60,\allowbreak{}78,\allowbreak{}25)\quad /\quad (80,\allowbreak{}54,\allowbreak{}6)}
\RouteRow{\(\kappa_D\)}{59.43358\allowbreak{}64101631\allowbreak{}67023563\allowbreak{}04}
\RouteRow{\(\kappa_\Omega\)}{79.60661\allowbreak{}01397948\allowbreak{}27586470\allowbreak{}91}
\RouteRow{\(R_{\mathrm{act},D}\)}{1.000000\allowbreak{}04120422\allowbreak{}24053449\allowbreak{}05}
\RouteRow{\(R_{\mathrm{act},\Omega}\)}{1.000000\allowbreak{}05518981\allowbreak{}25318591\allowbreak{}40}
\RouteGroup{Separación de prefijos}
\RouteRow{Área de ambigüedad; área \(\log_2\)}{54; 28.52932\allowbreak{}50129808}
\RouteRow{Primer bloque de separación}{órbita: \textemdash{}; celda: \textemdash{}; ruta: \textemdash{}}
\RouteRow{Ambigüedad final}{3}
\RouteGroup{Secuencias completas}
\RouteRow{\(w_6\)}{507 504 585 507 510 510 666 672 510 507 504 585 507 510 510 666 672 510}
\RouteRow{Tríadas estables por bloque}{0 1 2 3 4 5 6 7 8 9 10 11 12 13 14 15 16 17}
\end{tabularx}\endgroup\par\vspace{5pt}\noindent{\fontsize{8.4}{10}\selectfont\hyperref[trace-26]{Procedencia y códigos originales}\hfill Ficha completa · 026/324}
\clearpage\phantomsection\label{nav:vi-ruta-27}
% VI-CATALOG-ROW rutas 27
\pdfbookmark[3]{r1c7\_h+1v-1}{route-27}
\RouteTitle{Ruta 027}{r1c7\_\allowbreak{}h+1v-1}
\noindent Celda 7\quad (1, 7)\qquad Orientación \(h=1,\ v=-1\)\hfill\hyperref[sec:vi-indice-rutas]{Índice}\par\vspace{4pt}
\begingroup\fontsize{10.2}{12.5}\selectfont\setlength{\tabcolsep}{5pt}\renewcommand{\arraystretch}{1.1}\rowcolors{1}{routepale}{white}\noindent\begin{tabularx}{\linewidth}{@{}>{\raggedright\arraybackslash}p{.345\linewidth}>{\raggedright\arraybackslash}X@{}}
\RouteGroup{Identidad estructural}
\RouteRow{Clase y multisección}{charged\_\allowbreak{}lepton\_\allowbreak{}G4\quad /\quad charged\_\allowbreak{}lepton\_\allowbreak{}G4}
\RouteRow{Colección y sector}{leptón cargado; leptón}
\RouteRow{Clasificación física}{leptón cargado}
\RouteRow{Generación, profundidad y color}{4; G4; \textemdash{}}
\RouteRow{Ruta primaria y órbita de inversión}{\(B_1\): no\qquad 1,\allowbreak{}7;\allowbreak{}9,\allowbreak{}3}
\RouteRow{Firma de clasificación}{30|\allowbreak{}6|\allowbreak{}4|\allowbreak{}0|\allowbreak{}-12|\allowbreak{}+++-\kern0pt{}-\kern0pt{}-+++-\kern0pt{}-\kern0pt{}-}
\RouteGroup{Retornos, memoria y transporte}
\RouteRow{Retornos con fase}{clase: 9; órbita: 27; celda: 27}
\RouteRow{Retorno cinemático}{en 108: 54; extensión: 216}
\RouteRow{Giros y cambios de hoja}{71; 107}
\RouteRow{Acarreo y diferimiento}{deuda total: 2; final: 0; bloques diferidos: 1}
\RouteRow{Tríadas finales; persistencia de Pareto}{17; sí}
\RouteRow{\(K\longrightarrow K+9\)}{repeticiones: 9; cambios de memoria: 9}
\RouteGroup{Firma, fase y diferencias}
\RouteRow{\(n_A,\ s_C,\ \kappa_0,\ \nu_{120},\ \nu_{270}\)}{14, 4, 0, 0, -12}
\RouteRow{\(\tau_{12}\); firma histórica \(\mathbb Z^5\)}{+++-\kern0pt{}-\kern0pt{}-+++-\kern0pt{}-\kern0pt{}-; 14|\allowbreak{}4|\allowbreak{}0|\allowbreak{}0|\allowbreak{}-12}
\RouteRow{\(D_1,\ D_3,\ D_4,\ D_{27},\ D_{36}\)}{78, 28, 78, 36, 24}
\RouteRow{\(\Omega_{90/120},\ P_6\)}{80, 6}
\RouteGroup{Lectores y factores de acción}
\RouteRow{\(K_D\quad /\quad K_\Omega\)}{(60,\allowbreak{}78,\allowbreak{}25)\quad /\quad (80,\allowbreak{}54,\allowbreak{}6)}
\RouteRow{\(\kappa_D\)}{59.43358\allowbreak{}64101631\allowbreak{}67023563\allowbreak{}04}
\RouteRow{\(\kappa_\Omega\)}{79.60661\allowbreak{}01397948\allowbreak{}27586470\allowbreak{}91}
\RouteRow{\(R_{\mathrm{act},D}\)}{1.000000\allowbreak{}04120422\allowbreak{}24053449\allowbreak{}05}
\RouteRow{\(R_{\mathrm{act},\Omega}\)}{1.000000\allowbreak{}05518981\allowbreak{}25318591\allowbreak{}40}
\RouteGroup{Separación de prefijos}
\RouteRow{Área de ambigüedad; área \(\log_2\)}{54; 28.52932\allowbreak{}50129808}
\RouteRow{Primer bloque de separación}{órbita: \textemdash{}; celda: \textemdash{}; ruta: \textemdash{}}
\RouteRow{Ambigüedad final}{3}
\RouteGroup{Secuencias completas}
\RouteRow{\(w_6\)}{666 672 510 666 666 585 507 504 585 666 672 510 666 666 585 507 504 585}
\RouteRow{Tríadas estables por bloque}{0 1 2 3 4 5 6 7 8 9 10 11 12 12 13 15 16 17}
\end{tabularx}\endgroup\par\vspace{5pt}\noindent{\fontsize{8.4}{10}\selectfont\hyperref[trace-27]{Procedencia y códigos originales}\hfill Ficha completa · 027/324}
\clearpage\phantomsection\label{nav:vi-ruta-28}
% VI-CATALOG-ROW rutas 28
\pdfbookmark[3]{r1c7\_h+1v+1}{route-28}
\RouteTitle{Ruta 028}{r1c7\_\allowbreak{}h+1v+1}
\noindent Celda 7\quad (1, 7)\qquad Orientación \(h=1,\ v=1\)\hfill\hyperref[sec:vi-indice-rutas]{Índice}\par\vspace{4pt}
\begingroup\fontsize{10.2}{12.5}\selectfont\setlength{\tabcolsep}{5pt}\renewcommand{\arraystretch}{1.1}\rowcolors{1}{routepale}{white}\noindent\begin{tabularx}{\linewidth}{@{}>{\raggedright\arraybackslash}p{.345\linewidth}>{\raggedright\arraybackslash}X@{}}
\RouteGroup{Identidad estructural}
\RouteRow{Clase y multisección}{charged\_\allowbreak{}lepton\_\allowbreak{}G4\quad /\quad charged\_\allowbreak{}lepton\_\allowbreak{}G4}
\RouteRow{Colección y sector}{leptón cargado; leptón}
\RouteRow{Clasificación física}{leptón cargado}
\RouteRow{Generación, profundidad y color}{4; G4; \textemdash{}}
\RouteRow{Ruta primaria y órbita de inversión}{\(B_1\): no\qquad 1,\allowbreak{}7;\allowbreak{}9,\allowbreak{}3}
\RouteRow{Firma de clasificación}{30|\allowbreak{}6|\allowbreak{}4|\allowbreak{}0|\allowbreak{}-12|\allowbreak{}+++-\kern0pt{}-\kern0pt{}-+++-\kern0pt{}-\kern0pt{}-}
\RouteGroup{Retornos, memoria y transporte}
\RouteRow{Retornos con fase}{clase: 27; órbita: 27; celda: 27}
\RouteRow{Retorno cinemático}{en 108: 54; extensión: 216}
\RouteRow{Giros y cambios de hoja}{71; 107}
\RouteRow{Acarreo y diferimiento}{deuda total: 2; final: 0; bloques diferidos: 1}
\RouteRow{Tríadas finales; persistencia de Pareto}{17; sí}
\RouteRow{\(K\longrightarrow K+9\)}{repeticiones: 9; cambios de memoria: 9}
\RouteGroup{Firma, fase y diferencias}
\RouteRow{\(n_A,\ s_C,\ \kappa_0,\ \nu_{120},\ \nu_{270}\)}{30, 6, 0, 0, -12}
\RouteRow{\(\tau_{12}\); firma histórica \(\mathbb Z^5\)}{+++-\kern0pt{}-\kern0pt{}-+++-\kern0pt{}-\kern0pt{}-; 30|\allowbreak{}6|\allowbreak{}0|\allowbreak{}0|\allowbreak{}-12}
\RouteRow{\(D_1,\ D_3,\ D_4,\ D_{27},\ D_{36}\)}{54, 56, 70, 48, 32}
\RouteRow{\(\Omega_{90/120},\ P_6\)}{80, 6}
\RouteGroup{Lectores y factores de acción}
\RouteRow{\(K_D\quad /\quad K_\Omega\)}{(80,\allowbreak{}54,\allowbreak{}7)\quad /\quad (80,\allowbreak{}54,\allowbreak{}6)}
\RouteRow{\(\kappa_D\)}{79.60672\allowbreak{}13362175\allowbreak{}19726524\allowbreak{}96}
\RouteRow{\(\kappa_\Omega\)}{79.60661\allowbreak{}01397948\allowbreak{}27586470\allowbreak{}91}
\RouteRow{\(R_{\mathrm{act},D}\)}{1.000000\allowbreak{}05518988\allowbreak{}96223153\allowbreak{}37}
\RouteRow{\(R_{\mathrm{act},\Omega}\)}{1.000000\allowbreak{}05518981\allowbreak{}25318591\allowbreak{}40}
\RouteGroup{Separación de prefijos}
\RouteRow{Área de ambigüedad; área \(\log_2\)}{58; 30.18947\allowbreak{}50100961\allowbreak{}74}
\RouteRow{Primer bloque de separación}{órbita: \textemdash{}; celda: \textemdash{}; ruta: \textemdash{}}
\RouteRow{Ambigüedad final}{3}
\RouteGroup{Secuencias completas}
\RouteRow{\(w_6\)}{657 24 243 657 18 15 486 423 15 657 24 243 657 18 15 486 423 15}
\RouteRow{Tríadas estables por bloque}{0 1 2 3 4 5 6 7 8 9 10 11 12 13 14 14 15 17}
\end{tabularx}\endgroup\par\vspace{5pt}\noindent{\fontsize{8.4}{10}\selectfont\hyperref[trace-28]{Procedencia y códigos originales}\hfill Ficha completa · 028/324}
\clearpage\phantomsection\label{nav:vi-ruta-29}
% VI-CATALOG-ROW rutas 29
\pdfbookmark[2]{Celda 8}{cell-8}
\pdfbookmark[3]{r1c8\_h-1v-1}{route-29}
\RouteTitle{Ruta 029}{r1c8\_\allowbreak{}h-1v-1}
\noindent Celda 8\quad (1, 8)\qquad Orientación \(h=-1,\ v=-1\)\hfill\hyperref[sec:vi-indice-rutas]{Índice}\par\vspace{4pt}
\begingroup\fontsize{10.2}{12.5}\selectfont\setlength{\tabcolsep}{5pt}\renewcommand{\arraystretch}{1.1}\rowcolors{1}{routepale}{white}\noindent\begin{tabularx}{\linewidth}{@{}>{\raggedright\arraybackslash}p{.345\linewidth}>{\raggedright\arraybackslash}X@{}}
\RouteGroup{Identidad estructural}
\RouteRow{Clase y multisección}{neutrino\_\allowbreak{}G4\quad /\quad neutrino\_\allowbreak{}G4}
\RouteRow{Colección y sector}{neutrino; leptón}
\RouteRow{Clasificación física}{neutrino}
\RouteRow{Generación, profundidad y color}{4; G4; \textemdash{}}
\RouteRow{Ruta primaria y órbita de inversión}{\(B_1\): no\qquad 1,\allowbreak{}8;\allowbreak{}9,\allowbreak{}2}
\RouteRow{Firma de clasificación}{28|\allowbreak{}6|\allowbreak{}2|\allowbreak{}0|\allowbreak{}-12|\allowbreak{}+++-\kern0pt{}-\kern0pt{}-+++-\kern0pt{}-\kern0pt{}-}
\RouteGroup{Retornos, memoria y transporte}
\RouteRow{Retornos con fase}{clase: 27; órbita: 27; celda: 27}
\RouteRow{Retorno cinemático}{en 108: 54; extensión: 216}
\RouteRow{Giros y cambios de hoja}{71; 107}
\RouteRow{Acarreo y diferimiento}{deuda total: 1; final: 0; bloques diferidos: 1}
\RouteRow{Tríadas finales; persistencia de Pareto}{17; sí}
\RouteRow{\(K\longrightarrow K+9\)}{repeticiones: 9; cambios de memoria: 9}
\RouteGroup{Firma, fase y diferencias}
\RouteRow{\(n_A,\ s_C,\ \kappa_0,\ \nu_{120},\ \nu_{270}\)}{28, -6, 0, 0, -12}
\RouteRow{\(\tau_{12}\); firma histórica \(\mathbb Z^5\)}{+++-\kern0pt{}-\kern0pt{}-+++-\kern0pt{}-\kern0pt{}-; 28|\allowbreak{}-6|\allowbreak{}0|\allowbreak{}0|\allowbreak{}-12}
\RouteRow{\(D_1,\ D_3,\ D_4,\ D_{27},\ D_{36}\)}{66, 64, 64, 60, 40}
\RouteRow{\(\Omega_{90/120},\ P_6\)}{80, 6}
\RouteGroup{Lectores y factores de acción}
\RouteRow{\(K_D\quad /\quad K_\Omega\)}{(100,\allowbreak{}66,\allowbreak{}0)\quad /\quad (80,\allowbreak{}54,\allowbreak{}6)}
\RouteRow{\(\kappa_D\)}{99.51837\allowbreak{}47304272\allowbreak{}69134179\allowbreak{}18}
\RouteRow{\(\kappa_\Omega\)}{79.60661\allowbreak{}01397948\allowbreak{}27586470\allowbreak{}91}
\RouteRow{\(R_{\mathrm{act},D}\)}{1.000000\allowbreak{}06899427\allowbreak{}66450218\allowbreak{}95}
\RouteRow{\(R_{\mathrm{act},\Omega}\)}{1.000000\allowbreak{}05518981\allowbreak{}25318591\allowbreak{}40}
\RouteGroup{Separación de prefijos}
\RouteRow{Área de ambigüedad; área \(\log_2\)}{36; 18.0}
\RouteRow{Primer bloque de separación}{órbita: \textemdash{}; celda: \textemdash{}; ruta: \textemdash{}}
\RouteRow{Ambigüedad final}{2}
\RouteGroup{Secuencias completas}
\RouteRow{\(w_6\)}{90 570 264 90 567 657 24 243 657 90 570 264 90 567 657 24 243 657}
\RouteRow{Tríadas estables por bloque}{0 1 2 3 4 5 6 7 8 9 10 11 12 13 13 15 16 17}
\end{tabularx}\endgroup\par\vspace{5pt}\noindent{\fontsize{8.4}{10}\selectfont\hyperref[trace-29]{Procedencia y códigos originales}\hfill Ficha completa · 029/324}
\clearpage\phantomsection\label{nav:vi-ruta-30}
% VI-CATALOG-ROW rutas 30
\pdfbookmark[3]{r1c8\_h-1v+1}{route-30}
\RouteTitle{Ruta 030}{r1c8\_\allowbreak{}h-1v+1}
\noindent Celda 8\quad (1, 8)\qquad Orientación \(h=-1,\ v=1\)\hfill\hyperref[sec:vi-indice-rutas]{Índice}\par\vspace{4pt}
\begingroup\fontsize{10.2}{12.5}\selectfont\setlength{\tabcolsep}{5pt}\renewcommand{\arraystretch}{1.1}\rowcolors{1}{routepale}{white}\noindent\begin{tabularx}{\linewidth}{@{}>{\raggedright\arraybackslash}p{.345\linewidth}>{\raggedright\arraybackslash}X@{}}
\RouteGroup{Identidad estructural}
\RouteRow{Clase y multisección}{neutrino\_\allowbreak{}G4\quad /\quad neutrino\_\allowbreak{}G4}
\RouteRow{Colección y sector}{neutrino; leptón}
\RouteRow{Clasificación física}{neutrino}
\RouteRow{Generación, profundidad y color}{4; G4; \textemdash{}}
\RouteRow{Ruta primaria y órbita de inversión}{\(B_1\): no\qquad 1,\allowbreak{}8;\allowbreak{}9,\allowbreak{}2}
\RouteRow{Firma de clasificación}{28|\allowbreak{}6|\allowbreak{}2|\allowbreak{}0|\allowbreak{}-12|\allowbreak{}+++-\kern0pt{}-\kern0pt{}-+++-\kern0pt{}-\kern0pt{}-}
\RouteGroup{Retornos, memoria y transporte}
\RouteRow{Retornos con fase}{clase: 27; órbita: 27; celda: 27}
\RouteRow{Retorno cinemático}{en 108: 54; extensión: 216}
\RouteRow{Giros y cambios de hoja}{71; 107}
\RouteRow{Acarreo y diferimiento}{deuda total: 2; final: 0; bloques diferidos: 2}
\RouteRow{Tríadas finales; persistencia de Pareto}{17; sí}
\RouteRow{\(K\longrightarrow K+9\)}{repeticiones: 9; cambios de memoria: 9}
\RouteGroup{Firma, fase y diferencias}
\RouteRow{\(n_A,\ s_C,\ \kappa_0,\ \nu_{120},\ \nu_{270}\)}{20, 0, 0, 0, -4}
\RouteRow{\(\tau_{12}\); firma histórica \(\mathbb Z^5\)}{++00-\kern0pt{}-++00-\kern0pt{}-; 20|\allowbreak{}0|\allowbreak{}0|\allowbreak{}0|\allowbreak{}-4}
\RouteRow{\(D_1,\ D_3,\ D_4,\ D_{27},\ D_{36}\)}{24, 20, 24, 24, 16}
\RouteRow{\(\Omega_{90/120},\ P_6\)}{36, 4}
\RouteGroup{Lectores y factores de acción}
\RouteRow{\(K_D\quad /\quad K_\Omega\)}{(40,\allowbreak{}24,\allowbreak{}2)\quad /\quad (36,\allowbreak{}54,\allowbreak{}4)}
\RouteRow{\(\kappa_D\)}{39.82508\allowbreak{}59311825\allowbreak{}73056173\allowbreak{}26}
\RouteRow{\(\kappa_\Omega\)}{35.60638\allowbreak{}77469494\allowbreak{}43306362\allowbreak{}81}
\RouteRow{\(R_{\mathrm{act},D}\)}{1.000000\allowbreak{}02761000\allowbreak{}61595142\allowbreak{}15}
\RouteRow{\(R_{\mathrm{act},\Omega}\)}{1.000000\allowbreak{}02468525\allowbreak{}95691181\allowbreak{}52}
\RouteGroup{Separación de prefijos}
\RouteRow{Área de ambigüedad; área \(\log_2\)}{36; 18.0}
\RouteRow{Primer bloque de separación}{órbita: \textemdash{}; celda: \textemdash{}; ruta: \textemdash{}}
\RouteRow{Ambigüedad final}{2}
\RouteGroup{Secuencias completas}
\RouteRow{\(w_6\)}{81 3 0 81 0 3 0 81 3 81 3 0 81 0 3 0 81 3}
\RouteRow{Tríadas estables por bloque}{0 1 2 3 4 5 6 6 8 9 10 11 12 13 14 15 15 17}
\end{tabularx}\endgroup\par\vspace{5pt}\noindent{\fontsize{8.4}{10}\selectfont\hyperref[trace-30]{Procedencia y códigos originales}\hfill Ficha completa · 030/324}
\clearpage\phantomsection\label{nav:vi-ruta-31}
% VI-CATALOG-ROW rutas 31
\pdfbookmark[3]{r1c8\_h+1v-1}{route-31}
\RouteTitle{Ruta 031}{r1c8\_\allowbreak{}h+1v-1}
\noindent Celda 8\quad (1, 8)\qquad Orientación \(h=1,\ v=-1\)\hfill\hyperref[sec:vi-indice-rutas]{Índice}\par\vspace{4pt}
\begingroup\fontsize{10.2}{12.5}\selectfont\setlength{\tabcolsep}{5pt}\renewcommand{\arraystretch}{1.1}\rowcolors{1}{routepale}{white}\noindent\begin{tabularx}{\linewidth}{@{}>{\raggedright\arraybackslash}p{.345\linewidth}>{\raggedright\arraybackslash}X@{}}
\RouteGroup{Identidad estructural}
\RouteRow{Clase y multisección}{neutrino\_\allowbreak{}G4\quad /\quad neutrino\_\allowbreak{}G4}
\RouteRow{Colección y sector}{neutrino; leptón}
\RouteRow{Clasificación física}{neutrino}
\RouteRow{Generación, profundidad y color}{4; G4; \textemdash{}}
\RouteRow{Ruta primaria y órbita de inversión}{\(B_1\): sí\qquad 1,\allowbreak{}8;\allowbreak{}9,\allowbreak{}2}
\RouteRow{Firma de clasificación}{28|\allowbreak{}6|\allowbreak{}2|\allowbreak{}0|\allowbreak{}-12|\allowbreak{}+++-\kern0pt{}-\kern0pt{}-+++-\kern0pt{}-\kern0pt{}-}
\RouteGroup{Retornos, memoria y transporte}
\RouteRow{Retornos con fase}{clase: 27; órbita: 27; celda: 27}
\RouteRow{Retorno cinemático}{en 108: 54; extensión: 216}
\RouteRow{Giros y cambios de hoja}{71; 107}
\RouteRow{Acarreo y diferimiento}{deuda total: 0; final: 0; bloques diferidos: 0}
\RouteRow{Tríadas finales; persistencia de Pareto}{17; sí}
\RouteRow{\(K\longrightarrow K+9\)}{repeticiones: 9; cambios de memoria: 9}
\RouteGroup{Firma, fase y diferencias}
\RouteRow{\(n_A,\ s_C,\ \kappa_0,\ \nu_{120},\ \nu_{270}\)}{20, 0, 0, 0, -4}
\RouteRow{\(\tau_{12}\); firma histórica \(\mathbb Z^5\)}{0++-\kern0pt{}-00++-\kern0pt{}-0; 20|\allowbreak{}0|\allowbreak{}0|\allowbreak{}0|\allowbreak{}-4}
\RouteRow{\(D_1,\ D_3,\ D_4,\ D_{27},\ D_{36}\)}{24, 20, 24, 24, 16}
\RouteRow{\(\Omega_{90/120},\ P_6\)}{36, 4}
\RouteGroup{Lectores y factores de acción}
\RouteRow{\(K_D\quad /\quad K_\Omega\)}{(40,\allowbreak{}24,\allowbreak{}2)\quad /\quad (36,\allowbreak{}54,\allowbreak{}4)}
\RouteRow{\(\kappa_D\)}{39.82508\allowbreak{}59311825\allowbreak{}73056173\allowbreak{}26}
\RouteRow{\(\kappa_\Omega\)}{35.60638\allowbreak{}77469494\allowbreak{}43306362\allowbreak{}81}
\RouteRow{\(R_{\mathrm{act},D}\)}{1.000000\allowbreak{}02761000\allowbreak{}61595142\allowbreak{}15}
\RouteRow{\(R_{\mathrm{act},\Omega}\)}{1.000000\allowbreak{}02468525\allowbreak{}95691181\allowbreak{}52}
\RouteGroup{Separación de prefijos}
\RouteRow{Área de ambigüedad; área \(\log_2\)}{36; 18.0}
\RouteRow{Primer bloque de separación}{órbita: \textemdash{}; celda: \textemdash{}; ruta: \textemdash{}}
\RouteRow{Ambigüedad final}{2}
\RouteGroup{Secuencias completas}
\RouteRow{\(w_6\)}{0 81 3 0 84 0 81 3 0 0 81 3 0 84 0 81 3 0}
\RouteRow{Tríadas estables por bloque}{0 1 2 3 4 5 6 7 8 9 10 11 12 13 14 15 16 17}
\end{tabularx}\endgroup\par\vspace{5pt}\noindent{\fontsize{8.4}{10}\selectfont\hyperref[trace-31]{Procedencia y códigos originales}\hfill Ficha completa · 031/324}
\clearpage\phantomsection\label{nav:vi-ruta-32}
% VI-CATALOG-ROW rutas 32
\pdfbookmark[3]{r1c8\_h+1v+1}{route-32}
\RouteTitle{Ruta 032}{r1c8\_\allowbreak{}h+1v+1}
\noindent Celda 8\quad (1, 8)\qquad Orientación \(h=1,\ v=1\)\hfill\hyperref[sec:vi-indice-rutas]{Índice}\par\vspace{4pt}
\begingroup\fontsize{10.2}{12.5}\selectfont\setlength{\tabcolsep}{5pt}\renewcommand{\arraystretch}{1.1}\rowcolors{1}{routepale}{white}\noindent\begin{tabularx}{\linewidth}{@{}>{\raggedright\arraybackslash}p{.345\linewidth}>{\raggedright\arraybackslash}X@{}}
\RouteGroup{Identidad estructural}
\RouteRow{Clase y multisección}{neutrino\_\allowbreak{}G4\quad /\quad neutrino\_\allowbreak{}G4}
\RouteRow{Colección y sector}{neutrino; leptón}
\RouteRow{Clasificación física}{neutrino}
\RouteRow{Generación, profundidad y color}{4; G4; \textemdash{}}
\RouteRow{Ruta primaria y órbita de inversión}{\(B_1\): no\qquad 1,\allowbreak{}8;\allowbreak{}9,\allowbreak{}2}
\RouteRow{Firma de clasificación}{28|\allowbreak{}6|\allowbreak{}2|\allowbreak{}0|\allowbreak{}-12|\allowbreak{}+++-\kern0pt{}-\kern0pt{}-+++-\kern0pt{}-\kern0pt{}-}
\RouteGroup{Retornos, memoria y transporte}
\RouteRow{Retornos con fase}{clase: 27; órbita: 27; celda: 27}
\RouteRow{Retorno cinemático}{en 108: 54; extensión: 216}
\RouteRow{Giros y cambios de hoja}{71; 107}
\RouteRow{Acarreo y diferimiento}{deuda total: 0; final: 0; bloques diferidos: 0}
\RouteRow{Tríadas finales; persistencia de Pareto}{17; sí}
\RouteRow{\(K\longrightarrow K+9\)}{repeticiones: 9; cambios de memoria: 9}
\RouteGroup{Firma, fase y diferencias}
\RouteRow{\(n_A,\ s_C,\ \kappa_0,\ \nu_{120},\ \nu_{270}\)}{28, 6, 0, 0, -12}
\RouteRow{\(\tau_{12}\); firma histórica \(\mathbb Z^5\)}{+++-\kern0pt{}-\kern0pt{}-+++-\kern0pt{}-\kern0pt{}-; 28|\allowbreak{}6|\allowbreak{}0|\allowbreak{}0|\allowbreak{}-12}
\RouteRow{\(D_1,\ D_3,\ D_4,\ D_{27},\ D_{36}\)}{66, 64, 64, 60, 40}
\RouteRow{\(\Omega_{90/120},\ P_6\)}{80, 6}
\RouteGroup{Lectores y factores de acción}
\RouteRow{\(K_D\quad /\quad K_\Omega\)}{(100,\allowbreak{}66,\allowbreak{}0)\quad /\quad (80,\allowbreak{}54,\allowbreak{}6)}
\RouteRow{\(\kappa_D\)}{99.51837\allowbreak{}47304272\allowbreak{}69134179\allowbreak{}18}
\RouteRow{\(\kappa_\Omega\)}{79.60661\allowbreak{}01397948\allowbreak{}27586470\allowbreak{}91}
\RouteRow{\(R_{\mathrm{act},D}\)}{1.000000\allowbreak{}06899427\allowbreak{}66450218\allowbreak{}95}
\RouteRow{\(R_{\mathrm{act},\Omega}\)}{1.000000\allowbreak{}05518981\allowbreak{}25318591\allowbreak{}40}
\RouteGroup{Separación de prefijos}
\RouteRow{Área de ambigüedad; área \(\log_2\)}{40; 20.33985\allowbreak{}00028846\allowbreak{}24}
\RouteRow{Primer bloque de separación}{órbita: \textemdash{}; celda: \textemdash{}; ruta: \textemdash{}}
\RouteRow{Ambigüedad final}{2}
\RouteGroup{Secuencias completas}
\RouteRow{\(w_6\)}{24 243 657 24 246 264 90 570 264 24 243 657 24 246 264 90 570 264}
\RouteRow{Tríadas estables por bloque}{0 1 2 3 4 5 6 7 8 9 10 11 12 13 14 15 16 17}
\end{tabularx}\endgroup\par\vspace{5pt}\noindent{\fontsize{8.4}{10}\selectfont\hyperref[trace-32]{Procedencia y códigos originales}\hfill Ficha completa · 032/324}
\clearpage\phantomsection\label{nav:vi-ruta-33}
% VI-CATALOG-ROW rutas 33
\pdfbookmark[2]{Celda 9}{cell-9}
\pdfbookmark[3]{r1c9\_h-1v-1}{route-33}
\RouteTitle{Ruta 033}{r1c9\_\allowbreak{}h-1v-1}
\noindent Celda 9\quad (1, 9)\qquad Orientación \(h=-1,\ v=-1\)\hfill\hyperref[sec:vi-indice-rutas]{Índice}\par\vspace{4pt}
\begingroup\fontsize{10.2}{12.5}\selectfont\setlength{\tabcolsep}{5pt}\renewcommand{\arraystretch}{1.1}\rowcolors{1}{routepale}{white}\noindent\begin{tabularx}{\linewidth}{@{}>{\raggedright\arraybackslash}p{.345\linewidth}>{\raggedright\arraybackslash}X@{}}
\RouteGroup{Identidad estructural}
\RouteRow{Clase y multisección}{charged\_\allowbreak{}lepton\_\allowbreak{}G4\quad /\quad charged\_\allowbreak{}lepton\_\allowbreak{}G4}
\RouteRow{Colección y sector}{leptón cargado; leptón}
\RouteRow{Clasificación física}{leptón cargado}
\RouteRow{Generación, profundidad y color}{4; G4; \textemdash{}}
\RouteRow{Ruta primaria y órbita de inversión}{\(B_1\): no\qquad 1,\allowbreak{}9;\allowbreak{}9,\allowbreak{}1}
\RouteRow{Firma de clasificación}{28|\allowbreak{}-6|\allowbreak{}4|\allowbreak{}0|\allowbreak{}-12|\allowbreak{}+++-\kern0pt{}-\kern0pt{}-+++-\kern0pt{}-\kern0pt{}-}
\RouteGroup{Retornos, memoria y transporte}
\RouteRow{Retornos con fase}{clase: 27; órbita: 27; celda: 27}
\RouteRow{Retorno cinemático}{en 108: 54; extensión: 216}
\RouteRow{Giros y cambios de hoja}{71; 107}
\RouteRow{Acarreo y diferimiento}{deuda total: 1; final: 1; bloques diferidos: 1}
\RouteRow{Tríadas finales; persistencia de Pareto}{16; no}
\RouteRow{\(K\longrightarrow K+9\)}{repeticiones: 9; cambios de memoria: 9}
\RouteGroup{Firma, fase y diferencias}
\RouteRow{\(n_A,\ s_C,\ \kappa_0,\ \nu_{120},\ \nu_{270}\)}{28, 6, 0, 2, -12}
\RouteRow{\(\tau_{12}\); firma histórica \(\mathbb Z^5\)}{+++-\kern0pt{}-\kern0pt{}-+++-\kern0pt{}-\kern0pt{}-; 28|\allowbreak{}6|\allowbreak{}0|\allowbreak{}2|\allowbreak{}-12}
\RouteRow{\(D_1,\ D_3,\ D_4,\ D_{27},\ D_{36}\)}{66, 60, 64, 48, 32}
\RouteRow{\(\Omega_{90/120},\ P_6\)}{80, 6}
\RouteGroup{Lectores y factores de acción}
\RouteRow{\(K_D\quad /\quad K_\Omega\)}{(80,\allowbreak{}66,\allowbreak{}2)\quad /\quad (80,\allowbreak{}54,\allowbreak{}6)}
\RouteRow{\(\kappa_D\)}{79.51859\allowbreak{}71232726\allowbreak{}53414287\allowbreak{}29}
\RouteRow{\(\kappa_\Omega\)}{79.60661\allowbreak{}01397948\allowbreak{}27586470\allowbreak{}91}
\RouteRow{\(R_{\mathrm{act},D}\)}{1.000000\allowbreak{}05512879\allowbreak{}47092488\allowbreak{}80}
\RouteRow{\(R_{\mathrm{act},\Omega}\)}{1.000000\allowbreak{}05518981\allowbreak{}25318591\allowbreak{}40}
\RouteGroup{Separación de prefijos}
\RouteRow{Área de ambigüedad; área \(\log_2\)}{54; 28.52932\allowbreak{}50129808}
\RouteRow{Primer bloque de separación}{órbita: \textemdash{}; celda: \textemdash{}; ruta: \textemdash{}}
\RouteRow{Ambigüedad final}{3}
\RouteGroup{Secuencias completas}
\RouteRow{\(w_6\)}{423 15 486 423 12 99 327 498 99 423 15 486 423 12 99 327 498 99}
\RouteRow{Tríadas estables por bloque}{0 1 2 3 4 5 6 7 8 9 10 11 12 13 14 15 16 16}
\end{tabularx}\endgroup\par\vspace{5pt}\noindent{\fontsize{8.4}{10}\selectfont\hyperref[trace-33]{Procedencia y códigos originales}\hfill Ficha completa · 033/324}
\clearpage\phantomsection\label{nav:vi-ruta-34}
% VI-CATALOG-ROW rutas 34
\pdfbookmark[3]{r1c9\_h-1v+1}{route-34}
\RouteTitle{Ruta 034}{r1c9\_\allowbreak{}h-1v+1}
\noindent Celda 9\quad (1, 9)\qquad Orientación \(h=-1,\ v=1\)\hfill\hyperref[sec:vi-indice-rutas]{Índice}\par\vspace{4pt}
\begingroup\fontsize{10.2}{12.5}\selectfont\setlength{\tabcolsep}{5pt}\renewcommand{\arraystretch}{1.1}\rowcolors{1}{routepale}{white}\noindent\begin{tabularx}{\linewidth}{@{}>{\raggedright\arraybackslash}p{.345\linewidth}>{\raggedright\arraybackslash}X@{}}
\RouteGroup{Identidad estructural}
\RouteRow{Clase y multisección}{charged\_\allowbreak{}lepton\_\allowbreak{}G4\quad /\quad charged\_\allowbreak{}lepton\_\allowbreak{}G4}
\RouteRow{Colección y sector}{leptón cargado; leptón}
\RouteRow{Clasificación física}{leptón cargado}
\RouteRow{Generación, profundidad y color}{4; G4; \textemdash{}}
\RouteRow{Ruta primaria y órbita de inversión}{\(B_1\): no\qquad 1,\allowbreak{}9;\allowbreak{}9,\allowbreak{}1}
\RouteRow{Firma de clasificación}{28|\allowbreak{}-6|\allowbreak{}4|\allowbreak{}0|\allowbreak{}-12|\allowbreak{}+++-\kern0pt{}-\kern0pt{}-+++-\kern0pt{}-\kern0pt{}-}
\RouteGroup{Retornos, memoria y transporte}
\RouteRow{Retornos con fase}{clase: 27; órbita: 27; celda: 27}
\RouteRow{Retorno cinemático}{en 108: 54; extensión: 216}
\RouteRow{Giros y cambios de hoja}{71; 107}
\RouteRow{Acarreo y diferimiento}{deuda total: 1; final: 0; bloques diferidos: 1}
\RouteRow{Tríadas finales; persistencia de Pareto}{17; no}
\RouteRow{\(K\longrightarrow K+9\)}{repeticiones: 9; cambios de memoria: 9}
\RouteGroup{Firma, fase y diferencias}
\RouteRow{\(n_A,\ s_C,\ \kappa_0,\ \nu_{120},\ \nu_{270}\)}{50, 4, -2, 2, -8}
\RouteRow{\(\tau_{12}\); firma histórica \(\mathbb Z^5\)}{+++-0-+++-0-; 50|\allowbreak{}4|\allowbreak{}-2|\allowbreak{}2|\allowbreak{}-8}
\RouteRow{\(D_1,\ D_3,\ D_4,\ D_{27},\ D_{36}\)}{84, 24, 84, 24, 16}
\RouteRow{\(\Omega_{90/120},\ P_6\)}{48, 5}
\RouteGroup{Lectores y factores de acción}
\RouteRow{\(K_D\quad /\quad K_\Omega\)}{(40,\allowbreak{}84,\allowbreak{}30)\quad /\quad (48,\allowbreak{}54,\allowbreak{}5)}
\RouteRow{\(\kappa_D\)}{39.39035\allowbreak{}82768609\allowbreak{}24917849\allowbreak{}59}
\RouteRow{\(\kappa_\Omega\)}{47.60649\allowbreak{}89433721\allowbreak{}35446416\allowbreak{}86}
\RouteRow{\(R_{\mathrm{act},D}\)}{1.000000\allowbreak{}02730861\allowbreak{}73967087\allowbreak{}05}
\RouteRow{\(R_{\mathrm{act},\Omega}\)}{1.000000\allowbreak{}03300471\allowbreak{}80532430\allowbreak{}84}
\RouteGroup{Separación de prefijos}
\RouteRow{Área de ambigüedad; área \(\log_2\)}{54; 28.52932\allowbreak{}50129808}
\RouteRow{Primer bloque de separación}{órbita: \textemdash{}; celda: \textemdash{}; ruta: \textemdash{}}
\RouteRow{Ambigüedad final}{3}
\RouteGroup{Secuencias completas}
\RouteRow{\(w_6\)}{420 258 414 420 255 252 333 255 252 420 258 414 420 255 252 333 255 252}
\RouteRow{Tríadas estables por bloque}{0 1 2 3 4 5 6 7 8 9 10 11 12 13 14 15 15 17}
\end{tabularx}\endgroup\par\vspace{5pt}\noindent{\fontsize{8.4}{10}\selectfont\hyperref[trace-34]{Procedencia y códigos originales}\hfill Ficha completa · 034/324}
\clearpage\phantomsection\label{nav:vi-ruta-35}
% VI-CATALOG-ROW rutas 35
\pdfbookmark[3]{r1c9\_h+1v-1}{route-35}
\RouteTitle{Ruta 035}{r1c9\_\allowbreak{}h+1v-1}
\noindent Celda 9\quad (1, 9)\qquad Orientación \(h=1,\ v=-1\)\hfill\hyperref[sec:vi-indice-rutas]{Índice}\par\vspace{4pt}
\begingroup\fontsize{10.2}{12.5}\selectfont\setlength{\tabcolsep}{5pt}\renewcommand{\arraystretch}{1.1}\rowcolors{1}{routepale}{white}\noindent\begin{tabularx}{\linewidth}{@{}>{\raggedright\arraybackslash}p{.345\linewidth}>{\raggedright\arraybackslash}X@{}}
\RouteGroup{Identidad estructural}
\RouteRow{Clase y multisección}{charged\_\allowbreak{}lepton\_\allowbreak{}G4\quad /\quad charged\_\allowbreak{}lepton\_\allowbreak{}G4}
\RouteRow{Colección y sector}{leptón cargado; leptón}
\RouteRow{Clasificación física}{leptón cargado}
\RouteRow{Generación, profundidad y color}{4; G4; \textemdash{}}
\RouteRow{Ruta primaria y órbita de inversión}{\(B_1\): no\qquad 1,\allowbreak{}9;\allowbreak{}9,\allowbreak{}1}
\RouteRow{Firma de clasificación}{28|\allowbreak{}-6|\allowbreak{}4|\allowbreak{}0|\allowbreak{}-12|\allowbreak{}+++-\kern0pt{}-\kern0pt{}-+++-\kern0pt{}-\kern0pt{}-}
\RouteGroup{Retornos, memoria y transporte}
\RouteRow{Retornos con fase}{clase: 27; órbita: 27; celda: 27}
\RouteRow{Retorno cinemático}{en 108: 54; extensión: 216}
\RouteRow{Giros y cambios de hoja}{71; 107}
\RouteRow{Acarreo y diferimiento}{deuda total: 0; final: 0; bloques diferidos: 0}
\RouteRow{Tríadas finales; persistencia de Pareto}{17; no}
\RouteRow{\(K\longrightarrow K+9\)}{repeticiones: 9; cambios de memoria: 9}
\RouteGroup{Firma, fase y diferencias}
\RouteRow{\(n_A,\ s_C,\ \kappa_0,\ \nu_{120},\ \nu_{270}\)}{50, -4, 2, -2, -8}
\RouteRow{\(\tau_{12}\); firma histórica \(\mathbb Z^5\)}{+0+-\kern0pt{}-\kern0pt{}-+0+-\kern0pt{}-\kern0pt{}-; 50|\allowbreak{}-4|\allowbreak{}2|\allowbreak{}-2|\allowbreak{}-8}
\RouteRow{\(D_1,\ D_3,\ D_4,\ D_{27},\ D_{36}\)}{84, 24, 84, 24, 16}
\RouteRow{\(\Omega_{90/120},\ P_6\)}{48, 5}
\RouteGroup{Lectores y factores de acción}
\RouteRow{\(K_D\quad /\quad K_\Omega\)}{(40,\allowbreak{}84,\allowbreak{}30)\quad /\quad (48,\allowbreak{}54,\allowbreak{}5)}
\RouteRow{\(\kappa_D\)}{39.39035\allowbreak{}82768609\allowbreak{}24917849\allowbreak{}59}
\RouteRow{\(\kappa_\Omega\)}{47.60649\allowbreak{}89433721\allowbreak{}35446416\allowbreak{}86}
\RouteRow{\(R_{\mathrm{act},D}\)}{1.000000\allowbreak{}02730861\allowbreak{}73967087\allowbreak{}05}
\RouteRow{\(R_{\mathrm{act},\Omega}\)}{1.000000\allowbreak{}03300471\allowbreak{}80532430\allowbreak{}84}
\RouteGroup{Separación de prefijos}
\RouteRow{Área de ambigüedad; área \(\log_2\)}{54; 28.52932\allowbreak{}50129808}
\RouteRow{Primer bloque de separación}{órbita: \textemdash{}; celda: \textemdash{}; ruta: \textemdash{}}
\RouteRow{Ambigüedad final}{3}
\RouteGroup{Secuencias completas}
\RouteRow{\(w_6\)}{333 255 252 333 258 414 420 258 414 333 255 252 333 258 414 420 258 414}
\RouteRow{Tríadas estables por bloque}{0 1 2 3 4 5 6 7 8 9 10 11 12 13 14 15 16 17}
\end{tabularx}\endgroup\par\vspace{5pt}\noindent{\fontsize{8.4}{10}\selectfont\hyperref[trace-35]{Procedencia y códigos originales}\hfill Ficha completa · 035/324}
\clearpage\phantomsection\label{nav:vi-ruta-36}
% VI-CATALOG-ROW rutas 36
\pdfbookmark[3]{r1c9\_h+1v+1}{route-36}
\RouteTitle{Ruta 036}{r1c9\_\allowbreak{}h+1v+1}
\noindent Celda 9\quad (1, 9)\qquad Orientación \(h=1,\ v=1\)\hfill\hyperref[sec:vi-indice-rutas]{Índice}\par\vspace{4pt}
\begingroup\fontsize{10.2}{12.5}\selectfont\setlength{\tabcolsep}{5pt}\renewcommand{\arraystretch}{1.1}\rowcolors{1}{routepale}{white}\noindent\begin{tabularx}{\linewidth}{@{}>{\raggedright\arraybackslash}p{.345\linewidth}>{\raggedright\arraybackslash}X@{}}
\RouteGroup{Identidad estructural}
\RouteRow{Clase y multisección}{charged\_\allowbreak{}lepton\_\allowbreak{}G4\quad /\quad charged\_\allowbreak{}lepton\_\allowbreak{}G4}
\RouteRow{Colección y sector}{leptón cargado; leptón}
\RouteRow{Clasificación física}{leptón cargado}
\RouteRow{Generación, profundidad y color}{4; G4; \textemdash{}}
\RouteRow{Ruta primaria y órbita de inversión}{\(B_1\): sí\qquad 1,\allowbreak{}9;\allowbreak{}9,\allowbreak{}1}
\RouteRow{Firma de clasificación}{28|\allowbreak{}-6|\allowbreak{}4|\allowbreak{}0|\allowbreak{}-12|\allowbreak{}+++-\kern0pt{}-\kern0pt{}-+++-\kern0pt{}-\kern0pt{}-}
\RouteGroup{Retornos, memoria y transporte}
\RouteRow{Retornos con fase}{clase: 27; órbita: 27; celda: 27}
\RouteRow{Retorno cinemático}{en 108: 54; extensión: 216}
\RouteRow{Giros y cambios de hoja}{71; 107}
\RouteRow{Acarreo y diferimiento}{deuda total: 2; final: 1; bloques diferidos: 2}
\RouteRow{Tríadas finales; persistencia de Pareto}{16; no}
\RouteRow{\(K\longrightarrow K+9\)}{repeticiones: 9; cambios de memoria: 9}
\RouteGroup{Firma, fase y diferencias}
\RouteRow{\(n_A,\ s_C,\ \kappa_0,\ \nu_{120},\ \nu_{270}\)}{28, -6, 0, -2, -12}
\RouteRow{\(\tau_{12}\); firma histórica \(\mathbb Z^5\)}{+++-\kern0pt{}-\kern0pt{}-+++-\kern0pt{}-\kern0pt{}-; 28|\allowbreak{}-6|\allowbreak{}0|\allowbreak{}-2|\allowbreak{}-12}
\RouteRow{\(D_1,\ D_3,\ D_4,\ D_{27},\ D_{36}\)}{66, 60, 64, 48, 32}
\RouteRow{\(\Omega_{90/120},\ P_6\)}{80, 6}
\RouteGroup{Lectores y factores de acción}
\RouteRow{\(K_D\quad /\quad K_\Omega\)}{(80,\allowbreak{}66,\allowbreak{}2)\quad /\quad (80,\allowbreak{}54,\allowbreak{}6)}
\RouteRow{\(\kappa_D\)}{79.51859\allowbreak{}71232726\allowbreak{}53414287\allowbreak{}29}
\RouteRow{\(\kappa_\Omega\)}{79.60661\allowbreak{}01397948\allowbreak{}27586470\allowbreak{}91}
\RouteRow{\(R_{\mathrm{act},D}\)}{1.000000\allowbreak{}05512879\allowbreak{}47092488\allowbreak{}80}
\RouteRow{\(R_{\mathrm{act},\Omega}\)}{1.000000\allowbreak{}05518981\allowbreak{}25318591\allowbreak{}40}
\RouteGroup{Separación de prefijos}
\RouteRow{Área de ambigüedad; área \(\log_2\)}{54; 28.52932\allowbreak{}50129808}
\RouteRow{Primer bloque de separación}{órbita: \textemdash{}; celda: \textemdash{}; ruta: \textemdash{}}
\RouteRow{Ambigüedad final}{3}
\RouteGroup{Secuencias completas}
\RouteRow{\(w_6\)}{327 498 99 327 501 486 423 15 486 327 498 99 327 501 486 423 15 486}
\RouteRow{Tríadas estables por bloque}{0 1 2 3 4 5 6 7 7 9 10 11 12 13 14 15 16 16}
\end{tabularx}\endgroup\par\vspace{5pt}\noindent{\fontsize{8.4}{10}\selectfont\hyperref[trace-36]{Procedencia y códigos originales}\hfill Ficha completa · 036/324}
\clearpage\phantomsection\label{nav:vi-ruta-37}
% VI-CATALOG-ROW rutas 37
\pdfbookmark[2]{Celda 10}{cell-10}
\pdfbookmark[3]{r2c1\_h-1v-1}{route-37}
\RouteTitle{Ruta 037}{r2c1\_\allowbreak{}h-1v-1}
\noindent Celda 10\quad (2, 1)\qquad Orientación \(h=-1,\ v=-1\)\hfill\hyperref[sec:vi-indice-rutas]{Índice}\par\vspace{4pt}
\begingroup\fontsize{10.2}{12.5}\selectfont\setlength{\tabcolsep}{5pt}\renewcommand{\arraystretch}{1.1}\rowcolors{1}{routepale}{white}\noindent\begin{tabularx}{\linewidth}{@{}>{\raggedright\arraybackslash}p{.345\linewidth}>{\raggedright\arraybackslash}X@{}}
\RouteGroup{Identidad estructural}
\RouteRow{Clase y multisección}{quark\_\allowbreak{}down\_\allowbreak{}G4\quad /\quad quark\_\allowbreak{}down\_\allowbreak{}G4}
\RouteRow{Colección y sector}{quark de tipo abajo; quark}
\RouteRow{Clasificación física}{quark de tipo abajo}
\RouteRow{Generación, profundidad y color}{4; G4; G}
\RouteRow{Ruta primaria y órbita de inversión}{\(B_1\): no\qquad 2,\allowbreak{}1;\allowbreak{}8,\allowbreak{}9}
\RouteRow{Firma de clasificación}{28|\allowbreak{}-6|\allowbreak{}12|\allowbreak{}0|\allowbreak{}-12|\allowbreak{}+++-\kern0pt{}-\kern0pt{}-+++-\kern0pt{}-\kern0pt{}-}
\RouteGroup{Retornos, memoria y transporte}
\RouteRow{Retornos con fase}{clase: 27; órbita: 27; celda: 27}
\RouteRow{Retorno cinemático}{en 108: 54; extensión: 216}
\RouteRow{Giros y cambios de hoja}{71; 107}
\RouteRow{Acarreo y diferimiento}{deuda total: 1; final: 1; bloques diferidos: 1}
\RouteRow{Tríadas finales; persistencia de Pareto}{16; sí}
\RouteRow{\(K\longrightarrow K+9\)}{repeticiones: 9; cambios de memoria: 9}
\RouteGroup{Firma, fase y diferencias}
\RouteRow{\(n_A,\ s_C,\ \kappa_0,\ \nu_{120},\ \nu_{270}\)}{28, 6, 0, 4, -12}
\RouteRow{\(\tau_{12}\); firma histórica \(\mathbb Z^5\)}{+++-\kern0pt{}-\kern0pt{}-+++-\kern0pt{}-\kern0pt{}-; 28|\allowbreak{}6|\allowbreak{}0|\allowbreak{}4|\allowbreak{}-12}
\RouteRow{\(D_1,\ D_3,\ D_4,\ D_{27},\ D_{36}\)}{66, 64, 68, 60, 40}
\RouteRow{\(\Omega_{90/120},\ P_6\)}{80, 6}
\RouteGroup{Lectores y factores de acción}
\RouteRow{\(K_D\quad /\quad K_\Omega\)}{(100,\allowbreak{}66,\allowbreak{}2)\quad /\quad (80,\allowbreak{}54,\allowbreak{}6)}
\RouteRow{\(\kappa_D\)}{99.51859\allowbreak{}71232726\allowbreak{}53414287\allowbreak{}29}
\RouteRow{\(\kappa_\Omega\)}{79.60661\allowbreak{}01397948\allowbreak{}27586470\allowbreak{}91}
\RouteRow{\(R_{\mathrm{act},D}\)}{1.000000\allowbreak{}06899443\allowbreak{}08259364\allowbreak{}17}
\RouteRow{\(R_{\mathrm{act},\Omega}\)}{1.000000\allowbreak{}05518981\allowbreak{}25318591\allowbreak{}40}
\RouteGroup{Separación de prefijos}
\RouteRow{Área de ambigüedad; área \(\log_2\)}{40; 20.33985\allowbreak{}00028846\allowbreak{}24}
\RouteRow{Primer bloque de separación}{órbita: \textemdash{}; celda: \textemdash{}; ruta: \textemdash{}}
\RouteRow{Ambigüedad final}{2}
\RouteGroup{Secuencias completas}
\RouteRow{\(w_6\)}{15 486 423 15 489 498 99 327 498 15 486 423 15 489 498 99 327 498}
\RouteRow{Tríadas estables por bloque}{0 1 2 3 4 5 6 7 8 9 10 11 12 13 14 15 16 16}
\end{tabularx}\endgroup\par\vspace{5pt}\noindent{\fontsize{8.4}{10}\selectfont\hyperref[trace-37]{Procedencia y códigos originales}\hfill Ficha completa · 037/324}
\clearpage\phantomsection\label{nav:vi-ruta-38}
% VI-CATALOG-ROW rutas 38
\pdfbookmark[3]{r2c1\_h-1v+1}{route-38}
\RouteTitle{Ruta 038}{r2c1\_\allowbreak{}h-1v+1}
\noindent Celda 10\quad (2, 1)\qquad Orientación \(h=-1,\ v=1\)\hfill\hyperref[sec:vi-indice-rutas]{Índice}\par\vspace{4pt}
\begingroup\fontsize{10.2}{12.5}\selectfont\setlength{\tabcolsep}{5pt}\renewcommand{\arraystretch}{1.1}\rowcolors{1}{routepale}{white}\noindent\begin{tabularx}{\linewidth}{@{}>{\raggedright\arraybackslash}p{.345\linewidth}>{\raggedright\arraybackslash}X@{}}
\RouteGroup{Identidad estructural}
\RouteRow{Clase y multisección}{quark\_\allowbreak{}down\_\allowbreak{}G4\quad /\quad quark\_\allowbreak{}down\_\allowbreak{}G4}
\RouteRow{Colección y sector}{quark de tipo abajo; quark}
\RouteRow{Clasificación física}{quark de tipo abajo}
\RouteRow{Generación, profundidad y color}{4; G4; G}
\RouteRow{Ruta primaria y órbita de inversión}{\(B_1\): no\qquad 2,\allowbreak{}1;\allowbreak{}8,\allowbreak{}9}
\RouteRow{Firma de clasificación}{28|\allowbreak{}-6|\allowbreak{}12|\allowbreak{}0|\allowbreak{}-12|\allowbreak{}+++-\kern0pt{}-\kern0pt{}-+++-\kern0pt{}-\kern0pt{}-}
\RouteGroup{Retornos, memoria y transporte}
\RouteRow{Retornos con fase}{clase: 27; órbita: 27; celda: 27}
\RouteRow{Retorno cinemático}{en 108: 54; extensión: 216}
\RouteRow{Giros y cambios de hoja}{71; 107}
\RouteRow{Acarreo y diferimiento}{deuda total: 0; final: 0; bloques diferidos: 0}
\RouteRow{Tríadas finales; persistencia de Pareto}{17; sí}
\RouteRow{\(K\longrightarrow K+9\)}{repeticiones: 9; cambios de memoria: 9}
\RouteGroup{Firma, fase y diferencias}
\RouteRow{\(n_A,\ s_C,\ \kappa_0,\ \nu_{120},\ \nu_{270}\)}{50, -6, 2, -4, -8}
\RouteRow{\(\tau_{12}\); firma histórica \(\mathbb Z^5\)}{+0+-\kern0pt{}-\kern0pt{}-+0+-\kern0pt{}-\kern0pt{}-; 50|\allowbreak{}-6|\allowbreak{}2|\allowbreak{}-4|\allowbreak{}-8}
\RouteRow{\(D_1,\ D_3,\ D_4,\ D_{27},\ D_{36}\)}{24, 20, 24, 24, 16}
\RouteRow{\(\Omega_{90/120},\ P_6\)}{48, 5}
\RouteGroup{Lectores y factores de acción}
\RouteRow{\(K_D\quad /\quad K_\Omega\)}{(40,\allowbreak{}24,\allowbreak{}2)\quad /\quad (48,\allowbreak{}54,\allowbreak{}5)}
\RouteRow{\(\kappa_D\)}{39.82508\allowbreak{}59311825\allowbreak{}73056173\allowbreak{}26}
\RouteRow{\(\kappa_\Omega\)}{47.60649\allowbreak{}89433721\allowbreak{}35446416\allowbreak{}86}
\RouteRow{\(R_{\mathrm{act},D}\)}{1.000000\allowbreak{}02761000\allowbreak{}61595142\allowbreak{}15}
\RouteRow{\(R_{\mathrm{act},\Omega}\)}{1.000000\allowbreak{}03300471\allowbreak{}80532430\allowbreak{}84}
\RouteGroup{Separación de prefijos}
\RouteRow{Área de ambigüedad; área \(\log_2\)}{36; 18.0}
\RouteRow{Primer bloque de separación}{órbita: \textemdash{}; celda: \textemdash{}; ruta: \textemdash{}}
\RouteRow{Ambigüedad final}{2}
\RouteGroup{Secuencias completas}
\RouteRow{\(w_6\)}{0 81 3 0 84 0 81 3 0 0 81 3 0 84 0 81 3 0}
\RouteRow{Tríadas estables por bloque}{0 1 2 3 4 5 6 7 8 9 10 11 12 13 14 15 16 17}
\end{tabularx}\endgroup\par\vspace{5pt}\noindent{\fontsize{8.4}{10}\selectfont\hyperref[trace-38]{Procedencia y códigos originales}\hfill Ficha completa · 038/324}
\clearpage\phantomsection\label{nav:vi-ruta-39}
% VI-CATALOG-ROW rutas 39
\pdfbookmark[3]{r2c1\_h+1v-1}{route-39}
\RouteTitle{Ruta 039}{r2c1\_\allowbreak{}h+1v-1}
\noindent Celda 10\quad (2, 1)\qquad Orientación \(h=1,\ v=-1\)\hfill\hyperref[sec:vi-indice-rutas]{Índice}\par\vspace{4pt}
\begingroup\fontsize{10.2}{12.5}\selectfont\setlength{\tabcolsep}{5pt}\renewcommand{\arraystretch}{1.1}\rowcolors{1}{routepale}{white}\noindent\begin{tabularx}{\linewidth}{@{}>{\raggedright\arraybackslash}p{.345\linewidth}>{\raggedright\arraybackslash}X@{}}
\RouteGroup{Identidad estructural}
\RouteRow{Clase y multisección}{quark\_\allowbreak{}down\_\allowbreak{}G4\quad /\quad quark\_\allowbreak{}down\_\allowbreak{}G4}
\RouteRow{Colección y sector}{quark de tipo abajo; quark}
\RouteRow{Clasificación física}{quark de tipo abajo}
\RouteRow{Generación, profundidad y color}{4; G4; G}
\RouteRow{Ruta primaria y órbita de inversión}{\(B_1\): no\qquad 2,\allowbreak{}1;\allowbreak{}8,\allowbreak{}9}
\RouteRow{Firma de clasificación}{28|\allowbreak{}-6|\allowbreak{}12|\allowbreak{}0|\allowbreak{}-12|\allowbreak{}+++-\kern0pt{}-\kern0pt{}-+++-\kern0pt{}-\kern0pt{}-}
\RouteGroup{Retornos, memoria y transporte}
\RouteRow{Retornos con fase}{clase: 27; órbita: 27; celda: 27}
\RouteRow{Retorno cinemático}{en 108: 54; extensión: 216}
\RouteRow{Giros y cambios de hoja}{71; 107}
\RouteRow{Acarreo y diferimiento}{deuda total: 2; final: 0; bloques diferidos: 2}
\RouteRow{Tríadas finales; persistencia de Pareto}{17; sí}
\RouteRow{\(K\longrightarrow K+9\)}{repeticiones: 9; cambios de memoria: 9}
\RouteGroup{Firma, fase y diferencias}
\RouteRow{\(n_A,\ s_C,\ \kappa_0,\ \nu_{120},\ \nu_{270}\)}{50, 6, -2, 4, -8}
\RouteRow{\(\tau_{12}\); firma histórica \(\mathbb Z^5\)}{+++-0-+++-0-; 50|\allowbreak{}6|\allowbreak{}-2|\allowbreak{}4|\allowbreak{}-8}
\RouteRow{\(D_1,\ D_3,\ D_4,\ D_{27},\ D_{36}\)}{24, 20, 24, 24, 16}
\RouteRow{\(\Omega_{90/120},\ P_6\)}{48, 5}
\RouteGroup{Lectores y factores de acción}
\RouteRow{\(K_D\quad /\quad K_\Omega\)}{(40,\allowbreak{}24,\allowbreak{}2)\quad /\quad (48,\allowbreak{}54,\allowbreak{}5)}
\RouteRow{\(\kappa_D\)}{39.82508\allowbreak{}59311825\allowbreak{}73056173\allowbreak{}26}
\RouteRow{\(\kappa_\Omega\)}{47.60649\allowbreak{}89433721\allowbreak{}35446416\allowbreak{}86}
\RouteRow{\(R_{\mathrm{act},D}\)}{1.000000\allowbreak{}02761000\allowbreak{}61595142\allowbreak{}15}
\RouteRow{\(R_{\mathrm{act},\Omega}\)}{1.000000\allowbreak{}03300471\allowbreak{}80532430\allowbreak{}84}
\RouteGroup{Separación de prefijos}
\RouteRow{Área de ambigüedad; área \(\log_2\)}{36; 18.0}
\RouteRow{Primer bloque de separación}{órbita: \textemdash{}; celda: \textemdash{}; ruta: \textemdash{}}
\RouteRow{Ambigüedad final}{2}
\RouteGroup{Secuencias completas}
\RouteRow{\(w_6\)}{81 3 0 81 0 3 0 81 3 81 3 0 81 0 3 0 81 3}
\RouteRow{Tríadas estables por bloque}{0 1 2 3 4 5 6 6 8 9 10 11 12 13 14 15 15 17}
\end{tabularx}\endgroup\par\vspace{5pt}\noindent{\fontsize{8.4}{10}\selectfont\hyperref[trace-39]{Procedencia y códigos originales}\hfill Ficha completa · 039/324}
\clearpage\phantomsection\label{nav:vi-ruta-40}
% VI-CATALOG-ROW rutas 40
\pdfbookmark[3]{r2c1\_h+1v+1}{route-40}
\RouteTitle{Ruta 040}{r2c1\_\allowbreak{}h+1v+1}
\noindent Celda 10\quad (2, 1)\qquad Orientación \(h=1,\ v=1\)\hfill\hyperref[sec:vi-indice-rutas]{Índice}\par\vspace{4pt}
\begingroup\fontsize{10.2}{12.5}\selectfont\setlength{\tabcolsep}{5pt}\renewcommand{\arraystretch}{1.1}\rowcolors{1}{routepale}{white}\noindent\begin{tabularx}{\linewidth}{@{}>{\raggedright\arraybackslash}p{.345\linewidth}>{\raggedright\arraybackslash}X@{}}
\RouteGroup{Identidad estructural}
\RouteRow{Clase y multisección}{quark\_\allowbreak{}down\_\allowbreak{}G4\quad /\quad quark\_\allowbreak{}down\_\allowbreak{}G4}
\RouteRow{Colección y sector}{quark de tipo abajo; quark}
\RouteRow{Clasificación física}{quark de tipo abajo}
\RouteRow{Generación, profundidad y color}{4; G4; G}
\RouteRow{Ruta primaria y órbita de inversión}{\(B_1\): sí\qquad 2,\allowbreak{}1;\allowbreak{}8,\allowbreak{}9}
\RouteRow{Firma de clasificación}{28|\allowbreak{}-6|\allowbreak{}12|\allowbreak{}0|\allowbreak{}-12|\allowbreak{}+++-\kern0pt{}-\kern0pt{}-+++-\kern0pt{}-\kern0pt{}-}
\RouteGroup{Retornos, memoria y transporte}
\RouteRow{Retornos con fase}{clase: 27; órbita: 27; celda: 27}
\RouteRow{Retorno cinemático}{en 108: 54; extensión: 216}
\RouteRow{Giros y cambios de hoja}{71; 107}
\RouteRow{Acarreo y diferimiento}{deuda total: 1; final: 0; bloques diferidos: 1}
\RouteRow{Tríadas finales; persistencia de Pareto}{17; sí}
\RouteRow{\(K\longrightarrow K+9\)}{repeticiones: 9; cambios de memoria: 9}
\RouteGroup{Firma, fase y diferencias}
\RouteRow{\(n_A,\ s_C,\ \kappa_0,\ \nu_{120},\ \nu_{270}\)}{28, -6, 0, -4, -12}
\RouteRow{\(\tau_{12}\); firma histórica \(\mathbb Z^5\)}{+++-\kern0pt{}-\kern0pt{}-+++-\kern0pt{}-\kern0pt{}-; 28|\allowbreak{}-6|\allowbreak{}0|\allowbreak{}-4|\allowbreak{}-12}
\RouteRow{\(D_1,\ D_3,\ D_4,\ D_{27},\ D_{36}\)}{66, 64, 68, 60, 40}
\RouteRow{\(\Omega_{90/120},\ P_6\)}{80, 6}
\RouteGroup{Lectores y factores de acción}
\RouteRow{\(K_D\quad /\quad K_\Omega\)}{(100,\allowbreak{}66,\allowbreak{}2)\quad /\quad (80,\allowbreak{}54,\allowbreak{}6)}
\RouteRow{\(\kappa_D\)}{99.51859\allowbreak{}71232726\allowbreak{}53414287\allowbreak{}29}
\RouteRow{\(\kappa_\Omega\)}{79.60661\allowbreak{}01397948\allowbreak{}27586470\allowbreak{}91}
\RouteRow{\(R_{\mathrm{act},D}\)}{1.000000\allowbreak{}06899443\allowbreak{}08259364\allowbreak{}17}
\RouteRow{\(R_{\mathrm{act},\Omega}\)}{1.000000\allowbreak{}05518981\allowbreak{}25318591\allowbreak{}40}
\RouteGroup{Separación de prefijos}
\RouteRow{Área de ambigüedad; área \(\log_2\)}{36; 18.0}
\RouteRow{Primer bloque de separación}{órbita: \textemdash{}; celda: \textemdash{}; ruta: \textemdash{}}
\RouteRow{Ambigüedad final}{2}
\RouteGroup{Secuencias completas}
\RouteRow{\(w_6\)}{99 327 498 99 324 423 15 486 423 99 327 498 99 324 423 15 486 423}
\RouteRow{Tríadas estables por bloque}{0 1 2 3 4 5 6 7 8 9 10 11 12 13 14 14 16 17}
\end{tabularx}\endgroup\par\vspace{5pt}\noindent{\fontsize{8.4}{10}\selectfont\hyperref[trace-40]{Procedencia y códigos originales}\hfill Ficha completa · 040/324}
\clearpage\phantomsection\label{nav:vi-ruta-41}
% VI-CATALOG-ROW rutas 41
\pdfbookmark[2]{Celda 11}{cell-11}
\pdfbookmark[3]{r2c2\_h-1v-1}{route-41}
\RouteTitle{Ruta 041}{r2c2\_\allowbreak{}h-1v-1}
\noindent Celda 11\quad (2, 2)\qquad Orientación \(h=-1,\ v=-1\)\hfill\hyperref[sec:vi-indice-rutas]{Índice}\par\vspace{4pt}
\begingroup\fontsize{10.2}{12.5}\selectfont\setlength{\tabcolsep}{5pt}\renewcommand{\arraystretch}{1.1}\rowcolors{1}{routepale}{white}\noindent\begin{tabularx}{\linewidth}{@{}>{\raggedright\arraybackslash}p{.345\linewidth}>{\raggedright\arraybackslash}X@{}}
\RouteGroup{Identidad estructural}
\RouteRow{Clase y multisección}{quark\_\allowbreak{}up\_\allowbreak{}G3\quad /\quad quark\_\allowbreak{}up\_\allowbreak{}G3}
\RouteRow{Colección y sector}{quark de tipo arriba; quark}
\RouteRow{Clasificación física}{quark de tipo arriba}
\RouteRow{Generación, profundidad y color}{3; G3; R}
\RouteRow{Ruta primaria y órbita de inversión}{\(B_1\): no\qquad 2,\allowbreak{}2;\allowbreak{}8,\allowbreak{}8}
\RouteRow{Firma de clasificación}{24|\allowbreak{}0|\allowbreak{}4|\allowbreak{}0|\allowbreak{}-12|\allowbreak{}+++-\kern0pt{}-\kern0pt{}-+++-\kern0pt{}-\kern0pt{}-}
\RouteGroup{Retornos, memoria y transporte}
\RouteRow{Retornos con fase}{clase: 9; órbita: 9; celda: 27}
\RouteRow{Retorno cinemático}{en 108: 54; extensión: 216}
\RouteRow{Giros y cambios de hoja}{71; 107}
\RouteRow{Acarreo y diferimiento}{deuda total: 3; final: 0; bloques diferidos: 2}
\RouteRow{Tríadas finales; persistencia de Pareto}{17; sí}
\RouteRow{\(K\longrightarrow K+9\)}{repeticiones: 9; cambios de memoria: 9}
\RouteGroup{Firma, fase y diferencias}
\RouteRow{\(n_A,\ s_C,\ \kappa_0,\ \nu_{120},\ \nu_{270}\)}{24, 0, 0, 0, -12}
\RouteRow{\(\tau_{12}\); firma histórica \(\mathbb Z^5\)}{+++-\kern0pt{}-\kern0pt{}-+++-\kern0pt{}-\kern0pt{}-; 24|\allowbreak{}0|\allowbreak{}0|\allowbreak{}0|\allowbreak{}-12}
\RouteRow{\(D_1,\ D_3,\ D_4,\ D_{27},\ D_{36}\)}{54, 56, 68, 48, 32}
\RouteRow{\(\Omega_{90/120},\ P_6\)}{80, 6}
\RouteGroup{Lectores y factores de acción}
\RouteRow{\(K_D\quad /\quad K_\Omega\)}{(80,\allowbreak{}54,\allowbreak{}6)\quad /\quad (80,\allowbreak{}54,\allowbreak{}6)}
\RouteRow{\(\kappa_D\)}{79.60661\allowbreak{}01397948\allowbreak{}27586470\allowbreak{}91}
\RouteRow{\(\kappa_\Omega\)}{79.60661\allowbreak{}01397948\allowbreak{}27586470\allowbreak{}91}
\RouteRow{\(R_{\mathrm{act},D}\)}{1.000000\allowbreak{}05518981\allowbreak{}25318591\allowbreak{}40}
\RouteRow{\(R_{\mathrm{act},\Omega}\)}{1.000000\allowbreak{}05518981\allowbreak{}25318591\allowbreak{}40}
\RouteGroup{Separación de prefijos}
\RouteRow{Área de ambigüedad; área \(\log_2\)}{72; 36.0}
\RouteRow{Primer bloque de separación}{órbita: \textemdash{}; celda: \textemdash{}; ruta: \textemdash{}}
\RouteRow{Ambigüedad final}{4}
\RouteGroup{Secuencias completas}
\RouteRow{\(w_6\)}{423 15 486 423 9 24 243 657 24 423 15 486 423 9 24 243 657 24}
\RouteRow{Tríadas estables por bloque}{0 1 2 3 4 5 6 7 8 9 10 11 11 13 14 14 15 17}
\end{tabularx}\endgroup\par\vspace{5pt}\noindent{\fontsize{8.4}{10}\selectfont\hyperref[trace-41]{Procedencia y códigos originales}\hfill Ficha completa · 041/324}
\clearpage\phantomsection\label{nav:vi-ruta-42}
% VI-CATALOG-ROW rutas 42
\pdfbookmark[3]{r2c2\_h-1v+1}{route-42}
\RouteTitle{Ruta 042}{r2c2\_\allowbreak{}h-1v+1}
\noindent Celda 11\quad (2, 2)\qquad Orientación \(h=-1,\ v=1\)\hfill\hyperref[sec:vi-indice-rutas]{Índice}\par\vspace{4pt}
\begingroup\fontsize{10.2}{12.5}\selectfont\setlength{\tabcolsep}{5pt}\renewcommand{\arraystretch}{1.1}\rowcolors{1}{routepale}{white}\noindent\begin{tabularx}{\linewidth}{@{}>{\raggedright\arraybackslash}p{.345\linewidth}>{\raggedright\arraybackslash}X@{}}
\RouteGroup{Identidad estructural}
\RouteRow{Clase y multisección}{quark\_\allowbreak{}up\_\allowbreak{}G3\quad /\quad quark\_\allowbreak{}up\_\allowbreak{}G3}
\RouteRow{Colección y sector}{quark de tipo arriba; quark}
\RouteRow{Clasificación física}{quark de tipo arriba}
\RouteRow{Generación, profundidad y color}{3; G3; R}
\RouteRow{Ruta primaria y órbita de inversión}{\(B_1\): sí\qquad 2,\allowbreak{}2;\allowbreak{}8,\allowbreak{}8}
\RouteRow{Firma de clasificación}{24|\allowbreak{}0|\allowbreak{}4|\allowbreak{}0|\allowbreak{}-12|\allowbreak{}+++-\kern0pt{}-\kern0pt{}-+++-\kern0pt{}-\kern0pt{}-}
\RouteGroup{Retornos, memoria y transporte}
\RouteRow{Retornos con fase}{clase: 27; órbita: 27; celda: 27}
\RouteRow{Retorno cinemático}{en 108: 54; extensión: 216}
\RouteRow{Giros y cambios de hoja}{71; 107}
\RouteRow{Acarreo y diferimiento}{deuda total: 2; final: 1; bloques diferidos: 2}
\RouteRow{Tríadas finales; persistencia de Pareto}{16; sí}
\RouteRow{\(K\longrightarrow K+9\)}{repeticiones: 9; cambios de memoria: 9}
\RouteGroup{Firma, fase y diferencias}
\RouteRow{\(n_A,\ s_C,\ \kappa_0,\ \nu_{120},\ \nu_{270}\)}{8, -2, 0, 0, -12}
\RouteRow{\(\tau_{12}\); firma histórica \(\mathbb Z^5\)}{+++-\kern0pt{}-\kern0pt{}-+++-\kern0pt{}-\kern0pt{}-; 8|\allowbreak{}-2|\allowbreak{}0|\allowbreak{}0|\allowbreak{}-12}
\RouteRow{\(D_1,\ D_3,\ D_4,\ D_{27},\ D_{36}\)}{84, 28, 84, 36, 24}
\RouteRow{\(\Omega_{90/120},\ P_6\)}{80, 6}
\RouteGroup{Lectores y factores de acción}
\RouteRow{\(K_D\quad /\quad K_\Omega\)}{(60,\allowbreak{}84,\allowbreak{}28)\quad /\quad (80,\allowbreak{}54,\allowbreak{}6)}
\RouteRow{\(\kappa_D\)}{59.39013\allowbreak{}58840155\allowbreak{}40637741\allowbreak{}48}
\RouteRow{\(\kappa_\Omega\)}{79.60661\allowbreak{}01397948\allowbreak{}27586470\allowbreak{}91}
\RouteRow{\(R_{\mathrm{act},D}\)}{1.000000\allowbreak{}04117409\allowbreak{}89467415\allowbreak{}75}
\RouteRow{\(R_{\mathrm{act},\Omega}\)}{1.000000\allowbreak{}05518981\allowbreak{}25318591\allowbreak{}40}
\RouteGroup{Separación de prefijos}
\RouteRow{Área de ambigüedad; área \(\log_2\)}{72; 36.0}
\RouteRow{Primer bloque de separación}{órbita: \textemdash{}; celda: \textemdash{}; ruta: \textemdash{}}
\RouteRow{Ambigüedad final}{4}
\RouteGroup{Secuencias completas}
\RouteRow{\(w_6\)}{414 420 258 414 414 333 255 252 333 414 420 258 414 414 333 255 252 333}
\RouteRow{Tríadas estables por bloque}{0 1 2 3 4 5 6 7 8 9 10 11 12 13 14 14 16 16}
\end{tabularx}\endgroup\par\vspace{5pt}\noindent{\fontsize{8.4}{10}\selectfont\hyperref[trace-42]{Procedencia y códigos originales}\hfill Ficha completa · 042/324}
\clearpage\phantomsection\label{nav:vi-ruta-43}
% VI-CATALOG-ROW rutas 43
\pdfbookmark[3]{r2c2\_h+1v-1}{route-43}
\RouteTitle{Ruta 043}{r2c2\_\allowbreak{}h+1v-1}
\noindent Celda 11\quad (2, 2)\qquad Orientación \(h=1,\ v=-1\)\hfill\hyperref[sec:vi-indice-rutas]{Índice}\par\vspace{4pt}
\begingroup\fontsize{10.2}{12.5}\selectfont\setlength{\tabcolsep}{5pt}\renewcommand{\arraystretch}{1.1}\rowcolors{1}{routepale}{white}\noindent\begin{tabularx}{\linewidth}{@{}>{\raggedright\arraybackslash}p{.345\linewidth}>{\raggedright\arraybackslash}X@{}}
\RouteGroup{Identidad estructural}
\RouteRow{Clase y multisección}{quark\_\allowbreak{}up\_\allowbreak{}G3\quad /\quad quark\_\allowbreak{}up\_\allowbreak{}G3}
\RouteRow{Colección y sector}{quark de tipo arriba; quark}
\RouteRow{Clasificación física}{quark de tipo arriba}
\RouteRow{Generación, profundidad y color}{3; G3; R}
\RouteRow{Ruta primaria y órbita de inversión}{\(B_1\): no\qquad 2,\allowbreak{}2;\allowbreak{}8,\allowbreak{}8}
\RouteRow{Firma de clasificación}{24|\allowbreak{}0|\allowbreak{}4|\allowbreak{}0|\allowbreak{}-12|\allowbreak{}+++-\kern0pt{}-\kern0pt{}-+++-\kern0pt{}-\kern0pt{}-}
\RouteGroup{Retornos, memoria y transporte}
\RouteRow{Retornos con fase}{clase: 27; órbita: 27; celda: 27}
\RouteRow{Retorno cinemático}{en 108: 54; extensión: 216}
\RouteRow{Giros y cambios de hoja}{71; 107}
\RouteRow{Acarreo y diferimiento}{deuda total: 3; final: 0; bloques diferidos: 2}
\RouteRow{Tríadas finales; persistencia de Pareto}{17; sí}
\RouteRow{\(K\longrightarrow K+9\)}{repeticiones: 9; cambios de memoria: 9}
\RouteGroup{Firma, fase y diferencias}
\RouteRow{\(n_A,\ s_C,\ \kappa_0,\ \nu_{120},\ \nu_{270}\)}{8, 2, 0, 0, -12}
\RouteRow{\(\tau_{12}\); firma histórica \(\mathbb Z^5\)}{+++-\kern0pt{}-\kern0pt{}-+++-\kern0pt{}-\kern0pt{}-; 8|\allowbreak{}2|\allowbreak{}0|\allowbreak{}0|\allowbreak{}-12}
\RouteRow{\(D_1,\ D_3,\ D_4,\ D_{27},\ D_{36}\)}{84, 28, 84, 36, 24}
\RouteRow{\(\Omega_{90/120},\ P_6\)}{80, 6}
\RouteGroup{Lectores y factores de acción}
\RouteRow{\(K_D\quad /\quad K_\Omega\)}{(60,\allowbreak{}84,\allowbreak{}28)\quad /\quad (80,\allowbreak{}54,\allowbreak{}6)}
\RouteRow{\(\kappa_D\)}{59.39013\allowbreak{}58840155\allowbreak{}40637741\allowbreak{}48}
\RouteRow{\(\kappa_\Omega\)}{79.60661\allowbreak{}01397948\allowbreak{}27586470\allowbreak{}91}
\RouteRow{\(R_{\mathrm{act},D}\)}{1.000000\allowbreak{}04117409\allowbreak{}89467415\allowbreak{}75}
\RouteRow{\(R_{\mathrm{act},\Omega}\)}{1.000000\allowbreak{}05518981\allowbreak{}25318591\allowbreak{}40}
\RouteGroup{Separación de prefijos}
\RouteRow{Área de ambigüedad; área \(\log_2\)}{72; 36.0}
\RouteRow{Primer bloque de separación}{órbita: \textemdash{}; celda: \textemdash{}; ruta: \textemdash{}}
\RouteRow{Ambigüedad final}{4}
\RouteGroup{Secuencias completas}
\RouteRow{\(w_6\)}{255 252 333 255 258 258 414 420 258 255 252 333 255 258 258 414 420 258}
\RouteRow{Tríadas estables por bloque}{0 1 2 3 4 5 6 7 8 9 10 11 11 12 14 15 15 17}
\end{tabularx}\endgroup\par\vspace{5pt}\noindent{\fontsize{8.4}{10}\selectfont\hyperref[trace-43]{Procedencia y códigos originales}\hfill Ficha completa · 043/324}
\clearpage\phantomsection\label{nav:vi-ruta-44}
% VI-CATALOG-ROW rutas 44
\pdfbookmark[3]{r2c2\_h+1v+1}{route-44}
\RouteTitle{Ruta 044}{r2c2\_\allowbreak{}h+1v+1}
\noindent Celda 11\quad (2, 2)\qquad Orientación \(h=1,\ v=1\)\hfill\hyperref[sec:vi-indice-rutas]{Índice}\par\vspace{4pt}
\begingroup\fontsize{10.2}{12.5}\selectfont\setlength{\tabcolsep}{5pt}\renewcommand{\arraystretch}{1.1}\rowcolors{1}{routepale}{white}\noindent\begin{tabularx}{\linewidth}{@{}>{\raggedright\arraybackslash}p{.345\linewidth}>{\raggedright\arraybackslash}X@{}}
\RouteGroup{Identidad estructural}
\RouteRow{Clase y multisección}{quark\_\allowbreak{}up\_\allowbreak{}G3\quad /\quad quark\_\allowbreak{}up\_\allowbreak{}G3}
\RouteRow{Colección y sector}{quark de tipo arriba; quark}
\RouteRow{Clasificación física}{quark de tipo arriba}
\RouteRow{Generación, profundidad y color}{3; G3; R}
\RouteRow{Ruta primaria y órbita de inversión}{\(B_1\): no\qquad 2,\allowbreak{}2;\allowbreak{}8,\allowbreak{}8}
\RouteRow{Firma de clasificación}{24|\allowbreak{}0|\allowbreak{}4|\allowbreak{}0|\allowbreak{}-12|\allowbreak{}+++-\kern0pt{}-\kern0pt{}-+++-\kern0pt{}-\kern0pt{}-}
\RouteGroup{Retornos, memoria y transporte}
\RouteRow{Retornos con fase}{clase: 18; órbita: 18; celda: 27}
\RouteRow{Retorno cinemático}{en 108: 54; extensión: 216}
\RouteRow{Giros y cambios de hoja}{71; 107}
\RouteRow{Acarreo y diferimiento}{deuda total: 2; final: 1; bloques diferidos: 2}
\RouteRow{Tríadas finales; persistencia de Pareto}{16; sí}
\RouteRow{\(K\longrightarrow K+9\)}{repeticiones: 9; cambios de memoria: 9}
\RouteGroup{Firma, fase y diferencias}
\RouteRow{\(n_A,\ s_C,\ \kappa_0,\ \nu_{120},\ \nu_{270}\)}{24, 0, 0, 0, -12}
\RouteRow{\(\tau_{12}\); firma histórica \(\mathbb Z^5\)}{+++-\kern0pt{}-\kern0pt{}-+++-\kern0pt{}-\kern0pt{}-; 24|\allowbreak{}0|\allowbreak{}0|\allowbreak{}0|\allowbreak{}-12}
\RouteRow{\(D_1,\ D_3,\ D_4,\ D_{27},\ D_{36}\)}{54, 56, 68, 48, 32}
\RouteRow{\(\Omega_{90/120},\ P_6\)}{80, 6}
\RouteGroup{Lectores y factores de acción}
\RouteRow{\(K_D\quad /\quad K_\Omega\)}{(80,\allowbreak{}54,\allowbreak{}6)\quad /\quad (80,\allowbreak{}54,\allowbreak{}6)}
\RouteRow{\(\kappa_D\)}{79.60661\allowbreak{}01397948\allowbreak{}27586470\allowbreak{}91}
\RouteRow{\(\kappa_\Omega\)}{79.60661\allowbreak{}01397948\allowbreak{}27586470\allowbreak{}91}
\RouteRow{\(R_{\mathrm{act},D}\)}{1.000000\allowbreak{}05518981\allowbreak{}25318591\allowbreak{}40}
\RouteRow{\(R_{\mathrm{act},\Omega}\)}{1.000000\allowbreak{}05518981\allowbreak{}25318591\allowbreak{}40}
\RouteGroup{Separación de prefijos}
\RouteRow{Área de ambigüedad; área \(\log_2\)}{72; 36.0}
\RouteRow{Primer bloque de separación}{órbita: \textemdash{}; celda: \textemdash{}; ruta: \textemdash{}}
\RouteRow{Ambigüedad final}{4}
\RouteGroup{Secuencias completas}
\RouteRow{\(w_6\)}{243 657 24 243 663 486 423 15 486 243 657 24 243 663 486 423 15 486}
\RouteRow{Tríadas estables por bloque}{0 1 2 3 4 5 6 7 8 9 10 11 12 13 14 14 16 16}
\end{tabularx}\endgroup\par\vspace{5pt}\noindent{\fontsize{8.4}{10}\selectfont\hyperref[trace-44]{Procedencia y códigos originales}\hfill Ficha completa · 044/324}
\clearpage\phantomsection\label{nav:vi-ruta-45}
% VI-CATALOG-ROW rutas 45
\pdfbookmark[2]{Celda 12}{cell-12}
\pdfbookmark[3]{r2c3\_h-1v-1}{route-45}
\RouteTitle{Ruta 045}{r2c3\_\allowbreak{}h-1v-1}
\noindent Celda 12\quad (2, 3)\qquad Orientación \(h=-1,\ v=-1\)\hfill\hyperref[sec:vi-indice-rutas]{Índice}\par\vspace{4pt}
\begingroup\fontsize{10.2}{12.5}\selectfont\setlength{\tabcolsep}{5pt}\renewcommand{\arraystretch}{1.1}\rowcolors{1}{routepale}{white}\noindent\begin{tabularx}{\linewidth}{@{}>{\raggedright\arraybackslash}p{.345\linewidth}>{\raggedright\arraybackslash}X@{}}
\RouteGroup{Identidad estructural}
\RouteRow{Clase y multisección}{quark\_\allowbreak{}down\_\allowbreak{}G3\quad /\quad quark\_\allowbreak{}down\_\allowbreak{}G3}
\RouteRow{Colección y sector}{quark de tipo abajo; quark}
\RouteRow{Clasificación física}{quark de tipo abajo}
\RouteRow{Generación, profundidad y color}{3; G3; B}
\RouteRow{Ruta primaria y órbita de inversión}{\(B_1\): sí\qquad 2,\allowbreak{}3;\allowbreak{}8,\allowbreak{}7}
\RouteRow{Firma de clasificación}{28|\allowbreak{}6|\allowbreak{}4|\allowbreak{}0|\allowbreak{}-12|\allowbreak{}+++-\kern0pt{}-\kern0pt{}-+++-\kern0pt{}-\kern0pt{}-}
\RouteGroup{Retornos, memoria y transporte}
\RouteRow{Retornos con fase}{clase: 27; órbita: 27; celda: 27}
\RouteRow{Retorno cinemático}{en 108: 54; extensión: 216}
\RouteRow{Giros y cambios de hoja}{71; 107}
\RouteRow{Acarreo y diferimiento}{deuda total: 1; final: 0; bloques diferidos: 1}
\RouteRow{Tríadas finales; persistencia de Pareto}{17; sí}
\RouteRow{\(K\longrightarrow K+9\)}{repeticiones: 9; cambios de memoria: 9}
\RouteGroup{Firma, fase y diferencias}
\RouteRow{\(n_A,\ s_C,\ \kappa_0,\ \nu_{120},\ \nu_{270}\)}{28, -6, 0, -4, -12}
\RouteRow{\(\tau_{12}\); firma histórica \(\mathbb Z^5\)}{+++-\kern0pt{}-\kern0pt{}-+++-\kern0pt{}-\kern0pt{}-; 28|\allowbreak{}-6|\allowbreak{}0|\allowbreak{}-4|\allowbreak{}-12}
\RouteRow{\(D_1,\ D_3,\ D_4,\ D_{27},\ D_{36}\)}{66, 60, 68, 48, 32}
\RouteRow{\(\Omega_{90/120},\ P_6\)}{80, 6}
\RouteGroup{Lectores y factores de acción}
\RouteRow{\(K_D\quad /\quad K_\Omega\)}{(80,\allowbreak{}66,\allowbreak{}4)\quad /\quad (80,\allowbreak{}54,\allowbreak{}6)}
\RouteRow{\(\kappa_D\)}{79.51881\allowbreak{}95161180\allowbreak{}37694395\allowbreak{}39}
\RouteRow{\(\kappa_\Omega\)}{79.60661\allowbreak{}01397948\allowbreak{}27586470\allowbreak{}91}
\RouteRow{\(R_{\mathrm{act},D}\)}{1.000000\allowbreak{}05512894\allowbreak{}88901612\allowbreak{}64}
\RouteRow{\(R_{\mathrm{act},\Omega}\)}{1.000000\allowbreak{}05518981\allowbreak{}25318591\allowbreak{}40}
\RouteGroup{Separación de prefijos}
\RouteRow{Área de ambigüedad; área \(\log_2\)}{36; 18.0}
\RouteRow{Primer bloque de separación}{órbita: \textemdash{}; celda: \textemdash{}; ruta: \textemdash{}}
\RouteRow{Ambigüedad final}{2}
\RouteGroup{Secuencias completas}
\RouteRow{\(w_6\)}{570 264 90 570 267 243 657 24 243 570 264 90 570 267 243 657 24 243}
\RouteRow{Tríadas estables por bloque}{0 1 1 3 4 5 6 7 8 9 10 11 12 13 14 15 16 17}
\end{tabularx}\endgroup\par\vspace{5pt}\noindent{\fontsize{8.4}{10}\selectfont\hyperref[trace-45]{Procedencia y códigos originales}\hfill Ficha completa · 045/324}
\clearpage\phantomsection\label{nav:vi-ruta-46}
% VI-CATALOG-ROW rutas 46
\pdfbookmark[3]{r2c3\_h-1v+1}{route-46}
\RouteTitle{Ruta 046}{r2c3\_\allowbreak{}h-1v+1}
\noindent Celda 12\quad (2, 3)\qquad Orientación \(h=-1,\ v=1\)\hfill\hyperref[sec:vi-indice-rutas]{Índice}\par\vspace{4pt}
\begingroup\fontsize{10.2}{12.5}\selectfont\setlength{\tabcolsep}{5pt}\renewcommand{\arraystretch}{1.1}\rowcolors{1}{routepale}{white}\noindent\begin{tabularx}{\linewidth}{@{}>{\raggedright\arraybackslash}p{.345\linewidth}>{\raggedright\arraybackslash}X@{}}
\RouteGroup{Identidad estructural}
\RouteRow{Clase y multisección}{quark\_\allowbreak{}down\_\allowbreak{}G3\quad /\quad quark\_\allowbreak{}down\_\allowbreak{}G3}
\RouteRow{Colección y sector}{quark de tipo abajo; quark}
\RouteRow{Clasificación física}{quark de tipo abajo}
\RouteRow{Generación, profundidad y color}{3; G3; B}
\RouteRow{Ruta primaria y órbita de inversión}{\(B_1\): no\qquad 2,\allowbreak{}3;\allowbreak{}8,\allowbreak{}7}
\RouteRow{Firma de clasificación}{28|\allowbreak{}6|\allowbreak{}4|\allowbreak{}0|\allowbreak{}-12|\allowbreak{}+++-\kern0pt{}-\kern0pt{}-+++-\kern0pt{}-\kern0pt{}-}
\RouteGroup{Retornos, memoria y transporte}
\RouteRow{Retornos con fase}{clase: 27; órbita: 27; celda: 27}
\RouteRow{Retorno cinemático}{en 108: 54; extensión: 216}
\RouteRow{Giros y cambios de hoja}{71; 107}
\RouteRow{Acarreo y diferimiento}{deuda total: 1; final: 0; bloques diferidos: 1}
\RouteRow{Tríadas finales; persistencia de Pareto}{17; sí}
\RouteRow{\(K\longrightarrow K+9\)}{repeticiones: 9; cambios de memoria: 9}
\RouteGroup{Firma, fase y diferencias}
\RouteRow{\(n_A,\ s_C,\ \kappa_0,\ \nu_{120},\ \nu_{270}\)}{44, 2, 0, 2, -12}
\RouteRow{\(\tau_{12}\); firma histórica \(\mathbb Z^5\)}{+++-\kern0pt{}-\kern0pt{}-+++-\kern0pt{}-\kern0pt{}-; 44|\allowbreak{}2|\allowbreak{}0|\allowbreak{}2|\allowbreak{}-12}
\RouteRow{\(D_1,\ D_3,\ D_4,\ D_{27},\ D_{36}\)}{78, 24, 78, 24, 16}
\RouteRow{\(\Omega_{90/120},\ P_6\)}{80, 6}
\RouteGroup{Lectores y factores de acción}
\RouteRow{\(K_D\quad /\quad K_\Omega\)}{(40,\allowbreak{}78,\allowbreak{}27)\quad /\quad (80,\allowbreak{}54,\allowbreak{}6)}
\RouteRow{\(\kappa_D\)}{39.43380\allowbreak{}88030085\allowbreak{}51303671\allowbreak{}14}
\RouteRow{\(\kappa_\Omega\)}{79.60661\allowbreak{}01397948\allowbreak{}27586470\allowbreak{}91}
\RouteRow{\(R_{\mathrm{act},D}\)}{1.000000\allowbreak{}02733874\allowbreak{}08548943\allowbreak{}58}
\RouteRow{\(R_{\mathrm{act},\Omega}\)}{1.000000\allowbreak{}05518981\allowbreak{}25318591\allowbreak{}40}
\RouteGroup{Separación de prefijos}
\RouteRow{Área de ambigüedad; área \(\log_2\)}{36; 18.0}
\RouteRow{Primer bloque de separación}{órbita: \textemdash{}; celda: \textemdash{}; ruta: \textemdash{}}
\RouteRow{Ambigüedad final}{2}
\RouteGroup{Secuencias completas}
\RouteRow{\(w_6\)}{585 507 504 585 510 666 672 510 666 585 507 504 585 510 666 672 510 666}
\RouteRow{Tríadas estables por bloque}{0 1 2 3 4 5 6 7 8 9 10 11 12 13 14 14 16 17}
\end{tabularx}\endgroup\par\vspace{5pt}\noindent{\fontsize{8.4}{10}\selectfont\hyperref[trace-46]{Procedencia y códigos originales}\hfill Ficha completa · 046/324}
\clearpage\phantomsection\label{nav:vi-ruta-47}
% VI-CATALOG-ROW rutas 47
\pdfbookmark[3]{r2c3\_h+1v-1}{route-47}
\RouteTitle{Ruta 047}{r2c3\_\allowbreak{}h+1v-1}
\noindent Celda 12\quad (2, 3)\qquad Orientación \(h=1,\ v=-1\)\hfill\hyperref[sec:vi-indice-rutas]{Índice}\par\vspace{4pt}
\begingroup\fontsize{10.2}{12.5}\selectfont\setlength{\tabcolsep}{5pt}\renewcommand{\arraystretch}{1.1}\rowcolors{1}{routepale}{white}\noindent\begin{tabularx}{\linewidth}{@{}>{\raggedright\arraybackslash}p{.345\linewidth}>{\raggedright\arraybackslash}X@{}}
\RouteGroup{Identidad estructural}
\RouteRow{Clase y multisección}{quark\_\allowbreak{}down\_\allowbreak{}G3\quad /\quad quark\_\allowbreak{}down\_\allowbreak{}G3}
\RouteRow{Colección y sector}{quark de tipo abajo; quark}
\RouteRow{Clasificación física}{quark de tipo abajo}
\RouteRow{Generación, profundidad y color}{3; G3; B}
\RouteRow{Ruta primaria y órbita de inversión}{\(B_1\): no\qquad 2,\allowbreak{}3;\allowbreak{}8,\allowbreak{}7}
\RouteRow{Firma de clasificación}{28|\allowbreak{}6|\allowbreak{}4|\allowbreak{}0|\allowbreak{}-12|\allowbreak{}+++-\kern0pt{}-\kern0pt{}-+++-\kern0pt{}-\kern0pt{}-}
\RouteGroup{Retornos, memoria y transporte}
\RouteRow{Retornos con fase}{clase: 27; órbita: 27; celda: 27}
\RouteRow{Retorno cinemático}{en 108: 54; extensión: 216}
\RouteRow{Giros y cambios de hoja}{71; 107}
\RouteRow{Acarreo y diferimiento}{deuda total: 1; final: 0; bloques diferidos: 1}
\RouteRow{Tríadas finales; persistencia de Pareto}{17; sí}
\RouteRow{\(K\longrightarrow K+9\)}{repeticiones: 9; cambios de memoria: 9}
\RouteGroup{Firma, fase y diferencias}
\RouteRow{\(n_A,\ s_C,\ \kappa_0,\ \nu_{120},\ \nu_{270}\)}{44, -2, 0, -2, -12}
\RouteRow{\(\tau_{12}\); firma histórica \(\mathbb Z^5\)}{+++-\kern0pt{}-\kern0pt{}-+++-\kern0pt{}-\kern0pt{}-; 44|\allowbreak{}-2|\allowbreak{}0|\allowbreak{}-2|\allowbreak{}-12}
\RouteRow{\(D_1,\ D_3,\ D_4,\ D_{27},\ D_{36}\)}{78, 24, 78, 24, 16}
\RouteRow{\(\Omega_{90/120},\ P_6\)}{80, 6}
\RouteGroup{Lectores y factores de acción}
\RouteRow{\(K_D\quad /\quad K_\Omega\)}{(40,\allowbreak{}78,\allowbreak{}27)\quad /\quad (80,\allowbreak{}54,\allowbreak{}6)}
\RouteRow{\(\kappa_D\)}{39.43380\allowbreak{}88030085\allowbreak{}51303671\allowbreak{}14}
\RouteRow{\(\kappa_\Omega\)}{79.60661\allowbreak{}01397948\allowbreak{}27586470\allowbreak{}91}
\RouteRow{\(R_{\mathrm{act},D}\)}{1.000000\allowbreak{}02733874\allowbreak{}08548943\allowbreak{}58}
\RouteRow{\(R_{\mathrm{act},\Omega}\)}{1.000000\allowbreak{}05518981\allowbreak{}25318591\allowbreak{}40}
\RouteGroup{Separación de prefijos}
\RouteRow{Área de ambigüedad; área \(\log_2\)}{36; 18.0}
\RouteRow{Primer bloque de separación}{órbita: \textemdash{}; celda: \textemdash{}; ruta: \textemdash{}}
\RouteRow{Ambigüedad final}{2}
\RouteGroup{Secuencias completas}
\RouteRow{\(w_6\)}{672 510 666 672 507 504 585 507 504 672 510 666 672 507 504 585 507 504}
\RouteRow{Tríadas estables por bloque}{0 1 2 3 4 5 6 7 8 9 10 11 12 13 14 14 16 17}
\end{tabularx}\endgroup\par\vspace{5pt}\noindent{\fontsize{8.4}{10}\selectfont\hyperref[trace-47]{Procedencia y códigos originales}\hfill Ficha completa · 047/324}
\clearpage\phantomsection\label{nav:vi-ruta-48}
% VI-CATALOG-ROW rutas 48
\pdfbookmark[3]{r2c3\_h+1v+1}{route-48}
\RouteTitle{Ruta 048}{r2c3\_\allowbreak{}h+1v+1}
\noindent Celda 12\quad (2, 3)\qquad Orientación \(h=1,\ v=1\)\hfill\hyperref[sec:vi-indice-rutas]{Índice}\par\vspace{4pt}
\begingroup\fontsize{10.2}{12.5}\selectfont\setlength{\tabcolsep}{5pt}\renewcommand{\arraystretch}{1.1}\rowcolors{1}{routepale}{white}\noindent\begin{tabularx}{\linewidth}{@{}>{\raggedright\arraybackslash}p{.345\linewidth}>{\raggedright\arraybackslash}X@{}}
\RouteGroup{Identidad estructural}
\RouteRow{Clase y multisección}{quark\_\allowbreak{}down\_\allowbreak{}G3\quad /\quad quark\_\allowbreak{}down\_\allowbreak{}G3}
\RouteRow{Colección y sector}{quark de tipo abajo; quark}
\RouteRow{Clasificación física}{quark de tipo abajo}
\RouteRow{Generación, profundidad y color}{3; G3; B}
\RouteRow{Ruta primaria y órbita de inversión}{\(B_1\): no\qquad 2,\allowbreak{}3;\allowbreak{}8,\allowbreak{}7}
\RouteRow{Firma de clasificación}{28|\allowbreak{}6|\allowbreak{}4|\allowbreak{}0|\allowbreak{}-12|\allowbreak{}+++-\kern0pt{}-\kern0pt{}-+++-\kern0pt{}-\kern0pt{}-}
\RouteGroup{Retornos, memoria y transporte}
\RouteRow{Retornos con fase}{clase: 27; órbita: 27; celda: 27}
\RouteRow{Retorno cinemático}{en 108: 54; extensión: 216}
\RouteRow{Giros y cambios de hoja}{71; 107}
\RouteRow{Acarreo y diferimiento}{deuda total: 0; final: 0; bloques diferidos: 0}
\RouteRow{Tríadas finales; persistencia de Pareto}{17; sí}
\RouteRow{\(K\longrightarrow K+9\)}{repeticiones: 9; cambios de memoria: 9}
\RouteGroup{Firma, fase y diferencias}
\RouteRow{\(n_A,\ s_C,\ \kappa_0,\ \nu_{120},\ \nu_{270}\)}{28, 6, 0, 4, -12}
\RouteRow{\(\tau_{12}\); firma histórica \(\mathbb Z^5\)}{+++-\kern0pt{}-\kern0pt{}-+++-\kern0pt{}-\kern0pt{}-; 28|\allowbreak{}6|\allowbreak{}0|\allowbreak{}4|\allowbreak{}-12}
\RouteRow{\(D_1,\ D_3,\ D_4,\ D_{27},\ D_{36}\)}{66, 60, 68, 48, 32}
\RouteRow{\(\Omega_{90/120},\ P_6\)}{80, 6}
\RouteGroup{Lectores y factores de acción}
\RouteRow{\(K_D\quad /\quad K_\Omega\)}{(80,\allowbreak{}66,\allowbreak{}4)\quad /\quad (80,\allowbreak{}54,\allowbreak{}6)}
\RouteRow{\(\kappa_D\)}{79.51881\allowbreak{}95161180\allowbreak{}37694395\allowbreak{}39}
\RouteRow{\(\kappa_\Omega\)}{79.60661\allowbreak{}01397948\allowbreak{}27586470\allowbreak{}91}
\RouteRow{\(R_{\mathrm{act},D}\)}{1.000000\allowbreak{}05512894\allowbreak{}88901612\allowbreak{}64}
\RouteRow{\(R_{\mathrm{act},\Omega}\)}{1.000000\allowbreak{}05518981\allowbreak{}25318591\allowbreak{}40}
\RouteGroup{Separación de prefijos}
\RouteRow{Área de ambigüedad; área \(\log_2\)}{36; 18.0}
\RouteRow{Primer bloque de separación}{órbita: \textemdash{}; celda: \textemdash{}; ruta: \textemdash{}}
\RouteRow{Ambigüedad final}{2}
\RouteGroup{Secuencias completas}
\RouteRow{\(w_6\)}{657 24 243 657 21 90 570 264 90 657 24 243 657 21 90 570 264 90}
\RouteRow{Tríadas estables por bloque}{0 1 2 3 4 5 6 7 8 9 10 11 12 13 14 15 16 17}
\end{tabularx}\endgroup\par\vspace{5pt}\noindent{\fontsize{8.4}{10}\selectfont\hyperref[trace-48]{Procedencia y códigos originales}\hfill Ficha completa · 048/324}
\clearpage\phantomsection\label{nav:vi-ruta-49}
% VI-CATALOG-ROW rutas 49
\pdfbookmark[2]{Celda 13}{cell-13}
\pdfbookmark[3]{r2c4\_h-1v-1}{route-49}
\RouteTitle{Ruta 049}{r2c4\_\allowbreak{}h-1v-1}
\noindent Celda 13\quad (2, 4)\qquad Orientación \(h=-1,\ v=-1\)\hfill\hyperref[sec:vi-indice-rutas]{Índice}\par\vspace{4pt}
\begingroup\fontsize{10.2}{12.5}\selectfont\setlength{\tabcolsep}{5pt}\renewcommand{\arraystretch}{1.1}\rowcolors{1}{routepale}{white}\noindent\begin{tabularx}{\linewidth}{@{}>{\raggedright\arraybackslash}p{.345\linewidth}>{\raggedright\arraybackslash}X@{}}
\RouteGroup{Identidad estructural}
\RouteRow{Clase y multisección}{quark\_\allowbreak{}up\_\allowbreak{}G3\quad /\quad quark\_\allowbreak{}up\_\allowbreak{}G3}
\RouteRow{Colección y sector}{quark de tipo arriba; quark}
\RouteRow{Clasificación física}{quark de tipo arriba}
\RouteRow{Generación, profundidad y color}{3; G3; G}
\RouteRow{Ruta primaria y órbita de inversión}{\(B_1\): no\qquad 2,\allowbreak{}4;\allowbreak{}8,\allowbreak{}6}
\RouteRow{Firma de clasificación}{28|\allowbreak{}-6|\allowbreak{}6|\allowbreak{}0|\allowbreak{}-12|\allowbreak{}+++-\kern0pt{}-\kern0pt{}-+++-\kern0pt{}-\kern0pt{}-}
\RouteGroup{Retornos, memoria y transporte}
\RouteRow{Retornos con fase}{clase: 27; órbita: 27; celda: 27}
\RouteRow{Retorno cinemático}{en 108: 54; extensión: 216}
\RouteRow{Giros y cambios de hoja}{71; 107}
\RouteRow{Acarreo y diferimiento}{deuda total: 1; final: 1; bloques diferidos: 1}
\RouteRow{Tríadas finales; persistencia de Pareto}{16; no}
\RouteRow{\(K\longrightarrow K+9\)}{repeticiones: 9; cambios de memoria: 9}
\RouteGroup{Firma, fase y diferencias}
\RouteRow{\(n_A,\ s_C,\ \kappa_0,\ \nu_{120},\ \nu_{270}\)}{28, 6, 0, 0, -12}
\RouteRow{\(\tau_{12}\); firma histórica \(\mathbb Z^5\)}{+++-\kern0pt{}-\kern0pt{}-+++-\kern0pt{}-\kern0pt{}-; 28|\allowbreak{}6|\allowbreak{}0|\allowbreak{}0|\allowbreak{}-12}
\RouteRow{\(D_1,\ D_3,\ D_4,\ D_{27},\ D_{36}\)}{66, 64, 68, 60, 40}
\RouteRow{\(\Omega_{90/120},\ P_6\)}{80, 6}
\RouteGroup{Lectores y factores de acción}
\RouteRow{\(K_D\quad /\quad K_\Omega\)}{(100,\allowbreak{}66,\allowbreak{}2)\quad /\quad (80,\allowbreak{}54,\allowbreak{}6)}
\RouteRow{\(\kappa_D\)}{99.51859\allowbreak{}71232726\allowbreak{}53414287\allowbreak{}29}
\RouteRow{\(\kappa_\Omega\)}{79.60661\allowbreak{}01397948\allowbreak{}27586470\allowbreak{}91}
\RouteRow{\(R_{\mathrm{act},D}\)}{1.000000\allowbreak{}06899443\allowbreak{}08259364\allowbreak{}17}
\RouteRow{\(R_{\mathrm{act},\Omega}\)}{1.000000\allowbreak{}05518981\allowbreak{}25318591\allowbreak{}40}
\RouteGroup{Separación de prefijos}
\RouteRow{Área de ambigüedad; área \(\log_2\)}{36; 18.0}
\RouteRow{Primer bloque de separación}{órbita: \textemdash{}; celda: \textemdash{}; ruta: \textemdash{}}
\RouteRow{Ambigüedad final}{2}
\RouteGroup{Secuencias completas}
\RouteRow{\(w_6\)}{15 486 423 15 489 498 99 327 498 15 486 423 15 489 498 99 327 498}
\RouteRow{Tríadas estables por bloque}{0 1 2 3 4 5 6 7 8 9 10 11 12 13 14 15 16 16}
\end{tabularx}\endgroup\par\vspace{5pt}\noindent{\fontsize{8.4}{10}\selectfont\hyperref[trace-49]{Procedencia y códigos originales}\hfill Ficha completa · 049/324}
\clearpage\phantomsection\label{nav:vi-ruta-50}
% VI-CATALOG-ROW rutas 50
\pdfbookmark[3]{r2c4\_h-1v+1}{route-50}
\RouteTitle{Ruta 050}{r2c4\_\allowbreak{}h-1v+1}
\noindent Celda 13\quad (2, 4)\qquad Orientación \(h=-1,\ v=1\)\hfill\hyperref[sec:vi-indice-rutas]{Índice}\par\vspace{4pt}
\begingroup\fontsize{10.2}{12.5}\selectfont\setlength{\tabcolsep}{5pt}\renewcommand{\arraystretch}{1.1}\rowcolors{1}{routepale}{white}\noindent\begin{tabularx}{\linewidth}{@{}>{\raggedright\arraybackslash}p{.345\linewidth}>{\raggedright\arraybackslash}X@{}}
\RouteGroup{Identidad estructural}
\RouteRow{Clase y multisección}{quark\_\allowbreak{}up\_\allowbreak{}G3\quad /\quad quark\_\allowbreak{}up\_\allowbreak{}G3}
\RouteRow{Colección y sector}{quark de tipo arriba; quark}
\RouteRow{Clasificación física}{quark de tipo arriba}
\RouteRow{Generación, profundidad y color}{3; G3; G}
\RouteRow{Ruta primaria y órbita de inversión}{\(B_1\): no\qquad 2,\allowbreak{}4;\allowbreak{}8,\allowbreak{}6}
\RouteRow{Firma de clasificación}{28|\allowbreak{}-6|\allowbreak{}6|\allowbreak{}0|\allowbreak{}-12|\allowbreak{}+++-\kern0pt{}-\kern0pt{}-+++-\kern0pt{}-\kern0pt{}-}
\RouteGroup{Retornos, memoria y transporte}
\RouteRow{Retornos con fase}{clase: 27; órbita: 27; celda: 27}
\RouteRow{Retorno cinemático}{en 108: 54; extensión: 216}
\RouteRow{Giros y cambios de hoja}{71; 107}
\RouteRow{Acarreo y diferimiento}{deuda total: 0; final: 0; bloques diferidos: 0}
\RouteRow{Tríadas finales; persistencia de Pareto}{17; no}
\RouteRow{\(K\longrightarrow K+9\)}{repeticiones: 9; cambios de memoria: 9}
\RouteGroup{Firma, fase y diferencias}
\RouteRow{\(n_A,\ s_C,\ \kappa_0,\ \nu_{120},\ \nu_{270}\)}{14, 6, 0, 0, -12}
\RouteRow{\(\tau_{12}\); firma histórica \(\mathbb Z^5\)}{+++-\kern0pt{}-\kern0pt{}-+++-\kern0pt{}-\kern0pt{}-; 14|\allowbreak{}6|\allowbreak{}0|\allowbreak{}0|\allowbreak{}-12}
\RouteRow{\(D_1,\ D_3,\ D_4,\ D_{27},\ D_{36}\)}{24, 20, 24, 24, 16}
\RouteRow{\(\Omega_{90/120},\ P_6\)}{80, 6}
\RouteGroup{Lectores y factores de acción}
\RouteRow{\(K_D\quad /\quad K_\Omega\)}{(40,\allowbreak{}24,\allowbreak{}2)\quad /\quad (80,\allowbreak{}54,\allowbreak{}6)}
\RouteRow{\(\kappa_D\)}{39.82508\allowbreak{}59311825\allowbreak{}73056173\allowbreak{}26}
\RouteRow{\(\kappa_\Omega\)}{79.60661\allowbreak{}01397948\allowbreak{}27586470\allowbreak{}91}
\RouteRow{\(R_{\mathrm{act},D}\)}{1.000000\allowbreak{}02761000\allowbreak{}61595142\allowbreak{}15}
\RouteRow{\(R_{\mathrm{act},\Omega}\)}{1.000000\allowbreak{}05518981\allowbreak{}25318591\allowbreak{}40}
\RouteGroup{Separación de prefijos}
\RouteRow{Área de ambigüedad; área \(\log_2\)}{72; 36.0}
\RouteRow{Primer bloque de separación}{órbita: \textemdash{}; celda: \textemdash{}; ruta: \textemdash{}}
\RouteRow{Ambigüedad final}{4}
\RouteGroup{Secuencias completas}
\RouteRow{\(w_6\)}{0 81 3 0 84 0 81 3 0 0 81 3 0 84 0 81 3 0}
\RouteRow{Tríadas estables por bloque}{0 1 2 3 4 5 6 7 8 9 10 11 12 13 14 15 16 17}
\end{tabularx}\endgroup\par\vspace{5pt}\noindent{\fontsize{8.4}{10}\selectfont\hyperref[trace-50]{Procedencia y códigos originales}\hfill Ficha completa · 050/324}
\clearpage\phantomsection\label{nav:vi-ruta-51}
% VI-CATALOG-ROW rutas 51
\pdfbookmark[3]{r2c4\_h+1v-1}{route-51}
\RouteTitle{Ruta 051}{r2c4\_\allowbreak{}h+1v-1}
\noindent Celda 13\quad (2, 4)\qquad Orientación \(h=1,\ v=-1\)\hfill\hyperref[sec:vi-indice-rutas]{Índice}\par\vspace{4pt}
\begingroup\fontsize{10.2}{12.5}\selectfont\setlength{\tabcolsep}{5pt}\renewcommand{\arraystretch}{1.1}\rowcolors{1}{routepale}{white}\noindent\begin{tabularx}{\linewidth}{@{}>{\raggedright\arraybackslash}p{.345\linewidth}>{\raggedright\arraybackslash}X@{}}
\RouteGroup{Identidad estructural}
\RouteRow{Clase y multisección}{quark\_\allowbreak{}up\_\allowbreak{}G3\quad /\quad quark\_\allowbreak{}up\_\allowbreak{}G3}
\RouteRow{Colección y sector}{quark de tipo arriba; quark}
\RouteRow{Clasificación física}{quark de tipo arriba}
\RouteRow{Generación, profundidad y color}{3; G3; G}
\RouteRow{Ruta primaria y órbita de inversión}{\(B_1\): sí\qquad 2,\allowbreak{}4;\allowbreak{}8,\allowbreak{}6}
\RouteRow{Firma de clasificación}{28|\allowbreak{}-6|\allowbreak{}6|\allowbreak{}0|\allowbreak{}-12|\allowbreak{}+++-\kern0pt{}-\kern0pt{}-+++-\kern0pt{}-\kern0pt{}-}
\RouteGroup{Retornos, memoria y transporte}
\RouteRow{Retornos con fase}{clase: 27; órbita: 27; celda: 27}
\RouteRow{Retorno cinemático}{en 108: 54; extensión: 216}
\RouteRow{Giros y cambios de hoja}{71; 107}
\RouteRow{Acarreo y diferimiento}{deuda total: 2; final: 0; bloques diferidos: 2}
\RouteRow{Tríadas finales; persistencia de Pareto}{17; no}
\RouteRow{\(K\longrightarrow K+9\)}{repeticiones: 9; cambios de memoria: 9}
\RouteGroup{Firma, fase y diferencias}
\RouteRow{\(n_A,\ s_C,\ \kappa_0,\ \nu_{120},\ \nu_{270}\)}{14, -6, 0, 0, -12}
\RouteRow{\(\tau_{12}\); firma histórica \(\mathbb Z^5\)}{+++-\kern0pt{}-\kern0pt{}-+++-\kern0pt{}-\kern0pt{}-; 14|\allowbreak{}-6|\allowbreak{}0|\allowbreak{}0|\allowbreak{}-12}
\RouteRow{\(D_1,\ D_3,\ D_4,\ D_{27},\ D_{36}\)}{24, 20, 24, 24, 16}
\RouteRow{\(\Omega_{90/120},\ P_6\)}{80, 6}
\RouteGroup{Lectores y factores de acción}
\RouteRow{\(K_D\quad /\quad K_\Omega\)}{(40,\allowbreak{}24,\allowbreak{}2)\quad /\quad (80,\allowbreak{}54,\allowbreak{}6)}
\RouteRow{\(\kappa_D\)}{39.82508\allowbreak{}59311825\allowbreak{}73056173\allowbreak{}26}
\RouteRow{\(\kappa_\Omega\)}{79.60661\allowbreak{}01397948\allowbreak{}27586470\allowbreak{}91}
\RouteRow{\(R_{\mathrm{act},D}\)}{1.000000\allowbreak{}02761000\allowbreak{}61595142\allowbreak{}15}
\RouteRow{\(R_{\mathrm{act},\Omega}\)}{1.000000\allowbreak{}05518981\allowbreak{}25318591\allowbreak{}40}
\RouteGroup{Separación de prefijos}
\RouteRow{Área de ambigüedad; área \(\log_2\)}{72; 36.0}
\RouteRow{Primer bloque de separación}{órbita: \textemdash{}; celda: \textemdash{}; ruta: \textemdash{}}
\RouteRow{Ambigüedad final}{4}
\RouteGroup{Secuencias completas}
\RouteRow{\(w_6\)}{81 3 0 81 0 3 0 81 3 81 3 0 81 0 3 0 81 3}
\RouteRow{Tríadas estables por bloque}{0 1 2 3 4 5 6 6 8 9 10 11 12 13 14 15 15 17}
\end{tabularx}\endgroup\par\vspace{5pt}\noindent{\fontsize{8.4}{10}\selectfont\hyperref[trace-51]{Procedencia y códigos originales}\hfill Ficha completa · 051/324}
\clearpage\phantomsection\label{nav:vi-ruta-52}
% VI-CATALOG-ROW rutas 52
\pdfbookmark[3]{r2c4\_h+1v+1}{route-52}
\RouteTitle{Ruta 052}{r2c4\_\allowbreak{}h+1v+1}
\noindent Celda 13\quad (2, 4)\qquad Orientación \(h=1,\ v=1\)\hfill\hyperref[sec:vi-indice-rutas]{Índice}\par\vspace{4pt}
\begingroup\fontsize{10.2}{12.5}\selectfont\setlength{\tabcolsep}{5pt}\renewcommand{\arraystretch}{1.1}\rowcolors{1}{routepale}{white}\noindent\begin{tabularx}{\linewidth}{@{}>{\raggedright\arraybackslash}p{.345\linewidth}>{\raggedright\arraybackslash}X@{}}
\RouteGroup{Identidad estructural}
\RouteRow{Clase y multisección}{quark\_\allowbreak{}up\_\allowbreak{}G3\quad /\quad quark\_\allowbreak{}up\_\allowbreak{}G3}
\RouteRow{Colección y sector}{quark de tipo arriba; quark}
\RouteRow{Clasificación física}{quark de tipo arriba}
\RouteRow{Generación, profundidad y color}{3; G3; G}
\RouteRow{Ruta primaria y órbita de inversión}{\(B_1\): no\qquad 2,\allowbreak{}4;\allowbreak{}8,\allowbreak{}6}
\RouteRow{Firma de clasificación}{28|\allowbreak{}-6|\allowbreak{}6|\allowbreak{}0|\allowbreak{}-12|\allowbreak{}+++-\kern0pt{}-\kern0pt{}-+++-\kern0pt{}-\kern0pt{}-}
\RouteGroup{Retornos, memoria y transporte}
\RouteRow{Retornos con fase}{clase: 27; órbita: 27; celda: 27}
\RouteRow{Retorno cinemático}{en 108: 54; extensión: 216}
\RouteRow{Giros y cambios de hoja}{71; 107}
\RouteRow{Acarreo y diferimiento}{deuda total: 1; final: 0; bloques diferidos: 1}
\RouteRow{Tríadas finales; persistencia de Pareto}{17; no}
\RouteRow{\(K\longrightarrow K+9\)}{repeticiones: 9; cambios de memoria: 9}
\RouteGroup{Firma, fase y diferencias}
\RouteRow{\(n_A,\ s_C,\ \kappa_0,\ \nu_{120},\ \nu_{270}\)}{28, -6, 0, 0, -12}
\RouteRow{\(\tau_{12}\); firma histórica \(\mathbb Z^5\)}{+++-\kern0pt{}-\kern0pt{}-+++-\kern0pt{}-\kern0pt{}-; 28|\allowbreak{}-6|\allowbreak{}0|\allowbreak{}0|\allowbreak{}-12}
\RouteRow{\(D_1,\ D_3,\ D_4,\ D_{27},\ D_{36}\)}{66, 64, 68, 60, 40}
\RouteRow{\(\Omega_{90/120},\ P_6\)}{80, 6}
\RouteGroup{Lectores y factores de acción}
\RouteRow{\(K_D\quad /\quad K_\Omega\)}{(100,\allowbreak{}66,\allowbreak{}2)\quad /\quad (80,\allowbreak{}54,\allowbreak{}6)}
\RouteRow{\(\kappa_D\)}{99.51859\allowbreak{}71232726\allowbreak{}53414287\allowbreak{}29}
\RouteRow{\(\kappa_\Omega\)}{79.60661\allowbreak{}01397948\allowbreak{}27586470\allowbreak{}91}
\RouteRow{\(R_{\mathrm{act},D}\)}{1.000000\allowbreak{}06899443\allowbreak{}08259364\allowbreak{}17}
\RouteRow{\(R_{\mathrm{act},\Omega}\)}{1.000000\allowbreak{}05518981\allowbreak{}25318591\allowbreak{}40}
\RouteGroup{Separación de prefijos}
\RouteRow{Área de ambigüedad; área \(\log_2\)}{36; 18.0}
\RouteRow{Primer bloque de separación}{órbita: \textemdash{}; celda: \textemdash{}; ruta: \textemdash{}}
\RouteRow{Ambigüedad final}{2}
\RouteGroup{Secuencias completas}
\RouteRow{\(w_6\)}{99 327 498 99 324 423 15 486 423 99 327 498 99 324 423 15 486 423}
\RouteRow{Tríadas estables por bloque}{0 1 2 3 4 5 6 7 8 9 10 11 12 13 14 14 16 17}
\end{tabularx}\endgroup\par\vspace{5pt}\noindent{\fontsize{8.4}{10}\selectfont\hyperref[trace-52]{Procedencia y códigos originales}\hfill Ficha completa · 052/324}
\clearpage\phantomsection\label{nav:vi-ruta-53}
% VI-CATALOG-ROW rutas 53
\pdfbookmark[2]{Celda 14}{cell-14}
\pdfbookmark[3]{r2c5\_h-1v-1}{route-53}
\RouteTitle{Ruta 053}{r2c5\_\allowbreak{}h-1v-1}
\noindent Celda 14\quad (2, 5)\qquad Orientación \(h=-1,\ v=-1\)\hfill\hyperref[sec:vi-indice-rutas]{Índice}\par\vspace{4pt}
\begingroup\fontsize{10.2}{12.5}\selectfont\setlength{\tabcolsep}{5pt}\renewcommand{\arraystretch}{1.1}\rowcolors{1}{routepale}{white}\noindent\begin{tabularx}{\linewidth}{@{}>{\raggedright\arraybackslash}p{.345\linewidth}>{\raggedright\arraybackslash}X@{}}
\RouteGroup{Identidad estructural}
\RouteRow{Clase y multisección}{quark\_\allowbreak{}down\_\allowbreak{}G3\quad /\quad quark\_\allowbreak{}down\_\allowbreak{}G3}
\RouteRow{Colección y sector}{quark de tipo abajo; quark}
\RouteRow{Clasificación física}{quark de tipo abajo}
\RouteRow{Generación, profundidad y color}{3; G3; R}
\RouteRow{Ruta primaria y órbita de inversión}{\(B_1\): no\qquad 2,\allowbreak{}5;\allowbreak{}8,\allowbreak{}5}
\RouteRow{Firma de clasificación}{30|\allowbreak{}-6|\allowbreak{}0|\allowbreak{}0|\allowbreak{}-12|\allowbreak{}+++-\kern0pt{}-\kern0pt{}-+++-\kern0pt{}-\kern0pt{}-}
\RouteGroup{Retornos, memoria y transporte}
\RouteRow{Retornos con fase}{clase: 27; órbita: 27; celda: 27}
\RouteRow{Retorno cinemático}{en 108: 54; extensión: 216}
\RouteRow{Giros y cambios de hoja}{71; 107}
\RouteRow{Acarreo y diferimiento}{deuda total: 3; final: 0; bloques diferidos: 2}
\RouteRow{Tríadas finales; persistencia de Pareto}{17; sí}
\RouteRow{\(K\longrightarrow K+9\)}{repeticiones: 9; cambios de memoria: 9}
\RouteGroup{Firma, fase y diferencias}
\RouteRow{\(n_A,\ s_C,\ \kappa_0,\ \nu_{120},\ \nu_{270}\)}{30, 6, 0, 4, -12}
\RouteRow{\(\tau_{12}\); firma histórica \(\mathbb Z^5\)}{+++-\kern0pt{}-\kern0pt{}-+++-\kern0pt{}-\kern0pt{}-; 30|\allowbreak{}6|\allowbreak{}0|\allowbreak{}4|\allowbreak{}-12}
\RouteRow{\(D_1,\ D_3,\ D_4,\ D_{27},\ D_{36}\)}{54, 56, 68, 48, 32}
\RouteRow{\(\Omega_{90/120},\ P_6\)}{80, 6}
\RouteGroup{Lectores y factores de acción}
\RouteRow{\(K_D\quad /\quad K_\Omega\)}{(80,\allowbreak{}54,\allowbreak{}6)\quad /\quad (80,\allowbreak{}54,\allowbreak{}6)}
\RouteRow{\(\kappa_D\)}{79.60661\allowbreak{}01397948\allowbreak{}27586470\allowbreak{}91}
\RouteRow{\(\kappa_\Omega\)}{79.60661\allowbreak{}01397948\allowbreak{}27586470\allowbreak{}91}
\RouteRow{\(R_{\mathrm{act},D}\)}{1.000000\allowbreak{}05518981\allowbreak{}25318591\allowbreak{}40}
\RouteRow{\(R_{\mathrm{act},\Omega}\)}{1.000000\allowbreak{}05518981\allowbreak{}25318591\allowbreak{}40}
\RouteGroup{Separación de prefijos}
\RouteRow{Área de ambigüedad; área \(\log_2\)}{36; 18.0}
\RouteRow{Primer bloque de separación}{órbita: 1; celda: \textemdash{}; ruta: \textemdash{}}
\RouteRow{Ambigüedad final}{2}
\RouteGroup{Secuencias completas}
\RouteRow{\(w_6\)}{423 15 486 423 9 24 243 657 24 423 15 486 423 9 24 243 657 24}
\RouteRow{Tríadas estables por bloque}{0 1 2 3 4 5 6 7 8 9 10 11 11 13 14 14 15 17}
\end{tabularx}\endgroup\par\vspace{5pt}\noindent{\fontsize{8.4}{10}\selectfont\hyperref[trace-53]{Procedencia y códigos originales}\hfill Ficha completa · 053/324}
\clearpage\phantomsection\label{nav:vi-ruta-54}
% VI-CATALOG-ROW rutas 54
\pdfbookmark[3]{r2c5\_h-1v+1}{route-54}
\RouteTitle{Ruta 054}{r2c5\_\allowbreak{}h-1v+1}
\noindent Celda 14\quad (2, 5)\qquad Orientación \(h=-1,\ v=1\)\hfill\hyperref[sec:vi-indice-rutas]{Índice}\par\vspace{4pt}
\begingroup\fontsize{10.2}{12.5}\selectfont\setlength{\tabcolsep}{5pt}\renewcommand{\arraystretch}{1.1}\rowcolors{1}{routepale}{white}\noindent\begin{tabularx}{\linewidth}{@{}>{\raggedright\arraybackslash}p{.345\linewidth}>{\raggedright\arraybackslash}X@{}}
\RouteGroup{Identidad estructural}
\RouteRow{Clase y multisección}{quark\_\allowbreak{}down\_\allowbreak{}G3\quad /\quad quark\_\allowbreak{}down\_\allowbreak{}G3}
\RouteRow{Colección y sector}{quark de tipo abajo; quark}
\RouteRow{Clasificación física}{quark de tipo abajo}
\RouteRow{Generación, profundidad y color}{3; G3; R}
\RouteRow{Ruta primaria y órbita de inversión}{\(B_1\): no\qquad 2,\allowbreak{}5;\allowbreak{}8,\allowbreak{}5}
\RouteRow{Firma de clasificación}{30|\allowbreak{}-6|\allowbreak{}0|\allowbreak{}0|\allowbreak{}-12|\allowbreak{}+++-\kern0pt{}-\kern0pt{}-+++-\kern0pt{}-\kern0pt{}-}
\RouteGroup{Retornos, memoria y transporte}
\RouteRow{Retornos con fase}{clase: 27; órbita: 27; celda: 27}
\RouteRow{Retorno cinemático}{en 108: 54; extensión: 216}
\RouteRow{Giros y cambios de hoja}{71; 107}
\RouteRow{Acarreo y diferimiento}{deuda total: 2; final: 1; bloques diferidos: 2}
\RouteRow{Tríadas finales; persistencia de Pareto}{16; sí}
\RouteRow{\(K\longrightarrow K+9\)}{repeticiones: 9; cambios de memoria: 9}
\RouteGroup{Firma, fase y diferencias}
\RouteRow{\(n_A,\ s_C,\ \kappa_0,\ \nu_{120},\ \nu_{270}\)}{50, 4, -4, 0, -4}
\RouteRow{\(\tau_{12}\); firma histórica \(\mathbb Z^5\)}{+++00-+++00-; 50|\allowbreak{}4|\allowbreak{}-4|\allowbreak{}0|\allowbreak{}-4}
\RouteRow{\(D_1,\ D_3,\ D_4,\ D_{27},\ D_{36}\)}{84, 28, 84, 36, 24}
\RouteRow{\(\Omega_{90/120},\ P_6\)}{20, 4}
\RouteGroup{Lectores y factores de acción}
\RouteRow{\(K_D\quad /\quad K_\Omega\)}{(60,\allowbreak{}84,\allowbreak{}28)\quad /\quad (20,\allowbreak{}54,\allowbreak{}4)}
\RouteRow{\(\kappa_D\)}{59.39013\allowbreak{}58840155\allowbreak{}40637741\allowbreak{}48}
\RouteRow{\(\kappa_\Omega\)}{19.60638\allowbreak{}77469494\allowbreak{}43306362\allowbreak{}81}
\RouteRow{\(R_{\mathrm{act},D}\)}{1.000000\allowbreak{}04117409\allowbreak{}89467415\allowbreak{}75}
\RouteRow{\(R_{\mathrm{act},\Omega}\)}{1.000000\allowbreak{}01359275\allowbreak{}11518874\allowbreak{}98}
\RouteGroup{Separación de prefijos}
\RouteRow{Área de ambigüedad; área \(\log_2\)}{36; 18.0}
\RouteRow{Primer bloque de separación}{órbita: 1; celda: \textemdash{}; ruta: \textemdash{}}
\RouteRow{Ambigüedad final}{2}
\RouteGroup{Secuencias completas}
\RouteRow{\(w_6\)}{414 420 258 414 414 333 255 252 333 414 420 258 414 414 333 255 252 333}
\RouteRow{Tríadas estables por bloque}{0 1 2 3 4 5 6 7 8 9 10 11 12 13 14 14 16 16}
\end{tabularx}\endgroup\par\vspace{5pt}\noindent{\fontsize{8.4}{10}\selectfont\hyperref[trace-54]{Procedencia y códigos originales}\hfill Ficha completa · 054/324}
\clearpage\phantomsection\label{nav:vi-ruta-55}
% VI-CATALOG-ROW rutas 55
\pdfbookmark[3]{r2c5\_h+1v-1}{route-55}
\RouteTitle{Ruta 055}{r2c5\_\allowbreak{}h+1v-1}
\noindent Celda 14\quad (2, 5)\qquad Orientación \(h=1,\ v=-1\)\hfill\hyperref[sec:vi-indice-rutas]{Índice}\par\vspace{4pt}
\begingroup\fontsize{10.2}{12.5}\selectfont\setlength{\tabcolsep}{5pt}\renewcommand{\arraystretch}{1.1}\rowcolors{1}{routepale}{white}\noindent\begin{tabularx}{\linewidth}{@{}>{\raggedright\arraybackslash}p{.345\linewidth}>{\raggedright\arraybackslash}X@{}}
\RouteGroup{Identidad estructural}
\RouteRow{Clase y multisección}{quark\_\allowbreak{}down\_\allowbreak{}G3\quad /\quad quark\_\allowbreak{}down\_\allowbreak{}G3}
\RouteRow{Colección y sector}{quark de tipo abajo; quark}
\RouteRow{Clasificación física}{quark de tipo abajo}
\RouteRow{Generación, profundidad y color}{3; G3; R}
\RouteRow{Ruta primaria y órbita de inversión}{\(B_1\): no\qquad 2,\allowbreak{}5;\allowbreak{}8,\allowbreak{}5}
\RouteRow{Firma de clasificación}{30|\allowbreak{}-6|\allowbreak{}0|\allowbreak{}0|\allowbreak{}-12|\allowbreak{}+++-\kern0pt{}-\kern0pt{}-+++-\kern0pt{}-\kern0pt{}-}
\RouteGroup{Retornos, memoria y transporte}
\RouteRow{Retornos con fase}{clase: 27; órbita: 27; celda: 27}
\RouteRow{Retorno cinemático}{en 108: 54; extensión: 216}
\RouteRow{Giros y cambios de hoja}{71; 107}
\RouteRow{Acarreo y diferimiento}{deuda total: 3; final: 0; bloques diferidos: 2}
\RouteRow{Tríadas finales; persistencia de Pareto}{17; sí}
\RouteRow{\(K\longrightarrow K+9\)}{repeticiones: 9; cambios de memoria: 9}
\RouteGroup{Firma, fase y diferencias}
\RouteRow{\(n_A,\ s_C,\ \kappa_0,\ \nu_{120},\ \nu_{270}\)}{50, -4, 4, 0, -4}
\RouteRow{\(\tau_{12}\); firma histórica \(\mathbb Z^5\)}{00+-\kern0pt{}-\kern0pt{}-00+-\kern0pt{}-\kern0pt{}-; 50|\allowbreak{}-4|\allowbreak{}4|\allowbreak{}0|\allowbreak{}-4}
\RouteRow{\(D_1,\ D_3,\ D_4,\ D_{27},\ D_{36}\)}{84, 28, 84, 36, 24}
\RouteRow{\(\Omega_{90/120},\ P_6\)}{20, 4}
\RouteGroup{Lectores y factores de acción}
\RouteRow{\(K_D\quad /\quad K_\Omega\)}{(60,\allowbreak{}84,\allowbreak{}28)\quad /\quad (20,\allowbreak{}54,\allowbreak{}4)}
\RouteRow{\(\kappa_D\)}{59.39013\allowbreak{}58840155\allowbreak{}40637741\allowbreak{}48}
\RouteRow{\(\kappa_\Omega\)}{19.60638\allowbreak{}77469494\allowbreak{}43306362\allowbreak{}81}
\RouteRow{\(R_{\mathrm{act},D}\)}{1.000000\allowbreak{}04117409\allowbreak{}89467415\allowbreak{}75}
\RouteRow{\(R_{\mathrm{act},\Omega}\)}{1.000000\allowbreak{}01359275\allowbreak{}11518874\allowbreak{}98}
\RouteGroup{Separación de prefijos}
\RouteRow{Área de ambigüedad; área \(\log_2\)}{36; 18.0}
\RouteRow{Primer bloque de separación}{órbita: 1; celda: \textemdash{}; ruta: \textemdash{}}
\RouteRow{Ambigüedad final}{2}
\RouteGroup{Secuencias completas}
\RouteRow{\(w_6\)}{255 252 333 255 258 258 414 420 258 255 252 333 255 258 258 414 420 258}
\RouteRow{Tríadas estables por bloque}{0 1 2 3 4 5 6 7 8 9 10 11 11 12 14 15 15 17}
\end{tabularx}\endgroup\par\vspace{5pt}\noindent{\fontsize{8.4}{10}\selectfont\hyperref[trace-55]{Procedencia y códigos originales}\hfill Ficha completa · 055/324}
\clearpage\phantomsection\label{nav:vi-ruta-56}
% VI-CATALOG-ROW rutas 56
\pdfbookmark[3]{r2c5\_h+1v+1}{route-56}
\RouteTitle{Ruta 056}{r2c5\_\allowbreak{}h+1v+1}
\noindent Celda 14\quad (2, 5)\qquad Orientación \(h=1,\ v=1\)\hfill\hyperref[sec:vi-indice-rutas]{Índice}\par\vspace{4pt}
\begingroup\fontsize{10.2}{12.5}\selectfont\setlength{\tabcolsep}{5pt}\renewcommand{\arraystretch}{1.1}\rowcolors{1}{routepale}{white}\noindent\begin{tabularx}{\linewidth}{@{}>{\raggedright\arraybackslash}p{.345\linewidth}>{\raggedright\arraybackslash}X@{}}
\RouteGroup{Identidad estructural}
\RouteRow{Clase y multisección}{quark\_\allowbreak{}down\_\allowbreak{}G3\quad /\quad quark\_\allowbreak{}down\_\allowbreak{}G3}
\RouteRow{Colección y sector}{quark de tipo abajo; quark}
\RouteRow{Clasificación física}{quark de tipo abajo}
\RouteRow{Generación, profundidad y color}{3; G3; R}
\RouteRow{Ruta primaria y órbita de inversión}{\(B_1\): sí\qquad 2,\allowbreak{}5;\allowbreak{}8,\allowbreak{}5}
\RouteRow{Firma de clasificación}{30|\allowbreak{}-6|\allowbreak{}0|\allowbreak{}0|\allowbreak{}-12|\allowbreak{}+++-\kern0pt{}-\kern0pt{}-+++-\kern0pt{}-\kern0pt{}-}
\RouteGroup{Retornos, memoria y transporte}
\RouteRow{Retornos con fase}{clase: 27; órbita: 27; celda: 27}
\RouteRow{Retorno cinemático}{en 108: 54; extensión: 216}
\RouteRow{Giros y cambios de hoja}{71; 107}
\RouteRow{Acarreo y diferimiento}{deuda total: 2; final: 1; bloques diferidos: 2}
\RouteRow{Tríadas finales; persistencia de Pareto}{16; sí}
\RouteRow{\(K\longrightarrow K+9\)}{repeticiones: 9; cambios de memoria: 9}
\RouteGroup{Firma, fase y diferencias}
\RouteRow{\(n_A,\ s_C,\ \kappa_0,\ \nu_{120},\ \nu_{270}\)}{30, -6, 0, -4, -12}
\RouteRow{\(\tau_{12}\); firma histórica \(\mathbb Z^5\)}{+++-\kern0pt{}-\kern0pt{}-+++-\kern0pt{}-\kern0pt{}-; 30|\allowbreak{}-6|\allowbreak{}0|\allowbreak{}-4|\allowbreak{}-12}
\RouteRow{\(D_1,\ D_3,\ D_4,\ D_{27},\ D_{36}\)}{54, 56, 68, 48, 32}
\RouteRow{\(\Omega_{90/120},\ P_6\)}{80, 6}
\RouteGroup{Lectores y factores de acción}
\RouteRow{\(K_D\quad /\quad K_\Omega\)}{(80,\allowbreak{}54,\allowbreak{}6)\quad /\quad (80,\allowbreak{}54,\allowbreak{}6)}
\RouteRow{\(\kappa_D\)}{79.60661\allowbreak{}01397948\allowbreak{}27586470\allowbreak{}91}
\RouteRow{\(\kappa_\Omega\)}{79.60661\allowbreak{}01397948\allowbreak{}27586470\allowbreak{}91}
\RouteRow{\(R_{\mathrm{act},D}\)}{1.000000\allowbreak{}05518981\allowbreak{}25318591\allowbreak{}40}
\RouteRow{\(R_{\mathrm{act},\Omega}\)}{1.000000\allowbreak{}05518981\allowbreak{}25318591\allowbreak{}40}
\RouteGroup{Separación de prefijos}
\RouteRow{Área de ambigüedad; área \(\log_2\)}{36; 18.0}
\RouteRow{Primer bloque de separación}{órbita: 1; celda: \textemdash{}; ruta: \textemdash{}}
\RouteRow{Ambigüedad final}{2}
\RouteGroup{Secuencias completas}
\RouteRow{\(w_6\)}{243 657 24 243 663 486 423 15 486 243 657 24 243 663 486 423 15 486}
\RouteRow{Tríadas estables por bloque}{0 1 2 3 4 5 6 7 8 9 10 11 12 13 14 14 16 16}
\end{tabularx}\endgroup\par\vspace{5pt}\noindent{\fontsize{8.4}{10}\selectfont\hyperref[trace-56]{Procedencia y códigos originales}\hfill Ficha completa · 056/324}
\clearpage\phantomsection\label{nav:vi-ruta-57}
% VI-CATALOG-ROW rutas 57
\pdfbookmark[2]{Celda 15}{cell-15}
\pdfbookmark[3]{r2c6\_h-1v-1}{route-57}
\RouteTitle{Ruta 057}{r2c6\_\allowbreak{}h-1v-1}
\noindent Celda 15\quad (2, 6)\qquad Orientación \(h=-1,\ v=-1\)\hfill\hyperref[sec:vi-indice-rutas]{Índice}\par\vspace{4pt}
\begingroup\fontsize{10.2}{12.5}\selectfont\setlength{\tabcolsep}{5pt}\renewcommand{\arraystretch}{1.1}\rowcolors{1}{routepale}{white}\noindent\begin{tabularx}{\linewidth}{@{}>{\raggedright\arraybackslash}p{.345\linewidth}>{\raggedright\arraybackslash}X@{}}
\RouteGroup{Identidad estructural}
\RouteRow{Clase y multisección}{quark\_\allowbreak{}up\_\allowbreak{}G3\quad /\quad quark\_\allowbreak{}up\_\allowbreak{}G3}
\RouteRow{Colección y sector}{quark de tipo arriba; quark}
\RouteRow{Clasificación física}{quark de tipo arriba}
\RouteRow{Generación, profundidad y color}{3; G3; B}
\RouteRow{Ruta primaria y órbita de inversión}{\(B_1\): no\qquad 2,\allowbreak{}6;\allowbreak{}8,\allowbreak{}4}
\RouteRow{Firma de clasificación}{34|\allowbreak{}0|\allowbreak{}-16|\allowbreak{}0|\allowbreak{}-8|\allowbreak{}+0+-\kern0pt{}-\kern0pt{}-+0+-\kern0pt{}-\kern0pt{}-}
\RouteGroup{Retornos, memoria y transporte}
\RouteRow{Retornos con fase}{clase: 27; órbita: 27; celda: 27}
\RouteRow{Retorno cinemático}{en 108: 54; extensión: 216}
\RouteRow{Giros y cambios de hoja}{71; 107}
\RouteRow{Acarreo y diferimiento}{deuda total: 1; final: 0; bloques diferidos: 1}
\RouteRow{Tríadas finales; persistencia de Pareto}{17; no}
\RouteRow{\(K\longrightarrow K+9\)}{repeticiones: 9; cambios de memoria: 9}
\RouteGroup{Firma, fase y diferencias}
\RouteRow{\(n_A,\ s_C,\ \kappa_0,\ \nu_{120},\ \nu_{270}\)}{34, 0, -2, 0, -8}
\RouteRow{\(\tau_{12}\); firma histórica \(\mathbb Z^5\)}{+++-0-+++-0-; 34|\allowbreak{}0|\allowbreak{}-2|\allowbreak{}0|\allowbreak{}-8}
\RouteRow{\(D_1,\ D_3,\ D_4,\ D_{27},\ D_{36}\)}{66, 60, 68, 48, 32}
\RouteRow{\(\Omega_{90/120},\ P_6\)}{48, 5}
\RouteGroup{Lectores y factores de acción}
\RouteRow{\(K_D\quad /\quad K_\Omega\)}{(80,\allowbreak{}66,\allowbreak{}4)\quad /\quad (48,\allowbreak{}54,\allowbreak{}5)}
\RouteRow{\(\kappa_D\)}{79.51881\allowbreak{}95161180\allowbreak{}37694395\allowbreak{}39}
\RouteRow{\(\kappa_\Omega\)}{47.60649\allowbreak{}89433721\allowbreak{}35446416\allowbreak{}86}
\RouteRow{\(R_{\mathrm{act},D}\)}{1.000000\allowbreak{}05512894\allowbreak{}88901612\allowbreak{}64}
\RouteRow{\(R_{\mathrm{act},\Omega}\)}{1.000000\allowbreak{}03300471\allowbreak{}80532430\allowbreak{}84}
\RouteGroup{Separación de prefijos}
\RouteRow{Área de ambigüedad; área \(\log_2\)}{36; 18.0}
\RouteRow{Primer bloque de separación}{órbita: \textemdash{}; celda: \textemdash{}; ruta: \textemdash{}}
\RouteRow{Ambigüedad final}{2}
\RouteGroup{Secuencias completas}
\RouteRow{\(w_6\)}{570 264 90 570 267 243 657 24 243 570 264 90 570 267 243 657 24 243}
\RouteRow{Tríadas estables por bloque}{0 1 1 3 4 5 6 7 8 9 10 11 12 13 14 15 16 17}
\end{tabularx}\endgroup\par\vspace{5pt}\noindent{\fontsize{8.4}{10}\selectfont\hyperref[trace-57]{Procedencia y códigos originales}\hfill Ficha completa · 057/324}
\clearpage\phantomsection\label{nav:vi-ruta-58}
% VI-CATALOG-ROW rutas 58
\pdfbookmark[3]{r2c6\_h-1v+1}{route-58}
\RouteTitle{Ruta 058}{r2c6\_\allowbreak{}h-1v+1}
\noindent Celda 15\quad (2, 6)\qquad Orientación \(h=-1,\ v=1\)\hfill\hyperref[sec:vi-indice-rutas]{Índice}\par\vspace{4pt}
\begingroup\fontsize{10.2}{12.5}\selectfont\setlength{\tabcolsep}{5pt}\renewcommand{\arraystretch}{1.1}\rowcolors{1}{routepale}{white}\noindent\begin{tabularx}{\linewidth}{@{}>{\raggedright\arraybackslash}p{.345\linewidth}>{\raggedright\arraybackslash}X@{}}
\RouteGroup{Identidad estructural}
\RouteRow{Clase y multisección}{quark\_\allowbreak{}up\_\allowbreak{}G3\quad /\quad quark\_\allowbreak{}up\_\allowbreak{}G3}
\RouteRow{Colección y sector}{quark de tipo arriba; quark}
\RouteRow{Clasificación física}{quark de tipo arriba}
\RouteRow{Generación, profundidad y color}{3; G3; B}
\RouteRow{Ruta primaria y órbita de inversión}{\(B_1\): sí\qquad 2,\allowbreak{}6;\allowbreak{}8,\allowbreak{}4}
\RouteRow{Firma de clasificación}{34|\allowbreak{}0|\allowbreak{}-16|\allowbreak{}0|\allowbreak{}-8|\allowbreak{}+0+-\kern0pt{}-\kern0pt{}-+0+-\kern0pt{}-\kern0pt{}-}
\RouteGroup{Retornos, memoria y transporte}
\RouteRow{Retornos con fase}{clase: 27; órbita: 27; celda: 27}
\RouteRow{Retorno cinemático}{en 108: 54; extensión: 216}
\RouteRow{Giros y cambios de hoja}{71; 107}
\RouteRow{Acarreo y diferimiento}{deuda total: 1; final: 0; bloques diferidos: 1}
\RouteRow{Tríadas finales; persistencia de Pareto}{17; no}
\RouteRow{\(K\longrightarrow K+9\)}{repeticiones: 9; cambios de memoria: 9}
\RouteGroup{Firma, fase y diferencias}
\RouteRow{\(n_A,\ s_C,\ \kappa_0,\ \nu_{120},\ \nu_{270}\)}{14, -4, 0, 0, -12}
\RouteRow{\(\tau_{12}\); firma histórica \(\mathbb Z^5\)}{+++-\kern0pt{}-\kern0pt{}-+++-\kern0pt{}-\kern0pt{}-; 14|\allowbreak{}-4|\allowbreak{}0|\allowbreak{}0|\allowbreak{}-12}
\RouteRow{\(D_1,\ D_3,\ D_4,\ D_{27},\ D_{36}\)}{78, 24, 78, 24, 16}
\RouteRow{\(\Omega_{90/120},\ P_6\)}{80, 6}
\RouteGroup{Lectores y factores de acción}
\RouteRow{\(K_D\quad /\quad K_\Omega\)}{(40,\allowbreak{}78,\allowbreak{}27)\quad /\quad (80,\allowbreak{}54,\allowbreak{}6)}
\RouteRow{\(\kappa_D\)}{39.43380\allowbreak{}88030085\allowbreak{}51303671\allowbreak{}14}
\RouteRow{\(\kappa_\Omega\)}{79.60661\allowbreak{}01397948\allowbreak{}27586470\allowbreak{}91}
\RouteRow{\(R_{\mathrm{act},D}\)}{1.000000\allowbreak{}02733874\allowbreak{}08548943\allowbreak{}58}
\RouteRow{\(R_{\mathrm{act},\Omega}\)}{1.000000\allowbreak{}05518981\allowbreak{}25318591\allowbreak{}40}
\RouteGroup{Separación de prefijos}
\RouteRow{Área de ambigüedad; área \(\log_2\)}{36; 18.0}
\RouteRow{Primer bloque de separación}{órbita: \textemdash{}; celda: \textemdash{}; ruta: \textemdash{}}
\RouteRow{Ambigüedad final}{2}
\RouteGroup{Secuencias completas}
\RouteRow{\(w_6\)}{585 507 504 585 510 666 672 510 666 585 507 504 585 510 666 672 510 666}
\RouteRow{Tríadas estables por bloque}{0 1 2 3 4 5 6 7 8 9 10 11 12 13 14 14 16 17}
\end{tabularx}\endgroup\par\vspace{5pt}\noindent{\fontsize{8.4}{10}\selectfont\hyperref[trace-58]{Procedencia y códigos originales}\hfill Ficha completa · 058/324}
\clearpage\phantomsection\label{nav:vi-ruta-59}
% VI-CATALOG-ROW rutas 59
\pdfbookmark[3]{r2c6\_h+1v-1}{route-59}
\RouteTitle{Ruta 059}{r2c6\_\allowbreak{}h+1v-1}
\noindent Celda 15\quad (2, 6)\qquad Orientación \(h=1,\ v=-1\)\hfill\hyperref[sec:vi-indice-rutas]{Índice}\par\vspace{4pt}
\begingroup\fontsize{10.2}{12.5}\selectfont\setlength{\tabcolsep}{5pt}\renewcommand{\arraystretch}{1.1}\rowcolors{1}{routepale}{white}\noindent\begin{tabularx}{\linewidth}{@{}>{\raggedright\arraybackslash}p{.345\linewidth}>{\raggedright\arraybackslash}X@{}}
\RouteGroup{Identidad estructural}
\RouteRow{Clase y multisección}{quark\_\allowbreak{}up\_\allowbreak{}G3\quad /\quad quark\_\allowbreak{}up\_\allowbreak{}G3}
\RouteRow{Colección y sector}{quark de tipo arriba; quark}
\RouteRow{Clasificación física}{quark de tipo arriba}
\RouteRow{Generación, profundidad y color}{3; G3; B}
\RouteRow{Ruta primaria y órbita de inversión}{\(B_1\): no\qquad 2,\allowbreak{}6;\allowbreak{}8,\allowbreak{}4}
\RouteRow{Firma de clasificación}{34|\allowbreak{}0|\allowbreak{}-16|\allowbreak{}0|\allowbreak{}-8|\allowbreak{}+0+-\kern0pt{}-\kern0pt{}-+0+-\kern0pt{}-\kern0pt{}-}
\RouteGroup{Retornos, memoria y transporte}
\RouteRow{Retornos con fase}{clase: 27; órbita: 27; celda: 27}
\RouteRow{Retorno cinemático}{en 108: 54; extensión: 216}
\RouteRow{Giros y cambios de hoja}{71; 107}
\RouteRow{Acarreo y diferimiento}{deuda total: 1; final: 0; bloques diferidos: 1}
\RouteRow{Tríadas finales; persistencia de Pareto}{17; no}
\RouteRow{\(K\longrightarrow K+9\)}{repeticiones: 9; cambios de memoria: 9}
\RouteGroup{Firma, fase y diferencias}
\RouteRow{\(n_A,\ s_C,\ \kappa_0,\ \nu_{120},\ \nu_{270}\)}{14, 4, 0, 0, -12}
\RouteRow{\(\tau_{12}\); firma histórica \(\mathbb Z^5\)}{+++-\kern0pt{}-\kern0pt{}-+++-\kern0pt{}-\kern0pt{}-; 14|\allowbreak{}4|\allowbreak{}0|\allowbreak{}0|\allowbreak{}-12}
\RouteRow{\(D_1,\ D_3,\ D_4,\ D_{27},\ D_{36}\)}{78, 24, 78, 24, 16}
\RouteRow{\(\Omega_{90/120},\ P_6\)}{80, 6}
\RouteGroup{Lectores y factores de acción}
\RouteRow{\(K_D\quad /\quad K_\Omega\)}{(40,\allowbreak{}78,\allowbreak{}27)\quad /\quad (80,\allowbreak{}54,\allowbreak{}6)}
\RouteRow{\(\kappa_D\)}{39.43380\allowbreak{}88030085\allowbreak{}51303671\allowbreak{}14}
\RouteRow{\(\kappa_\Omega\)}{79.60661\allowbreak{}01397948\allowbreak{}27586470\allowbreak{}91}
\RouteRow{\(R_{\mathrm{act},D}\)}{1.000000\allowbreak{}02733874\allowbreak{}08548943\allowbreak{}58}
\RouteRow{\(R_{\mathrm{act},\Omega}\)}{1.000000\allowbreak{}05518981\allowbreak{}25318591\allowbreak{}40}
\RouteGroup{Separación de prefijos}
\RouteRow{Área de ambigüedad; área \(\log_2\)}{36; 18.0}
\RouteRow{Primer bloque de separación}{órbita: \textemdash{}; celda: \textemdash{}; ruta: \textemdash{}}
\RouteRow{Ambigüedad final}{2}
\RouteGroup{Secuencias completas}
\RouteRow{\(w_6\)}{672 510 666 672 507 504 585 507 504 672 510 666 672 507 504 585 507 504}
\RouteRow{Tríadas estables por bloque}{0 1 2 3 4 5 6 7 8 9 10 11 12 13 14 14 16 17}
\end{tabularx}\endgroup\par\vspace{5pt}\noindent{\fontsize{8.4}{10}\selectfont\hyperref[trace-59]{Procedencia y códigos originales}\hfill Ficha completa · 059/324}
\clearpage\phantomsection\label{nav:vi-ruta-60}
% VI-CATALOG-ROW rutas 60
\pdfbookmark[3]{r2c6\_h+1v+1}{route-60}
\RouteTitle{Ruta 060}{r2c6\_\allowbreak{}h+1v+1}
\noindent Celda 15\quad (2, 6)\qquad Orientación \(h=1,\ v=1\)\hfill\hyperref[sec:vi-indice-rutas]{Índice}\par\vspace{4pt}
\begingroup\fontsize{10.2}{12.5}\selectfont\setlength{\tabcolsep}{5pt}\renewcommand{\arraystretch}{1.1}\rowcolors{1}{routepale}{white}\noindent\begin{tabularx}{\linewidth}{@{}>{\raggedright\arraybackslash}p{.345\linewidth}>{\raggedright\arraybackslash}X@{}}
\RouteGroup{Identidad estructural}
\RouteRow{Clase y multisección}{quark\_\allowbreak{}up\_\allowbreak{}G3\quad /\quad quark\_\allowbreak{}up\_\allowbreak{}G3}
\RouteRow{Colección y sector}{quark de tipo arriba; quark}
\RouteRow{Clasificación física}{quark de tipo arriba}
\RouteRow{Generación, profundidad y color}{3; G3; B}
\RouteRow{Ruta primaria y órbita de inversión}{\(B_1\): no\qquad 2,\allowbreak{}6;\allowbreak{}8,\allowbreak{}4}
\RouteRow{Firma de clasificación}{34|\allowbreak{}0|\allowbreak{}-16|\allowbreak{}0|\allowbreak{}-8|\allowbreak{}+0+-\kern0pt{}-\kern0pt{}-+0+-\kern0pt{}-\kern0pt{}-}
\RouteGroup{Retornos, memoria y transporte}
\RouteRow{Retornos con fase}{clase: 27; órbita: 27; celda: 27}
\RouteRow{Retorno cinemático}{en 108: 54; extensión: 216}
\RouteRow{Giros y cambios de hoja}{71; 107}
\RouteRow{Acarreo y diferimiento}{deuda total: 0; final: 0; bloques diferidos: 0}
\RouteRow{Tríadas finales; persistencia de Pareto}{17; no}
\RouteRow{\(K\longrightarrow K+9\)}{repeticiones: 9; cambios de memoria: 9}
\RouteGroup{Firma, fase y diferencias}
\RouteRow{\(n_A,\ s_C,\ \kappa_0,\ \nu_{120},\ \nu_{270}\)}{34, 0, 2, 0, -8}
\RouteRow{\(\tau_{12}\); firma histórica \(\mathbb Z^5\)}{+0+-\kern0pt{}-\kern0pt{}-+0+-\kern0pt{}-\kern0pt{}-; 34|\allowbreak{}0|\allowbreak{}2|\allowbreak{}0|\allowbreak{}-8}
\RouteRow{\(D_1,\ D_3,\ D_4,\ D_{27},\ D_{36}\)}{66, 60, 68, 48, 32}
\RouteRow{\(\Omega_{90/120},\ P_6\)}{48, 5}
\RouteGroup{Lectores y factores de acción}
\RouteRow{\(K_D\quad /\quad K_\Omega\)}{(80,\allowbreak{}66,\allowbreak{}4)\quad /\quad (48,\allowbreak{}54,\allowbreak{}5)}
\RouteRow{\(\kappa_D\)}{79.51881\allowbreak{}95161180\allowbreak{}37694395\allowbreak{}39}
\RouteRow{\(\kappa_\Omega\)}{47.60649\allowbreak{}89433721\allowbreak{}35446416\allowbreak{}86}
\RouteRow{\(R_{\mathrm{act},D}\)}{1.000000\allowbreak{}05512894\allowbreak{}88901612\allowbreak{}64}
\RouteRow{\(R_{\mathrm{act},\Omega}\)}{1.000000\allowbreak{}03300471\allowbreak{}80532430\allowbreak{}84}
\RouteGroup{Separación de prefijos}
\RouteRow{Área de ambigüedad; área \(\log_2\)}{36; 18.0}
\RouteRow{Primer bloque de separación}{órbita: \textemdash{}; celda: \textemdash{}; ruta: \textemdash{}}
\RouteRow{Ambigüedad final}{2}
\RouteGroup{Secuencias completas}
\RouteRow{\(w_6\)}{657 24 243 657 21 90 570 264 90 657 24 243 657 21 90 570 264 90}
\RouteRow{Tríadas estables por bloque}{0 1 2 3 4 5 6 7 8 9 10 11 12 13 14 15 16 17}
\end{tabularx}\endgroup\par\vspace{5pt}\noindent{\fontsize{8.4}{10}\selectfont\hyperref[trace-60]{Procedencia y códigos originales}\hfill Ficha completa · 060/324}
\clearpage\phantomsection\label{nav:vi-ruta-61}
% VI-CATALOG-ROW rutas 61
\pdfbookmark[2]{Celda 16}{cell-16}
\pdfbookmark[3]{r2c7\_h-1v-1}{route-61}
\RouteTitle{Ruta 061}{r2c7\_\allowbreak{}h-1v-1}
\noindent Celda 16\quad (2, 7)\qquad Orientación \(h=-1,\ v=-1\)\hfill\hyperref[sec:vi-indice-rutas]{Índice}\par\vspace{4pt}
\begingroup\fontsize{10.2}{12.5}\selectfont\setlength{\tabcolsep}{5pt}\renewcommand{\arraystretch}{1.1}\rowcolors{1}{routepale}{white}\noindent\begin{tabularx}{\linewidth}{@{}>{\raggedright\arraybackslash}p{.345\linewidth}>{\raggedright\arraybackslash}X@{}}
\RouteGroup{Identidad estructural}
\RouteRow{Clase y multisección}{quark\_\allowbreak{}down\_\allowbreak{}G3\quad /\quad quark\_\allowbreak{}down\_\allowbreak{}G3}
\RouteRow{Colección y sector}{quark de tipo abajo; quark}
\RouteRow{Clasificación física}{quark de tipo abajo}
\RouteRow{Generación, profundidad y color}{3; G3; G}
\RouteRow{Ruta primaria y órbita de inversión}{\(B_1\): no\qquad 2,\allowbreak{}7;\allowbreak{}8,\allowbreak{}3}
\RouteRow{Firma de clasificación}{34|\allowbreak{}0|\allowbreak{}0|\allowbreak{}2|\allowbreak{}-8|\allowbreak{}+++0-\kern0pt{}-+++0-\kern0pt{}-}
\RouteGroup{Retornos, memoria y transporte}
\RouteRow{Retornos con fase}{clase: 27; órbita: 27; celda: 27}
\RouteRow{Retorno cinemático}{en 108: 54; extensión: 216}
\RouteRow{Giros y cambios de hoja}{71; 107}
\RouteRow{Acarreo y diferimiento}{deuda total: 1; final: 1; bloques diferidos: 1}
\RouteRow{Tríadas finales; persistencia de Pareto}{16; sí}
\RouteRow{\(K\longrightarrow K+9\)}{repeticiones: 9; cambios de memoria: 9}
\RouteGroup{Firma, fase y diferencias}
\RouteRow{\(n_A,\ s_C,\ \kappa_0,\ \nu_{120},\ \nu_{270}\)}{34, 0, 2, 0, -8}
\RouteRow{\(\tau_{12}\); firma histórica \(\mathbb Z^5\)}{0++-\kern0pt{}-\kern0pt{}-0++-\kern0pt{}-\kern0pt{}-; 34|\allowbreak{}0|\allowbreak{}2|\allowbreak{}0|\allowbreak{}-8}
\RouteRow{\(D_1,\ D_3,\ D_4,\ D_{27},\ D_{36}\)}{66, 64, 68, 60, 40}
\RouteRow{\(\Omega_{90/120},\ P_6\)}{56, 5}
\RouteGroup{Lectores y factores de acción}
\RouteRow{\(K_D\quad /\quad K_\Omega\)}{(100,\allowbreak{}66,\allowbreak{}2)\quad /\quad (56,\allowbreak{}54,\allowbreak{}5)}
\RouteRow{\(\kappa_D\)}{99.51859\allowbreak{}71232726\allowbreak{}53414287\allowbreak{}29}
\RouteRow{\(\kappa_\Omega\)}{55.60649\allowbreak{}89433721\allowbreak{}35446416\allowbreak{}86}
\RouteRow{\(R_{\mathrm{act},D}\)}{1.000000\allowbreak{}06899443\allowbreak{}08259364\allowbreak{}17}
\RouteRow{\(R_{\mathrm{act},\Omega}\)}{1.000000\allowbreak{}03855097\allowbreak{}23541416\allowbreak{}45}
\RouteGroup{Separación de prefijos}
\RouteRow{Área de ambigüedad; área \(\log_2\)}{36; 18.0}
\RouteRow{Primer bloque de separación}{órbita: \textemdash{}; celda: \textemdash{}; ruta: \textemdash{}}
\RouteRow{Ambigüedad final}{2}
\RouteGroup{Secuencias completas}
\RouteRow{\(w_6\)}{15 486 423 15 489 498 99 327 498 15 486 423 15 489 498 99 327 498}
\RouteRow{Tríadas estables por bloque}{0 1 2 3 4 5 6 7 8 9 10 11 12 13 14 15 16 16}
\end{tabularx}\endgroup\par\vspace{5pt}\noindent{\fontsize{8.4}{10}\selectfont\hyperref[trace-61]{Procedencia y códigos originales}\hfill Ficha completa · 061/324}
\clearpage\phantomsection\label{nav:vi-ruta-62}
% VI-CATALOG-ROW rutas 62
\pdfbookmark[3]{r2c7\_h-1v+1}{route-62}
\RouteTitle{Ruta 062}{r2c7\_\allowbreak{}h-1v+1}
\noindent Celda 16\quad (2, 7)\qquad Orientación \(h=-1,\ v=1\)\hfill\hyperref[sec:vi-indice-rutas]{Índice}\par\vspace{4pt}
\begingroup\fontsize{10.2}{12.5}\selectfont\setlength{\tabcolsep}{5pt}\renewcommand{\arraystretch}{1.1}\rowcolors{1}{routepale}{white}\noindent\begin{tabularx}{\linewidth}{@{}>{\raggedright\arraybackslash}p{.345\linewidth}>{\raggedright\arraybackslash}X@{}}
\RouteGroup{Identidad estructural}
\RouteRow{Clase y multisección}{quark\_\allowbreak{}down\_\allowbreak{}G3\quad /\quad quark\_\allowbreak{}down\_\allowbreak{}G3}
\RouteRow{Colección y sector}{quark de tipo abajo; quark}
\RouteRow{Clasificación física}{quark de tipo abajo}
\RouteRow{Generación, profundidad y color}{3; G3; G}
\RouteRow{Ruta primaria y órbita de inversión}{\(B_1\): sí\qquad 2,\allowbreak{}7;\allowbreak{}8,\allowbreak{}3}
\RouteRow{Firma de clasificación}{34|\allowbreak{}0|\allowbreak{}0|\allowbreak{}2|\allowbreak{}-8|\allowbreak{}+++0-\kern0pt{}-+++0-\kern0pt{}-}
\RouteGroup{Retornos, memoria y transporte}
\RouteRow{Retornos con fase}{clase: 27; órbita: 27; celda: 27}
\RouteRow{Retorno cinemático}{en 108: 54; extensión: 216}
\RouteRow{Giros y cambios de hoja}{71; 107}
\RouteRow{Acarreo y diferimiento}{deuda total: 0; final: 0; bloques diferidos: 0}
\RouteRow{Tríadas finales; persistencia de Pareto}{17; sí}
\RouteRow{\(K\longrightarrow K+9\)}{repeticiones: 9; cambios de memoria: 9}
\RouteGroup{Firma, fase y diferencias}
\RouteRow{\(n_A,\ s_C,\ \kappa_0,\ \nu_{120},\ \nu_{270}\)}{20, 0, 0, 0, -4}
\RouteRow{\(\tau_{12}\); firma histórica \(\mathbb Z^5\)}{++00-\kern0pt{}-++00-\kern0pt{}-; 20|\allowbreak{}0|\allowbreak{}0|\allowbreak{}0|\allowbreak{}-4}
\RouteRow{\(D_1,\ D_3,\ D_4,\ D_{27},\ D_{36}\)}{24, 20, 24, 24, 16}
\RouteRow{\(\Omega_{90/120},\ P_6\)}{36, 4}
\RouteGroup{Lectores y factores de acción}
\RouteRow{\(K_D\quad /\quad K_\Omega\)}{(40,\allowbreak{}24,\allowbreak{}2)\quad /\quad (36,\allowbreak{}54,\allowbreak{}4)}
\RouteRow{\(\kappa_D\)}{39.82508\allowbreak{}59311825\allowbreak{}73056173\allowbreak{}26}
\RouteRow{\(\kappa_\Omega\)}{35.60638\allowbreak{}77469494\allowbreak{}43306362\allowbreak{}81}
\RouteRow{\(R_{\mathrm{act},D}\)}{1.000000\allowbreak{}02761000\allowbreak{}61595142\allowbreak{}15}
\RouteRow{\(R_{\mathrm{act},\Omega}\)}{1.000000\allowbreak{}02468525\allowbreak{}95691181\allowbreak{}52}
\RouteGroup{Separación de prefijos}
\RouteRow{Área de ambigüedad; área \(\log_2\)}{36; 18.0}
\RouteRow{Primer bloque de separación}{órbita: \textemdash{}; celda: \textemdash{}; ruta: \textemdash{}}
\RouteRow{Ambigüedad final}{2}
\RouteGroup{Secuencias completas}
\RouteRow{\(w_6\)}{0 81 3 0 84 0 81 3 0 0 81 3 0 84 0 81 3 0}
\RouteRow{Tríadas estables por bloque}{0 1 2 3 4 5 6 7 8 9 10 11 12 13 14 15 16 17}
\end{tabularx}\endgroup\par\vspace{5pt}\noindent{\fontsize{8.4}{10}\selectfont\hyperref[trace-62]{Procedencia y códigos originales}\hfill Ficha completa · 062/324}
\clearpage\phantomsection\label{nav:vi-ruta-63}
% VI-CATALOG-ROW rutas 63
\pdfbookmark[3]{r2c7\_h+1v-1}{route-63}
\RouteTitle{Ruta 063}{r2c7\_\allowbreak{}h+1v-1}
\noindent Celda 16\quad (2, 7)\qquad Orientación \(h=1,\ v=-1\)\hfill\hyperref[sec:vi-indice-rutas]{Índice}\par\vspace{4pt}
\begingroup\fontsize{10.2}{12.5}\selectfont\setlength{\tabcolsep}{5pt}\renewcommand{\arraystretch}{1.1}\rowcolors{1}{routepale}{white}\noindent\begin{tabularx}{\linewidth}{@{}>{\raggedright\arraybackslash}p{.345\linewidth}>{\raggedright\arraybackslash}X@{}}
\RouteGroup{Identidad estructural}
\RouteRow{Clase y multisección}{quark\_\allowbreak{}down\_\allowbreak{}G3\quad /\quad quark\_\allowbreak{}down\_\allowbreak{}G3}
\RouteRow{Colección y sector}{quark de tipo abajo; quark}
\RouteRow{Clasificación física}{quark de tipo abajo}
\RouteRow{Generación, profundidad y color}{3; G3; G}
\RouteRow{Ruta primaria y órbita de inversión}{\(B_1\): no\qquad 2,\allowbreak{}7;\allowbreak{}8,\allowbreak{}3}
\RouteRow{Firma de clasificación}{34|\allowbreak{}0|\allowbreak{}0|\allowbreak{}2|\allowbreak{}-8|\allowbreak{}+++0-\kern0pt{}-+++0-\kern0pt{}-}
\RouteGroup{Retornos, memoria y transporte}
\RouteRow{Retornos con fase}{clase: 27; órbita: 27; celda: 27}
\RouteRow{Retorno cinemático}{en 108: 54; extensión: 216}
\RouteRow{Giros y cambios de hoja}{71; 107}
\RouteRow{Acarreo y diferimiento}{deuda total: 2; final: 0; bloques diferidos: 2}
\RouteRow{Tríadas finales; persistencia de Pareto}{17; sí}
\RouteRow{\(K\longrightarrow K+9\)}{repeticiones: 9; cambios de memoria: 9}
\RouteGroup{Firma, fase y diferencias}
\RouteRow{\(n_A,\ s_C,\ \kappa_0,\ \nu_{120},\ \nu_{270}\)}{20, 0, 0, 0, -4}
\RouteRow{\(\tau_{12}\); firma histórica \(\mathbb Z^5\)}{0++-\kern0pt{}-00++-\kern0pt{}-0; 20|\allowbreak{}0|\allowbreak{}0|\allowbreak{}0|\allowbreak{}-4}
\RouteRow{\(D_1,\ D_3,\ D_4,\ D_{27},\ D_{36}\)}{24, 20, 24, 24, 16}
\RouteRow{\(\Omega_{90/120},\ P_6\)}{36, 4}
\RouteGroup{Lectores y factores de acción}
\RouteRow{\(K_D\quad /\quad K_\Omega\)}{(40,\allowbreak{}24,\allowbreak{}2)\quad /\quad (36,\allowbreak{}54,\allowbreak{}4)}
\RouteRow{\(\kappa_D\)}{39.82508\allowbreak{}59311825\allowbreak{}73056173\allowbreak{}26}
\RouteRow{\(\kappa_\Omega\)}{35.60638\allowbreak{}77469494\allowbreak{}43306362\allowbreak{}81}
\RouteRow{\(R_{\mathrm{act},D}\)}{1.000000\allowbreak{}02761000\allowbreak{}61595142\allowbreak{}15}
\RouteRow{\(R_{\mathrm{act},\Omega}\)}{1.000000\allowbreak{}02468525\allowbreak{}95691181\allowbreak{}52}
\RouteGroup{Separación de prefijos}
\RouteRow{Área de ambigüedad; área \(\log_2\)}{36; 18.0}
\RouteRow{Primer bloque de separación}{órbita: \textemdash{}; celda: \textemdash{}; ruta: \textemdash{}}
\RouteRow{Ambigüedad final}{2}
\RouteGroup{Secuencias completas}
\RouteRow{\(w_6\)}{81 3 0 81 0 3 0 81 3 81 3 0 81 0 3 0 81 3}
\RouteRow{Tríadas estables por bloque}{0 1 2 3 4 5 6 6 8 9 10 11 12 13 14 15 15 17}
\end{tabularx}\endgroup\par\vspace{5pt}\noindent{\fontsize{8.4}{10}\selectfont\hyperref[trace-63]{Procedencia y códigos originales}\hfill Ficha completa · 063/324}
\clearpage\phantomsection\label{nav:vi-ruta-64}
% VI-CATALOG-ROW rutas 64
\pdfbookmark[3]{r2c7\_h+1v+1}{route-64}
\RouteTitle{Ruta 064}{r2c7\_\allowbreak{}h+1v+1}
\noindent Celda 16\quad (2, 7)\qquad Orientación \(h=1,\ v=1\)\hfill\hyperref[sec:vi-indice-rutas]{Índice}\par\vspace{4pt}
\begingroup\fontsize{10.2}{12.5}\selectfont\setlength{\tabcolsep}{5pt}\renewcommand{\arraystretch}{1.1}\rowcolors{1}{routepale}{white}\noindent\begin{tabularx}{\linewidth}{@{}>{\raggedright\arraybackslash}p{.345\linewidth}>{\raggedright\arraybackslash}X@{}}
\RouteGroup{Identidad estructural}
\RouteRow{Clase y multisección}{quark\_\allowbreak{}down\_\allowbreak{}G3\quad /\quad quark\_\allowbreak{}down\_\allowbreak{}G3}
\RouteRow{Colección y sector}{quark de tipo abajo; quark}
\RouteRow{Clasificación física}{quark de tipo abajo}
\RouteRow{Generación, profundidad y color}{3; G3; G}
\RouteRow{Ruta primaria y órbita de inversión}{\(B_1\): no\qquad 2,\allowbreak{}7;\allowbreak{}8,\allowbreak{}3}
\RouteRow{Firma de clasificación}{34|\allowbreak{}0|\allowbreak{}0|\allowbreak{}2|\allowbreak{}-8|\allowbreak{}+++0-\kern0pt{}-+++0-\kern0pt{}-}
\RouteGroup{Retornos, memoria y transporte}
\RouteRow{Retornos con fase}{clase: 27; órbita: 27; celda: 27}
\RouteRow{Retorno cinemático}{en 108: 54; extensión: 216}
\RouteRow{Giros y cambios de hoja}{71; 107}
\RouteRow{Acarreo y diferimiento}{deuda total: 1; final: 0; bloques diferidos: 1}
\RouteRow{Tríadas finales; persistencia de Pareto}{17; sí}
\RouteRow{\(K\longrightarrow K+9\)}{repeticiones: 9; cambios de memoria: 9}
\RouteGroup{Firma, fase y diferencias}
\RouteRow{\(n_A,\ s_C,\ \kappa_0,\ \nu_{120},\ \nu_{270}\)}{34, 0, -2, 0, -8}
\RouteRow{\(\tau_{12}\); firma histórica \(\mathbb Z^5\)}{+++0-\kern0pt{}-+++0-\kern0pt{}-; 34|\allowbreak{}0|\allowbreak{}-2|\allowbreak{}0|\allowbreak{}-8}
\RouteRow{\(D_1,\ D_3,\ D_4,\ D_{27},\ D_{36}\)}{66, 64, 68, 60, 40}
\RouteRow{\(\Omega_{90/120},\ P_6\)}{56, 5}
\RouteGroup{Lectores y factores de acción}
\RouteRow{\(K_D\quad /\quad K_\Omega\)}{(100,\allowbreak{}66,\allowbreak{}2)\quad /\quad (56,\allowbreak{}54,\allowbreak{}5)}
\RouteRow{\(\kappa_D\)}{99.51859\allowbreak{}71232726\allowbreak{}53414287\allowbreak{}29}
\RouteRow{\(\kappa_\Omega\)}{55.60649\allowbreak{}89433721\allowbreak{}35446416\allowbreak{}86}
\RouteRow{\(R_{\mathrm{act},D}\)}{1.000000\allowbreak{}06899443\allowbreak{}08259364\allowbreak{}17}
\RouteRow{\(R_{\mathrm{act},\Omega}\)}{1.000000\allowbreak{}03855097\allowbreak{}23541416\allowbreak{}45}
\RouteGroup{Separación de prefijos}
\RouteRow{Área de ambigüedad; área \(\log_2\)}{36; 18.0}
\RouteRow{Primer bloque de separación}{órbita: \textemdash{}; celda: \textemdash{}; ruta: \textemdash{}}
\RouteRow{Ambigüedad final}{2}
\RouteGroup{Secuencias completas}
\RouteRow{\(w_6\)}{99 327 498 99 324 423 15 486 423 99 327 498 99 324 423 15 486 423}
\RouteRow{Tríadas estables por bloque}{0 1 2 3 4 5 6 7 8 9 10 11 12 13 14 14 16 17}
\end{tabularx}\endgroup\par\vspace{5pt}\noindent{\fontsize{8.4}{10}\selectfont\hyperref[trace-64]{Procedencia y códigos originales}\hfill Ficha completa · 064/324}
\clearpage\phantomsection\label{nav:vi-ruta-65}
% VI-CATALOG-ROW rutas 65
\pdfbookmark[2]{Celda 17}{cell-17}
\pdfbookmark[3]{r2c8\_h-1v-1}{route-65}
\RouteTitle{Ruta 065}{r2c8\_\allowbreak{}h-1v-1}
\noindent Celda 17\quad (2, 8)\qquad Orientación \(h=-1,\ v=-1\)\hfill\hyperref[sec:vi-indice-rutas]{Índice}\par\vspace{4pt}
\begingroup\fontsize{10.2}{12.5}\selectfont\setlength{\tabcolsep}{5pt}\renewcommand{\arraystretch}{1.1}\rowcolors{1}{routepale}{white}\noindent\begin{tabularx}{\linewidth}{@{}>{\raggedright\arraybackslash}p{.345\linewidth}>{\raggedright\arraybackslash}X@{}}
\RouteGroup{Identidad estructural}
\RouteRow{Clase y multisección}{quark\_\allowbreak{}up\_\allowbreak{}G3\quad /\quad quark\_\allowbreak{}up\_\allowbreak{}G3}
\RouteRow{Colección y sector}{quark de tipo arriba; quark}
\RouteRow{Clasificación física}{quark de tipo arriba}
\RouteRow{Generación, profundidad y color}{3; G3; R}
\RouteRow{Ruta primaria y órbita de inversión}{\(B_1\): sí\qquad 2,\allowbreak{}8;\allowbreak{}8,\allowbreak{}2}
\RouteRow{Firma de clasificación}{30|\allowbreak{}6|\allowbreak{}-8|\allowbreak{}0|\allowbreak{}-12|\allowbreak{}+++-\kern0pt{}-\kern0pt{}-+++-\kern0pt{}-\kern0pt{}-}
\RouteGroup{Retornos, memoria y transporte}
\RouteRow{Retornos con fase}{clase: 27; órbita: 27; celda: 27}
\RouteRow{Retorno cinemático}{en 108: 54; extensión: 216}
\RouteRow{Giros y cambios de hoja}{71; 107}
\RouteRow{Acarreo y diferimiento}{deuda total: 3; final: 0; bloques diferidos: 2}
\RouteRow{Tríadas finales; persistencia de Pareto}{17; sí}
\RouteRow{\(K\longrightarrow K+9\)}{repeticiones: 9; cambios de memoria: 9}
\RouteGroup{Firma, fase y diferencias}
\RouteRow{\(n_A,\ s_C,\ \kappa_0,\ \nu_{120},\ \nu_{270}\)}{30, -6, 0, -2, -12}
\RouteRow{\(\tau_{12}\); firma histórica \(\mathbb Z^5\)}{+++-\kern0pt{}-\kern0pt{}-+++-\kern0pt{}-\kern0pt{}-; 30|\allowbreak{}-6|\allowbreak{}0|\allowbreak{}-2|\allowbreak{}-12}
\RouteRow{\(D_1,\ D_3,\ D_4,\ D_{27},\ D_{36}\)}{54, 56, 68, 48, 32}
\RouteRow{\(\Omega_{90/120},\ P_6\)}{80, 6}
\RouteGroup{Lectores y factores de acción}
\RouteRow{\(K_D\quad /\quad K_\Omega\)}{(80,\allowbreak{}54,\allowbreak{}6)\quad /\quad (80,\allowbreak{}54,\allowbreak{}6)}
\RouteRow{\(\kappa_D\)}{79.60661\allowbreak{}01397948\allowbreak{}27586470\allowbreak{}91}
\RouteRow{\(\kappa_\Omega\)}{79.60661\allowbreak{}01397948\allowbreak{}27586470\allowbreak{}91}
\RouteRow{\(R_{\mathrm{act},D}\)}{1.000000\allowbreak{}05518981\allowbreak{}25318591\allowbreak{}40}
\RouteRow{\(R_{\mathrm{act},\Omega}\)}{1.000000\allowbreak{}05518981\allowbreak{}25318591\allowbreak{}40}
\RouteGroup{Separación de prefijos}
\RouteRow{Área de ambigüedad; área \(\log_2\)}{72; 36.0}
\RouteRow{Primer bloque de separación}{órbita: \textemdash{}; celda: \textemdash{}; ruta: \textemdash{}}
\RouteRow{Ambigüedad final}{4}
\RouteGroup{Secuencias completas}
\RouteRow{\(w_6\)}{423 15 486 423 9 24 243 657 24 423 15 486 423 9 24 243 657 24}
\RouteRow{Tríadas estables por bloque}{0 1 2 3 4 5 6 7 8 9 10 11 11 13 14 14 15 17}
\end{tabularx}\endgroup\par\vspace{5pt}\noindent{\fontsize{8.4}{10}\selectfont\hyperref[trace-65]{Procedencia y códigos originales}\hfill Ficha completa · 065/324}
\clearpage\phantomsection\label{nav:vi-ruta-66}
% VI-CATALOG-ROW rutas 66
\pdfbookmark[3]{r2c8\_h-1v+1}{route-66}
\RouteTitle{Ruta 066}{r2c8\_\allowbreak{}h-1v+1}
\noindent Celda 17\quad (2, 8)\qquad Orientación \(h=-1,\ v=1\)\hfill\hyperref[sec:vi-indice-rutas]{Índice}\par\vspace{4pt}
\begingroup\fontsize{10.2}{12.5}\selectfont\setlength{\tabcolsep}{5pt}\renewcommand{\arraystretch}{1.1}\rowcolors{1}{routepale}{white}\noindent\begin{tabularx}{\linewidth}{@{}>{\raggedright\arraybackslash}p{.345\linewidth}>{\raggedright\arraybackslash}X@{}}
\RouteGroup{Identidad estructural}
\RouteRow{Clase y multisección}{quark\_\allowbreak{}up\_\allowbreak{}G3\quad /\quad quark\_\allowbreak{}up\_\allowbreak{}G3}
\RouteRow{Colección y sector}{quark de tipo arriba; quark}
\RouteRow{Clasificación física}{quark de tipo arriba}
\RouteRow{Generación, profundidad y color}{3; G3; R}
\RouteRow{Ruta primaria y órbita de inversión}{\(B_1\): no\qquad 2,\allowbreak{}8;\allowbreak{}8,\allowbreak{}2}
\RouteRow{Firma de clasificación}{30|\allowbreak{}6|\allowbreak{}-8|\allowbreak{}0|\allowbreak{}-12|\allowbreak{}+++-\kern0pt{}-\kern0pt{}-+++-\kern0pt{}-\kern0pt{}-}
\RouteGroup{Retornos, memoria y transporte}
\RouteRow{Retornos con fase}{clase: 18; órbita: 18; celda: 27}
\RouteRow{Retorno cinemático}{en 108: 54; extensión: 216}
\RouteRow{Giros y cambios de hoja}{71; 107}
\RouteRow{Acarreo y diferimiento}{deuda total: 2; final: 1; bloques diferidos: 2}
\RouteRow{Tríadas finales; persistencia de Pareto}{16; sí}
\RouteRow{\(K\longrightarrow K+9\)}{repeticiones: 9; cambios de memoria: 9}
\RouteGroup{Firma, fase y diferencias}
\RouteRow{\(n_A,\ s_C,\ \kappa_0,\ \nu_{120},\ \nu_{270}\)}{50, 4, -2, 2, -8}
\RouteRow{\(\tau_{12}\); firma histórica \(\mathbb Z^5\)}{+++0-\kern0pt{}-+++0-\kern0pt{}-; 50|\allowbreak{}4|\allowbreak{}-2|\allowbreak{}2|\allowbreak{}-8}
\RouteRow{\(D_1,\ D_3,\ D_4,\ D_{27},\ D_{36}\)}{84, 28, 84, 36, 24}
\RouteRow{\(\Omega_{90/120},\ P_6\)}{56, 5}
\RouteGroup{Lectores y factores de acción}
\RouteRow{\(K_D\quad /\quad K_\Omega\)}{(60,\allowbreak{}84,\allowbreak{}28)\quad /\quad (56,\allowbreak{}54,\allowbreak{}5)}
\RouteRow{\(\kappa_D\)}{59.39013\allowbreak{}58840155\allowbreak{}40637741\allowbreak{}48}
\RouteRow{\(\kappa_\Omega\)}{55.60649\allowbreak{}89433721\allowbreak{}35446416\allowbreak{}86}
\RouteRow{\(R_{\mathrm{act},D}\)}{1.000000\allowbreak{}04117409\allowbreak{}89467415\allowbreak{}75}
\RouteRow{\(R_{\mathrm{act},\Omega}\)}{1.000000\allowbreak{}03855097\allowbreak{}23541416\allowbreak{}45}
\RouteGroup{Separación de prefijos}
\RouteRow{Área de ambigüedad; área \(\log_2\)}{72; 36.0}
\RouteRow{Primer bloque de separación}{órbita: \textemdash{}; celda: \textemdash{}; ruta: \textemdash{}}
\RouteRow{Ambigüedad final}{4}
\RouteGroup{Secuencias completas}
\RouteRow{\(w_6\)}{414 420 258 414 414 333 255 252 333 414 420 258 414 414 333 255 252 333}
\RouteRow{Tríadas estables por bloque}{0 1 2 3 4 5 6 7 8 9 10 11 12 13 14 14 16 16}
\end{tabularx}\endgroup\par\vspace{5pt}\noindent{\fontsize{8.4}{10}\selectfont\hyperref[trace-66]{Procedencia y códigos originales}\hfill Ficha completa · 066/324}
\clearpage\phantomsection\label{nav:vi-ruta-67}
% VI-CATALOG-ROW rutas 67
\pdfbookmark[3]{r2c8\_h+1v-1}{route-67}
\RouteTitle{Ruta 067}{r2c8\_\allowbreak{}h+1v-1}
\noindent Celda 17\quad (2, 8)\qquad Orientación \(h=1,\ v=-1\)\hfill\hyperref[sec:vi-indice-rutas]{Índice}\par\vspace{4pt}
\begingroup\fontsize{10.2}{12.5}\selectfont\setlength{\tabcolsep}{5pt}\renewcommand{\arraystretch}{1.1}\rowcolors{1}{routepale}{white}\noindent\begin{tabularx}{\linewidth}{@{}>{\raggedright\arraybackslash}p{.345\linewidth}>{\raggedright\arraybackslash}X@{}}
\RouteGroup{Identidad estructural}
\RouteRow{Clase y multisección}{quark\_\allowbreak{}up\_\allowbreak{}G3\quad /\quad quark\_\allowbreak{}up\_\allowbreak{}G3}
\RouteRow{Colección y sector}{quark de tipo arriba; quark}
\RouteRow{Clasificación física}{quark de tipo arriba}
\RouteRow{Generación, profundidad y color}{3; G3; R}
\RouteRow{Ruta primaria y órbita de inversión}{\(B_1\): no\qquad 2,\allowbreak{}8;\allowbreak{}8,\allowbreak{}2}
\RouteRow{Firma de clasificación}{30|\allowbreak{}6|\allowbreak{}-8|\allowbreak{}0|\allowbreak{}-12|\allowbreak{}+++-\kern0pt{}-\kern0pt{}-+++-\kern0pt{}-\kern0pt{}-}
\RouteGroup{Retornos, memoria y transporte}
\RouteRow{Retornos con fase}{clase: 9; órbita: 9; celda: 27}
\RouteRow{Retorno cinemático}{en 108: 54; extensión: 216}
\RouteRow{Giros y cambios de hoja}{71; 107}
\RouteRow{Acarreo y diferimiento}{deuda total: 3; final: 0; bloques diferidos: 2}
\RouteRow{Tríadas finales; persistencia de Pareto}{17; sí}
\RouteRow{\(K\longrightarrow K+9\)}{repeticiones: 9; cambios de memoria: 9}
\RouteGroup{Firma, fase y diferencias}
\RouteRow{\(n_A,\ s_C,\ \kappa_0,\ \nu_{120},\ \nu_{270}\)}{50, -4, 2, -2, -8}
\RouteRow{\(\tau_{12}\); firma histórica \(\mathbb Z^5\)}{0++-\kern0pt{}-\kern0pt{}-0++-\kern0pt{}-\kern0pt{}-; 50|\allowbreak{}-4|\allowbreak{}2|\allowbreak{}-2|\allowbreak{}-8}
\RouteRow{\(D_1,\ D_3,\ D_4,\ D_{27},\ D_{36}\)}{84, 28, 84, 36, 24}
\RouteRow{\(\Omega_{90/120},\ P_6\)}{56, 5}
\RouteGroup{Lectores y factores de acción}
\RouteRow{\(K_D\quad /\quad K_\Omega\)}{(60,\allowbreak{}84,\allowbreak{}28)\quad /\quad (56,\allowbreak{}54,\allowbreak{}5)}
\RouteRow{\(\kappa_D\)}{59.39013\allowbreak{}58840155\allowbreak{}40637741\allowbreak{}48}
\RouteRow{\(\kappa_\Omega\)}{55.60649\allowbreak{}89433721\allowbreak{}35446416\allowbreak{}86}
\RouteRow{\(R_{\mathrm{act},D}\)}{1.000000\allowbreak{}04117409\allowbreak{}89467415\allowbreak{}75}
\RouteRow{\(R_{\mathrm{act},\Omega}\)}{1.000000\allowbreak{}03855097\allowbreak{}23541416\allowbreak{}45}
\RouteGroup{Separación de prefijos}
\RouteRow{Área de ambigüedad; área \(\log_2\)}{72; 36.0}
\RouteRow{Primer bloque de separación}{órbita: \textemdash{}; celda: \textemdash{}; ruta: \textemdash{}}
\RouteRow{Ambigüedad final}{4}
\RouteGroup{Secuencias completas}
\RouteRow{\(w_6\)}{255 252 333 255 258 258 414 420 258 255 252 333 255 258 258 414 420 258}
\RouteRow{Tríadas estables por bloque}{0 1 2 3 4 5 6 7 8 9 10 11 11 12 14 15 15 17}
\end{tabularx}\endgroup\par\vspace{5pt}\noindent{\fontsize{8.4}{10}\selectfont\hyperref[trace-67]{Procedencia y códigos originales}\hfill Ficha completa · 067/324}
\clearpage\phantomsection\label{nav:vi-ruta-68}
% VI-CATALOG-ROW rutas 68
\pdfbookmark[3]{r2c8\_h+1v+1}{route-68}
\RouteTitle{Ruta 068}{r2c8\_\allowbreak{}h+1v+1}
\noindent Celda 17\quad (2, 8)\qquad Orientación \(h=1,\ v=1\)\hfill\hyperref[sec:vi-indice-rutas]{Índice}\par\vspace{4pt}
\begingroup\fontsize{10.2}{12.5}\selectfont\setlength{\tabcolsep}{5pt}\renewcommand{\arraystretch}{1.1}\rowcolors{1}{routepale}{white}\noindent\begin{tabularx}{\linewidth}{@{}>{\raggedright\arraybackslash}p{.345\linewidth}>{\raggedright\arraybackslash}X@{}}
\RouteGroup{Identidad estructural}
\RouteRow{Clase y multisección}{quark\_\allowbreak{}up\_\allowbreak{}G3\quad /\quad quark\_\allowbreak{}up\_\allowbreak{}G3}
\RouteRow{Colección y sector}{quark de tipo arriba; quark}
\RouteRow{Clasificación física}{quark de tipo arriba}
\RouteRow{Generación, profundidad y color}{3; G3; R}
\RouteRow{Ruta primaria y órbita de inversión}{\(B_1\): no\qquad 2,\allowbreak{}8;\allowbreak{}8,\allowbreak{}2}
\RouteRow{Firma de clasificación}{30|\allowbreak{}6|\allowbreak{}-8|\allowbreak{}0|\allowbreak{}-12|\allowbreak{}+++-\kern0pt{}-\kern0pt{}-+++-\kern0pt{}-\kern0pt{}-}
\RouteGroup{Retornos, memoria y transporte}
\RouteRow{Retornos con fase}{clase: 27; órbita: 27; celda: 27}
\RouteRow{Retorno cinemático}{en 108: 54; extensión: 216}
\RouteRow{Giros y cambios de hoja}{71; 107}
\RouteRow{Acarreo y diferimiento}{deuda total: 2; final: 1; bloques diferidos: 2}
\RouteRow{Tríadas finales; persistencia de Pareto}{16; sí}
\RouteRow{\(K\longrightarrow K+9\)}{repeticiones: 9; cambios de memoria: 9}
\RouteGroup{Firma, fase y diferencias}
\RouteRow{\(n_A,\ s_C,\ \kappa_0,\ \nu_{120},\ \nu_{270}\)}{30, 6, 0, 2, -12}
\RouteRow{\(\tau_{12}\); firma histórica \(\mathbb Z^5\)}{+++-\kern0pt{}-\kern0pt{}-+++-\kern0pt{}-\kern0pt{}-; 30|\allowbreak{}6|\allowbreak{}0|\allowbreak{}2|\allowbreak{}-12}
\RouteRow{\(D_1,\ D_3,\ D_4,\ D_{27},\ D_{36}\)}{54, 56, 68, 48, 32}
\RouteRow{\(\Omega_{90/120},\ P_6\)}{80, 6}
\RouteGroup{Lectores y factores de acción}
\RouteRow{\(K_D\quad /\quad K_\Omega\)}{(80,\allowbreak{}54,\allowbreak{}6)\quad /\quad (80,\allowbreak{}54,\allowbreak{}6)}
\RouteRow{\(\kappa_D\)}{79.60661\allowbreak{}01397948\allowbreak{}27586470\allowbreak{}91}
\RouteRow{\(\kappa_\Omega\)}{79.60661\allowbreak{}01397948\allowbreak{}27586470\allowbreak{}91}
\RouteRow{\(R_{\mathrm{act},D}\)}{1.000000\allowbreak{}05518981\allowbreak{}25318591\allowbreak{}40}
\RouteRow{\(R_{\mathrm{act},\Omega}\)}{1.000000\allowbreak{}05518981\allowbreak{}25318591\allowbreak{}40}
\RouteGroup{Separación de prefijos}
\RouteRow{Área de ambigüedad; área \(\log_2\)}{72; 36.0}
\RouteRow{Primer bloque de separación}{órbita: \textemdash{}; celda: \textemdash{}; ruta: \textemdash{}}
\RouteRow{Ambigüedad final}{4}
\RouteGroup{Secuencias completas}
\RouteRow{\(w_6\)}{243 657 24 243 663 486 423 15 486 243 657 24 243 663 486 423 15 486}
\RouteRow{Tríadas estables por bloque}{0 1 2 3 4 5 6 7 8 9 10 11 12 13 14 14 16 16}
\end{tabularx}\endgroup\par\vspace{5pt}\noindent{\fontsize{8.4}{10}\selectfont\hyperref[trace-68]{Procedencia y códigos originales}\hfill Ficha completa · 068/324}
\clearpage\phantomsection\label{nav:vi-ruta-69}
% VI-CATALOG-ROW rutas 69
\pdfbookmark[2]{Celda 18}{cell-18}
\pdfbookmark[3]{r2c9\_h-1v-1}{route-69}
\RouteTitle{Ruta 069}{r2c9\_\allowbreak{}h-1v-1}
\noindent Celda 18\quad (2, 9)\qquad Orientación \(h=-1,\ v=-1\)\hfill\hyperref[sec:vi-indice-rutas]{Índice}\par\vspace{4pt}
\begingroup\fontsize{10.2}{12.5}\selectfont\setlength{\tabcolsep}{5pt}\renewcommand{\arraystretch}{1.1}\rowcolors{1}{routepale}{white}\noindent\begin{tabularx}{\linewidth}{@{}>{\raggedright\arraybackslash}p{.345\linewidth}>{\raggedright\arraybackslash}X@{}}
\RouteGroup{Identidad estructural}
\RouteRow{Clase y multisección}{quark\_\allowbreak{}down\_\allowbreak{}G4\quad /\quad quark\_\allowbreak{}down\_\allowbreak{}G4}
\RouteRow{Colección y sector}{quark de tipo abajo; quark}
\RouteRow{Clasificación física}{quark de tipo abajo}
\RouteRow{Generación, profundidad y color}{4; G4; B}
\RouteRow{Ruta primaria y órbita de inversión}{\(B_1\): no\qquad 2,\allowbreak{}9;\allowbreak{}8,\allowbreak{}1}
\RouteRow{Firma de clasificación}{28|\allowbreak{}6|\allowbreak{}-2|\allowbreak{}0|\allowbreak{}-12|\allowbreak{}+++-\kern0pt{}-\kern0pt{}-+++-\kern0pt{}-\kern0pt{}-}
\RouteGroup{Retornos, memoria y transporte}
\RouteRow{Retornos con fase}{clase: 27; órbita: 27; celda: 27}
\RouteRow{Retorno cinemático}{en 108: 54; extensión: 216}
\RouteRow{Giros y cambios de hoja}{71; 107}
\RouteRow{Acarreo y diferimiento}{deuda total: 1; final: 0; bloques diferidos: 1}
\RouteRow{Tríadas finales; persistencia de Pareto}{17; sí}
\RouteRow{\(K\longrightarrow K+9\)}{repeticiones: 9; cambios de memoria: 9}
\RouteGroup{Firma, fase y diferencias}
\RouteRow{\(n_A,\ s_C,\ \kappa_0,\ \nu_{120},\ \nu_{270}\)}{28, -6, 0, 0, -12}
\RouteRow{\(\tau_{12}\); firma histórica \(\mathbb Z^5\)}{+++-\kern0pt{}-\kern0pt{}-+++-\kern0pt{}-\kern0pt{}-; 28|\allowbreak{}-6|\allowbreak{}0|\allowbreak{}0|\allowbreak{}-12}
\RouteRow{\(D_1,\ D_3,\ D_4,\ D_{27},\ D_{36}\)}{66, 60, 68, 48, 32}
\RouteRow{\(\Omega_{90/120},\ P_6\)}{80, 6}
\RouteGroup{Lectores y factores de acción}
\RouteRow{\(K_D\quad /\quad K_\Omega\)}{(80,\allowbreak{}66,\allowbreak{}4)\quad /\quad (80,\allowbreak{}54,\allowbreak{}6)}
\RouteRow{\(\kappa_D\)}{79.51881\allowbreak{}95161180\allowbreak{}37694395\allowbreak{}39}
\RouteRow{\(\kappa_\Omega\)}{79.60661\allowbreak{}01397948\allowbreak{}27586470\allowbreak{}91}
\RouteRow{\(R_{\mathrm{act},D}\)}{1.000000\allowbreak{}05512894\allowbreak{}88901612\allowbreak{}64}
\RouteRow{\(R_{\mathrm{act},\Omega}\)}{1.000000\allowbreak{}05518981\allowbreak{}25318591\allowbreak{}40}
\RouteGroup{Separación de prefijos}
\RouteRow{Área de ambigüedad; área \(\log_2\)}{36; 18.0}
\RouteRow{Primer bloque de separación}{órbita: \textemdash{}; celda: \textemdash{}; ruta: \textemdash{}}
\RouteRow{Ambigüedad final}{2}
\RouteGroup{Secuencias completas}
\RouteRow{\(w_6\)}{570 264 90 570 267 243 657 24 243 570 264 90 570 267 243 657 24 243}
\RouteRow{Tríadas estables por bloque}{0 1 1 3 4 5 6 7 8 9 10 11 12 13 14 15 16 17}
\end{tabularx}\endgroup\par\vspace{5pt}\noindent{\fontsize{8.4}{10}\selectfont\hyperref[trace-69]{Procedencia y códigos originales}\hfill Ficha completa · 069/324}
\clearpage\phantomsection\label{nav:vi-ruta-70}
% VI-CATALOG-ROW rutas 70
\pdfbookmark[3]{r2c9\_h-1v+1}{route-70}
\RouteTitle{Ruta 070}{r2c9\_\allowbreak{}h-1v+1}
\noindent Celda 18\quad (2, 9)\qquad Orientación \(h=-1,\ v=1\)\hfill\hyperref[sec:vi-indice-rutas]{Índice}\par\vspace{4pt}
\begingroup\fontsize{10.2}{12.5}\selectfont\setlength{\tabcolsep}{5pt}\renewcommand{\arraystretch}{1.1}\rowcolors{1}{routepale}{white}\noindent\begin{tabularx}{\linewidth}{@{}>{\raggedright\arraybackslash}p{.345\linewidth}>{\raggedright\arraybackslash}X@{}}
\RouteGroup{Identidad estructural}
\RouteRow{Clase y multisección}{quark\_\allowbreak{}down\_\allowbreak{}G4\quad /\quad quark\_\allowbreak{}down\_\allowbreak{}G4}
\RouteRow{Colección y sector}{quark de tipo abajo; quark}
\RouteRow{Clasificación física}{quark de tipo abajo}
\RouteRow{Generación, profundidad y color}{4; G4; B}
\RouteRow{Ruta primaria y órbita de inversión}{\(B_1\): sí\qquad 2,\allowbreak{}9;\allowbreak{}8,\allowbreak{}1}
\RouteRow{Firma de clasificación}{28|\allowbreak{}6|\allowbreak{}-2|\allowbreak{}0|\allowbreak{}-12|\allowbreak{}+++-\kern0pt{}-\kern0pt{}-+++-\kern0pt{}-\kern0pt{}-}
\RouteGroup{Retornos, memoria y transporte}
\RouteRow{Retornos con fase}{clase: 27; órbita: 27; celda: 27}
\RouteRow{Retorno cinemático}{en 108: 54; extensión: 216}
\RouteRow{Giros y cambios de hoja}{71; 107}
\RouteRow{Acarreo y diferimiento}{deuda total: 1; final: 0; bloques diferidos: 1}
\RouteRow{Tríadas finales; persistencia de Pareto}{17; sí}
\RouteRow{\(K\longrightarrow K+9\)}{repeticiones: 9; cambios de memoria: 9}
\RouteGroup{Firma, fase y diferencias}
\RouteRow{\(n_A,\ s_C,\ \kappa_0,\ \nu_{120},\ \nu_{270}\)}{14, -4, 0, 0, -12}
\RouteRow{\(\tau_{12}\); firma histórica \(\mathbb Z^5\)}{+++-\kern0pt{}-\kern0pt{}-+++-\kern0pt{}-\kern0pt{}-; 14|\allowbreak{}-4|\allowbreak{}0|\allowbreak{}0|\allowbreak{}-12}
\RouteRow{\(D_1,\ D_3,\ D_4,\ D_{27},\ D_{36}\)}{78, 24, 78, 24, 16}
\RouteRow{\(\Omega_{90/120},\ P_6\)}{80, 6}
\RouteGroup{Lectores y factores de acción}
\RouteRow{\(K_D\quad /\quad K_\Omega\)}{(40,\allowbreak{}78,\allowbreak{}27)\quad /\quad (80,\allowbreak{}54,\allowbreak{}6)}
\RouteRow{\(\kappa_D\)}{39.43380\allowbreak{}88030085\allowbreak{}51303671\allowbreak{}14}
\RouteRow{\(\kappa_\Omega\)}{79.60661\allowbreak{}01397948\allowbreak{}27586470\allowbreak{}91}
\RouteRow{\(R_{\mathrm{act},D}\)}{1.000000\allowbreak{}02733874\allowbreak{}08548943\allowbreak{}58}
\RouteRow{\(R_{\mathrm{act},\Omega}\)}{1.000000\allowbreak{}05518981\allowbreak{}25318591\allowbreak{}40}
\RouteGroup{Separación de prefijos}
\RouteRow{Área de ambigüedad; área \(\log_2\)}{36; 18.0}
\RouteRow{Primer bloque de separación}{órbita: \textemdash{}; celda: \textemdash{}; ruta: \textemdash{}}
\RouteRow{Ambigüedad final}{2}
\RouteGroup{Secuencias completas}
\RouteRow{\(w_6\)}{585 507 504 585 510 666 672 510 666 585 507 504 585 510 666 672 510 666}
\RouteRow{Tríadas estables por bloque}{0 1 2 3 4 5 6 7 8 9 10 11 12 13 14 14 16 17}
\end{tabularx}\endgroup\par\vspace{5pt}\noindent{\fontsize{8.4}{10}\selectfont\hyperref[trace-70]{Procedencia y códigos originales}\hfill Ficha completa · 070/324}
\clearpage\phantomsection\label{nav:vi-ruta-71}
% VI-CATALOG-ROW rutas 71
\pdfbookmark[3]{r2c9\_h+1v-1}{route-71}
\RouteTitle{Ruta 071}{r2c9\_\allowbreak{}h+1v-1}
\noindent Celda 18\quad (2, 9)\qquad Orientación \(h=1,\ v=-1\)\hfill\hyperref[sec:vi-indice-rutas]{Índice}\par\vspace{4pt}
\begingroup\fontsize{10.2}{12.5}\selectfont\setlength{\tabcolsep}{5pt}\renewcommand{\arraystretch}{1.1}\rowcolors{1}{routepale}{white}\noindent\begin{tabularx}{\linewidth}{@{}>{\raggedright\arraybackslash}p{.345\linewidth}>{\raggedright\arraybackslash}X@{}}
\RouteGroup{Identidad estructural}
\RouteRow{Clase y multisección}{quark\_\allowbreak{}down\_\allowbreak{}G4\quad /\quad quark\_\allowbreak{}down\_\allowbreak{}G4}
\RouteRow{Colección y sector}{quark de tipo abajo; quark}
\RouteRow{Clasificación física}{quark de tipo abajo}
\RouteRow{Generación, profundidad y color}{4; G4; B}
\RouteRow{Ruta primaria y órbita de inversión}{\(B_1\): no\qquad 2,\allowbreak{}9;\allowbreak{}8,\allowbreak{}1}
\RouteRow{Firma de clasificación}{28|\allowbreak{}6|\allowbreak{}-2|\allowbreak{}0|\allowbreak{}-12|\allowbreak{}+++-\kern0pt{}-\kern0pt{}-+++-\kern0pt{}-\kern0pt{}-}
\RouteGroup{Retornos, memoria y transporte}
\RouteRow{Retornos con fase}{clase: 27; órbita: 27; celda: 27}
\RouteRow{Retorno cinemático}{en 108: 54; extensión: 216}
\RouteRow{Giros y cambios de hoja}{71; 107}
\RouteRow{Acarreo y diferimiento}{deuda total: 1; final: 0; bloques diferidos: 1}
\RouteRow{Tríadas finales; persistencia de Pareto}{17; sí}
\RouteRow{\(K\longrightarrow K+9\)}{repeticiones: 9; cambios de memoria: 9}
\RouteGroup{Firma, fase y diferencias}
\RouteRow{\(n_A,\ s_C,\ \kappa_0,\ \nu_{120},\ \nu_{270}\)}{14, 4, 0, 0, -12}
\RouteRow{\(\tau_{12}\); firma histórica \(\mathbb Z^5\)}{+++-\kern0pt{}-\kern0pt{}-+++-\kern0pt{}-\kern0pt{}-; 14|\allowbreak{}4|\allowbreak{}0|\allowbreak{}0|\allowbreak{}-12}
\RouteRow{\(D_1,\ D_3,\ D_4,\ D_{27},\ D_{36}\)}{78, 24, 78, 24, 16}
\RouteRow{\(\Omega_{90/120},\ P_6\)}{80, 6}
\RouteGroup{Lectores y factores de acción}
\RouteRow{\(K_D\quad /\quad K_\Omega\)}{(40,\allowbreak{}78,\allowbreak{}27)\quad /\quad (80,\allowbreak{}54,\allowbreak{}6)}
\RouteRow{\(\kappa_D\)}{39.43380\allowbreak{}88030085\allowbreak{}51303671\allowbreak{}14}
\RouteRow{\(\kappa_\Omega\)}{79.60661\allowbreak{}01397948\allowbreak{}27586470\allowbreak{}91}
\RouteRow{\(R_{\mathrm{act},D}\)}{1.000000\allowbreak{}02733874\allowbreak{}08548943\allowbreak{}58}
\RouteRow{\(R_{\mathrm{act},\Omega}\)}{1.000000\allowbreak{}05518981\allowbreak{}25318591\allowbreak{}40}
\RouteGroup{Separación de prefijos}
\RouteRow{Área de ambigüedad; área \(\log_2\)}{36; 18.0}
\RouteRow{Primer bloque de separación}{órbita: \textemdash{}; celda: \textemdash{}; ruta: \textemdash{}}
\RouteRow{Ambigüedad final}{2}
\RouteGroup{Secuencias completas}
\RouteRow{\(w_6\)}{672 510 666 672 507 504 585 507 504 672 510 666 672 507 504 585 507 504}
\RouteRow{Tríadas estables por bloque}{0 1 2 3 4 5 6 7 8 9 10 11 12 13 14 14 16 17}
\end{tabularx}\endgroup\par\vspace{5pt}\noindent{\fontsize{8.4}{10}\selectfont\hyperref[trace-71]{Procedencia y códigos originales}\hfill Ficha completa · 071/324}
\clearpage\phantomsection\label{nav:vi-ruta-72}
% VI-CATALOG-ROW rutas 72
\pdfbookmark[3]{r2c9\_h+1v+1}{route-72}
\RouteTitle{Ruta 072}{r2c9\_\allowbreak{}h+1v+1}
\noindent Celda 18\quad (2, 9)\qquad Orientación \(h=1,\ v=1\)\hfill\hyperref[sec:vi-indice-rutas]{Índice}\par\vspace{4pt}
\begingroup\fontsize{10.2}{12.5}\selectfont\setlength{\tabcolsep}{5pt}\renewcommand{\arraystretch}{1.1}\rowcolors{1}{routepale}{white}\noindent\begin{tabularx}{\linewidth}{@{}>{\raggedright\arraybackslash}p{.345\linewidth}>{\raggedright\arraybackslash}X@{}}
\RouteGroup{Identidad estructural}
\RouteRow{Clase y multisección}{quark\_\allowbreak{}down\_\allowbreak{}G4\quad /\quad quark\_\allowbreak{}down\_\allowbreak{}G4}
\RouteRow{Colección y sector}{quark de tipo abajo; quark}
\RouteRow{Clasificación física}{quark de tipo abajo}
\RouteRow{Generación, profundidad y color}{4; G4; B}
\RouteRow{Ruta primaria y órbita de inversión}{\(B_1\): no\qquad 2,\allowbreak{}9;\allowbreak{}8,\allowbreak{}1}
\RouteRow{Firma de clasificación}{28|\allowbreak{}6|\allowbreak{}-2|\allowbreak{}0|\allowbreak{}-12|\allowbreak{}+++-\kern0pt{}-\kern0pt{}-+++-\kern0pt{}-\kern0pt{}-}
\RouteGroup{Retornos, memoria y transporte}
\RouteRow{Retornos con fase}{clase: 27; órbita: 27; celda: 27}
\RouteRow{Retorno cinemático}{en 108: 54; extensión: 216}
\RouteRow{Giros y cambios de hoja}{71; 107}
\RouteRow{Acarreo y diferimiento}{deuda total: 0; final: 0; bloques diferidos: 0}
\RouteRow{Tríadas finales; persistencia de Pareto}{17; sí}
\RouteRow{\(K\longrightarrow K+9\)}{repeticiones: 9; cambios de memoria: 9}
\RouteGroup{Firma, fase y diferencias}
\RouteRow{\(n_A,\ s_C,\ \kappa_0,\ \nu_{120},\ \nu_{270}\)}{28, 6, 0, 0, -12}
\RouteRow{\(\tau_{12}\); firma histórica \(\mathbb Z^5\)}{+++-\kern0pt{}-\kern0pt{}-+++-\kern0pt{}-\kern0pt{}-; 28|\allowbreak{}6|\allowbreak{}0|\allowbreak{}0|\allowbreak{}-12}
\RouteRow{\(D_1,\ D_3,\ D_4,\ D_{27},\ D_{36}\)}{66, 60, 68, 48, 32}
\RouteRow{\(\Omega_{90/120},\ P_6\)}{80, 6}
\RouteGroup{Lectores y factores de acción}
\RouteRow{\(K_D\quad /\quad K_\Omega\)}{(80,\allowbreak{}66,\allowbreak{}4)\quad /\quad (80,\allowbreak{}54,\allowbreak{}6)}
\RouteRow{\(\kappa_D\)}{79.51881\allowbreak{}95161180\allowbreak{}37694395\allowbreak{}39}
\RouteRow{\(\kappa_\Omega\)}{79.60661\allowbreak{}01397948\allowbreak{}27586470\allowbreak{}91}
\RouteRow{\(R_{\mathrm{act},D}\)}{1.000000\allowbreak{}05512894\allowbreak{}88901612\allowbreak{}64}
\RouteRow{\(R_{\mathrm{act},\Omega}\)}{1.000000\allowbreak{}05518981\allowbreak{}25318591\allowbreak{}40}
\RouteGroup{Separación de prefijos}
\RouteRow{Área de ambigüedad; área \(\log_2\)}{40; 20.33985\allowbreak{}00028846\allowbreak{}24}
\RouteRow{Primer bloque de separación}{órbita: \textemdash{}; celda: \textemdash{}; ruta: \textemdash{}}
\RouteRow{Ambigüedad final}{2}
\RouteGroup{Secuencias completas}
\RouteRow{\(w_6\)}{657 24 243 657 21 90 570 264 90 657 24 243 657 21 90 570 264 90}
\RouteRow{Tríadas estables por bloque}{0 1 2 3 4 5 6 7 8 9 10 11 12 13 14 15 16 17}
\end{tabularx}\endgroup\par\vspace{5pt}\noindent{\fontsize{8.4}{10}\selectfont\hyperref[trace-72]{Procedencia y códigos originales}\hfill Ficha completa · 072/324}
\clearpage\phantomsection\label{nav:vi-ruta-73}
% VI-CATALOG-ROW rutas 73
\pdfbookmark[2]{Celda 19}{cell-19}
\pdfbookmark[3]{r3c1\_h-1v-1}{route-73}
\RouteTitle{Ruta 073}{r3c1\_\allowbreak{}h-1v-1}
\noindent Celda 19\quad (3, 1)\qquad Orientación \(h=-1,\ v=-1\)\hfill\hyperref[sec:vi-indice-rutas]{Índice}\par\vspace{4pt}
\begingroup\fontsize{10.2}{12.5}\selectfont\setlength{\tabcolsep}{5pt}\renewcommand{\arraystretch}{1.1}\rowcolors{1}{routepale}{white}\noindent\begin{tabularx}{\linewidth}{@{}>{\raggedright\arraybackslash}p{.345\linewidth}>{\raggedright\arraybackslash}X@{}}
\RouteGroup{Identidad estructural}
\RouteRow{Clase y multisección}{charged\_\allowbreak{}lepton\_\allowbreak{}G4\quad /\quad charged\_\allowbreak{}lepton\_\allowbreak{}G4}
\RouteRow{Colección y sector}{leptón cargado; leptón}
\RouteRow{Clasificación física}{leptón cargado}
\RouteRow{Generación, profundidad y color}{4; G4; \textemdash{}}
\RouteRow{Ruta primaria y órbita de inversión}{\(B_1\): no\qquad 3,\allowbreak{}1;\allowbreak{}7,\allowbreak{}9}
\RouteRow{Firma de clasificación}{28|\allowbreak{}6|\allowbreak{}-10|\allowbreak{}0|\allowbreak{}-12|\allowbreak{}+++-\kern0pt{}-\kern0pt{}-+++-\kern0pt{}-\kern0pt{}-}
\RouteGroup{Retornos, memoria y transporte}
\RouteRow{Retornos con fase}{clase: 9; órbita: 27; celda: 27}
\RouteRow{Retorno cinemático}{en 108: 54; extensión: 216}
\RouteRow{Giros y cambios de hoja}{71; 107}
\RouteRow{Acarreo y diferimiento}{deuda total: 2; final: 1; bloques diferidos: 1}
\RouteRow{Tríadas finales; persistencia de Pareto}{16; no}
\RouteRow{\(K\longrightarrow K+9\)}{repeticiones: 9; cambios de memoria: 9}
\RouteGroup{Firma, fase y diferencias}
\RouteRow{\(n_A,\ s_C,\ \kappa_0,\ \nu_{120},\ \nu_{270}\)}{28, -6, 0, 0, -12}
\RouteRow{\(\tau_{12}\); firma histórica \(\mathbb Z^5\)}{+++-\kern0pt{}-\kern0pt{}-+++-\kern0pt{}-\kern0pt{}-; 28|\allowbreak{}-6|\allowbreak{}0|\allowbreak{}0|\allowbreak{}-12}
\RouteRow{\(D_1,\ D_3,\ D_4,\ D_{27},\ D_{36}\)}{66, 60, 66, 48, 32}
\RouteRow{\(\Omega_{90/120},\ P_6\)}{80, 6}
\RouteGroup{Lectores y factores de acción}
\RouteRow{\(K_D\quad /\quad K_\Omega\)}{(80,\allowbreak{}66,\allowbreak{}3)\quad /\quad (80,\allowbreak{}54,\allowbreak{}6)}
\RouteRow{\(\kappa_D\)}{79.51870\allowbreak{}83196953\allowbreak{}45554341\allowbreak{}34}
\RouteRow{\(\kappa_\Omega\)}{79.60661\allowbreak{}01397948\allowbreak{}27586470\allowbreak{}91}
\RouteRow{\(R_{\mathrm{act},D}\)}{1.000000\allowbreak{}05512887\allowbreak{}17997050\allowbreak{}72}
\RouteRow{\(R_{\mathrm{act},\Omega}\)}{1.000000\allowbreak{}05518981\allowbreak{}25318591\allowbreak{}40}
\RouteGroup{Separación de prefijos}
\RouteRow{Área de ambigüedad; área \(\log_2\)}{54; 28.52932\allowbreak{}50129808}
\RouteRow{Primer bloque de separación}{órbita: \textemdash{}; celda: \textemdash{}; ruta: \textemdash{}}
\RouteRow{Ambigüedad final}{3}
\RouteGroup{Secuencias completas}
\RouteRow{\(w_6\)}{264 90 570 264 90 24 243 657 24 264 90 570 264 90 24 243 657 24}
\RouteRow{Tríadas estables por bloque}{0 1 2 3 4 5 6 7 8 9 10 11 12 13 14 15 15 16}
\end{tabularx}\endgroup\par\vspace{5pt}\noindent{\fontsize{8.4}{10}\selectfont\hyperref[trace-73]{Procedencia y códigos originales}\hfill Ficha completa · 073/324}
\clearpage\phantomsection\label{nav:vi-ruta-74}
% VI-CATALOG-ROW rutas 74
\pdfbookmark[3]{r3c1\_h-1v+1}{route-74}
\RouteTitle{Ruta 074}{r3c1\_\allowbreak{}h-1v+1}
\noindent Celda 19\quad (3, 1)\qquad Orientación \(h=-1,\ v=1\)\hfill\hyperref[sec:vi-indice-rutas]{Índice}\par\vspace{4pt}
\begingroup\fontsize{10.2}{12.5}\selectfont\setlength{\tabcolsep}{5pt}\renewcommand{\arraystretch}{1.1}\rowcolors{1}{routepale}{white}\noindent\begin{tabularx}{\linewidth}{@{}>{\raggedright\arraybackslash}p{.345\linewidth}>{\raggedright\arraybackslash}X@{}}
\RouteGroup{Identidad estructural}
\RouteRow{Clase y multisección}{charged\_\allowbreak{}lepton\_\allowbreak{}G4\quad /\quad charged\_\allowbreak{}lepton\_\allowbreak{}G4}
\RouteRow{Colección y sector}{leptón cargado; leptón}
\RouteRow{Clasificación física}{leptón cargado}
\RouteRow{Generación, profundidad y color}{4; G4; \textemdash{}}
\RouteRow{Ruta primaria y órbita de inversión}{\(B_1\): sí\qquad 3,\allowbreak{}1;\allowbreak{}7,\allowbreak{}9}
\RouteRow{Firma de clasificación}{28|\allowbreak{}6|\allowbreak{}-10|\allowbreak{}0|\allowbreak{}-12|\allowbreak{}+++-\kern0pt{}-\kern0pt{}-+++-\kern0pt{}-\kern0pt{}-}
\RouteGroup{Retornos, memoria y transporte}
\RouteRow{Retornos con fase}{clase: 27; órbita: 27; celda: 27}
\RouteRow{Retorno cinemático}{en 108: 54; extensión: 216}
\RouteRow{Giros y cambios de hoja}{71; 107}
\RouteRow{Acarreo y diferimiento}{deuda total: 2; final: 0; bloques diferidos: 1}
\RouteRow{Tríadas finales; persistencia de Pareto}{17; no}
\RouteRow{\(K\longrightarrow K+9\)}{repeticiones: 9; cambios de memoria: 9}
\RouteGroup{Firma, fase y diferencias}
\RouteRow{\(n_A,\ s_C,\ \kappa_0,\ \nu_{120},\ \nu_{270}\)}{8, -2, 0, 0, -12}
\RouteRow{\(\tau_{12}\); firma histórica \(\mathbb Z^5\)}{+++-\kern0pt{}-\kern0pt{}-+++-\kern0pt{}-\kern0pt{}-; 8|\allowbreak{}-2|\allowbreak{}0|\allowbreak{}0|\allowbreak{}-12}
\RouteRow{\(D_1,\ D_3,\ D_4,\ D_{27},\ D_{36}\)}{84, 24, 84, 24, 16}
\RouteRow{\(\Omega_{90/120},\ P_6\)}{80, 6}
\RouteGroup{Lectores y factores de acción}
\RouteRow{\(K_D\quad /\quad K_\Omega\)}{(40,\allowbreak{}84,\allowbreak{}30)\quad /\quad (80,\allowbreak{}54,\allowbreak{}6)}
\RouteRow{\(\kappa_D\)}{39.39035\allowbreak{}82768609\allowbreak{}24917849\allowbreak{}59}
\RouteRow{\(\kappa_\Omega\)}{79.60661\allowbreak{}01397948\allowbreak{}27586470\allowbreak{}91}
\RouteRow{\(R_{\mathrm{act},D}\)}{1.000000\allowbreak{}02730861\allowbreak{}73967087\allowbreak{}05}
\RouteRow{\(R_{\mathrm{act},\Omega}\)}{1.000000\allowbreak{}05518981\allowbreak{}25318591\allowbreak{}40}
\RouteGroup{Separación de prefijos}
\RouteRow{Área de ambigüedad; área \(\log_2\)}{54; 28.52932\allowbreak{}50129808}
\RouteRow{Primer bloque de separación}{órbita: \textemdash{}; celda: \textemdash{}; ruta: \textemdash{}}
\RouteRow{Ambigüedad final}{3}
\RouteGroup{Secuencias completas}
\RouteRow{\(w_6\)}{258 414 420 258 414 255 252 333 255 258 414 420 258 414 255 252 333 255}
\RouteRow{Tríadas estables por bloque}{0 1 2 3 4 5 6 7 8 9 10 11 11 12 14 15 16 17}
\end{tabularx}\endgroup\par\vspace{5pt}\noindent{\fontsize{8.4}{10}\selectfont\hyperref[trace-74]{Procedencia y códigos originales}\hfill Ficha completa · 074/324}
\clearpage\phantomsection\label{nav:vi-ruta-75}
% VI-CATALOG-ROW rutas 75
\pdfbookmark[3]{r3c1\_h+1v-1}{route-75}
\RouteTitle{Ruta 075}{r3c1\_\allowbreak{}h+1v-1}
\noindent Celda 19\quad (3, 1)\qquad Orientación \(h=1,\ v=-1\)\hfill\hyperref[sec:vi-indice-rutas]{Índice}\par\vspace{4pt}
\begingroup\fontsize{10.2}{12.5}\selectfont\setlength{\tabcolsep}{5pt}\renewcommand{\arraystretch}{1.1}\rowcolors{1}{routepale}{white}\noindent\begin{tabularx}{\linewidth}{@{}>{\raggedright\arraybackslash}p{.345\linewidth}>{\raggedright\arraybackslash}X@{}}
\RouteGroup{Identidad estructural}
\RouteRow{Clase y multisección}{charged\_\allowbreak{}lepton\_\allowbreak{}G4\quad /\quad charged\_\allowbreak{}lepton\_\allowbreak{}G4}
\RouteRow{Colección y sector}{leptón cargado; leptón}
\RouteRow{Clasificación física}{leptón cargado}
\RouteRow{Generación, profundidad y color}{4; G4; \textemdash{}}
\RouteRow{Ruta primaria y órbita de inversión}{\(B_1\): no\qquad 3,\allowbreak{}1;\allowbreak{}7,\allowbreak{}9}
\RouteRow{Firma de clasificación}{28|\allowbreak{}6|\allowbreak{}-10|\allowbreak{}0|\allowbreak{}-12|\allowbreak{}+++-\kern0pt{}-\kern0pt{}-+++-\kern0pt{}-\kern0pt{}-}
\RouteGroup{Retornos, memoria y transporte}
\RouteRow{Retornos con fase}{clase: 27; órbita: 27; celda: 27}
\RouteRow{Retorno cinemático}{en 108: 54; extensión: 216}
\RouteRow{Giros y cambios de hoja}{71; 107}
\RouteRow{Acarreo y diferimiento}{deuda total: 3; final: 1; bloques diferidos: 3}
\RouteRow{Tríadas finales; persistencia de Pareto}{16; no}
\RouteRow{\(K\longrightarrow K+9\)}{repeticiones: 9; cambios de memoria: 9}
\RouteGroup{Firma, fase y diferencias}
\RouteRow{\(n_A,\ s_C,\ \kappa_0,\ \nu_{120},\ \nu_{270}\)}{8, 2, 0, 0, -12}
\RouteRow{\(\tau_{12}\); firma histórica \(\mathbb Z^5\)}{+++-\kern0pt{}-\kern0pt{}-+++-\kern0pt{}-\kern0pt{}-; 8|\allowbreak{}2|\allowbreak{}0|\allowbreak{}0|\allowbreak{}-12}
\RouteRow{\(D_1,\ D_3,\ D_4,\ D_{27},\ D_{36}\)}{84, 24, 84, 24, 16}
\RouteRow{\(\Omega_{90/120},\ P_6\)}{80, 6}
\RouteGroup{Lectores y factores de acción}
\RouteRow{\(K_D\quad /\quad K_\Omega\)}{(40,\allowbreak{}84,\allowbreak{}30)\quad /\quad (80,\allowbreak{}54,\allowbreak{}6)}
\RouteRow{\(\kappa_D\)}{39.39035\allowbreak{}82768609\allowbreak{}24917849\allowbreak{}59}
\RouteRow{\(\kappa_\Omega\)}{79.60661\allowbreak{}01397948\allowbreak{}27586470\allowbreak{}91}
\RouteRow{\(R_{\mathrm{act},D}\)}{1.000000\allowbreak{}02730861\allowbreak{}73967087\allowbreak{}05}
\RouteRow{\(R_{\mathrm{act},\Omega}\)}{1.000000\allowbreak{}05518981\allowbreak{}25318591\allowbreak{}40}
\RouteGroup{Separación de prefijos}
\RouteRow{Área de ambigüedad; área \(\log_2\)}{54; 28.52932\allowbreak{}50129808}
\RouteRow{Primer bloque de separación}{órbita: \textemdash{}; celda: \textemdash{}; ruta: \textemdash{}}
\RouteRow{Ambigüedad final}{3}
\RouteGroup{Secuencias completas}
\RouteRow{\(w_6\)}{252 333 255 252 333 420 258 414 420 252 333 255 252 333 420 258 414 420}
\RouteRow{Tríadas estables por bloque}{0 1 2 3 4 5 6 7 8 9 10 11 12 12 14 14 16 16}
\end{tabularx}\endgroup\par\vspace{5pt}\noindent{\fontsize{8.4}{10}\selectfont\hyperref[trace-75]{Procedencia y códigos originales}\hfill Ficha completa · 075/324}
\clearpage\phantomsection\label{nav:vi-ruta-76}
% VI-CATALOG-ROW rutas 76
\pdfbookmark[3]{r3c1\_h+1v+1}{route-76}
\RouteTitle{Ruta 076}{r3c1\_\allowbreak{}h+1v+1}
\noindent Celda 19\quad (3, 1)\qquad Orientación \(h=1,\ v=1\)\hfill\hyperref[sec:vi-indice-rutas]{Índice}\par\vspace{4pt}
\begingroup\fontsize{10.2}{12.5}\selectfont\setlength{\tabcolsep}{5pt}\renewcommand{\arraystretch}{1.1}\rowcolors{1}{routepale}{white}\noindent\begin{tabularx}{\linewidth}{@{}>{\raggedright\arraybackslash}p{.345\linewidth}>{\raggedright\arraybackslash}X@{}}
\RouteGroup{Identidad estructural}
\RouteRow{Clase y multisección}{charged\_\allowbreak{}lepton\_\allowbreak{}G4\quad /\quad charged\_\allowbreak{}lepton\_\allowbreak{}G4}
\RouteRow{Colección y sector}{leptón cargado; leptón}
\RouteRow{Clasificación física}{leptón cargado}
\RouteRow{Generación, profundidad y color}{4; G4; \textemdash{}}
\RouteRow{Ruta primaria y órbita de inversión}{\(B_1\): no\qquad 3,\allowbreak{}1;\allowbreak{}7,\allowbreak{}9}
\RouteRow{Firma de clasificación}{28|\allowbreak{}6|\allowbreak{}-10|\allowbreak{}0|\allowbreak{}-12|\allowbreak{}+++-\kern0pt{}-\kern0pt{}-+++-\kern0pt{}-\kern0pt{}-}
\RouteGroup{Retornos, memoria y transporte}
\RouteRow{Retornos con fase}{clase: 18; órbita: 27; celda: 27}
\RouteRow{Retorno cinemático}{en 108: 54; extensión: 216}
\RouteRow{Giros y cambios de hoja}{71; 107}
\RouteRow{Acarreo y diferimiento}{deuda total: 0; final: 0; bloques diferidos: 0}
\RouteRow{Tríadas finales; persistencia de Pareto}{17; no}
\RouteRow{\(K\longrightarrow K+9\)}{repeticiones: 9; cambios de memoria: 9}
\RouteGroup{Firma, fase y diferencias}
\RouteRow{\(n_A,\ s_C,\ \kappa_0,\ \nu_{120},\ \nu_{270}\)}{28, 6, 0, 0, -12}
\RouteRow{\(\tau_{12}\); firma histórica \(\mathbb Z^5\)}{+++-\kern0pt{}-\kern0pt{}-+++-\kern0pt{}-\kern0pt{}-; 28|\allowbreak{}6|\allowbreak{}0|\allowbreak{}0|\allowbreak{}-12}
\RouteRow{\(D_1,\ D_3,\ D_4,\ D_{27},\ D_{36}\)}{66, 60, 66, 48, 32}
\RouteRow{\(\Omega_{90/120},\ P_6\)}{80, 6}
\RouteGroup{Lectores y factores de acción}
\RouteRow{\(K_D\quad /\quad K_\Omega\)}{(80,\allowbreak{}66,\allowbreak{}3)\quad /\quad (80,\allowbreak{}54,\allowbreak{}6)}
\RouteRow{\(\kappa_D\)}{79.51870\allowbreak{}83196953\allowbreak{}45554341\allowbreak{}34}
\RouteRow{\(\kappa_\Omega\)}{79.60661\allowbreak{}01397948\allowbreak{}27586470\allowbreak{}91}
\RouteRow{\(R_{\mathrm{act},D}\)}{1.000000\allowbreak{}05512887\allowbreak{}17997050\allowbreak{}72}
\RouteRow{\(R_{\mathrm{act},\Omega}\)}{1.000000\allowbreak{}05518981\allowbreak{}25318591\allowbreak{}40}
\RouteGroup{Separación de prefijos}
\RouteRow{Área de ambigüedad; área \(\log_2\)}{54; 28.52932\allowbreak{}50129808}
\RouteRow{Primer bloque de separación}{órbita: \textemdash{}; celda: \textemdash{}; ruta: \textemdash{}}
\RouteRow{Ambigüedad final}{3}
\RouteGroup{Secuencias completas}
\RouteRow{\(w_6\)}{243 657 24 243 657 570 264 90 570 243 657 24 243 657 570 264 90 570}
\RouteRow{Tríadas estables por bloque}{0 1 2 3 4 5 6 7 8 9 10 11 12 13 14 15 16 17}
\end{tabularx}\endgroup\par\vspace{5pt}\noindent{\fontsize{8.4}{10}\selectfont\hyperref[trace-76]{Procedencia y códigos originales}\hfill Ficha completa · 076/324}
\clearpage\phantomsection\label{nav:vi-ruta-77}
% VI-CATALOG-ROW rutas 77
\pdfbookmark[2]{Celda 20}{cell-20}
\pdfbookmark[3]{r3c2\_h-1v-1}{route-77}
\RouteTitle{Ruta 077}{r3c2\_\allowbreak{}h-1v-1}
\noindent Celda 20\quad (3, 2)\qquad Orientación \(h=-1,\ v=-1\)\hfill\hyperref[sec:vi-indice-rutas]{Índice}\par\vspace{4pt}
\begingroup\fontsize{10.2}{12.5}\selectfont\setlength{\tabcolsep}{5pt}\renewcommand{\arraystretch}{1.1}\rowcolors{1}{routepale}{white}\noindent\begin{tabularx}{\linewidth}{@{}>{\raggedright\arraybackslash}p{.345\linewidth}>{\raggedright\arraybackslash}X@{}}
\RouteGroup{Identidad estructural}
\RouteRow{Clase y multisección}{neutrino\_\allowbreak{}G3\quad /\quad neutrino\_\allowbreak{}G3}
\RouteRow{Colección y sector}{neutrino; leptón}
\RouteRow{Clasificación física}{neutrino}
\RouteRow{Generación, profundidad y color}{3; G3; \textemdash{}}
\RouteRow{Ruta primaria y órbita de inversión}{\(B_1\): no\qquad 3,\allowbreak{}2;\allowbreak{}7,\allowbreak{}8}
\RouteRow{Firma de clasificación}{28|\allowbreak{}-6|\allowbreak{}-4|\allowbreak{}-2|\allowbreak{}-8|\allowbreak{}0++-\kern0pt{}-\kern0pt{}-0++-\kern0pt{}-\kern0pt{}-}
\RouteGroup{Retornos, memoria y transporte}
\RouteRow{Retornos con fase}{clase: 27; órbita: 27; celda: 27}
\RouteRow{Retorno cinemático}{en 108: 54; extensión: 216}
\RouteRow{Giros y cambios de hoja}{71; 107}
\RouteRow{Acarreo y diferimiento}{deuda total: 1; final: 0; bloques diferidos: 1}
\RouteRow{Tríadas finales; persistencia de Pareto}{17; sí}
\RouteRow{\(K\longrightarrow K+9\)}{repeticiones: 9; cambios de memoria: 9}
\RouteGroup{Firma, fase y diferencias}
\RouteRow{\(n_A,\ s_C,\ \kappa_0,\ \nu_{120},\ \nu_{270}\)}{28, 6, -2, 4, -8}
\RouteRow{\(\tau_{12}\); firma histórica \(\mathbb Z^5\)}{+++0-\kern0pt{}-+++0-\kern0pt{}-; 28|\allowbreak{}6|\allowbreak{}-2|\allowbreak{}4|\allowbreak{}-8}
\RouteRow{\(D_1,\ D_3,\ D_4,\ D_{27},\ D_{36}\)}{66, 60, 66, 48, 32}
\RouteRow{\(\Omega_{90/120},\ P_6\)}{56, 5}
\RouteGroup{Lectores y factores de acción}
\RouteRow{\(K_D\quad /\quad K_\Omega\)}{(80,\allowbreak{}66,\allowbreak{}3)\quad /\quad (56,\allowbreak{}54,\allowbreak{}5)}
\RouteRow{\(\kappa_D\)}{79.51870\allowbreak{}83196953\allowbreak{}45554341\allowbreak{}34}
\RouteRow{\(\kappa_\Omega\)}{55.60649\allowbreak{}89433721\allowbreak{}35446416\allowbreak{}86}
\RouteRow{\(R_{\mathrm{act},D}\)}{1.000000\allowbreak{}05512887\allowbreak{}17997050\allowbreak{}72}
\RouteRow{\(R_{\mathrm{act},\Omega}\)}{1.000000\allowbreak{}03855097\allowbreak{}23541416\allowbreak{}45}
\RouteGroup{Separación de prefijos}
\RouteRow{Área de ambigüedad; área \(\log_2\)}{36; 18.0}
\RouteRow{Primer bloque de separación}{órbita: \textemdash{}; celda: \textemdash{}; ruta: \textemdash{}}
\RouteRow{Ambigüedad final}{2}
\RouteGroup{Secuencias completas}
\RouteRow{\(w_6\)}{486 423 15 486 423 327 498 99 327 486 423 15 486 423 327 498 99 327}
\RouteRow{Tríadas estables por bloque}{0 1 2 3 4 5 6 7 8 9 10 10 12 13 14 15 16 17}
\end{tabularx}\endgroup\par\vspace{5pt}\noindent{\fontsize{8.4}{10}\selectfont\hyperref[trace-77]{Procedencia y códigos originales}\hfill Ficha completa · 077/324}
\clearpage\phantomsection\label{nav:vi-ruta-78}
% VI-CATALOG-ROW rutas 78
\pdfbookmark[3]{r3c2\_h-1v+1}{route-78}
\RouteTitle{Ruta 078}{r3c2\_\allowbreak{}h-1v+1}
\noindent Celda 20\quad (3, 2)\qquad Orientación \(h=-1,\ v=1\)\hfill\hyperref[sec:vi-indice-rutas]{Índice}\par\vspace{4pt}
\begingroup\fontsize{10.2}{12.5}\selectfont\setlength{\tabcolsep}{5pt}\renewcommand{\arraystretch}{1.1}\rowcolors{1}{routepale}{white}\noindent\begin{tabularx}{\linewidth}{@{}>{\raggedright\arraybackslash}p{.345\linewidth}>{\raggedright\arraybackslash}X@{}}
\RouteGroup{Identidad estructural}
\RouteRow{Clase y multisección}{neutrino\_\allowbreak{}G3\quad /\quad neutrino\_\allowbreak{}G3}
\RouteRow{Colección y sector}{neutrino; leptón}
\RouteRow{Clasificación física}{neutrino}
\RouteRow{Generación, profundidad y color}{3; G3; \textemdash{}}
\RouteRow{Ruta primaria y órbita de inversión}{\(B_1\): no\qquad 3,\allowbreak{}2;\allowbreak{}7,\allowbreak{}8}
\RouteRow{Firma de clasificación}{28|\allowbreak{}-6|\allowbreak{}-4|\allowbreak{}-2|\allowbreak{}-8|\allowbreak{}0++-\kern0pt{}-\kern0pt{}-0++-\kern0pt{}-\kern0pt{}-}
\RouteGroup{Retornos, memoria y transporte}
\RouteRow{Retornos con fase}{clase: 27; órbita: 27; celda: 27}
\RouteRow{Retorno cinemático}{en 108: 54; extensión: 216}
\RouteRow{Giros y cambios de hoja}{71; 107}
\RouteRow{Acarreo y diferimiento}{deuda total: 1; final: 0; bloques diferidos: 1}
\RouteRow{Tríadas finales; persistencia de Pareto}{17; sí}
\RouteRow{\(K\longrightarrow K+9\)}{repeticiones: 9; cambios de memoria: 9}
\RouteGroup{Firma, fase y diferencias}
\RouteRow{\(n_A,\ s_C,\ \kappa_0,\ \nu_{120},\ \nu_{270}\)}{44, 2, 0, 2, -12}
\RouteRow{\(\tau_{12}\); firma histórica \(\mathbb Z^5\)}{+++-\kern0pt{}-\kern0pt{}-+++-\kern0pt{}-\kern0pt{}-; 44|\allowbreak{}2|\allowbreak{}0|\allowbreak{}2|\allowbreak{}-12}
\RouteRow{\(D_1,\ D_3,\ D_4,\ D_{27},\ D_{36}\)}{78, 24, 78, 24, 16}
\RouteRow{\(\Omega_{90/120},\ P_6\)}{80, 6}
\RouteGroup{Lectores y factores de acción}
\RouteRow{\(K_D\quad /\quad K_\Omega\)}{(40,\allowbreak{}78,\allowbreak{}27)\quad /\quad (80,\allowbreak{}54,\allowbreak{}6)}
\RouteRow{\(\kappa_D\)}{39.43380\allowbreak{}88030085\allowbreak{}51303671\allowbreak{}14}
\RouteRow{\(\kappa_\Omega\)}{79.60661\allowbreak{}01397948\allowbreak{}27586470\allowbreak{}91}
\RouteRow{\(R_{\mathrm{act},D}\)}{1.000000\allowbreak{}02733874\allowbreak{}08548943\allowbreak{}58}
\RouteRow{\(R_{\mathrm{act},\Omega}\)}{1.000000\allowbreak{}05518981\allowbreak{}25318591\allowbreak{}40}
\RouteGroup{Separación de prefijos}
\RouteRow{Área de ambigüedad; área \(\log_2\)}{36; 18.0}
\RouteRow{Primer bloque de separación}{órbita: \textemdash{}; celda: \textemdash{}; ruta: \textemdash{}}
\RouteRow{Ambigüedad final}{2}
\RouteGroup{Secuencias completas}
\RouteRow{\(w_6\)}{504 585 507 504 585 672 510 666 672 504 585 507 504 585 672 510 666 672}
\RouteRow{Tríadas estables por bloque}{0 1 2 3 4 5 6 7 8 9 10 11 12 13 14 14 16 17}
\end{tabularx}\endgroup\par\vspace{5pt}\noindent{\fontsize{8.4}{10}\selectfont\hyperref[trace-78]{Procedencia y códigos originales}\hfill Ficha completa · 078/324}
\clearpage\phantomsection\label{nav:vi-ruta-79}
% VI-CATALOG-ROW rutas 79
\pdfbookmark[3]{r3c2\_h+1v-1}{route-79}
\RouteTitle{Ruta 079}{r3c2\_\allowbreak{}h+1v-1}
\noindent Celda 20\quad (3, 2)\qquad Orientación \(h=1,\ v=-1\)\hfill\hyperref[sec:vi-indice-rutas]{Índice}\par\vspace{4pt}
\begingroup\fontsize{10.2}{12.5}\selectfont\setlength{\tabcolsep}{5pt}\renewcommand{\arraystretch}{1.1}\rowcolors{1}{routepale}{white}\noindent\begin{tabularx}{\linewidth}{@{}>{\raggedright\arraybackslash}p{.345\linewidth}>{\raggedright\arraybackslash}X@{}}
\RouteGroup{Identidad estructural}
\RouteRow{Clase y multisección}{neutrino\_\allowbreak{}G3\quad /\quad neutrino\_\allowbreak{}G3}
\RouteRow{Colección y sector}{neutrino; leptón}
\RouteRow{Clasificación física}{neutrino}
\RouteRow{Generación, profundidad y color}{3; G3; \textemdash{}}
\RouteRow{Ruta primaria y órbita de inversión}{\(B_1\): no\qquad 3,\allowbreak{}2;\allowbreak{}7,\allowbreak{}8}
\RouteRow{Firma de clasificación}{28|\allowbreak{}-6|\allowbreak{}-4|\allowbreak{}-2|\allowbreak{}-8|\allowbreak{}0++-\kern0pt{}-\kern0pt{}-0++-\kern0pt{}-\kern0pt{}-}
\RouteGroup{Retornos, memoria y transporte}
\RouteRow{Retornos con fase}{clase: 27; órbita: 27; celda: 27}
\RouteRow{Retorno cinemático}{en 108: 54; extensión: 216}
\RouteRow{Giros y cambios de hoja}{71; 107}
\RouteRow{Acarreo y diferimiento}{deuda total: 1; final: 0; bloques diferidos: 1}
\RouteRow{Tríadas finales; persistencia de Pareto}{17; sí}
\RouteRow{\(K\longrightarrow K+9\)}{repeticiones: 9; cambios de memoria: 9}
\RouteGroup{Firma, fase y diferencias}
\RouteRow{\(n_A,\ s_C,\ \kappa_0,\ \nu_{120},\ \nu_{270}\)}{44, -2, 0, -2, -12}
\RouteRow{\(\tau_{12}\); firma histórica \(\mathbb Z^5\)}{+++-\kern0pt{}-\kern0pt{}-+++-\kern0pt{}-\kern0pt{}-; 44|\allowbreak{}-2|\allowbreak{}0|\allowbreak{}-2|\allowbreak{}-12}
\RouteRow{\(D_1,\ D_3,\ D_4,\ D_{27},\ D_{36}\)}{78, 24, 78, 24, 16}
\RouteRow{\(\Omega_{90/120},\ P_6\)}{80, 6}
\RouteGroup{Lectores y factores de acción}
\RouteRow{\(K_D\quad /\quad K_\Omega\)}{(40,\allowbreak{}78,\allowbreak{}27)\quad /\quad (80,\allowbreak{}54,\allowbreak{}6)}
\RouteRow{\(\kappa_D\)}{39.43380\allowbreak{}88030085\allowbreak{}51303671\allowbreak{}14}
\RouteRow{\(\kappa_\Omega\)}{79.60661\allowbreak{}01397948\allowbreak{}27586470\allowbreak{}91}
\RouteRow{\(R_{\mathrm{act},D}\)}{1.000000\allowbreak{}02733874\allowbreak{}08548943\allowbreak{}58}
\RouteRow{\(R_{\mathrm{act},\Omega}\)}{1.000000\allowbreak{}05518981\allowbreak{}25318591\allowbreak{}40}
\RouteGroup{Separación de prefijos}
\RouteRow{Área de ambigüedad; área \(\log_2\)}{36; 18.0}
\RouteRow{Primer bloque de separación}{órbita: \textemdash{}; celda: \textemdash{}; ruta: \textemdash{}}
\RouteRow{Ambigüedad final}{2}
\RouteGroup{Secuencias completas}
\RouteRow{\(w_6\)}{510 666 672 510 666 507 504 585 507 510 666 672 510 666 507 504 585 507}
\RouteRow{Tríadas estables por bloque}{0 1 2 3 4 5 6 7 8 9 10 11 12 13 14 15 15 17}
\end{tabularx}\endgroup\par\vspace{5pt}\noindent{\fontsize{8.4}{10}\selectfont\hyperref[trace-79]{Procedencia y códigos originales}\hfill Ficha completa · 079/324}
\clearpage\phantomsection\label{nav:vi-ruta-80}
% VI-CATALOG-ROW rutas 80
\pdfbookmark[3]{r3c2\_h+1v+1}{route-80}
\RouteTitle{Ruta 080}{r3c2\_\allowbreak{}h+1v+1}
\noindent Celda 20\quad (3, 2)\qquad Orientación \(h=1,\ v=1\)\hfill\hyperref[sec:vi-indice-rutas]{Índice}\par\vspace{4pt}
\begingroup\fontsize{10.2}{12.5}\selectfont\setlength{\tabcolsep}{5pt}\renewcommand{\arraystretch}{1.1}\rowcolors{1}{routepale}{white}\noindent\begin{tabularx}{\linewidth}{@{}>{\raggedright\arraybackslash}p{.345\linewidth}>{\raggedright\arraybackslash}X@{}}
\RouteGroup{Identidad estructural}
\RouteRow{Clase y multisección}{neutrino\_\allowbreak{}G3\quad /\quad neutrino\_\allowbreak{}G3}
\RouteRow{Colección y sector}{neutrino; leptón}
\RouteRow{Clasificación física}{neutrino}
\RouteRow{Generación, profundidad y color}{3; G3; \textemdash{}}
\RouteRow{Ruta primaria y órbita de inversión}{\(B_1\): sí\qquad 3,\allowbreak{}2;\allowbreak{}7,\allowbreak{}8}
\RouteRow{Firma de clasificación}{28|\allowbreak{}-6|\allowbreak{}-4|\allowbreak{}-2|\allowbreak{}-8|\allowbreak{}0++-\kern0pt{}-\kern0pt{}-0++-\kern0pt{}-\kern0pt{}-}
\RouteGroup{Retornos, memoria y transporte}
\RouteRow{Retornos con fase}{clase: 27; órbita: 27; celda: 27}
\RouteRow{Retorno cinemático}{en 108: 54; extensión: 216}
\RouteRow{Giros y cambios de hoja}{71; 107}
\RouteRow{Acarreo y diferimiento}{deuda total: 0; final: 0; bloques diferidos: 0}
\RouteRow{Tríadas finales; persistencia de Pareto}{17; sí}
\RouteRow{\(K\longrightarrow K+9\)}{repeticiones: 9; cambios de memoria: 9}
\RouteGroup{Firma, fase y diferencias}
\RouteRow{\(n_A,\ s_C,\ \kappa_0,\ \nu_{120},\ \nu_{270}\)}{28, -6, 2, -4, -8}
\RouteRow{\(\tau_{12}\); firma histórica \(\mathbb Z^5\)}{0++-\kern0pt{}-\kern0pt{}-0++-\kern0pt{}-\kern0pt{}-; 28|\allowbreak{}-6|\allowbreak{}2|\allowbreak{}-4|\allowbreak{}-8}
\RouteRow{\(D_1,\ D_3,\ D_4,\ D_{27},\ D_{36}\)}{66, 60, 66, 48, 32}
\RouteRow{\(\Omega_{90/120},\ P_6\)}{56, 5}
\RouteGroup{Lectores y factores de acción}
\RouteRow{\(K_D\quad /\quad K_\Omega\)}{(80,\allowbreak{}66,\allowbreak{}3)\quad /\quad (56,\allowbreak{}54,\allowbreak{}5)}
\RouteRow{\(\kappa_D\)}{79.51870\allowbreak{}83196953\allowbreak{}45554341\allowbreak{}34}
\RouteRow{\(\kappa_\Omega\)}{55.60649\allowbreak{}89433721\allowbreak{}35446416\allowbreak{}86}
\RouteRow{\(R_{\mathrm{act},D}\)}{1.000000\allowbreak{}05512887\allowbreak{}17997050\allowbreak{}72}
\RouteRow{\(R_{\mathrm{act},\Omega}\)}{1.000000\allowbreak{}03855097\allowbreak{}23541416\allowbreak{}45}
\RouteGroup{Separación de prefijos}
\RouteRow{Área de ambigüedad; área \(\log_2\)}{36; 18.0}
\RouteRow{Primer bloque de separación}{órbita: \textemdash{}; celda: \textemdash{}; ruta: \textemdash{}}
\RouteRow{Ambigüedad final}{2}
\RouteGroup{Secuencias completas}
\RouteRow{\(w_6\)}{498 99 327 498 99 15 486 423 15 498 99 327 498 99 15 486 423 15}
\RouteRow{Tríadas estables por bloque}{0 1 2 3 4 5 6 7 8 9 10 11 12 13 14 15 16 17}
\end{tabularx}\endgroup\par\vspace{5pt}\noindent{\fontsize{8.4}{10}\selectfont\hyperref[trace-80]{Procedencia y códigos originales}\hfill Ficha completa · 080/324}
\clearpage\phantomsection\label{nav:vi-ruta-81}
% VI-CATALOG-ROW rutas 81
\pdfbookmark[2]{Celda 21}{cell-21}
\pdfbookmark[3]{r3c3\_h-1v-1}{route-81}
\RouteTitle{Ruta 081}{r3c3\_\allowbreak{}h-1v-1}
\noindent Celda 21\quad (3, 3)\qquad Orientación \(h=-1,\ v=-1\)\hfill\hyperref[sec:vi-indice-rutas]{Índice}\par\vspace{4pt}
\begingroup\fontsize{10.2}{12.5}\selectfont\setlength{\tabcolsep}{5pt}\renewcommand{\arraystretch}{1.1}\rowcolors{1}{routepale}{white}\noindent\begin{tabularx}{\linewidth}{@{}>{\raggedright\arraybackslash}p{.345\linewidth}>{\raggedright\arraybackslash}X@{}}
\RouteGroup{Identidad estructural}
\RouteRow{Clase y multisección}{charged\_\allowbreak{}lepton\_\allowbreak{}G2\quad /\quad charged\_\allowbreak{}lepton\_\allowbreak{}G2}
\RouteRow{Colección y sector}{leptón cargado; leptón}
\RouteRow{Clasificación física}{leptón cargado}
\RouteRow{Generación, profundidad y color}{2; G2; \textemdash{}}
\RouteRow{Ruta primaria y órbita de inversión}{\(B_1\): no\qquad 3,\allowbreak{}3;\allowbreak{}7,\allowbreak{}7}
\RouteRow{Firma de clasificación}{24|\allowbreak{}0|\allowbreak{}12|\allowbreak{}0|\allowbreak{}-12|\allowbreak{}+++-\kern0pt{}-\kern0pt{}-+++-\kern0pt{}-\kern0pt{}-}
\RouteGroup{Retornos, memoria y transporte}
\RouteRow{Retornos con fase}{clase: 27; órbita: 27; celda: 27}
\RouteRow{Retorno cinemático}{en 108: 54; extensión: 216}
\RouteRow{Giros y cambios de hoja}{71; 107}
\RouteRow{Acarreo y diferimiento}{deuda total: 1; final: 0; bloques diferidos: 1}
\RouteRow{Tríadas finales; persistencia de Pareto}{17; no}
\RouteRow{\(K\longrightarrow K+9\)}{repeticiones: 9; cambios de memoria: 9}
\RouteGroup{Firma, fase y diferencias}
\RouteRow{\(n_A,\ s_C,\ \kappa_0,\ \nu_{120},\ \nu_{270}\)}{24, 0, 0, 0, -12}
\RouteRow{\(\tau_{12}\); firma histórica \(\mathbb Z^5\)}{+++-\kern0pt{}-\kern0pt{}-+++-\kern0pt{}-\kern0pt{}-; 24|\allowbreak{}0|\allowbreak{}0|\allowbreak{}0|\allowbreak{}-12}
\RouteRow{\(D_1,\ D_3,\ D_4,\ D_{27},\ D_{36}\)}{54, 60, 72, 24, 16}
\RouteRow{\(\Omega_{90/120},\ P_6\)}{80, 6}
\RouteGroup{Lectores y factores de acción}
\RouteRow{\(K_D\quad /\quad K_\Omega\)}{(40,\allowbreak{}54,\allowbreak{}6)\quad /\quad (80,\allowbreak{}54,\allowbreak{}6)}
\RouteRow{\(\kappa_D\)}{39.60661\allowbreak{}01397948\allowbreak{}27586470\allowbreak{}91}
\RouteRow{\(\kappa_\Omega\)}{79.60661\allowbreak{}01397948\allowbreak{}27586470\allowbreak{}91}
\RouteRow{\(R_{\mathrm{act},D}\)}{1.000000\allowbreak{}02745854\allowbreak{}08735595\allowbreak{}18}
\RouteRow{\(R_{\mathrm{act},\Omega}\)}{1.000000\allowbreak{}05518981\allowbreak{}25318591\allowbreak{}40}
\RouteGroup{Separación de prefijos}
\RouteRow{Área de ambigüedad; área \(\log_2\)}{22; 4.0}
\RouteRow{Primer bloque de separación}{órbita: 5; celda: 5; ruta: 5}
\RouteRow{Ambigüedad final}{1}
\RouteGroup{Secuencias completas}
\RouteRow{\(w_6\)}{15 486 423 15 486 657 24 243 657 15 486 423 15 486 657 24 243 657}
\RouteRow{Tríadas estables por bloque}{0 1 2 3 4 5 6 7 8 9 10 11 12 13 14 15 15 17}
\end{tabularx}\endgroup\par\vspace{5pt}\noindent{\fontsize{8.4}{10}\selectfont\hyperref[trace-81]{Procedencia y códigos originales}\hfill Ficha completa · 081/324}
\clearpage\phantomsection\label{nav:vi-ruta-82}
% VI-CATALOG-ROW rutas 82
\pdfbookmark[3]{r3c3\_h-1v+1}{route-82}
\RouteTitle{Ruta 082}{r3c3\_\allowbreak{}h-1v+1}
\noindent Celda 21\quad (3, 3)\qquad Orientación \(h=-1,\ v=1\)\hfill\hyperref[sec:vi-indice-rutas]{Índice}\par\vspace{4pt}
\begingroup\fontsize{10.2}{12.5}\selectfont\setlength{\tabcolsep}{5pt}\renewcommand{\arraystretch}{1.1}\rowcolors{1}{routepale}{white}\noindent\begin{tabularx}{\linewidth}{@{}>{\raggedright\arraybackslash}p{.345\linewidth}>{\raggedright\arraybackslash}X@{}}
\RouteGroup{Identidad estructural}
\RouteRow{Clase y multisección}{charged\_\allowbreak{}lepton\_\allowbreak{}G2\quad /\quad charged\_\allowbreak{}lepton\_\allowbreak{}G2}
\RouteRow{Colección y sector}{leptón cargado; leptón}
\RouteRow{Clasificación física}{leptón cargado}
\RouteRow{Generación, profundidad y color}{2; G2; \textemdash{}}
\RouteRow{Ruta primaria y órbita de inversión}{\(B_1\): no\qquad 3,\allowbreak{}3;\allowbreak{}7,\allowbreak{}7}
\RouteRow{Firma de clasificación}{24|\allowbreak{}0|\allowbreak{}12|\allowbreak{}0|\allowbreak{}-12|\allowbreak{}+++-\kern0pt{}-\kern0pt{}-+++-\kern0pt{}-\kern0pt{}-}
\RouteGroup{Retornos, memoria y transporte}
\RouteRow{Retornos con fase}{clase: 27; órbita: 27; celda: 27}
\RouteRow{Retorno cinemático}{en 108: 54; extensión: 216}
\RouteRow{Giros y cambios de hoja}{71; 107}
\RouteRow{Acarreo y diferimiento}{deuda total: 0; final: 0; bloques diferidos: 0}
\RouteRow{Tríadas finales; persistencia de Pareto}{17; no}
\RouteRow{\(K\longrightarrow K+9\)}{repeticiones: 9; cambios de memoria: 9}
\RouteGroup{Firma, fase y diferencias}
\RouteRow{\(n_A,\ s_C,\ \kappa_0,\ \nu_{120},\ \nu_{270}\)}{14, 6, 0, 0, -12}
\RouteRow{\(\tau_{12}\); firma histórica \(\mathbb Z^5\)}{+++-\kern0pt{}-\kern0pt{}-+++-\kern0pt{}-\kern0pt{}-; 14|\allowbreak{}6|\allowbreak{}0|\allowbreak{}0|\allowbreak{}-12}
\RouteRow{\(D_1,\ D_3,\ D_4,\ D_{27},\ D_{36}\)}{24, 24, 24, 0, 0}
\RouteRow{\(\Omega_{90/120},\ P_6\)}{80, 6}
\RouteGroup{Lectores y factores de acción}
\RouteRow{\(K_D\quad /\quad K_\Omega\)}{(0,\allowbreak{}24,\allowbreak{}0)\quad /\quad (80,\allowbreak{}54,\allowbreak{}6)}
\RouteRow{\(\kappa_D\)}{-0.175136\allowbreak{}46166281\allowbreak{}12239348\allowbreak{}425}
\RouteRow{\(\kappa_\Omega\)}{79.60661\allowbreak{}01397948\allowbreak{}27586470\allowbreak{}91}
\RouteRow{\(R_{\mathrm{act},D}\)}{0.999999\allowbreak{}99987858\allowbreak{}10851337\allowbreak{}881}
\RouteRow{\(R_{\mathrm{act},\Omega}\)}{1.000000\allowbreak{}05518981\allowbreak{}25318591\allowbreak{}40}
\RouteGroup{Separación de prefijos}
\RouteRow{Área de ambigüedad; área \(\log_2\)}{36; 18.0}
\RouteRow{Primer bloque de separación}{órbita: 1; celda: 1; ruta: \textemdash{}}
\RouteRow{Ambigüedad final}{2}
\RouteGroup{Secuencias completas}
\RouteRow{\(w_6\)}{3 0 81 3 0 81 3 0 81 3 0 81 3 0 81 3 0 81}
\RouteRow{Tríadas estables por bloque}{0 1 2 3 4 5 6 7 8 9 10 11 12 13 14 15 16 17}
\end{tabularx}\endgroup\par\vspace{5pt}\noindent{\fontsize{8.4}{10}\selectfont\hyperref[trace-82]{Procedencia y códigos originales}\hfill Ficha completa · 082/324}
\clearpage\phantomsection\label{nav:vi-ruta-83}
% VI-CATALOG-ROW rutas 83
\pdfbookmark[3]{r3c3\_h+1v-1}{route-83}
\RouteTitle{Ruta 083}{r3c3\_\allowbreak{}h+1v-1}
\noindent Celda 21\quad (3, 3)\qquad Orientación \(h=1,\ v=-1\)\hfill\hyperref[sec:vi-indice-rutas]{Índice}\par\vspace{4pt}
\begingroup\fontsize{10.2}{12.5}\selectfont\setlength{\tabcolsep}{5pt}\renewcommand{\arraystretch}{1.1}\rowcolors{1}{routepale}{white}\noindent\begin{tabularx}{\linewidth}{@{}>{\raggedright\arraybackslash}p{.345\linewidth}>{\raggedright\arraybackslash}X@{}}
\RouteGroup{Identidad estructural}
\RouteRow{Clase y multisección}{charged\_\allowbreak{}lepton\_\allowbreak{}G2\quad /\quad charged\_\allowbreak{}lepton\_\allowbreak{}G2}
\RouteRow{Colección y sector}{leptón cargado; leptón}
\RouteRow{Clasificación física}{leptón cargado}
\RouteRow{Generación, profundidad y color}{2; G2; \textemdash{}}
\RouteRow{Ruta primaria y órbita de inversión}{\(B_1\): sí\qquad 3,\allowbreak{}3;\allowbreak{}7,\allowbreak{}7}
\RouteRow{Firma de clasificación}{24|\allowbreak{}0|\allowbreak{}12|\allowbreak{}0|\allowbreak{}-12|\allowbreak{}+++-\kern0pt{}-\kern0pt{}-+++-\kern0pt{}-\kern0pt{}-}
\RouteGroup{Retornos, memoria y transporte}
\RouteRow{Retornos con fase}{clase: 27; órbita: 27; celda: 27}
\RouteRow{Retorno cinemático}{en 108: 54; extensión: 216}
\RouteRow{Giros y cambios de hoja}{71; 107}
\RouteRow{Acarreo y diferimiento}{deuda total: 0; final: 0; bloques diferidos: 0}
\RouteRow{Tríadas finales; persistencia de Pareto}{17; no}
\RouteRow{\(K\longrightarrow K+9\)}{repeticiones: 9; cambios de memoria: 9}
\RouteGroup{Firma, fase y diferencias}
\RouteRow{\(n_A,\ s_C,\ \kappa_0,\ \nu_{120},\ \nu_{270}\)}{14, -6, 0, 0, -12}
\RouteRow{\(\tau_{12}\); firma histórica \(\mathbb Z^5\)}{+++-\kern0pt{}-\kern0pt{}-+++-\kern0pt{}-\kern0pt{}-; 14|\allowbreak{}-6|\allowbreak{}0|\allowbreak{}0|\allowbreak{}-12}
\RouteRow{\(D_1,\ D_3,\ D_4,\ D_{27},\ D_{36}\)}{24, 24, 24, 0, 0}
\RouteRow{\(\Omega_{90/120},\ P_6\)}{80, 6}
\RouteGroup{Lectores y factores de acción}
\RouteRow{\(K_D\quad /\quad K_\Omega\)}{(0,\allowbreak{}24,\allowbreak{}0)\quad /\quad (80,\allowbreak{}54,\allowbreak{}6)}
\RouteRow{\(\kappa_D\)}{-0.175136\allowbreak{}46166281\allowbreak{}12239348\allowbreak{}425}
\RouteRow{\(\kappa_\Omega\)}{79.60661\allowbreak{}01397948\allowbreak{}27586470\allowbreak{}91}
\RouteRow{\(R_{\mathrm{act},D}\)}{0.999999\allowbreak{}99987858\allowbreak{}10851337\allowbreak{}881}
\RouteRow{\(R_{\mathrm{act},\Omega}\)}{1.000000\allowbreak{}05518981\allowbreak{}25318591\allowbreak{}40}
\RouteGroup{Separación de prefijos}
\RouteRow{Área de ambigüedad; área \(\log_2\)}{36; 18.0}
\RouteRow{Primer bloque de separación}{órbita: 1; celda: 1; ruta: \textemdash{}}
\RouteRow{Ambigüedad final}{2}
\RouteGroup{Secuencias completas}
\RouteRow{\(w_6\)}{3 0 81 3 0 81 3 0 81 3 0 81 3 0 81 3 0 81}
\RouteRow{Tríadas estables por bloque}{0 1 2 3 4 5 6 7 8 9 10 11 12 13 14 15 16 17}
\end{tabularx}\endgroup\par\vspace{5pt}\noindent{\fontsize{8.4}{10}\selectfont\hyperref[trace-83]{Procedencia y códigos originales}\hfill Ficha completa · 083/324}
\clearpage\phantomsection\label{nav:vi-ruta-84}
% VI-CATALOG-ROW rutas 84
\pdfbookmark[3]{r3c3\_h+1v+1}{route-84}
\RouteTitle{Ruta 084}{r3c3\_\allowbreak{}h+1v+1}
\noindent Celda 21\quad (3, 3)\qquad Orientación \(h=1,\ v=1\)\hfill\hyperref[sec:vi-indice-rutas]{Índice}\par\vspace{4pt}
\begingroup\fontsize{10.2}{12.5}\selectfont\setlength{\tabcolsep}{5pt}\renewcommand{\arraystretch}{1.1}\rowcolors{1}{routepale}{white}\noindent\begin{tabularx}{\linewidth}{@{}>{\raggedright\arraybackslash}p{.345\linewidth}>{\raggedright\arraybackslash}X@{}}
\RouteGroup{Identidad estructural}
\RouteRow{Clase y multisección}{charged\_\allowbreak{}lepton\_\allowbreak{}G2\quad /\quad charged\_\allowbreak{}lepton\_\allowbreak{}G2}
\RouteRow{Colección y sector}{leptón cargado; leptón}
\RouteRow{Clasificación física}{leptón cargado}
\RouteRow{Generación, profundidad y color}{2; G2; \textemdash{}}
\RouteRow{Ruta primaria y órbita de inversión}{\(B_1\): no\qquad 3,\allowbreak{}3;\allowbreak{}7,\allowbreak{}7}
\RouteRow{Firma de clasificación}{24|\allowbreak{}0|\allowbreak{}12|\allowbreak{}0|\allowbreak{}-12|\allowbreak{}+++-\kern0pt{}-\kern0pt{}-+++-\kern0pt{}-\kern0pt{}-}
\RouteGroup{Retornos, memoria y transporte}
\RouteRow{Retornos con fase}{clase: 27; órbita: 27; celda: 27}
\RouteRow{Retorno cinemático}{en 108: 54; extensión: 216}
\RouteRow{Giros y cambios de hoja}{71; 107}
\RouteRow{Acarreo y diferimiento}{deuda total: 2; final: 1; bloques diferidos: 2}
\RouteRow{Tríadas finales; persistencia de Pareto}{16; no}
\RouteRow{\(K\longrightarrow K+9\)}{repeticiones: 9; cambios de memoria: 9}
\RouteGroup{Firma, fase y diferencias}
\RouteRow{\(n_A,\ s_C,\ \kappa_0,\ \nu_{120},\ \nu_{270}\)}{24, 0, 0, 0, -12}
\RouteRow{\(\tau_{12}\); firma histórica \(\mathbb Z^5\)}{+++-\kern0pt{}-\kern0pt{}-+++-\kern0pt{}-\kern0pt{}-; 24|\allowbreak{}0|\allowbreak{}0|\allowbreak{}0|\allowbreak{}-12}
\RouteRow{\(D_1,\ D_3,\ D_4,\ D_{27},\ D_{36}\)}{54, 60, 72, 24, 16}
\RouteRow{\(\Omega_{90/120},\ P_6\)}{80, 6}
\RouteGroup{Lectores y factores de acción}
\RouteRow{\(K_D\quad /\quad K_\Omega\)}{(40,\allowbreak{}54,\allowbreak{}6)\quad /\quad (80,\allowbreak{}54,\allowbreak{}6)}
\RouteRow{\(\kappa_D\)}{39.60661\allowbreak{}01397948\allowbreak{}27586470\allowbreak{}91}
\RouteRow{\(\kappa_\Omega\)}{79.60661\allowbreak{}01397948\allowbreak{}27586470\allowbreak{}91}
\RouteRow{\(R_{\mathrm{act},D}\)}{1.000000\allowbreak{}02745854\allowbreak{}08735595\allowbreak{}18}
\RouteRow{\(R_{\mathrm{act},\Omega}\)}{1.000000\allowbreak{}05518981\allowbreak{}25318591\allowbreak{}40}
\RouteGroup{Separación de prefijos}
\RouteRow{Área de ambigüedad; área \(\log_2\)}{22; 4.0}
\RouteRow{Primer bloque de separación}{órbita: 5; celda: 5; ruta: 5}
\RouteRow{Ambigüedad final}{1}
\RouteGroup{Secuencias completas}
\RouteRow{\(w_6\)}{24 243 657 24 243 423 15 486 423 24 243 657 24 243 423 15 486 423}
\RouteRow{Tríadas estables por bloque}{0 1 2 3 4 5 6 7 8 9 10 11 12 13 13 15 16 16}
\end{tabularx}\endgroup\par\vspace{5pt}\noindent{\fontsize{8.4}{10}\selectfont\hyperref[trace-84]{Procedencia y códigos originales}\hfill Ficha completa · 084/324}
\clearpage\phantomsection\label{nav:vi-ruta-85}
% VI-CATALOG-ROW rutas 85
\pdfbookmark[2]{Celda 22}{cell-22}
\pdfbookmark[3]{r3c4\_h-1v-1}{route-85}
\RouteTitle{Ruta 085}{r3c4\_\allowbreak{}h-1v-1}
\noindent Celda 22\quad (3, 4)\qquad Orientación \(h=-1,\ v=-1\)\hfill\hyperref[sec:vi-indice-rutas]{Índice}\par\vspace{4pt}
\begingroup\fontsize{10.2}{12.5}\selectfont\setlength{\tabcolsep}{5pt}\renewcommand{\arraystretch}{1.1}\rowcolors{1}{routepale}{white}\noindent\begin{tabularx}{\linewidth}{@{}>{\raggedright\arraybackslash}p{.345\linewidth}>{\raggedright\arraybackslash}X@{}}
\RouteGroup{Identidad estructural}
\RouteRow{Clase y multisección}{neutrino\_\allowbreak{}G2\quad /\quad neutrino\_\allowbreak{}G2}
\RouteRow{Colección y sector}{neutrino; leptón}
\RouteRow{Clasificación física}{neutrino}
\RouteRow{Generación, profundidad y color}{2; G2; \textemdash{}}
\RouteRow{Ruta primaria y órbita de inversión}{\(B_1\): sí\qquad 3,\allowbreak{}4;\allowbreak{}7,\allowbreak{}6}
\RouteRow{Firma de clasificación}{28|\allowbreak{}6|\allowbreak{}0|\allowbreak{}0|\allowbreak{}-12|\allowbreak{}+++-\kern0pt{}-\kern0pt{}-+++-\kern0pt{}-\kern0pt{}-}
\RouteGroup{Retornos, memoria y transporte}
\RouteRow{Retornos con fase}{clase: 27; órbita: 27; celda: 27}
\RouteRow{Retorno cinemático}{en 108: 54; extensión: 216}
\RouteRow{Giros y cambios de hoja}{71; 107}
\RouteRow{Acarreo y diferimiento}{deuda total: 2; final: 1; bloques diferidos: 1}
\RouteRow{Tríadas finales; persistencia de Pareto}{16; no}
\RouteRow{\(K\longrightarrow K+9\)}{repeticiones: 9; cambios de memoria: 9}
\RouteGroup{Firma, fase y diferencias}
\RouteRow{\(n_A,\ s_C,\ \kappa_0,\ \nu_{120},\ \nu_{270}\)}{28, -6, 0, -4, -12}
\RouteRow{\(\tau_{12}\); firma histórica \(\mathbb Z^5\)}{+++-\kern0pt{}-\kern0pt{}-+++-\kern0pt{}-\kern0pt{}-; 28|\allowbreak{}-6|\allowbreak{}0|\allowbreak{}-4|\allowbreak{}-12}
\RouteRow{\(D_1,\ D_3,\ D_4,\ D_{27},\ D_{36}\)}{66, 60, 66, 48, 32}
\RouteRow{\(\Omega_{90/120},\ P_6\)}{80, 6}
\RouteGroup{Lectores y factores de acción}
\RouteRow{\(K_D\quad /\quad K_\Omega\)}{(80,\allowbreak{}66,\allowbreak{}3)\quad /\quad (80,\allowbreak{}54,\allowbreak{}6)}
\RouteRow{\(\kappa_D\)}{79.51870\allowbreak{}83196953\allowbreak{}45554341\allowbreak{}34}
\RouteRow{\(\kappa_\Omega\)}{79.60661\allowbreak{}01397948\allowbreak{}27586470\allowbreak{}91}
\RouteRow{\(R_{\mathrm{act},D}\)}{1.000000\allowbreak{}05512887\allowbreak{}17997050\allowbreak{}72}
\RouteRow{\(R_{\mathrm{act},\Omega}\)}{1.000000\allowbreak{}05518981\allowbreak{}25318591\allowbreak{}40}
\RouteGroup{Separación de prefijos}
\RouteRow{Área de ambigüedad; área \(\log_2\)}{18; 0.0}
\RouteRow{Primer bloque de separación}{órbita: 1; celda: 1; ruta: 1}
\RouteRow{Ambigüedad final}{1}
\RouteGroup{Secuencias completas}
\RouteRow{\(w_6\)}{264 90 570 264 90 24 243 657 24 264 90 570 264 90 24 243 657 24}
\RouteRow{Tríadas estables por bloque}{0 1 2 3 4 5 6 7 8 9 10 11 12 13 14 15 15 16}
\end{tabularx}\endgroup\par\vspace{5pt}\noindent{\fontsize{8.4}{10}\selectfont\hyperref[trace-85]{Procedencia y códigos originales}\hfill Ficha completa · 085/324}
\clearpage\phantomsection\label{nav:vi-ruta-86}
% VI-CATALOG-ROW rutas 86
\pdfbookmark[3]{r3c4\_h-1v+1}{route-86}
\RouteTitle{Ruta 086}{r3c4\_\allowbreak{}h-1v+1}
\noindent Celda 22\quad (3, 4)\qquad Orientación \(h=-1,\ v=1\)\hfill\hyperref[sec:vi-indice-rutas]{Índice}\par\vspace{4pt}
\begingroup\fontsize{10.2}{12.5}\selectfont\setlength{\tabcolsep}{5pt}\renewcommand{\arraystretch}{1.1}\rowcolors{1}{routepale}{white}\noindent\begin{tabularx}{\linewidth}{@{}>{\raggedright\arraybackslash}p{.345\linewidth}>{\raggedright\arraybackslash}X@{}}
\RouteGroup{Identidad estructural}
\RouteRow{Clase y multisección}{neutrino\_\allowbreak{}G2\quad /\quad neutrino\_\allowbreak{}G2}
\RouteRow{Colección y sector}{neutrino; leptón}
\RouteRow{Clasificación física}{neutrino}
\RouteRow{Generación, profundidad y color}{2; G2; \textemdash{}}
\RouteRow{Ruta primaria y órbita de inversión}{\(B_1\): no\qquad 3,\allowbreak{}4;\allowbreak{}7,\allowbreak{}6}
\RouteRow{Firma de clasificación}{28|\allowbreak{}6|\allowbreak{}0|\allowbreak{}0|\allowbreak{}-12|\allowbreak{}+++-\kern0pt{}-\kern0pt{}-+++-\kern0pt{}-\kern0pt{}-}
\RouteGroup{Retornos, memoria y transporte}
\RouteRow{Retornos con fase}{clase: 27; órbita: 27; celda: 27}
\RouteRow{Retorno cinemático}{en 108: 54; extensión: 216}
\RouteRow{Giros y cambios de hoja}{71; 107}
\RouteRow{Acarreo y diferimiento}{deuda total: 2; final: 0; bloques diferidos: 1}
\RouteRow{Tríadas finales; persistencia de Pareto}{17; no}
\RouteRow{\(K\longrightarrow K+9\)}{repeticiones: 9; cambios de memoria: 9}
\RouteGroup{Firma, fase y diferencias}
\RouteRow{\(n_A,\ s_C,\ \kappa_0,\ \nu_{120},\ \nu_{270}\)}{50, 4, -2, 2, -8}
\RouteRow{\(\tau_{12}\); firma histórica \(\mathbb Z^5\)}{+++-0-+++-0-; 50|\allowbreak{}4|\allowbreak{}-2|\allowbreak{}2|\allowbreak{}-8}
\RouteRow{\(D_1,\ D_3,\ D_4,\ D_{27},\ D_{36}\)}{84, 24, 84, 24, 16}
\RouteRow{\(\Omega_{90/120},\ P_6\)}{48, 5}
\RouteGroup{Lectores y factores de acción}
\RouteRow{\(K_D\quad /\quad K_\Omega\)}{(40,\allowbreak{}84,\allowbreak{}30)\quad /\quad (48,\allowbreak{}54,\allowbreak{}5)}
\RouteRow{\(\kappa_D\)}{39.39035\allowbreak{}82768609\allowbreak{}24917849\allowbreak{}59}
\RouteRow{\(\kappa_\Omega\)}{47.60649\allowbreak{}89433721\allowbreak{}35446416\allowbreak{}86}
\RouteRow{\(R_{\mathrm{act},D}\)}{1.000000\allowbreak{}02730861\allowbreak{}73967087\allowbreak{}05}
\RouteRow{\(R_{\mathrm{act},\Omega}\)}{1.000000\allowbreak{}03300471\allowbreak{}80532430\allowbreak{}84}
\RouteGroup{Separación de prefijos}
\RouteRow{Área de ambigüedad; área \(\log_2\)}{18; 0.0}
\RouteRow{Primer bloque de separación}{órbita: 1; celda: 1; ruta: 1}
\RouteRow{Ambigüedad final}{1}
\RouteGroup{Secuencias completas}
\RouteRow{\(w_6\)}{258 414 420 258 414 255 252 333 255 258 414 420 258 414 255 252 333 255}
\RouteRow{Tríadas estables por bloque}{0 1 2 3 4 5 6 7 8 9 10 11 11 12 14 15 16 17}
\end{tabularx}\endgroup\par\vspace{5pt}\noindent{\fontsize{8.4}{10}\selectfont\hyperref[trace-86]{Procedencia y códigos originales}\hfill Ficha completa · 086/324}
\clearpage\phantomsection\label{nav:vi-ruta-87}
% VI-CATALOG-ROW rutas 87
\pdfbookmark[3]{r3c4\_h+1v-1}{route-87}
\RouteTitle{Ruta 087}{r3c4\_\allowbreak{}h+1v-1}
\noindent Celda 22\quad (3, 4)\qquad Orientación \(h=1,\ v=-1\)\hfill\hyperref[sec:vi-indice-rutas]{Índice}\par\vspace{4pt}
\begingroup\fontsize{10.2}{12.5}\selectfont\setlength{\tabcolsep}{5pt}\renewcommand{\arraystretch}{1.1}\rowcolors{1}{routepale}{white}\noindent\begin{tabularx}{\linewidth}{@{}>{\raggedright\arraybackslash}p{.345\linewidth}>{\raggedright\arraybackslash}X@{}}
\RouteGroup{Identidad estructural}
\RouteRow{Clase y multisección}{neutrino\_\allowbreak{}G2\quad /\quad neutrino\_\allowbreak{}G2}
\RouteRow{Colección y sector}{neutrino; leptón}
\RouteRow{Clasificación física}{neutrino}
\RouteRow{Generación, profundidad y color}{2; G2; \textemdash{}}
\RouteRow{Ruta primaria y órbita de inversión}{\(B_1\): no\qquad 3,\allowbreak{}4;\allowbreak{}7,\allowbreak{}6}
\RouteRow{Firma de clasificación}{28|\allowbreak{}6|\allowbreak{}0|\allowbreak{}0|\allowbreak{}-12|\allowbreak{}+++-\kern0pt{}-\kern0pt{}-+++-\kern0pt{}-\kern0pt{}-}
\RouteGroup{Retornos, memoria y transporte}
\RouteRow{Retornos con fase}{clase: 27; órbita: 27; celda: 27}
\RouteRow{Retorno cinemático}{en 108: 54; extensión: 216}
\RouteRow{Giros y cambios de hoja}{71; 107}
\RouteRow{Acarreo y diferimiento}{deuda total: 3; final: 1; bloques diferidos: 3}
\RouteRow{Tríadas finales; persistencia de Pareto}{16; no}
\RouteRow{\(K\longrightarrow K+9\)}{repeticiones: 9; cambios de memoria: 9}
\RouteGroup{Firma, fase y diferencias}
\RouteRow{\(n_A,\ s_C,\ \kappa_0,\ \nu_{120},\ \nu_{270}\)}{50, -4, 2, -2, -8}
\RouteRow{\(\tau_{12}\); firma histórica \(\mathbb Z^5\)}{+0+-\kern0pt{}-\kern0pt{}-+0+-\kern0pt{}-\kern0pt{}-; 50|\allowbreak{}-4|\allowbreak{}2|\allowbreak{}-2|\allowbreak{}-8}
\RouteRow{\(D_1,\ D_3,\ D_4,\ D_{27},\ D_{36}\)}{84, 24, 84, 24, 16}
\RouteRow{\(\Omega_{90/120},\ P_6\)}{48, 5}
\RouteGroup{Lectores y factores de acción}
\RouteRow{\(K_D\quad /\quad K_\Omega\)}{(40,\allowbreak{}84,\allowbreak{}30)\quad /\quad (48,\allowbreak{}54,\allowbreak{}5)}
\RouteRow{\(\kappa_D\)}{39.39035\allowbreak{}82768609\allowbreak{}24917849\allowbreak{}59}
\RouteRow{\(\kappa_\Omega\)}{47.60649\allowbreak{}89433721\allowbreak{}35446416\allowbreak{}86}
\RouteRow{\(R_{\mathrm{act},D}\)}{1.000000\allowbreak{}02730861\allowbreak{}73967087\allowbreak{}05}
\RouteRow{\(R_{\mathrm{act},\Omega}\)}{1.000000\allowbreak{}03300471\allowbreak{}80532430\allowbreak{}84}
\RouteGroup{Separación de prefijos}
\RouteRow{Área de ambigüedad; área \(\log_2\)}{18; 0.0}
\RouteRow{Primer bloque de separación}{órbita: 1; celda: 1; ruta: 1}
\RouteRow{Ambigüedad final}{1}
\RouteGroup{Secuencias completas}
\RouteRow{\(w_6\)}{252 333 255 252 333 420 258 414 420 252 333 255 252 333 420 258 414 420}
\RouteRow{Tríadas estables por bloque}{0 1 2 3 4 5 6 7 8 9 10 11 12 12 14 14 16 16}
\end{tabularx}\endgroup\par\vspace{5pt}\noindent{\fontsize{8.4}{10}\selectfont\hyperref[trace-87]{Procedencia y códigos originales}\hfill Ficha completa · 087/324}
\clearpage\phantomsection\label{nav:vi-ruta-88}
% VI-CATALOG-ROW rutas 88
\pdfbookmark[3]{r3c4\_h+1v+1}{route-88}
\RouteTitle{Ruta 088}{r3c4\_\allowbreak{}h+1v+1}
\noindent Celda 22\quad (3, 4)\qquad Orientación \(h=1,\ v=1\)\hfill\hyperref[sec:vi-indice-rutas]{Índice}\par\vspace{4pt}
\begingroup\fontsize{10.2}{12.5}\selectfont\setlength{\tabcolsep}{5pt}\renewcommand{\arraystretch}{1.1}\rowcolors{1}{routepale}{white}\noindent\begin{tabularx}{\linewidth}{@{}>{\raggedright\arraybackslash}p{.345\linewidth}>{\raggedright\arraybackslash}X@{}}
\RouteGroup{Identidad estructural}
\RouteRow{Clase y multisección}{neutrino\_\allowbreak{}G2\quad /\quad neutrino\_\allowbreak{}G2}
\RouteRow{Colección y sector}{neutrino; leptón}
\RouteRow{Clasificación física}{neutrino}
\RouteRow{Generación, profundidad y color}{2; G2; \textemdash{}}
\RouteRow{Ruta primaria y órbita de inversión}{\(B_1\): no\qquad 3,\allowbreak{}4;\allowbreak{}7,\allowbreak{}6}
\RouteRow{Firma de clasificación}{28|\allowbreak{}6|\allowbreak{}0|\allowbreak{}0|\allowbreak{}-12|\allowbreak{}+++-\kern0pt{}-\kern0pt{}-+++-\kern0pt{}-\kern0pt{}-}
\RouteGroup{Retornos, memoria y transporte}
\RouteRow{Retornos con fase}{clase: 27; órbita: 27; celda: 27}
\RouteRow{Retorno cinemático}{en 108: 54; extensión: 216}
\RouteRow{Giros y cambios de hoja}{71; 107}
\RouteRow{Acarreo y diferimiento}{deuda total: 0; final: 0; bloques diferidos: 0}
\RouteRow{Tríadas finales; persistencia de Pareto}{17; no}
\RouteRow{\(K\longrightarrow K+9\)}{repeticiones: 9; cambios de memoria: 9}
\RouteGroup{Firma, fase y diferencias}
\RouteRow{\(n_A,\ s_C,\ \kappa_0,\ \nu_{120},\ \nu_{270}\)}{28, 6, 0, 4, -12}
\RouteRow{\(\tau_{12}\); firma histórica \(\mathbb Z^5\)}{+++-\kern0pt{}-\kern0pt{}-+++-\kern0pt{}-\kern0pt{}-; 28|\allowbreak{}6|\allowbreak{}0|\allowbreak{}4|\allowbreak{}-12}
\RouteRow{\(D_1,\ D_3,\ D_4,\ D_{27},\ D_{36}\)}{66, 60, 66, 48, 32}
\RouteRow{\(\Omega_{90/120},\ P_6\)}{80, 6}
\RouteGroup{Lectores y factores de acción}
\RouteRow{\(K_D\quad /\quad K_\Omega\)}{(80,\allowbreak{}66,\allowbreak{}3)\quad /\quad (80,\allowbreak{}54,\allowbreak{}6)}
\RouteRow{\(\kappa_D\)}{79.51870\allowbreak{}83196953\allowbreak{}45554341\allowbreak{}34}
\RouteRow{\(\kappa_\Omega\)}{79.60661\allowbreak{}01397948\allowbreak{}27586470\allowbreak{}91}
\RouteRow{\(R_{\mathrm{act},D}\)}{1.000000\allowbreak{}05512887\allowbreak{}17997050\allowbreak{}72}
\RouteRow{\(R_{\mathrm{act},\Omega}\)}{1.000000\allowbreak{}05518981\allowbreak{}25318591\allowbreak{}40}
\RouteGroup{Separación de prefijos}
\RouteRow{Área de ambigüedad; área \(\log_2\)}{18; 0.0}
\RouteRow{Primer bloque de separación}{órbita: 1; celda: 1; ruta: 1}
\RouteRow{Ambigüedad final}{1}
\RouteGroup{Secuencias completas}
\RouteRow{\(w_6\)}{243 657 24 243 657 570 264 90 570 243 657 24 243 657 570 264 90 570}
\RouteRow{Tríadas estables por bloque}{0 1 2 3 4 5 6 7 8 9 10 11 12 13 14 15 16 17}
\end{tabularx}\endgroup\par\vspace{5pt}\noindent{\fontsize{8.4}{10}\selectfont\hyperref[trace-88]{Procedencia y códigos originales}\hfill Ficha completa · 088/324}
\clearpage\phantomsection\label{nav:vi-ruta-89}
% VI-CATALOG-ROW rutas 89
\pdfbookmark[2]{Celda 23}{cell-23}
\pdfbookmark[3]{r3c5\_h-1v-1}{route-89}
\RouteTitle{Ruta 089}{r3c5\_\allowbreak{}h-1v-1}
\noindent Celda 23\quad (3, 5)\qquad Orientación \(h=-1,\ v=-1\)\hfill\hyperref[sec:vi-indice-rutas]{Índice}\par\vspace{4pt}
\begingroup\fontsize{10.2}{12.5}\selectfont\setlength{\tabcolsep}{5pt}\renewcommand{\arraystretch}{1.1}\rowcolors{1}{routepale}{white}\noindent\begin{tabularx}{\linewidth}{@{}>{\raggedright\arraybackslash}p{.345\linewidth}>{\raggedright\arraybackslash}X@{}}
\RouteGroup{Identidad estructural}
\RouteRow{Clase y multisección}{charged\_\allowbreak{}lepton\_\allowbreak{}G2\quad /\quad charged\_\allowbreak{}lepton\_\allowbreak{}G2}
\RouteRow{Colección y sector}{leptón cargado; leptón}
\RouteRow{Clasificación física}{leptón cargado}
\RouteRow{Generación, profundidad y color}{2; G2; \textemdash{}}
\RouteRow{Ruta primaria y órbita de inversión}{\(B_1\): no\qquad 3,\allowbreak{}5;\allowbreak{}7,\allowbreak{}5}
\RouteRow{Firma de clasificación}{28|\allowbreak{}-6|\allowbreak{}6|\allowbreak{}0|\allowbreak{}-8|\allowbreak{}+0+-\kern0pt{}-\kern0pt{}-+0+-\kern0pt{}-\kern0pt{}-}
\RouteGroup{Retornos, memoria y transporte}
\RouteRow{Retornos con fase}{clase: 27; órbita: 27; celda: 27}
\RouteRow{Retorno cinemático}{en 108: 54; extensión: 216}
\RouteRow{Giros y cambios de hoja}{71; 107}
\RouteRow{Acarreo y diferimiento}{deuda total: 1; final: 0; bloques diferidos: 1}
\RouteRow{Tríadas finales; persistencia de Pareto}{17; sí}
\RouteRow{\(K\longrightarrow K+9\)}{repeticiones: 9; cambios de memoria: 9}
\RouteGroup{Firma, fase y diferencias}
\RouteRow{\(n_A,\ s_C,\ \kappa_0,\ \nu_{120},\ \nu_{270}\)}{28, 6, -2, 0, -8}
\RouteRow{\(\tau_{12}\); firma histórica \(\mathbb Z^5\)}{+++-0-+++-0-; 28|\allowbreak{}6|\allowbreak{}-2|\allowbreak{}0|\allowbreak{}-8}
\RouteRow{\(D_1,\ D_3,\ D_4,\ D_{27},\ D_{36}\)}{66, 60, 66, 48, 32}
\RouteRow{\(\Omega_{90/120},\ P_6\)}{48, 5}
\RouteGroup{Lectores y factores de acción}
\RouteRow{\(K_D\quad /\quad K_\Omega\)}{(80,\allowbreak{}66,\allowbreak{}3)\quad /\quad (48,\allowbreak{}54,\allowbreak{}5)}
\RouteRow{\(\kappa_D\)}{79.51870\allowbreak{}83196953\allowbreak{}45554341\allowbreak{}34}
\RouteRow{\(\kappa_\Omega\)}{47.60649\allowbreak{}89433721\allowbreak{}35446416\allowbreak{}86}
\RouteRow{\(R_{\mathrm{act},D}\)}{1.000000\allowbreak{}05512887\allowbreak{}17997050\allowbreak{}72}
\RouteRow{\(R_{\mathrm{act},\Omega}\)}{1.000000\allowbreak{}03300471\allowbreak{}80532430\allowbreak{}84}
\RouteGroup{Separación de prefijos}
\RouteRow{Área de ambigüedad; área \(\log_2\)}{22; 4.0}
\RouteRow{Primer bloque de separación}{órbita: 5; celda: 5; ruta: 5}
\RouteRow{Ambigüedad final}{1}
\RouteGroup{Secuencias completas}
\RouteRow{\(w_6\)}{486 423 15 486 423 327 498 99 327 486 423 15 486 423 327 498 99 327}
\RouteRow{Tríadas estables por bloque}{0 1 2 3 4 5 6 7 8 9 10 10 12 13 14 15 16 17}
\end{tabularx}\endgroup\par\vspace{5pt}\noindent{\fontsize{8.4}{10}\selectfont\hyperref[trace-89]{Procedencia y códigos originales}\hfill Ficha completa · 089/324}
\clearpage\phantomsection\label{nav:vi-ruta-90}
% VI-CATALOG-ROW rutas 90
\pdfbookmark[3]{r3c5\_h-1v+1}{route-90}
\RouteTitle{Ruta 090}{r3c5\_\allowbreak{}h-1v+1}
\noindent Celda 23\quad (3, 5)\qquad Orientación \(h=-1,\ v=1\)\hfill\hyperref[sec:vi-indice-rutas]{Índice}\par\vspace{4pt}
\begingroup\fontsize{10.2}{12.5}\selectfont\setlength{\tabcolsep}{5pt}\renewcommand{\arraystretch}{1.1}\rowcolors{1}{routepale}{white}\noindent\begin{tabularx}{\linewidth}{@{}>{\raggedright\arraybackslash}p{.345\linewidth}>{\raggedright\arraybackslash}X@{}}
\RouteGroup{Identidad estructural}
\RouteRow{Clase y multisección}{charged\_\allowbreak{}lepton\_\allowbreak{}G2\quad /\quad charged\_\allowbreak{}lepton\_\allowbreak{}G2}
\RouteRow{Colección y sector}{leptón cargado; leptón}
\RouteRow{Clasificación física}{leptón cargado}
\RouteRow{Generación, profundidad y color}{2; G2; \textemdash{}}
\RouteRow{Ruta primaria y órbita de inversión}{\(B_1\): sí\qquad 3,\allowbreak{}5;\allowbreak{}7,\allowbreak{}5}
\RouteRow{Firma de clasificación}{28|\allowbreak{}-6|\allowbreak{}6|\allowbreak{}0|\allowbreak{}-8|\allowbreak{}+0+-\kern0pt{}-\kern0pt{}-+0+-\kern0pt{}-\kern0pt{}-}
\RouteGroup{Retornos, memoria y transporte}
\RouteRow{Retornos con fase}{clase: 27; órbita: 27; celda: 27}
\RouteRow{Retorno cinemático}{en 108: 54; extensión: 216}
\RouteRow{Giros y cambios de hoja}{71; 107}
\RouteRow{Acarreo y diferimiento}{deuda total: 1; final: 0; bloques diferidos: 1}
\RouteRow{Tríadas finales; persistencia de Pareto}{17; sí}
\RouteRow{\(K\longrightarrow K+9\)}{repeticiones: 9; cambios de memoria: 9}
\RouteGroup{Firma, fase y diferencias}
\RouteRow{\(n_A,\ s_C,\ \kappa_0,\ \nu_{120},\ \nu_{270}\)}{14, -4, 0, 0, -12}
\RouteRow{\(\tau_{12}\); firma histórica \(\mathbb Z^5\)}{+++-\kern0pt{}-\kern0pt{}-+++-\kern0pt{}-\kern0pt{}-; 14|\allowbreak{}-4|\allowbreak{}0|\allowbreak{}0|\allowbreak{}-12}
\RouteRow{\(D_1,\ D_3,\ D_4,\ D_{27},\ D_{36}\)}{78, 24, 78, 24, 16}
\RouteRow{\(\Omega_{90/120},\ P_6\)}{80, 6}
\RouteGroup{Lectores y factores de acción}
\RouteRow{\(K_D\quad /\quad K_\Omega\)}{(40,\allowbreak{}78,\allowbreak{}27)\quad /\quad (80,\allowbreak{}54,\allowbreak{}6)}
\RouteRow{\(\kappa_D\)}{39.43380\allowbreak{}88030085\allowbreak{}51303671\allowbreak{}14}
\RouteRow{\(\kappa_\Omega\)}{79.60661\allowbreak{}01397948\allowbreak{}27586470\allowbreak{}91}
\RouteRow{\(R_{\mathrm{act},D}\)}{1.000000\allowbreak{}02733874\allowbreak{}08548943\allowbreak{}58}
\RouteRow{\(R_{\mathrm{act},\Omega}\)}{1.000000\allowbreak{}05518981\allowbreak{}25318591\allowbreak{}40}
\RouteGroup{Separación de prefijos}
\RouteRow{Área de ambigüedad; área \(\log_2\)}{18; 0.0}
\RouteRow{Primer bloque de separación}{órbita: 1; celda: 1; ruta: 1}
\RouteRow{Ambigüedad final}{1}
\RouteGroup{Secuencias completas}
\RouteRow{\(w_6\)}{504 585 507 504 585 672 510 666 672 504 585 507 504 585 672 510 666 672}
\RouteRow{Tríadas estables por bloque}{0 1 2 3 4 5 6 7 8 9 10 11 12 13 14 14 16 17}
\end{tabularx}\endgroup\par\vspace{5pt}\noindent{\fontsize{8.4}{10}\selectfont\hyperref[trace-90]{Procedencia y códigos originales}\hfill Ficha completa · 090/324}
\clearpage\phantomsection\label{nav:vi-ruta-91}
% VI-CATALOG-ROW rutas 91
\pdfbookmark[3]{r3c5\_h+1v-1}{route-91}
\RouteTitle{Ruta 091}{r3c5\_\allowbreak{}h+1v-1}
\noindent Celda 23\quad (3, 5)\qquad Orientación \(h=1,\ v=-1\)\hfill\hyperref[sec:vi-indice-rutas]{Índice}\par\vspace{4pt}
\begingroup\fontsize{10.2}{12.5}\selectfont\setlength{\tabcolsep}{5pt}\renewcommand{\arraystretch}{1.1}\rowcolors{1}{routepale}{white}\noindent\begin{tabularx}{\linewidth}{@{}>{\raggedright\arraybackslash}p{.345\linewidth}>{\raggedright\arraybackslash}X@{}}
\RouteGroup{Identidad estructural}
\RouteRow{Clase y multisección}{charged\_\allowbreak{}lepton\_\allowbreak{}G2\quad /\quad charged\_\allowbreak{}lepton\_\allowbreak{}G2}
\RouteRow{Colección y sector}{leptón cargado; leptón}
\RouteRow{Clasificación física}{leptón cargado}
\RouteRow{Generación, profundidad y color}{2; G2; \textemdash{}}
\RouteRow{Ruta primaria y órbita de inversión}{\(B_1\): no\qquad 3,\allowbreak{}5;\allowbreak{}7,\allowbreak{}5}
\RouteRow{Firma de clasificación}{28|\allowbreak{}-6|\allowbreak{}6|\allowbreak{}0|\allowbreak{}-8|\allowbreak{}+0+-\kern0pt{}-\kern0pt{}-+0+-\kern0pt{}-\kern0pt{}-}
\RouteGroup{Retornos, memoria y transporte}
\RouteRow{Retornos con fase}{clase: 27; órbita: 27; celda: 27}
\RouteRow{Retorno cinemático}{en 108: 54; extensión: 216}
\RouteRow{Giros y cambios de hoja}{71; 107}
\RouteRow{Acarreo y diferimiento}{deuda total: 1; final: 0; bloques diferidos: 1}
\RouteRow{Tríadas finales; persistencia de Pareto}{17; sí}
\RouteRow{\(K\longrightarrow K+9\)}{repeticiones: 9; cambios de memoria: 9}
\RouteGroup{Firma, fase y diferencias}
\RouteRow{\(n_A,\ s_C,\ \kappa_0,\ \nu_{120},\ \nu_{270}\)}{14, 4, 0, 0, -12}
\RouteRow{\(\tau_{12}\); firma histórica \(\mathbb Z^5\)}{+++-\kern0pt{}-\kern0pt{}-+++-\kern0pt{}-\kern0pt{}-; 14|\allowbreak{}4|\allowbreak{}0|\allowbreak{}0|\allowbreak{}-12}
\RouteRow{\(D_1,\ D_3,\ D_4,\ D_{27},\ D_{36}\)}{78, 24, 78, 24, 16}
\RouteRow{\(\Omega_{90/120},\ P_6\)}{80, 6}
\RouteGroup{Lectores y factores de acción}
\RouteRow{\(K_D\quad /\quad K_\Omega\)}{(40,\allowbreak{}78,\allowbreak{}27)\quad /\quad (80,\allowbreak{}54,\allowbreak{}6)}
\RouteRow{\(\kappa_D\)}{39.43380\allowbreak{}88030085\allowbreak{}51303671\allowbreak{}14}
\RouteRow{\(\kappa_\Omega\)}{79.60661\allowbreak{}01397948\allowbreak{}27586470\allowbreak{}91}
\RouteRow{\(R_{\mathrm{act},D}\)}{1.000000\allowbreak{}02733874\allowbreak{}08548943\allowbreak{}58}
\RouteRow{\(R_{\mathrm{act},\Omega}\)}{1.000000\allowbreak{}05518981\allowbreak{}25318591\allowbreak{}40}
\RouteGroup{Separación de prefijos}
\RouteRow{Área de ambigüedad; área \(\log_2\)}{18; 0.0}
\RouteRow{Primer bloque de separación}{órbita: 1; celda: 1; ruta: 1}
\RouteRow{Ambigüedad final}{1}
\RouteGroup{Secuencias completas}
\RouteRow{\(w_6\)}{510 666 672 510 666 507 504 585 507 510 666 672 510 666 507 504 585 507}
\RouteRow{Tríadas estables por bloque}{0 1 2 3 4 5 6 7 8 9 10 11 12 13 14 15 15 17}
\end{tabularx}\endgroup\par\vspace{5pt}\noindent{\fontsize{8.4}{10}\selectfont\hyperref[trace-91]{Procedencia y códigos originales}\hfill Ficha completa · 091/324}
\clearpage\phantomsection\label{nav:vi-ruta-92}
% VI-CATALOG-ROW rutas 92
\pdfbookmark[3]{r3c5\_h+1v+1}{route-92}
\RouteTitle{Ruta 092}{r3c5\_\allowbreak{}h+1v+1}
\noindent Celda 23\quad (3, 5)\qquad Orientación \(h=1,\ v=1\)\hfill\hyperref[sec:vi-indice-rutas]{Índice}\par\vspace{4pt}
\begingroup\fontsize{10.2}{12.5}\selectfont\setlength{\tabcolsep}{5pt}\renewcommand{\arraystretch}{1.1}\rowcolors{1}{routepale}{white}\noindent\begin{tabularx}{\linewidth}{@{}>{\raggedright\arraybackslash}p{.345\linewidth}>{\raggedright\arraybackslash}X@{}}
\RouteGroup{Identidad estructural}
\RouteRow{Clase y multisección}{charged\_\allowbreak{}lepton\_\allowbreak{}G2\quad /\quad charged\_\allowbreak{}lepton\_\allowbreak{}G2}
\RouteRow{Colección y sector}{leptón cargado; leptón}
\RouteRow{Clasificación física}{leptón cargado}
\RouteRow{Generación, profundidad y color}{2; G2; \textemdash{}}
\RouteRow{Ruta primaria y órbita de inversión}{\(B_1\): no\qquad 3,\allowbreak{}5;\allowbreak{}7,\allowbreak{}5}
\RouteRow{Firma de clasificación}{28|\allowbreak{}-6|\allowbreak{}6|\allowbreak{}0|\allowbreak{}-8|\allowbreak{}+0+-\kern0pt{}-\kern0pt{}-+0+-\kern0pt{}-\kern0pt{}-}
\RouteGroup{Retornos, memoria y transporte}
\RouteRow{Retornos con fase}{clase: 27; órbita: 27; celda: 27}
\RouteRow{Retorno cinemático}{en 108: 54; extensión: 216}
\RouteRow{Giros y cambios de hoja}{71; 107}
\RouteRow{Acarreo y diferimiento}{deuda total: 0; final: 0; bloques diferidos: 0}
\RouteRow{Tríadas finales; persistencia de Pareto}{17; sí}
\RouteRow{\(K\longrightarrow K+9\)}{repeticiones: 9; cambios de memoria: 9}
\RouteGroup{Firma, fase y diferencias}
\RouteRow{\(n_A,\ s_C,\ \kappa_0,\ \nu_{120},\ \nu_{270}\)}{28, -6, 2, 0, -8}
\RouteRow{\(\tau_{12}\); firma histórica \(\mathbb Z^5\)}{+0+-\kern0pt{}-\kern0pt{}-+0+-\kern0pt{}-\kern0pt{}-; 28|\allowbreak{}-6|\allowbreak{}2|\allowbreak{}0|\allowbreak{}-8}
\RouteRow{\(D_1,\ D_3,\ D_4,\ D_{27},\ D_{36}\)}{66, 60, 66, 48, 32}
\RouteRow{\(\Omega_{90/120},\ P_6\)}{48, 5}
\RouteGroup{Lectores y factores de acción}
\RouteRow{\(K_D\quad /\quad K_\Omega\)}{(80,\allowbreak{}66,\allowbreak{}3)\quad /\quad (48,\allowbreak{}54,\allowbreak{}5)}
\RouteRow{\(\kappa_D\)}{79.51870\allowbreak{}83196953\allowbreak{}45554341\allowbreak{}34}
\RouteRow{\(\kappa_\Omega\)}{47.60649\allowbreak{}89433721\allowbreak{}35446416\allowbreak{}86}
\RouteRow{\(R_{\mathrm{act},D}\)}{1.000000\allowbreak{}05512887\allowbreak{}17997050\allowbreak{}72}
\RouteRow{\(R_{\mathrm{act},\Omega}\)}{1.000000\allowbreak{}03300471\allowbreak{}80532430\allowbreak{}84}
\RouteGroup{Separación de prefijos}
\RouteRow{Área de ambigüedad; área \(\log_2\)}{18; 0.0}
\RouteRow{Primer bloque de separación}{órbita: 1; celda: 1; ruta: 1}
\RouteRow{Ambigüedad final}{1}
\RouteGroup{Secuencias completas}
\RouteRow{\(w_6\)}{498 99 327 498 99 15 486 423 15 498 99 327 498 99 15 486 423 15}
\RouteRow{Tríadas estables por bloque}{0 1 2 3 4 5 6 7 8 9 10 11 12 13 14 15 16 17}
\end{tabularx}\endgroup\par\vspace{5pt}\noindent{\fontsize{8.4}{10}\selectfont\hyperref[trace-92]{Procedencia y códigos originales}\hfill Ficha completa · 092/324}
\clearpage\phantomsection\label{nav:vi-ruta-93}
% VI-CATALOG-ROW rutas 93
\pdfbookmark[2]{Celda 24}{cell-24}
\pdfbookmark[3]{r3c6\_h-1v-1}{route-93}
\RouteTitle{Ruta 093}{r3c6\_\allowbreak{}h-1v-1}
\noindent Celda 24\quad (3, 6)\qquad Orientación \(h=-1,\ v=-1\)\hfill\hyperref[sec:vi-indice-rutas]{Índice}\par\vspace{4pt}
\begingroup\fontsize{10.2}{12.5}\selectfont\setlength{\tabcolsep}{5pt}\renewcommand{\arraystretch}{1.1}\rowcolors{1}{routepale}{white}\noindent\begin{tabularx}{\linewidth}{@{}>{\raggedright\arraybackslash}p{.345\linewidth}>{\raggedright\arraybackslash}X@{}}
\RouteGroup{Identidad estructural}
\RouteRow{Clase y multisección}{neutrino\_\allowbreak{}G2\quad /\quad neutrino\_\allowbreak{}G2}
\RouteRow{Colección y sector}{neutrino; leptón}
\RouteRow{Clasificación física}{neutrino}
\RouteRow{Generación, profundidad y color}{2; G2; \textemdash{}}
\RouteRow{Ruta primaria y órbita de inversión}{\(B_1\): no\qquad 3,\allowbreak{}6;\allowbreak{}7,\allowbreak{}4}
\RouteRow{Firma de clasificación}{30|\allowbreak{}-6|\allowbreak{}0|\allowbreak{}-2|\allowbreak{}-8|\allowbreak{}++0-\kern0pt{}-\kern0pt{}-++0-\kern0pt{}-\kern0pt{}-}
\RouteGroup{Retornos, memoria y transporte}
\RouteRow{Retornos con fase}{clase: 27; órbita: 27; celda: 27}
\RouteRow{Retorno cinemático}{en 108: 54; extensión: 216}
\RouteRow{Giros y cambios de hoja}{71; 107}
\RouteRow{Acarreo y diferimiento}{deuda total: 1; final: 0; bloques diferidos: 1}
\RouteRow{Tríadas finales; persistencia de Pareto}{17; sí}
\RouteRow{\(K\longrightarrow K+9\)}{repeticiones: 9; cambios de memoria: 9}
\RouteGroup{Firma, fase y diferencias}
\RouteRow{\(n_A,\ s_C,\ \kappa_0,\ \nu_{120},\ \nu_{270}\)}{30, 6, -2, 2, -8}
\RouteRow{\(\tau_{12}\); firma histórica \(\mathbb Z^5\)}{+++-\kern0pt{}-0+++-\kern0pt{}-0; 30|\allowbreak{}6|\allowbreak{}-2|\allowbreak{}2|\allowbreak{}-8}
\RouteRow{\(D_1,\ D_3,\ D_4,\ D_{27},\ D_{36}\)}{54, 60, 72, 24, 16}
\RouteRow{\(\Omega_{90/120},\ P_6\)}{56, 5}
\RouteGroup{Lectores y factores de acción}
\RouteRow{\(K_D\quad /\quad K_\Omega\)}{(40,\allowbreak{}54,\allowbreak{}6)\quad /\quad (56,\allowbreak{}54,\allowbreak{}5)}
\RouteRow{\(\kappa_D\)}{39.60661\allowbreak{}01397948\allowbreak{}27586470\allowbreak{}91}
\RouteRow{\(\kappa_\Omega\)}{55.60649\allowbreak{}89433721\allowbreak{}35446416\allowbreak{}86}
\RouteRow{\(R_{\mathrm{act},D}\)}{1.000000\allowbreak{}02745854\allowbreak{}08735595\allowbreak{}18}
\RouteRow{\(R_{\mathrm{act},\Omega}\)}{1.000000\allowbreak{}03855097\allowbreak{}23541416\allowbreak{}45}
\RouteGroup{Separación de prefijos}
\RouteRow{Área de ambigüedad; área \(\log_2\)}{18; 0.0}
\RouteRow{Primer bloque de separación}{órbita: 1; celda: 1; ruta: 1}
\RouteRow{Ambigüedad final}{1}
\RouteGroup{Secuencias completas}
\RouteRow{\(w_6\)}{15 486 423 15 486 657 24 243 657 15 486 423 15 486 657 24 243 657}
\RouteRow{Tríadas estables por bloque}{0 1 2 3 4 5 6 7 8 9 10 11 12 13 14 15 15 17}
\end{tabularx}\endgroup\par\vspace{5pt}\noindent{\fontsize{8.4}{10}\selectfont\hyperref[trace-93]{Procedencia y códigos originales}\hfill Ficha completa · 093/324}
\clearpage\phantomsection\label{nav:vi-ruta-94}
% VI-CATALOG-ROW rutas 94
\pdfbookmark[3]{r3c6\_h-1v+1}{route-94}
\RouteTitle{Ruta 094}{r3c6\_\allowbreak{}h-1v+1}
\noindent Celda 24\quad (3, 6)\qquad Orientación \(h=-1,\ v=1\)\hfill\hyperref[sec:vi-indice-rutas]{Índice}\par\vspace{4pt}
\begingroup\fontsize{10.2}{12.5}\selectfont\setlength{\tabcolsep}{5pt}\renewcommand{\arraystretch}{1.1}\rowcolors{1}{routepale}{white}\noindent\begin{tabularx}{\linewidth}{@{}>{\raggedright\arraybackslash}p{.345\linewidth}>{\raggedright\arraybackslash}X@{}}
\RouteGroup{Identidad estructural}
\RouteRow{Clase y multisección}{neutrino\_\allowbreak{}G2\quad /\quad neutrino\_\allowbreak{}G2}
\RouteRow{Colección y sector}{neutrino; leptón}
\RouteRow{Clasificación física}{neutrino}
\RouteRow{Generación, profundidad y color}{2; G2; \textemdash{}}
\RouteRow{Ruta primaria y órbita de inversión}{\(B_1\): sí\qquad 3,\allowbreak{}6;\allowbreak{}7,\allowbreak{}4}
\RouteRow{Firma de clasificación}{30|\allowbreak{}-6|\allowbreak{}0|\allowbreak{}-2|\allowbreak{}-8|\allowbreak{}++0-\kern0pt{}-\kern0pt{}-++0-\kern0pt{}-\kern0pt{}-}
\RouteGroup{Retornos, memoria y transporte}
\RouteRow{Retornos con fase}{clase: 27; órbita: 27; celda: 27}
\RouteRow{Retorno cinemático}{en 108: 54; extensión: 216}
\RouteRow{Giros y cambios de hoja}{71; 107}
\RouteRow{Acarreo y diferimiento}{deuda total: 0; final: 0; bloques diferidos: 0}
\RouteRow{Tríadas finales; persistencia de Pareto}{17; sí}
\RouteRow{\(K\longrightarrow K+9\)}{repeticiones: 9; cambios de memoria: 9}
\RouteGroup{Firma, fase y diferencias}
\RouteRow{\(n_A,\ s_C,\ \kappa_0,\ \nu_{120},\ \nu_{270}\)}{20, 0, 0, 0, 0}
\RouteRow{\(\tau_{12}\); firma histórica \(\mathbb Z^5\)}{+0000-+0000-; 20|\allowbreak{}0|\allowbreak{}0|\allowbreak{}0|\allowbreak{}0}
\RouteRow{\(D_1,\ D_3,\ D_4,\ D_{27},\ D_{36}\)}{24, 24, 24, 0, 0}
\RouteRow{\(\Omega_{90/120},\ P_6\)}{4, 2}
\RouteGroup{Lectores y factores de acción}
\RouteRow{\(K_D\quad /\quad K_\Omega\)}{(0,\allowbreak{}24,\allowbreak{}0)\quad /\quad (4,\allowbreak{}54,\allowbreak{}2)}
\RouteRow{\(\kappa_D\)}{-0.175136\allowbreak{}46166281\allowbreak{}12239348\allowbreak{}425}
\RouteRow{\(\kappa_\Omega\)}{3.606165\allowbreak{}35410405\allowbreak{}90262547\allowbreak{}07}
\RouteRow{\(R_{\mathrm{act},D}\)}{0.999999\allowbreak{}99987858\allowbreak{}10851337\allowbreak{}881}
\RouteRow{\(R_{\mathrm{act},\Omega}\)}{1.000000\allowbreak{}00250008\allowbreak{}86767963\allowbreak{}14}
\RouteGroup{Separación de prefijos}
\RouteRow{Área de ambigüedad; área \(\log_2\)}{36; 18.0}
\RouteRow{Primer bloque de separación}{órbita: 1; celda: 1; ruta: \textemdash{}}
\RouteRow{Ambigüedad final}{2}
\RouteGroup{Secuencias completas}
\RouteRow{\(w_6\)}{3 0 81 3 0 81 3 0 81 3 0 81 3 0 81 3 0 81}
\RouteRow{Tríadas estables por bloque}{0 1 2 3 4 5 6 7 8 9 10 11 12 13 14 15 16 17}
\end{tabularx}\endgroup\par\vspace{5pt}\noindent{\fontsize{8.4}{10}\selectfont\hyperref[trace-94]{Procedencia y códigos originales}\hfill Ficha completa · 094/324}
\clearpage\phantomsection\label{nav:vi-ruta-95}
% VI-CATALOG-ROW rutas 95
\pdfbookmark[3]{r3c6\_h+1v-1}{route-95}
\RouteTitle{Ruta 095}{r3c6\_\allowbreak{}h+1v-1}
\noindent Celda 24\quad (3, 6)\qquad Orientación \(h=1,\ v=-1\)\hfill\hyperref[sec:vi-indice-rutas]{Índice}\par\vspace{4pt}
\begingroup\fontsize{10.2}{12.5}\selectfont\setlength{\tabcolsep}{5pt}\renewcommand{\arraystretch}{1.1}\rowcolors{1}{routepale}{white}\noindent\begin{tabularx}{\linewidth}{@{}>{\raggedright\arraybackslash}p{.345\linewidth}>{\raggedright\arraybackslash}X@{}}
\RouteGroup{Identidad estructural}
\RouteRow{Clase y multisección}{neutrino\_\allowbreak{}G2\quad /\quad neutrino\_\allowbreak{}G2}
\RouteRow{Colección y sector}{neutrino; leptón}
\RouteRow{Clasificación física}{neutrino}
\RouteRow{Generación, profundidad y color}{2; G2; \textemdash{}}
\RouteRow{Ruta primaria y órbita de inversión}{\(B_1\): no\qquad 3,\allowbreak{}6;\allowbreak{}7,\allowbreak{}4}
\RouteRow{Firma de clasificación}{30|\allowbreak{}-6|\allowbreak{}0|\allowbreak{}-2|\allowbreak{}-8|\allowbreak{}++0-\kern0pt{}-\kern0pt{}-++0-\kern0pt{}-\kern0pt{}-}
\RouteGroup{Retornos, memoria y transporte}
\RouteRow{Retornos con fase}{clase: 27; órbita: 27; celda: 27}
\RouteRow{Retorno cinemático}{en 108: 54; extensión: 216}
\RouteRow{Giros y cambios de hoja}{71; 107}
\RouteRow{Acarreo y diferimiento}{deuda total: 0; final: 0; bloques diferidos: 0}
\RouteRow{Tríadas finales; persistencia de Pareto}{17; sí}
\RouteRow{\(K\longrightarrow K+9\)}{repeticiones: 9; cambios de memoria: 9}
\RouteGroup{Firma, fase y diferencias}
\RouteRow{\(n_A,\ s_C,\ \kappa_0,\ \nu_{120},\ \nu_{270}\)}{20, 0, 0, 0, 0}
\RouteRow{\(\tau_{12}\); firma histórica \(\mathbb Z^5\)}{00+-0000+-00; 20|\allowbreak{}0|\allowbreak{}0|\allowbreak{}0|\allowbreak{}0}
\RouteRow{\(D_1,\ D_3,\ D_4,\ D_{27},\ D_{36}\)}{24, 24, 24, 0, 0}
\RouteRow{\(\Omega_{90/120},\ P_6\)}{4, 2}
\RouteGroup{Lectores y factores de acción}
\RouteRow{\(K_D\quad /\quad K_\Omega\)}{(0,\allowbreak{}24,\allowbreak{}0)\quad /\quad (4,\allowbreak{}54,\allowbreak{}2)}
\RouteRow{\(\kappa_D\)}{-0.175136\allowbreak{}46166281\allowbreak{}12239348\allowbreak{}425}
\RouteRow{\(\kappa_\Omega\)}{3.606165\allowbreak{}35410405\allowbreak{}90262547\allowbreak{}07}
\RouteRow{\(R_{\mathrm{act},D}\)}{0.999999\allowbreak{}99987858\allowbreak{}10851337\allowbreak{}881}
\RouteRow{\(R_{\mathrm{act},\Omega}\)}{1.000000\allowbreak{}00250008\allowbreak{}86767963\allowbreak{}14}
\RouteGroup{Separación de prefijos}
\RouteRow{Área de ambigüedad; área \(\log_2\)}{36; 18.0}
\RouteRow{Primer bloque de separación}{órbita: 1; celda: 1; ruta: \textemdash{}}
\RouteRow{Ambigüedad final}{2}
\RouteGroup{Secuencias completas}
\RouteRow{\(w_6\)}{3 0 81 3 0 81 3 0 81 3 0 81 3 0 81 3 0 81}
\RouteRow{Tríadas estables por bloque}{0 1 2 3 4 5 6 7 8 9 10 11 12 13 14 15 16 17}
\end{tabularx}\endgroup\par\vspace{5pt}\noindent{\fontsize{8.4}{10}\selectfont\hyperref[trace-95]{Procedencia y códigos originales}\hfill Ficha completa · 095/324}
\clearpage\phantomsection\label{nav:vi-ruta-96}
% VI-CATALOG-ROW rutas 96
\pdfbookmark[3]{r3c6\_h+1v+1}{route-96}
\RouteTitle{Ruta 096}{r3c6\_\allowbreak{}h+1v+1}
\noindent Celda 24\quad (3, 6)\qquad Orientación \(h=1,\ v=1\)\hfill\hyperref[sec:vi-indice-rutas]{Índice}\par\vspace{4pt}
\begingroup\fontsize{10.2}{12.5}\selectfont\setlength{\tabcolsep}{5pt}\renewcommand{\arraystretch}{1.1}\rowcolors{1}{routepale}{white}\noindent\begin{tabularx}{\linewidth}{@{}>{\raggedright\arraybackslash}p{.345\linewidth}>{\raggedright\arraybackslash}X@{}}
\RouteGroup{Identidad estructural}
\RouteRow{Clase y multisección}{neutrino\_\allowbreak{}G2\quad /\quad neutrino\_\allowbreak{}G2}
\RouteRow{Colección y sector}{neutrino; leptón}
\RouteRow{Clasificación física}{neutrino}
\RouteRow{Generación, profundidad y color}{2; G2; \textemdash{}}
\RouteRow{Ruta primaria y órbita de inversión}{\(B_1\): no\qquad 3,\allowbreak{}6;\allowbreak{}7,\allowbreak{}4}
\RouteRow{Firma de clasificación}{30|\allowbreak{}-6|\allowbreak{}0|\allowbreak{}-2|\allowbreak{}-8|\allowbreak{}++0-\kern0pt{}-\kern0pt{}-++0-\kern0pt{}-\kern0pt{}-}
\RouteGroup{Retornos, memoria y transporte}
\RouteRow{Retornos con fase}{clase: 27; órbita: 27; celda: 27}
\RouteRow{Retorno cinemático}{en 108: 54; extensión: 216}
\RouteRow{Giros y cambios de hoja}{71; 107}
\RouteRow{Acarreo y diferimiento}{deuda total: 2; final: 1; bloques diferidos: 2}
\RouteRow{Tríadas finales; persistencia de Pareto}{16; sí}
\RouteRow{\(K\longrightarrow K+9\)}{repeticiones: 9; cambios de memoria: 9}
\RouteGroup{Firma, fase y diferencias}
\RouteRow{\(n_A,\ s_C,\ \kappa_0,\ \nu_{120},\ \nu_{270}\)}{30, -6, 2, -2, -8}
\RouteRow{\(\tau_{12}\); firma histórica \(\mathbb Z^5\)}{++0-\kern0pt{}-\kern0pt{}-++0-\kern0pt{}-\kern0pt{}-; 30|\allowbreak{}-6|\allowbreak{}2|\allowbreak{}-2|\allowbreak{}-8}
\RouteRow{\(D_1,\ D_3,\ D_4,\ D_{27},\ D_{36}\)}{54, 60, 72, 24, 16}
\RouteRow{\(\Omega_{90/120},\ P_6\)}{56, 5}
\RouteGroup{Lectores y factores de acción}
\RouteRow{\(K_D\quad /\quad K_\Omega\)}{(40,\allowbreak{}54,\allowbreak{}6)\quad /\quad (56,\allowbreak{}54,\allowbreak{}5)}
\RouteRow{\(\kappa_D\)}{39.60661\allowbreak{}01397948\allowbreak{}27586470\allowbreak{}91}
\RouteRow{\(\kappa_\Omega\)}{55.60649\allowbreak{}89433721\allowbreak{}35446416\allowbreak{}86}
\RouteRow{\(R_{\mathrm{act},D}\)}{1.000000\allowbreak{}02745854\allowbreak{}08735595\allowbreak{}18}
\RouteRow{\(R_{\mathrm{act},\Omega}\)}{1.000000\allowbreak{}03855097\allowbreak{}23541416\allowbreak{}45}
\RouteGroup{Separación de prefijos}
\RouteRow{Área de ambigüedad; área \(\log_2\)}{18; 0.0}
\RouteRow{Primer bloque de separación}{órbita: 1; celda: 1; ruta: 1}
\RouteRow{Ambigüedad final}{1}
\RouteGroup{Secuencias completas}
\RouteRow{\(w_6\)}{24 243 657 24 243 423 15 486 423 24 243 657 24 243 423 15 486 423}
\RouteRow{Tríadas estables por bloque}{0 1 2 3 4 5 6 7 8 9 10 11 12 13 13 15 16 16}
\end{tabularx}\endgroup\par\vspace{5pt}\noindent{\fontsize{8.4}{10}\selectfont\hyperref[trace-96]{Procedencia y códigos originales}\hfill Ficha completa · 096/324}
\clearpage\phantomsection\label{nav:vi-ruta-97}
% VI-CATALOG-ROW rutas 97
\pdfbookmark[2]{Celda 25}{cell-25}
\pdfbookmark[3]{r3c7\_h-1v-1}{route-97}
\RouteTitle{Ruta 097}{r3c7\_\allowbreak{}h-1v-1}
\noindent Celda 25\quad (3, 7)\qquad Orientación \(h=-1,\ v=-1\)\hfill\hyperref[sec:vi-indice-rutas]{Índice}\par\vspace{4pt}
\begingroup\fontsize{10.2}{12.5}\selectfont\setlength{\tabcolsep}{5pt}\renewcommand{\arraystretch}{1.1}\rowcolors{1}{routepale}{white}\noindent\begin{tabularx}{\linewidth}{@{}>{\raggedright\arraybackslash}p{.345\linewidth}>{\raggedright\arraybackslash}X@{}}
\RouteGroup{Identidad estructural}
\RouteRow{Clase y multisección}{charged\_\allowbreak{}lepton\_\allowbreak{}G2\quad /\quad charged\_\allowbreak{}lepton\_\allowbreak{}G2}
\RouteRow{Colección y sector}{leptón cargado; leptón}
\RouteRow{Clasificación física}{leptón cargado}
\RouteRow{Generación, profundidad y color}{2; G2; \textemdash{}}
\RouteRow{Ruta primaria y órbita de inversión}{\(B_1\): no\qquad 3,\allowbreak{}7;\allowbreak{}7,\allowbreak{}3}
\RouteRow{Firma de clasificación}{34|\allowbreak{}0|\allowbreak{}-4|\allowbreak{}0|\allowbreak{}-12|\allowbreak{}+++-\kern0pt{}-\kern0pt{}-+++-\kern0pt{}-\kern0pt{}-}
\RouteGroup{Retornos, memoria y transporte}
\RouteRow{Retornos con fase}{clase: 27; órbita: 27; celda: 27}
\RouteRow{Retorno cinemático}{en 108: 54; extensión: 216}
\RouteRow{Giros y cambios de hoja}{71; 107}
\RouteRow{Acarreo y diferimiento}{deuda total: 2; final: 1; bloques diferidos: 1}
\RouteRow{Tríadas finales; persistencia de Pareto}{16; sí}
\RouteRow{\(K\longrightarrow K+9\)}{repeticiones: 9; cambios de memoria: 9}
\RouteGroup{Firma, fase y diferencias}
\RouteRow{\(n_A,\ s_C,\ \kappa_0,\ \nu_{120},\ \nu_{270}\)}{34, 0, 0, 0, -12}
\RouteRow{\(\tau_{12}\); firma histórica \(\mathbb Z^5\)}{+++-\kern0pt{}-\kern0pt{}-+++-\kern0pt{}-\kern0pt{}-; 34|\allowbreak{}0|\allowbreak{}0|\allowbreak{}0|\allowbreak{}-12}
\RouteRow{\(D_1,\ D_3,\ D_4,\ D_{27},\ D_{36}\)}{66, 60, 66, 48, 32}
\RouteRow{\(\Omega_{90/120},\ P_6\)}{80, 6}
\RouteGroup{Lectores y factores de acción}
\RouteRow{\(K_D\quad /\quad K_\Omega\)}{(80,\allowbreak{}66,\allowbreak{}3)\quad /\quad (80,\allowbreak{}54,\allowbreak{}6)}
\RouteRow{\(\kappa_D\)}{79.51870\allowbreak{}83196953\allowbreak{}45554341\allowbreak{}34}
\RouteRow{\(\kappa_\Omega\)}{79.60661\allowbreak{}01397948\allowbreak{}27586470\allowbreak{}91}
\RouteRow{\(R_{\mathrm{act},D}\)}{1.000000\allowbreak{}05512887\allowbreak{}17997050\allowbreak{}72}
\RouteRow{\(R_{\mathrm{act},\Omega}\)}{1.000000\allowbreak{}05518981\allowbreak{}25318591\allowbreak{}40}
\RouteGroup{Separación de prefijos}
\RouteRow{Área de ambigüedad; área \(\log_2\)}{18; 0.0}
\RouteRow{Primer bloque de separación}{órbita: 1; celda: 1; ruta: 1}
\RouteRow{Ambigüedad final}{1}
\RouteGroup{Secuencias completas}
\RouteRow{\(w_6\)}{264 90 570 264 90 24 243 657 24 264 90 570 264 90 24 243 657 24}
\RouteRow{Tríadas estables por bloque}{0 1 2 3 4 5 6 7 8 9 10 11 12 13 14 15 15 16}
\end{tabularx}\endgroup\par\vspace{5pt}\noindent{\fontsize{8.4}{10}\selectfont\hyperref[trace-97]{Procedencia y códigos originales}\hfill Ficha completa · 097/324}
\clearpage\phantomsection\label{nav:vi-ruta-98}
% VI-CATALOG-ROW rutas 98
\pdfbookmark[3]{r3c7\_h-1v+1}{route-98}
\RouteTitle{Ruta 098}{r3c7\_\allowbreak{}h-1v+1}
\noindent Celda 25\quad (3, 7)\qquad Orientación \(h=-1,\ v=1\)\hfill\hyperref[sec:vi-indice-rutas]{Índice}\par\vspace{4pt}
\begingroup\fontsize{10.2}{12.5}\selectfont\setlength{\tabcolsep}{5pt}\renewcommand{\arraystretch}{1.1}\rowcolors{1}{routepale}{white}\noindent\begin{tabularx}{\linewidth}{@{}>{\raggedright\arraybackslash}p{.345\linewidth}>{\raggedright\arraybackslash}X@{}}
\RouteGroup{Identidad estructural}
\RouteRow{Clase y multisección}{charged\_\allowbreak{}lepton\_\allowbreak{}G2\quad /\quad charged\_\allowbreak{}lepton\_\allowbreak{}G2}
\RouteRow{Colección y sector}{leptón cargado; leptón}
\RouteRow{Clasificación física}{leptón cargado}
\RouteRow{Generación, profundidad y color}{2; G2; \textemdash{}}
\RouteRow{Ruta primaria y órbita de inversión}{\(B_1\): no\qquad 3,\allowbreak{}7;\allowbreak{}7,\allowbreak{}3}
\RouteRow{Firma de clasificación}{34|\allowbreak{}0|\allowbreak{}-4|\allowbreak{}0|\allowbreak{}-12|\allowbreak{}+++-\kern0pt{}-\kern0pt{}-+++-\kern0pt{}-\kern0pt{}-}
\RouteGroup{Retornos, memoria y transporte}
\RouteRow{Retornos con fase}{clase: 27; órbita: 27; celda: 27}
\RouteRow{Retorno cinemático}{en 108: 54; extensión: 216}
\RouteRow{Giros y cambios de hoja}{71; 107}
\RouteRow{Acarreo y diferimiento}{deuda total: 2; final: 0; bloques diferidos: 1}
\RouteRow{Tríadas finales; persistencia de Pareto}{17; sí}
\RouteRow{\(K\longrightarrow K+9\)}{repeticiones: 9; cambios de memoria: 9}
\RouteGroup{Firma, fase y diferencias}
\RouteRow{\(n_A,\ s_C,\ \kappa_0,\ \nu_{120},\ \nu_{270}\)}{50, 4, -2, 2, -8}
\RouteRow{\(\tau_{12}\); firma histórica \(\mathbb Z^5\)}{+++-\kern0pt{}-0+++-\kern0pt{}-0; 50|\allowbreak{}4|\allowbreak{}-2|\allowbreak{}2|\allowbreak{}-8}
\RouteRow{\(D_1,\ D_3,\ D_4,\ D_{27},\ D_{36}\)}{84, 24, 84, 24, 16}
\RouteRow{\(\Omega_{90/120},\ P_6\)}{56, 5}
\RouteGroup{Lectores y factores de acción}
\RouteRow{\(K_D\quad /\quad K_\Omega\)}{(40,\allowbreak{}84,\allowbreak{}30)\quad /\quad (56,\allowbreak{}54,\allowbreak{}5)}
\RouteRow{\(\kappa_D\)}{39.39035\allowbreak{}82768609\allowbreak{}24917849\allowbreak{}59}
\RouteRow{\(\kappa_\Omega\)}{55.60649\allowbreak{}89433721\allowbreak{}35446416\allowbreak{}86}
\RouteRow{\(R_{\mathrm{act},D}\)}{1.000000\allowbreak{}02730861\allowbreak{}73967087\allowbreak{}05}
\RouteRow{\(R_{\mathrm{act},\Omega}\)}{1.000000\allowbreak{}03855097\allowbreak{}23541416\allowbreak{}45}
\RouteGroup{Separación de prefijos}
\RouteRow{Área de ambigüedad; área \(\log_2\)}{18; 0.0}
\RouteRow{Primer bloque de separación}{órbita: 1; celda: 1; ruta: 1}
\RouteRow{Ambigüedad final}{1}
\RouteGroup{Secuencias completas}
\RouteRow{\(w_6\)}{258 414 420 258 414 255 252 333 255 258 414 420 258 414 255 252 333 255}
\RouteRow{Tríadas estables por bloque}{0 1 2 3 4 5 6 7 8 9 10 11 11 12 14 15 16 17}
\end{tabularx}\endgroup\par\vspace{5pt}\noindent{\fontsize{8.4}{10}\selectfont\hyperref[trace-98]{Procedencia y códigos originales}\hfill Ficha completa · 098/324}
\clearpage\phantomsection\label{nav:vi-ruta-99}
% VI-CATALOG-ROW rutas 99
\pdfbookmark[3]{r3c7\_h+1v-1}{route-99}
\RouteTitle{Ruta 099}{r3c7\_\allowbreak{}h+1v-1}
\noindent Celda 25\quad (3, 7)\qquad Orientación \(h=1,\ v=-1\)\hfill\hyperref[sec:vi-indice-rutas]{Índice}\par\vspace{4pt}
\begingroup\fontsize{10.2}{12.5}\selectfont\setlength{\tabcolsep}{5pt}\renewcommand{\arraystretch}{1.1}\rowcolors{1}{routepale}{white}\noindent\begin{tabularx}{\linewidth}{@{}>{\raggedright\arraybackslash}p{.345\linewidth}>{\raggedright\arraybackslash}X@{}}
\RouteGroup{Identidad estructural}
\RouteRow{Clase y multisección}{charged\_\allowbreak{}lepton\_\allowbreak{}G2\quad /\quad charged\_\allowbreak{}lepton\_\allowbreak{}G2}
\RouteRow{Colección y sector}{leptón cargado; leptón}
\RouteRow{Clasificación física}{leptón cargado}
\RouteRow{Generación, profundidad y color}{2; G2; \textemdash{}}
\RouteRow{Ruta primaria y órbita de inversión}{\(B_1\): no\qquad 3,\allowbreak{}7;\allowbreak{}7,\allowbreak{}3}
\RouteRow{Firma de clasificación}{34|\allowbreak{}0|\allowbreak{}-4|\allowbreak{}0|\allowbreak{}-12|\allowbreak{}+++-\kern0pt{}-\kern0pt{}-+++-\kern0pt{}-\kern0pt{}-}
\RouteGroup{Retornos, memoria y transporte}
\RouteRow{Retornos con fase}{clase: 27; órbita: 27; celda: 27}
\RouteRow{Retorno cinemático}{en 108: 54; extensión: 216}
\RouteRow{Giros y cambios de hoja}{71; 107}
\RouteRow{Acarreo y diferimiento}{deuda total: 3; final: 1; bloques diferidos: 3}
\RouteRow{Tríadas finales; persistencia de Pareto}{16; sí}
\RouteRow{\(K\longrightarrow K+9\)}{repeticiones: 9; cambios de memoria: 9}
\RouteGroup{Firma, fase y diferencias}
\RouteRow{\(n_A,\ s_C,\ \kappa_0,\ \nu_{120},\ \nu_{270}\)}{50, -4, 2, -2, -8}
\RouteRow{\(\tau_{12}\); firma histórica \(\mathbb Z^5\)}{++0-\kern0pt{}-\kern0pt{}-++0-\kern0pt{}-\kern0pt{}-; 50|\allowbreak{}-4|\allowbreak{}2|\allowbreak{}-2|\allowbreak{}-8}
\RouteRow{\(D_1,\ D_3,\ D_4,\ D_{27},\ D_{36}\)}{84, 24, 84, 24, 16}
\RouteRow{\(\Omega_{90/120},\ P_6\)}{56, 5}
\RouteGroup{Lectores y factores de acción}
\RouteRow{\(K_D\quad /\quad K_\Omega\)}{(40,\allowbreak{}84,\allowbreak{}30)\quad /\quad (56,\allowbreak{}54,\allowbreak{}5)}
\RouteRow{\(\kappa_D\)}{39.39035\allowbreak{}82768609\allowbreak{}24917849\allowbreak{}59}
\RouteRow{\(\kappa_\Omega\)}{55.60649\allowbreak{}89433721\allowbreak{}35446416\allowbreak{}86}
\RouteRow{\(R_{\mathrm{act},D}\)}{1.000000\allowbreak{}02730861\allowbreak{}73967087\allowbreak{}05}
\RouteRow{\(R_{\mathrm{act},\Omega}\)}{1.000000\allowbreak{}03855097\allowbreak{}23541416\allowbreak{}45}
\RouteGroup{Separación de prefijos}
\RouteRow{Área de ambigüedad; área \(\log_2\)}{18; 0.0}
\RouteRow{Primer bloque de separación}{órbita: 1; celda: 1; ruta: 1}
\RouteRow{Ambigüedad final}{1}
\RouteGroup{Secuencias completas}
\RouteRow{\(w_6\)}{252 333 255 252 333 420 258 414 420 252 333 255 252 333 420 258 414 420}
\RouteRow{Tríadas estables por bloque}{0 1 2 3 4 5 6 7 8 9 10 11 12 12 14 14 16 16}
\end{tabularx}\endgroup\par\vspace{5pt}\noindent{\fontsize{8.4}{10}\selectfont\hyperref[trace-99]{Procedencia y códigos originales}\hfill Ficha completa · 099/324}
\clearpage\phantomsection\label{nav:vi-ruta-100}
% VI-CATALOG-ROW rutas 100
\pdfbookmark[3]{r3c7\_h+1v+1}{route-100}
\RouteTitle{Ruta 100}{r3c7\_\allowbreak{}h+1v+1}
\noindent Celda 25\quad (3, 7)\qquad Orientación \(h=1,\ v=1\)\hfill\hyperref[sec:vi-indice-rutas]{Índice}\par\vspace{4pt}
\begingroup\fontsize{10.2}{12.5}\selectfont\setlength{\tabcolsep}{5pt}\renewcommand{\arraystretch}{1.1}\rowcolors{1}{routepale}{white}\noindent\begin{tabularx}{\linewidth}{@{}>{\raggedright\arraybackslash}p{.345\linewidth}>{\raggedright\arraybackslash}X@{}}
\RouteGroup{Identidad estructural}
\RouteRow{Clase y multisección}{charged\_\allowbreak{}lepton\_\allowbreak{}G2\quad /\quad charged\_\allowbreak{}lepton\_\allowbreak{}G2}
\RouteRow{Colección y sector}{leptón cargado; leptón}
\RouteRow{Clasificación física}{leptón cargado}
\RouteRow{Generación, profundidad y color}{2; G2; \textemdash{}}
\RouteRow{Ruta primaria y órbita de inversión}{\(B_1\): sí\qquad 3,\allowbreak{}7;\allowbreak{}7,\allowbreak{}3}
\RouteRow{Firma de clasificación}{34|\allowbreak{}0|\allowbreak{}-4|\allowbreak{}0|\allowbreak{}-12|\allowbreak{}+++-\kern0pt{}-\kern0pt{}-+++-\kern0pt{}-\kern0pt{}-}
\RouteGroup{Retornos, memoria y transporte}
\RouteRow{Retornos con fase}{clase: 27; órbita: 27; celda: 27}
\RouteRow{Retorno cinemático}{en 108: 54; extensión: 216}
\RouteRow{Giros y cambios de hoja}{71; 107}
\RouteRow{Acarreo y diferimiento}{deuda total: 0; final: 0; bloques diferidos: 0}
\RouteRow{Tríadas finales; persistencia de Pareto}{17; sí}
\RouteRow{\(K\longrightarrow K+9\)}{repeticiones: 9; cambios de memoria: 9}
\RouteGroup{Firma, fase y diferencias}
\RouteRow{\(n_A,\ s_C,\ \kappa_0,\ \nu_{120},\ \nu_{270}\)}{34, 0, 0, 0, -12}
\RouteRow{\(\tau_{12}\); firma histórica \(\mathbb Z^5\)}{+++-\kern0pt{}-\kern0pt{}-+++-\kern0pt{}-\kern0pt{}-; 34|\allowbreak{}0|\allowbreak{}0|\allowbreak{}0|\allowbreak{}-12}
\RouteRow{\(D_1,\ D_3,\ D_4,\ D_{27},\ D_{36}\)}{66, 60, 66, 48, 32}
\RouteRow{\(\Omega_{90/120},\ P_6\)}{80, 6}
\RouteGroup{Lectores y factores de acción}
\RouteRow{\(K_D\quad /\quad K_\Omega\)}{(80,\allowbreak{}66,\allowbreak{}3)\quad /\quad (80,\allowbreak{}54,\allowbreak{}6)}
\RouteRow{\(\kappa_D\)}{79.51870\allowbreak{}83196953\allowbreak{}45554341\allowbreak{}34}
\RouteRow{\(\kappa_\Omega\)}{79.60661\allowbreak{}01397948\allowbreak{}27586470\allowbreak{}91}
\RouteRow{\(R_{\mathrm{act},D}\)}{1.000000\allowbreak{}05512887\allowbreak{}17997050\allowbreak{}72}
\RouteRow{\(R_{\mathrm{act},\Omega}\)}{1.000000\allowbreak{}05518981\allowbreak{}25318591\allowbreak{}40}
\RouteGroup{Separación de prefijos}
\RouteRow{Área de ambigüedad; área \(\log_2\)}{18; 0.0}
\RouteRow{Primer bloque de separación}{órbita: 1; celda: 1; ruta: 1}
\RouteRow{Ambigüedad final}{1}
\RouteGroup{Secuencias completas}
\RouteRow{\(w_6\)}{243 657 24 243 657 570 264 90 570 243 657 24 243 657 570 264 90 570}
\RouteRow{Tríadas estables por bloque}{0 1 2 3 4 5 6 7 8 9 10 11 12 13 14 15 16 17}
\end{tabularx}\endgroup\par\vspace{5pt}\noindent{\fontsize{8.4}{10}\selectfont\hyperref[trace-100]{Procedencia y códigos originales}\hfill Ficha completa · 100/324}
\clearpage\phantomsection\label{nav:vi-ruta-101}
% VI-CATALOG-ROW rutas 101
\pdfbookmark[2]{Celda 26}{cell-26}
\pdfbookmark[3]{r3c8\_h-1v-1}{route-101}
\RouteTitle{Ruta 101}{r3c8\_\allowbreak{}h-1v-1}
\noindent Celda 26\quad (3, 8)\qquad Orientación \(h=-1,\ v=-1\)\hfill\hyperref[sec:vi-indice-rutas]{Índice}\par\vspace{4pt}
\begingroup\fontsize{10.2}{12.5}\selectfont\setlength{\tabcolsep}{5pt}\renewcommand{\arraystretch}{1.1}\rowcolors{1}{routepale}{white}\noindent\begin{tabularx}{\linewidth}{@{}>{\raggedright\arraybackslash}p{.345\linewidth}>{\raggedright\arraybackslash}X@{}}
\RouteGroup{Identidad estructural}
\RouteRow{Clase y multisección}{neutrino\_\allowbreak{}G3\quad /\quad neutrino\_\allowbreak{}G3}
\RouteRow{Colección y sector}{neutrino; leptón}
\RouteRow{Clasificación física}{neutrino}
\RouteRow{Generación, profundidad y color}{3; G3; \textemdash{}}
\RouteRow{Ruta primaria y órbita de inversión}{\(B_1\): no\qquad 3,\allowbreak{}8;\allowbreak{}7,\allowbreak{}2}
\RouteRow{Firma de clasificación}{34|\allowbreak{}0|\allowbreak{}0|\allowbreak{}-2|\allowbreak{}-8|\allowbreak{}++0-\kern0pt{}-\kern0pt{}-++0-\kern0pt{}-\kern0pt{}-}
\RouteGroup{Retornos, memoria y transporte}
\RouteRow{Retornos con fase}{clase: 27; órbita: 27; celda: 27}
\RouteRow{Retorno cinemático}{en 108: 54; extensión: 216}
\RouteRow{Giros y cambios de hoja}{71; 107}
\RouteRow{Acarreo y diferimiento}{deuda total: 1; final: 0; bloques diferidos: 1}
\RouteRow{Tríadas finales; persistencia de Pareto}{17; sí}
\RouteRow{\(K\longrightarrow K+9\)}{repeticiones: 9; cambios de memoria: 9}
\RouteGroup{Firma, fase y diferencias}
\RouteRow{\(n_A,\ s_C,\ \kappa_0,\ \nu_{120},\ \nu_{270}\)}{34, 0, -2, 0, -8}
\RouteRow{\(\tau_{12}\); firma histórica \(\mathbb Z^5\)}{+++-\kern0pt{}-0+++-\kern0pt{}-0; 34|\allowbreak{}0|\allowbreak{}-2|\allowbreak{}0|\allowbreak{}-8}
\RouteRow{\(D_1,\ D_3,\ D_4,\ D_{27},\ D_{36}\)}{66, 60, 66, 48, 32}
\RouteRow{\(\Omega_{90/120},\ P_6\)}{56, 5}
\RouteGroup{Lectores y factores de acción}
\RouteRow{\(K_D\quad /\quad K_\Omega\)}{(80,\allowbreak{}66,\allowbreak{}3)\quad /\quad (56,\allowbreak{}54,\allowbreak{}5)}
\RouteRow{\(\kappa_D\)}{79.51870\allowbreak{}83196953\allowbreak{}45554341\allowbreak{}34}
\RouteRow{\(\kappa_\Omega\)}{55.60649\allowbreak{}89433721\allowbreak{}35446416\allowbreak{}86}
\RouteRow{\(R_{\mathrm{act},D}\)}{1.000000\allowbreak{}05512887\allowbreak{}17997050\allowbreak{}72}
\RouteRow{\(R_{\mathrm{act},\Omega}\)}{1.000000\allowbreak{}03855097\allowbreak{}23541416\allowbreak{}45}
\RouteGroup{Separación de prefijos}
\RouteRow{Área de ambigüedad; área \(\log_2\)}{36; 18.0}
\RouteRow{Primer bloque de separación}{órbita: \textemdash{}; celda: \textemdash{}; ruta: \textemdash{}}
\RouteRow{Ambigüedad final}{2}
\RouteGroup{Secuencias completas}
\RouteRow{\(w_6\)}{486 423 15 486 423 327 498 99 327 486 423 15 486 423 327 498 99 327}
\RouteRow{Tríadas estables por bloque}{0 1 2 3 4 5 6 7 8 9 10 10 12 13 14 15 16 17}
\end{tabularx}\endgroup\par\vspace{5pt}\noindent{\fontsize{8.4}{10}\selectfont\hyperref[trace-101]{Procedencia y códigos originales}\hfill Ficha completa · 101/324}
\clearpage\phantomsection\label{nav:vi-ruta-102}
% VI-CATALOG-ROW rutas 102
\pdfbookmark[3]{r3c8\_h-1v+1}{route-102}
\RouteTitle{Ruta 102}{r3c8\_\allowbreak{}h-1v+1}
\noindent Celda 26\quad (3, 8)\qquad Orientación \(h=-1,\ v=1\)\hfill\hyperref[sec:vi-indice-rutas]{Índice}\par\vspace{4pt}
\begingroup\fontsize{10.2}{12.5}\selectfont\setlength{\tabcolsep}{5pt}\renewcommand{\arraystretch}{1.1}\rowcolors{1}{routepale}{white}\noindent\begin{tabularx}{\linewidth}{@{}>{\raggedright\arraybackslash}p{.345\linewidth}>{\raggedright\arraybackslash}X@{}}
\RouteGroup{Identidad estructural}
\RouteRow{Clase y multisección}{neutrino\_\allowbreak{}G3\quad /\quad neutrino\_\allowbreak{}G3}
\RouteRow{Colección y sector}{neutrino; leptón}
\RouteRow{Clasificación física}{neutrino}
\RouteRow{Generación, profundidad y color}{3; G3; \textemdash{}}
\RouteRow{Ruta primaria y órbita de inversión}{\(B_1\): sí\qquad 3,\allowbreak{}8;\allowbreak{}7,\allowbreak{}2}
\RouteRow{Firma de clasificación}{34|\allowbreak{}0|\allowbreak{}0|\allowbreak{}-2|\allowbreak{}-8|\allowbreak{}++0-\kern0pt{}-\kern0pt{}-++0-\kern0pt{}-\kern0pt{}-}
\RouteGroup{Retornos, memoria y transporte}
\RouteRow{Retornos con fase}{clase: 27; órbita: 27; celda: 27}
\RouteRow{Retorno cinemático}{en 108: 54; extensión: 216}
\RouteRow{Giros y cambios de hoja}{71; 107}
\RouteRow{Acarreo y diferimiento}{deuda total: 1; final: 0; bloques diferidos: 1}
\RouteRow{Tríadas finales; persistencia de Pareto}{17; sí}
\RouteRow{\(K\longrightarrow K+9\)}{repeticiones: 9; cambios de memoria: 9}
\RouteGroup{Firma, fase y diferencias}
\RouteRow{\(n_A,\ s_C,\ \kappa_0,\ \nu_{120},\ \nu_{270}\)}{14, -4, 0, 0, -12}
\RouteRow{\(\tau_{12}\); firma histórica \(\mathbb Z^5\)}{+++-\kern0pt{}-\kern0pt{}-+++-\kern0pt{}-\kern0pt{}-; 14|\allowbreak{}-4|\allowbreak{}0|\allowbreak{}0|\allowbreak{}-12}
\RouteRow{\(D_1,\ D_3,\ D_4,\ D_{27},\ D_{36}\)}{78, 24, 78, 24, 16}
\RouteRow{\(\Omega_{90/120},\ P_6\)}{80, 6}
\RouteGroup{Lectores y factores de acción}
\RouteRow{\(K_D\quad /\quad K_\Omega\)}{(40,\allowbreak{}78,\allowbreak{}27)\quad /\quad (80,\allowbreak{}54,\allowbreak{}6)}
\RouteRow{\(\kappa_D\)}{39.43380\allowbreak{}88030085\allowbreak{}51303671\allowbreak{}14}
\RouteRow{\(\kappa_\Omega\)}{79.60661\allowbreak{}01397948\allowbreak{}27586470\allowbreak{}91}
\RouteRow{\(R_{\mathrm{act},D}\)}{1.000000\allowbreak{}02733874\allowbreak{}08548943\allowbreak{}58}
\RouteRow{\(R_{\mathrm{act},\Omega}\)}{1.000000\allowbreak{}05518981\allowbreak{}25318591\allowbreak{}40}
\RouteGroup{Separación de prefijos}
\RouteRow{Área de ambigüedad; área \(\log_2\)}{36; 18.0}
\RouteRow{Primer bloque de separación}{órbita: \textemdash{}; celda: \textemdash{}; ruta: \textemdash{}}
\RouteRow{Ambigüedad final}{2}
\RouteGroup{Secuencias completas}
\RouteRow{\(w_6\)}{504 585 507 504 585 672 510 666 672 504 585 507 504 585 672 510 666 672}
\RouteRow{Tríadas estables por bloque}{0 1 2 3 4 5 6 7 8 9 10 11 12 13 14 14 16 17}
\end{tabularx}\endgroup\par\vspace{5pt}\noindent{\fontsize{8.4}{10}\selectfont\hyperref[trace-102]{Procedencia y códigos originales}\hfill Ficha completa · 102/324}
\clearpage\phantomsection\label{nav:vi-ruta-103}
% VI-CATALOG-ROW rutas 103
\pdfbookmark[3]{r3c8\_h+1v-1}{route-103}
\RouteTitle{Ruta 103}{r3c8\_\allowbreak{}h+1v-1}
\noindent Celda 26\quad (3, 8)\qquad Orientación \(h=1,\ v=-1\)\hfill\hyperref[sec:vi-indice-rutas]{Índice}\par\vspace{4pt}
\begingroup\fontsize{10.2}{12.5}\selectfont\setlength{\tabcolsep}{5pt}\renewcommand{\arraystretch}{1.1}\rowcolors{1}{routepale}{white}\noindent\begin{tabularx}{\linewidth}{@{}>{\raggedright\arraybackslash}p{.345\linewidth}>{\raggedright\arraybackslash}X@{}}
\RouteGroup{Identidad estructural}
\RouteRow{Clase y multisección}{neutrino\_\allowbreak{}G3\quad /\quad neutrino\_\allowbreak{}G3}
\RouteRow{Colección y sector}{neutrino; leptón}
\RouteRow{Clasificación física}{neutrino}
\RouteRow{Generación, profundidad y color}{3; G3; \textemdash{}}
\RouteRow{Ruta primaria y órbita de inversión}{\(B_1\): no\qquad 3,\allowbreak{}8;\allowbreak{}7,\allowbreak{}2}
\RouteRow{Firma de clasificación}{34|\allowbreak{}0|\allowbreak{}0|\allowbreak{}-2|\allowbreak{}-8|\allowbreak{}++0-\kern0pt{}-\kern0pt{}-++0-\kern0pt{}-\kern0pt{}-}
\RouteGroup{Retornos, memoria y transporte}
\RouteRow{Retornos con fase}{clase: 27; órbita: 27; celda: 27}
\RouteRow{Retorno cinemático}{en 108: 54; extensión: 216}
\RouteRow{Giros y cambios de hoja}{71; 107}
\RouteRow{Acarreo y diferimiento}{deuda total: 1; final: 0; bloques diferidos: 1}
\RouteRow{Tríadas finales; persistencia de Pareto}{17; sí}
\RouteRow{\(K\longrightarrow K+9\)}{repeticiones: 9; cambios de memoria: 9}
\RouteGroup{Firma, fase y diferencias}
\RouteRow{\(n_A,\ s_C,\ \kappa_0,\ \nu_{120},\ \nu_{270}\)}{14, 4, 0, 0, -12}
\RouteRow{\(\tau_{12}\); firma histórica \(\mathbb Z^5\)}{+++-\kern0pt{}-\kern0pt{}-+++-\kern0pt{}-\kern0pt{}-; 14|\allowbreak{}4|\allowbreak{}0|\allowbreak{}0|\allowbreak{}-12}
\RouteRow{\(D_1,\ D_3,\ D_4,\ D_{27},\ D_{36}\)}{78, 24, 78, 24, 16}
\RouteRow{\(\Omega_{90/120},\ P_6\)}{80, 6}
\RouteGroup{Lectores y factores de acción}
\RouteRow{\(K_D\quad /\quad K_\Omega\)}{(40,\allowbreak{}78,\allowbreak{}27)\quad /\quad (80,\allowbreak{}54,\allowbreak{}6)}
\RouteRow{\(\kappa_D\)}{39.43380\allowbreak{}88030085\allowbreak{}51303671\allowbreak{}14}
\RouteRow{\(\kappa_\Omega\)}{79.60661\allowbreak{}01397948\allowbreak{}27586470\allowbreak{}91}
\RouteRow{\(R_{\mathrm{act},D}\)}{1.000000\allowbreak{}02733874\allowbreak{}08548943\allowbreak{}58}
\RouteRow{\(R_{\mathrm{act},\Omega}\)}{1.000000\allowbreak{}05518981\allowbreak{}25318591\allowbreak{}40}
\RouteGroup{Separación de prefijos}
\RouteRow{Área de ambigüedad; área \(\log_2\)}{36; 18.0}
\RouteRow{Primer bloque de separación}{órbita: \textemdash{}; celda: \textemdash{}; ruta: \textemdash{}}
\RouteRow{Ambigüedad final}{2}
\RouteGroup{Secuencias completas}
\RouteRow{\(w_6\)}{510 666 672 510 666 507 504 585 507 510 666 672 510 666 507 504 585 507}
\RouteRow{Tríadas estables por bloque}{0 1 2 3 4 5 6 7 8 9 10 11 12 13 14 15 15 17}
\end{tabularx}\endgroup\par\vspace{5pt}\noindent{\fontsize{8.4}{10}\selectfont\hyperref[trace-103]{Procedencia y códigos originales}\hfill Ficha completa · 103/324}
\clearpage\phantomsection\label{nav:vi-ruta-104}
% VI-CATALOG-ROW rutas 104
\pdfbookmark[3]{r3c8\_h+1v+1}{route-104}
\RouteTitle{Ruta 104}{r3c8\_\allowbreak{}h+1v+1}
\noindent Celda 26\quad (3, 8)\qquad Orientación \(h=1,\ v=1\)\hfill\hyperref[sec:vi-indice-rutas]{Índice}\par\vspace{4pt}
\begingroup\fontsize{10.2}{12.5}\selectfont\setlength{\tabcolsep}{5pt}\renewcommand{\arraystretch}{1.1}\rowcolors{1}{routepale}{white}\noindent\begin{tabularx}{\linewidth}{@{}>{\raggedright\arraybackslash}p{.345\linewidth}>{\raggedright\arraybackslash}X@{}}
\RouteGroup{Identidad estructural}
\RouteRow{Clase y multisección}{neutrino\_\allowbreak{}G3\quad /\quad neutrino\_\allowbreak{}G3}
\RouteRow{Colección y sector}{neutrino; leptón}
\RouteRow{Clasificación física}{neutrino}
\RouteRow{Generación, profundidad y color}{3; G3; \textemdash{}}
\RouteRow{Ruta primaria y órbita de inversión}{\(B_1\): no\qquad 3,\allowbreak{}8;\allowbreak{}7,\allowbreak{}2}
\RouteRow{Firma de clasificación}{34|\allowbreak{}0|\allowbreak{}0|\allowbreak{}-2|\allowbreak{}-8|\allowbreak{}++0-\kern0pt{}-\kern0pt{}-++0-\kern0pt{}-\kern0pt{}-}
\RouteGroup{Retornos, memoria y transporte}
\RouteRow{Retornos con fase}{clase: 27; órbita: 27; celda: 27}
\RouteRow{Retorno cinemático}{en 108: 54; extensión: 216}
\RouteRow{Giros y cambios de hoja}{71; 107}
\RouteRow{Acarreo y diferimiento}{deuda total: 0; final: 0; bloques diferidos: 0}
\RouteRow{Tríadas finales; persistencia de Pareto}{17; sí}
\RouteRow{\(K\longrightarrow K+9\)}{repeticiones: 9; cambios de memoria: 9}
\RouteGroup{Firma, fase y diferencias}
\RouteRow{\(n_A,\ s_C,\ \kappa_0,\ \nu_{120},\ \nu_{270}\)}{34, 0, 2, 0, -8}
\RouteRow{\(\tau_{12}\); firma histórica \(\mathbb Z^5\)}{++0-\kern0pt{}-\kern0pt{}-++0-\kern0pt{}-\kern0pt{}-; 34|\allowbreak{}0|\allowbreak{}2|\allowbreak{}0|\allowbreak{}-8}
\RouteRow{\(D_1,\ D_3,\ D_4,\ D_{27},\ D_{36}\)}{66, 60, 66, 48, 32}
\RouteRow{\(\Omega_{90/120},\ P_6\)}{56, 5}
\RouteGroup{Lectores y factores de acción}
\RouteRow{\(K_D\quad /\quad K_\Omega\)}{(80,\allowbreak{}66,\allowbreak{}3)\quad /\quad (56,\allowbreak{}54,\allowbreak{}5)}
\RouteRow{\(\kappa_D\)}{79.51870\allowbreak{}83196953\allowbreak{}45554341\allowbreak{}34}
\RouteRow{\(\kappa_\Omega\)}{55.60649\allowbreak{}89433721\allowbreak{}35446416\allowbreak{}86}
\RouteRow{\(R_{\mathrm{act},D}\)}{1.000000\allowbreak{}05512887\allowbreak{}17997050\allowbreak{}72}
\RouteRow{\(R_{\mathrm{act},\Omega}\)}{1.000000\allowbreak{}03855097\allowbreak{}23541416\allowbreak{}45}
\RouteGroup{Separación de prefijos}
\RouteRow{Área de ambigüedad; área \(\log_2\)}{36; 18.0}
\RouteRow{Primer bloque de separación}{órbita: \textemdash{}; celda: \textemdash{}; ruta: \textemdash{}}
\RouteRow{Ambigüedad final}{2}
\RouteGroup{Secuencias completas}
\RouteRow{\(w_6\)}{498 99 327 498 99 15 486 423 15 498 99 327 498 99 15 486 423 15}
\RouteRow{Tríadas estables por bloque}{0 1 2 3 4 5 6 7 8 9 10 11 12 13 14 15 16 17}
\end{tabularx}\endgroup\par\vspace{5pt}\noindent{\fontsize{8.4}{10}\selectfont\hyperref[trace-104]{Procedencia y códigos originales}\hfill Ficha completa · 104/324}
\clearpage\phantomsection\label{nav:vi-ruta-105}
% VI-CATALOG-ROW rutas 105
\pdfbookmark[2]{Celda 27}{cell-27}
\pdfbookmark[3]{r3c9\_h-1v-1}{route-105}
\RouteTitle{Ruta 105}{r3c9\_\allowbreak{}h-1v-1}
\noindent Celda 27\quad (3, 9)\qquad Orientación \(h=-1,\ v=-1\)\hfill\hyperref[sec:vi-indice-rutas]{Índice}\par\vspace{4pt}
\begingroup\fontsize{10.2}{12.5}\selectfont\setlength{\tabcolsep}{5pt}\renewcommand{\arraystretch}{1.1}\rowcolors{1}{routepale}{white}\noindent\begin{tabularx}{\linewidth}{@{}>{\raggedright\arraybackslash}p{.345\linewidth}>{\raggedright\arraybackslash}X@{}}
\RouteGroup{Identidad estructural}
\RouteRow{Clase y multisección}{charged\_\allowbreak{}lepton\_\allowbreak{}G4\quad /\quad charged\_\allowbreak{}lepton\_\allowbreak{}G4}
\RouteRow{Colección y sector}{leptón cargado; leptón}
\RouteRow{Clasificación física}{leptón cargado}
\RouteRow{Generación, profundidad y color}{4; G4; \textemdash{}}
\RouteRow{Ruta primaria y órbita de inversión}{\(B_1\): sí\qquad 3,\allowbreak{}9;\allowbreak{}7,\allowbreak{}1}
\RouteRow{Firma de clasificación}{30|\allowbreak{}6|\allowbreak{}0|\allowbreak{}2|\allowbreak{}-8|\allowbreak{}+++-\kern0pt{}-0+++-\kern0pt{}-0}
\RouteGroup{Retornos, memoria y transporte}
\RouteRow{Retornos con fase}{clase: 27; órbita: 27; celda: 27}
\RouteRow{Retorno cinemático}{en 108: 54; extensión: 216}
\RouteRow{Giros y cambios de hoja}{71; 107}
\RouteRow{Acarreo y diferimiento}{deuda total: 1; final: 0; bloques diferidos: 1}
\RouteRow{Tríadas finales; persistencia de Pareto}{17; sí}
\RouteRow{\(K\longrightarrow K+9\)}{repeticiones: 9; cambios de memoria: 9}
\RouteGroup{Firma, fase y diferencias}
\RouteRow{\(n_A,\ s_C,\ \kappa_0,\ \nu_{120},\ \nu_{270}\)}{30, -6, 2, -4, -8}
\RouteRow{\(\tau_{12}\); firma histórica \(\mathbb Z^5\)}{++0-\kern0pt{}-\kern0pt{}-++0-\kern0pt{}-\kern0pt{}-; 30|\allowbreak{}-6|\allowbreak{}2|\allowbreak{}-4|\allowbreak{}-8}
\RouteRow{\(D_1,\ D_3,\ D_4,\ D_{27},\ D_{36}\)}{54, 60, 72, 24, 16}
\RouteRow{\(\Omega_{90/120},\ P_6\)}{56, 5}
\RouteGroup{Lectores y factores de acción}
\RouteRow{\(K_D\quad /\quad K_\Omega\)}{(40,\allowbreak{}54,\allowbreak{}6)\quad /\quad (56,\allowbreak{}54,\allowbreak{}5)}
\RouteRow{\(\kappa_D\)}{39.60661\allowbreak{}01397948\allowbreak{}27586470\allowbreak{}91}
\RouteRow{\(\kappa_\Omega\)}{55.60649\allowbreak{}89433721\allowbreak{}35446416\allowbreak{}86}
\RouteRow{\(R_{\mathrm{act},D}\)}{1.000000\allowbreak{}02745854\allowbreak{}08735595\allowbreak{}18}
\RouteRow{\(R_{\mathrm{act},\Omega}\)}{1.000000\allowbreak{}03855097\allowbreak{}23541416\allowbreak{}45}
\RouteGroup{Separación de prefijos}
\RouteRow{Área de ambigüedad; área \(\log_2\)}{58; 30.18947\allowbreak{}50100961\allowbreak{}74}
\RouteRow{Primer bloque de separación}{órbita: \textemdash{}; celda: \textemdash{}; ruta: \textemdash{}}
\RouteRow{Ambigüedad final}{3}
\RouteGroup{Secuencias completas}
\RouteRow{\(w_6\)}{15 486 423 15 486 657 24 243 657 15 486 423 15 486 657 24 243 657}
\RouteRow{Tríadas estables por bloque}{0 1 2 3 4 5 6 7 8 9 10 11 12 13 14 15 15 17}
\end{tabularx}\endgroup\par\vspace{5pt}\noindent{\fontsize{8.4}{10}\selectfont\hyperref[trace-105]{Procedencia y códigos originales}\hfill Ficha completa · 105/324}
\clearpage\phantomsection\label{nav:vi-ruta-106}
% VI-CATALOG-ROW rutas 106
\pdfbookmark[3]{r3c9\_h-1v+1}{route-106}
\RouteTitle{Ruta 106}{r3c9\_\allowbreak{}h-1v+1}
\noindent Celda 27\quad (3, 9)\qquad Orientación \(h=-1,\ v=1\)\hfill\hyperref[sec:vi-indice-rutas]{Índice}\par\vspace{4pt}
\begingroup\fontsize{10.2}{12.5}\selectfont\setlength{\tabcolsep}{5pt}\renewcommand{\arraystretch}{1.1}\rowcolors{1}{routepale}{white}\noindent\begin{tabularx}{\linewidth}{@{}>{\raggedright\arraybackslash}p{.345\linewidth}>{\raggedright\arraybackslash}X@{}}
\RouteGroup{Identidad estructural}
\RouteRow{Clase y multisección}{charged\_\allowbreak{}lepton\_\allowbreak{}G4\quad /\quad charged\_\allowbreak{}lepton\_\allowbreak{}G4}
\RouteRow{Colección y sector}{leptón cargado; leptón}
\RouteRow{Clasificación física}{leptón cargado}
\RouteRow{Generación, profundidad y color}{4; G4; \textemdash{}}
\RouteRow{Ruta primaria y órbita de inversión}{\(B_1\): no\qquad 3,\allowbreak{}9;\allowbreak{}7,\allowbreak{}1}
\RouteRow{Firma de clasificación}{30|\allowbreak{}6|\allowbreak{}0|\allowbreak{}2|\allowbreak{}-8|\allowbreak{}+++-\kern0pt{}-0+++-\kern0pt{}-0}
\RouteGroup{Retornos, memoria y transporte}
\RouteRow{Retornos con fase}{clase: 18; órbita: 27; celda: 27}
\RouteRow{Retorno cinemático}{en 108: 54; extensión: 216}
\RouteRow{Giros y cambios de hoja}{71; 107}
\RouteRow{Acarreo y diferimiento}{deuda total: 0; final: 0; bloques diferidos: 0}
\RouteRow{Tríadas finales; persistencia de Pareto}{17; sí}
\RouteRow{\(K\longrightarrow K+9\)}{repeticiones: 9; cambios de memoria: 9}
\RouteGroup{Firma, fase y diferencias}
\RouteRow{\(n_A,\ s_C,\ \kappa_0,\ \nu_{120},\ \nu_{270}\)}{50, -6, 2, -4, -8}
\RouteRow{\(\tau_{12}\); firma histórica \(\mathbb Z^5\)}{++0-\kern0pt{}-\kern0pt{}-++0-\kern0pt{}-\kern0pt{}-; 50|\allowbreak{}-6|\allowbreak{}2|\allowbreak{}-4|\allowbreak{}-8}
\RouteRow{\(D_1,\ D_3,\ D_4,\ D_{27},\ D_{36}\)}{24, 24, 24, 0, 0}
\RouteRow{\(\Omega_{90/120},\ P_6\)}{56, 5}
\RouteGroup{Lectores y factores de acción}
\RouteRow{\(K_D\quad /\quad K_\Omega\)}{(0,\allowbreak{}24,\allowbreak{}0)\quad /\quad (56,\allowbreak{}54,\allowbreak{}5)}
\RouteRow{\(\kappa_D\)}{-0.175136\allowbreak{}46166281\allowbreak{}12239348\allowbreak{}425}
\RouteRow{\(\kappa_\Omega\)}{55.60649\allowbreak{}89433721\allowbreak{}35446416\allowbreak{}86}
\RouteRow{\(R_{\mathrm{act},D}\)}{0.999999\allowbreak{}99987858\allowbreak{}10851337\allowbreak{}881}
\RouteRow{\(R_{\mathrm{act},\Omega}\)}{1.000000\allowbreak{}03855097\allowbreak{}23541416\allowbreak{}45}
\RouteGroup{Separación de prefijos}
\RouteRow{Área de ambigüedad; área \(\log_2\)}{108; 46.52932\allowbreak{}50129807\allowbreak{}9}
\RouteRow{Primer bloque de separación}{órbita: \textemdash{}; celda: \textemdash{}; ruta: \textemdash{}}
\RouteRow{Ambigüedad final}{6}
\RouteGroup{Secuencias completas}
\RouteRow{\(w_6\)}{3 0 81 3 0 81 3 0 81 3 0 81 3 0 81 3 0 81}
\RouteRow{Tríadas estables por bloque}{0 1 2 3 4 5 6 7 8 9 10 11 12 13 14 15 16 17}
\end{tabularx}\endgroup\par\vspace{5pt}\noindent{\fontsize{8.4}{10}\selectfont\hyperref[trace-106]{Procedencia y códigos originales}\hfill Ficha completa · 106/324}
\clearpage\phantomsection\label{nav:vi-ruta-107}
% VI-CATALOG-ROW rutas 107
\pdfbookmark[3]{r3c9\_h+1v-1}{route-107}
\RouteTitle{Ruta 107}{r3c9\_\allowbreak{}h+1v-1}
\noindent Celda 27\quad (3, 9)\qquad Orientación \(h=1,\ v=-1\)\hfill\hyperref[sec:vi-indice-rutas]{Índice}\par\vspace{4pt}
\begingroup\fontsize{10.2}{12.5}\selectfont\setlength{\tabcolsep}{5pt}\renewcommand{\arraystretch}{1.1}\rowcolors{1}{routepale}{white}\noindent\begin{tabularx}{\linewidth}{@{}>{\raggedright\arraybackslash}p{.345\linewidth}>{\raggedright\arraybackslash}X@{}}
\RouteGroup{Identidad estructural}
\RouteRow{Clase y multisección}{charged\_\allowbreak{}lepton\_\allowbreak{}G4\quad /\quad charged\_\allowbreak{}lepton\_\allowbreak{}G4}
\RouteRow{Colección y sector}{leptón cargado; leptón}
\RouteRow{Clasificación física}{leptón cargado}
\RouteRow{Generación, profundidad y color}{4; G4; \textemdash{}}
\RouteRow{Ruta primaria y órbita de inversión}{\(B_1\): no\qquad 3,\allowbreak{}9;\allowbreak{}7,\allowbreak{}1}
\RouteRow{Firma de clasificación}{30|\allowbreak{}6|\allowbreak{}0|\allowbreak{}2|\allowbreak{}-8|\allowbreak{}+++-\kern0pt{}-0+++-\kern0pt{}-0}
\RouteGroup{Retornos, memoria y transporte}
\RouteRow{Retornos con fase}{clase: 9; órbita: 27; celda: 27}
\RouteRow{Retorno cinemático}{en 108: 54; extensión: 216}
\RouteRow{Giros y cambios de hoja}{71; 107}
\RouteRow{Acarreo y diferimiento}{deuda total: 0; final: 0; bloques diferidos: 0}
\RouteRow{Tríadas finales; persistencia de Pareto}{17; sí}
\RouteRow{\(K\longrightarrow K+9\)}{repeticiones: 9; cambios de memoria: 9}
\RouteGroup{Firma, fase y diferencias}
\RouteRow{\(n_A,\ s_C,\ \kappa_0,\ \nu_{120},\ \nu_{270}\)}{50, 6, -2, 4, -8}
\RouteRow{\(\tau_{12}\); firma histórica \(\mathbb Z^5\)}{+++-\kern0pt{}-0+++-\kern0pt{}-0; 50|\allowbreak{}6|\allowbreak{}-2|\allowbreak{}4|\allowbreak{}-8}
\RouteRow{\(D_1,\ D_3,\ D_4,\ D_{27},\ D_{36}\)}{24, 24, 24, 0, 0}
\RouteRow{\(\Omega_{90/120},\ P_6\)}{56, 5}
\RouteGroup{Lectores y factores de acción}
\RouteRow{\(K_D\quad /\quad K_\Omega\)}{(0,\allowbreak{}24,\allowbreak{}0)\quad /\quad (56,\allowbreak{}54,\allowbreak{}5)}
\RouteRow{\(\kappa_D\)}{-0.175136\allowbreak{}46166281\allowbreak{}12239348\allowbreak{}425}
\RouteRow{\(\kappa_\Omega\)}{55.60649\allowbreak{}89433721\allowbreak{}35446416\allowbreak{}86}
\RouteRow{\(R_{\mathrm{act},D}\)}{0.999999\allowbreak{}99987858\allowbreak{}10851337\allowbreak{}881}
\RouteRow{\(R_{\mathrm{act},\Omega}\)}{1.000000\allowbreak{}03855097\allowbreak{}23541416\allowbreak{}45}
\RouteGroup{Separación de prefijos}
\RouteRow{Área de ambigüedad; área \(\log_2\)}{108; 46.52932\allowbreak{}50129807\allowbreak{}9}
\RouteRow{Primer bloque de separación}{órbita: \textemdash{}; celda: \textemdash{}; ruta: \textemdash{}}
\RouteRow{Ambigüedad final}{6}
\RouteGroup{Secuencias completas}
\RouteRow{\(w_6\)}{3 0 81 3 0 81 3 0 81 3 0 81 3 0 81 3 0 81}
\RouteRow{Tríadas estables por bloque}{0 1 2 3 4 5 6 7 8 9 10 11 12 13 14 15 16 17}
\end{tabularx}\endgroup\par\vspace{5pt}\noindent{\fontsize{8.4}{10}\selectfont\hyperref[trace-107]{Procedencia y códigos originales}\hfill Ficha completa · 107/324}
\clearpage\phantomsection\label{nav:vi-ruta-108}
% VI-CATALOG-ROW rutas 108
\pdfbookmark[3]{r3c9\_h+1v+1}{route-108}
\RouteTitle{Ruta 108}{r3c9\_\allowbreak{}h+1v+1}
\noindent Celda 27\quad (3, 9)\qquad Orientación \(h=1,\ v=1\)\hfill\hyperref[sec:vi-indice-rutas]{Índice}\par\vspace{4pt}
\begingroup\fontsize{10.2}{12.5}\selectfont\setlength{\tabcolsep}{5pt}\renewcommand{\arraystretch}{1.1}\rowcolors{1}{routepale}{white}\noindent\begin{tabularx}{\linewidth}{@{}>{\raggedright\arraybackslash}p{.345\linewidth}>{\raggedright\arraybackslash}X@{}}
\RouteGroup{Identidad estructural}
\RouteRow{Clase y multisección}{charged\_\allowbreak{}lepton\_\allowbreak{}G4\quad /\quad charged\_\allowbreak{}lepton\_\allowbreak{}G4}
\RouteRow{Colección y sector}{leptón cargado; leptón}
\RouteRow{Clasificación física}{leptón cargado}
\RouteRow{Generación, profundidad y color}{4; G4; \textemdash{}}
\RouteRow{Ruta primaria y órbita de inversión}{\(B_1\): no\qquad 3,\allowbreak{}9;\allowbreak{}7,\allowbreak{}1}
\RouteRow{Firma de clasificación}{30|\allowbreak{}6|\allowbreak{}0|\allowbreak{}2|\allowbreak{}-8|\allowbreak{}+++-\kern0pt{}-0+++-\kern0pt{}-0}
\RouteGroup{Retornos, memoria y transporte}
\RouteRow{Retornos con fase}{clase: 27; órbita: 27; celda: 27}
\RouteRow{Retorno cinemático}{en 108: 54; extensión: 216}
\RouteRow{Giros y cambios de hoja}{71; 107}
\RouteRow{Acarreo y diferimiento}{deuda total: 2; final: 1; bloques diferidos: 2}
\RouteRow{Tríadas finales; persistencia de Pareto}{16; sí}
\RouteRow{\(K\longrightarrow K+9\)}{repeticiones: 9; cambios de memoria: 9}
\RouteGroup{Firma, fase y diferencias}
\RouteRow{\(n_A,\ s_C,\ \kappa_0,\ \nu_{120},\ \nu_{270}\)}{30, 6, -2, 4, -8}
\RouteRow{\(\tau_{12}\); firma histórica \(\mathbb Z^5\)}{+++-\kern0pt{}-0+++-\kern0pt{}-0; 30|\allowbreak{}6|\allowbreak{}-2|\allowbreak{}4|\allowbreak{}-8}
\RouteRow{\(D_1,\ D_3,\ D_4,\ D_{27},\ D_{36}\)}{54, 60, 72, 24, 16}
\RouteRow{\(\Omega_{90/120},\ P_6\)}{56, 5}
\RouteGroup{Lectores y factores de acción}
\RouteRow{\(K_D\quad /\quad K_\Omega\)}{(40,\allowbreak{}54,\allowbreak{}6)\quad /\quad (56,\allowbreak{}54,\allowbreak{}5)}
\RouteRow{\(\kappa_D\)}{39.60661\allowbreak{}01397948\allowbreak{}27586470\allowbreak{}91}
\RouteRow{\(\kappa_\Omega\)}{55.60649\allowbreak{}89433721\allowbreak{}35446416\allowbreak{}86}
\RouteRow{\(R_{\mathrm{act},D}\)}{1.000000\allowbreak{}02745854\allowbreak{}08735595\allowbreak{}18}
\RouteRow{\(R_{\mathrm{act},\Omega}\)}{1.000000\allowbreak{}03855097\allowbreak{}23541416\allowbreak{}45}
\RouteGroup{Separación de prefijos}
\RouteRow{Área de ambigüedad; área \(\log_2\)}{58; 30.18947\allowbreak{}50100961\allowbreak{}74}
\RouteRow{Primer bloque de separación}{órbita: \textemdash{}; celda: \textemdash{}; ruta: \textemdash{}}
\RouteRow{Ambigüedad final}{3}
\RouteGroup{Secuencias completas}
\RouteRow{\(w_6\)}{24 243 657 24 243 423 15 486 423 24 243 657 24 243 423 15 486 423}
\RouteRow{Tríadas estables por bloque}{0 1 2 3 4 5 6 7 8 9 10 11 12 13 13 15 16 16}
\end{tabularx}\endgroup\par\vspace{5pt}\noindent{\fontsize{8.4}{10}\selectfont\hyperref[trace-108]{Procedencia y códigos originales}\hfill Ficha completa · 108/324}
\clearpage\phantomsection\label{nav:vi-ruta-109}
% VI-CATALOG-ROW rutas 109
\pdfbookmark[2]{Celda 28}{cell-28}
\pdfbookmark[3]{r4c1\_h-1v-1}{route-109}
\RouteTitle{Ruta 109}{r4c1\_\allowbreak{}h-1v-1}
\noindent Celda 28\quad (4, 1)\qquad Orientación \(h=-1,\ v=-1\)\hfill\hyperref[sec:vi-indice-rutas]{Índice}\par\vspace{4pt}
\begingroup\fontsize{10.2}{12.5}\selectfont\setlength{\tabcolsep}{5pt}\renewcommand{\arraystretch}{1.1}\rowcolors{1}{routepale}{white}\noindent\begin{tabularx}{\linewidth}{@{}>{\raggedright\arraybackslash}p{.345\linewidth}>{\raggedright\arraybackslash}X@{}}
\RouteGroup{Identidad estructural}
\RouteRow{Clase y multisección}{quark\_\allowbreak{}down\_\allowbreak{}G4\quad /\quad quark\_\allowbreak{}down\_\allowbreak{}G4}
\RouteRow{Colección y sector}{quark de tipo abajo; quark}
\RouteRow{Clasificación física}{quark de tipo abajo}
\RouteRow{Generación, profundidad y color}{4; G4; R}
\RouteRow{Ruta primaria y órbita de inversión}{\(B_1\): sí\qquad 4,\allowbreak{}1;\allowbreak{}6,\allowbreak{}9}
\RouteRow{Firma de clasificación}{30|\allowbreak{}6|\allowbreak{}8|\allowbreak{}0|\allowbreak{}-12|\allowbreak{}+++-\kern0pt{}-\kern0pt{}-+++-\kern0pt{}-\kern0pt{}-}
\RouteGroup{Retornos, memoria y transporte}
\RouteRow{Retornos con fase}{clase: 27; órbita: 27; celda: 27}
\RouteRow{Retorno cinemático}{en 108: 54; extensión: 216}
\RouteRow{Giros y cambios de hoja}{71; 107}
\RouteRow{Acarreo y diferimiento}{deuda total: 2; final: 1; bloques diferidos: 2}
\RouteRow{Tríadas finales; persistencia de Pareto}{16; sí}
\RouteRow{\(K\longrightarrow K+9\)}{repeticiones: 9; cambios de memoria: 9}
\RouteGroup{Firma, fase y diferencias}
\RouteRow{\(n_A,\ s_C,\ \kappa_0,\ \nu_{120},\ \nu_{270}\)}{30, -6, 0, -4, -12}
\RouteRow{\(\tau_{12}\); firma histórica \(\mathbb Z^5\)}{+++-\kern0pt{}-\kern0pt{}-+++-\kern0pt{}-\kern0pt{}-; 30|\allowbreak{}-6|\allowbreak{}0|\allowbreak{}-4|\allowbreak{}-12}
\RouteRow{\(D_1,\ D_3,\ D_4,\ D_{27},\ D_{36}\)}{54, 56, 70, 48, 32}
\RouteRow{\(\Omega_{90/120},\ P_6\)}{80, 6}
\RouteGroup{Lectores y factores de acción}
\RouteRow{\(K_D\quad /\quad K_\Omega\)}{(80,\allowbreak{}54,\allowbreak{}7)\quad /\quad (80,\allowbreak{}54,\allowbreak{}6)}
\RouteRow{\(\kappa_D\)}{79.60672\allowbreak{}13362175\allowbreak{}19726524\allowbreak{}96}
\RouteRow{\(\kappa_\Omega\)}{79.60661\allowbreak{}01397948\allowbreak{}27586470\allowbreak{}91}
\RouteRow{\(R_{\mathrm{act},D}\)}{1.000000\allowbreak{}05518988\allowbreak{}96223153\allowbreak{}37}
\RouteRow{\(R_{\mathrm{act},\Omega}\)}{1.000000\allowbreak{}05518981\allowbreak{}25318591\allowbreak{}40}
\RouteGroup{Separación de prefijos}
\RouteRow{Área de ambigüedad; área \(\log_2\)}{18; 0.0}
\RouteRow{Primer bloque de separación}{órbita: 1; celda: 1; ruta: 1}
\RouteRow{Ambigüedad final}{1}
\RouteGroup{Secuencias completas}
\RouteRow{\(w_6\)}{486 423 15 486 429 243 657 24 243 486 423 15 486 429 243 657 24 243}
\RouteRow{Tríadas estables por bloque}{0 1 2 3 4 5 6 7 8 9 10 11 12 13 13 15 16 16}
\end{tabularx}\endgroup\par\vspace{5pt}\noindent{\fontsize{8.4}{10}\selectfont\hyperref[trace-109]{Procedencia y códigos originales}\hfill Ficha completa · 109/324}
\clearpage\phantomsection\label{nav:vi-ruta-110}
% VI-CATALOG-ROW rutas 110
\pdfbookmark[3]{r4c1\_h-1v+1}{route-110}
\RouteTitle{Ruta 110}{r4c1\_\allowbreak{}h-1v+1}
\noindent Celda 28\quad (4, 1)\qquad Orientación \(h=-1,\ v=1\)\hfill\hyperref[sec:vi-indice-rutas]{Índice}\par\vspace{4pt}
\begingroup\fontsize{10.2}{12.5}\selectfont\setlength{\tabcolsep}{5pt}\renewcommand{\arraystretch}{1.1}\rowcolors{1}{routepale}{white}\noindent\begin{tabularx}{\linewidth}{@{}>{\raggedright\arraybackslash}p{.345\linewidth}>{\raggedright\arraybackslash}X@{}}
\RouteGroup{Identidad estructural}
\RouteRow{Clase y multisección}{quark\_\allowbreak{}down\_\allowbreak{}G4\quad /\quad quark\_\allowbreak{}down\_\allowbreak{}G4}
\RouteRow{Colección y sector}{quark de tipo abajo; quark}
\RouteRow{Clasificación física}{quark de tipo abajo}
\RouteRow{Generación, profundidad y color}{4; G4; R}
\RouteRow{Ruta primaria y órbita de inversión}{\(B_1\): no\qquad 4,\allowbreak{}1;\allowbreak{}6,\allowbreak{}9}
\RouteRow{Firma de clasificación}{30|\allowbreak{}6|\allowbreak{}8|\allowbreak{}0|\allowbreak{}-12|\allowbreak{}+++-\kern0pt{}-\kern0pt{}-+++-\kern0pt{}-\kern0pt{}-}
\RouteGroup{Retornos, memoria y transporte}
\RouteRow{Retornos con fase}{clase: 27; órbita: 27; celda: 27}
\RouteRow{Retorno cinemático}{en 108: 54; extensión: 216}
\RouteRow{Giros y cambios de hoja}{71; 107}
\RouteRow{Acarreo y diferimiento}{deuda total: 0; final: 0; bloques diferidos: 0}
\RouteRow{Tríadas finales; persistencia de Pareto}{17; sí}
\RouteRow{\(K\longrightarrow K+9\)}{repeticiones: 9; cambios de memoria: 9}
\RouteGroup{Firma, fase y diferencias}
\RouteRow{\(n_A,\ s_C,\ \kappa_0,\ \nu_{120},\ \nu_{270}\)}{44, 2, 0, 0, -12}
\RouteRow{\(\tau_{12}\); firma histórica \(\mathbb Z^5\)}{+++-\kern0pt{}-\kern0pt{}-+++-\kern0pt{}-\kern0pt{}-; 44|\allowbreak{}2|\allowbreak{}0|\allowbreak{}0|\allowbreak{}-12}
\RouteRow{\(D_1,\ D_3,\ D_4,\ D_{27},\ D_{36}\)}{78, 28, 78, 36, 24}
\RouteRow{\(\Omega_{90/120},\ P_6\)}{80, 6}
\RouteGroup{Lectores y factores de acción}
\RouteRow{\(K_D\quad /\quad K_\Omega\)}{(60,\allowbreak{}78,\allowbreak{}25)\quad /\quad (80,\allowbreak{}54,\allowbreak{}6)}
\RouteRow{\(\kappa_D\)}{59.43358\allowbreak{}64101631\allowbreak{}67023563\allowbreak{}04}
\RouteRow{\(\kappa_\Omega\)}{79.60661\allowbreak{}01397948\allowbreak{}27586470\allowbreak{}91}
\RouteRow{\(R_{\mathrm{act},D}\)}{1.000000\allowbreak{}04120422\allowbreak{}24053449\allowbreak{}05}
\RouteRow{\(R_{\mathrm{act},\Omega}\)}{1.000000\allowbreak{}05518981\allowbreak{}25318591\allowbreak{}40}
\RouteGroup{Separación de prefijos}
\RouteRow{Área de ambigüedad; área \(\log_2\)}{18; 0.0}
\RouteRow{Primer bloque de separación}{órbita: 1; celda: 1; ruta: 1}
\RouteRow{Ambigüedad final}{1}
\RouteGroup{Secuencias completas}
\RouteRow{\(w_6\)}{507 504 585 507 510 510 666 672 510 507 504 585 507 510 510 666 672 510}
\RouteRow{Tríadas estables por bloque}{0 1 2 3 4 5 6 7 8 9 10 11 12 13 14 15 16 17}
\end{tabularx}\endgroup\par\vspace{5pt}\noindent{\fontsize{8.4}{10}\selectfont\hyperref[trace-110]{Procedencia y códigos originales}\hfill Ficha completa · 110/324}
\clearpage\phantomsection\label{nav:vi-ruta-111}
% VI-CATALOG-ROW rutas 111
\pdfbookmark[3]{r4c1\_h+1v-1}{route-111}
\RouteTitle{Ruta 111}{r4c1\_\allowbreak{}h+1v-1}
\noindent Celda 28\quad (4, 1)\qquad Orientación \(h=1,\ v=-1\)\hfill\hyperref[sec:vi-indice-rutas]{Índice}\par\vspace{4pt}
\begingroup\fontsize{10.2}{12.5}\selectfont\setlength{\tabcolsep}{5pt}\renewcommand{\arraystretch}{1.1}\rowcolors{1}{routepale}{white}\noindent\begin{tabularx}{\linewidth}{@{}>{\raggedright\arraybackslash}p{.345\linewidth}>{\raggedright\arraybackslash}X@{}}
\RouteGroup{Identidad estructural}
\RouteRow{Clase y multisección}{quark\_\allowbreak{}down\_\allowbreak{}G4\quad /\quad quark\_\allowbreak{}down\_\allowbreak{}G4}
\RouteRow{Colección y sector}{quark de tipo abajo; quark}
\RouteRow{Clasificación física}{quark de tipo abajo}
\RouteRow{Generación, profundidad y color}{4; G4; R}
\RouteRow{Ruta primaria y órbita de inversión}{\(B_1\): no\qquad 4,\allowbreak{}1;\allowbreak{}6,\allowbreak{}9}
\RouteRow{Firma de clasificación}{30|\allowbreak{}6|\allowbreak{}8|\allowbreak{}0|\allowbreak{}-12|\allowbreak{}+++-\kern0pt{}-\kern0pt{}-+++-\kern0pt{}-\kern0pt{}-}
\RouteGroup{Retornos, memoria y transporte}
\RouteRow{Retornos con fase}{clase: 27; órbita: 27; celda: 27}
\RouteRow{Retorno cinemático}{en 108: 54; extensión: 216}
\RouteRow{Giros y cambios de hoja}{71; 107}
\RouteRow{Acarreo y diferimiento}{deuda total: 2; final: 0; bloques diferidos: 1}
\RouteRow{Tríadas finales; persistencia de Pareto}{17; sí}
\RouteRow{\(K\longrightarrow K+9\)}{repeticiones: 9; cambios de memoria: 9}
\RouteGroup{Firma, fase y diferencias}
\RouteRow{\(n_A,\ s_C,\ \kappa_0,\ \nu_{120},\ \nu_{270}\)}{44, -2, 0, 0, -12}
\RouteRow{\(\tau_{12}\); firma histórica \(\mathbb Z^5\)}{+++-\kern0pt{}-\kern0pt{}-+++-\kern0pt{}-\kern0pt{}-; 44|\allowbreak{}-2|\allowbreak{}0|\allowbreak{}0|\allowbreak{}-12}
\RouteRow{\(D_1,\ D_3,\ D_4,\ D_{27},\ D_{36}\)}{78, 28, 78, 36, 24}
\RouteRow{\(\Omega_{90/120},\ P_6\)}{80, 6}
\RouteGroup{Lectores y factores de acción}
\RouteRow{\(K_D\quad /\quad K_\Omega\)}{(60,\allowbreak{}78,\allowbreak{}25)\quad /\quad (80,\allowbreak{}54,\allowbreak{}6)}
\RouteRow{\(\kappa_D\)}{59.43358\allowbreak{}64101631\allowbreak{}67023563\allowbreak{}04}
\RouteRow{\(\kappa_\Omega\)}{79.60661\allowbreak{}01397948\allowbreak{}27586470\allowbreak{}91}
\RouteRow{\(R_{\mathrm{act},D}\)}{1.000000\allowbreak{}04120422\allowbreak{}24053449\allowbreak{}05}
\RouteRow{\(R_{\mathrm{act},\Omega}\)}{1.000000\allowbreak{}05518981\allowbreak{}25318591\allowbreak{}40}
\RouteGroup{Separación de prefijos}
\RouteRow{Área de ambigüedad; área \(\log_2\)}{18; 0.0}
\RouteRow{Primer bloque de separación}{órbita: 1; celda: 1; ruta: 1}
\RouteRow{Ambigüedad final}{1}
\RouteGroup{Secuencias completas}
\RouteRow{\(w_6\)}{666 672 510 666 666 585 507 504 585 666 672 510 666 666 585 507 504 585}
\RouteRow{Tríadas estables por bloque}{0 1 2 3 4 5 6 7 8 9 10 11 12 12 13 15 16 17}
\end{tabularx}\endgroup\par\vspace{5pt}\noindent{\fontsize{8.4}{10}\selectfont\hyperref[trace-111]{Procedencia y códigos originales}\hfill Ficha completa · 111/324}
\clearpage\phantomsection\label{nav:vi-ruta-112}
% VI-CATALOG-ROW rutas 112
\pdfbookmark[3]{r4c1\_h+1v+1}{route-112}
\RouteTitle{Ruta 112}{r4c1\_\allowbreak{}h+1v+1}
\noindent Celda 28\quad (4, 1)\qquad Orientación \(h=1,\ v=1\)\hfill\hyperref[sec:vi-indice-rutas]{Índice}\par\vspace{4pt}
\begingroup\fontsize{10.2}{12.5}\selectfont\setlength{\tabcolsep}{5pt}\renewcommand{\arraystretch}{1.1}\rowcolors{1}{routepale}{white}\noindent\begin{tabularx}{\linewidth}{@{}>{\raggedright\arraybackslash}p{.345\linewidth}>{\raggedright\arraybackslash}X@{}}
\RouteGroup{Identidad estructural}
\RouteRow{Clase y multisección}{quark\_\allowbreak{}down\_\allowbreak{}G4\quad /\quad quark\_\allowbreak{}down\_\allowbreak{}G4}
\RouteRow{Colección y sector}{quark de tipo abajo; quark}
\RouteRow{Clasificación física}{quark de tipo abajo}
\RouteRow{Generación, profundidad y color}{4; G4; R}
\RouteRow{Ruta primaria y órbita de inversión}{\(B_1\): no\qquad 4,\allowbreak{}1;\allowbreak{}6,\allowbreak{}9}
\RouteRow{Firma de clasificación}{30|\allowbreak{}6|\allowbreak{}8|\allowbreak{}0|\allowbreak{}-12|\allowbreak{}+++-\kern0pt{}-\kern0pt{}-+++-\kern0pt{}-\kern0pt{}-}
\RouteGroup{Retornos, memoria y transporte}
\RouteRow{Retornos con fase}{clase: 27; órbita: 27; celda: 27}
\RouteRow{Retorno cinemático}{en 108: 54; extensión: 216}
\RouteRow{Giros y cambios de hoja}{71; 107}
\RouteRow{Acarreo y diferimiento}{deuda total: 2; final: 0; bloques diferidos: 1}
\RouteRow{Tríadas finales; persistencia de Pareto}{17; sí}
\RouteRow{\(K\longrightarrow K+9\)}{repeticiones: 9; cambios de memoria: 9}
\RouteGroup{Firma, fase y diferencias}
\RouteRow{\(n_A,\ s_C,\ \kappa_0,\ \nu_{120},\ \nu_{270}\)}{30, 6, 0, 4, -12}
\RouteRow{\(\tau_{12}\); firma histórica \(\mathbb Z^5\)}{+++-\kern0pt{}-\kern0pt{}-+++-\kern0pt{}-\kern0pt{}-; 30|\allowbreak{}6|\allowbreak{}0|\allowbreak{}4|\allowbreak{}-12}
\RouteRow{\(D_1,\ D_3,\ D_4,\ D_{27},\ D_{36}\)}{54, 56, 70, 48, 32}
\RouteRow{\(\Omega_{90/120},\ P_6\)}{80, 6}
\RouteGroup{Lectores y factores de acción}
\RouteRow{\(K_D\quad /\quad K_\Omega\)}{(80,\allowbreak{}54,\allowbreak{}7)\quad /\quad (80,\allowbreak{}54,\allowbreak{}6)}
\RouteRow{\(\kappa_D\)}{79.60672\allowbreak{}13362175\allowbreak{}19726524\allowbreak{}96}
\RouteRow{\(\kappa_\Omega\)}{79.60661\allowbreak{}01397948\allowbreak{}27586470\allowbreak{}91}
\RouteRow{\(R_{\mathrm{act},D}\)}{1.000000\allowbreak{}05518988\allowbreak{}96223153\allowbreak{}37}
\RouteRow{\(R_{\mathrm{act},\Omega}\)}{1.000000\allowbreak{}05518981\allowbreak{}25318591\allowbreak{}40}
\RouteGroup{Separación de prefijos}
\RouteRow{Área de ambigüedad; área \(\log_2\)}{26; 6.339850\allowbreak{}00288462\allowbreak{}4}
\RouteRow{Primer bloque de separación}{órbita: 5; celda: 5; ruta: 5}
\RouteRow{Ambigüedad final}{1}
\RouteGroup{Secuencias completas}
\RouteRow{\(w_6\)}{657 24 243 657 18 15 486 423 15 657 24 243 657 18 15 486 423 15}
\RouteRow{Tríadas estables por bloque}{0 1 2 3 4 5 6 7 8 9 10 11 12 13 14 14 15 17}
\end{tabularx}\endgroup\par\vspace{5pt}\noindent{\fontsize{8.4}{10}\selectfont\hyperref[trace-112]{Procedencia y códigos originales}\hfill Ficha completa · 112/324}
\clearpage\phantomsection\label{nav:vi-ruta-113}
% VI-CATALOG-ROW rutas 113
\pdfbookmark[2]{Celda 29}{cell-29}
\pdfbookmark[3]{r4c2\_h-1v-1}{route-113}
\RouteTitle{Ruta 113}{r4c2\_\allowbreak{}h-1v-1}
\noindent Celda 29\quad (4, 2)\qquad Orientación \(h=-1,\ v=-1\)\hfill\hyperref[sec:vi-indice-rutas]{Índice}\par\vspace{4pt}
\begingroup\fontsize{10.2}{12.5}\selectfont\setlength{\tabcolsep}{5pt}\renewcommand{\arraystretch}{1.1}\rowcolors{1}{routepale}{white}\noindent\begin{tabularx}{\linewidth}{@{}>{\raggedright\arraybackslash}p{.345\linewidth}>{\raggedright\arraybackslash}X@{}}
\RouteGroup{Identidad estructural}
\RouteRow{Clase y multisección}{quark\_\allowbreak{}up\_\allowbreak{}G3\quad /\quad quark\_\allowbreak{}up\_\allowbreak{}G3}
\RouteRow{Colección y sector}{quark de tipo arriba; quark}
\RouteRow{Clasificación física}{quark de tipo arriba}
\RouteRow{Generación, profundidad y color}{3; G3; B}
\RouteRow{Ruta primaria y órbita de inversión}{\(B_1\): no\qquad 4,\allowbreak{}2;\allowbreak{}6,\allowbreak{}8}
\RouteRow{Firma de clasificación}{28|\allowbreak{}6|\allowbreak{}-2|\allowbreak{}0|\allowbreak{}-12|\allowbreak{}+++-\kern0pt{}-\kern0pt{}-+++-\kern0pt{}-\kern0pt{}-}
\RouteGroup{Retornos, memoria y transporte}
\RouteRow{Retornos con fase}{clase: 27; órbita: 27; celda: 27}
\RouteRow{Retorno cinemático}{en 108: 54; extensión: 216}
\RouteRow{Giros y cambios de hoja}{71; 107}
\RouteRow{Acarreo y diferimiento}{deuda total: 1; final: 0; bloques diferidos: 1}
\RouteRow{Tríadas finales; persistencia de Pareto}{17; sí}
\RouteRow{\(K\longrightarrow K+9\)}{repeticiones: 9; cambios de memoria: 9}
\RouteGroup{Firma, fase y diferencias}
\RouteRow{\(n_A,\ s_C,\ \kappa_0,\ \nu_{120},\ \nu_{270}\)}{28, -6, 0, 0, -12}
\RouteRow{\(\tau_{12}\); firma histórica \(\mathbb Z^5\)}{+++-\kern0pt{}-\kern0pt{}-+++-\kern0pt{}-\kern0pt{}-; 28|\allowbreak{}-6|\allowbreak{}0|\allowbreak{}0|\allowbreak{}-12}
\RouteRow{\(D_1,\ D_3,\ D_4,\ D_{27},\ D_{36}\)}{66, 64, 64, 60, 40}
\RouteRow{\(\Omega_{90/120},\ P_6\)}{80, 6}
\RouteGroup{Lectores y factores de acción}
\RouteRow{\(K_D\quad /\quad K_\Omega\)}{(100,\allowbreak{}66,\allowbreak{}0)\quad /\quad (80,\allowbreak{}54,\allowbreak{}6)}
\RouteRow{\(\kappa_D\)}{99.51837\allowbreak{}47304272\allowbreak{}69134179\allowbreak{}18}
\RouteRow{\(\kappa_\Omega\)}{79.60661\allowbreak{}01397948\allowbreak{}27586470\allowbreak{}91}
\RouteRow{\(R_{\mathrm{act},D}\)}{1.000000\allowbreak{}06899427\allowbreak{}66450218\allowbreak{}95}
\RouteRow{\(R_{\mathrm{act},\Omega}\)}{1.000000\allowbreak{}05518981\allowbreak{}25318591\allowbreak{}40}
\RouteGroup{Separación de prefijos}
\RouteRow{Área de ambigüedad; área \(\log_2\)}{36; 18.0}
\RouteRow{Primer bloque de separación}{órbita: \textemdash{}; celda: \textemdash{}; ruta: \textemdash{}}
\RouteRow{Ambigüedad final}{2}
\RouteGroup{Secuencias completas}
\RouteRow{\(w_6\)}{90 570 264 90 567 657 24 243 657 90 570 264 90 567 657 24 243 657}
\RouteRow{Tríadas estables por bloque}{0 1 2 3 4 5 6 7 8 9 10 11 12 13 13 15 16 17}
\end{tabularx}\endgroup\par\vspace{5pt}\noindent{\fontsize{8.4}{10}\selectfont\hyperref[trace-113]{Procedencia y códigos originales}\hfill Ficha completa · 113/324}
\clearpage\phantomsection\label{nav:vi-ruta-114}
% VI-CATALOG-ROW rutas 114
\pdfbookmark[3]{r4c2\_h-1v+1}{route-114}
\RouteTitle{Ruta 114}{r4c2\_\allowbreak{}h-1v+1}
\noindent Celda 29\quad (4, 2)\qquad Orientación \(h=-1,\ v=1\)\hfill\hyperref[sec:vi-indice-rutas]{Índice}\par\vspace{4pt}
\begingroup\fontsize{10.2}{12.5}\selectfont\setlength{\tabcolsep}{5pt}\renewcommand{\arraystretch}{1.1}\rowcolors{1}{routepale}{white}\noindent\begin{tabularx}{\linewidth}{@{}>{\raggedright\arraybackslash}p{.345\linewidth}>{\raggedright\arraybackslash}X@{}}
\RouteGroup{Identidad estructural}
\RouteRow{Clase y multisección}{quark\_\allowbreak{}up\_\allowbreak{}G3\quad /\quad quark\_\allowbreak{}up\_\allowbreak{}G3}
\RouteRow{Colección y sector}{quark de tipo arriba; quark}
\RouteRow{Clasificación física}{quark de tipo arriba}
\RouteRow{Generación, profundidad y color}{3; G3; B}
\RouteRow{Ruta primaria y órbita de inversión}{\(B_1\): no\qquad 4,\allowbreak{}2;\allowbreak{}6,\allowbreak{}8}
\RouteRow{Firma de clasificación}{28|\allowbreak{}6|\allowbreak{}-2|\allowbreak{}0|\allowbreak{}-12|\allowbreak{}+++-\kern0pt{}-\kern0pt{}-+++-\kern0pt{}-\kern0pt{}-}
\RouteGroup{Retornos, memoria y transporte}
\RouteRow{Retornos con fase}{clase: 27; órbita: 27; celda: 27}
\RouteRow{Retorno cinemático}{en 108: 54; extensión: 216}
\RouteRow{Giros y cambios de hoja}{71; 107}
\RouteRow{Acarreo y diferimiento}{deuda total: 2; final: 0; bloques diferidos: 2}
\RouteRow{Tríadas finales; persistencia de Pareto}{17; sí}
\RouteRow{\(K\longrightarrow K+9\)}{repeticiones: 9; cambios de memoria: 9}
\RouteGroup{Firma, fase y diferencias}
\RouteRow{\(n_A,\ s_C,\ \kappa_0,\ \nu_{120},\ \nu_{270}\)}{14, 6, 0, 0, -12}
\RouteRow{\(\tau_{12}\); firma histórica \(\mathbb Z^5\)}{+++-\kern0pt{}-\kern0pt{}-+++-\kern0pt{}-\kern0pt{}-; 14|\allowbreak{}6|\allowbreak{}0|\allowbreak{}0|\allowbreak{}-12}
\RouteRow{\(D_1,\ D_3,\ D_4,\ D_{27},\ D_{36}\)}{24, 20, 24, 24, 16}
\RouteRow{\(\Omega_{90/120},\ P_6\)}{80, 6}
\RouteGroup{Lectores y factores de acción}
\RouteRow{\(K_D\quad /\quad K_\Omega\)}{(40,\allowbreak{}24,\allowbreak{}2)\quad /\quad (80,\allowbreak{}54,\allowbreak{}6)}
\RouteRow{\(\kappa_D\)}{39.82508\allowbreak{}59311825\allowbreak{}73056173\allowbreak{}26}
\RouteRow{\(\kappa_\Omega\)}{79.60661\allowbreak{}01397948\allowbreak{}27586470\allowbreak{}91}
\RouteRow{\(R_{\mathrm{act},D}\)}{1.000000\allowbreak{}02761000\allowbreak{}61595142\allowbreak{}15}
\RouteRow{\(R_{\mathrm{act},\Omega}\)}{1.000000\allowbreak{}05518981\allowbreak{}25318591\allowbreak{}40}
\RouteGroup{Separación de prefijos}
\RouteRow{Área de ambigüedad; área \(\log_2\)}{72; 36.0}
\RouteRow{Primer bloque de separación}{órbita: \textemdash{}; celda: \textemdash{}; ruta: \textemdash{}}
\RouteRow{Ambigüedad final}{4}
\RouteGroup{Secuencias completas}
\RouteRow{\(w_6\)}{81 3 0 81 0 3 0 81 3 81 3 0 81 0 3 0 81 3}
\RouteRow{Tríadas estables por bloque}{0 1 2 3 4 5 6 6 8 9 10 11 12 13 14 15 15 17}
\end{tabularx}\endgroup\par\vspace{5pt}\noindent{\fontsize{8.4}{10}\selectfont\hyperref[trace-114]{Procedencia y códigos originales}\hfill Ficha completa · 114/324}
\clearpage\phantomsection\label{nav:vi-ruta-115}
% VI-CATALOG-ROW rutas 115
\pdfbookmark[3]{r4c2\_h+1v-1}{route-115}
\RouteTitle{Ruta 115}{r4c2\_\allowbreak{}h+1v-1}
\noindent Celda 29\quad (4, 2)\qquad Orientación \(h=1,\ v=-1\)\hfill\hyperref[sec:vi-indice-rutas]{Índice}\par\vspace{4pt}
\begingroup\fontsize{10.2}{12.5}\selectfont\setlength{\tabcolsep}{5pt}\renewcommand{\arraystretch}{1.1}\rowcolors{1}{routepale}{white}\noindent\begin{tabularx}{\linewidth}{@{}>{\raggedright\arraybackslash}p{.345\linewidth}>{\raggedright\arraybackslash}X@{}}
\RouteGroup{Identidad estructural}
\RouteRow{Clase y multisección}{quark\_\allowbreak{}up\_\allowbreak{}G3\quad /\quad quark\_\allowbreak{}up\_\allowbreak{}G3}
\RouteRow{Colección y sector}{quark de tipo arriba; quark}
\RouteRow{Clasificación física}{quark de tipo arriba}
\RouteRow{Generación, profundidad y color}{3; G3; B}
\RouteRow{Ruta primaria y órbita de inversión}{\(B_1\): sí\qquad 4,\allowbreak{}2;\allowbreak{}6,\allowbreak{}8}
\RouteRow{Firma de clasificación}{28|\allowbreak{}6|\allowbreak{}-2|\allowbreak{}0|\allowbreak{}-12|\allowbreak{}+++-\kern0pt{}-\kern0pt{}-+++-\kern0pt{}-\kern0pt{}-}
\RouteGroup{Retornos, memoria y transporte}
\RouteRow{Retornos con fase}{clase: 27; órbita: 27; celda: 27}
\RouteRow{Retorno cinemático}{en 108: 54; extensión: 216}
\RouteRow{Giros y cambios de hoja}{71; 107}
\RouteRow{Acarreo y diferimiento}{deuda total: 0; final: 0; bloques diferidos: 0}
\RouteRow{Tríadas finales; persistencia de Pareto}{17; sí}
\RouteRow{\(K\longrightarrow K+9\)}{repeticiones: 9; cambios de memoria: 9}
\RouteGroup{Firma, fase y diferencias}
\RouteRow{\(n_A,\ s_C,\ \kappa_0,\ \nu_{120},\ \nu_{270}\)}{14, -6, 0, 0, -12}
\RouteRow{\(\tau_{12}\); firma histórica \(\mathbb Z^5\)}{+++-\kern0pt{}-\kern0pt{}-+++-\kern0pt{}-\kern0pt{}-; 14|\allowbreak{}-6|\allowbreak{}0|\allowbreak{}0|\allowbreak{}-12}
\RouteRow{\(D_1,\ D_3,\ D_4,\ D_{27},\ D_{36}\)}{24, 20, 24, 24, 16}
\RouteRow{\(\Omega_{90/120},\ P_6\)}{80, 6}
\RouteGroup{Lectores y factores de acción}
\RouteRow{\(K_D\quad /\quad K_\Omega\)}{(40,\allowbreak{}24,\allowbreak{}2)\quad /\quad (80,\allowbreak{}54,\allowbreak{}6)}
\RouteRow{\(\kappa_D\)}{39.82508\allowbreak{}59311825\allowbreak{}73056173\allowbreak{}26}
\RouteRow{\(\kappa_\Omega\)}{79.60661\allowbreak{}01397948\allowbreak{}27586470\allowbreak{}91}
\RouteRow{\(R_{\mathrm{act},D}\)}{1.000000\allowbreak{}02761000\allowbreak{}61595142\allowbreak{}15}
\RouteRow{\(R_{\mathrm{act},\Omega}\)}{1.000000\allowbreak{}05518981\allowbreak{}25318591\allowbreak{}40}
\RouteGroup{Separación de prefijos}
\RouteRow{Área de ambigüedad; área \(\log_2\)}{72; 36.0}
\RouteRow{Primer bloque de separación}{órbita: \textemdash{}; celda: \textemdash{}; ruta: \textemdash{}}
\RouteRow{Ambigüedad final}{4}
\RouteGroup{Secuencias completas}
\RouteRow{\(w_6\)}{0 81 3 0 84 0 81 3 0 0 81 3 0 84 0 81 3 0}
\RouteRow{Tríadas estables por bloque}{0 1 2 3 4 5 6 7 8 9 10 11 12 13 14 15 16 17}
\end{tabularx}\endgroup\par\vspace{5pt}\noindent{\fontsize{8.4}{10}\selectfont\hyperref[trace-115]{Procedencia y códigos originales}\hfill Ficha completa · 115/324}
\clearpage\phantomsection\label{nav:vi-ruta-116}
% VI-CATALOG-ROW rutas 116
\pdfbookmark[3]{r4c2\_h+1v+1}{route-116}
\RouteTitle{Ruta 116}{r4c2\_\allowbreak{}h+1v+1}
\noindent Celda 29\quad (4, 2)\qquad Orientación \(h=1,\ v=1\)\hfill\hyperref[sec:vi-indice-rutas]{Índice}\par\vspace{4pt}
\begingroup\fontsize{10.2}{12.5}\selectfont\setlength{\tabcolsep}{5pt}\renewcommand{\arraystretch}{1.1}\rowcolors{1}{routepale}{white}\noindent\begin{tabularx}{\linewidth}{@{}>{\raggedright\arraybackslash}p{.345\linewidth}>{\raggedright\arraybackslash}X@{}}
\RouteGroup{Identidad estructural}
\RouteRow{Clase y multisección}{quark\_\allowbreak{}up\_\allowbreak{}G3\quad /\quad quark\_\allowbreak{}up\_\allowbreak{}G3}
\RouteRow{Colección y sector}{quark de tipo arriba; quark}
\RouteRow{Clasificación física}{quark de tipo arriba}
\RouteRow{Generación, profundidad y color}{3; G3; B}
\RouteRow{Ruta primaria y órbita de inversión}{\(B_1\): no\qquad 4,\allowbreak{}2;\allowbreak{}6,\allowbreak{}8}
\RouteRow{Firma de clasificación}{28|\allowbreak{}6|\allowbreak{}-2|\allowbreak{}0|\allowbreak{}-12|\allowbreak{}+++-\kern0pt{}-\kern0pt{}-+++-\kern0pt{}-\kern0pt{}-}
\RouteGroup{Retornos, memoria y transporte}
\RouteRow{Retornos con fase}{clase: 27; órbita: 27; celda: 27}
\RouteRow{Retorno cinemático}{en 108: 54; extensión: 216}
\RouteRow{Giros y cambios de hoja}{71; 107}
\RouteRow{Acarreo y diferimiento}{deuda total: 0; final: 0; bloques diferidos: 0}
\RouteRow{Tríadas finales; persistencia de Pareto}{17; sí}
\RouteRow{\(K\longrightarrow K+9\)}{repeticiones: 9; cambios de memoria: 9}
\RouteGroup{Firma, fase y diferencias}
\RouteRow{\(n_A,\ s_C,\ \kappa_0,\ \nu_{120},\ \nu_{270}\)}{28, 6, 0, 0, -12}
\RouteRow{\(\tau_{12}\); firma histórica \(\mathbb Z^5\)}{+++-\kern0pt{}-\kern0pt{}-+++-\kern0pt{}-\kern0pt{}-; 28|\allowbreak{}6|\allowbreak{}0|\allowbreak{}0|\allowbreak{}-12}
\RouteRow{\(D_1,\ D_3,\ D_4,\ D_{27},\ D_{36}\)}{66, 64, 64, 60, 40}
\RouteRow{\(\Omega_{90/120},\ P_6\)}{80, 6}
\RouteGroup{Lectores y factores de acción}
\RouteRow{\(K_D\quad /\quad K_\Omega\)}{(100,\allowbreak{}66,\allowbreak{}0)\quad /\quad (80,\allowbreak{}54,\allowbreak{}6)}
\RouteRow{\(\kappa_D\)}{99.51837\allowbreak{}47304272\allowbreak{}69134179\allowbreak{}18}
\RouteRow{\(\kappa_\Omega\)}{79.60661\allowbreak{}01397948\allowbreak{}27586470\allowbreak{}91}
\RouteRow{\(R_{\mathrm{act},D}\)}{1.000000\allowbreak{}06899427\allowbreak{}66450218\allowbreak{}95}
\RouteRow{\(R_{\mathrm{act},\Omega}\)}{1.000000\allowbreak{}05518981\allowbreak{}25318591\allowbreak{}40}
\RouteGroup{Separación de prefijos}
\RouteRow{Área de ambigüedad; área \(\log_2\)}{36; 18.0}
\RouteRow{Primer bloque de separación}{órbita: \textemdash{}; celda: \textemdash{}; ruta: \textemdash{}}
\RouteRow{Ambigüedad final}{2}
\RouteGroup{Secuencias completas}
\RouteRow{\(w_6\)}{24 243 657 24 246 264 90 570 264 24 243 657 24 246 264 90 570 264}
\RouteRow{Tríadas estables por bloque}{0 1 2 3 4 5 6 7 8 9 10 11 12 13 14 15 16 17}
\end{tabularx}\endgroup\par\vspace{5pt}\noindent{\fontsize{8.4}{10}\selectfont\hyperref[trace-116]{Procedencia y códigos originales}\hfill Ficha completa · 116/324}
\clearpage\phantomsection\label{nav:vi-ruta-117}
% VI-CATALOG-ROW rutas 117
\pdfbookmark[2]{Celda 30}{cell-30}
\pdfbookmark[3]{r4c3\_h-1v-1}{route-117}
\RouteTitle{Ruta 117}{r4c3\_\allowbreak{}h-1v-1}
\noindent Celda 30\quad (4, 3)\qquad Orientación \(h=-1,\ v=-1\)\hfill\hyperref[sec:vi-indice-rutas]{Índice}\par\vspace{4pt}
\begingroup\fontsize{10.2}{12.5}\selectfont\setlength{\tabcolsep}{5pt}\renewcommand{\arraystretch}{1.1}\rowcolors{1}{routepale}{white}\noindent\begin{tabularx}{\linewidth}{@{}>{\raggedright\arraybackslash}p{.345\linewidth}>{\raggedright\arraybackslash}X@{}}
\RouteGroup{Identidad estructural}
\RouteRow{Clase y multisección}{quark\_\allowbreak{}down\_\allowbreak{}G2\quad /\quad quark\_\allowbreak{}down\_\allowbreak{}G2}
\RouteRow{Colección y sector}{quark de tipo abajo; quark}
\RouteRow{Clasificación física}{quark de tipo abajo}
\RouteRow{Generación, profundidad y color}{2; G2; G}
\RouteRow{Ruta primaria y órbita de inversión}{\(B_1\): no\qquad 4,\allowbreak{}3;\allowbreak{}6,\allowbreak{}7}
\RouteRow{Firma de clasificación}{28|\allowbreak{}-6|\allowbreak{}0|\allowbreak{}0|\allowbreak{}-12|\allowbreak{}+++-\kern0pt{}-\kern0pt{}-+++-\kern0pt{}-\kern0pt{}-}
\RouteGroup{Retornos, memoria y transporte}
\RouteRow{Retornos con fase}{clase: 27; órbita: 27; celda: 27}
\RouteRow{Retorno cinemático}{en 108: 54; extensión: 216}
\RouteRow{Giros y cambios de hoja}{71; 107}
\RouteRow{Acarreo y diferimiento}{deuda total: 1; final: 1; bloques diferidos: 1}
\RouteRow{Tríadas finales; persistencia de Pareto}{16; no}
\RouteRow{\(K\longrightarrow K+9\)}{repeticiones: 9; cambios de memoria: 9}
\RouteGroup{Firma, fase y diferencias}
\RouteRow{\(n_A,\ s_C,\ \kappa_0,\ \nu_{120},\ \nu_{270}\)}{28, 6, 0, 4, -12}
\RouteRow{\(\tau_{12}\); firma histórica \(\mathbb Z^5\)}{+++-\kern0pt{}-\kern0pt{}-+++-\kern0pt{}-\kern0pt{}-; 28|\allowbreak{}6|\allowbreak{}0|\allowbreak{}4|\allowbreak{}-12}
\RouteRow{\(D_1,\ D_3,\ D_4,\ D_{27},\ D_{36}\)}{66, 60, 64, 48, 32}
\RouteRow{\(\Omega_{90/120},\ P_6\)}{80, 6}
\RouteGroup{Lectores y factores de acción}
\RouteRow{\(K_D\quad /\quad K_\Omega\)}{(80,\allowbreak{}66,\allowbreak{}2)\quad /\quad (80,\allowbreak{}54,\allowbreak{}6)}
\RouteRow{\(\kappa_D\)}{79.51859\allowbreak{}71232726\allowbreak{}53414287\allowbreak{}29}
\RouteRow{\(\kappa_\Omega\)}{79.60661\allowbreak{}01397948\allowbreak{}27586470\allowbreak{}91}
\RouteRow{\(R_{\mathrm{act},D}\)}{1.000000\allowbreak{}05512879\allowbreak{}47092488\allowbreak{}80}
\RouteRow{\(R_{\mathrm{act},\Omega}\)}{1.000000\allowbreak{}05518981\allowbreak{}25318591\allowbreak{}40}
\RouteGroup{Separación de prefijos}
\RouteRow{Área de ambigüedad; área \(\log_2\)}{18; 0.0}
\RouteRow{Primer bloque de separación}{órbita: 1; celda: 1; ruta: 1}
\RouteRow{Ambigüedad final}{1}
\RouteGroup{Secuencias completas}
\RouteRow{\(w_6\)}{423 15 486 423 12 99 327 498 99 423 15 486 423 12 99 327 498 99}
\RouteRow{Tríadas estables por bloque}{0 1 2 3 4 5 6 7 8 9 10 11 12 13 14 15 16 16}
\end{tabularx}\endgroup\par\vspace{5pt}\noindent{\fontsize{8.4}{10}\selectfont\hyperref[trace-117]{Procedencia y códigos originales}\hfill Ficha completa · 117/324}
\clearpage\phantomsection\label{nav:vi-ruta-118}
% VI-CATALOG-ROW rutas 118
\pdfbookmark[3]{r4c3\_h-1v+1}{route-118}
\RouteTitle{Ruta 118}{r4c3\_\allowbreak{}h-1v+1}
\noindent Celda 30\quad (4, 3)\qquad Orientación \(h=-1,\ v=1\)\hfill\hyperref[sec:vi-indice-rutas]{Índice}\par\vspace{4pt}
\begingroup\fontsize{10.2}{12.5}\selectfont\setlength{\tabcolsep}{5pt}\renewcommand{\arraystretch}{1.1}\rowcolors{1}{routepale}{white}\noindent\begin{tabularx}{\linewidth}{@{}>{\raggedright\arraybackslash}p{.345\linewidth}>{\raggedright\arraybackslash}X@{}}
\RouteGroup{Identidad estructural}
\RouteRow{Clase y multisección}{quark\_\allowbreak{}down\_\allowbreak{}G2\quad /\quad quark\_\allowbreak{}down\_\allowbreak{}G2}
\RouteRow{Colección y sector}{quark de tipo abajo; quark}
\RouteRow{Clasificación física}{quark de tipo abajo}
\RouteRow{Generación, profundidad y color}{2; G2; G}
\RouteRow{Ruta primaria y órbita de inversión}{\(B_1\): no\qquad 4,\allowbreak{}3;\allowbreak{}6,\allowbreak{}7}
\RouteRow{Firma de clasificación}{28|\allowbreak{}-6|\allowbreak{}0|\allowbreak{}0|\allowbreak{}-12|\allowbreak{}+++-\kern0pt{}-\kern0pt{}-+++-\kern0pt{}-\kern0pt{}-}
\RouteGroup{Retornos, memoria y transporte}
\RouteRow{Retornos con fase}{clase: 27; órbita: 27; celda: 27}
\RouteRow{Retorno cinemático}{en 108: 54; extensión: 216}
\RouteRow{Giros y cambios de hoja}{71; 107}
\RouteRow{Acarreo y diferimiento}{deuda total: 1; final: 0; bloques diferidos: 1}
\RouteRow{Tríadas finales; persistencia de Pareto}{17; no}
\RouteRow{\(K\longrightarrow K+9\)}{repeticiones: 9; cambios de memoria: 9}
\RouteGroup{Firma, fase y diferencias}
\RouteRow{\(n_A,\ s_C,\ \kappa_0,\ \nu_{120},\ \nu_{270}\)}{50, 4, -2, 2, -8}
\RouteRow{\(\tau_{12}\); firma histórica \(\mathbb Z^5\)}{+++-0-+++-0-; 50|\allowbreak{}4|\allowbreak{}-2|\allowbreak{}2|\allowbreak{}-8}
\RouteRow{\(D_1,\ D_3,\ D_4,\ D_{27},\ D_{36}\)}{84, 24, 84, 24, 16}
\RouteRow{\(\Omega_{90/120},\ P_6\)}{48, 5}
\RouteGroup{Lectores y factores de acción}
\RouteRow{\(K_D\quad /\quad K_\Omega\)}{(40,\allowbreak{}84,\allowbreak{}30)\quad /\quad (48,\allowbreak{}54,\allowbreak{}5)}
\RouteRow{\(\kappa_D\)}{39.39035\allowbreak{}82768609\allowbreak{}24917849\allowbreak{}59}
\RouteRow{\(\kappa_\Omega\)}{47.60649\allowbreak{}89433721\allowbreak{}35446416\allowbreak{}86}
\RouteRow{\(R_{\mathrm{act},D}\)}{1.000000\allowbreak{}02730861\allowbreak{}73967087\allowbreak{}05}
\RouteRow{\(R_{\mathrm{act},\Omega}\)}{1.000000\allowbreak{}03300471\allowbreak{}80532430\allowbreak{}84}
\RouteGroup{Separación de prefijos}
\RouteRow{Área de ambigüedad; área \(\log_2\)}{18; 0.0}
\RouteRow{Primer bloque de separación}{órbita: 1; celda: 1; ruta: 1}
\RouteRow{Ambigüedad final}{1}
\RouteGroup{Secuencias completas}
\RouteRow{\(w_6\)}{420 258 414 420 255 252 333 255 252 420 258 414 420 255 252 333 255 252}
\RouteRow{Tríadas estables por bloque}{0 1 2 3 4 5 6 7 8 9 10 11 12 13 14 15 15 17}
\end{tabularx}\endgroup\par\vspace{5pt}\noindent{\fontsize{8.4}{10}\selectfont\hyperref[trace-118]{Procedencia y códigos originales}\hfill Ficha completa · 118/324}
\clearpage\phantomsection\label{nav:vi-ruta-119}
% VI-CATALOG-ROW rutas 119
\pdfbookmark[3]{r4c3\_h+1v-1}{route-119}
\RouteTitle{Ruta 119}{r4c3\_\allowbreak{}h+1v-1}
\noindent Celda 30\quad (4, 3)\qquad Orientación \(h=1,\ v=-1\)\hfill\hyperref[sec:vi-indice-rutas]{Índice}\par\vspace{4pt}
\begingroup\fontsize{10.2}{12.5}\selectfont\setlength{\tabcolsep}{5pt}\renewcommand{\arraystretch}{1.1}\rowcolors{1}{routepale}{white}\noindent\begin{tabularx}{\linewidth}{@{}>{\raggedright\arraybackslash}p{.345\linewidth}>{\raggedright\arraybackslash}X@{}}
\RouteGroup{Identidad estructural}
\RouteRow{Clase y multisección}{quark\_\allowbreak{}down\_\allowbreak{}G2\quad /\quad quark\_\allowbreak{}down\_\allowbreak{}G2}
\RouteRow{Colección y sector}{quark de tipo abajo; quark}
\RouteRow{Clasificación física}{quark de tipo abajo}
\RouteRow{Generación, profundidad y color}{2; G2; G}
\RouteRow{Ruta primaria y órbita de inversión}{\(B_1\): no\qquad 4,\allowbreak{}3;\allowbreak{}6,\allowbreak{}7}
\RouteRow{Firma de clasificación}{28|\allowbreak{}-6|\allowbreak{}0|\allowbreak{}0|\allowbreak{}-12|\allowbreak{}+++-\kern0pt{}-\kern0pt{}-+++-\kern0pt{}-\kern0pt{}-}
\RouteGroup{Retornos, memoria y transporte}
\RouteRow{Retornos con fase}{clase: 27; órbita: 27; celda: 27}
\RouteRow{Retorno cinemático}{en 108: 54; extensión: 216}
\RouteRow{Giros y cambios de hoja}{71; 107}
\RouteRow{Acarreo y diferimiento}{deuda total: 0; final: 0; bloques diferidos: 0}
\RouteRow{Tríadas finales; persistencia de Pareto}{17; no}
\RouteRow{\(K\longrightarrow K+9\)}{repeticiones: 9; cambios de memoria: 9}
\RouteGroup{Firma, fase y diferencias}
\RouteRow{\(n_A,\ s_C,\ \kappa_0,\ \nu_{120},\ \nu_{270}\)}{50, -4, 2, -2, -8}
\RouteRow{\(\tau_{12}\); firma histórica \(\mathbb Z^5\)}{+0+-\kern0pt{}-\kern0pt{}-+0+-\kern0pt{}-\kern0pt{}-; 50|\allowbreak{}-4|\allowbreak{}2|\allowbreak{}-2|\allowbreak{}-8}
\RouteRow{\(D_1,\ D_3,\ D_4,\ D_{27},\ D_{36}\)}{84, 24, 84, 24, 16}
\RouteRow{\(\Omega_{90/120},\ P_6\)}{48, 5}
\RouteGroup{Lectores y factores de acción}
\RouteRow{\(K_D\quad /\quad K_\Omega\)}{(40,\allowbreak{}84,\allowbreak{}30)\quad /\quad (48,\allowbreak{}54,\allowbreak{}5)}
\RouteRow{\(\kappa_D\)}{39.39035\allowbreak{}82768609\allowbreak{}24917849\allowbreak{}59}
\RouteRow{\(\kappa_\Omega\)}{47.60649\allowbreak{}89433721\allowbreak{}35446416\allowbreak{}86}
\RouteRow{\(R_{\mathrm{act},D}\)}{1.000000\allowbreak{}02730861\allowbreak{}73967087\allowbreak{}05}
\RouteRow{\(R_{\mathrm{act},\Omega}\)}{1.000000\allowbreak{}03300471\allowbreak{}80532430\allowbreak{}84}
\RouteGroup{Separación de prefijos}
\RouteRow{Área de ambigüedad; área \(\log_2\)}{18; 0.0}
\RouteRow{Primer bloque de separación}{órbita: 1; celda: 1; ruta: 1}
\RouteRow{Ambigüedad final}{1}
\RouteGroup{Secuencias completas}
\RouteRow{\(w_6\)}{333 255 252 333 258 414 420 258 414 333 255 252 333 258 414 420 258 414}
\RouteRow{Tríadas estables por bloque}{0 1 2 3 4 5 6 7 8 9 10 11 12 13 14 15 16 17}
\end{tabularx}\endgroup\par\vspace{5pt}\noindent{\fontsize{8.4}{10}\selectfont\hyperref[trace-119]{Procedencia y códigos originales}\hfill Ficha completa · 119/324}
\clearpage\phantomsection\label{nav:vi-ruta-120}
% VI-CATALOG-ROW rutas 120
\pdfbookmark[3]{r4c3\_h+1v+1}{route-120}
\RouteTitle{Ruta 120}{r4c3\_\allowbreak{}h+1v+1}
\noindent Celda 30\quad (4, 3)\qquad Orientación \(h=1,\ v=1\)\hfill\hyperref[sec:vi-indice-rutas]{Índice}\par\vspace{4pt}
\begingroup\fontsize{10.2}{12.5}\selectfont\setlength{\tabcolsep}{5pt}\renewcommand{\arraystretch}{1.1}\rowcolors{1}{routepale}{white}\noindent\begin{tabularx}{\linewidth}{@{}>{\raggedright\arraybackslash}p{.345\linewidth}>{\raggedright\arraybackslash}X@{}}
\RouteGroup{Identidad estructural}
\RouteRow{Clase y multisección}{quark\_\allowbreak{}down\_\allowbreak{}G2\quad /\quad quark\_\allowbreak{}down\_\allowbreak{}G2}
\RouteRow{Colección y sector}{quark de tipo abajo; quark}
\RouteRow{Clasificación física}{quark de tipo abajo}
\RouteRow{Generación, profundidad y color}{2; G2; G}
\RouteRow{Ruta primaria y órbita de inversión}{\(B_1\): sí\qquad 4,\allowbreak{}3;\allowbreak{}6,\allowbreak{}7}
\RouteRow{Firma de clasificación}{28|\allowbreak{}-6|\allowbreak{}0|\allowbreak{}0|\allowbreak{}-12|\allowbreak{}+++-\kern0pt{}-\kern0pt{}-+++-\kern0pt{}-\kern0pt{}-}
\RouteGroup{Retornos, memoria y transporte}
\RouteRow{Retornos con fase}{clase: 27; órbita: 27; celda: 27}
\RouteRow{Retorno cinemático}{en 108: 54; extensión: 216}
\RouteRow{Giros y cambios de hoja}{71; 107}
\RouteRow{Acarreo y diferimiento}{deuda total: 2; final: 1; bloques diferidos: 2}
\RouteRow{Tríadas finales; persistencia de Pareto}{16; no}
\RouteRow{\(K\longrightarrow K+9\)}{repeticiones: 9; cambios de memoria: 9}
\RouteGroup{Firma, fase y diferencias}
\RouteRow{\(n_A,\ s_C,\ \kappa_0,\ \nu_{120},\ \nu_{270}\)}{28, -6, 0, -4, -12}
\RouteRow{\(\tau_{12}\); firma histórica \(\mathbb Z^5\)}{+++-\kern0pt{}-\kern0pt{}-+++-\kern0pt{}-\kern0pt{}-; 28|\allowbreak{}-6|\allowbreak{}0|\allowbreak{}-4|\allowbreak{}-12}
\RouteRow{\(D_1,\ D_3,\ D_4,\ D_{27},\ D_{36}\)}{66, 60, 64, 48, 32}
\RouteRow{\(\Omega_{90/120},\ P_6\)}{80, 6}
\RouteGroup{Lectores y factores de acción}
\RouteRow{\(K_D\quad /\quad K_\Omega\)}{(80,\allowbreak{}66,\allowbreak{}2)\quad /\quad (80,\allowbreak{}54,\allowbreak{}6)}
\RouteRow{\(\kappa_D\)}{79.51859\allowbreak{}71232726\allowbreak{}53414287\allowbreak{}29}
\RouteRow{\(\kappa_\Omega\)}{79.60661\allowbreak{}01397948\allowbreak{}27586470\allowbreak{}91}
\RouteRow{\(R_{\mathrm{act},D}\)}{1.000000\allowbreak{}05512879\allowbreak{}47092488\allowbreak{}80}
\RouteRow{\(R_{\mathrm{act},\Omega}\)}{1.000000\allowbreak{}05518981\allowbreak{}25318591\allowbreak{}40}
\RouteGroup{Separación de prefijos}
\RouteRow{Área de ambigüedad; área \(\log_2\)}{18; 0.0}
\RouteRow{Primer bloque de separación}{órbita: 1; celda: 1; ruta: 1}
\RouteRow{Ambigüedad final}{1}
\RouteGroup{Secuencias completas}
\RouteRow{\(w_6\)}{327 498 99 327 501 486 423 15 486 327 498 99 327 501 486 423 15 486}
\RouteRow{Tríadas estables por bloque}{0 1 2 3 4 5 6 7 7 9 10 11 12 13 14 15 16 16}
\end{tabularx}\endgroup\par\vspace{5pt}\noindent{\fontsize{8.4}{10}\selectfont\hyperref[trace-120]{Procedencia y códigos originales}\hfill Ficha completa · 120/324}
\clearpage\phantomsection\label{nav:vi-ruta-121}
% VI-CATALOG-ROW rutas 121
\pdfbookmark[2]{Celda 31}{cell-31}
\pdfbookmark[3]{r4c4\_h-1v-1}{route-121}
\RouteTitle{Ruta 121}{r4c4\_\allowbreak{}h-1v-1}
\noindent Celda 31\quad (4, 4)\qquad Orientación \(h=-1,\ v=-1\)\hfill\hyperref[sec:vi-indice-rutas]{Índice}\par\vspace{4pt}
\begingroup\fontsize{10.2}{12.5}\selectfont\setlength{\tabcolsep}{5pt}\renewcommand{\arraystretch}{1.1}\rowcolors{1}{routepale}{white}\noindent\begin{tabularx}{\linewidth}{@{}>{\raggedright\arraybackslash}p{.345\linewidth}>{\raggedright\arraybackslash}X@{}}
\RouteGroup{Identidad estructural}
\RouteRow{Clase y multisección}{quark\_\allowbreak{}up\_\allowbreak{}G1\quad /\quad quark\_\allowbreak{}up\_\allowbreak{}G1}
\RouteRow{Colección y sector}{quark de tipo arriba; quark}
\RouteRow{Clasificación física}{quark de tipo arriba}
\RouteRow{Generación, profundidad y color}{1; G1; R}
\RouteRow{Ruta primaria y órbita de inversión}{\(B_1\): no\qquad 4,\allowbreak{}4;\allowbreak{}6,\allowbreak{}6}
\RouteRow{Firma de clasificación}{24|\allowbreak{}0|\allowbreak{}12|\allowbreak{}0|\allowbreak{}-12|\allowbreak{}+++-\kern0pt{}-\kern0pt{}-+++-\kern0pt{}-\kern0pt{}-}
\RouteGroup{Retornos, memoria y transporte}
\RouteRow{Retornos con fase}{clase: 27; órbita: 27; celda: 27}
\RouteRow{Retorno cinemático}{en 108: 54; extensión: 216}
\RouteRow{Giros y cambios de hoja}{71; 107}
\RouteRow{Acarreo y diferimiento}{deuda total: 2; final: 1; bloques diferidos: 2}
\RouteRow{Tríadas finales; persistencia de Pareto}{16; sí}
\RouteRow{\(K\longrightarrow K+9\)}{repeticiones: 9; cambios de memoria: 9}
\RouteGroup{Firma, fase y diferencias}
\RouteRow{\(n_A,\ s_C,\ \kappa_0,\ \nu_{120},\ \nu_{270}\)}{24, 0, 0, 0, -12}
\RouteRow{\(\tau_{12}\); firma histórica \(\mathbb Z^5\)}{+++-\kern0pt{}-\kern0pt{}-+++-\kern0pt{}-\kern0pt{}-; 24|\allowbreak{}0|\allowbreak{}0|\allowbreak{}0|\allowbreak{}-12}
\RouteRow{\(D_1,\ D_3,\ D_4,\ D_{27},\ D_{36}\)}{54, 56, 70, 48, 32}
\RouteRow{\(\Omega_{90/120},\ P_6\)}{80, 6}
\RouteGroup{Lectores y factores de acción}
\RouteRow{\(K_D\quad /\quad K_\Omega\)}{(80,\allowbreak{}54,\allowbreak{}7)\quad /\quad (80,\allowbreak{}54,\allowbreak{}6)}
\RouteRow{\(\kappa_D\)}{79.60672\allowbreak{}13362175\allowbreak{}19726524\allowbreak{}96}
\RouteRow{\(\kappa_\Omega\)}{79.60661\allowbreak{}01397948\allowbreak{}27586470\allowbreak{}91}
\RouteRow{\(R_{\mathrm{act},D}\)}{1.000000\allowbreak{}05518988\allowbreak{}96223153\allowbreak{}37}
\RouteRow{\(R_{\mathrm{act},\Omega}\)}{1.000000\allowbreak{}05518981\allowbreak{}25318591\allowbreak{}40}
\RouteGroup{Separación de prefijos}
\RouteRow{Área de ambigüedad; área \(\log_2\)}{18; 0.0}
\RouteRow{Primer bloque de separación}{órbita: 1; celda: 1; ruta: 1}
\RouteRow{Ambigüedad final}{1}
\RouteGroup{Secuencias completas}
\RouteRow{\(w_6\)}{486 423 15 486 429 243 657 24 243 486 423 15 486 429 243 657 24 243}
\RouteRow{Tríadas estables por bloque}{0 1 2 3 4 5 6 7 8 9 10 11 12 13 13 15 16 16}
\end{tabularx}\endgroup\par\vspace{5pt}\noindent{\fontsize{8.4}{10}\selectfont\hyperref[trace-121]{Procedencia y códigos originales}\hfill Ficha completa · 121/324}
\clearpage\phantomsection\label{nav:vi-ruta-122}
% VI-CATALOG-ROW rutas 122
\pdfbookmark[3]{r4c4\_h-1v+1}{route-122}
\RouteTitle{Ruta 122}{r4c4\_\allowbreak{}h-1v+1}
\noindent Celda 31\quad (4, 4)\qquad Orientación \(h=-1,\ v=1\)\hfill\hyperref[sec:vi-indice-rutas]{Índice}\par\vspace{4pt}
\begingroup\fontsize{10.2}{12.5}\selectfont\setlength{\tabcolsep}{5pt}\renewcommand{\arraystretch}{1.1}\rowcolors{1}{routepale}{white}\noindent\begin{tabularx}{\linewidth}{@{}>{\raggedright\arraybackslash}p{.345\linewidth}>{\raggedright\arraybackslash}X@{}}
\RouteGroup{Identidad estructural}
\RouteRow{Clase y multisección}{quark\_\allowbreak{}up\_\allowbreak{}G1\quad /\quad quark\_\allowbreak{}up\_\allowbreak{}G1}
\RouteRow{Colección y sector}{quark de tipo arriba; quark}
\RouteRow{Clasificación física}{quark de tipo arriba}
\RouteRow{Generación, profundidad y color}{1; G1; R}
\RouteRow{Ruta primaria y órbita de inversión}{\(B_1\): sí\qquad 4,\allowbreak{}4;\allowbreak{}6,\allowbreak{}6}
\RouteRow{Firma de clasificación}{24|\allowbreak{}0|\allowbreak{}12|\allowbreak{}0|\allowbreak{}-12|\allowbreak{}+++-\kern0pt{}-\kern0pt{}-+++-\kern0pt{}-\kern0pt{}-}
\RouteGroup{Retornos, memoria y transporte}
\RouteRow{Retornos con fase}{clase: 27; órbita: 27; celda: 27}
\RouteRow{Retorno cinemático}{en 108: 54; extensión: 216}
\RouteRow{Giros y cambios de hoja}{71; 107}
\RouteRow{Acarreo y diferimiento}{deuda total: 0; final: 0; bloques diferidos: 0}
\RouteRow{Tríadas finales; persistencia de Pareto}{17; sí}
\RouteRow{\(K\longrightarrow K+9\)}{repeticiones: 9; cambios de memoria: 9}
\RouteGroup{Firma, fase y diferencias}
\RouteRow{\(n_A,\ s_C,\ \kappa_0,\ \nu_{120},\ \nu_{270}\)}{14, -4, 0, 0, -12}
\RouteRow{\(\tau_{12}\); firma histórica \(\mathbb Z^5\)}{+++-\kern0pt{}-\kern0pt{}-+++-\kern0pt{}-\kern0pt{}-; 14|\allowbreak{}-4|\allowbreak{}0|\allowbreak{}0|\allowbreak{}-12}
\RouteRow{\(D_1,\ D_3,\ D_4,\ D_{27},\ D_{36}\)}{78, 28, 78, 36, 24}
\RouteRow{\(\Omega_{90/120},\ P_6\)}{80, 6}
\RouteGroup{Lectores y factores de acción}
\RouteRow{\(K_D\quad /\quad K_\Omega\)}{(60,\allowbreak{}78,\allowbreak{}25)\quad /\quad (80,\allowbreak{}54,\allowbreak{}6)}
\RouteRow{\(\kappa_D\)}{59.43358\allowbreak{}64101631\allowbreak{}67023563\allowbreak{}04}
\RouteRow{\(\kappa_\Omega\)}{79.60661\allowbreak{}01397948\allowbreak{}27586470\allowbreak{}91}
\RouteRow{\(R_{\mathrm{act},D}\)}{1.000000\allowbreak{}04120422\allowbreak{}24053449\allowbreak{}05}
\RouteRow{\(R_{\mathrm{act},\Omega}\)}{1.000000\allowbreak{}05518981\allowbreak{}25318591\allowbreak{}40}
\RouteGroup{Separación de prefijos}
\RouteRow{Área de ambigüedad; área \(\log_2\)}{18; 0.0}
\RouteRow{Primer bloque de separación}{órbita: 1; celda: 1; ruta: 1}
\RouteRow{Ambigüedad final}{1}
\RouteGroup{Secuencias completas}
\RouteRow{\(w_6\)}{507 504 585 507 510 510 666 672 510 507 504 585 507 510 510 666 672 510}
\RouteRow{Tríadas estables por bloque}{0 1 2 3 4 5 6 7 8 9 10 11 12 13 14 15 16 17}
\end{tabularx}\endgroup\par\vspace{5pt}\noindent{\fontsize{8.4}{10}\selectfont\hyperref[trace-122]{Procedencia y códigos originales}\hfill Ficha completa · 122/324}
\clearpage\phantomsection\label{nav:vi-ruta-123}
% VI-CATALOG-ROW rutas 123
\pdfbookmark[3]{r4c4\_h+1v-1}{route-123}
\RouteTitle{Ruta 123}{r4c4\_\allowbreak{}h+1v-1}
\noindent Celda 31\quad (4, 4)\qquad Orientación \(h=1,\ v=-1\)\hfill\hyperref[sec:vi-indice-rutas]{Índice}\par\vspace{4pt}
\begingroup\fontsize{10.2}{12.5}\selectfont\setlength{\tabcolsep}{5pt}\renewcommand{\arraystretch}{1.1}\rowcolors{1}{routepale}{white}\noindent\begin{tabularx}{\linewidth}{@{}>{\raggedright\arraybackslash}p{.345\linewidth}>{\raggedright\arraybackslash}X@{}}
\RouteGroup{Identidad estructural}
\RouteRow{Clase y multisección}{quark\_\allowbreak{}up\_\allowbreak{}G1\quad /\quad quark\_\allowbreak{}up\_\allowbreak{}G1}
\RouteRow{Colección y sector}{quark de tipo arriba; quark}
\RouteRow{Clasificación física}{quark de tipo arriba}
\RouteRow{Generación, profundidad y color}{1; G1; R}
\RouteRow{Ruta primaria y órbita de inversión}{\(B_1\): no\qquad 4,\allowbreak{}4;\allowbreak{}6,\allowbreak{}6}
\RouteRow{Firma de clasificación}{24|\allowbreak{}0|\allowbreak{}12|\allowbreak{}0|\allowbreak{}-12|\allowbreak{}+++-\kern0pt{}-\kern0pt{}-+++-\kern0pt{}-\kern0pt{}-}
\RouteGroup{Retornos, memoria y transporte}
\RouteRow{Retornos con fase}{clase: 27; órbita: 27; celda: 27}
\RouteRow{Retorno cinemático}{en 108: 54; extensión: 216}
\RouteRow{Giros y cambios de hoja}{71; 107}
\RouteRow{Acarreo y diferimiento}{deuda total: 2; final: 0; bloques diferidos: 1}
\RouteRow{Tríadas finales; persistencia de Pareto}{17; sí}
\RouteRow{\(K\longrightarrow K+9\)}{repeticiones: 9; cambios de memoria: 9}
\RouteGroup{Firma, fase y diferencias}
\RouteRow{\(n_A,\ s_C,\ \kappa_0,\ \nu_{120},\ \nu_{270}\)}{14, 4, 0, 0, -12}
\RouteRow{\(\tau_{12}\); firma histórica \(\mathbb Z^5\)}{+++-\kern0pt{}-\kern0pt{}-+++-\kern0pt{}-\kern0pt{}-; 14|\allowbreak{}4|\allowbreak{}0|\allowbreak{}0|\allowbreak{}-12}
\RouteRow{\(D_1,\ D_3,\ D_4,\ D_{27},\ D_{36}\)}{78, 28, 78, 36, 24}
\RouteRow{\(\Omega_{90/120},\ P_6\)}{80, 6}
\RouteGroup{Lectores y factores de acción}
\RouteRow{\(K_D\quad /\quad K_\Omega\)}{(60,\allowbreak{}78,\allowbreak{}25)\quad /\quad (80,\allowbreak{}54,\allowbreak{}6)}
\RouteRow{\(\kappa_D\)}{59.43358\allowbreak{}64101631\allowbreak{}67023563\allowbreak{}04}
\RouteRow{\(\kappa_\Omega\)}{79.60661\allowbreak{}01397948\allowbreak{}27586470\allowbreak{}91}
\RouteRow{\(R_{\mathrm{act},D}\)}{1.000000\allowbreak{}04120422\allowbreak{}24053449\allowbreak{}05}
\RouteRow{\(R_{\mathrm{act},\Omega}\)}{1.000000\allowbreak{}05518981\allowbreak{}25318591\allowbreak{}40}
\RouteGroup{Separación de prefijos}
\RouteRow{Área de ambigüedad; área \(\log_2\)}{18; 0.0}
\RouteRow{Primer bloque de separación}{órbita: 1; celda: 1; ruta: 1}
\RouteRow{Ambigüedad final}{1}
\RouteGroup{Secuencias completas}
\RouteRow{\(w_6\)}{666 672 510 666 666 585 507 504 585 666 672 510 666 666 585 507 504 585}
\RouteRow{Tríadas estables por bloque}{0 1 2 3 4 5 6 7 8 9 10 11 12 12 13 15 16 17}
\end{tabularx}\endgroup\par\vspace{5pt}\noindent{\fontsize{8.4}{10}\selectfont\hyperref[trace-123]{Procedencia y códigos originales}\hfill Ficha completa · 123/324}
\clearpage\phantomsection\label{nav:vi-ruta-124}
% VI-CATALOG-ROW rutas 124
\pdfbookmark[3]{r4c4\_h+1v+1}{route-124}
\RouteTitle{Ruta 124}{r4c4\_\allowbreak{}h+1v+1}
\noindent Celda 31\quad (4, 4)\qquad Orientación \(h=1,\ v=1\)\hfill\hyperref[sec:vi-indice-rutas]{Índice}\par\vspace{4pt}
\begingroup\fontsize{10.2}{12.5}\selectfont\setlength{\tabcolsep}{5pt}\renewcommand{\arraystretch}{1.1}\rowcolors{1}{routepale}{white}\noindent\begin{tabularx}{\linewidth}{@{}>{\raggedright\arraybackslash}p{.345\linewidth}>{\raggedright\arraybackslash}X@{}}
\RouteGroup{Identidad estructural}
\RouteRow{Clase y multisección}{quark\_\allowbreak{}up\_\allowbreak{}G1\quad /\quad quark\_\allowbreak{}up\_\allowbreak{}G1}
\RouteRow{Colección y sector}{quark de tipo arriba; quark}
\RouteRow{Clasificación física}{quark de tipo arriba}
\RouteRow{Generación, profundidad y color}{1; G1; R}
\RouteRow{Ruta primaria y órbita de inversión}{\(B_1\): no\qquad 4,\allowbreak{}4;\allowbreak{}6,\allowbreak{}6}
\RouteRow{Firma de clasificación}{24|\allowbreak{}0|\allowbreak{}12|\allowbreak{}0|\allowbreak{}-12|\allowbreak{}+++-\kern0pt{}-\kern0pt{}-+++-\kern0pt{}-\kern0pt{}-}
\RouteGroup{Retornos, memoria y transporte}
\RouteRow{Retornos con fase}{clase: 27; órbita: 27; celda: 27}
\RouteRow{Retorno cinemático}{en 108: 54; extensión: 216}
\RouteRow{Giros y cambios de hoja}{71; 107}
\RouteRow{Acarreo y diferimiento}{deuda total: 2; final: 0; bloques diferidos: 1}
\RouteRow{Tríadas finales; persistencia de Pareto}{17; sí}
\RouteRow{\(K\longrightarrow K+9\)}{repeticiones: 9; cambios de memoria: 9}
\RouteGroup{Firma, fase y diferencias}
\RouteRow{\(n_A,\ s_C,\ \kappa_0,\ \nu_{120},\ \nu_{270}\)}{24, 0, 0, 0, -12}
\RouteRow{\(\tau_{12}\); firma histórica \(\mathbb Z^5\)}{+++-\kern0pt{}-\kern0pt{}-+++-\kern0pt{}-\kern0pt{}-; 24|\allowbreak{}0|\allowbreak{}0|\allowbreak{}0|\allowbreak{}-12}
\RouteRow{\(D_1,\ D_3,\ D_4,\ D_{27},\ D_{36}\)}{54, 56, 70, 48, 32}
\RouteRow{\(\Omega_{90/120},\ P_6\)}{80, 6}
\RouteGroup{Lectores y factores de acción}
\RouteRow{\(K_D\quad /\quad K_\Omega\)}{(80,\allowbreak{}54,\allowbreak{}7)\quad /\quad (80,\allowbreak{}54,\allowbreak{}6)}
\RouteRow{\(\kappa_D\)}{79.60672\allowbreak{}13362175\allowbreak{}19726524\allowbreak{}96}
\RouteRow{\(\kappa_\Omega\)}{79.60661\allowbreak{}01397948\allowbreak{}27586470\allowbreak{}91}
\RouteRow{\(R_{\mathrm{act},D}\)}{1.000000\allowbreak{}05518988\allowbreak{}96223153\allowbreak{}37}
\RouteRow{\(R_{\mathrm{act},\Omega}\)}{1.000000\allowbreak{}05518981\allowbreak{}25318591\allowbreak{}40}
\RouteGroup{Separación de prefijos}
\RouteRow{Área de ambigüedad; área \(\log_2\)}{18; 0.0}
\RouteRow{Primer bloque de separación}{órbita: 1; celda: 1; ruta: 1}
\RouteRow{Ambigüedad final}{1}
\RouteGroup{Secuencias completas}
\RouteRow{\(w_6\)}{657 24 243 657 18 15 486 423 15 657 24 243 657 18 15 486 423 15}
\RouteRow{Tríadas estables por bloque}{0 1 2 3 4 5 6 7 8 9 10 11 12 13 14 14 15 17}
\end{tabularx}\endgroup\par\vspace{5pt}\noindent{\fontsize{8.4}{10}\selectfont\hyperref[trace-124]{Procedencia y códigos originales}\hfill Ficha completa · 124/324}
\clearpage\phantomsection\label{nav:vi-ruta-125}
% VI-CATALOG-ROW rutas 125
\pdfbookmark[2]{Celda 32}{cell-32}
\pdfbookmark[3]{r4c5\_h-1v-1}{route-125}
\RouteTitle{Ruta 125}{r4c5\_\allowbreak{}h-1v-1}
\noindent Celda 32\quad (4, 5)\qquad Orientación \(h=-1,\ v=-1\)\hfill\hyperref[sec:vi-indice-rutas]{Índice}\par\vspace{4pt}
\begingroup\fontsize{10.2}{12.5}\selectfont\setlength{\tabcolsep}{5pt}\renewcommand{\arraystretch}{1.1}\rowcolors{1}{routepale}{white}\noindent\begin{tabularx}{\linewidth}{@{}>{\raggedright\arraybackslash}p{.345\linewidth}>{\raggedright\arraybackslash}X@{}}
\RouteGroup{Identidad estructural}
\RouteRow{Clase y multisección}{quark\_\allowbreak{}down\_\allowbreak{}G1\quad /\quad quark\_\allowbreak{}down\_\allowbreak{}G1}
\RouteRow{Colección y sector}{quark de tipo abajo; quark}
\RouteRow{Clasificación física}{quark de tipo abajo}
\RouteRow{Generación, profundidad y color}{1; G1; B}
\RouteRow{Ruta primaria y órbita de inversión}{\(B_1\): no\qquad 4,\allowbreak{}5;\allowbreak{}6,\allowbreak{}5}
\RouteRow{Firma de clasificación}{28|\allowbreak{}6|\allowbreak{}0|\allowbreak{}0|\allowbreak{}-12|\allowbreak{}+++-\kern0pt{}-\kern0pt{}-+++-\kern0pt{}-\kern0pt{}-}
\RouteGroup{Retornos, memoria y transporte}
\RouteRow{Retornos con fase}{clase: 27; órbita: 27; celda: 27}
\RouteRow{Retorno cinemático}{en 108: 54; extensión: 216}
\RouteRow{Giros y cambios de hoja}{71; 107}
\RouteRow{Acarreo y diferimiento}{deuda total: 1; final: 0; bloques diferidos: 1}
\RouteRow{Tríadas finales; persistencia de Pareto}{17; sí}
\RouteRow{\(K\longrightarrow K+9\)}{repeticiones: 9; cambios de memoria: 9}
\RouteGroup{Firma, fase y diferencias}
\RouteRow{\(n_A,\ s_C,\ \kappa_0,\ \nu_{120},\ \nu_{270}\)}{28, -6, 0, 0, -12}
\RouteRow{\(\tau_{12}\); firma histórica \(\mathbb Z^5\)}{+++-\kern0pt{}-\kern0pt{}-+++-\kern0pt{}-\kern0pt{}-; 28|\allowbreak{}-6|\allowbreak{}0|\allowbreak{}0|\allowbreak{}-12}
\RouteRow{\(D_1,\ D_3,\ D_4,\ D_{27},\ D_{36}\)}{66, 64, 64, 60, 40}
\RouteRow{\(\Omega_{90/120},\ P_6\)}{80, 6}
\RouteGroup{Lectores y factores de acción}
\RouteRow{\(K_D\quad /\quad K_\Omega\)}{(100,\allowbreak{}66,\allowbreak{}0)\quad /\quad (80,\allowbreak{}54,\allowbreak{}6)}
\RouteRow{\(\kappa_D\)}{99.51837\allowbreak{}47304272\allowbreak{}69134179\allowbreak{}18}
\RouteRow{\(\kappa_\Omega\)}{79.60661\allowbreak{}01397948\allowbreak{}27586470\allowbreak{}91}
\RouteRow{\(R_{\mathrm{act},D}\)}{1.000000\allowbreak{}06899427\allowbreak{}66450218\allowbreak{}95}
\RouteRow{\(R_{\mathrm{act},\Omega}\)}{1.000000\allowbreak{}05518981\allowbreak{}25318591\allowbreak{}40}
\RouteGroup{Separación de prefijos}
\RouteRow{Área de ambigüedad; área \(\log_2\)}{18; 0.0}
\RouteRow{Primer bloque de separación}{órbita: 1; celda: 1; ruta: 1}
\RouteRow{Ambigüedad final}{1}
\RouteGroup{Secuencias completas}
\RouteRow{\(w_6\)}{90 570 264 90 567 657 24 243 657 90 570 264 90 567 657 24 243 657}
\RouteRow{Tríadas estables por bloque}{0 1 2 3 4 5 6 7 8 9 10 11 12 13 13 15 16 17}
\end{tabularx}\endgroup\par\vspace{5pt}\noindent{\fontsize{8.4}{10}\selectfont\hyperref[trace-125]{Procedencia y códigos originales}\hfill Ficha completa · 125/324}
\clearpage\phantomsection\label{nav:vi-ruta-126}
% VI-CATALOG-ROW rutas 126
\pdfbookmark[3]{r4c5\_h-1v+1}{route-126}
\RouteTitle{Ruta 126}{r4c5\_\allowbreak{}h-1v+1}
\noindent Celda 32\quad (4, 5)\qquad Orientación \(h=-1,\ v=1\)\hfill\hyperref[sec:vi-indice-rutas]{Índice}\par\vspace{4pt}
\begingroup\fontsize{10.2}{12.5}\selectfont\setlength{\tabcolsep}{5pt}\renewcommand{\arraystretch}{1.1}\rowcolors{1}{routepale}{white}\noindent\begin{tabularx}{\linewidth}{@{}>{\raggedright\arraybackslash}p{.345\linewidth}>{\raggedright\arraybackslash}X@{}}
\RouteGroup{Identidad estructural}
\RouteRow{Clase y multisección}{quark\_\allowbreak{}down\_\allowbreak{}G1\quad /\quad quark\_\allowbreak{}down\_\allowbreak{}G1}
\RouteRow{Colección y sector}{quark de tipo abajo; quark}
\RouteRow{Clasificación física}{quark de tipo abajo}
\RouteRow{Generación, profundidad y color}{1; G1; B}
\RouteRow{Ruta primaria y órbita de inversión}{\(B_1\): no\qquad 4,\allowbreak{}5;\allowbreak{}6,\allowbreak{}5}
\RouteRow{Firma de clasificación}{28|\allowbreak{}6|\allowbreak{}0|\allowbreak{}0|\allowbreak{}-12|\allowbreak{}+++-\kern0pt{}-\kern0pt{}-+++-\kern0pt{}-\kern0pt{}-}
\RouteGroup{Retornos, memoria y transporte}
\RouteRow{Retornos con fase}{clase: 27; órbita: 27; celda: 27}
\RouteRow{Retorno cinemático}{en 108: 54; extensión: 216}
\RouteRow{Giros y cambios de hoja}{71; 107}
\RouteRow{Acarreo y diferimiento}{deuda total: 2; final: 0; bloques diferidos: 2}
\RouteRow{Tríadas finales; persistencia de Pareto}{17; sí}
\RouteRow{\(K\longrightarrow K+9\)}{repeticiones: 9; cambios de memoria: 9}
\RouteGroup{Firma, fase y diferencias}
\RouteRow{\(n_A,\ s_C,\ \kappa_0,\ \nu_{120},\ \nu_{270}\)}{20, 0, 0, 0, -8}
\RouteRow{\(\tau_{12}\); firma histórica \(\mathbb Z^5\)}{+0+-0-+0+-0-; 20|\allowbreak{}0|\allowbreak{}0|\allowbreak{}0|\allowbreak{}-8}
\RouteRow{\(D_1,\ D_3,\ D_4,\ D_{27},\ D_{36}\)}{24, 20, 24, 24, 16}
\RouteRow{\(\Omega_{90/120},\ P_6\)}{40, 4}
\RouteGroup{Lectores y factores de acción}
\RouteRow{\(K_D\quad /\quad K_\Omega\)}{(40,\allowbreak{}24,\allowbreak{}2)\quad /\quad (40,\allowbreak{}54,\allowbreak{}4)}
\RouteRow{\(\kappa_D\)}{39.82508\allowbreak{}59311825\allowbreak{}73056173\allowbreak{}26}
\RouteRow{\(\kappa_\Omega\)}{39.60638\allowbreak{}77469494\allowbreak{}43306362\allowbreak{}81}
\RouteRow{\(R_{\mathrm{act},D}\)}{1.000000\allowbreak{}02761000\allowbreak{}61595142\allowbreak{}15}
\RouteRow{\(R_{\mathrm{act},\Omega}\)}{1.000000\allowbreak{}02745838\allowbreak{}66926514\allowbreak{}00}
\RouteGroup{Separación de prefijos}
\RouteRow{Área de ambigüedad; área \(\log_2\)}{18; 0.0}
\RouteRow{Primer bloque de separación}{órbita: 1; celda: 1; ruta: 1}
\RouteRow{Ambigüedad final}{1}
\RouteGroup{Secuencias completas}
\RouteRow{\(w_6\)}{81 3 0 81 0 3 0 81 3 81 3 0 81 0 3 0 81 3}
\RouteRow{Tríadas estables por bloque}{0 1 2 3 4 5 6 6 8 9 10 11 12 13 14 15 15 17}
\end{tabularx}\endgroup\par\vspace{5pt}\noindent{\fontsize{8.4}{10}\selectfont\hyperref[trace-126]{Procedencia y códigos originales}\hfill Ficha completa · 126/324}
\clearpage\phantomsection\label{nav:vi-ruta-127}
% VI-CATALOG-ROW rutas 127
\pdfbookmark[3]{r4c5\_h+1v-1}{route-127}
\RouteTitle{Ruta 127}{r4c5\_\allowbreak{}h+1v-1}
\noindent Celda 32\quad (4, 5)\qquad Orientación \(h=1,\ v=-1\)\hfill\hyperref[sec:vi-indice-rutas]{Índice}\par\vspace{4pt}
\begingroup\fontsize{10.2}{12.5}\selectfont\setlength{\tabcolsep}{5pt}\renewcommand{\arraystretch}{1.1}\rowcolors{1}{routepale}{white}\noindent\begin{tabularx}{\linewidth}{@{}>{\raggedright\arraybackslash}p{.345\linewidth}>{\raggedright\arraybackslash}X@{}}
\RouteGroup{Identidad estructural}
\RouteRow{Clase y multisección}{quark\_\allowbreak{}down\_\allowbreak{}G1\quad /\quad quark\_\allowbreak{}down\_\allowbreak{}G1}
\RouteRow{Colección y sector}{quark de tipo abajo; quark}
\RouteRow{Clasificación física}{quark de tipo abajo}
\RouteRow{Generación, profundidad y color}{1; G1; B}
\RouteRow{Ruta primaria y órbita de inversión}{\(B_1\): sí\qquad 4,\allowbreak{}5;\allowbreak{}6,\allowbreak{}5}
\RouteRow{Firma de clasificación}{28|\allowbreak{}6|\allowbreak{}0|\allowbreak{}0|\allowbreak{}-12|\allowbreak{}+++-\kern0pt{}-\kern0pt{}-+++-\kern0pt{}-\kern0pt{}-}
\RouteGroup{Retornos, memoria y transporte}
\RouteRow{Retornos con fase}{clase: 27; órbita: 27; celda: 27}
\RouteRow{Retorno cinemático}{en 108: 54; extensión: 216}
\RouteRow{Giros y cambios de hoja}{71; 107}
\RouteRow{Acarreo y diferimiento}{deuda total: 0; final: 0; bloques diferidos: 0}
\RouteRow{Tríadas finales; persistencia de Pareto}{17; sí}
\RouteRow{\(K\longrightarrow K+9\)}{repeticiones: 9; cambios de memoria: 9}
\RouteGroup{Firma, fase y diferencias}
\RouteRow{\(n_A,\ s_C,\ \kappa_0,\ \nu_{120},\ \nu_{270}\)}{20, 0, 0, 0, -8}
\RouteRow{\(\tau_{12}\); firma histórica \(\mathbb Z^5\)}{+0+-0-+0+-0-; 20|\allowbreak{}0|\allowbreak{}0|\allowbreak{}0|\allowbreak{}-8}
\RouteRow{\(D_1,\ D_3,\ D_4,\ D_{27},\ D_{36}\)}{24, 20, 24, 24, 16}
\RouteRow{\(\Omega_{90/120},\ P_6\)}{40, 4}
\RouteGroup{Lectores y factores de acción}
\RouteRow{\(K_D\quad /\quad K_\Omega\)}{(40,\allowbreak{}24,\allowbreak{}2)\quad /\quad (40,\allowbreak{}54,\allowbreak{}4)}
\RouteRow{\(\kappa_D\)}{39.82508\allowbreak{}59311825\allowbreak{}73056173\allowbreak{}26}
\RouteRow{\(\kappa_\Omega\)}{39.60638\allowbreak{}77469494\allowbreak{}43306362\allowbreak{}81}
\RouteRow{\(R_{\mathrm{act},D}\)}{1.000000\allowbreak{}02761000\allowbreak{}61595142\allowbreak{}15}
\RouteRow{\(R_{\mathrm{act},\Omega}\)}{1.000000\allowbreak{}02745838\allowbreak{}66926514\allowbreak{}00}
\RouteGroup{Separación de prefijos}
\RouteRow{Área de ambigüedad; área \(\log_2\)}{18; 0.0}
\RouteRow{Primer bloque de separación}{órbita: 1; celda: 1; ruta: 1}
\RouteRow{Ambigüedad final}{1}
\RouteGroup{Secuencias completas}
\RouteRow{\(w_6\)}{0 81 3 0 84 0 81 3 0 0 81 3 0 84 0 81 3 0}
\RouteRow{Tríadas estables por bloque}{0 1 2 3 4 5 6 7 8 9 10 11 12 13 14 15 16 17}
\end{tabularx}\endgroup\par\vspace{5pt}\noindent{\fontsize{8.4}{10}\selectfont\hyperref[trace-127]{Procedencia y códigos originales}\hfill Ficha completa · 127/324}
\clearpage\phantomsection\label{nav:vi-ruta-128}
% VI-CATALOG-ROW rutas 128
\pdfbookmark[3]{r4c5\_h+1v+1}{route-128}
\RouteTitle{Ruta 128}{r4c5\_\allowbreak{}h+1v+1}
\noindent Celda 32\quad (4, 5)\qquad Orientación \(h=1,\ v=1\)\hfill\hyperref[sec:vi-indice-rutas]{Índice}\par\vspace{4pt}
\begingroup\fontsize{10.2}{12.5}\selectfont\setlength{\tabcolsep}{5pt}\renewcommand{\arraystretch}{1.1}\rowcolors{1}{routepale}{white}\noindent\begin{tabularx}{\linewidth}{@{}>{\raggedright\arraybackslash}p{.345\linewidth}>{\raggedright\arraybackslash}X@{}}
\RouteGroup{Identidad estructural}
\RouteRow{Clase y multisección}{quark\_\allowbreak{}down\_\allowbreak{}G1\quad /\quad quark\_\allowbreak{}down\_\allowbreak{}G1}
\RouteRow{Colección y sector}{quark de tipo abajo; quark}
\RouteRow{Clasificación física}{quark de tipo abajo}
\RouteRow{Generación, profundidad y color}{1; G1; B}
\RouteRow{Ruta primaria y órbita de inversión}{\(B_1\): no\qquad 4,\allowbreak{}5;\allowbreak{}6,\allowbreak{}5}
\RouteRow{Firma de clasificación}{28|\allowbreak{}6|\allowbreak{}0|\allowbreak{}0|\allowbreak{}-12|\allowbreak{}+++-\kern0pt{}-\kern0pt{}-+++-\kern0pt{}-\kern0pt{}-}
\RouteGroup{Retornos, memoria y transporte}
\RouteRow{Retornos con fase}{clase: 27; órbita: 27; celda: 27}
\RouteRow{Retorno cinemático}{en 108: 54; extensión: 216}
\RouteRow{Giros y cambios de hoja}{71; 107}
\RouteRow{Acarreo y diferimiento}{deuda total: 0; final: 0; bloques diferidos: 0}
\RouteRow{Tríadas finales; persistencia de Pareto}{17; sí}
\RouteRow{\(K\longrightarrow K+9\)}{repeticiones: 9; cambios de memoria: 9}
\RouteGroup{Firma, fase y diferencias}
\RouteRow{\(n_A,\ s_C,\ \kappa_0,\ \nu_{120},\ \nu_{270}\)}{28, 6, 0, 0, -12}
\RouteRow{\(\tau_{12}\); firma histórica \(\mathbb Z^5\)}{+++-\kern0pt{}-\kern0pt{}-+++-\kern0pt{}-\kern0pt{}-; 28|\allowbreak{}6|\allowbreak{}0|\allowbreak{}0|\allowbreak{}-12}
\RouteRow{\(D_1,\ D_3,\ D_4,\ D_{27},\ D_{36}\)}{66, 64, 64, 60, 40}
\RouteRow{\(\Omega_{90/120},\ P_6\)}{80, 6}
\RouteGroup{Lectores y factores de acción}
\RouteRow{\(K_D\quad /\quad K_\Omega\)}{(100,\allowbreak{}66,\allowbreak{}0)\quad /\quad (80,\allowbreak{}54,\allowbreak{}6)}
\RouteRow{\(\kappa_D\)}{99.51837\allowbreak{}47304272\allowbreak{}69134179\allowbreak{}18}
\RouteRow{\(\kappa_\Omega\)}{79.60661\allowbreak{}01397948\allowbreak{}27586470\allowbreak{}91}
\RouteRow{\(R_{\mathrm{act},D}\)}{1.000000\allowbreak{}06899427\allowbreak{}66450218\allowbreak{}95}
\RouteRow{\(R_{\mathrm{act},\Omega}\)}{1.000000\allowbreak{}05518981\allowbreak{}25318591\allowbreak{}40}
\RouteGroup{Separación de prefijos}
\RouteRow{Área de ambigüedad; área \(\log_2\)}{18; 0.0}
\RouteRow{Primer bloque de separación}{órbita: 1; celda: 1; ruta: 1}
\RouteRow{Ambigüedad final}{1}
\RouteGroup{Secuencias completas}
\RouteRow{\(w_6\)}{24 243 657 24 246 264 90 570 264 24 243 657 24 246 264 90 570 264}
\RouteRow{Tríadas estables por bloque}{0 1 2 3 4 5 6 7 8 9 10 11 12 13 14 15 16 17}
\end{tabularx}\endgroup\par\vspace{5pt}\noindent{\fontsize{8.4}{10}\selectfont\hyperref[trace-128]{Procedencia y códigos originales}\hfill Ficha completa · 128/324}
\clearpage\phantomsection\label{nav:vi-ruta-129}
% VI-CATALOG-ROW rutas 129
\pdfbookmark[2]{Celda 33}{cell-33}
\pdfbookmark[3]{r4c6\_h-1v-1}{route-129}
\RouteTitle{Ruta 129}{r4c6\_\allowbreak{}h-1v-1}
\noindent Celda 33\quad (4, 6)\qquad Orientación \(h=-1,\ v=-1\)\hfill\hyperref[sec:vi-indice-rutas]{Índice}\par\vspace{4pt}
\begingroup\fontsize{10.2}{12.5}\selectfont\setlength{\tabcolsep}{5pt}\renewcommand{\arraystretch}{1.1}\rowcolors{1}{routepale}{white}\noindent\begin{tabularx}{\linewidth}{@{}>{\raggedright\arraybackslash}p{.345\linewidth}>{\raggedright\arraybackslash}X@{}}
\RouteGroup{Identidad estructural}
\RouteRow{Clase y multisección}{quark\_\allowbreak{}up\_\allowbreak{}G1\quad /\quad quark\_\allowbreak{}up\_\allowbreak{}G1}
\RouteRow{Colección y sector}{quark de tipo arriba; quark}
\RouteRow{Clasificación física}{quark de tipo arriba}
\RouteRow{Generación, profundidad y color}{1; G1; G}
\RouteRow{Ruta primaria y órbita de inversión}{\(B_1\): no\qquad 4,\allowbreak{}6;\allowbreak{}6,\allowbreak{}4}
\RouteRow{Firma de clasificación}{28|\allowbreak{}-6|\allowbreak{}10|\allowbreak{}0|\allowbreak{}-12|\allowbreak{}+++-\kern0pt{}-\kern0pt{}-+++-\kern0pt{}-\kern0pt{}-}
\RouteGroup{Retornos, memoria y transporte}
\RouteRow{Retornos con fase}{clase: 27; órbita: 27; celda: 27}
\RouteRow{Retorno cinemático}{en 108: 54; extensión: 216}
\RouteRow{Giros y cambios de hoja}{71; 107}
\RouteRow{Acarreo y diferimiento}{deuda total: 1; final: 1; bloques diferidos: 1}
\RouteRow{Tríadas finales; persistencia de Pareto}{16; no}
\RouteRow{\(K\longrightarrow K+9\)}{repeticiones: 9; cambios de memoria: 9}
\RouteGroup{Firma, fase y diferencias}
\RouteRow{\(n_A,\ s_C,\ \kappa_0,\ \nu_{120},\ \nu_{270}\)}{28, 6, 0, 0, -12}
\RouteRow{\(\tau_{12}\); firma histórica \(\mathbb Z^5\)}{+++-\kern0pt{}-\kern0pt{}-+++-\kern0pt{}-\kern0pt{}-; 28|\allowbreak{}6|\allowbreak{}0|\allowbreak{}0|\allowbreak{}-12}
\RouteRow{\(D_1,\ D_3,\ D_4,\ D_{27},\ D_{36}\)}{66, 60, 64, 48, 32}
\RouteRow{\(\Omega_{90/120},\ P_6\)}{80, 6}
\RouteGroup{Lectores y factores de acción}
\RouteRow{\(K_D\quad /\quad K_\Omega\)}{(80,\allowbreak{}66,\allowbreak{}2)\quad /\quad (80,\allowbreak{}54,\allowbreak{}6)}
\RouteRow{\(\kappa_D\)}{79.51859\allowbreak{}71232726\allowbreak{}53414287\allowbreak{}29}
\RouteRow{\(\kappa_\Omega\)}{79.60661\allowbreak{}01397948\allowbreak{}27586470\allowbreak{}91}
\RouteRow{\(R_{\mathrm{act},D}\)}{1.000000\allowbreak{}05512879\allowbreak{}47092488\allowbreak{}80}
\RouteRow{\(R_{\mathrm{act},\Omega}\)}{1.000000\allowbreak{}05518981\allowbreak{}25318591\allowbreak{}40}
\RouteGroup{Separación de prefijos}
\RouteRow{Área de ambigüedad; área \(\log_2\)}{18; 0.0}
\RouteRow{Primer bloque de separación}{órbita: 1; celda: 1; ruta: 1}
\RouteRow{Ambigüedad final}{1}
\RouteGroup{Secuencias completas}
\RouteRow{\(w_6\)}{423 15 486 423 12 99 327 498 99 423 15 486 423 12 99 327 498 99}
\RouteRow{Tríadas estables por bloque}{0 1 2 3 4 5 6 7 8 9 10 11 12 13 14 15 16 16}
\end{tabularx}\endgroup\par\vspace{5pt}\noindent{\fontsize{8.4}{10}\selectfont\hyperref[trace-129]{Procedencia y códigos originales}\hfill Ficha completa · 129/324}
\clearpage\phantomsection\label{nav:vi-ruta-130}
% VI-CATALOG-ROW rutas 130
\pdfbookmark[3]{r4c6\_h-1v+1}{route-130}
\RouteTitle{Ruta 130}{r4c6\_\allowbreak{}h-1v+1}
\noindent Celda 33\quad (4, 6)\qquad Orientación \(h=-1,\ v=1\)\hfill\hyperref[sec:vi-indice-rutas]{Índice}\par\vspace{4pt}
\begingroup\fontsize{10.2}{12.5}\selectfont\setlength{\tabcolsep}{5pt}\renewcommand{\arraystretch}{1.1}\rowcolors{1}{routepale}{white}\noindent\begin{tabularx}{\linewidth}{@{}>{\raggedright\arraybackslash}p{.345\linewidth}>{\raggedright\arraybackslash}X@{}}
\RouteGroup{Identidad estructural}
\RouteRow{Clase y multisección}{quark\_\allowbreak{}up\_\allowbreak{}G1\quad /\quad quark\_\allowbreak{}up\_\allowbreak{}G1}
\RouteRow{Colección y sector}{quark de tipo arriba; quark}
\RouteRow{Clasificación física}{quark de tipo arriba}
\RouteRow{Generación, profundidad y color}{1; G1; G}
\RouteRow{Ruta primaria y órbita de inversión}{\(B_1\): no\qquad 4,\allowbreak{}6;\allowbreak{}6,\allowbreak{}4}
\RouteRow{Firma de clasificación}{28|\allowbreak{}-6|\allowbreak{}10|\allowbreak{}0|\allowbreak{}-12|\allowbreak{}+++-\kern0pt{}-\kern0pt{}-+++-\kern0pt{}-\kern0pt{}-}
\RouteGroup{Retornos, memoria y transporte}
\RouteRow{Retornos con fase}{clase: 27; órbita: 27; celda: 27}
\RouteRow{Retorno cinemático}{en 108: 54; extensión: 216}
\RouteRow{Giros y cambios de hoja}{71; 107}
\RouteRow{Acarreo y diferimiento}{deuda total: 1; final: 0; bloques diferidos: 1}
\RouteRow{Tríadas finales; persistencia de Pareto}{17; no}
\RouteRow{\(K\longrightarrow K+9\)}{repeticiones: 9; cambios de memoria: 9}
\RouteGroup{Firma, fase y diferencias}
\RouteRow{\(n_A,\ s_C,\ \kappa_0,\ \nu_{120},\ \nu_{270}\)}{50, 4, -2, 2, -8}
\RouteRow{\(\tau_{12}\); firma histórica \(\mathbb Z^5\)}{+++-\kern0pt{}-0+++-\kern0pt{}-0; 50|\allowbreak{}4|\allowbreak{}-2|\allowbreak{}2|\allowbreak{}-8}
\RouteRow{\(D_1,\ D_3,\ D_4,\ D_{27},\ D_{36}\)}{84, 24, 84, 24, 16}
\RouteRow{\(\Omega_{90/120},\ P_6\)}{56, 5}
\RouteGroup{Lectores y factores de acción}
\RouteRow{\(K_D\quad /\quad K_\Omega\)}{(40,\allowbreak{}84,\allowbreak{}30)\quad /\quad (56,\allowbreak{}54,\allowbreak{}5)}
\RouteRow{\(\kappa_D\)}{39.39035\allowbreak{}82768609\allowbreak{}24917849\allowbreak{}59}
\RouteRow{\(\kappa_\Omega\)}{55.60649\allowbreak{}89433721\allowbreak{}35446416\allowbreak{}86}
\RouteRow{\(R_{\mathrm{act},D}\)}{1.000000\allowbreak{}02730861\allowbreak{}73967087\allowbreak{}05}
\RouteRow{\(R_{\mathrm{act},\Omega}\)}{1.000000\allowbreak{}03855097\allowbreak{}23541416\allowbreak{}45}
\RouteGroup{Separación de prefijos}
\RouteRow{Área de ambigüedad; área \(\log_2\)}{18; 0.0}
\RouteRow{Primer bloque de separación}{órbita: 1; celda: 1; ruta: 1}
\RouteRow{Ambigüedad final}{1}
\RouteGroup{Secuencias completas}
\RouteRow{\(w_6\)}{420 258 414 420 255 252 333 255 252 420 258 414 420 255 252 333 255 252}
\RouteRow{Tríadas estables por bloque}{0 1 2 3 4 5 6 7 8 9 10 11 12 13 14 15 15 17}
\end{tabularx}\endgroup\par\vspace{5pt}\noindent{\fontsize{8.4}{10}\selectfont\hyperref[trace-130]{Procedencia y códigos originales}\hfill Ficha completa · 130/324}
\clearpage\phantomsection\label{nav:vi-ruta-131}
% VI-CATALOG-ROW rutas 131
\pdfbookmark[3]{r4c6\_h+1v-1}{route-131}
\RouteTitle{Ruta 131}{r4c6\_\allowbreak{}h+1v-1}
\noindent Celda 33\quad (4, 6)\qquad Orientación \(h=1,\ v=-1\)\hfill\hyperref[sec:vi-indice-rutas]{Índice}\par\vspace{4pt}
\begingroup\fontsize{10.2}{12.5}\selectfont\setlength{\tabcolsep}{5pt}\renewcommand{\arraystretch}{1.1}\rowcolors{1}{routepale}{white}\noindent\begin{tabularx}{\linewidth}{@{}>{\raggedright\arraybackslash}p{.345\linewidth}>{\raggedright\arraybackslash}X@{}}
\RouteGroup{Identidad estructural}
\RouteRow{Clase y multisección}{quark\_\allowbreak{}up\_\allowbreak{}G1\quad /\quad quark\_\allowbreak{}up\_\allowbreak{}G1}
\RouteRow{Colección y sector}{quark de tipo arriba; quark}
\RouteRow{Clasificación física}{quark de tipo arriba}
\RouteRow{Generación, profundidad y color}{1; G1; G}
\RouteRow{Ruta primaria y órbita de inversión}{\(B_1\): no\qquad 4,\allowbreak{}6;\allowbreak{}6,\allowbreak{}4}
\RouteRow{Firma de clasificación}{28|\allowbreak{}-6|\allowbreak{}10|\allowbreak{}0|\allowbreak{}-12|\allowbreak{}+++-\kern0pt{}-\kern0pt{}-+++-\kern0pt{}-\kern0pt{}-}
\RouteGroup{Retornos, memoria y transporte}
\RouteRow{Retornos con fase}{clase: 27; órbita: 27; celda: 27}
\RouteRow{Retorno cinemático}{en 108: 54; extensión: 216}
\RouteRow{Giros y cambios de hoja}{71; 107}
\RouteRow{Acarreo y diferimiento}{deuda total: 0; final: 0; bloques diferidos: 0}
\RouteRow{Tríadas finales; persistencia de Pareto}{17; no}
\RouteRow{\(K\longrightarrow K+9\)}{repeticiones: 9; cambios de memoria: 9}
\RouteGroup{Firma, fase y diferencias}
\RouteRow{\(n_A,\ s_C,\ \kappa_0,\ \nu_{120},\ \nu_{270}\)}{50, -4, 2, -2, -8}
\RouteRow{\(\tau_{12}\); firma histórica \(\mathbb Z^5\)}{++0-\kern0pt{}-\kern0pt{}-++0-\kern0pt{}-\kern0pt{}-; 50|\allowbreak{}-4|\allowbreak{}2|\allowbreak{}-2|\allowbreak{}-8}
\RouteRow{\(D_1,\ D_3,\ D_4,\ D_{27},\ D_{36}\)}{84, 24, 84, 24, 16}
\RouteRow{\(\Omega_{90/120},\ P_6\)}{56, 5}
\RouteGroup{Lectores y factores de acción}
\RouteRow{\(K_D\quad /\quad K_\Omega\)}{(40,\allowbreak{}84,\allowbreak{}30)\quad /\quad (56,\allowbreak{}54,\allowbreak{}5)}
\RouteRow{\(\kappa_D\)}{39.39035\allowbreak{}82768609\allowbreak{}24917849\allowbreak{}59}
\RouteRow{\(\kappa_\Omega\)}{55.60649\allowbreak{}89433721\allowbreak{}35446416\allowbreak{}86}
\RouteRow{\(R_{\mathrm{act},D}\)}{1.000000\allowbreak{}02730861\allowbreak{}73967087\allowbreak{}05}
\RouteRow{\(R_{\mathrm{act},\Omega}\)}{1.000000\allowbreak{}03855097\allowbreak{}23541416\allowbreak{}45}
\RouteGroup{Separación de prefijos}
\RouteRow{Área de ambigüedad; área \(\log_2\)}{18; 0.0}
\RouteRow{Primer bloque de separación}{órbita: 1; celda: 1; ruta: 1}
\RouteRow{Ambigüedad final}{1}
\RouteGroup{Secuencias completas}
\RouteRow{\(w_6\)}{333 255 252 333 258 414 420 258 414 333 255 252 333 258 414 420 258 414}
\RouteRow{Tríadas estables por bloque}{0 1 2 3 4 5 6 7 8 9 10 11 12 13 14 15 16 17}
\end{tabularx}\endgroup\par\vspace{5pt}\noindent{\fontsize{8.4}{10}\selectfont\hyperref[trace-131]{Procedencia y códigos originales}\hfill Ficha completa · 131/324}
\clearpage\phantomsection\label{nav:vi-ruta-132}
% VI-CATALOG-ROW rutas 132
\pdfbookmark[3]{r4c6\_h+1v+1}{route-132}
\RouteTitle{Ruta 132}{r4c6\_\allowbreak{}h+1v+1}
\noindent Celda 33\quad (4, 6)\qquad Orientación \(h=1,\ v=1\)\hfill\hyperref[sec:vi-indice-rutas]{Índice}\par\vspace{4pt}
\begingroup\fontsize{10.2}{12.5}\selectfont\setlength{\tabcolsep}{5pt}\renewcommand{\arraystretch}{1.1}\rowcolors{1}{routepale}{white}\noindent\begin{tabularx}{\linewidth}{@{}>{\raggedright\arraybackslash}p{.345\linewidth}>{\raggedright\arraybackslash}X@{}}
\RouteGroup{Identidad estructural}
\RouteRow{Clase y multisección}{quark\_\allowbreak{}up\_\allowbreak{}G1\quad /\quad quark\_\allowbreak{}up\_\allowbreak{}G1}
\RouteRow{Colección y sector}{quark de tipo arriba; quark}
\RouteRow{Clasificación física}{quark de tipo arriba}
\RouteRow{Generación, profundidad y color}{1; G1; G}
\RouteRow{Ruta primaria y órbita de inversión}{\(B_1\): sí\qquad 4,\allowbreak{}6;\allowbreak{}6,\allowbreak{}4}
\RouteRow{Firma de clasificación}{28|\allowbreak{}-6|\allowbreak{}10|\allowbreak{}0|\allowbreak{}-12|\allowbreak{}+++-\kern0pt{}-\kern0pt{}-+++-\kern0pt{}-\kern0pt{}-}
\RouteGroup{Retornos, memoria y transporte}
\RouteRow{Retornos con fase}{clase: 27; órbita: 27; celda: 27}
\RouteRow{Retorno cinemático}{en 108: 54; extensión: 216}
\RouteRow{Giros y cambios de hoja}{71; 107}
\RouteRow{Acarreo y diferimiento}{deuda total: 2; final: 1; bloques diferidos: 2}
\RouteRow{Tríadas finales; persistencia de Pareto}{16; no}
\RouteRow{\(K\longrightarrow K+9\)}{repeticiones: 9; cambios de memoria: 9}
\RouteGroup{Firma, fase y diferencias}
\RouteRow{\(n_A,\ s_C,\ \kappa_0,\ \nu_{120},\ \nu_{270}\)}{28, -6, 0, 0, -12}
\RouteRow{\(\tau_{12}\); firma histórica \(\mathbb Z^5\)}{+++-\kern0pt{}-\kern0pt{}-+++-\kern0pt{}-\kern0pt{}-; 28|\allowbreak{}-6|\allowbreak{}0|\allowbreak{}0|\allowbreak{}-12}
\RouteRow{\(D_1,\ D_3,\ D_4,\ D_{27},\ D_{36}\)}{66, 60, 64, 48, 32}
\RouteRow{\(\Omega_{90/120},\ P_6\)}{80, 6}
\RouteGroup{Lectores y factores de acción}
\RouteRow{\(K_D\quad /\quad K_\Omega\)}{(80,\allowbreak{}66,\allowbreak{}2)\quad /\quad (80,\allowbreak{}54,\allowbreak{}6)}
\RouteRow{\(\kappa_D\)}{79.51859\allowbreak{}71232726\allowbreak{}53414287\allowbreak{}29}
\RouteRow{\(\kappa_\Omega\)}{79.60661\allowbreak{}01397948\allowbreak{}27586470\allowbreak{}91}
\RouteRow{\(R_{\mathrm{act},D}\)}{1.000000\allowbreak{}05512879\allowbreak{}47092488\allowbreak{}80}
\RouteRow{\(R_{\mathrm{act},\Omega}\)}{1.000000\allowbreak{}05518981\allowbreak{}25318591\allowbreak{}40}
\RouteGroup{Separación de prefijos}
\RouteRow{Área de ambigüedad; área \(\log_2\)}{18; 0.0}
\RouteRow{Primer bloque de separación}{órbita: 1; celda: 1; ruta: 1}
\RouteRow{Ambigüedad final}{1}
\RouteGroup{Secuencias completas}
\RouteRow{\(w_6\)}{327 498 99 327 501 486 423 15 486 327 498 99 327 501 486 423 15 486}
\RouteRow{Tríadas estables por bloque}{0 1 2 3 4 5 6 7 7 9 10 11 12 13 14 15 16 16}
\end{tabularx}\endgroup\par\vspace{5pt}\noindent{\fontsize{8.4}{10}\selectfont\hyperref[trace-132]{Procedencia y códigos originales}\hfill Ficha completa · 132/324}
\clearpage\phantomsection\label{nav:vi-ruta-133}
% VI-CATALOG-ROW rutas 133
\pdfbookmark[2]{Celda 34}{cell-34}
\pdfbookmark[3]{r4c7\_h-1v-1}{route-133}
\RouteTitle{Ruta 133}{r4c7\_\allowbreak{}h-1v-1}
\noindent Celda 34\quad (4, 7)\qquad Orientación \(h=-1,\ v=-1\)\hfill\hyperref[sec:vi-indice-rutas]{Índice}\par\vspace{4pt}
\begingroup\fontsize{10.2}{12.5}\selectfont\setlength{\tabcolsep}{5pt}\renewcommand{\arraystretch}{1.1}\rowcolors{1}{routepale}{white}\noindent\begin{tabularx}{\linewidth}{@{}>{\raggedright\arraybackslash}p{.345\linewidth}>{\raggedright\arraybackslash}X@{}}
\RouteGroup{Identidad estructural}
\RouteRow{Clase y multisección}{quark\_\allowbreak{}down\_\allowbreak{}G2\quad /\quad quark\_\allowbreak{}down\_\allowbreak{}G2}
\RouteRow{Colección y sector}{quark de tipo abajo; quark}
\RouteRow{Clasificación física}{quark de tipo abajo}
\RouteRow{Generación, profundidad y color}{2; G2; R}
\RouteRow{Ruta primaria y órbita de inversión}{\(B_1\): no\qquad 4,\allowbreak{}7;\allowbreak{}6,\allowbreak{}3}
\RouteRow{Firma de clasificación}{30|\allowbreak{}-6|\allowbreak{}0|\allowbreak{}-2|\allowbreak{}-8|\allowbreak{}++0-\kern0pt{}-\kern0pt{}-++0-\kern0pt{}-\kern0pt{}-}
\RouteGroup{Retornos, memoria y transporte}
\RouteRow{Retornos con fase}{clase: 27; órbita: 27; celda: 27}
\RouteRow{Retorno cinemático}{en 108: 54; extensión: 216}
\RouteRow{Giros y cambios de hoja}{71; 107}
\RouteRow{Acarreo y diferimiento}{deuda total: 2; final: 1; bloques diferidos: 2}
\RouteRow{Tríadas finales; persistencia de Pareto}{16; sí}
\RouteRow{\(K\longrightarrow K+9\)}{repeticiones: 9; cambios de memoria: 9}
\RouteGroup{Firma, fase y diferencias}
\RouteRow{\(n_A,\ s_C,\ \kappa_0,\ \nu_{120},\ \nu_{270}\)}{30, 6, -2, 0, -8}
\RouteRow{\(\tau_{12}\); firma histórica \(\mathbb Z^5\)}{+++-\kern0pt{}-0+++-\kern0pt{}-0; 30|\allowbreak{}6|\allowbreak{}-2|\allowbreak{}0|\allowbreak{}-8}
\RouteRow{\(D_1,\ D_3,\ D_4,\ D_{27},\ D_{36}\)}{54, 56, 70, 48, 32}
\RouteRow{\(\Omega_{90/120},\ P_6\)}{56, 5}
\RouteGroup{Lectores y factores de acción}
\RouteRow{\(K_D\quad /\quad K_\Omega\)}{(80,\allowbreak{}54,\allowbreak{}7)\quad /\quad (56,\allowbreak{}54,\allowbreak{}5)}
\RouteRow{\(\kappa_D\)}{79.60672\allowbreak{}13362175\allowbreak{}19726524\allowbreak{}96}
\RouteRow{\(\kappa_\Omega\)}{55.60649\allowbreak{}89433721\allowbreak{}35446416\allowbreak{}86}
\RouteRow{\(R_{\mathrm{act},D}\)}{1.000000\allowbreak{}05518988\allowbreak{}96223153\allowbreak{}37}
\RouteRow{\(R_{\mathrm{act},\Omega}\)}{1.000000\allowbreak{}03855097\allowbreak{}23541416\allowbreak{}45}
\RouteGroup{Separación de prefijos}
\RouteRow{Área de ambigüedad; área \(\log_2\)}{18; 0.0}
\RouteRow{Primer bloque de separación}{órbita: 1; celda: 1; ruta: 1}
\RouteRow{Ambigüedad final}{1}
\RouteGroup{Secuencias completas}
\RouteRow{\(w_6\)}{486 423 15 486 429 243 657 24 243 486 423 15 486 429 243 657 24 243}
\RouteRow{Tríadas estables por bloque}{0 1 2 3 4 5 6 7 8 9 10 11 12 13 13 15 16 16}
\end{tabularx}\endgroup\par\vspace{5pt}\noindent{\fontsize{8.4}{10}\selectfont\hyperref[trace-133]{Procedencia y códigos originales}\hfill Ficha completa · 133/324}
\clearpage\phantomsection\label{nav:vi-ruta-134}
% VI-CATALOG-ROW rutas 134
\pdfbookmark[3]{r4c7\_h-1v+1}{route-134}
\RouteTitle{Ruta 134}{r4c7\_\allowbreak{}h-1v+1}
\noindent Celda 34\quad (4, 7)\qquad Orientación \(h=-1,\ v=1\)\hfill\hyperref[sec:vi-indice-rutas]{Índice}\par\vspace{4pt}
\begingroup\fontsize{10.2}{12.5}\selectfont\setlength{\tabcolsep}{5pt}\renewcommand{\arraystretch}{1.1}\rowcolors{1}{routepale}{white}\noindent\begin{tabularx}{\linewidth}{@{}>{\raggedright\arraybackslash}p{.345\linewidth}>{\raggedright\arraybackslash}X@{}}
\RouteGroup{Identidad estructural}
\RouteRow{Clase y multisección}{quark\_\allowbreak{}down\_\allowbreak{}G2\quad /\quad quark\_\allowbreak{}down\_\allowbreak{}G2}
\RouteRow{Colección y sector}{quark de tipo abajo; quark}
\RouteRow{Clasificación física}{quark de tipo abajo}
\RouteRow{Generación, profundidad y color}{2; G2; R}
\RouteRow{Ruta primaria y órbita de inversión}{\(B_1\): sí\qquad 4,\allowbreak{}7;\allowbreak{}6,\allowbreak{}3}
\RouteRow{Firma de clasificación}{30|\allowbreak{}-6|\allowbreak{}0|\allowbreak{}-2|\allowbreak{}-8|\allowbreak{}++0-\kern0pt{}-\kern0pt{}-++0-\kern0pt{}-\kern0pt{}-}
\RouteGroup{Retornos, memoria y transporte}
\RouteRow{Retornos con fase}{clase: 27; órbita: 27; celda: 27}
\RouteRow{Retorno cinemático}{en 108: 54; extensión: 216}
\RouteRow{Giros y cambios de hoja}{71; 107}
\RouteRow{Acarreo y diferimiento}{deuda total: 0; final: 0; bloques diferidos: 0}
\RouteRow{Tríadas finales; persistencia de Pareto}{17; sí}
\RouteRow{\(K\longrightarrow K+9\)}{repeticiones: 9; cambios de memoria: 9}
\RouteGroup{Firma, fase y diferencias}
\RouteRow{\(n_A,\ s_C,\ \kappa_0,\ \nu_{120},\ \nu_{270}\)}{14, -4, 0, 0, -12}
\RouteRow{\(\tau_{12}\); firma histórica \(\mathbb Z^5\)}{+++-\kern0pt{}-\kern0pt{}-+++-\kern0pt{}-\kern0pt{}-; 14|\allowbreak{}-4|\allowbreak{}0|\allowbreak{}0|\allowbreak{}-12}
\RouteRow{\(D_1,\ D_3,\ D_4,\ D_{27},\ D_{36}\)}{78, 28, 78, 36, 24}
\RouteRow{\(\Omega_{90/120},\ P_6\)}{80, 6}
\RouteGroup{Lectores y factores de acción}
\RouteRow{\(K_D\quad /\quad K_\Omega\)}{(60,\allowbreak{}78,\allowbreak{}25)\quad /\quad (80,\allowbreak{}54,\allowbreak{}6)}
\RouteRow{\(\kappa_D\)}{59.43358\allowbreak{}64101631\allowbreak{}67023563\allowbreak{}04}
\RouteRow{\(\kappa_\Omega\)}{79.60661\allowbreak{}01397948\allowbreak{}27586470\allowbreak{}91}
\RouteRow{\(R_{\mathrm{act},D}\)}{1.000000\allowbreak{}04120422\allowbreak{}24053449\allowbreak{}05}
\RouteRow{\(R_{\mathrm{act},\Omega}\)}{1.000000\allowbreak{}05518981\allowbreak{}25318591\allowbreak{}40}
\RouteGroup{Separación de prefijos}
\RouteRow{Área de ambigüedad; área \(\log_2\)}{18; 0.0}
\RouteRow{Primer bloque de separación}{órbita: 1; celda: 1; ruta: 1}
\RouteRow{Ambigüedad final}{1}
\RouteGroup{Secuencias completas}
\RouteRow{\(w_6\)}{507 504 585 507 510 510 666 672 510 507 504 585 507 510 510 666 672 510}
\RouteRow{Tríadas estables por bloque}{0 1 2 3 4 5 6 7 8 9 10 11 12 13 14 15 16 17}
\end{tabularx}\endgroup\par\vspace{5pt}\noindent{\fontsize{8.4}{10}\selectfont\hyperref[trace-134]{Procedencia y códigos originales}\hfill Ficha completa · 134/324}
\clearpage\phantomsection\label{nav:vi-ruta-135}
% VI-CATALOG-ROW rutas 135
\pdfbookmark[3]{r4c7\_h+1v-1}{route-135}
\RouteTitle{Ruta 135}{r4c7\_\allowbreak{}h+1v-1}
\noindent Celda 34\quad (4, 7)\qquad Orientación \(h=1,\ v=-1\)\hfill\hyperref[sec:vi-indice-rutas]{Índice}\par\vspace{4pt}
\begingroup\fontsize{10.2}{12.5}\selectfont\setlength{\tabcolsep}{5pt}\renewcommand{\arraystretch}{1.1}\rowcolors{1}{routepale}{white}\noindent\begin{tabularx}{\linewidth}{@{}>{\raggedright\arraybackslash}p{.345\linewidth}>{\raggedright\arraybackslash}X@{}}
\RouteGroup{Identidad estructural}
\RouteRow{Clase y multisección}{quark\_\allowbreak{}down\_\allowbreak{}G2\quad /\quad quark\_\allowbreak{}down\_\allowbreak{}G2}
\RouteRow{Colección y sector}{quark de tipo abajo; quark}
\RouteRow{Clasificación física}{quark de tipo abajo}
\RouteRow{Generación, profundidad y color}{2; G2; R}
\RouteRow{Ruta primaria y órbita de inversión}{\(B_1\): no\qquad 4,\allowbreak{}7;\allowbreak{}6,\allowbreak{}3}
\RouteRow{Firma de clasificación}{30|\allowbreak{}-6|\allowbreak{}0|\allowbreak{}-2|\allowbreak{}-8|\allowbreak{}++0-\kern0pt{}-\kern0pt{}-++0-\kern0pt{}-\kern0pt{}-}
\RouteGroup{Retornos, memoria y transporte}
\RouteRow{Retornos con fase}{clase: 27; órbita: 27; celda: 27}
\RouteRow{Retorno cinemático}{en 108: 54; extensión: 216}
\RouteRow{Giros y cambios de hoja}{71; 107}
\RouteRow{Acarreo y diferimiento}{deuda total: 2; final: 0; bloques diferidos: 1}
\RouteRow{Tríadas finales; persistencia de Pareto}{17; sí}
\RouteRow{\(K\longrightarrow K+9\)}{repeticiones: 9; cambios de memoria: 9}
\RouteGroup{Firma, fase y diferencias}
\RouteRow{\(n_A,\ s_C,\ \kappa_0,\ \nu_{120},\ \nu_{270}\)}{14, 4, 0, 0, -12}
\RouteRow{\(\tau_{12}\); firma histórica \(\mathbb Z^5\)}{+++-\kern0pt{}-\kern0pt{}-+++-\kern0pt{}-\kern0pt{}-; 14|\allowbreak{}4|\allowbreak{}0|\allowbreak{}0|\allowbreak{}-12}
\RouteRow{\(D_1,\ D_3,\ D_4,\ D_{27},\ D_{36}\)}{78, 28, 78, 36, 24}
\RouteRow{\(\Omega_{90/120},\ P_6\)}{80, 6}
\RouteGroup{Lectores y factores de acción}
\RouteRow{\(K_D\quad /\quad K_\Omega\)}{(60,\allowbreak{}78,\allowbreak{}25)\quad /\quad (80,\allowbreak{}54,\allowbreak{}6)}
\RouteRow{\(\kappa_D\)}{59.43358\allowbreak{}64101631\allowbreak{}67023563\allowbreak{}04}
\RouteRow{\(\kappa_\Omega\)}{79.60661\allowbreak{}01397948\allowbreak{}27586470\allowbreak{}91}
\RouteRow{\(R_{\mathrm{act},D}\)}{1.000000\allowbreak{}04120422\allowbreak{}24053449\allowbreak{}05}
\RouteRow{\(R_{\mathrm{act},\Omega}\)}{1.000000\allowbreak{}05518981\allowbreak{}25318591\allowbreak{}40}
\RouteGroup{Separación de prefijos}
\RouteRow{Área de ambigüedad; área \(\log_2\)}{18; 0.0}
\RouteRow{Primer bloque de separación}{órbita: 1; celda: 1; ruta: 1}
\RouteRow{Ambigüedad final}{1}
\RouteGroup{Secuencias completas}
\RouteRow{\(w_6\)}{666 672 510 666 666 585 507 504 585 666 672 510 666 666 585 507 504 585}
\RouteRow{Tríadas estables por bloque}{0 1 2 3 4 5 6 7 8 9 10 11 12 12 13 15 16 17}
\end{tabularx}\endgroup\par\vspace{5pt}\noindent{\fontsize{8.4}{10}\selectfont\hyperref[trace-135]{Procedencia y códigos originales}\hfill Ficha completa · 135/324}
\clearpage\phantomsection\label{nav:vi-ruta-136}
% VI-CATALOG-ROW rutas 136
\pdfbookmark[3]{r4c7\_h+1v+1}{route-136}
\RouteTitle{Ruta 136}{r4c7\_\allowbreak{}h+1v+1}
\noindent Celda 34\quad (4, 7)\qquad Orientación \(h=1,\ v=1\)\hfill\hyperref[sec:vi-indice-rutas]{Índice}\par\vspace{4pt}
\begingroup\fontsize{10.2}{12.5}\selectfont\setlength{\tabcolsep}{5pt}\renewcommand{\arraystretch}{1.1}\rowcolors{1}{routepale}{white}\noindent\begin{tabularx}{\linewidth}{@{}>{\raggedright\arraybackslash}p{.345\linewidth}>{\raggedright\arraybackslash}X@{}}
\RouteGroup{Identidad estructural}
\RouteRow{Clase y multisección}{quark\_\allowbreak{}down\_\allowbreak{}G2\quad /\quad quark\_\allowbreak{}down\_\allowbreak{}G2}
\RouteRow{Colección y sector}{quark de tipo abajo; quark}
\RouteRow{Clasificación física}{quark de tipo abajo}
\RouteRow{Generación, profundidad y color}{2; G2; R}
\RouteRow{Ruta primaria y órbita de inversión}{\(B_1\): no\qquad 4,\allowbreak{}7;\allowbreak{}6,\allowbreak{}3}
\RouteRow{Firma de clasificación}{30|\allowbreak{}-6|\allowbreak{}0|\allowbreak{}-2|\allowbreak{}-8|\allowbreak{}++0-\kern0pt{}-\kern0pt{}-++0-\kern0pt{}-\kern0pt{}-}
\RouteGroup{Retornos, memoria y transporte}
\RouteRow{Retornos con fase}{clase: 27; órbita: 27; celda: 27}
\RouteRow{Retorno cinemático}{en 108: 54; extensión: 216}
\RouteRow{Giros y cambios de hoja}{71; 107}
\RouteRow{Acarreo y diferimiento}{deuda total: 2; final: 0; bloques diferidos: 1}
\RouteRow{Tríadas finales; persistencia de Pareto}{17; sí}
\RouteRow{\(K\longrightarrow K+9\)}{repeticiones: 9; cambios de memoria: 9}
\RouteGroup{Firma, fase y diferencias}
\RouteRow{\(n_A,\ s_C,\ \kappa_0,\ \nu_{120},\ \nu_{270}\)}{30, -6, 2, 0, -8}
\RouteRow{\(\tau_{12}\); firma histórica \(\mathbb Z^5\)}{++0-\kern0pt{}-\kern0pt{}-++0-\kern0pt{}-\kern0pt{}-; 30|\allowbreak{}-6|\allowbreak{}2|\allowbreak{}0|\allowbreak{}-8}
\RouteRow{\(D_1,\ D_3,\ D_4,\ D_{27},\ D_{36}\)}{54, 56, 70, 48, 32}
\RouteRow{\(\Omega_{90/120},\ P_6\)}{56, 5}
\RouteGroup{Lectores y factores de acción}
\RouteRow{\(K_D\quad /\quad K_\Omega\)}{(80,\allowbreak{}54,\allowbreak{}7)\quad /\quad (56,\allowbreak{}54,\allowbreak{}5)}
\RouteRow{\(\kappa_D\)}{79.60672\allowbreak{}13362175\allowbreak{}19726524\allowbreak{}96}
\RouteRow{\(\kappa_\Omega\)}{55.60649\allowbreak{}89433721\allowbreak{}35446416\allowbreak{}86}
\RouteRow{\(R_{\mathrm{act},D}\)}{1.000000\allowbreak{}05518988\allowbreak{}96223153\allowbreak{}37}
\RouteRow{\(R_{\mathrm{act},\Omega}\)}{1.000000\allowbreak{}03855097\allowbreak{}23541416\allowbreak{}45}
\RouteGroup{Separación de prefijos}
\RouteRow{Área de ambigüedad; área \(\log_2\)}{18; 0.0}
\RouteRow{Primer bloque de separación}{órbita: 1; celda: 1; ruta: 1}
\RouteRow{Ambigüedad final}{1}
\RouteGroup{Secuencias completas}
\RouteRow{\(w_6\)}{657 24 243 657 18 15 486 423 15 657 24 243 657 18 15 486 423 15}
\RouteRow{Tríadas estables por bloque}{0 1 2 3 4 5 6 7 8 9 10 11 12 13 14 14 15 17}
\end{tabularx}\endgroup\par\vspace{5pt}\noindent{\fontsize{8.4}{10}\selectfont\hyperref[trace-136]{Procedencia y códigos originales}\hfill Ficha completa · 136/324}
\clearpage\phantomsection\label{nav:vi-ruta-137}
% VI-CATALOG-ROW rutas 137
\pdfbookmark[2]{Celda 35}{cell-35}
\pdfbookmark[3]{r4c8\_h-1v-1}{route-137}
\RouteTitle{Ruta 137}{r4c8\_\allowbreak{}h-1v-1}
\noindent Celda 35\quad (4, 8)\qquad Orientación \(h=-1,\ v=-1\)\hfill\hyperref[sec:vi-indice-rutas]{Índice}\par\vspace{4pt}
\begingroup\fontsize{10.2}{12.5}\selectfont\setlength{\tabcolsep}{5pt}\renewcommand{\arraystretch}{1.1}\rowcolors{1}{routepale}{white}\noindent\begin{tabularx}{\linewidth}{@{}>{\raggedright\arraybackslash}p{.345\linewidth}>{\raggedright\arraybackslash}X@{}}
\RouteGroup{Identidad estructural}
\RouteRow{Clase y multisección}{quark\_\allowbreak{}up\_\allowbreak{}G3\quad /\quad quark\_\allowbreak{}up\_\allowbreak{}G3}
\RouteRow{Colección y sector}{quark de tipo arriba; quark}
\RouteRow{Clasificación física}{quark de tipo arriba}
\RouteRow{Generación, profundidad y color}{3; G3; B}
\RouteRow{Ruta primaria y órbita de inversión}{\(B_1\): sí\qquad 4,\allowbreak{}8;\allowbreak{}6,\allowbreak{}2}
\RouteRow{Firma de clasificación}{34|\allowbreak{}0|\allowbreak{}-20|\allowbreak{}-2|\allowbreak{}-8|\allowbreak{}0++-\kern0pt{}-\kern0pt{}-0++-\kern0pt{}-\kern0pt{}-}
\RouteGroup{Retornos, memoria y transporte}
\RouteRow{Retornos con fase}{clase: 27; órbita: 27; celda: 27}
\RouteRow{Retorno cinemático}{en 108: 54; extensión: 216}
\RouteRow{Giros y cambios de hoja}{71; 107}
\RouteRow{Acarreo y diferimiento}{deuda total: 1; final: 0; bloques diferidos: 1}
\RouteRow{Tríadas finales; persistencia de Pareto}{17; sí}
\RouteRow{\(K\longrightarrow K+9\)}{repeticiones: 9; cambios de memoria: 9}
\RouteGroup{Firma, fase y diferencias}
\RouteRow{\(n_A,\ s_C,\ \kappa_0,\ \nu_{120},\ \nu_{270}\)}{34, 0, -2, 0, -8}
\RouteRow{\(\tau_{12}\); firma histórica \(\mathbb Z^5\)}{+++0-\kern0pt{}-+++0-\kern0pt{}-; 34|\allowbreak{}0|\allowbreak{}-2|\allowbreak{}0|\allowbreak{}-8}
\RouteRow{\(D_1,\ D_3,\ D_4,\ D_{27},\ D_{36}\)}{66, 64, 64, 60, 40}
\RouteRow{\(\Omega_{90/120},\ P_6\)}{56, 5}
\RouteGroup{Lectores y factores de acción}
\RouteRow{\(K_D\quad /\quad K_\Omega\)}{(100,\allowbreak{}66,\allowbreak{}0)\quad /\quad (56,\allowbreak{}54,\allowbreak{}5)}
\RouteRow{\(\kappa_D\)}{99.51837\allowbreak{}47304272\allowbreak{}69134179\allowbreak{}18}
\RouteRow{\(\kappa_\Omega\)}{55.60649\allowbreak{}89433721\allowbreak{}35446416\allowbreak{}86}
\RouteRow{\(R_{\mathrm{act},D}\)}{1.000000\allowbreak{}06899427\allowbreak{}66450218\allowbreak{}95}
\RouteRow{\(R_{\mathrm{act},\Omega}\)}{1.000000\allowbreak{}03855097\allowbreak{}23541416\allowbreak{}45}
\RouteGroup{Separación de prefijos}
\RouteRow{Área de ambigüedad; área \(\log_2\)}{36; 18.0}
\RouteRow{Primer bloque de separación}{órbita: \textemdash{}; celda: \textemdash{}; ruta: \textemdash{}}
\RouteRow{Ambigüedad final}{2}
\RouteGroup{Secuencias completas}
\RouteRow{\(w_6\)}{90 570 264 90 567 657 24 243 657 90 570 264 90 567 657 24 243 657}
\RouteRow{Tríadas estables por bloque}{0 1 2 3 4 5 6 7 8 9 10 11 12 13 13 15 16 17}
\end{tabularx}\endgroup\par\vspace{5pt}\noindent{\fontsize{8.4}{10}\selectfont\hyperref[trace-137]{Procedencia y códigos originales}\hfill Ficha completa · 137/324}
\clearpage\phantomsection\label{nav:vi-ruta-138}
% VI-CATALOG-ROW rutas 138
\pdfbookmark[3]{r4c8\_h-1v+1}{route-138}
\RouteTitle{Ruta 138}{r4c8\_\allowbreak{}h-1v+1}
\noindent Celda 35\quad (4, 8)\qquad Orientación \(h=-1,\ v=1\)\hfill\hyperref[sec:vi-indice-rutas]{Índice}\par\vspace{4pt}
\begingroup\fontsize{10.2}{12.5}\selectfont\setlength{\tabcolsep}{5pt}\renewcommand{\arraystretch}{1.1}\rowcolors{1}{routepale}{white}\noindent\begin{tabularx}{\linewidth}{@{}>{\raggedright\arraybackslash}p{.345\linewidth}>{\raggedright\arraybackslash}X@{}}
\RouteGroup{Identidad estructural}
\RouteRow{Clase y multisección}{quark\_\allowbreak{}up\_\allowbreak{}G3\quad /\quad quark\_\allowbreak{}up\_\allowbreak{}G3}
\RouteRow{Colección y sector}{quark de tipo arriba; quark}
\RouteRow{Clasificación física}{quark de tipo arriba}
\RouteRow{Generación, profundidad y color}{3; G3; B}
\RouteRow{Ruta primaria y órbita de inversión}{\(B_1\): no\qquad 4,\allowbreak{}8;\allowbreak{}6,\allowbreak{}2}
\RouteRow{Firma de clasificación}{34|\allowbreak{}0|\allowbreak{}-20|\allowbreak{}-2|\allowbreak{}-8|\allowbreak{}0++-\kern0pt{}-\kern0pt{}-0++-\kern0pt{}-\kern0pt{}-}
\RouteGroup{Retornos, memoria y transporte}
\RouteRow{Retornos con fase}{clase: 27; órbita: 27; celda: 27}
\RouteRow{Retorno cinemático}{en 108: 54; extensión: 216}
\RouteRow{Giros y cambios de hoja}{71; 107}
\RouteRow{Acarreo y diferimiento}{deuda total: 2; final: 0; bloques diferidos: 2}
\RouteRow{Tríadas finales; persistencia de Pareto}{17; sí}
\RouteRow{\(K\longrightarrow K+9\)}{repeticiones: 9; cambios de memoria: 9}
\RouteGroup{Firma, fase y diferencias}
\RouteRow{\(n_A,\ s_C,\ \kappa_0,\ \nu_{120},\ \nu_{270}\)}{50, -6, 2, -2, -8}
\RouteRow{\(\tau_{12}\); firma histórica \(\mathbb Z^5\)}{0++-\kern0pt{}-\kern0pt{}-0++-\kern0pt{}-\kern0pt{}-; 50|\allowbreak{}-6|\allowbreak{}2|\allowbreak{}-2|\allowbreak{}-8}
\RouteRow{\(D_1,\ D_3,\ D_4,\ D_{27},\ D_{36}\)}{24, 20, 24, 24, 16}
\RouteRow{\(\Omega_{90/120},\ P_6\)}{56, 5}
\RouteGroup{Lectores y factores de acción}
\RouteRow{\(K_D\quad /\quad K_\Omega\)}{(40,\allowbreak{}24,\allowbreak{}2)\quad /\quad (56,\allowbreak{}54,\allowbreak{}5)}
\RouteRow{\(\kappa_D\)}{39.82508\allowbreak{}59311825\allowbreak{}73056173\allowbreak{}26}
\RouteRow{\(\kappa_\Omega\)}{55.60649\allowbreak{}89433721\allowbreak{}35446416\allowbreak{}86}
\RouteRow{\(R_{\mathrm{act},D}\)}{1.000000\allowbreak{}02761000\allowbreak{}61595142\allowbreak{}15}
\RouteRow{\(R_{\mathrm{act},\Omega}\)}{1.000000\allowbreak{}03855097\allowbreak{}23541416\allowbreak{}45}
\RouteGroup{Separación de prefijos}
\RouteRow{Área de ambigüedad; área \(\log_2\)}{72; 36.0}
\RouteRow{Primer bloque de separación}{órbita: \textemdash{}; celda: \textemdash{}; ruta: \textemdash{}}
\RouteRow{Ambigüedad final}{4}
\RouteGroup{Secuencias completas}
\RouteRow{\(w_6\)}{81 3 0 81 0 3 0 81 3 81 3 0 81 0 3 0 81 3}
\RouteRow{Tríadas estables por bloque}{0 1 2 3 4 5 6 6 8 9 10 11 12 13 14 15 15 17}
\end{tabularx}\endgroup\par\vspace{5pt}\noindent{\fontsize{8.4}{10}\selectfont\hyperref[trace-138]{Procedencia y códigos originales}\hfill Ficha completa · 138/324}
\clearpage\phantomsection\label{nav:vi-ruta-139}
% VI-CATALOG-ROW rutas 139
\pdfbookmark[3]{r4c8\_h+1v-1}{route-139}
\RouteTitle{Ruta 139}{r4c8\_\allowbreak{}h+1v-1}
\noindent Celda 35\quad (4, 8)\qquad Orientación \(h=1,\ v=-1\)\hfill\hyperref[sec:vi-indice-rutas]{Índice}\par\vspace{4pt}
\begingroup\fontsize{10.2}{12.5}\selectfont\setlength{\tabcolsep}{5pt}\renewcommand{\arraystretch}{1.1}\rowcolors{1}{routepale}{white}\noindent\begin{tabularx}{\linewidth}{@{}>{\raggedright\arraybackslash}p{.345\linewidth}>{\raggedright\arraybackslash}X@{}}
\RouteGroup{Identidad estructural}
\RouteRow{Clase y multisección}{quark\_\allowbreak{}up\_\allowbreak{}G3\quad /\quad quark\_\allowbreak{}up\_\allowbreak{}G3}
\RouteRow{Colección y sector}{quark de tipo arriba; quark}
\RouteRow{Clasificación física}{quark de tipo arriba}
\RouteRow{Generación, profundidad y color}{3; G3; B}
\RouteRow{Ruta primaria y órbita de inversión}{\(B_1\): no\qquad 4,\allowbreak{}8;\allowbreak{}6,\allowbreak{}2}
\RouteRow{Firma de clasificación}{34|\allowbreak{}0|\allowbreak{}-20|\allowbreak{}-2|\allowbreak{}-8|\allowbreak{}0++-\kern0pt{}-\kern0pt{}-0++-\kern0pt{}-\kern0pt{}-}
\RouteGroup{Retornos, memoria y transporte}
\RouteRow{Retornos con fase}{clase: 27; órbita: 27; celda: 27}
\RouteRow{Retorno cinemático}{en 108: 54; extensión: 216}
\RouteRow{Giros y cambios de hoja}{71; 107}
\RouteRow{Acarreo y diferimiento}{deuda total: 0; final: 0; bloques diferidos: 0}
\RouteRow{Tríadas finales; persistencia de Pareto}{17; sí}
\RouteRow{\(K\longrightarrow K+9\)}{repeticiones: 9; cambios de memoria: 9}
\RouteGroup{Firma, fase y diferencias}
\RouteRow{\(n_A,\ s_C,\ \kappa_0,\ \nu_{120},\ \nu_{270}\)}{50, 6, -2, 2, -8}
\RouteRow{\(\tau_{12}\); firma histórica \(\mathbb Z^5\)}{+++0-\kern0pt{}-+++0-\kern0pt{}-; 50|\allowbreak{}6|\allowbreak{}-2|\allowbreak{}2|\allowbreak{}-8}
\RouteRow{\(D_1,\ D_3,\ D_4,\ D_{27},\ D_{36}\)}{24, 20, 24, 24, 16}
\RouteRow{\(\Omega_{90/120},\ P_6\)}{56, 5}
\RouteGroup{Lectores y factores de acción}
\RouteRow{\(K_D\quad /\quad K_\Omega\)}{(40,\allowbreak{}24,\allowbreak{}2)\quad /\quad (56,\allowbreak{}54,\allowbreak{}5)}
\RouteRow{\(\kappa_D\)}{39.82508\allowbreak{}59311825\allowbreak{}73056173\allowbreak{}26}
\RouteRow{\(\kappa_\Omega\)}{55.60649\allowbreak{}89433721\allowbreak{}35446416\allowbreak{}86}
\RouteRow{\(R_{\mathrm{act},D}\)}{1.000000\allowbreak{}02761000\allowbreak{}61595142\allowbreak{}15}
\RouteRow{\(R_{\mathrm{act},\Omega}\)}{1.000000\allowbreak{}03855097\allowbreak{}23541416\allowbreak{}45}
\RouteGroup{Separación de prefijos}
\RouteRow{Área de ambigüedad; área \(\log_2\)}{72; 36.0}
\RouteRow{Primer bloque de separación}{órbita: \textemdash{}; celda: \textemdash{}; ruta: \textemdash{}}
\RouteRow{Ambigüedad final}{4}
\RouteGroup{Secuencias completas}
\RouteRow{\(w_6\)}{0 81 3 0 84 0 81 3 0 0 81 3 0 84 0 81 3 0}
\RouteRow{Tríadas estables por bloque}{0 1 2 3 4 5 6 7 8 9 10 11 12 13 14 15 16 17}
\end{tabularx}\endgroup\par\vspace{5pt}\noindent{\fontsize{8.4}{10}\selectfont\hyperref[trace-139]{Procedencia y códigos originales}\hfill Ficha completa · 139/324}
\clearpage\phantomsection\label{nav:vi-ruta-140}
% VI-CATALOG-ROW rutas 140
\pdfbookmark[3]{r4c8\_h+1v+1}{route-140}
\RouteTitle{Ruta 140}{r4c8\_\allowbreak{}h+1v+1}
\noindent Celda 35\quad (4, 8)\qquad Orientación \(h=1,\ v=1\)\hfill\hyperref[sec:vi-indice-rutas]{Índice}\par\vspace{4pt}
\begingroup\fontsize{10.2}{12.5}\selectfont\setlength{\tabcolsep}{5pt}\renewcommand{\arraystretch}{1.1}\rowcolors{1}{routepale}{white}\noindent\begin{tabularx}{\linewidth}{@{}>{\raggedright\arraybackslash}p{.345\linewidth}>{\raggedright\arraybackslash}X@{}}
\RouteGroup{Identidad estructural}
\RouteRow{Clase y multisección}{quark\_\allowbreak{}up\_\allowbreak{}G3\quad /\quad quark\_\allowbreak{}up\_\allowbreak{}G3}
\RouteRow{Colección y sector}{quark de tipo arriba; quark}
\RouteRow{Clasificación física}{quark de tipo arriba}
\RouteRow{Generación, profundidad y color}{3; G3; B}
\RouteRow{Ruta primaria y órbita de inversión}{\(B_1\): no\qquad 4,\allowbreak{}8;\allowbreak{}6,\allowbreak{}2}
\RouteRow{Firma de clasificación}{34|\allowbreak{}0|\allowbreak{}-20|\allowbreak{}-2|\allowbreak{}-8|\allowbreak{}0++-\kern0pt{}-\kern0pt{}-0++-\kern0pt{}-\kern0pt{}-}
\RouteGroup{Retornos, memoria y transporte}
\RouteRow{Retornos con fase}{clase: 27; órbita: 27; celda: 27}
\RouteRow{Retorno cinemático}{en 108: 54; extensión: 216}
\RouteRow{Giros y cambios de hoja}{71; 107}
\RouteRow{Acarreo y diferimiento}{deuda total: 0; final: 0; bloques diferidos: 0}
\RouteRow{Tríadas finales; persistencia de Pareto}{17; sí}
\RouteRow{\(K\longrightarrow K+9\)}{repeticiones: 9; cambios de memoria: 9}
\RouteGroup{Firma, fase y diferencias}
\RouteRow{\(n_A,\ s_C,\ \kappa_0,\ \nu_{120},\ \nu_{270}\)}{34, 0, 2, 0, -8}
\RouteRow{\(\tau_{12}\); firma histórica \(\mathbb Z^5\)}{0++-\kern0pt{}-\kern0pt{}-0++-\kern0pt{}-\kern0pt{}-; 34|\allowbreak{}0|\allowbreak{}2|\allowbreak{}0|\allowbreak{}-8}
\RouteRow{\(D_1,\ D_3,\ D_4,\ D_{27},\ D_{36}\)}{66, 64, 64, 60, 40}
\RouteRow{\(\Omega_{90/120},\ P_6\)}{56, 5}
\RouteGroup{Lectores y factores de acción}
\RouteRow{\(K_D\quad /\quad K_\Omega\)}{(100,\allowbreak{}66,\allowbreak{}0)\quad /\quad (56,\allowbreak{}54,\allowbreak{}5)}
\RouteRow{\(\kappa_D\)}{99.51837\allowbreak{}47304272\allowbreak{}69134179\allowbreak{}18}
\RouteRow{\(\kappa_\Omega\)}{55.60649\allowbreak{}89433721\allowbreak{}35446416\allowbreak{}86}
\RouteRow{\(R_{\mathrm{act},D}\)}{1.000000\allowbreak{}06899427\allowbreak{}66450218\allowbreak{}95}
\RouteRow{\(R_{\mathrm{act},\Omega}\)}{1.000000\allowbreak{}03855097\allowbreak{}23541416\allowbreak{}45}
\RouteGroup{Separación de prefijos}
\RouteRow{Área de ambigüedad; área \(\log_2\)}{36; 18.0}
\RouteRow{Primer bloque de separación}{órbita: \textemdash{}; celda: \textemdash{}; ruta: \textemdash{}}
\RouteRow{Ambigüedad final}{2}
\RouteGroup{Secuencias completas}
\RouteRow{\(w_6\)}{24 243 657 24 246 264 90 570 264 24 243 657 24 246 264 90 570 264}
\RouteRow{Tríadas estables por bloque}{0 1 2 3 4 5 6 7 8 9 10 11 12 13 14 15 16 17}
\end{tabularx}\endgroup\par\vspace{5pt}\noindent{\fontsize{8.4}{10}\selectfont\hyperref[trace-140]{Procedencia y códigos originales}\hfill Ficha completa · 140/324}
\clearpage\phantomsection\label{nav:vi-ruta-141}
% VI-CATALOG-ROW rutas 141
\pdfbookmark[2]{Celda 36}{cell-36}
\pdfbookmark[3]{r4c9\_h-1v-1}{route-141}
\RouteTitle{Ruta 141}{r4c9\_\allowbreak{}h-1v-1}
\noindent Celda 36\quad (4, 9)\qquad Orientación \(h=-1,\ v=-1\)\hfill\hyperref[sec:vi-indice-rutas]{Índice}\par\vspace{4pt}
\begingroup\fontsize{10.2}{12.5}\selectfont\setlength{\tabcolsep}{5pt}\renewcommand{\arraystretch}{1.1}\rowcolors{1}{routepale}{white}\noindent\begin{tabularx}{\linewidth}{@{}>{\raggedright\arraybackslash}p{.345\linewidth}>{\raggedright\arraybackslash}X@{}}
\RouteGroup{Identidad estructural}
\RouteRow{Clase y multisección}{quark\_\allowbreak{}down\_\allowbreak{}G4\quad /\quad quark\_\allowbreak{}down\_\allowbreak{}G4}
\RouteRow{Colección y sector}{quark de tipo abajo; quark}
\RouteRow{Clasificación física}{quark de tipo abajo}
\RouteRow{Generación, profundidad y color}{4; G4; G}
\RouteRow{Ruta primaria y órbita de inversión}{\(B_1\): no\qquad 4,\allowbreak{}9;\allowbreak{}6,\allowbreak{}1}
\RouteRow{Firma de clasificación}{34|\allowbreak{}0|\allowbreak{}-8|\allowbreak{}0|\allowbreak{}-12|\allowbreak{}+++-\kern0pt{}-\kern0pt{}-+++-\kern0pt{}-\kern0pt{}-}
\RouteGroup{Retornos, memoria y transporte}
\RouteRow{Retornos con fase}{clase: 27; órbita: 27; celda: 27}
\RouteRow{Retorno cinemático}{en 108: 54; extensión: 216}
\RouteRow{Giros y cambios de hoja}{71; 107}
\RouteRow{Acarreo y diferimiento}{deuda total: 1; final: 1; bloques diferidos: 1}
\RouteRow{Tríadas finales; persistencia de Pareto}{16; sí}
\RouteRow{\(K\longrightarrow K+9\)}{repeticiones: 9; cambios de memoria: 9}
\RouteGroup{Firma, fase y diferencias}
\RouteRow{\(n_A,\ s_C,\ \kappa_0,\ \nu_{120},\ \nu_{270}\)}{34, 0, 0, 0, -12}
\RouteRow{\(\tau_{12}\); firma histórica \(\mathbb Z^5\)}{+++-\kern0pt{}-\kern0pt{}-+++-\kern0pt{}-\kern0pt{}-; 34|\allowbreak{}0|\allowbreak{}0|\allowbreak{}0|\allowbreak{}-12}
\RouteRow{\(D_1,\ D_3,\ D_4,\ D_{27},\ D_{36}\)}{66, 60, 64, 48, 32}
\RouteRow{\(\Omega_{90/120},\ P_6\)}{80, 6}
\RouteGroup{Lectores y factores de acción}
\RouteRow{\(K_D\quad /\quad K_\Omega\)}{(80,\allowbreak{}66,\allowbreak{}2)\quad /\quad (80,\allowbreak{}54,\allowbreak{}6)}
\RouteRow{\(\kappa_D\)}{79.51859\allowbreak{}71232726\allowbreak{}53414287\allowbreak{}29}
\RouteRow{\(\kappa_\Omega\)}{79.60661\allowbreak{}01397948\allowbreak{}27586470\allowbreak{}91}
\RouteRow{\(R_{\mathrm{act},D}\)}{1.000000\allowbreak{}05512879\allowbreak{}47092488\allowbreak{}80}
\RouteRow{\(R_{\mathrm{act},\Omega}\)}{1.000000\allowbreak{}05518981\allowbreak{}25318591\allowbreak{}40}
\RouteGroup{Separación de prefijos}
\RouteRow{Área de ambigüedad; área \(\log_2\)}{18; 0.0}
\RouteRow{Primer bloque de separación}{órbita: 1; celda: 1; ruta: 1}
\RouteRow{Ambigüedad final}{1}
\RouteGroup{Secuencias completas}
\RouteRow{\(w_6\)}{423 15 486 423 12 99 327 498 99 423 15 486 423 12 99 327 498 99}
\RouteRow{Tríadas estables por bloque}{0 1 2 3 4 5 6 7 8 9 10 11 12 13 14 15 16 16}
\end{tabularx}\endgroup\par\vspace{5pt}\noindent{\fontsize{8.4}{10}\selectfont\hyperref[trace-141]{Procedencia y códigos originales}\hfill Ficha completa · 141/324}
\clearpage\phantomsection\label{nav:vi-ruta-142}
% VI-CATALOG-ROW rutas 142
\pdfbookmark[3]{r4c9\_h-1v+1}{route-142}
\RouteTitle{Ruta 142}{r4c9\_\allowbreak{}h-1v+1}
\noindent Celda 36\quad (4, 9)\qquad Orientación \(h=-1,\ v=1\)\hfill\hyperref[sec:vi-indice-rutas]{Índice}\par\vspace{4pt}
\begingroup\fontsize{10.2}{12.5}\selectfont\setlength{\tabcolsep}{5pt}\renewcommand{\arraystretch}{1.1}\rowcolors{1}{routepale}{white}\noindent\begin{tabularx}{\linewidth}{@{}>{\raggedright\arraybackslash}p{.345\linewidth}>{\raggedright\arraybackslash}X@{}}
\RouteGroup{Identidad estructural}
\RouteRow{Clase y multisección}{quark\_\allowbreak{}down\_\allowbreak{}G4\quad /\quad quark\_\allowbreak{}down\_\allowbreak{}G4}
\RouteRow{Colección y sector}{quark de tipo abajo; quark}
\RouteRow{Clasificación física}{quark de tipo abajo}
\RouteRow{Generación, profundidad y color}{4; G4; G}
\RouteRow{Ruta primaria y órbita de inversión}{\(B_1\): sí\qquad 4,\allowbreak{}9;\allowbreak{}6,\allowbreak{}1}
\RouteRow{Firma de clasificación}{34|\allowbreak{}0|\allowbreak{}-8|\allowbreak{}0|\allowbreak{}-12|\allowbreak{}+++-\kern0pt{}-\kern0pt{}-+++-\kern0pt{}-\kern0pt{}-}
\RouteGroup{Retornos, memoria y transporte}
\RouteRow{Retornos con fase}{clase: 27; órbita: 27; celda: 27}
\RouteRow{Retorno cinemático}{en 108: 54; extensión: 216}
\RouteRow{Giros y cambios de hoja}{71; 107}
\RouteRow{Acarreo y diferimiento}{deuda total: 1; final: 0; bloques diferidos: 1}
\RouteRow{Tríadas finales; persistencia de Pareto}{17; sí}
\RouteRow{\(K\longrightarrow K+9\)}{repeticiones: 9; cambios de memoria: 9}
\RouteGroup{Firma, fase y diferencias}
\RouteRow{\(n_A,\ s_C,\ \kappa_0,\ \nu_{120},\ \nu_{270}\)}{8, -2, 0, 0, -12}
\RouteRow{\(\tau_{12}\); firma histórica \(\mathbb Z^5\)}{+++-\kern0pt{}-\kern0pt{}-+++-\kern0pt{}-\kern0pt{}-; 8|\allowbreak{}-2|\allowbreak{}0|\allowbreak{}0|\allowbreak{}-12}
\RouteRow{\(D_1,\ D_3,\ D_4,\ D_{27},\ D_{36}\)}{84, 24, 84, 24, 16}
\RouteRow{\(\Omega_{90/120},\ P_6\)}{80, 6}
\RouteGroup{Lectores y factores de acción}
\RouteRow{\(K_D\quad /\quad K_\Omega\)}{(40,\allowbreak{}84,\allowbreak{}30)\quad /\quad (80,\allowbreak{}54,\allowbreak{}6)}
\RouteRow{\(\kappa_D\)}{39.39035\allowbreak{}82768609\allowbreak{}24917849\allowbreak{}59}
\RouteRow{\(\kappa_\Omega\)}{79.60661\allowbreak{}01397948\allowbreak{}27586470\allowbreak{}91}
\RouteRow{\(R_{\mathrm{act},D}\)}{1.000000\allowbreak{}02730861\allowbreak{}73967087\allowbreak{}05}
\RouteRow{\(R_{\mathrm{act},\Omega}\)}{1.000000\allowbreak{}05518981\allowbreak{}25318591\allowbreak{}40}
\RouteGroup{Separación de prefijos}
\RouteRow{Área de ambigüedad; área \(\log_2\)}{18; 0.0}
\RouteRow{Primer bloque de separación}{órbita: 1; celda: 1; ruta: 1}
\RouteRow{Ambigüedad final}{1}
\RouteGroup{Secuencias completas}
\RouteRow{\(w_6\)}{420 258 414 420 255 252 333 255 252 420 258 414 420 255 252 333 255 252}
\RouteRow{Tríadas estables por bloque}{0 1 2 3 4 5 6 7 8 9 10 11 12 13 14 15 15 17}
\end{tabularx}\endgroup\par\vspace{5pt}\noindent{\fontsize{8.4}{10}\selectfont\hyperref[trace-142]{Procedencia y códigos originales}\hfill Ficha completa · 142/324}
\clearpage\phantomsection\label{nav:vi-ruta-143}
% VI-CATALOG-ROW rutas 143
\pdfbookmark[3]{r4c9\_h+1v-1}{route-143}
\RouteTitle{Ruta 143}{r4c9\_\allowbreak{}h+1v-1}
\noindent Celda 36\quad (4, 9)\qquad Orientación \(h=1,\ v=-1\)\hfill\hyperref[sec:vi-indice-rutas]{Índice}\par\vspace{4pt}
\begingroup\fontsize{10.2}{12.5}\selectfont\setlength{\tabcolsep}{5pt}\renewcommand{\arraystretch}{1.1}\rowcolors{1}{routepale}{white}\noindent\begin{tabularx}{\linewidth}{@{}>{\raggedright\arraybackslash}p{.345\linewidth}>{\raggedright\arraybackslash}X@{}}
\RouteGroup{Identidad estructural}
\RouteRow{Clase y multisección}{quark\_\allowbreak{}down\_\allowbreak{}G4\quad /\quad quark\_\allowbreak{}down\_\allowbreak{}G4}
\RouteRow{Colección y sector}{quark de tipo abajo; quark}
\RouteRow{Clasificación física}{quark de tipo abajo}
\RouteRow{Generación, profundidad y color}{4; G4; G}
\RouteRow{Ruta primaria y órbita de inversión}{\(B_1\): no\qquad 4,\allowbreak{}9;\allowbreak{}6,\allowbreak{}1}
\RouteRow{Firma de clasificación}{34|\allowbreak{}0|\allowbreak{}-8|\allowbreak{}0|\allowbreak{}-12|\allowbreak{}+++-\kern0pt{}-\kern0pt{}-+++-\kern0pt{}-\kern0pt{}-}
\RouteGroup{Retornos, memoria y transporte}
\RouteRow{Retornos con fase}{clase: 27; órbita: 27; celda: 27}
\RouteRow{Retorno cinemático}{en 108: 54; extensión: 216}
\RouteRow{Giros y cambios de hoja}{71; 107}
\RouteRow{Acarreo y diferimiento}{deuda total: 0; final: 0; bloques diferidos: 0}
\RouteRow{Tríadas finales; persistencia de Pareto}{17; sí}
\RouteRow{\(K\longrightarrow K+9\)}{repeticiones: 9; cambios de memoria: 9}
\RouteGroup{Firma, fase y diferencias}
\RouteRow{\(n_A,\ s_C,\ \kappa_0,\ \nu_{120},\ \nu_{270}\)}{8, 2, 0, 0, -12}
\RouteRow{\(\tau_{12}\); firma histórica \(\mathbb Z^5\)}{+++-\kern0pt{}-\kern0pt{}-+++-\kern0pt{}-\kern0pt{}-; 8|\allowbreak{}2|\allowbreak{}0|\allowbreak{}0|\allowbreak{}-12}
\RouteRow{\(D_1,\ D_3,\ D_4,\ D_{27},\ D_{36}\)}{84, 24, 84, 24, 16}
\RouteRow{\(\Omega_{90/120},\ P_6\)}{80, 6}
\RouteGroup{Lectores y factores de acción}
\RouteRow{\(K_D\quad /\quad K_\Omega\)}{(40,\allowbreak{}84,\allowbreak{}30)\quad /\quad (80,\allowbreak{}54,\allowbreak{}6)}
\RouteRow{\(\kappa_D\)}{39.39035\allowbreak{}82768609\allowbreak{}24917849\allowbreak{}59}
\RouteRow{\(\kappa_\Omega\)}{79.60661\allowbreak{}01397948\allowbreak{}27586470\allowbreak{}91}
\RouteRow{\(R_{\mathrm{act},D}\)}{1.000000\allowbreak{}02730861\allowbreak{}73967087\allowbreak{}05}
\RouteRow{\(R_{\mathrm{act},\Omega}\)}{1.000000\allowbreak{}05518981\allowbreak{}25318591\allowbreak{}40}
\RouteGroup{Separación de prefijos}
\RouteRow{Área de ambigüedad; área \(\log_2\)}{18; 0.0}
\RouteRow{Primer bloque de separación}{órbita: 1; celda: 1; ruta: 1}
\RouteRow{Ambigüedad final}{1}
\RouteGroup{Secuencias completas}
\RouteRow{\(w_6\)}{333 255 252 333 258 414 420 258 414 333 255 252 333 258 414 420 258 414}
\RouteRow{Tríadas estables por bloque}{0 1 2 3 4 5 6 7 8 9 10 11 12 13 14 15 16 17}
\end{tabularx}\endgroup\par\vspace{5pt}\noindent{\fontsize{8.4}{10}\selectfont\hyperref[trace-143]{Procedencia y códigos originales}\hfill Ficha completa · 143/324}
\clearpage\phantomsection\label{nav:vi-ruta-144}
% VI-CATALOG-ROW rutas 144
\pdfbookmark[3]{r4c9\_h+1v+1}{route-144}
\RouteTitle{Ruta 144}{r4c9\_\allowbreak{}h+1v+1}
\noindent Celda 36\quad (4, 9)\qquad Orientación \(h=1,\ v=1\)\hfill\hyperref[sec:vi-indice-rutas]{Índice}\par\vspace{4pt}
\begingroup\fontsize{10.2}{12.5}\selectfont\setlength{\tabcolsep}{5pt}\renewcommand{\arraystretch}{1.1}\rowcolors{1}{routepale}{white}\noindent\begin{tabularx}{\linewidth}{@{}>{\raggedright\arraybackslash}p{.345\linewidth}>{\raggedright\arraybackslash}X@{}}
\RouteGroup{Identidad estructural}
\RouteRow{Clase y multisección}{quark\_\allowbreak{}down\_\allowbreak{}G4\quad /\quad quark\_\allowbreak{}down\_\allowbreak{}G4}
\RouteRow{Colección y sector}{quark de tipo abajo; quark}
\RouteRow{Clasificación física}{quark de tipo abajo}
\RouteRow{Generación, profundidad y color}{4; G4; G}
\RouteRow{Ruta primaria y órbita de inversión}{\(B_1\): no\qquad 4,\allowbreak{}9;\allowbreak{}6,\allowbreak{}1}
\RouteRow{Firma de clasificación}{34|\allowbreak{}0|\allowbreak{}-8|\allowbreak{}0|\allowbreak{}-12|\allowbreak{}+++-\kern0pt{}-\kern0pt{}-+++-\kern0pt{}-\kern0pt{}-}
\RouteGroup{Retornos, memoria y transporte}
\RouteRow{Retornos con fase}{clase: 27; órbita: 27; celda: 27}
\RouteRow{Retorno cinemático}{en 108: 54; extensión: 216}
\RouteRow{Giros y cambios de hoja}{71; 107}
\RouteRow{Acarreo y diferimiento}{deuda total: 2; final: 1; bloques diferidos: 2}
\RouteRow{Tríadas finales; persistencia de Pareto}{16; sí}
\RouteRow{\(K\longrightarrow K+9\)}{repeticiones: 9; cambios de memoria: 9}
\RouteGroup{Firma, fase y diferencias}
\RouteRow{\(n_A,\ s_C,\ \kappa_0,\ \nu_{120},\ \nu_{270}\)}{34, 0, 0, 0, -12}
\RouteRow{\(\tau_{12}\); firma histórica \(\mathbb Z^5\)}{+++-\kern0pt{}-\kern0pt{}-+++-\kern0pt{}-\kern0pt{}-; 34|\allowbreak{}0|\allowbreak{}0|\allowbreak{}0|\allowbreak{}-12}
\RouteRow{\(D_1,\ D_3,\ D_4,\ D_{27},\ D_{36}\)}{66, 60, 64, 48, 32}
\RouteRow{\(\Omega_{90/120},\ P_6\)}{80, 6}
\RouteGroup{Lectores y factores de acción}
\RouteRow{\(K_D\quad /\quad K_\Omega\)}{(80,\allowbreak{}66,\allowbreak{}2)\quad /\quad (80,\allowbreak{}54,\allowbreak{}6)}
\RouteRow{\(\kappa_D\)}{79.51859\allowbreak{}71232726\allowbreak{}53414287\allowbreak{}29}
\RouteRow{\(\kappa_\Omega\)}{79.60661\allowbreak{}01397948\allowbreak{}27586470\allowbreak{}91}
\RouteRow{\(R_{\mathrm{act},D}\)}{1.000000\allowbreak{}05512879\allowbreak{}47092488\allowbreak{}80}
\RouteRow{\(R_{\mathrm{act},\Omega}\)}{1.000000\allowbreak{}05518981\allowbreak{}25318591\allowbreak{}40}
\RouteGroup{Separación de prefijos}
\RouteRow{Área de ambigüedad; área \(\log_2\)}{18; 0.0}
\RouteRow{Primer bloque de separación}{órbita: 1; celda: 1; ruta: 1}
\RouteRow{Ambigüedad final}{1}
\RouteGroup{Secuencias completas}
\RouteRow{\(w_6\)}{327 498 99 327 501 486 423 15 486 327 498 99 327 501 486 423 15 486}
\RouteRow{Tríadas estables por bloque}{0 1 2 3 4 5 6 7 7 9 10 11 12 13 14 15 16 16}
\end{tabularx}\endgroup\par\vspace{5pt}\noindent{\fontsize{8.4}{10}\selectfont\hyperref[trace-144]{Procedencia y códigos originales}\hfill Ficha completa · 144/324}
\clearpage\phantomsection\label{nav:vi-ruta-145}
% VI-CATALOG-ROW rutas 145
\pdfbookmark[2]{Celda 37}{cell-37}
\pdfbookmark[3]{r5c1\_h-1v-1}{route-145}
\RouteTitle{Ruta 145}{r5c1\_\allowbreak{}h-1v-1}
\noindent Celda 37\quad (5, 1)\qquad Orientación \(h=-1,\ v=-1\)\hfill\hyperref[sec:vi-indice-rutas]{Índice}\par\vspace{4pt}
\begingroup\fontsize{10.2}{12.5}\selectfont\setlength{\tabcolsep}{5pt}\renewcommand{\arraystretch}{1.1}\rowcolors{1}{routepale}{white}\noindent\begin{tabularx}{\linewidth}{@{}>{\raggedright\arraybackslash}p{.345\linewidth}>{\raggedright\arraybackslash}X@{}}
\RouteGroup{Identidad estructural}
\RouteRow{Clase y multisección}{charged\_\allowbreak{}lepton\_\allowbreak{}G4\quad /\quad charged\_\allowbreak{}lepton\_\allowbreak{}G4}
\RouteRow{Colección y sector}{leptón cargado; leptón}
\RouteRow{Clasificación física}{leptón cargado}
\RouteRow{Generación, profundidad y color}{4; G4; \textemdash{}}
\RouteRow{Ruta primaria y órbita de inversión}{\(B_1\): no\qquad 5,\allowbreak{}1;\allowbreak{}5,\allowbreak{}9}
\RouteRow{Firma de clasificación}{34|\allowbreak{}0|\allowbreak{}-16|\allowbreak{}0|\allowbreak{}-8|\allowbreak{}+++-0-+++-0-}
\RouteGroup{Retornos, memoria y transporte}
\RouteRow{Retornos con fase}{clase: 27; órbita: 27; celda: 27}
\RouteRow{Retorno cinemático}{en 108: 54; extensión: 216}
\RouteRow{Giros y cambios de hoja}{71; 107}
\RouteRow{Acarreo y diferimiento}{deuda total: 1; final: 1; bloques diferidos: 1}
\RouteRow{Tríadas finales; persistencia de Pareto}{16; no}
\RouteRow{\(K\longrightarrow K+9\)}{repeticiones: 9; cambios de memoria: 9}
\RouteGroup{Firma, fase y diferencias}
\RouteRow{\(n_A,\ s_C,\ \kappa_0,\ \nu_{120},\ \nu_{270}\)}{34, 0, 2, 0, -8}
\RouteRow{\(\tau_{12}\); firma histórica \(\mathbb Z^5\)}{+0+-\kern0pt{}-\kern0pt{}-+0+-\kern0pt{}-\kern0pt{}-; 34|\allowbreak{}0|\allowbreak{}2|\allowbreak{}0|\allowbreak{}-8}
\RouteRow{\(D_1,\ D_3,\ D_4,\ D_{27},\ D_{36}\)}{66, 64, 68, 60, 40}
\RouteRow{\(\Omega_{90/120},\ P_6\)}{48, 5}
\RouteGroup{Lectores y factores de acción}
\RouteRow{\(K_D\quad /\quad K_\Omega\)}{(100,\allowbreak{}66,\allowbreak{}2)\quad /\quad (48,\allowbreak{}54,\allowbreak{}5)}
\RouteRow{\(\kappa_D\)}{99.51859\allowbreak{}71232726\allowbreak{}53414287\allowbreak{}29}
\RouteRow{\(\kappa_\Omega\)}{47.60649\allowbreak{}89433721\allowbreak{}35446416\allowbreak{}86}
\RouteRow{\(R_{\mathrm{act},D}\)}{1.000000\allowbreak{}06899443\allowbreak{}08259364\allowbreak{}17}
\RouteRow{\(R_{\mathrm{act},\Omega}\)}{1.000000\allowbreak{}03300471\allowbreak{}80532430\allowbreak{}84}
\RouteGroup{Separación de prefijos}
\RouteRow{Área de ambigüedad; área \(\log_2\)}{30; 8.0}
\RouteRow{Primer bloque de separación}{órbita: 5; celda: 5; ruta: 5}
\RouteRow{Ambigüedad final}{1}
\RouteGroup{Secuencias completas}
\RouteRow{\(w_6\)}{15 486 423 15 489 498 99 327 498 15 486 423 15 489 498 99 327 498}
\RouteRow{Tríadas estables por bloque}{0 1 2 3 4 5 6 7 8 9 10 11 12 13 14 15 16 16}
\end{tabularx}\endgroup\par\vspace{5pt}\noindent{\fontsize{8.4}{10}\selectfont\hyperref[trace-145]{Procedencia y códigos originales}\hfill Ficha completa · 145/324}
\clearpage\phantomsection\label{nav:vi-ruta-146}
% VI-CATALOG-ROW rutas 146
\pdfbookmark[3]{r5c1\_h-1v+1}{route-146}
\RouteTitle{Ruta 146}{r5c1\_\allowbreak{}h-1v+1}
\noindent Celda 37\quad (5, 1)\qquad Orientación \(h=-1,\ v=1\)\hfill\hyperref[sec:vi-indice-rutas]{Índice}\par\vspace{4pt}
\begingroup\fontsize{10.2}{12.5}\selectfont\setlength{\tabcolsep}{5pt}\renewcommand{\arraystretch}{1.1}\rowcolors{1}{routepale}{white}\noindent\begin{tabularx}{\linewidth}{@{}>{\raggedright\arraybackslash}p{.345\linewidth}>{\raggedright\arraybackslash}X@{}}
\RouteGroup{Identidad estructural}
\RouteRow{Clase y multisección}{charged\_\allowbreak{}lepton\_\allowbreak{}G4\quad /\quad charged\_\allowbreak{}lepton\_\allowbreak{}G4}
\RouteRow{Colección y sector}{leptón cargado; leptón}
\RouteRow{Clasificación física}{leptón cargado}
\RouteRow{Generación, profundidad y color}{4; G4; \textemdash{}}
\RouteRow{Ruta primaria y órbita de inversión}{\(B_1\): no\qquad 5,\allowbreak{}1;\allowbreak{}5,\allowbreak{}9}
\RouteRow{Firma de clasificación}{34|\allowbreak{}0|\allowbreak{}-16|\allowbreak{}0|\allowbreak{}-8|\allowbreak{}+++-0-+++-0-}
\RouteGroup{Retornos, memoria y transporte}
\RouteRow{Retornos con fase}{clase: 27; órbita: 27; celda: 27}
\RouteRow{Retorno cinemático}{en 108: 54; extensión: 216}
\RouteRow{Giros y cambios de hoja}{71; 107}
\RouteRow{Acarreo y diferimiento}{deuda total: 0; final: 0; bloques diferidos: 0}
\RouteRow{Tríadas finales; persistencia de Pareto}{17; no}
\RouteRow{\(K\longrightarrow K+9\)}{repeticiones: 9; cambios de memoria: 9}
\RouteGroup{Firma, fase y diferencias}
\RouteRow{\(n_A,\ s_C,\ \kappa_0,\ \nu_{120},\ \nu_{270}\)}{14, 6, 0, 0, -12}
\RouteRow{\(\tau_{12}\); firma histórica \(\mathbb Z^5\)}{+++-\kern0pt{}-\kern0pt{}-+++-\kern0pt{}-\kern0pt{}-; 14|\allowbreak{}6|\allowbreak{}0|\allowbreak{}0|\allowbreak{}-12}
\RouteRow{\(D_1,\ D_3,\ D_4,\ D_{27},\ D_{36}\)}{24, 20, 24, 24, 16}
\RouteRow{\(\Omega_{90/120},\ P_6\)}{80, 6}
\RouteGroup{Lectores y factores de acción}
\RouteRow{\(K_D\quad /\quad K_\Omega\)}{(40,\allowbreak{}24,\allowbreak{}2)\quad /\quad (80,\allowbreak{}54,\allowbreak{}6)}
\RouteRow{\(\kappa_D\)}{39.82508\allowbreak{}59311825\allowbreak{}73056173\allowbreak{}26}
\RouteRow{\(\kappa_\Omega\)}{79.60661\allowbreak{}01397948\allowbreak{}27586470\allowbreak{}91}
\RouteRow{\(R_{\mathrm{act},D}\)}{1.000000\allowbreak{}02761000\allowbreak{}61595142\allowbreak{}15}
\RouteRow{\(R_{\mathrm{act},\Omega}\)}{1.000000\allowbreak{}05518981\allowbreak{}25318591\allowbreak{}40}
\RouteGroup{Separación de prefijos}
\RouteRow{Área de ambigüedad; área \(\log_2\)}{36; 18.0}
\RouteRow{Primer bloque de separación}{órbita: \textemdash{}; celda: \textemdash{}; ruta: \textemdash{}}
\RouteRow{Ambigüedad final}{2}
\RouteGroup{Secuencias completas}
\RouteRow{\(w_6\)}{0 81 3 0 84 0 81 3 0 0 81 3 0 84 0 81 3 0}
\RouteRow{Tríadas estables por bloque}{0 1 2 3 4 5 6 7 8 9 10 11 12 13 14 15 16 17}
\end{tabularx}\endgroup\par\vspace{5pt}\noindent{\fontsize{8.4}{10}\selectfont\hyperref[trace-146]{Procedencia y códigos originales}\hfill Ficha completa · 146/324}
\clearpage\phantomsection\label{nav:vi-ruta-147}
% VI-CATALOG-ROW rutas 147
\pdfbookmark[3]{r5c1\_h+1v-1}{route-147}
\RouteTitle{Ruta 147}{r5c1\_\allowbreak{}h+1v-1}
\noindent Celda 37\quad (5, 1)\qquad Orientación \(h=1,\ v=-1\)\hfill\hyperref[sec:vi-indice-rutas]{Índice}\par\vspace{4pt}
\begingroup\fontsize{10.2}{12.5}\selectfont\setlength{\tabcolsep}{5pt}\renewcommand{\arraystretch}{1.1}\rowcolors{1}{routepale}{white}\noindent\begin{tabularx}{\linewidth}{@{}>{\raggedright\arraybackslash}p{.345\linewidth}>{\raggedright\arraybackslash}X@{}}
\RouteGroup{Identidad estructural}
\RouteRow{Clase y multisección}{charged\_\allowbreak{}lepton\_\allowbreak{}G4\quad /\quad charged\_\allowbreak{}lepton\_\allowbreak{}G4}
\RouteRow{Colección y sector}{leptón cargado; leptón}
\RouteRow{Clasificación física}{leptón cargado}
\RouteRow{Generación, profundidad y color}{4; G4; \textemdash{}}
\RouteRow{Ruta primaria y órbita de inversión}{\(B_1\): sí\qquad 5,\allowbreak{}1;\allowbreak{}5,\allowbreak{}9}
\RouteRow{Firma de clasificación}{34|\allowbreak{}0|\allowbreak{}-16|\allowbreak{}0|\allowbreak{}-8|\allowbreak{}+++-0-+++-0-}
\RouteGroup{Retornos, memoria y transporte}
\RouteRow{Retornos con fase}{clase: 27; órbita: 27; celda: 27}
\RouteRow{Retorno cinemático}{en 108: 54; extensión: 216}
\RouteRow{Giros y cambios de hoja}{71; 107}
\RouteRow{Acarreo y diferimiento}{deuda total: 2; final: 0; bloques diferidos: 2}
\RouteRow{Tríadas finales; persistencia de Pareto}{17; no}
\RouteRow{\(K\longrightarrow K+9\)}{repeticiones: 9; cambios de memoria: 9}
\RouteGroup{Firma, fase y diferencias}
\RouteRow{\(n_A,\ s_C,\ \kappa_0,\ \nu_{120},\ \nu_{270}\)}{14, -6, 0, 0, -12}
\RouteRow{\(\tau_{12}\); firma histórica \(\mathbb Z^5\)}{+++-\kern0pt{}-\kern0pt{}-+++-\kern0pt{}-\kern0pt{}-; 14|\allowbreak{}-6|\allowbreak{}0|\allowbreak{}0|\allowbreak{}-12}
\RouteRow{\(D_1,\ D_3,\ D_4,\ D_{27},\ D_{36}\)}{24, 20, 24, 24, 16}
\RouteRow{\(\Omega_{90/120},\ P_6\)}{80, 6}
\RouteGroup{Lectores y factores de acción}
\RouteRow{\(K_D\quad /\quad K_\Omega\)}{(40,\allowbreak{}24,\allowbreak{}2)\quad /\quad (80,\allowbreak{}54,\allowbreak{}6)}
\RouteRow{\(\kappa_D\)}{39.82508\allowbreak{}59311825\allowbreak{}73056173\allowbreak{}26}
\RouteRow{\(\kappa_\Omega\)}{79.60661\allowbreak{}01397948\allowbreak{}27586470\allowbreak{}91}
\RouteRow{\(R_{\mathrm{act},D}\)}{1.000000\allowbreak{}02761000\allowbreak{}61595142\allowbreak{}15}
\RouteRow{\(R_{\mathrm{act},\Omega}\)}{1.000000\allowbreak{}05518981\allowbreak{}25318591\allowbreak{}40}
\RouteGroup{Separación de prefijos}
\RouteRow{Área de ambigüedad; área \(\log_2\)}{36; 18.0}
\RouteRow{Primer bloque de separación}{órbita: \textemdash{}; celda: \textemdash{}; ruta: \textemdash{}}
\RouteRow{Ambigüedad final}{2}
\RouteGroup{Secuencias completas}
\RouteRow{\(w_6\)}{81 3 0 81 0 3 0 81 3 81 3 0 81 0 3 0 81 3}
\RouteRow{Tríadas estables por bloque}{0 1 2 3 4 5 6 6 8 9 10 11 12 13 14 15 15 17}
\end{tabularx}\endgroup\par\vspace{5pt}\noindent{\fontsize{8.4}{10}\selectfont\hyperref[trace-147]{Procedencia y códigos originales}\hfill Ficha completa · 147/324}
\clearpage\phantomsection\label{nav:vi-ruta-148}
% VI-CATALOG-ROW rutas 148
\pdfbookmark[3]{r5c1\_h+1v+1}{route-148}
\RouteTitle{Ruta 148}{r5c1\_\allowbreak{}h+1v+1}
\noindent Celda 37\quad (5, 1)\qquad Orientación \(h=1,\ v=1\)\hfill\hyperref[sec:vi-indice-rutas]{Índice}\par\vspace{4pt}
\begingroup\fontsize{10.2}{12.5}\selectfont\setlength{\tabcolsep}{5pt}\renewcommand{\arraystretch}{1.1}\rowcolors{1}{routepale}{white}\noindent\begin{tabularx}{\linewidth}{@{}>{\raggedright\arraybackslash}p{.345\linewidth}>{\raggedright\arraybackslash}X@{}}
\RouteGroup{Identidad estructural}
\RouteRow{Clase y multisección}{charged\_\allowbreak{}lepton\_\allowbreak{}G4\quad /\quad charged\_\allowbreak{}lepton\_\allowbreak{}G4}
\RouteRow{Colección y sector}{leptón cargado; leptón}
\RouteRow{Clasificación física}{leptón cargado}
\RouteRow{Generación, profundidad y color}{4; G4; \textemdash{}}
\RouteRow{Ruta primaria y órbita de inversión}{\(B_1\): no\qquad 5,\allowbreak{}1;\allowbreak{}5,\allowbreak{}9}
\RouteRow{Firma de clasificación}{34|\allowbreak{}0|\allowbreak{}-16|\allowbreak{}0|\allowbreak{}-8|\allowbreak{}+++-0-+++-0-}
\RouteGroup{Retornos, memoria y transporte}
\RouteRow{Retornos con fase}{clase: 27; órbita: 27; celda: 27}
\RouteRow{Retorno cinemático}{en 108: 54; extensión: 216}
\RouteRow{Giros y cambios de hoja}{71; 107}
\RouteRow{Acarreo y diferimiento}{deuda total: 1; final: 0; bloques diferidos: 1}
\RouteRow{Tríadas finales; persistencia de Pareto}{17; no}
\RouteRow{\(K\longrightarrow K+9\)}{repeticiones: 9; cambios de memoria: 9}
\RouteGroup{Firma, fase y diferencias}
\RouteRow{\(n_A,\ s_C,\ \kappa_0,\ \nu_{120},\ \nu_{270}\)}{34, 0, -2, 0, -8}
\RouteRow{\(\tau_{12}\); firma histórica \(\mathbb Z^5\)}{+++-0-+++-0-; 34|\allowbreak{}0|\allowbreak{}-2|\allowbreak{}0|\allowbreak{}-8}
\RouteRow{\(D_1,\ D_3,\ D_4,\ D_{27},\ D_{36}\)}{66, 64, 68, 60, 40}
\RouteRow{\(\Omega_{90/120},\ P_6\)}{48, 5}
\RouteGroup{Lectores y factores de acción}
\RouteRow{\(K_D\quad /\quad K_\Omega\)}{(100,\allowbreak{}66,\allowbreak{}2)\quad /\quad (48,\allowbreak{}54,\allowbreak{}5)}
\RouteRow{\(\kappa_D\)}{99.51859\allowbreak{}71232726\allowbreak{}53414287\allowbreak{}29}
\RouteRow{\(\kappa_\Omega\)}{47.60649\allowbreak{}89433721\allowbreak{}35446416\allowbreak{}86}
\RouteRow{\(R_{\mathrm{act},D}\)}{1.000000\allowbreak{}06899443\allowbreak{}08259364\allowbreak{}17}
\RouteRow{\(R_{\mathrm{act},\Omega}\)}{1.000000\allowbreak{}03300471\allowbreak{}80532430\allowbreak{}84}
\RouteGroup{Separación de prefijos}
\RouteRow{Área de ambigüedad; área \(\log_2\)}{18; 0.0}
\RouteRow{Primer bloque de separación}{órbita: 1; celda: 1; ruta: 1}
\RouteRow{Ambigüedad final}{1}
\RouteGroup{Secuencias completas}
\RouteRow{\(w_6\)}{99 327 498 99 324 423 15 486 423 99 327 498 99 324 423 15 486 423}
\RouteRow{Tríadas estables por bloque}{0 1 2 3 4 5 6 7 8 9 10 11 12 13 14 14 16 17}
\end{tabularx}\endgroup\par\vspace{5pt}\noindent{\fontsize{8.4}{10}\selectfont\hyperref[trace-148]{Procedencia y códigos originales}\hfill Ficha completa · 148/324}
\clearpage\phantomsection\label{nav:vi-ruta-149}
% VI-CATALOG-ROW rutas 149
\pdfbookmark[2]{Celda 38}{cell-38}
\pdfbookmark[3]{r5c2\_h-1v-1}{route-149}
\RouteTitle{Ruta 149}{r5c2\_\allowbreak{}h-1v-1}
\noindent Celda 38\quad (5, 2)\qquad Orientación \(h=-1,\ v=-1\)\hfill\hyperref[sec:vi-indice-rutas]{Índice}\par\vspace{4pt}
\begingroup\fontsize{10.2}{12.5}\selectfont\setlength{\tabcolsep}{5pt}\renewcommand{\arraystretch}{1.1}\rowcolors{1}{routepale}{white}\noindent\begin{tabularx}{\linewidth}{@{}>{\raggedright\arraybackslash}p{.345\linewidth}>{\raggedright\arraybackslash}X@{}}
\RouteGroup{Identidad estructural}
\RouteRow{Clase y multisección}{neutrino\_\allowbreak{}G3\quad /\quad neutrino\_\allowbreak{}G3}
\RouteRow{Colección y sector}{neutrino; leptón}
\RouteRow{Clasificación física}{neutrino}
\RouteRow{Generación, profundidad y color}{3; G3; \textemdash{}}
\RouteRow{Ruta primaria y órbita de inversión}{\(B_1\): sí\qquad 5,\allowbreak{}2;\allowbreak{}5,\allowbreak{}8}
\RouteRow{Firma de clasificación}{30|\allowbreak{}6|\allowbreak{}0|\allowbreak{}0|\allowbreak{}-12|\allowbreak{}+++-\kern0pt{}-\kern0pt{}-+++-\kern0pt{}-\kern0pt{}-}
\RouteGroup{Retornos, memoria y transporte}
\RouteRow{Retornos con fase}{clase: 27; órbita: 27; celda: 27}
\RouteRow{Retorno cinemático}{en 108: 54; extensión: 216}
\RouteRow{Giros y cambios de hoja}{71; 107}
\RouteRow{Acarreo y diferimiento}{deuda total: 3; final: 0; bloques diferidos: 2}
\RouteRow{Tríadas finales; persistencia de Pareto}{17; sí}
\RouteRow{\(K\longrightarrow K+9\)}{repeticiones: 9; cambios de memoria: 9}
\RouteGroup{Firma, fase y diferencias}
\RouteRow{\(n_A,\ s_C,\ \kappa_0,\ \nu_{120},\ \nu_{270}\)}{30, -6, 0, -4, -12}
\RouteRow{\(\tau_{12}\); firma histórica \(\mathbb Z^5\)}{+++-\kern0pt{}-\kern0pt{}-+++-\kern0pt{}-\kern0pt{}-; 30|\allowbreak{}-6|\allowbreak{}0|\allowbreak{}-4|\allowbreak{}-12}
\RouteRow{\(D_1,\ D_3,\ D_4,\ D_{27},\ D_{36}\)}{54, 56, 68, 48, 32}
\RouteRow{\(\Omega_{90/120},\ P_6\)}{80, 6}
\RouteGroup{Lectores y factores de acción}
\RouteRow{\(K_D\quad /\quad K_\Omega\)}{(80,\allowbreak{}54,\allowbreak{}6)\quad /\quad (80,\allowbreak{}54,\allowbreak{}6)}
\RouteRow{\(\kappa_D\)}{79.60661\allowbreak{}01397948\allowbreak{}27586470\allowbreak{}91}
\RouteRow{\(\kappa_\Omega\)}{79.60661\allowbreak{}01397948\allowbreak{}27586470\allowbreak{}91}
\RouteRow{\(R_{\mathrm{act},D}\)}{1.000000\allowbreak{}05518981\allowbreak{}25318591\allowbreak{}40}
\RouteRow{\(R_{\mathrm{act},\Omega}\)}{1.000000\allowbreak{}05518981\allowbreak{}25318591\allowbreak{}40}
\RouteGroup{Separación de prefijos}
\RouteRow{Área de ambigüedad; área \(\log_2\)}{36; 18.0}
\RouteRow{Primer bloque de separación}{órbita: 1; celda: \textemdash{}; ruta: \textemdash{}}
\RouteRow{Ambigüedad final}{2}
\RouteGroup{Secuencias completas}
\RouteRow{\(w_6\)}{423 15 486 423 9 24 243 657 24 423 15 486 423 9 24 243 657 24}
\RouteRow{Tríadas estables por bloque}{0 1 2 3 4 5 6 7 8 9 10 11 11 13 14 14 15 17}
\end{tabularx}\endgroup\par\vspace{5pt}\noindent{\fontsize{8.4}{10}\selectfont\hyperref[trace-149]{Procedencia y códigos originales}\hfill Ficha completa · 149/324}
\clearpage\phantomsection\label{nav:vi-ruta-150}
% VI-CATALOG-ROW rutas 150
\pdfbookmark[3]{r5c2\_h-1v+1}{route-150}
\RouteTitle{Ruta 150}{r5c2\_\allowbreak{}h-1v+1}
\noindent Celda 38\quad (5, 2)\qquad Orientación \(h=-1,\ v=1\)\hfill\hyperref[sec:vi-indice-rutas]{Índice}\par\vspace{4pt}
\begingroup\fontsize{10.2}{12.5}\selectfont\setlength{\tabcolsep}{5pt}\renewcommand{\arraystretch}{1.1}\rowcolors{1}{routepale}{white}\noindent\begin{tabularx}{\linewidth}{@{}>{\raggedright\arraybackslash}p{.345\linewidth}>{\raggedright\arraybackslash}X@{}}
\RouteGroup{Identidad estructural}
\RouteRow{Clase y multisección}{neutrino\_\allowbreak{}G3\quad /\quad neutrino\_\allowbreak{}G3}
\RouteRow{Colección y sector}{neutrino; leptón}
\RouteRow{Clasificación física}{neutrino}
\RouteRow{Generación, profundidad y color}{3; G3; \textemdash{}}
\RouteRow{Ruta primaria y órbita de inversión}{\(B_1\): no\qquad 5,\allowbreak{}2;\allowbreak{}5,\allowbreak{}8}
\RouteRow{Firma de clasificación}{30|\allowbreak{}6|\allowbreak{}0|\allowbreak{}0|\allowbreak{}-12|\allowbreak{}+++-\kern0pt{}-\kern0pt{}-+++-\kern0pt{}-\kern0pt{}-}
\RouteGroup{Retornos, memoria y transporte}
\RouteRow{Retornos con fase}{clase: 27; órbita: 27; celda: 27}
\RouteRow{Retorno cinemático}{en 108: 54; extensión: 216}
\RouteRow{Giros y cambios de hoja}{71; 107}
\RouteRow{Acarreo y diferimiento}{deuda total: 2; final: 1; bloques diferidos: 2}
\RouteRow{Tríadas finales; persistencia de Pareto}{16; sí}
\RouteRow{\(K\longrightarrow K+9\)}{repeticiones: 9; cambios de memoria: 9}
\RouteGroup{Firma, fase y diferencias}
\RouteRow{\(n_A,\ s_C,\ \kappa_0,\ \nu_{120},\ \nu_{270}\)}{50, 4, -4, 0, -4}
\RouteRow{\(\tau_{12}\); firma histórica \(\mathbb Z^5\)}{+++-00+++-00; 50|\allowbreak{}4|\allowbreak{}-4|\allowbreak{}0|\allowbreak{}-4}
\RouteRow{\(D_1,\ D_3,\ D_4,\ D_{27},\ D_{36}\)}{84, 28, 84, 36, 24}
\RouteRow{\(\Omega_{90/120},\ P_6\)}{20, 4}
\RouteGroup{Lectores y factores de acción}
\RouteRow{\(K_D\quad /\quad K_\Omega\)}{(60,\allowbreak{}84,\allowbreak{}28)\quad /\quad (20,\allowbreak{}54,\allowbreak{}4)}
\RouteRow{\(\kappa_D\)}{59.39013\allowbreak{}58840155\allowbreak{}40637741\allowbreak{}48}
\RouteRow{\(\kappa_\Omega\)}{19.60638\allowbreak{}77469494\allowbreak{}43306362\allowbreak{}81}
\RouteRow{\(R_{\mathrm{act},D}\)}{1.000000\allowbreak{}04117409\allowbreak{}89467415\allowbreak{}75}
\RouteRow{\(R_{\mathrm{act},\Omega}\)}{1.000000\allowbreak{}01359275\allowbreak{}11518874\allowbreak{}98}
\RouteGroup{Separación de prefijos}
\RouteRow{Área de ambigüedad; área \(\log_2\)}{36; 18.0}
\RouteRow{Primer bloque de separación}{órbita: 1; celda: \textemdash{}; ruta: \textemdash{}}
\RouteRow{Ambigüedad final}{2}
\RouteGroup{Secuencias completas}
\RouteRow{\(w_6\)}{414 420 258 414 414 333 255 252 333 414 420 258 414 414 333 255 252 333}
\RouteRow{Tríadas estables por bloque}{0 1 2 3 4 5 6 7 8 9 10 11 12 13 14 14 16 16}
\end{tabularx}\endgroup\par\vspace{5pt}\noindent{\fontsize{8.4}{10}\selectfont\hyperref[trace-150]{Procedencia y códigos originales}\hfill Ficha completa · 150/324}
\clearpage\phantomsection\label{nav:vi-ruta-151}
% VI-CATALOG-ROW rutas 151
\pdfbookmark[3]{r5c2\_h+1v-1}{route-151}
\RouteTitle{Ruta 151}{r5c2\_\allowbreak{}h+1v-1}
\noindent Celda 38\quad (5, 2)\qquad Orientación \(h=1,\ v=-1\)\hfill\hyperref[sec:vi-indice-rutas]{Índice}\par\vspace{4pt}
\begingroup\fontsize{10.2}{12.5}\selectfont\setlength{\tabcolsep}{5pt}\renewcommand{\arraystretch}{1.1}\rowcolors{1}{routepale}{white}\noindent\begin{tabularx}{\linewidth}{@{}>{\raggedright\arraybackslash}p{.345\linewidth}>{\raggedright\arraybackslash}X@{}}
\RouteGroup{Identidad estructural}
\RouteRow{Clase y multisección}{neutrino\_\allowbreak{}G3\quad /\quad neutrino\_\allowbreak{}G3}
\RouteRow{Colección y sector}{neutrino; leptón}
\RouteRow{Clasificación física}{neutrino}
\RouteRow{Generación, profundidad y color}{3; G3; \textemdash{}}
\RouteRow{Ruta primaria y órbita de inversión}{\(B_1\): no\qquad 5,\allowbreak{}2;\allowbreak{}5,\allowbreak{}8}
\RouteRow{Firma de clasificación}{30|\allowbreak{}6|\allowbreak{}0|\allowbreak{}0|\allowbreak{}-12|\allowbreak{}+++-\kern0pt{}-\kern0pt{}-+++-\kern0pt{}-\kern0pt{}-}
\RouteGroup{Retornos, memoria y transporte}
\RouteRow{Retornos con fase}{clase: 27; órbita: 27; celda: 27}
\RouteRow{Retorno cinemático}{en 108: 54; extensión: 216}
\RouteRow{Giros y cambios de hoja}{71; 107}
\RouteRow{Acarreo y diferimiento}{deuda total: 3; final: 0; bloques diferidos: 2}
\RouteRow{Tríadas finales; persistencia de Pareto}{17; sí}
\RouteRow{\(K\longrightarrow K+9\)}{repeticiones: 9; cambios de memoria: 9}
\RouteGroup{Firma, fase y diferencias}
\RouteRow{\(n_A,\ s_C,\ \kappa_0,\ \nu_{120},\ \nu_{270}\)}{50, -4, 4, 0, -4}
\RouteRow{\(\tau_{12}\); firma histórica \(\mathbb Z^5\)}{+00-\kern0pt{}-\kern0pt{}-+00-\kern0pt{}-\kern0pt{}-; 50|\allowbreak{}-4|\allowbreak{}4|\allowbreak{}0|\allowbreak{}-4}
\RouteRow{\(D_1,\ D_3,\ D_4,\ D_{27},\ D_{36}\)}{84, 28, 84, 36, 24}
\RouteRow{\(\Omega_{90/120},\ P_6\)}{20, 4}
\RouteGroup{Lectores y factores de acción}
\RouteRow{\(K_D\quad /\quad K_\Omega\)}{(60,\allowbreak{}84,\allowbreak{}28)\quad /\quad (20,\allowbreak{}54,\allowbreak{}4)}
\RouteRow{\(\kappa_D\)}{59.39013\allowbreak{}58840155\allowbreak{}40637741\allowbreak{}48}
\RouteRow{\(\kappa_\Omega\)}{19.60638\allowbreak{}77469494\allowbreak{}43306362\allowbreak{}81}
\RouteRow{\(R_{\mathrm{act},D}\)}{1.000000\allowbreak{}04117409\allowbreak{}89467415\allowbreak{}75}
\RouteRow{\(R_{\mathrm{act},\Omega}\)}{1.000000\allowbreak{}01359275\allowbreak{}11518874\allowbreak{}98}
\RouteGroup{Separación de prefijos}
\RouteRow{Área de ambigüedad; área \(\log_2\)}{36; 18.0}
\RouteRow{Primer bloque de separación}{órbita: 1; celda: \textemdash{}; ruta: \textemdash{}}
\RouteRow{Ambigüedad final}{2}
\RouteGroup{Secuencias completas}
\RouteRow{\(w_6\)}{255 252 333 255 258 258 414 420 258 255 252 333 255 258 258 414 420 258}
\RouteRow{Tríadas estables por bloque}{0 1 2 3 4 5 6 7 8 9 10 11 11 12 14 15 15 17}
\end{tabularx}\endgroup\par\vspace{5pt}\noindent{\fontsize{8.4}{10}\selectfont\hyperref[trace-151]{Procedencia y códigos originales}\hfill Ficha completa · 151/324}
\clearpage\phantomsection\label{nav:vi-ruta-152}
% VI-CATALOG-ROW rutas 152
\pdfbookmark[3]{r5c2\_h+1v+1}{route-152}
\RouteTitle{Ruta 152}{r5c2\_\allowbreak{}h+1v+1}
\noindent Celda 38\quad (5, 2)\qquad Orientación \(h=1,\ v=1\)\hfill\hyperref[sec:vi-indice-rutas]{Índice}\par\vspace{4pt}
\begingroup\fontsize{10.2}{12.5}\selectfont\setlength{\tabcolsep}{5pt}\renewcommand{\arraystretch}{1.1}\rowcolors{1}{routepale}{white}\noindent\begin{tabularx}{\linewidth}{@{}>{\raggedright\arraybackslash}p{.345\linewidth}>{\raggedright\arraybackslash}X@{}}
\RouteGroup{Identidad estructural}
\RouteRow{Clase y multisección}{neutrino\_\allowbreak{}G3\quad /\quad neutrino\_\allowbreak{}G3}
\RouteRow{Colección y sector}{neutrino; leptón}
\RouteRow{Clasificación física}{neutrino}
\RouteRow{Generación, profundidad y color}{3; G3; \textemdash{}}
\RouteRow{Ruta primaria y órbita de inversión}{\(B_1\): no\qquad 5,\allowbreak{}2;\allowbreak{}5,\allowbreak{}8}
\RouteRow{Firma de clasificación}{30|\allowbreak{}6|\allowbreak{}0|\allowbreak{}0|\allowbreak{}-12|\allowbreak{}+++-\kern0pt{}-\kern0pt{}-+++-\kern0pt{}-\kern0pt{}-}
\RouteGroup{Retornos, memoria y transporte}
\RouteRow{Retornos con fase}{clase: 27; órbita: 27; celda: 27}
\RouteRow{Retorno cinemático}{en 108: 54; extensión: 216}
\RouteRow{Giros y cambios de hoja}{71; 107}
\RouteRow{Acarreo y diferimiento}{deuda total: 2; final: 1; bloques diferidos: 2}
\RouteRow{Tríadas finales; persistencia de Pareto}{16; sí}
\RouteRow{\(K\longrightarrow K+9\)}{repeticiones: 9; cambios de memoria: 9}
\RouteGroup{Firma, fase y diferencias}
\RouteRow{\(n_A,\ s_C,\ \kappa_0,\ \nu_{120},\ \nu_{270}\)}{30, 6, 0, 4, -12}
\RouteRow{\(\tau_{12}\); firma histórica \(\mathbb Z^5\)}{+++-\kern0pt{}-\kern0pt{}-+++-\kern0pt{}-\kern0pt{}-; 30|\allowbreak{}6|\allowbreak{}0|\allowbreak{}4|\allowbreak{}-12}
\RouteRow{\(D_1,\ D_3,\ D_4,\ D_{27},\ D_{36}\)}{54, 56, 68, 48, 32}
\RouteRow{\(\Omega_{90/120},\ P_6\)}{80, 6}
\RouteGroup{Lectores y factores de acción}
\RouteRow{\(K_D\quad /\quad K_\Omega\)}{(80,\allowbreak{}54,\allowbreak{}6)\quad /\quad (80,\allowbreak{}54,\allowbreak{}6)}
\RouteRow{\(\kappa_D\)}{79.60661\allowbreak{}01397948\allowbreak{}27586470\allowbreak{}91}
\RouteRow{\(\kappa_\Omega\)}{79.60661\allowbreak{}01397948\allowbreak{}27586470\allowbreak{}91}
\RouteRow{\(R_{\mathrm{act},D}\)}{1.000000\allowbreak{}05518981\allowbreak{}25318591\allowbreak{}40}
\RouteRow{\(R_{\mathrm{act},\Omega}\)}{1.000000\allowbreak{}05518981\allowbreak{}25318591\allowbreak{}40}
\RouteGroup{Separación de prefijos}
\RouteRow{Área de ambigüedad; área \(\log_2\)}{36; 18.0}
\RouteRow{Primer bloque de separación}{órbita: 1; celda: \textemdash{}; ruta: \textemdash{}}
\RouteRow{Ambigüedad final}{2}
\RouteGroup{Secuencias completas}
\RouteRow{\(w_6\)}{243 657 24 243 663 486 423 15 486 243 657 24 243 663 486 423 15 486}
\RouteRow{Tríadas estables por bloque}{0 1 2 3 4 5 6 7 8 9 10 11 12 13 14 14 16 16}
\end{tabularx}\endgroup\par\vspace{5pt}\noindent{\fontsize{8.4}{10}\selectfont\hyperref[trace-152]{Procedencia y códigos originales}\hfill Ficha completa · 152/324}
\clearpage\phantomsection\label{nav:vi-ruta-153}
% VI-CATALOG-ROW rutas 153
\pdfbookmark[2]{Celda 39}{cell-39}
\pdfbookmark[3]{r5c3\_h-1v-1}{route-153}
\RouteTitle{Ruta 153}{r5c3\_\allowbreak{}h-1v-1}
\noindent Celda 39\quad (5, 3)\qquad Orientación \(h=-1,\ v=-1\)\hfill\hyperref[sec:vi-indice-rutas]{Índice}\par\vspace{4pt}
\begingroup\fontsize{10.2}{12.5}\selectfont\setlength{\tabcolsep}{5pt}\renewcommand{\arraystretch}{1.1}\rowcolors{1}{routepale}{white}\noindent\begin{tabularx}{\linewidth}{@{}>{\raggedright\arraybackslash}p{.345\linewidth}>{\raggedright\arraybackslash}X@{}}
\RouteGroup{Identidad estructural}
\RouteRow{Clase y multisección}{charged\_\allowbreak{}lepton\_\allowbreak{}G2\quad /\quad charged\_\allowbreak{}lepton\_\allowbreak{}G2}
\RouteRow{Colección y sector}{leptón cargado; leptón}
\RouteRow{Clasificación física}{leptón cargado}
\RouteRow{Generación, profundidad y color}{2; G2; \textemdash{}}
\RouteRow{Ruta primaria y órbita de inversión}{\(B_1\): no\qquad 5,\allowbreak{}3;\allowbreak{}5,\allowbreak{}7}
\RouteRow{Firma de clasificación}{28|\allowbreak{}6|\allowbreak{}2|\allowbreak{}0|\allowbreak{}-12|\allowbreak{}+++-\kern0pt{}-\kern0pt{}-+++-\kern0pt{}-\kern0pt{}-}
\RouteGroup{Retornos, memoria y transporte}
\RouteRow{Retornos con fase}{clase: 27; órbita: 27; celda: 27}
\RouteRow{Retorno cinemático}{en 108: 54; extensión: 216}
\RouteRow{Giros y cambios de hoja}{71; 107}
\RouteRow{Acarreo y diferimiento}{deuda total: 1; final: 0; bloques diferidos: 1}
\RouteRow{Tríadas finales; persistencia de Pareto}{17; no}
\RouteRow{\(K\longrightarrow K+9\)}{repeticiones: 9; cambios de memoria: 9}
\RouteGroup{Firma, fase y diferencias}
\RouteRow{\(n_A,\ s_C,\ \kappa_0,\ \nu_{120},\ \nu_{270}\)}{28, -6, 0, 0, -12}
\RouteRow{\(\tau_{12}\); firma histórica \(\mathbb Z^5\)}{+++-\kern0pt{}-\kern0pt{}-+++-\kern0pt{}-\kern0pt{}-; 28|\allowbreak{}-6|\allowbreak{}0|\allowbreak{}0|\allowbreak{}-12}
\RouteRow{\(D_1,\ D_3,\ D_4,\ D_{27},\ D_{36}\)}{66, 60, 68, 48, 32}
\RouteRow{\(\Omega_{90/120},\ P_6\)}{80, 6}
\RouteGroup{Lectores y factores de acción}
\RouteRow{\(K_D\quad /\quad K_\Omega\)}{(80,\allowbreak{}66,\allowbreak{}4)\quad /\quad (80,\allowbreak{}54,\allowbreak{}6)}
\RouteRow{\(\kappa_D\)}{79.51881\allowbreak{}95161180\allowbreak{}37694395\allowbreak{}39}
\RouteRow{\(\kappa_\Omega\)}{79.60661\allowbreak{}01397948\allowbreak{}27586470\allowbreak{}91}
\RouteRow{\(R_{\mathrm{act},D}\)}{1.000000\allowbreak{}05512894\allowbreak{}88901612\allowbreak{}64}
\RouteRow{\(R_{\mathrm{act},\Omega}\)}{1.000000\allowbreak{}05518981\allowbreak{}25318591\allowbreak{}40}
\RouteGroup{Separación de prefijos}
\RouteRow{Área de ambigüedad; área \(\log_2\)}{18; 0.0}
\RouteRow{Primer bloque de separación}{órbita: 1; celda: 1; ruta: 1}
\RouteRow{Ambigüedad final}{1}
\RouteGroup{Secuencias completas}
\RouteRow{\(w_6\)}{570 264 90 570 267 243 657 24 243 570 264 90 570 267 243 657 24 243}
\RouteRow{Tríadas estables por bloque}{0 1 1 3 4 5 6 7 8 9 10 11 12 13 14 15 16 17}
\end{tabularx}\endgroup\par\vspace{5pt}\noindent{\fontsize{8.4}{10}\selectfont\hyperref[trace-153]{Procedencia y códigos originales}\hfill Ficha completa · 153/324}
\clearpage\phantomsection\label{nav:vi-ruta-154}
% VI-CATALOG-ROW rutas 154
\pdfbookmark[3]{r5c3\_h-1v+1}{route-154}
\RouteTitle{Ruta 154}{r5c3\_\allowbreak{}h-1v+1}
\noindent Celda 39\quad (5, 3)\qquad Orientación \(h=-1,\ v=1\)\hfill\hyperref[sec:vi-indice-rutas]{Índice}\par\vspace{4pt}
\begingroup\fontsize{10.2}{12.5}\selectfont\setlength{\tabcolsep}{5pt}\renewcommand{\arraystretch}{1.1}\rowcolors{1}{routepale}{white}\noindent\begin{tabularx}{\linewidth}{@{}>{\raggedright\arraybackslash}p{.345\linewidth}>{\raggedright\arraybackslash}X@{}}
\RouteGroup{Identidad estructural}
\RouteRow{Clase y multisección}{charged\_\allowbreak{}lepton\_\allowbreak{}G2\quad /\quad charged\_\allowbreak{}lepton\_\allowbreak{}G2}
\RouteRow{Colección y sector}{leptón cargado; leptón}
\RouteRow{Clasificación física}{leptón cargado}
\RouteRow{Generación, profundidad y color}{2; G2; \textemdash{}}
\RouteRow{Ruta primaria y órbita de inversión}{\(B_1\): sí\qquad 5,\allowbreak{}3;\allowbreak{}5,\allowbreak{}7}
\RouteRow{Firma de clasificación}{28|\allowbreak{}6|\allowbreak{}2|\allowbreak{}0|\allowbreak{}-12|\allowbreak{}+++-\kern0pt{}-\kern0pt{}-+++-\kern0pt{}-\kern0pt{}-}
\RouteGroup{Retornos, memoria y transporte}
\RouteRow{Retornos con fase}{clase: 27; órbita: 27; celda: 27}
\RouteRow{Retorno cinemático}{en 108: 54; extensión: 216}
\RouteRow{Giros y cambios de hoja}{71; 107}
\RouteRow{Acarreo y diferimiento}{deuda total: 1; final: 0; bloques diferidos: 1}
\RouteRow{Tríadas finales; persistencia de Pareto}{17; no}
\RouteRow{\(K\longrightarrow K+9\)}{repeticiones: 9; cambios de memoria: 9}
\RouteGroup{Firma, fase y diferencias}
\RouteRow{\(n_A,\ s_C,\ \kappa_0,\ \nu_{120},\ \nu_{270}\)}{14, -4, 0, 0, -12}
\RouteRow{\(\tau_{12}\); firma histórica \(\mathbb Z^5\)}{+++-\kern0pt{}-\kern0pt{}-+++-\kern0pt{}-\kern0pt{}-; 14|\allowbreak{}-4|\allowbreak{}0|\allowbreak{}0|\allowbreak{}-12}
\RouteRow{\(D_1,\ D_3,\ D_4,\ D_{27},\ D_{36}\)}{78, 24, 78, 24, 16}
\RouteRow{\(\Omega_{90/120},\ P_6\)}{80, 6}
\RouteGroup{Lectores y factores de acción}
\RouteRow{\(K_D\quad /\quad K_\Omega\)}{(40,\allowbreak{}78,\allowbreak{}27)\quad /\quad (80,\allowbreak{}54,\allowbreak{}6)}
\RouteRow{\(\kappa_D\)}{39.43380\allowbreak{}88030085\allowbreak{}51303671\allowbreak{}14}
\RouteRow{\(\kappa_\Omega\)}{79.60661\allowbreak{}01397948\allowbreak{}27586470\allowbreak{}91}
\RouteRow{\(R_{\mathrm{act},D}\)}{1.000000\allowbreak{}02733874\allowbreak{}08548943\allowbreak{}58}
\RouteRow{\(R_{\mathrm{act},\Omega}\)}{1.000000\allowbreak{}05518981\allowbreak{}25318591\allowbreak{}40}
\RouteGroup{Separación de prefijos}
\RouteRow{Área de ambigüedad; área \(\log_2\)}{18; 0.0}
\RouteRow{Primer bloque de separación}{órbita: 1; celda: 1; ruta: 1}
\RouteRow{Ambigüedad final}{1}
\RouteGroup{Secuencias completas}
\RouteRow{\(w_6\)}{585 507 504 585 510 666 672 510 666 585 507 504 585 510 666 672 510 666}
\RouteRow{Tríadas estables por bloque}{0 1 2 3 4 5 6 7 8 9 10 11 12 13 14 14 16 17}
\end{tabularx}\endgroup\par\vspace{5pt}\noindent{\fontsize{8.4}{10}\selectfont\hyperref[trace-154]{Procedencia y códigos originales}\hfill Ficha completa · 154/324}
\clearpage\phantomsection\label{nav:vi-ruta-155}
% VI-CATALOG-ROW rutas 155
\pdfbookmark[3]{r5c3\_h+1v-1}{route-155}
\RouteTitle{Ruta 155}{r5c3\_\allowbreak{}h+1v-1}
\noindent Celda 39\quad (5, 3)\qquad Orientación \(h=1,\ v=-1\)\hfill\hyperref[sec:vi-indice-rutas]{Índice}\par\vspace{4pt}
\begingroup\fontsize{10.2}{12.5}\selectfont\setlength{\tabcolsep}{5pt}\renewcommand{\arraystretch}{1.1}\rowcolors{1}{routepale}{white}\noindent\begin{tabularx}{\linewidth}{@{}>{\raggedright\arraybackslash}p{.345\linewidth}>{\raggedright\arraybackslash}X@{}}
\RouteGroup{Identidad estructural}
\RouteRow{Clase y multisección}{charged\_\allowbreak{}lepton\_\allowbreak{}G2\quad /\quad charged\_\allowbreak{}lepton\_\allowbreak{}G2}
\RouteRow{Colección y sector}{leptón cargado; leptón}
\RouteRow{Clasificación física}{leptón cargado}
\RouteRow{Generación, profundidad y color}{2; G2; \textemdash{}}
\RouteRow{Ruta primaria y órbita de inversión}{\(B_1\): no\qquad 5,\allowbreak{}3;\allowbreak{}5,\allowbreak{}7}
\RouteRow{Firma de clasificación}{28|\allowbreak{}6|\allowbreak{}2|\allowbreak{}0|\allowbreak{}-12|\allowbreak{}+++-\kern0pt{}-\kern0pt{}-+++-\kern0pt{}-\kern0pt{}-}
\RouteGroup{Retornos, memoria y transporte}
\RouteRow{Retornos con fase}{clase: 27; órbita: 27; celda: 27}
\RouteRow{Retorno cinemático}{en 108: 54; extensión: 216}
\RouteRow{Giros y cambios de hoja}{71; 107}
\RouteRow{Acarreo y diferimiento}{deuda total: 1; final: 0; bloques diferidos: 1}
\RouteRow{Tríadas finales; persistencia de Pareto}{17; no}
\RouteRow{\(K\longrightarrow K+9\)}{repeticiones: 9; cambios de memoria: 9}
\RouteGroup{Firma, fase y diferencias}
\RouteRow{\(n_A,\ s_C,\ \kappa_0,\ \nu_{120},\ \nu_{270}\)}{14, 4, 0, 0, -12}
\RouteRow{\(\tau_{12}\); firma histórica \(\mathbb Z^5\)}{+++-\kern0pt{}-\kern0pt{}-+++-\kern0pt{}-\kern0pt{}-; 14|\allowbreak{}4|\allowbreak{}0|\allowbreak{}0|\allowbreak{}-12}
\RouteRow{\(D_1,\ D_3,\ D_4,\ D_{27},\ D_{36}\)}{78, 24, 78, 24, 16}
\RouteRow{\(\Omega_{90/120},\ P_6\)}{80, 6}
\RouteGroup{Lectores y factores de acción}
\RouteRow{\(K_D\quad /\quad K_\Omega\)}{(40,\allowbreak{}78,\allowbreak{}27)\quad /\quad (80,\allowbreak{}54,\allowbreak{}6)}
\RouteRow{\(\kappa_D\)}{39.43380\allowbreak{}88030085\allowbreak{}51303671\allowbreak{}14}
\RouteRow{\(\kappa_\Omega\)}{79.60661\allowbreak{}01397948\allowbreak{}27586470\allowbreak{}91}
\RouteRow{\(R_{\mathrm{act},D}\)}{1.000000\allowbreak{}02733874\allowbreak{}08548943\allowbreak{}58}
\RouteRow{\(R_{\mathrm{act},\Omega}\)}{1.000000\allowbreak{}05518981\allowbreak{}25318591\allowbreak{}40}
\RouteGroup{Separación de prefijos}
\RouteRow{Área de ambigüedad; área \(\log_2\)}{18; 0.0}
\RouteRow{Primer bloque de separación}{órbita: 1; celda: 1; ruta: 1}
\RouteRow{Ambigüedad final}{1}
\RouteGroup{Secuencias completas}
\RouteRow{\(w_6\)}{672 510 666 672 507 504 585 507 504 672 510 666 672 507 504 585 507 504}
\RouteRow{Tríadas estables por bloque}{0 1 2 3 4 5 6 7 8 9 10 11 12 13 14 14 16 17}
\end{tabularx}\endgroup\par\vspace{5pt}\noindent{\fontsize{8.4}{10}\selectfont\hyperref[trace-155]{Procedencia y códigos originales}\hfill Ficha completa · 155/324}
\clearpage\phantomsection\label{nav:vi-ruta-156}
% VI-CATALOG-ROW rutas 156
\pdfbookmark[3]{r5c3\_h+1v+1}{route-156}
\RouteTitle{Ruta 156}{r5c3\_\allowbreak{}h+1v+1}
\noindent Celda 39\quad (5, 3)\qquad Orientación \(h=1,\ v=1\)\hfill\hyperref[sec:vi-indice-rutas]{Índice}\par\vspace{4pt}
\begingroup\fontsize{10.2}{12.5}\selectfont\setlength{\tabcolsep}{5pt}\renewcommand{\arraystretch}{1.1}\rowcolors{1}{routepale}{white}\noindent\begin{tabularx}{\linewidth}{@{}>{\raggedright\arraybackslash}p{.345\linewidth}>{\raggedright\arraybackslash}X@{}}
\RouteGroup{Identidad estructural}
\RouteRow{Clase y multisección}{charged\_\allowbreak{}lepton\_\allowbreak{}G2\quad /\quad charged\_\allowbreak{}lepton\_\allowbreak{}G2}
\RouteRow{Colección y sector}{leptón cargado; leptón}
\RouteRow{Clasificación física}{leptón cargado}
\RouteRow{Generación, profundidad y color}{2; G2; \textemdash{}}
\RouteRow{Ruta primaria y órbita de inversión}{\(B_1\): no\qquad 5,\allowbreak{}3;\allowbreak{}5,\allowbreak{}7}
\RouteRow{Firma de clasificación}{28|\allowbreak{}6|\allowbreak{}2|\allowbreak{}0|\allowbreak{}-12|\allowbreak{}+++-\kern0pt{}-\kern0pt{}-+++-\kern0pt{}-\kern0pt{}-}
\RouteGroup{Retornos, memoria y transporte}
\RouteRow{Retornos con fase}{clase: 27; órbita: 27; celda: 27}
\RouteRow{Retorno cinemático}{en 108: 54; extensión: 216}
\RouteRow{Giros y cambios de hoja}{71; 107}
\RouteRow{Acarreo y diferimiento}{deuda total: 0; final: 0; bloques diferidos: 0}
\RouteRow{Tríadas finales; persistencia de Pareto}{17; no}
\RouteRow{\(K\longrightarrow K+9\)}{repeticiones: 9; cambios de memoria: 9}
\RouteGroup{Firma, fase y diferencias}
\RouteRow{\(n_A,\ s_C,\ \kappa_0,\ \nu_{120},\ \nu_{270}\)}{28, 6, 0, 0, -12}
\RouteRow{\(\tau_{12}\); firma histórica \(\mathbb Z^5\)}{+++-\kern0pt{}-\kern0pt{}-+++-\kern0pt{}-\kern0pt{}-; 28|\allowbreak{}6|\allowbreak{}0|\allowbreak{}0|\allowbreak{}-12}
\RouteRow{\(D_1,\ D_3,\ D_4,\ D_{27},\ D_{36}\)}{66, 60, 68, 48, 32}
\RouteRow{\(\Omega_{90/120},\ P_6\)}{80, 6}
\RouteGroup{Lectores y factores de acción}
\RouteRow{\(K_D\quad /\quad K_\Omega\)}{(80,\allowbreak{}66,\allowbreak{}4)\quad /\quad (80,\allowbreak{}54,\allowbreak{}6)}
\RouteRow{\(\kappa_D\)}{79.51881\allowbreak{}95161180\allowbreak{}37694395\allowbreak{}39}
\RouteRow{\(\kappa_\Omega\)}{79.60661\allowbreak{}01397948\allowbreak{}27586470\allowbreak{}91}
\RouteRow{\(R_{\mathrm{act},D}\)}{1.000000\allowbreak{}05512894\allowbreak{}88901612\allowbreak{}64}
\RouteRow{\(R_{\mathrm{act},\Omega}\)}{1.000000\allowbreak{}05518981\allowbreak{}25318591\allowbreak{}40}
\RouteGroup{Separación de prefijos}
\RouteRow{Área de ambigüedad; área \(\log_2\)}{22; 4.0}
\RouteRow{Primer bloque de separación}{órbita: 5; celda: 5; ruta: 5}
\RouteRow{Ambigüedad final}{1}
\RouteGroup{Secuencias completas}
\RouteRow{\(w_6\)}{657 24 243 657 21 90 570 264 90 657 24 243 657 21 90 570 264 90}
\RouteRow{Tríadas estables por bloque}{0 1 2 3 4 5 6 7 8 9 10 11 12 13 14 15 16 17}
\end{tabularx}\endgroup\par\vspace{5pt}\noindent{\fontsize{8.4}{10}\selectfont\hyperref[trace-156]{Procedencia y códigos originales}\hfill Ficha completa · 156/324}
\clearpage\phantomsection\label{nav:vi-ruta-157}
% VI-CATALOG-ROW rutas 157
\pdfbookmark[2]{Celda 40}{cell-40}
\pdfbookmark[3]{r5c4\_h-1v-1}{route-157}
\RouteTitle{Ruta 157}{r5c4\_\allowbreak{}h-1v-1}
\noindent Celda 40\quad (5, 4)\qquad Orientación \(h=-1,\ v=-1\)\hfill\hyperref[sec:vi-indice-rutas]{Índice}\par\vspace{4pt}
\begingroup\fontsize{10.2}{12.5}\selectfont\setlength{\tabcolsep}{5pt}\renewcommand{\arraystretch}{1.1}\rowcolors{1}{routepale}{white}\noindent\begin{tabularx}{\linewidth}{@{}>{\raggedright\arraybackslash}p{.345\linewidth}>{\raggedright\arraybackslash}X@{}}
\RouteGroup{Identidad estructural}
\RouteRow{Clase y multisección}{neutrino\_\allowbreak{}G1\quad /\quad neutrino\_\allowbreak{}G1}
\RouteRow{Colección y sector}{neutrino; leptón}
\RouteRow{Clasificación física}{neutrino}
\RouteRow{Generación, profundidad y color}{1; G1; \textemdash{}}
\RouteRow{Ruta primaria y órbita de inversión}{\(B_1\): no\qquad 5,\allowbreak{}4;\allowbreak{}5,\allowbreak{}6}
\RouteRow{Firma de clasificación}{28|\allowbreak{}-6|\allowbreak{}4|\allowbreak{}0|\allowbreak{}-12|\allowbreak{}+++-\kern0pt{}-\kern0pt{}-+++-\kern0pt{}-\kern0pt{}-}
\RouteGroup{Retornos, memoria y transporte}
\RouteRow{Retornos con fase}{clase: 27; órbita: 27; celda: 27}
\RouteRow{Retorno cinemático}{en 108: 54; extensión: 216}
\RouteRow{Giros y cambios de hoja}{71; 107}
\RouteRow{Acarreo y diferimiento}{deuda total: 1; final: 1; bloques diferidos: 1}
\RouteRow{Tríadas finales; persistencia de Pareto}{16; sí}
\RouteRow{\(K\longrightarrow K+9\)}{repeticiones: 9; cambios de memoria: 9}
\RouteGroup{Firma, fase y diferencias}
\RouteRow{\(n_A,\ s_C,\ \kappa_0,\ \nu_{120},\ \nu_{270}\)}{28, 6, 0, 0, -12}
\RouteRow{\(\tau_{12}\); firma histórica \(\mathbb Z^5\)}{+++-\kern0pt{}-\kern0pt{}-+++-\kern0pt{}-\kern0pt{}-; 28|\allowbreak{}6|\allowbreak{}0|\allowbreak{}0|\allowbreak{}-12}
\RouteRow{\(D_1,\ D_3,\ D_4,\ D_{27},\ D_{36}\)}{66, 64, 68, 60, 40}
\RouteRow{\(\Omega_{90/120},\ P_6\)}{80, 6}
\RouteGroup{Lectores y factores de acción}
\RouteRow{\(K_D\quad /\quad K_\Omega\)}{(100,\allowbreak{}66,\allowbreak{}2)\quad /\quad (80,\allowbreak{}54,\allowbreak{}6)}
\RouteRow{\(\kappa_D\)}{99.51859\allowbreak{}71232726\allowbreak{}53414287\allowbreak{}29}
\RouteRow{\(\kappa_\Omega\)}{79.60661\allowbreak{}01397948\allowbreak{}27586470\allowbreak{}91}
\RouteRow{\(R_{\mathrm{act},D}\)}{1.000000\allowbreak{}06899443\allowbreak{}08259364\allowbreak{}17}
\RouteRow{\(R_{\mathrm{act},\Omega}\)}{1.000000\allowbreak{}05518981\allowbreak{}25318591\allowbreak{}40}
\RouteGroup{Separación de prefijos}
\RouteRow{Área de ambigüedad; área \(\log_2\)}{18; 0.0}
\RouteRow{Primer bloque de separación}{órbita: 1; celda: 1; ruta: 1}
\RouteRow{Ambigüedad final}{1}
\RouteGroup{Secuencias completas}
\RouteRow{\(w_6\)}{15 486 423 15 489 498 99 327 498 15 486 423 15 489 498 99 327 498}
\RouteRow{Tríadas estables por bloque}{0 1 2 3 4 5 6 7 8 9 10 11 12 13 14 15 16 16}
\end{tabularx}\endgroup\par\vspace{5pt}\noindent{\fontsize{8.4}{10}\selectfont\hyperref[trace-157]{Procedencia y códigos originales}\hfill Ficha completa · 157/324}
\clearpage\phantomsection\label{nav:vi-ruta-158}
% VI-CATALOG-ROW rutas 158
\pdfbookmark[3]{r5c4\_h-1v+1}{route-158}
\RouteTitle{Ruta 158}{r5c4\_\allowbreak{}h-1v+1}
\noindent Celda 40\quad (5, 4)\qquad Orientación \(h=-1,\ v=1\)\hfill\hyperref[sec:vi-indice-rutas]{Índice}\par\vspace{4pt}
\begingroup\fontsize{10.2}{12.5}\selectfont\setlength{\tabcolsep}{5pt}\renewcommand{\arraystretch}{1.1}\rowcolors{1}{routepale}{white}\noindent\begin{tabularx}{\linewidth}{@{}>{\raggedright\arraybackslash}p{.345\linewidth}>{\raggedright\arraybackslash}X@{}}
\RouteGroup{Identidad estructural}
\RouteRow{Clase y multisección}{neutrino\_\allowbreak{}G1\quad /\quad neutrino\_\allowbreak{}G1}
\RouteRow{Colección y sector}{neutrino; leptón}
\RouteRow{Clasificación física}{neutrino}
\RouteRow{Generación, profundidad y color}{1; G1; \textemdash{}}
\RouteRow{Ruta primaria y órbita de inversión}{\(B_1\): no\qquad 5,\allowbreak{}4;\allowbreak{}5,\allowbreak{}6}
\RouteRow{Firma de clasificación}{28|\allowbreak{}-6|\allowbreak{}4|\allowbreak{}0|\allowbreak{}-12|\allowbreak{}+++-\kern0pt{}-\kern0pt{}-+++-\kern0pt{}-\kern0pt{}-}
\RouteGroup{Retornos, memoria y transporte}
\RouteRow{Retornos con fase}{clase: 27; órbita: 27; celda: 27}
\RouteRow{Retorno cinemático}{en 108: 54; extensión: 216}
\RouteRow{Giros y cambios de hoja}{71; 107}
\RouteRow{Acarreo y diferimiento}{deuda total: 0; final: 0; bloques diferidos: 0}
\RouteRow{Tríadas finales; persistencia de Pareto}{17; sí}
\RouteRow{\(K\longrightarrow K+9\)}{repeticiones: 9; cambios de memoria: 9}
\RouteGroup{Firma, fase y diferencias}
\RouteRow{\(n_A,\ s_C,\ \kappa_0,\ \nu_{120},\ \nu_{270}\)}{20, 0, 0, 0, -8}
\RouteRow{\(\tau_{12}\); firma histórica \(\mathbb Z^5\)}{+0+-0-+0+-0-; 20|\allowbreak{}0|\allowbreak{}0|\allowbreak{}0|\allowbreak{}-8}
\RouteRow{\(D_1,\ D_3,\ D_4,\ D_{27},\ D_{36}\)}{24, 20, 24, 24, 16}
\RouteRow{\(\Omega_{90/120},\ P_6\)}{40, 4}
\RouteGroup{Lectores y factores de acción}
\RouteRow{\(K_D\quad /\quad K_\Omega\)}{(40,\allowbreak{}24,\allowbreak{}2)\quad /\quad (40,\allowbreak{}54,\allowbreak{}4)}
\RouteRow{\(\kappa_D\)}{39.82508\allowbreak{}59311825\allowbreak{}73056173\allowbreak{}26}
\RouteRow{\(\kappa_\Omega\)}{39.60638\allowbreak{}77469494\allowbreak{}43306362\allowbreak{}81}
\RouteRow{\(R_{\mathrm{act},D}\)}{1.000000\allowbreak{}02761000\allowbreak{}61595142\allowbreak{}15}
\RouteRow{\(R_{\mathrm{act},\Omega}\)}{1.000000\allowbreak{}02745838\allowbreak{}66926514\allowbreak{}00}
\RouteGroup{Separación de prefijos}
\RouteRow{Área de ambigüedad; área \(\log_2\)}{18; 0.0}
\RouteRow{Primer bloque de separación}{órbita: 1; celda: 1; ruta: 1}
\RouteRow{Ambigüedad final}{1}
\RouteGroup{Secuencias completas}
\RouteRow{\(w_6\)}{0 81 3 0 84 0 81 3 0 0 81 3 0 84 0 81 3 0}
\RouteRow{Tríadas estables por bloque}{0 1 2 3 4 5 6 7 8 9 10 11 12 13 14 15 16 17}
\end{tabularx}\endgroup\par\vspace{5pt}\noindent{\fontsize{8.4}{10}\selectfont\hyperref[trace-158]{Procedencia y códigos originales}\hfill Ficha completa · 158/324}
\clearpage\phantomsection\label{nav:vi-ruta-159}
% VI-CATALOG-ROW rutas 159
\pdfbookmark[3]{r5c4\_h+1v-1}{route-159}
\RouteTitle{Ruta 159}{r5c4\_\allowbreak{}h+1v-1}
\noindent Celda 40\quad (5, 4)\qquad Orientación \(h=1,\ v=-1\)\hfill\hyperref[sec:vi-indice-rutas]{Índice}\par\vspace{4pt}
\begingroup\fontsize{10.2}{12.5}\selectfont\setlength{\tabcolsep}{5pt}\renewcommand{\arraystretch}{1.1}\rowcolors{1}{routepale}{white}\noindent\begin{tabularx}{\linewidth}{@{}>{\raggedright\arraybackslash}p{.345\linewidth}>{\raggedright\arraybackslash}X@{}}
\RouteGroup{Identidad estructural}
\RouteRow{Clase y multisección}{neutrino\_\allowbreak{}G1\quad /\quad neutrino\_\allowbreak{}G1}
\RouteRow{Colección y sector}{neutrino; leptón}
\RouteRow{Clasificación física}{neutrino}
\RouteRow{Generación, profundidad y color}{1; G1; \textemdash{}}
\RouteRow{Ruta primaria y órbita de inversión}{\(B_1\): sí\qquad 5,\allowbreak{}4;\allowbreak{}5,\allowbreak{}6}
\RouteRow{Firma de clasificación}{28|\allowbreak{}-6|\allowbreak{}4|\allowbreak{}0|\allowbreak{}-12|\allowbreak{}+++-\kern0pt{}-\kern0pt{}-+++-\kern0pt{}-\kern0pt{}-}
\RouteGroup{Retornos, memoria y transporte}
\RouteRow{Retornos con fase}{clase: 27; órbita: 27; celda: 27}
\RouteRow{Retorno cinemático}{en 108: 54; extensión: 216}
\RouteRow{Giros y cambios de hoja}{71; 107}
\RouteRow{Acarreo y diferimiento}{deuda total: 2; final: 0; bloques diferidos: 2}
\RouteRow{Tríadas finales; persistencia de Pareto}{17; sí}
\RouteRow{\(K\longrightarrow K+9\)}{repeticiones: 9; cambios de memoria: 9}
\RouteGroup{Firma, fase y diferencias}
\RouteRow{\(n_A,\ s_C,\ \kappa_0,\ \nu_{120},\ \nu_{270}\)}{20, 0, 0, 0, -8}
\RouteRow{\(\tau_{12}\); firma histórica \(\mathbb Z^5\)}{+0+-0-+0+-0-; 20|\allowbreak{}0|\allowbreak{}0|\allowbreak{}0|\allowbreak{}-8}
\RouteRow{\(D_1,\ D_3,\ D_4,\ D_{27},\ D_{36}\)}{24, 20, 24, 24, 16}
\RouteRow{\(\Omega_{90/120},\ P_6\)}{40, 4}
\RouteGroup{Lectores y factores de acción}
\RouteRow{\(K_D\quad /\quad K_\Omega\)}{(40,\allowbreak{}24,\allowbreak{}2)\quad /\quad (40,\allowbreak{}54,\allowbreak{}4)}
\RouteRow{\(\kappa_D\)}{39.82508\allowbreak{}59311825\allowbreak{}73056173\allowbreak{}26}
\RouteRow{\(\kappa_\Omega\)}{39.60638\allowbreak{}77469494\allowbreak{}43306362\allowbreak{}81}
\RouteRow{\(R_{\mathrm{act},D}\)}{1.000000\allowbreak{}02761000\allowbreak{}61595142\allowbreak{}15}
\RouteRow{\(R_{\mathrm{act},\Omega}\)}{1.000000\allowbreak{}02745838\allowbreak{}66926514\allowbreak{}00}
\RouteGroup{Separación de prefijos}
\RouteRow{Área de ambigüedad; área \(\log_2\)}{18; 0.0}
\RouteRow{Primer bloque de separación}{órbita: 1; celda: 1; ruta: 1}
\RouteRow{Ambigüedad final}{1}
\RouteGroup{Secuencias completas}
\RouteRow{\(w_6\)}{81 3 0 81 0 3 0 81 3 81 3 0 81 0 3 0 81 3}
\RouteRow{Tríadas estables por bloque}{0 1 2 3 4 5 6 6 8 9 10 11 12 13 14 15 15 17}
\end{tabularx}\endgroup\par\vspace{5pt}\noindent{\fontsize{8.4}{10}\selectfont\hyperref[trace-159]{Procedencia y códigos originales}\hfill Ficha completa · 159/324}
\clearpage\phantomsection\label{nav:vi-ruta-160}
% VI-CATALOG-ROW rutas 160
\pdfbookmark[3]{r5c4\_h+1v+1}{route-160}
\RouteTitle{Ruta 160}{r5c4\_\allowbreak{}h+1v+1}
\noindent Celda 40\quad (5, 4)\qquad Orientación \(h=1,\ v=1\)\hfill\hyperref[sec:vi-indice-rutas]{Índice}\par\vspace{4pt}
\begingroup\fontsize{10.2}{12.5}\selectfont\setlength{\tabcolsep}{5pt}\renewcommand{\arraystretch}{1.1}\rowcolors{1}{routepale}{white}\noindent\begin{tabularx}{\linewidth}{@{}>{\raggedright\arraybackslash}p{.345\linewidth}>{\raggedright\arraybackslash}X@{}}
\RouteGroup{Identidad estructural}
\RouteRow{Clase y multisección}{neutrino\_\allowbreak{}G1\quad /\quad neutrino\_\allowbreak{}G1}
\RouteRow{Colección y sector}{neutrino; leptón}
\RouteRow{Clasificación física}{neutrino}
\RouteRow{Generación, profundidad y color}{1; G1; \textemdash{}}
\RouteRow{Ruta primaria y órbita de inversión}{\(B_1\): no\qquad 5,\allowbreak{}4;\allowbreak{}5,\allowbreak{}6}
\RouteRow{Firma de clasificación}{28|\allowbreak{}-6|\allowbreak{}4|\allowbreak{}0|\allowbreak{}-12|\allowbreak{}+++-\kern0pt{}-\kern0pt{}-+++-\kern0pt{}-\kern0pt{}-}
\RouteGroup{Retornos, memoria y transporte}
\RouteRow{Retornos con fase}{clase: 27; órbita: 27; celda: 27}
\RouteRow{Retorno cinemático}{en 108: 54; extensión: 216}
\RouteRow{Giros y cambios de hoja}{71; 107}
\RouteRow{Acarreo y diferimiento}{deuda total: 1; final: 0; bloques diferidos: 1}
\RouteRow{Tríadas finales; persistencia de Pareto}{17; sí}
\RouteRow{\(K\longrightarrow K+9\)}{repeticiones: 9; cambios de memoria: 9}
\RouteGroup{Firma, fase y diferencias}
\RouteRow{\(n_A,\ s_C,\ \kappa_0,\ \nu_{120},\ \nu_{270}\)}{28, -6, 0, 0, -12}
\RouteRow{\(\tau_{12}\); firma histórica \(\mathbb Z^5\)}{+++-\kern0pt{}-\kern0pt{}-+++-\kern0pt{}-\kern0pt{}-; 28|\allowbreak{}-6|\allowbreak{}0|\allowbreak{}0|\allowbreak{}-12}
\RouteRow{\(D_1,\ D_3,\ D_4,\ D_{27},\ D_{36}\)}{66, 64, 68, 60, 40}
\RouteRow{\(\Omega_{90/120},\ P_6\)}{80, 6}
\RouteGroup{Lectores y factores de acción}
\RouteRow{\(K_D\quad /\quad K_\Omega\)}{(100,\allowbreak{}66,\allowbreak{}2)\quad /\quad (80,\allowbreak{}54,\allowbreak{}6)}
\RouteRow{\(\kappa_D\)}{99.51859\allowbreak{}71232726\allowbreak{}53414287\allowbreak{}29}
\RouteRow{\(\kappa_\Omega\)}{79.60661\allowbreak{}01397948\allowbreak{}27586470\allowbreak{}91}
\RouteRow{\(R_{\mathrm{act},D}\)}{1.000000\allowbreak{}06899443\allowbreak{}08259364\allowbreak{}17}
\RouteRow{\(R_{\mathrm{act},\Omega}\)}{1.000000\allowbreak{}05518981\allowbreak{}25318591\allowbreak{}40}
\RouteGroup{Separación de prefijos}
\RouteRow{Área de ambigüedad; área \(\log_2\)}{18; 0.0}
\RouteRow{Primer bloque de separación}{órbita: 1; celda: 1; ruta: 1}
\RouteRow{Ambigüedad final}{1}
\RouteGroup{Secuencias completas}
\RouteRow{\(w_6\)}{99 327 498 99 324 423 15 486 423 99 327 498 99 324 423 15 486 423}
\RouteRow{Tríadas estables por bloque}{0 1 2 3 4 5 6 7 8 9 10 11 12 13 14 14 16 17}
\end{tabularx}\endgroup\par\vspace{5pt}\noindent{\fontsize{8.4}{10}\selectfont\hyperref[trace-160]{Procedencia y códigos originales}\hfill Ficha completa · 160/324}
\clearpage\phantomsection\label{nav:vi-ruta-161}
% VI-CATALOG-ROW rutas 161
\pdfbookmark[2]{Celda 41}{cell-41}
\pdfbookmark[3]{r5c5\_h-1v-1}{route-161}
\RouteTitle{Ruta 161}{r5c5\_\allowbreak{}h-1v-1}
\noindent Celda 41\quad (5, 5)\qquad Orientación \(h=-1,\ v=-1\)\hfill\hyperref[sec:vi-indice-rutas]{Índice}\par\vspace{4pt}
\begingroup\fontsize{10.2}{12.5}\selectfont\setlength{\tabcolsep}{5pt}\renewcommand{\arraystretch}{1.1}\rowcolors{1}{routepale}{white}\noindent\begin{tabularx}{\linewidth}{@{}>{\raggedright\arraybackslash}p{.345\linewidth}>{\raggedright\arraybackslash}X@{}}
\RouteGroup{Identidad estructural}
\RouteRow{Clase y multisección}{charged\_\allowbreak{}lepton\_\allowbreak{}G0\quad /\quad charged\_\allowbreak{}lepton\_\allowbreak{}G0}
\RouteRow{Colección y sector}{leptón cargado; leptón}
\RouteRow{Clasificación física}{leptón cargado}
\RouteRow{Generación, profundidad y color}{0; G0, centro; \textemdash{}}
\RouteRow{Ruta primaria y órbita de inversión}{\(B_1\): no\qquad 5,\allowbreak{}5}
\RouteRow{Firma de clasificación}{24|\allowbreak{}0|\allowbreak{}12|\allowbreak{}0|\allowbreak{}-12|\allowbreak{}+++-\kern0pt{}-\kern0pt{}-+++-\kern0pt{}-\kern0pt{}-}
\RouteGroup{Retornos, memoria y transporte}
\RouteRow{Retornos con fase}{clase: 27; órbita: 27; celda: 27}
\RouteRow{Retorno cinemático}{en 108: 54; extensión: 216}
\RouteRow{Giros y cambios de hoja}{71; 107}
\RouteRow{Acarreo y diferimiento}{deuda total: 3; final: 0; bloques diferidos: 2}
\RouteRow{Tríadas finales; persistencia de Pareto}{17; sí}
\RouteRow{\(K\longrightarrow K+9\)}{repeticiones: 9; cambios de memoria: 9}
\RouteGroup{Firma, fase y diferencias}
\RouteRow{\(n_A,\ s_C,\ \kappa_0,\ \nu_{120},\ \nu_{270}\)}{24, 0, 0, 0, -12}
\RouteRow{\(\tau_{12}\); firma histórica \(\mathbb Z^5\)}{+++-\kern0pt{}-\kern0pt{}-+++-\kern0pt{}-\kern0pt{}-; 24|\allowbreak{}0|\allowbreak{}0|\allowbreak{}0|\allowbreak{}-12}
\RouteRow{\(D_1,\ D_3,\ D_4,\ D_{27},\ D_{36}\)}{54, 56, 68, 48, 32}
\RouteRow{\(\Omega_{90/120},\ P_6\)}{80, 6}
\RouteGroup{Lectores y factores de acción}
\RouteRow{\(K_D\quad /\quad K_\Omega\)}{(80,\allowbreak{}54,\allowbreak{}6)\quad /\quad (80,\allowbreak{}54,\allowbreak{}6)}
\RouteRow{\(\kappa_D\)}{79.60661\allowbreak{}01397948\allowbreak{}27586470\allowbreak{}91}
\RouteRow{\(\kappa_\Omega\)}{79.60661\allowbreak{}01397948\allowbreak{}27586470\allowbreak{}91}
\RouteRow{\(R_{\mathrm{act},D}\)}{1.000000\allowbreak{}05518981\allowbreak{}25318591\allowbreak{}40}
\RouteRow{\(R_{\mathrm{act},\Omega}\)}{1.000000\allowbreak{}05518981\allowbreak{}25318591\allowbreak{}40}
\RouteGroup{Separación de prefijos}
\RouteRow{Área de ambigüedad; área \(\log_2\)}{18; 0.0}
\RouteRow{Primer bloque de separación}{órbita: 1; celda: 1; ruta: 1}
\RouteRow{Ambigüedad final}{1}
\RouteGroup{Secuencias completas}
\RouteRow{\(w_6\)}{423 15 486 423 9 24 243 657 24 423 15 486 423 9 24 243 657 24}
\RouteRow{Tríadas estables por bloque}{0 1 2 3 4 5 6 7 8 9 10 11 11 13 14 14 15 17}
\end{tabularx}\endgroup\par\vspace{5pt}\noindent{\fontsize{8.4}{10}\selectfont\hyperref[trace-161]{Procedencia y códigos originales}\hfill Ficha completa · 161/324}
\clearpage\phantomsection\label{nav:vi-ruta-162}
% VI-CATALOG-ROW rutas 162
\pdfbookmark[3]{r5c5\_h-1v+1}{route-162}
\RouteTitle{Ruta 162}{r5c5\_\allowbreak{}h-1v+1}
\noindent Celda 41\quad (5, 5)\qquad Orientación \(h=-1,\ v=1\)\hfill\hyperref[sec:vi-indice-rutas]{Índice}\par\vspace{4pt}
\begingroup\fontsize{10.2}{12.5}\selectfont\setlength{\tabcolsep}{5pt}\renewcommand{\arraystretch}{1.1}\rowcolors{1}{routepale}{white}\noindent\begin{tabularx}{\linewidth}{@{}>{\raggedright\arraybackslash}p{.345\linewidth}>{\raggedright\arraybackslash}X@{}}
\RouteGroup{Identidad estructural}
\RouteRow{Clase y multisección}{charged\_\allowbreak{}lepton\_\allowbreak{}G0\quad /\quad charged\_\allowbreak{}lepton\_\allowbreak{}G0}
\RouteRow{Colección y sector}{leptón cargado; leptón}
\RouteRow{Clasificación física}{leptón cargado}
\RouteRow{Generación, profundidad y color}{0; G0, centro; \textemdash{}}
\RouteRow{Ruta primaria y órbita de inversión}{\(B_1\): no\qquad 5,\allowbreak{}5}
\RouteRow{Firma de clasificación}{24|\allowbreak{}0|\allowbreak{}12|\allowbreak{}0|\allowbreak{}-12|\allowbreak{}+++-\kern0pt{}-\kern0pt{}-+++-\kern0pt{}-\kern0pt{}-}
\RouteGroup{Retornos, memoria y transporte}
\RouteRow{Retornos con fase}{clase: 27; órbita: 27; celda: 27}
\RouteRow{Retorno cinemático}{en 108: 54; extensión: 216}
\RouteRow{Giros y cambios de hoja}{71; 107}
\RouteRow{Acarreo y diferimiento}{deuda total: 2; final: 1; bloques diferidos: 2}
\RouteRow{Tríadas finales; persistencia de Pareto}{16; sí}
\RouteRow{\(K\longrightarrow K+9\)}{repeticiones: 9; cambios de memoria: 9}
\RouteGroup{Firma, fase y diferencias}
\RouteRow{\(n_A,\ s_C,\ \kappa_0,\ \nu_{120},\ \nu_{270}\)}{50, 4, -2, 2, -8}
\RouteRow{\(\tau_{12}\); firma histórica \(\mathbb Z^5\)}{+++-0-+++-0-; 50|\allowbreak{}4|\allowbreak{}-2|\allowbreak{}2|\allowbreak{}-8}
\RouteRow{\(D_1,\ D_3,\ D_4,\ D_{27},\ D_{36}\)}{84, 28, 84, 36, 24}
\RouteRow{\(\Omega_{90/120},\ P_6\)}{48, 5}
\RouteGroup{Lectores y factores de acción}
\RouteRow{\(K_D\quad /\quad K_\Omega\)}{(60,\allowbreak{}84,\allowbreak{}28)\quad /\quad (48,\allowbreak{}54,\allowbreak{}5)}
\RouteRow{\(\kappa_D\)}{59.39013\allowbreak{}58840155\allowbreak{}40637741\allowbreak{}48}
\RouteRow{\(\kappa_\Omega\)}{47.60649\allowbreak{}89433721\allowbreak{}35446416\allowbreak{}86}
\RouteRow{\(R_{\mathrm{act},D}\)}{1.000000\allowbreak{}04117409\allowbreak{}89467415\allowbreak{}75}
\RouteRow{\(R_{\mathrm{act},\Omega}\)}{1.000000\allowbreak{}03300471\allowbreak{}80532430\allowbreak{}84}
\RouteGroup{Separación de prefijos}
\RouteRow{Área de ambigüedad; área \(\log_2\)}{18; 0.0}
\RouteRow{Primer bloque de separación}{órbita: 1; celda: 1; ruta: 1}
\RouteRow{Ambigüedad final}{1}
\RouteGroup{Secuencias completas}
\RouteRow{\(w_6\)}{414 420 258 414 414 333 255 252 333 414 420 258 414 414 333 255 252 333}
\RouteRow{Tríadas estables por bloque}{0 1 2 3 4 5 6 7 8 9 10 11 12 13 14 14 16 16}
\end{tabularx}\endgroup\par\vspace{5pt}\noindent{\fontsize{8.4}{10}\selectfont\hyperref[trace-162]{Procedencia y códigos originales}\hfill Ficha completa · 162/324}
\clearpage\phantomsection\label{nav:vi-ruta-163}
% VI-CATALOG-ROW rutas 163
\pdfbookmark[3]{r5c5\_h+1v-1}{route-163}
\RouteTitle{Ruta 163}{r5c5\_\allowbreak{}h+1v-1}
\noindent Celda 41\quad (5, 5)\qquad Orientación \(h=1,\ v=-1\)\hfill\hyperref[sec:vi-indice-rutas]{Índice}\par\vspace{4pt}
\begingroup\fontsize{10.2}{12.5}\selectfont\setlength{\tabcolsep}{5pt}\renewcommand{\arraystretch}{1.1}\rowcolors{1}{routepale}{white}\noindent\begin{tabularx}{\linewidth}{@{}>{\raggedright\arraybackslash}p{.345\linewidth}>{\raggedright\arraybackslash}X@{}}
\RouteGroup{Identidad estructural}
\RouteRow{Clase y multisección}{charged\_\allowbreak{}lepton\_\allowbreak{}G0\quad /\quad charged\_\allowbreak{}lepton\_\allowbreak{}G0}
\RouteRow{Colección y sector}{leptón cargado; leptón}
\RouteRow{Clasificación física}{leptón cargado}
\RouteRow{Generación, profundidad y color}{0; G0, centro; \textemdash{}}
\RouteRow{Ruta primaria y órbita de inversión}{\(B_1\): no\qquad 5,\allowbreak{}5}
\RouteRow{Firma de clasificación}{24|\allowbreak{}0|\allowbreak{}12|\allowbreak{}0|\allowbreak{}-12|\allowbreak{}+++-\kern0pt{}-\kern0pt{}-+++-\kern0pt{}-\kern0pt{}-}
\RouteGroup{Retornos, memoria y transporte}
\RouteRow{Retornos con fase}{clase: 27; órbita: 27; celda: 27}
\RouteRow{Retorno cinemático}{en 108: 54; extensión: 216}
\RouteRow{Giros y cambios de hoja}{71; 107}
\RouteRow{Acarreo y diferimiento}{deuda total: 3; final: 0; bloques diferidos: 2}
\RouteRow{Tríadas finales; persistencia de Pareto}{17; sí}
\RouteRow{\(K\longrightarrow K+9\)}{repeticiones: 9; cambios de memoria: 9}
\RouteGroup{Firma, fase y diferencias}
\RouteRow{\(n_A,\ s_C,\ \kappa_0,\ \nu_{120},\ \nu_{270}\)}{50, -4, 2, -2, -8}
\RouteRow{\(\tau_{12}\); firma histórica \(\mathbb Z^5\)}{+0+-\kern0pt{}-\kern0pt{}-+0+-\kern0pt{}-\kern0pt{}-; 50|\allowbreak{}-4|\allowbreak{}2|\allowbreak{}-2|\allowbreak{}-8}
\RouteRow{\(D_1,\ D_3,\ D_4,\ D_{27},\ D_{36}\)}{84, 28, 84, 36, 24}
\RouteRow{\(\Omega_{90/120},\ P_6\)}{48, 5}
\RouteGroup{Lectores y factores de acción}
\RouteRow{\(K_D\quad /\quad K_\Omega\)}{(60,\allowbreak{}84,\allowbreak{}28)\quad /\quad (48,\allowbreak{}54,\allowbreak{}5)}
\RouteRow{\(\kappa_D\)}{59.39013\allowbreak{}58840155\allowbreak{}40637741\allowbreak{}48}
\RouteRow{\(\kappa_\Omega\)}{47.60649\allowbreak{}89433721\allowbreak{}35446416\allowbreak{}86}
\RouteRow{\(R_{\mathrm{act},D}\)}{1.000000\allowbreak{}04117409\allowbreak{}89467415\allowbreak{}75}
\RouteRow{\(R_{\mathrm{act},\Omega}\)}{1.000000\allowbreak{}03300471\allowbreak{}80532430\allowbreak{}84}
\RouteGroup{Separación de prefijos}
\RouteRow{Área de ambigüedad; área \(\log_2\)}{18; 0.0}
\RouteRow{Primer bloque de separación}{órbita: 1; celda: 1; ruta: 1}
\RouteRow{Ambigüedad final}{1}
\RouteGroup{Secuencias completas}
\RouteRow{\(w_6\)}{255 252 333 255 258 258 414 420 258 255 252 333 255 258 258 414 420 258}
\RouteRow{Tríadas estables por bloque}{0 1 2 3 4 5 6 7 8 9 10 11 11 12 14 15 15 17}
\end{tabularx}\endgroup\par\vspace{5pt}\noindent{\fontsize{8.4}{10}\selectfont\hyperref[trace-163]{Procedencia y códigos originales}\hfill Ficha completa · 163/324}
\clearpage\phantomsection\label{nav:vi-ruta-164}
% VI-CATALOG-ROW rutas 164
\pdfbookmark[3]{r5c5\_h+1v+1}{route-164}
\RouteTitle{Ruta 164}{r5c5\_\allowbreak{}h+1v+1}
\noindent Celda 41\quad (5, 5)\qquad Orientación \(h=1,\ v=1\)\hfill\hyperref[sec:vi-indice-rutas]{Índice}\par\vspace{4pt}
\begingroup\fontsize{10.2}{12.5}\selectfont\setlength{\tabcolsep}{5pt}\renewcommand{\arraystretch}{1.1}\rowcolors{1}{routepale}{white}\noindent\begin{tabularx}{\linewidth}{@{}>{\raggedright\arraybackslash}p{.345\linewidth}>{\raggedright\arraybackslash}X@{}}
\RouteGroup{Identidad estructural}
\RouteRow{Clase y multisección}{charged\_\allowbreak{}lepton\_\allowbreak{}G0\quad /\quad charged\_\allowbreak{}lepton\_\allowbreak{}G0}
\RouteRow{Colección y sector}{leptón cargado; leptón}
\RouteRow{Clasificación física}{leptón cargado}
\RouteRow{Generación, profundidad y color}{0; G0, centro; \textemdash{}}
\RouteRow{Ruta primaria y órbita de inversión}{\(B_1\): sí\qquad 5,\allowbreak{}5}
\RouteRow{Firma de clasificación}{24|\allowbreak{}0|\allowbreak{}12|\allowbreak{}0|\allowbreak{}-12|\allowbreak{}+++-\kern0pt{}-\kern0pt{}-+++-\kern0pt{}-\kern0pt{}-}
\RouteGroup{Retornos, memoria y transporte}
\RouteRow{Retornos con fase}{clase: 27; órbita: 27; celda: 27}
\RouteRow{Retorno cinemático}{en 108: 54; extensión: 216}
\RouteRow{Giros y cambios de hoja}{71; 107}
\RouteRow{Acarreo y diferimiento}{deuda total: 2; final: 1; bloques diferidos: 2}
\RouteRow{Tríadas finales; persistencia de Pareto}{16; sí}
\RouteRow{\(K\longrightarrow K+9\)}{repeticiones: 9; cambios de memoria: 9}
\RouteGroup{Firma, fase y diferencias}
\RouteRow{\(n_A,\ s_C,\ \kappa_0,\ \nu_{120},\ \nu_{270}\)}{24, 0, 0, 0, -12}
\RouteRow{\(\tau_{12}\); firma histórica \(\mathbb Z^5\)}{+++-\kern0pt{}-\kern0pt{}-+++-\kern0pt{}-\kern0pt{}-; 24|\allowbreak{}0|\allowbreak{}0|\allowbreak{}0|\allowbreak{}-12}
\RouteRow{\(D_1,\ D_3,\ D_4,\ D_{27},\ D_{36}\)}{54, 56, 68, 48, 32}
\RouteRow{\(\Omega_{90/120},\ P_6\)}{80, 6}
\RouteGroup{Lectores y factores de acción}
\RouteRow{\(K_D\quad /\quad K_\Omega\)}{(80,\allowbreak{}54,\allowbreak{}6)\quad /\quad (80,\allowbreak{}54,\allowbreak{}6)}
\RouteRow{\(\kappa_D\)}{79.60661\allowbreak{}01397948\allowbreak{}27586470\allowbreak{}91}
\RouteRow{\(\kappa_\Omega\)}{79.60661\allowbreak{}01397948\allowbreak{}27586470\allowbreak{}91}
\RouteRow{\(R_{\mathrm{act},D}\)}{1.000000\allowbreak{}05518981\allowbreak{}25318591\allowbreak{}40}
\RouteRow{\(R_{\mathrm{act},\Omega}\)}{1.000000\allowbreak{}05518981\allowbreak{}25318591\allowbreak{}40}
\RouteGroup{Separación de prefijos}
\RouteRow{Área de ambigüedad; área \(\log_2\)}{18; 0.0}
\RouteRow{Primer bloque de separación}{órbita: 1; celda: 1; ruta: 1}
\RouteRow{Ambigüedad final}{1}
\RouteGroup{Secuencias completas}
\RouteRow{\(w_6\)}{243 657 24 243 663 486 423 15 486 243 657 24 243 663 486 423 15 486}
\RouteRow{Tríadas estables por bloque}{0 1 2 3 4 5 6 7 8 9 10 11 12 13 14 14 16 16}
\end{tabularx}\endgroup\par\vspace{5pt}\noindent{\fontsize{8.4}{10}\selectfont\hyperref[trace-164]{Procedencia y códigos originales}\hfill Ficha completa · 164/324}
\clearpage\phantomsection\label{nav:vi-ruta-165}
% VI-CATALOG-ROW rutas 165
\pdfbookmark[2]{Celda 42}{cell-42}
\pdfbookmark[3]{r5c6\_h-1v-1}{route-165}
\RouteTitle{Ruta 165}{r5c6\_\allowbreak{}h-1v-1}
\noindent Celda 42\quad (5, 6)\qquad Orientación \(h=-1,\ v=-1\)\hfill\hyperref[sec:vi-indice-rutas]{Índice}\par\vspace{4pt}
\begingroup\fontsize{10.2}{12.5}\selectfont\setlength{\tabcolsep}{5pt}\renewcommand{\arraystretch}{1.1}\rowcolors{1}{routepale}{white}\noindent\begin{tabularx}{\linewidth}{@{}>{\raggedright\arraybackslash}p{.345\linewidth}>{\raggedright\arraybackslash}X@{}}
\RouteGroup{Identidad estructural}
\RouteRow{Clase y multisección}{neutrino\_\allowbreak{}G1\quad /\quad neutrino\_\allowbreak{}G1}
\RouteRow{Colección y sector}{neutrino; leptón}
\RouteRow{Clasificación física}{neutrino}
\RouteRow{Generación, profundidad y color}{1; G1; \textemdash{}}
\RouteRow{Ruta primaria y órbita de inversión}{\(B_1\): no\qquad 5,\allowbreak{}4;\allowbreak{}5,\allowbreak{}6}
\RouteRow{Firma de clasificación}{28|\allowbreak{}6|\allowbreak{}-4|\allowbreak{}0|\allowbreak{}-12|\allowbreak{}+++-\kern0pt{}-\kern0pt{}-+++-\kern0pt{}-\kern0pt{}-}
\RouteGroup{Retornos, memoria y transporte}
\RouteRow{Retornos con fase}{clase: 27; órbita: 27; celda: 27}
\RouteRow{Retorno cinemático}{en 108: 54; extensión: 216}
\RouteRow{Giros y cambios de hoja}{71; 107}
\RouteRow{Acarreo y diferimiento}{deuda total: 1; final: 0; bloques diferidos: 1}
\RouteRow{Tríadas finales; persistencia de Pareto}{17; sí}
\RouteRow{\(K\longrightarrow K+9\)}{repeticiones: 9; cambios de memoria: 9}
\RouteGroup{Firma, fase y diferencias}
\RouteRow{\(n_A,\ s_C,\ \kappa_0,\ \nu_{120},\ \nu_{270}\)}{28, -6, 0, 0, -12}
\RouteRow{\(\tau_{12}\); firma histórica \(\mathbb Z^5\)}{+++-\kern0pt{}-\kern0pt{}-+++-\kern0pt{}-\kern0pt{}-; 28|\allowbreak{}-6|\allowbreak{}0|\allowbreak{}0|\allowbreak{}-12}
\RouteRow{\(D_1,\ D_3,\ D_4,\ D_{27},\ D_{36}\)}{66, 60, 68, 48, 32}
\RouteRow{\(\Omega_{90/120},\ P_6\)}{80, 6}
\RouteGroup{Lectores y factores de acción}
\RouteRow{\(K_D\quad /\quad K_\Omega\)}{(80,\allowbreak{}66,\allowbreak{}4)\quad /\quad (80,\allowbreak{}54,\allowbreak{}6)}
\RouteRow{\(\kappa_D\)}{79.51881\allowbreak{}95161180\allowbreak{}37694395\allowbreak{}39}
\RouteRow{\(\kappa_\Omega\)}{79.60661\allowbreak{}01397948\allowbreak{}27586470\allowbreak{}91}
\RouteRow{\(R_{\mathrm{act},D}\)}{1.000000\allowbreak{}05512894\allowbreak{}88901612\allowbreak{}64}
\RouteRow{\(R_{\mathrm{act},\Omega}\)}{1.000000\allowbreak{}05518981\allowbreak{}25318591\allowbreak{}40}
\RouteGroup{Separación de prefijos}
\RouteRow{Área de ambigüedad; área \(\log_2\)}{18; 0.0}
\RouteRow{Primer bloque de separación}{órbita: 1; celda: 1; ruta: 1}
\RouteRow{Ambigüedad final}{1}
\RouteGroup{Secuencias completas}
\RouteRow{\(w_6\)}{570 264 90 570 267 243 657 24 243 570 264 90 570 267 243 657 24 243}
\RouteRow{Tríadas estables por bloque}{0 1 1 3 4 5 6 7 8 9 10 11 12 13 14 15 16 17}
\end{tabularx}\endgroup\par\vspace{5pt}\noindent{\fontsize{8.4}{10}\selectfont\hyperref[trace-165]{Procedencia y códigos originales}\hfill Ficha completa · 165/324}
\clearpage\phantomsection\label{nav:vi-ruta-166}
% VI-CATALOG-ROW rutas 166
\pdfbookmark[3]{r5c6\_h-1v+1}{route-166}
\RouteTitle{Ruta 166}{r5c6\_\allowbreak{}h-1v+1}
\noindent Celda 42\quad (5, 6)\qquad Orientación \(h=-1,\ v=1\)\hfill\hyperref[sec:vi-indice-rutas]{Índice}\par\vspace{4pt}
\begingroup\fontsize{10.2}{12.5}\selectfont\setlength{\tabcolsep}{5pt}\renewcommand{\arraystretch}{1.1}\rowcolors{1}{routepale}{white}\noindent\begin{tabularx}{\linewidth}{@{}>{\raggedright\arraybackslash}p{.345\linewidth}>{\raggedright\arraybackslash}X@{}}
\RouteGroup{Identidad estructural}
\RouteRow{Clase y multisección}{neutrino\_\allowbreak{}G1\quad /\quad neutrino\_\allowbreak{}G1}
\RouteRow{Colección y sector}{neutrino; leptón}
\RouteRow{Clasificación física}{neutrino}
\RouteRow{Generación, profundidad y color}{1; G1; \textemdash{}}
\RouteRow{Ruta primaria y órbita de inversión}{\(B_1\): sí\qquad 5,\allowbreak{}4;\allowbreak{}5,\allowbreak{}6}
\RouteRow{Firma de clasificación}{28|\allowbreak{}6|\allowbreak{}-4|\allowbreak{}0|\allowbreak{}-12|\allowbreak{}+++-\kern0pt{}-\kern0pt{}-+++-\kern0pt{}-\kern0pt{}-}
\RouteGroup{Retornos, memoria y transporte}
\RouteRow{Retornos con fase}{clase: 27; órbita: 27; celda: 27}
\RouteRow{Retorno cinemático}{en 108: 54; extensión: 216}
\RouteRow{Giros y cambios de hoja}{71; 107}
\RouteRow{Acarreo y diferimiento}{deuda total: 1; final: 0; bloques diferidos: 1}
\RouteRow{Tríadas finales; persistencia de Pareto}{17; sí}
\RouteRow{\(K\longrightarrow K+9\)}{repeticiones: 9; cambios de memoria: 9}
\RouteGroup{Firma, fase y diferencias}
\RouteRow{\(n_A,\ s_C,\ \kappa_0,\ \nu_{120},\ \nu_{270}\)}{14, -4, 0, 0, -12}
\RouteRow{\(\tau_{12}\); firma histórica \(\mathbb Z^5\)}{+++-\kern0pt{}-\kern0pt{}-+++-\kern0pt{}-\kern0pt{}-; 14|\allowbreak{}-4|\allowbreak{}0|\allowbreak{}0|\allowbreak{}-12}
\RouteRow{\(D_1,\ D_3,\ D_4,\ D_{27},\ D_{36}\)}{78, 24, 78, 24, 16}
\RouteRow{\(\Omega_{90/120},\ P_6\)}{80, 6}
\RouteGroup{Lectores y factores de acción}
\RouteRow{\(K_D\quad /\quad K_\Omega\)}{(40,\allowbreak{}78,\allowbreak{}27)\quad /\quad (80,\allowbreak{}54,\allowbreak{}6)}
\RouteRow{\(\kappa_D\)}{39.43380\allowbreak{}88030085\allowbreak{}51303671\allowbreak{}14}
\RouteRow{\(\kappa_\Omega\)}{79.60661\allowbreak{}01397948\allowbreak{}27586470\allowbreak{}91}
\RouteRow{\(R_{\mathrm{act},D}\)}{1.000000\allowbreak{}02733874\allowbreak{}08548943\allowbreak{}58}
\RouteRow{\(R_{\mathrm{act},\Omega}\)}{1.000000\allowbreak{}05518981\allowbreak{}25318591\allowbreak{}40}
\RouteGroup{Separación de prefijos}
\RouteRow{Área de ambigüedad; área \(\log_2\)}{18; 0.0}
\RouteRow{Primer bloque de separación}{órbita: 1; celda: 1; ruta: 1}
\RouteRow{Ambigüedad final}{1}
\RouteGroup{Secuencias completas}
\RouteRow{\(w_6\)}{585 507 504 585 510 666 672 510 666 585 507 504 585 510 666 672 510 666}
\RouteRow{Tríadas estables por bloque}{0 1 2 3 4 5 6 7 8 9 10 11 12 13 14 14 16 17}
\end{tabularx}\endgroup\par\vspace{5pt}\noindent{\fontsize{8.4}{10}\selectfont\hyperref[trace-166]{Procedencia y códigos originales}\hfill Ficha completa · 166/324}
\clearpage\phantomsection\label{nav:vi-ruta-167}
% VI-CATALOG-ROW rutas 167
\pdfbookmark[3]{r5c6\_h+1v-1}{route-167}
\RouteTitle{Ruta 167}{r5c6\_\allowbreak{}h+1v-1}
\noindent Celda 42\quad (5, 6)\qquad Orientación \(h=1,\ v=-1\)\hfill\hyperref[sec:vi-indice-rutas]{Índice}\par\vspace{4pt}
\begingroup\fontsize{10.2}{12.5}\selectfont\setlength{\tabcolsep}{5pt}\renewcommand{\arraystretch}{1.1}\rowcolors{1}{routepale}{white}\noindent\begin{tabularx}{\linewidth}{@{}>{\raggedright\arraybackslash}p{.345\linewidth}>{\raggedright\arraybackslash}X@{}}
\RouteGroup{Identidad estructural}
\RouteRow{Clase y multisección}{neutrino\_\allowbreak{}G1\quad /\quad neutrino\_\allowbreak{}G1}
\RouteRow{Colección y sector}{neutrino; leptón}
\RouteRow{Clasificación física}{neutrino}
\RouteRow{Generación, profundidad y color}{1; G1; \textemdash{}}
\RouteRow{Ruta primaria y órbita de inversión}{\(B_1\): no\qquad 5,\allowbreak{}4;\allowbreak{}5,\allowbreak{}6}
\RouteRow{Firma de clasificación}{28|\allowbreak{}6|\allowbreak{}-4|\allowbreak{}0|\allowbreak{}-12|\allowbreak{}+++-\kern0pt{}-\kern0pt{}-+++-\kern0pt{}-\kern0pt{}-}
\RouteGroup{Retornos, memoria y transporte}
\RouteRow{Retornos con fase}{clase: 27; órbita: 27; celda: 27}
\RouteRow{Retorno cinemático}{en 108: 54; extensión: 216}
\RouteRow{Giros y cambios de hoja}{71; 107}
\RouteRow{Acarreo y diferimiento}{deuda total: 1; final: 0; bloques diferidos: 1}
\RouteRow{Tríadas finales; persistencia de Pareto}{17; sí}
\RouteRow{\(K\longrightarrow K+9\)}{repeticiones: 9; cambios de memoria: 9}
\RouteGroup{Firma, fase y diferencias}
\RouteRow{\(n_A,\ s_C,\ \kappa_0,\ \nu_{120},\ \nu_{270}\)}{14, 4, 0, 0, -12}
\RouteRow{\(\tau_{12}\); firma histórica \(\mathbb Z^5\)}{+++-\kern0pt{}-\kern0pt{}-+++-\kern0pt{}-\kern0pt{}-; 14|\allowbreak{}4|\allowbreak{}0|\allowbreak{}0|\allowbreak{}-12}
\RouteRow{\(D_1,\ D_3,\ D_4,\ D_{27},\ D_{36}\)}{78, 24, 78, 24, 16}
\RouteRow{\(\Omega_{90/120},\ P_6\)}{80, 6}
\RouteGroup{Lectores y factores de acción}
\RouteRow{\(K_D\quad /\quad K_\Omega\)}{(40,\allowbreak{}78,\allowbreak{}27)\quad /\quad (80,\allowbreak{}54,\allowbreak{}6)}
\RouteRow{\(\kappa_D\)}{39.43380\allowbreak{}88030085\allowbreak{}51303671\allowbreak{}14}
\RouteRow{\(\kappa_\Omega\)}{79.60661\allowbreak{}01397948\allowbreak{}27586470\allowbreak{}91}
\RouteRow{\(R_{\mathrm{act},D}\)}{1.000000\allowbreak{}02733874\allowbreak{}08548943\allowbreak{}58}
\RouteRow{\(R_{\mathrm{act},\Omega}\)}{1.000000\allowbreak{}05518981\allowbreak{}25318591\allowbreak{}40}
\RouteGroup{Separación de prefijos}
\RouteRow{Área de ambigüedad; área \(\log_2\)}{18; 0.0}
\RouteRow{Primer bloque de separación}{órbita: 1; celda: 1; ruta: 1}
\RouteRow{Ambigüedad final}{1}
\RouteGroup{Secuencias completas}
\RouteRow{\(w_6\)}{672 510 666 672 507 504 585 507 504 672 510 666 672 507 504 585 507 504}
\RouteRow{Tríadas estables por bloque}{0 1 2 3 4 5 6 7 8 9 10 11 12 13 14 14 16 17}
\end{tabularx}\endgroup\par\vspace{5pt}\noindent{\fontsize{8.4}{10}\selectfont\hyperref[trace-167]{Procedencia y códigos originales}\hfill Ficha completa · 167/324}
\clearpage\phantomsection\label{nav:vi-ruta-168}
% VI-CATALOG-ROW rutas 168
\pdfbookmark[3]{r5c6\_h+1v+1}{route-168}
\RouteTitle{Ruta 168}{r5c6\_\allowbreak{}h+1v+1}
\noindent Celda 42\quad (5, 6)\qquad Orientación \(h=1,\ v=1\)\hfill\hyperref[sec:vi-indice-rutas]{Índice}\par\vspace{4pt}
\begingroup\fontsize{10.2}{12.5}\selectfont\setlength{\tabcolsep}{5pt}\renewcommand{\arraystretch}{1.1}\rowcolors{1}{routepale}{white}\noindent\begin{tabularx}{\linewidth}{@{}>{\raggedright\arraybackslash}p{.345\linewidth}>{\raggedright\arraybackslash}X@{}}
\RouteGroup{Identidad estructural}
\RouteRow{Clase y multisección}{neutrino\_\allowbreak{}G1\quad /\quad neutrino\_\allowbreak{}G1}
\RouteRow{Colección y sector}{neutrino; leptón}
\RouteRow{Clasificación física}{neutrino}
\RouteRow{Generación, profundidad y color}{1; G1; \textemdash{}}
\RouteRow{Ruta primaria y órbita de inversión}{\(B_1\): no\qquad 5,\allowbreak{}4;\allowbreak{}5,\allowbreak{}6}
\RouteRow{Firma de clasificación}{28|\allowbreak{}6|\allowbreak{}-4|\allowbreak{}0|\allowbreak{}-12|\allowbreak{}+++-\kern0pt{}-\kern0pt{}-+++-\kern0pt{}-\kern0pt{}-}
\RouteGroup{Retornos, memoria y transporte}
\RouteRow{Retornos con fase}{clase: 27; órbita: 27; celda: 27}
\RouteRow{Retorno cinemático}{en 108: 54; extensión: 216}
\RouteRow{Giros y cambios de hoja}{71; 107}
\RouteRow{Acarreo y diferimiento}{deuda total: 0; final: 0; bloques diferidos: 0}
\RouteRow{Tríadas finales; persistencia de Pareto}{17; sí}
\RouteRow{\(K\longrightarrow K+9\)}{repeticiones: 9; cambios de memoria: 9}
\RouteGroup{Firma, fase y diferencias}
\RouteRow{\(n_A,\ s_C,\ \kappa_0,\ \nu_{120},\ \nu_{270}\)}{28, 6, 0, 0, -12}
\RouteRow{\(\tau_{12}\); firma histórica \(\mathbb Z^5\)}{+++-\kern0pt{}-\kern0pt{}-+++-\kern0pt{}-\kern0pt{}-; 28|\allowbreak{}6|\allowbreak{}0|\allowbreak{}0|\allowbreak{}-12}
\RouteRow{\(D_1,\ D_3,\ D_4,\ D_{27},\ D_{36}\)}{66, 60, 68, 48, 32}
\RouteRow{\(\Omega_{90/120},\ P_6\)}{80, 6}
\RouteGroup{Lectores y factores de acción}
\RouteRow{\(K_D\quad /\quad K_\Omega\)}{(80,\allowbreak{}66,\allowbreak{}4)\quad /\quad (80,\allowbreak{}54,\allowbreak{}6)}
\RouteRow{\(\kappa_D\)}{79.51881\allowbreak{}95161180\allowbreak{}37694395\allowbreak{}39}
\RouteRow{\(\kappa_\Omega\)}{79.60661\allowbreak{}01397948\allowbreak{}27586470\allowbreak{}91}
\RouteRow{\(R_{\mathrm{act},D}\)}{1.000000\allowbreak{}05512894\allowbreak{}88901612\allowbreak{}64}
\RouteRow{\(R_{\mathrm{act},\Omega}\)}{1.000000\allowbreak{}05518981\allowbreak{}25318591\allowbreak{}40}
\RouteGroup{Separación de prefijos}
\RouteRow{Área de ambigüedad; área \(\log_2\)}{18; 0.0}
\RouteRow{Primer bloque de separación}{órbita: 1; celda: 1; ruta: 1}
\RouteRow{Ambigüedad final}{1}
\RouteGroup{Secuencias completas}
\RouteRow{\(w_6\)}{657 24 243 657 21 90 570 264 90 657 24 243 657 21 90 570 264 90}
\RouteRow{Tríadas estables por bloque}{0 1 2 3 4 5 6 7 8 9 10 11 12 13 14 15 16 17}
\end{tabularx}\endgroup\par\vspace{5pt}\noindent{\fontsize{8.4}{10}\selectfont\hyperref[trace-168]{Procedencia y códigos originales}\hfill Ficha completa · 168/324}
\clearpage\phantomsection\label{nav:vi-ruta-169}
% VI-CATALOG-ROW rutas 169
\pdfbookmark[2]{Celda 43}{cell-43}
\pdfbookmark[3]{r5c7\_h-1v-1}{route-169}
\RouteTitle{Ruta 169}{r5c7\_\allowbreak{}h-1v-1}
\noindent Celda 43\quad (5, 7)\qquad Orientación \(h=-1,\ v=-1\)\hfill\hyperref[sec:vi-indice-rutas]{Índice}\par\vspace{4pt}
\begingroup\fontsize{10.2}{12.5}\selectfont\setlength{\tabcolsep}{5pt}\renewcommand{\arraystretch}{1.1}\rowcolors{1}{routepale}{white}\noindent\begin{tabularx}{\linewidth}{@{}>{\raggedright\arraybackslash}p{.345\linewidth}>{\raggedright\arraybackslash}X@{}}
\RouteGroup{Identidad estructural}
\RouteRow{Clase y multisección}{charged\_\allowbreak{}lepton\_\allowbreak{}G2\quad /\quad charged\_\allowbreak{}lepton\_\allowbreak{}G2}
\RouteRow{Colección y sector}{leptón cargado; leptón}
\RouteRow{Clasificación física}{leptón cargado}
\RouteRow{Generación, profundidad y color}{2; G2; \textemdash{}}
\RouteRow{Ruta primaria y órbita de inversión}{\(B_1\): no\qquad 5,\allowbreak{}3;\allowbreak{}5,\allowbreak{}7}
\RouteRow{Firma de clasificación}{28|\allowbreak{}-6|\allowbreak{}6|\allowbreak{}0|\allowbreak{}-12|\allowbreak{}+++-\kern0pt{}-\kern0pt{}-+++-\kern0pt{}-\kern0pt{}-}
\RouteGroup{Retornos, memoria y transporte}
\RouteRow{Retornos con fase}{clase: 27; órbita: 27; celda: 27}
\RouteRow{Retorno cinemático}{en 108: 54; extensión: 216}
\RouteRow{Giros y cambios de hoja}{71; 107}
\RouteRow{Acarreo y diferimiento}{deuda total: 1; final: 1; bloques diferidos: 1}
\RouteRow{Tríadas finales; persistencia de Pareto}{16; no}
\RouteRow{\(K\longrightarrow K+9\)}{repeticiones: 9; cambios de memoria: 9}
\RouteGroup{Firma, fase y diferencias}
\RouteRow{\(n_A,\ s_C,\ \kappa_0,\ \nu_{120},\ \nu_{270}\)}{28, 6, 0, 4, -12}
\RouteRow{\(\tau_{12}\); firma histórica \(\mathbb Z^5\)}{+++-\kern0pt{}-\kern0pt{}-+++-\kern0pt{}-\kern0pt{}-; 28|\allowbreak{}6|\allowbreak{}0|\allowbreak{}4|\allowbreak{}-12}
\RouteRow{\(D_1,\ D_3,\ D_4,\ D_{27},\ D_{36}\)}{66, 64, 68, 60, 40}
\RouteRow{\(\Omega_{90/120},\ P_6\)}{80, 6}
\RouteGroup{Lectores y factores de acción}
\RouteRow{\(K_D\quad /\quad K_\Omega\)}{(100,\allowbreak{}66,\allowbreak{}2)\quad /\quad (80,\allowbreak{}54,\allowbreak{}6)}
\RouteRow{\(\kappa_D\)}{99.51859\allowbreak{}71232726\allowbreak{}53414287\allowbreak{}29}
\RouteRow{\(\kappa_\Omega\)}{79.60661\allowbreak{}01397948\allowbreak{}27586470\allowbreak{}91}
\RouteRow{\(R_{\mathrm{act},D}\)}{1.000000\allowbreak{}06899443\allowbreak{}08259364\allowbreak{}17}
\RouteRow{\(R_{\mathrm{act},\Omega}\)}{1.000000\allowbreak{}05518981\allowbreak{}25318591\allowbreak{}40}
\RouteGroup{Separación de prefijos}
\RouteRow{Área de ambigüedad; área \(\log_2\)}{22; 4.0}
\RouteRow{Primer bloque de separación}{órbita: 5; celda: 5; ruta: 5}
\RouteRow{Ambigüedad final}{1}
\RouteGroup{Secuencias completas}
\RouteRow{\(w_6\)}{15 486 423 15 489 498 99 327 498 15 486 423 15 489 498 99 327 498}
\RouteRow{Tríadas estables por bloque}{0 1 2 3 4 5 6 7 8 9 10 11 12 13 14 15 16 16}
\end{tabularx}\endgroup\par\vspace{5pt}\noindent{\fontsize{8.4}{10}\selectfont\hyperref[trace-169]{Procedencia y códigos originales}\hfill Ficha completa · 169/324}
\clearpage\phantomsection\label{nav:vi-ruta-170}
% VI-CATALOG-ROW rutas 170
\pdfbookmark[3]{r5c7\_h-1v+1}{route-170}
\RouteTitle{Ruta 170}{r5c7\_\allowbreak{}h-1v+1}
\noindent Celda 43\quad (5, 7)\qquad Orientación \(h=-1,\ v=1\)\hfill\hyperref[sec:vi-indice-rutas]{Índice}\par\vspace{4pt}
\begingroup\fontsize{10.2}{12.5}\selectfont\setlength{\tabcolsep}{5pt}\renewcommand{\arraystretch}{1.1}\rowcolors{1}{routepale}{white}\noindent\begin{tabularx}{\linewidth}{@{}>{\raggedright\arraybackslash}p{.345\linewidth}>{\raggedright\arraybackslash}X@{}}
\RouteGroup{Identidad estructural}
\RouteRow{Clase y multisección}{charged\_\allowbreak{}lepton\_\allowbreak{}G2\quad /\quad charged\_\allowbreak{}lepton\_\allowbreak{}G2}
\RouteRow{Colección y sector}{leptón cargado; leptón}
\RouteRow{Clasificación física}{leptón cargado}
\RouteRow{Generación, profundidad y color}{2; G2; \textemdash{}}
\RouteRow{Ruta primaria y órbita de inversión}{\(B_1\): no\qquad 5,\allowbreak{}3;\allowbreak{}5,\allowbreak{}7}
\RouteRow{Firma de clasificación}{28|\allowbreak{}-6|\allowbreak{}6|\allowbreak{}0|\allowbreak{}-12|\allowbreak{}+++-\kern0pt{}-\kern0pt{}-+++-\kern0pt{}-\kern0pt{}-}
\RouteGroup{Retornos, memoria y transporte}
\RouteRow{Retornos con fase}{clase: 27; órbita: 27; celda: 27}
\RouteRow{Retorno cinemático}{en 108: 54; extensión: 216}
\RouteRow{Giros y cambios de hoja}{71; 107}
\RouteRow{Acarreo y diferimiento}{deuda total: 0; final: 0; bloques diferidos: 0}
\RouteRow{Tríadas finales; persistencia de Pareto}{17; no}
\RouteRow{\(K\longrightarrow K+9\)}{repeticiones: 9; cambios de memoria: 9}
\RouteGroup{Firma, fase y diferencias}
\RouteRow{\(n_A,\ s_C,\ \kappa_0,\ \nu_{120},\ \nu_{270}\)}{50, -6, 2, -2, -8}
\RouteRow{\(\tau_{12}\); firma histórica \(\mathbb Z^5\)}{0++-\kern0pt{}-\kern0pt{}-0++-\kern0pt{}-\kern0pt{}-; 50|\allowbreak{}-6|\allowbreak{}2|\allowbreak{}-2|\allowbreak{}-8}
\RouteRow{\(D_1,\ D_3,\ D_4,\ D_{27},\ D_{36}\)}{24, 20, 24, 24, 16}
\RouteRow{\(\Omega_{90/120},\ P_6\)}{56, 5}
\RouteGroup{Lectores y factores de acción}
\RouteRow{\(K_D\quad /\quad K_\Omega\)}{(40,\allowbreak{}24,\allowbreak{}2)\quad /\quad (56,\allowbreak{}54,\allowbreak{}5)}
\RouteRow{\(\kappa_D\)}{39.82508\allowbreak{}59311825\allowbreak{}73056173\allowbreak{}26}
\RouteRow{\(\kappa_\Omega\)}{55.60649\allowbreak{}89433721\allowbreak{}35446416\allowbreak{}86}
\RouteRow{\(R_{\mathrm{act},D}\)}{1.000000\allowbreak{}02761000\allowbreak{}61595142\allowbreak{}15}
\RouteRow{\(R_{\mathrm{act},\Omega}\)}{1.000000\allowbreak{}03855097\allowbreak{}23541416\allowbreak{}45}
\RouteGroup{Separación de prefijos}
\RouteRow{Área de ambigüedad; área \(\log_2\)}{36; 18.0}
\RouteRow{Primer bloque de separación}{órbita: \textemdash{}; celda: \textemdash{}; ruta: \textemdash{}}
\RouteRow{Ambigüedad final}{2}
\RouteGroup{Secuencias completas}
\RouteRow{\(w_6\)}{0 81 3 0 84 0 81 3 0 0 81 3 0 84 0 81 3 0}
\RouteRow{Tríadas estables por bloque}{0 1 2 3 4 5 6 7 8 9 10 11 12 13 14 15 16 17}
\end{tabularx}\endgroup\par\vspace{5pt}\noindent{\fontsize{8.4}{10}\selectfont\hyperref[trace-170]{Procedencia y códigos originales}\hfill Ficha completa · 170/324}
\clearpage\phantomsection\label{nav:vi-ruta-171}
% VI-CATALOG-ROW rutas 171
\pdfbookmark[3]{r5c7\_h+1v-1}{route-171}
\RouteTitle{Ruta 171}{r5c7\_\allowbreak{}h+1v-1}
\noindent Celda 43\quad (5, 7)\qquad Orientación \(h=1,\ v=-1\)\hfill\hyperref[sec:vi-indice-rutas]{Índice}\par\vspace{4pt}
\begingroup\fontsize{10.2}{12.5}\selectfont\setlength{\tabcolsep}{5pt}\renewcommand{\arraystretch}{1.1}\rowcolors{1}{routepale}{white}\noindent\begin{tabularx}{\linewidth}{@{}>{\raggedright\arraybackslash}p{.345\linewidth}>{\raggedright\arraybackslash}X@{}}
\RouteGroup{Identidad estructural}
\RouteRow{Clase y multisección}{charged\_\allowbreak{}lepton\_\allowbreak{}G2\quad /\quad charged\_\allowbreak{}lepton\_\allowbreak{}G2}
\RouteRow{Colección y sector}{leptón cargado; leptón}
\RouteRow{Clasificación física}{leptón cargado}
\RouteRow{Generación, profundidad y color}{2; G2; \textemdash{}}
\RouteRow{Ruta primaria y órbita de inversión}{\(B_1\): no\qquad 5,\allowbreak{}3;\allowbreak{}5,\allowbreak{}7}
\RouteRow{Firma de clasificación}{28|\allowbreak{}-6|\allowbreak{}6|\allowbreak{}0|\allowbreak{}-12|\allowbreak{}+++-\kern0pt{}-\kern0pt{}-+++-\kern0pt{}-\kern0pt{}-}
\RouteGroup{Retornos, memoria y transporte}
\RouteRow{Retornos con fase}{clase: 27; órbita: 27; celda: 27}
\RouteRow{Retorno cinemático}{en 108: 54; extensión: 216}
\RouteRow{Giros y cambios de hoja}{71; 107}
\RouteRow{Acarreo y diferimiento}{deuda total: 2; final: 0; bloques diferidos: 2}
\RouteRow{Tríadas finales; persistencia de Pareto}{17; no}
\RouteRow{\(K\longrightarrow K+9\)}{repeticiones: 9; cambios de memoria: 9}
\RouteGroup{Firma, fase y diferencias}
\RouteRow{\(n_A,\ s_C,\ \kappa_0,\ \nu_{120},\ \nu_{270}\)}{50, 6, -2, 2, -8}
\RouteRow{\(\tau_{12}\); firma histórica \(\mathbb Z^5\)}{+++0-\kern0pt{}-+++0-\kern0pt{}-; 50|\allowbreak{}6|\allowbreak{}-2|\allowbreak{}2|\allowbreak{}-8}
\RouteRow{\(D_1,\ D_3,\ D_4,\ D_{27},\ D_{36}\)}{24, 20, 24, 24, 16}
\RouteRow{\(\Omega_{90/120},\ P_6\)}{56, 5}
\RouteGroup{Lectores y factores de acción}
\RouteRow{\(K_D\quad /\quad K_\Omega\)}{(40,\allowbreak{}24,\allowbreak{}2)\quad /\quad (56,\allowbreak{}54,\allowbreak{}5)}
\RouteRow{\(\kappa_D\)}{39.82508\allowbreak{}59311825\allowbreak{}73056173\allowbreak{}26}
\RouteRow{\(\kappa_\Omega\)}{55.60649\allowbreak{}89433721\allowbreak{}35446416\allowbreak{}86}
\RouteRow{\(R_{\mathrm{act},D}\)}{1.000000\allowbreak{}02761000\allowbreak{}61595142\allowbreak{}15}
\RouteRow{\(R_{\mathrm{act},\Omega}\)}{1.000000\allowbreak{}03855097\allowbreak{}23541416\allowbreak{}45}
\RouteGroup{Separación de prefijos}
\RouteRow{Área de ambigüedad; área \(\log_2\)}{36; 18.0}
\RouteRow{Primer bloque de separación}{órbita: \textemdash{}; celda: \textemdash{}; ruta: \textemdash{}}
\RouteRow{Ambigüedad final}{2}
\RouteGroup{Secuencias completas}
\RouteRow{\(w_6\)}{81 3 0 81 0 3 0 81 3 81 3 0 81 0 3 0 81 3}
\RouteRow{Tríadas estables por bloque}{0 1 2 3 4 5 6 6 8 9 10 11 12 13 14 15 15 17}
\end{tabularx}\endgroup\par\vspace{5pt}\noindent{\fontsize{8.4}{10}\selectfont\hyperref[trace-171]{Procedencia y códigos originales}\hfill Ficha completa · 171/324}
\clearpage\phantomsection\label{nav:vi-ruta-172}
% VI-CATALOG-ROW rutas 172
\pdfbookmark[3]{r5c7\_h+1v+1}{route-172}
\RouteTitle{Ruta 172}{r5c7\_\allowbreak{}h+1v+1}
\noindent Celda 43\quad (5, 7)\qquad Orientación \(h=1,\ v=1\)\hfill\hyperref[sec:vi-indice-rutas]{Índice}\par\vspace{4pt}
\begingroup\fontsize{10.2}{12.5}\selectfont\setlength{\tabcolsep}{5pt}\renewcommand{\arraystretch}{1.1}\rowcolors{1}{routepale}{white}\noindent\begin{tabularx}{\linewidth}{@{}>{\raggedright\arraybackslash}p{.345\linewidth}>{\raggedright\arraybackslash}X@{}}
\RouteGroup{Identidad estructural}
\RouteRow{Clase y multisección}{charged\_\allowbreak{}lepton\_\allowbreak{}G2\quad /\quad charged\_\allowbreak{}lepton\_\allowbreak{}G2}
\RouteRow{Colección y sector}{leptón cargado; leptón}
\RouteRow{Clasificación física}{leptón cargado}
\RouteRow{Generación, profundidad y color}{2; G2; \textemdash{}}
\RouteRow{Ruta primaria y órbita de inversión}{\(B_1\): sí\qquad 5,\allowbreak{}3;\allowbreak{}5,\allowbreak{}7}
\RouteRow{Firma de clasificación}{28|\allowbreak{}-6|\allowbreak{}6|\allowbreak{}0|\allowbreak{}-12|\allowbreak{}+++-\kern0pt{}-\kern0pt{}-+++-\kern0pt{}-\kern0pt{}-}
\RouteGroup{Retornos, memoria y transporte}
\RouteRow{Retornos con fase}{clase: 27; órbita: 27; celda: 27}
\RouteRow{Retorno cinemático}{en 108: 54; extensión: 216}
\RouteRow{Giros y cambios de hoja}{71; 107}
\RouteRow{Acarreo y diferimiento}{deuda total: 1; final: 0; bloques diferidos: 1}
\RouteRow{Tríadas finales; persistencia de Pareto}{17; no}
\RouteRow{\(K\longrightarrow K+9\)}{repeticiones: 9; cambios de memoria: 9}
\RouteGroup{Firma, fase y diferencias}
\RouteRow{\(n_A,\ s_C,\ \kappa_0,\ \nu_{120},\ \nu_{270}\)}{28, -6, 0, -4, -12}
\RouteRow{\(\tau_{12}\); firma histórica \(\mathbb Z^5\)}{+++-\kern0pt{}-\kern0pt{}-+++-\kern0pt{}-\kern0pt{}-; 28|\allowbreak{}-6|\allowbreak{}0|\allowbreak{}-4|\allowbreak{}-12}
\RouteRow{\(D_1,\ D_3,\ D_4,\ D_{27},\ D_{36}\)}{66, 64, 68, 60, 40}
\RouteRow{\(\Omega_{90/120},\ P_6\)}{80, 6}
\RouteGroup{Lectores y factores de acción}
\RouteRow{\(K_D\quad /\quad K_\Omega\)}{(100,\allowbreak{}66,\allowbreak{}2)\quad /\quad (80,\allowbreak{}54,\allowbreak{}6)}
\RouteRow{\(\kappa_D\)}{99.51859\allowbreak{}71232726\allowbreak{}53414287\allowbreak{}29}
\RouteRow{\(\kappa_\Omega\)}{79.60661\allowbreak{}01397948\allowbreak{}27586470\allowbreak{}91}
\RouteRow{\(R_{\mathrm{act},D}\)}{1.000000\allowbreak{}06899443\allowbreak{}08259364\allowbreak{}17}
\RouteRow{\(R_{\mathrm{act},\Omega}\)}{1.000000\allowbreak{}05518981\allowbreak{}25318591\allowbreak{}40}
\RouteGroup{Separación de prefijos}
\RouteRow{Área de ambigüedad; área \(\log_2\)}{18; 0.0}
\RouteRow{Primer bloque de separación}{órbita: 1; celda: 1; ruta: 1}
\RouteRow{Ambigüedad final}{1}
\RouteGroup{Secuencias completas}
\RouteRow{\(w_6\)}{99 327 498 99 324 423 15 486 423 99 327 498 99 324 423 15 486 423}
\RouteRow{Tríadas estables por bloque}{0 1 2 3 4 5 6 7 8 9 10 11 12 13 14 14 16 17}
\end{tabularx}\endgroup\par\vspace{5pt}\noindent{\fontsize{8.4}{10}\selectfont\hyperref[trace-172]{Procedencia y códigos originales}\hfill Ficha completa · 172/324}
\clearpage\phantomsection\label{nav:vi-ruta-173}
% VI-CATALOG-ROW rutas 173
\pdfbookmark[2]{Celda 44}{cell-44}
\pdfbookmark[3]{r5c8\_h-1v-1}{route-173}
\RouteTitle{Ruta 173}{r5c8\_\allowbreak{}h-1v-1}
\noindent Celda 44\quad (5, 8)\qquad Orientación \(h=-1,\ v=-1\)\hfill\hyperref[sec:vi-indice-rutas]{Índice}\par\vspace{4pt}
\begingroup\fontsize{10.2}{12.5}\selectfont\setlength{\tabcolsep}{5pt}\renewcommand{\arraystretch}{1.1}\rowcolors{1}{routepale}{white}\noindent\begin{tabularx}{\linewidth}{@{}>{\raggedright\arraybackslash}p{.345\linewidth}>{\raggedright\arraybackslash}X@{}}
\RouteGroup{Identidad estructural}
\RouteRow{Clase y multisección}{neutrino\_\allowbreak{}G3\quad /\quad neutrino\_\allowbreak{}G3}
\RouteRow{Colección y sector}{neutrino; leptón}
\RouteRow{Clasificación física}{neutrino}
\RouteRow{Generación, profundidad y color}{3; G3; \textemdash{}}
\RouteRow{Ruta primaria y órbita de inversión}{\(B_1\): no\qquad 5,\allowbreak{}2;\allowbreak{}5,\allowbreak{}8}
\RouteRow{Firma de clasificación}{30|\allowbreak{}-6|\allowbreak{}8|\allowbreak{}0|\allowbreak{}-12|\allowbreak{}+++-\kern0pt{}-\kern0pt{}-+++-\kern0pt{}-\kern0pt{}-}
\RouteGroup{Retornos, memoria y transporte}
\RouteRow{Retornos con fase}{clase: 27; órbita: 27; celda: 27}
\RouteRow{Retorno cinemático}{en 108: 54; extensión: 216}
\RouteRow{Giros y cambios de hoja}{71; 107}
\RouteRow{Acarreo y diferimiento}{deuda total: 3; final: 0; bloques diferidos: 2}
\RouteRow{Tríadas finales; persistencia de Pareto}{17; sí}
\RouteRow{\(K\longrightarrow K+9\)}{repeticiones: 9; cambios de memoria: 9}
\RouteGroup{Firma, fase y diferencias}
\RouteRow{\(n_A,\ s_C,\ \kappa_0,\ \nu_{120},\ \nu_{270}\)}{30, 6, 0, 0, -12}
\RouteRow{\(\tau_{12}\); firma histórica \(\mathbb Z^5\)}{+++-\kern0pt{}-\kern0pt{}-+++-\kern0pt{}-\kern0pt{}-; 30|\allowbreak{}6|\allowbreak{}0|\allowbreak{}0|\allowbreak{}-12}
\RouteRow{\(D_1,\ D_3,\ D_4,\ D_{27},\ D_{36}\)}{54, 56, 68, 48, 32}
\RouteRow{\(\Omega_{90/120},\ P_6\)}{80, 6}
\RouteGroup{Lectores y factores de acción}
\RouteRow{\(K_D\quad /\quad K_\Omega\)}{(80,\allowbreak{}54,\allowbreak{}6)\quad /\quad (80,\allowbreak{}54,\allowbreak{}6)}
\RouteRow{\(\kappa_D\)}{79.60661\allowbreak{}01397948\allowbreak{}27586470\allowbreak{}91}
\RouteRow{\(\kappa_\Omega\)}{79.60661\allowbreak{}01397948\allowbreak{}27586470\allowbreak{}91}
\RouteRow{\(R_{\mathrm{act},D}\)}{1.000000\allowbreak{}05518981\allowbreak{}25318591\allowbreak{}40}
\RouteRow{\(R_{\mathrm{act},\Omega}\)}{1.000000\allowbreak{}05518981\allowbreak{}25318591\allowbreak{}40}
\RouteGroup{Separación de prefijos}
\RouteRow{Área de ambigüedad; área \(\log_2\)}{36; 18.0}
\RouteRow{Primer bloque de separación}{órbita: 1; celda: \textemdash{}; ruta: \textemdash{}}
\RouteRow{Ambigüedad final}{2}
\RouteGroup{Secuencias completas}
\RouteRow{\(w_6\)}{423 15 486 423 9 24 243 657 24 423 15 486 423 9 24 243 657 24}
\RouteRow{Tríadas estables por bloque}{0 1 2 3 4 5 6 7 8 9 10 11 11 13 14 14 15 17}
\end{tabularx}\endgroup\par\vspace{5pt}\noindent{\fontsize{8.4}{10}\selectfont\hyperref[trace-173]{Procedencia y códigos originales}\hfill Ficha completa · 173/324}
\clearpage\phantomsection\label{nav:vi-ruta-174}
% VI-CATALOG-ROW rutas 174
\pdfbookmark[3]{r5c8\_h-1v+1}{route-174}
\RouteTitle{Ruta 174}{r5c8\_\allowbreak{}h-1v+1}
\noindent Celda 44\quad (5, 8)\qquad Orientación \(h=-1,\ v=1\)\hfill\hyperref[sec:vi-indice-rutas]{Índice}\par\vspace{4pt}
\begingroup\fontsize{10.2}{12.5}\selectfont\setlength{\tabcolsep}{5pt}\renewcommand{\arraystretch}{1.1}\rowcolors{1}{routepale}{white}\noindent\begin{tabularx}{\linewidth}{@{}>{\raggedright\arraybackslash}p{.345\linewidth}>{\raggedright\arraybackslash}X@{}}
\RouteGroup{Identidad estructural}
\RouteRow{Clase y multisección}{neutrino\_\allowbreak{}G3\quad /\quad neutrino\_\allowbreak{}G3}
\RouteRow{Colección y sector}{neutrino; leptón}
\RouteRow{Clasificación física}{neutrino}
\RouteRow{Generación, profundidad y color}{3; G3; \textemdash{}}
\RouteRow{Ruta primaria y órbita de inversión}{\(B_1\): sí\qquad 5,\allowbreak{}2;\allowbreak{}5,\allowbreak{}8}
\RouteRow{Firma de clasificación}{30|\allowbreak{}-6|\allowbreak{}8|\allowbreak{}0|\allowbreak{}-12|\allowbreak{}+++-\kern0pt{}-\kern0pt{}-+++-\kern0pt{}-\kern0pt{}-}
\RouteGroup{Retornos, memoria y transporte}
\RouteRow{Retornos con fase}{clase: 27; órbita: 27; celda: 27}
\RouteRow{Retorno cinemático}{en 108: 54; extensión: 216}
\RouteRow{Giros y cambios de hoja}{71; 107}
\RouteRow{Acarreo y diferimiento}{deuda total: 2; final: 1; bloques diferidos: 2}
\RouteRow{Tríadas finales; persistencia de Pareto}{16; sí}
\RouteRow{\(K\longrightarrow K+9\)}{repeticiones: 9; cambios de memoria: 9}
\RouteGroup{Firma, fase y diferencias}
\RouteRow{\(n_A,\ s_C,\ \kappa_0,\ \nu_{120},\ \nu_{270}\)}{8, -2, 0, 0, -12}
\RouteRow{\(\tau_{12}\); firma histórica \(\mathbb Z^5\)}{+++-\kern0pt{}-\kern0pt{}-+++-\kern0pt{}-\kern0pt{}-; 8|\allowbreak{}-2|\allowbreak{}0|\allowbreak{}0|\allowbreak{}-12}
\RouteRow{\(D_1,\ D_3,\ D_4,\ D_{27},\ D_{36}\)}{84, 28, 84, 36, 24}
\RouteRow{\(\Omega_{90/120},\ P_6\)}{80, 6}
\RouteGroup{Lectores y factores de acción}
\RouteRow{\(K_D\quad /\quad K_\Omega\)}{(60,\allowbreak{}84,\allowbreak{}28)\quad /\quad (80,\allowbreak{}54,\allowbreak{}6)}
\RouteRow{\(\kappa_D\)}{59.39013\allowbreak{}58840155\allowbreak{}40637741\allowbreak{}48}
\RouteRow{\(\kappa_\Omega\)}{79.60661\allowbreak{}01397948\allowbreak{}27586470\allowbreak{}91}
\RouteRow{\(R_{\mathrm{act},D}\)}{1.000000\allowbreak{}04117409\allowbreak{}89467415\allowbreak{}75}
\RouteRow{\(R_{\mathrm{act},\Omega}\)}{1.000000\allowbreak{}05518981\allowbreak{}25318591\allowbreak{}40}
\RouteGroup{Separación de prefijos}
\RouteRow{Área de ambigüedad; área \(\log_2\)}{36; 18.0}
\RouteRow{Primer bloque de separación}{órbita: 1; celda: \textemdash{}; ruta: \textemdash{}}
\RouteRow{Ambigüedad final}{2}
\RouteGroup{Secuencias completas}
\RouteRow{\(w_6\)}{414 420 258 414 414 333 255 252 333 414 420 258 414 414 333 255 252 333}
\RouteRow{Tríadas estables por bloque}{0 1 2 3 4 5 6 7 8 9 10 11 12 13 14 14 16 16}
\end{tabularx}\endgroup\par\vspace{5pt}\noindent{\fontsize{8.4}{10}\selectfont\hyperref[trace-174]{Procedencia y códigos originales}\hfill Ficha completa · 174/324}
\clearpage\phantomsection\label{nav:vi-ruta-175}
% VI-CATALOG-ROW rutas 175
\pdfbookmark[3]{r5c8\_h+1v-1}{route-175}
\RouteTitle{Ruta 175}{r5c8\_\allowbreak{}h+1v-1}
\noindent Celda 44\quad (5, 8)\qquad Orientación \(h=1,\ v=-1\)\hfill\hyperref[sec:vi-indice-rutas]{Índice}\par\vspace{4pt}
\begingroup\fontsize{10.2}{12.5}\selectfont\setlength{\tabcolsep}{5pt}\renewcommand{\arraystretch}{1.1}\rowcolors{1}{routepale}{white}\noindent\begin{tabularx}{\linewidth}{@{}>{\raggedright\arraybackslash}p{.345\linewidth}>{\raggedright\arraybackslash}X@{}}
\RouteGroup{Identidad estructural}
\RouteRow{Clase y multisección}{neutrino\_\allowbreak{}G3\quad /\quad neutrino\_\allowbreak{}G3}
\RouteRow{Colección y sector}{neutrino; leptón}
\RouteRow{Clasificación física}{neutrino}
\RouteRow{Generación, profundidad y color}{3; G3; \textemdash{}}
\RouteRow{Ruta primaria y órbita de inversión}{\(B_1\): no\qquad 5,\allowbreak{}2;\allowbreak{}5,\allowbreak{}8}
\RouteRow{Firma de clasificación}{30|\allowbreak{}-6|\allowbreak{}8|\allowbreak{}0|\allowbreak{}-12|\allowbreak{}+++-\kern0pt{}-\kern0pt{}-+++-\kern0pt{}-\kern0pt{}-}
\RouteGroup{Retornos, memoria y transporte}
\RouteRow{Retornos con fase}{clase: 27; órbita: 27; celda: 27}
\RouteRow{Retorno cinemático}{en 108: 54; extensión: 216}
\RouteRow{Giros y cambios de hoja}{71; 107}
\RouteRow{Acarreo y diferimiento}{deuda total: 3; final: 0; bloques diferidos: 2}
\RouteRow{Tríadas finales; persistencia de Pareto}{17; sí}
\RouteRow{\(K\longrightarrow K+9\)}{repeticiones: 9; cambios de memoria: 9}
\RouteGroup{Firma, fase y diferencias}
\RouteRow{\(n_A,\ s_C,\ \kappa_0,\ \nu_{120},\ \nu_{270}\)}{8, 2, 0, 0, -12}
\RouteRow{\(\tau_{12}\); firma histórica \(\mathbb Z^5\)}{+++-\kern0pt{}-\kern0pt{}-+++-\kern0pt{}-\kern0pt{}-; 8|\allowbreak{}2|\allowbreak{}0|\allowbreak{}0|\allowbreak{}-12}
\RouteRow{\(D_1,\ D_3,\ D_4,\ D_{27},\ D_{36}\)}{84, 28, 84, 36, 24}
\RouteRow{\(\Omega_{90/120},\ P_6\)}{80, 6}
\RouteGroup{Lectores y factores de acción}
\RouteRow{\(K_D\quad /\quad K_\Omega\)}{(60,\allowbreak{}84,\allowbreak{}28)\quad /\quad (80,\allowbreak{}54,\allowbreak{}6)}
\RouteRow{\(\kappa_D\)}{59.39013\allowbreak{}58840155\allowbreak{}40637741\allowbreak{}48}
\RouteRow{\(\kappa_\Omega\)}{79.60661\allowbreak{}01397948\allowbreak{}27586470\allowbreak{}91}
\RouteRow{\(R_{\mathrm{act},D}\)}{1.000000\allowbreak{}04117409\allowbreak{}89467415\allowbreak{}75}
\RouteRow{\(R_{\mathrm{act},\Omega}\)}{1.000000\allowbreak{}05518981\allowbreak{}25318591\allowbreak{}40}
\RouteGroup{Separación de prefijos}
\RouteRow{Área de ambigüedad; área \(\log_2\)}{36; 18.0}
\RouteRow{Primer bloque de separación}{órbita: 1; celda: \textemdash{}; ruta: \textemdash{}}
\RouteRow{Ambigüedad final}{2}
\RouteGroup{Secuencias completas}
\RouteRow{\(w_6\)}{255 252 333 255 258 258 414 420 258 255 252 333 255 258 258 414 420 258}
\RouteRow{Tríadas estables por bloque}{0 1 2 3 4 5 6 7 8 9 10 11 11 12 14 15 15 17}
\end{tabularx}\endgroup\par\vspace{5pt}\noindent{\fontsize{8.4}{10}\selectfont\hyperref[trace-175]{Procedencia y códigos originales}\hfill Ficha completa · 175/324}
\clearpage\phantomsection\label{nav:vi-ruta-176}
% VI-CATALOG-ROW rutas 176
\pdfbookmark[3]{r5c8\_h+1v+1}{route-176}
\RouteTitle{Ruta 176}{r5c8\_\allowbreak{}h+1v+1}
\noindent Celda 44\quad (5, 8)\qquad Orientación \(h=1,\ v=1\)\hfill\hyperref[sec:vi-indice-rutas]{Índice}\par\vspace{4pt}
\begingroup\fontsize{10.2}{12.5}\selectfont\setlength{\tabcolsep}{5pt}\renewcommand{\arraystretch}{1.1}\rowcolors{1}{routepale}{white}\noindent\begin{tabularx}{\linewidth}{@{}>{\raggedright\arraybackslash}p{.345\linewidth}>{\raggedright\arraybackslash}X@{}}
\RouteGroup{Identidad estructural}
\RouteRow{Clase y multisección}{neutrino\_\allowbreak{}G3\quad /\quad neutrino\_\allowbreak{}G3}
\RouteRow{Colección y sector}{neutrino; leptón}
\RouteRow{Clasificación física}{neutrino}
\RouteRow{Generación, profundidad y color}{3; G3; \textemdash{}}
\RouteRow{Ruta primaria y órbita de inversión}{\(B_1\): no\qquad 5,\allowbreak{}2;\allowbreak{}5,\allowbreak{}8}
\RouteRow{Firma de clasificación}{30|\allowbreak{}-6|\allowbreak{}8|\allowbreak{}0|\allowbreak{}-12|\allowbreak{}+++-\kern0pt{}-\kern0pt{}-+++-\kern0pt{}-\kern0pt{}-}
\RouteGroup{Retornos, memoria y transporte}
\RouteRow{Retornos con fase}{clase: 27; órbita: 27; celda: 27}
\RouteRow{Retorno cinemático}{en 108: 54; extensión: 216}
\RouteRow{Giros y cambios de hoja}{71; 107}
\RouteRow{Acarreo y diferimiento}{deuda total: 2; final: 1; bloques diferidos: 2}
\RouteRow{Tríadas finales; persistencia de Pareto}{16; sí}
\RouteRow{\(K\longrightarrow K+9\)}{repeticiones: 9; cambios de memoria: 9}
\RouteGroup{Firma, fase y diferencias}
\RouteRow{\(n_A,\ s_C,\ \kappa_0,\ \nu_{120},\ \nu_{270}\)}{30, -6, 0, 0, -12}
\RouteRow{\(\tau_{12}\); firma histórica \(\mathbb Z^5\)}{+++-\kern0pt{}-\kern0pt{}-+++-\kern0pt{}-\kern0pt{}-; 30|\allowbreak{}-6|\allowbreak{}0|\allowbreak{}0|\allowbreak{}-12}
\RouteRow{\(D_1,\ D_3,\ D_4,\ D_{27},\ D_{36}\)}{54, 56, 68, 48, 32}
\RouteRow{\(\Omega_{90/120},\ P_6\)}{80, 6}
\RouteGroup{Lectores y factores de acción}
\RouteRow{\(K_D\quad /\quad K_\Omega\)}{(80,\allowbreak{}54,\allowbreak{}6)\quad /\quad (80,\allowbreak{}54,\allowbreak{}6)}
\RouteRow{\(\kappa_D\)}{79.60661\allowbreak{}01397948\allowbreak{}27586470\allowbreak{}91}
\RouteRow{\(\kappa_\Omega\)}{79.60661\allowbreak{}01397948\allowbreak{}27586470\allowbreak{}91}
\RouteRow{\(R_{\mathrm{act},D}\)}{1.000000\allowbreak{}05518981\allowbreak{}25318591\allowbreak{}40}
\RouteRow{\(R_{\mathrm{act},\Omega}\)}{1.000000\allowbreak{}05518981\allowbreak{}25318591\allowbreak{}40}
\RouteGroup{Separación de prefijos}
\RouteRow{Área de ambigüedad; área \(\log_2\)}{36; 18.0}
\RouteRow{Primer bloque de separación}{órbita: 1; celda: \textemdash{}; ruta: \textemdash{}}
\RouteRow{Ambigüedad final}{2}
\RouteGroup{Secuencias completas}
\RouteRow{\(w_6\)}{243 657 24 243 663 486 423 15 486 243 657 24 243 663 486 423 15 486}
\RouteRow{Tríadas estables por bloque}{0 1 2 3 4 5 6 7 8 9 10 11 12 13 14 14 16 16}
\end{tabularx}\endgroup\par\vspace{5pt}\noindent{\fontsize{8.4}{10}\selectfont\hyperref[trace-176]{Procedencia y códigos originales}\hfill Ficha completa · 176/324}
\clearpage\phantomsection\label{nav:vi-ruta-177}
% VI-CATALOG-ROW rutas 177
\pdfbookmark[2]{Celda 45}{cell-45}
\pdfbookmark[3]{r5c9\_h-1v-1}{route-177}
\RouteTitle{Ruta 177}{r5c9\_\allowbreak{}h-1v-1}
\noindent Celda 45\quad (5, 9)\qquad Orientación \(h=-1,\ v=-1\)\hfill\hyperref[sec:vi-indice-rutas]{Índice}\par\vspace{4pt}
\begingroup\fontsize{10.2}{12.5}\selectfont\setlength{\tabcolsep}{5pt}\renewcommand{\arraystretch}{1.1}\rowcolors{1}{routepale}{white}\noindent\begin{tabularx}{\linewidth}{@{}>{\raggedright\arraybackslash}p{.345\linewidth}>{\raggedright\arraybackslash}X@{}}
\RouteGroup{Identidad estructural}
\RouteRow{Clase y multisección}{charged\_\allowbreak{}lepton\_\allowbreak{}G4\quad /\quad charged\_\allowbreak{}lepton\_\allowbreak{}G4}
\RouteRow{Colección y sector}{leptón cargado; leptón}
\RouteRow{Clasificación física}{leptón cargado}
\RouteRow{Generación, profundidad y color}{4; G4; \textemdash{}}
\RouteRow{Ruta primaria y órbita de inversión}{\(B_1\): sí\qquad 5,\allowbreak{}1;\allowbreak{}5,\allowbreak{}9}
\RouteRow{Firma de clasificación}{34|\allowbreak{}0|\allowbreak{}-12|\allowbreak{}-2|\allowbreak{}-8|\allowbreak{}0++-\kern0pt{}-\kern0pt{}-0++-\kern0pt{}-\kern0pt{}-}
\RouteGroup{Retornos, memoria y transporte}
\RouteRow{Retornos con fase}{clase: 27; órbita: 27; celda: 27}
\RouteRow{Retorno cinemático}{en 108: 54; extensión: 216}
\RouteRow{Giros y cambios de hoja}{71; 107}
\RouteRow{Acarreo y diferimiento}{deuda total: 1; final: 0; bloques diferidos: 1}
\RouteRow{Tríadas finales; persistencia de Pareto}{17; no}
\RouteRow{\(K\longrightarrow K+9\)}{repeticiones: 9; cambios de memoria: 9}
\RouteGroup{Firma, fase y diferencias}
\RouteRow{\(n_A,\ s_C,\ \kappa_0,\ \nu_{120},\ \nu_{270}\)}{34, 0, -2, 0, -8}
\RouteRow{\(\tau_{12}\); firma histórica \(\mathbb Z^5\)}{+++0-\kern0pt{}-+++0-\kern0pt{}-; 34|\allowbreak{}0|\allowbreak{}-2|\allowbreak{}0|\allowbreak{}-8}
\RouteRow{\(D_1,\ D_3,\ D_4,\ D_{27},\ D_{36}\)}{66, 60, 68, 48, 32}
\RouteRow{\(\Omega_{90/120},\ P_6\)}{56, 5}
\RouteGroup{Lectores y factores de acción}
\RouteRow{\(K_D\quad /\quad K_\Omega\)}{(80,\allowbreak{}66,\allowbreak{}4)\quad /\quad (56,\allowbreak{}54,\allowbreak{}5)}
\RouteRow{\(\kappa_D\)}{79.51881\allowbreak{}95161180\allowbreak{}37694395\allowbreak{}39}
\RouteRow{\(\kappa_\Omega\)}{55.60649\allowbreak{}89433721\allowbreak{}35446416\allowbreak{}86}
\RouteRow{\(R_{\mathrm{act},D}\)}{1.000000\allowbreak{}05512894\allowbreak{}88901612\allowbreak{}64}
\RouteRow{\(R_{\mathrm{act},\Omega}\)}{1.000000\allowbreak{}03855097\allowbreak{}23541416\allowbreak{}45}
\RouteGroup{Separación de prefijos}
\RouteRow{Área de ambigüedad; área \(\log_2\)}{18; 0.0}
\RouteRow{Primer bloque de separación}{órbita: 1; celda: 1; ruta: 1}
\RouteRow{Ambigüedad final}{1}
\RouteGroup{Secuencias completas}
\RouteRow{\(w_6\)}{570 264 90 570 267 243 657 24 243 570 264 90 570 267 243 657 24 243}
\RouteRow{Tríadas estables por bloque}{0 1 1 3 4 5 6 7 8 9 10 11 12 13 14 15 16 17}
\end{tabularx}\endgroup\par\vspace{5pt}\noindent{\fontsize{8.4}{10}\selectfont\hyperref[trace-177]{Procedencia y códigos originales}\hfill Ficha completa · 177/324}
\clearpage\phantomsection\label{nav:vi-ruta-178}
% VI-CATALOG-ROW rutas 178
\pdfbookmark[3]{r5c9\_h-1v+1}{route-178}
\RouteTitle{Ruta 178}{r5c9\_\allowbreak{}h-1v+1}
\noindent Celda 45\quad (5, 9)\qquad Orientación \(h=-1,\ v=1\)\hfill\hyperref[sec:vi-indice-rutas]{Índice}\par\vspace{4pt}
\begingroup\fontsize{10.2}{12.5}\selectfont\setlength{\tabcolsep}{5pt}\renewcommand{\arraystretch}{1.1}\rowcolors{1}{routepale}{white}\noindent\begin{tabularx}{\linewidth}{@{}>{\raggedright\arraybackslash}p{.345\linewidth}>{\raggedright\arraybackslash}X@{}}
\RouteGroup{Identidad estructural}
\RouteRow{Clase y multisección}{charged\_\allowbreak{}lepton\_\allowbreak{}G4\quad /\quad charged\_\allowbreak{}lepton\_\allowbreak{}G4}
\RouteRow{Colección y sector}{leptón cargado; leptón}
\RouteRow{Clasificación física}{leptón cargado}
\RouteRow{Generación, profundidad y color}{4; G4; \textemdash{}}
\RouteRow{Ruta primaria y órbita de inversión}{\(B_1\): no\qquad 5,\allowbreak{}1;\allowbreak{}5,\allowbreak{}9}
\RouteRow{Firma de clasificación}{34|\allowbreak{}0|\allowbreak{}-12|\allowbreak{}-2|\allowbreak{}-8|\allowbreak{}0++-\kern0pt{}-\kern0pt{}-0++-\kern0pt{}-\kern0pt{}-}
\RouteGroup{Retornos, memoria y transporte}
\RouteRow{Retornos con fase}{clase: 27; órbita: 27; celda: 27}
\RouteRow{Retorno cinemático}{en 108: 54; extensión: 216}
\RouteRow{Giros y cambios de hoja}{71; 107}
\RouteRow{Acarreo y diferimiento}{deuda total: 1; final: 0; bloques diferidos: 1}
\RouteRow{Tríadas finales; persistencia de Pareto}{17; no}
\RouteRow{\(K\longrightarrow K+9\)}{repeticiones: 9; cambios de memoria: 9}
\RouteGroup{Firma, fase y diferencias}
\RouteRow{\(n_A,\ s_C,\ \kappa_0,\ \nu_{120},\ \nu_{270}\)}{44, 2, 0, 0, -12}
\RouteRow{\(\tau_{12}\); firma histórica \(\mathbb Z^5\)}{+++-\kern0pt{}-\kern0pt{}-+++-\kern0pt{}-\kern0pt{}-; 44|\allowbreak{}2|\allowbreak{}0|\allowbreak{}0|\allowbreak{}-12}
\RouteRow{\(D_1,\ D_3,\ D_4,\ D_{27},\ D_{36}\)}{78, 24, 78, 24, 16}
\RouteRow{\(\Omega_{90/120},\ P_6\)}{80, 6}
\RouteGroup{Lectores y factores de acción}
\RouteRow{\(K_D\quad /\quad K_\Omega\)}{(40,\allowbreak{}78,\allowbreak{}27)\quad /\quad (80,\allowbreak{}54,\allowbreak{}6)}
\RouteRow{\(\kappa_D\)}{39.43380\allowbreak{}88030085\allowbreak{}51303671\allowbreak{}14}
\RouteRow{\(\kappa_\Omega\)}{79.60661\allowbreak{}01397948\allowbreak{}27586470\allowbreak{}91}
\RouteRow{\(R_{\mathrm{act},D}\)}{1.000000\allowbreak{}02733874\allowbreak{}08548943\allowbreak{}58}
\RouteRow{\(R_{\mathrm{act},\Omega}\)}{1.000000\allowbreak{}05518981\allowbreak{}25318591\allowbreak{}40}
\RouteGroup{Separación de prefijos}
\RouteRow{Área de ambigüedad; área \(\log_2\)}{18; 0.0}
\RouteRow{Primer bloque de separación}{órbita: 1; celda: 1; ruta: 1}
\RouteRow{Ambigüedad final}{1}
\RouteGroup{Secuencias completas}
\RouteRow{\(w_6\)}{585 507 504 585 510 666 672 510 666 585 507 504 585 510 666 672 510 666}
\RouteRow{Tríadas estables por bloque}{0 1 2 3 4 5 6 7 8 9 10 11 12 13 14 14 16 17}
\end{tabularx}\endgroup\par\vspace{5pt}\noindent{\fontsize{8.4}{10}\selectfont\hyperref[trace-178]{Procedencia y códigos originales}\hfill Ficha completa · 178/324}
\clearpage\phantomsection\label{nav:vi-ruta-179}
% VI-CATALOG-ROW rutas 179
\pdfbookmark[3]{r5c9\_h+1v-1}{route-179}
\RouteTitle{Ruta 179}{r5c9\_\allowbreak{}h+1v-1}
\noindent Celda 45\quad (5, 9)\qquad Orientación \(h=1,\ v=-1\)\hfill\hyperref[sec:vi-indice-rutas]{Índice}\par\vspace{4pt}
\begingroup\fontsize{10.2}{12.5}\selectfont\setlength{\tabcolsep}{5pt}\renewcommand{\arraystretch}{1.1}\rowcolors{1}{routepale}{white}\noindent\begin{tabularx}{\linewidth}{@{}>{\raggedright\arraybackslash}p{.345\linewidth}>{\raggedright\arraybackslash}X@{}}
\RouteGroup{Identidad estructural}
\RouteRow{Clase y multisección}{charged\_\allowbreak{}lepton\_\allowbreak{}G4\quad /\quad charged\_\allowbreak{}lepton\_\allowbreak{}G4}
\RouteRow{Colección y sector}{leptón cargado; leptón}
\RouteRow{Clasificación física}{leptón cargado}
\RouteRow{Generación, profundidad y color}{4; G4; \textemdash{}}
\RouteRow{Ruta primaria y órbita de inversión}{\(B_1\): no\qquad 5,\allowbreak{}1;\allowbreak{}5,\allowbreak{}9}
\RouteRow{Firma de clasificación}{34|\allowbreak{}0|\allowbreak{}-12|\allowbreak{}-2|\allowbreak{}-8|\allowbreak{}0++-\kern0pt{}-\kern0pt{}-0++-\kern0pt{}-\kern0pt{}-}
\RouteGroup{Retornos, memoria y transporte}
\RouteRow{Retornos con fase}{clase: 27; órbita: 27; celda: 27}
\RouteRow{Retorno cinemático}{en 108: 54; extensión: 216}
\RouteRow{Giros y cambios de hoja}{71; 107}
\RouteRow{Acarreo y diferimiento}{deuda total: 1; final: 0; bloques diferidos: 1}
\RouteRow{Tríadas finales; persistencia de Pareto}{17; no}
\RouteRow{\(K\longrightarrow K+9\)}{repeticiones: 9; cambios de memoria: 9}
\RouteGroup{Firma, fase y diferencias}
\RouteRow{\(n_A,\ s_C,\ \kappa_0,\ \nu_{120},\ \nu_{270}\)}{44, -2, 0, 0, -12}
\RouteRow{\(\tau_{12}\); firma histórica \(\mathbb Z^5\)}{+++-\kern0pt{}-\kern0pt{}-+++-\kern0pt{}-\kern0pt{}-; 44|\allowbreak{}-2|\allowbreak{}0|\allowbreak{}0|\allowbreak{}-12}
\RouteRow{\(D_1,\ D_3,\ D_4,\ D_{27},\ D_{36}\)}{78, 24, 78, 24, 16}
\RouteRow{\(\Omega_{90/120},\ P_6\)}{80, 6}
\RouteGroup{Lectores y factores de acción}
\RouteRow{\(K_D\quad /\quad K_\Omega\)}{(40,\allowbreak{}78,\allowbreak{}27)\quad /\quad (80,\allowbreak{}54,\allowbreak{}6)}
\RouteRow{\(\kappa_D\)}{39.43380\allowbreak{}88030085\allowbreak{}51303671\allowbreak{}14}
\RouteRow{\(\kappa_\Omega\)}{79.60661\allowbreak{}01397948\allowbreak{}27586470\allowbreak{}91}
\RouteRow{\(R_{\mathrm{act},D}\)}{1.000000\allowbreak{}02733874\allowbreak{}08548943\allowbreak{}58}
\RouteRow{\(R_{\mathrm{act},\Omega}\)}{1.000000\allowbreak{}05518981\allowbreak{}25318591\allowbreak{}40}
\RouteGroup{Separación de prefijos}
\RouteRow{Área de ambigüedad; área \(\log_2\)}{18; 0.0}
\RouteRow{Primer bloque de separación}{órbita: 1; celda: 1; ruta: 1}
\RouteRow{Ambigüedad final}{1}
\RouteGroup{Secuencias completas}
\RouteRow{\(w_6\)}{672 510 666 672 507 504 585 507 504 672 510 666 672 507 504 585 507 504}
\RouteRow{Tríadas estables por bloque}{0 1 2 3 4 5 6 7 8 9 10 11 12 13 14 14 16 17}
\end{tabularx}\endgroup\par\vspace{5pt}\noindent{\fontsize{8.4}{10}\selectfont\hyperref[trace-179]{Procedencia y códigos originales}\hfill Ficha completa · 179/324}
\clearpage\phantomsection\label{nav:vi-ruta-180}
% VI-CATALOG-ROW rutas 180
\pdfbookmark[3]{r5c9\_h+1v+1}{route-180}
\RouteTitle{Ruta 180}{r5c9\_\allowbreak{}h+1v+1}
\noindent Celda 45\quad (5, 9)\qquad Orientación \(h=1,\ v=1\)\hfill\hyperref[sec:vi-indice-rutas]{Índice}\par\vspace{4pt}
\begingroup\fontsize{10.2}{12.5}\selectfont\setlength{\tabcolsep}{5pt}\renewcommand{\arraystretch}{1.1}\rowcolors{1}{routepale}{white}\noindent\begin{tabularx}{\linewidth}{@{}>{\raggedright\arraybackslash}p{.345\linewidth}>{\raggedright\arraybackslash}X@{}}
\RouteGroup{Identidad estructural}
\RouteRow{Clase y multisección}{charged\_\allowbreak{}lepton\_\allowbreak{}G4\quad /\quad charged\_\allowbreak{}lepton\_\allowbreak{}G4}
\RouteRow{Colección y sector}{leptón cargado; leptón}
\RouteRow{Clasificación física}{leptón cargado}
\RouteRow{Generación, profundidad y color}{4; G4; \textemdash{}}
\RouteRow{Ruta primaria y órbita de inversión}{\(B_1\): no\qquad 5,\allowbreak{}1;\allowbreak{}5,\allowbreak{}9}
\RouteRow{Firma de clasificación}{34|\allowbreak{}0|\allowbreak{}-12|\allowbreak{}-2|\allowbreak{}-8|\allowbreak{}0++-\kern0pt{}-\kern0pt{}-0++-\kern0pt{}-\kern0pt{}-}
\RouteGroup{Retornos, memoria y transporte}
\RouteRow{Retornos con fase}{clase: 27; órbita: 27; celda: 27}
\RouteRow{Retorno cinemático}{en 108: 54; extensión: 216}
\RouteRow{Giros y cambios de hoja}{71; 107}
\RouteRow{Acarreo y diferimiento}{deuda total: 0; final: 0; bloques diferidos: 0}
\RouteRow{Tríadas finales; persistencia de Pareto}{17; no}
\RouteRow{\(K\longrightarrow K+9\)}{repeticiones: 9; cambios de memoria: 9}
\RouteGroup{Firma, fase y diferencias}
\RouteRow{\(n_A,\ s_C,\ \kappa_0,\ \nu_{120},\ \nu_{270}\)}{34, 0, 2, 0, -8}
\RouteRow{\(\tau_{12}\); firma histórica \(\mathbb Z^5\)}{0++-\kern0pt{}-\kern0pt{}-0++-\kern0pt{}-\kern0pt{}-; 34|\allowbreak{}0|\allowbreak{}2|\allowbreak{}0|\allowbreak{}-8}
\RouteRow{\(D_1,\ D_3,\ D_4,\ D_{27},\ D_{36}\)}{66, 60, 68, 48, 32}
\RouteRow{\(\Omega_{90/120},\ P_6\)}{56, 5}
\RouteGroup{Lectores y factores de acción}
\RouteRow{\(K_D\quad /\quad K_\Omega\)}{(80,\allowbreak{}66,\allowbreak{}4)\quad /\quad (56,\allowbreak{}54,\allowbreak{}5)}
\RouteRow{\(\kappa_D\)}{79.51881\allowbreak{}95161180\allowbreak{}37694395\allowbreak{}39}
\RouteRow{\(\kappa_\Omega\)}{55.60649\allowbreak{}89433721\allowbreak{}35446416\allowbreak{}86}
\RouteRow{\(R_{\mathrm{act},D}\)}{1.000000\allowbreak{}05512894\allowbreak{}88901612\allowbreak{}64}
\RouteRow{\(R_{\mathrm{act},\Omega}\)}{1.000000\allowbreak{}03855097\allowbreak{}23541416\allowbreak{}45}
\RouteGroup{Separación de prefijos}
\RouteRow{Área de ambigüedad; área \(\log_2\)}{30; 8.0}
\RouteRow{Primer bloque de separación}{órbita: 5; celda: 5; ruta: 5}
\RouteRow{Ambigüedad final}{1}
\RouteGroup{Secuencias completas}
\RouteRow{\(w_6\)}{657 24 243 657 21 90 570 264 90 657 24 243 657 21 90 570 264 90}
\RouteRow{Tríadas estables por bloque}{0 1 2 3 4 5 6 7 8 9 10 11 12 13 14 15 16 17}
\end{tabularx}\endgroup\par\vspace{5pt}\noindent{\fontsize{8.4}{10}\selectfont\hyperref[trace-180]{Procedencia y códigos originales}\hfill Ficha completa · 180/324}
\clearpage\phantomsection\label{nav:vi-ruta-181}
% VI-CATALOG-ROW rutas 181
\pdfbookmark[2]{Celda 46}{cell-46}
\pdfbookmark[3]{r6c1\_h-1v-1}{route-181}
\RouteTitle{Ruta 181}{r6c1\_\allowbreak{}h-1v-1}
\noindent Celda 46\quad (6, 1)\qquad Orientación \(h=-1,\ v=-1\)\hfill\hyperref[sec:vi-indice-rutas]{Índice}\par\vspace{4pt}
\begingroup\fontsize{10.2}{12.5}\selectfont\setlength{\tabcolsep}{5pt}\renewcommand{\arraystretch}{1.1}\rowcolors{1}{routepale}{white}\noindent\begin{tabularx}{\linewidth}{@{}>{\raggedright\arraybackslash}p{.345\linewidth}>{\raggedright\arraybackslash}X@{}}
\RouteGroup{Identidad estructural}
\RouteRow{Clase y multisección}{quark\_\allowbreak{}down\_\allowbreak{}G4\quad /\quad quark\_\allowbreak{}down\_\allowbreak{}G4}
\RouteRow{Colección y sector}{quark de tipo abajo; quark}
\RouteRow{Clasificación física}{quark de tipo abajo}
\RouteRow{Generación, profundidad y color}{4; G4; B}
\RouteRow{Ruta primaria y órbita de inversión}{\(B_1\): sí\qquad 4,\allowbreak{}9;\allowbreak{}6,\allowbreak{}1}
\RouteRow{Firma de clasificación}{34|\allowbreak{}0|\allowbreak{}-4|\allowbreak{}0|\allowbreak{}-12|\allowbreak{}+++-\kern0pt{}-\kern0pt{}-+++-\kern0pt{}-\kern0pt{}-}
\RouteGroup{Retornos, memoria y transporte}
\RouteRow{Retornos con fase}{clase: 27; órbita: 27; celda: 27}
\RouteRow{Retorno cinemático}{en 108: 54; extensión: 216}
\RouteRow{Giros y cambios de hoja}{71; 107}
\RouteRow{Acarreo y diferimiento}{deuda total: 2; final: 1; bloques diferidos: 1}
\RouteRow{Tríadas finales; persistencia de Pareto}{16; sí}
\RouteRow{\(K\longrightarrow K+9\)}{repeticiones: 9; cambios de memoria: 9}
\RouteGroup{Firma, fase y diferencias}
\RouteRow{\(n_A,\ s_C,\ \kappa_0,\ \nu_{120},\ \nu_{270}\)}{34, 0, 0, 0, -12}
\RouteRow{\(\tau_{12}\); firma histórica \(\mathbb Z^5\)}{+++-\kern0pt{}-\kern0pt{}-+++-\kern0pt{}-\kern0pt{}-; 34|\allowbreak{}0|\allowbreak{}0|\allowbreak{}0|\allowbreak{}-12}
\RouteRow{\(D_1,\ D_3,\ D_4,\ D_{27},\ D_{36}\)}{66, 60, 66, 48, 32}
\RouteRow{\(\Omega_{90/120},\ P_6\)}{80, 6}
\RouteGroup{Lectores y factores de acción}
\RouteRow{\(K_D\quad /\quad K_\Omega\)}{(80,\allowbreak{}66,\allowbreak{}3)\quad /\quad (80,\allowbreak{}54,\allowbreak{}6)}
\RouteRow{\(\kappa_D\)}{79.51870\allowbreak{}83196953\allowbreak{}45554341\allowbreak{}34}
\RouteRow{\(\kappa_\Omega\)}{79.60661\allowbreak{}01397948\allowbreak{}27586470\allowbreak{}91}
\RouteRow{\(R_{\mathrm{act},D}\)}{1.000000\allowbreak{}05512887\allowbreak{}17997050\allowbreak{}72}
\RouteRow{\(R_{\mathrm{act},\Omega}\)}{1.000000\allowbreak{}05518981\allowbreak{}25318591\allowbreak{}40}
\RouteGroup{Separación de prefijos}
\RouteRow{Área de ambigüedad; área \(\log_2\)}{18; 0.0}
\RouteRow{Primer bloque de separación}{órbita: 1; celda: 1; ruta: 1}
\RouteRow{Ambigüedad final}{1}
\RouteGroup{Secuencias completas}
\RouteRow{\(w_6\)}{264 90 570 264 90 24 243 657 24 264 90 570 264 90 24 243 657 24}
\RouteRow{Tríadas estables por bloque}{0 1 2 3 4 5 6 7 8 9 10 11 12 13 14 15 15 16}
\end{tabularx}\endgroup\par\vspace{5pt}\noindent{\fontsize{8.4}{10}\selectfont\hyperref[trace-181]{Procedencia y códigos originales}\hfill Ficha completa · 181/324}
\clearpage\phantomsection\label{nav:vi-ruta-182}
% VI-CATALOG-ROW rutas 182
\pdfbookmark[3]{r6c1\_h-1v+1}{route-182}
\RouteTitle{Ruta 182}{r6c1\_\allowbreak{}h-1v+1}
\noindent Celda 46\quad (6, 1)\qquad Orientación \(h=-1,\ v=1\)\hfill\hyperref[sec:vi-indice-rutas]{Índice}\par\vspace{4pt}
\begingroup\fontsize{10.2}{12.5}\selectfont\setlength{\tabcolsep}{5pt}\renewcommand{\arraystretch}{1.1}\rowcolors{1}{routepale}{white}\noindent\begin{tabularx}{\linewidth}{@{}>{\raggedright\arraybackslash}p{.345\linewidth}>{\raggedright\arraybackslash}X@{}}
\RouteGroup{Identidad estructural}
\RouteRow{Clase y multisección}{quark\_\allowbreak{}down\_\allowbreak{}G4\quad /\quad quark\_\allowbreak{}down\_\allowbreak{}G4}
\RouteRow{Colección y sector}{quark de tipo abajo; quark}
\RouteRow{Clasificación física}{quark de tipo abajo}
\RouteRow{Generación, profundidad y color}{4; G4; B}
\RouteRow{Ruta primaria y órbita de inversión}{\(B_1\): no\qquad 4,\allowbreak{}9;\allowbreak{}6,\allowbreak{}1}
\RouteRow{Firma de clasificación}{34|\allowbreak{}0|\allowbreak{}-4|\allowbreak{}0|\allowbreak{}-12|\allowbreak{}+++-\kern0pt{}-\kern0pt{}-+++-\kern0pt{}-\kern0pt{}-}
\RouteGroup{Retornos, memoria y transporte}
\RouteRow{Retornos con fase}{clase: 27; órbita: 27; celda: 27}
\RouteRow{Retorno cinemático}{en 108: 54; extensión: 216}
\RouteRow{Giros y cambios de hoja}{71; 107}
\RouteRow{Acarreo y diferimiento}{deuda total: 2; final: 0; bloques diferidos: 1}
\RouteRow{Tríadas finales; persistencia de Pareto}{17; sí}
\RouteRow{\(K\longrightarrow K+9\)}{repeticiones: 9; cambios de memoria: 9}
\RouteGroup{Firma, fase y diferencias}
\RouteRow{\(n_A,\ s_C,\ \kappa_0,\ \nu_{120},\ \nu_{270}\)}{50, 4, -2, 2, -8}
\RouteRow{\(\tau_{12}\); firma histórica \(\mathbb Z^5\)}{+++-\kern0pt{}-0+++-\kern0pt{}-0; 50|\allowbreak{}4|\allowbreak{}-2|\allowbreak{}2|\allowbreak{}-8}
\RouteRow{\(D_1,\ D_3,\ D_4,\ D_{27},\ D_{36}\)}{84, 24, 84, 24, 16}
\RouteRow{\(\Omega_{90/120},\ P_6\)}{56, 5}
\RouteGroup{Lectores y factores de acción}
\RouteRow{\(K_D\quad /\quad K_\Omega\)}{(40,\allowbreak{}84,\allowbreak{}30)\quad /\quad (56,\allowbreak{}54,\allowbreak{}5)}
\RouteRow{\(\kappa_D\)}{39.39035\allowbreak{}82768609\allowbreak{}24917849\allowbreak{}59}
\RouteRow{\(\kappa_\Omega\)}{55.60649\allowbreak{}89433721\allowbreak{}35446416\allowbreak{}86}
\RouteRow{\(R_{\mathrm{act},D}\)}{1.000000\allowbreak{}02730861\allowbreak{}73967087\allowbreak{}05}
\RouteRow{\(R_{\mathrm{act},\Omega}\)}{1.000000\allowbreak{}03855097\allowbreak{}23541416\allowbreak{}45}
\RouteGroup{Separación de prefijos}
\RouteRow{Área de ambigüedad; área \(\log_2\)}{18; 0.0}
\RouteRow{Primer bloque de separación}{órbita: 1; celda: 1; ruta: 1}
\RouteRow{Ambigüedad final}{1}
\RouteGroup{Secuencias completas}
\RouteRow{\(w_6\)}{258 414 420 258 414 255 252 333 255 258 414 420 258 414 255 252 333 255}
\RouteRow{Tríadas estables por bloque}{0 1 2 3 4 5 6 7 8 9 10 11 11 12 14 15 16 17}
\end{tabularx}\endgroup\par\vspace{5pt}\noindent{\fontsize{8.4}{10}\selectfont\hyperref[trace-182]{Procedencia y códigos originales}\hfill Ficha completa · 182/324}
\clearpage\phantomsection\label{nav:vi-ruta-183}
% VI-CATALOG-ROW rutas 183
\pdfbookmark[3]{r6c1\_h+1v-1}{route-183}
\RouteTitle{Ruta 183}{r6c1\_\allowbreak{}h+1v-1}
\noindent Celda 46\quad (6, 1)\qquad Orientación \(h=1,\ v=-1\)\hfill\hyperref[sec:vi-indice-rutas]{Índice}\par\vspace{4pt}
\begingroup\fontsize{10.2}{12.5}\selectfont\setlength{\tabcolsep}{5pt}\renewcommand{\arraystretch}{1.1}\rowcolors{1}{routepale}{white}\noindent\begin{tabularx}{\linewidth}{@{}>{\raggedright\arraybackslash}p{.345\linewidth}>{\raggedright\arraybackslash}X@{}}
\RouteGroup{Identidad estructural}
\RouteRow{Clase y multisección}{quark\_\allowbreak{}down\_\allowbreak{}G4\quad /\quad quark\_\allowbreak{}down\_\allowbreak{}G4}
\RouteRow{Colección y sector}{quark de tipo abajo; quark}
\RouteRow{Clasificación física}{quark de tipo abajo}
\RouteRow{Generación, profundidad y color}{4; G4; B}
\RouteRow{Ruta primaria y órbita de inversión}{\(B_1\): no\qquad 4,\allowbreak{}9;\allowbreak{}6,\allowbreak{}1}
\RouteRow{Firma de clasificación}{34|\allowbreak{}0|\allowbreak{}-4|\allowbreak{}0|\allowbreak{}-12|\allowbreak{}+++-\kern0pt{}-\kern0pt{}-+++-\kern0pt{}-\kern0pt{}-}
\RouteGroup{Retornos, memoria y transporte}
\RouteRow{Retornos con fase}{clase: 27; órbita: 27; celda: 27}
\RouteRow{Retorno cinemático}{en 108: 54; extensión: 216}
\RouteRow{Giros y cambios de hoja}{71; 107}
\RouteRow{Acarreo y diferimiento}{deuda total: 3; final: 1; bloques diferidos: 3}
\RouteRow{Tríadas finales; persistencia de Pareto}{16; sí}
\RouteRow{\(K\longrightarrow K+9\)}{repeticiones: 9; cambios de memoria: 9}
\RouteGroup{Firma, fase y diferencias}
\RouteRow{\(n_A,\ s_C,\ \kappa_0,\ \nu_{120},\ \nu_{270}\)}{50, -4, 2, -2, -8}
\RouteRow{\(\tau_{12}\); firma histórica \(\mathbb Z^5\)}{++0-\kern0pt{}-\kern0pt{}-++0-\kern0pt{}-\kern0pt{}-; 50|\allowbreak{}-4|\allowbreak{}2|\allowbreak{}-2|\allowbreak{}-8}
\RouteRow{\(D_1,\ D_3,\ D_4,\ D_{27},\ D_{36}\)}{84, 24, 84, 24, 16}
\RouteRow{\(\Omega_{90/120},\ P_6\)}{56, 5}
\RouteGroup{Lectores y factores de acción}
\RouteRow{\(K_D\quad /\quad K_\Omega\)}{(40,\allowbreak{}84,\allowbreak{}30)\quad /\quad (56,\allowbreak{}54,\allowbreak{}5)}
\RouteRow{\(\kappa_D\)}{39.39035\allowbreak{}82768609\allowbreak{}24917849\allowbreak{}59}
\RouteRow{\(\kappa_\Omega\)}{55.60649\allowbreak{}89433721\allowbreak{}35446416\allowbreak{}86}
\RouteRow{\(R_{\mathrm{act},D}\)}{1.000000\allowbreak{}02730861\allowbreak{}73967087\allowbreak{}05}
\RouteRow{\(R_{\mathrm{act},\Omega}\)}{1.000000\allowbreak{}03855097\allowbreak{}23541416\allowbreak{}45}
\RouteGroup{Separación de prefijos}
\RouteRow{Área de ambigüedad; área \(\log_2\)}{18; 0.0}
\RouteRow{Primer bloque de separación}{órbita: 1; celda: 1; ruta: 1}
\RouteRow{Ambigüedad final}{1}
\RouteGroup{Secuencias completas}
\RouteRow{\(w_6\)}{252 333 255 252 333 420 258 414 420 252 333 255 252 333 420 258 414 420}
\RouteRow{Tríadas estables por bloque}{0 1 2 3 4 5 6 7 8 9 10 11 12 12 14 14 16 16}
\end{tabularx}\endgroup\par\vspace{5pt}\noindent{\fontsize{8.4}{10}\selectfont\hyperref[trace-183]{Procedencia y códigos originales}\hfill Ficha completa · 183/324}
\clearpage\phantomsection\label{nav:vi-ruta-184}
% VI-CATALOG-ROW rutas 184
\pdfbookmark[3]{r6c1\_h+1v+1}{route-184}
\RouteTitle{Ruta 184}{r6c1\_\allowbreak{}h+1v+1}
\noindent Celda 46\quad (6, 1)\qquad Orientación \(h=1,\ v=1\)\hfill\hyperref[sec:vi-indice-rutas]{Índice}\par\vspace{4pt}
\begingroup\fontsize{10.2}{12.5}\selectfont\setlength{\tabcolsep}{5pt}\renewcommand{\arraystretch}{1.1}\rowcolors{1}{routepale}{white}\noindent\begin{tabularx}{\linewidth}{@{}>{\raggedright\arraybackslash}p{.345\linewidth}>{\raggedright\arraybackslash}X@{}}
\RouteGroup{Identidad estructural}
\RouteRow{Clase y multisección}{quark\_\allowbreak{}down\_\allowbreak{}G4\quad /\quad quark\_\allowbreak{}down\_\allowbreak{}G4}
\RouteRow{Colección y sector}{quark de tipo abajo; quark}
\RouteRow{Clasificación física}{quark de tipo abajo}
\RouteRow{Generación, profundidad y color}{4; G4; B}
\RouteRow{Ruta primaria y órbita de inversión}{\(B_1\): no\qquad 4,\allowbreak{}9;\allowbreak{}6,\allowbreak{}1}
\RouteRow{Firma de clasificación}{34|\allowbreak{}0|\allowbreak{}-4|\allowbreak{}0|\allowbreak{}-12|\allowbreak{}+++-\kern0pt{}-\kern0pt{}-+++-\kern0pt{}-\kern0pt{}-}
\RouteGroup{Retornos, memoria y transporte}
\RouteRow{Retornos con fase}{clase: 27; órbita: 27; celda: 27}
\RouteRow{Retorno cinemático}{en 108: 54; extensión: 216}
\RouteRow{Giros y cambios de hoja}{71; 107}
\RouteRow{Acarreo y diferimiento}{deuda total: 0; final: 0; bloques diferidos: 0}
\RouteRow{Tríadas finales; persistencia de Pareto}{17; sí}
\RouteRow{\(K\longrightarrow K+9\)}{repeticiones: 9; cambios de memoria: 9}
\RouteGroup{Firma, fase y diferencias}
\RouteRow{\(n_A,\ s_C,\ \kappa_0,\ \nu_{120},\ \nu_{270}\)}{34, 0, 0, 0, -12}
\RouteRow{\(\tau_{12}\); firma histórica \(\mathbb Z^5\)}{+++-\kern0pt{}-\kern0pt{}-+++-\kern0pt{}-\kern0pt{}-; 34|\allowbreak{}0|\allowbreak{}0|\allowbreak{}0|\allowbreak{}-12}
\RouteRow{\(D_1,\ D_3,\ D_4,\ D_{27},\ D_{36}\)}{66, 60, 66, 48, 32}
\RouteRow{\(\Omega_{90/120},\ P_6\)}{80, 6}
\RouteGroup{Lectores y factores de acción}
\RouteRow{\(K_D\quad /\quad K_\Omega\)}{(80,\allowbreak{}66,\allowbreak{}3)\quad /\quad (80,\allowbreak{}54,\allowbreak{}6)}
\RouteRow{\(\kappa_D\)}{79.51870\allowbreak{}83196953\allowbreak{}45554341\allowbreak{}34}
\RouteRow{\(\kappa_\Omega\)}{79.60661\allowbreak{}01397948\allowbreak{}27586470\allowbreak{}91}
\RouteRow{\(R_{\mathrm{act},D}\)}{1.000000\allowbreak{}05512887\allowbreak{}17997050\allowbreak{}72}
\RouteRow{\(R_{\mathrm{act},\Omega}\)}{1.000000\allowbreak{}05518981\allowbreak{}25318591\allowbreak{}40}
\RouteGroup{Separación de prefijos}
\RouteRow{Área de ambigüedad; área \(\log_2\)}{18; 0.0}
\RouteRow{Primer bloque de separación}{órbita: 1; celda: 1; ruta: 1}
\RouteRow{Ambigüedad final}{1}
\RouteGroup{Secuencias completas}
\RouteRow{\(w_6\)}{243 657 24 243 657 570 264 90 570 243 657 24 243 657 570 264 90 570}
\RouteRow{Tríadas estables por bloque}{0 1 2 3 4 5 6 7 8 9 10 11 12 13 14 15 16 17}
\end{tabularx}\endgroup\par\vspace{5pt}\noindent{\fontsize{8.4}{10}\selectfont\hyperref[trace-184]{Procedencia y códigos originales}\hfill Ficha completa · 184/324}
\clearpage\phantomsection\label{nav:vi-ruta-185}
% VI-CATALOG-ROW rutas 185
\pdfbookmark[2]{Celda 47}{cell-47}
\pdfbookmark[3]{r6c2\_h-1v-1}{route-185}
\RouteTitle{Ruta 185}{r6c2\_\allowbreak{}h-1v-1}
\noindent Celda 47\quad (6, 2)\qquad Orientación \(h=-1,\ v=-1\)\hfill\hyperref[sec:vi-indice-rutas]{Índice}\par\vspace{4pt}
\begingroup\fontsize{10.2}{12.5}\selectfont\setlength{\tabcolsep}{5pt}\renewcommand{\arraystretch}{1.1}\rowcolors{1}{routepale}{white}\noindent\begin{tabularx}{\linewidth}{@{}>{\raggedright\arraybackslash}p{.345\linewidth}>{\raggedright\arraybackslash}X@{}}
\RouteGroup{Identidad estructural}
\RouteRow{Clase y multisección}{quark\_\allowbreak{}up\_\allowbreak{}G3\quad /\quad quark\_\allowbreak{}up\_\allowbreak{}G3}
\RouteRow{Colección y sector}{quark de tipo arriba; quark}
\RouteRow{Clasificación física}{quark de tipo arriba}
\RouteRow{Generación, profundidad y color}{3; G3; G}
\RouteRow{Ruta primaria y órbita de inversión}{\(B_1\): no\qquad 4,\allowbreak{}8;\allowbreak{}6,\allowbreak{}2}
\RouteRow{Firma de clasificación}{34|\allowbreak{}0|\allowbreak{}0|\allowbreak{}0|\allowbreak{}-8|\allowbreak{}+0+-\kern0pt{}-\kern0pt{}-+0+-\kern0pt{}-\kern0pt{}-}
\RouteGroup{Retornos, memoria y transporte}
\RouteRow{Retornos con fase}{clase: 27; órbita: 27; celda: 27}
\RouteRow{Retorno cinemático}{en 108: 54; extensión: 216}
\RouteRow{Giros y cambios de hoja}{71; 107}
\RouteRow{Acarreo y diferimiento}{deuda total: 1; final: 0; bloques diferidos: 1}
\RouteRow{Tríadas finales; persistencia de Pareto}{17; sí}
\RouteRow{\(K\longrightarrow K+9\)}{repeticiones: 9; cambios de memoria: 9}
\RouteGroup{Firma, fase y diferencias}
\RouteRow{\(n_A,\ s_C,\ \kappa_0,\ \nu_{120},\ \nu_{270}\)}{34, 0, -2, 0, -8}
\RouteRow{\(\tau_{12}\); firma histórica \(\mathbb Z^5\)}{+++-0-+++-0-; 34|\allowbreak{}0|\allowbreak{}-2|\allowbreak{}0|\allowbreak{}-8}
\RouteRow{\(D_1,\ D_3,\ D_4,\ D_{27},\ D_{36}\)}{66, 60, 66, 48, 32}
\RouteRow{\(\Omega_{90/120},\ P_6\)}{48, 5}
\RouteGroup{Lectores y factores de acción}
\RouteRow{\(K_D\quad /\quad K_\Omega\)}{(80,\allowbreak{}66,\allowbreak{}3)\quad /\quad (48,\allowbreak{}54,\allowbreak{}5)}
\RouteRow{\(\kappa_D\)}{79.51870\allowbreak{}83196953\allowbreak{}45554341\allowbreak{}34}
\RouteRow{\(\kappa_\Omega\)}{47.60649\allowbreak{}89433721\allowbreak{}35446416\allowbreak{}86}
\RouteRow{\(R_{\mathrm{act},D}\)}{1.000000\allowbreak{}05512887\allowbreak{}17997050\allowbreak{}72}
\RouteRow{\(R_{\mathrm{act},\Omega}\)}{1.000000\allowbreak{}03300471\allowbreak{}80532430\allowbreak{}84}
\RouteGroup{Separación de prefijos}
\RouteRow{Área de ambigüedad; área \(\log_2\)}{36; 18.0}
\RouteRow{Primer bloque de separación}{órbita: \textemdash{}; celda: \textemdash{}; ruta: \textemdash{}}
\RouteRow{Ambigüedad final}{2}
\RouteGroup{Secuencias completas}
\RouteRow{\(w_6\)}{486 423 15 486 423 327 498 99 327 486 423 15 486 423 327 498 99 327}
\RouteRow{Tríadas estables por bloque}{0 1 2 3 4 5 6 7 8 9 10 10 12 13 14 15 16 17}
\end{tabularx}\endgroup\par\vspace{5pt}\noindent{\fontsize{8.4}{10}\selectfont\hyperref[trace-185]{Procedencia y códigos originales}\hfill Ficha completa · 185/324}
\clearpage\phantomsection\label{nav:vi-ruta-186}
% VI-CATALOG-ROW rutas 186
\pdfbookmark[3]{r6c2\_h-1v+1}{route-186}
\RouteTitle{Ruta 186}{r6c2\_\allowbreak{}h-1v+1}
\noindent Celda 47\quad (6, 2)\qquad Orientación \(h=-1,\ v=1\)\hfill\hyperref[sec:vi-indice-rutas]{Índice}\par\vspace{4pt}
\begingroup\fontsize{10.2}{12.5}\selectfont\setlength{\tabcolsep}{5pt}\renewcommand{\arraystretch}{1.1}\rowcolors{1}{routepale}{white}\noindent\begin{tabularx}{\linewidth}{@{}>{\raggedright\arraybackslash}p{.345\linewidth}>{\raggedright\arraybackslash}X@{}}
\RouteGroup{Identidad estructural}
\RouteRow{Clase y multisección}{quark\_\allowbreak{}up\_\allowbreak{}G3\quad /\quad quark\_\allowbreak{}up\_\allowbreak{}G3}
\RouteRow{Colección y sector}{quark de tipo arriba; quark}
\RouteRow{Clasificación física}{quark de tipo arriba}
\RouteRow{Generación, profundidad y color}{3; G3; G}
\RouteRow{Ruta primaria y órbita de inversión}{\(B_1\): sí\qquad 4,\allowbreak{}8;\allowbreak{}6,\allowbreak{}2}
\RouteRow{Firma de clasificación}{34|\allowbreak{}0|\allowbreak{}0|\allowbreak{}0|\allowbreak{}-8|\allowbreak{}+0+-\kern0pt{}-\kern0pt{}-+0+-\kern0pt{}-\kern0pt{}-}
\RouteGroup{Retornos, memoria y transporte}
\RouteRow{Retornos con fase}{clase: 27; órbita: 27; celda: 27}
\RouteRow{Retorno cinemático}{en 108: 54; extensión: 216}
\RouteRow{Giros y cambios de hoja}{71; 107}
\RouteRow{Acarreo y diferimiento}{deuda total: 1; final: 0; bloques diferidos: 1}
\RouteRow{Tríadas finales; persistencia de Pareto}{17; sí}
\RouteRow{\(K\longrightarrow K+9\)}{repeticiones: 9; cambios de memoria: 9}
\RouteGroup{Firma, fase y diferencias}
\RouteRow{\(n_A,\ s_C,\ \kappa_0,\ \nu_{120},\ \nu_{270}\)}{14, -4, 0, 0, -12}
\RouteRow{\(\tau_{12}\); firma histórica \(\mathbb Z^5\)}{+++-\kern0pt{}-\kern0pt{}-+++-\kern0pt{}-\kern0pt{}-; 14|\allowbreak{}-4|\allowbreak{}0|\allowbreak{}0|\allowbreak{}-12}
\RouteRow{\(D_1,\ D_3,\ D_4,\ D_{27},\ D_{36}\)}{78, 24, 78, 24, 16}
\RouteRow{\(\Omega_{90/120},\ P_6\)}{80, 6}
\RouteGroup{Lectores y factores de acción}
\RouteRow{\(K_D\quad /\quad K_\Omega\)}{(40,\allowbreak{}78,\allowbreak{}27)\quad /\quad (80,\allowbreak{}54,\allowbreak{}6)}
\RouteRow{\(\kappa_D\)}{39.43380\allowbreak{}88030085\allowbreak{}51303671\allowbreak{}14}
\RouteRow{\(\kappa_\Omega\)}{79.60661\allowbreak{}01397948\allowbreak{}27586470\allowbreak{}91}
\RouteRow{\(R_{\mathrm{act},D}\)}{1.000000\allowbreak{}02733874\allowbreak{}08548943\allowbreak{}58}
\RouteRow{\(R_{\mathrm{act},\Omega}\)}{1.000000\allowbreak{}05518981\allowbreak{}25318591\allowbreak{}40}
\RouteGroup{Separación de prefijos}
\RouteRow{Área de ambigüedad; área \(\log_2\)}{36; 18.0}
\RouteRow{Primer bloque de separación}{órbita: \textemdash{}; celda: \textemdash{}; ruta: \textemdash{}}
\RouteRow{Ambigüedad final}{2}
\RouteGroup{Secuencias completas}
\RouteRow{\(w_6\)}{504 585 507 504 585 672 510 666 672 504 585 507 504 585 672 510 666 672}
\RouteRow{Tríadas estables por bloque}{0 1 2 3 4 5 6 7 8 9 10 11 12 13 14 14 16 17}
\end{tabularx}\endgroup\par\vspace{5pt}\noindent{\fontsize{8.4}{10}\selectfont\hyperref[trace-186]{Procedencia y códigos originales}\hfill Ficha completa · 186/324}
\clearpage\phantomsection\label{nav:vi-ruta-187}
% VI-CATALOG-ROW rutas 187
\pdfbookmark[3]{r6c2\_h+1v-1}{route-187}
\RouteTitle{Ruta 187}{r6c2\_\allowbreak{}h+1v-1}
\noindent Celda 47\quad (6, 2)\qquad Orientación \(h=1,\ v=-1\)\hfill\hyperref[sec:vi-indice-rutas]{Índice}\par\vspace{4pt}
\begingroup\fontsize{10.2}{12.5}\selectfont\setlength{\tabcolsep}{5pt}\renewcommand{\arraystretch}{1.1}\rowcolors{1}{routepale}{white}\noindent\begin{tabularx}{\linewidth}{@{}>{\raggedright\arraybackslash}p{.345\linewidth}>{\raggedright\arraybackslash}X@{}}
\RouteGroup{Identidad estructural}
\RouteRow{Clase y multisección}{quark\_\allowbreak{}up\_\allowbreak{}G3\quad /\quad quark\_\allowbreak{}up\_\allowbreak{}G3}
\RouteRow{Colección y sector}{quark de tipo arriba; quark}
\RouteRow{Clasificación física}{quark de tipo arriba}
\RouteRow{Generación, profundidad y color}{3; G3; G}
\RouteRow{Ruta primaria y órbita de inversión}{\(B_1\): no\qquad 4,\allowbreak{}8;\allowbreak{}6,\allowbreak{}2}
\RouteRow{Firma de clasificación}{34|\allowbreak{}0|\allowbreak{}0|\allowbreak{}0|\allowbreak{}-8|\allowbreak{}+0+-\kern0pt{}-\kern0pt{}-+0+-\kern0pt{}-\kern0pt{}-}
\RouteGroup{Retornos, memoria y transporte}
\RouteRow{Retornos con fase}{clase: 27; órbita: 27; celda: 27}
\RouteRow{Retorno cinemático}{en 108: 54; extensión: 216}
\RouteRow{Giros y cambios de hoja}{71; 107}
\RouteRow{Acarreo y diferimiento}{deuda total: 1; final: 0; bloques diferidos: 1}
\RouteRow{Tríadas finales; persistencia de Pareto}{17; sí}
\RouteRow{\(K\longrightarrow K+9\)}{repeticiones: 9; cambios de memoria: 9}
\RouteGroup{Firma, fase y diferencias}
\RouteRow{\(n_A,\ s_C,\ \kappa_0,\ \nu_{120},\ \nu_{270}\)}{14, 4, 0, 0, -12}
\RouteRow{\(\tau_{12}\); firma histórica \(\mathbb Z^5\)}{+++-\kern0pt{}-\kern0pt{}-+++-\kern0pt{}-\kern0pt{}-; 14|\allowbreak{}4|\allowbreak{}0|\allowbreak{}0|\allowbreak{}-12}
\RouteRow{\(D_1,\ D_3,\ D_4,\ D_{27},\ D_{36}\)}{78, 24, 78, 24, 16}
\RouteRow{\(\Omega_{90/120},\ P_6\)}{80, 6}
\RouteGroup{Lectores y factores de acción}
\RouteRow{\(K_D\quad /\quad K_\Omega\)}{(40,\allowbreak{}78,\allowbreak{}27)\quad /\quad (80,\allowbreak{}54,\allowbreak{}6)}
\RouteRow{\(\kappa_D\)}{39.43380\allowbreak{}88030085\allowbreak{}51303671\allowbreak{}14}
\RouteRow{\(\kappa_\Omega\)}{79.60661\allowbreak{}01397948\allowbreak{}27586470\allowbreak{}91}
\RouteRow{\(R_{\mathrm{act},D}\)}{1.000000\allowbreak{}02733874\allowbreak{}08548943\allowbreak{}58}
\RouteRow{\(R_{\mathrm{act},\Omega}\)}{1.000000\allowbreak{}05518981\allowbreak{}25318591\allowbreak{}40}
\RouteGroup{Separación de prefijos}
\RouteRow{Área de ambigüedad; área \(\log_2\)}{36; 18.0}
\RouteRow{Primer bloque de separación}{órbita: \textemdash{}; celda: \textemdash{}; ruta: \textemdash{}}
\RouteRow{Ambigüedad final}{2}
\RouteGroup{Secuencias completas}
\RouteRow{\(w_6\)}{510 666 672 510 666 507 504 585 507 510 666 672 510 666 507 504 585 507}
\RouteRow{Tríadas estables por bloque}{0 1 2 3 4 5 6 7 8 9 10 11 12 13 14 15 15 17}
\end{tabularx}\endgroup\par\vspace{5pt}\noindent{\fontsize{8.4}{10}\selectfont\hyperref[trace-187]{Procedencia y códigos originales}\hfill Ficha completa · 187/324}
\clearpage\phantomsection\label{nav:vi-ruta-188}
% VI-CATALOG-ROW rutas 188
\pdfbookmark[3]{r6c2\_h+1v+1}{route-188}
\RouteTitle{Ruta 188}{r6c2\_\allowbreak{}h+1v+1}
\noindent Celda 47\quad (6, 2)\qquad Orientación \(h=1,\ v=1\)\hfill\hyperref[sec:vi-indice-rutas]{Índice}\par\vspace{4pt}
\begingroup\fontsize{10.2}{12.5}\selectfont\setlength{\tabcolsep}{5pt}\renewcommand{\arraystretch}{1.1}\rowcolors{1}{routepale}{white}\noindent\begin{tabularx}{\linewidth}{@{}>{\raggedright\arraybackslash}p{.345\linewidth}>{\raggedright\arraybackslash}X@{}}
\RouteGroup{Identidad estructural}
\RouteRow{Clase y multisección}{quark\_\allowbreak{}up\_\allowbreak{}G3\quad /\quad quark\_\allowbreak{}up\_\allowbreak{}G3}
\RouteRow{Colección y sector}{quark de tipo arriba; quark}
\RouteRow{Clasificación física}{quark de tipo arriba}
\RouteRow{Generación, profundidad y color}{3; G3; G}
\RouteRow{Ruta primaria y órbita de inversión}{\(B_1\): no\qquad 4,\allowbreak{}8;\allowbreak{}6,\allowbreak{}2}
\RouteRow{Firma de clasificación}{34|\allowbreak{}0|\allowbreak{}0|\allowbreak{}0|\allowbreak{}-8|\allowbreak{}+0+-\kern0pt{}-\kern0pt{}-+0+-\kern0pt{}-\kern0pt{}-}
\RouteGroup{Retornos, memoria y transporte}
\RouteRow{Retornos con fase}{clase: 27; órbita: 27; celda: 27}
\RouteRow{Retorno cinemático}{en 108: 54; extensión: 216}
\RouteRow{Giros y cambios de hoja}{71; 107}
\RouteRow{Acarreo y diferimiento}{deuda total: 0; final: 0; bloques diferidos: 0}
\RouteRow{Tríadas finales; persistencia de Pareto}{17; sí}
\RouteRow{\(K\longrightarrow K+9\)}{repeticiones: 9; cambios de memoria: 9}
\RouteGroup{Firma, fase y diferencias}
\RouteRow{\(n_A,\ s_C,\ \kappa_0,\ \nu_{120},\ \nu_{270}\)}{34, 0, 2, 0, -8}
\RouteRow{\(\tau_{12}\); firma histórica \(\mathbb Z^5\)}{+0+-\kern0pt{}-\kern0pt{}-+0+-\kern0pt{}-\kern0pt{}-; 34|\allowbreak{}0|\allowbreak{}2|\allowbreak{}0|\allowbreak{}-8}
\RouteRow{\(D_1,\ D_3,\ D_4,\ D_{27},\ D_{36}\)}{66, 60, 66, 48, 32}
\RouteRow{\(\Omega_{90/120},\ P_6\)}{48, 5}
\RouteGroup{Lectores y factores de acción}
\RouteRow{\(K_D\quad /\quad K_\Omega\)}{(80,\allowbreak{}66,\allowbreak{}3)\quad /\quad (48,\allowbreak{}54,\allowbreak{}5)}
\RouteRow{\(\kappa_D\)}{79.51870\allowbreak{}83196953\allowbreak{}45554341\allowbreak{}34}
\RouteRow{\(\kappa_\Omega\)}{47.60649\allowbreak{}89433721\allowbreak{}35446416\allowbreak{}86}
\RouteRow{\(R_{\mathrm{act},D}\)}{1.000000\allowbreak{}05512887\allowbreak{}17997050\allowbreak{}72}
\RouteRow{\(R_{\mathrm{act},\Omega}\)}{1.000000\allowbreak{}03300471\allowbreak{}80532430\allowbreak{}84}
\RouteGroup{Separación de prefijos}
\RouteRow{Área de ambigüedad; área \(\log_2\)}{36; 18.0}
\RouteRow{Primer bloque de separación}{órbita: \textemdash{}; celda: \textemdash{}; ruta: \textemdash{}}
\RouteRow{Ambigüedad final}{2}
\RouteGroup{Secuencias completas}
\RouteRow{\(w_6\)}{498 99 327 498 99 15 486 423 15 498 99 327 498 99 15 486 423 15}
\RouteRow{Tríadas estables por bloque}{0 1 2 3 4 5 6 7 8 9 10 11 12 13 14 15 16 17}
\end{tabularx}\endgroup\par\vspace{5pt}\noindent{\fontsize{8.4}{10}\selectfont\hyperref[trace-188]{Procedencia y códigos originales}\hfill Ficha completa · 188/324}
\clearpage\phantomsection\label{nav:vi-ruta-189}
% VI-CATALOG-ROW rutas 189
\pdfbookmark[2]{Celda 48}{cell-48}
\pdfbookmark[3]{r6c3\_h-1v-1}{route-189}
\RouteTitle{Ruta 189}{r6c3\_\allowbreak{}h-1v-1}
\noindent Celda 48\quad (6, 3)\qquad Orientación \(h=-1,\ v=-1\)\hfill\hyperref[sec:vi-indice-rutas]{Índice}\par\vspace{4pt}
\begingroup\fontsize{10.2}{12.5}\selectfont\setlength{\tabcolsep}{5pt}\renewcommand{\arraystretch}{1.1}\rowcolors{1}{routepale}{white}\noindent\begin{tabularx}{\linewidth}{@{}>{\raggedright\arraybackslash}p{.345\linewidth}>{\raggedright\arraybackslash}X@{}}
\RouteGroup{Identidad estructural}
\RouteRow{Clase y multisección}{quark\_\allowbreak{}down\_\allowbreak{}G2\quad /\quad quark\_\allowbreak{}down\_\allowbreak{}G2}
\RouteRow{Colección y sector}{quark de tipo abajo; quark}
\RouteRow{Clasificación física}{quark de tipo abajo}
\RouteRow{Generación, profundidad y color}{2; G2; R}
\RouteRow{Ruta primaria y órbita de inversión}{\(B_1\): no\qquad 4,\allowbreak{}7;\allowbreak{}6,\allowbreak{}3}
\RouteRow{Firma de clasificación}{30|\allowbreak{}6|\allowbreak{}0|\allowbreak{}2|\allowbreak{}-8|\allowbreak{}+++0-\kern0pt{}-+++0-\kern0pt{}-}
\RouteGroup{Retornos, memoria y transporte}
\RouteRow{Retornos con fase}{clase: 27; órbita: 27; celda: 27}
\RouteRow{Retorno cinemático}{en 108: 54; extensión: 216}
\RouteRow{Giros y cambios de hoja}{71; 107}
\RouteRow{Acarreo y diferimiento}{deuda total: 1; final: 0; bloques diferidos: 1}
\RouteRow{Tríadas finales; persistencia de Pareto}{17; sí}
\RouteRow{\(K\longrightarrow K+9\)}{repeticiones: 9; cambios de memoria: 9}
\RouteGroup{Firma, fase y diferencias}
\RouteRow{\(n_A,\ s_C,\ \kappa_0,\ \nu_{120},\ \nu_{270}\)}{30, -6, 2, -2, -8}
\RouteRow{\(\tau_{12}\); firma histórica \(\mathbb Z^5\)}{0++-\kern0pt{}-\kern0pt{}-0++-\kern0pt{}-\kern0pt{}-; 30|\allowbreak{}-6|\allowbreak{}2|\allowbreak{}-2|\allowbreak{}-8}
\RouteRow{\(D_1,\ D_3,\ D_4,\ D_{27},\ D_{36}\)}{54, 60, 72, 24, 16}
\RouteRow{\(\Omega_{90/120},\ P_6\)}{56, 5}
\RouteGroup{Lectores y factores de acción}
\RouteRow{\(K_D\quad /\quad K_\Omega\)}{(40,\allowbreak{}54,\allowbreak{}6)\quad /\quad (56,\allowbreak{}54,\allowbreak{}5)}
\RouteRow{\(\kappa_D\)}{39.60661\allowbreak{}01397948\allowbreak{}27586470\allowbreak{}91}
\RouteRow{\(\kappa_\Omega\)}{55.60649\allowbreak{}89433721\allowbreak{}35446416\allowbreak{}86}
\RouteRow{\(R_{\mathrm{act},D}\)}{1.000000\allowbreak{}02745854\allowbreak{}08735595\allowbreak{}18}
\RouteRow{\(R_{\mathrm{act},\Omega}\)}{1.000000\allowbreak{}03855097\allowbreak{}23541416\allowbreak{}45}
\RouteGroup{Separación de prefijos}
\RouteRow{Área de ambigüedad; área \(\log_2\)}{18; 0.0}
\RouteRow{Primer bloque de separación}{órbita: 1; celda: 1; ruta: 1}
\RouteRow{Ambigüedad final}{1}
\RouteGroup{Secuencias completas}
\RouteRow{\(w_6\)}{15 486 423 15 486 657 24 243 657 15 486 423 15 486 657 24 243 657}
\RouteRow{Tríadas estables por bloque}{0 1 2 3 4 5 6 7 8 9 10 11 12 13 14 15 15 17}
\end{tabularx}\endgroup\par\vspace{5pt}\noindent{\fontsize{8.4}{10}\selectfont\hyperref[trace-189]{Procedencia y códigos originales}\hfill Ficha completa · 189/324}
\clearpage\phantomsection\label{nav:vi-ruta-190}
% VI-CATALOG-ROW rutas 190
\pdfbookmark[3]{r6c3\_h-1v+1}{route-190}
\RouteTitle{Ruta 190}{r6c3\_\allowbreak{}h-1v+1}
\noindent Celda 48\quad (6, 3)\qquad Orientación \(h=-1,\ v=1\)\hfill\hyperref[sec:vi-indice-rutas]{Índice}\par\vspace{4pt}
\begingroup\fontsize{10.2}{12.5}\selectfont\setlength{\tabcolsep}{5pt}\renewcommand{\arraystretch}{1.1}\rowcolors{1}{routepale}{white}\noindent\begin{tabularx}{\linewidth}{@{}>{\raggedright\arraybackslash}p{.345\linewidth}>{\raggedright\arraybackslash}X@{}}
\RouteGroup{Identidad estructural}
\RouteRow{Clase y multisección}{quark\_\allowbreak{}down\_\allowbreak{}G2\quad /\quad quark\_\allowbreak{}down\_\allowbreak{}G2}
\RouteRow{Colección y sector}{quark de tipo abajo; quark}
\RouteRow{Clasificación física}{quark de tipo abajo}
\RouteRow{Generación, profundidad y color}{2; G2; R}
\RouteRow{Ruta primaria y órbita de inversión}{\(B_1\): sí\qquad 4,\allowbreak{}7;\allowbreak{}6,\allowbreak{}3}
\RouteRow{Firma de clasificación}{30|\allowbreak{}6|\allowbreak{}0|\allowbreak{}2|\allowbreak{}-8|\allowbreak{}+++0-\kern0pt{}-+++0-\kern0pt{}-}
\RouteGroup{Retornos, memoria y transporte}
\RouteRow{Retornos con fase}{clase: 27; órbita: 27; celda: 27}
\RouteRow{Retorno cinemático}{en 108: 54; extensión: 216}
\RouteRow{Giros y cambios de hoja}{71; 107}
\RouteRow{Acarreo y diferimiento}{deuda total: 0; final: 0; bloques diferidos: 0}
\RouteRow{Tríadas finales; persistencia de Pareto}{17; sí}
\RouteRow{\(K\longrightarrow K+9\)}{repeticiones: 9; cambios de memoria: 9}
\RouteGroup{Firma, fase y diferencias}
\RouteRow{\(n_A,\ s_C,\ \kappa_0,\ \nu_{120},\ \nu_{270}\)}{20, 0, 0, 0, 0}
\RouteRow{\(\tau_{12}\); firma histórica \(\mathbb Z^5\)}{00+-0000+-00; 20|\allowbreak{}0|\allowbreak{}0|\allowbreak{}0|\allowbreak{}0}
\RouteRow{\(D_1,\ D_3,\ D_4,\ D_{27},\ D_{36}\)}{24, 24, 24, 0, 0}
\RouteRow{\(\Omega_{90/120},\ P_6\)}{4, 2}
\RouteGroup{Lectores y factores de acción}
\RouteRow{\(K_D\quad /\quad K_\Omega\)}{(0,\allowbreak{}24,\allowbreak{}0)\quad /\quad (4,\allowbreak{}54,\allowbreak{}2)}
\RouteRow{\(\kappa_D\)}{-0.175136\allowbreak{}46166281\allowbreak{}12239348\allowbreak{}425}
\RouteRow{\(\kappa_\Omega\)}{3.606165\allowbreak{}35410405\allowbreak{}90262547\allowbreak{}07}
\RouteRow{\(R_{\mathrm{act},D}\)}{0.999999\allowbreak{}99987858\allowbreak{}10851337\allowbreak{}881}
\RouteRow{\(R_{\mathrm{act},\Omega}\)}{1.000000\allowbreak{}00250008\allowbreak{}86767963\allowbreak{}14}
\RouteGroup{Separación de prefijos}
\RouteRow{Área de ambigüedad; área \(\log_2\)}{36; 18.0}
\RouteRow{Primer bloque de separación}{órbita: 1; celda: 1; ruta: \textemdash{}}
\RouteRow{Ambigüedad final}{2}
\RouteGroup{Secuencias completas}
\RouteRow{\(w_6\)}{3 0 81 3 0 81 3 0 81 3 0 81 3 0 81 3 0 81}
\RouteRow{Tríadas estables por bloque}{0 1 2 3 4 5 6 7 8 9 10 11 12 13 14 15 16 17}
\end{tabularx}\endgroup\par\vspace{5pt}\noindent{\fontsize{8.4}{10}\selectfont\hyperref[trace-190]{Procedencia y códigos originales}\hfill Ficha completa · 190/324}
\clearpage\phantomsection\label{nav:vi-ruta-191}
% VI-CATALOG-ROW rutas 191
\pdfbookmark[3]{r6c3\_h+1v-1}{route-191}
\RouteTitle{Ruta 191}{r6c3\_\allowbreak{}h+1v-1}
\noindent Celda 48\quad (6, 3)\qquad Orientación \(h=1,\ v=-1\)\hfill\hyperref[sec:vi-indice-rutas]{Índice}\par\vspace{4pt}
\begingroup\fontsize{10.2}{12.5}\selectfont\setlength{\tabcolsep}{5pt}\renewcommand{\arraystretch}{1.1}\rowcolors{1}{routepale}{white}\noindent\begin{tabularx}{\linewidth}{@{}>{\raggedright\arraybackslash}p{.345\linewidth}>{\raggedright\arraybackslash}X@{}}
\RouteGroup{Identidad estructural}
\RouteRow{Clase y multisección}{quark\_\allowbreak{}down\_\allowbreak{}G2\quad /\quad quark\_\allowbreak{}down\_\allowbreak{}G2}
\RouteRow{Colección y sector}{quark de tipo abajo; quark}
\RouteRow{Clasificación física}{quark de tipo abajo}
\RouteRow{Generación, profundidad y color}{2; G2; R}
\RouteRow{Ruta primaria y órbita de inversión}{\(B_1\): no\qquad 4,\allowbreak{}7;\allowbreak{}6,\allowbreak{}3}
\RouteRow{Firma de clasificación}{30|\allowbreak{}6|\allowbreak{}0|\allowbreak{}2|\allowbreak{}-8|\allowbreak{}+++0-\kern0pt{}-+++0-\kern0pt{}-}
\RouteGroup{Retornos, memoria y transporte}
\RouteRow{Retornos con fase}{clase: 27; órbita: 27; celda: 27}
\RouteRow{Retorno cinemático}{en 108: 54; extensión: 216}
\RouteRow{Giros y cambios de hoja}{71; 107}
\RouteRow{Acarreo y diferimiento}{deuda total: 0; final: 0; bloques diferidos: 0}
\RouteRow{Tríadas finales; persistencia de Pareto}{17; sí}
\RouteRow{\(K\longrightarrow K+9\)}{repeticiones: 9; cambios de memoria: 9}
\RouteGroup{Firma, fase y diferencias}
\RouteRow{\(n_A,\ s_C,\ \kappa_0,\ \nu_{120},\ \nu_{270}\)}{20, 0, 0, 0, 0}
\RouteRow{\(\tau_{12}\); firma histórica \(\mathbb Z^5\)}{+0000-+0000-; 20|\allowbreak{}0|\allowbreak{}0|\allowbreak{}0|\allowbreak{}0}
\RouteRow{\(D_1,\ D_3,\ D_4,\ D_{27},\ D_{36}\)}{24, 24, 24, 0, 0}
\RouteRow{\(\Omega_{90/120},\ P_6\)}{4, 2}
\RouteGroup{Lectores y factores de acción}
\RouteRow{\(K_D\quad /\quad K_\Omega\)}{(0,\allowbreak{}24,\allowbreak{}0)\quad /\quad (4,\allowbreak{}54,\allowbreak{}2)}
\RouteRow{\(\kappa_D\)}{-0.175136\allowbreak{}46166281\allowbreak{}12239348\allowbreak{}425}
\RouteRow{\(\kappa_\Omega\)}{3.606165\allowbreak{}35410405\allowbreak{}90262547\allowbreak{}07}
\RouteRow{\(R_{\mathrm{act},D}\)}{0.999999\allowbreak{}99987858\allowbreak{}10851337\allowbreak{}881}
\RouteRow{\(R_{\mathrm{act},\Omega}\)}{1.000000\allowbreak{}00250008\allowbreak{}86767963\allowbreak{}14}
\RouteGroup{Separación de prefijos}
\RouteRow{Área de ambigüedad; área \(\log_2\)}{36; 18.0}
\RouteRow{Primer bloque de separación}{órbita: 1; celda: 1; ruta: \textemdash{}}
\RouteRow{Ambigüedad final}{2}
\RouteGroup{Secuencias completas}
\RouteRow{\(w_6\)}{3 0 81 3 0 81 3 0 81 3 0 81 3 0 81 3 0 81}
\RouteRow{Tríadas estables por bloque}{0 1 2 3 4 5 6 7 8 9 10 11 12 13 14 15 16 17}
\end{tabularx}\endgroup\par\vspace{5pt}\noindent{\fontsize{8.4}{10}\selectfont\hyperref[trace-191]{Procedencia y códigos originales}\hfill Ficha completa · 191/324}
\clearpage\phantomsection\label{nav:vi-ruta-192}
% VI-CATALOG-ROW rutas 192
\pdfbookmark[3]{r6c3\_h+1v+1}{route-192}
\RouteTitle{Ruta 192}{r6c3\_\allowbreak{}h+1v+1}
\noindent Celda 48\quad (6, 3)\qquad Orientación \(h=1,\ v=1\)\hfill\hyperref[sec:vi-indice-rutas]{Índice}\par\vspace{4pt}
\begingroup\fontsize{10.2}{12.5}\selectfont\setlength{\tabcolsep}{5pt}\renewcommand{\arraystretch}{1.1}\rowcolors{1}{routepale}{white}\noindent\begin{tabularx}{\linewidth}{@{}>{\raggedright\arraybackslash}p{.345\linewidth}>{\raggedright\arraybackslash}X@{}}
\RouteGroup{Identidad estructural}
\RouteRow{Clase y multisección}{quark\_\allowbreak{}down\_\allowbreak{}G2\quad /\quad quark\_\allowbreak{}down\_\allowbreak{}G2}
\RouteRow{Colección y sector}{quark de tipo abajo; quark}
\RouteRow{Clasificación física}{quark de tipo abajo}
\RouteRow{Generación, profundidad y color}{2; G2; R}
\RouteRow{Ruta primaria y órbita de inversión}{\(B_1\): no\qquad 4,\allowbreak{}7;\allowbreak{}6,\allowbreak{}3}
\RouteRow{Firma de clasificación}{30|\allowbreak{}6|\allowbreak{}0|\allowbreak{}2|\allowbreak{}-8|\allowbreak{}+++0-\kern0pt{}-+++0-\kern0pt{}-}
\RouteGroup{Retornos, memoria y transporte}
\RouteRow{Retornos con fase}{clase: 27; órbita: 27; celda: 27}
\RouteRow{Retorno cinemático}{en 108: 54; extensión: 216}
\RouteRow{Giros y cambios de hoja}{71; 107}
\RouteRow{Acarreo y diferimiento}{deuda total: 2; final: 1; bloques diferidos: 2}
\RouteRow{Tríadas finales; persistencia de Pareto}{16; sí}
\RouteRow{\(K\longrightarrow K+9\)}{repeticiones: 9; cambios de memoria: 9}
\RouteGroup{Firma, fase y diferencias}
\RouteRow{\(n_A,\ s_C,\ \kappa_0,\ \nu_{120},\ \nu_{270}\)}{30, 6, -2, 2, -8}
\RouteRow{\(\tau_{12}\); firma histórica \(\mathbb Z^5\)}{+++0-\kern0pt{}-+++0-\kern0pt{}-; 30|\allowbreak{}6|\allowbreak{}-2|\allowbreak{}2|\allowbreak{}-8}
\RouteRow{\(D_1,\ D_3,\ D_4,\ D_{27},\ D_{36}\)}{54, 60, 72, 24, 16}
\RouteRow{\(\Omega_{90/120},\ P_6\)}{56, 5}
\RouteGroup{Lectores y factores de acción}
\RouteRow{\(K_D\quad /\quad K_\Omega\)}{(40,\allowbreak{}54,\allowbreak{}6)\quad /\quad (56,\allowbreak{}54,\allowbreak{}5)}
\RouteRow{\(\kappa_D\)}{39.60661\allowbreak{}01397948\allowbreak{}27586470\allowbreak{}91}
\RouteRow{\(\kappa_\Omega\)}{55.60649\allowbreak{}89433721\allowbreak{}35446416\allowbreak{}86}
\RouteRow{\(R_{\mathrm{act},D}\)}{1.000000\allowbreak{}02745854\allowbreak{}08735595\allowbreak{}18}
\RouteRow{\(R_{\mathrm{act},\Omega}\)}{1.000000\allowbreak{}03855097\allowbreak{}23541416\allowbreak{}45}
\RouteGroup{Separación de prefijos}
\RouteRow{Área de ambigüedad; área \(\log_2\)}{18; 0.0}
\RouteRow{Primer bloque de separación}{órbita: 1; celda: 1; ruta: 1}
\RouteRow{Ambigüedad final}{1}
\RouteGroup{Secuencias completas}
\RouteRow{\(w_6\)}{24 243 657 24 243 423 15 486 423 24 243 657 24 243 423 15 486 423}
\RouteRow{Tríadas estables por bloque}{0 1 2 3 4 5 6 7 8 9 10 11 12 13 13 15 16 16}
\end{tabularx}\endgroup\par\vspace{5pt}\noindent{\fontsize{8.4}{10}\selectfont\hyperref[trace-192]{Procedencia y códigos originales}\hfill Ficha completa · 192/324}
\clearpage\phantomsection\label{nav:vi-ruta-193}
% VI-CATALOG-ROW rutas 193
\pdfbookmark[2]{Celda 49}{cell-49}
\pdfbookmark[3]{r6c4\_h-1v-1}{route-193}
\RouteTitle{Ruta 193}{r6c4\_\allowbreak{}h-1v-1}
\noindent Celda 49\quad (6, 4)\qquad Orientación \(h=-1,\ v=-1\)\hfill\hyperref[sec:vi-indice-rutas]{Índice}\par\vspace{4pt}
\begingroup\fontsize{10.2}{12.5}\selectfont\setlength{\tabcolsep}{5pt}\renewcommand{\arraystretch}{1.1}\rowcolors{1}{routepale}{white}\noindent\begin{tabularx}{\linewidth}{@{}>{\raggedright\arraybackslash}p{.345\linewidth}>{\raggedright\arraybackslash}X@{}}
\RouteGroup{Identidad estructural}
\RouteRow{Clase y multisección}{quark\_\allowbreak{}up\_\allowbreak{}G1\quad /\quad quark\_\allowbreak{}up\_\allowbreak{}G1}
\RouteRow{Colección y sector}{quark de tipo arriba; quark}
\RouteRow{Clasificación física}{quark de tipo arriba}
\RouteRow{Generación, profundidad y color}{1; G1; B}
\RouteRow{Ruta primaria y órbita de inversión}{\(B_1\): sí\qquad 4,\allowbreak{}6;\allowbreak{}6,\allowbreak{}4}
\RouteRow{Firma de clasificación}{28|\allowbreak{}6|\allowbreak{}-2|\allowbreak{}0|\allowbreak{}-12|\allowbreak{}+++-\kern0pt{}-\kern0pt{}-+++-\kern0pt{}-\kern0pt{}-}
\RouteGroup{Retornos, memoria y transporte}
\RouteRow{Retornos con fase}{clase: 27; órbita: 27; celda: 27}
\RouteRow{Retorno cinemático}{en 108: 54; extensión: 216}
\RouteRow{Giros y cambios de hoja}{71; 107}
\RouteRow{Acarreo y diferimiento}{deuda total: 2; final: 1; bloques diferidos: 1}
\RouteRow{Tríadas finales; persistencia de Pareto}{16; no}
\RouteRow{\(K\longrightarrow K+9\)}{repeticiones: 9; cambios de memoria: 9}
\RouteGroup{Firma, fase y diferencias}
\RouteRow{\(n_A,\ s_C,\ \kappa_0,\ \nu_{120},\ \nu_{270}\)}{28, -6, 0, 0, -12}
\RouteRow{\(\tau_{12}\); firma histórica \(\mathbb Z^5\)}{+++-\kern0pt{}-\kern0pt{}-+++-\kern0pt{}-\kern0pt{}-; 28|\allowbreak{}-6|\allowbreak{}0|\allowbreak{}0|\allowbreak{}-12}
\RouteRow{\(D_1,\ D_3,\ D_4,\ D_{27},\ D_{36}\)}{66, 60, 66, 48, 32}
\RouteRow{\(\Omega_{90/120},\ P_6\)}{80, 6}
\RouteGroup{Lectores y factores de acción}
\RouteRow{\(K_D\quad /\quad K_\Omega\)}{(80,\allowbreak{}66,\allowbreak{}3)\quad /\quad (80,\allowbreak{}54,\allowbreak{}6)}
\RouteRow{\(\kappa_D\)}{79.51870\allowbreak{}83196953\allowbreak{}45554341\allowbreak{}34}
\RouteRow{\(\kappa_\Omega\)}{79.60661\allowbreak{}01397948\allowbreak{}27586470\allowbreak{}91}
\RouteRow{\(R_{\mathrm{act},D}\)}{1.000000\allowbreak{}05512887\allowbreak{}17997050\allowbreak{}72}
\RouteRow{\(R_{\mathrm{act},\Omega}\)}{1.000000\allowbreak{}05518981\allowbreak{}25318591\allowbreak{}40}
\RouteGroup{Separación de prefijos}
\RouteRow{Área de ambigüedad; área \(\log_2\)}{18; 0.0}
\RouteRow{Primer bloque de separación}{órbita: 1; celda: 1; ruta: 1}
\RouteRow{Ambigüedad final}{1}
\RouteGroup{Secuencias completas}
\RouteRow{\(w_6\)}{264 90 570 264 90 24 243 657 24 264 90 570 264 90 24 243 657 24}
\RouteRow{Tríadas estables por bloque}{0 1 2 3 4 5 6 7 8 9 10 11 12 13 14 15 15 16}
\end{tabularx}\endgroup\par\vspace{5pt}\noindent{\fontsize{8.4}{10}\selectfont\hyperref[trace-193]{Procedencia y códigos originales}\hfill Ficha completa · 193/324}
\clearpage\phantomsection\label{nav:vi-ruta-194}
% VI-CATALOG-ROW rutas 194
\pdfbookmark[3]{r6c4\_h-1v+1}{route-194}
\RouteTitle{Ruta 194}{r6c4\_\allowbreak{}h-1v+1}
\noindent Celda 49\quad (6, 4)\qquad Orientación \(h=-1,\ v=1\)\hfill\hyperref[sec:vi-indice-rutas]{Índice}\par\vspace{4pt}
\begingroup\fontsize{10.2}{12.5}\selectfont\setlength{\tabcolsep}{5pt}\renewcommand{\arraystretch}{1.1}\rowcolors{1}{routepale}{white}\noindent\begin{tabularx}{\linewidth}{@{}>{\raggedright\arraybackslash}p{.345\linewidth}>{\raggedright\arraybackslash}X@{}}
\RouteGroup{Identidad estructural}
\RouteRow{Clase y multisección}{quark\_\allowbreak{}up\_\allowbreak{}G1\quad /\quad quark\_\allowbreak{}up\_\allowbreak{}G1}
\RouteRow{Colección y sector}{quark de tipo arriba; quark}
\RouteRow{Clasificación física}{quark de tipo arriba}
\RouteRow{Generación, profundidad y color}{1; G1; B}
\RouteRow{Ruta primaria y órbita de inversión}{\(B_1\): no\qquad 4,\allowbreak{}6;\allowbreak{}6,\allowbreak{}4}
\RouteRow{Firma de clasificación}{28|\allowbreak{}6|\allowbreak{}-2|\allowbreak{}0|\allowbreak{}-12|\allowbreak{}+++-\kern0pt{}-\kern0pt{}-+++-\kern0pt{}-\kern0pt{}-}
\RouteGroup{Retornos, memoria y transporte}
\RouteRow{Retornos con fase}{clase: 27; órbita: 27; celda: 27}
\RouteRow{Retorno cinemático}{en 108: 54; extensión: 216}
\RouteRow{Giros y cambios de hoja}{71; 107}
\RouteRow{Acarreo y diferimiento}{deuda total: 2; final: 0; bloques diferidos: 1}
\RouteRow{Tríadas finales; persistencia de Pareto}{17; no}
\RouteRow{\(K\longrightarrow K+9\)}{repeticiones: 9; cambios de memoria: 9}
\RouteGroup{Firma, fase y diferencias}
\RouteRow{\(n_A,\ s_C,\ \kappa_0,\ \nu_{120},\ \nu_{270}\)}{50, 4, -2, 2, -8}
\RouteRow{\(\tau_{12}\); firma histórica \(\mathbb Z^5\)}{+++0-\kern0pt{}-+++0-\kern0pt{}-; 50|\allowbreak{}4|\allowbreak{}-2|\allowbreak{}2|\allowbreak{}-8}
\RouteRow{\(D_1,\ D_3,\ D_4,\ D_{27},\ D_{36}\)}{84, 24, 84, 24, 16}
\RouteRow{\(\Omega_{90/120},\ P_6\)}{56, 5}
\RouteGroup{Lectores y factores de acción}
\RouteRow{\(K_D\quad /\quad K_\Omega\)}{(40,\allowbreak{}84,\allowbreak{}30)\quad /\quad (56,\allowbreak{}54,\allowbreak{}5)}
\RouteRow{\(\kappa_D\)}{39.39035\allowbreak{}82768609\allowbreak{}24917849\allowbreak{}59}
\RouteRow{\(\kappa_\Omega\)}{55.60649\allowbreak{}89433721\allowbreak{}35446416\allowbreak{}86}
\RouteRow{\(R_{\mathrm{act},D}\)}{1.000000\allowbreak{}02730861\allowbreak{}73967087\allowbreak{}05}
\RouteRow{\(R_{\mathrm{act},\Omega}\)}{1.000000\allowbreak{}03855097\allowbreak{}23541416\allowbreak{}45}
\RouteGroup{Separación de prefijos}
\RouteRow{Área de ambigüedad; área \(\log_2\)}{18; 0.0}
\RouteRow{Primer bloque de separación}{órbita: 1; celda: 1; ruta: 1}
\RouteRow{Ambigüedad final}{1}
\RouteGroup{Secuencias completas}
\RouteRow{\(w_6\)}{258 414 420 258 414 255 252 333 255 258 414 420 258 414 255 252 333 255}
\RouteRow{Tríadas estables por bloque}{0 1 2 3 4 5 6 7 8 9 10 11 11 12 14 15 16 17}
\end{tabularx}\endgroup\par\vspace{5pt}\noindent{\fontsize{8.4}{10}\selectfont\hyperref[trace-194]{Procedencia y códigos originales}\hfill Ficha completa · 194/324}
\clearpage\phantomsection\label{nav:vi-ruta-195}
% VI-CATALOG-ROW rutas 195
\pdfbookmark[3]{r6c4\_h+1v-1}{route-195}
\RouteTitle{Ruta 195}{r6c4\_\allowbreak{}h+1v-1}
\noindent Celda 49\quad (6, 4)\qquad Orientación \(h=1,\ v=-1\)\hfill\hyperref[sec:vi-indice-rutas]{Índice}\par\vspace{4pt}
\begingroup\fontsize{10.2}{12.5}\selectfont\setlength{\tabcolsep}{5pt}\renewcommand{\arraystretch}{1.1}\rowcolors{1}{routepale}{white}\noindent\begin{tabularx}{\linewidth}{@{}>{\raggedright\arraybackslash}p{.345\linewidth}>{\raggedright\arraybackslash}X@{}}
\RouteGroup{Identidad estructural}
\RouteRow{Clase y multisección}{quark\_\allowbreak{}up\_\allowbreak{}G1\quad /\quad quark\_\allowbreak{}up\_\allowbreak{}G1}
\RouteRow{Colección y sector}{quark de tipo arriba; quark}
\RouteRow{Clasificación física}{quark de tipo arriba}
\RouteRow{Generación, profundidad y color}{1; G1; B}
\RouteRow{Ruta primaria y órbita de inversión}{\(B_1\): no\qquad 4,\allowbreak{}6;\allowbreak{}6,\allowbreak{}4}
\RouteRow{Firma de clasificación}{28|\allowbreak{}6|\allowbreak{}-2|\allowbreak{}0|\allowbreak{}-12|\allowbreak{}+++-\kern0pt{}-\kern0pt{}-+++-\kern0pt{}-\kern0pt{}-}
\RouteGroup{Retornos, memoria y transporte}
\RouteRow{Retornos con fase}{clase: 27; órbita: 27; celda: 27}
\RouteRow{Retorno cinemático}{en 108: 54; extensión: 216}
\RouteRow{Giros y cambios de hoja}{71; 107}
\RouteRow{Acarreo y diferimiento}{deuda total: 3; final: 1; bloques diferidos: 3}
\RouteRow{Tríadas finales; persistencia de Pareto}{16; no}
\RouteRow{\(K\longrightarrow K+9\)}{repeticiones: 9; cambios de memoria: 9}
\RouteGroup{Firma, fase y diferencias}
\RouteRow{\(n_A,\ s_C,\ \kappa_0,\ \nu_{120},\ \nu_{270}\)}{50, -4, 2, -2, -8}
\RouteRow{\(\tau_{12}\); firma histórica \(\mathbb Z^5\)}{0++-\kern0pt{}-\kern0pt{}-0++-\kern0pt{}-\kern0pt{}-; 50|\allowbreak{}-4|\allowbreak{}2|\allowbreak{}-2|\allowbreak{}-8}
\RouteRow{\(D_1,\ D_3,\ D_4,\ D_{27},\ D_{36}\)}{84, 24, 84, 24, 16}
\RouteRow{\(\Omega_{90/120},\ P_6\)}{56, 5}
\RouteGroup{Lectores y factores de acción}
\RouteRow{\(K_D\quad /\quad K_\Omega\)}{(40,\allowbreak{}84,\allowbreak{}30)\quad /\quad (56,\allowbreak{}54,\allowbreak{}5)}
\RouteRow{\(\kappa_D\)}{39.39035\allowbreak{}82768609\allowbreak{}24917849\allowbreak{}59}
\RouteRow{\(\kappa_\Omega\)}{55.60649\allowbreak{}89433721\allowbreak{}35446416\allowbreak{}86}
\RouteRow{\(R_{\mathrm{act},D}\)}{1.000000\allowbreak{}02730861\allowbreak{}73967087\allowbreak{}05}
\RouteRow{\(R_{\mathrm{act},\Omega}\)}{1.000000\allowbreak{}03855097\allowbreak{}23541416\allowbreak{}45}
\RouteGroup{Separación de prefijos}
\RouteRow{Área de ambigüedad; área \(\log_2\)}{18; 0.0}
\RouteRow{Primer bloque de separación}{órbita: 1; celda: 1; ruta: 1}
\RouteRow{Ambigüedad final}{1}
\RouteGroup{Secuencias completas}
\RouteRow{\(w_6\)}{252 333 255 252 333 420 258 414 420 252 333 255 252 333 420 258 414 420}
\RouteRow{Tríadas estables por bloque}{0 1 2 3 4 5 6 7 8 9 10 11 12 12 14 14 16 16}
\end{tabularx}\endgroup\par\vspace{5pt}\noindent{\fontsize{8.4}{10}\selectfont\hyperref[trace-195]{Procedencia y códigos originales}\hfill Ficha completa · 195/324}
\clearpage\phantomsection\label{nav:vi-ruta-196}
% VI-CATALOG-ROW rutas 196
\pdfbookmark[3]{r6c4\_h+1v+1}{route-196}
\RouteTitle{Ruta 196}{r6c4\_\allowbreak{}h+1v+1}
\noindent Celda 49\quad (6, 4)\qquad Orientación \(h=1,\ v=1\)\hfill\hyperref[sec:vi-indice-rutas]{Índice}\par\vspace{4pt}
\begingroup\fontsize{10.2}{12.5}\selectfont\setlength{\tabcolsep}{5pt}\renewcommand{\arraystretch}{1.1}\rowcolors{1}{routepale}{white}\noindent\begin{tabularx}{\linewidth}{@{}>{\raggedright\arraybackslash}p{.345\linewidth}>{\raggedright\arraybackslash}X@{}}
\RouteGroup{Identidad estructural}
\RouteRow{Clase y multisección}{quark\_\allowbreak{}up\_\allowbreak{}G1\quad /\quad quark\_\allowbreak{}up\_\allowbreak{}G1}
\RouteRow{Colección y sector}{quark de tipo arriba; quark}
\RouteRow{Clasificación física}{quark de tipo arriba}
\RouteRow{Generación, profundidad y color}{1; G1; B}
\RouteRow{Ruta primaria y órbita de inversión}{\(B_1\): no\qquad 4,\allowbreak{}6;\allowbreak{}6,\allowbreak{}4}
\RouteRow{Firma de clasificación}{28|\allowbreak{}6|\allowbreak{}-2|\allowbreak{}0|\allowbreak{}-12|\allowbreak{}+++-\kern0pt{}-\kern0pt{}-+++-\kern0pt{}-\kern0pt{}-}
\RouteGroup{Retornos, memoria y transporte}
\RouteRow{Retornos con fase}{clase: 27; órbita: 27; celda: 27}
\RouteRow{Retorno cinemático}{en 108: 54; extensión: 216}
\RouteRow{Giros y cambios de hoja}{71; 107}
\RouteRow{Acarreo y diferimiento}{deuda total: 0; final: 0; bloques diferidos: 0}
\RouteRow{Tríadas finales; persistencia de Pareto}{17; no}
\RouteRow{\(K\longrightarrow K+9\)}{repeticiones: 9; cambios de memoria: 9}
\RouteGroup{Firma, fase y diferencias}
\RouteRow{\(n_A,\ s_C,\ \kappa_0,\ \nu_{120},\ \nu_{270}\)}{28, 6, 0, 0, -12}
\RouteRow{\(\tau_{12}\); firma histórica \(\mathbb Z^5\)}{+++-\kern0pt{}-\kern0pt{}-+++-\kern0pt{}-\kern0pt{}-; 28|\allowbreak{}6|\allowbreak{}0|\allowbreak{}0|\allowbreak{}-12}
\RouteRow{\(D_1,\ D_3,\ D_4,\ D_{27},\ D_{36}\)}{66, 60, 66, 48, 32}
\RouteRow{\(\Omega_{90/120},\ P_6\)}{80, 6}
\RouteGroup{Lectores y factores de acción}
\RouteRow{\(K_D\quad /\quad K_\Omega\)}{(80,\allowbreak{}66,\allowbreak{}3)\quad /\quad (80,\allowbreak{}54,\allowbreak{}6)}
\RouteRow{\(\kappa_D\)}{79.51870\allowbreak{}83196953\allowbreak{}45554341\allowbreak{}34}
\RouteRow{\(\kappa_\Omega\)}{79.60661\allowbreak{}01397948\allowbreak{}27586470\allowbreak{}91}
\RouteRow{\(R_{\mathrm{act},D}\)}{1.000000\allowbreak{}05512887\allowbreak{}17997050\allowbreak{}72}
\RouteRow{\(R_{\mathrm{act},\Omega}\)}{1.000000\allowbreak{}05518981\allowbreak{}25318591\allowbreak{}40}
\RouteGroup{Separación de prefijos}
\RouteRow{Área de ambigüedad; área \(\log_2\)}{18; 0.0}
\RouteRow{Primer bloque de separación}{órbita: 1; celda: 1; ruta: 1}
\RouteRow{Ambigüedad final}{1}
\RouteGroup{Secuencias completas}
\RouteRow{\(w_6\)}{243 657 24 243 657 570 264 90 570 243 657 24 243 657 570 264 90 570}
\RouteRow{Tríadas estables por bloque}{0 1 2 3 4 5 6 7 8 9 10 11 12 13 14 15 16 17}
\end{tabularx}\endgroup\par\vspace{5pt}\noindent{\fontsize{8.4}{10}\selectfont\hyperref[trace-196]{Procedencia y códigos originales}\hfill Ficha completa · 196/324}
\clearpage\phantomsection\label{nav:vi-ruta-197}
% VI-CATALOG-ROW rutas 197
\pdfbookmark[2]{Celda 50}{cell-50}
\pdfbookmark[3]{r6c5\_h-1v-1}{route-197}
\RouteTitle{Ruta 197}{r6c5\_\allowbreak{}h-1v-1}
\noindent Celda 50\quad (6, 5)\qquad Orientación \(h=-1,\ v=-1\)\hfill\hyperref[sec:vi-indice-rutas]{Índice}\par\vspace{4pt}
\begingroup\fontsize{10.2}{12.5}\selectfont\setlength{\tabcolsep}{5pt}\renewcommand{\arraystretch}{1.1}\rowcolors{1}{routepale}{white}\noindent\begin{tabularx}{\linewidth}{@{}>{\raggedright\arraybackslash}p{.345\linewidth}>{\raggedright\arraybackslash}X@{}}
\RouteGroup{Identidad estructural}
\RouteRow{Clase y multisección}{quark\_\allowbreak{}down\_\allowbreak{}G1\quad /\quad quark\_\allowbreak{}down\_\allowbreak{}G1}
\RouteRow{Colección y sector}{quark de tipo abajo; quark}
\RouteRow{Clasificación física}{quark de tipo abajo}
\RouteRow{Generación, profundidad y color}{1; G1; G}
\RouteRow{Ruta primaria y órbita de inversión}{\(B_1\): no\qquad 4,\allowbreak{}5;\allowbreak{}6,\allowbreak{}5}
\RouteRow{Firma de clasificación}{28|\allowbreak{}-6|\allowbreak{}4|\allowbreak{}-2|\allowbreak{}-8|\allowbreak{}++0-\kern0pt{}-\kern0pt{}-++0-\kern0pt{}-\kern0pt{}-}
\RouteGroup{Retornos, memoria y transporte}
\RouteRow{Retornos con fase}{clase: 27; órbita: 27; celda: 27}
\RouteRow{Retorno cinemático}{en 108: 54; extensión: 216}
\RouteRow{Giros y cambios de hoja}{71; 107}
\RouteRow{Acarreo y diferimiento}{deuda total: 1; final: 0; bloques diferidos: 1}
\RouteRow{Tríadas finales; persistencia de Pareto}{17; sí}
\RouteRow{\(K\longrightarrow K+9\)}{repeticiones: 9; cambios de memoria: 9}
\RouteGroup{Firma, fase y diferencias}
\RouteRow{\(n_A,\ s_C,\ \kappa_0,\ \nu_{120},\ \nu_{270}\)}{28, 6, -2, 0, -8}
\RouteRow{\(\tau_{12}\); firma histórica \(\mathbb Z^5\)}{+++-\kern0pt{}-0+++-\kern0pt{}-0; 28|\allowbreak{}6|\allowbreak{}-2|\allowbreak{}0|\allowbreak{}-8}
\RouteRow{\(D_1,\ D_3,\ D_4,\ D_{27},\ D_{36}\)}{66, 60, 66, 48, 32}
\RouteRow{\(\Omega_{90/120},\ P_6\)}{56, 5}
\RouteGroup{Lectores y factores de acción}
\RouteRow{\(K_D\quad /\quad K_\Omega\)}{(80,\allowbreak{}66,\allowbreak{}3)\quad /\quad (56,\allowbreak{}54,\allowbreak{}5)}
\RouteRow{\(\kappa_D\)}{79.51870\allowbreak{}83196953\allowbreak{}45554341\allowbreak{}34}
\RouteRow{\(\kappa_\Omega\)}{55.60649\allowbreak{}89433721\allowbreak{}35446416\allowbreak{}86}
\RouteRow{\(R_{\mathrm{act},D}\)}{1.000000\allowbreak{}05512887\allowbreak{}17997050\allowbreak{}72}
\RouteRow{\(R_{\mathrm{act},\Omega}\)}{1.000000\allowbreak{}03855097\allowbreak{}23541416\allowbreak{}45}
\RouteGroup{Separación de prefijos}
\RouteRow{Área de ambigüedad; área \(\log_2\)}{18; 0.0}
\RouteRow{Primer bloque de separación}{órbita: 1; celda: 1; ruta: 1}
\RouteRow{Ambigüedad final}{1}
\RouteGroup{Secuencias completas}
\RouteRow{\(w_6\)}{486 423 15 486 423 327 498 99 327 486 423 15 486 423 327 498 99 327}
\RouteRow{Tríadas estables por bloque}{0 1 2 3 4 5 6 7 8 9 10 10 12 13 14 15 16 17}
\end{tabularx}\endgroup\par\vspace{5pt}\noindent{\fontsize{8.4}{10}\selectfont\hyperref[trace-197]{Procedencia y códigos originales}\hfill Ficha completa · 197/324}
\clearpage\phantomsection\label{nav:vi-ruta-198}
% VI-CATALOG-ROW rutas 198
\pdfbookmark[3]{r6c5\_h-1v+1}{route-198}
\RouteTitle{Ruta 198}{r6c5\_\allowbreak{}h-1v+1}
\noindent Celda 50\quad (6, 5)\qquad Orientación \(h=-1,\ v=1\)\hfill\hyperref[sec:vi-indice-rutas]{Índice}\par\vspace{4pt}
\begingroup\fontsize{10.2}{12.5}\selectfont\setlength{\tabcolsep}{5pt}\renewcommand{\arraystretch}{1.1}\rowcolors{1}{routepale}{white}\noindent\begin{tabularx}{\linewidth}{@{}>{\raggedright\arraybackslash}p{.345\linewidth}>{\raggedright\arraybackslash}X@{}}
\RouteGroup{Identidad estructural}
\RouteRow{Clase y multisección}{quark\_\allowbreak{}down\_\allowbreak{}G1\quad /\quad quark\_\allowbreak{}down\_\allowbreak{}G1}
\RouteRow{Colección y sector}{quark de tipo abajo; quark}
\RouteRow{Clasificación física}{quark de tipo abajo}
\RouteRow{Generación, profundidad y color}{1; G1; G}
\RouteRow{Ruta primaria y órbita de inversión}{\(B_1\): sí\qquad 4,\allowbreak{}5;\allowbreak{}6,\allowbreak{}5}
\RouteRow{Firma de clasificación}{28|\allowbreak{}-6|\allowbreak{}4|\allowbreak{}-2|\allowbreak{}-8|\allowbreak{}++0-\kern0pt{}-\kern0pt{}-++0-\kern0pt{}-\kern0pt{}-}
\RouteGroup{Retornos, memoria y transporte}
\RouteRow{Retornos con fase}{clase: 27; órbita: 27; celda: 27}
\RouteRow{Retorno cinemático}{en 108: 54; extensión: 216}
\RouteRow{Giros y cambios de hoja}{71; 107}
\RouteRow{Acarreo y diferimiento}{deuda total: 1; final: 0; bloques diferidos: 1}
\RouteRow{Tríadas finales; persistencia de Pareto}{17; sí}
\RouteRow{\(K\longrightarrow K+9\)}{repeticiones: 9; cambios de memoria: 9}
\RouteGroup{Firma, fase y diferencias}
\RouteRow{\(n_A,\ s_C,\ \kappa_0,\ \nu_{120},\ \nu_{270}\)}{14, -4, 0, 0, -12}
\RouteRow{\(\tau_{12}\); firma histórica \(\mathbb Z^5\)}{+++-\kern0pt{}-\kern0pt{}-+++-\kern0pt{}-\kern0pt{}-; 14|\allowbreak{}-4|\allowbreak{}0|\allowbreak{}0|\allowbreak{}-12}
\RouteRow{\(D_1,\ D_3,\ D_4,\ D_{27},\ D_{36}\)}{78, 24, 78, 24, 16}
\RouteRow{\(\Omega_{90/120},\ P_6\)}{80, 6}
\RouteGroup{Lectores y factores de acción}
\RouteRow{\(K_D\quad /\quad K_\Omega\)}{(40,\allowbreak{}78,\allowbreak{}27)\quad /\quad (80,\allowbreak{}54,\allowbreak{}6)}
\RouteRow{\(\kappa_D\)}{39.43380\allowbreak{}88030085\allowbreak{}51303671\allowbreak{}14}
\RouteRow{\(\kappa_\Omega\)}{79.60661\allowbreak{}01397948\allowbreak{}27586470\allowbreak{}91}
\RouteRow{\(R_{\mathrm{act},D}\)}{1.000000\allowbreak{}02733874\allowbreak{}08548943\allowbreak{}58}
\RouteRow{\(R_{\mathrm{act},\Omega}\)}{1.000000\allowbreak{}05518981\allowbreak{}25318591\allowbreak{}40}
\RouteGroup{Separación de prefijos}
\RouteRow{Área de ambigüedad; área \(\log_2\)}{18; 0.0}
\RouteRow{Primer bloque de separación}{órbita: 1; celda: 1; ruta: 1}
\RouteRow{Ambigüedad final}{1}
\RouteGroup{Secuencias completas}
\RouteRow{\(w_6\)}{504 585 507 504 585 672 510 666 672 504 585 507 504 585 672 510 666 672}
\RouteRow{Tríadas estables por bloque}{0 1 2 3 4 5 6 7 8 9 10 11 12 13 14 14 16 17}
\end{tabularx}\endgroup\par\vspace{5pt}\noindent{\fontsize{8.4}{10}\selectfont\hyperref[trace-198]{Procedencia y códigos originales}\hfill Ficha completa · 198/324}
\clearpage\phantomsection\label{nav:vi-ruta-199}
% VI-CATALOG-ROW rutas 199
\pdfbookmark[3]{r6c5\_h+1v-1}{route-199}
\RouteTitle{Ruta 199}{r6c5\_\allowbreak{}h+1v-1}
\noindent Celda 50\quad (6, 5)\qquad Orientación \(h=1,\ v=-1\)\hfill\hyperref[sec:vi-indice-rutas]{Índice}\par\vspace{4pt}
\begingroup\fontsize{10.2}{12.5}\selectfont\setlength{\tabcolsep}{5pt}\renewcommand{\arraystretch}{1.1}\rowcolors{1}{routepale}{white}\noindent\begin{tabularx}{\linewidth}{@{}>{\raggedright\arraybackslash}p{.345\linewidth}>{\raggedright\arraybackslash}X@{}}
\RouteGroup{Identidad estructural}
\RouteRow{Clase y multisección}{quark\_\allowbreak{}down\_\allowbreak{}G1\quad /\quad quark\_\allowbreak{}down\_\allowbreak{}G1}
\RouteRow{Colección y sector}{quark de tipo abajo; quark}
\RouteRow{Clasificación física}{quark de tipo abajo}
\RouteRow{Generación, profundidad y color}{1; G1; G}
\RouteRow{Ruta primaria y órbita de inversión}{\(B_1\): no\qquad 4,\allowbreak{}5;\allowbreak{}6,\allowbreak{}5}
\RouteRow{Firma de clasificación}{28|\allowbreak{}-6|\allowbreak{}4|\allowbreak{}-2|\allowbreak{}-8|\allowbreak{}++0-\kern0pt{}-\kern0pt{}-++0-\kern0pt{}-\kern0pt{}-}
\RouteGroup{Retornos, memoria y transporte}
\RouteRow{Retornos con fase}{clase: 27; órbita: 27; celda: 27}
\RouteRow{Retorno cinemático}{en 108: 54; extensión: 216}
\RouteRow{Giros y cambios de hoja}{71; 107}
\RouteRow{Acarreo y diferimiento}{deuda total: 1; final: 0; bloques diferidos: 1}
\RouteRow{Tríadas finales; persistencia de Pareto}{17; sí}
\RouteRow{\(K\longrightarrow K+9\)}{repeticiones: 9; cambios de memoria: 9}
\RouteGroup{Firma, fase y diferencias}
\RouteRow{\(n_A,\ s_C,\ \kappa_0,\ \nu_{120},\ \nu_{270}\)}{14, 4, 0, 0, -12}
\RouteRow{\(\tau_{12}\); firma histórica \(\mathbb Z^5\)}{+++-\kern0pt{}-\kern0pt{}-+++-\kern0pt{}-\kern0pt{}-; 14|\allowbreak{}4|\allowbreak{}0|\allowbreak{}0|\allowbreak{}-12}
\RouteRow{\(D_1,\ D_3,\ D_4,\ D_{27},\ D_{36}\)}{78, 24, 78, 24, 16}
\RouteRow{\(\Omega_{90/120},\ P_6\)}{80, 6}
\RouteGroup{Lectores y factores de acción}
\RouteRow{\(K_D\quad /\quad K_\Omega\)}{(40,\allowbreak{}78,\allowbreak{}27)\quad /\quad (80,\allowbreak{}54,\allowbreak{}6)}
\RouteRow{\(\kappa_D\)}{39.43380\allowbreak{}88030085\allowbreak{}51303671\allowbreak{}14}
\RouteRow{\(\kappa_\Omega\)}{79.60661\allowbreak{}01397948\allowbreak{}27586470\allowbreak{}91}
\RouteRow{\(R_{\mathrm{act},D}\)}{1.000000\allowbreak{}02733874\allowbreak{}08548943\allowbreak{}58}
\RouteRow{\(R_{\mathrm{act},\Omega}\)}{1.000000\allowbreak{}05518981\allowbreak{}25318591\allowbreak{}40}
\RouteGroup{Separación de prefijos}
\RouteRow{Área de ambigüedad; área \(\log_2\)}{18; 0.0}
\RouteRow{Primer bloque de separación}{órbita: 1; celda: 1; ruta: 1}
\RouteRow{Ambigüedad final}{1}
\RouteGroup{Secuencias completas}
\RouteRow{\(w_6\)}{510 666 672 510 666 507 504 585 507 510 666 672 510 666 507 504 585 507}
\RouteRow{Tríadas estables por bloque}{0 1 2 3 4 5 6 7 8 9 10 11 12 13 14 15 15 17}
\end{tabularx}\endgroup\par\vspace{5pt}\noindent{\fontsize{8.4}{10}\selectfont\hyperref[trace-199]{Procedencia y códigos originales}\hfill Ficha completa · 199/324}
\clearpage\phantomsection\label{nav:vi-ruta-200}
% VI-CATALOG-ROW rutas 200
\pdfbookmark[3]{r6c5\_h+1v+1}{route-200}
\RouteTitle{Ruta 200}{r6c5\_\allowbreak{}h+1v+1}
\noindent Celda 50\quad (6, 5)\qquad Orientación \(h=1,\ v=1\)\hfill\hyperref[sec:vi-indice-rutas]{Índice}\par\vspace{4pt}
\begingroup\fontsize{10.2}{12.5}\selectfont\setlength{\tabcolsep}{5pt}\renewcommand{\arraystretch}{1.1}\rowcolors{1}{routepale}{white}\noindent\begin{tabularx}{\linewidth}{@{}>{\raggedright\arraybackslash}p{.345\linewidth}>{\raggedright\arraybackslash}X@{}}
\RouteGroup{Identidad estructural}
\RouteRow{Clase y multisección}{quark\_\allowbreak{}down\_\allowbreak{}G1\quad /\quad quark\_\allowbreak{}down\_\allowbreak{}G1}
\RouteRow{Colección y sector}{quark de tipo abajo; quark}
\RouteRow{Clasificación física}{quark de tipo abajo}
\RouteRow{Generación, profundidad y color}{1; G1; G}
\RouteRow{Ruta primaria y órbita de inversión}{\(B_1\): no\qquad 4,\allowbreak{}5;\allowbreak{}6,\allowbreak{}5}
\RouteRow{Firma de clasificación}{28|\allowbreak{}-6|\allowbreak{}4|\allowbreak{}-2|\allowbreak{}-8|\allowbreak{}++0-\kern0pt{}-\kern0pt{}-++0-\kern0pt{}-\kern0pt{}-}
\RouteGroup{Retornos, memoria y transporte}
\RouteRow{Retornos con fase}{clase: 27; órbita: 27; celda: 27}
\RouteRow{Retorno cinemático}{en 108: 54; extensión: 216}
\RouteRow{Giros y cambios de hoja}{71; 107}
\RouteRow{Acarreo y diferimiento}{deuda total: 0; final: 0; bloques diferidos: 0}
\RouteRow{Tríadas finales; persistencia de Pareto}{17; sí}
\RouteRow{\(K\longrightarrow K+9\)}{repeticiones: 9; cambios de memoria: 9}
\RouteGroup{Firma, fase y diferencias}
\RouteRow{\(n_A,\ s_C,\ \kappa_0,\ \nu_{120},\ \nu_{270}\)}{28, -6, 2, 0, -8}
\RouteRow{\(\tau_{12}\); firma histórica \(\mathbb Z^5\)}{++0-\kern0pt{}-\kern0pt{}-++0-\kern0pt{}-\kern0pt{}-; 28|\allowbreak{}-6|\allowbreak{}2|\allowbreak{}0|\allowbreak{}-8}
\RouteRow{\(D_1,\ D_3,\ D_4,\ D_{27},\ D_{36}\)}{66, 60, 66, 48, 32}
\RouteRow{\(\Omega_{90/120},\ P_6\)}{56, 5}
\RouteGroup{Lectores y factores de acción}
\RouteRow{\(K_D\quad /\quad K_\Omega\)}{(80,\allowbreak{}66,\allowbreak{}3)\quad /\quad (56,\allowbreak{}54,\allowbreak{}5)}
\RouteRow{\(\kappa_D\)}{79.51870\allowbreak{}83196953\allowbreak{}45554341\allowbreak{}34}
\RouteRow{\(\kappa_\Omega\)}{55.60649\allowbreak{}89433721\allowbreak{}35446416\allowbreak{}86}
\RouteRow{\(R_{\mathrm{act},D}\)}{1.000000\allowbreak{}05512887\allowbreak{}17997050\allowbreak{}72}
\RouteRow{\(R_{\mathrm{act},\Omega}\)}{1.000000\allowbreak{}03855097\allowbreak{}23541416\allowbreak{}45}
\RouteGroup{Separación de prefijos}
\RouteRow{Área de ambigüedad; área \(\log_2\)}{18; 0.0}
\RouteRow{Primer bloque de separación}{órbita: 1; celda: 1; ruta: 1}
\RouteRow{Ambigüedad final}{1}
\RouteGroup{Secuencias completas}
\RouteRow{\(w_6\)}{498 99 327 498 99 15 486 423 15 498 99 327 498 99 15 486 423 15}
\RouteRow{Tríadas estables por bloque}{0 1 2 3 4 5 6 7 8 9 10 11 12 13 14 15 16 17}
\end{tabularx}\endgroup\par\vspace{5pt}\noindent{\fontsize{8.4}{10}\selectfont\hyperref[trace-200]{Procedencia y códigos originales}\hfill Ficha completa · 200/324}
\clearpage\phantomsection\label{nav:vi-ruta-201}
% VI-CATALOG-ROW rutas 201
\pdfbookmark[2]{Celda 51}{cell-51}
\pdfbookmark[3]{r6c6\_h-1v-1}{route-201}
\RouteTitle{Ruta 201}{r6c6\_\allowbreak{}h-1v-1}
\noindent Celda 51\quad (6, 6)\qquad Orientación \(h=-1,\ v=-1\)\hfill\hyperref[sec:vi-indice-rutas]{Índice}\par\vspace{4pt}
\begingroup\fontsize{10.2}{12.5}\selectfont\setlength{\tabcolsep}{5pt}\renewcommand{\arraystretch}{1.1}\rowcolors{1}{routepale}{white}\noindent\begin{tabularx}{\linewidth}{@{}>{\raggedright\arraybackslash}p{.345\linewidth}>{\raggedright\arraybackslash}X@{}}
\RouteGroup{Identidad estructural}
\RouteRow{Clase y multisección}{quark\_\allowbreak{}up\_\allowbreak{}G1\quad /\quad quark\_\allowbreak{}up\_\allowbreak{}G1}
\RouteRow{Colección y sector}{quark de tipo arriba; quark}
\RouteRow{Clasificación física}{quark de tipo arriba}
\RouteRow{Generación, profundidad y color}{1; G1; R}
\RouteRow{Ruta primaria y órbita de inversión}{\(B_1\): no\qquad 4,\allowbreak{}4;\allowbreak{}6,\allowbreak{}6}
\RouteRow{Firma de clasificación}{24|\allowbreak{}0|\allowbreak{}12|\allowbreak{}0|\allowbreak{}-12|\allowbreak{}+++-\kern0pt{}-\kern0pt{}-+++-\kern0pt{}-\kern0pt{}-}
\RouteGroup{Retornos, memoria y transporte}
\RouteRow{Retornos con fase}{clase: 27; órbita: 27; celda: 27}
\RouteRow{Retorno cinemático}{en 108: 54; extensión: 216}
\RouteRow{Giros y cambios de hoja}{71; 107}
\RouteRow{Acarreo y diferimiento}{deuda total: 1; final: 0; bloques diferidos: 1}
\RouteRow{Tríadas finales; persistencia de Pareto}{17; sí}
\RouteRow{\(K\longrightarrow K+9\)}{repeticiones: 9; cambios de memoria: 9}
\RouteGroup{Firma, fase y diferencias}
\RouteRow{\(n_A,\ s_C,\ \kappa_0,\ \nu_{120},\ \nu_{270}\)}{24, 0, 0, 0, -12}
\RouteRow{\(\tau_{12}\); firma histórica \(\mathbb Z^5\)}{+++-\kern0pt{}-\kern0pt{}-+++-\kern0pt{}-\kern0pt{}-; 24|\allowbreak{}0|\allowbreak{}0|\allowbreak{}0|\allowbreak{}-12}
\RouteRow{\(D_1,\ D_3,\ D_4,\ D_{27},\ D_{36}\)}{54, 60, 72, 24, 16}
\RouteRow{\(\Omega_{90/120},\ P_6\)}{80, 6}
\RouteGroup{Lectores y factores de acción}
\RouteRow{\(K_D\quad /\quad K_\Omega\)}{(40,\allowbreak{}54,\allowbreak{}6)\quad /\quad (80,\allowbreak{}54,\allowbreak{}6)}
\RouteRow{\(\kappa_D\)}{39.60661\allowbreak{}01397948\allowbreak{}27586470\allowbreak{}91}
\RouteRow{\(\kappa_\Omega\)}{79.60661\allowbreak{}01397948\allowbreak{}27586470\allowbreak{}91}
\RouteRow{\(R_{\mathrm{act},D}\)}{1.000000\allowbreak{}02745854\allowbreak{}08735595\allowbreak{}18}
\RouteRow{\(R_{\mathrm{act},\Omega}\)}{1.000000\allowbreak{}05518981\allowbreak{}25318591\allowbreak{}40}
\RouteGroup{Separación de prefijos}
\RouteRow{Área de ambigüedad; área \(\log_2\)}{18; 0.0}
\RouteRow{Primer bloque de separación}{órbita: 1; celda: 1; ruta: 1}
\RouteRow{Ambigüedad final}{1}
\RouteGroup{Secuencias completas}
\RouteRow{\(w_6\)}{15 486 423 15 486 657 24 243 657 15 486 423 15 486 657 24 243 657}
\RouteRow{Tríadas estables por bloque}{0 1 2 3 4 5 6 7 8 9 10 11 12 13 14 15 15 17}
\end{tabularx}\endgroup\par\vspace{5pt}\noindent{\fontsize{8.4}{10}\selectfont\hyperref[trace-201]{Procedencia y códigos originales}\hfill Ficha completa · 201/324}
\clearpage\phantomsection\label{nav:vi-ruta-202}
% VI-CATALOG-ROW rutas 202
\pdfbookmark[3]{r6c6\_h-1v+1}{route-202}
\RouteTitle{Ruta 202}{r6c6\_\allowbreak{}h-1v+1}
\noindent Celda 51\quad (6, 6)\qquad Orientación \(h=-1,\ v=1\)\hfill\hyperref[sec:vi-indice-rutas]{Índice}\par\vspace{4pt}
\begingroup\fontsize{10.2}{12.5}\selectfont\setlength{\tabcolsep}{5pt}\renewcommand{\arraystretch}{1.1}\rowcolors{1}{routepale}{white}\noindent\begin{tabularx}{\linewidth}{@{}>{\raggedright\arraybackslash}p{.345\linewidth}>{\raggedright\arraybackslash}X@{}}
\RouteGroup{Identidad estructural}
\RouteRow{Clase y multisección}{quark\_\allowbreak{}up\_\allowbreak{}G1\quad /\quad quark\_\allowbreak{}up\_\allowbreak{}G1}
\RouteRow{Colección y sector}{quark de tipo arriba; quark}
\RouteRow{Clasificación física}{quark de tipo arriba}
\RouteRow{Generación, profundidad y color}{1; G1; R}
\RouteRow{Ruta primaria y órbita de inversión}{\(B_1\): no\qquad 4,\allowbreak{}4;\allowbreak{}6,\allowbreak{}6}
\RouteRow{Firma de clasificación}{24|\allowbreak{}0|\allowbreak{}12|\allowbreak{}0|\allowbreak{}-12|\allowbreak{}+++-\kern0pt{}-\kern0pt{}-+++-\kern0pt{}-\kern0pt{}-}
\RouteGroup{Retornos, memoria y transporte}
\RouteRow{Retornos con fase}{clase: 27; órbita: 27; celda: 27}
\RouteRow{Retorno cinemático}{en 108: 54; extensión: 216}
\RouteRow{Giros y cambios de hoja}{71; 107}
\RouteRow{Acarreo y diferimiento}{deuda total: 0; final: 0; bloques diferidos: 0}
\RouteRow{Tríadas finales; persistencia de Pareto}{17; sí}
\RouteRow{\(K\longrightarrow K+9\)}{repeticiones: 9; cambios de memoria: 9}
\RouteGroup{Firma, fase y diferencias}
\RouteRow{\(n_A,\ s_C,\ \kappa_0,\ \nu_{120},\ \nu_{270}\)}{50, -6, 2, 0, -8}
\RouteRow{\(\tau_{12}\); firma histórica \(\mathbb Z^5\)}{+0+-\kern0pt{}-\kern0pt{}-+0+-\kern0pt{}-\kern0pt{}-; 50|\allowbreak{}-6|\allowbreak{}2|\allowbreak{}0|\allowbreak{}-8}
\RouteRow{\(D_1,\ D_3,\ D_4,\ D_{27},\ D_{36}\)}{24, 24, 24, 0, 0}
\RouteRow{\(\Omega_{90/120},\ P_6\)}{48, 5}
\RouteGroup{Lectores y factores de acción}
\RouteRow{\(K_D\quad /\quad K_\Omega\)}{(0,\allowbreak{}24,\allowbreak{}0)\quad /\quad (48,\allowbreak{}54,\allowbreak{}5)}
\RouteRow{\(\kappa_D\)}{-0.175136\allowbreak{}46166281\allowbreak{}12239348\allowbreak{}425}
\RouteRow{\(\kappa_\Omega\)}{47.60649\allowbreak{}89433721\allowbreak{}35446416\allowbreak{}86}
\RouteRow{\(R_{\mathrm{act},D}\)}{0.999999\allowbreak{}99987858\allowbreak{}10851337\allowbreak{}881}
\RouteRow{\(R_{\mathrm{act},\Omega}\)}{1.000000\allowbreak{}03300471\allowbreak{}80532430\allowbreak{}84}
\RouteGroup{Separación de prefijos}
\RouteRow{Área de ambigüedad; área \(\log_2\)}{36; 18.0}
\RouteRow{Primer bloque de separación}{órbita: 1; celda: 1; ruta: \textemdash{}}
\RouteRow{Ambigüedad final}{2}
\RouteGroup{Secuencias completas}
\RouteRow{\(w_6\)}{3 0 81 3 0 81 3 0 81 3 0 81 3 0 81 3 0 81}
\RouteRow{Tríadas estables por bloque}{0 1 2 3 4 5 6 7 8 9 10 11 12 13 14 15 16 17}
\end{tabularx}\endgroup\par\vspace{5pt}\noindent{\fontsize{8.4}{10}\selectfont\hyperref[trace-202]{Procedencia y códigos originales}\hfill Ficha completa · 202/324}
\clearpage\phantomsection\label{nav:vi-ruta-203}
% VI-CATALOG-ROW rutas 203
\pdfbookmark[3]{r6c6\_h+1v-1}{route-203}
\RouteTitle{Ruta 203}{r6c6\_\allowbreak{}h+1v-1}
\noindent Celda 51\quad (6, 6)\qquad Orientación \(h=1,\ v=-1\)\hfill\hyperref[sec:vi-indice-rutas]{Índice}\par\vspace{4pt}
\begingroup\fontsize{10.2}{12.5}\selectfont\setlength{\tabcolsep}{5pt}\renewcommand{\arraystretch}{1.1}\rowcolors{1}{routepale}{white}\noindent\begin{tabularx}{\linewidth}{@{}>{\raggedright\arraybackslash}p{.345\linewidth}>{\raggedright\arraybackslash}X@{}}
\RouteGroup{Identidad estructural}
\RouteRow{Clase y multisección}{quark\_\allowbreak{}up\_\allowbreak{}G1\quad /\quad quark\_\allowbreak{}up\_\allowbreak{}G1}
\RouteRow{Colección y sector}{quark de tipo arriba; quark}
\RouteRow{Clasificación física}{quark de tipo arriba}
\RouteRow{Generación, profundidad y color}{1; G1; R}
\RouteRow{Ruta primaria y órbita de inversión}{\(B_1\): no\qquad 4,\allowbreak{}4;\allowbreak{}6,\allowbreak{}6}
\RouteRow{Firma de clasificación}{24|\allowbreak{}0|\allowbreak{}12|\allowbreak{}0|\allowbreak{}-12|\allowbreak{}+++-\kern0pt{}-\kern0pt{}-+++-\kern0pt{}-\kern0pt{}-}
\RouteGroup{Retornos, memoria y transporte}
\RouteRow{Retornos con fase}{clase: 27; órbita: 27; celda: 27}
\RouteRow{Retorno cinemático}{en 108: 54; extensión: 216}
\RouteRow{Giros y cambios de hoja}{71; 107}
\RouteRow{Acarreo y diferimiento}{deuda total: 0; final: 0; bloques diferidos: 0}
\RouteRow{Tríadas finales; persistencia de Pareto}{17; sí}
\RouteRow{\(K\longrightarrow K+9\)}{repeticiones: 9; cambios de memoria: 9}
\RouteGroup{Firma, fase y diferencias}
\RouteRow{\(n_A,\ s_C,\ \kappa_0,\ \nu_{120},\ \nu_{270}\)}{50, 6, -2, 0, -8}
\RouteRow{\(\tau_{12}\); firma histórica \(\mathbb Z^5\)}{+++-0-+++-0-; 50|\allowbreak{}6|\allowbreak{}-2|\allowbreak{}0|\allowbreak{}-8}
\RouteRow{\(D_1,\ D_3,\ D_4,\ D_{27},\ D_{36}\)}{24, 24, 24, 0, 0}
\RouteRow{\(\Omega_{90/120},\ P_6\)}{48, 5}
\RouteGroup{Lectores y factores de acción}
\RouteRow{\(K_D\quad /\quad K_\Omega\)}{(0,\allowbreak{}24,\allowbreak{}0)\quad /\quad (48,\allowbreak{}54,\allowbreak{}5)}
\RouteRow{\(\kappa_D\)}{-0.175136\allowbreak{}46166281\allowbreak{}12239348\allowbreak{}425}
\RouteRow{\(\kappa_\Omega\)}{47.60649\allowbreak{}89433721\allowbreak{}35446416\allowbreak{}86}
\RouteRow{\(R_{\mathrm{act},D}\)}{0.999999\allowbreak{}99987858\allowbreak{}10851337\allowbreak{}881}
\RouteRow{\(R_{\mathrm{act},\Omega}\)}{1.000000\allowbreak{}03300471\allowbreak{}80532430\allowbreak{}84}
\RouteGroup{Separación de prefijos}
\RouteRow{Área de ambigüedad; área \(\log_2\)}{36; 18.0}
\RouteRow{Primer bloque de separación}{órbita: 1; celda: 1; ruta: \textemdash{}}
\RouteRow{Ambigüedad final}{2}
\RouteGroup{Secuencias completas}
\RouteRow{\(w_6\)}{3 0 81 3 0 81 3 0 81 3 0 81 3 0 81 3 0 81}
\RouteRow{Tríadas estables por bloque}{0 1 2 3 4 5 6 7 8 9 10 11 12 13 14 15 16 17}
\end{tabularx}\endgroup\par\vspace{5pt}\noindent{\fontsize{8.4}{10}\selectfont\hyperref[trace-203]{Procedencia y códigos originales}\hfill Ficha completa · 203/324}
\clearpage\phantomsection\label{nav:vi-ruta-204}
% VI-CATALOG-ROW rutas 204
\pdfbookmark[3]{r6c6\_h+1v+1}{route-204}
\RouteTitle{Ruta 204}{r6c6\_\allowbreak{}h+1v+1}
\noindent Celda 51\quad (6, 6)\qquad Orientación \(h=1,\ v=1\)\hfill\hyperref[sec:vi-indice-rutas]{Índice}\par\vspace{4pt}
\begingroup\fontsize{10.2}{12.5}\selectfont\setlength{\tabcolsep}{5pt}\renewcommand{\arraystretch}{1.1}\rowcolors{1}{routepale}{white}\noindent\begin{tabularx}{\linewidth}{@{}>{\raggedright\arraybackslash}p{.345\linewidth}>{\raggedright\arraybackslash}X@{}}
\RouteGroup{Identidad estructural}
\RouteRow{Clase y multisección}{quark\_\allowbreak{}up\_\allowbreak{}G1\quad /\quad quark\_\allowbreak{}up\_\allowbreak{}G1}
\RouteRow{Colección y sector}{quark de tipo arriba; quark}
\RouteRow{Clasificación física}{quark de tipo arriba}
\RouteRow{Generación, profundidad y color}{1; G1; R}
\RouteRow{Ruta primaria y órbita de inversión}{\(B_1\): sí\qquad 4,\allowbreak{}4;\allowbreak{}6,\allowbreak{}6}
\RouteRow{Firma de clasificación}{24|\allowbreak{}0|\allowbreak{}12|\allowbreak{}0|\allowbreak{}-12|\allowbreak{}+++-\kern0pt{}-\kern0pt{}-+++-\kern0pt{}-\kern0pt{}-}
\RouteGroup{Retornos, memoria y transporte}
\RouteRow{Retornos con fase}{clase: 27; órbita: 27; celda: 27}
\RouteRow{Retorno cinemático}{en 108: 54; extensión: 216}
\RouteRow{Giros y cambios de hoja}{71; 107}
\RouteRow{Acarreo y diferimiento}{deuda total: 2; final: 1; bloques diferidos: 2}
\RouteRow{Tríadas finales; persistencia de Pareto}{16; sí}
\RouteRow{\(K\longrightarrow K+9\)}{repeticiones: 9; cambios de memoria: 9}
\RouteGroup{Firma, fase y diferencias}
\RouteRow{\(n_A,\ s_C,\ \kappa_0,\ \nu_{120},\ \nu_{270}\)}{24, 0, 0, 0, -12}
\RouteRow{\(\tau_{12}\); firma histórica \(\mathbb Z^5\)}{+++-\kern0pt{}-\kern0pt{}-+++-\kern0pt{}-\kern0pt{}-; 24|\allowbreak{}0|\allowbreak{}0|\allowbreak{}0|\allowbreak{}-12}
\RouteRow{\(D_1,\ D_3,\ D_4,\ D_{27},\ D_{36}\)}{54, 60, 72, 24, 16}
\RouteRow{\(\Omega_{90/120},\ P_6\)}{80, 6}
\RouteGroup{Lectores y factores de acción}
\RouteRow{\(K_D\quad /\quad K_\Omega\)}{(40,\allowbreak{}54,\allowbreak{}6)\quad /\quad (80,\allowbreak{}54,\allowbreak{}6)}
\RouteRow{\(\kappa_D\)}{39.60661\allowbreak{}01397948\allowbreak{}27586470\allowbreak{}91}
\RouteRow{\(\kappa_\Omega\)}{79.60661\allowbreak{}01397948\allowbreak{}27586470\allowbreak{}91}
\RouteRow{\(R_{\mathrm{act},D}\)}{1.000000\allowbreak{}02745854\allowbreak{}08735595\allowbreak{}18}
\RouteRow{\(R_{\mathrm{act},\Omega}\)}{1.000000\allowbreak{}05518981\allowbreak{}25318591\allowbreak{}40}
\RouteGroup{Separación de prefijos}
\RouteRow{Área de ambigüedad; área \(\log_2\)}{18; 0.0}
\RouteRow{Primer bloque de separación}{órbita: 1; celda: 1; ruta: 1}
\RouteRow{Ambigüedad final}{1}
\RouteGroup{Secuencias completas}
\RouteRow{\(w_6\)}{24 243 657 24 243 423 15 486 423 24 243 657 24 243 423 15 486 423}
\RouteRow{Tríadas estables por bloque}{0 1 2 3 4 5 6 7 8 9 10 11 12 13 13 15 16 16}
\end{tabularx}\endgroup\par\vspace{5pt}\noindent{\fontsize{8.4}{10}\selectfont\hyperref[trace-204]{Procedencia y códigos originales}\hfill Ficha completa · 204/324}
\clearpage\phantomsection\label{nav:vi-ruta-205}
% VI-CATALOG-ROW rutas 205
\pdfbookmark[2]{Celda 52}{cell-52}
\pdfbookmark[3]{r6c7\_h-1v-1}{route-205}
\RouteTitle{Ruta 205}{r6c7\_\allowbreak{}h-1v-1}
\noindent Celda 52\quad (6, 7)\qquad Orientación \(h=-1,\ v=-1\)\hfill\hyperref[sec:vi-indice-rutas]{Índice}\par\vspace{4pt}
\begingroup\fontsize{10.2}{12.5}\selectfont\setlength{\tabcolsep}{5pt}\renewcommand{\arraystretch}{1.1}\rowcolors{1}{routepale}{white}\noindent\begin{tabularx}{\linewidth}{@{}>{\raggedright\arraybackslash}p{.345\linewidth}>{\raggedright\arraybackslash}X@{}}
\RouteGroup{Identidad estructural}
\RouteRow{Clase y multisección}{quark\_\allowbreak{}down\_\allowbreak{}G2\quad /\quad quark\_\allowbreak{}down\_\allowbreak{}G2}
\RouteRow{Colección y sector}{quark de tipo abajo; quark}
\RouteRow{Clasificación física}{quark de tipo abajo}
\RouteRow{Generación, profundidad y color}{2; G2; B}
\RouteRow{Ruta primaria y órbita de inversión}{\(B_1\): no\qquad 4,\allowbreak{}3;\allowbreak{}6,\allowbreak{}7}
\RouteRow{Firma de clasificación}{28|\allowbreak{}6|\allowbreak{}-8|\allowbreak{}0|\allowbreak{}-12|\allowbreak{}+++-\kern0pt{}-\kern0pt{}-+++-\kern0pt{}-\kern0pt{}-}
\RouteGroup{Retornos, memoria y transporte}
\RouteRow{Retornos con fase}{clase: 27; órbita: 27; celda: 27}
\RouteRow{Retorno cinemático}{en 108: 54; extensión: 216}
\RouteRow{Giros y cambios de hoja}{71; 107}
\RouteRow{Acarreo y diferimiento}{deuda total: 2; final: 1; bloques diferidos: 1}
\RouteRow{Tríadas finales; persistencia de Pareto}{16; no}
\RouteRow{\(K\longrightarrow K+9\)}{repeticiones: 9; cambios de memoria: 9}
\RouteGroup{Firma, fase y diferencias}
\RouteRow{\(n_A,\ s_C,\ \kappa_0,\ \nu_{120},\ \nu_{270}\)}{28, -6, 0, 0, -12}
\RouteRow{\(\tau_{12}\); firma histórica \(\mathbb Z^5\)}{+++-\kern0pt{}-\kern0pt{}-+++-\kern0pt{}-\kern0pt{}-; 28|\allowbreak{}-6|\allowbreak{}0|\allowbreak{}0|\allowbreak{}-12}
\RouteRow{\(D_1,\ D_3,\ D_4,\ D_{27},\ D_{36}\)}{66, 60, 66, 48, 32}
\RouteRow{\(\Omega_{90/120},\ P_6\)}{80, 6}
\RouteGroup{Lectores y factores de acción}
\RouteRow{\(K_D\quad /\quad K_\Omega\)}{(80,\allowbreak{}66,\allowbreak{}3)\quad /\quad (80,\allowbreak{}54,\allowbreak{}6)}
\RouteRow{\(\kappa_D\)}{79.51870\allowbreak{}83196953\allowbreak{}45554341\allowbreak{}34}
\RouteRow{\(\kappa_\Omega\)}{79.60661\allowbreak{}01397948\allowbreak{}27586470\allowbreak{}91}
\RouteRow{\(R_{\mathrm{act},D}\)}{1.000000\allowbreak{}05512887\allowbreak{}17997050\allowbreak{}72}
\RouteRow{\(R_{\mathrm{act},\Omega}\)}{1.000000\allowbreak{}05518981\allowbreak{}25318591\allowbreak{}40}
\RouteGroup{Separación de prefijos}
\RouteRow{Área de ambigüedad; área \(\log_2\)}{18; 0.0}
\RouteRow{Primer bloque de separación}{órbita: 1; celda: 1; ruta: 1}
\RouteRow{Ambigüedad final}{1}
\RouteGroup{Secuencias completas}
\RouteRow{\(w_6\)}{264 90 570 264 90 24 243 657 24 264 90 570 264 90 24 243 657 24}
\RouteRow{Tríadas estables por bloque}{0 1 2 3 4 5 6 7 8 9 10 11 12 13 14 15 15 16}
\end{tabularx}\endgroup\par\vspace{5pt}\noindent{\fontsize{8.4}{10}\selectfont\hyperref[trace-205]{Procedencia y códigos originales}\hfill Ficha completa · 205/324}
\clearpage\phantomsection\label{nav:vi-ruta-206}
% VI-CATALOG-ROW rutas 206
\pdfbookmark[3]{r6c7\_h-1v+1}{route-206}
\RouteTitle{Ruta 206}{r6c7\_\allowbreak{}h-1v+1}
\noindent Celda 52\quad (6, 7)\qquad Orientación \(h=-1,\ v=1\)\hfill\hyperref[sec:vi-indice-rutas]{Índice}\par\vspace{4pt}
\begingroup\fontsize{10.2}{12.5}\selectfont\setlength{\tabcolsep}{5pt}\renewcommand{\arraystretch}{1.1}\rowcolors{1}{routepale}{white}\noindent\begin{tabularx}{\linewidth}{@{}>{\raggedright\arraybackslash}p{.345\linewidth}>{\raggedright\arraybackslash}X@{}}
\RouteGroup{Identidad estructural}
\RouteRow{Clase y multisección}{quark\_\allowbreak{}down\_\allowbreak{}G2\quad /\quad quark\_\allowbreak{}down\_\allowbreak{}G2}
\RouteRow{Colección y sector}{quark de tipo abajo; quark}
\RouteRow{Clasificación física}{quark de tipo abajo}
\RouteRow{Generación, profundidad y color}{2; G2; B}
\RouteRow{Ruta primaria y órbita de inversión}{\(B_1\): sí\qquad 4,\allowbreak{}3;\allowbreak{}6,\allowbreak{}7}
\RouteRow{Firma de clasificación}{28|\allowbreak{}6|\allowbreak{}-8|\allowbreak{}0|\allowbreak{}-12|\allowbreak{}+++-\kern0pt{}-\kern0pt{}-+++-\kern0pt{}-\kern0pt{}-}
\RouteGroup{Retornos, memoria y transporte}
\RouteRow{Retornos con fase}{clase: 27; órbita: 27; celda: 27}
\RouteRow{Retorno cinemático}{en 108: 54; extensión: 216}
\RouteRow{Giros y cambios de hoja}{71; 107}
\RouteRow{Acarreo y diferimiento}{deuda total: 2; final: 0; bloques diferidos: 1}
\RouteRow{Tríadas finales; persistencia de Pareto}{17; no}
\RouteRow{\(K\longrightarrow K+9\)}{repeticiones: 9; cambios de memoria: 9}
\RouteGroup{Firma, fase y diferencias}
\RouteRow{\(n_A,\ s_C,\ \kappa_0,\ \nu_{120},\ \nu_{270}\)}{8, -2, 0, 0, -12}
\RouteRow{\(\tau_{12}\); firma histórica \(\mathbb Z^5\)}{+++-\kern0pt{}-\kern0pt{}-+++-\kern0pt{}-\kern0pt{}-; 8|\allowbreak{}-2|\allowbreak{}0|\allowbreak{}0|\allowbreak{}-12}
\RouteRow{\(D_1,\ D_3,\ D_4,\ D_{27},\ D_{36}\)}{84, 24, 84, 24, 16}
\RouteRow{\(\Omega_{90/120},\ P_6\)}{80, 6}
\RouteGroup{Lectores y factores de acción}
\RouteRow{\(K_D\quad /\quad K_\Omega\)}{(40,\allowbreak{}84,\allowbreak{}30)\quad /\quad (80,\allowbreak{}54,\allowbreak{}6)}
\RouteRow{\(\kappa_D\)}{39.39035\allowbreak{}82768609\allowbreak{}24917849\allowbreak{}59}
\RouteRow{\(\kappa_\Omega\)}{79.60661\allowbreak{}01397948\allowbreak{}27586470\allowbreak{}91}
\RouteRow{\(R_{\mathrm{act},D}\)}{1.000000\allowbreak{}02730861\allowbreak{}73967087\allowbreak{}05}
\RouteRow{\(R_{\mathrm{act},\Omega}\)}{1.000000\allowbreak{}05518981\allowbreak{}25318591\allowbreak{}40}
\RouteGroup{Separación de prefijos}
\RouteRow{Área de ambigüedad; área \(\log_2\)}{18; 0.0}
\RouteRow{Primer bloque de separación}{órbita: 1; celda: 1; ruta: 1}
\RouteRow{Ambigüedad final}{1}
\RouteGroup{Secuencias completas}
\RouteRow{\(w_6\)}{258 414 420 258 414 255 252 333 255 258 414 420 258 414 255 252 333 255}
\RouteRow{Tríadas estables por bloque}{0 1 2 3 4 5 6 7 8 9 10 11 11 12 14 15 16 17}
\end{tabularx}\endgroup\par\vspace{5pt}\noindent{\fontsize{8.4}{10}\selectfont\hyperref[trace-206]{Procedencia y códigos originales}\hfill Ficha completa · 206/324}
\clearpage\phantomsection\label{nav:vi-ruta-207}
% VI-CATALOG-ROW rutas 207
\pdfbookmark[3]{r6c7\_h+1v-1}{route-207}
\RouteTitle{Ruta 207}{r6c7\_\allowbreak{}h+1v-1}
\noindent Celda 52\quad (6, 7)\qquad Orientación \(h=1,\ v=-1\)\hfill\hyperref[sec:vi-indice-rutas]{Índice}\par\vspace{4pt}
\begingroup\fontsize{10.2}{12.5}\selectfont\setlength{\tabcolsep}{5pt}\renewcommand{\arraystretch}{1.1}\rowcolors{1}{routepale}{white}\noindent\begin{tabularx}{\linewidth}{@{}>{\raggedright\arraybackslash}p{.345\linewidth}>{\raggedright\arraybackslash}X@{}}
\RouteGroup{Identidad estructural}
\RouteRow{Clase y multisección}{quark\_\allowbreak{}down\_\allowbreak{}G2\quad /\quad quark\_\allowbreak{}down\_\allowbreak{}G2}
\RouteRow{Colección y sector}{quark de tipo abajo; quark}
\RouteRow{Clasificación física}{quark de tipo abajo}
\RouteRow{Generación, profundidad y color}{2; G2; B}
\RouteRow{Ruta primaria y órbita de inversión}{\(B_1\): no\qquad 4,\allowbreak{}3;\allowbreak{}6,\allowbreak{}7}
\RouteRow{Firma de clasificación}{28|\allowbreak{}6|\allowbreak{}-8|\allowbreak{}0|\allowbreak{}-12|\allowbreak{}+++-\kern0pt{}-\kern0pt{}-+++-\kern0pt{}-\kern0pt{}-}
\RouteGroup{Retornos, memoria y transporte}
\RouteRow{Retornos con fase}{clase: 27; órbita: 27; celda: 27}
\RouteRow{Retorno cinemático}{en 108: 54; extensión: 216}
\RouteRow{Giros y cambios de hoja}{71; 107}
\RouteRow{Acarreo y diferimiento}{deuda total: 3; final: 1; bloques diferidos: 3}
\RouteRow{Tríadas finales; persistencia de Pareto}{16; no}
\RouteRow{\(K\longrightarrow K+9\)}{repeticiones: 9; cambios de memoria: 9}
\RouteGroup{Firma, fase y diferencias}
\RouteRow{\(n_A,\ s_C,\ \kappa_0,\ \nu_{120},\ \nu_{270}\)}{8, 2, 0, 0, -12}
\RouteRow{\(\tau_{12}\); firma histórica \(\mathbb Z^5\)}{+++-\kern0pt{}-\kern0pt{}-+++-\kern0pt{}-\kern0pt{}-; 8|\allowbreak{}2|\allowbreak{}0|\allowbreak{}0|\allowbreak{}-12}
\RouteRow{\(D_1,\ D_3,\ D_4,\ D_{27},\ D_{36}\)}{84, 24, 84, 24, 16}
\RouteRow{\(\Omega_{90/120},\ P_6\)}{80, 6}
\RouteGroup{Lectores y factores de acción}
\RouteRow{\(K_D\quad /\quad K_\Omega\)}{(40,\allowbreak{}84,\allowbreak{}30)\quad /\quad (80,\allowbreak{}54,\allowbreak{}6)}
\RouteRow{\(\kappa_D\)}{39.39035\allowbreak{}82768609\allowbreak{}24917849\allowbreak{}59}
\RouteRow{\(\kappa_\Omega\)}{79.60661\allowbreak{}01397948\allowbreak{}27586470\allowbreak{}91}
\RouteRow{\(R_{\mathrm{act},D}\)}{1.000000\allowbreak{}02730861\allowbreak{}73967087\allowbreak{}05}
\RouteRow{\(R_{\mathrm{act},\Omega}\)}{1.000000\allowbreak{}05518981\allowbreak{}25318591\allowbreak{}40}
\RouteGroup{Separación de prefijos}
\RouteRow{Área de ambigüedad; área \(\log_2\)}{18; 0.0}
\RouteRow{Primer bloque de separación}{órbita: 1; celda: 1; ruta: 1}
\RouteRow{Ambigüedad final}{1}
\RouteGroup{Secuencias completas}
\RouteRow{\(w_6\)}{252 333 255 252 333 420 258 414 420 252 333 255 252 333 420 258 414 420}
\RouteRow{Tríadas estables por bloque}{0 1 2 3 4 5 6 7 8 9 10 11 12 12 14 14 16 16}
\end{tabularx}\endgroup\par\vspace{5pt}\noindent{\fontsize{8.4}{10}\selectfont\hyperref[trace-207]{Procedencia y códigos originales}\hfill Ficha completa · 207/324}
\clearpage\phantomsection\label{nav:vi-ruta-208}
% VI-CATALOG-ROW rutas 208
\pdfbookmark[3]{r6c7\_h+1v+1}{route-208}
\RouteTitle{Ruta 208}{r6c7\_\allowbreak{}h+1v+1}
\noindent Celda 52\quad (6, 7)\qquad Orientación \(h=1,\ v=1\)\hfill\hyperref[sec:vi-indice-rutas]{Índice}\par\vspace{4pt}
\begingroup\fontsize{10.2}{12.5}\selectfont\setlength{\tabcolsep}{5pt}\renewcommand{\arraystretch}{1.1}\rowcolors{1}{routepale}{white}\noindent\begin{tabularx}{\linewidth}{@{}>{\raggedright\arraybackslash}p{.345\linewidth}>{\raggedright\arraybackslash}X@{}}
\RouteGroup{Identidad estructural}
\RouteRow{Clase y multisección}{quark\_\allowbreak{}down\_\allowbreak{}G2\quad /\quad quark\_\allowbreak{}down\_\allowbreak{}G2}
\RouteRow{Colección y sector}{quark de tipo abajo; quark}
\RouteRow{Clasificación física}{quark de tipo abajo}
\RouteRow{Generación, profundidad y color}{2; G2; B}
\RouteRow{Ruta primaria y órbita de inversión}{\(B_1\): no\qquad 4,\allowbreak{}3;\allowbreak{}6,\allowbreak{}7}
\RouteRow{Firma de clasificación}{28|\allowbreak{}6|\allowbreak{}-8|\allowbreak{}0|\allowbreak{}-12|\allowbreak{}+++-\kern0pt{}-\kern0pt{}-+++-\kern0pt{}-\kern0pt{}-}
\RouteGroup{Retornos, memoria y transporte}
\RouteRow{Retornos con fase}{clase: 27; órbita: 27; celda: 27}
\RouteRow{Retorno cinemático}{en 108: 54; extensión: 216}
\RouteRow{Giros y cambios de hoja}{71; 107}
\RouteRow{Acarreo y diferimiento}{deuda total: 0; final: 0; bloques diferidos: 0}
\RouteRow{Tríadas finales; persistencia de Pareto}{17; no}
\RouteRow{\(K\longrightarrow K+9\)}{repeticiones: 9; cambios de memoria: 9}
\RouteGroup{Firma, fase y diferencias}
\RouteRow{\(n_A,\ s_C,\ \kappa_0,\ \nu_{120},\ \nu_{270}\)}{28, 6, 0, 0, -12}
\RouteRow{\(\tau_{12}\); firma histórica \(\mathbb Z^5\)}{+++-\kern0pt{}-\kern0pt{}-+++-\kern0pt{}-\kern0pt{}-; 28|\allowbreak{}6|\allowbreak{}0|\allowbreak{}0|\allowbreak{}-12}
\RouteRow{\(D_1,\ D_3,\ D_4,\ D_{27},\ D_{36}\)}{66, 60, 66, 48, 32}
\RouteRow{\(\Omega_{90/120},\ P_6\)}{80, 6}
\RouteGroup{Lectores y factores de acción}
\RouteRow{\(K_D\quad /\quad K_\Omega\)}{(80,\allowbreak{}66,\allowbreak{}3)\quad /\quad (80,\allowbreak{}54,\allowbreak{}6)}
\RouteRow{\(\kappa_D\)}{79.51870\allowbreak{}83196953\allowbreak{}45554341\allowbreak{}34}
\RouteRow{\(\kappa_\Omega\)}{79.60661\allowbreak{}01397948\allowbreak{}27586470\allowbreak{}91}
\RouteRow{\(R_{\mathrm{act},D}\)}{1.000000\allowbreak{}05512887\allowbreak{}17997050\allowbreak{}72}
\RouteRow{\(R_{\mathrm{act},\Omega}\)}{1.000000\allowbreak{}05518981\allowbreak{}25318591\allowbreak{}40}
\RouteGroup{Separación de prefijos}
\RouteRow{Área de ambigüedad; área \(\log_2\)}{18; 0.0}
\RouteRow{Primer bloque de separación}{órbita: 1; celda: 1; ruta: 1}
\RouteRow{Ambigüedad final}{1}
\RouteGroup{Secuencias completas}
\RouteRow{\(w_6\)}{243 657 24 243 657 570 264 90 570 243 657 24 243 657 570 264 90 570}
\RouteRow{Tríadas estables por bloque}{0 1 2 3 4 5 6 7 8 9 10 11 12 13 14 15 16 17}
\end{tabularx}\endgroup\par\vspace{5pt}\noindent{\fontsize{8.4}{10}\selectfont\hyperref[trace-208]{Procedencia y códigos originales}\hfill Ficha completa · 208/324}
\clearpage\phantomsection\label{nav:vi-ruta-209}
% VI-CATALOG-ROW rutas 209
\pdfbookmark[2]{Celda 53}{cell-53}
\pdfbookmark[3]{r6c8\_h-1v-1}{route-209}
\RouteTitle{Ruta 209}{r6c8\_\allowbreak{}h-1v-1}
\noindent Celda 53\quad (6, 8)\qquad Orientación \(h=-1,\ v=-1\)\hfill\hyperref[sec:vi-indice-rutas]{Índice}\par\vspace{4pt}
\begingroup\fontsize{10.2}{12.5}\selectfont\setlength{\tabcolsep}{5pt}\renewcommand{\arraystretch}{1.1}\rowcolors{1}{routepale}{white}\noindent\begin{tabularx}{\linewidth}{@{}>{\raggedright\arraybackslash}p{.345\linewidth}>{\raggedright\arraybackslash}X@{}}
\RouteGroup{Identidad estructural}
\RouteRow{Clase y multisección}{quark\_\allowbreak{}up\_\allowbreak{}G3\quad /\quad quark\_\allowbreak{}up\_\allowbreak{}G3}
\RouteRow{Colección y sector}{quark de tipo arriba; quark}
\RouteRow{Clasificación física}{quark de tipo arriba}
\RouteRow{Generación, profundidad y color}{3; G3; G}
\RouteRow{Ruta primaria y órbita de inversión}{\(B_1\): no\qquad 4,\allowbreak{}2;\allowbreak{}6,\allowbreak{}8}
\RouteRow{Firma de clasificación}{28|\allowbreak{}-6|\allowbreak{}-2|\allowbreak{}-2|\allowbreak{}-8|\allowbreak{}0++-\kern0pt{}-\kern0pt{}-0++-\kern0pt{}-\kern0pt{}-}
\RouteGroup{Retornos, memoria y transporte}
\RouteRow{Retornos con fase}{clase: 27; órbita: 27; celda: 27}
\RouteRow{Retorno cinemático}{en 108: 54; extensión: 216}
\RouteRow{Giros y cambios de hoja}{71; 107}
\RouteRow{Acarreo y diferimiento}{deuda total: 1; final: 0; bloques diferidos: 1}
\RouteRow{Tríadas finales; persistencia de Pareto}{17; sí}
\RouteRow{\(K\longrightarrow K+9\)}{repeticiones: 9; cambios de memoria: 9}
\RouteGroup{Firma, fase y diferencias}
\RouteRow{\(n_A,\ s_C,\ \kappa_0,\ \nu_{120},\ \nu_{270}\)}{28, 6, -2, 2, -8}
\RouteRow{\(\tau_{12}\); firma histórica \(\mathbb Z^5\)}{+++0-\kern0pt{}-+++0-\kern0pt{}-; 28|\allowbreak{}6|\allowbreak{}-2|\allowbreak{}2|\allowbreak{}-8}
\RouteRow{\(D_1,\ D_3,\ D_4,\ D_{27},\ D_{36}\)}{66, 60, 66, 48, 32}
\RouteRow{\(\Omega_{90/120},\ P_6\)}{56, 5}
\RouteGroup{Lectores y factores de acción}
\RouteRow{\(K_D\quad /\quad K_\Omega\)}{(80,\allowbreak{}66,\allowbreak{}3)\quad /\quad (56,\allowbreak{}54,\allowbreak{}5)}
\RouteRow{\(\kappa_D\)}{79.51870\allowbreak{}83196953\allowbreak{}45554341\allowbreak{}34}
\RouteRow{\(\kappa_\Omega\)}{55.60649\allowbreak{}89433721\allowbreak{}35446416\allowbreak{}86}
\RouteRow{\(R_{\mathrm{act},D}\)}{1.000000\allowbreak{}05512887\allowbreak{}17997050\allowbreak{}72}
\RouteRow{\(R_{\mathrm{act},\Omega}\)}{1.000000\allowbreak{}03855097\allowbreak{}23541416\allowbreak{}45}
\RouteGroup{Separación de prefijos}
\RouteRow{Área de ambigüedad; área \(\log_2\)}{36; 18.0}
\RouteRow{Primer bloque de separación}{órbita: \textemdash{}; celda: \textemdash{}; ruta: \textemdash{}}
\RouteRow{Ambigüedad final}{2}
\RouteGroup{Secuencias completas}
\RouteRow{\(w_6\)}{486 423 15 486 423 327 498 99 327 486 423 15 486 423 327 498 99 327}
\RouteRow{Tríadas estables por bloque}{0 1 2 3 4 5 6 7 8 9 10 10 12 13 14 15 16 17}
\end{tabularx}\endgroup\par\vspace{5pt}\noindent{\fontsize{8.4}{10}\selectfont\hyperref[trace-209]{Procedencia y códigos originales}\hfill Ficha completa · 209/324}
\clearpage\phantomsection\label{nav:vi-ruta-210}
% VI-CATALOG-ROW rutas 210
\pdfbookmark[3]{r6c8\_h-1v+1}{route-210}
\RouteTitle{Ruta 210}{r6c8\_\allowbreak{}h-1v+1}
\noindent Celda 53\quad (6, 8)\qquad Orientación \(h=-1,\ v=1\)\hfill\hyperref[sec:vi-indice-rutas]{Índice}\par\vspace{4pt}
\begingroup\fontsize{10.2}{12.5}\selectfont\setlength{\tabcolsep}{5pt}\renewcommand{\arraystretch}{1.1}\rowcolors{1}{routepale}{white}\noindent\begin{tabularx}{\linewidth}{@{}>{\raggedright\arraybackslash}p{.345\linewidth}>{\raggedright\arraybackslash}X@{}}
\RouteGroup{Identidad estructural}
\RouteRow{Clase y multisección}{quark\_\allowbreak{}up\_\allowbreak{}G3\quad /\quad quark\_\allowbreak{}up\_\allowbreak{}G3}
\RouteRow{Colección y sector}{quark de tipo arriba; quark}
\RouteRow{Clasificación física}{quark de tipo arriba}
\RouteRow{Generación, profundidad y color}{3; G3; G}
\RouteRow{Ruta primaria y órbita de inversión}{\(B_1\): no\qquad 4,\allowbreak{}2;\allowbreak{}6,\allowbreak{}8}
\RouteRow{Firma de clasificación}{28|\allowbreak{}-6|\allowbreak{}-2|\allowbreak{}-2|\allowbreak{}-8|\allowbreak{}0++-\kern0pt{}-\kern0pt{}-0++-\kern0pt{}-\kern0pt{}-}
\RouteGroup{Retornos, memoria y transporte}
\RouteRow{Retornos con fase}{clase: 27; órbita: 27; celda: 27}
\RouteRow{Retorno cinemático}{en 108: 54; extensión: 216}
\RouteRow{Giros y cambios de hoja}{71; 107}
\RouteRow{Acarreo y diferimiento}{deuda total: 1; final: 0; bloques diferidos: 1}
\RouteRow{Tríadas finales; persistencia de Pareto}{17; sí}
\RouteRow{\(K\longrightarrow K+9\)}{repeticiones: 9; cambios de memoria: 9}
\RouteGroup{Firma, fase y diferencias}
\RouteRow{\(n_A,\ s_C,\ \kappa_0,\ \nu_{120},\ \nu_{270}\)}{44, 2, 0, 0, -12}
\RouteRow{\(\tau_{12}\); firma histórica \(\mathbb Z^5\)}{+++-\kern0pt{}-\kern0pt{}-+++-\kern0pt{}-\kern0pt{}-; 44|\allowbreak{}2|\allowbreak{}0|\allowbreak{}0|\allowbreak{}-12}
\RouteRow{\(D_1,\ D_3,\ D_4,\ D_{27},\ D_{36}\)}{78, 24, 78, 24, 16}
\RouteRow{\(\Omega_{90/120},\ P_6\)}{80, 6}
\RouteGroup{Lectores y factores de acción}
\RouteRow{\(K_D\quad /\quad K_\Omega\)}{(40,\allowbreak{}78,\allowbreak{}27)\quad /\quad (80,\allowbreak{}54,\allowbreak{}6)}
\RouteRow{\(\kappa_D\)}{39.43380\allowbreak{}88030085\allowbreak{}51303671\allowbreak{}14}
\RouteRow{\(\kappa_\Omega\)}{79.60661\allowbreak{}01397948\allowbreak{}27586470\allowbreak{}91}
\RouteRow{\(R_{\mathrm{act},D}\)}{1.000000\allowbreak{}02733874\allowbreak{}08548943\allowbreak{}58}
\RouteRow{\(R_{\mathrm{act},\Omega}\)}{1.000000\allowbreak{}05518981\allowbreak{}25318591\allowbreak{}40}
\RouteGroup{Separación de prefijos}
\RouteRow{Área de ambigüedad; área \(\log_2\)}{36; 18.0}
\RouteRow{Primer bloque de separación}{órbita: \textemdash{}; celda: \textemdash{}; ruta: \textemdash{}}
\RouteRow{Ambigüedad final}{2}
\RouteGroup{Secuencias completas}
\RouteRow{\(w_6\)}{504 585 507 504 585 672 510 666 672 504 585 507 504 585 672 510 666 672}
\RouteRow{Tríadas estables por bloque}{0 1 2 3 4 5 6 7 8 9 10 11 12 13 14 14 16 17}
\end{tabularx}\endgroup\par\vspace{5pt}\noindent{\fontsize{8.4}{10}\selectfont\hyperref[trace-210]{Procedencia y códigos originales}\hfill Ficha completa · 210/324}
\clearpage\phantomsection\label{nav:vi-ruta-211}
% VI-CATALOG-ROW rutas 211
\pdfbookmark[3]{r6c8\_h+1v-1}{route-211}
\RouteTitle{Ruta 211}{r6c8\_\allowbreak{}h+1v-1}
\noindent Celda 53\quad (6, 8)\qquad Orientación \(h=1,\ v=-1\)\hfill\hyperref[sec:vi-indice-rutas]{Índice}\par\vspace{4pt}
\begingroup\fontsize{10.2}{12.5}\selectfont\setlength{\tabcolsep}{5pt}\renewcommand{\arraystretch}{1.1}\rowcolors{1}{routepale}{white}\noindent\begin{tabularx}{\linewidth}{@{}>{\raggedright\arraybackslash}p{.345\linewidth}>{\raggedright\arraybackslash}X@{}}
\RouteGroup{Identidad estructural}
\RouteRow{Clase y multisección}{quark\_\allowbreak{}up\_\allowbreak{}G3\quad /\quad quark\_\allowbreak{}up\_\allowbreak{}G3}
\RouteRow{Colección y sector}{quark de tipo arriba; quark}
\RouteRow{Clasificación física}{quark de tipo arriba}
\RouteRow{Generación, profundidad y color}{3; G3; G}
\RouteRow{Ruta primaria y órbita de inversión}{\(B_1\): no\qquad 4,\allowbreak{}2;\allowbreak{}6,\allowbreak{}8}
\RouteRow{Firma de clasificación}{28|\allowbreak{}-6|\allowbreak{}-2|\allowbreak{}-2|\allowbreak{}-8|\allowbreak{}0++-\kern0pt{}-\kern0pt{}-0++-\kern0pt{}-\kern0pt{}-}
\RouteGroup{Retornos, memoria y transporte}
\RouteRow{Retornos con fase}{clase: 27; órbita: 27; celda: 27}
\RouteRow{Retorno cinemático}{en 108: 54; extensión: 216}
\RouteRow{Giros y cambios de hoja}{71; 107}
\RouteRow{Acarreo y diferimiento}{deuda total: 1; final: 0; bloques diferidos: 1}
\RouteRow{Tríadas finales; persistencia de Pareto}{17; sí}
\RouteRow{\(K\longrightarrow K+9\)}{repeticiones: 9; cambios de memoria: 9}
\RouteGroup{Firma, fase y diferencias}
\RouteRow{\(n_A,\ s_C,\ \kappa_0,\ \nu_{120},\ \nu_{270}\)}{44, -2, 0, 0, -12}
\RouteRow{\(\tau_{12}\); firma histórica \(\mathbb Z^5\)}{+++-\kern0pt{}-\kern0pt{}-+++-\kern0pt{}-\kern0pt{}-; 44|\allowbreak{}-2|\allowbreak{}0|\allowbreak{}0|\allowbreak{}-12}
\RouteRow{\(D_1,\ D_3,\ D_4,\ D_{27},\ D_{36}\)}{78, 24, 78, 24, 16}
\RouteRow{\(\Omega_{90/120},\ P_6\)}{80, 6}
\RouteGroup{Lectores y factores de acción}
\RouteRow{\(K_D\quad /\quad K_\Omega\)}{(40,\allowbreak{}78,\allowbreak{}27)\quad /\quad (80,\allowbreak{}54,\allowbreak{}6)}
\RouteRow{\(\kappa_D\)}{39.43380\allowbreak{}88030085\allowbreak{}51303671\allowbreak{}14}
\RouteRow{\(\kappa_\Omega\)}{79.60661\allowbreak{}01397948\allowbreak{}27586470\allowbreak{}91}
\RouteRow{\(R_{\mathrm{act},D}\)}{1.000000\allowbreak{}02733874\allowbreak{}08548943\allowbreak{}58}
\RouteRow{\(R_{\mathrm{act},\Omega}\)}{1.000000\allowbreak{}05518981\allowbreak{}25318591\allowbreak{}40}
\RouteGroup{Separación de prefijos}
\RouteRow{Área de ambigüedad; área \(\log_2\)}{36; 18.0}
\RouteRow{Primer bloque de separación}{órbita: \textemdash{}; celda: \textemdash{}; ruta: \textemdash{}}
\RouteRow{Ambigüedad final}{2}
\RouteGroup{Secuencias completas}
\RouteRow{\(w_6\)}{510 666 672 510 666 507 504 585 507 510 666 672 510 666 507 504 585 507}
\RouteRow{Tríadas estables por bloque}{0 1 2 3 4 5 6 7 8 9 10 11 12 13 14 15 15 17}
\end{tabularx}\endgroup\par\vspace{5pt}\noindent{\fontsize{8.4}{10}\selectfont\hyperref[trace-211]{Procedencia y códigos originales}\hfill Ficha completa · 211/324}
\clearpage\phantomsection\label{nav:vi-ruta-212}
% VI-CATALOG-ROW rutas 212
\pdfbookmark[3]{r6c8\_h+1v+1}{route-212}
\RouteTitle{Ruta 212}{r6c8\_\allowbreak{}h+1v+1}
\noindent Celda 53\quad (6, 8)\qquad Orientación \(h=1,\ v=1\)\hfill\hyperref[sec:vi-indice-rutas]{Índice}\par\vspace{4pt}
\begingroup\fontsize{10.2}{12.5}\selectfont\setlength{\tabcolsep}{5pt}\renewcommand{\arraystretch}{1.1}\rowcolors{1}{routepale}{white}\noindent\begin{tabularx}{\linewidth}{@{}>{\raggedright\arraybackslash}p{.345\linewidth}>{\raggedright\arraybackslash}X@{}}
\RouteGroup{Identidad estructural}
\RouteRow{Clase y multisección}{quark\_\allowbreak{}up\_\allowbreak{}G3\quad /\quad quark\_\allowbreak{}up\_\allowbreak{}G3}
\RouteRow{Colección y sector}{quark de tipo arriba; quark}
\RouteRow{Clasificación física}{quark de tipo arriba}
\RouteRow{Generación, profundidad y color}{3; G3; G}
\RouteRow{Ruta primaria y órbita de inversión}{\(B_1\): sí\qquad 4,\allowbreak{}2;\allowbreak{}6,\allowbreak{}8}
\RouteRow{Firma de clasificación}{28|\allowbreak{}-6|\allowbreak{}-2|\allowbreak{}-2|\allowbreak{}-8|\allowbreak{}0++-\kern0pt{}-\kern0pt{}-0++-\kern0pt{}-\kern0pt{}-}
\RouteGroup{Retornos, memoria y transporte}
\RouteRow{Retornos con fase}{clase: 27; órbita: 27; celda: 27}
\RouteRow{Retorno cinemático}{en 108: 54; extensión: 216}
\RouteRow{Giros y cambios de hoja}{71; 107}
\RouteRow{Acarreo y diferimiento}{deuda total: 0; final: 0; bloques diferidos: 0}
\RouteRow{Tríadas finales; persistencia de Pareto}{17; sí}
\RouteRow{\(K\longrightarrow K+9\)}{repeticiones: 9; cambios de memoria: 9}
\RouteGroup{Firma, fase y diferencias}
\RouteRow{\(n_A,\ s_C,\ \kappa_0,\ \nu_{120},\ \nu_{270}\)}{28, -6, 2, -2, -8}
\RouteRow{\(\tau_{12}\); firma histórica \(\mathbb Z^5\)}{0++-\kern0pt{}-\kern0pt{}-0++-\kern0pt{}-\kern0pt{}-; 28|\allowbreak{}-6|\allowbreak{}2|\allowbreak{}-2|\allowbreak{}-8}
\RouteRow{\(D_1,\ D_3,\ D_4,\ D_{27},\ D_{36}\)}{66, 60, 66, 48, 32}
\RouteRow{\(\Omega_{90/120},\ P_6\)}{56, 5}
\RouteGroup{Lectores y factores de acción}
\RouteRow{\(K_D\quad /\quad K_\Omega\)}{(80,\allowbreak{}66,\allowbreak{}3)\quad /\quad (56,\allowbreak{}54,\allowbreak{}5)}
\RouteRow{\(\kappa_D\)}{79.51870\allowbreak{}83196953\allowbreak{}45554341\allowbreak{}34}
\RouteRow{\(\kappa_\Omega\)}{55.60649\allowbreak{}89433721\allowbreak{}35446416\allowbreak{}86}
\RouteRow{\(R_{\mathrm{act},D}\)}{1.000000\allowbreak{}05512887\allowbreak{}17997050\allowbreak{}72}
\RouteRow{\(R_{\mathrm{act},\Omega}\)}{1.000000\allowbreak{}03855097\allowbreak{}23541416\allowbreak{}45}
\RouteGroup{Separación de prefijos}
\RouteRow{Área de ambigüedad; área \(\log_2\)}{36; 18.0}
\RouteRow{Primer bloque de separación}{órbita: \textemdash{}; celda: \textemdash{}; ruta: \textemdash{}}
\RouteRow{Ambigüedad final}{2}
\RouteGroup{Secuencias completas}
\RouteRow{\(w_6\)}{498 99 327 498 99 15 486 423 15 498 99 327 498 99 15 486 423 15}
\RouteRow{Tríadas estables por bloque}{0 1 2 3 4 5 6 7 8 9 10 11 12 13 14 15 16 17}
\end{tabularx}\endgroup\par\vspace{5pt}\noindent{\fontsize{8.4}{10}\selectfont\hyperref[trace-212]{Procedencia y códigos originales}\hfill Ficha completa · 212/324}
\clearpage\phantomsection\label{nav:vi-ruta-213}
% VI-CATALOG-ROW rutas 213
\pdfbookmark[2]{Celda 54}{cell-54}
\pdfbookmark[3]{r6c9\_h-1v-1}{route-213}
\RouteTitle{Ruta 213}{r6c9\_\allowbreak{}h-1v-1}
\noindent Celda 54\quad (6, 9)\qquad Orientación \(h=-1,\ v=-1\)\hfill\hyperref[sec:vi-indice-rutas]{Índice}\par\vspace{4pt}
\begingroup\fontsize{10.2}{12.5}\selectfont\setlength{\tabcolsep}{5pt}\renewcommand{\arraystretch}{1.1}\rowcolors{1}{routepale}{white}\noindent\begin{tabularx}{\linewidth}{@{}>{\raggedright\arraybackslash}p{.345\linewidth}>{\raggedright\arraybackslash}X@{}}
\RouteGroup{Identidad estructural}
\RouteRow{Clase y multisección}{quark\_\allowbreak{}down\_\allowbreak{}G4\quad /\quad quark\_\allowbreak{}down\_\allowbreak{}G4}
\RouteRow{Colección y sector}{quark de tipo abajo; quark}
\RouteRow{Clasificación física}{quark de tipo abajo}
\RouteRow{Generación, profundidad y color}{4; G4; R}
\RouteRow{Ruta primaria y órbita de inversión}{\(B_1\): no\qquad 4,\allowbreak{}1;\allowbreak{}6,\allowbreak{}9}
\RouteRow{Firma de clasificación}{30|\allowbreak{}-6|\allowbreak{}0|\allowbreak{}0|\allowbreak{}-8|\allowbreak{}+0+-\kern0pt{}-\kern0pt{}-+0+-\kern0pt{}-\kern0pt{}-}
\RouteGroup{Retornos, memoria y transporte}
\RouteRow{Retornos con fase}{clase: 27; órbita: 27; celda: 27}
\RouteRow{Retorno cinemático}{en 108: 54; extensión: 216}
\RouteRow{Giros y cambios de hoja}{71; 107}
\RouteRow{Acarreo y diferimiento}{deuda total: 1; final: 0; bloques diferidos: 1}
\RouteRow{Tríadas finales; persistencia de Pareto}{17; sí}
\RouteRow{\(K\longrightarrow K+9\)}{repeticiones: 9; cambios de memoria: 9}
\RouteGroup{Firma, fase y diferencias}
\RouteRow{\(n_A,\ s_C,\ \kappa_0,\ \nu_{120},\ \nu_{270}\)}{30, 6, -2, 0, -8}
\RouteRow{\(\tau_{12}\); firma histórica \(\mathbb Z^5\)}{+++-0-+++-0-; 30|\allowbreak{}6|\allowbreak{}-2|\allowbreak{}0|\allowbreak{}-8}
\RouteRow{\(D_1,\ D_3,\ D_4,\ D_{27},\ D_{36}\)}{54, 60, 72, 24, 16}
\RouteRow{\(\Omega_{90/120},\ P_6\)}{48, 5}
\RouteGroup{Lectores y factores de acción}
\RouteRow{\(K_D\quad /\quad K_\Omega\)}{(40,\allowbreak{}54,\allowbreak{}6)\quad /\quad (48,\allowbreak{}54,\allowbreak{}5)}
\RouteRow{\(\kappa_D\)}{39.60661\allowbreak{}01397948\allowbreak{}27586470\allowbreak{}91}
\RouteRow{\(\kappa_\Omega\)}{47.60649\allowbreak{}89433721\allowbreak{}35446416\allowbreak{}86}
\RouteRow{\(R_{\mathrm{act},D}\)}{1.000000\allowbreak{}02745854\allowbreak{}08735595\allowbreak{}18}
\RouteRow{\(R_{\mathrm{act},\Omega}\)}{1.000000\allowbreak{}03300471\allowbreak{}80532430\allowbreak{}84}
\RouteGroup{Separación de prefijos}
\RouteRow{Área de ambigüedad; área \(\log_2\)}{26; 6.339850\allowbreak{}00288462\allowbreak{}4}
\RouteRow{Primer bloque de separación}{órbita: 5; celda: 5; ruta: 5}
\RouteRow{Ambigüedad final}{1}
\RouteGroup{Secuencias completas}
\RouteRow{\(w_6\)}{15 486 423 15 486 657 24 243 657 15 486 423 15 486 657 24 243 657}
\RouteRow{Tríadas estables por bloque}{0 1 2 3 4 5 6 7 8 9 10 11 12 13 14 15 15 17}
\end{tabularx}\endgroup\par\vspace{5pt}\noindent{\fontsize{8.4}{10}\selectfont\hyperref[trace-213]{Procedencia y códigos originales}\hfill Ficha completa · 213/324}
\clearpage\phantomsection\label{nav:vi-ruta-214}
% VI-CATALOG-ROW rutas 214
\pdfbookmark[3]{r6c9\_h-1v+1}{route-214}
\RouteTitle{Ruta 214}{r6c9\_\allowbreak{}h-1v+1}
\noindent Celda 54\quad (6, 9)\qquad Orientación \(h=-1,\ v=1\)\hfill\hyperref[sec:vi-indice-rutas]{Índice}\par\vspace{4pt}
\begingroup\fontsize{10.2}{12.5}\selectfont\setlength{\tabcolsep}{5pt}\renewcommand{\arraystretch}{1.1}\rowcolors{1}{routepale}{white}\noindent\begin{tabularx}{\linewidth}{@{}>{\raggedright\arraybackslash}p{.345\linewidth}>{\raggedright\arraybackslash}X@{}}
\RouteGroup{Identidad estructural}
\RouteRow{Clase y multisección}{quark\_\allowbreak{}down\_\allowbreak{}G4\quad /\quad quark\_\allowbreak{}down\_\allowbreak{}G4}
\RouteRow{Colección y sector}{quark de tipo abajo; quark}
\RouteRow{Clasificación física}{quark de tipo abajo}
\RouteRow{Generación, profundidad y color}{4; G4; R}
\RouteRow{Ruta primaria y órbita de inversión}{\(B_1\): no\qquad 4,\allowbreak{}1;\allowbreak{}6,\allowbreak{}9}
\RouteRow{Firma de clasificación}{30|\allowbreak{}-6|\allowbreak{}0|\allowbreak{}0|\allowbreak{}-8|\allowbreak{}+0+-\kern0pt{}-\kern0pt{}-+0+-\kern0pt{}-\kern0pt{}-}
\RouteGroup{Retornos, memoria y transporte}
\RouteRow{Retornos con fase}{clase: 27; órbita: 27; celda: 27}
\RouteRow{Retorno cinemático}{en 108: 54; extensión: 216}
\RouteRow{Giros y cambios de hoja}{71; 107}
\RouteRow{Acarreo y diferimiento}{deuda total: 0; final: 0; bloques diferidos: 0}
\RouteRow{Tríadas finales; persistencia de Pareto}{17; sí}
\RouteRow{\(K\longrightarrow K+9\)}{repeticiones: 9; cambios de memoria: 9}
\RouteGroup{Firma, fase y diferencias}
\RouteRow{\(n_A,\ s_C,\ \kappa_0,\ \nu_{120},\ \nu_{270}\)}{14, 6, 0, 0, -12}
\RouteRow{\(\tau_{12}\); firma histórica \(\mathbb Z^5\)}{+++-\kern0pt{}-\kern0pt{}-+++-\kern0pt{}-\kern0pt{}-; 14|\allowbreak{}6|\allowbreak{}0|\allowbreak{}0|\allowbreak{}-12}
\RouteRow{\(D_1,\ D_3,\ D_4,\ D_{27},\ D_{36}\)}{24, 24, 24, 0, 0}
\RouteRow{\(\Omega_{90/120},\ P_6\)}{80, 6}
\RouteGroup{Lectores y factores de acción}
\RouteRow{\(K_D\quad /\quad K_\Omega\)}{(0,\allowbreak{}24,\allowbreak{}0)\quad /\quad (80,\allowbreak{}54,\allowbreak{}6)}
\RouteRow{\(\kappa_D\)}{-0.175136\allowbreak{}46166281\allowbreak{}12239348\allowbreak{}425}
\RouteRow{\(\kappa_\Omega\)}{79.60661\allowbreak{}01397948\allowbreak{}27586470\allowbreak{}91}
\RouteRow{\(R_{\mathrm{act},D}\)}{0.999999\allowbreak{}99987858\allowbreak{}10851337\allowbreak{}881}
\RouteRow{\(R_{\mathrm{act},\Omega}\)}{1.000000\allowbreak{}05518981\allowbreak{}25318591\allowbreak{}40}
\RouteGroup{Separación de prefijos}
\RouteRow{Área de ambigüedad; área \(\log_2\)}{36; 18.0}
\RouteRow{Primer bloque de separación}{órbita: 1; celda: 1; ruta: \textemdash{}}
\RouteRow{Ambigüedad final}{2}
\RouteGroup{Secuencias completas}
\RouteRow{\(w_6\)}{3 0 81 3 0 81 3 0 81 3 0 81 3 0 81 3 0 81}
\RouteRow{Tríadas estables por bloque}{0 1 2 3 4 5 6 7 8 9 10 11 12 13 14 15 16 17}
\end{tabularx}\endgroup\par\vspace{5pt}\noindent{\fontsize{8.4}{10}\selectfont\hyperref[trace-214]{Procedencia y códigos originales}\hfill Ficha completa · 214/324}
\clearpage\phantomsection\label{nav:vi-ruta-215}
% VI-CATALOG-ROW rutas 215
\pdfbookmark[3]{r6c9\_h+1v-1}{route-215}
\RouteTitle{Ruta 215}{r6c9\_\allowbreak{}h+1v-1}
\noindent Celda 54\quad (6, 9)\qquad Orientación \(h=1,\ v=-1\)\hfill\hyperref[sec:vi-indice-rutas]{Índice}\par\vspace{4pt}
\begingroup\fontsize{10.2}{12.5}\selectfont\setlength{\tabcolsep}{5pt}\renewcommand{\arraystretch}{1.1}\rowcolors{1}{routepale}{white}\noindent\begin{tabularx}{\linewidth}{@{}>{\raggedright\arraybackslash}p{.345\linewidth}>{\raggedright\arraybackslash}X@{}}
\RouteGroup{Identidad estructural}
\RouteRow{Clase y multisección}{quark\_\allowbreak{}down\_\allowbreak{}G4\quad /\quad quark\_\allowbreak{}down\_\allowbreak{}G4}
\RouteRow{Colección y sector}{quark de tipo abajo; quark}
\RouteRow{Clasificación física}{quark de tipo abajo}
\RouteRow{Generación, profundidad y color}{4; G4; R}
\RouteRow{Ruta primaria y órbita de inversión}{\(B_1\): sí\qquad 4,\allowbreak{}1;\allowbreak{}6,\allowbreak{}9}
\RouteRow{Firma de clasificación}{30|\allowbreak{}-6|\allowbreak{}0|\allowbreak{}0|\allowbreak{}-8|\allowbreak{}+0+-\kern0pt{}-\kern0pt{}-+0+-\kern0pt{}-\kern0pt{}-}
\RouteGroup{Retornos, memoria y transporte}
\RouteRow{Retornos con fase}{clase: 27; órbita: 27; celda: 27}
\RouteRow{Retorno cinemático}{en 108: 54; extensión: 216}
\RouteRow{Giros y cambios de hoja}{71; 107}
\RouteRow{Acarreo y diferimiento}{deuda total: 0; final: 0; bloques diferidos: 0}
\RouteRow{Tríadas finales; persistencia de Pareto}{17; sí}
\RouteRow{\(K\longrightarrow K+9\)}{repeticiones: 9; cambios de memoria: 9}
\RouteGroup{Firma, fase y diferencias}
\RouteRow{\(n_A,\ s_C,\ \kappa_0,\ \nu_{120},\ \nu_{270}\)}{14, -6, 0, 0, -12}
\RouteRow{\(\tau_{12}\); firma histórica \(\mathbb Z^5\)}{+++-\kern0pt{}-\kern0pt{}-+++-\kern0pt{}-\kern0pt{}-; 14|\allowbreak{}-6|\allowbreak{}0|\allowbreak{}0|\allowbreak{}-12}
\RouteRow{\(D_1,\ D_3,\ D_4,\ D_{27},\ D_{36}\)}{24, 24, 24, 0, 0}
\RouteRow{\(\Omega_{90/120},\ P_6\)}{80, 6}
\RouteGroup{Lectores y factores de acción}
\RouteRow{\(K_D\quad /\quad K_\Omega\)}{(0,\allowbreak{}24,\allowbreak{}0)\quad /\quad (80,\allowbreak{}54,\allowbreak{}6)}
\RouteRow{\(\kappa_D\)}{-0.175136\allowbreak{}46166281\allowbreak{}12239348\allowbreak{}425}
\RouteRow{\(\kappa_\Omega\)}{79.60661\allowbreak{}01397948\allowbreak{}27586470\allowbreak{}91}
\RouteRow{\(R_{\mathrm{act},D}\)}{0.999999\allowbreak{}99987858\allowbreak{}10851337\allowbreak{}881}
\RouteRow{\(R_{\mathrm{act},\Omega}\)}{1.000000\allowbreak{}05518981\allowbreak{}25318591\allowbreak{}40}
\RouteGroup{Separación de prefijos}
\RouteRow{Área de ambigüedad; área \(\log_2\)}{36; 18.0}
\RouteRow{Primer bloque de separación}{órbita: 1; celda: 1; ruta: \textemdash{}}
\RouteRow{Ambigüedad final}{2}
\RouteGroup{Secuencias completas}
\RouteRow{\(w_6\)}{3 0 81 3 0 81 3 0 81 3 0 81 3 0 81 3 0 81}
\RouteRow{Tríadas estables por bloque}{0 1 2 3 4 5 6 7 8 9 10 11 12 13 14 15 16 17}
\end{tabularx}\endgroup\par\vspace{5pt}\noindent{\fontsize{8.4}{10}\selectfont\hyperref[trace-215]{Procedencia y códigos originales}\hfill Ficha completa · 215/324}
\clearpage\phantomsection\label{nav:vi-ruta-216}
% VI-CATALOG-ROW rutas 216
\pdfbookmark[3]{r6c9\_h+1v+1}{route-216}
\RouteTitle{Ruta 216}{r6c9\_\allowbreak{}h+1v+1}
\noindent Celda 54\quad (6, 9)\qquad Orientación \(h=1,\ v=1\)\hfill\hyperref[sec:vi-indice-rutas]{Índice}\par\vspace{4pt}
\begingroup\fontsize{10.2}{12.5}\selectfont\setlength{\tabcolsep}{5pt}\renewcommand{\arraystretch}{1.1}\rowcolors{1}{routepale}{white}\noindent\begin{tabularx}{\linewidth}{@{}>{\raggedright\arraybackslash}p{.345\linewidth}>{\raggedright\arraybackslash}X@{}}
\RouteGroup{Identidad estructural}
\RouteRow{Clase y multisección}{quark\_\allowbreak{}down\_\allowbreak{}G4\quad /\quad quark\_\allowbreak{}down\_\allowbreak{}G4}
\RouteRow{Colección y sector}{quark de tipo abajo; quark}
\RouteRow{Clasificación física}{quark de tipo abajo}
\RouteRow{Generación, profundidad y color}{4; G4; R}
\RouteRow{Ruta primaria y órbita de inversión}{\(B_1\): no\qquad 4,\allowbreak{}1;\allowbreak{}6,\allowbreak{}9}
\RouteRow{Firma de clasificación}{30|\allowbreak{}-6|\allowbreak{}0|\allowbreak{}0|\allowbreak{}-8|\allowbreak{}+0+-\kern0pt{}-\kern0pt{}-+0+-\kern0pt{}-\kern0pt{}-}
\RouteGroup{Retornos, memoria y transporte}
\RouteRow{Retornos con fase}{clase: 27; órbita: 27; celda: 27}
\RouteRow{Retorno cinemático}{en 108: 54; extensión: 216}
\RouteRow{Giros y cambios de hoja}{71; 107}
\RouteRow{Acarreo y diferimiento}{deuda total: 2; final: 1; bloques diferidos: 2}
\RouteRow{Tríadas finales; persistencia de Pareto}{16; sí}
\RouteRow{\(K\longrightarrow K+9\)}{repeticiones: 9; cambios de memoria: 9}
\RouteGroup{Firma, fase y diferencias}
\RouteRow{\(n_A,\ s_C,\ \kappa_0,\ \nu_{120},\ \nu_{270}\)}{30, -6, 2, 0, -8}
\RouteRow{\(\tau_{12}\); firma histórica \(\mathbb Z^5\)}{+0+-\kern0pt{}-\kern0pt{}-+0+-\kern0pt{}-\kern0pt{}-; 30|\allowbreak{}-6|\allowbreak{}2|\allowbreak{}0|\allowbreak{}-8}
\RouteRow{\(D_1,\ D_3,\ D_4,\ D_{27},\ D_{36}\)}{54, 60, 72, 24, 16}
\RouteRow{\(\Omega_{90/120},\ P_6\)}{48, 5}
\RouteGroup{Lectores y factores de acción}
\RouteRow{\(K_D\quad /\quad K_\Omega\)}{(40,\allowbreak{}54,\allowbreak{}6)\quad /\quad (48,\allowbreak{}54,\allowbreak{}5)}
\RouteRow{\(\kappa_D\)}{39.60661\allowbreak{}01397948\allowbreak{}27586470\allowbreak{}91}
\RouteRow{\(\kappa_\Omega\)}{47.60649\allowbreak{}89433721\allowbreak{}35446416\allowbreak{}86}
\RouteRow{\(R_{\mathrm{act},D}\)}{1.000000\allowbreak{}02745854\allowbreak{}08735595\allowbreak{}18}
\RouteRow{\(R_{\mathrm{act},\Omega}\)}{1.000000\allowbreak{}03300471\allowbreak{}80532430\allowbreak{}84}
\RouteGroup{Separación de prefijos}
\RouteRow{Área de ambigüedad; área \(\log_2\)}{18; 0.0}
\RouteRow{Primer bloque de separación}{órbita: 1; celda: 1; ruta: 1}
\RouteRow{Ambigüedad final}{1}
\RouteGroup{Secuencias completas}
\RouteRow{\(w_6\)}{24 243 657 24 243 423 15 486 423 24 243 657 24 243 423 15 486 423}
\RouteRow{Tríadas estables por bloque}{0 1 2 3 4 5 6 7 8 9 10 11 12 13 13 15 16 16}
\end{tabularx}\endgroup\par\vspace{5pt}\noindent{\fontsize{8.4}{10}\selectfont\hyperref[trace-216]{Procedencia y códigos originales}\hfill Ficha completa · 216/324}
\clearpage\phantomsection\label{nav:vi-ruta-217}
% VI-CATALOG-ROW rutas 217
\pdfbookmark[2]{Celda 55}{cell-55}
\pdfbookmark[3]{r7c1\_h-1v-1}{route-217}
\RouteTitle{Ruta 217}{r7c1\_\allowbreak{}h-1v-1}
\noindent Celda 55\quad (7, 1)\qquad Orientación \(h=-1,\ v=-1\)\hfill\hyperref[sec:vi-indice-rutas]{Índice}\par\vspace{4pt}
\begingroup\fontsize{10.2}{12.5}\selectfont\setlength{\tabcolsep}{5pt}\renewcommand{\arraystretch}{1.1}\rowcolors{1}{routepale}{white}\noindent\begin{tabularx}{\linewidth}{@{}>{\raggedright\arraybackslash}p{.345\linewidth}>{\raggedright\arraybackslash}X@{}}
\RouteGroup{Identidad estructural}
\RouteRow{Clase y multisección}{charged\_\allowbreak{}lepton\_\allowbreak{}G4\quad /\quad charged\_\allowbreak{}lepton\_\allowbreak{}G4}
\RouteRow{Colección y sector}{leptón cargado; leptón}
\RouteRow{Clasificación física}{leptón cargado}
\RouteRow{Generación, profundidad y color}{4; G4; \textemdash{}}
\RouteRow{Ruta primaria y órbita de inversión}{\(B_1\): no\qquad 3,\allowbreak{}9;\allowbreak{}7,\allowbreak{}1}
\RouteRow{Firma de clasificación}{30|\allowbreak{}-6|\allowbreak{}-8|\allowbreak{}0|\allowbreak{}-8|\allowbreak{}+0+-\kern0pt{}-\kern0pt{}-+0+-\kern0pt{}-\kern0pt{}-}
\RouteGroup{Retornos, memoria y transporte}
\RouteRow{Retornos con fase}{clase: 27; órbita: 27; celda: 27}
\RouteRow{Retorno cinemático}{en 108: 54; extensión: 216}
\RouteRow{Giros y cambios de hoja}{71; 107}
\RouteRow{Acarreo y diferimiento}{deuda total: 2; final: 1; bloques diferidos: 2}
\RouteRow{Tríadas finales; persistencia de Pareto}{16; sí}
\RouteRow{\(K\longrightarrow K+9\)}{repeticiones: 9; cambios de memoria: 9}
\RouteGroup{Firma, fase y diferencias}
\RouteRow{\(n_A,\ s_C,\ \kappa_0,\ \nu_{120},\ \nu_{270}\)}{30, 6, -2, 0, -8}
\RouteRow{\(\tau_{12}\); firma histórica \(\mathbb Z^5\)}{+++-0-+++-0-; 30|\allowbreak{}6|\allowbreak{}-2|\allowbreak{}0|\allowbreak{}-8}
\RouteRow{\(D_1,\ D_3,\ D_4,\ D_{27},\ D_{36}\)}{54, 56, 70, 48, 32}
\RouteRow{\(\Omega_{90/120},\ P_6\)}{48, 5}
\RouteGroup{Lectores y factores de acción}
\RouteRow{\(K_D\quad /\quad K_\Omega\)}{(80,\allowbreak{}54,\allowbreak{}7)\quad /\quad (48,\allowbreak{}54,\allowbreak{}5)}
\RouteRow{\(\kappa_D\)}{79.60672\allowbreak{}13362175\allowbreak{}19726524\allowbreak{}96}
\RouteRow{\(\kappa_\Omega\)}{47.60649\allowbreak{}89433721\allowbreak{}35446416\allowbreak{}86}
\RouteRow{\(R_{\mathrm{act},D}\)}{1.000000\allowbreak{}05518988\allowbreak{}96223153\allowbreak{}37}
\RouteRow{\(R_{\mathrm{act},\Omega}\)}{1.000000\allowbreak{}03300471\allowbreak{}80532430\allowbreak{}84}
\RouteGroup{Separación de prefijos}
\RouteRow{Área de ambigüedad; área \(\log_2\)}{58; 30.18947\allowbreak{}50100961\allowbreak{}74}
\RouteRow{Primer bloque de separación}{órbita: \textemdash{}; celda: \textemdash{}; ruta: \textemdash{}}
\RouteRow{Ambigüedad final}{3}
\RouteGroup{Secuencias completas}
\RouteRow{\(w_6\)}{486 423 15 486 429 243 657 24 243 486 423 15 486 429 243 657 24 243}
\RouteRow{Tríadas estables por bloque}{0 1 2 3 4 5 6 7 8 9 10 11 12 13 13 15 16 16}
\end{tabularx}\endgroup\par\vspace{5pt}\noindent{\fontsize{8.4}{10}\selectfont\hyperref[trace-217]{Procedencia y códigos originales}\hfill Ficha completa · 217/324}
\clearpage\phantomsection\label{nav:vi-ruta-218}
% VI-CATALOG-ROW rutas 218
\pdfbookmark[3]{r7c1\_h-1v+1}{route-218}
\RouteTitle{Ruta 218}{r7c1\_\allowbreak{}h-1v+1}
\noindent Celda 55\quad (7, 1)\qquad Orientación \(h=-1,\ v=1\)\hfill\hyperref[sec:vi-indice-rutas]{Índice}\par\vspace{4pt}
\begingroup\fontsize{10.2}{12.5}\selectfont\setlength{\tabcolsep}{5pt}\renewcommand{\arraystretch}{1.1}\rowcolors{1}{routepale}{white}\noindent\begin{tabularx}{\linewidth}{@{}>{\raggedright\arraybackslash}p{.345\linewidth}>{\raggedright\arraybackslash}X@{}}
\RouteGroup{Identidad estructural}
\RouteRow{Clase y multisección}{charged\_\allowbreak{}lepton\_\allowbreak{}G4\quad /\quad charged\_\allowbreak{}lepton\_\allowbreak{}G4}
\RouteRow{Colección y sector}{leptón cargado; leptón}
\RouteRow{Clasificación física}{leptón cargado}
\RouteRow{Generación, profundidad y color}{4; G4; \textemdash{}}
\RouteRow{Ruta primaria y órbita de inversión}{\(B_1\): sí\qquad 3,\allowbreak{}9;\allowbreak{}7,\allowbreak{}1}
\RouteRow{Firma de clasificación}{30|\allowbreak{}-6|\allowbreak{}-8|\allowbreak{}0|\allowbreak{}-8|\allowbreak{}+0+-\kern0pt{}-\kern0pt{}-+0+-\kern0pt{}-\kern0pt{}-}
\RouteGroup{Retornos, memoria y transporte}
\RouteRow{Retornos con fase}{clase: 9; órbita: 27; celda: 27}
\RouteRow{Retorno cinemático}{en 108: 54; extensión: 216}
\RouteRow{Giros y cambios de hoja}{71; 107}
\RouteRow{Acarreo y diferimiento}{deuda total: 0; final: 0; bloques diferidos: 0}
\RouteRow{Tríadas finales; persistencia de Pareto}{17; sí}
\RouteRow{\(K\longrightarrow K+9\)}{repeticiones: 9; cambios de memoria: 9}
\RouteGroup{Firma, fase y diferencias}
\RouteRow{\(n_A,\ s_C,\ \kappa_0,\ \nu_{120},\ \nu_{270}\)}{14, -4, 0, 0, -12}
\RouteRow{\(\tau_{12}\); firma histórica \(\mathbb Z^5\)}{+++-\kern0pt{}-\kern0pt{}-+++-\kern0pt{}-\kern0pt{}-; 14|\allowbreak{}-4|\allowbreak{}0|\allowbreak{}0|\allowbreak{}-12}
\RouteRow{\(D_1,\ D_3,\ D_4,\ D_{27},\ D_{36}\)}{78, 28, 78, 36, 24}
\RouteRow{\(\Omega_{90/120},\ P_6\)}{80, 6}
\RouteGroup{Lectores y factores de acción}
\RouteRow{\(K_D\quad /\quad K_\Omega\)}{(60,\allowbreak{}78,\allowbreak{}25)\quad /\quad (80,\allowbreak{}54,\allowbreak{}6)}
\RouteRow{\(\kappa_D\)}{59.43358\allowbreak{}64101631\allowbreak{}67023563\allowbreak{}04}
\RouteRow{\(\kappa_\Omega\)}{79.60661\allowbreak{}01397948\allowbreak{}27586470\allowbreak{}91}
\RouteRow{\(R_{\mathrm{act},D}\)}{1.000000\allowbreak{}04120422\allowbreak{}24053449\allowbreak{}05}
\RouteRow{\(R_{\mathrm{act},\Omega}\)}{1.000000\allowbreak{}05518981\allowbreak{}25318591\allowbreak{}40}
\RouteGroup{Separación de prefijos}
\RouteRow{Área de ambigüedad; área \(\log_2\)}{54; 28.52932\allowbreak{}50129808}
\RouteRow{Primer bloque de separación}{órbita: \textemdash{}; celda: \textemdash{}; ruta: \textemdash{}}
\RouteRow{Ambigüedad final}{3}
\RouteGroup{Secuencias completas}
\RouteRow{\(w_6\)}{507 504 585 507 510 510 666 672 510 507 504 585 507 510 510 666 672 510}
\RouteRow{Tríadas estables por bloque}{0 1 2 3 4 5 6 7 8 9 10 11 12 13 14 15 16 17}
\end{tabularx}\endgroup\par\vspace{5pt}\noindent{\fontsize{8.4}{10}\selectfont\hyperref[trace-218]{Procedencia y códigos originales}\hfill Ficha completa · 218/324}
\clearpage\phantomsection\label{nav:vi-ruta-219}
% VI-CATALOG-ROW rutas 219
\pdfbookmark[3]{r7c1\_h+1v-1}{route-219}
\RouteTitle{Ruta 219}{r7c1\_\allowbreak{}h+1v-1}
\noindent Celda 55\quad (7, 1)\qquad Orientación \(h=1,\ v=-1\)\hfill\hyperref[sec:vi-indice-rutas]{Índice}\par\vspace{4pt}
\begingroup\fontsize{10.2}{12.5}\selectfont\setlength{\tabcolsep}{5pt}\renewcommand{\arraystretch}{1.1}\rowcolors{1}{routepale}{white}\noindent\begin{tabularx}{\linewidth}{@{}>{\raggedright\arraybackslash}p{.345\linewidth}>{\raggedright\arraybackslash}X@{}}
\RouteGroup{Identidad estructural}
\RouteRow{Clase y multisección}{charged\_\allowbreak{}lepton\_\allowbreak{}G4\quad /\quad charged\_\allowbreak{}lepton\_\allowbreak{}G4}
\RouteRow{Colección y sector}{leptón cargado; leptón}
\RouteRow{Clasificación física}{leptón cargado}
\RouteRow{Generación, profundidad y color}{4; G4; \textemdash{}}
\RouteRow{Ruta primaria y órbita de inversión}{\(B_1\): no\qquad 3,\allowbreak{}9;\allowbreak{}7,\allowbreak{}1}
\RouteRow{Firma de clasificación}{30|\allowbreak{}-6|\allowbreak{}-8|\allowbreak{}0|\allowbreak{}-8|\allowbreak{}+0+-\kern0pt{}-\kern0pt{}-+0+-\kern0pt{}-\kern0pt{}-}
\RouteGroup{Retornos, memoria y transporte}
\RouteRow{Retornos con fase}{clase: 18; órbita: 27; celda: 27}
\RouteRow{Retorno cinemático}{en 108: 54; extensión: 216}
\RouteRow{Giros y cambios de hoja}{71; 107}
\RouteRow{Acarreo y diferimiento}{deuda total: 2; final: 0; bloques diferidos: 1}
\RouteRow{Tríadas finales; persistencia de Pareto}{17; sí}
\RouteRow{\(K\longrightarrow K+9\)}{repeticiones: 9; cambios de memoria: 9}
\RouteGroup{Firma, fase y diferencias}
\RouteRow{\(n_A,\ s_C,\ \kappa_0,\ \nu_{120},\ \nu_{270}\)}{14, 4, 0, 0, -12}
\RouteRow{\(\tau_{12}\); firma histórica \(\mathbb Z^5\)}{+++-\kern0pt{}-\kern0pt{}-+++-\kern0pt{}-\kern0pt{}-; 14|\allowbreak{}4|\allowbreak{}0|\allowbreak{}0|\allowbreak{}-12}
\RouteRow{\(D_1,\ D_3,\ D_4,\ D_{27},\ D_{36}\)}{78, 28, 78, 36, 24}
\RouteRow{\(\Omega_{90/120},\ P_6\)}{80, 6}
\RouteGroup{Lectores y factores de acción}
\RouteRow{\(K_D\quad /\quad K_\Omega\)}{(60,\allowbreak{}78,\allowbreak{}25)\quad /\quad (80,\allowbreak{}54,\allowbreak{}6)}
\RouteRow{\(\kappa_D\)}{59.43358\allowbreak{}64101631\allowbreak{}67023563\allowbreak{}04}
\RouteRow{\(\kappa_\Omega\)}{79.60661\allowbreak{}01397948\allowbreak{}27586470\allowbreak{}91}
\RouteRow{\(R_{\mathrm{act},D}\)}{1.000000\allowbreak{}04120422\allowbreak{}24053449\allowbreak{}05}
\RouteRow{\(R_{\mathrm{act},\Omega}\)}{1.000000\allowbreak{}05518981\allowbreak{}25318591\allowbreak{}40}
\RouteGroup{Separación de prefijos}
\RouteRow{Área de ambigüedad; área \(\log_2\)}{54; 28.52932\allowbreak{}50129808}
\RouteRow{Primer bloque de separación}{órbita: \textemdash{}; celda: \textemdash{}; ruta: \textemdash{}}
\RouteRow{Ambigüedad final}{3}
\RouteGroup{Secuencias completas}
\RouteRow{\(w_6\)}{666 672 510 666 666 585 507 504 585 666 672 510 666 666 585 507 504 585}
\RouteRow{Tríadas estables por bloque}{0 1 2 3 4 5 6 7 8 9 10 11 12 12 13 15 16 17}
\end{tabularx}\endgroup\par\vspace{5pt}\noindent{\fontsize{8.4}{10}\selectfont\hyperref[trace-219]{Procedencia y códigos originales}\hfill Ficha completa · 219/324}
\clearpage\phantomsection\label{nav:vi-ruta-220}
% VI-CATALOG-ROW rutas 220
\pdfbookmark[3]{r7c1\_h+1v+1}{route-220}
\RouteTitle{Ruta 220}{r7c1\_\allowbreak{}h+1v+1}
\noindent Celda 55\quad (7, 1)\qquad Orientación \(h=1,\ v=1\)\hfill\hyperref[sec:vi-indice-rutas]{Índice}\par\vspace{4pt}
\begingroup\fontsize{10.2}{12.5}\selectfont\setlength{\tabcolsep}{5pt}\renewcommand{\arraystretch}{1.1}\rowcolors{1}{routepale}{white}\noindent\begin{tabularx}{\linewidth}{@{}>{\raggedright\arraybackslash}p{.345\linewidth}>{\raggedright\arraybackslash}X@{}}
\RouteGroup{Identidad estructural}
\RouteRow{Clase y multisección}{charged\_\allowbreak{}lepton\_\allowbreak{}G4\quad /\quad charged\_\allowbreak{}lepton\_\allowbreak{}G4}
\RouteRow{Colección y sector}{leptón cargado; leptón}
\RouteRow{Clasificación física}{leptón cargado}
\RouteRow{Generación, profundidad y color}{4; G4; \textemdash{}}
\RouteRow{Ruta primaria y órbita de inversión}{\(B_1\): no\qquad 3,\allowbreak{}9;\allowbreak{}7,\allowbreak{}1}
\RouteRow{Firma de clasificación}{30|\allowbreak{}-6|\allowbreak{}-8|\allowbreak{}0|\allowbreak{}-8|\allowbreak{}+0+-\kern0pt{}-\kern0pt{}-+0+-\kern0pt{}-\kern0pt{}-}
\RouteGroup{Retornos, memoria y transporte}
\RouteRow{Retornos con fase}{clase: 27; órbita: 27; celda: 27}
\RouteRow{Retorno cinemático}{en 108: 54; extensión: 216}
\RouteRow{Giros y cambios de hoja}{71; 107}
\RouteRow{Acarreo y diferimiento}{deuda total: 2; final: 0; bloques diferidos: 1}
\RouteRow{Tríadas finales; persistencia de Pareto}{17; sí}
\RouteRow{\(K\longrightarrow K+9\)}{repeticiones: 9; cambios de memoria: 9}
\RouteGroup{Firma, fase y diferencias}
\RouteRow{\(n_A,\ s_C,\ \kappa_0,\ \nu_{120},\ \nu_{270}\)}{30, -6, 2, 0, -8}
\RouteRow{\(\tau_{12}\); firma histórica \(\mathbb Z^5\)}{+0+-\kern0pt{}-\kern0pt{}-+0+-\kern0pt{}-\kern0pt{}-; 30|\allowbreak{}-6|\allowbreak{}2|\allowbreak{}0|\allowbreak{}-8}
\RouteRow{\(D_1,\ D_3,\ D_4,\ D_{27},\ D_{36}\)}{54, 56, 70, 48, 32}
\RouteRow{\(\Omega_{90/120},\ P_6\)}{48, 5}
\RouteGroup{Lectores y factores de acción}
\RouteRow{\(K_D\quad /\quad K_\Omega\)}{(80,\allowbreak{}54,\allowbreak{}7)\quad /\quad (48,\allowbreak{}54,\allowbreak{}5)}
\RouteRow{\(\kappa_D\)}{79.60672\allowbreak{}13362175\allowbreak{}19726524\allowbreak{}96}
\RouteRow{\(\kappa_\Omega\)}{47.60649\allowbreak{}89433721\allowbreak{}35446416\allowbreak{}86}
\RouteRow{\(R_{\mathrm{act},D}\)}{1.000000\allowbreak{}05518988\allowbreak{}96223153\allowbreak{}37}
\RouteRow{\(R_{\mathrm{act},\Omega}\)}{1.000000\allowbreak{}03300471\allowbreak{}80532430\allowbreak{}84}
\RouteGroup{Separación de prefijos}
\RouteRow{Área de ambigüedad; área \(\log_2\)}{58; 30.18947\allowbreak{}50100961\allowbreak{}74}
\RouteRow{Primer bloque de separación}{órbita: \textemdash{}; celda: \textemdash{}; ruta: \textemdash{}}
\RouteRow{Ambigüedad final}{3}
\RouteGroup{Secuencias completas}
\RouteRow{\(w_6\)}{657 24 243 657 18 15 486 423 15 657 24 243 657 18 15 486 423 15}
\RouteRow{Tríadas estables por bloque}{0 1 2 3 4 5 6 7 8 9 10 11 12 13 14 14 15 17}
\end{tabularx}\endgroup\par\vspace{5pt}\noindent{\fontsize{8.4}{10}\selectfont\hyperref[trace-220]{Procedencia y códigos originales}\hfill Ficha completa · 220/324}
\clearpage\phantomsection\label{nav:vi-ruta-221}
% VI-CATALOG-ROW rutas 221
\pdfbookmark[2]{Celda 56}{cell-56}
\pdfbookmark[3]{r7c2\_h-1v-1}{route-221}
\RouteTitle{Ruta 221}{r7c2\_\allowbreak{}h-1v-1}
\noindent Celda 56\quad (7, 2)\qquad Orientación \(h=-1,\ v=-1\)\hfill\hyperref[sec:vi-indice-rutas]{Índice}\par\vspace{4pt}
\begingroup\fontsize{10.2}{12.5}\selectfont\setlength{\tabcolsep}{5pt}\renewcommand{\arraystretch}{1.1}\rowcolors{1}{routepale}{white}\noindent\begin{tabularx}{\linewidth}{@{}>{\raggedright\arraybackslash}p{.345\linewidth}>{\raggedright\arraybackslash}X@{}}
\RouteGroup{Identidad estructural}
\RouteRow{Clase y multisección}{neutrino\_\allowbreak{}G3\quad /\quad neutrino\_\allowbreak{}G3}
\RouteRow{Colección y sector}{neutrino; leptón}
\RouteRow{Clasificación física}{neutrino}
\RouteRow{Generación, profundidad y color}{3; G3; \textemdash{}}
\RouteRow{Ruta primaria y órbita de inversión}{\(B_1\): no\qquad 3,\allowbreak{}8;\allowbreak{}7,\allowbreak{}2}
\RouteRow{Firma de clasificación}{34|\allowbreak{}0|\allowbreak{}-4|\allowbreak{}-2|\allowbreak{}-8|\allowbreak{}++0-\kern0pt{}-\kern0pt{}-++0-\kern0pt{}-\kern0pt{}-}
\RouteGroup{Retornos, memoria y transporte}
\RouteRow{Retornos con fase}{clase: 27; órbita: 27; celda: 27}
\RouteRow{Retorno cinemático}{en 108: 54; extensión: 216}
\RouteRow{Giros y cambios de hoja}{71; 107}
\RouteRow{Acarreo y diferimiento}{deuda total: 1; final: 0; bloques diferidos: 1}
\RouteRow{Tríadas finales; persistencia de Pareto}{17; sí}
\RouteRow{\(K\longrightarrow K+9\)}{repeticiones: 9; cambios de memoria: 9}
\RouteGroup{Firma, fase y diferencias}
\RouteRow{\(n_A,\ s_C,\ \kappa_0,\ \nu_{120},\ \nu_{270}\)}{34, 0, -2, 0, -8}
\RouteRow{\(\tau_{12}\); firma histórica \(\mathbb Z^5\)}{+++-\kern0pt{}-0+++-\kern0pt{}-0; 34|\allowbreak{}0|\allowbreak{}-2|\allowbreak{}0|\allowbreak{}-8}
\RouteRow{\(D_1,\ D_3,\ D_4,\ D_{27},\ D_{36}\)}{66, 64, 64, 60, 40}
\RouteRow{\(\Omega_{90/120},\ P_6\)}{56, 5}
\RouteGroup{Lectores y factores de acción}
\RouteRow{\(K_D\quad /\quad K_\Omega\)}{(100,\allowbreak{}66,\allowbreak{}0)\quad /\quad (56,\allowbreak{}54,\allowbreak{}5)}
\RouteRow{\(\kappa_D\)}{99.51837\allowbreak{}47304272\allowbreak{}69134179\allowbreak{}18}
\RouteRow{\(\kappa_\Omega\)}{55.60649\allowbreak{}89433721\allowbreak{}35446416\allowbreak{}86}
\RouteRow{\(R_{\mathrm{act},D}\)}{1.000000\allowbreak{}06899427\allowbreak{}66450218\allowbreak{}95}
\RouteRow{\(R_{\mathrm{act},\Omega}\)}{1.000000\allowbreak{}03855097\allowbreak{}23541416\allowbreak{}45}
\RouteGroup{Separación de prefijos}
\RouteRow{Área de ambigüedad; área \(\log_2\)}{36; 18.0}
\RouteRow{Primer bloque de separación}{órbita: \textemdash{}; celda: \textemdash{}; ruta: \textemdash{}}
\RouteRow{Ambigüedad final}{2}
\RouteGroup{Secuencias completas}
\RouteRow{\(w_6\)}{90 570 264 90 567 657 24 243 657 90 570 264 90 567 657 24 243 657}
\RouteRow{Tríadas estables por bloque}{0 1 2 3 4 5 6 7 8 9 10 11 12 13 13 15 16 17}
\end{tabularx}\endgroup\par\vspace{5pt}\noindent{\fontsize{8.4}{10}\selectfont\hyperref[trace-221]{Procedencia y códigos originales}\hfill Ficha completa · 221/324}
\clearpage\phantomsection\label{nav:vi-ruta-222}
% VI-CATALOG-ROW rutas 222
\pdfbookmark[3]{r7c2\_h-1v+1}{route-222}
\RouteTitle{Ruta 222}{r7c2\_\allowbreak{}h-1v+1}
\noindent Celda 56\quad (7, 2)\qquad Orientación \(h=-1,\ v=1\)\hfill\hyperref[sec:vi-indice-rutas]{Índice}\par\vspace{4pt}
\begingroup\fontsize{10.2}{12.5}\selectfont\setlength{\tabcolsep}{5pt}\renewcommand{\arraystretch}{1.1}\rowcolors{1}{routepale}{white}\noindent\begin{tabularx}{\linewidth}{@{}>{\raggedright\arraybackslash}p{.345\linewidth}>{\raggedright\arraybackslash}X@{}}
\RouteGroup{Identidad estructural}
\RouteRow{Clase y multisección}{neutrino\_\allowbreak{}G3\quad /\quad neutrino\_\allowbreak{}G3}
\RouteRow{Colección y sector}{neutrino; leptón}
\RouteRow{Clasificación física}{neutrino}
\RouteRow{Generación, profundidad y color}{3; G3; \textemdash{}}
\RouteRow{Ruta primaria y órbita de inversión}{\(B_1\): sí\qquad 3,\allowbreak{}8;\allowbreak{}7,\allowbreak{}2}
\RouteRow{Firma de clasificación}{34|\allowbreak{}0|\allowbreak{}-4|\allowbreak{}-2|\allowbreak{}-8|\allowbreak{}++0-\kern0pt{}-\kern0pt{}-++0-\kern0pt{}-\kern0pt{}-}
\RouteGroup{Retornos, memoria y transporte}
\RouteRow{Retornos con fase}{clase: 27; órbita: 27; celda: 27}
\RouteRow{Retorno cinemático}{en 108: 54; extensión: 216}
\RouteRow{Giros y cambios de hoja}{71; 107}
\RouteRow{Acarreo y diferimiento}{deuda total: 2; final: 0; bloques diferidos: 2}
\RouteRow{Tríadas finales; persistencia de Pareto}{17; sí}
\RouteRow{\(K\longrightarrow K+9\)}{repeticiones: 9; cambios de memoria: 9}
\RouteGroup{Firma, fase y diferencias}
\RouteRow{\(n_A,\ s_C,\ \kappa_0,\ \nu_{120},\ \nu_{270}\)}{20, 0, 0, 0, -4}
\RouteRow{\(\tau_{12}\); firma histórica \(\mathbb Z^5\)}{0++-\kern0pt{}-00++-\kern0pt{}-0; 20|\allowbreak{}0|\allowbreak{}0|\allowbreak{}0|\allowbreak{}-4}
\RouteRow{\(D_1,\ D_3,\ D_4,\ D_{27},\ D_{36}\)}{24, 20, 24, 24, 16}
\RouteRow{\(\Omega_{90/120},\ P_6\)}{36, 4}
\RouteGroup{Lectores y factores de acción}
\RouteRow{\(K_D\quad /\quad K_\Omega\)}{(40,\allowbreak{}24,\allowbreak{}2)\quad /\quad (36,\allowbreak{}54,\allowbreak{}4)}
\RouteRow{\(\kappa_D\)}{39.82508\allowbreak{}59311825\allowbreak{}73056173\allowbreak{}26}
\RouteRow{\(\kappa_\Omega\)}{35.60638\allowbreak{}77469494\allowbreak{}43306362\allowbreak{}81}
\RouteRow{\(R_{\mathrm{act},D}\)}{1.000000\allowbreak{}02761000\allowbreak{}61595142\allowbreak{}15}
\RouteRow{\(R_{\mathrm{act},\Omega}\)}{1.000000\allowbreak{}02468525\allowbreak{}95691181\allowbreak{}52}
\RouteGroup{Separación de prefijos}
\RouteRow{Área de ambigüedad; área \(\log_2\)}{36; 18.0}
\RouteRow{Primer bloque de separación}{órbita: \textemdash{}; celda: \textemdash{}; ruta: \textemdash{}}
\RouteRow{Ambigüedad final}{2}
\RouteGroup{Secuencias completas}
\RouteRow{\(w_6\)}{81 3 0 81 0 3 0 81 3 81 3 0 81 0 3 0 81 3}
\RouteRow{Tríadas estables por bloque}{0 1 2 3 4 5 6 6 8 9 10 11 12 13 14 15 15 17}
\end{tabularx}\endgroup\par\vspace{5pt}\noindent{\fontsize{8.4}{10}\selectfont\hyperref[trace-222]{Procedencia y códigos originales}\hfill Ficha completa · 222/324}
\clearpage\phantomsection\label{nav:vi-ruta-223}
% VI-CATALOG-ROW rutas 223
\pdfbookmark[3]{r7c2\_h+1v-1}{route-223}
\RouteTitle{Ruta 223}{r7c2\_\allowbreak{}h+1v-1}
\noindent Celda 56\quad (7, 2)\qquad Orientación \(h=1,\ v=-1\)\hfill\hyperref[sec:vi-indice-rutas]{Índice}\par\vspace{4pt}
\begingroup\fontsize{10.2}{12.5}\selectfont\setlength{\tabcolsep}{5pt}\renewcommand{\arraystretch}{1.1}\rowcolors{1}{routepale}{white}\noindent\begin{tabularx}{\linewidth}{@{}>{\raggedright\arraybackslash}p{.345\linewidth}>{\raggedright\arraybackslash}X@{}}
\RouteGroup{Identidad estructural}
\RouteRow{Clase y multisección}{neutrino\_\allowbreak{}G3\quad /\quad neutrino\_\allowbreak{}G3}
\RouteRow{Colección y sector}{neutrino; leptón}
\RouteRow{Clasificación física}{neutrino}
\RouteRow{Generación, profundidad y color}{3; G3; \textemdash{}}
\RouteRow{Ruta primaria y órbita de inversión}{\(B_1\): no\qquad 3,\allowbreak{}8;\allowbreak{}7,\allowbreak{}2}
\RouteRow{Firma de clasificación}{34|\allowbreak{}0|\allowbreak{}-4|\allowbreak{}-2|\allowbreak{}-8|\allowbreak{}++0-\kern0pt{}-\kern0pt{}-++0-\kern0pt{}-\kern0pt{}-}
\RouteGroup{Retornos, memoria y transporte}
\RouteRow{Retornos con fase}{clase: 27; órbita: 27; celda: 27}
\RouteRow{Retorno cinemático}{en 108: 54; extensión: 216}
\RouteRow{Giros y cambios de hoja}{71; 107}
\RouteRow{Acarreo y diferimiento}{deuda total: 0; final: 0; bloques diferidos: 0}
\RouteRow{Tríadas finales; persistencia de Pareto}{17; sí}
\RouteRow{\(K\longrightarrow K+9\)}{repeticiones: 9; cambios de memoria: 9}
\RouteGroup{Firma, fase y diferencias}
\RouteRow{\(n_A,\ s_C,\ \kappa_0,\ \nu_{120},\ \nu_{270}\)}{20, 0, 0, 0, -4}
\RouteRow{\(\tau_{12}\); firma histórica \(\mathbb Z^5\)}{++00-\kern0pt{}-++00-\kern0pt{}-; 20|\allowbreak{}0|\allowbreak{}0|\allowbreak{}0|\allowbreak{}-4}
\RouteRow{\(D_1,\ D_3,\ D_4,\ D_{27},\ D_{36}\)}{24, 20, 24, 24, 16}
\RouteRow{\(\Omega_{90/120},\ P_6\)}{36, 4}
\RouteGroup{Lectores y factores de acción}
\RouteRow{\(K_D\quad /\quad K_\Omega\)}{(40,\allowbreak{}24,\allowbreak{}2)\quad /\quad (36,\allowbreak{}54,\allowbreak{}4)}
\RouteRow{\(\kappa_D\)}{39.82508\allowbreak{}59311825\allowbreak{}73056173\allowbreak{}26}
\RouteRow{\(\kappa_\Omega\)}{35.60638\allowbreak{}77469494\allowbreak{}43306362\allowbreak{}81}
\RouteRow{\(R_{\mathrm{act},D}\)}{1.000000\allowbreak{}02761000\allowbreak{}61595142\allowbreak{}15}
\RouteRow{\(R_{\mathrm{act},\Omega}\)}{1.000000\allowbreak{}02468525\allowbreak{}95691181\allowbreak{}52}
\RouteGroup{Separación de prefijos}
\RouteRow{Área de ambigüedad; área \(\log_2\)}{36; 18.0}
\RouteRow{Primer bloque de separación}{órbita: \textemdash{}; celda: \textemdash{}; ruta: \textemdash{}}
\RouteRow{Ambigüedad final}{2}
\RouteGroup{Secuencias completas}
\RouteRow{\(w_6\)}{0 81 3 0 84 0 81 3 0 0 81 3 0 84 0 81 3 0}
\RouteRow{Tríadas estables por bloque}{0 1 2 3 4 5 6 7 8 9 10 11 12 13 14 15 16 17}
\end{tabularx}\endgroup\par\vspace{5pt}\noindent{\fontsize{8.4}{10}\selectfont\hyperref[trace-223]{Procedencia y códigos originales}\hfill Ficha completa · 223/324}
\clearpage\phantomsection\label{nav:vi-ruta-224}
% VI-CATALOG-ROW rutas 224
\pdfbookmark[3]{r7c2\_h+1v+1}{route-224}
\RouteTitle{Ruta 224}{r7c2\_\allowbreak{}h+1v+1}
\noindent Celda 56\quad (7, 2)\qquad Orientación \(h=1,\ v=1\)\hfill\hyperref[sec:vi-indice-rutas]{Índice}\par\vspace{4pt}
\begingroup\fontsize{10.2}{12.5}\selectfont\setlength{\tabcolsep}{5pt}\renewcommand{\arraystretch}{1.1}\rowcolors{1}{routepale}{white}\noindent\begin{tabularx}{\linewidth}{@{}>{\raggedright\arraybackslash}p{.345\linewidth}>{\raggedright\arraybackslash}X@{}}
\RouteGroup{Identidad estructural}
\RouteRow{Clase y multisección}{neutrino\_\allowbreak{}G3\quad /\quad neutrino\_\allowbreak{}G3}
\RouteRow{Colección y sector}{neutrino; leptón}
\RouteRow{Clasificación física}{neutrino}
\RouteRow{Generación, profundidad y color}{3; G3; \textemdash{}}
\RouteRow{Ruta primaria y órbita de inversión}{\(B_1\): no\qquad 3,\allowbreak{}8;\allowbreak{}7,\allowbreak{}2}
\RouteRow{Firma de clasificación}{34|\allowbreak{}0|\allowbreak{}-4|\allowbreak{}-2|\allowbreak{}-8|\allowbreak{}++0-\kern0pt{}-\kern0pt{}-++0-\kern0pt{}-\kern0pt{}-}
\RouteGroup{Retornos, memoria y transporte}
\RouteRow{Retornos con fase}{clase: 27; órbita: 27; celda: 27}
\RouteRow{Retorno cinemático}{en 108: 54; extensión: 216}
\RouteRow{Giros y cambios de hoja}{71; 107}
\RouteRow{Acarreo y diferimiento}{deuda total: 0; final: 0; bloques diferidos: 0}
\RouteRow{Tríadas finales; persistencia de Pareto}{17; sí}
\RouteRow{\(K\longrightarrow K+9\)}{repeticiones: 9; cambios de memoria: 9}
\RouteGroup{Firma, fase y diferencias}
\RouteRow{\(n_A,\ s_C,\ \kappa_0,\ \nu_{120},\ \nu_{270}\)}{34, 0, 2, 0, -8}
\RouteRow{\(\tau_{12}\); firma histórica \(\mathbb Z^5\)}{++0-\kern0pt{}-\kern0pt{}-++0-\kern0pt{}-\kern0pt{}-; 34|\allowbreak{}0|\allowbreak{}2|\allowbreak{}0|\allowbreak{}-8}
\RouteRow{\(D_1,\ D_3,\ D_4,\ D_{27},\ D_{36}\)}{66, 64, 64, 60, 40}
\RouteRow{\(\Omega_{90/120},\ P_6\)}{56, 5}
\RouteGroup{Lectores y factores de acción}
\RouteRow{\(K_D\quad /\quad K_\Omega\)}{(100,\allowbreak{}66,\allowbreak{}0)\quad /\quad (56,\allowbreak{}54,\allowbreak{}5)}
\RouteRow{\(\kappa_D\)}{99.51837\allowbreak{}47304272\allowbreak{}69134179\allowbreak{}18}
\RouteRow{\(\kappa_\Omega\)}{55.60649\allowbreak{}89433721\allowbreak{}35446416\allowbreak{}86}
\RouteRow{\(R_{\mathrm{act},D}\)}{1.000000\allowbreak{}06899427\allowbreak{}66450218\allowbreak{}95}
\RouteRow{\(R_{\mathrm{act},\Omega}\)}{1.000000\allowbreak{}03855097\allowbreak{}23541416\allowbreak{}45}
\RouteGroup{Separación de prefijos}
\RouteRow{Área de ambigüedad; área \(\log_2\)}{36; 18.0}
\RouteRow{Primer bloque de separación}{órbita: \textemdash{}; celda: \textemdash{}; ruta: \textemdash{}}
\RouteRow{Ambigüedad final}{2}
\RouteGroup{Secuencias completas}
\RouteRow{\(w_6\)}{24 243 657 24 246 264 90 570 264 24 243 657 24 246 264 90 570 264}
\RouteRow{Tríadas estables por bloque}{0 1 2 3 4 5 6 7 8 9 10 11 12 13 14 15 16 17}
\end{tabularx}\endgroup\par\vspace{5pt}\noindent{\fontsize{8.4}{10}\selectfont\hyperref[trace-224]{Procedencia y códigos originales}\hfill Ficha completa · 224/324}
\clearpage\phantomsection\label{nav:vi-ruta-225}
% VI-CATALOG-ROW rutas 225
\pdfbookmark[2]{Celda 57}{cell-57}
\pdfbookmark[3]{r7c3\_h-1v-1}{route-225}
\RouteTitle{Ruta 225}{r7c3\_\allowbreak{}h-1v-1}
\noindent Celda 57\quad (7, 3)\qquad Orientación \(h=-1,\ v=-1\)\hfill\hyperref[sec:vi-indice-rutas]{Índice}\par\vspace{4pt}
\begingroup\fontsize{10.2}{12.5}\selectfont\setlength{\tabcolsep}{5pt}\renewcommand{\arraystretch}{1.1}\rowcolors{1}{routepale}{white}\noindent\begin{tabularx}{\linewidth}{@{}>{\raggedright\arraybackslash}p{.345\linewidth}>{\raggedright\arraybackslash}X@{}}
\RouteGroup{Identidad estructural}
\RouteRow{Clase y multisección}{charged\_\allowbreak{}lepton\_\allowbreak{}G2\quad /\quad charged\_\allowbreak{}lepton\_\allowbreak{}G2}
\RouteRow{Colección y sector}{leptón cargado; leptón}
\RouteRow{Clasificación física}{leptón cargado}
\RouteRow{Generación, profundidad y color}{2; G2; \textemdash{}}
\RouteRow{Ruta primaria y órbita de inversión}{\(B_1\): no\qquad 3,\allowbreak{}7;\allowbreak{}7,\allowbreak{}3}
\RouteRow{Firma de clasificación}{34|\allowbreak{}0|\allowbreak{}-12|\allowbreak{}0|\allowbreak{}-12|\allowbreak{}+++-\kern0pt{}-\kern0pt{}-+++-\kern0pt{}-\kern0pt{}-}
\RouteGroup{Retornos, memoria y transporte}
\RouteRow{Retornos con fase}{clase: 27; órbita: 27; celda: 27}
\RouteRow{Retorno cinemático}{en 108: 54; extensión: 216}
\RouteRow{Giros y cambios de hoja}{71; 107}
\RouteRow{Acarreo y diferimiento}{deuda total: 1; final: 1; bloques diferidos: 1}
\RouteRow{Tríadas finales; persistencia de Pareto}{16; sí}
\RouteRow{\(K\longrightarrow K+9\)}{repeticiones: 9; cambios de memoria: 9}
\RouteGroup{Firma, fase y diferencias}
\RouteRow{\(n_A,\ s_C,\ \kappa_0,\ \nu_{120},\ \nu_{270}\)}{34, 0, 0, 0, -12}
\RouteRow{\(\tau_{12}\); firma histórica \(\mathbb Z^5\)}{+++-\kern0pt{}-\kern0pt{}-+++-\kern0pt{}-\kern0pt{}-; 34|\allowbreak{}0|\allowbreak{}0|\allowbreak{}0|\allowbreak{}-12}
\RouteRow{\(D_1,\ D_3,\ D_4,\ D_{27},\ D_{36}\)}{66, 60, 64, 48, 32}
\RouteRow{\(\Omega_{90/120},\ P_6\)}{80, 6}
\RouteGroup{Lectores y factores de acción}
\RouteRow{\(K_D\quad /\quad K_\Omega\)}{(80,\allowbreak{}66,\allowbreak{}2)\quad /\quad (80,\allowbreak{}54,\allowbreak{}6)}
\RouteRow{\(\kappa_D\)}{79.51859\allowbreak{}71232726\allowbreak{}53414287\allowbreak{}29}
\RouteRow{\(\kappa_\Omega\)}{79.60661\allowbreak{}01397948\allowbreak{}27586470\allowbreak{}91}
\RouteRow{\(R_{\mathrm{act},D}\)}{1.000000\allowbreak{}05512879\allowbreak{}47092488\allowbreak{}80}
\RouteRow{\(R_{\mathrm{act},\Omega}\)}{1.000000\allowbreak{}05518981\allowbreak{}25318591\allowbreak{}40}
\RouteGroup{Separación de prefijos}
\RouteRow{Área de ambigüedad; área \(\log_2\)}{18; 0.0}
\RouteRow{Primer bloque de separación}{órbita: 1; celda: 1; ruta: 1}
\RouteRow{Ambigüedad final}{1}
\RouteGroup{Secuencias completas}
\RouteRow{\(w_6\)}{423 15 486 423 12 99 327 498 99 423 15 486 423 12 99 327 498 99}
\RouteRow{Tríadas estables por bloque}{0 1 2 3 4 5 6 7 8 9 10 11 12 13 14 15 16 16}
\end{tabularx}\endgroup\par\vspace{5pt}\noindent{\fontsize{8.4}{10}\selectfont\hyperref[trace-225]{Procedencia y códigos originales}\hfill Ficha completa · 225/324}
\clearpage\phantomsection\label{nav:vi-ruta-226}
% VI-CATALOG-ROW rutas 226
\pdfbookmark[3]{r7c3\_h-1v+1}{route-226}
\RouteTitle{Ruta 226}{r7c3\_\allowbreak{}h-1v+1}
\noindent Celda 57\quad (7, 3)\qquad Orientación \(h=-1,\ v=1\)\hfill\hyperref[sec:vi-indice-rutas]{Índice}\par\vspace{4pt}
\begingroup\fontsize{10.2}{12.5}\selectfont\setlength{\tabcolsep}{5pt}\renewcommand{\arraystretch}{1.1}\rowcolors{1}{routepale}{white}\noindent\begin{tabularx}{\linewidth}{@{}>{\raggedright\arraybackslash}p{.345\linewidth}>{\raggedright\arraybackslash}X@{}}
\RouteGroup{Identidad estructural}
\RouteRow{Clase y multisección}{charged\_\allowbreak{}lepton\_\allowbreak{}G2\quad /\quad charged\_\allowbreak{}lepton\_\allowbreak{}G2}
\RouteRow{Colección y sector}{leptón cargado; leptón}
\RouteRow{Clasificación física}{leptón cargado}
\RouteRow{Generación, profundidad y color}{2; G2; \textemdash{}}
\RouteRow{Ruta primaria y órbita de inversión}{\(B_1\): no\qquad 3,\allowbreak{}7;\allowbreak{}7,\allowbreak{}3}
\RouteRow{Firma de clasificación}{34|\allowbreak{}0|\allowbreak{}-12|\allowbreak{}0|\allowbreak{}-12|\allowbreak{}+++-\kern0pt{}-\kern0pt{}-+++-\kern0pt{}-\kern0pt{}-}
\RouteGroup{Retornos, memoria y transporte}
\RouteRow{Retornos con fase}{clase: 27; órbita: 27; celda: 27}
\RouteRow{Retorno cinemático}{en 108: 54; extensión: 216}
\RouteRow{Giros y cambios de hoja}{71; 107}
\RouteRow{Acarreo y diferimiento}{deuda total: 1; final: 0; bloques diferidos: 1}
\RouteRow{Tríadas finales; persistencia de Pareto}{17; sí}
\RouteRow{\(K\longrightarrow K+9\)}{repeticiones: 9; cambios de memoria: 9}
\RouteGroup{Firma, fase y diferencias}
\RouteRow{\(n_A,\ s_C,\ \kappa_0,\ \nu_{120},\ \nu_{270}\)}{50, 4, -2, 2, -8}
\RouteRow{\(\tau_{12}\); firma histórica \(\mathbb Z^5\)}{+++0-\kern0pt{}-+++0-\kern0pt{}-; 50|\allowbreak{}4|\allowbreak{}-2|\allowbreak{}2|\allowbreak{}-8}
\RouteRow{\(D_1,\ D_3,\ D_4,\ D_{27},\ D_{36}\)}{84, 24, 84, 24, 16}
\RouteRow{\(\Omega_{90/120},\ P_6\)}{56, 5}
\RouteGroup{Lectores y factores de acción}
\RouteRow{\(K_D\quad /\quad K_\Omega\)}{(40,\allowbreak{}84,\allowbreak{}30)\quad /\quad (56,\allowbreak{}54,\allowbreak{}5)}
\RouteRow{\(\kappa_D\)}{39.39035\allowbreak{}82768609\allowbreak{}24917849\allowbreak{}59}
\RouteRow{\(\kappa_\Omega\)}{55.60649\allowbreak{}89433721\allowbreak{}35446416\allowbreak{}86}
\RouteRow{\(R_{\mathrm{act},D}\)}{1.000000\allowbreak{}02730861\allowbreak{}73967087\allowbreak{}05}
\RouteRow{\(R_{\mathrm{act},\Omega}\)}{1.000000\allowbreak{}03855097\allowbreak{}23541416\allowbreak{}45}
\RouteGroup{Separación de prefijos}
\RouteRow{Área de ambigüedad; área \(\log_2\)}{18; 0.0}
\RouteRow{Primer bloque de separación}{órbita: 1; celda: 1; ruta: 1}
\RouteRow{Ambigüedad final}{1}
\RouteGroup{Secuencias completas}
\RouteRow{\(w_6\)}{420 258 414 420 255 252 333 255 252 420 258 414 420 255 252 333 255 252}
\RouteRow{Tríadas estables por bloque}{0 1 2 3 4 5 6 7 8 9 10 11 12 13 14 15 15 17}
\end{tabularx}\endgroup\par\vspace{5pt}\noindent{\fontsize{8.4}{10}\selectfont\hyperref[trace-226]{Procedencia y códigos originales}\hfill Ficha completa · 226/324}
\clearpage\phantomsection\label{nav:vi-ruta-227}
% VI-CATALOG-ROW rutas 227
\pdfbookmark[3]{r7c3\_h+1v-1}{route-227}
\RouteTitle{Ruta 227}{r7c3\_\allowbreak{}h+1v-1}
\noindent Celda 57\quad (7, 3)\qquad Orientación \(h=1,\ v=-1\)\hfill\hyperref[sec:vi-indice-rutas]{Índice}\par\vspace{4pt}
\begingroup\fontsize{10.2}{12.5}\selectfont\setlength{\tabcolsep}{5pt}\renewcommand{\arraystretch}{1.1}\rowcolors{1}{routepale}{white}\noindent\begin{tabularx}{\linewidth}{@{}>{\raggedright\arraybackslash}p{.345\linewidth}>{\raggedright\arraybackslash}X@{}}
\RouteGroup{Identidad estructural}
\RouteRow{Clase y multisección}{charged\_\allowbreak{}lepton\_\allowbreak{}G2\quad /\quad charged\_\allowbreak{}lepton\_\allowbreak{}G2}
\RouteRow{Colección y sector}{leptón cargado; leptón}
\RouteRow{Clasificación física}{leptón cargado}
\RouteRow{Generación, profundidad y color}{2; G2; \textemdash{}}
\RouteRow{Ruta primaria y órbita de inversión}{\(B_1\): no\qquad 3,\allowbreak{}7;\allowbreak{}7,\allowbreak{}3}
\RouteRow{Firma de clasificación}{34|\allowbreak{}0|\allowbreak{}-12|\allowbreak{}0|\allowbreak{}-12|\allowbreak{}+++-\kern0pt{}-\kern0pt{}-+++-\kern0pt{}-\kern0pt{}-}
\RouteGroup{Retornos, memoria y transporte}
\RouteRow{Retornos con fase}{clase: 27; órbita: 27; celda: 27}
\RouteRow{Retorno cinemático}{en 108: 54; extensión: 216}
\RouteRow{Giros y cambios de hoja}{71; 107}
\RouteRow{Acarreo y diferimiento}{deuda total: 0; final: 0; bloques diferidos: 0}
\RouteRow{Tríadas finales; persistencia de Pareto}{17; sí}
\RouteRow{\(K\longrightarrow K+9\)}{repeticiones: 9; cambios de memoria: 9}
\RouteGroup{Firma, fase y diferencias}
\RouteRow{\(n_A,\ s_C,\ \kappa_0,\ \nu_{120},\ \nu_{270}\)}{50, -4, 2, -2, -8}
\RouteRow{\(\tau_{12}\); firma histórica \(\mathbb Z^5\)}{0++-\kern0pt{}-\kern0pt{}-0++-\kern0pt{}-\kern0pt{}-; 50|\allowbreak{}-4|\allowbreak{}2|\allowbreak{}-2|\allowbreak{}-8}
\RouteRow{\(D_1,\ D_3,\ D_4,\ D_{27},\ D_{36}\)}{84, 24, 84, 24, 16}
\RouteRow{\(\Omega_{90/120},\ P_6\)}{56, 5}
\RouteGroup{Lectores y factores de acción}
\RouteRow{\(K_D\quad /\quad K_\Omega\)}{(40,\allowbreak{}84,\allowbreak{}30)\quad /\quad (56,\allowbreak{}54,\allowbreak{}5)}
\RouteRow{\(\kappa_D\)}{39.39035\allowbreak{}82768609\allowbreak{}24917849\allowbreak{}59}
\RouteRow{\(\kappa_\Omega\)}{55.60649\allowbreak{}89433721\allowbreak{}35446416\allowbreak{}86}
\RouteRow{\(R_{\mathrm{act},D}\)}{1.000000\allowbreak{}02730861\allowbreak{}73967087\allowbreak{}05}
\RouteRow{\(R_{\mathrm{act},\Omega}\)}{1.000000\allowbreak{}03855097\allowbreak{}23541416\allowbreak{}45}
\RouteGroup{Separación de prefijos}
\RouteRow{Área de ambigüedad; área \(\log_2\)}{18; 0.0}
\RouteRow{Primer bloque de separación}{órbita: 1; celda: 1; ruta: 1}
\RouteRow{Ambigüedad final}{1}
\RouteGroup{Secuencias completas}
\RouteRow{\(w_6\)}{333 255 252 333 258 414 420 258 414 333 255 252 333 258 414 420 258 414}
\RouteRow{Tríadas estables por bloque}{0 1 2 3 4 5 6 7 8 9 10 11 12 13 14 15 16 17}
\end{tabularx}\endgroup\par\vspace{5pt}\noindent{\fontsize{8.4}{10}\selectfont\hyperref[trace-227]{Procedencia y códigos originales}\hfill Ficha completa · 227/324}
\clearpage\phantomsection\label{nav:vi-ruta-228}
% VI-CATALOG-ROW rutas 228
\pdfbookmark[3]{r7c3\_h+1v+1}{route-228}
\RouteTitle{Ruta 228}{r7c3\_\allowbreak{}h+1v+1}
\noindent Celda 57\quad (7, 3)\qquad Orientación \(h=1,\ v=1\)\hfill\hyperref[sec:vi-indice-rutas]{Índice}\par\vspace{4pt}
\begingroup\fontsize{10.2}{12.5}\selectfont\setlength{\tabcolsep}{5pt}\renewcommand{\arraystretch}{1.1}\rowcolors{1}{routepale}{white}\noindent\begin{tabularx}{\linewidth}{@{}>{\raggedright\arraybackslash}p{.345\linewidth}>{\raggedright\arraybackslash}X@{}}
\RouteGroup{Identidad estructural}
\RouteRow{Clase y multisección}{charged\_\allowbreak{}lepton\_\allowbreak{}G2\quad /\quad charged\_\allowbreak{}lepton\_\allowbreak{}G2}
\RouteRow{Colección y sector}{leptón cargado; leptón}
\RouteRow{Clasificación física}{leptón cargado}
\RouteRow{Generación, profundidad y color}{2; G2; \textemdash{}}
\RouteRow{Ruta primaria y órbita de inversión}{\(B_1\): sí\qquad 3,\allowbreak{}7;\allowbreak{}7,\allowbreak{}3}
\RouteRow{Firma de clasificación}{34|\allowbreak{}0|\allowbreak{}-12|\allowbreak{}0|\allowbreak{}-12|\allowbreak{}+++-\kern0pt{}-\kern0pt{}-+++-\kern0pt{}-\kern0pt{}-}
\RouteGroup{Retornos, memoria y transporte}
\RouteRow{Retornos con fase}{clase: 27; órbita: 27; celda: 27}
\RouteRow{Retorno cinemático}{en 108: 54; extensión: 216}
\RouteRow{Giros y cambios de hoja}{71; 107}
\RouteRow{Acarreo y diferimiento}{deuda total: 2; final: 1; bloques diferidos: 2}
\RouteRow{Tríadas finales; persistencia de Pareto}{16; sí}
\RouteRow{\(K\longrightarrow K+9\)}{repeticiones: 9; cambios de memoria: 9}
\RouteGroup{Firma, fase y diferencias}
\RouteRow{\(n_A,\ s_C,\ \kappa_0,\ \nu_{120},\ \nu_{270}\)}{34, 0, 0, 0, -12}
\RouteRow{\(\tau_{12}\); firma histórica \(\mathbb Z^5\)}{+++-\kern0pt{}-\kern0pt{}-+++-\kern0pt{}-\kern0pt{}-; 34|\allowbreak{}0|\allowbreak{}0|\allowbreak{}0|\allowbreak{}-12}
\RouteRow{\(D_1,\ D_3,\ D_4,\ D_{27},\ D_{36}\)}{66, 60, 64, 48, 32}
\RouteRow{\(\Omega_{90/120},\ P_6\)}{80, 6}
\RouteGroup{Lectores y factores de acción}
\RouteRow{\(K_D\quad /\quad K_\Omega\)}{(80,\allowbreak{}66,\allowbreak{}2)\quad /\quad (80,\allowbreak{}54,\allowbreak{}6)}
\RouteRow{\(\kappa_D\)}{79.51859\allowbreak{}71232726\allowbreak{}53414287\allowbreak{}29}
\RouteRow{\(\kappa_\Omega\)}{79.60661\allowbreak{}01397948\allowbreak{}27586470\allowbreak{}91}
\RouteRow{\(R_{\mathrm{act},D}\)}{1.000000\allowbreak{}05512879\allowbreak{}47092488\allowbreak{}80}
\RouteRow{\(R_{\mathrm{act},\Omega}\)}{1.000000\allowbreak{}05518981\allowbreak{}25318591\allowbreak{}40}
\RouteGroup{Separación de prefijos}
\RouteRow{Área de ambigüedad; área \(\log_2\)}{18; 0.0}
\RouteRow{Primer bloque de separación}{órbita: 1; celda: 1; ruta: 1}
\RouteRow{Ambigüedad final}{1}
\RouteGroup{Secuencias completas}
\RouteRow{\(w_6\)}{327 498 99 327 501 486 423 15 486 327 498 99 327 501 486 423 15 486}
\RouteRow{Tríadas estables por bloque}{0 1 2 3 4 5 6 7 7 9 10 11 12 13 14 15 16 16}
\end{tabularx}\endgroup\par\vspace{5pt}\noindent{\fontsize{8.4}{10}\selectfont\hyperref[trace-228]{Procedencia y códigos originales}\hfill Ficha completa · 228/324}
\clearpage\phantomsection\label{nav:vi-ruta-229}
% VI-CATALOG-ROW rutas 229
\pdfbookmark[2]{Celda 58}{cell-58}
\pdfbookmark[3]{r7c4\_h-1v-1}{route-229}
\RouteTitle{Ruta 229}{r7c4\_\allowbreak{}h-1v-1}
\noindent Celda 58\quad (7, 4)\qquad Orientación \(h=-1,\ v=-1\)\hfill\hyperref[sec:vi-indice-rutas]{Índice}\par\vspace{4pt}
\begingroup\fontsize{10.2}{12.5}\selectfont\setlength{\tabcolsep}{5pt}\renewcommand{\arraystretch}{1.1}\rowcolors{1}{routepale}{white}\noindent\begin{tabularx}{\linewidth}{@{}>{\raggedright\arraybackslash}p{.345\linewidth}>{\raggedright\arraybackslash}X@{}}
\RouteGroup{Identidad estructural}
\RouteRow{Clase y multisección}{neutrino\_\allowbreak{}G2\quad /\quad neutrino\_\allowbreak{}G2}
\RouteRow{Colección y sector}{neutrino; leptón}
\RouteRow{Clasificación física}{neutrino}
\RouteRow{Generación, profundidad y color}{2; G2; \textemdash{}}
\RouteRow{Ruta primaria y órbita de inversión}{\(B_1\): no\qquad 3,\allowbreak{}6;\allowbreak{}7,\allowbreak{}4}
\RouteRow{Firma de clasificación}{30|\allowbreak{}6|\allowbreak{}0|\allowbreak{}0|\allowbreak{}-12|\allowbreak{}+++-\kern0pt{}-\kern0pt{}-+++-\kern0pt{}-\kern0pt{}-}
\RouteGroup{Retornos, memoria y transporte}
\RouteRow{Retornos con fase}{clase: 27; órbita: 27; celda: 27}
\RouteRow{Retorno cinemático}{en 108: 54; extensión: 216}
\RouteRow{Giros y cambios de hoja}{71; 107}
\RouteRow{Acarreo y diferimiento}{deuda total: 2; final: 1; bloques diferidos: 2}
\RouteRow{Tríadas finales; persistencia de Pareto}{16; sí}
\RouteRow{\(K\longrightarrow K+9\)}{repeticiones: 9; cambios de memoria: 9}
\RouteGroup{Firma, fase y diferencias}
\RouteRow{\(n_A,\ s_C,\ \kappa_0,\ \nu_{120},\ \nu_{270}\)}{30, -6, 0, 0, -12}
\RouteRow{\(\tau_{12}\); firma histórica \(\mathbb Z^5\)}{+++-\kern0pt{}-\kern0pt{}-+++-\kern0pt{}-\kern0pt{}-; 30|\allowbreak{}-6|\allowbreak{}0|\allowbreak{}0|\allowbreak{}-12}
\RouteRow{\(D_1,\ D_3,\ D_4,\ D_{27},\ D_{36}\)}{54, 56, 70, 48, 32}
\RouteRow{\(\Omega_{90/120},\ P_6\)}{80, 6}
\RouteGroup{Lectores y factores de acción}
\RouteRow{\(K_D\quad /\quad K_\Omega\)}{(80,\allowbreak{}54,\allowbreak{}7)\quad /\quad (80,\allowbreak{}54,\allowbreak{}6)}
\RouteRow{\(\kappa_D\)}{79.60672\allowbreak{}13362175\allowbreak{}19726524\allowbreak{}96}
\RouteRow{\(\kappa_\Omega\)}{79.60661\allowbreak{}01397948\allowbreak{}27586470\allowbreak{}91}
\RouteRow{\(R_{\mathrm{act},D}\)}{1.000000\allowbreak{}05518988\allowbreak{}96223153\allowbreak{}37}
\RouteRow{\(R_{\mathrm{act},\Omega}\)}{1.000000\allowbreak{}05518981\allowbreak{}25318591\allowbreak{}40}
\RouteGroup{Separación de prefijos}
\RouteRow{Área de ambigüedad; área \(\log_2\)}{18; 0.0}
\RouteRow{Primer bloque de separación}{órbita: 1; celda: 1; ruta: 1}
\RouteRow{Ambigüedad final}{1}
\RouteGroup{Secuencias completas}
\RouteRow{\(w_6\)}{486 423 15 486 429 243 657 24 243 486 423 15 486 429 243 657 24 243}
\RouteRow{Tríadas estables por bloque}{0 1 2 3 4 5 6 7 8 9 10 11 12 13 13 15 16 16}
\end{tabularx}\endgroup\par\vspace{5pt}\noindent{\fontsize{8.4}{10}\selectfont\hyperref[trace-229]{Procedencia y códigos originales}\hfill Ficha completa · 229/324}
\clearpage\phantomsection\label{nav:vi-ruta-230}
% VI-CATALOG-ROW rutas 230
\pdfbookmark[3]{r7c4\_h-1v+1}{route-230}
\RouteTitle{Ruta 230}{r7c4\_\allowbreak{}h-1v+1}
\noindent Celda 58\quad (7, 4)\qquad Orientación \(h=-1,\ v=1\)\hfill\hyperref[sec:vi-indice-rutas]{Índice}\par\vspace{4pt}
\begingroup\fontsize{10.2}{12.5}\selectfont\setlength{\tabcolsep}{5pt}\renewcommand{\arraystretch}{1.1}\rowcolors{1}{routepale}{white}\noindent\begin{tabularx}{\linewidth}{@{}>{\raggedright\arraybackslash}p{.345\linewidth}>{\raggedright\arraybackslash}X@{}}
\RouteGroup{Identidad estructural}
\RouteRow{Clase y multisección}{neutrino\_\allowbreak{}G2\quad /\quad neutrino\_\allowbreak{}G2}
\RouteRow{Colección y sector}{neutrino; leptón}
\RouteRow{Clasificación física}{neutrino}
\RouteRow{Generación, profundidad y color}{2; G2; \textemdash{}}
\RouteRow{Ruta primaria y órbita de inversión}{\(B_1\): sí\qquad 3,\allowbreak{}6;\allowbreak{}7,\allowbreak{}4}
\RouteRow{Firma de clasificación}{30|\allowbreak{}6|\allowbreak{}0|\allowbreak{}0|\allowbreak{}-12|\allowbreak{}+++-\kern0pt{}-\kern0pt{}-+++-\kern0pt{}-\kern0pt{}-}
\RouteGroup{Retornos, memoria y transporte}
\RouteRow{Retornos con fase}{clase: 27; órbita: 27; celda: 27}
\RouteRow{Retorno cinemático}{en 108: 54; extensión: 216}
\RouteRow{Giros y cambios de hoja}{71; 107}
\RouteRow{Acarreo y diferimiento}{deuda total: 0; final: 0; bloques diferidos: 0}
\RouteRow{Tríadas finales; persistencia de Pareto}{17; sí}
\RouteRow{\(K\longrightarrow K+9\)}{repeticiones: 9; cambios de memoria: 9}
\RouteGroup{Firma, fase y diferencias}
\RouteRow{\(n_A,\ s_C,\ \kappa_0,\ \nu_{120},\ \nu_{270}\)}{14, -4, 0, 0, -12}
\RouteRow{\(\tau_{12}\); firma histórica \(\mathbb Z^5\)}{+++-\kern0pt{}-\kern0pt{}-+++-\kern0pt{}-\kern0pt{}-; 14|\allowbreak{}-4|\allowbreak{}0|\allowbreak{}0|\allowbreak{}-12}
\RouteRow{\(D_1,\ D_3,\ D_4,\ D_{27},\ D_{36}\)}{78, 28, 78, 36, 24}
\RouteRow{\(\Omega_{90/120},\ P_6\)}{80, 6}
\RouteGroup{Lectores y factores de acción}
\RouteRow{\(K_D\quad /\quad K_\Omega\)}{(60,\allowbreak{}78,\allowbreak{}25)\quad /\quad (80,\allowbreak{}54,\allowbreak{}6)}
\RouteRow{\(\kappa_D\)}{59.43358\allowbreak{}64101631\allowbreak{}67023563\allowbreak{}04}
\RouteRow{\(\kappa_\Omega\)}{79.60661\allowbreak{}01397948\allowbreak{}27586470\allowbreak{}91}
\RouteRow{\(R_{\mathrm{act},D}\)}{1.000000\allowbreak{}04120422\allowbreak{}24053449\allowbreak{}05}
\RouteRow{\(R_{\mathrm{act},\Omega}\)}{1.000000\allowbreak{}05518981\allowbreak{}25318591\allowbreak{}40}
\RouteGroup{Separación de prefijos}
\RouteRow{Área de ambigüedad; área \(\log_2\)}{18; 0.0}
\RouteRow{Primer bloque de separación}{órbita: 1; celda: 1; ruta: 1}
\RouteRow{Ambigüedad final}{1}
\RouteGroup{Secuencias completas}
\RouteRow{\(w_6\)}{507 504 585 507 510 510 666 672 510 507 504 585 507 510 510 666 672 510}
\RouteRow{Tríadas estables por bloque}{0 1 2 3 4 5 6 7 8 9 10 11 12 13 14 15 16 17}
\end{tabularx}\endgroup\par\vspace{5pt}\noindent{\fontsize{8.4}{10}\selectfont\hyperref[trace-230]{Procedencia y códigos originales}\hfill Ficha completa · 230/324}
\clearpage\phantomsection\label{nav:vi-ruta-231}
% VI-CATALOG-ROW rutas 231
\pdfbookmark[3]{r7c4\_h+1v-1}{route-231}
\RouteTitle{Ruta 231}{r7c4\_\allowbreak{}h+1v-1}
\noindent Celda 58\quad (7, 4)\qquad Orientación \(h=1,\ v=-1\)\hfill\hyperref[sec:vi-indice-rutas]{Índice}\par\vspace{4pt}
\begingroup\fontsize{10.2}{12.5}\selectfont\setlength{\tabcolsep}{5pt}\renewcommand{\arraystretch}{1.1}\rowcolors{1}{routepale}{white}\noindent\begin{tabularx}{\linewidth}{@{}>{\raggedright\arraybackslash}p{.345\linewidth}>{\raggedright\arraybackslash}X@{}}
\RouteGroup{Identidad estructural}
\RouteRow{Clase y multisección}{neutrino\_\allowbreak{}G2\quad /\quad neutrino\_\allowbreak{}G2}
\RouteRow{Colección y sector}{neutrino; leptón}
\RouteRow{Clasificación física}{neutrino}
\RouteRow{Generación, profundidad y color}{2; G2; \textemdash{}}
\RouteRow{Ruta primaria y órbita de inversión}{\(B_1\): no\qquad 3,\allowbreak{}6;\allowbreak{}7,\allowbreak{}4}
\RouteRow{Firma de clasificación}{30|\allowbreak{}6|\allowbreak{}0|\allowbreak{}0|\allowbreak{}-12|\allowbreak{}+++-\kern0pt{}-\kern0pt{}-+++-\kern0pt{}-\kern0pt{}-}
\RouteGroup{Retornos, memoria y transporte}
\RouteRow{Retornos con fase}{clase: 27; órbita: 27; celda: 27}
\RouteRow{Retorno cinemático}{en 108: 54; extensión: 216}
\RouteRow{Giros y cambios de hoja}{71; 107}
\RouteRow{Acarreo y diferimiento}{deuda total: 2; final: 0; bloques diferidos: 1}
\RouteRow{Tríadas finales; persistencia de Pareto}{17; sí}
\RouteRow{\(K\longrightarrow K+9\)}{repeticiones: 9; cambios de memoria: 9}
\RouteGroup{Firma, fase y diferencias}
\RouteRow{\(n_A,\ s_C,\ \kappa_0,\ \nu_{120},\ \nu_{270}\)}{14, 4, 0, 0, -12}
\RouteRow{\(\tau_{12}\); firma histórica \(\mathbb Z^5\)}{+++-\kern0pt{}-\kern0pt{}-+++-\kern0pt{}-\kern0pt{}-; 14|\allowbreak{}4|\allowbreak{}0|\allowbreak{}0|\allowbreak{}-12}
\RouteRow{\(D_1,\ D_3,\ D_4,\ D_{27},\ D_{36}\)}{78, 28, 78, 36, 24}
\RouteRow{\(\Omega_{90/120},\ P_6\)}{80, 6}
\RouteGroup{Lectores y factores de acción}
\RouteRow{\(K_D\quad /\quad K_\Omega\)}{(60,\allowbreak{}78,\allowbreak{}25)\quad /\quad (80,\allowbreak{}54,\allowbreak{}6)}
\RouteRow{\(\kappa_D\)}{59.43358\allowbreak{}64101631\allowbreak{}67023563\allowbreak{}04}
\RouteRow{\(\kappa_\Omega\)}{79.60661\allowbreak{}01397948\allowbreak{}27586470\allowbreak{}91}
\RouteRow{\(R_{\mathrm{act},D}\)}{1.000000\allowbreak{}04120422\allowbreak{}24053449\allowbreak{}05}
\RouteRow{\(R_{\mathrm{act},\Omega}\)}{1.000000\allowbreak{}05518981\allowbreak{}25318591\allowbreak{}40}
\RouteGroup{Separación de prefijos}
\RouteRow{Área de ambigüedad; área \(\log_2\)}{18; 0.0}
\RouteRow{Primer bloque de separación}{órbita: 1; celda: 1; ruta: 1}
\RouteRow{Ambigüedad final}{1}
\RouteGroup{Secuencias completas}
\RouteRow{\(w_6\)}{666 672 510 666 666 585 507 504 585 666 672 510 666 666 585 507 504 585}
\RouteRow{Tríadas estables por bloque}{0 1 2 3 4 5 6 7 8 9 10 11 12 12 13 15 16 17}
\end{tabularx}\endgroup\par\vspace{5pt}\noindent{\fontsize{8.4}{10}\selectfont\hyperref[trace-231]{Procedencia y códigos originales}\hfill Ficha completa · 231/324}
\clearpage\phantomsection\label{nav:vi-ruta-232}
% VI-CATALOG-ROW rutas 232
\pdfbookmark[3]{r7c4\_h+1v+1}{route-232}
\RouteTitle{Ruta 232}{r7c4\_\allowbreak{}h+1v+1}
\noindent Celda 58\quad (7, 4)\qquad Orientación \(h=1,\ v=1\)\hfill\hyperref[sec:vi-indice-rutas]{Índice}\par\vspace{4pt}
\begingroup\fontsize{10.2}{12.5}\selectfont\setlength{\tabcolsep}{5pt}\renewcommand{\arraystretch}{1.1}\rowcolors{1}{routepale}{white}\noindent\begin{tabularx}{\linewidth}{@{}>{\raggedright\arraybackslash}p{.345\linewidth}>{\raggedright\arraybackslash}X@{}}
\RouteGroup{Identidad estructural}
\RouteRow{Clase y multisección}{neutrino\_\allowbreak{}G2\quad /\quad neutrino\_\allowbreak{}G2}
\RouteRow{Colección y sector}{neutrino; leptón}
\RouteRow{Clasificación física}{neutrino}
\RouteRow{Generación, profundidad y color}{2; G2; \textemdash{}}
\RouteRow{Ruta primaria y órbita de inversión}{\(B_1\): no\qquad 3,\allowbreak{}6;\allowbreak{}7,\allowbreak{}4}
\RouteRow{Firma de clasificación}{30|\allowbreak{}6|\allowbreak{}0|\allowbreak{}0|\allowbreak{}-12|\allowbreak{}+++-\kern0pt{}-\kern0pt{}-+++-\kern0pt{}-\kern0pt{}-}
\RouteGroup{Retornos, memoria y transporte}
\RouteRow{Retornos con fase}{clase: 27; órbita: 27; celda: 27}
\RouteRow{Retorno cinemático}{en 108: 54; extensión: 216}
\RouteRow{Giros y cambios de hoja}{71; 107}
\RouteRow{Acarreo y diferimiento}{deuda total: 2; final: 0; bloques diferidos: 1}
\RouteRow{Tríadas finales; persistencia de Pareto}{17; sí}
\RouteRow{\(K\longrightarrow K+9\)}{repeticiones: 9; cambios de memoria: 9}
\RouteGroup{Firma, fase y diferencias}
\RouteRow{\(n_A,\ s_C,\ \kappa_0,\ \nu_{120},\ \nu_{270}\)}{30, 6, 0, 0, -12}
\RouteRow{\(\tau_{12}\); firma histórica \(\mathbb Z^5\)}{+++-\kern0pt{}-\kern0pt{}-+++-\kern0pt{}-\kern0pt{}-; 30|\allowbreak{}6|\allowbreak{}0|\allowbreak{}0|\allowbreak{}-12}
\RouteRow{\(D_1,\ D_3,\ D_4,\ D_{27},\ D_{36}\)}{54, 56, 70, 48, 32}
\RouteRow{\(\Omega_{90/120},\ P_6\)}{80, 6}
\RouteGroup{Lectores y factores de acción}
\RouteRow{\(K_D\quad /\quad K_\Omega\)}{(80,\allowbreak{}54,\allowbreak{}7)\quad /\quad (80,\allowbreak{}54,\allowbreak{}6)}
\RouteRow{\(\kappa_D\)}{79.60672\allowbreak{}13362175\allowbreak{}19726524\allowbreak{}96}
\RouteRow{\(\kappa_\Omega\)}{79.60661\allowbreak{}01397948\allowbreak{}27586470\allowbreak{}91}
\RouteRow{\(R_{\mathrm{act},D}\)}{1.000000\allowbreak{}05518988\allowbreak{}96223153\allowbreak{}37}
\RouteRow{\(R_{\mathrm{act},\Omega}\)}{1.000000\allowbreak{}05518981\allowbreak{}25318591\allowbreak{}40}
\RouteGroup{Separación de prefijos}
\RouteRow{Área de ambigüedad; área \(\log_2\)}{18; 0.0}
\RouteRow{Primer bloque de separación}{órbita: 1; celda: 1; ruta: 1}
\RouteRow{Ambigüedad final}{1}
\RouteGroup{Secuencias completas}
\RouteRow{\(w_6\)}{657 24 243 657 18 15 486 423 15 657 24 243 657 18 15 486 423 15}
\RouteRow{Tríadas estables por bloque}{0 1 2 3 4 5 6 7 8 9 10 11 12 13 14 14 15 17}
\end{tabularx}\endgroup\par\vspace{5pt}\noindent{\fontsize{8.4}{10}\selectfont\hyperref[trace-232]{Procedencia y códigos originales}\hfill Ficha completa · 232/324}
\clearpage\phantomsection\label{nav:vi-ruta-233}
% VI-CATALOG-ROW rutas 233
\pdfbookmark[2]{Celda 59}{cell-59}
\pdfbookmark[3]{r7c5\_h-1v-1}{route-233}
\RouteTitle{Ruta 233}{r7c5\_\allowbreak{}h-1v-1}
\noindent Celda 59\quad (7, 5)\qquad Orientación \(h=-1,\ v=-1\)\hfill\hyperref[sec:vi-indice-rutas]{Índice}\par\vspace{4pt}
\begingroup\fontsize{10.2}{12.5}\selectfont\setlength{\tabcolsep}{5pt}\renewcommand{\arraystretch}{1.1}\rowcolors{1}{routepale}{white}\noindent\begin{tabularx}{\linewidth}{@{}>{\raggedright\arraybackslash}p{.345\linewidth}>{\raggedright\arraybackslash}X@{}}
\RouteGroup{Identidad estructural}
\RouteRow{Clase y multisección}{charged\_\allowbreak{}lepton\_\allowbreak{}G2\quad /\quad charged\_\allowbreak{}lepton\_\allowbreak{}G2}
\RouteRow{Colección y sector}{leptón cargado; leptón}
\RouteRow{Clasificación física}{leptón cargado}
\RouteRow{Generación, profundidad y color}{2; G2; \textemdash{}}
\RouteRow{Ruta primaria y órbita de inversión}{\(B_1\): sí\qquad 3,\allowbreak{}5;\allowbreak{}7,\allowbreak{}5}
\RouteRow{Firma de clasificación}{28|\allowbreak{}6|\allowbreak{}-2|\allowbreak{}0|\allowbreak{}-12|\allowbreak{}+++-\kern0pt{}-\kern0pt{}-+++-\kern0pt{}-\kern0pt{}-}
\RouteGroup{Retornos, memoria y transporte}
\RouteRow{Retornos con fase}{clase: 27; órbita: 27; celda: 27}
\RouteRow{Retorno cinemático}{en 108: 54; extensión: 216}
\RouteRow{Giros y cambios de hoja}{71; 107}
\RouteRow{Acarreo y diferimiento}{deuda total: 1; final: 0; bloques diferidos: 1}
\RouteRow{Tríadas finales; persistencia de Pareto}{17; sí}
\RouteRow{\(K\longrightarrow K+9\)}{repeticiones: 9; cambios de memoria: 9}
\RouteGroup{Firma, fase y diferencias}
\RouteRow{\(n_A,\ s_C,\ \kappa_0,\ \nu_{120},\ \nu_{270}\)}{28, -6, 0, -4, -12}
\RouteRow{\(\tau_{12}\); firma histórica \(\mathbb Z^5\)}{+++-\kern0pt{}-\kern0pt{}-+++-\kern0pt{}-\kern0pt{}-; 28|\allowbreak{}-6|\allowbreak{}0|\allowbreak{}-4|\allowbreak{}-12}
\RouteRow{\(D_1,\ D_3,\ D_4,\ D_{27},\ D_{36}\)}{66, 64, 64, 60, 40}
\RouteRow{\(\Omega_{90/120},\ P_6\)}{80, 6}
\RouteGroup{Lectores y factores de acción}
\RouteRow{\(K_D\quad /\quad K_\Omega\)}{(100,\allowbreak{}66,\allowbreak{}0)\quad /\quad (80,\allowbreak{}54,\allowbreak{}6)}
\RouteRow{\(\kappa_D\)}{99.51837\allowbreak{}47304272\allowbreak{}69134179\allowbreak{}18}
\RouteRow{\(\kappa_\Omega\)}{79.60661\allowbreak{}01397948\allowbreak{}27586470\allowbreak{}91}
\RouteRow{\(R_{\mathrm{act},D}\)}{1.000000\allowbreak{}06899427\allowbreak{}66450218\allowbreak{}95}
\RouteRow{\(R_{\mathrm{act},\Omega}\)}{1.000000\allowbreak{}05518981\allowbreak{}25318591\allowbreak{}40}
\RouteGroup{Separación de prefijos}
\RouteRow{Área de ambigüedad; área \(\log_2\)}{18; 0.0}
\RouteRow{Primer bloque de separación}{órbita: 1; celda: 1; ruta: 1}
\RouteRow{Ambigüedad final}{1}
\RouteGroup{Secuencias completas}
\RouteRow{\(w_6\)}{90 570 264 90 567 657 24 243 657 90 570 264 90 567 657 24 243 657}
\RouteRow{Tríadas estables por bloque}{0 1 2 3 4 5 6 7 8 9 10 11 12 13 13 15 16 17}
\end{tabularx}\endgroup\par\vspace{5pt}\noindent{\fontsize{8.4}{10}\selectfont\hyperref[trace-233]{Procedencia y códigos originales}\hfill Ficha completa · 233/324}
\clearpage\phantomsection\label{nav:vi-ruta-234}
% VI-CATALOG-ROW rutas 234
\pdfbookmark[3]{r7c5\_h-1v+1}{route-234}
\RouteTitle{Ruta 234}{r7c5\_\allowbreak{}h-1v+1}
\noindent Celda 59\quad (7, 5)\qquad Orientación \(h=-1,\ v=1\)\hfill\hyperref[sec:vi-indice-rutas]{Índice}\par\vspace{4pt}
\begingroup\fontsize{10.2}{12.5}\selectfont\setlength{\tabcolsep}{5pt}\renewcommand{\arraystretch}{1.1}\rowcolors{1}{routepale}{white}\noindent\begin{tabularx}{\linewidth}{@{}>{\raggedright\arraybackslash}p{.345\linewidth}>{\raggedright\arraybackslash}X@{}}
\RouteGroup{Identidad estructural}
\RouteRow{Clase y multisección}{charged\_\allowbreak{}lepton\_\allowbreak{}G2\quad /\quad charged\_\allowbreak{}lepton\_\allowbreak{}G2}
\RouteRow{Colección y sector}{leptón cargado; leptón}
\RouteRow{Clasificación física}{leptón cargado}
\RouteRow{Generación, profundidad y color}{2; G2; \textemdash{}}
\RouteRow{Ruta primaria y órbita de inversión}{\(B_1\): no\qquad 3,\allowbreak{}5;\allowbreak{}7,\allowbreak{}5}
\RouteRow{Firma de clasificación}{28|\allowbreak{}6|\allowbreak{}-2|\allowbreak{}0|\allowbreak{}-12|\allowbreak{}+++-\kern0pt{}-\kern0pt{}-+++-\kern0pt{}-\kern0pt{}-}
\RouteGroup{Retornos, memoria y transporte}
\RouteRow{Retornos con fase}{clase: 27; órbita: 27; celda: 27}
\RouteRow{Retorno cinemático}{en 108: 54; extensión: 216}
\RouteRow{Giros y cambios de hoja}{71; 107}
\RouteRow{Acarreo y diferimiento}{deuda total: 2; final: 0; bloques diferidos: 2}
\RouteRow{Tríadas finales; persistencia de Pareto}{17; sí}
\RouteRow{\(K\longrightarrow K+9\)}{repeticiones: 9; cambios de memoria: 9}
\RouteGroup{Firma, fase y diferencias}
\RouteRow{\(n_A,\ s_C,\ \kappa_0,\ \nu_{120},\ \nu_{270}\)}{50, -6, 2, -2, -8}
\RouteRow{\(\tau_{12}\); firma histórica \(\mathbb Z^5\)}{++0-\kern0pt{}-\kern0pt{}-++0-\kern0pt{}-\kern0pt{}-; 50|\allowbreak{}-6|\allowbreak{}2|\allowbreak{}-2|\allowbreak{}-8}
\RouteRow{\(D_1,\ D_3,\ D_4,\ D_{27},\ D_{36}\)}{24, 20, 24, 24, 16}
\RouteRow{\(\Omega_{90/120},\ P_6\)}{56, 5}
\RouteGroup{Lectores y factores de acción}
\RouteRow{\(K_D\quad /\quad K_\Omega\)}{(40,\allowbreak{}24,\allowbreak{}2)\quad /\quad (56,\allowbreak{}54,\allowbreak{}5)}
\RouteRow{\(\kappa_D\)}{39.82508\allowbreak{}59311825\allowbreak{}73056173\allowbreak{}26}
\RouteRow{\(\kappa_\Omega\)}{55.60649\allowbreak{}89433721\allowbreak{}35446416\allowbreak{}86}
\RouteRow{\(R_{\mathrm{act},D}\)}{1.000000\allowbreak{}02761000\allowbreak{}61595142\allowbreak{}15}
\RouteRow{\(R_{\mathrm{act},\Omega}\)}{1.000000\allowbreak{}03855097\allowbreak{}23541416\allowbreak{}45}
\RouteGroup{Separación de prefijos}
\RouteRow{Área de ambigüedad; área \(\log_2\)}{36; 18.0}
\RouteRow{Primer bloque de separación}{órbita: \textemdash{}; celda: \textemdash{}; ruta: \textemdash{}}
\RouteRow{Ambigüedad final}{2}
\RouteGroup{Secuencias completas}
\RouteRow{\(w_6\)}{81 3 0 81 0 3 0 81 3 81 3 0 81 0 3 0 81 3}
\RouteRow{Tríadas estables por bloque}{0 1 2 3 4 5 6 6 8 9 10 11 12 13 14 15 15 17}
\end{tabularx}\endgroup\par\vspace{5pt}\noindent{\fontsize{8.4}{10}\selectfont\hyperref[trace-234]{Procedencia y códigos originales}\hfill Ficha completa · 234/324}
\clearpage\phantomsection\label{nav:vi-ruta-235}
% VI-CATALOG-ROW rutas 235
\pdfbookmark[3]{r7c5\_h+1v-1}{route-235}
\RouteTitle{Ruta 235}{r7c5\_\allowbreak{}h+1v-1}
\noindent Celda 59\quad (7, 5)\qquad Orientación \(h=1,\ v=-1\)\hfill\hyperref[sec:vi-indice-rutas]{Índice}\par\vspace{4pt}
\begingroup\fontsize{10.2}{12.5}\selectfont\setlength{\tabcolsep}{5pt}\renewcommand{\arraystretch}{1.1}\rowcolors{1}{routepale}{white}\noindent\begin{tabularx}{\linewidth}{@{}>{\raggedright\arraybackslash}p{.345\linewidth}>{\raggedright\arraybackslash}X@{}}
\RouteGroup{Identidad estructural}
\RouteRow{Clase y multisección}{charged\_\allowbreak{}lepton\_\allowbreak{}G2\quad /\quad charged\_\allowbreak{}lepton\_\allowbreak{}G2}
\RouteRow{Colección y sector}{leptón cargado; leptón}
\RouteRow{Clasificación física}{leptón cargado}
\RouteRow{Generación, profundidad y color}{2; G2; \textemdash{}}
\RouteRow{Ruta primaria y órbita de inversión}{\(B_1\): no\qquad 3,\allowbreak{}5;\allowbreak{}7,\allowbreak{}5}
\RouteRow{Firma de clasificación}{28|\allowbreak{}6|\allowbreak{}-2|\allowbreak{}0|\allowbreak{}-12|\allowbreak{}+++-\kern0pt{}-\kern0pt{}-+++-\kern0pt{}-\kern0pt{}-}
\RouteGroup{Retornos, memoria y transporte}
\RouteRow{Retornos con fase}{clase: 27; órbita: 27; celda: 27}
\RouteRow{Retorno cinemático}{en 108: 54; extensión: 216}
\RouteRow{Giros y cambios de hoja}{71; 107}
\RouteRow{Acarreo y diferimiento}{deuda total: 0; final: 0; bloques diferidos: 0}
\RouteRow{Tríadas finales; persistencia de Pareto}{17; sí}
\RouteRow{\(K\longrightarrow K+9\)}{repeticiones: 9; cambios de memoria: 9}
\RouteGroup{Firma, fase y diferencias}
\RouteRow{\(n_A,\ s_C,\ \kappa_0,\ \nu_{120},\ \nu_{270}\)}{50, 6, -2, 2, -8}
\RouteRow{\(\tau_{12}\); firma histórica \(\mathbb Z^5\)}{+++-\kern0pt{}-0+++-\kern0pt{}-0; 50|\allowbreak{}6|\allowbreak{}-2|\allowbreak{}2|\allowbreak{}-8}
\RouteRow{\(D_1,\ D_3,\ D_4,\ D_{27},\ D_{36}\)}{24, 20, 24, 24, 16}
\RouteRow{\(\Omega_{90/120},\ P_6\)}{56, 5}
\RouteGroup{Lectores y factores de acción}
\RouteRow{\(K_D\quad /\quad K_\Omega\)}{(40,\allowbreak{}24,\allowbreak{}2)\quad /\quad (56,\allowbreak{}54,\allowbreak{}5)}
\RouteRow{\(\kappa_D\)}{39.82508\allowbreak{}59311825\allowbreak{}73056173\allowbreak{}26}
\RouteRow{\(\kappa_\Omega\)}{55.60649\allowbreak{}89433721\allowbreak{}35446416\allowbreak{}86}
\RouteRow{\(R_{\mathrm{act},D}\)}{1.000000\allowbreak{}02761000\allowbreak{}61595142\allowbreak{}15}
\RouteRow{\(R_{\mathrm{act},\Omega}\)}{1.000000\allowbreak{}03855097\allowbreak{}23541416\allowbreak{}45}
\RouteGroup{Separación de prefijos}
\RouteRow{Área de ambigüedad; área \(\log_2\)}{36; 18.0}
\RouteRow{Primer bloque de separación}{órbita: \textemdash{}; celda: \textemdash{}; ruta: \textemdash{}}
\RouteRow{Ambigüedad final}{2}
\RouteGroup{Secuencias completas}
\RouteRow{\(w_6\)}{0 81 3 0 84 0 81 3 0 0 81 3 0 84 0 81 3 0}
\RouteRow{Tríadas estables por bloque}{0 1 2 3 4 5 6 7 8 9 10 11 12 13 14 15 16 17}
\end{tabularx}\endgroup\par\vspace{5pt}\noindent{\fontsize{8.4}{10}\selectfont\hyperref[trace-235]{Procedencia y códigos originales}\hfill Ficha completa · 235/324}
\clearpage\phantomsection\label{nav:vi-ruta-236}
% VI-CATALOG-ROW rutas 236
\pdfbookmark[3]{r7c5\_h+1v+1}{route-236}
\RouteTitle{Ruta 236}{r7c5\_\allowbreak{}h+1v+1}
\noindent Celda 59\quad (7, 5)\qquad Orientación \(h=1,\ v=1\)\hfill\hyperref[sec:vi-indice-rutas]{Índice}\par\vspace{4pt}
\begingroup\fontsize{10.2}{12.5}\selectfont\setlength{\tabcolsep}{5pt}\renewcommand{\arraystretch}{1.1}\rowcolors{1}{routepale}{white}\noindent\begin{tabularx}{\linewidth}{@{}>{\raggedright\arraybackslash}p{.345\linewidth}>{\raggedright\arraybackslash}X@{}}
\RouteGroup{Identidad estructural}
\RouteRow{Clase y multisección}{charged\_\allowbreak{}lepton\_\allowbreak{}G2\quad /\quad charged\_\allowbreak{}lepton\_\allowbreak{}G2}
\RouteRow{Colección y sector}{leptón cargado; leptón}
\RouteRow{Clasificación física}{leptón cargado}
\RouteRow{Generación, profundidad y color}{2; G2; \textemdash{}}
\RouteRow{Ruta primaria y órbita de inversión}{\(B_1\): no\qquad 3,\allowbreak{}5;\allowbreak{}7,\allowbreak{}5}
\RouteRow{Firma de clasificación}{28|\allowbreak{}6|\allowbreak{}-2|\allowbreak{}0|\allowbreak{}-12|\allowbreak{}+++-\kern0pt{}-\kern0pt{}-+++-\kern0pt{}-\kern0pt{}-}
\RouteGroup{Retornos, memoria y transporte}
\RouteRow{Retornos con fase}{clase: 27; órbita: 27; celda: 27}
\RouteRow{Retorno cinemático}{en 108: 54; extensión: 216}
\RouteRow{Giros y cambios de hoja}{71; 107}
\RouteRow{Acarreo y diferimiento}{deuda total: 0; final: 0; bloques diferidos: 0}
\RouteRow{Tríadas finales; persistencia de Pareto}{17; sí}
\RouteRow{\(K\longrightarrow K+9\)}{repeticiones: 9; cambios de memoria: 9}
\RouteGroup{Firma, fase y diferencias}
\RouteRow{\(n_A,\ s_C,\ \kappa_0,\ \nu_{120},\ \nu_{270}\)}{28, 6, 0, 4, -12}
\RouteRow{\(\tau_{12}\); firma histórica \(\mathbb Z^5\)}{+++-\kern0pt{}-\kern0pt{}-+++-\kern0pt{}-\kern0pt{}-; 28|\allowbreak{}6|\allowbreak{}0|\allowbreak{}4|\allowbreak{}-12}
\RouteRow{\(D_1,\ D_3,\ D_4,\ D_{27},\ D_{36}\)}{66, 64, 64, 60, 40}
\RouteRow{\(\Omega_{90/120},\ P_6\)}{80, 6}
\RouteGroup{Lectores y factores de acción}
\RouteRow{\(K_D\quad /\quad K_\Omega\)}{(100,\allowbreak{}66,\allowbreak{}0)\quad /\quad (80,\allowbreak{}54,\allowbreak{}6)}
\RouteRow{\(\kappa_D\)}{99.51837\allowbreak{}47304272\allowbreak{}69134179\allowbreak{}18}
\RouteRow{\(\kappa_\Omega\)}{79.60661\allowbreak{}01397948\allowbreak{}27586470\allowbreak{}91}
\RouteRow{\(R_{\mathrm{act},D}\)}{1.000000\allowbreak{}06899427\allowbreak{}66450218\allowbreak{}95}
\RouteRow{\(R_{\mathrm{act},\Omega}\)}{1.000000\allowbreak{}05518981\allowbreak{}25318591\allowbreak{}40}
\RouteGroup{Separación de prefijos}
\RouteRow{Área de ambigüedad; área \(\log_2\)}{22; 4.0}
\RouteRow{Primer bloque de separación}{órbita: 5; celda: 5; ruta: 5}
\RouteRow{Ambigüedad final}{1}
\RouteGroup{Secuencias completas}
\RouteRow{\(w_6\)}{24 243 657 24 246 264 90 570 264 24 243 657 24 246 264 90 570 264}
\RouteRow{Tríadas estables por bloque}{0 1 2 3 4 5 6 7 8 9 10 11 12 13 14 15 16 17}
\end{tabularx}\endgroup\par\vspace{5pt}\noindent{\fontsize{8.4}{10}\selectfont\hyperref[trace-236]{Procedencia y códigos originales}\hfill Ficha completa · 236/324}
\clearpage\phantomsection\label{nav:vi-ruta-237}
% VI-CATALOG-ROW rutas 237
\pdfbookmark[2]{Celda 60}{cell-60}
\pdfbookmark[3]{r7c6\_h-1v-1}{route-237}
\RouteTitle{Ruta 237}{r7c6\_\allowbreak{}h-1v-1}
\noindent Celda 60\quad (7, 6)\qquad Orientación \(h=-1,\ v=-1\)\hfill\hyperref[sec:vi-indice-rutas]{Índice}\par\vspace{4pt}
\begingroup\fontsize{10.2}{12.5}\selectfont\setlength{\tabcolsep}{5pt}\renewcommand{\arraystretch}{1.1}\rowcolors{1}{routepale}{white}\noindent\begin{tabularx}{\linewidth}{@{}>{\raggedright\arraybackslash}p{.345\linewidth}>{\raggedright\arraybackslash}X@{}}
\RouteGroup{Identidad estructural}
\RouteRow{Clase y multisección}{neutrino\_\allowbreak{}G2\quad /\quad neutrino\_\allowbreak{}G2}
\RouteRow{Colección y sector}{neutrino; leptón}
\RouteRow{Clasificación física}{neutrino}
\RouteRow{Generación, profundidad y color}{2; G2; \textemdash{}}
\RouteRow{Ruta primaria y órbita de inversión}{\(B_1\): no\qquad 3,\allowbreak{}4;\allowbreak{}7,\allowbreak{}6}
\RouteRow{Firma de clasificación}{28|\allowbreak{}-6|\allowbreak{}8|\allowbreak{}0|\allowbreak{}-12|\allowbreak{}+++-\kern0pt{}-\kern0pt{}-+++-\kern0pt{}-\kern0pt{}-}
\RouteGroup{Retornos, memoria y transporte}
\RouteRow{Retornos con fase}{clase: 27; órbita: 27; celda: 27}
\RouteRow{Retorno cinemático}{en 108: 54; extensión: 216}
\RouteRow{Giros y cambios de hoja}{71; 107}
\RouteRow{Acarreo y diferimiento}{deuda total: 1; final: 1; bloques diferidos: 1}
\RouteRow{Tríadas finales; persistencia de Pareto}{16; no}
\RouteRow{\(K\longrightarrow K+9\)}{repeticiones: 9; cambios de memoria: 9}
\RouteGroup{Firma, fase y diferencias}
\RouteRow{\(n_A,\ s_C,\ \kappa_0,\ \nu_{120},\ \nu_{270}\)}{28, 6, 0, 0, -12}
\RouteRow{\(\tau_{12}\); firma histórica \(\mathbb Z^5\)}{+++-\kern0pt{}-\kern0pt{}-+++-\kern0pt{}-\kern0pt{}-; 28|\allowbreak{}6|\allowbreak{}0|\allowbreak{}0|\allowbreak{}-12}
\RouteRow{\(D_1,\ D_3,\ D_4,\ D_{27},\ D_{36}\)}{66, 60, 64, 48, 32}
\RouteRow{\(\Omega_{90/120},\ P_6\)}{80, 6}
\RouteGroup{Lectores y factores de acción}
\RouteRow{\(K_D\quad /\quad K_\Omega\)}{(80,\allowbreak{}66,\allowbreak{}2)\quad /\quad (80,\allowbreak{}54,\allowbreak{}6)}
\RouteRow{\(\kappa_D\)}{79.51859\allowbreak{}71232726\allowbreak{}53414287\allowbreak{}29}
\RouteRow{\(\kappa_\Omega\)}{79.60661\allowbreak{}01397948\allowbreak{}27586470\allowbreak{}91}
\RouteRow{\(R_{\mathrm{act},D}\)}{1.000000\allowbreak{}05512879\allowbreak{}47092488\allowbreak{}80}
\RouteRow{\(R_{\mathrm{act},\Omega}\)}{1.000000\allowbreak{}05518981\allowbreak{}25318591\allowbreak{}40}
\RouteGroup{Separación de prefijos}
\RouteRow{Área de ambigüedad; área \(\log_2\)}{18; 0.0}
\RouteRow{Primer bloque de separación}{órbita: 1; celda: 1; ruta: 1}
\RouteRow{Ambigüedad final}{1}
\RouteGroup{Secuencias completas}
\RouteRow{\(w_6\)}{423 15 486 423 12 99 327 498 99 423 15 486 423 12 99 327 498 99}
\RouteRow{Tríadas estables por bloque}{0 1 2 3 4 5 6 7 8 9 10 11 12 13 14 15 16 16}
\end{tabularx}\endgroup\par\vspace{5pt}\noindent{\fontsize{8.4}{10}\selectfont\hyperref[trace-237]{Procedencia y códigos originales}\hfill Ficha completa · 237/324}
\clearpage\phantomsection\label{nav:vi-ruta-238}
% VI-CATALOG-ROW rutas 238
\pdfbookmark[3]{r7c6\_h-1v+1}{route-238}
\RouteTitle{Ruta 238}{r7c6\_\allowbreak{}h-1v+1}
\noindent Celda 60\quad (7, 6)\qquad Orientación \(h=-1,\ v=1\)\hfill\hyperref[sec:vi-indice-rutas]{Índice}\par\vspace{4pt}
\begingroup\fontsize{10.2}{12.5}\selectfont\setlength{\tabcolsep}{5pt}\renewcommand{\arraystretch}{1.1}\rowcolors{1}{routepale}{white}\noindent\begin{tabularx}{\linewidth}{@{}>{\raggedright\arraybackslash}p{.345\linewidth}>{\raggedright\arraybackslash}X@{}}
\RouteGroup{Identidad estructural}
\RouteRow{Clase y multisección}{neutrino\_\allowbreak{}G2\quad /\quad neutrino\_\allowbreak{}G2}
\RouteRow{Colección y sector}{neutrino; leptón}
\RouteRow{Clasificación física}{neutrino}
\RouteRow{Generación, profundidad y color}{2; G2; \textemdash{}}
\RouteRow{Ruta primaria y órbita de inversión}{\(B_1\): sí\qquad 3,\allowbreak{}4;\allowbreak{}7,\allowbreak{}6}
\RouteRow{Firma de clasificación}{28|\allowbreak{}-6|\allowbreak{}8|\allowbreak{}0|\allowbreak{}-12|\allowbreak{}+++-\kern0pt{}-\kern0pt{}-+++-\kern0pt{}-\kern0pt{}-}
\RouteGroup{Retornos, memoria y transporte}
\RouteRow{Retornos con fase}{clase: 27; órbita: 27; celda: 27}
\RouteRow{Retorno cinemático}{en 108: 54; extensión: 216}
\RouteRow{Giros y cambios de hoja}{71; 107}
\RouteRow{Acarreo y diferimiento}{deuda total: 1; final: 0; bloques diferidos: 1}
\RouteRow{Tríadas finales; persistencia de Pareto}{17; no}
\RouteRow{\(K\longrightarrow K+9\)}{repeticiones: 9; cambios de memoria: 9}
\RouteGroup{Firma, fase y diferencias}
\RouteRow{\(n_A,\ s_C,\ \kappa_0,\ \nu_{120},\ \nu_{270}\)}{8, -2, 0, 0, -12}
\RouteRow{\(\tau_{12}\); firma histórica \(\mathbb Z^5\)}{+++-\kern0pt{}-\kern0pt{}-+++-\kern0pt{}-\kern0pt{}-; 8|\allowbreak{}-2|\allowbreak{}0|\allowbreak{}0|\allowbreak{}-12}
\RouteRow{\(D_1,\ D_3,\ D_4,\ D_{27},\ D_{36}\)}{84, 24, 84, 24, 16}
\RouteRow{\(\Omega_{90/120},\ P_6\)}{80, 6}
\RouteGroup{Lectores y factores de acción}
\RouteRow{\(K_D\quad /\quad K_\Omega\)}{(40,\allowbreak{}84,\allowbreak{}30)\quad /\quad (80,\allowbreak{}54,\allowbreak{}6)}
\RouteRow{\(\kappa_D\)}{39.39035\allowbreak{}82768609\allowbreak{}24917849\allowbreak{}59}
\RouteRow{\(\kappa_\Omega\)}{79.60661\allowbreak{}01397948\allowbreak{}27586470\allowbreak{}91}
\RouteRow{\(R_{\mathrm{act},D}\)}{1.000000\allowbreak{}02730861\allowbreak{}73967087\allowbreak{}05}
\RouteRow{\(R_{\mathrm{act},\Omega}\)}{1.000000\allowbreak{}05518981\allowbreak{}25318591\allowbreak{}40}
\RouteGroup{Separación de prefijos}
\RouteRow{Área de ambigüedad; área \(\log_2\)}{18; 0.0}
\RouteRow{Primer bloque de separación}{órbita: 1; celda: 1; ruta: 1}
\RouteRow{Ambigüedad final}{1}
\RouteGroup{Secuencias completas}
\RouteRow{\(w_6\)}{420 258 414 420 255 252 333 255 252 420 258 414 420 255 252 333 255 252}
\RouteRow{Tríadas estables por bloque}{0 1 2 3 4 5 6 7 8 9 10 11 12 13 14 15 15 17}
\end{tabularx}\endgroup\par\vspace{5pt}\noindent{\fontsize{8.4}{10}\selectfont\hyperref[trace-238]{Procedencia y códigos originales}\hfill Ficha completa · 238/324}
\clearpage\phantomsection\label{nav:vi-ruta-239}
% VI-CATALOG-ROW rutas 239
\pdfbookmark[3]{r7c6\_h+1v-1}{route-239}
\RouteTitle{Ruta 239}{r7c6\_\allowbreak{}h+1v-1}
\noindent Celda 60\quad (7, 6)\qquad Orientación \(h=1,\ v=-1\)\hfill\hyperref[sec:vi-indice-rutas]{Índice}\par\vspace{4pt}
\begingroup\fontsize{10.2}{12.5}\selectfont\setlength{\tabcolsep}{5pt}\renewcommand{\arraystretch}{1.1}\rowcolors{1}{routepale}{white}\noindent\begin{tabularx}{\linewidth}{@{}>{\raggedright\arraybackslash}p{.345\linewidth}>{\raggedright\arraybackslash}X@{}}
\RouteGroup{Identidad estructural}
\RouteRow{Clase y multisección}{neutrino\_\allowbreak{}G2\quad /\quad neutrino\_\allowbreak{}G2}
\RouteRow{Colección y sector}{neutrino; leptón}
\RouteRow{Clasificación física}{neutrino}
\RouteRow{Generación, profundidad y color}{2; G2; \textemdash{}}
\RouteRow{Ruta primaria y órbita de inversión}{\(B_1\): no\qquad 3,\allowbreak{}4;\allowbreak{}7,\allowbreak{}6}
\RouteRow{Firma de clasificación}{28|\allowbreak{}-6|\allowbreak{}8|\allowbreak{}0|\allowbreak{}-12|\allowbreak{}+++-\kern0pt{}-\kern0pt{}-+++-\kern0pt{}-\kern0pt{}-}
\RouteGroup{Retornos, memoria y transporte}
\RouteRow{Retornos con fase}{clase: 27; órbita: 27; celda: 27}
\RouteRow{Retorno cinemático}{en 108: 54; extensión: 216}
\RouteRow{Giros y cambios de hoja}{71; 107}
\RouteRow{Acarreo y diferimiento}{deuda total: 0; final: 0; bloques diferidos: 0}
\RouteRow{Tríadas finales; persistencia de Pareto}{17; no}
\RouteRow{\(K\longrightarrow K+9\)}{repeticiones: 9; cambios de memoria: 9}
\RouteGroup{Firma, fase y diferencias}
\RouteRow{\(n_A,\ s_C,\ \kappa_0,\ \nu_{120},\ \nu_{270}\)}{8, 2, 0, 0, -12}
\RouteRow{\(\tau_{12}\); firma histórica \(\mathbb Z^5\)}{+++-\kern0pt{}-\kern0pt{}-+++-\kern0pt{}-\kern0pt{}-; 8|\allowbreak{}2|\allowbreak{}0|\allowbreak{}0|\allowbreak{}-12}
\RouteRow{\(D_1,\ D_3,\ D_4,\ D_{27},\ D_{36}\)}{84, 24, 84, 24, 16}
\RouteRow{\(\Omega_{90/120},\ P_6\)}{80, 6}
\RouteGroup{Lectores y factores de acción}
\RouteRow{\(K_D\quad /\quad K_\Omega\)}{(40,\allowbreak{}84,\allowbreak{}30)\quad /\quad (80,\allowbreak{}54,\allowbreak{}6)}
\RouteRow{\(\kappa_D\)}{39.39035\allowbreak{}82768609\allowbreak{}24917849\allowbreak{}59}
\RouteRow{\(\kappa_\Omega\)}{79.60661\allowbreak{}01397948\allowbreak{}27586470\allowbreak{}91}
\RouteRow{\(R_{\mathrm{act},D}\)}{1.000000\allowbreak{}02730861\allowbreak{}73967087\allowbreak{}05}
\RouteRow{\(R_{\mathrm{act},\Omega}\)}{1.000000\allowbreak{}05518981\allowbreak{}25318591\allowbreak{}40}
\RouteGroup{Separación de prefijos}
\RouteRow{Área de ambigüedad; área \(\log_2\)}{18; 0.0}
\RouteRow{Primer bloque de separación}{órbita: 1; celda: 1; ruta: 1}
\RouteRow{Ambigüedad final}{1}
\RouteGroup{Secuencias completas}
\RouteRow{\(w_6\)}{333 255 252 333 258 414 420 258 414 333 255 252 333 258 414 420 258 414}
\RouteRow{Tríadas estables por bloque}{0 1 2 3 4 5 6 7 8 9 10 11 12 13 14 15 16 17}
\end{tabularx}\endgroup\par\vspace{5pt}\noindent{\fontsize{8.4}{10}\selectfont\hyperref[trace-239]{Procedencia y códigos originales}\hfill Ficha completa · 239/324}
\clearpage\phantomsection\label{nav:vi-ruta-240}
% VI-CATALOG-ROW rutas 240
\pdfbookmark[3]{r7c6\_h+1v+1}{route-240}
\RouteTitle{Ruta 240}{r7c6\_\allowbreak{}h+1v+1}
\noindent Celda 60\quad (7, 6)\qquad Orientación \(h=1,\ v=1\)\hfill\hyperref[sec:vi-indice-rutas]{Índice}\par\vspace{4pt}
\begingroup\fontsize{10.2}{12.5}\selectfont\setlength{\tabcolsep}{5pt}\renewcommand{\arraystretch}{1.1}\rowcolors{1}{routepale}{white}\noindent\begin{tabularx}{\linewidth}{@{}>{\raggedright\arraybackslash}p{.345\linewidth}>{\raggedright\arraybackslash}X@{}}
\RouteGroup{Identidad estructural}
\RouteRow{Clase y multisección}{neutrino\_\allowbreak{}G2\quad /\quad neutrino\_\allowbreak{}G2}
\RouteRow{Colección y sector}{neutrino; leptón}
\RouteRow{Clasificación física}{neutrino}
\RouteRow{Generación, profundidad y color}{2; G2; \textemdash{}}
\RouteRow{Ruta primaria y órbita de inversión}{\(B_1\): no\qquad 3,\allowbreak{}4;\allowbreak{}7,\allowbreak{}6}
\RouteRow{Firma de clasificación}{28|\allowbreak{}-6|\allowbreak{}8|\allowbreak{}0|\allowbreak{}-12|\allowbreak{}+++-\kern0pt{}-\kern0pt{}-+++-\kern0pt{}-\kern0pt{}-}
\RouteGroup{Retornos, memoria y transporte}
\RouteRow{Retornos con fase}{clase: 27; órbita: 27; celda: 27}
\RouteRow{Retorno cinemático}{en 108: 54; extensión: 216}
\RouteRow{Giros y cambios de hoja}{71; 107}
\RouteRow{Acarreo y diferimiento}{deuda total: 2; final: 1; bloques diferidos: 2}
\RouteRow{Tríadas finales; persistencia de Pareto}{16; no}
\RouteRow{\(K\longrightarrow K+9\)}{repeticiones: 9; cambios de memoria: 9}
\RouteGroup{Firma, fase y diferencias}
\RouteRow{\(n_A,\ s_C,\ \kappa_0,\ \nu_{120},\ \nu_{270}\)}{28, -6, 0, 0, -12}
\RouteRow{\(\tau_{12}\); firma histórica \(\mathbb Z^5\)}{+++-\kern0pt{}-\kern0pt{}-+++-\kern0pt{}-\kern0pt{}-; 28|\allowbreak{}-6|\allowbreak{}0|\allowbreak{}0|\allowbreak{}-12}
\RouteRow{\(D_1,\ D_3,\ D_4,\ D_{27},\ D_{36}\)}{66, 60, 64, 48, 32}
\RouteRow{\(\Omega_{90/120},\ P_6\)}{80, 6}
\RouteGroup{Lectores y factores de acción}
\RouteRow{\(K_D\quad /\quad K_\Omega\)}{(80,\allowbreak{}66,\allowbreak{}2)\quad /\quad (80,\allowbreak{}54,\allowbreak{}6)}
\RouteRow{\(\kappa_D\)}{79.51859\allowbreak{}71232726\allowbreak{}53414287\allowbreak{}29}
\RouteRow{\(\kappa_\Omega\)}{79.60661\allowbreak{}01397948\allowbreak{}27586470\allowbreak{}91}
\RouteRow{\(R_{\mathrm{act},D}\)}{1.000000\allowbreak{}05512879\allowbreak{}47092488\allowbreak{}80}
\RouteRow{\(R_{\mathrm{act},\Omega}\)}{1.000000\allowbreak{}05518981\allowbreak{}25318591\allowbreak{}40}
\RouteGroup{Separación de prefijos}
\RouteRow{Área de ambigüedad; área \(\log_2\)}{18; 0.0}
\RouteRow{Primer bloque de separación}{órbita: 1; celda: 1; ruta: 1}
\RouteRow{Ambigüedad final}{1}
\RouteGroup{Secuencias completas}
\RouteRow{\(w_6\)}{327 498 99 327 501 486 423 15 486 327 498 99 327 501 486 423 15 486}
\RouteRow{Tríadas estables por bloque}{0 1 2 3 4 5 6 7 7 9 10 11 12 13 14 15 16 16}
\end{tabularx}\endgroup\par\vspace{5pt}\noindent{\fontsize{8.4}{10}\selectfont\hyperref[trace-240]{Procedencia y códigos originales}\hfill Ficha completa · 240/324}
\clearpage\phantomsection\label{nav:vi-ruta-241}
% VI-CATALOG-ROW rutas 241
\pdfbookmark[2]{Celda 61}{cell-61}
\pdfbookmark[3]{r7c7\_h-1v-1}{route-241}
\RouteTitle{Ruta 241}{r7c7\_\allowbreak{}h-1v-1}
\noindent Celda 61\quad (7, 7)\qquad Orientación \(h=-1,\ v=-1\)\hfill\hyperref[sec:vi-indice-rutas]{Índice}\par\vspace{4pt}
\begingroup\fontsize{10.2}{12.5}\selectfont\setlength{\tabcolsep}{5pt}\renewcommand{\arraystretch}{1.1}\rowcolors{1}{routepale}{white}\noindent\begin{tabularx}{\linewidth}{@{}>{\raggedright\arraybackslash}p{.345\linewidth}>{\raggedright\arraybackslash}X@{}}
\RouteGroup{Identidad estructural}
\RouteRow{Clase y multisección}{charged\_\allowbreak{}lepton\_\allowbreak{}G2\quad /\quad charged\_\allowbreak{}lepton\_\allowbreak{}G2}
\RouteRow{Colección y sector}{leptón cargado; leptón}
\RouteRow{Clasificación física}{leptón cargado}
\RouteRow{Generación, profundidad y color}{2; G2; \textemdash{}}
\RouteRow{Ruta primaria y órbita de inversión}{\(B_1\): sí\qquad 3,\allowbreak{}3;\allowbreak{}7,\allowbreak{}7}
\RouteRow{Firma de clasificación}{24|\allowbreak{}0|\allowbreak{}4|\allowbreak{}0|\allowbreak{}-12|\allowbreak{}+++-\kern0pt{}-\kern0pt{}-+++-\kern0pt{}-\kern0pt{}-}
\RouteGroup{Retornos, memoria y transporte}
\RouteRow{Retornos con fase}{clase: 27; órbita: 27; celda: 27}
\RouteRow{Retorno cinemático}{en 108: 54; extensión: 216}
\RouteRow{Giros y cambios de hoja}{71; 107}
\RouteRow{Acarreo y diferimiento}{deuda total: 2; final: 1; bloques diferidos: 2}
\RouteRow{Tríadas finales; persistencia de Pareto}{16; no}
\RouteRow{\(K\longrightarrow K+9\)}{repeticiones: 9; cambios de memoria: 9}
\RouteGroup{Firma, fase y diferencias}
\RouteRow{\(n_A,\ s_C,\ \kappa_0,\ \nu_{120},\ \nu_{270}\)}{24, 0, 0, 0, -12}
\RouteRow{\(\tau_{12}\); firma histórica \(\mathbb Z^5\)}{+++-\kern0pt{}-\kern0pt{}-+++-\kern0pt{}-\kern0pt{}-; 24|\allowbreak{}0|\allowbreak{}0|\allowbreak{}0|\allowbreak{}-12}
\RouteRow{\(D_1,\ D_3,\ D_4,\ D_{27},\ D_{36}\)}{54, 56, 70, 48, 32}
\RouteRow{\(\Omega_{90/120},\ P_6\)}{80, 6}
\RouteGroup{Lectores y factores de acción}
\RouteRow{\(K_D\quad /\quad K_\Omega\)}{(80,\allowbreak{}54,\allowbreak{}7)\quad /\quad (80,\allowbreak{}54,\allowbreak{}6)}
\RouteRow{\(\kappa_D\)}{79.60672\allowbreak{}13362175\allowbreak{}19726524\allowbreak{}96}
\RouteRow{\(\kappa_\Omega\)}{79.60661\allowbreak{}01397948\allowbreak{}27586470\allowbreak{}91}
\RouteRow{\(R_{\mathrm{act},D}\)}{1.000000\allowbreak{}05518988\allowbreak{}96223153\allowbreak{}37}
\RouteRow{\(R_{\mathrm{act},\Omega}\)}{1.000000\allowbreak{}05518981\allowbreak{}25318591\allowbreak{}40}
\RouteGroup{Separación de prefijos}
\RouteRow{Área de ambigüedad; área \(\log_2\)}{22; 4.0}
\RouteRow{Primer bloque de separación}{órbita: 5; celda: 5; ruta: 5}
\RouteRow{Ambigüedad final}{1}
\RouteGroup{Secuencias completas}
\RouteRow{\(w_6\)}{486 423 15 486 429 243 657 24 243 486 423 15 486 429 243 657 24 243}
\RouteRow{Tríadas estables por bloque}{0 1 2 3 4 5 6 7 8 9 10 11 12 13 13 15 16 16}
\end{tabularx}\endgroup\par\vspace{5pt}\noindent{\fontsize{8.4}{10}\selectfont\hyperref[trace-241]{Procedencia y códigos originales}\hfill Ficha completa · 241/324}
\clearpage\phantomsection\label{nav:vi-ruta-242}
% VI-CATALOG-ROW rutas 242
\pdfbookmark[3]{r7c7\_h-1v+1}{route-242}
\RouteTitle{Ruta 242}{r7c7\_\allowbreak{}h-1v+1}
\noindent Celda 61\quad (7, 7)\qquad Orientación \(h=-1,\ v=1\)\hfill\hyperref[sec:vi-indice-rutas]{Índice}\par\vspace{4pt}
\begingroup\fontsize{10.2}{12.5}\selectfont\setlength{\tabcolsep}{5pt}\renewcommand{\arraystretch}{1.1}\rowcolors{1}{routepale}{white}\noindent\begin{tabularx}{\linewidth}{@{}>{\raggedright\arraybackslash}p{.345\linewidth}>{\raggedright\arraybackslash}X@{}}
\RouteGroup{Identidad estructural}
\RouteRow{Clase y multisección}{charged\_\allowbreak{}lepton\_\allowbreak{}G2\quad /\quad charged\_\allowbreak{}lepton\_\allowbreak{}G2}
\RouteRow{Colección y sector}{leptón cargado; leptón}
\RouteRow{Clasificación física}{leptón cargado}
\RouteRow{Generación, profundidad y color}{2; G2; \textemdash{}}
\RouteRow{Ruta primaria y órbita de inversión}{\(B_1\): no\qquad 3,\allowbreak{}3;\allowbreak{}7,\allowbreak{}7}
\RouteRow{Firma de clasificación}{24|\allowbreak{}0|\allowbreak{}4|\allowbreak{}0|\allowbreak{}-12|\allowbreak{}+++-\kern0pt{}-\kern0pt{}-+++-\kern0pt{}-\kern0pt{}-}
\RouteGroup{Retornos, memoria y transporte}
\RouteRow{Retornos con fase}{clase: 27; órbita: 27; celda: 27}
\RouteRow{Retorno cinemático}{en 108: 54; extensión: 216}
\RouteRow{Giros y cambios de hoja}{71; 107}
\RouteRow{Acarreo y diferimiento}{deuda total: 0; final: 0; bloques diferidos: 0}
\RouteRow{Tríadas finales; persistencia de Pareto}{17; no}
\RouteRow{\(K\longrightarrow K+9\)}{repeticiones: 9; cambios de memoria: 9}
\RouteGroup{Firma, fase y diferencias}
\RouteRow{\(n_A,\ s_C,\ \kappa_0,\ \nu_{120},\ \nu_{270}\)}{44, 2, 0, 0, -12}
\RouteRow{\(\tau_{12}\); firma histórica \(\mathbb Z^5\)}{+++-\kern0pt{}-\kern0pt{}-+++-\kern0pt{}-\kern0pt{}-; 44|\allowbreak{}2|\allowbreak{}0|\allowbreak{}0|\allowbreak{}-12}
\RouteRow{\(D_1,\ D_3,\ D_4,\ D_{27},\ D_{36}\)}{78, 28, 78, 36, 24}
\RouteRow{\(\Omega_{90/120},\ P_6\)}{80, 6}
\RouteGroup{Lectores y factores de acción}
\RouteRow{\(K_D\quad /\quad K_\Omega\)}{(60,\allowbreak{}78,\allowbreak{}25)\quad /\quad (80,\allowbreak{}54,\allowbreak{}6)}
\RouteRow{\(\kappa_D\)}{59.43358\allowbreak{}64101631\allowbreak{}67023563\allowbreak{}04}
\RouteRow{\(\kappa_\Omega\)}{79.60661\allowbreak{}01397948\allowbreak{}27586470\allowbreak{}91}
\RouteRow{\(R_{\mathrm{act},D}\)}{1.000000\allowbreak{}04120422\allowbreak{}24053449\allowbreak{}05}
\RouteRow{\(R_{\mathrm{act},\Omega}\)}{1.000000\allowbreak{}05518981\allowbreak{}25318591\allowbreak{}40}
\RouteGroup{Separación de prefijos}
\RouteRow{Área de ambigüedad; área \(\log_2\)}{18; 0.0}
\RouteRow{Primer bloque de separación}{órbita: 1; celda: 1; ruta: 1}
\RouteRow{Ambigüedad final}{1}
\RouteGroup{Secuencias completas}
\RouteRow{\(w_6\)}{507 504 585 507 510 510 666 672 510 507 504 585 507 510 510 666 672 510}
\RouteRow{Tríadas estables por bloque}{0 1 2 3 4 5 6 7 8 9 10 11 12 13 14 15 16 17}
\end{tabularx}\endgroup\par\vspace{5pt}\noindent{\fontsize{8.4}{10}\selectfont\hyperref[trace-242]{Procedencia y códigos originales}\hfill Ficha completa · 242/324}
\clearpage\phantomsection\label{nav:vi-ruta-243}
% VI-CATALOG-ROW rutas 243
\pdfbookmark[3]{r7c7\_h+1v-1}{route-243}
\RouteTitle{Ruta 243}{r7c7\_\allowbreak{}h+1v-1}
\noindent Celda 61\quad (7, 7)\qquad Orientación \(h=1,\ v=-1\)\hfill\hyperref[sec:vi-indice-rutas]{Índice}\par\vspace{4pt}
\begingroup\fontsize{10.2}{12.5}\selectfont\setlength{\tabcolsep}{5pt}\renewcommand{\arraystretch}{1.1}\rowcolors{1}{routepale}{white}\noindent\begin{tabularx}{\linewidth}{@{}>{\raggedright\arraybackslash}p{.345\linewidth}>{\raggedright\arraybackslash}X@{}}
\RouteGroup{Identidad estructural}
\RouteRow{Clase y multisección}{charged\_\allowbreak{}lepton\_\allowbreak{}G2\quad /\quad charged\_\allowbreak{}lepton\_\allowbreak{}G2}
\RouteRow{Colección y sector}{leptón cargado; leptón}
\RouteRow{Clasificación física}{leptón cargado}
\RouteRow{Generación, profundidad y color}{2; G2; \textemdash{}}
\RouteRow{Ruta primaria y órbita de inversión}{\(B_1\): no\qquad 3,\allowbreak{}3;\allowbreak{}7,\allowbreak{}7}
\RouteRow{Firma de clasificación}{24|\allowbreak{}0|\allowbreak{}4|\allowbreak{}0|\allowbreak{}-12|\allowbreak{}+++-\kern0pt{}-\kern0pt{}-+++-\kern0pt{}-\kern0pt{}-}
\RouteGroup{Retornos, memoria y transporte}
\RouteRow{Retornos con fase}{clase: 27; órbita: 27; celda: 27}
\RouteRow{Retorno cinemático}{en 108: 54; extensión: 216}
\RouteRow{Giros y cambios de hoja}{71; 107}
\RouteRow{Acarreo y diferimiento}{deuda total: 2; final: 0; bloques diferidos: 1}
\RouteRow{Tríadas finales; persistencia de Pareto}{17; no}
\RouteRow{\(K\longrightarrow K+9\)}{repeticiones: 9; cambios de memoria: 9}
\RouteGroup{Firma, fase y diferencias}
\RouteRow{\(n_A,\ s_C,\ \kappa_0,\ \nu_{120},\ \nu_{270}\)}{44, -2, 0, 0, -12}
\RouteRow{\(\tau_{12}\); firma histórica \(\mathbb Z^5\)}{+++-\kern0pt{}-\kern0pt{}-+++-\kern0pt{}-\kern0pt{}-; 44|\allowbreak{}-2|\allowbreak{}0|\allowbreak{}0|\allowbreak{}-12}
\RouteRow{\(D_1,\ D_3,\ D_4,\ D_{27},\ D_{36}\)}{78, 28, 78, 36, 24}
\RouteRow{\(\Omega_{90/120},\ P_6\)}{80, 6}
\RouteGroup{Lectores y factores de acción}
\RouteRow{\(K_D\quad /\quad K_\Omega\)}{(60,\allowbreak{}78,\allowbreak{}25)\quad /\quad (80,\allowbreak{}54,\allowbreak{}6)}
\RouteRow{\(\kappa_D\)}{59.43358\allowbreak{}64101631\allowbreak{}67023563\allowbreak{}04}
\RouteRow{\(\kappa_\Omega\)}{79.60661\allowbreak{}01397948\allowbreak{}27586470\allowbreak{}91}
\RouteRow{\(R_{\mathrm{act},D}\)}{1.000000\allowbreak{}04120422\allowbreak{}24053449\allowbreak{}05}
\RouteRow{\(R_{\mathrm{act},\Omega}\)}{1.000000\allowbreak{}05518981\allowbreak{}25318591\allowbreak{}40}
\RouteGroup{Separación de prefijos}
\RouteRow{Área de ambigüedad; área \(\log_2\)}{18; 0.0}
\RouteRow{Primer bloque de separación}{órbita: 1; celda: 1; ruta: 1}
\RouteRow{Ambigüedad final}{1}
\RouteGroup{Secuencias completas}
\RouteRow{\(w_6\)}{666 672 510 666 666 585 507 504 585 666 672 510 666 666 585 507 504 585}
\RouteRow{Tríadas estables por bloque}{0 1 2 3 4 5 6 7 8 9 10 11 12 12 13 15 16 17}
\end{tabularx}\endgroup\par\vspace{5pt}\noindent{\fontsize{8.4}{10}\selectfont\hyperref[trace-243]{Procedencia y códigos originales}\hfill Ficha completa · 243/324}
\clearpage\phantomsection\label{nav:vi-ruta-244}
% VI-CATALOG-ROW rutas 244
\pdfbookmark[3]{r7c7\_h+1v+1}{route-244}
\RouteTitle{Ruta 244}{r7c7\_\allowbreak{}h+1v+1}
\noindent Celda 61\quad (7, 7)\qquad Orientación \(h=1,\ v=1\)\hfill\hyperref[sec:vi-indice-rutas]{Índice}\par\vspace{4pt}
\begingroup\fontsize{10.2}{12.5}\selectfont\setlength{\tabcolsep}{5pt}\renewcommand{\arraystretch}{1.1}\rowcolors{1}{routepale}{white}\noindent\begin{tabularx}{\linewidth}{@{}>{\raggedright\arraybackslash}p{.345\linewidth}>{\raggedright\arraybackslash}X@{}}
\RouteGroup{Identidad estructural}
\RouteRow{Clase y multisección}{charged\_\allowbreak{}lepton\_\allowbreak{}G2\quad /\quad charged\_\allowbreak{}lepton\_\allowbreak{}G2}
\RouteRow{Colección y sector}{leptón cargado; leptón}
\RouteRow{Clasificación física}{leptón cargado}
\RouteRow{Generación, profundidad y color}{2; G2; \textemdash{}}
\RouteRow{Ruta primaria y órbita de inversión}{\(B_1\): no\qquad 3,\allowbreak{}3;\allowbreak{}7,\allowbreak{}7}
\RouteRow{Firma de clasificación}{24|\allowbreak{}0|\allowbreak{}4|\allowbreak{}0|\allowbreak{}-12|\allowbreak{}+++-\kern0pt{}-\kern0pt{}-+++-\kern0pt{}-\kern0pt{}-}
\RouteGroup{Retornos, memoria y transporte}
\RouteRow{Retornos con fase}{clase: 27; órbita: 27; celda: 27}
\RouteRow{Retorno cinemático}{en 108: 54; extensión: 216}
\RouteRow{Giros y cambios de hoja}{71; 107}
\RouteRow{Acarreo y diferimiento}{deuda total: 2; final: 0; bloques diferidos: 1}
\RouteRow{Tríadas finales; persistencia de Pareto}{17; no}
\RouteRow{\(K\longrightarrow K+9\)}{repeticiones: 9; cambios de memoria: 9}
\RouteGroup{Firma, fase y diferencias}
\RouteRow{\(n_A,\ s_C,\ \kappa_0,\ \nu_{120},\ \nu_{270}\)}{24, 0, 0, 0, -12}
\RouteRow{\(\tau_{12}\); firma histórica \(\mathbb Z^5\)}{+++-\kern0pt{}-\kern0pt{}-+++-\kern0pt{}-\kern0pt{}-; 24|\allowbreak{}0|\allowbreak{}0|\allowbreak{}0|\allowbreak{}-12}
\RouteRow{\(D_1,\ D_3,\ D_4,\ D_{27},\ D_{36}\)}{54, 56, 70, 48, 32}
\RouteRow{\(\Omega_{90/120},\ P_6\)}{80, 6}
\RouteGroup{Lectores y factores de acción}
\RouteRow{\(K_D\quad /\quad K_\Omega\)}{(80,\allowbreak{}54,\allowbreak{}7)\quad /\quad (80,\allowbreak{}54,\allowbreak{}6)}
\RouteRow{\(\kappa_D\)}{79.60672\allowbreak{}13362175\allowbreak{}19726524\allowbreak{}96}
\RouteRow{\(\kappa_\Omega\)}{79.60661\allowbreak{}01397948\allowbreak{}27586470\allowbreak{}91}
\RouteRow{\(R_{\mathrm{act},D}\)}{1.000000\allowbreak{}05518988\allowbreak{}96223153\allowbreak{}37}
\RouteRow{\(R_{\mathrm{act},\Omega}\)}{1.000000\allowbreak{}05518981\allowbreak{}25318591\allowbreak{}40}
\RouteGroup{Separación de prefijos}
\RouteRow{Área de ambigüedad; área \(\log_2\)}{22; 4.0}
\RouteRow{Primer bloque de separación}{órbita: 5; celda: 5; ruta: 5}
\RouteRow{Ambigüedad final}{1}
\RouteGroup{Secuencias completas}
\RouteRow{\(w_6\)}{657 24 243 657 18 15 486 423 15 657 24 243 657 18 15 486 423 15}
\RouteRow{Tríadas estables por bloque}{0 1 2 3 4 5 6 7 8 9 10 11 12 13 14 14 15 17}
\end{tabularx}\endgroup\par\vspace{5pt}\noindent{\fontsize{8.4}{10}\selectfont\hyperref[trace-244]{Procedencia y códigos originales}\hfill Ficha completa · 244/324}
\clearpage\phantomsection\label{nav:vi-ruta-245}
% VI-CATALOG-ROW rutas 245
\pdfbookmark[2]{Celda 62}{cell-62}
\pdfbookmark[3]{r7c8\_h-1v-1}{route-245}
\RouteTitle{Ruta 245}{r7c8\_\allowbreak{}h-1v-1}
\noindent Celda 62\quad (7, 8)\qquad Orientación \(h=-1,\ v=-1\)\hfill\hyperref[sec:vi-indice-rutas]{Índice}\par\vspace{4pt}
\begingroup\fontsize{10.2}{12.5}\selectfont\setlength{\tabcolsep}{5pt}\renewcommand{\arraystretch}{1.1}\rowcolors{1}{routepale}{white}\noindent\begin{tabularx}{\linewidth}{@{}>{\raggedright\arraybackslash}p{.345\linewidth}>{\raggedright\arraybackslash}X@{}}
\RouteGroup{Identidad estructural}
\RouteRow{Clase y multisección}{neutrino\_\allowbreak{}G3\quad /\quad neutrino\_\allowbreak{}G3}
\RouteRow{Colección y sector}{neutrino; leptón}
\RouteRow{Clasificación física}{neutrino}
\RouteRow{Generación, profundidad y color}{3; G3; \textemdash{}}
\RouteRow{Ruta primaria y órbita de inversión}{\(B_1\): no\qquad 3,\allowbreak{}2;\allowbreak{}7,\allowbreak{}8}
\RouteRow{Firma de clasificación}{28|\allowbreak{}6|\allowbreak{}8|\allowbreak{}0|\allowbreak{}-12|\allowbreak{}+++-\kern0pt{}-\kern0pt{}-+++-\kern0pt{}-\kern0pt{}-}
\RouteGroup{Retornos, memoria y transporte}
\RouteRow{Retornos con fase}{clase: 27; órbita: 27; celda: 27}
\RouteRow{Retorno cinemático}{en 108: 54; extensión: 216}
\RouteRow{Giros y cambios de hoja}{71; 107}
\RouteRow{Acarreo y diferimiento}{deuda total: 1; final: 0; bloques diferidos: 1}
\RouteRow{Tríadas finales; persistencia de Pareto}{17; sí}
\RouteRow{\(K\longrightarrow K+9\)}{repeticiones: 9; cambios de memoria: 9}
\RouteGroup{Firma, fase y diferencias}
\RouteRow{\(n_A,\ s_C,\ \kappa_0,\ \nu_{120},\ \nu_{270}\)}{28, -6, 0, 0, -12}
\RouteRow{\(\tau_{12}\); firma histórica \(\mathbb Z^5\)}{+++-\kern0pt{}-\kern0pt{}-+++-\kern0pt{}-\kern0pt{}-; 28|\allowbreak{}-6|\allowbreak{}0|\allowbreak{}0|\allowbreak{}-12}
\RouteRow{\(D_1,\ D_3,\ D_4,\ D_{27},\ D_{36}\)}{66, 64, 64, 60, 40}
\RouteRow{\(\Omega_{90/120},\ P_6\)}{80, 6}
\RouteGroup{Lectores y factores de acción}
\RouteRow{\(K_D\quad /\quad K_\Omega\)}{(100,\allowbreak{}66,\allowbreak{}0)\quad /\quad (80,\allowbreak{}54,\allowbreak{}6)}
\RouteRow{\(\kappa_D\)}{99.51837\allowbreak{}47304272\allowbreak{}69134179\allowbreak{}18}
\RouteRow{\(\kappa_\Omega\)}{79.60661\allowbreak{}01397948\allowbreak{}27586470\allowbreak{}91}
\RouteRow{\(R_{\mathrm{act},D}\)}{1.000000\allowbreak{}06899427\allowbreak{}66450218\allowbreak{}95}
\RouteRow{\(R_{\mathrm{act},\Omega}\)}{1.000000\allowbreak{}05518981\allowbreak{}25318591\allowbreak{}40}
\RouteGroup{Separación de prefijos}
\RouteRow{Área de ambigüedad; área \(\log_2\)}{36; 18.0}
\RouteRow{Primer bloque de separación}{órbita: \textemdash{}; celda: \textemdash{}; ruta: \textemdash{}}
\RouteRow{Ambigüedad final}{2}
\RouteGroup{Secuencias completas}
\RouteRow{\(w_6\)}{90 570 264 90 567 657 24 243 657 90 570 264 90 567 657 24 243 657}
\RouteRow{Tríadas estables por bloque}{0 1 2 3 4 5 6 7 8 9 10 11 12 13 13 15 16 17}
\end{tabularx}\endgroup\par\vspace{5pt}\noindent{\fontsize{8.4}{10}\selectfont\hyperref[trace-245]{Procedencia y códigos originales}\hfill Ficha completa · 245/324}
\clearpage\phantomsection\label{nav:vi-ruta-246}
% VI-CATALOG-ROW rutas 246
\pdfbookmark[3]{r7c8\_h-1v+1}{route-246}
\RouteTitle{Ruta 246}{r7c8\_\allowbreak{}h-1v+1}
\noindent Celda 62\quad (7, 8)\qquad Orientación \(h=-1,\ v=1\)\hfill\hyperref[sec:vi-indice-rutas]{Índice}\par\vspace{4pt}
\begingroup\fontsize{10.2}{12.5}\selectfont\setlength{\tabcolsep}{5pt}\renewcommand{\arraystretch}{1.1}\rowcolors{1}{routepale}{white}\noindent\begin{tabularx}{\linewidth}{@{}>{\raggedright\arraybackslash}p{.345\linewidth}>{\raggedright\arraybackslash}X@{}}
\RouteGroup{Identidad estructural}
\RouteRow{Clase y multisección}{neutrino\_\allowbreak{}G3\quad /\quad neutrino\_\allowbreak{}G3}
\RouteRow{Colección y sector}{neutrino; leptón}
\RouteRow{Clasificación física}{neutrino}
\RouteRow{Generación, profundidad y color}{3; G3; \textemdash{}}
\RouteRow{Ruta primaria y órbita de inversión}{\(B_1\): no\qquad 3,\allowbreak{}2;\allowbreak{}7,\allowbreak{}8}
\RouteRow{Firma de clasificación}{28|\allowbreak{}6|\allowbreak{}8|\allowbreak{}0|\allowbreak{}-12|\allowbreak{}+++-\kern0pt{}-\kern0pt{}-+++-\kern0pt{}-\kern0pt{}-}
\RouteGroup{Retornos, memoria y transporte}
\RouteRow{Retornos con fase}{clase: 27; órbita: 27; celda: 27}
\RouteRow{Retorno cinemático}{en 108: 54; extensión: 216}
\RouteRow{Giros y cambios de hoja}{71; 107}
\RouteRow{Acarreo y diferimiento}{deuda total: 2; final: 0; bloques diferidos: 2}
\RouteRow{Tríadas finales; persistencia de Pareto}{17; sí}
\RouteRow{\(K\longrightarrow K+9\)}{repeticiones: 9; cambios de memoria: 9}
\RouteGroup{Firma, fase y diferencias}
\RouteRow{\(n_A,\ s_C,\ \kappa_0,\ \nu_{120},\ \nu_{270}\)}{14, 6, 0, 0, -12}
\RouteRow{\(\tau_{12}\); firma histórica \(\mathbb Z^5\)}{+++-\kern0pt{}-\kern0pt{}-+++-\kern0pt{}-\kern0pt{}-; 14|\allowbreak{}6|\allowbreak{}0|\allowbreak{}0|\allowbreak{}-12}
\RouteRow{\(D_1,\ D_3,\ D_4,\ D_{27},\ D_{36}\)}{24, 20, 24, 24, 16}
\RouteRow{\(\Omega_{90/120},\ P_6\)}{80, 6}
\RouteGroup{Lectores y factores de acción}
\RouteRow{\(K_D\quad /\quad K_\Omega\)}{(40,\allowbreak{}24,\allowbreak{}2)\quad /\quad (80,\allowbreak{}54,\allowbreak{}6)}
\RouteRow{\(\kappa_D\)}{39.82508\allowbreak{}59311825\allowbreak{}73056173\allowbreak{}26}
\RouteRow{\(\kappa_\Omega\)}{79.60661\allowbreak{}01397948\allowbreak{}27586470\allowbreak{}91}
\RouteRow{\(R_{\mathrm{act},D}\)}{1.000000\allowbreak{}02761000\allowbreak{}61595142\allowbreak{}15}
\RouteRow{\(R_{\mathrm{act},\Omega}\)}{1.000000\allowbreak{}05518981\allowbreak{}25318591\allowbreak{}40}
\RouteGroup{Separación de prefijos}
\RouteRow{Área de ambigüedad; área \(\log_2\)}{36; 18.0}
\RouteRow{Primer bloque de separación}{órbita: \textemdash{}; celda: \textemdash{}; ruta: \textemdash{}}
\RouteRow{Ambigüedad final}{2}
\RouteGroup{Secuencias completas}
\RouteRow{\(w_6\)}{81 3 0 81 0 3 0 81 3 81 3 0 81 0 3 0 81 3}
\RouteRow{Tríadas estables por bloque}{0 1 2 3 4 5 6 6 8 9 10 11 12 13 14 15 15 17}
\end{tabularx}\endgroup\par\vspace{5pt}\noindent{\fontsize{8.4}{10}\selectfont\hyperref[trace-246]{Procedencia y códigos originales}\hfill Ficha completa · 246/324}
\clearpage\phantomsection\label{nav:vi-ruta-247}
% VI-CATALOG-ROW rutas 247
\pdfbookmark[3]{r7c8\_h+1v-1}{route-247}
\RouteTitle{Ruta 247}{r7c8\_\allowbreak{}h+1v-1}
\noindent Celda 62\quad (7, 8)\qquad Orientación \(h=1,\ v=-1\)\hfill\hyperref[sec:vi-indice-rutas]{Índice}\par\vspace{4pt}
\begingroup\fontsize{10.2}{12.5}\selectfont\setlength{\tabcolsep}{5pt}\renewcommand{\arraystretch}{1.1}\rowcolors{1}{routepale}{white}\noindent\begin{tabularx}{\linewidth}{@{}>{\raggedright\arraybackslash}p{.345\linewidth}>{\raggedright\arraybackslash}X@{}}
\RouteGroup{Identidad estructural}
\RouteRow{Clase y multisección}{neutrino\_\allowbreak{}G3\quad /\quad neutrino\_\allowbreak{}G3}
\RouteRow{Colección y sector}{neutrino; leptón}
\RouteRow{Clasificación física}{neutrino}
\RouteRow{Generación, profundidad y color}{3; G3; \textemdash{}}
\RouteRow{Ruta primaria y órbita de inversión}{\(B_1\): sí\qquad 3,\allowbreak{}2;\allowbreak{}7,\allowbreak{}8}
\RouteRow{Firma de clasificación}{28|\allowbreak{}6|\allowbreak{}8|\allowbreak{}0|\allowbreak{}-12|\allowbreak{}+++-\kern0pt{}-\kern0pt{}-+++-\kern0pt{}-\kern0pt{}-}
\RouteGroup{Retornos, memoria y transporte}
\RouteRow{Retornos con fase}{clase: 27; órbita: 27; celda: 27}
\RouteRow{Retorno cinemático}{en 108: 54; extensión: 216}
\RouteRow{Giros y cambios de hoja}{71; 107}
\RouteRow{Acarreo y diferimiento}{deuda total: 0; final: 0; bloques diferidos: 0}
\RouteRow{Tríadas finales; persistencia de Pareto}{17; sí}
\RouteRow{\(K\longrightarrow K+9\)}{repeticiones: 9; cambios de memoria: 9}
\RouteGroup{Firma, fase y diferencias}
\RouteRow{\(n_A,\ s_C,\ \kappa_0,\ \nu_{120},\ \nu_{270}\)}{14, -6, 0, 0, -12}
\RouteRow{\(\tau_{12}\); firma histórica \(\mathbb Z^5\)}{+++-\kern0pt{}-\kern0pt{}-+++-\kern0pt{}-\kern0pt{}-; 14|\allowbreak{}-6|\allowbreak{}0|\allowbreak{}0|\allowbreak{}-12}
\RouteRow{\(D_1,\ D_3,\ D_4,\ D_{27},\ D_{36}\)}{24, 20, 24, 24, 16}
\RouteRow{\(\Omega_{90/120},\ P_6\)}{80, 6}
\RouteGroup{Lectores y factores de acción}
\RouteRow{\(K_D\quad /\quad K_\Omega\)}{(40,\allowbreak{}24,\allowbreak{}2)\quad /\quad (80,\allowbreak{}54,\allowbreak{}6)}
\RouteRow{\(\kappa_D\)}{39.82508\allowbreak{}59311825\allowbreak{}73056173\allowbreak{}26}
\RouteRow{\(\kappa_\Omega\)}{79.60661\allowbreak{}01397948\allowbreak{}27586470\allowbreak{}91}
\RouteRow{\(R_{\mathrm{act},D}\)}{1.000000\allowbreak{}02761000\allowbreak{}61595142\allowbreak{}15}
\RouteRow{\(R_{\mathrm{act},\Omega}\)}{1.000000\allowbreak{}05518981\allowbreak{}25318591\allowbreak{}40}
\RouteGroup{Separación de prefijos}
\RouteRow{Área de ambigüedad; área \(\log_2\)}{36; 18.0}
\RouteRow{Primer bloque de separación}{órbita: \textemdash{}; celda: \textemdash{}; ruta: \textemdash{}}
\RouteRow{Ambigüedad final}{2}
\RouteGroup{Secuencias completas}
\RouteRow{\(w_6\)}{0 81 3 0 84 0 81 3 0 0 81 3 0 84 0 81 3 0}
\RouteRow{Tríadas estables por bloque}{0 1 2 3 4 5 6 7 8 9 10 11 12 13 14 15 16 17}
\end{tabularx}\endgroup\par\vspace{5pt}\noindent{\fontsize{8.4}{10}\selectfont\hyperref[trace-247]{Procedencia y códigos originales}\hfill Ficha completa · 247/324}
\clearpage\phantomsection\label{nav:vi-ruta-248}
% VI-CATALOG-ROW rutas 248
\pdfbookmark[3]{r7c8\_h+1v+1}{route-248}
\RouteTitle{Ruta 248}{r7c8\_\allowbreak{}h+1v+1}
\noindent Celda 62\quad (7, 8)\qquad Orientación \(h=1,\ v=1\)\hfill\hyperref[sec:vi-indice-rutas]{Índice}\par\vspace{4pt}
\begingroup\fontsize{10.2}{12.5}\selectfont\setlength{\tabcolsep}{5pt}\renewcommand{\arraystretch}{1.1}\rowcolors{1}{routepale}{white}\noindent\begin{tabularx}{\linewidth}{@{}>{\raggedright\arraybackslash}p{.345\linewidth}>{\raggedright\arraybackslash}X@{}}
\RouteGroup{Identidad estructural}
\RouteRow{Clase y multisección}{neutrino\_\allowbreak{}G3\quad /\quad neutrino\_\allowbreak{}G3}
\RouteRow{Colección y sector}{neutrino; leptón}
\RouteRow{Clasificación física}{neutrino}
\RouteRow{Generación, profundidad y color}{3; G3; \textemdash{}}
\RouteRow{Ruta primaria y órbita de inversión}{\(B_1\): no\qquad 3,\allowbreak{}2;\allowbreak{}7,\allowbreak{}8}
\RouteRow{Firma de clasificación}{28|\allowbreak{}6|\allowbreak{}8|\allowbreak{}0|\allowbreak{}-12|\allowbreak{}+++-\kern0pt{}-\kern0pt{}-+++-\kern0pt{}-\kern0pt{}-}
\RouteGroup{Retornos, memoria y transporte}
\RouteRow{Retornos con fase}{clase: 27; órbita: 27; celda: 27}
\RouteRow{Retorno cinemático}{en 108: 54; extensión: 216}
\RouteRow{Giros y cambios de hoja}{71; 107}
\RouteRow{Acarreo y diferimiento}{deuda total: 0; final: 0; bloques diferidos: 0}
\RouteRow{Tríadas finales; persistencia de Pareto}{17; sí}
\RouteRow{\(K\longrightarrow K+9\)}{repeticiones: 9; cambios de memoria: 9}
\RouteGroup{Firma, fase y diferencias}
\RouteRow{\(n_A,\ s_C,\ \kappa_0,\ \nu_{120},\ \nu_{270}\)}{28, 6, 0, 0, -12}
\RouteRow{\(\tau_{12}\); firma histórica \(\mathbb Z^5\)}{+++-\kern0pt{}-\kern0pt{}-+++-\kern0pt{}-\kern0pt{}-; 28|\allowbreak{}6|\allowbreak{}0|\allowbreak{}0|\allowbreak{}-12}
\RouteRow{\(D_1,\ D_3,\ D_4,\ D_{27},\ D_{36}\)}{66, 64, 64, 60, 40}
\RouteRow{\(\Omega_{90/120},\ P_6\)}{80, 6}
\RouteGroup{Lectores y factores de acción}
\RouteRow{\(K_D\quad /\quad K_\Omega\)}{(100,\allowbreak{}66,\allowbreak{}0)\quad /\quad (80,\allowbreak{}54,\allowbreak{}6)}
\RouteRow{\(\kappa_D\)}{99.51837\allowbreak{}47304272\allowbreak{}69134179\allowbreak{}18}
\RouteRow{\(\kappa_\Omega\)}{79.60661\allowbreak{}01397948\allowbreak{}27586470\allowbreak{}91}
\RouteRow{\(R_{\mathrm{act},D}\)}{1.000000\allowbreak{}06899427\allowbreak{}66450218\allowbreak{}95}
\RouteRow{\(R_{\mathrm{act},\Omega}\)}{1.000000\allowbreak{}05518981\allowbreak{}25318591\allowbreak{}40}
\RouteGroup{Separación de prefijos}
\RouteRow{Área de ambigüedad; área \(\log_2\)}{36; 18.0}
\RouteRow{Primer bloque de separación}{órbita: \textemdash{}; celda: \textemdash{}; ruta: \textemdash{}}
\RouteRow{Ambigüedad final}{2}
\RouteGroup{Secuencias completas}
\RouteRow{\(w_6\)}{24 243 657 24 246 264 90 570 264 24 243 657 24 246 264 90 570 264}
\RouteRow{Tríadas estables por bloque}{0 1 2 3 4 5 6 7 8 9 10 11 12 13 14 15 16 17}
\end{tabularx}\endgroup\par\vspace{5pt}\noindent{\fontsize{8.4}{10}\selectfont\hyperref[trace-248]{Procedencia y códigos originales}\hfill Ficha completa · 248/324}
\clearpage\phantomsection\label{nav:vi-ruta-249}
% VI-CATALOG-ROW rutas 249
\pdfbookmark[2]{Celda 63}{cell-63}
\pdfbookmark[3]{r7c9\_h-1v-1}{route-249}
\RouteTitle{Ruta 249}{r7c9\_\allowbreak{}h-1v-1}
\noindent Celda 63\quad (7, 9)\qquad Orientación \(h=-1,\ v=-1\)\hfill\hyperref[sec:vi-indice-rutas]{Índice}\par\vspace{4pt}
\begingroup\fontsize{10.2}{12.5}\selectfont\setlength{\tabcolsep}{5pt}\renewcommand{\arraystretch}{1.1}\rowcolors{1}{routepale}{white}\noindent\begin{tabularx}{\linewidth}{@{}>{\raggedright\arraybackslash}p{.345\linewidth}>{\raggedright\arraybackslash}X@{}}
\RouteGroup{Identidad estructural}
\RouteRow{Clase y multisección}{charged\_\allowbreak{}lepton\_\allowbreak{}G4\quad /\quad charged\_\allowbreak{}lepton\_\allowbreak{}G4}
\RouteRow{Colección y sector}{leptón cargado; leptón}
\RouteRow{Clasificación física}{leptón cargado}
\RouteRow{Generación, profundidad y color}{4; G4; \textemdash{}}
\RouteRow{Ruta primaria y órbita de inversión}{\(B_1\): no\qquad 3,\allowbreak{}1;\allowbreak{}7,\allowbreak{}9}
\RouteRow{Firma de clasificación}{28|\allowbreak{}-6|\allowbreak{}6|\allowbreak{}0|\allowbreak{}-12|\allowbreak{}+++-\kern0pt{}-\kern0pt{}-+++-\kern0pt{}-\kern0pt{}-}
\RouteGroup{Retornos, memoria y transporte}
\RouteRow{Retornos con fase}{clase: 18; órbita: 27; celda: 27}
\RouteRow{Retorno cinemático}{en 108: 54; extensión: 216}
\RouteRow{Giros y cambios de hoja}{71; 107}
\RouteRow{Acarreo y diferimiento}{deuda total: 1; final: 1; bloques diferidos: 1}
\RouteRow{Tríadas finales; persistencia de Pareto}{16; no}
\RouteRow{\(K\longrightarrow K+9\)}{repeticiones: 9; cambios de memoria: 9}
\RouteGroup{Firma, fase y diferencias}
\RouteRow{\(n_A,\ s_C,\ \kappa_0,\ \nu_{120},\ \nu_{270}\)}{28, 6, 0, 4, -12}
\RouteRow{\(\tau_{12}\); firma histórica \(\mathbb Z^5\)}{+++-\kern0pt{}-\kern0pt{}-+++-\kern0pt{}-\kern0pt{}-; 28|\allowbreak{}6|\allowbreak{}0|\allowbreak{}4|\allowbreak{}-12}
\RouteRow{\(D_1,\ D_3,\ D_4,\ D_{27},\ D_{36}\)}{66, 60, 64, 48, 32}
\RouteRow{\(\Omega_{90/120},\ P_6\)}{80, 6}
\RouteGroup{Lectores y factores de acción}
\RouteRow{\(K_D\quad /\quad K_\Omega\)}{(80,\allowbreak{}66,\allowbreak{}2)\quad /\quad (80,\allowbreak{}54,\allowbreak{}6)}
\RouteRow{\(\kappa_D\)}{79.51859\allowbreak{}71232726\allowbreak{}53414287\allowbreak{}29}
\RouteRow{\(\kappa_\Omega\)}{79.60661\allowbreak{}01397948\allowbreak{}27586470\allowbreak{}91}
\RouteRow{\(R_{\mathrm{act},D}\)}{1.000000\allowbreak{}05512879\allowbreak{}47092488\allowbreak{}80}
\RouteRow{\(R_{\mathrm{act},\Omega}\)}{1.000000\allowbreak{}05518981\allowbreak{}25318591\allowbreak{}40}
\RouteGroup{Separación de prefijos}
\RouteRow{Área de ambigüedad; área \(\log_2\)}{54; 28.52932\allowbreak{}50129808}
\RouteRow{Primer bloque de separación}{órbita: \textemdash{}; celda: \textemdash{}; ruta: \textemdash{}}
\RouteRow{Ambigüedad final}{3}
\RouteGroup{Secuencias completas}
\RouteRow{\(w_6\)}{423 15 486 423 12 99 327 498 99 423 15 486 423 12 99 327 498 99}
\RouteRow{Tríadas estables por bloque}{0 1 2 3 4 5 6 7 8 9 10 11 12 13 14 15 16 16}
\end{tabularx}\endgroup\par\vspace{5pt}\noindent{\fontsize{8.4}{10}\selectfont\hyperref[trace-249]{Procedencia y códigos originales}\hfill Ficha completa · 249/324}
\clearpage\phantomsection\label{nav:vi-ruta-250}
% VI-CATALOG-ROW rutas 250
\pdfbookmark[3]{r7c9\_h-1v+1}{route-250}
\RouteTitle{Ruta 250}{r7c9\_\allowbreak{}h-1v+1}
\noindent Celda 63\quad (7, 9)\qquad Orientación \(h=-1,\ v=1\)\hfill\hyperref[sec:vi-indice-rutas]{Índice}\par\vspace{4pt}
\begingroup\fontsize{10.2}{12.5}\selectfont\setlength{\tabcolsep}{5pt}\renewcommand{\arraystretch}{1.1}\rowcolors{1}{routepale}{white}\noindent\begin{tabularx}{\linewidth}{@{}>{\raggedright\arraybackslash}p{.345\linewidth}>{\raggedright\arraybackslash}X@{}}
\RouteGroup{Identidad estructural}
\RouteRow{Clase y multisección}{charged\_\allowbreak{}lepton\_\allowbreak{}G4\quad /\quad charged\_\allowbreak{}lepton\_\allowbreak{}G4}
\RouteRow{Colección y sector}{leptón cargado; leptón}
\RouteRow{Clasificación física}{leptón cargado}
\RouteRow{Generación, profundidad y color}{4; G4; \textemdash{}}
\RouteRow{Ruta primaria y órbita de inversión}{\(B_1\): no\qquad 3,\allowbreak{}1;\allowbreak{}7,\allowbreak{}9}
\RouteRow{Firma de clasificación}{28|\allowbreak{}-6|\allowbreak{}6|\allowbreak{}0|\allowbreak{}-12|\allowbreak{}+++-\kern0pt{}-\kern0pt{}-+++-\kern0pt{}-\kern0pt{}-}
\RouteGroup{Retornos, memoria y transporte}
\RouteRow{Retornos con fase}{clase: 27; órbita: 27; celda: 27}
\RouteRow{Retorno cinemático}{en 108: 54; extensión: 216}
\RouteRow{Giros y cambios de hoja}{71; 107}
\RouteRow{Acarreo y diferimiento}{deuda total: 1; final: 0; bloques diferidos: 1}
\RouteRow{Tríadas finales; persistencia de Pareto}{17; no}
\RouteRow{\(K\longrightarrow K+9\)}{repeticiones: 9; cambios de memoria: 9}
\RouteGroup{Firma, fase y diferencias}
\RouteRow{\(n_A,\ s_C,\ \kappa_0,\ \nu_{120},\ \nu_{270}\)}{50, 4, -2, 2, -8}
\RouteRow{\(\tau_{12}\); firma histórica \(\mathbb Z^5\)}{+++-\kern0pt{}-0+++-\kern0pt{}-0; 50|\allowbreak{}4|\allowbreak{}-2|\allowbreak{}2|\allowbreak{}-8}
\RouteRow{\(D_1,\ D_3,\ D_4,\ D_{27},\ D_{36}\)}{84, 24, 84, 24, 16}
\RouteRow{\(\Omega_{90/120},\ P_6\)}{56, 5}
\RouteGroup{Lectores y factores de acción}
\RouteRow{\(K_D\quad /\quad K_\Omega\)}{(40,\allowbreak{}84,\allowbreak{}30)\quad /\quad (56,\allowbreak{}54,\allowbreak{}5)}
\RouteRow{\(\kappa_D\)}{39.39035\allowbreak{}82768609\allowbreak{}24917849\allowbreak{}59}
\RouteRow{\(\kappa_\Omega\)}{55.60649\allowbreak{}89433721\allowbreak{}35446416\allowbreak{}86}
\RouteRow{\(R_{\mathrm{act},D}\)}{1.000000\allowbreak{}02730861\allowbreak{}73967087\allowbreak{}05}
\RouteRow{\(R_{\mathrm{act},\Omega}\)}{1.000000\allowbreak{}03855097\allowbreak{}23541416\allowbreak{}45}
\RouteGroup{Separación de prefijos}
\RouteRow{Área de ambigüedad; área \(\log_2\)}{54; 28.52932\allowbreak{}50129808}
\RouteRow{Primer bloque de separación}{órbita: \textemdash{}; celda: \textemdash{}; ruta: \textemdash{}}
\RouteRow{Ambigüedad final}{3}
\RouteGroup{Secuencias completas}
\RouteRow{\(w_6\)}{420 258 414 420 255 252 333 255 252 420 258 414 420 255 252 333 255 252}
\RouteRow{Tríadas estables por bloque}{0 1 2 3 4 5 6 7 8 9 10 11 12 13 14 15 15 17}
\end{tabularx}\endgroup\par\vspace{5pt}\noindent{\fontsize{8.4}{10}\selectfont\hyperref[trace-250]{Procedencia y códigos originales}\hfill Ficha completa · 250/324}
\clearpage\phantomsection\label{nav:vi-ruta-251}
% VI-CATALOG-ROW rutas 251
\pdfbookmark[3]{r7c9\_h+1v-1}{route-251}
\RouteTitle{Ruta 251}{r7c9\_\allowbreak{}h+1v-1}
\noindent Celda 63\quad (7, 9)\qquad Orientación \(h=1,\ v=-1\)\hfill\hyperref[sec:vi-indice-rutas]{Índice}\par\vspace{4pt}
\begingroup\fontsize{10.2}{12.5}\selectfont\setlength{\tabcolsep}{5pt}\renewcommand{\arraystretch}{1.1}\rowcolors{1}{routepale}{white}\noindent\begin{tabularx}{\linewidth}{@{}>{\raggedright\arraybackslash}p{.345\linewidth}>{\raggedright\arraybackslash}X@{}}
\RouteGroup{Identidad estructural}
\RouteRow{Clase y multisección}{charged\_\allowbreak{}lepton\_\allowbreak{}G4\quad /\quad charged\_\allowbreak{}lepton\_\allowbreak{}G4}
\RouteRow{Colección y sector}{leptón cargado; leptón}
\RouteRow{Clasificación física}{leptón cargado}
\RouteRow{Generación, profundidad y color}{4; G4; \textemdash{}}
\RouteRow{Ruta primaria y órbita de inversión}{\(B_1\): no\qquad 3,\allowbreak{}1;\allowbreak{}7,\allowbreak{}9}
\RouteRow{Firma de clasificación}{28|\allowbreak{}-6|\allowbreak{}6|\allowbreak{}0|\allowbreak{}-12|\allowbreak{}+++-\kern0pt{}-\kern0pt{}-+++-\kern0pt{}-\kern0pt{}-}
\RouteGroup{Retornos, memoria y transporte}
\RouteRow{Retornos con fase}{clase: 27; órbita: 27; celda: 27}
\RouteRow{Retorno cinemático}{en 108: 54; extensión: 216}
\RouteRow{Giros y cambios de hoja}{71; 107}
\RouteRow{Acarreo y diferimiento}{deuda total: 0; final: 0; bloques diferidos: 0}
\RouteRow{Tríadas finales; persistencia de Pareto}{17; no}
\RouteRow{\(K\longrightarrow K+9\)}{repeticiones: 9; cambios de memoria: 9}
\RouteGroup{Firma, fase y diferencias}
\RouteRow{\(n_A,\ s_C,\ \kappa_0,\ \nu_{120},\ \nu_{270}\)}{50, -4, 2, -2, -8}
\RouteRow{\(\tau_{12}\); firma histórica \(\mathbb Z^5\)}{++0-\kern0pt{}-\kern0pt{}-++0-\kern0pt{}-\kern0pt{}-; 50|\allowbreak{}-4|\allowbreak{}2|\allowbreak{}-2|\allowbreak{}-8}
\RouteRow{\(D_1,\ D_3,\ D_4,\ D_{27},\ D_{36}\)}{84, 24, 84, 24, 16}
\RouteRow{\(\Omega_{90/120},\ P_6\)}{56, 5}
\RouteGroup{Lectores y factores de acción}
\RouteRow{\(K_D\quad /\quad K_\Omega\)}{(40,\allowbreak{}84,\allowbreak{}30)\quad /\quad (56,\allowbreak{}54,\allowbreak{}5)}
\RouteRow{\(\kappa_D\)}{39.39035\allowbreak{}82768609\allowbreak{}24917849\allowbreak{}59}
\RouteRow{\(\kappa_\Omega\)}{55.60649\allowbreak{}89433721\allowbreak{}35446416\allowbreak{}86}
\RouteRow{\(R_{\mathrm{act},D}\)}{1.000000\allowbreak{}02730861\allowbreak{}73967087\allowbreak{}05}
\RouteRow{\(R_{\mathrm{act},\Omega}\)}{1.000000\allowbreak{}03855097\allowbreak{}23541416\allowbreak{}45}
\RouteGroup{Separación de prefijos}
\RouteRow{Área de ambigüedad; área \(\log_2\)}{54; 28.52932\allowbreak{}50129808}
\RouteRow{Primer bloque de separación}{órbita: \textemdash{}; celda: \textemdash{}; ruta: \textemdash{}}
\RouteRow{Ambigüedad final}{3}
\RouteGroup{Secuencias completas}
\RouteRow{\(w_6\)}{333 255 252 333 258 414 420 258 414 333 255 252 333 258 414 420 258 414}
\RouteRow{Tríadas estables por bloque}{0 1 2 3 4 5 6 7 8 9 10 11 12 13 14 15 16 17}
\end{tabularx}\endgroup\par\vspace{5pt}\noindent{\fontsize{8.4}{10}\selectfont\hyperref[trace-251]{Procedencia y códigos originales}\hfill Ficha completa · 251/324}
\clearpage\phantomsection\label{nav:vi-ruta-252}
% VI-CATALOG-ROW rutas 252
\pdfbookmark[3]{r7c9\_h+1v+1}{route-252}
\RouteTitle{Ruta 252}{r7c9\_\allowbreak{}h+1v+1}
\noindent Celda 63\quad (7, 9)\qquad Orientación \(h=1,\ v=1\)\hfill\hyperref[sec:vi-indice-rutas]{Índice}\par\vspace{4pt}
\begingroup\fontsize{10.2}{12.5}\selectfont\setlength{\tabcolsep}{5pt}\renewcommand{\arraystretch}{1.1}\rowcolors{1}{routepale}{white}\noindent\begin{tabularx}{\linewidth}{@{}>{\raggedright\arraybackslash}p{.345\linewidth}>{\raggedright\arraybackslash}X@{}}
\RouteGroup{Identidad estructural}
\RouteRow{Clase y multisección}{charged\_\allowbreak{}lepton\_\allowbreak{}G4\quad /\quad charged\_\allowbreak{}lepton\_\allowbreak{}G4}
\RouteRow{Colección y sector}{leptón cargado; leptón}
\RouteRow{Clasificación física}{leptón cargado}
\RouteRow{Generación, profundidad y color}{4; G4; \textemdash{}}
\RouteRow{Ruta primaria y órbita de inversión}{\(B_1\): sí\qquad 3,\allowbreak{}1;\allowbreak{}7,\allowbreak{}9}
\RouteRow{Firma de clasificación}{28|\allowbreak{}-6|\allowbreak{}6|\allowbreak{}0|\allowbreak{}-12|\allowbreak{}+++-\kern0pt{}-\kern0pt{}-+++-\kern0pt{}-\kern0pt{}-}
\RouteGroup{Retornos, memoria y transporte}
\RouteRow{Retornos con fase}{clase: 9; órbita: 27; celda: 27}
\RouteRow{Retorno cinemático}{en 108: 54; extensión: 216}
\RouteRow{Giros y cambios de hoja}{71; 107}
\RouteRow{Acarreo y diferimiento}{deuda total: 2; final: 1; bloques diferidos: 2}
\RouteRow{Tríadas finales; persistencia de Pareto}{16; no}
\RouteRow{\(K\longrightarrow K+9\)}{repeticiones: 9; cambios de memoria: 9}
\RouteGroup{Firma, fase y diferencias}
\RouteRow{\(n_A,\ s_C,\ \kappa_0,\ \nu_{120},\ \nu_{270}\)}{28, -6, 0, -4, -12}
\RouteRow{\(\tau_{12}\); firma histórica \(\mathbb Z^5\)}{+++-\kern0pt{}-\kern0pt{}-+++-\kern0pt{}-\kern0pt{}-; 28|\allowbreak{}-6|\allowbreak{}0|\allowbreak{}-4|\allowbreak{}-12}
\RouteRow{\(D_1,\ D_3,\ D_4,\ D_{27},\ D_{36}\)}{66, 60, 64, 48, 32}
\RouteRow{\(\Omega_{90/120},\ P_6\)}{80, 6}
\RouteGroup{Lectores y factores de acción}
\RouteRow{\(K_D\quad /\quad K_\Omega\)}{(80,\allowbreak{}66,\allowbreak{}2)\quad /\quad (80,\allowbreak{}54,\allowbreak{}6)}
\RouteRow{\(\kappa_D\)}{79.51859\allowbreak{}71232726\allowbreak{}53414287\allowbreak{}29}
\RouteRow{\(\kappa_\Omega\)}{79.60661\allowbreak{}01397948\allowbreak{}27586470\allowbreak{}91}
\RouteRow{\(R_{\mathrm{act},D}\)}{1.000000\allowbreak{}05512879\allowbreak{}47092488\allowbreak{}80}
\RouteRow{\(R_{\mathrm{act},\Omega}\)}{1.000000\allowbreak{}05518981\allowbreak{}25318591\allowbreak{}40}
\RouteGroup{Separación de prefijos}
\RouteRow{Área de ambigüedad; área \(\log_2\)}{54; 28.52932\allowbreak{}50129808}
\RouteRow{Primer bloque de separación}{órbita: \textemdash{}; celda: \textemdash{}; ruta: \textemdash{}}
\RouteRow{Ambigüedad final}{3}
\RouteGroup{Secuencias completas}
\RouteRow{\(w_6\)}{327 498 99 327 501 486 423 15 486 327 498 99 327 501 486 423 15 486}
\RouteRow{Tríadas estables por bloque}{0 1 2 3 4 5 6 7 7 9 10 11 12 13 14 15 16 16}
\end{tabularx}\endgroup\par\vspace{5pt}\noindent{\fontsize{8.4}{10}\selectfont\hyperref[trace-252]{Procedencia y códigos originales}\hfill Ficha completa · 252/324}
\clearpage\phantomsection\label{nav:vi-ruta-253}
% VI-CATALOG-ROW rutas 253
\pdfbookmark[2]{Celda 64}{cell-64}
\pdfbookmark[3]{r8c1\_h-1v-1}{route-253}
\RouteTitle{Ruta 253}{r8c1\_\allowbreak{}h-1v-1}
\noindent Celda 64\quad (8, 1)\qquad Orientación \(h=-1,\ v=-1\)\hfill\hyperref[sec:vi-indice-rutas]{Índice}\par\vspace{4pt}
\begingroup\fontsize{10.2}{12.5}\selectfont\setlength{\tabcolsep}{5pt}\renewcommand{\arraystretch}{1.1}\rowcolors{1}{routepale}{white}\noindent\begin{tabularx}{\linewidth}{@{}>{\raggedright\arraybackslash}p{.345\linewidth}>{\raggedright\arraybackslash}X@{}}
\RouteGroup{Identidad estructural}
\RouteRow{Clase y multisección}{quark\_\allowbreak{}down\_\allowbreak{}G4\quad /\quad quark\_\allowbreak{}down\_\allowbreak{}G4}
\RouteRow{Colección y sector}{quark de tipo abajo; quark}
\RouteRow{Clasificación física}{quark de tipo abajo}
\RouteRow{Generación, profundidad y color}{4; G4; G}
\RouteRow{Ruta primaria y órbita de inversión}{\(B_1\): no\qquad 2,\allowbreak{}9;\allowbreak{}8,\allowbreak{}1}
\RouteRow{Firma de clasificación}{28|\allowbreak{}-6|\allowbreak{}10|\allowbreak{}0|\allowbreak{}-12|\allowbreak{}+++-\kern0pt{}-\kern0pt{}-+++-\kern0pt{}-\kern0pt{}-}
\RouteGroup{Retornos, memoria y transporte}
\RouteRow{Retornos con fase}{clase: 27; órbita: 27; celda: 27}
\RouteRow{Retorno cinemático}{en 108: 54; extensión: 216}
\RouteRow{Giros y cambios de hoja}{71; 107}
\RouteRow{Acarreo y diferimiento}{deuda total: 1; final: 1; bloques diferidos: 1}
\RouteRow{Tríadas finales; persistencia de Pareto}{16; sí}
\RouteRow{\(K\longrightarrow K+9\)}{repeticiones: 9; cambios de memoria: 9}
\RouteGroup{Firma, fase y diferencias}
\RouteRow{\(n_A,\ s_C,\ \kappa_0,\ \nu_{120},\ \nu_{270}\)}{28, 6, 0, 0, -12}
\RouteRow{\(\tau_{12}\); firma histórica \(\mathbb Z^5\)}{+++-\kern0pt{}-\kern0pt{}-+++-\kern0pt{}-\kern0pt{}-; 28|\allowbreak{}6|\allowbreak{}0|\allowbreak{}0|\allowbreak{}-12}
\RouteRow{\(D_1,\ D_3,\ D_4,\ D_{27},\ D_{36}\)}{66, 64, 68, 60, 40}
\RouteRow{\(\Omega_{90/120},\ P_6\)}{80, 6}
\RouteGroup{Lectores y factores de acción}
\RouteRow{\(K_D\quad /\quad K_\Omega\)}{(100,\allowbreak{}66,\allowbreak{}2)\quad /\quad (80,\allowbreak{}54,\allowbreak{}6)}
\RouteRow{\(\kappa_D\)}{99.51859\allowbreak{}71232726\allowbreak{}53414287\allowbreak{}29}
\RouteRow{\(\kappa_\Omega\)}{79.60661\allowbreak{}01397948\allowbreak{}27586470\allowbreak{}91}
\RouteRow{\(R_{\mathrm{act},D}\)}{1.000000\allowbreak{}06899443\allowbreak{}08259364\allowbreak{}17}
\RouteRow{\(R_{\mathrm{act},\Omega}\)}{1.000000\allowbreak{}05518981\allowbreak{}25318591\allowbreak{}40}
\RouteGroup{Separación de prefijos}
\RouteRow{Área de ambigüedad; área \(\log_2\)}{40; 20.33985\allowbreak{}00028846\allowbreak{}24}
\RouteRow{Primer bloque de separación}{órbita: \textemdash{}; celda: \textemdash{}; ruta: \textemdash{}}
\RouteRow{Ambigüedad final}{2}
\RouteGroup{Secuencias completas}
\RouteRow{\(w_6\)}{15 486 423 15 489 498 99 327 498 15 486 423 15 489 498 99 327 498}
\RouteRow{Tríadas estables por bloque}{0 1 2 3 4 5 6 7 8 9 10 11 12 13 14 15 16 16}
\end{tabularx}\endgroup\par\vspace{5pt}\noindent{\fontsize{8.4}{10}\selectfont\hyperref[trace-253]{Procedencia y códigos originales}\hfill Ficha completa · 253/324}
\clearpage\phantomsection\label{nav:vi-ruta-254}
% VI-CATALOG-ROW rutas 254
\pdfbookmark[3]{r8c1\_h-1v+1}{route-254}
\RouteTitle{Ruta 254}{r8c1\_\allowbreak{}h-1v+1}
\noindent Celda 64\quad (8, 1)\qquad Orientación \(h=-1,\ v=1\)\hfill\hyperref[sec:vi-indice-rutas]{Índice}\par\vspace{4pt}
\begingroup\fontsize{10.2}{12.5}\selectfont\setlength{\tabcolsep}{5pt}\renewcommand{\arraystretch}{1.1}\rowcolors{1}{routepale}{white}\noindent\begin{tabularx}{\linewidth}{@{}>{\raggedright\arraybackslash}p{.345\linewidth}>{\raggedright\arraybackslash}X@{}}
\RouteGroup{Identidad estructural}
\RouteRow{Clase y multisección}{quark\_\allowbreak{}down\_\allowbreak{}G4\quad /\quad quark\_\allowbreak{}down\_\allowbreak{}G4}
\RouteRow{Colección y sector}{quark de tipo abajo; quark}
\RouteRow{Clasificación física}{quark de tipo abajo}
\RouteRow{Generación, profundidad y color}{4; G4; G}
\RouteRow{Ruta primaria y órbita de inversión}{\(B_1\): no\qquad 2,\allowbreak{}9;\allowbreak{}8,\allowbreak{}1}
\RouteRow{Firma de clasificación}{28|\allowbreak{}-6|\allowbreak{}10|\allowbreak{}0|\allowbreak{}-12|\allowbreak{}+++-\kern0pt{}-\kern0pt{}-+++-\kern0pt{}-\kern0pt{}-}
\RouteGroup{Retornos, memoria y transporte}
\RouteRow{Retornos con fase}{clase: 27; órbita: 27; celda: 27}
\RouteRow{Retorno cinemático}{en 108: 54; extensión: 216}
\RouteRow{Giros y cambios de hoja}{71; 107}
\RouteRow{Acarreo y diferimiento}{deuda total: 0; final: 0; bloques diferidos: 0}
\RouteRow{Tríadas finales; persistencia de Pareto}{17; sí}
\RouteRow{\(K\longrightarrow K+9\)}{repeticiones: 9; cambios de memoria: 9}
\RouteGroup{Firma, fase y diferencias}
\RouteRow{\(n_A,\ s_C,\ \kappa_0,\ \nu_{120},\ \nu_{270}\)}{20, 0, 0, 0, -4}
\RouteRow{\(\tau_{12}\); firma histórica \(\mathbb Z^5\)}{0++-\kern0pt{}-00++-\kern0pt{}-0; 20|\allowbreak{}0|\allowbreak{}0|\allowbreak{}0|\allowbreak{}-4}
\RouteRow{\(D_1,\ D_3,\ D_4,\ D_{27},\ D_{36}\)}{24, 20, 24, 24, 16}
\RouteRow{\(\Omega_{90/120},\ P_6\)}{36, 4}
\RouteGroup{Lectores y factores de acción}
\RouteRow{\(K_D\quad /\quad K_\Omega\)}{(40,\allowbreak{}24,\allowbreak{}2)\quad /\quad (36,\allowbreak{}54,\allowbreak{}4)}
\RouteRow{\(\kappa_D\)}{39.82508\allowbreak{}59311825\allowbreak{}73056173\allowbreak{}26}
\RouteRow{\(\kappa_\Omega\)}{35.60638\allowbreak{}77469494\allowbreak{}43306362\allowbreak{}81}
\RouteRow{\(R_{\mathrm{act},D}\)}{1.000000\allowbreak{}02761000\allowbreak{}61595142\allowbreak{}15}
\RouteRow{\(R_{\mathrm{act},\Omega}\)}{1.000000\allowbreak{}02468525\allowbreak{}95691181\allowbreak{}52}
\RouteGroup{Separación de prefijos}
\RouteRow{Área de ambigüedad; área \(\log_2\)}{36; 18.0}
\RouteRow{Primer bloque de separación}{órbita: \textemdash{}; celda: \textemdash{}; ruta: \textemdash{}}
\RouteRow{Ambigüedad final}{2}
\RouteGroup{Secuencias completas}
\RouteRow{\(w_6\)}{0 81 3 0 84 0 81 3 0 0 81 3 0 84 0 81 3 0}
\RouteRow{Tríadas estables por bloque}{0 1 2 3 4 5 6 7 8 9 10 11 12 13 14 15 16 17}
\end{tabularx}\endgroup\par\vspace{5pt}\noindent{\fontsize{8.4}{10}\selectfont\hyperref[trace-254]{Procedencia y códigos originales}\hfill Ficha completa · 254/324}
\clearpage\phantomsection\label{nav:vi-ruta-255}
% VI-CATALOG-ROW rutas 255
\pdfbookmark[3]{r8c1\_h+1v-1}{route-255}
\RouteTitle{Ruta 255}{r8c1\_\allowbreak{}h+1v-1}
\noindent Celda 64\quad (8, 1)\qquad Orientación \(h=1,\ v=-1\)\hfill\hyperref[sec:vi-indice-rutas]{Índice}\par\vspace{4pt}
\begingroup\fontsize{10.2}{12.5}\selectfont\setlength{\tabcolsep}{5pt}\renewcommand{\arraystretch}{1.1}\rowcolors{1}{routepale}{white}\noindent\begin{tabularx}{\linewidth}{@{}>{\raggedright\arraybackslash}p{.345\linewidth}>{\raggedright\arraybackslash}X@{}}
\RouteGroup{Identidad estructural}
\RouteRow{Clase y multisección}{quark\_\allowbreak{}down\_\allowbreak{}G4\quad /\quad quark\_\allowbreak{}down\_\allowbreak{}G4}
\RouteRow{Colección y sector}{quark de tipo abajo; quark}
\RouteRow{Clasificación física}{quark de tipo abajo}
\RouteRow{Generación, profundidad y color}{4; G4; G}
\RouteRow{Ruta primaria y órbita de inversión}{\(B_1\): sí\qquad 2,\allowbreak{}9;\allowbreak{}8,\allowbreak{}1}
\RouteRow{Firma de clasificación}{28|\allowbreak{}-6|\allowbreak{}10|\allowbreak{}0|\allowbreak{}-12|\allowbreak{}+++-\kern0pt{}-\kern0pt{}-+++-\kern0pt{}-\kern0pt{}-}
\RouteGroup{Retornos, memoria y transporte}
\RouteRow{Retornos con fase}{clase: 27; órbita: 27; celda: 27}
\RouteRow{Retorno cinemático}{en 108: 54; extensión: 216}
\RouteRow{Giros y cambios de hoja}{71; 107}
\RouteRow{Acarreo y diferimiento}{deuda total: 2; final: 0; bloques diferidos: 2}
\RouteRow{Tríadas finales; persistencia de Pareto}{17; sí}
\RouteRow{\(K\longrightarrow K+9\)}{repeticiones: 9; cambios de memoria: 9}
\RouteGroup{Firma, fase y diferencias}
\RouteRow{\(n_A,\ s_C,\ \kappa_0,\ \nu_{120},\ \nu_{270}\)}{20, 0, 0, 0, -4}
\RouteRow{\(\tau_{12}\); firma histórica \(\mathbb Z^5\)}{++00-\kern0pt{}-++00-\kern0pt{}-; 20|\allowbreak{}0|\allowbreak{}0|\allowbreak{}0|\allowbreak{}-4}
\RouteRow{\(D_1,\ D_3,\ D_4,\ D_{27},\ D_{36}\)}{24, 20, 24, 24, 16}
\RouteRow{\(\Omega_{90/120},\ P_6\)}{36, 4}
\RouteGroup{Lectores y factores de acción}
\RouteRow{\(K_D\quad /\quad K_\Omega\)}{(40,\allowbreak{}24,\allowbreak{}2)\quad /\quad (36,\allowbreak{}54,\allowbreak{}4)}
\RouteRow{\(\kappa_D\)}{39.82508\allowbreak{}59311825\allowbreak{}73056173\allowbreak{}26}
\RouteRow{\(\kappa_\Omega\)}{35.60638\allowbreak{}77469494\allowbreak{}43306362\allowbreak{}81}
\RouteRow{\(R_{\mathrm{act},D}\)}{1.000000\allowbreak{}02761000\allowbreak{}61595142\allowbreak{}15}
\RouteRow{\(R_{\mathrm{act},\Omega}\)}{1.000000\allowbreak{}02468525\allowbreak{}95691181\allowbreak{}52}
\RouteGroup{Separación de prefijos}
\RouteRow{Área de ambigüedad; área \(\log_2\)}{36; 18.0}
\RouteRow{Primer bloque de separación}{órbita: \textemdash{}; celda: \textemdash{}; ruta: \textemdash{}}
\RouteRow{Ambigüedad final}{2}
\RouteGroup{Secuencias completas}
\RouteRow{\(w_6\)}{81 3 0 81 0 3 0 81 3 81 3 0 81 0 3 0 81 3}
\RouteRow{Tríadas estables por bloque}{0 1 2 3 4 5 6 6 8 9 10 11 12 13 14 15 15 17}
\end{tabularx}\endgroup\par\vspace{5pt}\noindent{\fontsize{8.4}{10}\selectfont\hyperref[trace-255]{Procedencia y códigos originales}\hfill Ficha completa · 255/324}
\clearpage\phantomsection\label{nav:vi-ruta-256}
% VI-CATALOG-ROW rutas 256
\pdfbookmark[3]{r8c1\_h+1v+1}{route-256}
\RouteTitle{Ruta 256}{r8c1\_\allowbreak{}h+1v+1}
\noindent Celda 64\quad (8, 1)\qquad Orientación \(h=1,\ v=1\)\hfill\hyperref[sec:vi-indice-rutas]{Índice}\par\vspace{4pt}
\begingroup\fontsize{10.2}{12.5}\selectfont\setlength{\tabcolsep}{5pt}\renewcommand{\arraystretch}{1.1}\rowcolors{1}{routepale}{white}\noindent\begin{tabularx}{\linewidth}{@{}>{\raggedright\arraybackslash}p{.345\linewidth}>{\raggedright\arraybackslash}X@{}}
\RouteGroup{Identidad estructural}
\RouteRow{Clase y multisección}{quark\_\allowbreak{}down\_\allowbreak{}G4\quad /\quad quark\_\allowbreak{}down\_\allowbreak{}G4}
\RouteRow{Colección y sector}{quark de tipo abajo; quark}
\RouteRow{Clasificación física}{quark de tipo abajo}
\RouteRow{Generación, profundidad y color}{4; G4; G}
\RouteRow{Ruta primaria y órbita de inversión}{\(B_1\): no\qquad 2,\allowbreak{}9;\allowbreak{}8,\allowbreak{}1}
\RouteRow{Firma de clasificación}{28|\allowbreak{}-6|\allowbreak{}10|\allowbreak{}0|\allowbreak{}-12|\allowbreak{}+++-\kern0pt{}-\kern0pt{}-+++-\kern0pt{}-\kern0pt{}-}
\RouteGroup{Retornos, memoria y transporte}
\RouteRow{Retornos con fase}{clase: 27; órbita: 27; celda: 27}
\RouteRow{Retorno cinemático}{en 108: 54; extensión: 216}
\RouteRow{Giros y cambios de hoja}{71; 107}
\RouteRow{Acarreo y diferimiento}{deuda total: 1; final: 0; bloques diferidos: 1}
\RouteRow{Tríadas finales; persistencia de Pareto}{17; sí}
\RouteRow{\(K\longrightarrow K+9\)}{repeticiones: 9; cambios de memoria: 9}
\RouteGroup{Firma, fase y diferencias}
\RouteRow{\(n_A,\ s_C,\ \kappa_0,\ \nu_{120},\ \nu_{270}\)}{28, -6, 0, 0, -12}
\RouteRow{\(\tau_{12}\); firma histórica \(\mathbb Z^5\)}{+++-\kern0pt{}-\kern0pt{}-+++-\kern0pt{}-\kern0pt{}-; 28|\allowbreak{}-6|\allowbreak{}0|\allowbreak{}0|\allowbreak{}-12}
\RouteRow{\(D_1,\ D_3,\ D_4,\ D_{27},\ D_{36}\)}{66, 64, 68, 60, 40}
\RouteRow{\(\Omega_{90/120},\ P_6\)}{80, 6}
\RouteGroup{Lectores y factores de acción}
\RouteRow{\(K_D\quad /\quad K_\Omega\)}{(100,\allowbreak{}66,\allowbreak{}2)\quad /\quad (80,\allowbreak{}54,\allowbreak{}6)}
\RouteRow{\(\kappa_D\)}{99.51859\allowbreak{}71232726\allowbreak{}53414287\allowbreak{}29}
\RouteRow{\(\kappa_\Omega\)}{79.60661\allowbreak{}01397948\allowbreak{}27586470\allowbreak{}91}
\RouteRow{\(R_{\mathrm{act},D}\)}{1.000000\allowbreak{}06899443\allowbreak{}08259364\allowbreak{}17}
\RouteRow{\(R_{\mathrm{act},\Omega}\)}{1.000000\allowbreak{}05518981\allowbreak{}25318591\allowbreak{}40}
\RouteGroup{Separación de prefijos}
\RouteRow{Área de ambigüedad; área \(\log_2\)}{36; 18.0}
\RouteRow{Primer bloque de separación}{órbita: \textemdash{}; celda: \textemdash{}; ruta: \textemdash{}}
\RouteRow{Ambigüedad final}{2}
\RouteGroup{Secuencias completas}
\RouteRow{\(w_6\)}{99 327 498 99 324 423 15 486 423 99 327 498 99 324 423 15 486 423}
\RouteRow{Tríadas estables por bloque}{0 1 2 3 4 5 6 7 8 9 10 11 12 13 14 14 16 17}
\end{tabularx}\endgroup\par\vspace{5pt}\noindent{\fontsize{8.4}{10}\selectfont\hyperref[trace-256]{Procedencia y códigos originales}\hfill Ficha completa · 256/324}
\clearpage\phantomsection\label{nav:vi-ruta-257}
% VI-CATALOG-ROW rutas 257
\pdfbookmark[2]{Celda 65}{cell-65}
\pdfbookmark[3]{r8c2\_h-1v-1}{route-257}
\RouteTitle{Ruta 257}{r8c2\_\allowbreak{}h-1v-1}
\noindent Celda 65\quad (8, 2)\qquad Orientación \(h=-1,\ v=-1\)\hfill\hyperref[sec:vi-indice-rutas]{Índice}\par\vspace{4pt}
\begingroup\fontsize{10.2}{12.5}\selectfont\setlength{\tabcolsep}{5pt}\renewcommand{\arraystretch}{1.1}\rowcolors{1}{routepale}{white}\noindent\begin{tabularx}{\linewidth}{@{}>{\raggedright\arraybackslash}p{.345\linewidth}>{\raggedright\arraybackslash}X@{}}
\RouteGroup{Identidad estructural}
\RouteRow{Clase y multisección}{quark\_\allowbreak{}up\_\allowbreak{}G3\quad /\quad quark\_\allowbreak{}up\_\allowbreak{}G3}
\RouteRow{Colección y sector}{quark de tipo arriba; quark}
\RouteRow{Clasificación física}{quark de tipo arriba}
\RouteRow{Generación, profundidad y color}{3; G3; R}
\RouteRow{Ruta primaria y órbita de inversión}{\(B_1\): no\qquad 2,\allowbreak{}8;\allowbreak{}8,\allowbreak{}2}
\RouteRow{Firma de clasificación}{30|\allowbreak{}-6|\allowbreak{}4|\allowbreak{}0|\allowbreak{}-12|\allowbreak{}+++-\kern0pt{}-\kern0pt{}-+++-\kern0pt{}-\kern0pt{}-}
\RouteGroup{Retornos, memoria y transporte}
\RouteRow{Retornos con fase}{clase: 27; órbita: 27; celda: 27}
\RouteRow{Retorno cinemático}{en 108: 54; extensión: 216}
\RouteRow{Giros y cambios de hoja}{71; 107}
\RouteRow{Acarreo y diferimiento}{deuda total: 3; final: 0; bloques diferidos: 2}
\RouteRow{Tríadas finales; persistencia de Pareto}{17; sí}
\RouteRow{\(K\longrightarrow K+9\)}{repeticiones: 9; cambios de memoria: 9}
\RouteGroup{Firma, fase y diferencias}
\RouteRow{\(n_A,\ s_C,\ \kappa_0,\ \nu_{120},\ \nu_{270}\)}{30, 6, 0, 2, -12}
\RouteRow{\(\tau_{12}\); firma histórica \(\mathbb Z^5\)}{+++-\kern0pt{}-\kern0pt{}-+++-\kern0pt{}-\kern0pt{}-; 30|\allowbreak{}6|\allowbreak{}0|\allowbreak{}2|\allowbreak{}-12}
\RouteRow{\(D_1,\ D_3,\ D_4,\ D_{27},\ D_{36}\)}{54, 56, 68, 48, 32}
\RouteRow{\(\Omega_{90/120},\ P_6\)}{80, 6}
\RouteGroup{Lectores y factores de acción}
\RouteRow{\(K_D\quad /\quad K_\Omega\)}{(80,\allowbreak{}54,\allowbreak{}6)\quad /\quad (80,\allowbreak{}54,\allowbreak{}6)}
\RouteRow{\(\kappa_D\)}{79.60661\allowbreak{}01397948\allowbreak{}27586470\allowbreak{}91}
\RouteRow{\(\kappa_\Omega\)}{79.60661\allowbreak{}01397948\allowbreak{}27586470\allowbreak{}91}
\RouteRow{\(R_{\mathrm{act},D}\)}{1.000000\allowbreak{}05518981\allowbreak{}25318591\allowbreak{}40}
\RouteRow{\(R_{\mathrm{act},\Omega}\)}{1.000000\allowbreak{}05518981\allowbreak{}25318591\allowbreak{}40}
\RouteGroup{Separación de prefijos}
\RouteRow{Área de ambigüedad; área \(\log_2\)}{72; 36.0}
\RouteRow{Primer bloque de separación}{órbita: \textemdash{}; celda: \textemdash{}; ruta: \textemdash{}}
\RouteRow{Ambigüedad final}{4}
\RouteGroup{Secuencias completas}
\RouteRow{\(w_6\)}{423 15 486 423 9 24 243 657 24 423 15 486 423 9 24 243 657 24}
\RouteRow{Tríadas estables por bloque}{0 1 2 3 4 5 6 7 8 9 10 11 11 13 14 14 15 17}
\end{tabularx}\endgroup\par\vspace{5pt}\noindent{\fontsize{8.4}{10}\selectfont\hyperref[trace-257]{Procedencia y códigos originales}\hfill Ficha completa · 257/324}
\clearpage\phantomsection\label{nav:vi-ruta-258}
% VI-CATALOG-ROW rutas 258
\pdfbookmark[3]{r8c2\_h-1v+1}{route-258}
\RouteTitle{Ruta 258}{r8c2\_\allowbreak{}h-1v+1}
\noindent Celda 65\quad (8, 2)\qquad Orientación \(h=-1,\ v=1\)\hfill\hyperref[sec:vi-indice-rutas]{Índice}\par\vspace{4pt}
\begingroup\fontsize{10.2}{12.5}\selectfont\setlength{\tabcolsep}{5pt}\renewcommand{\arraystretch}{1.1}\rowcolors{1}{routepale}{white}\noindent\begin{tabularx}{\linewidth}{@{}>{\raggedright\arraybackslash}p{.345\linewidth}>{\raggedright\arraybackslash}X@{}}
\RouteGroup{Identidad estructural}
\RouteRow{Clase y multisección}{quark\_\allowbreak{}up\_\allowbreak{}G3\quad /\quad quark\_\allowbreak{}up\_\allowbreak{}G3}
\RouteRow{Colección y sector}{quark de tipo arriba; quark}
\RouteRow{Clasificación física}{quark de tipo arriba}
\RouteRow{Generación, profundidad y color}{3; G3; R}
\RouteRow{Ruta primaria y órbita de inversión}{\(B_1\): no\qquad 2,\allowbreak{}8;\allowbreak{}8,\allowbreak{}2}
\RouteRow{Firma de clasificación}{30|\allowbreak{}-6|\allowbreak{}4|\allowbreak{}0|\allowbreak{}-12|\allowbreak{}+++-\kern0pt{}-\kern0pt{}-+++-\kern0pt{}-\kern0pt{}-}
\RouteGroup{Retornos, memoria y transporte}
\RouteRow{Retornos con fase}{clase: 9; órbita: 9; celda: 27}
\RouteRow{Retorno cinemático}{en 108: 54; extensión: 216}
\RouteRow{Giros y cambios de hoja}{71; 107}
\RouteRow{Acarreo y diferimiento}{deuda total: 2; final: 1; bloques diferidos: 2}
\RouteRow{Tríadas finales; persistencia de Pareto}{16; sí}
\RouteRow{\(K\longrightarrow K+9\)}{repeticiones: 9; cambios de memoria: 9}
\RouteGroup{Firma, fase y diferencias}
\RouteRow{\(n_A,\ s_C,\ \kappa_0,\ \nu_{120},\ \nu_{270}\)}{50, 4, -2, 2, -8}
\RouteRow{\(\tau_{12}\); firma histórica \(\mathbb Z^5\)}{+++-\kern0pt{}-0+++-\kern0pt{}-0; 50|\allowbreak{}4|\allowbreak{}-2|\allowbreak{}2|\allowbreak{}-8}
\RouteRow{\(D_1,\ D_3,\ D_4,\ D_{27},\ D_{36}\)}{84, 28, 84, 36, 24}
\RouteRow{\(\Omega_{90/120},\ P_6\)}{56, 5}
\RouteGroup{Lectores y factores de acción}
\RouteRow{\(K_D\quad /\quad K_\Omega\)}{(60,\allowbreak{}84,\allowbreak{}28)\quad /\quad (56,\allowbreak{}54,\allowbreak{}5)}
\RouteRow{\(\kappa_D\)}{59.39013\allowbreak{}58840155\allowbreak{}40637741\allowbreak{}48}
\RouteRow{\(\kappa_\Omega\)}{55.60649\allowbreak{}89433721\allowbreak{}35446416\allowbreak{}86}
\RouteRow{\(R_{\mathrm{act},D}\)}{1.000000\allowbreak{}04117409\allowbreak{}89467415\allowbreak{}75}
\RouteRow{\(R_{\mathrm{act},\Omega}\)}{1.000000\allowbreak{}03855097\allowbreak{}23541416\allowbreak{}45}
\RouteGroup{Separación de prefijos}
\RouteRow{Área de ambigüedad; área \(\log_2\)}{72; 36.0}
\RouteRow{Primer bloque de separación}{órbita: \textemdash{}; celda: \textemdash{}; ruta: \textemdash{}}
\RouteRow{Ambigüedad final}{4}
\RouteGroup{Secuencias completas}
\RouteRow{\(w_6\)}{414 420 258 414 414 333 255 252 333 414 420 258 414 414 333 255 252 333}
\RouteRow{Tríadas estables por bloque}{0 1 2 3 4 5 6 7 8 9 10 11 12 13 14 14 16 16}
\end{tabularx}\endgroup\par\vspace{5pt}\noindent{\fontsize{8.4}{10}\selectfont\hyperref[trace-258]{Procedencia y códigos originales}\hfill Ficha completa · 258/324}
\clearpage\phantomsection\label{nav:vi-ruta-259}
% VI-CATALOG-ROW rutas 259
\pdfbookmark[3]{r8c2\_h+1v-1}{route-259}
\RouteTitle{Ruta 259}{r8c2\_\allowbreak{}h+1v-1}
\noindent Celda 65\quad (8, 2)\qquad Orientación \(h=1,\ v=-1\)\hfill\hyperref[sec:vi-indice-rutas]{Índice}\par\vspace{4pt}
\begingroup\fontsize{10.2}{12.5}\selectfont\setlength{\tabcolsep}{5pt}\renewcommand{\arraystretch}{1.1}\rowcolors{1}{routepale}{white}\noindent\begin{tabularx}{\linewidth}{@{}>{\raggedright\arraybackslash}p{.345\linewidth}>{\raggedright\arraybackslash}X@{}}
\RouteGroup{Identidad estructural}
\RouteRow{Clase y multisección}{quark\_\allowbreak{}up\_\allowbreak{}G3\quad /\quad quark\_\allowbreak{}up\_\allowbreak{}G3}
\RouteRow{Colección y sector}{quark de tipo arriba; quark}
\RouteRow{Clasificación física}{quark de tipo arriba}
\RouteRow{Generación, profundidad y color}{3; G3; R}
\RouteRow{Ruta primaria y órbita de inversión}{\(B_1\): no\qquad 2,\allowbreak{}8;\allowbreak{}8,\allowbreak{}2}
\RouteRow{Firma de clasificación}{30|\allowbreak{}-6|\allowbreak{}4|\allowbreak{}0|\allowbreak{}-12|\allowbreak{}+++-\kern0pt{}-\kern0pt{}-+++-\kern0pt{}-\kern0pt{}-}
\RouteGroup{Retornos, memoria y transporte}
\RouteRow{Retornos con fase}{clase: 18; órbita: 18; celda: 27}
\RouteRow{Retorno cinemático}{en 108: 54; extensión: 216}
\RouteRow{Giros y cambios de hoja}{71; 107}
\RouteRow{Acarreo y diferimiento}{deuda total: 3; final: 0; bloques diferidos: 2}
\RouteRow{Tríadas finales; persistencia de Pareto}{17; sí}
\RouteRow{\(K\longrightarrow K+9\)}{repeticiones: 9; cambios de memoria: 9}
\RouteGroup{Firma, fase y diferencias}
\RouteRow{\(n_A,\ s_C,\ \kappa_0,\ \nu_{120},\ \nu_{270}\)}{50, -4, 2, -2, -8}
\RouteRow{\(\tau_{12}\); firma histórica \(\mathbb Z^5\)}{++0-\kern0pt{}-\kern0pt{}-++0-\kern0pt{}-\kern0pt{}-; 50|\allowbreak{}-4|\allowbreak{}2|\allowbreak{}-2|\allowbreak{}-8}
\RouteRow{\(D_1,\ D_3,\ D_4,\ D_{27},\ D_{36}\)}{84, 28, 84, 36, 24}
\RouteRow{\(\Omega_{90/120},\ P_6\)}{56, 5}
\RouteGroup{Lectores y factores de acción}
\RouteRow{\(K_D\quad /\quad K_\Omega\)}{(60,\allowbreak{}84,\allowbreak{}28)\quad /\quad (56,\allowbreak{}54,\allowbreak{}5)}
\RouteRow{\(\kappa_D\)}{59.39013\allowbreak{}58840155\allowbreak{}40637741\allowbreak{}48}
\RouteRow{\(\kappa_\Omega\)}{55.60649\allowbreak{}89433721\allowbreak{}35446416\allowbreak{}86}
\RouteRow{\(R_{\mathrm{act},D}\)}{1.000000\allowbreak{}04117409\allowbreak{}89467415\allowbreak{}75}
\RouteRow{\(R_{\mathrm{act},\Omega}\)}{1.000000\allowbreak{}03855097\allowbreak{}23541416\allowbreak{}45}
\RouteGroup{Separación de prefijos}
\RouteRow{Área de ambigüedad; área \(\log_2\)}{72; 36.0}
\RouteRow{Primer bloque de separación}{órbita: \textemdash{}; celda: \textemdash{}; ruta: \textemdash{}}
\RouteRow{Ambigüedad final}{4}
\RouteGroup{Secuencias completas}
\RouteRow{\(w_6\)}{255 252 333 255 258 258 414 420 258 255 252 333 255 258 258 414 420 258}
\RouteRow{Tríadas estables por bloque}{0 1 2 3 4 5 6 7 8 9 10 11 11 12 14 15 15 17}
\end{tabularx}\endgroup\par\vspace{5pt}\noindent{\fontsize{8.4}{10}\selectfont\hyperref[trace-259]{Procedencia y códigos originales}\hfill Ficha completa · 259/324}
\clearpage\phantomsection\label{nav:vi-ruta-260}
% VI-CATALOG-ROW rutas 260
\pdfbookmark[3]{r8c2\_h+1v+1}{route-260}
\RouteTitle{Ruta 260}{r8c2\_\allowbreak{}h+1v+1}
\noindent Celda 65\quad (8, 2)\qquad Orientación \(h=1,\ v=1\)\hfill\hyperref[sec:vi-indice-rutas]{Índice}\par\vspace{4pt}
\begingroup\fontsize{10.2}{12.5}\selectfont\setlength{\tabcolsep}{5pt}\renewcommand{\arraystretch}{1.1}\rowcolors{1}{routepale}{white}\noindent\begin{tabularx}{\linewidth}{@{}>{\raggedright\arraybackslash}p{.345\linewidth}>{\raggedright\arraybackslash}X@{}}
\RouteGroup{Identidad estructural}
\RouteRow{Clase y multisección}{quark\_\allowbreak{}up\_\allowbreak{}G3\quad /\quad quark\_\allowbreak{}up\_\allowbreak{}G3}
\RouteRow{Colección y sector}{quark de tipo arriba; quark}
\RouteRow{Clasificación física}{quark de tipo arriba}
\RouteRow{Generación, profundidad y color}{3; G3; R}
\RouteRow{Ruta primaria y órbita de inversión}{\(B_1\): sí\qquad 2,\allowbreak{}8;\allowbreak{}8,\allowbreak{}2}
\RouteRow{Firma de clasificación}{30|\allowbreak{}-6|\allowbreak{}4|\allowbreak{}0|\allowbreak{}-12|\allowbreak{}+++-\kern0pt{}-\kern0pt{}-+++-\kern0pt{}-\kern0pt{}-}
\RouteGroup{Retornos, memoria y transporte}
\RouteRow{Retornos con fase}{clase: 27; órbita: 27; celda: 27}
\RouteRow{Retorno cinemático}{en 108: 54; extensión: 216}
\RouteRow{Giros y cambios de hoja}{71; 107}
\RouteRow{Acarreo y diferimiento}{deuda total: 2; final: 1; bloques diferidos: 2}
\RouteRow{Tríadas finales; persistencia de Pareto}{16; sí}
\RouteRow{\(K\longrightarrow K+9\)}{repeticiones: 9; cambios de memoria: 9}
\RouteGroup{Firma, fase y diferencias}
\RouteRow{\(n_A,\ s_C,\ \kappa_0,\ \nu_{120},\ \nu_{270}\)}{30, -6, 0, -2, -12}
\RouteRow{\(\tau_{12}\); firma histórica \(\mathbb Z^5\)}{+++-\kern0pt{}-\kern0pt{}-+++-\kern0pt{}-\kern0pt{}-; 30|\allowbreak{}-6|\allowbreak{}0|\allowbreak{}-2|\allowbreak{}-12}
\RouteRow{\(D_1,\ D_3,\ D_4,\ D_{27},\ D_{36}\)}{54, 56, 68, 48, 32}
\RouteRow{\(\Omega_{90/120},\ P_6\)}{80, 6}
\RouteGroup{Lectores y factores de acción}
\RouteRow{\(K_D\quad /\quad K_\Omega\)}{(80,\allowbreak{}54,\allowbreak{}6)\quad /\quad (80,\allowbreak{}54,\allowbreak{}6)}
\RouteRow{\(\kappa_D\)}{79.60661\allowbreak{}01397948\allowbreak{}27586470\allowbreak{}91}
\RouteRow{\(\kappa_\Omega\)}{79.60661\allowbreak{}01397948\allowbreak{}27586470\allowbreak{}91}
\RouteRow{\(R_{\mathrm{act},D}\)}{1.000000\allowbreak{}05518981\allowbreak{}25318591\allowbreak{}40}
\RouteRow{\(R_{\mathrm{act},\Omega}\)}{1.000000\allowbreak{}05518981\allowbreak{}25318591\allowbreak{}40}
\RouteGroup{Separación de prefijos}
\RouteRow{Área de ambigüedad; área \(\log_2\)}{72; 36.0}
\RouteRow{Primer bloque de separación}{órbita: \textemdash{}; celda: \textemdash{}; ruta: \textemdash{}}
\RouteRow{Ambigüedad final}{4}
\RouteGroup{Secuencias completas}
\RouteRow{\(w_6\)}{243 657 24 243 663 486 423 15 486 243 657 24 243 663 486 423 15 486}
\RouteRow{Tríadas estables por bloque}{0 1 2 3 4 5 6 7 8 9 10 11 12 13 14 14 16 16}
\end{tabularx}\endgroup\par\vspace{5pt}\noindent{\fontsize{8.4}{10}\selectfont\hyperref[trace-260]{Procedencia y códigos originales}\hfill Ficha completa · 260/324}
\clearpage\phantomsection\label{nav:vi-ruta-261}
% VI-CATALOG-ROW rutas 261
\pdfbookmark[2]{Celda 66}{cell-66}
\pdfbookmark[3]{r8c3\_h-1v-1}{route-261}
\RouteTitle{Ruta 261}{r8c3\_\allowbreak{}h-1v-1}
\noindent Celda 66\quad (8, 3)\qquad Orientación \(h=-1,\ v=-1\)\hfill\hyperref[sec:vi-indice-rutas]{Índice}\par\vspace{4pt}
\begingroup\fontsize{10.2}{12.5}\selectfont\setlength{\tabcolsep}{5pt}\renewcommand{\arraystretch}{1.1}\rowcolors{1}{routepale}{white}\noindent\begin{tabularx}{\linewidth}{@{}>{\raggedright\arraybackslash}p{.345\linewidth}>{\raggedright\arraybackslash}X@{}}
\RouteGroup{Identidad estructural}
\RouteRow{Clase y multisección}{quark\_\allowbreak{}down\_\allowbreak{}G3\quad /\quad quark\_\allowbreak{}down\_\allowbreak{}G3}
\RouteRow{Colección y sector}{quark de tipo abajo; quark}
\RouteRow{Clasificación física}{quark de tipo abajo}
\RouteRow{Generación, profundidad y color}{3; G3; B}
\RouteRow{Ruta primaria y órbita de inversión}{\(B_1\): no\qquad 2,\allowbreak{}7;\allowbreak{}8,\allowbreak{}3}
\RouteRow{Firma de clasificación}{34|\allowbreak{}0|\allowbreak{}-8|\allowbreak{}-2|\allowbreak{}-8|\allowbreak{}++0-\kern0pt{}-\kern0pt{}-++0-\kern0pt{}-\kern0pt{}-}
\RouteGroup{Retornos, memoria y transporte}
\RouteRow{Retornos con fase}{clase: 27; órbita: 27; celda: 27}
\RouteRow{Retorno cinemático}{en 108: 54; extensión: 216}
\RouteRow{Giros y cambios de hoja}{71; 107}
\RouteRow{Acarreo y diferimiento}{deuda total: 1; final: 0; bloques diferidos: 1}
\RouteRow{Tríadas finales; persistencia de Pareto}{17; sí}
\RouteRow{\(K\longrightarrow K+9\)}{repeticiones: 9; cambios de memoria: 9}
\RouteGroup{Firma, fase y diferencias}
\RouteRow{\(n_A,\ s_C,\ \kappa_0,\ \nu_{120},\ \nu_{270}\)}{34, 0, -2, 0, -8}
\RouteRow{\(\tau_{12}\); firma histórica \(\mathbb Z^5\)}{+++-\kern0pt{}-0+++-\kern0pt{}-0; 34|\allowbreak{}0|\allowbreak{}-2|\allowbreak{}0|\allowbreak{}-8}
\RouteRow{\(D_1,\ D_3,\ D_4,\ D_{27},\ D_{36}\)}{66, 60, 68, 48, 32}
\RouteRow{\(\Omega_{90/120},\ P_6\)}{56, 5}
\RouteGroup{Lectores y factores de acción}
\RouteRow{\(K_D\quad /\quad K_\Omega\)}{(80,\allowbreak{}66,\allowbreak{}4)\quad /\quad (56,\allowbreak{}54,\allowbreak{}5)}
\RouteRow{\(\kappa_D\)}{79.51881\allowbreak{}95161180\allowbreak{}37694395\allowbreak{}39}
\RouteRow{\(\kappa_\Omega\)}{55.60649\allowbreak{}89433721\allowbreak{}35446416\allowbreak{}86}
\RouteRow{\(R_{\mathrm{act},D}\)}{1.000000\allowbreak{}05512894\allowbreak{}88901612\allowbreak{}64}
\RouteRow{\(R_{\mathrm{act},\Omega}\)}{1.000000\allowbreak{}03855097\allowbreak{}23541416\allowbreak{}45}
\RouteGroup{Separación de prefijos}
\RouteRow{Área de ambigüedad; área \(\log_2\)}{36; 18.0}
\RouteRow{Primer bloque de separación}{órbita: \textemdash{}; celda: \textemdash{}; ruta: \textemdash{}}
\RouteRow{Ambigüedad final}{2}
\RouteGroup{Secuencias completas}
\RouteRow{\(w_6\)}{570 264 90 570 267 243 657 24 243 570 264 90 570 267 243 657 24 243}
\RouteRow{Tríadas estables por bloque}{0 1 1 3 4 5 6 7 8 9 10 11 12 13 14 15 16 17}
\end{tabularx}\endgroup\par\vspace{5pt}\noindent{\fontsize{8.4}{10}\selectfont\hyperref[trace-261]{Procedencia y códigos originales}\hfill Ficha completa · 261/324}
\clearpage\phantomsection\label{nav:vi-ruta-262}
% VI-CATALOG-ROW rutas 262
\pdfbookmark[3]{r8c3\_h-1v+1}{route-262}
\RouteTitle{Ruta 262}{r8c3\_\allowbreak{}h-1v+1}
\noindent Celda 66\quad (8, 3)\qquad Orientación \(h=-1,\ v=1\)\hfill\hyperref[sec:vi-indice-rutas]{Índice}\par\vspace{4pt}
\begingroup\fontsize{10.2}{12.5}\selectfont\setlength{\tabcolsep}{5pt}\renewcommand{\arraystretch}{1.1}\rowcolors{1}{routepale}{white}\noindent\begin{tabularx}{\linewidth}{@{}>{\raggedright\arraybackslash}p{.345\linewidth}>{\raggedright\arraybackslash}X@{}}
\RouteGroup{Identidad estructural}
\RouteRow{Clase y multisección}{quark\_\allowbreak{}down\_\allowbreak{}G3\quad /\quad quark\_\allowbreak{}down\_\allowbreak{}G3}
\RouteRow{Colección y sector}{quark de tipo abajo; quark}
\RouteRow{Clasificación física}{quark de tipo abajo}
\RouteRow{Generación, profundidad y color}{3; G3; B}
\RouteRow{Ruta primaria y órbita de inversión}{\(B_1\): sí\qquad 2,\allowbreak{}7;\allowbreak{}8,\allowbreak{}3}
\RouteRow{Firma de clasificación}{34|\allowbreak{}0|\allowbreak{}-8|\allowbreak{}-2|\allowbreak{}-8|\allowbreak{}++0-\kern0pt{}-\kern0pt{}-++0-\kern0pt{}-\kern0pt{}-}
\RouteGroup{Retornos, memoria y transporte}
\RouteRow{Retornos con fase}{clase: 27; órbita: 27; celda: 27}
\RouteRow{Retorno cinemático}{en 108: 54; extensión: 216}
\RouteRow{Giros y cambios de hoja}{71; 107}
\RouteRow{Acarreo y diferimiento}{deuda total: 1; final: 0; bloques diferidos: 1}
\RouteRow{Tríadas finales; persistencia de Pareto}{17; sí}
\RouteRow{\(K\longrightarrow K+9\)}{repeticiones: 9; cambios de memoria: 9}
\RouteGroup{Firma, fase y diferencias}
\RouteRow{\(n_A,\ s_C,\ \kappa_0,\ \nu_{120},\ \nu_{270}\)}{14, -4, 0, 0, -12}
\RouteRow{\(\tau_{12}\); firma histórica \(\mathbb Z^5\)}{+++-\kern0pt{}-\kern0pt{}-+++-\kern0pt{}-\kern0pt{}-; 14|\allowbreak{}-4|\allowbreak{}0|\allowbreak{}0|\allowbreak{}-12}
\RouteRow{\(D_1,\ D_3,\ D_4,\ D_{27},\ D_{36}\)}{78, 24, 78, 24, 16}
\RouteRow{\(\Omega_{90/120},\ P_6\)}{80, 6}
\RouteGroup{Lectores y factores de acción}
\RouteRow{\(K_D\quad /\quad K_\Omega\)}{(40,\allowbreak{}78,\allowbreak{}27)\quad /\quad (80,\allowbreak{}54,\allowbreak{}6)}
\RouteRow{\(\kappa_D\)}{39.43380\allowbreak{}88030085\allowbreak{}51303671\allowbreak{}14}
\RouteRow{\(\kappa_\Omega\)}{79.60661\allowbreak{}01397948\allowbreak{}27586470\allowbreak{}91}
\RouteRow{\(R_{\mathrm{act},D}\)}{1.000000\allowbreak{}02733874\allowbreak{}08548943\allowbreak{}58}
\RouteRow{\(R_{\mathrm{act},\Omega}\)}{1.000000\allowbreak{}05518981\allowbreak{}25318591\allowbreak{}40}
\RouteGroup{Separación de prefijos}
\RouteRow{Área de ambigüedad; área \(\log_2\)}{36; 18.0}
\RouteRow{Primer bloque de separación}{órbita: \textemdash{}; celda: \textemdash{}; ruta: \textemdash{}}
\RouteRow{Ambigüedad final}{2}
\RouteGroup{Secuencias completas}
\RouteRow{\(w_6\)}{585 507 504 585 510 666 672 510 666 585 507 504 585 510 666 672 510 666}
\RouteRow{Tríadas estables por bloque}{0 1 2 3 4 5 6 7 8 9 10 11 12 13 14 14 16 17}
\end{tabularx}\endgroup\par\vspace{5pt}\noindent{\fontsize{8.4}{10}\selectfont\hyperref[trace-262]{Procedencia y códigos originales}\hfill Ficha completa · 262/324}
\clearpage\phantomsection\label{nav:vi-ruta-263}
% VI-CATALOG-ROW rutas 263
\pdfbookmark[3]{r8c3\_h+1v-1}{route-263}
\RouteTitle{Ruta 263}{r8c3\_\allowbreak{}h+1v-1}
\noindent Celda 66\quad (8, 3)\qquad Orientación \(h=1,\ v=-1\)\hfill\hyperref[sec:vi-indice-rutas]{Índice}\par\vspace{4pt}
\begingroup\fontsize{10.2}{12.5}\selectfont\setlength{\tabcolsep}{5pt}\renewcommand{\arraystretch}{1.1}\rowcolors{1}{routepale}{white}\noindent\begin{tabularx}{\linewidth}{@{}>{\raggedright\arraybackslash}p{.345\linewidth}>{\raggedright\arraybackslash}X@{}}
\RouteGroup{Identidad estructural}
\RouteRow{Clase y multisección}{quark\_\allowbreak{}down\_\allowbreak{}G3\quad /\quad quark\_\allowbreak{}down\_\allowbreak{}G3}
\RouteRow{Colección y sector}{quark de tipo abajo; quark}
\RouteRow{Clasificación física}{quark de tipo abajo}
\RouteRow{Generación, profundidad y color}{3; G3; B}
\RouteRow{Ruta primaria y órbita de inversión}{\(B_1\): no\qquad 2,\allowbreak{}7;\allowbreak{}8,\allowbreak{}3}
\RouteRow{Firma de clasificación}{34|\allowbreak{}0|\allowbreak{}-8|\allowbreak{}-2|\allowbreak{}-8|\allowbreak{}++0-\kern0pt{}-\kern0pt{}-++0-\kern0pt{}-\kern0pt{}-}
\RouteGroup{Retornos, memoria y transporte}
\RouteRow{Retornos con fase}{clase: 27; órbita: 27; celda: 27}
\RouteRow{Retorno cinemático}{en 108: 54; extensión: 216}
\RouteRow{Giros y cambios de hoja}{71; 107}
\RouteRow{Acarreo y diferimiento}{deuda total: 1; final: 0; bloques diferidos: 1}
\RouteRow{Tríadas finales; persistencia de Pareto}{17; sí}
\RouteRow{\(K\longrightarrow K+9\)}{repeticiones: 9; cambios de memoria: 9}
\RouteGroup{Firma, fase y diferencias}
\RouteRow{\(n_A,\ s_C,\ \kappa_0,\ \nu_{120},\ \nu_{270}\)}{14, 4, 0, 0, -12}
\RouteRow{\(\tau_{12}\); firma histórica \(\mathbb Z^5\)}{+++-\kern0pt{}-\kern0pt{}-+++-\kern0pt{}-\kern0pt{}-; 14|\allowbreak{}4|\allowbreak{}0|\allowbreak{}0|\allowbreak{}-12}
\RouteRow{\(D_1,\ D_3,\ D_4,\ D_{27},\ D_{36}\)}{78, 24, 78, 24, 16}
\RouteRow{\(\Omega_{90/120},\ P_6\)}{80, 6}
\RouteGroup{Lectores y factores de acción}
\RouteRow{\(K_D\quad /\quad K_\Omega\)}{(40,\allowbreak{}78,\allowbreak{}27)\quad /\quad (80,\allowbreak{}54,\allowbreak{}6)}
\RouteRow{\(\kappa_D\)}{39.43380\allowbreak{}88030085\allowbreak{}51303671\allowbreak{}14}
\RouteRow{\(\kappa_\Omega\)}{79.60661\allowbreak{}01397948\allowbreak{}27586470\allowbreak{}91}
\RouteRow{\(R_{\mathrm{act},D}\)}{1.000000\allowbreak{}02733874\allowbreak{}08548943\allowbreak{}58}
\RouteRow{\(R_{\mathrm{act},\Omega}\)}{1.000000\allowbreak{}05518981\allowbreak{}25318591\allowbreak{}40}
\RouteGroup{Separación de prefijos}
\RouteRow{Área de ambigüedad; área \(\log_2\)}{36; 18.0}
\RouteRow{Primer bloque de separación}{órbita: \textemdash{}; celda: \textemdash{}; ruta: \textemdash{}}
\RouteRow{Ambigüedad final}{2}
\RouteGroup{Secuencias completas}
\RouteRow{\(w_6\)}{672 510 666 672 507 504 585 507 504 672 510 666 672 507 504 585 507 504}
\RouteRow{Tríadas estables por bloque}{0 1 2 3 4 5 6 7 8 9 10 11 12 13 14 14 16 17}
\end{tabularx}\endgroup\par\vspace{5pt}\noindent{\fontsize{8.4}{10}\selectfont\hyperref[trace-263]{Procedencia y códigos originales}\hfill Ficha completa · 263/324}
\clearpage\phantomsection\label{nav:vi-ruta-264}
% VI-CATALOG-ROW rutas 264
\pdfbookmark[3]{r8c3\_h+1v+1}{route-264}
\RouteTitle{Ruta 264}{r8c3\_\allowbreak{}h+1v+1}
\noindent Celda 66\quad (8, 3)\qquad Orientación \(h=1,\ v=1\)\hfill\hyperref[sec:vi-indice-rutas]{Índice}\par\vspace{4pt}
\begingroup\fontsize{10.2}{12.5}\selectfont\setlength{\tabcolsep}{5pt}\renewcommand{\arraystretch}{1.1}\rowcolors{1}{routepale}{white}\noindent\begin{tabularx}{\linewidth}{@{}>{\raggedright\arraybackslash}p{.345\linewidth}>{\raggedright\arraybackslash}X@{}}
\RouteGroup{Identidad estructural}
\RouteRow{Clase y multisección}{quark\_\allowbreak{}down\_\allowbreak{}G3\quad /\quad quark\_\allowbreak{}down\_\allowbreak{}G3}
\RouteRow{Colección y sector}{quark de tipo abajo; quark}
\RouteRow{Clasificación física}{quark de tipo abajo}
\RouteRow{Generación, profundidad y color}{3; G3; B}
\RouteRow{Ruta primaria y órbita de inversión}{\(B_1\): no\qquad 2,\allowbreak{}7;\allowbreak{}8,\allowbreak{}3}
\RouteRow{Firma de clasificación}{34|\allowbreak{}0|\allowbreak{}-8|\allowbreak{}-2|\allowbreak{}-8|\allowbreak{}++0-\kern0pt{}-\kern0pt{}-++0-\kern0pt{}-\kern0pt{}-}
\RouteGroup{Retornos, memoria y transporte}
\RouteRow{Retornos con fase}{clase: 27; órbita: 27; celda: 27}
\RouteRow{Retorno cinemático}{en 108: 54; extensión: 216}
\RouteRow{Giros y cambios de hoja}{71; 107}
\RouteRow{Acarreo y diferimiento}{deuda total: 0; final: 0; bloques diferidos: 0}
\RouteRow{Tríadas finales; persistencia de Pareto}{17; sí}
\RouteRow{\(K\longrightarrow K+9\)}{repeticiones: 9; cambios de memoria: 9}
\RouteGroup{Firma, fase y diferencias}
\RouteRow{\(n_A,\ s_C,\ \kappa_0,\ \nu_{120},\ \nu_{270}\)}{34, 0, 2, 0, -8}
\RouteRow{\(\tau_{12}\); firma histórica \(\mathbb Z^5\)}{++0-\kern0pt{}-\kern0pt{}-++0-\kern0pt{}-\kern0pt{}-; 34|\allowbreak{}0|\allowbreak{}2|\allowbreak{}0|\allowbreak{}-8}
\RouteRow{\(D_1,\ D_3,\ D_4,\ D_{27},\ D_{36}\)}{66, 60, 68, 48, 32}
\RouteRow{\(\Omega_{90/120},\ P_6\)}{56, 5}
\RouteGroup{Lectores y factores de acción}
\RouteRow{\(K_D\quad /\quad K_\Omega\)}{(80,\allowbreak{}66,\allowbreak{}4)\quad /\quad (56,\allowbreak{}54,\allowbreak{}5)}
\RouteRow{\(\kappa_D\)}{79.51881\allowbreak{}95161180\allowbreak{}37694395\allowbreak{}39}
\RouteRow{\(\kappa_\Omega\)}{55.60649\allowbreak{}89433721\allowbreak{}35446416\allowbreak{}86}
\RouteRow{\(R_{\mathrm{act},D}\)}{1.000000\allowbreak{}05512894\allowbreak{}88901612\allowbreak{}64}
\RouteRow{\(R_{\mathrm{act},\Omega}\)}{1.000000\allowbreak{}03855097\allowbreak{}23541416\allowbreak{}45}
\RouteGroup{Separación de prefijos}
\RouteRow{Área de ambigüedad; área \(\log_2\)}{36; 18.0}
\RouteRow{Primer bloque de separación}{órbita: \textemdash{}; celda: \textemdash{}; ruta: \textemdash{}}
\RouteRow{Ambigüedad final}{2}
\RouteGroup{Secuencias completas}
\RouteRow{\(w_6\)}{657 24 243 657 21 90 570 264 90 657 24 243 657 21 90 570 264 90}
\RouteRow{Tríadas estables por bloque}{0 1 2 3 4 5 6 7 8 9 10 11 12 13 14 15 16 17}
\end{tabularx}\endgroup\par\vspace{5pt}\noindent{\fontsize{8.4}{10}\selectfont\hyperref[trace-264]{Procedencia y códigos originales}\hfill Ficha completa · 264/324}
\clearpage\phantomsection\label{nav:vi-ruta-265}
% VI-CATALOG-ROW rutas 265
\pdfbookmark[2]{Celda 67}{cell-67}
\pdfbookmark[3]{r8c4\_h-1v-1}{route-265}
\RouteTitle{Ruta 265}{r8c4\_\allowbreak{}h-1v-1}
\noindent Celda 67\quad (8, 4)\qquad Orientación \(h=-1,\ v=-1\)\hfill\hyperref[sec:vi-indice-rutas]{Índice}\par\vspace{4pt}
\begingroup\fontsize{10.2}{12.5}\selectfont\setlength{\tabcolsep}{5pt}\renewcommand{\arraystretch}{1.1}\rowcolors{1}{routepale}{white}\noindent\begin{tabularx}{\linewidth}{@{}>{\raggedright\arraybackslash}p{.345\linewidth}>{\raggedright\arraybackslash}X@{}}
\RouteGroup{Identidad estructural}
\RouteRow{Clase y multisección}{quark\_\allowbreak{}up\_\allowbreak{}G3\quad /\quad quark\_\allowbreak{}up\_\allowbreak{}G3}
\RouteRow{Colección y sector}{quark de tipo arriba; quark}
\RouteRow{Clasificación física}{quark de tipo arriba}
\RouteRow{Generación, profundidad y color}{3; G3; G}
\RouteRow{Ruta primaria y órbita de inversión}{\(B_1\): no\qquad 2,\allowbreak{}6;\allowbreak{}8,\allowbreak{}4}
\RouteRow{Firma de clasificación}{34|\allowbreak{}0|\allowbreak{}-4|\allowbreak{}2|\allowbreak{}-8|\allowbreak{}+++-\kern0pt{}-0+++-\kern0pt{}-0}
\RouteGroup{Retornos, memoria y transporte}
\RouteRow{Retornos con fase}{clase: 27; órbita: 27; celda: 27}
\RouteRow{Retorno cinemático}{en 108: 54; extensión: 216}
\RouteRow{Giros y cambios de hoja}{71; 107}
\RouteRow{Acarreo y diferimiento}{deuda total: 1; final: 1; bloques diferidos: 1}
\RouteRow{Tríadas finales; persistencia de Pareto}{16; no}
\RouteRow{\(K\longrightarrow K+9\)}{repeticiones: 9; cambios de memoria: 9}
\RouteGroup{Firma, fase y diferencias}
\RouteRow{\(n_A,\ s_C,\ \kappa_0,\ \nu_{120},\ \nu_{270}\)}{34, 0, 2, 0, -8}
\RouteRow{\(\tau_{12}\); firma histórica \(\mathbb Z^5\)}{++0-\kern0pt{}-\kern0pt{}-++0-\kern0pt{}-\kern0pt{}-; 34|\allowbreak{}0|\allowbreak{}2|\allowbreak{}0|\allowbreak{}-8}
\RouteRow{\(D_1,\ D_3,\ D_4,\ D_{27},\ D_{36}\)}{66, 64, 68, 60, 40}
\RouteRow{\(\Omega_{90/120},\ P_6\)}{56, 5}
\RouteGroup{Lectores y factores de acción}
\RouteRow{\(K_D\quad /\quad K_\Omega\)}{(100,\allowbreak{}66,\allowbreak{}2)\quad /\quad (56,\allowbreak{}54,\allowbreak{}5)}
\RouteRow{\(\kappa_D\)}{99.51859\allowbreak{}71232726\allowbreak{}53414287\allowbreak{}29}
\RouteRow{\(\kappa_\Omega\)}{55.60649\allowbreak{}89433721\allowbreak{}35446416\allowbreak{}86}
\RouteRow{\(R_{\mathrm{act},D}\)}{1.000000\allowbreak{}06899443\allowbreak{}08259364\allowbreak{}17}
\RouteRow{\(R_{\mathrm{act},\Omega}\)}{1.000000\allowbreak{}03855097\allowbreak{}23541416\allowbreak{}45}
\RouteGroup{Separación de prefijos}
\RouteRow{Área de ambigüedad; área \(\log_2\)}{36; 18.0}
\RouteRow{Primer bloque de separación}{órbita: \textemdash{}; celda: \textemdash{}; ruta: \textemdash{}}
\RouteRow{Ambigüedad final}{2}
\RouteGroup{Secuencias completas}
\RouteRow{\(w_6\)}{15 486 423 15 489 498 99 327 498 15 486 423 15 489 498 99 327 498}
\RouteRow{Tríadas estables por bloque}{0 1 2 3 4 5 6 7 8 9 10 11 12 13 14 15 16 16}
\end{tabularx}\endgroup\par\vspace{5pt}\noindent{\fontsize{8.4}{10}\selectfont\hyperref[trace-265]{Procedencia y códigos originales}\hfill Ficha completa · 265/324}
\clearpage\phantomsection\label{nav:vi-ruta-266}
% VI-CATALOG-ROW rutas 266
\pdfbookmark[3]{r8c4\_h-1v+1}{route-266}
\RouteTitle{Ruta 266}{r8c4\_\allowbreak{}h-1v+1}
\noindent Celda 67\quad (8, 4)\qquad Orientación \(h=-1,\ v=1\)\hfill\hyperref[sec:vi-indice-rutas]{Índice}\par\vspace{4pt}
\begingroup\fontsize{10.2}{12.5}\selectfont\setlength{\tabcolsep}{5pt}\renewcommand{\arraystretch}{1.1}\rowcolors{1}{routepale}{white}\noindent\begin{tabularx}{\linewidth}{@{}>{\raggedright\arraybackslash}p{.345\linewidth}>{\raggedright\arraybackslash}X@{}}
\RouteGroup{Identidad estructural}
\RouteRow{Clase y multisección}{quark\_\allowbreak{}up\_\allowbreak{}G3\quad /\quad quark\_\allowbreak{}up\_\allowbreak{}G3}
\RouteRow{Colección y sector}{quark de tipo arriba; quark}
\RouteRow{Clasificación física}{quark de tipo arriba}
\RouteRow{Generación, profundidad y color}{3; G3; G}
\RouteRow{Ruta primaria y órbita de inversión}{\(B_1\): no\qquad 2,\allowbreak{}6;\allowbreak{}8,\allowbreak{}4}
\RouteRow{Firma de clasificación}{34|\allowbreak{}0|\allowbreak{}-4|\allowbreak{}2|\allowbreak{}-8|\allowbreak{}+++-\kern0pt{}-0+++-\kern0pt{}-0}
\RouteGroup{Retornos, memoria y transporte}
\RouteRow{Retornos con fase}{clase: 27; órbita: 27; celda: 27}
\RouteRow{Retorno cinemático}{en 108: 54; extensión: 216}
\RouteRow{Giros y cambios de hoja}{71; 107}
\RouteRow{Acarreo y diferimiento}{deuda total: 0; final: 0; bloques diferidos: 0}
\RouteRow{Tríadas finales; persistencia de Pareto}{17; no}
\RouteRow{\(K\longrightarrow K+9\)}{repeticiones: 9; cambios de memoria: 9}
\RouteGroup{Firma, fase y diferencias}
\RouteRow{\(n_A,\ s_C,\ \kappa_0,\ \nu_{120},\ \nu_{270}\)}{50, -6, 2, -2, -8}
\RouteRow{\(\tau_{12}\); firma histórica \(\mathbb Z^5\)}{++0-\kern0pt{}-\kern0pt{}-++0-\kern0pt{}-\kern0pt{}-; 50|\allowbreak{}-6|\allowbreak{}2|\allowbreak{}-2|\allowbreak{}-8}
\RouteRow{\(D_1,\ D_3,\ D_4,\ D_{27},\ D_{36}\)}{24, 20, 24, 24, 16}
\RouteRow{\(\Omega_{90/120},\ P_6\)}{56, 5}
\RouteGroup{Lectores y factores de acción}
\RouteRow{\(K_D\quad /\quad K_\Omega\)}{(40,\allowbreak{}24,\allowbreak{}2)\quad /\quad (56,\allowbreak{}54,\allowbreak{}5)}
\RouteRow{\(\kappa_D\)}{39.82508\allowbreak{}59311825\allowbreak{}73056173\allowbreak{}26}
\RouteRow{\(\kappa_\Omega\)}{55.60649\allowbreak{}89433721\allowbreak{}35446416\allowbreak{}86}
\RouteRow{\(R_{\mathrm{act},D}\)}{1.000000\allowbreak{}02761000\allowbreak{}61595142\allowbreak{}15}
\RouteRow{\(R_{\mathrm{act},\Omega}\)}{1.000000\allowbreak{}03855097\allowbreak{}23541416\allowbreak{}45}
\RouteGroup{Separación de prefijos}
\RouteRow{Área de ambigüedad; área \(\log_2\)}{72; 36.0}
\RouteRow{Primer bloque de separación}{órbita: \textemdash{}; celda: \textemdash{}; ruta: \textemdash{}}
\RouteRow{Ambigüedad final}{4}
\RouteGroup{Secuencias completas}
\RouteRow{\(w_6\)}{0 81 3 0 84 0 81 3 0 0 81 3 0 84 0 81 3 0}
\RouteRow{Tríadas estables por bloque}{0 1 2 3 4 5 6 7 8 9 10 11 12 13 14 15 16 17}
\end{tabularx}\endgroup\par\vspace{5pt}\noindent{\fontsize{8.4}{10}\selectfont\hyperref[trace-266]{Procedencia y códigos originales}\hfill Ficha completa · 266/324}
\clearpage\phantomsection\label{nav:vi-ruta-267}
% VI-CATALOG-ROW rutas 267
\pdfbookmark[3]{r8c4\_h+1v-1}{route-267}
\RouteTitle{Ruta 267}{r8c4\_\allowbreak{}h+1v-1}
\noindent Celda 67\quad (8, 4)\qquad Orientación \(h=1,\ v=-1\)\hfill\hyperref[sec:vi-indice-rutas]{Índice}\par\vspace{4pt}
\begingroup\fontsize{10.2}{12.5}\selectfont\setlength{\tabcolsep}{5pt}\renewcommand{\arraystretch}{1.1}\rowcolors{1}{routepale}{white}\noindent\begin{tabularx}{\linewidth}{@{}>{\raggedright\arraybackslash}p{.345\linewidth}>{\raggedright\arraybackslash}X@{}}
\RouteGroup{Identidad estructural}
\RouteRow{Clase y multisección}{quark\_\allowbreak{}up\_\allowbreak{}G3\quad /\quad quark\_\allowbreak{}up\_\allowbreak{}G3}
\RouteRow{Colección y sector}{quark de tipo arriba; quark}
\RouteRow{Clasificación física}{quark de tipo arriba}
\RouteRow{Generación, profundidad y color}{3; G3; G}
\RouteRow{Ruta primaria y órbita de inversión}{\(B_1\): no\qquad 2,\allowbreak{}6;\allowbreak{}8,\allowbreak{}4}
\RouteRow{Firma de clasificación}{34|\allowbreak{}0|\allowbreak{}-4|\allowbreak{}2|\allowbreak{}-8|\allowbreak{}+++-\kern0pt{}-0+++-\kern0pt{}-0}
\RouteGroup{Retornos, memoria y transporte}
\RouteRow{Retornos con fase}{clase: 27; órbita: 27; celda: 27}
\RouteRow{Retorno cinemático}{en 108: 54; extensión: 216}
\RouteRow{Giros y cambios de hoja}{71; 107}
\RouteRow{Acarreo y diferimiento}{deuda total: 2; final: 0; bloques diferidos: 2}
\RouteRow{Tríadas finales; persistencia de Pareto}{17; no}
\RouteRow{\(K\longrightarrow K+9\)}{repeticiones: 9; cambios de memoria: 9}
\RouteGroup{Firma, fase y diferencias}
\RouteRow{\(n_A,\ s_C,\ \kappa_0,\ \nu_{120},\ \nu_{270}\)}{50, 6, -2, 2, -8}
\RouteRow{\(\tau_{12}\); firma histórica \(\mathbb Z^5\)}{+++-\kern0pt{}-0+++-\kern0pt{}-0; 50|\allowbreak{}6|\allowbreak{}-2|\allowbreak{}2|\allowbreak{}-8}
\RouteRow{\(D_1,\ D_3,\ D_4,\ D_{27},\ D_{36}\)}{24, 20, 24, 24, 16}
\RouteRow{\(\Omega_{90/120},\ P_6\)}{56, 5}
\RouteGroup{Lectores y factores de acción}
\RouteRow{\(K_D\quad /\quad K_\Omega\)}{(40,\allowbreak{}24,\allowbreak{}2)\quad /\quad (56,\allowbreak{}54,\allowbreak{}5)}
\RouteRow{\(\kappa_D\)}{39.82508\allowbreak{}59311825\allowbreak{}73056173\allowbreak{}26}
\RouteRow{\(\kappa_\Omega\)}{55.60649\allowbreak{}89433721\allowbreak{}35446416\allowbreak{}86}
\RouteRow{\(R_{\mathrm{act},D}\)}{1.000000\allowbreak{}02761000\allowbreak{}61595142\allowbreak{}15}
\RouteRow{\(R_{\mathrm{act},\Omega}\)}{1.000000\allowbreak{}03855097\allowbreak{}23541416\allowbreak{}45}
\RouteGroup{Separación de prefijos}
\RouteRow{Área de ambigüedad; área \(\log_2\)}{72; 36.0}
\RouteRow{Primer bloque de separación}{órbita: \textemdash{}; celda: \textemdash{}; ruta: \textemdash{}}
\RouteRow{Ambigüedad final}{4}
\RouteGroup{Secuencias completas}
\RouteRow{\(w_6\)}{81 3 0 81 0 3 0 81 3 81 3 0 81 0 3 0 81 3}
\RouteRow{Tríadas estables por bloque}{0 1 2 3 4 5 6 6 8 9 10 11 12 13 14 15 15 17}
\end{tabularx}\endgroup\par\vspace{5pt}\noindent{\fontsize{8.4}{10}\selectfont\hyperref[trace-267]{Procedencia y códigos originales}\hfill Ficha completa · 267/324}
\clearpage\phantomsection\label{nav:vi-ruta-268}
% VI-CATALOG-ROW rutas 268
\pdfbookmark[3]{r8c4\_h+1v+1}{route-268}
\RouteTitle{Ruta 268}{r8c4\_\allowbreak{}h+1v+1}
\noindent Celda 67\quad (8, 4)\qquad Orientación \(h=1,\ v=1\)\hfill\hyperref[sec:vi-indice-rutas]{Índice}\par\vspace{4pt}
\begingroup\fontsize{10.2}{12.5}\selectfont\setlength{\tabcolsep}{5pt}\renewcommand{\arraystretch}{1.1}\rowcolors{1}{routepale}{white}\noindent\begin{tabularx}{\linewidth}{@{}>{\raggedright\arraybackslash}p{.345\linewidth}>{\raggedright\arraybackslash}X@{}}
\RouteGroup{Identidad estructural}
\RouteRow{Clase y multisección}{quark\_\allowbreak{}up\_\allowbreak{}G3\quad /\quad quark\_\allowbreak{}up\_\allowbreak{}G3}
\RouteRow{Colección y sector}{quark de tipo arriba; quark}
\RouteRow{Clasificación física}{quark de tipo arriba}
\RouteRow{Generación, profundidad y color}{3; G3; G}
\RouteRow{Ruta primaria y órbita de inversión}{\(B_1\): sí\qquad 2,\allowbreak{}6;\allowbreak{}8,\allowbreak{}4}
\RouteRow{Firma de clasificación}{34|\allowbreak{}0|\allowbreak{}-4|\allowbreak{}2|\allowbreak{}-8|\allowbreak{}+++-\kern0pt{}-0+++-\kern0pt{}-0}
\RouteGroup{Retornos, memoria y transporte}
\RouteRow{Retornos con fase}{clase: 27; órbita: 27; celda: 27}
\RouteRow{Retorno cinemático}{en 108: 54; extensión: 216}
\RouteRow{Giros y cambios de hoja}{71; 107}
\RouteRow{Acarreo y diferimiento}{deuda total: 1; final: 0; bloques diferidos: 1}
\RouteRow{Tríadas finales; persistencia de Pareto}{17; no}
\RouteRow{\(K\longrightarrow K+9\)}{repeticiones: 9; cambios de memoria: 9}
\RouteGroup{Firma, fase y diferencias}
\RouteRow{\(n_A,\ s_C,\ \kappa_0,\ \nu_{120},\ \nu_{270}\)}{34, 0, -2, 0, -8}
\RouteRow{\(\tau_{12}\); firma histórica \(\mathbb Z^5\)}{+++-\kern0pt{}-0+++-\kern0pt{}-0; 34|\allowbreak{}0|\allowbreak{}-2|\allowbreak{}0|\allowbreak{}-8}
\RouteRow{\(D_1,\ D_3,\ D_4,\ D_{27},\ D_{36}\)}{66, 64, 68, 60, 40}
\RouteRow{\(\Omega_{90/120},\ P_6\)}{56, 5}
\RouteGroup{Lectores y factores de acción}
\RouteRow{\(K_D\quad /\quad K_\Omega\)}{(100,\allowbreak{}66,\allowbreak{}2)\quad /\quad (56,\allowbreak{}54,\allowbreak{}5)}
\RouteRow{\(\kappa_D\)}{99.51859\allowbreak{}71232726\allowbreak{}53414287\allowbreak{}29}
\RouteRow{\(\kappa_\Omega\)}{55.60649\allowbreak{}89433721\allowbreak{}35446416\allowbreak{}86}
\RouteRow{\(R_{\mathrm{act},D}\)}{1.000000\allowbreak{}06899443\allowbreak{}08259364\allowbreak{}17}
\RouteRow{\(R_{\mathrm{act},\Omega}\)}{1.000000\allowbreak{}03855097\allowbreak{}23541416\allowbreak{}45}
\RouteGroup{Separación de prefijos}
\RouteRow{Área de ambigüedad; área \(\log_2\)}{36; 18.0}
\RouteRow{Primer bloque de separación}{órbita: \textemdash{}; celda: \textemdash{}; ruta: \textemdash{}}
\RouteRow{Ambigüedad final}{2}
\RouteGroup{Secuencias completas}
\RouteRow{\(w_6\)}{99 327 498 99 324 423 15 486 423 99 327 498 99 324 423 15 486 423}
\RouteRow{Tríadas estables por bloque}{0 1 2 3 4 5 6 7 8 9 10 11 12 13 14 14 16 17}
\end{tabularx}\endgroup\par\vspace{5pt}\noindent{\fontsize{8.4}{10}\selectfont\hyperref[trace-268]{Procedencia y códigos originales}\hfill Ficha completa · 268/324}
\clearpage\phantomsection\label{nav:vi-ruta-269}
% VI-CATALOG-ROW rutas 269
\pdfbookmark[2]{Celda 68}{cell-68}
\pdfbookmark[3]{r8c5\_h-1v-1}{route-269}
\RouteTitle{Ruta 269}{r8c5\_\allowbreak{}h-1v-1}
\noindent Celda 68\quad (8, 5)\qquad Orientación \(h=-1,\ v=-1\)\hfill\hyperref[sec:vi-indice-rutas]{Índice}\par\vspace{4pt}
\begingroup\fontsize{10.2}{12.5}\selectfont\setlength{\tabcolsep}{5pt}\renewcommand{\arraystretch}{1.1}\rowcolors{1}{routepale}{white}\noindent\begin{tabularx}{\linewidth}{@{}>{\raggedright\arraybackslash}p{.345\linewidth}>{\raggedright\arraybackslash}X@{}}
\RouteGroup{Identidad estructural}
\RouteRow{Clase y multisección}{quark\_\allowbreak{}down\_\allowbreak{}G3\quad /\quad quark\_\allowbreak{}down\_\allowbreak{}G3}
\RouteRow{Colección y sector}{quark de tipo abajo; quark}
\RouteRow{Clasificación física}{quark de tipo abajo}
\RouteRow{Generación, profundidad y color}{3; G3; R}
\RouteRow{Ruta primaria y órbita de inversión}{\(B_1\): no\qquad 2,\allowbreak{}5;\allowbreak{}8,\allowbreak{}5}
\RouteRow{Firma de clasificación}{30|\allowbreak{}6|\allowbreak{}-4|\allowbreak{}0|\allowbreak{}-12|\allowbreak{}+++-\kern0pt{}-\kern0pt{}-+++-\kern0pt{}-\kern0pt{}-}
\RouteGroup{Retornos, memoria y transporte}
\RouteRow{Retornos con fase}{clase: 27; órbita: 27; celda: 27}
\RouteRow{Retorno cinemático}{en 108: 54; extensión: 216}
\RouteRow{Giros y cambios de hoja}{71; 107}
\RouteRow{Acarreo y diferimiento}{deuda total: 3; final: 0; bloques diferidos: 2}
\RouteRow{Tríadas finales; persistencia de Pareto}{17; sí}
\RouteRow{\(K\longrightarrow K+9\)}{repeticiones: 9; cambios de memoria: 9}
\RouteGroup{Firma, fase y diferencias}
\RouteRow{\(n_A,\ s_C,\ \kappa_0,\ \nu_{120},\ \nu_{270}\)}{30, -6, 0, 0, -12}
\RouteRow{\(\tau_{12}\); firma histórica \(\mathbb Z^5\)}{+++-\kern0pt{}-\kern0pt{}-+++-\kern0pt{}-\kern0pt{}-; 30|\allowbreak{}-6|\allowbreak{}0|\allowbreak{}0|\allowbreak{}-12}
\RouteRow{\(D_1,\ D_3,\ D_4,\ D_{27},\ D_{36}\)}{54, 56, 68, 48, 32}
\RouteRow{\(\Omega_{90/120},\ P_6\)}{80, 6}
\RouteGroup{Lectores y factores de acción}
\RouteRow{\(K_D\quad /\quad K_\Omega\)}{(80,\allowbreak{}54,\allowbreak{}6)\quad /\quad (80,\allowbreak{}54,\allowbreak{}6)}
\RouteRow{\(\kappa_D\)}{79.60661\allowbreak{}01397948\allowbreak{}27586470\allowbreak{}91}
\RouteRow{\(\kappa_\Omega\)}{79.60661\allowbreak{}01397948\allowbreak{}27586470\allowbreak{}91}
\RouteRow{\(R_{\mathrm{act},D}\)}{1.000000\allowbreak{}05518981\allowbreak{}25318591\allowbreak{}40}
\RouteRow{\(R_{\mathrm{act},\Omega}\)}{1.000000\allowbreak{}05518981\allowbreak{}25318591\allowbreak{}40}
\RouteGroup{Separación de prefijos}
\RouteRow{Área de ambigüedad; área \(\log_2\)}{36; 18.0}
\RouteRow{Primer bloque de separación}{órbita: 1; celda: \textemdash{}; ruta: \textemdash{}}
\RouteRow{Ambigüedad final}{2}
\RouteGroup{Secuencias completas}
\RouteRow{\(w_6\)}{423 15 486 423 9 24 243 657 24 423 15 486 423 9 24 243 657 24}
\RouteRow{Tríadas estables por bloque}{0 1 2 3 4 5 6 7 8 9 10 11 11 13 14 14 15 17}
\end{tabularx}\endgroup\par\vspace{5pt}\noindent{\fontsize{8.4}{10}\selectfont\hyperref[trace-269]{Procedencia y códigos originales}\hfill Ficha completa · 269/324}
\clearpage\phantomsection\label{nav:vi-ruta-270}
% VI-CATALOG-ROW rutas 270
\pdfbookmark[3]{r8c5\_h-1v+1}{route-270}
\RouteTitle{Ruta 270}{r8c5\_\allowbreak{}h-1v+1}
\noindent Celda 68\quad (8, 5)\qquad Orientación \(h=-1,\ v=1\)\hfill\hyperref[sec:vi-indice-rutas]{Índice}\par\vspace{4pt}
\begingroup\fontsize{10.2}{12.5}\selectfont\setlength{\tabcolsep}{5pt}\renewcommand{\arraystretch}{1.1}\rowcolors{1}{routepale}{white}\noindent\begin{tabularx}{\linewidth}{@{}>{\raggedright\arraybackslash}p{.345\linewidth}>{\raggedright\arraybackslash}X@{}}
\RouteGroup{Identidad estructural}
\RouteRow{Clase y multisección}{quark\_\allowbreak{}down\_\allowbreak{}G3\quad /\quad quark\_\allowbreak{}down\_\allowbreak{}G3}
\RouteRow{Colección y sector}{quark de tipo abajo; quark}
\RouteRow{Clasificación física}{quark de tipo abajo}
\RouteRow{Generación, profundidad y color}{3; G3; R}
\RouteRow{Ruta primaria y órbita de inversión}{\(B_1\): sí\qquad 2,\allowbreak{}5;\allowbreak{}8,\allowbreak{}5}
\RouteRow{Firma de clasificación}{30|\allowbreak{}6|\allowbreak{}-4|\allowbreak{}0|\allowbreak{}-12|\allowbreak{}+++-\kern0pt{}-\kern0pt{}-+++-\kern0pt{}-\kern0pt{}-}
\RouteGroup{Retornos, memoria y transporte}
\RouteRow{Retornos con fase}{clase: 27; órbita: 27; celda: 27}
\RouteRow{Retorno cinemático}{en 108: 54; extensión: 216}
\RouteRow{Giros y cambios de hoja}{71; 107}
\RouteRow{Acarreo y diferimiento}{deuda total: 2; final: 1; bloques diferidos: 2}
\RouteRow{Tríadas finales; persistencia de Pareto}{16; sí}
\RouteRow{\(K\longrightarrow K+9\)}{repeticiones: 9; cambios de memoria: 9}
\RouteGroup{Firma, fase y diferencias}
\RouteRow{\(n_A,\ s_C,\ \kappa_0,\ \nu_{120},\ \nu_{270}\)}{8, -2, 0, 0, -12}
\RouteRow{\(\tau_{12}\); firma histórica \(\mathbb Z^5\)}{+++-\kern0pt{}-\kern0pt{}-+++-\kern0pt{}-\kern0pt{}-; 8|\allowbreak{}-2|\allowbreak{}0|\allowbreak{}0|\allowbreak{}-12}
\RouteRow{\(D_1,\ D_3,\ D_4,\ D_{27},\ D_{36}\)}{84, 28, 84, 36, 24}
\RouteRow{\(\Omega_{90/120},\ P_6\)}{80, 6}
\RouteGroup{Lectores y factores de acción}
\RouteRow{\(K_D\quad /\quad K_\Omega\)}{(60,\allowbreak{}84,\allowbreak{}28)\quad /\quad (80,\allowbreak{}54,\allowbreak{}6)}
\RouteRow{\(\kappa_D\)}{59.39013\allowbreak{}58840155\allowbreak{}40637741\allowbreak{}48}
\RouteRow{\(\kappa_\Omega\)}{79.60661\allowbreak{}01397948\allowbreak{}27586470\allowbreak{}91}
\RouteRow{\(R_{\mathrm{act},D}\)}{1.000000\allowbreak{}04117409\allowbreak{}89467415\allowbreak{}75}
\RouteRow{\(R_{\mathrm{act},\Omega}\)}{1.000000\allowbreak{}05518981\allowbreak{}25318591\allowbreak{}40}
\RouteGroup{Separación de prefijos}
\RouteRow{Área de ambigüedad; área \(\log_2\)}{36; 18.0}
\RouteRow{Primer bloque de separación}{órbita: 1; celda: \textemdash{}; ruta: \textemdash{}}
\RouteRow{Ambigüedad final}{2}
\RouteGroup{Secuencias completas}
\RouteRow{\(w_6\)}{414 420 258 414 414 333 255 252 333 414 420 258 414 414 333 255 252 333}
\RouteRow{Tríadas estables por bloque}{0 1 2 3 4 5 6 7 8 9 10 11 12 13 14 14 16 16}
\end{tabularx}\endgroup\par\vspace{5pt}\noindent{\fontsize{8.4}{10}\selectfont\hyperref[trace-270]{Procedencia y códigos originales}\hfill Ficha completa · 270/324}
\clearpage\phantomsection\label{nav:vi-ruta-271}
% VI-CATALOG-ROW rutas 271
\pdfbookmark[3]{r8c5\_h+1v-1}{route-271}
\RouteTitle{Ruta 271}{r8c5\_\allowbreak{}h+1v-1}
\noindent Celda 68\quad (8, 5)\qquad Orientación \(h=1,\ v=-1\)\hfill\hyperref[sec:vi-indice-rutas]{Índice}\par\vspace{4pt}
\begingroup\fontsize{10.2}{12.5}\selectfont\setlength{\tabcolsep}{5pt}\renewcommand{\arraystretch}{1.1}\rowcolors{1}{routepale}{white}\noindent\begin{tabularx}{\linewidth}{@{}>{\raggedright\arraybackslash}p{.345\linewidth}>{\raggedright\arraybackslash}X@{}}
\RouteGroup{Identidad estructural}
\RouteRow{Clase y multisección}{quark\_\allowbreak{}down\_\allowbreak{}G3\quad /\quad quark\_\allowbreak{}down\_\allowbreak{}G3}
\RouteRow{Colección y sector}{quark de tipo abajo; quark}
\RouteRow{Clasificación física}{quark de tipo abajo}
\RouteRow{Generación, profundidad y color}{3; G3; R}
\RouteRow{Ruta primaria y órbita de inversión}{\(B_1\): no\qquad 2,\allowbreak{}5;\allowbreak{}8,\allowbreak{}5}
\RouteRow{Firma de clasificación}{30|\allowbreak{}6|\allowbreak{}-4|\allowbreak{}0|\allowbreak{}-12|\allowbreak{}+++-\kern0pt{}-\kern0pt{}-+++-\kern0pt{}-\kern0pt{}-}
\RouteGroup{Retornos, memoria y transporte}
\RouteRow{Retornos con fase}{clase: 27; órbita: 27; celda: 27}
\RouteRow{Retorno cinemático}{en 108: 54; extensión: 216}
\RouteRow{Giros y cambios de hoja}{71; 107}
\RouteRow{Acarreo y diferimiento}{deuda total: 3; final: 0; bloques diferidos: 2}
\RouteRow{Tríadas finales; persistencia de Pareto}{17; sí}
\RouteRow{\(K\longrightarrow K+9\)}{repeticiones: 9; cambios de memoria: 9}
\RouteGroup{Firma, fase y diferencias}
\RouteRow{\(n_A,\ s_C,\ \kappa_0,\ \nu_{120},\ \nu_{270}\)}{8, 2, 0, 0, -12}
\RouteRow{\(\tau_{12}\); firma histórica \(\mathbb Z^5\)}{+++-\kern0pt{}-\kern0pt{}-+++-\kern0pt{}-\kern0pt{}-; 8|\allowbreak{}2|\allowbreak{}0|\allowbreak{}0|\allowbreak{}-12}
\RouteRow{\(D_1,\ D_3,\ D_4,\ D_{27},\ D_{36}\)}{84, 28, 84, 36, 24}
\RouteRow{\(\Omega_{90/120},\ P_6\)}{80, 6}
\RouteGroup{Lectores y factores de acción}
\RouteRow{\(K_D\quad /\quad K_\Omega\)}{(60,\allowbreak{}84,\allowbreak{}28)\quad /\quad (80,\allowbreak{}54,\allowbreak{}6)}
\RouteRow{\(\kappa_D\)}{59.39013\allowbreak{}58840155\allowbreak{}40637741\allowbreak{}48}
\RouteRow{\(\kappa_\Omega\)}{79.60661\allowbreak{}01397948\allowbreak{}27586470\allowbreak{}91}
\RouteRow{\(R_{\mathrm{act},D}\)}{1.000000\allowbreak{}04117409\allowbreak{}89467415\allowbreak{}75}
\RouteRow{\(R_{\mathrm{act},\Omega}\)}{1.000000\allowbreak{}05518981\allowbreak{}25318591\allowbreak{}40}
\RouteGroup{Separación de prefijos}
\RouteRow{Área de ambigüedad; área \(\log_2\)}{36; 18.0}
\RouteRow{Primer bloque de separación}{órbita: 1; celda: \textemdash{}; ruta: \textemdash{}}
\RouteRow{Ambigüedad final}{2}
\RouteGroup{Secuencias completas}
\RouteRow{\(w_6\)}{255 252 333 255 258 258 414 420 258 255 252 333 255 258 258 414 420 258}
\RouteRow{Tríadas estables por bloque}{0 1 2 3 4 5 6 7 8 9 10 11 11 12 14 15 15 17}
\end{tabularx}\endgroup\par\vspace{5pt}\noindent{\fontsize{8.4}{10}\selectfont\hyperref[trace-271]{Procedencia y códigos originales}\hfill Ficha completa · 271/324}
\clearpage\phantomsection\label{nav:vi-ruta-272}
% VI-CATALOG-ROW rutas 272
\pdfbookmark[3]{r8c5\_h+1v+1}{route-272}
\RouteTitle{Ruta 272}{r8c5\_\allowbreak{}h+1v+1}
\noindent Celda 68\quad (8, 5)\qquad Orientación \(h=1,\ v=1\)\hfill\hyperref[sec:vi-indice-rutas]{Índice}\par\vspace{4pt}
\begingroup\fontsize{10.2}{12.5}\selectfont\setlength{\tabcolsep}{5pt}\renewcommand{\arraystretch}{1.1}\rowcolors{1}{routepale}{white}\noindent\begin{tabularx}{\linewidth}{@{}>{\raggedright\arraybackslash}p{.345\linewidth}>{\raggedright\arraybackslash}X@{}}
\RouteGroup{Identidad estructural}
\RouteRow{Clase y multisección}{quark\_\allowbreak{}down\_\allowbreak{}G3\quad /\quad quark\_\allowbreak{}down\_\allowbreak{}G3}
\RouteRow{Colección y sector}{quark de tipo abajo; quark}
\RouteRow{Clasificación física}{quark de tipo abajo}
\RouteRow{Generación, profundidad y color}{3; G3; R}
\RouteRow{Ruta primaria y órbita de inversión}{\(B_1\): no\qquad 2,\allowbreak{}5;\allowbreak{}8,\allowbreak{}5}
\RouteRow{Firma de clasificación}{30|\allowbreak{}6|\allowbreak{}-4|\allowbreak{}0|\allowbreak{}-12|\allowbreak{}+++-\kern0pt{}-\kern0pt{}-+++-\kern0pt{}-\kern0pt{}-}
\RouteGroup{Retornos, memoria y transporte}
\RouteRow{Retornos con fase}{clase: 27; órbita: 27; celda: 27}
\RouteRow{Retorno cinemático}{en 108: 54; extensión: 216}
\RouteRow{Giros y cambios de hoja}{71; 107}
\RouteRow{Acarreo y diferimiento}{deuda total: 2; final: 1; bloques diferidos: 2}
\RouteRow{Tríadas finales; persistencia de Pareto}{16; sí}
\RouteRow{\(K\longrightarrow K+9\)}{repeticiones: 9; cambios de memoria: 9}
\RouteGroup{Firma, fase y diferencias}
\RouteRow{\(n_A,\ s_C,\ \kappa_0,\ \nu_{120},\ \nu_{270}\)}{30, 6, 0, 0, -12}
\RouteRow{\(\tau_{12}\); firma histórica \(\mathbb Z^5\)}{+++-\kern0pt{}-\kern0pt{}-+++-\kern0pt{}-\kern0pt{}-; 30|\allowbreak{}6|\allowbreak{}0|\allowbreak{}0|\allowbreak{}-12}
\RouteRow{\(D_1,\ D_3,\ D_4,\ D_{27},\ D_{36}\)}{54, 56, 68, 48, 32}
\RouteRow{\(\Omega_{90/120},\ P_6\)}{80, 6}
\RouteGroup{Lectores y factores de acción}
\RouteRow{\(K_D\quad /\quad K_\Omega\)}{(80,\allowbreak{}54,\allowbreak{}6)\quad /\quad (80,\allowbreak{}54,\allowbreak{}6)}
\RouteRow{\(\kappa_D\)}{79.60661\allowbreak{}01397948\allowbreak{}27586470\allowbreak{}91}
\RouteRow{\(\kappa_\Omega\)}{79.60661\allowbreak{}01397948\allowbreak{}27586470\allowbreak{}91}
\RouteRow{\(R_{\mathrm{act},D}\)}{1.000000\allowbreak{}05518981\allowbreak{}25318591\allowbreak{}40}
\RouteRow{\(R_{\mathrm{act},\Omega}\)}{1.000000\allowbreak{}05518981\allowbreak{}25318591\allowbreak{}40}
\RouteGroup{Separación de prefijos}
\RouteRow{Área de ambigüedad; área \(\log_2\)}{36; 18.0}
\RouteRow{Primer bloque de separación}{órbita: 1; celda: \textemdash{}; ruta: \textemdash{}}
\RouteRow{Ambigüedad final}{2}
\RouteGroup{Secuencias completas}
\RouteRow{\(w_6\)}{243 657 24 243 663 486 423 15 486 243 657 24 243 663 486 423 15 486}
\RouteRow{Tríadas estables por bloque}{0 1 2 3 4 5 6 7 8 9 10 11 12 13 14 14 16 16}
\end{tabularx}\endgroup\par\vspace{5pt}\noindent{\fontsize{8.4}{10}\selectfont\hyperref[trace-272]{Procedencia y códigos originales}\hfill Ficha completa · 272/324}
\clearpage\phantomsection\label{nav:vi-ruta-273}
% VI-CATALOG-ROW rutas 273
\pdfbookmark[2]{Celda 69}{cell-69}
\pdfbookmark[3]{r8c6\_h-1v-1}{route-273}
\RouteTitle{Ruta 273}{r8c6\_\allowbreak{}h-1v-1}
\noindent Celda 69\quad (8, 6)\qquad Orientación \(h=-1,\ v=-1\)\hfill\hyperref[sec:vi-indice-rutas]{Índice}\par\vspace{4pt}
\begingroup\fontsize{10.2}{12.5}\selectfont\setlength{\tabcolsep}{5pt}\renewcommand{\arraystretch}{1.1}\rowcolors{1}{routepale}{white}\noindent\begin{tabularx}{\linewidth}{@{}>{\raggedright\arraybackslash}p{.345\linewidth}>{\raggedright\arraybackslash}X@{}}
\RouteGroup{Identidad estructural}
\RouteRow{Clase y multisección}{quark\_\allowbreak{}up\_\allowbreak{}G3\quad /\quad quark\_\allowbreak{}up\_\allowbreak{}G3}
\RouteRow{Colección y sector}{quark de tipo arriba; quark}
\RouteRow{Clasificación física}{quark de tipo arriba}
\RouteRow{Generación, profundidad y color}{3; G3; B}
\RouteRow{Ruta primaria y órbita de inversión}{\(B_1\): sí\qquad 2,\allowbreak{}4;\allowbreak{}8,\allowbreak{}6}
\RouteRow{Firma de clasificación}{28|\allowbreak{}6|\allowbreak{}-6|\allowbreak{}0|\allowbreak{}-12|\allowbreak{}+++-\kern0pt{}-\kern0pt{}-+++-\kern0pt{}-\kern0pt{}-}
\RouteGroup{Retornos, memoria y transporte}
\RouteRow{Retornos con fase}{clase: 27; órbita: 27; celda: 27}
\RouteRow{Retorno cinemático}{en 108: 54; extensión: 216}
\RouteRow{Giros y cambios de hoja}{71; 107}
\RouteRow{Acarreo y diferimiento}{deuda total: 1; final: 0; bloques diferidos: 1}
\RouteRow{Tríadas finales; persistencia de Pareto}{17; no}
\RouteRow{\(K\longrightarrow K+9\)}{repeticiones: 9; cambios de memoria: 9}
\RouteGroup{Firma, fase y diferencias}
\RouteRow{\(n_A,\ s_C,\ \kappa_0,\ \nu_{120},\ \nu_{270}\)}{28, -6, 0, -2, -12}
\RouteRow{\(\tau_{12}\); firma histórica \(\mathbb Z^5\)}{+++-\kern0pt{}-\kern0pt{}-+++-\kern0pt{}-\kern0pt{}-; 28|\allowbreak{}-6|\allowbreak{}0|\allowbreak{}-2|\allowbreak{}-12}
\RouteRow{\(D_1,\ D_3,\ D_4,\ D_{27},\ D_{36}\)}{66, 60, 68, 48, 32}
\RouteRow{\(\Omega_{90/120},\ P_6\)}{80, 6}
\RouteGroup{Lectores y factores de acción}
\RouteRow{\(K_D\quad /\quad K_\Omega\)}{(80,\allowbreak{}66,\allowbreak{}4)\quad /\quad (80,\allowbreak{}54,\allowbreak{}6)}
\RouteRow{\(\kappa_D\)}{79.51881\allowbreak{}95161180\allowbreak{}37694395\allowbreak{}39}
\RouteRow{\(\kappa_\Omega\)}{79.60661\allowbreak{}01397948\allowbreak{}27586470\allowbreak{}91}
\RouteRow{\(R_{\mathrm{act},D}\)}{1.000000\allowbreak{}05512894\allowbreak{}88901612\allowbreak{}64}
\RouteRow{\(R_{\mathrm{act},\Omega}\)}{1.000000\allowbreak{}05518981\allowbreak{}25318591\allowbreak{}40}
\RouteGroup{Separación de prefijos}
\RouteRow{Área de ambigüedad; área \(\log_2\)}{36; 18.0}
\RouteRow{Primer bloque de separación}{órbita: \textemdash{}; celda: \textemdash{}; ruta: \textemdash{}}
\RouteRow{Ambigüedad final}{2}
\RouteGroup{Secuencias completas}
\RouteRow{\(w_6\)}{570 264 90 570 267 243 657 24 243 570 264 90 570 267 243 657 24 243}
\RouteRow{Tríadas estables por bloque}{0 1 1 3 4 5 6 7 8 9 10 11 12 13 14 15 16 17}
\end{tabularx}\endgroup\par\vspace{5pt}\noindent{\fontsize{8.4}{10}\selectfont\hyperref[trace-273]{Procedencia y códigos originales}\hfill Ficha completa · 273/324}
\clearpage\phantomsection\label{nav:vi-ruta-274}
% VI-CATALOG-ROW rutas 274
\pdfbookmark[3]{r8c6\_h-1v+1}{route-274}
\RouteTitle{Ruta 274}{r8c6\_\allowbreak{}h-1v+1}
\noindent Celda 69\quad (8, 6)\qquad Orientación \(h=-1,\ v=1\)\hfill\hyperref[sec:vi-indice-rutas]{Índice}\par\vspace{4pt}
\begingroup\fontsize{10.2}{12.5}\selectfont\setlength{\tabcolsep}{5pt}\renewcommand{\arraystretch}{1.1}\rowcolors{1}{routepale}{white}\noindent\begin{tabularx}{\linewidth}{@{}>{\raggedright\arraybackslash}p{.345\linewidth}>{\raggedright\arraybackslash}X@{}}
\RouteGroup{Identidad estructural}
\RouteRow{Clase y multisección}{quark\_\allowbreak{}up\_\allowbreak{}G3\quad /\quad quark\_\allowbreak{}up\_\allowbreak{}G3}
\RouteRow{Colección y sector}{quark de tipo arriba; quark}
\RouteRow{Clasificación física}{quark de tipo arriba}
\RouteRow{Generación, profundidad y color}{3; G3; B}
\RouteRow{Ruta primaria y órbita de inversión}{\(B_1\): no\qquad 2,\allowbreak{}4;\allowbreak{}8,\allowbreak{}6}
\RouteRow{Firma de clasificación}{28|\allowbreak{}6|\allowbreak{}-6|\allowbreak{}0|\allowbreak{}-12|\allowbreak{}+++-\kern0pt{}-\kern0pt{}-+++-\kern0pt{}-\kern0pt{}-}
\RouteGroup{Retornos, memoria y transporte}
\RouteRow{Retornos con fase}{clase: 27; órbita: 27; celda: 27}
\RouteRow{Retorno cinemático}{en 108: 54; extensión: 216}
\RouteRow{Giros y cambios de hoja}{71; 107}
\RouteRow{Acarreo y diferimiento}{deuda total: 1; final: 0; bloques diferidos: 1}
\RouteRow{Tríadas finales; persistencia de Pareto}{17; no}
\RouteRow{\(K\longrightarrow K+9\)}{repeticiones: 9; cambios de memoria: 9}
\RouteGroup{Firma, fase y diferencias}
\RouteRow{\(n_A,\ s_C,\ \kappa_0,\ \nu_{120},\ \nu_{270}\)}{44, 2, 0, 0, -12}
\RouteRow{\(\tau_{12}\); firma histórica \(\mathbb Z^5\)}{+++-\kern0pt{}-\kern0pt{}-+++-\kern0pt{}-\kern0pt{}-; 44|\allowbreak{}2|\allowbreak{}0|\allowbreak{}0|\allowbreak{}-12}
\RouteRow{\(D_1,\ D_3,\ D_4,\ D_{27},\ D_{36}\)}{78, 24, 78, 24, 16}
\RouteRow{\(\Omega_{90/120},\ P_6\)}{80, 6}
\RouteGroup{Lectores y factores de acción}
\RouteRow{\(K_D\quad /\quad K_\Omega\)}{(40,\allowbreak{}78,\allowbreak{}27)\quad /\quad (80,\allowbreak{}54,\allowbreak{}6)}
\RouteRow{\(\kappa_D\)}{39.43380\allowbreak{}88030085\allowbreak{}51303671\allowbreak{}14}
\RouteRow{\(\kappa_\Omega\)}{79.60661\allowbreak{}01397948\allowbreak{}27586470\allowbreak{}91}
\RouteRow{\(R_{\mathrm{act},D}\)}{1.000000\allowbreak{}02733874\allowbreak{}08548943\allowbreak{}58}
\RouteRow{\(R_{\mathrm{act},\Omega}\)}{1.000000\allowbreak{}05518981\allowbreak{}25318591\allowbreak{}40}
\RouteGroup{Separación de prefijos}
\RouteRow{Área de ambigüedad; área \(\log_2\)}{36; 18.0}
\RouteRow{Primer bloque de separación}{órbita: \textemdash{}; celda: \textemdash{}; ruta: \textemdash{}}
\RouteRow{Ambigüedad final}{2}
\RouteGroup{Secuencias completas}
\RouteRow{\(w_6\)}{585 507 504 585 510 666 672 510 666 585 507 504 585 510 666 672 510 666}
\RouteRow{Tríadas estables por bloque}{0 1 2 3 4 5 6 7 8 9 10 11 12 13 14 14 16 17}
\end{tabularx}\endgroup\par\vspace{5pt}\noindent{\fontsize{8.4}{10}\selectfont\hyperref[trace-274]{Procedencia y códigos originales}\hfill Ficha completa · 274/324}
\clearpage\phantomsection\label{nav:vi-ruta-275}
% VI-CATALOG-ROW rutas 275
\pdfbookmark[3]{r8c6\_h+1v-1}{route-275}
\RouteTitle{Ruta 275}{r8c6\_\allowbreak{}h+1v-1}
\noindent Celda 69\quad (8, 6)\qquad Orientación \(h=1,\ v=-1\)\hfill\hyperref[sec:vi-indice-rutas]{Índice}\par\vspace{4pt}
\begingroup\fontsize{10.2}{12.5}\selectfont\setlength{\tabcolsep}{5pt}\renewcommand{\arraystretch}{1.1}\rowcolors{1}{routepale}{white}\noindent\begin{tabularx}{\linewidth}{@{}>{\raggedright\arraybackslash}p{.345\linewidth}>{\raggedright\arraybackslash}X@{}}
\RouteGroup{Identidad estructural}
\RouteRow{Clase y multisección}{quark\_\allowbreak{}up\_\allowbreak{}G3\quad /\quad quark\_\allowbreak{}up\_\allowbreak{}G3}
\RouteRow{Colección y sector}{quark de tipo arriba; quark}
\RouteRow{Clasificación física}{quark de tipo arriba}
\RouteRow{Generación, profundidad y color}{3; G3; B}
\RouteRow{Ruta primaria y órbita de inversión}{\(B_1\): no\qquad 2,\allowbreak{}4;\allowbreak{}8,\allowbreak{}6}
\RouteRow{Firma de clasificación}{28|\allowbreak{}6|\allowbreak{}-6|\allowbreak{}0|\allowbreak{}-12|\allowbreak{}+++-\kern0pt{}-\kern0pt{}-+++-\kern0pt{}-\kern0pt{}-}
\RouteGroup{Retornos, memoria y transporte}
\RouteRow{Retornos con fase}{clase: 27; órbita: 27; celda: 27}
\RouteRow{Retorno cinemático}{en 108: 54; extensión: 216}
\RouteRow{Giros y cambios de hoja}{71; 107}
\RouteRow{Acarreo y diferimiento}{deuda total: 1; final: 0; bloques diferidos: 1}
\RouteRow{Tríadas finales; persistencia de Pareto}{17; no}
\RouteRow{\(K\longrightarrow K+9\)}{repeticiones: 9; cambios de memoria: 9}
\RouteGroup{Firma, fase y diferencias}
\RouteRow{\(n_A,\ s_C,\ \kappa_0,\ \nu_{120},\ \nu_{270}\)}{44, -2, 0, 0, -12}
\RouteRow{\(\tau_{12}\); firma histórica \(\mathbb Z^5\)}{+++-\kern0pt{}-\kern0pt{}-+++-\kern0pt{}-\kern0pt{}-; 44|\allowbreak{}-2|\allowbreak{}0|\allowbreak{}0|\allowbreak{}-12}
\RouteRow{\(D_1,\ D_3,\ D_4,\ D_{27},\ D_{36}\)}{78, 24, 78, 24, 16}
\RouteRow{\(\Omega_{90/120},\ P_6\)}{80, 6}
\RouteGroup{Lectores y factores de acción}
\RouteRow{\(K_D\quad /\quad K_\Omega\)}{(40,\allowbreak{}78,\allowbreak{}27)\quad /\quad (80,\allowbreak{}54,\allowbreak{}6)}
\RouteRow{\(\kappa_D\)}{39.43380\allowbreak{}88030085\allowbreak{}51303671\allowbreak{}14}
\RouteRow{\(\kappa_\Omega\)}{79.60661\allowbreak{}01397948\allowbreak{}27586470\allowbreak{}91}
\RouteRow{\(R_{\mathrm{act},D}\)}{1.000000\allowbreak{}02733874\allowbreak{}08548943\allowbreak{}58}
\RouteRow{\(R_{\mathrm{act},\Omega}\)}{1.000000\allowbreak{}05518981\allowbreak{}25318591\allowbreak{}40}
\RouteGroup{Separación de prefijos}
\RouteRow{Área de ambigüedad; área \(\log_2\)}{36; 18.0}
\RouteRow{Primer bloque de separación}{órbita: \textemdash{}; celda: \textemdash{}; ruta: \textemdash{}}
\RouteRow{Ambigüedad final}{2}
\RouteGroup{Secuencias completas}
\RouteRow{\(w_6\)}{672 510 666 672 507 504 585 507 504 672 510 666 672 507 504 585 507 504}
\RouteRow{Tríadas estables por bloque}{0 1 2 3 4 5 6 7 8 9 10 11 12 13 14 14 16 17}
\end{tabularx}\endgroup\par\vspace{5pt}\noindent{\fontsize{8.4}{10}\selectfont\hyperref[trace-275]{Procedencia y códigos originales}\hfill Ficha completa · 275/324}
\clearpage\phantomsection\label{nav:vi-ruta-276}
% VI-CATALOG-ROW rutas 276
\pdfbookmark[3]{r8c6\_h+1v+1}{route-276}
\RouteTitle{Ruta 276}{r8c6\_\allowbreak{}h+1v+1}
\noindent Celda 69\quad (8, 6)\qquad Orientación \(h=1,\ v=1\)\hfill\hyperref[sec:vi-indice-rutas]{Índice}\par\vspace{4pt}
\begingroup\fontsize{10.2}{12.5}\selectfont\setlength{\tabcolsep}{5pt}\renewcommand{\arraystretch}{1.1}\rowcolors{1}{routepale}{white}\noindent\begin{tabularx}{\linewidth}{@{}>{\raggedright\arraybackslash}p{.345\linewidth}>{\raggedright\arraybackslash}X@{}}
\RouteGroup{Identidad estructural}
\RouteRow{Clase y multisección}{quark\_\allowbreak{}up\_\allowbreak{}G3\quad /\quad quark\_\allowbreak{}up\_\allowbreak{}G3}
\RouteRow{Colección y sector}{quark de tipo arriba; quark}
\RouteRow{Clasificación física}{quark de tipo arriba}
\RouteRow{Generación, profundidad y color}{3; G3; B}
\RouteRow{Ruta primaria y órbita de inversión}{\(B_1\): no\qquad 2,\allowbreak{}4;\allowbreak{}8,\allowbreak{}6}
\RouteRow{Firma de clasificación}{28|\allowbreak{}6|\allowbreak{}-6|\allowbreak{}0|\allowbreak{}-12|\allowbreak{}+++-\kern0pt{}-\kern0pt{}-+++-\kern0pt{}-\kern0pt{}-}
\RouteGroup{Retornos, memoria y transporte}
\RouteRow{Retornos con fase}{clase: 27; órbita: 27; celda: 27}
\RouteRow{Retorno cinemático}{en 108: 54; extensión: 216}
\RouteRow{Giros y cambios de hoja}{71; 107}
\RouteRow{Acarreo y diferimiento}{deuda total: 0; final: 0; bloques diferidos: 0}
\RouteRow{Tríadas finales; persistencia de Pareto}{17; no}
\RouteRow{\(K\longrightarrow K+9\)}{repeticiones: 9; cambios de memoria: 9}
\RouteGroup{Firma, fase y diferencias}
\RouteRow{\(n_A,\ s_C,\ \kappa_0,\ \nu_{120},\ \nu_{270}\)}{28, 6, 0, 2, -12}
\RouteRow{\(\tau_{12}\); firma histórica \(\mathbb Z^5\)}{+++-\kern0pt{}-\kern0pt{}-+++-\kern0pt{}-\kern0pt{}-; 28|\allowbreak{}6|\allowbreak{}0|\allowbreak{}2|\allowbreak{}-12}
\RouteRow{\(D_1,\ D_3,\ D_4,\ D_{27},\ D_{36}\)}{66, 60, 68, 48, 32}
\RouteRow{\(\Omega_{90/120},\ P_6\)}{80, 6}
\RouteGroup{Lectores y factores de acción}
\RouteRow{\(K_D\quad /\quad K_\Omega\)}{(80,\allowbreak{}66,\allowbreak{}4)\quad /\quad (80,\allowbreak{}54,\allowbreak{}6)}
\RouteRow{\(\kappa_D\)}{79.51881\allowbreak{}95161180\allowbreak{}37694395\allowbreak{}39}
\RouteRow{\(\kappa_\Omega\)}{79.60661\allowbreak{}01397948\allowbreak{}27586470\allowbreak{}91}
\RouteRow{\(R_{\mathrm{act},D}\)}{1.000000\allowbreak{}05512894\allowbreak{}88901612\allowbreak{}64}
\RouteRow{\(R_{\mathrm{act},\Omega}\)}{1.000000\allowbreak{}05518981\allowbreak{}25318591\allowbreak{}40}
\RouteGroup{Separación de prefijos}
\RouteRow{Área de ambigüedad; área \(\log_2\)}{36; 18.0}
\RouteRow{Primer bloque de separación}{órbita: \textemdash{}; celda: \textemdash{}; ruta: \textemdash{}}
\RouteRow{Ambigüedad final}{2}
\RouteGroup{Secuencias completas}
\RouteRow{\(w_6\)}{657 24 243 657 21 90 570 264 90 657 24 243 657 21 90 570 264 90}
\RouteRow{Tríadas estables por bloque}{0 1 2 3 4 5 6 7 8 9 10 11 12 13 14 15 16 17}
\end{tabularx}\endgroup\par\vspace{5pt}\noindent{\fontsize{8.4}{10}\selectfont\hyperref[trace-276]{Procedencia y códigos originales}\hfill Ficha completa · 276/324}
\clearpage\phantomsection\label{nav:vi-ruta-277}
% VI-CATALOG-ROW rutas 277
\pdfbookmark[2]{Celda 70}{cell-70}
\pdfbookmark[3]{r8c7\_h-1v-1}{route-277}
\RouteTitle{Ruta 277}{r8c7\_\allowbreak{}h-1v-1}
\noindent Celda 70\quad (8, 7)\qquad Orientación \(h=-1,\ v=-1\)\hfill\hyperref[sec:vi-indice-rutas]{Índice}\par\vspace{4pt}
\begingroup\fontsize{10.2}{12.5}\selectfont\setlength{\tabcolsep}{5pt}\renewcommand{\arraystretch}{1.1}\rowcolors{1}{routepale}{white}\noindent\begin{tabularx}{\linewidth}{@{}>{\raggedright\arraybackslash}p{.345\linewidth}>{\raggedright\arraybackslash}X@{}}
\RouteGroup{Identidad estructural}
\RouteRow{Clase y multisección}{quark\_\allowbreak{}down\_\allowbreak{}G3\quad /\quad quark\_\allowbreak{}down\_\allowbreak{}G3}
\RouteRow{Colección y sector}{quark de tipo abajo; quark}
\RouteRow{Clasificación física}{quark de tipo abajo}
\RouteRow{Generación, profundidad y color}{3; G3; G}
\RouteRow{Ruta primaria y órbita de inversión}{\(B_1\): no\qquad 2,\allowbreak{}3;\allowbreak{}8,\allowbreak{}7}
\RouteRow{Firma de clasificación}{28|\allowbreak{}-6|\allowbreak{}0|\allowbreak{}0|\allowbreak{}-12|\allowbreak{}+++-\kern0pt{}-\kern0pt{}-+++-\kern0pt{}-\kern0pt{}-}
\RouteGroup{Retornos, memoria y transporte}
\RouteRow{Retornos con fase}{clase: 27; órbita: 27; celda: 27}
\RouteRow{Retorno cinemático}{en 108: 54; extensión: 216}
\RouteRow{Giros y cambios de hoja}{71; 107}
\RouteRow{Acarreo y diferimiento}{deuda total: 1; final: 1; bloques diferidos: 1}
\RouteRow{Tríadas finales; persistencia de Pareto}{16; sí}
\RouteRow{\(K\longrightarrow K+9\)}{repeticiones: 9; cambios de memoria: 9}
\RouteGroup{Firma, fase y diferencias}
\RouteRow{\(n_A,\ s_C,\ \kappa_0,\ \nu_{120},\ \nu_{270}\)}{28, 6, 0, 0, -12}
\RouteRow{\(\tau_{12}\); firma histórica \(\mathbb Z^5\)}{+++-\kern0pt{}-\kern0pt{}-+++-\kern0pt{}-\kern0pt{}-; 28|\allowbreak{}6|\allowbreak{}0|\allowbreak{}0|\allowbreak{}-12}
\RouteRow{\(D_1,\ D_3,\ D_4,\ D_{27},\ D_{36}\)}{66, 64, 68, 60, 40}
\RouteRow{\(\Omega_{90/120},\ P_6\)}{80, 6}
\RouteGroup{Lectores y factores de acción}
\RouteRow{\(K_D\quad /\quad K_\Omega\)}{(100,\allowbreak{}66,\allowbreak{}2)\quad /\quad (80,\allowbreak{}54,\allowbreak{}6)}
\RouteRow{\(\kappa_D\)}{99.51859\allowbreak{}71232726\allowbreak{}53414287\allowbreak{}29}
\RouteRow{\(\kappa_\Omega\)}{79.60661\allowbreak{}01397948\allowbreak{}27586470\allowbreak{}91}
\RouteRow{\(R_{\mathrm{act},D}\)}{1.000000\allowbreak{}06899443\allowbreak{}08259364\allowbreak{}17}
\RouteRow{\(R_{\mathrm{act},\Omega}\)}{1.000000\allowbreak{}05518981\allowbreak{}25318591\allowbreak{}40}
\RouteGroup{Separación de prefijos}
\RouteRow{Área de ambigüedad; área \(\log_2\)}{36; 18.0}
\RouteRow{Primer bloque de separación}{órbita: \textemdash{}; celda: \textemdash{}; ruta: \textemdash{}}
\RouteRow{Ambigüedad final}{2}
\RouteGroup{Secuencias completas}
\RouteRow{\(w_6\)}{15 486 423 15 489 498 99 327 498 15 486 423 15 489 498 99 327 498}
\RouteRow{Tríadas estables por bloque}{0 1 2 3 4 5 6 7 8 9 10 11 12 13 14 15 16 16}
\end{tabularx}\endgroup\par\vspace{5pt}\noindent{\fontsize{8.4}{10}\selectfont\hyperref[trace-277]{Procedencia y códigos originales}\hfill Ficha completa · 277/324}
\clearpage\phantomsection\label{nav:vi-ruta-278}
% VI-CATALOG-ROW rutas 278
\pdfbookmark[3]{r8c7\_h-1v+1}{route-278}
\RouteTitle{Ruta 278}{r8c7\_\allowbreak{}h-1v+1}
\noindent Celda 70\quad (8, 7)\qquad Orientación \(h=-1,\ v=1\)\hfill\hyperref[sec:vi-indice-rutas]{Índice}\par\vspace{4pt}
\begingroup\fontsize{10.2}{12.5}\selectfont\setlength{\tabcolsep}{5pt}\renewcommand{\arraystretch}{1.1}\rowcolors{1}{routepale}{white}\noindent\begin{tabularx}{\linewidth}{@{}>{\raggedright\arraybackslash}p{.345\linewidth}>{\raggedright\arraybackslash}X@{}}
\RouteGroup{Identidad estructural}
\RouteRow{Clase y multisección}{quark\_\allowbreak{}down\_\allowbreak{}G3\quad /\quad quark\_\allowbreak{}down\_\allowbreak{}G3}
\RouteRow{Colección y sector}{quark de tipo abajo; quark}
\RouteRow{Clasificación física}{quark de tipo abajo}
\RouteRow{Generación, profundidad y color}{3; G3; G}
\RouteRow{Ruta primaria y órbita de inversión}{\(B_1\): no\qquad 2,\allowbreak{}3;\allowbreak{}8,\allowbreak{}7}
\RouteRow{Firma de clasificación}{28|\allowbreak{}-6|\allowbreak{}0|\allowbreak{}0|\allowbreak{}-12|\allowbreak{}+++-\kern0pt{}-\kern0pt{}-+++-\kern0pt{}-\kern0pt{}-}
\RouteGroup{Retornos, memoria y transporte}
\RouteRow{Retornos con fase}{clase: 27; órbita: 27; celda: 27}
\RouteRow{Retorno cinemático}{en 108: 54; extensión: 216}
\RouteRow{Giros y cambios de hoja}{71; 107}
\RouteRow{Acarreo y diferimiento}{deuda total: 0; final: 0; bloques diferidos: 0}
\RouteRow{Tríadas finales; persistencia de Pareto}{17; sí}
\RouteRow{\(K\longrightarrow K+9\)}{repeticiones: 9; cambios de memoria: 9}
\RouteGroup{Firma, fase y diferencias}
\RouteRow{\(n_A,\ s_C,\ \kappa_0,\ \nu_{120},\ \nu_{270}\)}{14, 6, 0, 0, -12}
\RouteRow{\(\tau_{12}\); firma histórica \(\mathbb Z^5\)}{+++-\kern0pt{}-\kern0pt{}-+++-\kern0pt{}-\kern0pt{}-; 14|\allowbreak{}6|\allowbreak{}0|\allowbreak{}0|\allowbreak{}-12}
\RouteRow{\(D_1,\ D_3,\ D_4,\ D_{27},\ D_{36}\)}{24, 20, 24, 24, 16}
\RouteRow{\(\Omega_{90/120},\ P_6\)}{80, 6}
\RouteGroup{Lectores y factores de acción}
\RouteRow{\(K_D\quad /\quad K_\Omega\)}{(40,\allowbreak{}24,\allowbreak{}2)\quad /\quad (80,\allowbreak{}54,\allowbreak{}6)}
\RouteRow{\(\kappa_D\)}{39.82508\allowbreak{}59311825\allowbreak{}73056173\allowbreak{}26}
\RouteRow{\(\kappa_\Omega\)}{79.60661\allowbreak{}01397948\allowbreak{}27586470\allowbreak{}91}
\RouteRow{\(R_{\mathrm{act},D}\)}{1.000000\allowbreak{}02761000\allowbreak{}61595142\allowbreak{}15}
\RouteRow{\(R_{\mathrm{act},\Omega}\)}{1.000000\allowbreak{}05518981\allowbreak{}25318591\allowbreak{}40}
\RouteGroup{Separación de prefijos}
\RouteRow{Área de ambigüedad; área \(\log_2\)}{36; 18.0}
\RouteRow{Primer bloque de separación}{órbita: \textemdash{}; celda: \textemdash{}; ruta: \textemdash{}}
\RouteRow{Ambigüedad final}{2}
\RouteGroup{Secuencias completas}
\RouteRow{\(w_6\)}{0 81 3 0 84 0 81 3 0 0 81 3 0 84 0 81 3 0}
\RouteRow{Tríadas estables por bloque}{0 1 2 3 4 5 6 7 8 9 10 11 12 13 14 15 16 17}
\end{tabularx}\endgroup\par\vspace{5pt}\noindent{\fontsize{8.4}{10}\selectfont\hyperref[trace-278]{Procedencia y códigos originales}\hfill Ficha completa · 278/324}
\clearpage\phantomsection\label{nav:vi-ruta-279}
% VI-CATALOG-ROW rutas 279
\pdfbookmark[3]{r8c7\_h+1v-1}{route-279}
\RouteTitle{Ruta 279}{r8c7\_\allowbreak{}h+1v-1}
\noindent Celda 70\quad (8, 7)\qquad Orientación \(h=1,\ v=-1\)\hfill\hyperref[sec:vi-indice-rutas]{Índice}\par\vspace{4pt}
\begingroup\fontsize{10.2}{12.5}\selectfont\setlength{\tabcolsep}{5pt}\renewcommand{\arraystretch}{1.1}\rowcolors{1}{routepale}{white}\noindent\begin{tabularx}{\linewidth}{@{}>{\raggedright\arraybackslash}p{.345\linewidth}>{\raggedright\arraybackslash}X@{}}
\RouteGroup{Identidad estructural}
\RouteRow{Clase y multisección}{quark\_\allowbreak{}down\_\allowbreak{}G3\quad /\quad quark\_\allowbreak{}down\_\allowbreak{}G3}
\RouteRow{Colección y sector}{quark de tipo abajo; quark}
\RouteRow{Clasificación física}{quark de tipo abajo}
\RouteRow{Generación, profundidad y color}{3; G3; G}
\RouteRow{Ruta primaria y órbita de inversión}{\(B_1\): sí\qquad 2,\allowbreak{}3;\allowbreak{}8,\allowbreak{}7}
\RouteRow{Firma de clasificación}{28|\allowbreak{}-6|\allowbreak{}0|\allowbreak{}0|\allowbreak{}-12|\allowbreak{}+++-\kern0pt{}-\kern0pt{}-+++-\kern0pt{}-\kern0pt{}-}
\RouteGroup{Retornos, memoria y transporte}
\RouteRow{Retornos con fase}{clase: 27; órbita: 27; celda: 27}
\RouteRow{Retorno cinemático}{en 108: 54; extensión: 216}
\RouteRow{Giros y cambios de hoja}{71; 107}
\RouteRow{Acarreo y diferimiento}{deuda total: 2; final: 0; bloques diferidos: 2}
\RouteRow{Tríadas finales; persistencia de Pareto}{17; sí}
\RouteRow{\(K\longrightarrow K+9\)}{repeticiones: 9; cambios de memoria: 9}
\RouteGroup{Firma, fase y diferencias}
\RouteRow{\(n_A,\ s_C,\ \kappa_0,\ \nu_{120},\ \nu_{270}\)}{14, -6, 0, 0, -12}
\RouteRow{\(\tau_{12}\); firma histórica \(\mathbb Z^5\)}{+++-\kern0pt{}-\kern0pt{}-+++-\kern0pt{}-\kern0pt{}-; 14|\allowbreak{}-6|\allowbreak{}0|\allowbreak{}0|\allowbreak{}-12}
\RouteRow{\(D_1,\ D_3,\ D_4,\ D_{27},\ D_{36}\)}{24, 20, 24, 24, 16}
\RouteRow{\(\Omega_{90/120},\ P_6\)}{80, 6}
\RouteGroup{Lectores y factores de acción}
\RouteRow{\(K_D\quad /\quad K_\Omega\)}{(40,\allowbreak{}24,\allowbreak{}2)\quad /\quad (80,\allowbreak{}54,\allowbreak{}6)}
\RouteRow{\(\kappa_D\)}{39.82508\allowbreak{}59311825\allowbreak{}73056173\allowbreak{}26}
\RouteRow{\(\kappa_\Omega\)}{79.60661\allowbreak{}01397948\allowbreak{}27586470\allowbreak{}91}
\RouteRow{\(R_{\mathrm{act},D}\)}{1.000000\allowbreak{}02761000\allowbreak{}61595142\allowbreak{}15}
\RouteRow{\(R_{\mathrm{act},\Omega}\)}{1.000000\allowbreak{}05518981\allowbreak{}25318591\allowbreak{}40}
\RouteGroup{Separación de prefijos}
\RouteRow{Área de ambigüedad; área \(\log_2\)}{36; 18.0}
\RouteRow{Primer bloque de separación}{órbita: \textemdash{}; celda: \textemdash{}; ruta: \textemdash{}}
\RouteRow{Ambigüedad final}{2}
\RouteGroup{Secuencias completas}
\RouteRow{\(w_6\)}{81 3 0 81 0 3 0 81 3 81 3 0 81 0 3 0 81 3}
\RouteRow{Tríadas estables por bloque}{0 1 2 3 4 5 6 6 8 9 10 11 12 13 14 15 15 17}
\end{tabularx}\endgroup\par\vspace{5pt}\noindent{\fontsize{8.4}{10}\selectfont\hyperref[trace-279]{Procedencia y códigos originales}\hfill Ficha completa · 279/324}
\clearpage\phantomsection\label{nav:vi-ruta-280}
% VI-CATALOG-ROW rutas 280
\pdfbookmark[3]{r8c7\_h+1v+1}{route-280}
\RouteTitle{Ruta 280}{r8c7\_\allowbreak{}h+1v+1}
\noindent Celda 70\quad (8, 7)\qquad Orientación \(h=1,\ v=1\)\hfill\hyperref[sec:vi-indice-rutas]{Índice}\par\vspace{4pt}
\begingroup\fontsize{10.2}{12.5}\selectfont\setlength{\tabcolsep}{5pt}\renewcommand{\arraystretch}{1.1}\rowcolors{1}{routepale}{white}\noindent\begin{tabularx}{\linewidth}{@{}>{\raggedright\arraybackslash}p{.345\linewidth}>{\raggedright\arraybackslash}X@{}}
\RouteGroup{Identidad estructural}
\RouteRow{Clase y multisección}{quark\_\allowbreak{}down\_\allowbreak{}G3\quad /\quad quark\_\allowbreak{}down\_\allowbreak{}G3}
\RouteRow{Colección y sector}{quark de tipo abajo; quark}
\RouteRow{Clasificación física}{quark de tipo abajo}
\RouteRow{Generación, profundidad y color}{3; G3; G}
\RouteRow{Ruta primaria y órbita de inversión}{\(B_1\): no\qquad 2,\allowbreak{}3;\allowbreak{}8,\allowbreak{}7}
\RouteRow{Firma de clasificación}{28|\allowbreak{}-6|\allowbreak{}0|\allowbreak{}0|\allowbreak{}-12|\allowbreak{}+++-\kern0pt{}-\kern0pt{}-+++-\kern0pt{}-\kern0pt{}-}
\RouteGroup{Retornos, memoria y transporte}
\RouteRow{Retornos con fase}{clase: 27; órbita: 27; celda: 27}
\RouteRow{Retorno cinemático}{en 108: 54; extensión: 216}
\RouteRow{Giros y cambios de hoja}{71; 107}
\RouteRow{Acarreo y diferimiento}{deuda total: 1; final: 0; bloques diferidos: 1}
\RouteRow{Tríadas finales; persistencia de Pareto}{17; sí}
\RouteRow{\(K\longrightarrow K+9\)}{repeticiones: 9; cambios de memoria: 9}
\RouteGroup{Firma, fase y diferencias}
\RouteRow{\(n_A,\ s_C,\ \kappa_0,\ \nu_{120},\ \nu_{270}\)}{28, -6, 0, 0, -12}
\RouteRow{\(\tau_{12}\); firma histórica \(\mathbb Z^5\)}{+++-\kern0pt{}-\kern0pt{}-+++-\kern0pt{}-\kern0pt{}-; 28|\allowbreak{}-6|\allowbreak{}0|\allowbreak{}0|\allowbreak{}-12}
\RouteRow{\(D_1,\ D_3,\ D_4,\ D_{27},\ D_{36}\)}{66, 64, 68, 60, 40}
\RouteRow{\(\Omega_{90/120},\ P_6\)}{80, 6}
\RouteGroup{Lectores y factores de acción}
\RouteRow{\(K_D\quad /\quad K_\Omega\)}{(100,\allowbreak{}66,\allowbreak{}2)\quad /\quad (80,\allowbreak{}54,\allowbreak{}6)}
\RouteRow{\(\kappa_D\)}{99.51859\allowbreak{}71232726\allowbreak{}53414287\allowbreak{}29}
\RouteRow{\(\kappa_\Omega\)}{79.60661\allowbreak{}01397948\allowbreak{}27586470\allowbreak{}91}
\RouteRow{\(R_{\mathrm{act},D}\)}{1.000000\allowbreak{}06899443\allowbreak{}08259364\allowbreak{}17}
\RouteRow{\(R_{\mathrm{act},\Omega}\)}{1.000000\allowbreak{}05518981\allowbreak{}25318591\allowbreak{}40}
\RouteGroup{Separación de prefijos}
\RouteRow{Área de ambigüedad; área \(\log_2\)}{36; 18.0}
\RouteRow{Primer bloque de separación}{órbita: \textemdash{}; celda: \textemdash{}; ruta: \textemdash{}}
\RouteRow{Ambigüedad final}{2}
\RouteGroup{Secuencias completas}
\RouteRow{\(w_6\)}{99 327 498 99 324 423 15 486 423 99 327 498 99 324 423 15 486 423}
\RouteRow{Tríadas estables por bloque}{0 1 2 3 4 5 6 7 8 9 10 11 12 13 14 14 16 17}
\end{tabularx}\endgroup\par\vspace{5pt}\noindent{\fontsize{8.4}{10}\selectfont\hyperref[trace-280]{Procedencia y códigos originales}\hfill Ficha completa · 280/324}
\clearpage\phantomsection\label{nav:vi-ruta-281}
% VI-CATALOG-ROW rutas 281
\pdfbookmark[2]{Celda 71}{cell-71}
\pdfbookmark[3]{r8c8\_h-1v-1}{route-281}
\RouteTitle{Ruta 281}{r8c8\_\allowbreak{}h-1v-1}
\noindent Celda 71\quad (8, 8)\qquad Orientación \(h=-1,\ v=-1\)\hfill\hyperref[sec:vi-indice-rutas]{Índice}\par\vspace{4pt}
\begingroup\fontsize{10.2}{12.5}\selectfont\setlength{\tabcolsep}{5pt}\renewcommand{\arraystretch}{1.1}\rowcolors{1}{routepale}{white}\noindent\begin{tabularx}{\linewidth}{@{}>{\raggedright\arraybackslash}p{.345\linewidth}>{\raggedright\arraybackslash}X@{}}
\RouteGroup{Identidad estructural}
\RouteRow{Clase y multisección}{quark\_\allowbreak{}up\_\allowbreak{}G3\quad /\quad quark\_\allowbreak{}up\_\allowbreak{}G3}
\RouteRow{Colección y sector}{quark de tipo arriba; quark}
\RouteRow{Clasificación física}{quark de tipo arriba}
\RouteRow{Generación, profundidad y color}{3; G3; R}
\RouteRow{Ruta primaria y órbita de inversión}{\(B_1\): no\qquad 2,\allowbreak{}2;\allowbreak{}8,\allowbreak{}8}
\RouteRow{Firma de clasificación}{24|\allowbreak{}0|\allowbreak{}8|\allowbreak{}0|\allowbreak{}-12|\allowbreak{}+++-\kern0pt{}-\kern0pt{}-+++-\kern0pt{}-\kern0pt{}-}
\RouteGroup{Retornos, memoria y transporte}
\RouteRow{Retornos con fase}{clase: 18; órbita: 18; celda: 27}
\RouteRow{Retorno cinemático}{en 108: 54; extensión: 216}
\RouteRow{Giros y cambios de hoja}{71; 107}
\RouteRow{Acarreo y diferimiento}{deuda total: 3; final: 0; bloques diferidos: 2}
\RouteRow{Tríadas finales; persistencia de Pareto}{17; sí}
\RouteRow{\(K\longrightarrow K+9\)}{repeticiones: 9; cambios de memoria: 9}
\RouteGroup{Firma, fase y diferencias}
\RouteRow{\(n_A,\ s_C,\ \kappa_0,\ \nu_{120},\ \nu_{270}\)}{24, 0, 0, 0, -12}
\RouteRow{\(\tau_{12}\); firma histórica \(\mathbb Z^5\)}{+++-\kern0pt{}-\kern0pt{}-+++-\kern0pt{}-\kern0pt{}-; 24|\allowbreak{}0|\allowbreak{}0|\allowbreak{}0|\allowbreak{}-12}
\RouteRow{\(D_1,\ D_3,\ D_4,\ D_{27},\ D_{36}\)}{54, 56, 68, 48, 32}
\RouteRow{\(\Omega_{90/120},\ P_6\)}{80, 6}
\RouteGroup{Lectores y factores de acción}
\RouteRow{\(K_D\quad /\quad K_\Omega\)}{(80,\allowbreak{}54,\allowbreak{}6)\quad /\quad (80,\allowbreak{}54,\allowbreak{}6)}
\RouteRow{\(\kappa_D\)}{79.60661\allowbreak{}01397948\allowbreak{}27586470\allowbreak{}91}
\RouteRow{\(\kappa_\Omega\)}{79.60661\allowbreak{}01397948\allowbreak{}27586470\allowbreak{}91}
\RouteRow{\(R_{\mathrm{act},D}\)}{1.000000\allowbreak{}05518981\allowbreak{}25318591\allowbreak{}40}
\RouteRow{\(R_{\mathrm{act},\Omega}\)}{1.000000\allowbreak{}05518981\allowbreak{}25318591\allowbreak{}40}
\RouteGroup{Separación de prefijos}
\RouteRow{Área de ambigüedad; área \(\log_2\)}{72; 36.0}
\RouteRow{Primer bloque de separación}{órbita: \textemdash{}; celda: \textemdash{}; ruta: \textemdash{}}
\RouteRow{Ambigüedad final}{4}
\RouteGroup{Secuencias completas}
\RouteRow{\(w_6\)}{423 15 486 423 9 24 243 657 24 423 15 486 423 9 24 243 657 24}
\RouteRow{Tríadas estables por bloque}{0 1 2 3 4 5 6 7 8 9 10 11 11 13 14 14 15 17}
\end{tabularx}\endgroup\par\vspace{5pt}\noindent{\fontsize{8.4}{10}\selectfont\hyperref[trace-281]{Procedencia y códigos originales}\hfill Ficha completa · 281/324}
\clearpage\phantomsection\label{nav:vi-ruta-282}
% VI-CATALOG-ROW rutas 282
\pdfbookmark[3]{r8c8\_h-1v+1}{route-282}
\RouteTitle{Ruta 282}{r8c8\_\allowbreak{}h-1v+1}
\noindent Celda 71\quad (8, 8)\qquad Orientación \(h=-1,\ v=1\)\hfill\hyperref[sec:vi-indice-rutas]{Índice}\par\vspace{4pt}
\begingroup\fontsize{10.2}{12.5}\selectfont\setlength{\tabcolsep}{5pt}\renewcommand{\arraystretch}{1.1}\rowcolors{1}{routepale}{white}\noindent\begin{tabularx}{\linewidth}{@{}>{\raggedright\arraybackslash}p{.345\linewidth}>{\raggedright\arraybackslash}X@{}}
\RouteGroup{Identidad estructural}
\RouteRow{Clase y multisección}{quark\_\allowbreak{}up\_\allowbreak{}G3\quad /\quad quark\_\allowbreak{}up\_\allowbreak{}G3}
\RouteRow{Colección y sector}{quark de tipo arriba; quark}
\RouteRow{Clasificación física}{quark de tipo arriba}
\RouteRow{Generación, profundidad y color}{3; G3; R}
\RouteRow{Ruta primaria y órbita de inversión}{\(B_1\): no\qquad 2,\allowbreak{}2;\allowbreak{}8,\allowbreak{}8}
\RouteRow{Firma de clasificación}{24|\allowbreak{}0|\allowbreak{}8|\allowbreak{}0|\allowbreak{}-12|\allowbreak{}+++-\kern0pt{}-\kern0pt{}-+++-\kern0pt{}-\kern0pt{}-}
\RouteGroup{Retornos, memoria y transporte}
\RouteRow{Retornos con fase}{clase: 27; órbita: 27; celda: 27}
\RouteRow{Retorno cinemático}{en 108: 54; extensión: 216}
\RouteRow{Giros y cambios de hoja}{71; 107}
\RouteRow{Acarreo y diferimiento}{deuda total: 2; final: 1; bloques diferidos: 2}
\RouteRow{Tríadas finales; persistencia de Pareto}{16; sí}
\RouteRow{\(K\longrightarrow K+9\)}{repeticiones: 9; cambios de memoria: 9}
\RouteGroup{Firma, fase y diferencias}
\RouteRow{\(n_A,\ s_C,\ \kappa_0,\ \nu_{120},\ \nu_{270}\)}{50, 4, -4, 2, -4}
\RouteRow{\(\tau_{12}\); firma histórica \(\mathbb Z^5\)}{+++0-0+++0-0; 50|\allowbreak{}4|\allowbreak{}-4|\allowbreak{}2|\allowbreak{}-4}
\RouteRow{\(D_1,\ D_3,\ D_4,\ D_{27},\ D_{36}\)}{84, 28, 84, 36, 24}
\RouteRow{\(\Omega_{90/120},\ P_6\)}{24, 4}
\RouteGroup{Lectores y factores de acción}
\RouteRow{\(K_D\quad /\quad K_\Omega\)}{(60,\allowbreak{}84,\allowbreak{}28)\quad /\quad (24,\allowbreak{}54,\allowbreak{}4)}
\RouteRow{\(\kappa_D\)}{59.39013\allowbreak{}58840155\allowbreak{}40637741\allowbreak{}48}
\RouteRow{\(\kappa_\Omega\)}{23.60638\allowbreak{}77469494\allowbreak{}43306362\allowbreak{}81}
\RouteRow{\(R_{\mathrm{act},D}\)}{1.000000\allowbreak{}04117409\allowbreak{}89467415\allowbreak{}75}
\RouteRow{\(R_{\mathrm{act},\Omega}\)}{1.000000\allowbreak{}01636587\allowbreak{}82446598\allowbreak{}10}
\RouteGroup{Separación de prefijos}
\RouteRow{Área de ambigüedad; área \(\log_2\)}{72; 36.0}
\RouteRow{Primer bloque de separación}{órbita: \textemdash{}; celda: \textemdash{}; ruta: \textemdash{}}
\RouteRow{Ambigüedad final}{4}
\RouteGroup{Secuencias completas}
\RouteRow{\(w_6\)}{414 420 258 414 414 333 255 252 333 414 420 258 414 414 333 255 252 333}
\RouteRow{Tríadas estables por bloque}{0 1 2 3 4 5 6 7 8 9 10 11 12 13 14 14 16 16}
\end{tabularx}\endgroup\par\vspace{5pt}\noindent{\fontsize{8.4}{10}\selectfont\hyperref[trace-282]{Procedencia y códigos originales}\hfill Ficha completa · 282/324}
\clearpage\phantomsection\label{nav:vi-ruta-283}
% VI-CATALOG-ROW rutas 283
\pdfbookmark[3]{r8c8\_h+1v-1}{route-283}
\RouteTitle{Ruta 283}{r8c8\_\allowbreak{}h+1v-1}
\noindent Celda 71\quad (8, 8)\qquad Orientación \(h=1,\ v=-1\)\hfill\hyperref[sec:vi-indice-rutas]{Índice}\par\vspace{4pt}
\begingroup\fontsize{10.2}{12.5}\selectfont\setlength{\tabcolsep}{5pt}\renewcommand{\arraystretch}{1.1}\rowcolors{1}{routepale}{white}\noindent\begin{tabularx}{\linewidth}{@{}>{\raggedright\arraybackslash}p{.345\linewidth}>{\raggedright\arraybackslash}X@{}}
\RouteGroup{Identidad estructural}
\RouteRow{Clase y multisección}{quark\_\allowbreak{}up\_\allowbreak{}G3\quad /\quad quark\_\allowbreak{}up\_\allowbreak{}G3}
\RouteRow{Colección y sector}{quark de tipo arriba; quark}
\RouteRow{Clasificación física}{quark de tipo arriba}
\RouteRow{Generación, profundidad y color}{3; G3; R}
\RouteRow{Ruta primaria y órbita de inversión}{\(B_1\): no\qquad 2,\allowbreak{}2;\allowbreak{}8,\allowbreak{}8}
\RouteRow{Firma de clasificación}{24|\allowbreak{}0|\allowbreak{}8|\allowbreak{}0|\allowbreak{}-12|\allowbreak{}+++-\kern0pt{}-\kern0pt{}-+++-\kern0pt{}-\kern0pt{}-}
\RouteGroup{Retornos, memoria y transporte}
\RouteRow{Retornos con fase}{clase: 27; órbita: 27; celda: 27}
\RouteRow{Retorno cinemático}{en 108: 54; extensión: 216}
\RouteRow{Giros y cambios de hoja}{71; 107}
\RouteRow{Acarreo y diferimiento}{deuda total: 3; final: 0; bloques diferidos: 2}
\RouteRow{Tríadas finales; persistencia de Pareto}{17; sí}
\RouteRow{\(K\longrightarrow K+9\)}{repeticiones: 9; cambios de memoria: 9}
\RouteGroup{Firma, fase y diferencias}
\RouteRow{\(n_A,\ s_C,\ \kappa_0,\ \nu_{120},\ \nu_{270}\)}{50, -4, 4, -2, -4}
\RouteRow{\(\tau_{12}\); firma histórica \(\mathbb Z^5\)}{0+0-\kern0pt{}-\kern0pt{}-0+0-\kern0pt{}-\kern0pt{}-; 50|\allowbreak{}-4|\allowbreak{}4|\allowbreak{}-2|\allowbreak{}-4}
\RouteRow{\(D_1,\ D_3,\ D_4,\ D_{27},\ D_{36}\)}{84, 28, 84, 36, 24}
\RouteRow{\(\Omega_{90/120},\ P_6\)}{24, 4}
\RouteGroup{Lectores y factores de acción}
\RouteRow{\(K_D\quad /\quad K_\Omega\)}{(60,\allowbreak{}84,\allowbreak{}28)\quad /\quad (24,\allowbreak{}54,\allowbreak{}4)}
\RouteRow{\(\kappa_D\)}{59.39013\allowbreak{}58840155\allowbreak{}40637741\allowbreak{}48}
\RouteRow{\(\kappa_\Omega\)}{23.60638\allowbreak{}77469494\allowbreak{}43306362\allowbreak{}81}
\RouteRow{\(R_{\mathrm{act},D}\)}{1.000000\allowbreak{}04117409\allowbreak{}89467415\allowbreak{}75}
\RouteRow{\(R_{\mathrm{act},\Omega}\)}{1.000000\allowbreak{}01636587\allowbreak{}82446598\allowbreak{}10}
\RouteGroup{Separación de prefijos}
\RouteRow{Área de ambigüedad; área \(\log_2\)}{72; 36.0}
\RouteRow{Primer bloque de separación}{órbita: \textemdash{}; celda: \textemdash{}; ruta: \textemdash{}}
\RouteRow{Ambigüedad final}{4}
\RouteGroup{Secuencias completas}
\RouteRow{\(w_6\)}{255 252 333 255 258 258 414 420 258 255 252 333 255 258 258 414 420 258}
\RouteRow{Tríadas estables por bloque}{0 1 2 3 4 5 6 7 8 9 10 11 11 12 14 15 15 17}
\end{tabularx}\endgroup\par\vspace{5pt}\noindent{\fontsize{8.4}{10}\selectfont\hyperref[trace-283]{Procedencia y códigos originales}\hfill Ficha completa · 283/324}
\clearpage\phantomsection\label{nav:vi-ruta-284}
% VI-CATALOG-ROW rutas 284
\pdfbookmark[3]{r8c8\_h+1v+1}{route-284}
\RouteTitle{Ruta 284}{r8c8\_\allowbreak{}h+1v+1}
\noindent Celda 71\quad (8, 8)\qquad Orientación \(h=1,\ v=1\)\hfill\hyperref[sec:vi-indice-rutas]{Índice}\par\vspace{4pt}
\begingroup\fontsize{10.2}{12.5}\selectfont\setlength{\tabcolsep}{5pt}\renewcommand{\arraystretch}{1.1}\rowcolors{1}{routepale}{white}\noindent\begin{tabularx}{\linewidth}{@{}>{\raggedright\arraybackslash}p{.345\linewidth}>{\raggedright\arraybackslash}X@{}}
\RouteGroup{Identidad estructural}
\RouteRow{Clase y multisección}{quark\_\allowbreak{}up\_\allowbreak{}G3\quad /\quad quark\_\allowbreak{}up\_\allowbreak{}G3}
\RouteRow{Colección y sector}{quark de tipo arriba; quark}
\RouteRow{Clasificación física}{quark de tipo arriba}
\RouteRow{Generación, profundidad y color}{3; G3; R}
\RouteRow{Ruta primaria y órbita de inversión}{\(B_1\): sí\qquad 2,\allowbreak{}2;\allowbreak{}8,\allowbreak{}8}
\RouteRow{Firma de clasificación}{24|\allowbreak{}0|\allowbreak{}8|\allowbreak{}0|\allowbreak{}-12|\allowbreak{}+++-\kern0pt{}-\kern0pt{}-+++-\kern0pt{}-\kern0pt{}-}
\RouteGroup{Retornos, memoria y transporte}
\RouteRow{Retornos con fase}{clase: 9; órbita: 9; celda: 27}
\RouteRow{Retorno cinemático}{en 108: 54; extensión: 216}
\RouteRow{Giros y cambios de hoja}{71; 107}
\RouteRow{Acarreo y diferimiento}{deuda total: 2; final: 1; bloques diferidos: 2}
\RouteRow{Tríadas finales; persistencia de Pareto}{16; sí}
\RouteRow{\(K\longrightarrow K+9\)}{repeticiones: 9; cambios de memoria: 9}
\RouteGroup{Firma, fase y diferencias}
\RouteRow{\(n_A,\ s_C,\ \kappa_0,\ \nu_{120},\ \nu_{270}\)}{24, 0, 0, 0, -12}
\RouteRow{\(\tau_{12}\); firma histórica \(\mathbb Z^5\)}{+++-\kern0pt{}-\kern0pt{}-+++-\kern0pt{}-\kern0pt{}-; 24|\allowbreak{}0|\allowbreak{}0|\allowbreak{}0|\allowbreak{}-12}
\RouteRow{\(D_1,\ D_3,\ D_4,\ D_{27},\ D_{36}\)}{54, 56, 68, 48, 32}
\RouteRow{\(\Omega_{90/120},\ P_6\)}{80, 6}
\RouteGroup{Lectores y factores de acción}
\RouteRow{\(K_D\quad /\quad K_\Omega\)}{(80,\allowbreak{}54,\allowbreak{}6)\quad /\quad (80,\allowbreak{}54,\allowbreak{}6)}
\RouteRow{\(\kappa_D\)}{79.60661\allowbreak{}01397948\allowbreak{}27586470\allowbreak{}91}
\RouteRow{\(\kappa_\Omega\)}{79.60661\allowbreak{}01397948\allowbreak{}27586470\allowbreak{}91}
\RouteRow{\(R_{\mathrm{act},D}\)}{1.000000\allowbreak{}05518981\allowbreak{}25318591\allowbreak{}40}
\RouteRow{\(R_{\mathrm{act},\Omega}\)}{1.000000\allowbreak{}05518981\allowbreak{}25318591\allowbreak{}40}
\RouteGroup{Separación de prefijos}
\RouteRow{Área de ambigüedad; área \(\log_2\)}{72; 36.0}
\RouteRow{Primer bloque de separación}{órbita: \textemdash{}; celda: \textemdash{}; ruta: \textemdash{}}
\RouteRow{Ambigüedad final}{4}
\RouteGroup{Secuencias completas}
\RouteRow{\(w_6\)}{243 657 24 243 663 486 423 15 486 243 657 24 243 663 486 423 15 486}
\RouteRow{Tríadas estables por bloque}{0 1 2 3 4 5 6 7 8 9 10 11 12 13 14 14 16 16}
\end{tabularx}\endgroup\par\vspace{5pt}\noindent{\fontsize{8.4}{10}\selectfont\hyperref[trace-284]{Procedencia y códigos originales}\hfill Ficha completa · 284/324}
\clearpage\phantomsection\label{nav:vi-ruta-285}
% VI-CATALOG-ROW rutas 285
\pdfbookmark[2]{Celda 72}{cell-72}
\pdfbookmark[3]{r8c9\_h-1v-1}{route-285}
\RouteTitle{Ruta 285}{r8c9\_\allowbreak{}h-1v-1}
\noindent Celda 72\quad (8, 9)\qquad Orientación \(h=-1,\ v=-1\)\hfill\hyperref[sec:vi-indice-rutas]{Índice}\par\vspace{4pt}
\begingroup\fontsize{10.2}{12.5}\selectfont\setlength{\tabcolsep}{5pt}\renewcommand{\arraystretch}{1.1}\rowcolors{1}{routepale}{white}\noindent\begin{tabularx}{\linewidth}{@{}>{\raggedright\arraybackslash}p{.345\linewidth}>{\raggedright\arraybackslash}X@{}}
\RouteGroup{Identidad estructural}
\RouteRow{Clase y multisección}{quark\_\allowbreak{}down\_\allowbreak{}G4\quad /\quad quark\_\allowbreak{}down\_\allowbreak{}G4}
\RouteRow{Colección y sector}{quark de tipo abajo; quark}
\RouteRow{Clasificación física}{quark de tipo abajo}
\RouteRow{Generación, profundidad y color}{4; G4; B}
\RouteRow{Ruta primaria y órbita de inversión}{\(B_1\): no\qquad 2,\allowbreak{}1;\allowbreak{}8,\allowbreak{}9}
\RouteRow{Firma de clasificación}{28|\allowbreak{}6|\allowbreak{}0|\allowbreak{}0|\allowbreak{}-12|\allowbreak{}+++-\kern0pt{}-\kern0pt{}-+++-\kern0pt{}-\kern0pt{}-}
\RouteGroup{Retornos, memoria y transporte}
\RouteRow{Retornos con fase}{clase: 27; órbita: 27; celda: 27}
\RouteRow{Retorno cinemático}{en 108: 54; extensión: 216}
\RouteRow{Giros y cambios de hoja}{71; 107}
\RouteRow{Acarreo y diferimiento}{deuda total: 1; final: 0; bloques diferidos: 1}
\RouteRow{Tríadas finales; persistencia de Pareto}{17; sí}
\RouteRow{\(K\longrightarrow K+9\)}{repeticiones: 9; cambios de memoria: 9}
\RouteGroup{Firma, fase y diferencias}
\RouteRow{\(n_A,\ s_C,\ \kappa_0,\ \nu_{120},\ \nu_{270}\)}{28, -6, 0, 0, -12}
\RouteRow{\(\tau_{12}\); firma histórica \(\mathbb Z^5\)}{+++-\kern0pt{}-\kern0pt{}-+++-\kern0pt{}-\kern0pt{}-; 28|\allowbreak{}-6|\allowbreak{}0|\allowbreak{}0|\allowbreak{}-12}
\RouteRow{\(D_1,\ D_3,\ D_4,\ D_{27},\ D_{36}\)}{66, 60, 68, 48, 32}
\RouteRow{\(\Omega_{90/120},\ P_6\)}{80, 6}
\RouteGroup{Lectores y factores de acción}
\RouteRow{\(K_D\quad /\quad K_\Omega\)}{(80,\allowbreak{}66,\allowbreak{}4)\quad /\quad (80,\allowbreak{}54,\allowbreak{}6)}
\RouteRow{\(\kappa_D\)}{79.51881\allowbreak{}95161180\allowbreak{}37694395\allowbreak{}39}
\RouteRow{\(\kappa_\Omega\)}{79.60661\allowbreak{}01397948\allowbreak{}27586470\allowbreak{}91}
\RouteRow{\(R_{\mathrm{act},D}\)}{1.000000\allowbreak{}05512894\allowbreak{}88901612\allowbreak{}64}
\RouteRow{\(R_{\mathrm{act},\Omega}\)}{1.000000\allowbreak{}05518981\allowbreak{}25318591\allowbreak{}40}
\RouteGroup{Separación de prefijos}
\RouteRow{Área de ambigüedad; área \(\log_2\)}{36; 18.0}
\RouteRow{Primer bloque de separación}{órbita: \textemdash{}; celda: \textemdash{}; ruta: \textemdash{}}
\RouteRow{Ambigüedad final}{2}
\RouteGroup{Secuencias completas}
\RouteRow{\(w_6\)}{570 264 90 570 267 243 657 24 243 570 264 90 570 267 243 657 24 243}
\RouteRow{Tríadas estables por bloque}{0 1 1 3 4 5 6 7 8 9 10 11 12 13 14 15 16 17}
\end{tabularx}\endgroup\par\vspace{5pt}\noindent{\fontsize{8.4}{10}\selectfont\hyperref[trace-285]{Procedencia y códigos originales}\hfill Ficha completa · 285/324}
\clearpage\phantomsection\label{nav:vi-ruta-286}
% VI-CATALOG-ROW rutas 286
\pdfbookmark[3]{r8c9\_h-1v+1}{route-286}
\RouteTitle{Ruta 286}{r8c9\_\allowbreak{}h-1v+1}
\noindent Celda 72\quad (8, 9)\qquad Orientación \(h=-1,\ v=1\)\hfill\hyperref[sec:vi-indice-rutas]{Índice}\par\vspace{4pt}
\begingroup\fontsize{10.2}{12.5}\selectfont\setlength{\tabcolsep}{5pt}\renewcommand{\arraystretch}{1.1}\rowcolors{1}{routepale}{white}\noindent\begin{tabularx}{\linewidth}{@{}>{\raggedright\arraybackslash}p{.345\linewidth}>{\raggedright\arraybackslash}X@{}}
\RouteGroup{Identidad estructural}
\RouteRow{Clase y multisección}{quark\_\allowbreak{}down\_\allowbreak{}G4\quad /\quad quark\_\allowbreak{}down\_\allowbreak{}G4}
\RouteRow{Colección y sector}{quark de tipo abajo; quark}
\RouteRow{Clasificación física}{quark de tipo abajo}
\RouteRow{Generación, profundidad y color}{4; G4; B}
\RouteRow{Ruta primaria y órbita de inversión}{\(B_1\): sí\qquad 2,\allowbreak{}1;\allowbreak{}8,\allowbreak{}9}
\RouteRow{Firma de clasificación}{28|\allowbreak{}6|\allowbreak{}0|\allowbreak{}0|\allowbreak{}-12|\allowbreak{}+++-\kern0pt{}-\kern0pt{}-+++-\kern0pt{}-\kern0pt{}-}
\RouteGroup{Retornos, memoria y transporte}
\RouteRow{Retornos con fase}{clase: 27; órbita: 27; celda: 27}
\RouteRow{Retorno cinemático}{en 108: 54; extensión: 216}
\RouteRow{Giros y cambios de hoja}{71; 107}
\RouteRow{Acarreo y diferimiento}{deuda total: 1; final: 0; bloques diferidos: 1}
\RouteRow{Tríadas finales; persistencia de Pareto}{17; sí}
\RouteRow{\(K\longrightarrow K+9\)}{repeticiones: 9; cambios de memoria: 9}
\RouteGroup{Firma, fase y diferencias}
\RouteRow{\(n_A,\ s_C,\ \kappa_0,\ \nu_{120},\ \nu_{270}\)}{14, -4, 0, 0, -12}
\RouteRow{\(\tau_{12}\); firma histórica \(\mathbb Z^5\)}{+++-\kern0pt{}-\kern0pt{}-+++-\kern0pt{}-\kern0pt{}-; 14|\allowbreak{}-4|\allowbreak{}0|\allowbreak{}0|\allowbreak{}-12}
\RouteRow{\(D_1,\ D_3,\ D_4,\ D_{27},\ D_{36}\)}{78, 24, 78, 24, 16}
\RouteRow{\(\Omega_{90/120},\ P_6\)}{80, 6}
\RouteGroup{Lectores y factores de acción}
\RouteRow{\(K_D\quad /\quad K_\Omega\)}{(40,\allowbreak{}78,\allowbreak{}27)\quad /\quad (80,\allowbreak{}54,\allowbreak{}6)}
\RouteRow{\(\kappa_D\)}{39.43380\allowbreak{}88030085\allowbreak{}51303671\allowbreak{}14}
\RouteRow{\(\kappa_\Omega\)}{79.60661\allowbreak{}01397948\allowbreak{}27586470\allowbreak{}91}
\RouteRow{\(R_{\mathrm{act},D}\)}{1.000000\allowbreak{}02733874\allowbreak{}08548943\allowbreak{}58}
\RouteRow{\(R_{\mathrm{act},\Omega}\)}{1.000000\allowbreak{}05518981\allowbreak{}25318591\allowbreak{}40}
\RouteGroup{Separación de prefijos}
\RouteRow{Área de ambigüedad; área \(\log_2\)}{36; 18.0}
\RouteRow{Primer bloque de separación}{órbita: \textemdash{}; celda: \textemdash{}; ruta: \textemdash{}}
\RouteRow{Ambigüedad final}{2}
\RouteGroup{Secuencias completas}
\RouteRow{\(w_6\)}{585 507 504 585 510 666 672 510 666 585 507 504 585 510 666 672 510 666}
\RouteRow{Tríadas estables por bloque}{0 1 2 3 4 5 6 7 8 9 10 11 12 13 14 14 16 17}
\end{tabularx}\endgroup\par\vspace{5pt}\noindent{\fontsize{8.4}{10}\selectfont\hyperref[trace-286]{Procedencia y códigos originales}\hfill Ficha completa · 286/324}
\clearpage\phantomsection\label{nav:vi-ruta-287}
% VI-CATALOG-ROW rutas 287
\pdfbookmark[3]{r8c9\_h+1v-1}{route-287}
\RouteTitle{Ruta 287}{r8c9\_\allowbreak{}h+1v-1}
\noindent Celda 72\quad (8, 9)\qquad Orientación \(h=1,\ v=-1\)\hfill\hyperref[sec:vi-indice-rutas]{Índice}\par\vspace{4pt}
\begingroup\fontsize{10.2}{12.5}\selectfont\setlength{\tabcolsep}{5pt}\renewcommand{\arraystretch}{1.1}\rowcolors{1}{routepale}{white}\noindent\begin{tabularx}{\linewidth}{@{}>{\raggedright\arraybackslash}p{.345\linewidth}>{\raggedright\arraybackslash}X@{}}
\RouteGroup{Identidad estructural}
\RouteRow{Clase y multisección}{quark\_\allowbreak{}down\_\allowbreak{}G4\quad /\quad quark\_\allowbreak{}down\_\allowbreak{}G4}
\RouteRow{Colección y sector}{quark de tipo abajo; quark}
\RouteRow{Clasificación física}{quark de tipo abajo}
\RouteRow{Generación, profundidad y color}{4; G4; B}
\RouteRow{Ruta primaria y órbita de inversión}{\(B_1\): no\qquad 2,\allowbreak{}1;\allowbreak{}8,\allowbreak{}9}
\RouteRow{Firma de clasificación}{28|\allowbreak{}6|\allowbreak{}0|\allowbreak{}0|\allowbreak{}-12|\allowbreak{}+++-\kern0pt{}-\kern0pt{}-+++-\kern0pt{}-\kern0pt{}-}
\RouteGroup{Retornos, memoria y transporte}
\RouteRow{Retornos con fase}{clase: 27; órbita: 27; celda: 27}
\RouteRow{Retorno cinemático}{en 108: 54; extensión: 216}
\RouteRow{Giros y cambios de hoja}{71; 107}
\RouteRow{Acarreo y diferimiento}{deuda total: 1; final: 0; bloques diferidos: 1}
\RouteRow{Tríadas finales; persistencia de Pareto}{17; sí}
\RouteRow{\(K\longrightarrow K+9\)}{repeticiones: 9; cambios de memoria: 9}
\RouteGroup{Firma, fase y diferencias}
\RouteRow{\(n_A,\ s_C,\ \kappa_0,\ \nu_{120},\ \nu_{270}\)}{14, 4, 0, 0, -12}
\RouteRow{\(\tau_{12}\); firma histórica \(\mathbb Z^5\)}{+++-\kern0pt{}-\kern0pt{}-+++-\kern0pt{}-\kern0pt{}-; 14|\allowbreak{}4|\allowbreak{}0|\allowbreak{}0|\allowbreak{}-12}
\RouteRow{\(D_1,\ D_3,\ D_4,\ D_{27},\ D_{36}\)}{78, 24, 78, 24, 16}
\RouteRow{\(\Omega_{90/120},\ P_6\)}{80, 6}
\RouteGroup{Lectores y factores de acción}
\RouteRow{\(K_D\quad /\quad K_\Omega\)}{(40,\allowbreak{}78,\allowbreak{}27)\quad /\quad (80,\allowbreak{}54,\allowbreak{}6)}
\RouteRow{\(\kappa_D\)}{39.43380\allowbreak{}88030085\allowbreak{}51303671\allowbreak{}14}
\RouteRow{\(\kappa_\Omega\)}{79.60661\allowbreak{}01397948\allowbreak{}27586470\allowbreak{}91}
\RouteRow{\(R_{\mathrm{act},D}\)}{1.000000\allowbreak{}02733874\allowbreak{}08548943\allowbreak{}58}
\RouteRow{\(R_{\mathrm{act},\Omega}\)}{1.000000\allowbreak{}05518981\allowbreak{}25318591\allowbreak{}40}
\RouteGroup{Separación de prefijos}
\RouteRow{Área de ambigüedad; área \(\log_2\)}{36; 18.0}
\RouteRow{Primer bloque de separación}{órbita: \textemdash{}; celda: \textemdash{}; ruta: \textemdash{}}
\RouteRow{Ambigüedad final}{2}
\RouteGroup{Secuencias completas}
\RouteRow{\(w_6\)}{672 510 666 672 507 504 585 507 504 672 510 666 672 507 504 585 507 504}
\RouteRow{Tríadas estables por bloque}{0 1 2 3 4 5 6 7 8 9 10 11 12 13 14 14 16 17}
\end{tabularx}\endgroup\par\vspace{5pt}\noindent{\fontsize{8.4}{10}\selectfont\hyperref[trace-287]{Procedencia y códigos originales}\hfill Ficha completa · 287/324}
\clearpage\phantomsection\label{nav:vi-ruta-288}
% VI-CATALOG-ROW rutas 288
\pdfbookmark[3]{r8c9\_h+1v+1}{route-288}
\RouteTitle{Ruta 288}{r8c9\_\allowbreak{}h+1v+1}
\noindent Celda 72\quad (8, 9)\qquad Orientación \(h=1,\ v=1\)\hfill\hyperref[sec:vi-indice-rutas]{Índice}\par\vspace{4pt}
\begingroup\fontsize{10.2}{12.5}\selectfont\setlength{\tabcolsep}{5pt}\renewcommand{\arraystretch}{1.1}\rowcolors{1}{routepale}{white}\noindent\begin{tabularx}{\linewidth}{@{}>{\raggedright\arraybackslash}p{.345\linewidth}>{\raggedright\arraybackslash}X@{}}
\RouteGroup{Identidad estructural}
\RouteRow{Clase y multisección}{quark\_\allowbreak{}down\_\allowbreak{}G4\quad /\quad quark\_\allowbreak{}down\_\allowbreak{}G4}
\RouteRow{Colección y sector}{quark de tipo abajo; quark}
\RouteRow{Clasificación física}{quark de tipo abajo}
\RouteRow{Generación, profundidad y color}{4; G4; B}
\RouteRow{Ruta primaria y órbita de inversión}{\(B_1\): no\qquad 2,\allowbreak{}1;\allowbreak{}8,\allowbreak{}9}
\RouteRow{Firma de clasificación}{28|\allowbreak{}6|\allowbreak{}0|\allowbreak{}0|\allowbreak{}-12|\allowbreak{}+++-\kern0pt{}-\kern0pt{}-+++-\kern0pt{}-\kern0pt{}-}
\RouteGroup{Retornos, memoria y transporte}
\RouteRow{Retornos con fase}{clase: 27; órbita: 27; celda: 27}
\RouteRow{Retorno cinemático}{en 108: 54; extensión: 216}
\RouteRow{Giros y cambios de hoja}{71; 107}
\RouteRow{Acarreo y diferimiento}{deuda total: 0; final: 0; bloques diferidos: 0}
\RouteRow{Tríadas finales; persistencia de Pareto}{17; sí}
\RouteRow{\(K\longrightarrow K+9\)}{repeticiones: 9; cambios de memoria: 9}
\RouteGroup{Firma, fase y diferencias}
\RouteRow{\(n_A,\ s_C,\ \kappa_0,\ \nu_{120},\ \nu_{270}\)}{28, 6, 0, 0, -12}
\RouteRow{\(\tau_{12}\); firma histórica \(\mathbb Z^5\)}{+++-\kern0pt{}-\kern0pt{}-+++-\kern0pt{}-\kern0pt{}-; 28|\allowbreak{}6|\allowbreak{}0|\allowbreak{}0|\allowbreak{}-12}
\RouteRow{\(D_1,\ D_3,\ D_4,\ D_{27},\ D_{36}\)}{66, 60, 68, 48, 32}
\RouteRow{\(\Omega_{90/120},\ P_6\)}{80, 6}
\RouteGroup{Lectores y factores de acción}
\RouteRow{\(K_D\quad /\quad K_\Omega\)}{(80,\allowbreak{}66,\allowbreak{}4)\quad /\quad (80,\allowbreak{}54,\allowbreak{}6)}
\RouteRow{\(\kappa_D\)}{79.51881\allowbreak{}95161180\allowbreak{}37694395\allowbreak{}39}
\RouteRow{\(\kappa_\Omega\)}{79.60661\allowbreak{}01397948\allowbreak{}27586470\allowbreak{}91}
\RouteRow{\(R_{\mathrm{act},D}\)}{1.000000\allowbreak{}05512894\allowbreak{}88901612\allowbreak{}64}
\RouteRow{\(R_{\mathrm{act},\Omega}\)}{1.000000\allowbreak{}05518981\allowbreak{}25318591\allowbreak{}40}
\RouteGroup{Separación de prefijos}
\RouteRow{Área de ambigüedad; área \(\log_2\)}{40; 20.33985\allowbreak{}00028846\allowbreak{}24}
\RouteRow{Primer bloque de separación}{órbita: \textemdash{}; celda: \textemdash{}; ruta: \textemdash{}}
\RouteRow{Ambigüedad final}{2}
\RouteGroup{Secuencias completas}
\RouteRow{\(w_6\)}{657 24 243 657 21 90 570 264 90 657 24 243 657 21 90 570 264 90}
\RouteRow{Tríadas estables por bloque}{0 1 2 3 4 5 6 7 8 9 10 11 12 13 14 15 16 17}
\end{tabularx}\endgroup\par\vspace{5pt}\noindent{\fontsize{8.4}{10}\selectfont\hyperref[trace-288]{Procedencia y códigos originales}\hfill Ficha completa · 288/324}
\clearpage\phantomsection\label{nav:vi-ruta-289}
% VI-CATALOG-ROW rutas 289
\pdfbookmark[2]{Celda 73}{cell-73}
\pdfbookmark[3]{r9c1\_h-1v-1}{route-289}
\RouteTitle{Ruta 289}{r9c1\_\allowbreak{}h-1v-1}
\noindent Celda 73\quad (9, 1)\qquad Orientación \(h=-1,\ v=-1\)\hfill\hyperref[sec:vi-indice-rutas]{Índice}\par\vspace{4pt}
\begingroup\fontsize{10.2}{12.5}\selectfont\setlength{\tabcolsep}{5pt}\renewcommand{\arraystretch}{1.1}\rowcolors{1}{routepale}{white}\noindent\begin{tabularx}{\linewidth}{@{}>{\raggedright\arraybackslash}p{.345\linewidth}>{\raggedright\arraybackslash}X@{}}
\RouteGroup{Identidad estructural}
\RouteRow{Clase y multisección}{charged\_\allowbreak{}lepton\_\allowbreak{}G4\quad /\quad charged\_\allowbreak{}lepton\_\allowbreak{}G4}
\RouteRow{Colección y sector}{leptón cargado; leptón}
\RouteRow{Clasificación física}{leptón cargado}
\RouteRow{Generación, profundidad y color}{4; G4; \textemdash{}}
\RouteRow{Ruta primaria y órbita de inversión}{\(B_1\): sí\qquad 1,\allowbreak{}9;\allowbreak{}9,\allowbreak{}1}
\RouteRow{Firma de clasificación}{28|\allowbreak{}6|\allowbreak{}0|\allowbreak{}0|\allowbreak{}-12|\allowbreak{}+++-\kern0pt{}-\kern0pt{}-+++-\kern0pt{}-\kern0pt{}-}
\RouteGroup{Retornos, memoria y transporte}
\RouteRow{Retornos con fase}{clase: 27; órbita: 27; celda: 27}
\RouteRow{Retorno cinemático}{en 108: 54; extensión: 216}
\RouteRow{Giros y cambios de hoja}{71; 107}
\RouteRow{Acarreo y diferimiento}{deuda total: 2; final: 1; bloques diferidos: 1}
\RouteRow{Tríadas finales; persistencia de Pareto}{16; no}
\RouteRow{\(K\longrightarrow K+9\)}{repeticiones: 9; cambios de memoria: 9}
\RouteGroup{Firma, fase y diferencias}
\RouteRow{\(n_A,\ s_C,\ \kappa_0,\ \nu_{120},\ \nu_{270}\)}{28, -6, 0, -2, -12}
\RouteRow{\(\tau_{12}\); firma histórica \(\mathbb Z^5\)}{+++-\kern0pt{}-\kern0pt{}-+++-\kern0pt{}-\kern0pt{}-; 28|\allowbreak{}-6|\allowbreak{}0|\allowbreak{}-2|\allowbreak{}-12}
\RouteRow{\(D_1,\ D_3,\ D_4,\ D_{27},\ D_{36}\)}{66, 60, 66, 48, 32}
\RouteRow{\(\Omega_{90/120},\ P_6\)}{80, 6}
\RouteGroup{Lectores y factores de acción}
\RouteRow{\(K_D\quad /\quad K_\Omega\)}{(80,\allowbreak{}66,\allowbreak{}3)\quad /\quad (80,\allowbreak{}54,\allowbreak{}6)}
\RouteRow{\(\kappa_D\)}{79.51870\allowbreak{}83196953\allowbreak{}45554341\allowbreak{}34}
\RouteRow{\(\kappa_\Omega\)}{79.60661\allowbreak{}01397948\allowbreak{}27586470\allowbreak{}91}
\RouteRow{\(R_{\mathrm{act},D}\)}{1.000000\allowbreak{}05512887\allowbreak{}17997050\allowbreak{}72}
\RouteRow{\(R_{\mathrm{act},\Omega}\)}{1.000000\allowbreak{}05518981\allowbreak{}25318591\allowbreak{}40}
\RouteGroup{Separación de prefijos}
\RouteRow{Área de ambigüedad; área \(\log_2\)}{54; 28.52932\allowbreak{}50129808}
\RouteRow{Primer bloque de separación}{órbita: \textemdash{}; celda: \textemdash{}; ruta: \textemdash{}}
\RouteRow{Ambigüedad final}{3}
\RouteGroup{Secuencias completas}
\RouteRow{\(w_6\)}{264 90 570 264 90 24 243 657 24 264 90 570 264 90 24 243 657 24}
\RouteRow{Tríadas estables por bloque}{0 1 2 3 4 5 6 7 8 9 10 11 12 13 14 15 15 16}
\end{tabularx}\endgroup\par\vspace{5pt}\noindent{\fontsize{8.4}{10}\selectfont\hyperref[trace-289]{Procedencia y códigos originales}\hfill Ficha completa · 289/324}
\clearpage\phantomsection\label{nav:vi-ruta-290}
% VI-CATALOG-ROW rutas 290
\pdfbookmark[3]{r9c1\_h-1v+1}{route-290}
\RouteTitle{Ruta 290}{r9c1\_\allowbreak{}h-1v+1}
\noindent Celda 73\quad (9, 1)\qquad Orientación \(h=-1,\ v=1\)\hfill\hyperref[sec:vi-indice-rutas]{Índice}\par\vspace{4pt}
\begingroup\fontsize{10.2}{12.5}\selectfont\setlength{\tabcolsep}{5pt}\renewcommand{\arraystretch}{1.1}\rowcolors{1}{routepale}{white}\noindent\begin{tabularx}{\linewidth}{@{}>{\raggedright\arraybackslash}p{.345\linewidth}>{\raggedright\arraybackslash}X@{}}
\RouteGroup{Identidad estructural}
\RouteRow{Clase y multisección}{charged\_\allowbreak{}lepton\_\allowbreak{}G4\quad /\quad charged\_\allowbreak{}lepton\_\allowbreak{}G4}
\RouteRow{Colección y sector}{leptón cargado; leptón}
\RouteRow{Clasificación física}{leptón cargado}
\RouteRow{Generación, profundidad y color}{4; G4; \textemdash{}}
\RouteRow{Ruta primaria y órbita de inversión}{\(B_1\): no\qquad 1,\allowbreak{}9;\allowbreak{}9,\allowbreak{}1}
\RouteRow{Firma de clasificación}{28|\allowbreak{}6|\allowbreak{}0|\allowbreak{}0|\allowbreak{}-12|\allowbreak{}+++-\kern0pt{}-\kern0pt{}-+++-\kern0pt{}-\kern0pt{}-}
\RouteGroup{Retornos, memoria y transporte}
\RouteRow{Retornos con fase}{clase: 27; órbita: 27; celda: 27}
\RouteRow{Retorno cinemático}{en 108: 54; extensión: 216}
\RouteRow{Giros y cambios de hoja}{71; 107}
\RouteRow{Acarreo y diferimiento}{deuda total: 2; final: 0; bloques diferidos: 1}
\RouteRow{Tríadas finales; persistencia de Pareto}{17; no}
\RouteRow{\(K\longrightarrow K+9\)}{repeticiones: 9; cambios de memoria: 9}
\RouteGroup{Firma, fase y diferencias}
\RouteRow{\(n_A,\ s_C,\ \kappa_0,\ \nu_{120},\ \nu_{270}\)}{50, 4, -2, 2, -8}
\RouteRow{\(\tau_{12}\); firma histórica \(\mathbb Z^5\)}{+++-0-+++-0-; 50|\allowbreak{}4|\allowbreak{}-2|\allowbreak{}2|\allowbreak{}-8}
\RouteRow{\(D_1,\ D_3,\ D_4,\ D_{27},\ D_{36}\)}{84, 24, 84, 24, 16}
\RouteRow{\(\Omega_{90/120},\ P_6\)}{48, 5}
\RouteGroup{Lectores y factores de acción}
\RouteRow{\(K_D\quad /\quad K_\Omega\)}{(40,\allowbreak{}84,\allowbreak{}30)\quad /\quad (48,\allowbreak{}54,\allowbreak{}5)}
\RouteRow{\(\kappa_D\)}{39.39035\allowbreak{}82768609\allowbreak{}24917849\allowbreak{}59}
\RouteRow{\(\kappa_\Omega\)}{47.60649\allowbreak{}89433721\allowbreak{}35446416\allowbreak{}86}
\RouteRow{\(R_{\mathrm{act},D}\)}{1.000000\allowbreak{}02730861\allowbreak{}73967087\allowbreak{}05}
\RouteRow{\(R_{\mathrm{act},\Omega}\)}{1.000000\allowbreak{}03300471\allowbreak{}80532430\allowbreak{}84}
\RouteGroup{Separación de prefijos}
\RouteRow{Área de ambigüedad; área \(\log_2\)}{54; 28.52932\allowbreak{}50129808}
\RouteRow{Primer bloque de separación}{órbita: \textemdash{}; celda: \textemdash{}; ruta: \textemdash{}}
\RouteRow{Ambigüedad final}{3}
\RouteGroup{Secuencias completas}
\RouteRow{\(w_6\)}{258 414 420 258 414 255 252 333 255 258 414 420 258 414 255 252 333 255}
\RouteRow{Tríadas estables por bloque}{0 1 2 3 4 5 6 7 8 9 10 11 11 12 14 15 16 17}
\end{tabularx}\endgroup\par\vspace{5pt}\noindent{\fontsize{8.4}{10}\selectfont\hyperref[trace-290]{Procedencia y códigos originales}\hfill Ficha completa · 290/324}
\clearpage\phantomsection\label{nav:vi-ruta-291}
% VI-CATALOG-ROW rutas 291
\pdfbookmark[3]{r9c1\_h+1v-1}{route-291}
\RouteTitle{Ruta 291}{r9c1\_\allowbreak{}h+1v-1}
\noindent Celda 73\quad (9, 1)\qquad Orientación \(h=1,\ v=-1\)\hfill\hyperref[sec:vi-indice-rutas]{Índice}\par\vspace{4pt}
\begingroup\fontsize{10.2}{12.5}\selectfont\setlength{\tabcolsep}{5pt}\renewcommand{\arraystretch}{1.1}\rowcolors{1}{routepale}{white}\noindent\begin{tabularx}{\linewidth}{@{}>{\raggedright\arraybackslash}p{.345\linewidth}>{\raggedright\arraybackslash}X@{}}
\RouteGroup{Identidad estructural}
\RouteRow{Clase y multisección}{charged\_\allowbreak{}lepton\_\allowbreak{}G4\quad /\quad charged\_\allowbreak{}lepton\_\allowbreak{}G4}
\RouteRow{Colección y sector}{leptón cargado; leptón}
\RouteRow{Clasificación física}{leptón cargado}
\RouteRow{Generación, profundidad y color}{4; G4; \textemdash{}}
\RouteRow{Ruta primaria y órbita de inversión}{\(B_1\): no\qquad 1,\allowbreak{}9;\allowbreak{}9,\allowbreak{}1}
\RouteRow{Firma de clasificación}{28|\allowbreak{}6|\allowbreak{}0|\allowbreak{}0|\allowbreak{}-12|\allowbreak{}+++-\kern0pt{}-\kern0pt{}-+++-\kern0pt{}-\kern0pt{}-}
\RouteGroup{Retornos, memoria y transporte}
\RouteRow{Retornos con fase}{clase: 27; órbita: 27; celda: 27}
\RouteRow{Retorno cinemático}{en 108: 54; extensión: 216}
\RouteRow{Giros y cambios de hoja}{71; 107}
\RouteRow{Acarreo y diferimiento}{deuda total: 3; final: 1; bloques diferidos: 3}
\RouteRow{Tríadas finales; persistencia de Pareto}{16; no}
\RouteRow{\(K\longrightarrow K+9\)}{repeticiones: 9; cambios de memoria: 9}
\RouteGroup{Firma, fase y diferencias}
\RouteRow{\(n_A,\ s_C,\ \kappa_0,\ \nu_{120},\ \nu_{270}\)}{50, -4, 2, -2, -8}
\RouteRow{\(\tau_{12}\); firma histórica \(\mathbb Z^5\)}{+0+-\kern0pt{}-\kern0pt{}-+0+-\kern0pt{}-\kern0pt{}-; 50|\allowbreak{}-4|\allowbreak{}2|\allowbreak{}-2|\allowbreak{}-8}
\RouteRow{\(D_1,\ D_3,\ D_4,\ D_{27},\ D_{36}\)}{84, 24, 84, 24, 16}
\RouteRow{\(\Omega_{90/120},\ P_6\)}{48, 5}
\RouteGroup{Lectores y factores de acción}
\RouteRow{\(K_D\quad /\quad K_\Omega\)}{(40,\allowbreak{}84,\allowbreak{}30)\quad /\quad (48,\allowbreak{}54,\allowbreak{}5)}
\RouteRow{\(\kappa_D\)}{39.39035\allowbreak{}82768609\allowbreak{}24917849\allowbreak{}59}
\RouteRow{\(\kappa_\Omega\)}{47.60649\allowbreak{}89433721\allowbreak{}35446416\allowbreak{}86}
\RouteRow{\(R_{\mathrm{act},D}\)}{1.000000\allowbreak{}02730861\allowbreak{}73967087\allowbreak{}05}
\RouteRow{\(R_{\mathrm{act},\Omega}\)}{1.000000\allowbreak{}03300471\allowbreak{}80532430\allowbreak{}84}
\RouteGroup{Separación de prefijos}
\RouteRow{Área de ambigüedad; área \(\log_2\)}{54; 28.52932\allowbreak{}50129808}
\RouteRow{Primer bloque de separación}{órbita: \textemdash{}; celda: \textemdash{}; ruta: \textemdash{}}
\RouteRow{Ambigüedad final}{3}
\RouteGroup{Secuencias completas}
\RouteRow{\(w_6\)}{252 333 255 252 333 420 258 414 420 252 333 255 252 333 420 258 414 420}
\RouteRow{Tríadas estables por bloque}{0 1 2 3 4 5 6 7 8 9 10 11 12 12 14 14 16 16}
\end{tabularx}\endgroup\par\vspace{5pt}\noindent{\fontsize{8.4}{10}\selectfont\hyperref[trace-291]{Procedencia y códigos originales}\hfill Ficha completa · 291/324}
\clearpage\phantomsection\label{nav:vi-ruta-292}
% VI-CATALOG-ROW rutas 292
\pdfbookmark[3]{r9c1\_h+1v+1}{route-292}
\RouteTitle{Ruta 292}{r9c1\_\allowbreak{}h+1v+1}
\noindent Celda 73\quad (9, 1)\qquad Orientación \(h=1,\ v=1\)\hfill\hyperref[sec:vi-indice-rutas]{Índice}\par\vspace{4pt}
\begingroup\fontsize{10.2}{12.5}\selectfont\setlength{\tabcolsep}{5pt}\renewcommand{\arraystretch}{1.1}\rowcolors{1}{routepale}{white}\noindent\begin{tabularx}{\linewidth}{@{}>{\raggedright\arraybackslash}p{.345\linewidth}>{\raggedright\arraybackslash}X@{}}
\RouteGroup{Identidad estructural}
\RouteRow{Clase y multisección}{charged\_\allowbreak{}lepton\_\allowbreak{}G4\quad /\quad charged\_\allowbreak{}lepton\_\allowbreak{}G4}
\RouteRow{Colección y sector}{leptón cargado; leptón}
\RouteRow{Clasificación física}{leptón cargado}
\RouteRow{Generación, profundidad y color}{4; G4; \textemdash{}}
\RouteRow{Ruta primaria y órbita de inversión}{\(B_1\): no\qquad 1,\allowbreak{}9;\allowbreak{}9,\allowbreak{}1}
\RouteRow{Firma de clasificación}{28|\allowbreak{}6|\allowbreak{}0|\allowbreak{}0|\allowbreak{}-12|\allowbreak{}+++-\kern0pt{}-\kern0pt{}-+++-\kern0pt{}-\kern0pt{}-}
\RouteGroup{Retornos, memoria y transporte}
\RouteRow{Retornos con fase}{clase: 27; órbita: 27; celda: 27}
\RouteRow{Retorno cinemático}{en 108: 54; extensión: 216}
\RouteRow{Giros y cambios de hoja}{71; 107}
\RouteRow{Acarreo y diferimiento}{deuda total: 0; final: 0; bloques diferidos: 0}
\RouteRow{Tríadas finales; persistencia de Pareto}{17; no}
\RouteRow{\(K\longrightarrow K+9\)}{repeticiones: 9; cambios de memoria: 9}
\RouteGroup{Firma, fase y diferencias}
\RouteRow{\(n_A,\ s_C,\ \kappa_0,\ \nu_{120},\ \nu_{270}\)}{28, 6, 0, 2, -12}
\RouteRow{\(\tau_{12}\); firma histórica \(\mathbb Z^5\)}{+++-\kern0pt{}-\kern0pt{}-+++-\kern0pt{}-\kern0pt{}-; 28|\allowbreak{}6|\allowbreak{}0|\allowbreak{}2|\allowbreak{}-12}
\RouteRow{\(D_1,\ D_3,\ D_4,\ D_{27},\ D_{36}\)}{66, 60, 66, 48, 32}
\RouteRow{\(\Omega_{90/120},\ P_6\)}{80, 6}
\RouteGroup{Lectores y factores de acción}
\RouteRow{\(K_D\quad /\quad K_\Omega\)}{(80,\allowbreak{}66,\allowbreak{}3)\quad /\quad (80,\allowbreak{}54,\allowbreak{}6)}
\RouteRow{\(\kappa_D\)}{79.51870\allowbreak{}83196953\allowbreak{}45554341\allowbreak{}34}
\RouteRow{\(\kappa_\Omega\)}{79.60661\allowbreak{}01397948\allowbreak{}27586470\allowbreak{}91}
\RouteRow{\(R_{\mathrm{act},D}\)}{1.000000\allowbreak{}05512887\allowbreak{}17997050\allowbreak{}72}
\RouteRow{\(R_{\mathrm{act},\Omega}\)}{1.000000\allowbreak{}05518981\allowbreak{}25318591\allowbreak{}40}
\RouteGroup{Separación de prefijos}
\RouteRow{Área de ambigüedad; área \(\log_2\)}{54; 28.52932\allowbreak{}50129808}
\RouteRow{Primer bloque de separación}{órbita: \textemdash{}; celda: \textemdash{}; ruta: \textemdash{}}
\RouteRow{Ambigüedad final}{3}
\RouteGroup{Secuencias completas}
\RouteRow{\(w_6\)}{243 657 24 243 657 570 264 90 570 243 657 24 243 657 570 264 90 570}
\RouteRow{Tríadas estables por bloque}{0 1 2 3 4 5 6 7 8 9 10 11 12 13 14 15 16 17}
\end{tabularx}\endgroup\par\vspace{5pt}\noindent{\fontsize{8.4}{10}\selectfont\hyperref[trace-292]{Procedencia y códigos originales}\hfill Ficha completa · 292/324}
\clearpage\phantomsection\label{nav:vi-ruta-293}
% VI-CATALOG-ROW rutas 293
\pdfbookmark[2]{Celda 74}{cell-74}
\pdfbookmark[3]{r9c2\_h-1v-1}{route-293}
\RouteTitle{Ruta 293}{r9c2\_\allowbreak{}h-1v-1}
\noindent Celda 74\quad (9, 2)\qquad Orientación \(h=-1,\ v=-1\)\hfill\hyperref[sec:vi-indice-rutas]{Índice}\par\vspace{4pt}
\begingroup\fontsize{10.2}{12.5}\selectfont\setlength{\tabcolsep}{5pt}\renewcommand{\arraystretch}{1.1}\rowcolors{1}{routepale}{white}\noindent\begin{tabularx}{\linewidth}{@{}>{\raggedright\arraybackslash}p{.345\linewidth}>{\raggedright\arraybackslash}X@{}}
\RouteGroup{Identidad estructural}
\RouteRow{Clase y multisección}{neutrino\_\allowbreak{}G4\quad /\quad neutrino\_\allowbreak{}G4}
\RouteRow{Colección y sector}{neutrino; leptón}
\RouteRow{Clasificación física}{neutrino}
\RouteRow{Generación, profundidad y color}{4; G4; \textemdash{}}
\RouteRow{Ruta primaria y órbita de inversión}{\(B_1\): no\qquad 1,\allowbreak{}8;\allowbreak{}9,\allowbreak{}2}
\RouteRow{Firma de clasificación}{28|\allowbreak{}-6|\allowbreak{}6|\allowbreak{}-2|\allowbreak{}-8|\allowbreak{}++0-\kern0pt{}-\kern0pt{}-++0-\kern0pt{}-\kern0pt{}-}
\RouteGroup{Retornos, memoria y transporte}
\RouteRow{Retornos con fase}{clase: 27; órbita: 27; celda: 27}
\RouteRow{Retorno cinemático}{en 108: 54; extensión: 216}
\RouteRow{Giros y cambios de hoja}{71; 107}
\RouteRow{Acarreo y diferimiento}{deuda total: 1; final: 0; bloques diferidos: 1}
\RouteRow{Tríadas finales; persistencia de Pareto}{17; sí}
\RouteRow{\(K\longrightarrow K+9\)}{repeticiones: 9; cambios de memoria: 9}
\RouteGroup{Firma, fase y diferencias}
\RouteRow{\(n_A,\ s_C,\ \kappa_0,\ \nu_{120},\ \nu_{270}\)}{28, 6, -2, 0, -8}
\RouteRow{\(\tau_{12}\); firma histórica \(\mathbb Z^5\)}{+++-\kern0pt{}-0+++-\kern0pt{}-0; 28|\allowbreak{}6|\allowbreak{}-2|\allowbreak{}0|\allowbreak{}-8}
\RouteRow{\(D_1,\ D_3,\ D_4,\ D_{27},\ D_{36}\)}{66, 60, 66, 48, 32}
\RouteRow{\(\Omega_{90/120},\ P_6\)}{56, 5}
\RouteGroup{Lectores y factores de acción}
\RouteRow{\(K_D\quad /\quad K_\Omega\)}{(80,\allowbreak{}66,\allowbreak{}3)\quad /\quad (56,\allowbreak{}54,\allowbreak{}5)}
\RouteRow{\(\kappa_D\)}{79.51870\allowbreak{}83196953\allowbreak{}45554341\allowbreak{}34}
\RouteRow{\(\kappa_\Omega\)}{55.60649\allowbreak{}89433721\allowbreak{}35446416\allowbreak{}86}
\RouteRow{\(R_{\mathrm{act},D}\)}{1.000000\allowbreak{}05512887\allowbreak{}17997050\allowbreak{}72}
\RouteRow{\(R_{\mathrm{act},\Omega}\)}{1.000000\allowbreak{}03855097\allowbreak{}23541416\allowbreak{}45}
\RouteGroup{Separación de prefijos}
\RouteRow{Área de ambigüedad; área \(\log_2\)}{40; 20.33985\allowbreak{}00028846\allowbreak{}24}
\RouteRow{Primer bloque de separación}{órbita: \textemdash{}; celda: \textemdash{}; ruta: \textemdash{}}
\RouteRow{Ambigüedad final}{2}
\RouteGroup{Secuencias completas}
\RouteRow{\(w_6\)}{486 423 15 486 423 327 498 99 327 486 423 15 486 423 327 498 99 327}
\RouteRow{Tríadas estables por bloque}{0 1 2 3 4 5 6 7 8 9 10 10 12 13 14 15 16 17}
\end{tabularx}\endgroup\par\vspace{5pt}\noindent{\fontsize{8.4}{10}\selectfont\hyperref[trace-293]{Procedencia y códigos originales}\hfill Ficha completa · 293/324}
\clearpage\phantomsection\label{nav:vi-ruta-294}
% VI-CATALOG-ROW rutas 294
\pdfbookmark[3]{r9c2\_h-1v+1}{route-294}
\RouteTitle{Ruta 294}{r9c2\_\allowbreak{}h-1v+1}
\noindent Celda 74\quad (9, 2)\qquad Orientación \(h=-1,\ v=1\)\hfill\hyperref[sec:vi-indice-rutas]{Índice}\par\vspace{4pt}
\begingroup\fontsize{10.2}{12.5}\selectfont\setlength{\tabcolsep}{5pt}\renewcommand{\arraystretch}{1.1}\rowcolors{1}{routepale}{white}\noindent\begin{tabularx}{\linewidth}{@{}>{\raggedright\arraybackslash}p{.345\linewidth}>{\raggedright\arraybackslash}X@{}}
\RouteGroup{Identidad estructural}
\RouteRow{Clase y multisección}{neutrino\_\allowbreak{}G4\quad /\quad neutrino\_\allowbreak{}G4}
\RouteRow{Colección y sector}{neutrino; leptón}
\RouteRow{Clasificación física}{neutrino}
\RouteRow{Generación, profundidad y color}{4; G4; \textemdash{}}
\RouteRow{Ruta primaria y órbita de inversión}{\(B_1\): sí\qquad 1,\allowbreak{}8;\allowbreak{}9,\allowbreak{}2}
\RouteRow{Firma de clasificación}{28|\allowbreak{}-6|\allowbreak{}6|\allowbreak{}-2|\allowbreak{}-8|\allowbreak{}++0-\kern0pt{}-\kern0pt{}-++0-\kern0pt{}-\kern0pt{}-}
\RouteGroup{Retornos, memoria y transporte}
\RouteRow{Retornos con fase}{clase: 27; órbita: 27; celda: 27}
\RouteRow{Retorno cinemático}{en 108: 54; extensión: 216}
\RouteRow{Giros y cambios de hoja}{71; 107}
\RouteRow{Acarreo y diferimiento}{deuda total: 1; final: 0; bloques diferidos: 1}
\RouteRow{Tríadas finales; persistencia de Pareto}{17; sí}
\RouteRow{\(K\longrightarrow K+9\)}{repeticiones: 9; cambios de memoria: 9}
\RouteGroup{Firma, fase y diferencias}
\RouteRow{\(n_A,\ s_C,\ \kappa_0,\ \nu_{120},\ \nu_{270}\)}{14, -4, 0, 0, -12}
\RouteRow{\(\tau_{12}\); firma histórica \(\mathbb Z^5\)}{+++-\kern0pt{}-\kern0pt{}-+++-\kern0pt{}-\kern0pt{}-; 14|\allowbreak{}-4|\allowbreak{}0|\allowbreak{}0|\allowbreak{}-12}
\RouteRow{\(D_1,\ D_3,\ D_4,\ D_{27},\ D_{36}\)}{78, 24, 78, 24, 16}
\RouteRow{\(\Omega_{90/120},\ P_6\)}{80, 6}
\RouteGroup{Lectores y factores de acción}
\RouteRow{\(K_D\quad /\quad K_\Omega\)}{(40,\allowbreak{}78,\allowbreak{}27)\quad /\quad (80,\allowbreak{}54,\allowbreak{}6)}
\RouteRow{\(\kappa_D\)}{39.43380\allowbreak{}88030085\allowbreak{}51303671\allowbreak{}14}
\RouteRow{\(\kappa_\Omega\)}{79.60661\allowbreak{}01397948\allowbreak{}27586470\allowbreak{}91}
\RouteRow{\(R_{\mathrm{act},D}\)}{1.000000\allowbreak{}02733874\allowbreak{}08548943\allowbreak{}58}
\RouteRow{\(R_{\mathrm{act},\Omega}\)}{1.000000\allowbreak{}05518981\allowbreak{}25318591\allowbreak{}40}
\RouteGroup{Separación de prefijos}
\RouteRow{Área de ambigüedad; área \(\log_2\)}{36; 18.0}
\RouteRow{Primer bloque de separación}{órbita: \textemdash{}; celda: \textemdash{}; ruta: \textemdash{}}
\RouteRow{Ambigüedad final}{2}
\RouteGroup{Secuencias completas}
\RouteRow{\(w_6\)}{504 585 507 504 585 672 510 666 672 504 585 507 504 585 672 510 666 672}
\RouteRow{Tríadas estables por bloque}{0 1 2 3 4 5 6 7 8 9 10 11 12 13 14 14 16 17}
\end{tabularx}\endgroup\par\vspace{5pt}\noindent{\fontsize{8.4}{10}\selectfont\hyperref[trace-294]{Procedencia y códigos originales}\hfill Ficha completa · 294/324}
\clearpage\phantomsection\label{nav:vi-ruta-295}
% VI-CATALOG-ROW rutas 295
\pdfbookmark[3]{r9c2\_h+1v-1}{route-295}
\RouteTitle{Ruta 295}{r9c2\_\allowbreak{}h+1v-1}
\noindent Celda 74\quad (9, 2)\qquad Orientación \(h=1,\ v=-1\)\hfill\hyperref[sec:vi-indice-rutas]{Índice}\par\vspace{4pt}
\begingroup\fontsize{10.2}{12.5}\selectfont\setlength{\tabcolsep}{5pt}\renewcommand{\arraystretch}{1.1}\rowcolors{1}{routepale}{white}\noindent\begin{tabularx}{\linewidth}{@{}>{\raggedright\arraybackslash}p{.345\linewidth}>{\raggedright\arraybackslash}X@{}}
\RouteGroup{Identidad estructural}
\RouteRow{Clase y multisección}{neutrino\_\allowbreak{}G4\quad /\quad neutrino\_\allowbreak{}G4}
\RouteRow{Colección y sector}{neutrino; leptón}
\RouteRow{Clasificación física}{neutrino}
\RouteRow{Generación, profundidad y color}{4; G4; \textemdash{}}
\RouteRow{Ruta primaria y órbita de inversión}{\(B_1\): no\qquad 1,\allowbreak{}8;\allowbreak{}9,\allowbreak{}2}
\RouteRow{Firma de clasificación}{28|\allowbreak{}-6|\allowbreak{}6|\allowbreak{}-2|\allowbreak{}-8|\allowbreak{}++0-\kern0pt{}-\kern0pt{}-++0-\kern0pt{}-\kern0pt{}-}
\RouteGroup{Retornos, memoria y transporte}
\RouteRow{Retornos con fase}{clase: 27; órbita: 27; celda: 27}
\RouteRow{Retorno cinemático}{en 108: 54; extensión: 216}
\RouteRow{Giros y cambios de hoja}{71; 107}
\RouteRow{Acarreo y diferimiento}{deuda total: 1; final: 0; bloques diferidos: 1}
\RouteRow{Tríadas finales; persistencia de Pareto}{17; sí}
\RouteRow{\(K\longrightarrow K+9\)}{repeticiones: 9; cambios de memoria: 9}
\RouteGroup{Firma, fase y diferencias}
\RouteRow{\(n_A,\ s_C,\ \kappa_0,\ \nu_{120},\ \nu_{270}\)}{14, 4, 0, 0, -12}
\RouteRow{\(\tau_{12}\); firma histórica \(\mathbb Z^5\)}{+++-\kern0pt{}-\kern0pt{}-+++-\kern0pt{}-\kern0pt{}-; 14|\allowbreak{}4|\allowbreak{}0|\allowbreak{}0|\allowbreak{}-12}
\RouteRow{\(D_1,\ D_3,\ D_4,\ D_{27},\ D_{36}\)}{78, 24, 78, 24, 16}
\RouteRow{\(\Omega_{90/120},\ P_6\)}{80, 6}
\RouteGroup{Lectores y factores de acción}
\RouteRow{\(K_D\quad /\quad K_\Omega\)}{(40,\allowbreak{}78,\allowbreak{}27)\quad /\quad (80,\allowbreak{}54,\allowbreak{}6)}
\RouteRow{\(\kappa_D\)}{39.43380\allowbreak{}88030085\allowbreak{}51303671\allowbreak{}14}
\RouteRow{\(\kappa_\Omega\)}{79.60661\allowbreak{}01397948\allowbreak{}27586470\allowbreak{}91}
\RouteRow{\(R_{\mathrm{act},D}\)}{1.000000\allowbreak{}02733874\allowbreak{}08548943\allowbreak{}58}
\RouteRow{\(R_{\mathrm{act},\Omega}\)}{1.000000\allowbreak{}05518981\allowbreak{}25318591\allowbreak{}40}
\RouteGroup{Separación de prefijos}
\RouteRow{Área de ambigüedad; área \(\log_2\)}{36; 18.0}
\RouteRow{Primer bloque de separación}{órbita: \textemdash{}; celda: \textemdash{}; ruta: \textemdash{}}
\RouteRow{Ambigüedad final}{2}
\RouteGroup{Secuencias completas}
\RouteRow{\(w_6\)}{510 666 672 510 666 507 504 585 507 510 666 672 510 666 507 504 585 507}
\RouteRow{Tríadas estables por bloque}{0 1 2 3 4 5 6 7 8 9 10 11 12 13 14 15 15 17}
\end{tabularx}\endgroup\par\vspace{5pt}\noindent{\fontsize{8.4}{10}\selectfont\hyperref[trace-295]{Procedencia y códigos originales}\hfill Ficha completa · 295/324}
\clearpage\phantomsection\label{nav:vi-ruta-296}
% VI-CATALOG-ROW rutas 296
\pdfbookmark[3]{r9c2\_h+1v+1}{route-296}
\RouteTitle{Ruta 296}{r9c2\_\allowbreak{}h+1v+1}
\noindent Celda 74\quad (9, 2)\qquad Orientación \(h=1,\ v=1\)\hfill\hyperref[sec:vi-indice-rutas]{Índice}\par\vspace{4pt}
\begingroup\fontsize{10.2}{12.5}\selectfont\setlength{\tabcolsep}{5pt}\renewcommand{\arraystretch}{1.1}\rowcolors{1}{routepale}{white}\noindent\begin{tabularx}{\linewidth}{@{}>{\raggedright\arraybackslash}p{.345\linewidth}>{\raggedright\arraybackslash}X@{}}
\RouteGroup{Identidad estructural}
\RouteRow{Clase y multisección}{neutrino\_\allowbreak{}G4\quad /\quad neutrino\_\allowbreak{}G4}
\RouteRow{Colección y sector}{neutrino; leptón}
\RouteRow{Clasificación física}{neutrino}
\RouteRow{Generación, profundidad y color}{4; G4; \textemdash{}}
\RouteRow{Ruta primaria y órbita de inversión}{\(B_1\): no\qquad 1,\allowbreak{}8;\allowbreak{}9,\allowbreak{}2}
\RouteRow{Firma de clasificación}{28|\allowbreak{}-6|\allowbreak{}6|\allowbreak{}-2|\allowbreak{}-8|\allowbreak{}++0-\kern0pt{}-\kern0pt{}-++0-\kern0pt{}-\kern0pt{}-}
\RouteGroup{Retornos, memoria y transporte}
\RouteRow{Retornos con fase}{clase: 27; órbita: 27; celda: 27}
\RouteRow{Retorno cinemático}{en 108: 54; extensión: 216}
\RouteRow{Giros y cambios de hoja}{71; 107}
\RouteRow{Acarreo y diferimiento}{deuda total: 0; final: 0; bloques diferidos: 0}
\RouteRow{Tríadas finales; persistencia de Pareto}{17; sí}
\RouteRow{\(K\longrightarrow K+9\)}{repeticiones: 9; cambios de memoria: 9}
\RouteGroup{Firma, fase y diferencias}
\RouteRow{\(n_A,\ s_C,\ \kappa_0,\ \nu_{120},\ \nu_{270}\)}{28, -6, 2, 0, -8}
\RouteRow{\(\tau_{12}\); firma histórica \(\mathbb Z^5\)}{++0-\kern0pt{}-\kern0pt{}-++0-\kern0pt{}-\kern0pt{}-; 28|\allowbreak{}-6|\allowbreak{}2|\allowbreak{}0|\allowbreak{}-8}
\RouteRow{\(D_1,\ D_3,\ D_4,\ D_{27},\ D_{36}\)}{66, 60, 66, 48, 32}
\RouteRow{\(\Omega_{90/120},\ P_6\)}{56, 5}
\RouteGroup{Lectores y factores de acción}
\RouteRow{\(K_D\quad /\quad K_\Omega\)}{(80,\allowbreak{}66,\allowbreak{}3)\quad /\quad (56,\allowbreak{}54,\allowbreak{}5)}
\RouteRow{\(\kappa_D\)}{79.51870\allowbreak{}83196953\allowbreak{}45554341\allowbreak{}34}
\RouteRow{\(\kappa_\Omega\)}{55.60649\allowbreak{}89433721\allowbreak{}35446416\allowbreak{}86}
\RouteRow{\(R_{\mathrm{act},D}\)}{1.000000\allowbreak{}05512887\allowbreak{}17997050\allowbreak{}72}
\RouteRow{\(R_{\mathrm{act},\Omega}\)}{1.000000\allowbreak{}03855097\allowbreak{}23541416\allowbreak{}45}
\RouteGroup{Separación de prefijos}
\RouteRow{Área de ambigüedad; área \(\log_2\)}{36; 18.0}
\RouteRow{Primer bloque de separación}{órbita: \textemdash{}; celda: \textemdash{}; ruta: \textemdash{}}
\RouteRow{Ambigüedad final}{2}
\RouteGroup{Secuencias completas}
\RouteRow{\(w_6\)}{498 99 327 498 99 15 486 423 15 498 99 327 498 99 15 486 423 15}
\RouteRow{Tríadas estables por bloque}{0 1 2 3 4 5 6 7 8 9 10 11 12 13 14 15 16 17}
\end{tabularx}\endgroup\par\vspace{5pt}\noindent{\fontsize{8.4}{10}\selectfont\hyperref[trace-296]{Procedencia y códigos originales}\hfill Ficha completa · 296/324}
\clearpage\phantomsection\label{nav:vi-ruta-297}
% VI-CATALOG-ROW rutas 297
\pdfbookmark[2]{Celda 75}{cell-75}
\pdfbookmark[3]{r9c3\_h-1v-1}{route-297}
\RouteTitle{Ruta 297}{r9c3\_\allowbreak{}h-1v-1}
\noindent Celda 75\quad (9, 3)\qquad Orientación \(h=-1,\ v=-1\)\hfill\hyperref[sec:vi-indice-rutas]{Índice}\par\vspace{4pt}
\begingroup\fontsize{10.2}{12.5}\selectfont\setlength{\tabcolsep}{5pt}\renewcommand{\arraystretch}{1.1}\rowcolors{1}{routepale}{white}\noindent\begin{tabularx}{\linewidth}{@{}>{\raggedright\arraybackslash}p{.345\linewidth}>{\raggedright\arraybackslash}X@{}}
\RouteGroup{Identidad estructural}
\RouteRow{Clase y multisección}{charged\_\allowbreak{}lepton\_\allowbreak{}G4\quad /\quad charged\_\allowbreak{}lepton\_\allowbreak{}G4}
\RouteRow{Colección y sector}{leptón cargado; leptón}
\RouteRow{Clasificación física}{leptón cargado}
\RouteRow{Generación, profundidad y color}{4; G4; \textemdash{}}
\RouteRow{Ruta primaria y órbita de inversión}{\(B_1\): no\qquad 1,\allowbreak{}7;\allowbreak{}9,\allowbreak{}3}
\RouteRow{Firma de clasificación}{30|\allowbreak{}-6|\allowbreak{}0|\allowbreak{}-2|\allowbreak{}-8|\allowbreak{}0++-\kern0pt{}-\kern0pt{}-0++-\kern0pt{}-\kern0pt{}-}
\RouteGroup{Retornos, memoria y transporte}
\RouteRow{Retornos con fase}{clase: 27; órbita: 27; celda: 27}
\RouteRow{Retorno cinemático}{en 108: 54; extensión: 216}
\RouteRow{Giros y cambios de hoja}{71; 107}
\RouteRow{Acarreo y diferimiento}{deuda total: 1; final: 0; bloques diferidos: 1}
\RouteRow{Tríadas finales; persistencia de Pareto}{17; sí}
\RouteRow{\(K\longrightarrow K+9\)}{repeticiones: 9; cambios de memoria: 9}
\RouteGroup{Firma, fase y diferencias}
\RouteRow{\(n_A,\ s_C,\ \kappa_0,\ \nu_{120},\ \nu_{270}\)}{30, 6, -2, 4, -8}
\RouteRow{\(\tau_{12}\); firma histórica \(\mathbb Z^5\)}{+++0-\kern0pt{}-+++0-\kern0pt{}-; 30|\allowbreak{}6|\allowbreak{}-2|\allowbreak{}4|\allowbreak{}-8}
\RouteRow{\(D_1,\ D_3,\ D_4,\ D_{27},\ D_{36}\)}{54, 60, 72, 24, 16}
\RouteRow{\(\Omega_{90/120},\ P_6\)}{56, 5}
\RouteGroup{Lectores y factores de acción}
\RouteRow{\(K_D\quad /\quad K_\Omega\)}{(40,\allowbreak{}54,\allowbreak{}6)\quad /\quad (56,\allowbreak{}54,\allowbreak{}5)}
\RouteRow{\(\kappa_D\)}{39.60661\allowbreak{}01397948\allowbreak{}27586470\allowbreak{}91}
\RouteRow{\(\kappa_\Omega\)}{55.60649\allowbreak{}89433721\allowbreak{}35446416\allowbreak{}86}
\RouteRow{\(R_{\mathrm{act},D}\)}{1.000000\allowbreak{}02745854\allowbreak{}08735595\allowbreak{}18}
\RouteRow{\(R_{\mathrm{act},\Omega}\)}{1.000000\allowbreak{}03855097\allowbreak{}23541416\allowbreak{}45}
\RouteGroup{Separación de prefijos}
\RouteRow{Área de ambigüedad; área \(\log_2\)}{58; 30.18947\allowbreak{}50100961\allowbreak{}74}
\RouteRow{Primer bloque de separación}{órbita: \textemdash{}; celda: \textemdash{}; ruta: \textemdash{}}
\RouteRow{Ambigüedad final}{3}
\RouteGroup{Secuencias completas}
\RouteRow{\(w_6\)}{15 486 423 15 486 657 24 243 657 15 486 423 15 486 657 24 243 657}
\RouteRow{Tríadas estables por bloque}{0 1 2 3 4 5 6 7 8 9 10 11 12 13 14 15 15 17}
\end{tabularx}\endgroup\par\vspace{5pt}\noindent{\fontsize{8.4}{10}\selectfont\hyperref[trace-297]{Procedencia y códigos originales}\hfill Ficha completa · 297/324}
\clearpage\phantomsection\label{nav:vi-ruta-298}
% VI-CATALOG-ROW rutas 298
\pdfbookmark[3]{r9c3\_h-1v+1}{route-298}
\RouteTitle{Ruta 298}{r9c3\_\allowbreak{}h-1v+1}
\noindent Celda 75\quad (9, 3)\qquad Orientación \(h=-1,\ v=1\)\hfill\hyperref[sec:vi-indice-rutas]{Índice}\par\vspace{4pt}
\begingroup\fontsize{10.2}{12.5}\selectfont\setlength{\tabcolsep}{5pt}\renewcommand{\arraystretch}{1.1}\rowcolors{1}{routepale}{white}\noindent\begin{tabularx}{\linewidth}{@{}>{\raggedright\arraybackslash}p{.345\linewidth}>{\raggedright\arraybackslash}X@{}}
\RouteGroup{Identidad estructural}
\RouteRow{Clase y multisección}{charged\_\allowbreak{}lepton\_\allowbreak{}G4\quad /\quad charged\_\allowbreak{}lepton\_\allowbreak{}G4}
\RouteRow{Colección y sector}{leptón cargado; leptón}
\RouteRow{Clasificación física}{leptón cargado}
\RouteRow{Generación, profundidad y color}{4; G4; \textemdash{}}
\RouteRow{Ruta primaria y órbita de inversión}{\(B_1\): no\qquad 1,\allowbreak{}7;\allowbreak{}9,\allowbreak{}3}
\RouteRow{Firma de clasificación}{30|\allowbreak{}-6|\allowbreak{}0|\allowbreak{}-2|\allowbreak{}-8|\allowbreak{}0++-\kern0pt{}-\kern0pt{}-0++-\kern0pt{}-\kern0pt{}-}
\RouteGroup{Retornos, memoria y transporte}
\RouteRow{Retornos con fase}{clase: 9; órbita: 27; celda: 27}
\RouteRow{Retorno cinemático}{en 108: 54; extensión: 216}
\RouteRow{Giros y cambios de hoja}{71; 107}
\RouteRow{Acarreo y diferimiento}{deuda total: 0; final: 0; bloques diferidos: 0}
\RouteRow{Tríadas finales; persistencia de Pareto}{17; sí}
\RouteRow{\(K\longrightarrow K+9\)}{repeticiones: 9; cambios de memoria: 9}
\RouteGroup{Firma, fase y diferencias}
\RouteRow{\(n_A,\ s_C,\ \kappa_0,\ \nu_{120},\ \nu_{270}\)}{50, -6, 2, -4, -8}
\RouteRow{\(\tau_{12}\); firma histórica \(\mathbb Z^5\)}{0++-\kern0pt{}-\kern0pt{}-0++-\kern0pt{}-\kern0pt{}-; 50|\allowbreak{}-6|\allowbreak{}2|\allowbreak{}-4|\allowbreak{}-8}
\RouteRow{\(D_1,\ D_3,\ D_4,\ D_{27},\ D_{36}\)}{24, 24, 24, 0, 0}
\RouteRow{\(\Omega_{90/120},\ P_6\)}{56, 5}
\RouteGroup{Lectores y factores de acción}
\RouteRow{\(K_D\quad /\quad K_\Omega\)}{(0,\allowbreak{}24,\allowbreak{}0)\quad /\quad (56,\allowbreak{}54,\allowbreak{}5)}
\RouteRow{\(\kappa_D\)}{-0.175136\allowbreak{}46166281\allowbreak{}12239348\allowbreak{}425}
\RouteRow{\(\kappa_\Omega\)}{55.60649\allowbreak{}89433721\allowbreak{}35446416\allowbreak{}86}
\RouteRow{\(R_{\mathrm{act},D}\)}{0.999999\allowbreak{}99987858\allowbreak{}10851337\allowbreak{}881}
\RouteRow{\(R_{\mathrm{act},\Omega}\)}{1.000000\allowbreak{}03855097\allowbreak{}23541416\allowbreak{}45}
\RouteGroup{Separación de prefijos}
\RouteRow{Área de ambigüedad; área \(\log_2\)}{108; 46.52932\allowbreak{}50129807\allowbreak{}9}
\RouteRow{Primer bloque de separación}{órbita: \textemdash{}; celda: \textemdash{}; ruta: \textemdash{}}
\RouteRow{Ambigüedad final}{6}
\RouteGroup{Secuencias completas}
\RouteRow{\(w_6\)}{3 0 81 3 0 81 3 0 81 3 0 81 3 0 81 3 0 81}
\RouteRow{Tríadas estables por bloque}{0 1 2 3 4 5 6 7 8 9 10 11 12 13 14 15 16 17}
\end{tabularx}\endgroup\par\vspace{5pt}\noindent{\fontsize{8.4}{10}\selectfont\hyperref[trace-298]{Procedencia y códigos originales}\hfill Ficha completa · 298/324}
\clearpage\phantomsection\label{nav:vi-ruta-299}
% VI-CATALOG-ROW rutas 299
\pdfbookmark[3]{r9c3\_h+1v-1}{route-299}
\RouteTitle{Ruta 299}{r9c3\_\allowbreak{}h+1v-1}
\noindent Celda 75\quad (9, 3)\qquad Orientación \(h=1,\ v=-1\)\hfill\hyperref[sec:vi-indice-rutas]{Índice}\par\vspace{4pt}
\begingroup\fontsize{10.2}{12.5}\selectfont\setlength{\tabcolsep}{5pt}\renewcommand{\arraystretch}{1.1}\rowcolors{1}{routepale}{white}\noindent\begin{tabularx}{\linewidth}{@{}>{\raggedright\arraybackslash}p{.345\linewidth}>{\raggedright\arraybackslash}X@{}}
\RouteGroup{Identidad estructural}
\RouteRow{Clase y multisección}{charged\_\allowbreak{}lepton\_\allowbreak{}G4\quad /\quad charged\_\allowbreak{}lepton\_\allowbreak{}G4}
\RouteRow{Colección y sector}{leptón cargado; leptón}
\RouteRow{Clasificación física}{leptón cargado}
\RouteRow{Generación, profundidad y color}{4; G4; \textemdash{}}
\RouteRow{Ruta primaria y órbita de inversión}{\(B_1\): no\qquad 1,\allowbreak{}7;\allowbreak{}9,\allowbreak{}3}
\RouteRow{Firma de clasificación}{30|\allowbreak{}-6|\allowbreak{}0|\allowbreak{}-2|\allowbreak{}-8|\allowbreak{}0++-\kern0pt{}-\kern0pt{}-0++-\kern0pt{}-\kern0pt{}-}
\RouteGroup{Retornos, memoria y transporte}
\RouteRow{Retornos con fase}{clase: 18; órbita: 27; celda: 27}
\RouteRow{Retorno cinemático}{en 108: 54; extensión: 216}
\RouteRow{Giros y cambios de hoja}{71; 107}
\RouteRow{Acarreo y diferimiento}{deuda total: 0; final: 0; bloques diferidos: 0}
\RouteRow{Tríadas finales; persistencia de Pareto}{17; sí}
\RouteRow{\(K\longrightarrow K+9\)}{repeticiones: 9; cambios de memoria: 9}
\RouteGroup{Firma, fase y diferencias}
\RouteRow{\(n_A,\ s_C,\ \kappa_0,\ \nu_{120},\ \nu_{270}\)}{50, 6, -2, 4, -8}
\RouteRow{\(\tau_{12}\); firma histórica \(\mathbb Z^5\)}{+++0-\kern0pt{}-+++0-\kern0pt{}-; 50|\allowbreak{}6|\allowbreak{}-2|\allowbreak{}4|\allowbreak{}-8}
\RouteRow{\(D_1,\ D_3,\ D_4,\ D_{27},\ D_{36}\)}{24, 24, 24, 0, 0}
\RouteRow{\(\Omega_{90/120},\ P_6\)}{56, 5}
\RouteGroup{Lectores y factores de acción}
\RouteRow{\(K_D\quad /\quad K_\Omega\)}{(0,\allowbreak{}24,\allowbreak{}0)\quad /\quad (56,\allowbreak{}54,\allowbreak{}5)}
\RouteRow{\(\kappa_D\)}{-0.175136\allowbreak{}46166281\allowbreak{}12239348\allowbreak{}425}
\RouteRow{\(\kappa_\Omega\)}{55.60649\allowbreak{}89433721\allowbreak{}35446416\allowbreak{}86}
\RouteRow{\(R_{\mathrm{act},D}\)}{0.999999\allowbreak{}99987858\allowbreak{}10851337\allowbreak{}881}
\RouteRow{\(R_{\mathrm{act},\Omega}\)}{1.000000\allowbreak{}03855097\allowbreak{}23541416\allowbreak{}45}
\RouteGroup{Separación de prefijos}
\RouteRow{Área de ambigüedad; área \(\log_2\)}{108; 46.52932\allowbreak{}50129807\allowbreak{}9}
\RouteRow{Primer bloque de separación}{órbita: \textemdash{}; celda: \textemdash{}; ruta: \textemdash{}}
\RouteRow{Ambigüedad final}{6}
\RouteGroup{Secuencias completas}
\RouteRow{\(w_6\)}{3 0 81 3 0 81 3 0 81 3 0 81 3 0 81 3 0 81}
\RouteRow{Tríadas estables por bloque}{0 1 2 3 4 5 6 7 8 9 10 11 12 13 14 15 16 17}
\end{tabularx}\endgroup\par\vspace{5pt}\noindent{\fontsize{8.4}{10}\selectfont\hyperref[trace-299]{Procedencia y códigos originales}\hfill Ficha completa · 299/324}
\clearpage\phantomsection\label{nav:vi-ruta-300}
% VI-CATALOG-ROW rutas 300
\pdfbookmark[3]{r9c3\_h+1v+1}{route-300}
\RouteTitle{Ruta 300}{r9c3\_\allowbreak{}h+1v+1}
\noindent Celda 75\quad (9, 3)\qquad Orientación \(h=1,\ v=1\)\hfill\hyperref[sec:vi-indice-rutas]{Índice}\par\vspace{4pt}
\begingroup\fontsize{10.2}{12.5}\selectfont\setlength{\tabcolsep}{5pt}\renewcommand{\arraystretch}{1.1}\rowcolors{1}{routepale}{white}\noindent\begin{tabularx}{\linewidth}{@{}>{\raggedright\arraybackslash}p{.345\linewidth}>{\raggedright\arraybackslash}X@{}}
\RouteGroup{Identidad estructural}
\RouteRow{Clase y multisección}{charged\_\allowbreak{}lepton\_\allowbreak{}G4\quad /\quad charged\_\allowbreak{}lepton\_\allowbreak{}G4}
\RouteRow{Colección y sector}{leptón cargado; leptón}
\RouteRow{Clasificación física}{leptón cargado}
\RouteRow{Generación, profundidad y color}{4; G4; \textemdash{}}
\RouteRow{Ruta primaria y órbita de inversión}{\(B_1\): sí\qquad 1,\allowbreak{}7;\allowbreak{}9,\allowbreak{}3}
\RouteRow{Firma de clasificación}{30|\allowbreak{}-6|\allowbreak{}0|\allowbreak{}-2|\allowbreak{}-8|\allowbreak{}0++-\kern0pt{}-\kern0pt{}-0++-\kern0pt{}-\kern0pt{}-}
\RouteGroup{Retornos, memoria y transporte}
\RouteRow{Retornos con fase}{clase: 27; órbita: 27; celda: 27}
\RouteRow{Retorno cinemático}{en 108: 54; extensión: 216}
\RouteRow{Giros y cambios de hoja}{71; 107}
\RouteRow{Acarreo y diferimiento}{deuda total: 2; final: 1; bloques diferidos: 2}
\RouteRow{Tríadas finales; persistencia de Pareto}{16; sí}
\RouteRow{\(K\longrightarrow K+9\)}{repeticiones: 9; cambios de memoria: 9}
\RouteGroup{Firma, fase y diferencias}
\RouteRow{\(n_A,\ s_C,\ \kappa_0,\ \nu_{120},\ \nu_{270}\)}{30, -6, 2, -4, -8}
\RouteRow{\(\tau_{12}\); firma histórica \(\mathbb Z^5\)}{0++-\kern0pt{}-\kern0pt{}-0++-\kern0pt{}-\kern0pt{}-; 30|\allowbreak{}-6|\allowbreak{}2|\allowbreak{}-4|\allowbreak{}-8}
\RouteRow{\(D_1,\ D_3,\ D_4,\ D_{27},\ D_{36}\)}{54, 60, 72, 24, 16}
\RouteRow{\(\Omega_{90/120},\ P_6\)}{56, 5}
\RouteGroup{Lectores y factores de acción}
\RouteRow{\(K_D\quad /\quad K_\Omega\)}{(40,\allowbreak{}54,\allowbreak{}6)\quad /\quad (56,\allowbreak{}54,\allowbreak{}5)}
\RouteRow{\(\kappa_D\)}{39.60661\allowbreak{}01397948\allowbreak{}27586470\allowbreak{}91}
\RouteRow{\(\kappa_\Omega\)}{55.60649\allowbreak{}89433721\allowbreak{}35446416\allowbreak{}86}
\RouteRow{\(R_{\mathrm{act},D}\)}{1.000000\allowbreak{}02745854\allowbreak{}08735595\allowbreak{}18}
\RouteRow{\(R_{\mathrm{act},\Omega}\)}{1.000000\allowbreak{}03855097\allowbreak{}23541416\allowbreak{}45}
\RouteGroup{Separación de prefijos}
\RouteRow{Área de ambigüedad; área \(\log_2\)}{58; 30.18947\allowbreak{}50100961\allowbreak{}74}
\RouteRow{Primer bloque de separación}{órbita: \textemdash{}; celda: \textemdash{}; ruta: \textemdash{}}
\RouteRow{Ambigüedad final}{3}
\RouteGroup{Secuencias completas}
\RouteRow{\(w_6\)}{24 243 657 24 243 423 15 486 423 24 243 657 24 243 423 15 486 423}
\RouteRow{Tríadas estables por bloque}{0 1 2 3 4 5 6 7 8 9 10 11 12 13 13 15 16 16}
\end{tabularx}\endgroup\par\vspace{5pt}\noindent{\fontsize{8.4}{10}\selectfont\hyperref[trace-300]{Procedencia y códigos originales}\hfill Ficha completa · 300/324}
\clearpage\phantomsection\label{nav:vi-ruta-301}
% VI-CATALOG-ROW rutas 301
\pdfbookmark[2]{Celda 76}{cell-76}
\pdfbookmark[3]{r9c4\_h-1v-1}{route-301}
\RouteTitle{Ruta 301}{r9c4\_\allowbreak{}h-1v-1}
\noindent Celda 76\quad (9, 4)\qquad Orientación \(h=-1,\ v=-1\)\hfill\hyperref[sec:vi-indice-rutas]{Índice}\par\vspace{4pt}
\begingroup\fontsize{10.2}{12.5}\selectfont\setlength{\tabcolsep}{5pt}\renewcommand{\arraystretch}{1.1}\rowcolors{1}{routepale}{white}\noindent\begin{tabularx}{\linewidth}{@{}>{\raggedright\arraybackslash}p{.345\linewidth}>{\raggedright\arraybackslash}X@{}}
\RouteGroup{Identidad estructural}
\RouteRow{Clase y multisección}{neutrino\_\allowbreak{}G4\quad /\quad neutrino\_\allowbreak{}G4}
\RouteRow{Colección y sector}{neutrino; leptón}
\RouteRow{Clasificación física}{neutrino}
\RouteRow{Generación, profundidad y color}{4; G4; \textemdash{}}
\RouteRow{Ruta primaria y órbita de inversión}{\(B_1\): no\qquad 1,\allowbreak{}6;\allowbreak{}9,\allowbreak{}4}
\RouteRow{Firma de clasificación}{34|\allowbreak{}0|\allowbreak{}-12|\allowbreak{}0|\allowbreak{}-12|\allowbreak{}+++-\kern0pt{}-\kern0pt{}-+++-\kern0pt{}-\kern0pt{}-}
\RouteGroup{Retornos, memoria y transporte}
\RouteRow{Retornos con fase}{clase: 27; órbita: 27; celda: 27}
\RouteRow{Retorno cinemático}{en 108: 54; extensión: 216}
\RouteRow{Giros y cambios de hoja}{71; 107}
\RouteRow{Acarreo y diferimiento}{deuda total: 2; final: 1; bloques diferidos: 1}
\RouteRow{Tríadas finales; persistencia de Pareto}{16; sí}
\RouteRow{\(K\longrightarrow K+9\)}{repeticiones: 9; cambios de memoria: 9}
\RouteGroup{Firma, fase y diferencias}
\RouteRow{\(n_A,\ s_C,\ \kappa_0,\ \nu_{120},\ \nu_{270}\)}{34, 0, 0, 0, -12}
\RouteRow{\(\tau_{12}\); firma histórica \(\mathbb Z^5\)}{+++-\kern0pt{}-\kern0pt{}-+++-\kern0pt{}-\kern0pt{}-; 34|\allowbreak{}0|\allowbreak{}0|\allowbreak{}0|\allowbreak{}-12}
\RouteRow{\(D_1,\ D_3,\ D_4,\ D_{27},\ D_{36}\)}{66, 60, 66, 48, 32}
\RouteRow{\(\Omega_{90/120},\ P_6\)}{80, 6}
\RouteGroup{Lectores y factores de acción}
\RouteRow{\(K_D\quad /\quad K_\Omega\)}{(80,\allowbreak{}66,\allowbreak{}3)\quad /\quad (80,\allowbreak{}54,\allowbreak{}6)}
\RouteRow{\(\kappa_D\)}{79.51870\allowbreak{}83196953\allowbreak{}45554341\allowbreak{}34}
\RouteRow{\(\kappa_\Omega\)}{79.60661\allowbreak{}01397948\allowbreak{}27586470\allowbreak{}91}
\RouteRow{\(R_{\mathrm{act},D}\)}{1.000000\allowbreak{}05512887\allowbreak{}17997050\allowbreak{}72}
\RouteRow{\(R_{\mathrm{act},\Omega}\)}{1.000000\allowbreak{}05518981\allowbreak{}25318591\allowbreak{}40}
\RouteGroup{Separación de prefijos}
\RouteRow{Área de ambigüedad; área \(\log_2\)}{18; 0.0}
\RouteRow{Primer bloque de separación}{órbita: 1; celda: 1; ruta: 1}
\RouteRow{Ambigüedad final}{1}
\RouteGroup{Secuencias completas}
\RouteRow{\(w_6\)}{264 90 570 264 90 24 243 657 24 264 90 570 264 90 24 243 657 24}
\RouteRow{Tríadas estables por bloque}{0 1 2 3 4 5 6 7 8 9 10 11 12 13 14 15 15 16}
\end{tabularx}\endgroup\par\vspace{5pt}\noindent{\fontsize{8.4}{10}\selectfont\hyperref[trace-301]{Procedencia y códigos originales}\hfill Ficha completa · 301/324}
\clearpage\phantomsection\label{nav:vi-ruta-302}
% VI-CATALOG-ROW rutas 302
\pdfbookmark[3]{r9c4\_h-1v+1}{route-302}
\RouteTitle{Ruta 302}{r9c4\_\allowbreak{}h-1v+1}
\noindent Celda 76\quad (9, 4)\qquad Orientación \(h=-1,\ v=1\)\hfill\hyperref[sec:vi-indice-rutas]{Índice}\par\vspace{4pt}
\begingroup\fontsize{10.2}{12.5}\selectfont\setlength{\tabcolsep}{5pt}\renewcommand{\arraystretch}{1.1}\rowcolors{1}{routepale}{white}\noindent\begin{tabularx}{\linewidth}{@{}>{\raggedright\arraybackslash}p{.345\linewidth}>{\raggedright\arraybackslash}X@{}}
\RouteGroup{Identidad estructural}
\RouteRow{Clase y multisección}{neutrino\_\allowbreak{}G4\quad /\quad neutrino\_\allowbreak{}G4}
\RouteRow{Colección y sector}{neutrino; leptón}
\RouteRow{Clasificación física}{neutrino}
\RouteRow{Generación, profundidad y color}{4; G4; \textemdash{}}
\RouteRow{Ruta primaria y órbita de inversión}{\(B_1\): sí\qquad 1,\allowbreak{}6;\allowbreak{}9,\allowbreak{}4}
\RouteRow{Firma de clasificación}{34|\allowbreak{}0|\allowbreak{}-12|\allowbreak{}0|\allowbreak{}-12|\allowbreak{}+++-\kern0pt{}-\kern0pt{}-+++-\kern0pt{}-\kern0pt{}-}
\RouteGroup{Retornos, memoria y transporte}
\RouteRow{Retornos con fase}{clase: 27; órbita: 27; celda: 27}
\RouteRow{Retorno cinemático}{en 108: 54; extensión: 216}
\RouteRow{Giros y cambios de hoja}{71; 107}
\RouteRow{Acarreo y diferimiento}{deuda total: 2; final: 0; bloques diferidos: 1}
\RouteRow{Tríadas finales; persistencia de Pareto}{17; sí}
\RouteRow{\(K\longrightarrow K+9\)}{repeticiones: 9; cambios de memoria: 9}
\RouteGroup{Firma, fase y diferencias}
\RouteRow{\(n_A,\ s_C,\ \kappa_0,\ \nu_{120},\ \nu_{270}\)}{8, -2, 0, 0, -12}
\RouteRow{\(\tau_{12}\); firma histórica \(\mathbb Z^5\)}{+++-\kern0pt{}-\kern0pt{}-+++-\kern0pt{}-\kern0pt{}-; 8|\allowbreak{}-2|\allowbreak{}0|\allowbreak{}0|\allowbreak{}-12}
\RouteRow{\(D_1,\ D_3,\ D_4,\ D_{27},\ D_{36}\)}{84, 24, 84, 24, 16}
\RouteRow{\(\Omega_{90/120},\ P_6\)}{80, 6}
\RouteGroup{Lectores y factores de acción}
\RouteRow{\(K_D\quad /\quad K_\Omega\)}{(40,\allowbreak{}84,\allowbreak{}30)\quad /\quad (80,\allowbreak{}54,\allowbreak{}6)}
\RouteRow{\(\kappa_D\)}{39.39035\allowbreak{}82768609\allowbreak{}24917849\allowbreak{}59}
\RouteRow{\(\kappa_\Omega\)}{79.60661\allowbreak{}01397948\allowbreak{}27586470\allowbreak{}91}
\RouteRow{\(R_{\mathrm{act},D}\)}{1.000000\allowbreak{}02730861\allowbreak{}73967087\allowbreak{}05}
\RouteRow{\(R_{\mathrm{act},\Omega}\)}{1.000000\allowbreak{}05518981\allowbreak{}25318591\allowbreak{}40}
\RouteGroup{Separación de prefijos}
\RouteRow{Área de ambigüedad; área \(\log_2\)}{18; 0.0}
\RouteRow{Primer bloque de separación}{órbita: 1; celda: 1; ruta: 1}
\RouteRow{Ambigüedad final}{1}
\RouteGroup{Secuencias completas}
\RouteRow{\(w_6\)}{258 414 420 258 414 255 252 333 255 258 414 420 258 414 255 252 333 255}
\RouteRow{Tríadas estables por bloque}{0 1 2 3 4 5 6 7 8 9 10 11 11 12 14 15 16 17}
\end{tabularx}\endgroup\par\vspace{5pt}\noindent{\fontsize{8.4}{10}\selectfont\hyperref[trace-302]{Procedencia y códigos originales}\hfill Ficha completa · 302/324}
\clearpage\phantomsection\label{nav:vi-ruta-303}
% VI-CATALOG-ROW rutas 303
\pdfbookmark[3]{r9c4\_h+1v-1}{route-303}
\RouteTitle{Ruta 303}{r9c4\_\allowbreak{}h+1v-1}
\noindent Celda 76\quad (9, 4)\qquad Orientación \(h=1,\ v=-1\)\hfill\hyperref[sec:vi-indice-rutas]{Índice}\par\vspace{4pt}
\begingroup\fontsize{10.2}{12.5}\selectfont\setlength{\tabcolsep}{5pt}\renewcommand{\arraystretch}{1.1}\rowcolors{1}{routepale}{white}\noindent\begin{tabularx}{\linewidth}{@{}>{\raggedright\arraybackslash}p{.345\linewidth}>{\raggedright\arraybackslash}X@{}}
\RouteGroup{Identidad estructural}
\RouteRow{Clase y multisección}{neutrino\_\allowbreak{}G4\quad /\quad neutrino\_\allowbreak{}G4}
\RouteRow{Colección y sector}{neutrino; leptón}
\RouteRow{Clasificación física}{neutrino}
\RouteRow{Generación, profundidad y color}{4; G4; \textemdash{}}
\RouteRow{Ruta primaria y órbita de inversión}{\(B_1\): no\qquad 1,\allowbreak{}6;\allowbreak{}9,\allowbreak{}4}
\RouteRow{Firma de clasificación}{34|\allowbreak{}0|\allowbreak{}-12|\allowbreak{}0|\allowbreak{}-12|\allowbreak{}+++-\kern0pt{}-\kern0pt{}-+++-\kern0pt{}-\kern0pt{}-}
\RouteGroup{Retornos, memoria y transporte}
\RouteRow{Retornos con fase}{clase: 27; órbita: 27; celda: 27}
\RouteRow{Retorno cinemático}{en 108: 54; extensión: 216}
\RouteRow{Giros y cambios de hoja}{71; 107}
\RouteRow{Acarreo y diferimiento}{deuda total: 3; final: 1; bloques diferidos: 3}
\RouteRow{Tríadas finales; persistencia de Pareto}{16; sí}
\RouteRow{\(K\longrightarrow K+9\)}{repeticiones: 9; cambios de memoria: 9}
\RouteGroup{Firma, fase y diferencias}
\RouteRow{\(n_A,\ s_C,\ \kappa_0,\ \nu_{120},\ \nu_{270}\)}{8, 2, 0, 0, -12}
\RouteRow{\(\tau_{12}\); firma histórica \(\mathbb Z^5\)}{+++-\kern0pt{}-\kern0pt{}-+++-\kern0pt{}-\kern0pt{}-; 8|\allowbreak{}2|\allowbreak{}0|\allowbreak{}0|\allowbreak{}-12}
\RouteRow{\(D_1,\ D_3,\ D_4,\ D_{27},\ D_{36}\)}{84, 24, 84, 24, 16}
\RouteRow{\(\Omega_{90/120},\ P_6\)}{80, 6}
\RouteGroup{Lectores y factores de acción}
\RouteRow{\(K_D\quad /\quad K_\Omega\)}{(40,\allowbreak{}84,\allowbreak{}30)\quad /\quad (80,\allowbreak{}54,\allowbreak{}6)}
\RouteRow{\(\kappa_D\)}{39.39035\allowbreak{}82768609\allowbreak{}24917849\allowbreak{}59}
\RouteRow{\(\kappa_\Omega\)}{79.60661\allowbreak{}01397948\allowbreak{}27586470\allowbreak{}91}
\RouteRow{\(R_{\mathrm{act},D}\)}{1.000000\allowbreak{}02730861\allowbreak{}73967087\allowbreak{}05}
\RouteRow{\(R_{\mathrm{act},\Omega}\)}{1.000000\allowbreak{}05518981\allowbreak{}25318591\allowbreak{}40}
\RouteGroup{Separación de prefijos}
\RouteRow{Área de ambigüedad; área \(\log_2\)}{18; 0.0}
\RouteRow{Primer bloque de separación}{órbita: 1; celda: 1; ruta: 1}
\RouteRow{Ambigüedad final}{1}
\RouteGroup{Secuencias completas}
\RouteRow{\(w_6\)}{252 333 255 252 333 420 258 414 420 252 333 255 252 333 420 258 414 420}
\RouteRow{Tríadas estables por bloque}{0 1 2 3 4 5 6 7 8 9 10 11 12 12 14 14 16 16}
\end{tabularx}\endgroup\par\vspace{5pt}\noindent{\fontsize{8.4}{10}\selectfont\hyperref[trace-303]{Procedencia y códigos originales}\hfill Ficha completa · 303/324}
\clearpage\phantomsection\label{nav:vi-ruta-304}
% VI-CATALOG-ROW rutas 304
\pdfbookmark[3]{r9c4\_h+1v+1}{route-304}
\RouteTitle{Ruta 304}{r9c4\_\allowbreak{}h+1v+1}
\noindent Celda 76\quad (9, 4)\qquad Orientación \(h=1,\ v=1\)\hfill\hyperref[sec:vi-indice-rutas]{Índice}\par\vspace{4pt}
\begingroup\fontsize{10.2}{12.5}\selectfont\setlength{\tabcolsep}{5pt}\renewcommand{\arraystretch}{1.1}\rowcolors{1}{routepale}{white}\noindent\begin{tabularx}{\linewidth}{@{}>{\raggedright\arraybackslash}p{.345\linewidth}>{\raggedright\arraybackslash}X@{}}
\RouteGroup{Identidad estructural}
\RouteRow{Clase y multisección}{neutrino\_\allowbreak{}G4\quad /\quad neutrino\_\allowbreak{}G4}
\RouteRow{Colección y sector}{neutrino; leptón}
\RouteRow{Clasificación física}{neutrino}
\RouteRow{Generación, profundidad y color}{4; G4; \textemdash{}}
\RouteRow{Ruta primaria y órbita de inversión}{\(B_1\): no\qquad 1,\allowbreak{}6;\allowbreak{}9,\allowbreak{}4}
\RouteRow{Firma de clasificación}{34|\allowbreak{}0|\allowbreak{}-12|\allowbreak{}0|\allowbreak{}-12|\allowbreak{}+++-\kern0pt{}-\kern0pt{}-+++-\kern0pt{}-\kern0pt{}-}
\RouteGroup{Retornos, memoria y transporte}
\RouteRow{Retornos con fase}{clase: 27; órbita: 27; celda: 27}
\RouteRow{Retorno cinemático}{en 108: 54; extensión: 216}
\RouteRow{Giros y cambios de hoja}{71; 107}
\RouteRow{Acarreo y diferimiento}{deuda total: 0; final: 0; bloques diferidos: 0}
\RouteRow{Tríadas finales; persistencia de Pareto}{17; sí}
\RouteRow{\(K\longrightarrow K+9\)}{repeticiones: 9; cambios de memoria: 9}
\RouteGroup{Firma, fase y diferencias}
\RouteRow{\(n_A,\ s_C,\ \kappa_0,\ \nu_{120},\ \nu_{270}\)}{34, 0, 0, 0, -12}
\RouteRow{\(\tau_{12}\); firma histórica \(\mathbb Z^5\)}{+++-\kern0pt{}-\kern0pt{}-+++-\kern0pt{}-\kern0pt{}-; 34|\allowbreak{}0|\allowbreak{}0|\allowbreak{}0|\allowbreak{}-12}
\RouteRow{\(D_1,\ D_3,\ D_4,\ D_{27},\ D_{36}\)}{66, 60, 66, 48, 32}
\RouteRow{\(\Omega_{90/120},\ P_6\)}{80, 6}
\RouteGroup{Lectores y factores de acción}
\RouteRow{\(K_D\quad /\quad K_\Omega\)}{(80,\allowbreak{}66,\allowbreak{}3)\quad /\quad (80,\allowbreak{}54,\allowbreak{}6)}
\RouteRow{\(\kappa_D\)}{79.51870\allowbreak{}83196953\allowbreak{}45554341\allowbreak{}34}
\RouteRow{\(\kappa_\Omega\)}{79.60661\allowbreak{}01397948\allowbreak{}27586470\allowbreak{}91}
\RouteRow{\(R_{\mathrm{act},D}\)}{1.000000\allowbreak{}05512887\allowbreak{}17997050\allowbreak{}72}
\RouteRow{\(R_{\mathrm{act},\Omega}\)}{1.000000\allowbreak{}05518981\allowbreak{}25318591\allowbreak{}40}
\RouteGroup{Separación de prefijos}
\RouteRow{Área de ambigüedad; área \(\log_2\)}{18; 0.0}
\RouteRow{Primer bloque de separación}{órbita: 1; celda: 1; ruta: 1}
\RouteRow{Ambigüedad final}{1}
\RouteGroup{Secuencias completas}
\RouteRow{\(w_6\)}{243 657 24 243 657 570 264 90 570 243 657 24 243 657 570 264 90 570}
\RouteRow{Tríadas estables por bloque}{0 1 2 3 4 5 6 7 8 9 10 11 12 13 14 15 16 17}
\end{tabularx}\endgroup\par\vspace{5pt}\noindent{\fontsize{8.4}{10}\selectfont\hyperref[trace-304]{Procedencia y códigos originales}\hfill Ficha completa · 304/324}
\clearpage\phantomsection\label{nav:vi-ruta-305}
% VI-CATALOG-ROW rutas 305
\pdfbookmark[2]{Celda 77}{cell-77}
\pdfbookmark[3]{r9c5\_h-1v-1}{route-305}
\RouteTitle{Ruta 305}{r9c5\_\allowbreak{}h-1v-1}
\noindent Celda 77\quad (9, 5)\qquad Orientación \(h=-1,\ v=-1\)\hfill\hyperref[sec:vi-indice-rutas]{Índice}\par\vspace{4pt}
\begingroup\fontsize{10.2}{12.5}\selectfont\setlength{\tabcolsep}{5pt}\renewcommand{\arraystretch}{1.1}\rowcolors{1}{routepale}{white}\noindent\begin{tabularx}{\linewidth}{@{}>{\raggedright\arraybackslash}p{.345\linewidth}>{\raggedright\arraybackslash}X@{}}
\RouteGroup{Identidad estructural}
\RouteRow{Clase y multisección}{charged\_\allowbreak{}lepton\_\allowbreak{}G4\quad /\quad charged\_\allowbreak{}lepton\_\allowbreak{}G4}
\RouteRow{Colección y sector}{leptón cargado; leptón}
\RouteRow{Clasificación física}{leptón cargado}
\RouteRow{Generación, profundidad y color}{4; G4; \textemdash{}}
\RouteRow{Ruta primaria y órbita de inversión}{\(B_1\): sí\qquad 1,\allowbreak{}5;\allowbreak{}9,\allowbreak{}5}
\RouteRow{Firma de clasificación}{34|\allowbreak{}0|\allowbreak{}-8|\allowbreak{}-2|\allowbreak{}-8|\allowbreak{}0++-\kern0pt{}-\kern0pt{}-0++-\kern0pt{}-\kern0pt{}-}
\RouteGroup{Retornos, memoria y transporte}
\RouteRow{Retornos con fase}{clase: 27; órbita: 27; celda: 27}
\RouteRow{Retorno cinemático}{en 108: 54; extensión: 216}
\RouteRow{Giros y cambios de hoja}{71; 107}
\RouteRow{Acarreo y diferimiento}{deuda total: 1; final: 0; bloques diferidos: 1}
\RouteRow{Tríadas finales; persistencia de Pareto}{17; sí}
\RouteRow{\(K\longrightarrow K+9\)}{repeticiones: 9; cambios de memoria: 9}
\RouteGroup{Firma, fase y diferencias}
\RouteRow{\(n_A,\ s_C,\ \kappa_0,\ \nu_{120},\ \nu_{270}\)}{34, 0, -2, 0, -8}
\RouteRow{\(\tau_{12}\); firma histórica \(\mathbb Z^5\)}{+++0-\kern0pt{}-+++0-\kern0pt{}-; 34|\allowbreak{}0|\allowbreak{}-2|\allowbreak{}0|\allowbreak{}-8}
\RouteRow{\(D_1,\ D_3,\ D_4,\ D_{27},\ D_{36}\)}{66, 60, 66, 48, 32}
\RouteRow{\(\Omega_{90/120},\ P_6\)}{56, 5}
\RouteGroup{Lectores y factores de acción}
\RouteRow{\(K_D\quad /\quad K_\Omega\)}{(80,\allowbreak{}66,\allowbreak{}3)\quad /\quad (56,\allowbreak{}54,\allowbreak{}5)}
\RouteRow{\(\kappa_D\)}{79.51870\allowbreak{}83196953\allowbreak{}45554341\allowbreak{}34}
\RouteRow{\(\kappa_\Omega\)}{55.60649\allowbreak{}89433721\allowbreak{}35446416\allowbreak{}86}
\RouteRow{\(R_{\mathrm{act},D}\)}{1.000000\allowbreak{}05512887\allowbreak{}17997050\allowbreak{}72}
\RouteRow{\(R_{\mathrm{act},\Omega}\)}{1.000000\allowbreak{}03855097\allowbreak{}23541416\allowbreak{}45}
\RouteGroup{Separación de prefijos}
\RouteRow{Área de ambigüedad; área \(\log_2\)}{30; 8.0}
\RouteRow{Primer bloque de separación}{órbita: 5; celda: 5; ruta: 5}
\RouteRow{Ambigüedad final}{1}
\RouteGroup{Secuencias completas}
\RouteRow{\(w_6\)}{486 423 15 486 423 327 498 99 327 486 423 15 486 423 327 498 99 327}
\RouteRow{Tríadas estables por bloque}{0 1 2 3 4 5 6 7 8 9 10 10 12 13 14 15 16 17}
\end{tabularx}\endgroup\par\vspace{5pt}\noindent{\fontsize{8.4}{10}\selectfont\hyperref[trace-305]{Procedencia y códigos originales}\hfill Ficha completa · 305/324}
\clearpage\phantomsection\label{nav:vi-ruta-306}
% VI-CATALOG-ROW rutas 306
\pdfbookmark[3]{r9c5\_h-1v+1}{route-306}
\RouteTitle{Ruta 306}{r9c5\_\allowbreak{}h-1v+1}
\noindent Celda 77\quad (9, 5)\qquad Orientación \(h=-1,\ v=1\)\hfill\hyperref[sec:vi-indice-rutas]{Índice}\par\vspace{4pt}
\begingroup\fontsize{10.2}{12.5}\selectfont\setlength{\tabcolsep}{5pt}\renewcommand{\arraystretch}{1.1}\rowcolors{1}{routepale}{white}\noindent\begin{tabularx}{\linewidth}{@{}>{\raggedright\arraybackslash}p{.345\linewidth}>{\raggedright\arraybackslash}X@{}}
\RouteGroup{Identidad estructural}
\RouteRow{Clase y multisección}{charged\_\allowbreak{}lepton\_\allowbreak{}G4\quad /\quad charged\_\allowbreak{}lepton\_\allowbreak{}G4}
\RouteRow{Colección y sector}{leptón cargado; leptón}
\RouteRow{Clasificación física}{leptón cargado}
\RouteRow{Generación, profundidad y color}{4; G4; \textemdash{}}
\RouteRow{Ruta primaria y órbita de inversión}{\(B_1\): no\qquad 1,\allowbreak{}5;\allowbreak{}9,\allowbreak{}5}
\RouteRow{Firma de clasificación}{34|\allowbreak{}0|\allowbreak{}-8|\allowbreak{}-2|\allowbreak{}-8|\allowbreak{}0++-\kern0pt{}-\kern0pt{}-0++-\kern0pt{}-\kern0pt{}-}
\RouteGroup{Retornos, memoria y transporte}
\RouteRow{Retornos con fase}{clase: 27; órbita: 27; celda: 27}
\RouteRow{Retorno cinemático}{en 108: 54; extensión: 216}
\RouteRow{Giros y cambios de hoja}{71; 107}
\RouteRow{Acarreo y diferimiento}{deuda total: 1; final: 0; bloques diferidos: 1}
\RouteRow{Tríadas finales; persistencia de Pareto}{17; sí}
\RouteRow{\(K\longrightarrow K+9\)}{repeticiones: 9; cambios de memoria: 9}
\RouteGroup{Firma, fase y diferencias}
\RouteRow{\(n_A,\ s_C,\ \kappa_0,\ \nu_{120},\ \nu_{270}\)}{44, 2, 0, 0, -12}
\RouteRow{\(\tau_{12}\); firma histórica \(\mathbb Z^5\)}{+++-\kern0pt{}-\kern0pt{}-+++-\kern0pt{}-\kern0pt{}-; 44|\allowbreak{}2|\allowbreak{}0|\allowbreak{}0|\allowbreak{}-12}
\RouteRow{\(D_1,\ D_3,\ D_4,\ D_{27},\ D_{36}\)}{78, 24, 78, 24, 16}
\RouteRow{\(\Omega_{90/120},\ P_6\)}{80, 6}
\RouteGroup{Lectores y factores de acción}
\RouteRow{\(K_D\quad /\quad K_\Omega\)}{(40,\allowbreak{}78,\allowbreak{}27)\quad /\quad (80,\allowbreak{}54,\allowbreak{}6)}
\RouteRow{\(\kappa_D\)}{39.43380\allowbreak{}88030085\allowbreak{}51303671\allowbreak{}14}
\RouteRow{\(\kappa_\Omega\)}{79.60661\allowbreak{}01397948\allowbreak{}27586470\allowbreak{}91}
\RouteRow{\(R_{\mathrm{act},D}\)}{1.000000\allowbreak{}02733874\allowbreak{}08548943\allowbreak{}58}
\RouteRow{\(R_{\mathrm{act},\Omega}\)}{1.000000\allowbreak{}05518981\allowbreak{}25318591\allowbreak{}40}
\RouteGroup{Separación de prefijos}
\RouteRow{Área de ambigüedad; área \(\log_2\)}{18; 0.0}
\RouteRow{Primer bloque de separación}{órbita: 1; celda: 1; ruta: 1}
\RouteRow{Ambigüedad final}{1}
\RouteGroup{Secuencias completas}
\RouteRow{\(w_6\)}{504 585 507 504 585 672 510 666 672 504 585 507 504 585 672 510 666 672}
\RouteRow{Tríadas estables por bloque}{0 1 2 3 4 5 6 7 8 9 10 11 12 13 14 14 16 17}
\end{tabularx}\endgroup\par\vspace{5pt}\noindent{\fontsize{8.4}{10}\selectfont\hyperref[trace-306]{Procedencia y códigos originales}\hfill Ficha completa · 306/324}
\clearpage\phantomsection\label{nav:vi-ruta-307}
% VI-CATALOG-ROW rutas 307
\pdfbookmark[3]{r9c5\_h+1v-1}{route-307}
\RouteTitle{Ruta 307}{r9c5\_\allowbreak{}h+1v-1}
\noindent Celda 77\quad (9, 5)\qquad Orientación \(h=1,\ v=-1\)\hfill\hyperref[sec:vi-indice-rutas]{Índice}\par\vspace{4pt}
\begingroup\fontsize{10.2}{12.5}\selectfont\setlength{\tabcolsep}{5pt}\renewcommand{\arraystretch}{1.1}\rowcolors{1}{routepale}{white}\noindent\begin{tabularx}{\linewidth}{@{}>{\raggedright\arraybackslash}p{.345\linewidth}>{\raggedright\arraybackslash}X@{}}
\RouteGroup{Identidad estructural}
\RouteRow{Clase y multisección}{charged\_\allowbreak{}lepton\_\allowbreak{}G4\quad /\quad charged\_\allowbreak{}lepton\_\allowbreak{}G4}
\RouteRow{Colección y sector}{leptón cargado; leptón}
\RouteRow{Clasificación física}{leptón cargado}
\RouteRow{Generación, profundidad y color}{4; G4; \textemdash{}}
\RouteRow{Ruta primaria y órbita de inversión}{\(B_1\): no\qquad 1,\allowbreak{}5;\allowbreak{}9,\allowbreak{}5}
\RouteRow{Firma de clasificación}{34|\allowbreak{}0|\allowbreak{}-8|\allowbreak{}-2|\allowbreak{}-8|\allowbreak{}0++-\kern0pt{}-\kern0pt{}-0++-\kern0pt{}-\kern0pt{}-}
\RouteGroup{Retornos, memoria y transporte}
\RouteRow{Retornos con fase}{clase: 27; órbita: 27; celda: 27}
\RouteRow{Retorno cinemático}{en 108: 54; extensión: 216}
\RouteRow{Giros y cambios de hoja}{71; 107}
\RouteRow{Acarreo y diferimiento}{deuda total: 1; final: 0; bloques diferidos: 1}
\RouteRow{Tríadas finales; persistencia de Pareto}{17; sí}
\RouteRow{\(K\longrightarrow K+9\)}{repeticiones: 9; cambios de memoria: 9}
\RouteGroup{Firma, fase y diferencias}
\RouteRow{\(n_A,\ s_C,\ \kappa_0,\ \nu_{120},\ \nu_{270}\)}{44, -2, 0, 0, -12}
\RouteRow{\(\tau_{12}\); firma histórica \(\mathbb Z^5\)}{+++-\kern0pt{}-\kern0pt{}-+++-\kern0pt{}-\kern0pt{}-; 44|\allowbreak{}-2|\allowbreak{}0|\allowbreak{}0|\allowbreak{}-12}
\RouteRow{\(D_1,\ D_3,\ D_4,\ D_{27},\ D_{36}\)}{78, 24, 78, 24, 16}
\RouteRow{\(\Omega_{90/120},\ P_6\)}{80, 6}
\RouteGroup{Lectores y factores de acción}
\RouteRow{\(K_D\quad /\quad K_\Omega\)}{(40,\allowbreak{}78,\allowbreak{}27)\quad /\quad (80,\allowbreak{}54,\allowbreak{}6)}
\RouteRow{\(\kappa_D\)}{39.43380\allowbreak{}88030085\allowbreak{}51303671\allowbreak{}14}
\RouteRow{\(\kappa_\Omega\)}{79.60661\allowbreak{}01397948\allowbreak{}27586470\allowbreak{}91}
\RouteRow{\(R_{\mathrm{act},D}\)}{1.000000\allowbreak{}02733874\allowbreak{}08548943\allowbreak{}58}
\RouteRow{\(R_{\mathrm{act},\Omega}\)}{1.000000\allowbreak{}05518981\allowbreak{}25318591\allowbreak{}40}
\RouteGroup{Separación de prefijos}
\RouteRow{Área de ambigüedad; área \(\log_2\)}{18; 0.0}
\RouteRow{Primer bloque de separación}{órbita: 1; celda: 1; ruta: 1}
\RouteRow{Ambigüedad final}{1}
\RouteGroup{Secuencias completas}
\RouteRow{\(w_6\)}{510 666 672 510 666 507 504 585 507 510 666 672 510 666 507 504 585 507}
\RouteRow{Tríadas estables por bloque}{0 1 2 3 4 5 6 7 8 9 10 11 12 13 14 15 15 17}
\end{tabularx}\endgroup\par\vspace{5pt}\noindent{\fontsize{8.4}{10}\selectfont\hyperref[trace-307]{Procedencia y códigos originales}\hfill Ficha completa · 307/324}
\clearpage\phantomsection\label{nav:vi-ruta-308}
% VI-CATALOG-ROW rutas 308
\pdfbookmark[3]{r9c5\_h+1v+1}{route-308}
\RouteTitle{Ruta 308}{r9c5\_\allowbreak{}h+1v+1}
\noindent Celda 77\quad (9, 5)\qquad Orientación \(h=1,\ v=1\)\hfill\hyperref[sec:vi-indice-rutas]{Índice}\par\vspace{4pt}
\begingroup\fontsize{10.2}{12.5}\selectfont\setlength{\tabcolsep}{5pt}\renewcommand{\arraystretch}{1.1}\rowcolors{1}{routepale}{white}\noindent\begin{tabularx}{\linewidth}{@{}>{\raggedright\arraybackslash}p{.345\linewidth}>{\raggedright\arraybackslash}X@{}}
\RouteGroup{Identidad estructural}
\RouteRow{Clase y multisección}{charged\_\allowbreak{}lepton\_\allowbreak{}G4\quad /\quad charged\_\allowbreak{}lepton\_\allowbreak{}G4}
\RouteRow{Colección y sector}{leptón cargado; leptón}
\RouteRow{Clasificación física}{leptón cargado}
\RouteRow{Generación, profundidad y color}{4; G4; \textemdash{}}
\RouteRow{Ruta primaria y órbita de inversión}{\(B_1\): no\qquad 1,\allowbreak{}5;\allowbreak{}9,\allowbreak{}5}
\RouteRow{Firma de clasificación}{34|\allowbreak{}0|\allowbreak{}-8|\allowbreak{}-2|\allowbreak{}-8|\allowbreak{}0++-\kern0pt{}-\kern0pt{}-0++-\kern0pt{}-\kern0pt{}-}
\RouteGroup{Retornos, memoria y transporte}
\RouteRow{Retornos con fase}{clase: 27; órbita: 27; celda: 27}
\RouteRow{Retorno cinemático}{en 108: 54; extensión: 216}
\RouteRow{Giros y cambios de hoja}{71; 107}
\RouteRow{Acarreo y diferimiento}{deuda total: 0; final: 0; bloques diferidos: 0}
\RouteRow{Tríadas finales; persistencia de Pareto}{17; sí}
\RouteRow{\(K\longrightarrow K+9\)}{repeticiones: 9; cambios de memoria: 9}
\RouteGroup{Firma, fase y diferencias}
\RouteRow{\(n_A,\ s_C,\ \kappa_0,\ \nu_{120},\ \nu_{270}\)}{34, 0, 2, 0, -8}
\RouteRow{\(\tau_{12}\); firma histórica \(\mathbb Z^5\)}{0++-\kern0pt{}-\kern0pt{}-0++-\kern0pt{}-\kern0pt{}-; 34|\allowbreak{}0|\allowbreak{}2|\allowbreak{}0|\allowbreak{}-8}
\RouteRow{\(D_1,\ D_3,\ D_4,\ D_{27},\ D_{36}\)}{66, 60, 66, 48, 32}
\RouteRow{\(\Omega_{90/120},\ P_6\)}{56, 5}
\RouteGroup{Lectores y factores de acción}
\RouteRow{\(K_D\quad /\quad K_\Omega\)}{(80,\allowbreak{}66,\allowbreak{}3)\quad /\quad (56,\allowbreak{}54,\allowbreak{}5)}
\RouteRow{\(\kappa_D\)}{79.51870\allowbreak{}83196953\allowbreak{}45554341\allowbreak{}34}
\RouteRow{\(\kappa_\Omega\)}{55.60649\allowbreak{}89433721\allowbreak{}35446416\allowbreak{}86}
\RouteRow{\(R_{\mathrm{act},D}\)}{1.000000\allowbreak{}05512887\allowbreak{}17997050\allowbreak{}72}
\RouteRow{\(R_{\mathrm{act},\Omega}\)}{1.000000\allowbreak{}03855097\allowbreak{}23541416\allowbreak{}45}
\RouteGroup{Separación de prefijos}
\RouteRow{Área de ambigüedad; área \(\log_2\)}{18; 0.0}
\RouteRow{Primer bloque de separación}{órbita: 1; celda: 1; ruta: 1}
\RouteRow{Ambigüedad final}{1}
\RouteGroup{Secuencias completas}
\RouteRow{\(w_6\)}{498 99 327 498 99 15 486 423 15 498 99 327 498 99 15 486 423 15}
\RouteRow{Tríadas estables por bloque}{0 1 2 3 4 5 6 7 8 9 10 11 12 13 14 15 16 17}
\end{tabularx}\endgroup\par\vspace{5pt}\noindent{\fontsize{8.4}{10}\selectfont\hyperref[trace-308]{Procedencia y códigos originales}\hfill Ficha completa · 308/324}
\clearpage\phantomsection\label{nav:vi-ruta-309}
% VI-CATALOG-ROW rutas 309
\pdfbookmark[2]{Celda 78}{cell-78}
\pdfbookmark[3]{r9c6\_h-1v-1}{route-309}
\RouteTitle{Ruta 309}{r9c6\_\allowbreak{}h-1v-1}
\noindent Celda 78\quad (9, 6)\qquad Orientación \(h=-1,\ v=-1\)\hfill\hyperref[sec:vi-indice-rutas]{Índice}\par\vspace{4pt}
\begingroup\fontsize{10.2}{12.5}\selectfont\setlength{\tabcolsep}{5pt}\renewcommand{\arraystretch}{1.1}\rowcolors{1}{routepale}{white}\noindent\begin{tabularx}{\linewidth}{@{}>{\raggedright\arraybackslash}p{.345\linewidth}>{\raggedright\arraybackslash}X@{}}
\RouteGroup{Identidad estructural}
\RouteRow{Clase y multisección}{neutrino\_\allowbreak{}G4\quad /\quad neutrino\_\allowbreak{}G4}
\RouteRow{Colección y sector}{neutrino; leptón}
\RouteRow{Clasificación física}{neutrino}
\RouteRow{Generación, profundidad y color}{4; G4; \textemdash{}}
\RouteRow{Ruta primaria y órbita de inversión}{\(B_1\): no\qquad 1,\allowbreak{}4;\allowbreak{}9,\allowbreak{}6}
\RouteRow{Firma de clasificación}{30|\allowbreak{}6|\allowbreak{}0|\allowbreak{}0|\allowbreak{}-8|\allowbreak{}+++-0-+++-0-}
\RouteGroup{Retornos, memoria y transporte}
\RouteRow{Retornos con fase}{clase: 27; órbita: 27; celda: 27}
\RouteRow{Retorno cinemático}{en 108: 54; extensión: 216}
\RouteRow{Giros y cambios de hoja}{71; 107}
\RouteRow{Acarreo y diferimiento}{deuda total: 1; final: 0; bloques diferidos: 1}
\RouteRow{Tríadas finales; persistencia de Pareto}{17; sí}
\RouteRow{\(K\longrightarrow K+9\)}{repeticiones: 9; cambios de memoria: 9}
\RouteGroup{Firma, fase y diferencias}
\RouteRow{\(n_A,\ s_C,\ \kappa_0,\ \nu_{120},\ \nu_{270}\)}{30, -6, 2, 0, -8}
\RouteRow{\(\tau_{12}\); firma histórica \(\mathbb Z^5\)}{+0+-\kern0pt{}-\kern0pt{}-+0+-\kern0pt{}-\kern0pt{}-; 30|\allowbreak{}-6|\allowbreak{}2|\allowbreak{}0|\allowbreak{}-8}
\RouteRow{\(D_1,\ D_3,\ D_4,\ D_{27},\ D_{36}\)}{54, 60, 72, 24, 16}
\RouteRow{\(\Omega_{90/120},\ P_6\)}{48, 5}
\RouteGroup{Lectores y factores de acción}
\RouteRow{\(K_D\quad /\quad K_\Omega\)}{(40,\allowbreak{}54,\allowbreak{}6)\quad /\quad (48,\allowbreak{}54,\allowbreak{}5)}
\RouteRow{\(\kappa_D\)}{39.60661\allowbreak{}01397948\allowbreak{}27586470\allowbreak{}91}
\RouteRow{\(\kappa_\Omega\)}{47.60649\allowbreak{}89433721\allowbreak{}35446416\allowbreak{}86}
\RouteRow{\(R_{\mathrm{act},D}\)}{1.000000\allowbreak{}02745854\allowbreak{}08735595\allowbreak{}18}
\RouteRow{\(R_{\mathrm{act},\Omega}\)}{1.000000\allowbreak{}03300471\allowbreak{}80532430\allowbreak{}84}
\RouteGroup{Separación de prefijos}
\RouteRow{Área de ambigüedad; área \(\log_2\)}{18; 0.0}
\RouteRow{Primer bloque de separación}{órbita: 1; celda: 1; ruta: 1}
\RouteRow{Ambigüedad final}{1}
\RouteGroup{Secuencias completas}
\RouteRow{\(w_6\)}{15 486 423 15 486 657 24 243 657 15 486 423 15 486 657 24 243 657}
\RouteRow{Tríadas estables por bloque}{0 1 2 3 4 5 6 7 8 9 10 11 12 13 14 15 15 17}
\end{tabularx}\endgroup\par\vspace{5pt}\noindent{\fontsize{8.4}{10}\selectfont\hyperref[trace-309]{Procedencia y códigos originales}\hfill Ficha completa · 309/324}
\clearpage\phantomsection\label{nav:vi-ruta-310}
% VI-CATALOG-ROW rutas 310
\pdfbookmark[3]{r9c6\_h-1v+1}{route-310}
\RouteTitle{Ruta 310}{r9c6\_\allowbreak{}h-1v+1}
\noindent Celda 78\quad (9, 6)\qquad Orientación \(h=-1,\ v=1\)\hfill\hyperref[sec:vi-indice-rutas]{Índice}\par\vspace{4pt}
\begingroup\fontsize{10.2}{12.5}\selectfont\setlength{\tabcolsep}{5pt}\renewcommand{\arraystretch}{1.1}\rowcolors{1}{routepale}{white}\noindent\begin{tabularx}{\linewidth}{@{}>{\raggedright\arraybackslash}p{.345\linewidth}>{\raggedright\arraybackslash}X@{}}
\RouteGroup{Identidad estructural}
\RouteRow{Clase y multisección}{neutrino\_\allowbreak{}G4\quad /\quad neutrino\_\allowbreak{}G4}
\RouteRow{Colección y sector}{neutrino; leptón}
\RouteRow{Clasificación física}{neutrino}
\RouteRow{Generación, profundidad y color}{4; G4; \textemdash{}}
\RouteRow{Ruta primaria y órbita de inversión}{\(B_1\): no\qquad 1,\allowbreak{}4;\allowbreak{}9,\allowbreak{}6}
\RouteRow{Firma de clasificación}{30|\allowbreak{}6|\allowbreak{}0|\allowbreak{}0|\allowbreak{}-8|\allowbreak{}+++-0-+++-0-}
\RouteGroup{Retornos, memoria y transporte}
\RouteRow{Retornos con fase}{clase: 27; órbita: 27; celda: 27}
\RouteRow{Retorno cinemático}{en 108: 54; extensión: 216}
\RouteRow{Giros y cambios de hoja}{71; 107}
\RouteRow{Acarreo y diferimiento}{deuda total: 0; final: 0; bloques diferidos: 0}
\RouteRow{Tríadas finales; persistencia de Pareto}{17; sí}
\RouteRow{\(K\longrightarrow K+9\)}{repeticiones: 9; cambios de memoria: 9}
\RouteGroup{Firma, fase y diferencias}
\RouteRow{\(n_A,\ s_C,\ \kappa_0,\ \nu_{120},\ \nu_{270}\)}{14, 6, 0, 0, -12}
\RouteRow{\(\tau_{12}\); firma histórica \(\mathbb Z^5\)}{+++-\kern0pt{}-\kern0pt{}-+++-\kern0pt{}-\kern0pt{}-; 14|\allowbreak{}6|\allowbreak{}0|\allowbreak{}0|\allowbreak{}-12}
\RouteRow{\(D_1,\ D_3,\ D_4,\ D_{27},\ D_{36}\)}{24, 24, 24, 0, 0}
\RouteRow{\(\Omega_{90/120},\ P_6\)}{80, 6}
\RouteGroup{Lectores y factores de acción}
\RouteRow{\(K_D\quad /\quad K_\Omega\)}{(0,\allowbreak{}24,\allowbreak{}0)\quad /\quad (80,\allowbreak{}54,\allowbreak{}6)}
\RouteRow{\(\kappa_D\)}{-0.175136\allowbreak{}46166281\allowbreak{}12239348\allowbreak{}425}
\RouteRow{\(\kappa_\Omega\)}{79.60661\allowbreak{}01397948\allowbreak{}27586470\allowbreak{}91}
\RouteRow{\(R_{\mathrm{act},D}\)}{0.999999\allowbreak{}99987858\allowbreak{}10851337\allowbreak{}881}
\RouteRow{\(R_{\mathrm{act},\Omega}\)}{1.000000\allowbreak{}05518981\allowbreak{}25318591\allowbreak{}40}
\RouteGroup{Separación de prefijos}
\RouteRow{Área de ambigüedad; área \(\log_2\)}{36; 18.0}
\RouteRow{Primer bloque de separación}{órbita: 1; celda: 1; ruta: \textemdash{}}
\RouteRow{Ambigüedad final}{2}
\RouteGroup{Secuencias completas}
\RouteRow{\(w_6\)}{3 0 81 3 0 81 3 0 81 3 0 81 3 0 81 3 0 81}
\RouteRow{Tríadas estables por bloque}{0 1 2 3 4 5 6 7 8 9 10 11 12 13 14 15 16 17}
\end{tabularx}\endgroup\par\vspace{5pt}\noindent{\fontsize{8.4}{10}\selectfont\hyperref[trace-310]{Procedencia y códigos originales}\hfill Ficha completa · 310/324}
\clearpage\phantomsection\label{nav:vi-ruta-311}
% VI-CATALOG-ROW rutas 311
\pdfbookmark[3]{r9c6\_h+1v-1}{route-311}
\RouteTitle{Ruta 311}{r9c6\_\allowbreak{}h+1v-1}
\noindent Celda 78\quad (9, 6)\qquad Orientación \(h=1,\ v=-1\)\hfill\hyperref[sec:vi-indice-rutas]{Índice}\par\vspace{4pt}
\begingroup\fontsize{10.2}{12.5}\selectfont\setlength{\tabcolsep}{5pt}\renewcommand{\arraystretch}{1.1}\rowcolors{1}{routepale}{white}\noindent\begin{tabularx}{\linewidth}{@{}>{\raggedright\arraybackslash}p{.345\linewidth}>{\raggedright\arraybackslash}X@{}}
\RouteGroup{Identidad estructural}
\RouteRow{Clase y multisección}{neutrino\_\allowbreak{}G4\quad /\quad neutrino\_\allowbreak{}G4}
\RouteRow{Colección y sector}{neutrino; leptón}
\RouteRow{Clasificación física}{neutrino}
\RouteRow{Generación, profundidad y color}{4; G4; \textemdash{}}
\RouteRow{Ruta primaria y órbita de inversión}{\(B_1\): sí\qquad 1,\allowbreak{}4;\allowbreak{}9,\allowbreak{}6}
\RouteRow{Firma de clasificación}{30|\allowbreak{}6|\allowbreak{}0|\allowbreak{}0|\allowbreak{}-8|\allowbreak{}+++-0-+++-0-}
\RouteGroup{Retornos, memoria y transporte}
\RouteRow{Retornos con fase}{clase: 27; órbita: 27; celda: 27}
\RouteRow{Retorno cinemático}{en 108: 54; extensión: 216}
\RouteRow{Giros y cambios de hoja}{71; 107}
\RouteRow{Acarreo y diferimiento}{deuda total: 0; final: 0; bloques diferidos: 0}
\RouteRow{Tríadas finales; persistencia de Pareto}{17; sí}
\RouteRow{\(K\longrightarrow K+9\)}{repeticiones: 9; cambios de memoria: 9}
\RouteGroup{Firma, fase y diferencias}
\RouteRow{\(n_A,\ s_C,\ \kappa_0,\ \nu_{120},\ \nu_{270}\)}{14, -6, 0, 0, -12}
\RouteRow{\(\tau_{12}\); firma histórica \(\mathbb Z^5\)}{+++-\kern0pt{}-\kern0pt{}-+++-\kern0pt{}-\kern0pt{}-; 14|\allowbreak{}-6|\allowbreak{}0|\allowbreak{}0|\allowbreak{}-12}
\RouteRow{\(D_1,\ D_3,\ D_4,\ D_{27},\ D_{36}\)}{24, 24, 24, 0, 0}
\RouteRow{\(\Omega_{90/120},\ P_6\)}{80, 6}
\RouteGroup{Lectores y factores de acción}
\RouteRow{\(K_D\quad /\quad K_\Omega\)}{(0,\allowbreak{}24,\allowbreak{}0)\quad /\quad (80,\allowbreak{}54,\allowbreak{}6)}
\RouteRow{\(\kappa_D\)}{-0.175136\allowbreak{}46166281\allowbreak{}12239348\allowbreak{}425}
\RouteRow{\(\kappa_\Omega\)}{79.60661\allowbreak{}01397948\allowbreak{}27586470\allowbreak{}91}
\RouteRow{\(R_{\mathrm{act},D}\)}{0.999999\allowbreak{}99987858\allowbreak{}10851337\allowbreak{}881}
\RouteRow{\(R_{\mathrm{act},\Omega}\)}{1.000000\allowbreak{}05518981\allowbreak{}25318591\allowbreak{}40}
\RouteGroup{Separación de prefijos}
\RouteRow{Área de ambigüedad; área \(\log_2\)}{36; 18.0}
\RouteRow{Primer bloque de separación}{órbita: 1; celda: 1; ruta: \textemdash{}}
\RouteRow{Ambigüedad final}{2}
\RouteGroup{Secuencias completas}
\RouteRow{\(w_6\)}{3 0 81 3 0 81 3 0 81 3 0 81 3 0 81 3 0 81}
\RouteRow{Tríadas estables por bloque}{0 1 2 3 4 5 6 7 8 9 10 11 12 13 14 15 16 17}
\end{tabularx}\endgroup\par\vspace{5pt}\noindent{\fontsize{8.4}{10}\selectfont\hyperref[trace-311]{Procedencia y códigos originales}\hfill Ficha completa · 311/324}
\clearpage\phantomsection\label{nav:vi-ruta-312}
% VI-CATALOG-ROW rutas 312
\pdfbookmark[3]{r9c6\_h+1v+1}{route-312}
\RouteTitle{Ruta 312}{r9c6\_\allowbreak{}h+1v+1}
\noindent Celda 78\quad (9, 6)\qquad Orientación \(h=1,\ v=1\)\hfill\hyperref[sec:vi-indice-rutas]{Índice}\par\vspace{4pt}
\begingroup\fontsize{10.2}{12.5}\selectfont\setlength{\tabcolsep}{5pt}\renewcommand{\arraystretch}{1.1}\rowcolors{1}{routepale}{white}\noindent\begin{tabularx}{\linewidth}{@{}>{\raggedright\arraybackslash}p{.345\linewidth}>{\raggedright\arraybackslash}X@{}}
\RouteGroup{Identidad estructural}
\RouteRow{Clase y multisección}{neutrino\_\allowbreak{}G4\quad /\quad neutrino\_\allowbreak{}G4}
\RouteRow{Colección y sector}{neutrino; leptón}
\RouteRow{Clasificación física}{neutrino}
\RouteRow{Generación, profundidad y color}{4; G4; \textemdash{}}
\RouteRow{Ruta primaria y órbita de inversión}{\(B_1\): no\qquad 1,\allowbreak{}4;\allowbreak{}9,\allowbreak{}6}
\RouteRow{Firma de clasificación}{30|\allowbreak{}6|\allowbreak{}0|\allowbreak{}0|\allowbreak{}-8|\allowbreak{}+++-0-+++-0-}
\RouteGroup{Retornos, memoria y transporte}
\RouteRow{Retornos con fase}{clase: 27; órbita: 27; celda: 27}
\RouteRow{Retorno cinemático}{en 108: 54; extensión: 216}
\RouteRow{Giros y cambios de hoja}{71; 107}
\RouteRow{Acarreo y diferimiento}{deuda total: 2; final: 1; bloques diferidos: 2}
\RouteRow{Tríadas finales; persistencia de Pareto}{16; sí}
\RouteRow{\(K\longrightarrow K+9\)}{repeticiones: 9; cambios de memoria: 9}
\RouteGroup{Firma, fase y diferencias}
\RouteRow{\(n_A,\ s_C,\ \kappa_0,\ \nu_{120},\ \nu_{270}\)}{30, 6, -2, 0, -8}
\RouteRow{\(\tau_{12}\); firma histórica \(\mathbb Z^5\)}{+++-0-+++-0-; 30|\allowbreak{}6|\allowbreak{}-2|\allowbreak{}0|\allowbreak{}-8}
\RouteRow{\(D_1,\ D_3,\ D_4,\ D_{27},\ D_{36}\)}{54, 60, 72, 24, 16}
\RouteRow{\(\Omega_{90/120},\ P_6\)}{48, 5}
\RouteGroup{Lectores y factores de acción}
\RouteRow{\(K_D\quad /\quad K_\Omega\)}{(40,\allowbreak{}54,\allowbreak{}6)\quad /\quad (48,\allowbreak{}54,\allowbreak{}5)}
\RouteRow{\(\kappa_D\)}{39.60661\allowbreak{}01397948\allowbreak{}27586470\allowbreak{}91}
\RouteRow{\(\kappa_\Omega\)}{47.60649\allowbreak{}89433721\allowbreak{}35446416\allowbreak{}86}
\RouteRow{\(R_{\mathrm{act},D}\)}{1.000000\allowbreak{}02745854\allowbreak{}08735595\allowbreak{}18}
\RouteRow{\(R_{\mathrm{act},\Omega}\)}{1.000000\allowbreak{}03300471\allowbreak{}80532430\allowbreak{}84}
\RouteGroup{Separación de prefijos}
\RouteRow{Área de ambigüedad; área \(\log_2\)}{26; 6.339850\allowbreak{}00288462\allowbreak{}4}
\RouteRow{Primer bloque de separación}{órbita: 5; celda: 5; ruta: 5}
\RouteRow{Ambigüedad final}{1}
\RouteGroup{Secuencias completas}
\RouteRow{\(w_6\)}{24 243 657 24 243 423 15 486 423 24 243 657 24 243 423 15 486 423}
\RouteRow{Tríadas estables por bloque}{0 1 2 3 4 5 6 7 8 9 10 11 12 13 13 15 16 16}
\end{tabularx}\endgroup\par\vspace{5pt}\noindent{\fontsize{8.4}{10}\selectfont\hyperref[trace-312]{Procedencia y códigos originales}\hfill Ficha completa · 312/324}
\clearpage\phantomsection\label{nav:vi-ruta-313}
% VI-CATALOG-ROW rutas 313
\pdfbookmark[2]{Celda 79}{cell-79}
\pdfbookmark[3]{r9c7\_h-1v-1}{route-313}
\RouteTitle{Ruta 313}{r9c7\_\allowbreak{}h-1v-1}
\noindent Celda 79\quad (9, 7)\qquad Orientación \(h=-1,\ v=-1\)\hfill\hyperref[sec:vi-indice-rutas]{Índice}\par\vspace{4pt}
\begingroup\fontsize{10.2}{12.5}\selectfont\setlength{\tabcolsep}{5pt}\renewcommand{\arraystretch}{1.1}\rowcolors{1}{routepale}{white}\noindent\begin{tabularx}{\linewidth}{@{}>{\raggedright\arraybackslash}p{.345\linewidth}>{\raggedright\arraybackslash}X@{}}
\RouteGroup{Identidad estructural}
\RouteRow{Clase y multisección}{charged\_\allowbreak{}lepton\_\allowbreak{}G4\quad /\quad charged\_\allowbreak{}lepton\_\allowbreak{}G4}
\RouteRow{Colección y sector}{leptón cargado; leptón}
\RouteRow{Clasificación física}{leptón cargado}
\RouteRow{Generación, profundidad y color}{4; G4; \textemdash{}}
\RouteRow{Ruta primaria y órbita de inversión}{\(B_1\): sí\qquad 1,\allowbreak{}3;\allowbreak{}9,\allowbreak{}7}
\RouteRow{Firma de clasificación}{28|\allowbreak{}6|\allowbreak{}-2|\allowbreak{}0|\allowbreak{}-12|\allowbreak{}+++-\kern0pt{}-\kern0pt{}-+++-\kern0pt{}-\kern0pt{}-}
\RouteGroup{Retornos, memoria y transporte}
\RouteRow{Retornos con fase}{clase: 18; órbita: 27; celda: 27}
\RouteRow{Retorno cinemático}{en 108: 54; extensión: 216}
\RouteRow{Giros y cambios de hoja}{71; 107}
\RouteRow{Acarreo y diferimiento}{deuda total: 2; final: 1; bloques diferidos: 1}
\RouteRow{Tríadas finales; persistencia de Pareto}{16; no}
\RouteRow{\(K\longrightarrow K+9\)}{repeticiones: 9; cambios de memoria: 9}
\RouteGroup{Firma, fase y diferencias}
\RouteRow{\(n_A,\ s_C,\ \kappa_0,\ \nu_{120},\ \nu_{270}\)}{28, -6, 0, -4, -12}
\RouteRow{\(\tau_{12}\); firma histórica \(\mathbb Z^5\)}{+++-\kern0pt{}-\kern0pt{}-+++-\kern0pt{}-\kern0pt{}-; 28|\allowbreak{}-6|\allowbreak{}0|\allowbreak{}-4|\allowbreak{}-12}
\RouteRow{\(D_1,\ D_3,\ D_4,\ D_{27},\ D_{36}\)}{66, 60, 66, 48, 32}
\RouteRow{\(\Omega_{90/120},\ P_6\)}{80, 6}
\RouteGroup{Lectores y factores de acción}
\RouteRow{\(K_D\quad /\quad K_\Omega\)}{(80,\allowbreak{}66,\allowbreak{}3)\quad /\quad (80,\allowbreak{}54,\allowbreak{}6)}
\RouteRow{\(\kappa_D\)}{79.51870\allowbreak{}83196953\allowbreak{}45554341\allowbreak{}34}
\RouteRow{\(\kappa_\Omega\)}{79.60661\allowbreak{}01397948\allowbreak{}27586470\allowbreak{}91}
\RouteRow{\(R_{\mathrm{act},D}\)}{1.000000\allowbreak{}05512887\allowbreak{}17997050\allowbreak{}72}
\RouteRow{\(R_{\mathrm{act},\Omega}\)}{1.000000\allowbreak{}05518981\allowbreak{}25318591\allowbreak{}40}
\RouteGroup{Separación de prefijos}
\RouteRow{Área de ambigüedad; área \(\log_2\)}{54; 28.52932\allowbreak{}50129808}
\RouteRow{Primer bloque de separación}{órbita: \textemdash{}; celda: \textemdash{}; ruta: \textemdash{}}
\RouteRow{Ambigüedad final}{3}
\RouteGroup{Secuencias completas}
\RouteRow{\(w_6\)}{264 90 570 264 90 24 243 657 24 264 90 570 264 90 24 243 657 24}
\RouteRow{Tríadas estables por bloque}{0 1 2 3 4 5 6 7 8 9 10 11 12 13 14 15 15 16}
\end{tabularx}\endgroup\par\vspace{5pt}\noindent{\fontsize{8.4}{10}\selectfont\hyperref[trace-313]{Procedencia y códigos originales}\hfill Ficha completa · 313/324}
\clearpage\phantomsection\label{nav:vi-ruta-314}
% VI-CATALOG-ROW rutas 314
\pdfbookmark[3]{r9c7\_h-1v+1}{route-314}
\RouteTitle{Ruta 314}{r9c7\_\allowbreak{}h-1v+1}
\noindent Celda 79\quad (9, 7)\qquad Orientación \(h=-1,\ v=1\)\hfill\hyperref[sec:vi-indice-rutas]{Índice}\par\vspace{4pt}
\begingroup\fontsize{10.2}{12.5}\selectfont\setlength{\tabcolsep}{5pt}\renewcommand{\arraystretch}{1.1}\rowcolors{1}{routepale}{white}\noindent\begin{tabularx}{\linewidth}{@{}>{\raggedright\arraybackslash}p{.345\linewidth}>{\raggedright\arraybackslash}X@{}}
\RouteGroup{Identidad estructural}
\RouteRow{Clase y multisección}{charged\_\allowbreak{}lepton\_\allowbreak{}G4\quad /\quad charged\_\allowbreak{}lepton\_\allowbreak{}G4}
\RouteRow{Colección y sector}{leptón cargado; leptón}
\RouteRow{Clasificación física}{leptón cargado}
\RouteRow{Generación, profundidad y color}{4; G4; \textemdash{}}
\RouteRow{Ruta primaria y órbita de inversión}{\(B_1\): no\qquad 1,\allowbreak{}3;\allowbreak{}9,\allowbreak{}7}
\RouteRow{Firma de clasificación}{28|\allowbreak{}6|\allowbreak{}-2|\allowbreak{}0|\allowbreak{}-12|\allowbreak{}+++-\kern0pt{}-\kern0pt{}-+++-\kern0pt{}-\kern0pt{}-}
\RouteGroup{Retornos, memoria y transporte}
\RouteRow{Retornos con fase}{clase: 27; órbita: 27; celda: 27}
\RouteRow{Retorno cinemático}{en 108: 54; extensión: 216}
\RouteRow{Giros y cambios de hoja}{71; 107}
\RouteRow{Acarreo y diferimiento}{deuda total: 2; final: 0; bloques diferidos: 1}
\RouteRow{Tríadas finales; persistencia de Pareto}{17; no}
\RouteRow{\(K\longrightarrow K+9\)}{repeticiones: 9; cambios de memoria: 9}
\RouteGroup{Firma, fase y diferencias}
\RouteRow{\(n_A,\ s_C,\ \kappa_0,\ \nu_{120},\ \nu_{270}\)}{50, 4, -2, 2, -8}
\RouteRow{\(\tau_{12}\); firma histórica \(\mathbb Z^5\)}{+++0-\kern0pt{}-+++0-\kern0pt{}-; 50|\allowbreak{}4|\allowbreak{}-2|\allowbreak{}2|\allowbreak{}-8}
\RouteRow{\(D_1,\ D_3,\ D_4,\ D_{27},\ D_{36}\)}{84, 24, 84, 24, 16}
\RouteRow{\(\Omega_{90/120},\ P_6\)}{56, 5}
\RouteGroup{Lectores y factores de acción}
\RouteRow{\(K_D\quad /\quad K_\Omega\)}{(40,\allowbreak{}84,\allowbreak{}30)\quad /\quad (56,\allowbreak{}54,\allowbreak{}5)}
\RouteRow{\(\kappa_D\)}{39.39035\allowbreak{}82768609\allowbreak{}24917849\allowbreak{}59}
\RouteRow{\(\kappa_\Omega\)}{55.60649\allowbreak{}89433721\allowbreak{}35446416\allowbreak{}86}
\RouteRow{\(R_{\mathrm{act},D}\)}{1.000000\allowbreak{}02730861\allowbreak{}73967087\allowbreak{}05}
\RouteRow{\(R_{\mathrm{act},\Omega}\)}{1.000000\allowbreak{}03855097\allowbreak{}23541416\allowbreak{}45}
\RouteGroup{Separación de prefijos}
\RouteRow{Área de ambigüedad; área \(\log_2\)}{54; 28.52932\allowbreak{}50129808}
\RouteRow{Primer bloque de separación}{órbita: \textemdash{}; celda: \textemdash{}; ruta: \textemdash{}}
\RouteRow{Ambigüedad final}{3}
\RouteGroup{Secuencias completas}
\RouteRow{\(w_6\)}{258 414 420 258 414 255 252 333 255 258 414 420 258 414 255 252 333 255}
\RouteRow{Tríadas estables por bloque}{0 1 2 3 4 5 6 7 8 9 10 11 11 12 14 15 16 17}
\end{tabularx}\endgroup\par\vspace{5pt}\noindent{\fontsize{8.4}{10}\selectfont\hyperref[trace-314]{Procedencia y códigos originales}\hfill Ficha completa · 314/324}
\clearpage\phantomsection\label{nav:vi-ruta-315}
% VI-CATALOG-ROW rutas 315
\pdfbookmark[3]{r9c7\_h+1v-1}{route-315}
\RouteTitle{Ruta 315}{r9c7\_\allowbreak{}h+1v-1}
\noindent Celda 79\quad (9, 7)\qquad Orientación \(h=1,\ v=-1\)\hfill\hyperref[sec:vi-indice-rutas]{Índice}\par\vspace{4pt}
\begingroup\fontsize{10.2}{12.5}\selectfont\setlength{\tabcolsep}{5pt}\renewcommand{\arraystretch}{1.1}\rowcolors{1}{routepale}{white}\noindent\begin{tabularx}{\linewidth}{@{}>{\raggedright\arraybackslash}p{.345\linewidth}>{\raggedright\arraybackslash}X@{}}
\RouteGroup{Identidad estructural}
\RouteRow{Clase y multisección}{charged\_\allowbreak{}lepton\_\allowbreak{}G4\quad /\quad charged\_\allowbreak{}lepton\_\allowbreak{}G4}
\RouteRow{Colección y sector}{leptón cargado; leptón}
\RouteRow{Clasificación física}{leptón cargado}
\RouteRow{Generación, profundidad y color}{4; G4; \textemdash{}}
\RouteRow{Ruta primaria y órbita de inversión}{\(B_1\): no\qquad 1,\allowbreak{}3;\allowbreak{}9,\allowbreak{}7}
\RouteRow{Firma de clasificación}{28|\allowbreak{}6|\allowbreak{}-2|\allowbreak{}0|\allowbreak{}-12|\allowbreak{}+++-\kern0pt{}-\kern0pt{}-+++-\kern0pt{}-\kern0pt{}-}
\RouteGroup{Retornos, memoria y transporte}
\RouteRow{Retornos con fase}{clase: 27; órbita: 27; celda: 27}
\RouteRow{Retorno cinemático}{en 108: 54; extensión: 216}
\RouteRow{Giros y cambios de hoja}{71; 107}
\RouteRow{Acarreo y diferimiento}{deuda total: 3; final: 1; bloques diferidos: 3}
\RouteRow{Tríadas finales; persistencia de Pareto}{16; no}
\RouteRow{\(K\longrightarrow K+9\)}{repeticiones: 9; cambios de memoria: 9}
\RouteGroup{Firma, fase y diferencias}
\RouteRow{\(n_A,\ s_C,\ \kappa_0,\ \nu_{120},\ \nu_{270}\)}{50, -4, 2, -2, -8}
\RouteRow{\(\tau_{12}\); firma histórica \(\mathbb Z^5\)}{0++-\kern0pt{}-\kern0pt{}-0++-\kern0pt{}-\kern0pt{}-; 50|\allowbreak{}-4|\allowbreak{}2|\allowbreak{}-2|\allowbreak{}-8}
\RouteRow{\(D_1,\ D_3,\ D_4,\ D_{27},\ D_{36}\)}{84, 24, 84, 24, 16}
\RouteRow{\(\Omega_{90/120},\ P_6\)}{56, 5}
\RouteGroup{Lectores y factores de acción}
\RouteRow{\(K_D\quad /\quad K_\Omega\)}{(40,\allowbreak{}84,\allowbreak{}30)\quad /\quad (56,\allowbreak{}54,\allowbreak{}5)}
\RouteRow{\(\kappa_D\)}{39.39035\allowbreak{}82768609\allowbreak{}24917849\allowbreak{}59}
\RouteRow{\(\kappa_\Omega\)}{55.60649\allowbreak{}89433721\allowbreak{}35446416\allowbreak{}86}
\RouteRow{\(R_{\mathrm{act},D}\)}{1.000000\allowbreak{}02730861\allowbreak{}73967087\allowbreak{}05}
\RouteRow{\(R_{\mathrm{act},\Omega}\)}{1.000000\allowbreak{}03855097\allowbreak{}23541416\allowbreak{}45}
\RouteGroup{Separación de prefijos}
\RouteRow{Área de ambigüedad; área \(\log_2\)}{54; 28.52932\allowbreak{}50129808}
\RouteRow{Primer bloque de separación}{órbita: \textemdash{}; celda: \textemdash{}; ruta: \textemdash{}}
\RouteRow{Ambigüedad final}{3}
\RouteGroup{Secuencias completas}
\RouteRow{\(w_6\)}{252 333 255 252 333 420 258 414 420 252 333 255 252 333 420 258 414 420}
\RouteRow{Tríadas estables por bloque}{0 1 2 3 4 5 6 7 8 9 10 11 12 12 14 14 16 16}
\end{tabularx}\endgroup\par\vspace{5pt}\noindent{\fontsize{8.4}{10}\selectfont\hyperref[trace-315]{Procedencia y códigos originales}\hfill Ficha completa · 315/324}
\clearpage\phantomsection\label{nav:vi-ruta-316}
% VI-CATALOG-ROW rutas 316
\pdfbookmark[3]{r9c7\_h+1v+1}{route-316}
\RouteTitle{Ruta 316}{r9c7\_\allowbreak{}h+1v+1}
\noindent Celda 79\quad (9, 7)\qquad Orientación \(h=1,\ v=1\)\hfill\hyperref[sec:vi-indice-rutas]{Índice}\par\vspace{4pt}
\begingroup\fontsize{10.2}{12.5}\selectfont\setlength{\tabcolsep}{5pt}\renewcommand{\arraystretch}{1.1}\rowcolors{1}{routepale}{white}\noindent\begin{tabularx}{\linewidth}{@{}>{\raggedright\arraybackslash}p{.345\linewidth}>{\raggedright\arraybackslash}X@{}}
\RouteGroup{Identidad estructural}
\RouteRow{Clase y multisección}{charged\_\allowbreak{}lepton\_\allowbreak{}G4\quad /\quad charged\_\allowbreak{}lepton\_\allowbreak{}G4}
\RouteRow{Colección y sector}{leptón cargado; leptón}
\RouteRow{Clasificación física}{leptón cargado}
\RouteRow{Generación, profundidad y color}{4; G4; \textemdash{}}
\RouteRow{Ruta primaria y órbita de inversión}{\(B_1\): no\qquad 1,\allowbreak{}3;\allowbreak{}9,\allowbreak{}7}
\RouteRow{Firma de clasificación}{28|\allowbreak{}6|\allowbreak{}-2|\allowbreak{}0|\allowbreak{}-12|\allowbreak{}+++-\kern0pt{}-\kern0pt{}-+++-\kern0pt{}-\kern0pt{}-}
\RouteGroup{Retornos, memoria y transporte}
\RouteRow{Retornos con fase}{clase: 9; órbita: 27; celda: 27}
\RouteRow{Retorno cinemático}{en 108: 54; extensión: 216}
\RouteRow{Giros y cambios de hoja}{71; 107}
\RouteRow{Acarreo y diferimiento}{deuda total: 0; final: 0; bloques diferidos: 0}
\RouteRow{Tríadas finales; persistencia de Pareto}{17; no}
\RouteRow{\(K\longrightarrow K+9\)}{repeticiones: 9; cambios de memoria: 9}
\RouteGroup{Firma, fase y diferencias}
\RouteRow{\(n_A,\ s_C,\ \kappa_0,\ \nu_{120},\ \nu_{270}\)}{28, 6, 0, 4, -12}
\RouteRow{\(\tau_{12}\); firma histórica \(\mathbb Z^5\)}{+++-\kern0pt{}-\kern0pt{}-+++-\kern0pt{}-\kern0pt{}-; 28|\allowbreak{}6|\allowbreak{}0|\allowbreak{}4|\allowbreak{}-12}
\RouteRow{\(D_1,\ D_3,\ D_4,\ D_{27},\ D_{36}\)}{66, 60, 66, 48, 32}
\RouteRow{\(\Omega_{90/120},\ P_6\)}{80, 6}
\RouteGroup{Lectores y factores de acción}
\RouteRow{\(K_D\quad /\quad K_\Omega\)}{(80,\allowbreak{}66,\allowbreak{}3)\quad /\quad (80,\allowbreak{}54,\allowbreak{}6)}
\RouteRow{\(\kappa_D\)}{79.51870\allowbreak{}83196953\allowbreak{}45554341\allowbreak{}34}
\RouteRow{\(\kappa_\Omega\)}{79.60661\allowbreak{}01397948\allowbreak{}27586470\allowbreak{}91}
\RouteRow{\(R_{\mathrm{act},D}\)}{1.000000\allowbreak{}05512887\allowbreak{}17997050\allowbreak{}72}
\RouteRow{\(R_{\mathrm{act},\Omega}\)}{1.000000\allowbreak{}05518981\allowbreak{}25318591\allowbreak{}40}
\RouteGroup{Separación de prefijos}
\RouteRow{Área de ambigüedad; área \(\log_2\)}{54; 28.52932\allowbreak{}50129808}
\RouteRow{Primer bloque de separación}{órbita: \textemdash{}; celda: \textemdash{}; ruta: \textemdash{}}
\RouteRow{Ambigüedad final}{3}
\RouteGroup{Secuencias completas}
\RouteRow{\(w_6\)}{243 657 24 243 657 570 264 90 570 243 657 24 243 657 570 264 90 570}
\RouteRow{Tríadas estables por bloque}{0 1 2 3 4 5 6 7 8 9 10 11 12 13 14 15 16 17}
\end{tabularx}\endgroup\par\vspace{5pt}\noindent{\fontsize{8.4}{10}\selectfont\hyperref[trace-316]{Procedencia y códigos originales}\hfill Ficha completa · 316/324}
\clearpage\phantomsection\label{nav:vi-ruta-317}
% VI-CATALOG-ROW rutas 317
\pdfbookmark[2]{Celda 80}{cell-80}
\pdfbookmark[3]{r9c8\_h-1v-1}{route-317}
\RouteTitle{Ruta 317}{r9c8\_\allowbreak{}h-1v-1}
\noindent Celda 80\quad (9, 8)\qquad Orientación \(h=-1,\ v=-1\)\hfill\hyperref[sec:vi-indice-rutas]{Índice}\par\vspace{4pt}
\begingroup\fontsize{10.2}{12.5}\selectfont\setlength{\tabcolsep}{5pt}\renewcommand{\arraystretch}{1.1}\rowcolors{1}{routepale}{white}\noindent\begin{tabularx}{\linewidth}{@{}>{\raggedright\arraybackslash}p{.345\linewidth}>{\raggedright\arraybackslash}X@{}}
\RouteGroup{Identidad estructural}
\RouteRow{Clase y multisección}{neutrino\_\allowbreak{}G4\quad /\quad neutrino\_\allowbreak{}G4}
\RouteRow{Colección y sector}{neutrino; leptón}
\RouteRow{Clasificación física}{neutrino}
\RouteRow{Generación, profundidad y color}{4; G4; \textemdash{}}
\RouteRow{Ruta primaria y órbita de inversión}{\(B_1\): no\qquad 1,\allowbreak{}2;\allowbreak{}9,\allowbreak{}8}
\RouteRow{Firma de clasificación}{28|\allowbreak{}-6|\allowbreak{}4|\allowbreak{}0|\allowbreak{}-8|\allowbreak{}+0+-\kern0pt{}-\kern0pt{}-+0+-\kern0pt{}-\kern0pt{}-}
\RouteGroup{Retornos, memoria y transporte}
\RouteRow{Retornos con fase}{clase: 27; órbita: 27; celda: 27}
\RouteRow{Retorno cinemático}{en 108: 54; extensión: 216}
\RouteRow{Giros y cambios de hoja}{71; 107}
\RouteRow{Acarreo y diferimiento}{deuda total: 1; final: 0; bloques diferidos: 1}
\RouteRow{Tríadas finales; persistencia de Pareto}{17; sí}
\RouteRow{\(K\longrightarrow K+9\)}{repeticiones: 9; cambios de memoria: 9}
\RouteGroup{Firma, fase y diferencias}
\RouteRow{\(n_A,\ s_C,\ \kappa_0,\ \nu_{120},\ \nu_{270}\)}{28, 6, -2, 0, -8}
\RouteRow{\(\tau_{12}\); firma histórica \(\mathbb Z^5\)}{+++-0-+++-0-; 28|\allowbreak{}6|\allowbreak{}-2|\allowbreak{}0|\allowbreak{}-8}
\RouteRow{\(D_1,\ D_3,\ D_4,\ D_{27},\ D_{36}\)}{66, 60, 66, 48, 32}
\RouteRow{\(\Omega_{90/120},\ P_6\)}{48, 5}
\RouteGroup{Lectores y factores de acción}
\RouteRow{\(K_D\quad /\quad K_\Omega\)}{(80,\allowbreak{}66,\allowbreak{}3)\quad /\quad (48,\allowbreak{}54,\allowbreak{}5)}
\RouteRow{\(\kappa_D\)}{79.51870\allowbreak{}83196953\allowbreak{}45554341\allowbreak{}34}
\RouteRow{\(\kappa_\Omega\)}{47.60649\allowbreak{}89433721\allowbreak{}35446416\allowbreak{}86}
\RouteRow{\(R_{\mathrm{act},D}\)}{1.000000\allowbreak{}05512887\allowbreak{}17997050\allowbreak{}72}
\RouteRow{\(R_{\mathrm{act},\Omega}\)}{1.000000\allowbreak{}03300471\allowbreak{}80532430\allowbreak{}84}
\RouteGroup{Separación de prefijos}
\RouteRow{Área de ambigüedad; área \(\log_2\)}{40; 20.33985\allowbreak{}00028846\allowbreak{}24}
\RouteRow{Primer bloque de separación}{órbita: \textemdash{}; celda: \textemdash{}; ruta: \textemdash{}}
\RouteRow{Ambigüedad final}{2}
\RouteGroup{Secuencias completas}
\RouteRow{\(w_6\)}{486 423 15 486 423 327 498 99 327 486 423 15 486 423 327 498 99 327}
\RouteRow{Tríadas estables por bloque}{0 1 2 3 4 5 6 7 8 9 10 10 12 13 14 15 16 17}
\end{tabularx}\endgroup\par\vspace{5pt}\noindent{\fontsize{8.4}{10}\selectfont\hyperref[trace-317]{Procedencia y códigos originales}\hfill Ficha completa · 317/324}
\clearpage\phantomsection\label{nav:vi-ruta-318}
% VI-CATALOG-ROW rutas 318
\pdfbookmark[3]{r9c8\_h-1v+1}{route-318}
\RouteTitle{Ruta 318}{r9c8\_\allowbreak{}h-1v+1}
\noindent Celda 80\quad (9, 8)\qquad Orientación \(h=-1,\ v=1\)\hfill\hyperref[sec:vi-indice-rutas]{Índice}\par\vspace{4pt}
\begingroup\fontsize{10.2}{12.5}\selectfont\setlength{\tabcolsep}{5pt}\renewcommand{\arraystretch}{1.1}\rowcolors{1}{routepale}{white}\noindent\begin{tabularx}{\linewidth}{@{}>{\raggedright\arraybackslash}p{.345\linewidth}>{\raggedright\arraybackslash}X@{}}
\RouteGroup{Identidad estructural}
\RouteRow{Clase y multisección}{neutrino\_\allowbreak{}G4\quad /\quad neutrino\_\allowbreak{}G4}
\RouteRow{Colección y sector}{neutrino; leptón}
\RouteRow{Clasificación física}{neutrino}
\RouteRow{Generación, profundidad y color}{4; G4; \textemdash{}}
\RouteRow{Ruta primaria y órbita de inversión}{\(B_1\): sí\qquad 1,\allowbreak{}2;\allowbreak{}9,\allowbreak{}8}
\RouteRow{Firma de clasificación}{28|\allowbreak{}-6|\allowbreak{}4|\allowbreak{}0|\allowbreak{}-8|\allowbreak{}+0+-\kern0pt{}-\kern0pt{}-+0+-\kern0pt{}-\kern0pt{}-}
\RouteGroup{Retornos, memoria y transporte}
\RouteRow{Retornos con fase}{clase: 27; órbita: 27; celda: 27}
\RouteRow{Retorno cinemático}{en 108: 54; extensión: 216}
\RouteRow{Giros y cambios de hoja}{71; 107}
\RouteRow{Acarreo y diferimiento}{deuda total: 1; final: 0; bloques diferidos: 1}
\RouteRow{Tríadas finales; persistencia de Pareto}{17; sí}
\RouteRow{\(K\longrightarrow K+9\)}{repeticiones: 9; cambios de memoria: 9}
\RouteGroup{Firma, fase y diferencias}
\RouteRow{\(n_A,\ s_C,\ \kappa_0,\ \nu_{120},\ \nu_{270}\)}{14, -4, 0, 0, -12}
\RouteRow{\(\tau_{12}\); firma histórica \(\mathbb Z^5\)}{+++-\kern0pt{}-\kern0pt{}-+++-\kern0pt{}-\kern0pt{}-; 14|\allowbreak{}-4|\allowbreak{}0|\allowbreak{}0|\allowbreak{}-12}
\RouteRow{\(D_1,\ D_3,\ D_4,\ D_{27},\ D_{36}\)}{78, 24, 78, 24, 16}
\RouteRow{\(\Omega_{90/120},\ P_6\)}{80, 6}
\RouteGroup{Lectores y factores de acción}
\RouteRow{\(K_D\quad /\quad K_\Omega\)}{(40,\allowbreak{}78,\allowbreak{}27)\quad /\quad (80,\allowbreak{}54,\allowbreak{}6)}
\RouteRow{\(\kappa_D\)}{39.43380\allowbreak{}88030085\allowbreak{}51303671\allowbreak{}14}
\RouteRow{\(\kappa_\Omega\)}{79.60661\allowbreak{}01397948\allowbreak{}27586470\allowbreak{}91}
\RouteRow{\(R_{\mathrm{act},D}\)}{1.000000\allowbreak{}02733874\allowbreak{}08548943\allowbreak{}58}
\RouteRow{\(R_{\mathrm{act},\Omega}\)}{1.000000\allowbreak{}05518981\allowbreak{}25318591\allowbreak{}40}
\RouteGroup{Separación de prefijos}
\RouteRow{Área de ambigüedad; área \(\log_2\)}{36; 18.0}
\RouteRow{Primer bloque de separación}{órbita: \textemdash{}; celda: \textemdash{}; ruta: \textemdash{}}
\RouteRow{Ambigüedad final}{2}
\RouteGroup{Secuencias completas}
\RouteRow{\(w_6\)}{504 585 507 504 585 672 510 666 672 504 585 507 504 585 672 510 666 672}
\RouteRow{Tríadas estables por bloque}{0 1 2 3 4 5 6 7 8 9 10 11 12 13 14 14 16 17}
\end{tabularx}\endgroup\par\vspace{5pt}\noindent{\fontsize{8.4}{10}\selectfont\hyperref[trace-318]{Procedencia y códigos originales}\hfill Ficha completa · 318/324}
\clearpage\phantomsection\label{nav:vi-ruta-319}
% VI-CATALOG-ROW rutas 319
\pdfbookmark[3]{r9c8\_h+1v-1}{route-319}
\RouteTitle{Ruta 319}{r9c8\_\allowbreak{}h+1v-1}
\noindent Celda 80\quad (9, 8)\qquad Orientación \(h=1,\ v=-1\)\hfill\hyperref[sec:vi-indice-rutas]{Índice}\par\vspace{4pt}
\begingroup\fontsize{10.2}{12.5}\selectfont\setlength{\tabcolsep}{5pt}\renewcommand{\arraystretch}{1.1}\rowcolors{1}{routepale}{white}\noindent\begin{tabularx}{\linewidth}{@{}>{\raggedright\arraybackslash}p{.345\linewidth}>{\raggedright\arraybackslash}X@{}}
\RouteGroup{Identidad estructural}
\RouteRow{Clase y multisección}{neutrino\_\allowbreak{}G4\quad /\quad neutrino\_\allowbreak{}G4}
\RouteRow{Colección y sector}{neutrino; leptón}
\RouteRow{Clasificación física}{neutrino}
\RouteRow{Generación, profundidad y color}{4; G4; \textemdash{}}
\RouteRow{Ruta primaria y órbita de inversión}{\(B_1\): no\qquad 1,\allowbreak{}2;\allowbreak{}9,\allowbreak{}8}
\RouteRow{Firma de clasificación}{28|\allowbreak{}-6|\allowbreak{}4|\allowbreak{}0|\allowbreak{}-8|\allowbreak{}+0+-\kern0pt{}-\kern0pt{}-+0+-\kern0pt{}-\kern0pt{}-}
\RouteGroup{Retornos, memoria y transporte}
\RouteRow{Retornos con fase}{clase: 27; órbita: 27; celda: 27}
\RouteRow{Retorno cinemático}{en 108: 54; extensión: 216}
\RouteRow{Giros y cambios de hoja}{71; 107}
\RouteRow{Acarreo y diferimiento}{deuda total: 1; final: 0; bloques diferidos: 1}
\RouteRow{Tríadas finales; persistencia de Pareto}{17; sí}
\RouteRow{\(K\longrightarrow K+9\)}{repeticiones: 9; cambios de memoria: 9}
\RouteGroup{Firma, fase y diferencias}
\RouteRow{\(n_A,\ s_C,\ \kappa_0,\ \nu_{120},\ \nu_{270}\)}{14, 4, 0, 0, -12}
\RouteRow{\(\tau_{12}\); firma histórica \(\mathbb Z^5\)}{+++-\kern0pt{}-\kern0pt{}-+++-\kern0pt{}-\kern0pt{}-; 14|\allowbreak{}4|\allowbreak{}0|\allowbreak{}0|\allowbreak{}-12}
\RouteRow{\(D_1,\ D_3,\ D_4,\ D_{27},\ D_{36}\)}{78, 24, 78, 24, 16}
\RouteRow{\(\Omega_{90/120},\ P_6\)}{80, 6}
\RouteGroup{Lectores y factores de acción}
\RouteRow{\(K_D\quad /\quad K_\Omega\)}{(40,\allowbreak{}78,\allowbreak{}27)\quad /\quad (80,\allowbreak{}54,\allowbreak{}6)}
\RouteRow{\(\kappa_D\)}{39.43380\allowbreak{}88030085\allowbreak{}51303671\allowbreak{}14}
\RouteRow{\(\kappa_\Omega\)}{79.60661\allowbreak{}01397948\allowbreak{}27586470\allowbreak{}91}
\RouteRow{\(R_{\mathrm{act},D}\)}{1.000000\allowbreak{}02733874\allowbreak{}08548943\allowbreak{}58}
\RouteRow{\(R_{\mathrm{act},\Omega}\)}{1.000000\allowbreak{}05518981\allowbreak{}25318591\allowbreak{}40}
\RouteGroup{Separación de prefijos}
\RouteRow{Área de ambigüedad; área \(\log_2\)}{36; 18.0}
\RouteRow{Primer bloque de separación}{órbita: \textemdash{}; celda: \textemdash{}; ruta: \textemdash{}}
\RouteRow{Ambigüedad final}{2}
\RouteGroup{Secuencias completas}
\RouteRow{\(w_6\)}{510 666 672 510 666 507 504 585 507 510 666 672 510 666 507 504 585 507}
\RouteRow{Tríadas estables por bloque}{0 1 2 3 4 5 6 7 8 9 10 11 12 13 14 15 15 17}
\end{tabularx}\endgroup\par\vspace{5pt}\noindent{\fontsize{8.4}{10}\selectfont\hyperref[trace-319]{Procedencia y códigos originales}\hfill Ficha completa · 319/324}
\clearpage\phantomsection\label{nav:vi-ruta-320}
% VI-CATALOG-ROW rutas 320
\pdfbookmark[3]{r9c8\_h+1v+1}{route-320}
\RouteTitle{Ruta 320}{r9c8\_\allowbreak{}h+1v+1}
\noindent Celda 80\quad (9, 8)\qquad Orientación \(h=1,\ v=1\)\hfill\hyperref[sec:vi-indice-rutas]{Índice}\par\vspace{4pt}
\begingroup\fontsize{10.2}{12.5}\selectfont\setlength{\tabcolsep}{5pt}\renewcommand{\arraystretch}{1.1}\rowcolors{1}{routepale}{white}\noindent\begin{tabularx}{\linewidth}{@{}>{\raggedright\arraybackslash}p{.345\linewidth}>{\raggedright\arraybackslash}X@{}}
\RouteGroup{Identidad estructural}
\RouteRow{Clase y multisección}{neutrino\_\allowbreak{}G4\quad /\quad neutrino\_\allowbreak{}G4}
\RouteRow{Colección y sector}{neutrino; leptón}
\RouteRow{Clasificación física}{neutrino}
\RouteRow{Generación, profundidad y color}{4; G4; \textemdash{}}
\RouteRow{Ruta primaria y órbita de inversión}{\(B_1\): no\qquad 1,\allowbreak{}2;\allowbreak{}9,\allowbreak{}8}
\RouteRow{Firma de clasificación}{28|\allowbreak{}-6|\allowbreak{}4|\allowbreak{}0|\allowbreak{}-8|\allowbreak{}+0+-\kern0pt{}-\kern0pt{}-+0+-\kern0pt{}-\kern0pt{}-}
\RouteGroup{Retornos, memoria y transporte}
\RouteRow{Retornos con fase}{clase: 27; órbita: 27; celda: 27}
\RouteRow{Retorno cinemático}{en 108: 54; extensión: 216}
\RouteRow{Giros y cambios de hoja}{71; 107}
\RouteRow{Acarreo y diferimiento}{deuda total: 0; final: 0; bloques diferidos: 0}
\RouteRow{Tríadas finales; persistencia de Pareto}{17; sí}
\RouteRow{\(K\longrightarrow K+9\)}{repeticiones: 9; cambios de memoria: 9}
\RouteGroup{Firma, fase y diferencias}
\RouteRow{\(n_A,\ s_C,\ \kappa_0,\ \nu_{120},\ \nu_{270}\)}{28, -6, 2, 0, -8}
\RouteRow{\(\tau_{12}\); firma histórica \(\mathbb Z^5\)}{+0+-\kern0pt{}-\kern0pt{}-+0+-\kern0pt{}-\kern0pt{}-; 28|\allowbreak{}-6|\allowbreak{}2|\allowbreak{}0|\allowbreak{}-8}
\RouteRow{\(D_1,\ D_3,\ D_4,\ D_{27},\ D_{36}\)}{66, 60, 66, 48, 32}
\RouteRow{\(\Omega_{90/120},\ P_6\)}{48, 5}
\RouteGroup{Lectores y factores de acción}
\RouteRow{\(K_D\quad /\quad K_\Omega\)}{(80,\allowbreak{}66,\allowbreak{}3)\quad /\quad (48,\allowbreak{}54,\allowbreak{}5)}
\RouteRow{\(\kappa_D\)}{79.51870\allowbreak{}83196953\allowbreak{}45554341\allowbreak{}34}
\RouteRow{\(\kappa_\Omega\)}{47.60649\allowbreak{}89433721\allowbreak{}35446416\allowbreak{}86}
\RouteRow{\(R_{\mathrm{act},D}\)}{1.000000\allowbreak{}05512887\allowbreak{}17997050\allowbreak{}72}
\RouteRow{\(R_{\mathrm{act},\Omega}\)}{1.000000\allowbreak{}03300471\allowbreak{}80532430\allowbreak{}84}
\RouteGroup{Separación de prefijos}
\RouteRow{Área de ambigüedad; área \(\log_2\)}{36; 18.0}
\RouteRow{Primer bloque de separación}{órbita: \textemdash{}; celda: \textemdash{}; ruta: \textemdash{}}
\RouteRow{Ambigüedad final}{2}
\RouteGroup{Secuencias completas}
\RouteRow{\(w_6\)}{498 99 327 498 99 15 486 423 15 498 99 327 498 99 15 486 423 15}
\RouteRow{Tríadas estables por bloque}{0 1 2 3 4 5 6 7 8 9 10 11 12 13 14 15 16 17}
\end{tabularx}\endgroup\par\vspace{5pt}\noindent{\fontsize{8.4}{10}\selectfont\hyperref[trace-320]{Procedencia y códigos originales}\hfill Ficha completa · 320/324}
\clearpage\phantomsection\label{nav:vi-ruta-321}
% VI-CATALOG-ROW rutas 321
\pdfbookmark[2]{Celda 81}{cell-81}
\pdfbookmark[3]{r9c9\_h-1v-1}{route-321}
\RouteTitle{Ruta 321}{r9c9\_\allowbreak{}h-1v-1}
\noindent Celda 81\quad (9, 9)\qquad Orientación \(h=-1,\ v=-1\)\hfill\hyperref[sec:vi-indice-rutas]{Índice}\par\vspace{4pt}
\begingroup\fontsize{10.2}{12.5}\selectfont\setlength{\tabcolsep}{5pt}\renewcommand{\arraystretch}{1.1}\rowcolors{1}{routepale}{white}\noindent\begin{tabularx}{\linewidth}{@{}>{\raggedright\arraybackslash}p{.345\linewidth}>{\raggedright\arraybackslash}X@{}}
\RouteGroup{Identidad estructural}
\RouteRow{Clase y multisección}{charged\_\allowbreak{}lepton\_\allowbreak{}G4\quad /\quad charged\_\allowbreak{}lepton\_\allowbreak{}G4}
\RouteRow{Colección y sector}{leptón cargado; leptón}
\RouteRow{Clasificación física}{leptón cargado}
\RouteRow{Generación, profundidad y color}{4; G4; \textemdash{}}
\RouteRow{Ruta primaria y órbita de inversión}{\(B_1\): no\qquad 1,\allowbreak{}1;\allowbreak{}9,\allowbreak{}9}
\RouteRow{Firma de clasificación}{24|\allowbreak{}0|\allowbreak{}12|\allowbreak{}0|\allowbreak{}-12|\allowbreak{}+++-\kern0pt{}-\kern0pt{}-+++-\kern0pt{}-\kern0pt{}-}
\RouteGroup{Retornos, memoria y transporte}
\RouteRow{Retornos con fase}{clase: 27; órbita: 27; celda: 27}
\RouteRow{Retorno cinemático}{en 108: 54; extensión: 216}
\RouteRow{Giros y cambios de hoja}{71; 107}
\RouteRow{Acarreo y diferimiento}{deuda total: 1; final: 0; bloques diferidos: 1}
\RouteRow{Tríadas finales; persistencia de Pareto}{17; no}
\RouteRow{\(K\longrightarrow K+9\)}{repeticiones: 9; cambios de memoria: 9}
\RouteGroup{Firma, fase y diferencias}
\RouteRow{\(n_A,\ s_C,\ \kappa_0,\ \nu_{120},\ \nu_{270}\)}{24, 0, 0, 0, -12}
\RouteRow{\(\tau_{12}\); firma histórica \(\mathbb Z^5\)}{+++-\kern0pt{}-\kern0pt{}-+++-\kern0pt{}-\kern0pt{}-; 24|\allowbreak{}0|\allowbreak{}0|\allowbreak{}0|\allowbreak{}-12}
\RouteRow{\(D_1,\ D_3,\ D_4,\ D_{27},\ D_{36}\)}{54, 60, 72, 24, 16}
\RouteRow{\(\Omega_{90/120},\ P_6\)}{80, 6}
\RouteGroup{Lectores y factores de acción}
\RouteRow{\(K_D\quad /\quad K_\Omega\)}{(40,\allowbreak{}54,\allowbreak{}6)\quad /\quad (80,\allowbreak{}54,\allowbreak{}6)}
\RouteRow{\(\kappa_D\)}{39.60661\allowbreak{}01397948\allowbreak{}27586470\allowbreak{}91}
\RouteRow{\(\kappa_\Omega\)}{79.60661\allowbreak{}01397948\allowbreak{}27586470\allowbreak{}91}
\RouteRow{\(R_{\mathrm{act},D}\)}{1.000000\allowbreak{}02745854\allowbreak{}08735595\allowbreak{}18}
\RouteRow{\(R_{\mathrm{act},\Omega}\)}{1.000000\allowbreak{}05518981\allowbreak{}25318591\allowbreak{}40}
\RouteGroup{Separación de prefijos}
\RouteRow{Área de ambigüedad; área \(\log_2\)}{58; 30.18947\allowbreak{}50100961\allowbreak{}74}
\RouteRow{Primer bloque de separación}{órbita: \textemdash{}; celda: \textemdash{}; ruta: \textemdash{}}
\RouteRow{Ambigüedad final}{3}
\RouteGroup{Secuencias completas}
\RouteRow{\(w_6\)}{15 486 423 15 486 657 24 243 657 15 486 423 15 486 657 24 243 657}
\RouteRow{Tríadas estables por bloque}{0 1 2 3 4 5 6 7 8 9 10 11 12 13 14 15 15 17}
\end{tabularx}\endgroup\par\vspace{5pt}\noindent{\fontsize{8.4}{10}\selectfont\hyperref[trace-321]{Procedencia y códigos originales}\hfill Ficha completa · 321/324}
\clearpage\phantomsection\label{nav:vi-ruta-322}
% VI-CATALOG-ROW rutas 322
\pdfbookmark[3]{r9c9\_h-1v+1}{route-322}
\RouteTitle{Ruta 322}{r9c9\_\allowbreak{}h-1v+1}
\noindent Celda 81\quad (9, 9)\qquad Orientación \(h=-1,\ v=1\)\hfill\hyperref[sec:vi-indice-rutas]{Índice}\par\vspace{4pt}
\begingroup\fontsize{10.2}{12.5}\selectfont\setlength{\tabcolsep}{5pt}\renewcommand{\arraystretch}{1.1}\rowcolors{1}{routepale}{white}\noindent\begin{tabularx}{\linewidth}{@{}>{\raggedright\arraybackslash}p{.345\linewidth}>{\raggedright\arraybackslash}X@{}}
\RouteGroup{Identidad estructural}
\RouteRow{Clase y multisección}{charged\_\allowbreak{}lepton\_\allowbreak{}G4\quad /\quad charged\_\allowbreak{}lepton\_\allowbreak{}G4}
\RouteRow{Colección y sector}{leptón cargado; leptón}
\RouteRow{Clasificación física}{leptón cargado}
\RouteRow{Generación, profundidad y color}{4; G4; \textemdash{}}
\RouteRow{Ruta primaria y órbita de inversión}{\(B_1\): sí\qquad 1,\allowbreak{}1;\allowbreak{}9,\allowbreak{}9}
\RouteRow{Firma de clasificación}{24|\allowbreak{}0|\allowbreak{}12|\allowbreak{}0|\allowbreak{}-12|\allowbreak{}+++-\kern0pt{}-\kern0pt{}-+++-\kern0pt{}-\kern0pt{}-}
\RouteGroup{Retornos, memoria y transporte}
\RouteRow{Retornos con fase}{clase: 27; órbita: 27; celda: 27}
\RouteRow{Retorno cinemático}{en 108: 54; extensión: 216}
\RouteRow{Giros y cambios de hoja}{71; 107}
\RouteRow{Acarreo y diferimiento}{deuda total: 0; final: 0; bloques diferidos: 0}
\RouteRow{Tríadas finales; persistencia de Pareto}{17; no}
\RouteRow{\(K\longrightarrow K+9\)}{repeticiones: 9; cambios de memoria: 9}
\RouteGroup{Firma, fase y diferencias}
\RouteRow{\(n_A,\ s_C,\ \kappa_0,\ \nu_{120},\ \nu_{270}\)}{20, 0, 0, 0, -4}
\RouteRow{\(\tau_{12}\); firma histórica \(\mathbb Z^5\)}{0+00-00+00-0; 20|\allowbreak{}0|\allowbreak{}0|\allowbreak{}0|\allowbreak{}-4}
\RouteRow{\(D_1,\ D_3,\ D_4,\ D_{27},\ D_{36}\)}{24, 24, 24, 0, 0}
\RouteRow{\(\Omega_{90/120},\ P_6\)}{24, 2}
\RouteGroup{Lectores y factores de acción}
\RouteRow{\(K_D\quad /\quad K_\Omega\)}{(0,\allowbreak{}24,\allowbreak{}0)\quad /\quad (24,\allowbreak{}54,\allowbreak{}2)}
\RouteRow{\(\kappa_D\)}{-0.175136\allowbreak{}46166281\allowbreak{}12239348\allowbreak{}425}
\RouteRow{\(\kappa_\Omega\)}{23.60616\allowbreak{}53541040\allowbreak{}59026254\allowbreak{}71}
\RouteRow{\(R_{\mathrm{act},D}\)}{0.999999\allowbreak{}99987858\allowbreak{}10851337\allowbreak{}881}
\RouteRow{\(R_{\mathrm{act},\Omega}\)}{1.000000\allowbreak{}01636572\allowbreak{}40637534\allowbreak{}03}
\RouteGroup{Separación de prefijos}
\RouteRow{Área de ambigüedad; área \(\log_2\)}{108; 46.52932\allowbreak{}50129807\allowbreak{}9}
\RouteRow{Primer bloque de separación}{órbita: \textemdash{}; celda: \textemdash{}; ruta: \textemdash{}}
\RouteRow{Ambigüedad final}{6}
\RouteGroup{Secuencias completas}
\RouteRow{\(w_6\)}{3 0 81 3 0 81 3 0 81 3 0 81 3 0 81 3 0 81}
\RouteRow{Tríadas estables por bloque}{0 1 2 3 4 5 6 7 8 9 10 11 12 13 14 15 16 17}
\end{tabularx}\endgroup\par\vspace{5pt}\noindent{\fontsize{8.4}{10}\selectfont\hyperref[trace-322]{Procedencia y códigos originales}\hfill Ficha completa · 322/324}
\clearpage\phantomsection\label{nav:vi-ruta-323}
% VI-CATALOG-ROW rutas 323
\pdfbookmark[3]{r9c9\_h+1v-1}{route-323}
\RouteTitle{Ruta 323}{r9c9\_\allowbreak{}h+1v-1}
\noindent Celda 81\quad (9, 9)\qquad Orientación \(h=1,\ v=-1\)\hfill\hyperref[sec:vi-indice-rutas]{Índice}\par\vspace{4pt}
\begingroup\fontsize{10.2}{12.5}\selectfont\setlength{\tabcolsep}{5pt}\renewcommand{\arraystretch}{1.1}\rowcolors{1}{routepale}{white}\noindent\begin{tabularx}{\linewidth}{@{}>{\raggedright\arraybackslash}p{.345\linewidth}>{\raggedright\arraybackslash}X@{}}
\RouteGroup{Identidad estructural}
\RouteRow{Clase y multisección}{charged\_\allowbreak{}lepton\_\allowbreak{}G4\quad /\quad charged\_\allowbreak{}lepton\_\allowbreak{}G4}
\RouteRow{Colección y sector}{leptón cargado; leptón}
\RouteRow{Clasificación física}{leptón cargado}
\RouteRow{Generación, profundidad y color}{4; G4; \textemdash{}}
\RouteRow{Ruta primaria y órbita de inversión}{\(B_1\): no\qquad 1,\allowbreak{}1;\allowbreak{}9,\allowbreak{}9}
\RouteRow{Firma de clasificación}{24|\allowbreak{}0|\allowbreak{}12|\allowbreak{}0|\allowbreak{}-12|\allowbreak{}+++-\kern0pt{}-\kern0pt{}-+++-\kern0pt{}-\kern0pt{}-}
\RouteGroup{Retornos, memoria y transporte}
\RouteRow{Retornos con fase}{clase: 27; órbita: 27; celda: 27}
\RouteRow{Retorno cinemático}{en 108: 54; extensión: 216}
\RouteRow{Giros y cambios de hoja}{71; 107}
\RouteRow{Acarreo y diferimiento}{deuda total: 0; final: 0; bloques diferidos: 0}
\RouteRow{Tríadas finales; persistencia de Pareto}{17; no}
\RouteRow{\(K\longrightarrow K+9\)}{repeticiones: 9; cambios de memoria: 9}
\RouteGroup{Firma, fase y diferencias}
\RouteRow{\(n_A,\ s_C,\ \kappa_0,\ \nu_{120},\ \nu_{270}\)}{20, 0, 0, 0, -4}
\RouteRow{\(\tau_{12}\); firma histórica \(\mathbb Z^5\)}{0+00-00+00-0; 20|\allowbreak{}0|\allowbreak{}0|\allowbreak{}0|\allowbreak{}-4}
\RouteRow{\(D_1,\ D_3,\ D_4,\ D_{27},\ D_{36}\)}{24, 24, 24, 0, 0}
\RouteRow{\(\Omega_{90/120},\ P_6\)}{24, 2}
\RouteGroup{Lectores y factores de acción}
\RouteRow{\(K_D\quad /\quad K_\Omega\)}{(0,\allowbreak{}24,\allowbreak{}0)\quad /\quad (24,\allowbreak{}54,\allowbreak{}2)}
\RouteRow{\(\kappa_D\)}{-0.175136\allowbreak{}46166281\allowbreak{}12239348\allowbreak{}425}
\RouteRow{\(\kappa_\Omega\)}{23.60616\allowbreak{}53541040\allowbreak{}59026254\allowbreak{}71}
\RouteRow{\(R_{\mathrm{act},D}\)}{0.999999\allowbreak{}99987858\allowbreak{}10851337\allowbreak{}881}
\RouteRow{\(R_{\mathrm{act},\Omega}\)}{1.000000\allowbreak{}01636572\allowbreak{}40637534\allowbreak{}03}
\RouteGroup{Separación de prefijos}
\RouteRow{Área de ambigüedad; área \(\log_2\)}{108; 46.52932\allowbreak{}50129807\allowbreak{}9}
\RouteRow{Primer bloque de separación}{órbita: \textemdash{}; celda: \textemdash{}; ruta: \textemdash{}}
\RouteRow{Ambigüedad final}{6}
\RouteGroup{Secuencias completas}
\RouteRow{\(w_6\)}{3 0 81 3 0 81 3 0 81 3 0 81 3 0 81 3 0 81}
\RouteRow{Tríadas estables por bloque}{0 1 2 3 4 5 6 7 8 9 10 11 12 13 14 15 16 17}
\end{tabularx}\endgroup\par\vspace{5pt}\noindent{\fontsize{8.4}{10}\selectfont\hyperref[trace-323]{Procedencia y códigos originales}\hfill Ficha completa · 323/324}
\clearpage\phantomsection\label{nav:vi-ruta-324}
% VI-CATALOG-ROW rutas 324
\pdfbookmark[3]{r9c9\_h+1v+1}{route-324}
\RouteTitle{Ruta 324}{r9c9\_\allowbreak{}h+1v+1}
\noindent Celda 81\quad (9, 9)\qquad Orientación \(h=1,\ v=1\)\hfill\hyperref[sec:vi-indice-rutas]{Índice}\par\vspace{4pt}
\begingroup\fontsize{10.2}{12.5}\selectfont\setlength{\tabcolsep}{5pt}\renewcommand{\arraystretch}{1.1}\rowcolors{1}{routepale}{white}\noindent\begin{tabularx}{\linewidth}{@{}>{\raggedright\arraybackslash}p{.345\linewidth}>{\raggedright\arraybackslash}X@{}}
\RouteGroup{Identidad estructural}
\RouteRow{Clase y multisección}{charged\_\allowbreak{}lepton\_\allowbreak{}G4\quad /\quad charged\_\allowbreak{}lepton\_\allowbreak{}G4}
\RouteRow{Colección y sector}{leptón cargado; leptón}
\RouteRow{Clasificación física}{leptón cargado}
\RouteRow{Generación, profundidad y color}{4; G4; \textemdash{}}
\RouteRow{Ruta primaria y órbita de inversión}{\(B_1\): no\qquad 1,\allowbreak{}1;\allowbreak{}9,\allowbreak{}9}
\RouteRow{Firma de clasificación}{24|\allowbreak{}0|\allowbreak{}12|\allowbreak{}0|\allowbreak{}-12|\allowbreak{}+++-\kern0pt{}-\kern0pt{}-+++-\kern0pt{}-\kern0pt{}-}
\RouteGroup{Retornos, memoria y transporte}
\RouteRow{Retornos con fase}{clase: 27; órbita: 27; celda: 27}
\RouteRow{Retorno cinemático}{en 108: 54; extensión: 216}
\RouteRow{Giros y cambios de hoja}{71; 107}
\RouteRow{Acarreo y diferimiento}{deuda total: 2; final: 1; bloques diferidos: 2}
\RouteRow{Tríadas finales; persistencia de Pareto}{16; no}
\RouteRow{\(K\longrightarrow K+9\)}{repeticiones: 9; cambios de memoria: 9}
\RouteGroup{Firma, fase y diferencias}
\RouteRow{\(n_A,\ s_C,\ \kappa_0,\ \nu_{120},\ \nu_{270}\)}{24, 0, 0, 0, -12}
\RouteRow{\(\tau_{12}\); firma histórica \(\mathbb Z^5\)}{+++-\kern0pt{}-\kern0pt{}-+++-\kern0pt{}-\kern0pt{}-; 24|\allowbreak{}0|\allowbreak{}0|\allowbreak{}0|\allowbreak{}-12}
\RouteRow{\(D_1,\ D_3,\ D_4,\ D_{27},\ D_{36}\)}{54, 60, 72, 24, 16}
\RouteRow{\(\Omega_{90/120},\ P_6\)}{80, 6}
\RouteGroup{Lectores y factores de acción}
\RouteRow{\(K_D\quad /\quad K_\Omega\)}{(40,\allowbreak{}54,\allowbreak{}6)\quad /\quad (80,\allowbreak{}54,\allowbreak{}6)}
\RouteRow{\(\kappa_D\)}{39.60661\allowbreak{}01397948\allowbreak{}27586470\allowbreak{}91}
\RouteRow{\(\kappa_\Omega\)}{79.60661\allowbreak{}01397948\allowbreak{}27586470\allowbreak{}91}
\RouteRow{\(R_{\mathrm{act},D}\)}{1.000000\allowbreak{}02745854\allowbreak{}08735595\allowbreak{}18}
\RouteRow{\(R_{\mathrm{act},\Omega}\)}{1.000000\allowbreak{}05518981\allowbreak{}25318591\allowbreak{}40}
\RouteGroup{Separación de prefijos}
\RouteRow{Área de ambigüedad; área \(\log_2\)}{58; 30.18947\allowbreak{}50100961\allowbreak{}74}
\RouteRow{Primer bloque de separación}{órbita: \textemdash{}; celda: \textemdash{}; ruta: \textemdash{}}
\RouteRow{Ambigüedad final}{3}
\RouteGroup{Secuencias completas}
\RouteRow{\(w_6\)}{24 243 657 24 243 423 15 486 423 24 243 657 24 243 423 15 486 423}
\RouteRow{Tríadas estables por bloque}{0 1 2 3 4 5 6 7 8 9 10 11 12 13 13 15 16 16}
\end{tabularx}\endgroup\par\vspace{5pt}\noindent{\fontsize{8.4}{10}\selectfont\hyperref[trace-324]{Procedencia y códigos originales}\hfill Ficha completa · 324/324}
\clearpage
\section{Inventario metrológico: 471 masas y 384 anchuras}
\label{sec:vi-inventario-metrologico}

Las tablas siguientes son un catálogo de observaciones externas y no un
recuento de predicciones HMT. Conservan el corte documental PDG 2026 con fecha
\texttt{2026-01-15}, tal como figura en los archivos preservados. La fecha de
consulta consignada en el integral es \texttt{2026-07-26}. No se ha actualizado
la edición ni se han reinterpretado sus observables durante esta transcripción.
El identificador estable distingue registros diferentes de una misma especie;
no se eliminan repeticiones ni se funden masa, polo, anchura o estimación.

La columna de valor conserva tanto la expresión literal como, cuando difiere,
su presentación de origen. No se reconstruye una incertidumbre a partir de
decimales ni se transforma un intervalo en una desviación estándar.
Se mantienen las incertidumbres positiva y negativa separadas, la unidad,
el tipo de límite y el nivel de confianza disponible. \texttt{ND} abrevia
\texttt{UNAVAILABLE}; \texttt{NA} abrevia \texttt{NOT\_APPLICABLE}.
Una raya significa un campo vacío. Los tipos originales \texttt{NONE},
\texttt{R} y \texttt{U} se conservan como códigos de la fuente, junto a la
expresión del valor: respectivamente, ausencia de marca de límite, rango y
límite superior. No se infiere un nivel de confianza donde el registro no lo da.

En este corte, los 855 registros indican que esquema y escala no están
disponibles en la instantánea. Todos conservan la condición de comparación
posterior a la fijación del resultado interno y declaran que el valor externo
no se utilizó como selector. Las claves de estado \(E\) y enlace \(M\)
son abreviaturas editoriales reversibles de los estados y mapas de la fuente,
descritos en las dos leyendas siguientes. La palabra «abierto» en esas leyendas
reproduce el estado documental de ese corte; no es una revisión matemática nueva.

\begingroup
\setlength{\tabcolsep}{4pt}
\setlength{\extrarowheight}{4pt}
\renewcommand{\arraystretch}{1.13}
\begin{longtable}{@{}>{\raggedright\arraybackslash}p{\dimexpr 0.08860759\linewidth-2\tabcolsep\relax}>{\raggedright\arraybackslash}p{\dimexpr 0.41772152\linewidth-2\tabcolsep\relax}>{\raggedright\arraybackslash}p{\dimexpr 0.49367089\linewidth-2\tabcolsep\relax}@{}}
\caption{Estados conservados del inventario externo.}\label{tab:vi-estados-externos}\\
\toprule
Clave & Estado literal & Lectura documental \\
\midrule
\endfirsthead
\multicolumn{3}{l}{\textit{Continuación}}\\
\toprule
Clave & Estado literal & Lectura documental \\
\midrule
\endhead
\midrule
\multicolumn{3}{r}{Continúa en la página siguiente}\\
\endfoot
\bottomrule
\endlastfoot
E01 & CONDITIONED\_\allowbreak{}NOMINAL\_\allowbreak{}SIGNATURE\_\allowbreak{}\_\allowbreak{}ROUTE\_\allowbreak{}AND\_\allowbreak{}POLE\_\allowbreak{}REALIZATION\_\allowbreak{}OPEN & Firma nominal condicionada; realización de ruta y polo abierta. \\
E02 & EXACT\_\allowbreak{}COLOR\_\allowbreak{}FIBER\_\allowbreak{}OPERATOR\_\allowbreak{}\_\allowbreak{}NOMINAL\_\allowbreak{}SIGNATURE\_\allowbreak{}CONDITIONED\_\allowbreak{}\_\allowbreak{}SCHEME\_\allowbreak{}OPEN & Operador de fibra de color exacto; firma nominal condicionada; esquema abierto. \\
E03 & EXACT\_\allowbreak{}FIBER\_\allowbreak{}OPERATOR\_\allowbreak{}\_\allowbreak{}NOMINAL\_\allowbreak{}SIGNATURE\_\allowbreak{}CONDITIONED\_\allowbreak{}\_\allowbreak{}SCALARIZATION\_\allowbreak{}OPEN & Operador de fibra exacto; firma nominal condicionada; escalarización abierta. \\
E04 & EXTERNAL\_\allowbreak{}RECORD\_\allowbreak{}ACCOUNTED\_\allowbreak{}\_\allowbreak{}INDIVIDUAL\_\allowbreak{}HMT\_\allowbreak{}REALIZATION\_\allowbreak{}NOT\_\allowbreak{}YET\_\allowbreak{}JOINED & Registro externo censado; realización HMT individual todavía no enlazada en este corte. \\
E05 & NULL\_\allowbreak{}SECTION\_\allowbreak{}CLOSED & Sección nula cerrada. \\
E06 & SCALAR\_\allowbreak{}CLOSED\_\allowbreak{}ELECTRON\_\allowbreak{}BETA & Escalar electrónico beta cerrado. \\
\end{longtable}
\endgroup

\begingroup
\setlength{\tabcolsep}{4pt}
\setlength{\extrarowheight}{4pt}
\renewcommand{\arraystretch}{1.13}
\begin{longtable}{@{}>{\raggedright\arraybackslash}p{\dimexpr 0.08860759\linewidth-2\tabcolsep\relax}>{\raggedright\arraybackslash}p{\dimexpr 0.41772152\linewidth-2\tabcolsep\relax}>{\raggedright\arraybackslash}p{\dimexpr 0.49367089\linewidth-2\tabcolsep\relax}@{}}
\caption{Enlaces de realización consignados por el inventario.}\label{tab:vi-enlaces-externos}\\
\toprule
Clave & Estado literal & Lectura documental \\
\midrule
\endfirsthead
\multicolumn{3}{l}{\textit{Continuación}}\\
\toprule
Clave & Estado literal & Lectura documental \\
\midrule
\endhead
\midrule
\multicolumn{3}{r}{Continúa en la página siguiente}\\
\endfoot
\bottomrule
\endlastfoot
M01 & NONE & Ninguno, según el registro. \\
M02 & NONE\_\allowbreak{}FOR\_\allowbreak{}MASS\_\allowbreak{}NULLITY & Ninguno para la nulidad de masa, según el registro. \\
M03 & flavor\_\allowbreak{}color\_\allowbreak{}coupling\_\allowbreak{}to\_\allowbreak{}spectral\_\allowbreak{}projection\_\allowbreak{}then\_\allowbreak{}scheme\_\allowbreak{}scale\_\allowbreak{}transport & Acoplamiento sabor–color a la proyección espectral y transporte de esquema y escala. \\
M04 & generation\_\allowbreak{}coupling\_\allowbreak{}to\_\allowbreak{}spectral\_\allowbreak{}projection\_\allowbreak{}without\_\allowbreak{}mass\_\allowbreak{}target & Acoplamiento generacional a la proyección espectral sin masa objetivo. \\
M05 & identity\_\allowbreak{}to\_\allowbreak{}public\_\allowbreak{}sector\_\allowbreak{}or\_\allowbreak{}new\_\allowbreak{}typed\_\allowbreak{}sector & Identidad y sector identificado o nuevo sector tipado. \\
M06 & join\_\allowbreak{}named\_\allowbreak{}boson\_\allowbreak{}to\_\allowbreak{}enriched\_\allowbreak{}route\_\allowbreak{}then\_\allowbreak{}complex\_\allowbreak{}pole\_\allowbreak{}operator & Bosón individualizado, ruta con memoria y operador de polo complejo. \\
M07 & named\_\allowbreak{}baryon\_\allowbreak{}binding\_\allowbreak{}map & Enlace de unión del barión individualizado. \\
M08 & named\_\allowbreak{}meson\_\allowbreak{}binding\_\allowbreak{}map & Enlace de unión del mesón individualizado. \\
M09 & quark\_\allowbreak{}identity\_\allowbreak{}scheme\_\allowbreak{}scale\_\allowbreak{}map & Identidad del quark, esquema y escala. \\
M10 & typed\_\allowbreak{}gravitational\_\allowbreak{}representation\_\allowbreak{}map & Representación gravitatoria tipada. \\
\end{longtable}
\endgroup
\clearpage\subsection{Masas: 471 registros}
\begingroup\fontsize{9.2}{11.1}\selectfont\setlength{\tabcolsep}{5pt}\setlength{\extrarowheight}{3pt}\renewcommand{\arraystretch}{1.08}
\begin{longtable}{@{}>{\raggedright\arraybackslash}p{.445\linewidth}>{\raggedright\arraybackslash}p{.205\linewidth}>{\raggedright\arraybackslash}p{.18\linewidth}>{\raggedright\arraybackslash}p{\dimexpr.17\linewidth-6\tabcolsep\relax}@{}}
\toprule
Identidad y especies & Valor y unidad & Incertidumbre & Sector y estado\\\midrule
\endfirsthead
\toprule
Identidad y especies & Valor y unidad & Incertidumbre & Sector y estado\\\midrule
\endhead
\bottomrule
\endfoot
% VI-CATALOG-ROW masas 1
\textbf{1. Sigma(1580) MASS}\par Sigma(1580)\allowbreak{}- | Sigmabar(1580)\allowbreak{}+ | Sigma(1580)\allowbreak{}0 | Sigmabar(1580)\allowbreak{}0 | Sigma(1580)\allowbreak{}+ | Sigmabar(1580)\allowbreak{}-\par {\fontsize{8.1}{9.4}\selectfont PDG2026:\allowbreak{}observable:\allowbreak{}B000:\allowbreak{}B000M:\allowbreak{}349891} & \textbf{\textasciitilde{}1580} MeV & \(u_+\): ND\par \(u_-\): ND\par Tipo: NONE\par Confianza: ND & COMP-BARYON\par E04 / M07\\\addlinespace[3pt]
% VI-CATALOG-ROW masas 2
\textbf{2. N(1895) BREIT-WIGNER MASS}\par N(1895)\allowbreak{}0 | Nbar(1895)\allowbreak{}0 | N(1895)\allowbreak{}+ | Nbar(1895)\allowbreak{}-\par {\fontsize{8.1}{9.4}\selectfont PDG2026:\allowbreak{}observable:\allowbreak{}B004:\allowbreak{}B004M:\allowbreak{}348705} & \textbf{1870 to 1895 to 1920} MeV\par Presentación: 1870 to 1920 (\textasciitilde{}1895) & \(u_+\): NA\par \(u_-\): NA\par Tipo: R\par Confianza: ND & COMP-BARYON\par E04 / M07\\\addlinespace[3pt]
% VI-CATALOG-ROW masas 3
\textbf{3. N(2060) BREIT-WIGNER MASS}\par N(2060)\allowbreak{}0 | Nbar(2060)\allowbreak{}0 | N(2060)\allowbreak{}+ | Nbar(2060)\allowbreak{}-\par {\fontsize{8.1}{9.4}\selectfont PDG2026:\allowbreak{}observable:\allowbreak{}B005:\allowbreak{}B005M:\allowbreak{}348836} & \textbf{2030 to 2100 to 2200} MeV\par Presentación: 2030 to 2200 (\textasciitilde{}2100) & \(u_+\): NA\par \(u_-\): NA\par Tipo: R\par Confianza: ND & COMP-BARYON\par E04 / M07\\\addlinespace[3pt]
% VI-CATALOG-ROW masas 4
\textbf{4. Lambda(2585) MASS (BUMPS)}\par Lambda(2585)\allowbreak{}0 | Lambdabar(2585)\allowbreak{}0\par {\fontsize{8.1}{9.4}\selectfont PDG2026:\allowbreak{}observable:\allowbreak{}B007:\allowbreak{}B007M:\allowbreak{}349780} & \textbf{\textasciitilde{}2585} MeV & \(u_+\): ND\par \(u_-\): ND\par Tipo: NONE\par Confianza: ND & COMP-BARYON\par E04 / M07\\\addlinespace[3pt]
% VI-CATALOG-ROW masas 5
\textbf{5. Delta(1700) BREIT-WIGNER MASS}\par Delta(1700)\allowbreak{}- | Deltabar(1700)\allowbreak{}+ | Delta(1700)\allowbreak{}0 | Deltabar(1700)\allowbreak{}0 | Delta(1700)\allowbreak{}+ | Deltabar(1700)\allowbreak{}- | Delta(1700)\allowbreak{}++ | Deltabar(1700)\allowbreak{}-\kern0pt{}-\par {\fontsize{8.1}{9.4}\selectfont PDG2026:\allowbreak{}observable:\allowbreak{}B010:\allowbreak{}B010M:\allowbreak{}349097} & \textbf{1690 to 1710 to 1730} MeV\par Presentación: 1690 to 1730 (\textasciitilde{}1710) & \(u_+\): NA\par \(u_-\): NA\par Tipo: R\par Confianza: ND & COMP-BARYON\par E04 / M07\\\addlinespace[3pt]
% VI-CATALOG-ROW masas 6
\textbf{6. Delta(1905) BREIT-WIGNER MASS}\par Delta(1905)\allowbreak{}- | Deltabar(1905)\allowbreak{}+ | Delta(1905)\allowbreak{}0 | Deltabar(1905)\allowbreak{}0 | Delta(1905)\allowbreak{}+ | Deltabar(1905)\allowbreak{}- | Delta(1905)\allowbreak{}++ | Deltabar(1905)\allowbreak{}-\kern0pt{}-\par {\fontsize{8.1}{9.4}\selectfont PDG2026:\allowbreak{}observable:\allowbreak{}B011:\allowbreak{}B011M:\allowbreak{}349151} & \textbf{1855 to 1880 to 1910} MeV\par Presentación: 1855 to 1910 (\textasciitilde{}1880) & \(u_+\): NA\par \(u_-\): NA\par Tipo: R\par Confianza: ND & COMP-BARYON\par E04 / M07\\\addlinespace[3pt]
% VI-CATALOG-ROW masas 7
\textbf{7. Delta(1910) BREIT-WIGNER MASS}\par Delta(1910)\allowbreak{}- | Deltabar(1910)\allowbreak{}+ | Delta(1910)\allowbreak{}0 | Deltabar(1910)\allowbreak{}0 | Delta(1910)\allowbreak{}+ | Deltabar(1910)\allowbreak{}- | Delta(1910)\allowbreak{}++ | Deltabar(1910)\allowbreak{}-\kern0pt{}-\par {\fontsize{8.1}{9.4}\selectfont PDG2026:\allowbreak{}observable:\allowbreak{}B012:\allowbreak{}B012M:\allowbreak{}349182} & \textbf{1850 to 1900 to 1950} MeV\par Presentación: 1850 to 1950 (\textasciitilde{}1900) & \(u_+\): NA\par \(u_-\): NA\par Tipo: R\par Confianza: ND & COMP-BARYON\par E04 / M07\\\addlinespace[3pt]
% VI-CATALOG-ROW masas 8
\textbf{8. Delta(1930) BREIT-WIGNER MASS}\par Delta(1930)\allowbreak{}- | Deltabar(1930)\allowbreak{}+ | Delta(1930)\allowbreak{}0 | Deltabar(1930)\allowbreak{}0 | Delta(1930)\allowbreak{}+ | Deltabar(1930)\allowbreak{}- | Delta(1930)\allowbreak{}++ | Deltabar(1930)\allowbreak{}-\kern0pt{}-\par {\fontsize{8.1}{9.4}\selectfont PDG2026:\allowbreak{}observable:\allowbreak{}B013:\allowbreak{}B013M:\allowbreak{}349249} & \textbf{1900 to 1950 to 2000} MeV\par Presentación: 1900 to 2000 (\textasciitilde{}1950) & \(u_+\): NA\par \(u_-\): NA\par Tipo: R\par Confianza: ND & COMP-BARYON\par E04 / M07\\\addlinespace[3pt]
% VI-CATALOG-ROW masas 9
\textbf{9. N(1710) BREIT-WIGNER MASS}\par N(1710)\allowbreak{}0 | Nbar(1710)\allowbreak{}0 | N(1710)\allowbreak{}+ | Nbar(1710)\allowbreak{}-\par {\fontsize{8.1}{9.4}\selectfont PDG2026:\allowbreak{}observable:\allowbreak{}B014:\allowbreak{}B014M:\allowbreak{}348535} & \textbf{1680 TO 1710 TO 1740} MeV\par Presentación: 1680 to 1740 (\textasciitilde{}1710) & \(u_+\): NA\par \(u_-\): NA\par Tipo: R\par Confianza: ND & COMP-BARYON\par E04 / M07\\\addlinespace[3pt]
% VI-CATALOG-ROW masas 10
\textbf{10. N(1720) BREIT-WIGNER MASS}\par N(1720)\allowbreak{}0 | Nbar(1720)\allowbreak{}0 | N(1720)\allowbreak{}+ | Nbar(1720)\allowbreak{}-\par {\fontsize{8.1}{9.4}\selectfont PDG2026:\allowbreak{}observable:\allowbreak{}B015:\allowbreak{}B015M:\allowbreak{}348569} & \textbf{1680 to 1720 to 1750} MeV\par Presentación: 1680 to 1750 (\textasciitilde{}1720) & \(u_+\): NA\par \(u_-\): NA\par Tipo: R\par Confianza: ND & COMP-BARYON\par E04 / M07\\\addlinespace[3pt]
% VI-CATALOG-ROW masas 11
\textbf{11. N(1875) BREIT-WIGNER MASS}\par N(1875)\allowbreak{}0 | Nbar(1875)\allowbreak{}0 | N(1875)\allowbreak{}+ | Nbar(1875)\allowbreak{}-\par {\fontsize{8.1}{9.4}\selectfont PDG2026:\allowbreak{}observable:\allowbreak{}B016:\allowbreak{}B016M:\allowbreak{}348633} & \textbf{1850 to 1875 to 1920} MeV\par Presentación: 1850 to 1920 (\textasciitilde{}1875) & \(u_+\): NA\par \(u_-\): NA\par Tipo: R\par Confianza: ND & COMP-BARYON\par E04 / M07\\\addlinespace[3pt]
% VI-CATALOG-ROW masas 12
\textbf{12. N(1990) BREIT-WIGNER MASS}\par N(1990)\allowbreak{}0 | Nbar(1990)\allowbreak{}0 | N(1990)\allowbreak{}+ | Nbar(1990)\allowbreak{}-\par {\fontsize{8.1}{9.4}\selectfont PDG2026:\allowbreak{}observable:\allowbreak{}B017:\allowbreak{}B017M:\allowbreak{}348778} & \textbf{1950 to 2020 to 2100} MeV\par Presentación: 1950 to 2100 (\textasciitilde{}2020) & \(u_+\): NA\par \(u_-\): NA\par Tipo: R\par Confianza: ND & COMP-BARYON\par E04 / M07\\\addlinespace[3pt]
% VI-CATALOG-ROW masas 13
\textbf{13. N(1700) BREIT-WIGNER MASS}\par N(1700)\allowbreak{}0 | Nbar(1700)\allowbreak{}0 | N(1700)\allowbreak{}+ | Nbar(1700)\allowbreak{}-\par {\fontsize{8.1}{9.4}\selectfont PDG2026:\allowbreak{}observable:\allowbreak{}B018:\allowbreak{}B018M:\allowbreak{}348499} & \textbf{1650 to 1720 to 1800} MeV\par Presentación: 1650 to 1800 (\textasciitilde{}1720) & \(u_+\): NA\par \(u_-\): NA\par Tipo: R\par Confianza: ND & COMP-BARYON\par E04 / M07\\\addlinespace[3pt]
% VI-CATALOG-ROW masas 14
\textbf{14. Delta(1600) BREIT-WIGNER MASS}\par Delta(1600)\allowbreak{}- | Deltabar(1600)\allowbreak{}+ | Delta(1600)\allowbreak{}0 | Deltabar(1600)\allowbreak{}0 | Delta(1600)\allowbreak{}+ | Deltabar(1600)\allowbreak{}- | Delta(1600)\allowbreak{}++ | Deltabar(1600)\allowbreak{}-\kern0pt{}-\par {\fontsize{8.1}{9.4}\selectfont PDG2026:\allowbreak{}observable:\allowbreak{}B019:\allowbreak{}B019M:\allowbreak{}349050} & \textbf{1500 to 1570 to 1640} MeV\par Presentación: 1500 to 1640 (\textasciitilde{}1570) & \(u_+\): NA\par \(u_-\): NA\par Tipo: R\par Confianza: ND & COMP-BARYON\par E04 / M07\\\addlinespace[3pt]
% VI-CATALOG-ROW masas 15
\textbf{15. Xi(1620) MASS}\par Xi(1620)\allowbreak{}- | Xibar(1620)\allowbreak{}+ | Xi(1620)\allowbreak{}0 | Xibar(1620)\allowbreak{}0\par {\fontsize{8.1}{9.4}\selectfont PDG2026:\allowbreak{}observable:\allowbreak{}B021:\allowbreak{}B021M:\allowbreak{}350410} & \textbf{\textasciitilde{}1620} MeV & \(u_+\): ND\par \(u_-\): ND\par Tipo: NONE\par Confianza: ND & COMP-BARYON\par E04 / M07\\\addlinespace[3pt]
% VI-CATALOG-ROW masas 16
\textbf{16. Xi(2250) MASS}\par Xi(2250)\allowbreak{}- | Xibar(2250)\allowbreak{}+ | Xi(2250)\allowbreak{}0 | Xibar(2250)\allowbreak{}0\par {\fontsize{8.1}{9.4}\selectfont PDG2026:\allowbreak{}observable:\allowbreak{}B022:\allowbreak{}B022M:\allowbreak{}350457} & \textbf{\textasciitilde{}2250} MeV & \(u_+\): ND\par \(u_-\): ND\par Tipo: NONE\par Confianza: ND & COMP-BARYON\par E04 / M07\\\addlinespace[3pt]
% VI-CATALOG-ROW masas 17
\textbf{17. N(2120) BREIT-WIGNER MASS}\par N(2120)\allowbreak{}0 | Nbar(2120)\allowbreak{}0 | N(2120)\allowbreak{}+ | Nbar(2120)\allowbreak{}-\par {\fontsize{8.1}{9.4}\selectfont PDG2026:\allowbreak{}observable:\allowbreak{}B024:\allowbreak{}B024M:\allowbreak{}348906} & \textbf{2060 to 2120 to 2160} MeV\par Presentación: 2060 to 2160 (\textasciitilde{}2120) & \(u_+\): NA\par \(u_-\): NA\par Tipo: R\par Confianza: ND & COMP-BARYON\par E04 / M07\\\addlinespace[3pt]
% VI-CATALOG-ROW masas 18
\textbf{18. Sigma(2100) MASS}\par Sigma(2100)\allowbreak{}- | Sigmabar(2100)\allowbreak{}+ | Sigma(2100)\allowbreak{}0 | Sigmabar(2100)\allowbreak{}0 | Sigma(2100)\allowbreak{}+ | Sigmabar(2100)\allowbreak{}-\par {\fontsize{8.1}{9.4}\selectfont PDG2026:\allowbreak{}observable:\allowbreak{}B026:\allowbreak{}B026M:\allowbreak{}350202} & \textbf{\textasciitilde{}2100} MeV & \(u_+\): ND\par \(u_-\): ND\par Tipo: NONE\par Confianza: ND & COMP-BARYON\par E04 / M07\\\addlinespace[3pt]
% VI-CATALOG-ROW masas 19
\textbf{19. Lambda(2085) MASS}\par Lambda(2085)\allowbreak{}0 | Lambdabar(2085)\allowbreak{}0\par {\fontsize{8.1}{9.4}\selectfont PDG2026:\allowbreak{}observable:\allowbreak{}B027:\allowbreak{}B027M:\allowbreak{}349731} & \textbf{\textasciitilde{}2020} MeV & \(u_+\): ND\par \(u_-\): ND\par Tipo: NONE\par Confianza: ND & COMP-BARYON\par E04 / M07\\\addlinespace[3pt]
% VI-CATALOG-ROW masas 20
\textbf{20. Delta(1900) BREIT-WIGNER MASS}\par Delta(1900)\allowbreak{}- | Deltabar(1900)\allowbreak{}+ | Delta(1900)\allowbreak{}0 | Deltabar(1900)\allowbreak{}0 | Delta(1900)\allowbreak{}+ | Deltabar(1900)\allowbreak{}- | Delta(1900)\allowbreak{}++ | Deltabar(1900)\allowbreak{}-\kern0pt{}-\par {\fontsize{8.1}{9.4}\selectfont PDG2026:\allowbreak{}observable:\allowbreak{}B030:\allowbreak{}B030M:\allowbreak{}349126} & \textbf{1840 TO 1860 TO 1920} MeV\par Presentación: 1840 to 1920 (\textasciitilde{}1860) & \(u_+\): NA\par \(u_-\): NA\par Tipo: R\par Confianza: ND & COMP-BARYON\par E04 / M07\\\addlinespace[3pt]
% VI-CATALOG-ROW masas 21
\textbf{21. Sigma(1620) MASS}\par Sigma(1620)\allowbreak{}- | Sigmabar(1620)\allowbreak{}+ | Sigma(1620)\allowbreak{}0 | Sigmabar(1620)\allowbreak{}0 | Sigma(1620)\allowbreak{}+ | Sigmabar(1620)\allowbreak{}-\par {\fontsize{8.1}{9.4}\selectfont PDG2026:\allowbreak{}observable:\allowbreak{}B032:\allowbreak{}B032M:\allowbreak{}349902} & \textbf{1600 to 1620 to 1650} MeV\par Presentación: 1600 to 1650 (\textasciitilde{}1620) & \(u_+\): NA\par \(u_-\): NA\par Tipo: R\par Confianza: ND & COMP-BARYON\par E04 / M07\\\addlinespace[3pt]
% VI-CATALOG-ROW masas 22
\textbf{22. MIXED CHARGES}\par Delta(1232)\allowbreak{}- | Deltabar(1232)\allowbreak{}+ | Delta(1232)\allowbreak{}0 | Deltabar(1232)\allowbreak{}0 | Delta(1232)\allowbreak{}+ | Deltabar(1232)\allowbreak{}- | Delta(1232)\allowbreak{}++ | Deltabar(1232)\allowbreak{}-\kern0pt{}-\par {\fontsize{8.1}{9.4}\selectfont PDG2026:\allowbreak{}observable:\allowbreak{}B033:\allowbreak{}B033M:\allowbreak{}349034} & \textbf{1230 TO 1232 TO 1234} MeV\par Presentación: 1230 to 1234 (\textasciitilde{}1232) & \(u_+\): NA\par \(u_-\): NA\par Tipo: R\par Confianza: ND & COMP-BARYON\par E04 / M07\\\addlinespace[3pt]
% VI-CATALOG-ROW masas 23
\textbf{23. Sigma(2070) MASS}\par Sigma(2070)\allowbreak{}- | Sigmabar(2070)\allowbreak{}+ | Sigma(2070)\allowbreak{}0 | Sigmabar(2070)\allowbreak{}0 | Sigma(2070)\allowbreak{}+ | Sigmabar(2070)\allowbreak{}-\par {\fontsize{8.1}{9.4}\selectfont PDG2026:\allowbreak{}observable:\allowbreak{}B034:\allowbreak{}B034M:\allowbreak{}350186} & \textbf{2020 to 2060 to 2100} MeV\par Presentación: 2020 to 2100 (\textasciitilde{}2060) & \(u_+\): NA\par \(u_-\): NA\par Tipo: R\par Confianza: ND & COMP-BARYON\par E04 / M07\\\addlinespace[3pt]
% VI-CATALOG-ROW masas 24
\textbf{24. Lambda(2110) MASS}\par Lambda(2110)\allowbreak{}0 | Lambdabar(2110)\allowbreak{}0\par {\fontsize{8.1}{9.4}\selectfont PDG2026:\allowbreak{}observable:\allowbreak{}B035:\allowbreak{}B035M:\allowbreak{}349763} & \textbf{2050 to 2090 to 2130} MeV\par Presentación: 2050 to 2130 (\textasciitilde{}2090) & \(u_+\): NA\par \(u_-\): NA\par Tipo: R\par Confianza: ND & COMP-BARYON\par E04 / M07\\\addlinespace[3pt]
% VI-CATALOG-ROW masas 25
\textbf{25. Lambda(1800) MASS}\par Lambda(1800)\allowbreak{}0 | Lambdabar(1800)\allowbreak{}0\par {\fontsize{8.1}{9.4}\selectfont PDG2026:\allowbreak{}observable:\allowbreak{}B036:\allowbreak{}B036M:\allowbreak{}349555} & \textbf{1750 to 1800 to 1850} MeV\par Presentación: 1750 to 1850 (\textasciitilde{}1800) & \(u_+\): NA\par \(u_-\): NA\par Tipo: R\par Confianza: ND & COMP-BARYON\par E04 / M07\\\addlinespace[3pt]
% VI-CATALOG-ROW masas 26
\textbf{26. PRODUCTION EXPERIMENTS}\par Lambda(1405)\allowbreak{}0 | Lambdabar(1405)\allowbreak{}0\par {\fontsize{8.1}{9.4}\selectfont PDG2026:\allowbreak{}observable:\allowbreak{}B037:\allowbreak{}B037M1:\allowbreak{}349427} & \textbf{1405.1+1.3-1.0} MeV & \(u_+\): 1.311969\allowbreak{}35395758\allowbreak{}21\par \(u_-\): 0.964554\allowbreak{}87987763\allowbreak{}788\par Tipo: NONE\par Confianza: ND & COMP-BARYON\par E04 / M07\\\addlinespace[3pt]
% VI-CATALOG-ROW masas 27
\textbf{27. Lambda(1520) MASS}\par Lambda(1520)\allowbreak{}0 | Lambdabar(1520)\allowbreak{}0\par {\fontsize{8.1}{9.4}\selectfont PDG2026:\allowbreak{}observable:\allowbreak{}B038:\allowbreak{}B038M:\allowbreak{}349437} & \textbf{1518 to 1519 to 1520} MeV\par Presentación: 1518 to 1520 (\textasciitilde{}1519) & \(u_+\): NA\par \(u_-\): NA\par Tipo: R\par Confianza: ND & COMP-BARYON\par E04 / M07\\\addlinespace[3pt]
% VI-CATALOG-ROW masas 28
\textbf{28. Lambda(1820) MASS}\par Lambda(1820)\allowbreak{}0 | Lambdabar(1820)\allowbreak{}0\par {\fontsize{8.1}{9.4}\selectfont PDG2026:\allowbreak{}observable:\allowbreak{}B039:\allowbreak{}B039M:\allowbreak{}349594} & \textbf{1815 TO 1820 TO 1825} MeV\par Presentación: 1815 to 1825 (\textasciitilde{}1820) & \(u_+\): NA\par \(u_-\): NA\par Tipo: R\par Confianza: ND & COMP-BARYON\par E04 / M07\\\addlinespace[3pt]
% VI-CATALOG-ROW masas 29
\textbf{29. Lambda(1670) MASS}\par Lambda(1670)\allowbreak{}0 | Lambdabar(1670)\allowbreak{}0\par {\fontsize{8.1}{9.4}\selectfont PDG2026:\allowbreak{}observable:\allowbreak{}B040:\allowbreak{}B040M:\allowbreak{}349494} & \textbf{1670 to 1674 to 1678} MeV\par Presentación: 1670 to 1678 (\textasciitilde{}1674) & \(u_+\): NA\par \(u_-\): NA\par Tipo: R\par Confianza: ND & COMP-BARYON\par E04 / M07\\\addlinespace[3pt]
% VI-CATALOG-ROW masas 30
\textbf{30. Lambda(2100) MASS}\par Lambda(2100)\allowbreak{}0 | Lambdabar(2100)\allowbreak{}0\par {\fontsize{8.1}{9.4}\selectfont PDG2026:\allowbreak{}observable:\allowbreak{}B041:\allowbreak{}B041M:\allowbreak{}349741} & \textbf{2090 TO 2100 TO 2110} MeV\par Presentación: 2090 to 2110 (\textasciitilde{}2100) & \(u_+\): NA\par \(u_-\): NA\par Tipo: R\par Confianza: ND & COMP-BARYON\par E04 / M07\\\addlinespace[3pt]
% VI-CATALOG-ROW masas 31
\textbf{31. Lambda(2350) MASS}\par Lambda(2350)\allowbreak{}0 | Lambdabar(2350)\allowbreak{}0\par {\fontsize{8.1}{9.4}\selectfont PDG2026:\allowbreak{}observable:\allowbreak{}B042:\allowbreak{}B042M:\allowbreak{}349775} & \textbf{2340 TO 2350 TO 2370} MeV\par Presentación: 2340 to 2370 (\textasciitilde{}2350) & \(u_+\): NA\par \(u_-\): NA\par Tipo: R\par Confianza: ND & COMP-BARYON\par E04 / M07\\\addlinespace[3pt]
% VI-CATALOG-ROW masas 32
\textbf{32. Sigma(1385)\allowbreak{}+ MASS}\par Sigma(1385)\allowbreak{}- | Sigmabar(1385)\allowbreak{}+ | Sigma(1385)\allowbreak{}0 | Sigmabar(1385)\allowbreak{}0 | Sigma(1385)\allowbreak{}+ | Sigmabar(1385)\allowbreak{}-\par {\fontsize{8.1}{9.4}\selectfont PDG2026:\allowbreak{}observable:\allowbreak{}B043:\allowbreak{}B043M+:\allowbreak{}349876} & \textbf{1382.83+-0.34} MeV & \(u_+\): 0.339054\allowbreak{}35517601\allowbreak{}018\par \(u_-\): 0.344190\allowbreak{}99315609\allowbreak{}542\par Tipo: NONE\par Confianza: ND & COMP-BARYON\par E04 / M07\\\addlinespace[3pt]
% VI-CATALOG-ROW masas 33
\textbf{33. Sigma(1385)\allowbreak{}0 MASS}\par Sigma(1385)\allowbreak{}- | Sigmabar(1385)\allowbreak{}+ | Sigma(1385)\allowbreak{}0 | Sigmabar(1385)\allowbreak{}0 | Sigma(1385)\allowbreak{}+ | Sigmabar(1385)\allowbreak{}-\par {\fontsize{8.1}{9.4}\selectfont PDG2026:\allowbreak{}observable:\allowbreak{}B043:\allowbreak{}B043M0:\allowbreak{}349877} & \textbf{1383.8+-0.9} MeV & \(u_+\): 0.935144\allowbreak{}04403237\allowbreak{}228\par \(u_-\): 0.935144\allowbreak{}04403237\allowbreak{}228\par Tipo: NONE\par Confianza: ND & COMP-BARYON\par E04 / M07\\\addlinespace[3pt]
% VI-CATALOG-ROW masas 34
\textbf{34. Sigma(1385)\allowbreak{}- MASS}\par Sigma(1385)\allowbreak{}- | Sigmabar(1385)\allowbreak{}+ | Sigma(1385)\allowbreak{}0 | Sigmabar(1385)\allowbreak{}0 | Sigma(1385)\allowbreak{}+ | Sigmabar(1385)\allowbreak{}-\par {\fontsize{8.1}{9.4}\selectfont PDG2026:\allowbreak{}observable:\allowbreak{}B043:\allowbreak{}B043M-:\allowbreak{}349878} & \textbf{1387.2+-0.5} MeV & \(u_+\): 0.546537\allowbreak{}03280513\allowbreak{}61\par \(u_-\): 0.546537\allowbreak{}03280513\allowbreak{}61\par Tipo: NONE\par Confianza: ND & COMP-BARYON\par E04 / M07\\\addlinespace[3pt]
% VI-CATALOG-ROW masas 35
\textbf{35. Sigma(1670) MASS}\par Sigma(1670)\allowbreak{}- | Sigmabar(1670)\allowbreak{}+ | Sigma(1670)\allowbreak{}0 | Sigmabar(1670)\allowbreak{}0 | Sigma(1670)\allowbreak{}+ | Sigmabar(1670)\allowbreak{}-\par {\fontsize{8.1}{9.4}\selectfont PDG2026:\allowbreak{}observable:\allowbreak{}B044:\allowbreak{}B044M:\allowbreak{}349946} & \textbf{1665 to 1675 to 1685} MeV\par Presentación: 1665 to 1685 (\textasciitilde{}1675) & \(u_+\): NA\par \(u_-\): NA\par Tipo: R\par Confianza: ND & COMP-BARYON\par E04 / M07\\\addlinespace[3pt]
% VI-CATALOG-ROW masas 36
\textbf{36. Sigma(1775) MASS}\par Sigma(1775)\allowbreak{}- | Sigmabar(1775)\allowbreak{}+ | Sigma(1775)\allowbreak{}0 | Sigmabar(1775)\allowbreak{}0 | Sigma(1775)\allowbreak{}+ | Sigmabar(1775)\allowbreak{}-\par {\fontsize{8.1}{9.4}\selectfont PDG2026:\allowbreak{}observable:\allowbreak{}B045:\allowbreak{}B045M:\allowbreak{}349983} & \textbf{1770 TO 1775 TO 1780} MeV\par Presentación: 1770 to 1780 (\textasciitilde{}1775) & \(u_+\): NA\par \(u_-\): NA\par Tipo: R\par Confianza: ND & COMP-BARYON\par E04 / M07\\\addlinespace[3pt]
% VI-CATALOG-ROW masas 37
\textbf{37. Sigma(1915) MASS}\par Sigma(1915)\allowbreak{}- | Sigmabar(1915)\allowbreak{}+ | Sigma(1915)\allowbreak{}0 | Sigmabar(1915)\allowbreak{}0 | Sigma(1915)\allowbreak{}+ | Sigmabar(1915)\allowbreak{}-\par {\fontsize{8.1}{9.4}\selectfont PDG2026:\allowbreak{}observable:\allowbreak{}B046:\allowbreak{}B046M:\allowbreak{}350080} & \textbf{1900 TO 1915 TO 1935} MeV\par Presentación: 1900 to 1935 (\textasciitilde{}1915) & \(u_+\): NA\par \(u_-\): NA\par Tipo: R\par Confianza: ND & COMP-BARYON\par E04 / M07\\\addlinespace[3pt]
% VI-CATALOG-ROW masas 38
\textbf{38. Sigma(2030) MASS}\par Sigma(2030)\allowbreak{}- | Sigmabar(2030)\allowbreak{}+ | Sigma(2030)\allowbreak{}0 | Sigmabar(2030)\allowbreak{}0 | Sigma(2030)\allowbreak{}+ | Sigmabar(2030)\allowbreak{}-\par {\fontsize{8.1}{9.4}\selectfont PDG2026:\allowbreak{}observable:\allowbreak{}B047:\allowbreak{}B047M:\allowbreak{}350166} & \textbf{2025 TO 2030 TO 2040} MeV\par Presentación: 2025 to 2040 (\textasciitilde{}2030) & \(u_+\): NA\par \(u_-\): NA\par Tipo: R\par Confianza: ND & COMP-BARYON\par E04 / M07\\\addlinespace[3pt]
% VI-CATALOG-ROW masas 39
\textbf{39. Sigma(2250) MASS}\par Sigma(2250)\allowbreak{}- | Sigmabar(2250)\allowbreak{}+ | Sigma(2250)\allowbreak{}0 | Sigmabar(2250)\allowbreak{}0 | Sigma(2250)\allowbreak{}+ | Sigmabar(2250)\allowbreak{}-\par {\fontsize{8.1}{9.4}\selectfont PDG2026:\allowbreak{}observable:\allowbreak{}B048:\allowbreak{}B048M:\allowbreak{}350277} & \textbf{2210 TO 2250 TO 2280} MeV\par Presentación: 2210 to 2280 (\textasciitilde{}2250) & \(u_+\): NA\par \(u_-\): NA\par Tipo: R\par Confianza: ND & COMP-BARYON\par E04 / M07\\\addlinespace[3pt]
% VI-CATALOG-ROW masas 40
\textbf{40. Xi(1530)\allowbreak{}0 MASS}\par Xi(1530)\allowbreak{}- | Xibar(1530)\allowbreak{}+ | Xi(1530)\allowbreak{}0 | Xibar(1530)\allowbreak{}0\par {\fontsize{8.1}{9.4}\selectfont PDG2026:\allowbreak{}observable:\allowbreak{}B049:\allowbreak{}B049M0:\allowbreak{}350399} & \textbf{1531.80+-0.32} MeV & \(u_+\): 0.315752\allowbreak{}72391181\allowbreak{}958\par \(u_-\): 0.315752\allowbreak{}72391181\allowbreak{}958\par Tipo: NONE\par Confianza: ND & COMP-BARYON\par E04 / M07\\\addlinespace[3pt]
% VI-CATALOG-ROW masas 41
\textbf{41. Xi(1530)\allowbreak{}- MASS}\par Xi(1530)\allowbreak{}- | Xibar(1530)\allowbreak{}+ | Xi(1530)\allowbreak{}0 | Xibar(1530)\allowbreak{}0\par {\fontsize{8.1}{9.4}\selectfont PDG2026:\allowbreak{}observable:\allowbreak{}B049:\allowbreak{}B049M-:\allowbreak{}350401} & \textbf{1535.0+-0.6} MeV & \(u_+\): 0.632021\allowbreak{}66316087\allowbreak{}116\par \(u_-\): 0.632021\allowbreak{}66316087\allowbreak{}116\par Tipo: NONE\par Confianza: ND & COMP-BARYON\par E04 / M07\\\addlinespace[3pt]
% VI-CATALOG-ROW masas 42
\textbf{42. Xi(1820) MASS}\par Xi(1820)\allowbreak{}- | Xibar(1820)\allowbreak{}+ | Xi(1820)\allowbreak{}0 | Xibar(1820)\allowbreak{}0\par {\fontsize{8.1}{9.4}\selectfont PDG2026:\allowbreak{}observable:\allowbreak{}B050:\allowbreak{}B050M:\allowbreak{}350420} & \textbf{1823+-5} MeV & \(u_+\): 5\par \(u_-\): 5\par Tipo: NONE\par Confianza: ND & COMP-BARYON\par E04 / M07\\\addlinespace[3pt]
% VI-CATALOG-ROW masas 43
\textbf{43. Xi(1950) MASS}\par Xi(1950)\allowbreak{}- | Xibar(1950)\allowbreak{}+ | Xi(1950)\allowbreak{}0 | Xibar(1950)\allowbreak{}0\par {\fontsize{8.1}{9.4}\selectfont PDG2026:\allowbreak{}observable:\allowbreak{}B052:\allowbreak{}B052M:\allowbreak{}350437} & \textbf{1950+-15} MeV & \(u_+\): 15\par \(u_-\): 15\par Tipo: NONE\par Confianza: ND & COMP-BARYON\par E04 / M07\\\addlinespace[3pt]
% VI-CATALOG-ROW masas 44
\textbf{44. Sigma(2455) MASS}\par Sigma(2455)\allowbreak{}- | Sigmabar(2455)\allowbreak{}+ | Sigma(2455)\allowbreak{}0 | Sigmabar(2455)\allowbreak{}0 | Sigma(2455)\allowbreak{}+ | Sigmabar(2455)\allowbreak{}-\par {\fontsize{8.1}{9.4}\selectfont PDG2026:\allowbreak{}observable:\allowbreak{}B053:\allowbreak{}B053M:\allowbreak{}350283} & \textbf{\textasciitilde{}2455} MeV & \(u_+\): ND\par \(u_-\): ND\par Tipo: NONE\par Confianza: ND & COMP-BARYON\par E04 / M07\\\addlinespace[3pt]
% VI-CATALOG-ROW masas 45
\textbf{45. Sigma(2620) MASS}\par Sigma(2620)\allowbreak{}- | Sigmabar(2620)\allowbreak{}+ | Sigma(2620)\allowbreak{}0 | Sigmabar(2620)\allowbreak{}0 | Sigma(2620)\allowbreak{}+ | Sigmabar(2620)\allowbreak{}-\par {\fontsize{8.1}{9.4}\selectfont PDG2026:\allowbreak{}observable:\allowbreak{}B054:\allowbreak{}B054M:\allowbreak{}350284} & \textbf{\textasciitilde{}2620} MeV & \(u_+\): ND\par \(u_-\): ND\par Tipo: NONE\par Confianza: ND & COMP-BARYON\par E04 / M07\\\addlinespace[3pt]
% VI-CATALOG-ROW masas 46
\textbf{46. Lambda(1690) MASS}\par Lambda(1690)\allowbreak{}0 | Lambdabar(1690)\allowbreak{}0\par {\fontsize{8.1}{9.4}\selectfont PDG2026:\allowbreak{}observable:\allowbreak{}B055:\allowbreak{}B055M:\allowbreak{}349518} & \textbf{1685 to 1690 to 1695} MeV\par Presentación: 1685 to 1695 (\textasciitilde{}1690) & \(u_+\): NA\par \(u_-\): NA\par Tipo: R\par Confianza: ND & COMP-BARYON\par E04 / M07\\\addlinespace[3pt]
% VI-CATALOG-ROW masas 47
\textbf{47. Lambda(1830) MASS}\par Lambda(1830)\allowbreak{}0 | Lambdabar(1830)\allowbreak{}0\par {\fontsize{8.1}{9.4}\selectfont PDG2026:\allowbreak{}observable:\allowbreak{}B056:\allowbreak{}B056M:\allowbreak{}349614} & \textbf{1820 to 1825 to 1830} MeV\par Presentación: 1820 to 1830 (\textasciitilde{}1825) & \(u_+\): NA\par \(u_-\): NA\par Tipo: R\par Confianza: ND & COMP-BARYON\par E04 / M07\\\addlinespace[3pt]
% VI-CATALOG-ROW masas 48
\textbf{48. Sigma(1750) MASS}\par Sigma(1750)\allowbreak{}- | Sigmabar(1750)\allowbreak{}+ | Sigma(1750)\allowbreak{}0 | Sigmabar(1750)\allowbreak{}0 | Sigma(1750)\allowbreak{}+ | Sigmabar(1750)\allowbreak{}-\par {\fontsize{8.1}{9.4}\selectfont PDG2026:\allowbreak{}observable:\allowbreak{}B057:\allowbreak{}B057M:\allowbreak{}349966} & \textbf{1700 to 1750 to 1800} MeV\par Presentación: 1700 to 1800 (\textasciitilde{}1750) & \(u_+\): NA\par \(u_-\): NA\par Tipo: R\par Confianza: ND & COMP-BARYON\par E04 / M07\\\addlinespace[3pt]
% VI-CATALOG-ROW masas 49
\textbf{49. Sigma(3000) MASS}\par Sigma(3000)\allowbreak{}- | Sigmabar(3000)\allowbreak{}+ | Sigma(3000)\allowbreak{}0 | Sigmabar(3000)\allowbreak{}0 | Sigma(3000)\allowbreak{}+ | Sigmabar(3000)\allowbreak{}-\par {\fontsize{8.1}{9.4}\selectfont PDG2026:\allowbreak{}observable:\allowbreak{}B059:\allowbreak{}B059M:\allowbreak{}350285} & \textbf{\textasciitilde{}3000} MeV & \(u_+\): ND\par \(u_-\): ND\par Tipo: NONE\par Confianza: ND & COMP-BARYON\par E04 / M07\\\addlinespace[3pt]
% VI-CATALOG-ROW masas 50
\textbf{50. Lambda(1890) MASS}\par Lambda(1890)\allowbreak{}0 | Lambdabar(1890)\allowbreak{}0\par {\fontsize{8.1}{9.4}\selectfont PDG2026:\allowbreak{}observable:\allowbreak{}B060:\allowbreak{}B060M:\allowbreak{}349635} & \textbf{1870 to 1890 to 1910} MeV\par Presentación: 1870 to 1910 (\textasciitilde{}1890) & \(u_+\): NA\par \(u_-\): NA\par Tipo: R\par Confianza: ND & COMP-BARYON\par E04 / M07\\\addlinespace[3pt]
% VI-CATALOG-ROW masas 51
\textbf{51. N(1440) BREIT-WIGNER MASS}\par N(1440)\allowbreak{}0 | Nbar(1440)\allowbreak{}0 | N(1440)\allowbreak{}+ | Nbar(1440)\allowbreak{}-\par {\fontsize{8.1}{9.4}\selectfont PDG2026:\allowbreak{}observable:\allowbreak{}B061:\allowbreak{}B061M:\allowbreak{}348326} & \textbf{1410 to 1440 to 1470} MeV\par Presentación: 1410 to 1470 (\textasciitilde{}1440) & \(u_+\): NA\par \(u_-\): NA\par Tipo: R\par Confianza: ND & COMP-BARYON\par E04 / M07\\\addlinespace[3pt]
% VI-CATALOG-ROW masas 52
\textbf{52. N(1520) BREIT-WIGNER MASS}\par N(1520)\allowbreak{}0 | Nbar(1520)\allowbreak{}0 | N(1520)\allowbreak{}+ | Nbar(1520)\allowbreak{}-\par {\fontsize{8.1}{9.4}\selectfont PDG2026:\allowbreak{}observable:\allowbreak{}B062:\allowbreak{}B062M:\allowbreak{}348349} & \textbf{1510 TO 1515 TO 1520} MeV\par Presentación: 1510 to 1520 (\textasciitilde{}1515) & \(u_+\): NA\par \(u_-\): NA\par Tipo: R\par Confianza: ND & COMP-BARYON\par E04 / M07\\\addlinespace[3pt]
% VI-CATALOG-ROW masas 53
\textbf{53. N(1535) BREIT-WIGNER MASS}\par N(1535)\allowbreak{}0 | Nbar(1535)\allowbreak{}0 | N(1535)\allowbreak{}+ | Nbar(1535)\allowbreak{}-\par {\fontsize{8.1}{9.4}\selectfont PDG2026:\allowbreak{}observable:\allowbreak{}B063:\allowbreak{}B063M:\allowbreak{}348383} & \textbf{1515 to 1530 to 1545} MeV\par Presentación: 1515 to 1545 (\textasciitilde{}1530) & \(u_+\): NA\par \(u_-\): NA\par Tipo: R\par Confianza: ND & COMP-BARYON\par E04 / M07\\\addlinespace[3pt]
% VI-CATALOG-ROW masas 54
\textbf{54. N(1675) BREIT-WIGNER MASS}\par N(1675)\allowbreak{}0 | Nbar(1675)\allowbreak{}0 | N(1675)\allowbreak{}+ | Nbar(1675)\allowbreak{}-\par {\fontsize{8.1}{9.4}\selectfont PDG2026:\allowbreak{}observable:\allowbreak{}B064:\allowbreak{}B064M:\allowbreak{}348435} & \textbf{1665 to 1675 to 1680} MeV\par Presentación: 1665 to 1680 (\textasciitilde{}1675) & \(u_+\): NA\par \(u_-\): NA\par Tipo: R\par Confianza: ND & COMP-BARYON\par E04 / M07\\\addlinespace[3pt]
% VI-CATALOG-ROW masas 55
\textbf{55. N(1680) BREIT-WIGNER MASS}\par N(1680)\allowbreak{}0 | Nbar(1680)\allowbreak{}0 | N(1680)\allowbreak{}+ | Nbar(1680)\allowbreak{}-\par {\fontsize{8.1}{9.4}\selectfont PDG2026:\allowbreak{}observable:\allowbreak{}B065:\allowbreak{}B065M:\allowbreak{}348467} & \textbf{1680 TO 1685 TO 1690} MeV\par Presentación: 1680 to 1690 (\textasciitilde{}1685) & \(u_+\): NA\par \(u_-\): NA\par Tipo: R\par Confianza: ND & COMP-BARYON\par E04 / M07\\\addlinespace[3pt]
% VI-CATALOG-ROW masas 56
\textbf{56. N(1650) BREIT-WIGNER MASS}\par N(1650)\allowbreak{}0 | Nbar(1650)\allowbreak{}0 | N(1650)\allowbreak{}+ | Nbar(1650)\allowbreak{}-\par {\fontsize{8.1}{9.4}\selectfont PDG2026:\allowbreak{}observable:\allowbreak{}B066:\allowbreak{}B066M:\allowbreak{}348410} & \textbf{1635 to 1650 to 1665} MeV\par Presentación: 1635 to 1665 (\textasciitilde{}1650) & \(u_+\): NA\par \(u_-\): NA\par Tipo: R\par Confianza: ND & COMP-BARYON\par E04 / M07\\\addlinespace[3pt]
% VI-CATALOG-ROW masas 57
\textbf{57. Sigma(1880) MASS}\par Sigma(1880)\allowbreak{}- | Sigmabar(1880)\allowbreak{}+ | Sigma(1880)\allowbreak{}0 | Sigmabar(1880)\allowbreak{}0 | Sigma(1880)\allowbreak{}+ | Sigmabar(1880)\allowbreak{}-\par {\fontsize{8.1}{9.4}\selectfont PDG2026:\allowbreak{}observable:\allowbreak{}B067:\allowbreak{}B067M:\allowbreak{}349999} & \textbf{1820 to 1880 to 1940} MeV\par Presentación: 1820 to 1940 (\textasciitilde{}1880) & \(u_+\): NA\par \(u_-\): NA\par Tipo: R\par Confianza: ND & COMP-BARYON\par E04 / M07\\\addlinespace[3pt]
% VI-CATALOG-ROW masas 58
\textbf{58. Xi(2030) MASS}\par Xi(2030)\allowbreak{}- | Xibar(2030)\allowbreak{}+ | Xi(2030)\allowbreak{}0 | Xibar(2030)\allowbreak{}0\par {\fontsize{8.1}{9.4}\selectfont PDG2026:\allowbreak{}observable:\allowbreak{}B068:\allowbreak{}B068M:\allowbreak{}350443} & \textbf{2025+-5} MeV & \(u_+\): 5\par \(u_-\): 5\par Tipo: NONE\par Confianza: ND & COMP-BARYON\par E04 / M07\\\addlinespace[3pt]
% VI-CATALOG-ROW masas 59
\textbf{59. N(2190) BREIT-WIGNER MASS}\par N(2190)\allowbreak{}0 | Nbar(2190)\allowbreak{}0 | N(2190)\allowbreak{}+ | Nbar(2190)\allowbreak{}-\par {\fontsize{8.1}{9.4}\selectfont PDG2026:\allowbreak{}observable:\allowbreak{}B071:\allowbreak{}B071M:\allowbreak{}348953} & \textbf{2140 to 2180 to 2220} MeV\par Presentación: 2140 to 2220 (\textasciitilde{}2180) & \(u_+\): NA\par \(u_-\): NA\par Tipo: R\par Confianza: ND & COMP-BARYON\par E04 / M07\\\addlinespace[3pt]
% VI-CATALOG-ROW masas 60
\textbf{60. Lambda(1810) MASS}\par Lambda(1810)\allowbreak{}0 | Lambdabar(1810)\allowbreak{}0\par {\fontsize{8.1}{9.4}\selectfont PDG2026:\allowbreak{}observable:\allowbreak{}B077:\allowbreak{}B077M:\allowbreak{}349576} & \textbf{1740 to 1790 to 1840} MeV\par Presentación: 1740 to 1840 (\textasciitilde{}1790) & \(u_+\): NA\par \(u_-\): NA\par Tipo: R\par Confianza: ND & COMP-BARYON\par E04 / M07\\\addlinespace[3pt]
% VI-CATALOG-ROW masas 61
\textbf{61. Sigma(1660) MASS}\par Sigma(1660)\allowbreak{}- | Sigmabar(1660)\allowbreak{}+ | Sigma(1660)\allowbreak{}0 | Sigmabar(1660)\allowbreak{}0 | Sigma(1660)\allowbreak{}+ | Sigmabar(1660)\allowbreak{}-\par {\fontsize{8.1}{9.4}\selectfont PDG2026:\allowbreak{}observable:\allowbreak{}B079:\allowbreak{}B079M:\allowbreak{}349920} & \textbf{1640 to 1660 to 1680} MeV\par Presentación: 1640 to 1680 (\textasciitilde{}1660) & \(u_+\): NA\par \(u_-\): NA\par Tipo: R\par Confianza: ND & COMP-BARYON\par E04 / M07\\\addlinespace[3pt]
% VI-CATALOG-ROW masas 62
\textbf{62. Delta(1620) BREIT-WIGNER MASS}\par Delta(1620)\allowbreak{}- | Deltabar(1620)\allowbreak{}+ | Delta(1620)\allowbreak{}0 | Deltabar(1620)\allowbreak{}0 | Delta(1620)\allowbreak{}+ | Deltabar(1620)\allowbreak{}- | Delta(1620)\allowbreak{}++ | Deltabar(1620)\allowbreak{}-\kern0pt{}-\par {\fontsize{8.1}{9.4}\selectfont PDG2026:\allowbreak{}observable:\allowbreak{}B082:\allowbreak{}B082M:\allowbreak{}349078} & \textbf{1590 to 1610 to 1630} MeV\par Presentación: 1590 to 1630 (\textasciitilde{}1610) & \(u_+\): NA\par \(u_-\): NA\par Tipo: R\par Confianza: ND & COMP-BARYON\par E04 / M07\\\addlinespace[3pt]
% VI-CATALOG-ROW masas 63
\textbf{63. Delta(1950) BREIT-WIGNER MASS}\par Delta(1950)\allowbreak{}- | Deltabar(1950)\allowbreak{}+ | Delta(1950)\allowbreak{}0 | Deltabar(1950)\allowbreak{}0 | Delta(1950)\allowbreak{}+ | Deltabar(1950)\allowbreak{}- | Delta(1950)\allowbreak{}++ | Deltabar(1950)\allowbreak{}-\kern0pt{}-\par {\fontsize{8.1}{9.4}\selectfont PDG2026:\allowbreak{}observable:\allowbreak{}B083:\allowbreak{}B083M:\allowbreak{}349308} & \textbf{1915 to 1930 to 1950} MeV\par Presentación: 1915 to 1950 (\textasciitilde{}1930) & \(u_+\): NA\par \(u_-\): NA\par Tipo: R\par Confianza: ND & COMP-BARYON\par E04 / M07\\\addlinespace[3pt]
% VI-CATALOG-ROW masas 64
\textbf{64. Delta(2420) BREIT-WIGNER MASS}\par Delta(2420)\allowbreak{}- | Deltabar(2420)\allowbreak{}+ | Delta(2420)\allowbreak{}0 | Deltabar(2420)\allowbreak{}0 | Delta(2420)\allowbreak{}+ | Deltabar(2420)\allowbreak{}- | Delta(2420)\allowbreak{}++ | Deltabar(2420)\allowbreak{}-\kern0pt{}-\par {\fontsize{8.1}{9.4}\selectfont PDG2026:\allowbreak{}observable:\allowbreak{}B084:\allowbreak{}B084M:\allowbreak{}349359} & \textbf{2300 to 2450 to 2600} MeV\par Presentación: 2300 to 2600 (\textasciitilde{}2450) & \(u_+\): NA\par \(u_-\): NA\par Tipo: R\par Confianza: ND & COMP-BARYON\par E04 / M07\\\addlinespace[3pt]
% VI-CATALOG-ROW masas 65
\textbf{65. N(1880) BREIT-WIGNER MASS}\par N(1880)\allowbreak{}0 | Nbar(1880)\allowbreak{}0 | N(1880)\allowbreak{}+ | Nbar(1880)\allowbreak{}-\par {\fontsize{8.1}{9.4}\selectfont PDG2026:\allowbreak{}observable:\allowbreak{}B087:\allowbreak{}B087M:\allowbreak{}348671} & \textbf{1830 to 1880 to 1930} MeV\par Presentación: 1830 to 1930 (\textasciitilde{}1880) & \(u_+\): NA\par \(u_-\): NA\par Tipo: R\par Confianza: ND & COMP-BARYON\par E04 / M07\\\addlinespace[3pt]
% VI-CATALOG-ROW masas 66
\textbf{66. Sigma(2080) MASS}\par Sigma(2080)\allowbreak{}- | Sigmabar(2080)\allowbreak{}+ | Sigma(2080)\allowbreak{}0 | Sigmabar(2080)\allowbreak{}0 | Sigma(2080)\allowbreak{}+ | Sigmabar(2080)\allowbreak{}-\par {\fontsize{8.1}{9.4}\selectfont PDG2026:\allowbreak{}observable:\allowbreak{}B088:\allowbreak{}B088M:\allowbreak{}350188} & \textbf{2060 to 2090 to 2120} MeV\par Presentación: 2060 to 2120 (\textasciitilde{}2090) & \(u_+\): NA\par \(u_-\): NA\par Tipo: R\par Confianza: ND & COMP-BARYON\par E04 / M07\\\addlinespace[3pt]
% VI-CATALOG-ROW masas 67
\textbf{67. Lambda(2000) MASS}\par Lambda(2000)\allowbreak{}0 | Lambdabar(2000)\allowbreak{}0\par {\fontsize{8.1}{9.4}\selectfont PDG2026:\allowbreak{}observable:\allowbreak{}B089:\allowbreak{}B089M:\allowbreak{}349646} & \textbf{\textasciitilde{}2000} MeV & \(u_+\): ND\par \(u_-\): ND\par Tipo: NONE\par Confianza: ND & COMP-BARYON\par E04 / M07\\\addlinespace[3pt]
% VI-CATALOG-ROW masas 68
\textbf{68. N(2220) BREIT-WIGNER MASS}\par N(2220)\allowbreak{}0 | Nbar(2220)\allowbreak{}0 | N(2220)\allowbreak{}+ | Nbar(2220)\allowbreak{}-\par {\fontsize{8.1}{9.4}\selectfont PDG2026:\allowbreak{}observable:\allowbreak{}B090:\allowbreak{}B090M:\allowbreak{}348986} & \textbf{2200 TO 2250 TO 2300} MeV\par Presentación: 2200 to 2300 (\textasciitilde{}2250) & \(u_+\): NA\par \(u_-\): NA\par Tipo: R\par Confianza: ND & COMP-BARYON\par E04 / M07\\\addlinespace[3pt]
% VI-CATALOG-ROW masas 69
\textbf{69. Sigma(1910) MASS}\par Sigma(1910)\allowbreak{}- | Sigmabar(1910)\allowbreak{}+ | Sigma(1910)\allowbreak{}0 | Sigmabar(1910)\allowbreak{}0 | Sigma(1910)\allowbreak{}+ | Sigmabar(1910)\allowbreak{}-\par {\fontsize{8.1}{9.4}\selectfont PDG2026:\allowbreak{}observable:\allowbreak{}B098:\allowbreak{}B098M:\allowbreak{}350050} & \textbf{1870 to 1910 to 1950} MeV\par Presentación: 1870 to 1950 (\textasciitilde{}1910) & \(u_+\): NA\par \(u_-\): NA\par Tipo: R\par Confianza: ND & COMP-BARYON\par E04 / M07\\\addlinespace[3pt]
% VI-CATALOG-ROW masas 70
\textbf{70. Xi(2500) MASS}\par Xi(2500)\allowbreak{}- | Xibar(2500)\allowbreak{}+ | Xi(2500)\allowbreak{}0 | Xibar(2500)\allowbreak{}0\par {\fontsize{8.1}{9.4}\selectfont PDG2026:\allowbreak{}observable:\allowbreak{}B099:\allowbreak{}B099M:\allowbreak{}350464} & \textbf{\textasciitilde{}2500} MeV & \(u_+\): ND\par \(u_-\): ND\par Tipo: NONE\par Confianza: ND & COMP-BARYON\par E04 / M07\\\addlinespace[3pt]
% VI-CATALOG-ROW masas 71
\textbf{71. Lambda(1600) MASS}\par Lambda(1600)\allowbreak{}0 | Lambdabar(1600)\allowbreak{}0\par {\fontsize{8.1}{9.4}\selectfont PDG2026:\allowbreak{}observable:\allowbreak{}B101:\allowbreak{}B101M:\allowbreak{}349472} & \textbf{1570 to 1600 to 1630} MeV\par Presentación: 1570 to 1630 (\textasciitilde{}1600) & \(u_+\): NA\par \(u_-\): NA\par Tipo: R\par Confianza: ND & COMP-BARYON\par E04 / M07\\\addlinespace[3pt]
% VI-CATALOG-ROW masas 72
\textbf{72. Lambda\_\allowbreak{}c(2625)\allowbreak{}+ MASS}\par Lambda\_\allowbreak{}c(2625)\allowbreak{}+ | Lambdabar\_\allowbreak{}c(2625)\allowbreak{}-\par {\fontsize{8.1}{9.4}\selectfont PDG2026:\allowbreak{}observable:\allowbreak{}B102:\allowbreak{}B102M:\allowbreak{}350934} & \textbf{2628.00+-0.15} MeV & \(u_+\): 0.147359\allowbreak{}68814977\allowbreak{}221\par \(u_-\): 0.147359\allowbreak{}68814977\allowbreak{}221\par Tipo: NONE\par Confianza: ND & COMP-BARYON\par E04 / M07\\\addlinespace[3pt]
% VI-CATALOG-ROW masas 73
\textbf{73. Xi(2120) MASS}\par Xi(2120)\allowbreak{}- | Xibar(2120)\allowbreak{}+ | Xi(2120)\allowbreak{}0 | Xibar(2120)\allowbreak{}0\par {\fontsize{8.1}{9.4}\selectfont PDG2026:\allowbreak{}observable:\allowbreak{}B103:\allowbreak{}B103M:\allowbreak{}350454} & \textbf{\textasciitilde{}2120} MeV & \(u_+\): ND\par \(u_-\): ND\par Tipo: NONE\par Confianza: ND & COMP-BARYON\par E04 / M07\\\addlinespace[3pt]
% VI-CATALOG-ROW masas 74
\textbf{74. Sigma\_\allowbreak{}c(2455)\allowbreak{}++ MASS}\par Sigma\_\allowbreak{}c(2455)\allowbreak{}0 | Sigmabar\_\allowbreak{}c(2455)\allowbreak{}0 | Sigma\_\allowbreak{}c(2455)\allowbreak{}+ | Sigmabar\_\allowbreak{}c(2455)\allowbreak{}- | Sigma\_\allowbreak{}c(2455)\allowbreak{}++ | Sigmabar\_\allowbreak{}c(2455)\allowbreak{}-\kern0pt{}-\par {\fontsize{8.1}{9.4}\selectfont PDG2026:\allowbreak{}observable:\allowbreak{}B104:\allowbreak{}B104M++:\allowbreak{}350985} & \textbf{2453.97+-0.14} MeV & \(u_+\): 0.141009\allowbreak{}78554010\allowbreak{}379\par \(u_-\): 0.141009\allowbreak{}78554010\allowbreak{}379\par Tipo: NONE\par Confianza: ND & COMP-BARYON\par E04 / M07\\\addlinespace[3pt]
% VI-CATALOG-ROW masas 75
\textbf{75. Sigma\_\allowbreak{}c(2455)\allowbreak{}+ MASS}\par Sigma\_\allowbreak{}c(2455)\allowbreak{}0 | Sigmabar\_\allowbreak{}c(2455)\allowbreak{}0 | Sigma\_\allowbreak{}c(2455)\allowbreak{}+ | Sigmabar\_\allowbreak{}c(2455)\allowbreak{}- | Sigma\_\allowbreak{}c(2455)\allowbreak{}++ | Sigmabar\_\allowbreak{}c(2455)\allowbreak{}-\kern0pt{}-\par {\fontsize{8.1}{9.4}\selectfont PDG2026:\allowbreak{}observable:\allowbreak{}B104:\allowbreak{}B104M+:\allowbreak{}350986} & \textbf{2452.65+0.22-0.16} MeV & \(u_+\): 0.217932\allowbreak{}57206126\allowbreak{}77\par \(u_-\): 0.157506\allowbreak{}25023659\allowbreak{}159\par Tipo: NONE\par Confianza: ND & COMP-BARYON\par E04 / M07\\\addlinespace[3pt]
% VI-CATALOG-ROW masas 76
\textbf{76. Sigma\_\allowbreak{}c(2455)\allowbreak{}0 MASS}\par Sigma\_\allowbreak{}c(2455)\allowbreak{}0 | Sigmabar\_\allowbreak{}c(2455)\allowbreak{}0 | Sigma\_\allowbreak{}c(2455)\allowbreak{}+ | Sigmabar\_\allowbreak{}c(2455)\allowbreak{}- | Sigma\_\allowbreak{}c(2455)\allowbreak{}++ | Sigmabar\_\allowbreak{}c(2455)\allowbreak{}-\kern0pt{}-\par {\fontsize{8.1}{9.4}\selectfont PDG2026:\allowbreak{}observable:\allowbreak{}B104:\allowbreak{}B104M0:\allowbreak{}350987} & \textbf{2453.75+-0.14} MeV & \(u_+\): 0.141009\allowbreak{}37832457\allowbreak{}261\par \(u_-\): 0.141009\allowbreak{}37832457\allowbreak{}261\par Tipo: NONE\par Confianza: ND & COMP-BARYON\par E04 / M07\\\addlinespace[3pt]
% VI-CATALOG-ROW masas 77
\textbf{77. MIXED CHARGES}\par Xi(1690)\allowbreak{}- | Xibar(1690)\allowbreak{}+ | Xi(1690)\allowbreak{}0 | Xibar(1690)\allowbreak{}0\par {\fontsize{8.1}{9.4}\selectfont PDG2026:\allowbreak{}observable:\allowbreak{}B105:\allowbreak{}B105M:\allowbreak{}350412} & \textbf{1690+-10} MeV & \(u_+\): 10\par \(u_-\): 10\par Tipo: NONE\par Confianza: ND & COMP-BARYON\par E04 / M07\\\addlinespace[3pt]
% VI-CATALOG-ROW masas 78
\textbf{78. Lambda(2325) MASS}\par Lambda(2325)\allowbreak{}0 | Lambdabar(2325)\allowbreak{}0\par {\fontsize{8.1}{9.4}\selectfont PDG2026:\allowbreak{}observable:\allowbreak{}B112:\allowbreak{}B112M:\allowbreak{}349774} & \textbf{\textasciitilde{}2325} MeV & \(u_+\): ND\par \(u_-\): ND\par Tipo: NONE\par Confianza: ND & COMP-BARYON\par E04 / M07\\\addlinespace[3pt]
% VI-CATALOG-ROW masas 79
\textbf{79. N(2250) BREIT-WIGNER MASS}\par N(2250)\allowbreak{}0 | Nbar(2250)\allowbreak{}0 | N(2250)\allowbreak{}+ | Nbar(2250)\allowbreak{}-\par {\fontsize{8.1}{9.4}\selectfont PDG2026:\allowbreak{}observable:\allowbreak{}B113:\allowbreak{}B113M:\allowbreak{}349004} & \textbf{2250 to 2280 to 2320} MeV\par Presentación: 2250 to 2320 (\textasciitilde{}2280) & \(u_+\): NA\par \(u_-\): NA\par Tipo: R\par Confianza: ND & COMP-BARYON\par E04 / M07\\\addlinespace[3pt]
% VI-CATALOG-ROW masas 80
\textbf{80. Sigma\_\allowbreak{}c(2520)\allowbreak{}++ MASS}\par Sigma\_\allowbreak{}c(2520)\allowbreak{}0 | Sigmabar\_\allowbreak{}c(2520)\allowbreak{}0 | Sigma\_\allowbreak{}c(2520)\allowbreak{}+ | Sigmabar\_\allowbreak{}c(2520)\allowbreak{}- | Sigma\_\allowbreak{}c(2520)\allowbreak{}++ | Sigmabar\_\allowbreak{}c(2520)\allowbreak{}-\kern0pt{}-\par {\fontsize{8.1}{9.4}\selectfont PDG2026:\allowbreak{}observable:\allowbreak{}B115:\allowbreak{}B115M++:\allowbreak{}351001} & \textbf{2518.41+-0.22} MeV & \(u_+\): 0.224414\allowbreak{}20873039\allowbreak{}311\par \(u_-\): 0.224414\allowbreak{}20873039\allowbreak{}311\par Tipo: NONE\par Confianza: ND & COMP-BARYON\par E04 / M07\\\addlinespace[3pt]
% VI-CATALOG-ROW masas 81
\textbf{81. Sigma\_\allowbreak{}c(2520)\allowbreak{}+ MASS}\par Sigma\_\allowbreak{}c(2520)\allowbreak{}0 | Sigmabar\_\allowbreak{}c(2520)\allowbreak{}0 | Sigma\_\allowbreak{}c(2520)\allowbreak{}+ | Sigmabar\_\allowbreak{}c(2520)\allowbreak{}- | Sigma\_\allowbreak{}c(2520)\allowbreak{}++ | Sigmabar\_\allowbreak{}c(2520)\allowbreak{}-\kern0pt{}-\par {\fontsize{8.1}{9.4}\selectfont PDG2026:\allowbreak{}observable:\allowbreak{}B115:\allowbreak{}B115M+:\allowbreak{}351002} & \textbf{2517.4+0.7-0.5} MeV & \(u_+\): 0.685051\allowbreak{}02093303\allowbreak{}338\par \(u_-\): 0.524193\allowbreak{}84686487\allowbreak{}569\par Tipo: NONE\par Confianza: ND & COMP-BARYON\par E04 / M07\\\addlinespace[3pt]
% VI-CATALOG-ROW masas 82
\textbf{82. Sigma\_\allowbreak{}c(2520)\allowbreak{}0 MASS}\par Sigma\_\allowbreak{}c(2520)\allowbreak{}0 | Sigmabar\_\allowbreak{}c(2520)\allowbreak{}0 | Sigma\_\allowbreak{}c(2520)\allowbreak{}+ | Sigmabar\_\allowbreak{}c(2520)\allowbreak{}- | Sigma\_\allowbreak{}c(2520)\allowbreak{}++ | Sigmabar\_\allowbreak{}c(2520)\allowbreak{}-\kern0pt{}-\par {\fontsize{8.1}{9.4}\selectfont PDG2026:\allowbreak{}observable:\allowbreak{}B115:\allowbreak{}B115M0:\allowbreak{}351003} & \textbf{2518.48+-0.21} MeV & \(u_+\): 0.207957\allowbreak{}24066290\allowbreak{}969\par \(u_-\): 0.207957\allowbreak{}24066290\allowbreak{}969\par Tipo: NONE\par Confianza: ND & COMP-BARYON\par E04 / M07\\\addlinespace[3pt]
% VI-CATALOG-ROW masas 83
\textbf{83. Delta(1920) BREIT-WIGNER MASS}\par Delta(1920)\allowbreak{}- | Deltabar(1920)\allowbreak{}+ | Delta(1920)\allowbreak{}0 | Deltabar(1920)\allowbreak{}0 | Delta(1920)\allowbreak{}+ | Deltabar(1920)\allowbreak{}- | Delta(1920)\allowbreak{}++ | Deltabar(1920)\allowbreak{}-\kern0pt{}-\par {\fontsize{8.1}{9.4}\selectfont PDG2026:\allowbreak{}observable:\allowbreak{}B117:\allowbreak{}B117M:\allowbreak{}349208} & \textbf{1870 to 1920 to 1970} MeV\par Presentación: 1870 to 1970 (\textasciitilde{}1920) & \(u_+\): NA\par \(u_-\): NA\par Tipo: R\par Confianza: ND & COMP-BARYON\par E04 / M07\\\addlinespace[3pt]
% VI-CATALOG-ROW masas 84
\textbf{84. Sigma(3170) MASS (PRODUCTION EXPERIMENTS)}\par Sigma(3170)\allowbreak{}- | Sigmabar(3170)\allowbreak{}+ | Sigma(3170)\allowbreak{}0 | Sigmabar(3170)\allowbreak{}0 | Sigma(3170)\allowbreak{}+ | Sigmabar(3170)\allowbreak{}-\par {\fontsize{8.1}{9.4}\selectfont PDG2026:\allowbreak{}observable:\allowbreak{}B118:\allowbreak{}B118M:\allowbreak{}350286} & \textbf{\textasciitilde{}3170} MeV & \(u_+\): ND\par \(u_-\): ND\par Tipo: NONE\par Confianza: ND & COMP-BARYON\par E04 / M07\\\addlinespace[3pt]
% VI-CATALOG-ROW masas 85
\textbf{85. Lambda\_\allowbreak{}c(2595)\allowbreak{}+ MASS}\par Lambda\_\allowbreak{}c(2595)\allowbreak{}+ | Lambdabar\_\allowbreak{}c(2595)\allowbreak{}-\par {\fontsize{8.1}{9.4}\selectfont PDG2026:\allowbreak{}observable:\allowbreak{}B119:\allowbreak{}B119M:\allowbreak{}350917} & \textbf{2592.25+-0.28} MeV & \(u_+\): 0.281424\allowbreak{}94558940\allowbreak{}517\par \(u_-\): 0.281424\allowbreak{}94558940\allowbreak{}517\par Tipo: NONE\par Confianza: ND & COMP-BARYON\par E04 / M07\\\addlinespace[3pt]
% VI-CATALOG-ROW masas 86
\textbf{86. N(2600) BREIT-WIGNER MASS}\par N(2600)\allowbreak{}0 | Nbar(2600)\allowbreak{}0 | N(2600)\allowbreak{}+ | Nbar(2600)\allowbreak{}-\par {\fontsize{8.1}{9.4}\selectfont PDG2026:\allowbreak{}observable:\allowbreak{}B120:\allowbreak{}B120M:\allowbreak{}349022} & \textbf{2550 TO 2600 TO 2750} MeV\par Presentación: 2550 to 2750 (\textasciitilde{}2600) & \(u_+\): NA\par \(u_-\): NA\par Tipo: R\par Confianza: ND & COMP-BARYON\par E04 / M07\\\addlinespace[3pt]
% VI-CATALOG-ROW masas 87
\textbf{87. Lambda\_\allowbreak{}c(2940)\allowbreak{}+ MASS}\par Lambda\_\allowbreak{}c(2940)\allowbreak{}0 | Lambdabar\_\allowbreak{}c(2940)\allowbreak{}0\par {\fontsize{8.1}{9.4}\selectfont PDG2026:\allowbreak{}observable:\allowbreak{}B122:\allowbreak{}B122M:\allowbreak{}350979} & \textbf{2939.6+1.3-1.5} MeV & \(u_+\): 1.261876\allowbreak{}68447521\allowbreak{}3\par \(u_-\): 1.467184\allowbreak{}98507142\allowbreak{}8\par Tipo: NONE\par Confianza: ND & COMP-BARYON\par E04 / M07\\\addlinespace[3pt]
% VI-CATALOG-ROW masas 88
\textbf{88. N(\textasciitilde{}3000) BREIT-WIGNER MASS}\par N(\textasciitilde{}3000)\allowbreak{}0 | Nbar(\textasciitilde{}3000)\allowbreak{}0 | N(\textasciitilde{}3000)\allowbreak{}+ | Nbar(\textasciitilde{}3000)\allowbreak{}-\par {\fontsize{8.1}{9.4}\selectfont PDG2026:\allowbreak{}observable:\allowbreak{}B128:\allowbreak{}B128M:\allowbreak{}349027} & \textbf{\textasciitilde{}3000} MeV & \(u_+\): ND\par \(u_-\): ND\par Tipo: NONE\par Confianza: ND & COMP-BARYON\par E04 / M07\\\addlinespace[3pt]
% VI-CATALOG-ROW masas 89
\textbf{89. Xi\_\allowbreak{}c(2970)\allowbreak{}+ MASS}\par Xi\_\allowbreak{}c(2970)\allowbreak{}0 | Xibar\_\allowbreak{}c(2970)\allowbreak{}0\par {\fontsize{8.1}{9.4}\selectfont PDG2026:\allowbreak{}observable:\allowbreak{}B130:\allowbreak{}B130M+:\allowbreak{}351271} & \textbf{2965.2+-1.0} MeV & \(u_+\): 1.015681\allowbreak{}06038590\allowbreak{}41\par \(u_-\): 1.015681\allowbreak{}06038590\allowbreak{}41\par Tipo: NONE\par Confianza: ND & COMP-BARYON\par E04 / M07\\\addlinespace[3pt]
% VI-CATALOG-ROW masas 90
\textbf{90. Xi\_\allowbreak{}c(2970)\allowbreak{}0 MASS}\par Xi\_\allowbreak{}c(2970)\allowbreak{}0 | Xibar\_\allowbreak{}c(2970)\allowbreak{}0\par {\fontsize{8.1}{9.4}\selectfont PDG2026:\allowbreak{}observable:\allowbreak{}B130:\allowbreak{}B130M0:\allowbreak{}351273} & \textbf{2967.3+-1.9} MeV & \(u_+\): 1.850084\allowbreak{}39480852\allowbreak{}7\par \(u_-\): 1.850084\allowbreak{}39480852\allowbreak{}7\par Tipo: NONE\par Confianza: ND & COMP-BARYON\par E04 / M07\\\addlinespace[3pt]
% VI-CATALOG-ROW masas 91
\textbf{91. Xi(2370) MASS}\par Xi(2370)\allowbreak{}- | Xibar(2370)\allowbreak{}+ | Xi(2370)\allowbreak{}0 | Xibar(2370)\allowbreak{}0\par {\fontsize{8.1}{9.4}\selectfont PDG2026:\allowbreak{}observable:\allowbreak{}B131:\allowbreak{}B131M:\allowbreak{}350458} & \textbf{\textasciitilde{}2370} MeV & \(u_+\): ND\par \(u_-\): ND\par Tipo: NONE\par Confianza: ND & COMP-BARYON\par E04 / M07\\\addlinespace[3pt]
% VI-CATALOG-ROW masas 92
\textbf{92. N(2100) BREIT-WIGNER MASS}\par N(2100)\allowbreak{}0 | Nbar(2100)\allowbreak{}0 | N(2100)\allowbreak{}+ | Nbar(2100)\allowbreak{}-\par {\fontsize{8.1}{9.4}\selectfont PDG2026:\allowbreak{}observable:\allowbreak{}B132:\allowbreak{}B132M:\allowbreak{}348871} & \textbf{2050 to 2100 to 2150} MeV\par Presentación: 2050 to 2150 (\textasciitilde{}2100) & \(u_+\): NA\par \(u_-\): NA\par Tipo: R\par Confianza: ND & COMP-BARYON\par E04 / M07\\\addlinespace[3pt]
% VI-CATALOG-ROW masas 93
\textbf{93. Delta(2200) BREIT-WIGNER MASS}\par Delta(2200)\allowbreak{}- | Deltabar(2200)\allowbreak{}+ | Delta(2200)\allowbreak{}0 | Deltabar(2200)\allowbreak{}0 | Delta(2200)\allowbreak{}+ | Deltabar(2200)\allowbreak{}- | Delta(2200)\allowbreak{}++ | Deltabar(2200)\allowbreak{}-\kern0pt{}-\par {\fontsize{8.1}{9.4}\selectfont PDG2026:\allowbreak{}observable:\allowbreak{}B135:\allowbreak{}B135M:\allowbreak{}349339} & \textbf{2150 to 2200 to 2250} MeV\par Presentación: 2150 to 2250 (\textasciitilde{}2200) & \(u_+\): NA\par \(u_-\): NA\par Tipo: R\par Confianza: ND & COMP-BARYON\par E04 / M07\\\addlinespace[3pt]
% VI-CATALOG-ROW masas 94
\textbf{94. Delta(1940) BREIT-WIGNER MASS}\par Delta(1940)\allowbreak{}- | Deltabar(1940)\allowbreak{}+ | Delta(1940)\allowbreak{}0 | Deltabar(1940)\allowbreak{}0 | Delta(1940)\allowbreak{}+ | Deltabar(1940)\allowbreak{}- | Delta(1940)\allowbreak{}++ | Deltabar(1940)\allowbreak{}-\kern0pt{}-\par {\fontsize{8.1}{9.4}\selectfont PDG2026:\allowbreak{}observable:\allowbreak{}B136:\allowbreak{}B136M:\allowbreak{}349272} & \textbf{1940 to 2000 to 2060} MeV\par Presentación: 1940 to 2060 (\textasciitilde{}2000) & \(u_+\): NA\par \(u_-\): NA\par Tipo: R\par Confianza: ND & COMP-BARYON\par E04 / M07\\\addlinespace[3pt]
% VI-CATALOG-ROW masas 95
\textbf{95. Omega(2250)\allowbreak{}- MASS}\par Omega(2250)\allowbreak{}- | Omegabar(2250)\allowbreak{}+\par {\fontsize{8.1}{9.4}\selectfont PDG2026:\allowbreak{}observable:\allowbreak{}B141:\allowbreak{}B141M:\allowbreak{}350516} & \textbf{2252+-9} MeV & \(u_+\): 8.834104\allowbreak{}44096544\allowbreak{}19\par \(u_-\): 8.834104\allowbreak{}44096544\allowbreak{}19\par Tipo: NONE\par Confianza: ND & COMP-BARYON\par E04 / M07\\\addlinespace[3pt]
% VI-CATALOG-ROW masas 96
\textbf{96. Omega(2380)\allowbreak{}- MASS}\par Omega(2380)\allowbreak{}- | Omegabar(2380)\allowbreak{}+\par {\fontsize{8.1}{9.4}\selectfont PDG2026:\allowbreak{}observable:\allowbreak{}B142:\allowbreak{}B142M:\allowbreak{}350520} & \textbf{\textasciitilde{}2380} MeV & \(u_+\): ND\par \(u_-\): ND\par Tipo: NONE\par Confianza: ND & COMP-BARYON\par E04 / M07\\\addlinespace[3pt]
% VI-CATALOG-ROW masas 97
\textbf{97. Omega(2470)\allowbreak{}- MASS}\par Omega(2470)\allowbreak{}- | Omegabar(2470)\allowbreak{}+\par {\fontsize{8.1}{9.4}\selectfont PDG2026:\allowbreak{}observable:\allowbreak{}B143:\allowbreak{}B143M:\allowbreak{}350523} & \textbf{2474+-12} MeV & \(u_+\): 12\par \(u_-\): 12\par Tipo: NONE\par Confianza: ND & COMP-BARYON\par E04 / M07\\\addlinespace[3pt]
% VI-CATALOG-ROW masas 98
\textbf{98. N(1900) BREIT-WIGNER MASS}\par N(1900)\allowbreak{}0 | Nbar(1900)\allowbreak{}0 | N(1900)\allowbreak{}+ | Nbar(1900)\allowbreak{}-\par {\fontsize{8.1}{9.4}\selectfont PDG2026:\allowbreak{}observable:\allowbreak{}B144:\allowbreak{}B144M:\allowbreak{}348740} & \textbf{1890 to 1920 to 1950} MeV\par Presentación: 1890 to 1950 (\textasciitilde{}1920) & \(u_+\): NA\par \(u_-\): NA\par Tipo: R\par Confianza: ND & COMP-BARYON\par E04 / M07\\\addlinespace[3pt]
% VI-CATALOG-ROW masas 99
\textbf{99. Xi\_\allowbreak{}c(2645)\allowbreak{}+ MASS}\par Xi\_\allowbreak{}c(2645)\allowbreak{}0 | Xibar\_\allowbreak{}c(2645)\allowbreak{}0 | Xi\_\allowbreak{}c(2645)\allowbreak{}+ | Xibar\_\allowbreak{}c(2645)\allowbreak{}-\par {\fontsize{8.1}{9.4}\selectfont PDG2026:\allowbreak{}observable:\allowbreak{}B146:\allowbreak{}B146M+:\allowbreak{}351204} & \textbf{2645.17+-0.27} MeV & \(u_+\): 0.262812\allowbreak{}13522405\allowbreak{}311\par \(u_-\): 0.275224\allowbreak{}52658037\allowbreak{}4\par Tipo: NONE\par Confianza: ND & COMP-BARYON\par E04 / M07\\\addlinespace[3pt]
% VI-CATALOG-ROW masas 100
\textbf{100. Xi\_\allowbreak{}c(2645)\allowbreak{}0 MASS}\par Xi\_\allowbreak{}c(2645)\allowbreak{}0 | Xibar\_\allowbreak{}c(2645)\allowbreak{}0 | Xi\_\allowbreak{}c(2645)\allowbreak{}+ | Xibar\_\allowbreak{}c(2645)\allowbreak{}-\par {\fontsize{8.1}{9.4}\selectfont PDG2026:\allowbreak{}observable:\allowbreak{}B146:\allowbreak{}B146M0:\allowbreak{}351206} & \textbf{2646.24+-0.18} MeV & \(u_+\): 0.177927\allowbreak{}28963027\allowbreak{}479\par \(u_-\): 0.180979\allowbreak{}89449585\allowbreak{}48\par Tipo: NONE\par Confianza: ND & COMP-BARYON\par E04 / M07\\\addlinespace[3pt]
% VI-CATALOG-ROW masas 101
\textbf{101. Xi\_\allowbreak{}c(3080)\allowbreak{}+ MASS}\par Xi\_\allowbreak{}c(3080)\allowbreak{}0 | Xibar\_\allowbreak{}c(3080)\allowbreak{}0\par {\fontsize{8.1}{9.4}\selectfont PDG2026:\allowbreak{}observable:\allowbreak{}B147:\allowbreak{}B147M+:\allowbreak{}351307} & \textbf{3077.2+-0.4} MeV & \(u_+\): 0.386527\allowbreak{}49298648\allowbreak{}748\par \(u_-\): 0.386527\allowbreak{}49298648\allowbreak{}748\par Tipo: NONE\par Confianza: ND & COMP-BARYON\par E04 / M07\\\addlinespace[3pt]
% VI-CATALOG-ROW masas 102
\textbf{102. Xi\_\allowbreak{}c(3080)\allowbreak{}0 MASS}\par Xi\_\allowbreak{}c(3080)\allowbreak{}0 | Xibar\_\allowbreak{}c(3080)\allowbreak{}0\par {\fontsize{8.1}{9.4}\selectfont PDG2026:\allowbreak{}observable:\allowbreak{}B147:\allowbreak{}B147M0:\allowbreak{}351308} & \textbf{3079.9+-1.4} MeV & \(u_+\): 1.360345\allowbreak{}85012566\allowbreak{}09\par \(u_-\): 1.360345\allowbreak{}85012566\allowbreak{}09\par Tipo: NONE\par Confianza: ND & COMP-BARYON\par E04 / M07\\\addlinespace[3pt]
% VI-CATALOG-ROW masas 103
\textbf{103. Xi\_\allowbreak{}c(2815)\allowbreak{}+ MASS}\par Xi\_\allowbreak{}c(2815)\allowbreak{}0 | Xibar\_\allowbreak{}c(2815)\allowbreak{}0 | Xi\_\allowbreak{}c(2815)\allowbreak{}+ | Xibar\_\allowbreak{}c(2815)\allowbreak{}-\par {\fontsize{8.1}{9.4}\selectfont PDG2026:\allowbreak{}observable:\allowbreak{}B148:\allowbreak{}B148M+:\allowbreak{}351232} & \textbf{2816.60+-0.16} MeV & \(u_+\): 0.155171\allowbreak{}26784005\allowbreak{}561\par \(u_-\): 0.157672\allowbreak{}30758639\allowbreak{}7\par Tipo: NONE\par Confianza: ND & COMP-BARYON\par E04 / M07\\\addlinespace[3pt]
% VI-CATALOG-ROW masas 104
\textbf{104. Xi\_\allowbreak{}c(2815)\allowbreak{}0 MASS}\par Xi\_\allowbreak{}c(2815)\allowbreak{}0 | Xibar\_\allowbreak{}c(2815)\allowbreak{}0 | Xi\_\allowbreak{}c(2815)\allowbreak{}+ | Xibar\_\allowbreak{}c(2815)\allowbreak{}-\par {\fontsize{8.1}{9.4}\selectfont PDG2026:\allowbreak{}observable:\allowbreak{}B148:\allowbreak{}B148M0:\allowbreak{}351234} & \textbf{2819.85+-0.27} MeV & \(u_+\): 0.268091\allowbreak{}34636591\allowbreak{}098\par \(u_-\): 0.280254\allowbreak{}92882096\allowbreak{}131\par Tipo: NONE\par Confianza: ND & COMP-BARYON\par E04 / M07\\\addlinespace[3pt]
% VI-CATALOG-ROW masas 105
\textbf{105. Xi\_\allowbreak{}c(2790)\allowbreak{}+ MASS}\par Xi\_\allowbreak{}c(2790)\allowbreak{}0 | Xibar\_\allowbreak{}c(2790)\allowbreak{}0 | Xi\_\allowbreak{}c(2790)\allowbreak{}+ | Xibar\_\allowbreak{}c(2790)\allowbreak{}-\par {\fontsize{8.1}{9.4}\selectfont PDG2026:\allowbreak{}observable:\allowbreak{}B149:\allowbreak{}B149M+:\allowbreak{}351217} & \textbf{2792.0+-0.5} MeV & \(u_+\): 0.529738\allowbreak{}97333377\allowbreak{}062\par \(u_-\): 0.535852\allowbreak{}70924116\allowbreak{}923\par Tipo: NONE\par Confianza: ND & COMP-BARYON\par E04 / M07\\\addlinespace[3pt]
% VI-CATALOG-ROW masas 106
\textbf{106. Xi\_\allowbreak{}c(2790)\allowbreak{}0 MASS}\par Xi\_\allowbreak{}c(2790)\allowbreak{}0 | Xibar\_\allowbreak{}c(2790)\allowbreak{}0 | Xi\_\allowbreak{}c(2790)\allowbreak{}+ | Xibar\_\allowbreak{}c(2790)\allowbreak{}-\par {\fontsize{8.1}{9.4}\selectfont PDG2026:\allowbreak{}observable:\allowbreak{}B149:\allowbreak{}B149M0:\allowbreak{}351218} & \textbf{2794.0+-0.5} MeV & \(u_+\): 0.492344\allowbreak{}23994343\allowbreak{}779\par \(u_-\): 0.492344\allowbreak{}23994343\allowbreak{}779\par Tipo: NONE\par Confianza: ND & COMP-BARYON\par E04 / M07\\\addlinespace[3pt]
% VI-CATALOG-ROW masas 107
\textbf{107. Lambda\_\allowbreak{}c(2765)\allowbreak{}+ MASS}\par Lambda\_\allowbreak{}c(2765)\allowbreak{}0 | Lambdabar\_\allowbreak{}c(2765)\allowbreak{}0\par {\fontsize{8.1}{9.4}\selectfont PDG2026:\allowbreak{}observable:\allowbreak{}B150:\allowbreak{}B150M:\allowbreak{}350951} & \textbf{2766.6+-2.4} MeV & \(u_+\): 2.404079\allowbreak{}86556187\allowbreak{}4\par \(u_-\): 2.404079\allowbreak{}86556187\allowbreak{}4\par Tipo: NONE\par Confianza: ND & COMP-BARYON\par E04 / M07\\\addlinespace[3pt]
% VI-CATALOG-ROW masas 108
\textbf{108. Lambda\_\allowbreak{}c(2880)\allowbreak{}+ MASS}\par Lambda\_\allowbreak{}c(2880)\allowbreak{}+ | Lambdabar\_\allowbreak{}c(2880)\allowbreak{}-\par {\fontsize{8.1}{9.4}\selectfont PDG2026:\allowbreak{}observable:\allowbreak{}B151:\allowbreak{}B151M:\allowbreak{}350960} & \textbf{2881.63+-0.24} MeV & \(u_+\): 0.236132\allowbreak{}03643473\allowbreak{}499\par \(u_-\): 0.245558\allowbreak{}68414433\allowbreak{}36\par Tipo: NONE\par Confianza: ND & COMP-BARYON\par E04 / M07\\\addlinespace[3pt]
% VI-CATALOG-ROW masas 109
\textbf{109. Sigma\_\allowbreak{}c(2800)\allowbreak{}++ MASS}\par Sigma\_\allowbreak{}c(2800)\allowbreak{}0 | Sigmabar\_\allowbreak{}c(2800)\allowbreak{}0\par {\fontsize{8.1}{9.4}\selectfont PDG2026:\allowbreak{}observable:\allowbreak{}B155:\allowbreak{}B155M++:\allowbreak{}351015} & \textbf{2801+4-6} MeV & \(u_+\): 4.406767\allowbreak{}52279935\allowbreak{}26\par \(u_-\): 5.799965\allowbreak{}51713898\allowbreak{}52\par Tipo: NONE\par Confianza: ND & COMP-BARYON\par E04 / M07\\\addlinespace[3pt]
% VI-CATALOG-ROW masas 110
\textbf{110. Sigma\_\allowbreak{}c(2800)\allowbreak{}+ MASS}\par Sigma\_\allowbreak{}c(2800)\allowbreak{}0 | Sigmabar\_\allowbreak{}c(2800)\allowbreak{}0\par {\fontsize{8.1}{9.4}\selectfont PDG2026:\allowbreak{}observable:\allowbreak{}B155:\allowbreak{}B155M+:\allowbreak{}351016} & \textbf{2792+14-5} MeV & \(u_+\): 13.69012\allowbreak{}78299361\allowbreak{}1\par \(u_-\): 5.017927\allowbreak{}85918669\par Tipo: NONE\par Confianza: ND & COMP-BARYON\par E04 / M07\\\addlinespace[3pt]
% VI-CATALOG-ROW masas 111
\textbf{111. Sigma\_\allowbreak{}c(2800)\allowbreak{}0 MASS}\par Sigma\_\allowbreak{}c(2800)\allowbreak{}0 | Sigmabar\_\allowbreak{}c(2800)\allowbreak{}0\par {\fontsize{8.1}{9.4}\selectfont PDG2026:\allowbreak{}observable:\allowbreak{}B155:\allowbreak{}B155M0:\allowbreak{}351017} & \textbf{2806+5-7} MeV & \(u_+\): 4.859001\allowbreak{}99216405\allowbreak{}33\par \(u_-\): 7.387144\allowbreak{}26282947\allowbreak{}39\par Tipo: NONE\par Confianza: ND & COMP-BARYON\par E04 / M07\\\addlinespace[3pt]
% VI-CATALOG-ROW masas 112
\textbf{112. Xi\_\allowbreak{}c(2930)\allowbreak{}+ MASS}\par Xi\_\allowbreak{}c(2930)\allowbreak{}0 | Xibar\_\allowbreak{}c(2930)\allowbreak{}0\par {\fontsize{8.1}{9.4}\selectfont PDG2026:\allowbreak{}observable:\allowbreak{}B156:\allowbreak{}B156A03:\allowbreak{}351263} & \textbf{2942+-5} MeV & \(u_+\): 4.648655\allowbreak{}71966778\allowbreak{}38\par \(u_-\): 4.648655\allowbreak{}71966778\allowbreak{}38\par Tipo: NONE\par Confianza: ND & COMP-BARYON\par E04 / M07\\\addlinespace[3pt]
% VI-CATALOG-ROW masas 113
\textbf{113. Xi\_\allowbreak{}c(2930)\allowbreak{}0 MASS}\par Xi\_\allowbreak{}c(2930)\allowbreak{}0 | Xibar\_\allowbreak{}c(2930)\allowbreak{}0\par {\fontsize{8.1}{9.4}\selectfont PDG2026:\allowbreak{}observable:\allowbreak{}B156:\allowbreak{}B156A04:\allowbreak{}351264} & \textbf{2938.55+-0.30} MeV & \(u_+\): 0.301858\allowbreak{}05977630\allowbreak{}881\par \(u_-\): 0.301858\allowbreak{}05977630\allowbreak{}881\par Tipo: NONE\par Confianza: ND & COMP-BARYON\par E04 / M07\\\addlinespace[3pt]
% VI-CATALOG-ROW masas 114
\textbf{114. Xi\_\allowbreak{}c(3055)\allowbreak{}+ MASS}\par Xi\_\allowbreak{}c(3055)\allowbreak{}0 | Xibar\_\allowbreak{}c(3055)\allowbreak{}0\par {\fontsize{8.1}{9.4}\selectfont PDG2026:\allowbreak{}observable:\allowbreak{}B157:\allowbreak{}B157M+:\allowbreak{}351297} & \textbf{3055.2+-0.7} MeV & \(u_+\): 0.689992\allowbreak{}35826894\allowbreak{}623\par \(u_-\): 0.689992\allowbreak{}35826894\allowbreak{}623\par Tipo: NONE\par Confianza: ND & COMP-BARYON\par E04 / M07\\\addlinespace[3pt]
% VI-CATALOG-ROW masas 115
\textbf{115. Xi\_\allowbreak{}c(3055)\allowbreak{}0 MASS}\par Xi\_\allowbreak{}c(3055)\allowbreak{}0 | Xibar\_\allowbreak{}c(3055)\allowbreak{}0\par {\fontsize{8.1}{9.4}\selectfont PDG2026:\allowbreak{}observable:\allowbreak{}B157:\allowbreak{}B157M:\allowbreak{}351298} & \textbf{3061.0+-0.8} MeV & \(u_+\): 0.832406\allowbreak{}15086627\allowbreak{}034\par \(u_-\): 0.832406\allowbreak{}15086627\allowbreak{}034\par Tipo: NONE\par Confianza: ND & COMP-BARYON\par E04 / M07\\\addlinespace[3pt]
% VI-CATALOG-ROW masas 116
\textbf{116. Xi\_\allowbreak{}c(3123)\allowbreak{}+ MASS}\par Xi\_\allowbreak{}c(3123)\allowbreak{}0 | Xibar\_\allowbreak{}c(3123)\allowbreak{}0\par {\fontsize{8.1}{9.4}\selectfont PDG2026:\allowbreak{}observable:\allowbreak{}B158:\allowbreak{}B158M+:\allowbreak{}351325} & \textbf{3122.9+-1.3} MeV & \(u_+\): 1.334166\allowbreak{}40641263\allowbreak{}4\par \(u_-\): 1.334166\allowbreak{}40641263\allowbreak{}4\par Tipo: NONE\par Confianza: ND & COMP-BARYON\par E04 / M07\\\addlinespace[3pt]
% VI-CATALOG-ROW masas 117
\textbf{117. Xi\_\allowbreak{}b(5945)\allowbreak{}0 MASS}\par Xi\_\allowbreak{}b(5945)\allowbreak{}0 | Xibar\_\allowbreak{}b(5945)\allowbreak{}0\par {\fontsize{8.1}{9.4}\selectfont PDG2026:\allowbreak{}observable:\allowbreak{}B161:\allowbreak{}B161M:\allowbreak{}351844} & \textbf{5952.3+-0.6} MeV & \(u_+\): 0.577652\allowbreak{}72773906\allowbreak{}229\par \(u_-\): 0.577652\allowbreak{}72773906\allowbreak{}229\par Tipo: NONE\par Confianza: ND & COMP-BARYON\par E04 / M07\\\addlinespace[3pt]
% VI-CATALOG-ROW masas 118
\textbf{118. Lambda\_\allowbreak{}b(5912)\allowbreak{}0 MASS}\par Lambda\_\allowbreak{}b(5912)\allowbreak{}0 | Lambdabar\_\allowbreak{}b(5912)\allowbreak{}0\par {\fontsize{8.1}{9.4}\selectfont PDG2026:\allowbreak{}observable:\allowbreak{}B162:\allowbreak{}B162M:\allowbreak{}351694} & \textbf{5912.16+-0.16} MeV & \(u_+\): 0.160070\allowbreak{}03100330\allowbreak{}261\par \(u_-\): 0.160070\allowbreak{}03100330\allowbreak{}261\par Tipo: NONE\par Confianza: ND & COMP-BARYON\par E04 / M07\\\addlinespace[3pt]
% VI-CATALOG-ROW masas 119
\textbf{119. Lambda\_\allowbreak{}b(5920)\allowbreak{}0 MASS}\par Lambda\_\allowbreak{}b(5920)\allowbreak{}0 | Lambdabar\_\allowbreak{}b(5920)\allowbreak{}0\par {\fontsize{8.1}{9.4}\selectfont PDG2026:\allowbreak{}observable:\allowbreak{}B163:\allowbreak{}B163M:\allowbreak{}351698} & \textbf{5920.07+-0.16} MeV & \(u_+\): 0.158621\allowbreak{}70338385\allowbreak{}659\par \(u_-\): 0.158621\allowbreak{}70338385\allowbreak{}659\par Tipo: NONE\par Confianza: ND & COMP-BARYON\par E04 / M07\\\addlinespace[3pt]
% VI-CATALOG-ROW masas 120
\textbf{120. Lambda(1710) MASS}\par Lambda(1710)\allowbreak{}0 | Lambdabar(1710)\allowbreak{}0\par {\fontsize{8.1}{9.4}\selectfont PDG2026:\allowbreak{}observable:\allowbreak{}B164:\allowbreak{}B164M:\allowbreak{}349532} & \textbf{1713+-13} MeV & \(u_+\): 13\par \(u_-\): 13\par Tipo: NONE\par Confianza: ND & COMP-BARYON\par E04 / M07\\\addlinespace[3pt]
% VI-CATALOG-ROW masas 121
\textbf{121. Lambda(2050) MASS}\par Lambda(2050)\allowbreak{}0 | Lambdabar(2050)\allowbreak{}0\par {\fontsize{8.1}{9.4}\selectfont PDG2026:\allowbreak{}observable:\allowbreak{}B165:\allowbreak{}B165M:\allowbreak{}349651} & \textbf{2056+-22} MeV & \(u_+\): 22\par \(u_-\): 22\par Tipo: NONE\par Confianza: ND & COMP-BARYON\par E04 / M07\\\addlinespace[3pt]
% VI-CATALOG-ROW masas 122
\textbf{122. Sigma(1780) MASS}\par Sigma(1780)\allowbreak{}- | Sigmabar(1780)\allowbreak{}+ | Sigma(1780)\allowbreak{}0 | Sigmabar(1780)\allowbreak{}0 | Sigma(1780)\allowbreak{}+ | Sigmabar(1780)\allowbreak{}-\par {\fontsize{8.1}{9.4}\selectfont PDG2026:\allowbreak{}observable:\allowbreak{}B166:\allowbreak{}B166M:\allowbreak{}349991} & \textbf{1730 to 1780 to 1830} MeV\par Presentación: 1730 to 1830 (\textasciitilde{}1780) & \(u_+\): NA\par \(u_-\): NA\par Tipo: R\par Confianza: ND & COMP-BARYON\par E04 / M07\\\addlinespace[3pt]
% VI-CATALOG-ROW masas 123
\textbf{123. Sigma(1900) MASS}\par Sigma(1900)\allowbreak{}- | Sigmabar(1900)\allowbreak{}+ | Sigma(1900)\allowbreak{}0 | Sigmabar(1900)\allowbreak{}0 | Sigma(1900)\allowbreak{}+ | Sigmabar(1900)\allowbreak{}-\par {\fontsize{8.1}{9.4}\selectfont PDG2026:\allowbreak{}observable:\allowbreak{}B167:\allowbreak{}B167M:\allowbreak{}350019} & \textbf{1900 to 1925 to 1950} MeV\par Presentación: 1900 to 1950 (\textasciitilde{}1925) & \(u_+\): NA\par \(u_-\): NA\par Tipo: R\par Confianza: ND & COMP-BARYON\par E04 / M07\\\addlinespace[3pt]
% VI-CATALOG-ROW masas 124
\textbf{124. Sigma(1940) MASS}\par Sigma(1940)\allowbreak{}- | Sigmabar(1940)\allowbreak{}+ | Sigma(1940)\allowbreak{}0 | Sigmabar(1940)\allowbreak{}0 | Sigma(1940)\allowbreak{}+ | Sigmabar(1940)\allowbreak{}-\par {\fontsize{8.1}{9.4}\selectfont PDG2026:\allowbreak{}observable:\allowbreak{}B168:\allowbreak{}B168M:\allowbreak{}350102} & \textbf{1920 to 1940 to 1960} MeV\par Presentación: 1920 to 1960 (\textasciitilde{}1940) & \(u_+\): NA\par \(u_-\): NA\par Tipo: R\par Confianza: ND & COMP-BARYON\par E04 / M07\\\addlinespace[3pt]
% VI-CATALOG-ROW masas 125
\textbf{125. Xi\_\allowbreak{}b'(5935)\allowbreak{}- MASS}\par Xi\_\allowbreak{}b\textasciicircum{}'(5935)\allowbreak{}- | Xibar\_\allowbreak{}b\textasciicircum{}'(5935)\allowbreak{}+\par {\fontsize{8.1}{9.4}\selectfont PDG2026:\allowbreak{}observable:\allowbreak{}B169:\allowbreak{}B169M:\allowbreak{}351840} & \textbf{5934.9+-0.4} MeV & \(u_+\): 0.413524\allowbreak{}29707823\allowbreak{}54\par \(u_-\): 0.413524\allowbreak{}29707823\allowbreak{}54\par Tipo: NONE\par Confianza: ND & COMP-BARYON\par E04 / M07\\\addlinespace[3pt]
% VI-CATALOG-ROW masas 126
\textbf{126. Xi\_\allowbreak{}b(5955)\allowbreak{}- MASS}\par Xi\_\allowbreak{}b(5955)\allowbreak{}- | Xibar\_\allowbreak{}b(5955)\allowbreak{}+\par {\fontsize{8.1}{9.4}\selectfont PDG2026:\allowbreak{}observable:\allowbreak{}B170:\allowbreak{}B170M:\allowbreak{}351848} & \textbf{5955.5+-0.4} MeV & \(u_+\): 0.414573\allowbreak{}81951292\allowbreak{}599\par \(u_-\): 0.414573\allowbreak{}81951292\allowbreak{}599\par Tipo: NONE\par Confianza: ND & COMP-BARYON\par E04 / M07\\\addlinespace[3pt]
% VI-CATALOG-ROW masas 127
\textbf{127. P\_\allowbreak{}c\_\allowbreak{}cbar(4380)\allowbreak{}+ MASS}\par P\_\allowbreak{}c\_\allowbreak{}cbar(4380)\allowbreak{}+\par {\fontsize{8.1}{9.4}\selectfont PDG2026:\allowbreak{}observable:\allowbreak{}B171:\allowbreak{}B171M:\allowbreak{}351951} & \textbf{4380+-30} MeV & \(u_+\): 30.08321\allowbreak{}79129826\allowbreak{}51\par \(u_-\): 30.08321\allowbreak{}79129826\allowbreak{}51\par Tipo: NONE\par Confianza: ND & COMP-BARYON\par E04 / M07\\\addlinespace[3pt]
% VI-CATALOG-ROW masas 128
\textbf{128. P\_\allowbreak{}c\_\allowbreak{}cbar(4457)\allowbreak{}+ MASS}\par P\_\allowbreak{}c\_\allowbreak{}cbar(4457)\allowbreak{}+\par {\fontsize{8.1}{9.4}\selectfont PDG2026:\allowbreak{}observable:\allowbreak{}B172:\allowbreak{}B172M:\allowbreak{}351977} & \textbf{4457.3+4.0-1.8} MeV & \(u_+\): 4.143669\allowbreak{}87102013\allowbreak{}16\par \(u_-\): 1.802775\allowbreak{}63773199\allowbreak{}5\par Tipo: NONE\par Confianza: ND & COMP-BARYON\par E04 / M07\\\addlinespace[3pt]
% VI-CATALOG-ROW masas 129
\textbf{129. Omega\_\allowbreak{}c(3000)\allowbreak{}0 MASS}\par Omega\_\allowbreak{}c(3000)\allowbreak{}0 | Omegabar\_\allowbreak{}c(3000)\allowbreak{}0\par {\fontsize{8.1}{9.4}\selectfont PDG2026:\allowbreak{}observable:\allowbreak{}B173:\allowbreak{}B173M:\allowbreak{}351376} & \textbf{3000.46+-0.25} MeV & \(u_+\): 0.243176\allowbreak{}58361921\allowbreak{}44\par \(u_-\): 0.263996\allowbreak{}37894382\allowbreak{}499\par Tipo: NONE\par Confianza: ND & COMP-BARYON\par E04 / M07\\\addlinespace[3pt]
% VI-CATALOG-ROW masas 130
\textbf{130. Omega\_\allowbreak{}c(3050)\allowbreak{}0 MASS}\par Omega\_\allowbreak{}c(3050)\allowbreak{}0 | Omegabar\_\allowbreak{}c(3050)\allowbreak{}0\par {\fontsize{8.1}{9.4}\selectfont PDG2026:\allowbreak{}observable:\allowbreak{}B174:\allowbreak{}B174M:\allowbreak{}351380} & \textbf{3050.17+-0.19} MeV & \(u_+\): 0.188203\allowbreak{}56358563\allowbreak{}539\par \(u_-\): 0.190899\allowbreak{}50264325\allowbreak{}409\par Tipo: NONE\par Confianza: ND & COMP-BARYON\par E04 / M07\\\addlinespace[3pt]
% VI-CATALOG-ROW masas 131
\textbf{131. Omega\_\allowbreak{}c(3065)\allowbreak{}0 MASS}\par Omega\_\allowbreak{}c(3065)\allowbreak{}0 | Omegabar\_\allowbreak{}c(3065)\allowbreak{}0\par {\fontsize{8.1}{9.4}\selectfont PDG2026:\allowbreak{}observable:\allowbreak{}B175:\allowbreak{}B175M:\allowbreak{}351384} & \textbf{3065.58+-0.21} MeV & \(u_+\): 0.213457\allowbreak{}33449248\allowbreak{}469\par \(u_-\): 0.213914\allowbreak{}75011375\allowbreak{}18\par Tipo: NONE\par Confianza: ND & COMP-BARYON\par E04 / M07\\\addlinespace[3pt]
% VI-CATALOG-ROW masas 132
\textbf{132. Omega\_\allowbreak{}c(3090)\allowbreak{}0 MASS}\par Omega\_\allowbreak{}c(3090)\allowbreak{}0 | Omegabar\_\allowbreak{}c(3090)\allowbreak{}0\par {\fontsize{8.1}{9.4}\selectfont PDG2026:\allowbreak{}observable:\allowbreak{}B176:\allowbreak{}B176M:\allowbreak{}351388} & \textbf{3090.15+-0.26} MeV & \(u_+\): 0.252392\allowbreak{}09275859\allowbreak{}231\par \(u_-\): 0.263050\allowbreak{}43429981\allowbreak{}828\par Tipo: NONE\par Confianza: ND & COMP-BARYON\par E04 / M07\\\addlinespace[3pt]
% VI-CATALOG-ROW masas 133
\textbf{133. Omega\_\allowbreak{}c(3120)\allowbreak{}0 MASS}\par Omega\_\allowbreak{}c(3120)\allowbreak{}0 | Omegabar\_\allowbreak{}c(3120)\allowbreak{}0\par {\fontsize{8.1}{9.4}\selectfont PDG2026:\allowbreak{}observable:\allowbreak{}B177:\allowbreak{}B177M:\allowbreak{}351392} & \textbf{3118.98+0.27-0.35} MeV & \(u_+\): 0.274590\allowbreak{}60435491\allowbreak{}961\par \(u_-\): 0.346698\allowbreak{}71646719\allowbreak{}441\par Tipo: NONE\par Confianza: ND & COMP-BARYON\par E04 / M07\\\addlinespace[3pt]
% VI-CATALOG-ROW masas 134
\textbf{134. Lambda\_\allowbreak{}c(2860)\allowbreak{}+ MASS}\par Lambda\_\allowbreak{}c(2860)\allowbreak{}0 | Lambdabar\_\allowbreak{}c(2860)\allowbreak{}0\par {\fontsize{8.1}{9.4}\selectfont PDG2026:\allowbreak{}observable:\allowbreak{}B178:\allowbreak{}B178M:\allowbreak{}350956} & \textbf{2856.1+2.3-6.0} MeV & \(u_+\): 2.336664\allowbreak{}28910958\allowbreak{}5\par \(u_-\): 5.873670\allowbreak{}06223536\allowbreak{}5\par Tipo: NONE\par Confianza: ND & COMP-BARYON\par E04 / M07\\\addlinespace[3pt]
% VI-CATALOG-ROW masas 135
\textbf{135. Omega(2012)\allowbreak{}- MASS}\par Omega(2012)\allowbreak{}- | Omegabar(2012)\allowbreak{}+\par {\fontsize{8.1}{9.4}\selectfont PDG2026:\allowbreak{}observable:\allowbreak{}B179:\allowbreak{}B179M:\allowbreak{}350506} & \textbf{2012.9+-0.4} MeV & \(u_+\): 0.445362\allowbreak{}59664740\allowbreak{}679\par \(u_-\): 0.445362\allowbreak{}59664740\allowbreak{}679\par Tipo: NONE\par Confianza: ND & COMP-BARYON\par E04 / M07\\\addlinespace[3pt]
% VI-CATALOG-ROW masas 136
\textbf{136. Xi\_\allowbreak{}b(6227)\allowbreak{}- MASS}\par Xi\_\allowbreak{}b(6227)\allowbreak{}- | Xi\_\allowbreak{}b(6227)\allowbreak{}+\par {\fontsize{8.1}{9.4}\selectfont PDG2026:\allowbreak{}observable:\allowbreak{}B180:\allowbreak{}B180M:\allowbreak{}351864} & \textbf{6227.9+-0.9} MeV & \(u_+\): 0.908151\allowbreak{}28746094\allowbreak{}542\par \(u_-\): 0.908151\allowbreak{}28746094\allowbreak{}542\par Tipo: NONE\par Confianza: ND & COMP-BARYON\par E04 / M07\\\addlinespace[3pt]
% VI-CATALOG-ROW masas 137
\textbf{137. Sigma\_\allowbreak{}b(6097)\allowbreak{}+ MASS}\par Sigma\_\allowbreak{}b(6097)\allowbreak{}+ | Sigmabar\_\allowbreak{}b(6097)\allowbreak{}-\par {\fontsize{8.1}{9.4}\selectfont PDG2026:\allowbreak{}observable:\allowbreak{}B181:\allowbreak{}B181M:\allowbreak{}351734} & \textbf{6095.8+-1.7} MeV & \(u_+\): 1.746424\allowbreak{}91965729\allowbreak{}81\par \(u_-\): 1.746424\allowbreak{}91965729\allowbreak{}81\par Tipo: NONE\par Confianza: ND & COMP-BARYON\par E04 / M07\\\addlinespace[3pt]
% VI-CATALOG-ROW masas 138
\textbf{138. Sigma\_\allowbreak{}b(6097)\allowbreak{}- MASS}\par Sigma\_\allowbreak{}b(6097)\allowbreak{}- | Sigmabar\_\allowbreak{}b(6097)\allowbreak{}+\par {\fontsize{8.1}{9.4}\selectfont PDG2026:\allowbreak{}observable:\allowbreak{}B182:\allowbreak{}B182M:\allowbreak{}351738} & \textbf{6098.0+-1.8} MeV & \(u_+\): 1.772004\allowbreak{}51466693\allowbreak{}5\par \(u_-\): 1.772004\allowbreak{}51466693\allowbreak{}5\par Tipo: NONE\par Confianza: ND & COMP-BARYON\par E04 / M07\\\addlinespace[3pt]
% VI-CATALOG-ROW masas 139
\textbf{139. Sigma(2110) MASS}\par Sigma(2110)\allowbreak{}- | Sigmabar(2110)\allowbreak{}+ | Sigma(2110)\allowbreak{}0 | Sigmabar(2110)\allowbreak{}0 | Sigma(2110)\allowbreak{}+ | Sigmabar(2110)\allowbreak{}-\par {\fontsize{8.1}{9.4}\selectfont PDG2026:\allowbreak{}observable:\allowbreak{}B183:\allowbreak{}B183M:\allowbreak{}350228} & \textbf{2110+-50} MeV & \(u_+\): 51.99138\allowbreak{}75255207\allowbreak{}71\par \(u_-\): 51.99138\allowbreak{}75255207\allowbreak{}71\par Tipo: NONE\par Confianza: ND & COMP-BARYON\par E04 / M07\\\addlinespace[3pt]
% VI-CATALOG-ROW masas 140
\textbf{140. Sigma(2230) MASS}\par Sigma(2230)\allowbreak{}- | Sigmabar(2230)\allowbreak{}+ | Sigma(2230)\allowbreak{}0 | Sigmabar(2230)\allowbreak{}0 | Sigma(2230)\allowbreak{}+ | Sigmabar(2230)\allowbreak{}-\par {\fontsize{8.1}{9.4}\selectfont PDG2026:\allowbreak{}observable:\allowbreak{}B184:\allowbreak{}B184M:\allowbreak{}350253} & \textbf{2240+-27} MeV & \(u_+\): 27\par \(u_-\): 27\par Tipo: NONE\par Confianza: ND & COMP-BARYON\par E04 / M07\\\addlinespace[3pt]
% VI-CATALOG-ROW masas 141
\textbf{141. P\_\allowbreak{}c\_\allowbreak{}cbar(4312)\allowbreak{}+ MASS}\par P\_\allowbreak{}c\_\allowbreak{}cbar(4312)\allowbreak{}+\par {\fontsize{8.1}{9.4}\selectfont PDG2026:\allowbreak{}observable:\allowbreak{}B185:\allowbreak{}B185M:\allowbreak{}351935} & \textbf{4311.9+7.0-0.9} MeV & \(u_+\): 6.835934\allowbreak{}46428503\allowbreak{}97\par \(u_-\): 0.921954\allowbreak{}44572928\allowbreak{}864\par Tipo: NONE\par Confianza: ND & COMP-BARYON\par E04 / M07\\\addlinespace[3pt]
% VI-CATALOG-ROW masas 142
\textbf{142. P\_\allowbreak{}c\_\allowbreak{}cbar(4440)\allowbreak{}+ MASS}\par P\_\allowbreak{}c\_\allowbreak{}cbar(4440)\allowbreak{}+\par {\fontsize{8.1}{9.4}\selectfont PDG2026:\allowbreak{}observable:\allowbreak{}B186:\allowbreak{}B186M:\allowbreak{}351955} & \textbf{4440+4-5} MeV & \(u_+\): 4.301162\allowbreak{}63352131\allowbreak{}33\par \(u_-\): 4.876474\allowbreak{}13609464\allowbreak{}4\par Tipo: NONE\par Confianza: ND & COMP-BARYON\par E04 / M07\\\addlinespace[3pt]
% VI-CATALOG-ROW masas 143
\textbf{143. Lambda\_\allowbreak{}b(6146)\allowbreak{}0 MASS}\par Lambda\_\allowbreak{}b(6146)\allowbreak{}0 | Lambdabar\_\allowbreak{}b(6146)\allowbreak{}0\par {\fontsize{8.1}{9.4}\selectfont PDG2026:\allowbreak{}observable:\allowbreak{}B187:\allowbreak{}B187M:\allowbreak{}351706} & \textbf{6146.2+-0.4} MeV & \(u_+\): 0.417592\allowbreak{}68734261\allowbreak{}917\par \(u_-\): 0.417592\allowbreak{}68734261\allowbreak{}917\par Tipo: NONE\par Confianza: ND & COMP-BARYON\par E04 / M07\\\addlinespace[3pt]
% VI-CATALOG-ROW masas 144
\textbf{144. Lambda\_\allowbreak{}b(6152)\allowbreak{}0 MASS}\par Lambda\_\allowbreak{}b(6152)\allowbreak{}0 | Lambdabar\_\allowbreak{}b(6152)\allowbreak{}0\par {\fontsize{8.1}{9.4}\selectfont PDG2026:\allowbreak{}observable:\allowbreak{}B188:\allowbreak{}B188M:\allowbreak{}351711} & \textbf{6152.5+-0.4} MeV & \(u_+\): 0.357270\allowbreak{}06911778\allowbreak{}508\par \(u_-\): 0.357270\allowbreak{}06911778\allowbreak{}508\par Tipo: NONE\par Confianza: ND & COMP-BARYON\par E04 / M07\\\addlinespace[3pt]
% VI-CATALOG-ROW masas 145
\textbf{145. Lambda(2070) MASS}\par Lambda(2070)\allowbreak{}0 | Lambdabar(2070)\allowbreak{}0\par {\fontsize{8.1}{9.4}\selectfont PDG2026:\allowbreak{}observable:\allowbreak{}B190:\allowbreak{}B190M:\allowbreak{}349676} & \textbf{2070+-24} MeV & \(u_+\): 24\par \(u_-\): 24\par Tipo: NONE\par Confianza: ND & COMP-BARYON\par E04 / M07\\\addlinespace[3pt]
% VI-CATALOG-ROW masas 146
\textbf{146. Lambda(2080) MASS}\par Lambda(2080)\allowbreak{}0 | Lambdabar(2080)\allowbreak{}0\par {\fontsize{8.1}{9.4}\selectfont PDG2026:\allowbreak{}observable:\allowbreak{}B191:\allowbreak{}B191M:\allowbreak{}349711} & \textbf{2082+-13} MeV & \(u_+\): 13\par \(u_-\): 13\par Tipo: NONE\par Confianza: ND & COMP-BARYON\par E04 / M07\\\addlinespace[3pt]
% VI-CATALOG-ROW masas 147
\textbf{147. Xi\_\allowbreak{}c(2923)\allowbreak{}0 MASS}\par Xi\_\allowbreak{}c(2923)\allowbreak{}0 | Xibar\_\allowbreak{}c(2923)\allowbreak{}0\par {\fontsize{8.1}{9.4}\selectfont PDG2026:\allowbreak{}observable:\allowbreak{}B192:\allowbreak{}B192M0:\allowbreak{}351255} & \textbf{2923.2+-0.4} MeV & \(u_+\): 0.397437\allowbreak{}85839914\allowbreak{}991\par \(u_-\): 0.397437\allowbreak{}85839914\allowbreak{}991\par Tipo: NONE\par Confianza: ND & COMP-BARYON\par E04 / M07\\\addlinespace[3pt]
% VI-CATALOG-ROW masas 148
\textbf{148. Xi\_\allowbreak{}c(2923)\allowbreak{}+ MASS}\par Xi\_\allowbreak{}c(2923)\allowbreak{}0 | Xibar\_\allowbreak{}c(2923)\allowbreak{}0\par {\fontsize{8.1}{9.4}\selectfont PDG2026:\allowbreak{}observable:\allowbreak{}B192:\allowbreak{}B192MP:\allowbreak{}351256} & \textbf{2922.8+-0.6} MeV & \(u_+\): 0.616441\allowbreak{}40029689\allowbreak{}765\par \(u_-\): 0.616441\allowbreak{}40029689\allowbreak{}765\par Tipo: NONE\par Confianza: ND & COMP-BARYON\par E04 / M07\\\addlinespace[3pt]
% VI-CATALOG-ROW masas 149
\textbf{149. Omega\_\allowbreak{}b(6316)\allowbreak{}- MASS}\par Omega\_\allowbreak{}b(6316)\allowbreak{}- | Omegabar\_\allowbreak{}b(6316)\allowbreak{}+\par {\fontsize{8.1}{9.4}\selectfont PDG2026:\allowbreak{}observable:\allowbreak{}B193:\allowbreak{}B193M:\allowbreak{}351905} & \textbf{6315.6+-0.6} MeV & \(u_+\): 0.580487\allowbreak{}24929847\allowbreak{}628\par \(u_-\): 0.580487\allowbreak{}24929847\allowbreak{}628\par Tipo: NONE\par Confianza: ND & COMP-BARYON\par E04 / M07\\\addlinespace[3pt]
% VI-CATALOG-ROW masas 150
\textbf{150. Omega\_\allowbreak{}b(6330)\allowbreak{}- MASS}\par Omega\_\allowbreak{}b(6330)\allowbreak{}- | Omegabar\_\allowbreak{}b(6330)\allowbreak{}+\par {\fontsize{8.1}{9.4}\selectfont PDG2026:\allowbreak{}observable:\allowbreak{}B194:\allowbreak{}B194M:\allowbreak{}351909} & \textbf{6330.3+-0.6} MeV & \(u_+\): 0.565035\allowbreak{}79231594\allowbreak{}825\par \(u_-\): 0.565035\allowbreak{}79231594\allowbreak{}825\par Tipo: NONE\par Confianza: ND & COMP-BARYON\par E04 / M07\\\addlinespace[3pt]
% VI-CATALOG-ROW masas 151
\textbf{151. Omega\_\allowbreak{}b(6340)\allowbreak{}- MASS}\par Omega\_\allowbreak{}b(6340)\allowbreak{}- | Omegabar\_\allowbreak{}b(6340)\allowbreak{}+\par {\fontsize{8.1}{9.4}\selectfont PDG2026:\allowbreak{}observable:\allowbreak{}B195:\allowbreak{}B195M:\allowbreak{}351913} & \textbf{6339.7+-0.6} MeV & \(u_+\): 0.553231\allowbreak{}81994360\allowbreak{}315\par \(u_-\): 0.553231\allowbreak{}81994360\allowbreak{}315\par Tipo: NONE\par Confianza: ND & COMP-BARYON\par E04 / M07\\\addlinespace[3pt]
% VI-CATALOG-ROW masas 152
\textbf{152. Omega\_\allowbreak{}b(6350)\allowbreak{}- MASS}\par Omega\_\allowbreak{}b(6350)\allowbreak{}- | Omegabar\_\allowbreak{}b(6350)\allowbreak{}+\par {\fontsize{8.1}{9.4}\selectfont PDG2026:\allowbreak{}observable:\allowbreak{}B196:\allowbreak{}B196M:\allowbreak{}351917} & \textbf{6349.8+-0.6} MeV & \(u_+\): 0.600804\allowbreak{}00015155\allowbreak{}633\par \(u_-\): 0.600804\allowbreak{}00015155\allowbreak{}633\par Tipo: NONE\par Confianza: ND & COMP-BARYON\par E04 / M07\\\addlinespace[3pt]
% VI-CATALOG-ROW masas 153
\textbf{153. Lambda\_\allowbreak{}b(6070)\allowbreak{}0 MASS}\par Lambda\_\allowbreak{}b(6070)\allowbreak{}0 | Lambdabar\_\allowbreak{}b(6070)\allowbreak{}0\par {\fontsize{8.1}{9.4}\selectfont PDG2026:\allowbreak{}observable:\allowbreak{}B197:\allowbreak{}B197M:\allowbreak{}351702} & \textbf{6072.3+-2.9} MeV & \(u_+\): 2.946988\allowbreak{}08292754\allowbreak{}59\par \(u_-\): 2.946988\allowbreak{}08292754\allowbreak{}59\par Tipo: NONE\par Confianza: ND & COMP-BARYON\par E04 / M07\\\addlinespace[3pt]
% VI-CATALOG-ROW masas 154
\textbf{154. Xi\_\allowbreak{}b(6227)\allowbreak{}0 MASS}\par Xi\_\allowbreak{}b(6227)\allowbreak{}0 | Xibar\_\allowbreak{}b(6227)\allowbreak{}0\par {\fontsize{8.1}{9.4}\selectfont PDG2026:\allowbreak{}observable:\allowbreak{}B199:\allowbreak{}B199M:\allowbreak{}351870} & \textbf{6226.8+-1.6} MeV & \(u_+\): 1.543777\allowbreak{}84025284\allowbreak{}7\par \(u_-\): 1.635007\allowbreak{}65137529\allowbreak{}09\par Tipo: NONE\par Confianza: ND & COMP-BARYON\par E04 / M07\\\addlinespace[3pt]
% VI-CATALOG-ROW masas 155
\textbf{155. Xi\_\allowbreak{}b(6100)\allowbreak{}- MASS}\par Xi\_\allowbreak{}b(6100)\allowbreak{}- | Xibar\_\allowbreak{}b(6100)\allowbreak{}+\par {\fontsize{8.1}{9.4}\selectfont PDG2026:\allowbreak{}observable:\allowbreak{}B200:\allowbreak{}B200M:\allowbreak{}351860} & \textbf{6099.8+-0.4} MeV & \(u_+\): 0.448307\allowbreak{}69038414\allowbreak{}819\par \(u_-\): 0.448307\allowbreak{}69038414\allowbreak{}819\par Tipo: NONE\par Confianza: ND & COMP-BARYON\par E04 / M07\\\addlinespace[3pt]
% VI-CATALOG-ROW masas 156
\textbf{156. Xi\_\allowbreak{}b(6327)\allowbreak{}0 MASS}\par Xi\_\allowbreak{}b(6327)\allowbreak{}0 | Xibar\_\allowbreak{}b(6327)\allowbreak{}0\par {\fontsize{8.1}{9.4}\selectfont PDG2026:\allowbreak{}observable:\allowbreak{}B201:\allowbreak{}B201M:\allowbreak{}351874} & \textbf{6327.28+-0.35} MeV & \(u_+\): 0.354682\allowbreak{}95701936\allowbreak{}398\par \(u_-\): 0.342052\allowbreak{}62752974\allowbreak{}142\par Tipo: NONE\par Confianza: ND & COMP-BARYON\par E04 / M07\\\addlinespace[3pt]
% VI-CATALOG-ROW masas 157
\textbf{157. Xi\_\allowbreak{}b(6333)\allowbreak{}0 MASS}\par Xi\_\allowbreak{}b(6333)\allowbreak{}0 | Xibar\_\allowbreak{}b(6333)\allowbreak{}0\par {\fontsize{8.1}{9.4}\selectfont PDG2026:\allowbreak{}observable:\allowbreak{}B202:\allowbreak{}B202M:\allowbreak{}351878} & \textbf{6332.69+-0.28} MeV & \(u_+\): 0.278028\allowbreak{}77548915\allowbreak{}689\par \(u_-\): 0.284253\allowbreak{}40807103\allowbreak{}788\par Tipo: NONE\par Confianza: ND & COMP-BARYON\par E04 / M07\\\addlinespace[3pt]
% VI-CATALOG-ROW masas 158
\textbf{158. Xi\_\allowbreak{}c(2882) MASS}\par Xi\_\allowbreak{}c(2882)\allowbreak{}0 | Xibar\_\allowbreak{}c(2882)\allowbreak{}0\par {\fontsize{8.1}{9.4}\selectfont PDG2026:\allowbreak{}observable:\allowbreak{}B203:\allowbreak{}B203M0:\allowbreak{}351251} & \textbf{2882+-9} MeV & \(u_+\): 9.047651\allowbreak{}62901401\allowbreak{}79\par \(u_-\): 9.047651\allowbreak{}62901401\allowbreak{}79\par Tipo: NONE\par Confianza: ND & COMP-BARYON\par E04 / M07\\\addlinespace[3pt]
% VI-CATALOG-ROW masas 159
\textbf{159. P\_\allowbreak{}c\_\allowbreak{}cbar\_\allowbreak{}s(4338)\allowbreak{}0 MASS}\par P\_\allowbreak{}c\_\allowbreak{}cbar\_\allowbreak{}s(4338)\allowbreak{}+\par {\fontsize{8.1}{9.4}\selectfont PDG2026:\allowbreak{}observable:\allowbreak{}B204:\allowbreak{}B204M:\allowbreak{}351945} & \textbf{4338.2+-0.8} MeV & \(u_+\): 0.806225\allowbreak{}77482985\allowbreak{}491\par \(u_-\): 0.806225\allowbreak{}77482985\allowbreak{}491\par Tipo: NONE\par Confianza: ND & COMP-BARYON\par E04 / M07\\\addlinespace[3pt]
% VI-CATALOG-ROW masas 160
\textbf{160. P\_\allowbreak{}c\_\allowbreak{}cbar(4312)\allowbreak{}+ MASS}\par P\_\allowbreak{}c\_\allowbreak{}cbar(4312)\allowbreak{}+\par {\fontsize{8.1}{9.4}\selectfont PDG2026:\allowbreak{}observable:\allowbreak{}B205:\allowbreak{}B205M:\allowbreak{}351999} & \textbf{4466+6-5} MeV & \(u_+\): 6.393796\allowbreak{}24128860\allowbreak{}78\par \(u_-\): 4.587776\allowbreak{}54200301\allowbreak{}74\par Tipo: NONE\par Confianza: ND & COMP-BARYON\par E04 / M07\\\addlinespace[3pt]
% VI-CATALOG-ROW masas 161
\textbf{161. Lambda\_\allowbreak{}c(2910)\allowbreak{}+ MASS}\par Lambda\_\allowbreak{}c(2910)\allowbreak{}0 | Lambdabar\_\allowbreak{}c(2910)\allowbreak{}0\par {\fontsize{8.1}{9.4}\selectfont PDG2026:\allowbreak{}observable:\allowbreak{}B206:\allowbreak{}B206M:\allowbreak{}350973} & \textbf{2914+-7} MeV & \(u_+\): 6.767569\allowbreak{}72627545\allowbreak{}22\par \(u_-\): 6.767569\allowbreak{}72627545\allowbreak{}22\par Tipo: NONE\par Confianza: ND & COMP-BARYON\par E04 / M07\\\addlinespace[3pt]
% VI-CATALOG-ROW masas 162
\textbf{162. Xi\_\allowbreak{}b(6087)\allowbreak{}0 MASS}\par Xi\_\allowbreak{}b(6087)\par {\fontsize{8.1}{9.4}\selectfont PDG2026:\allowbreak{}observable:\allowbreak{}B207:\allowbreak{}B207M:\allowbreak{}351852} & \textbf{6087.0+-0.5} MeV & \(u_+\): 0.463167\allowbreak{}55368257\allowbreak{}7\par \(u_-\): 0.463167\allowbreak{}55368257\allowbreak{}7\par Tipo: NONE\par Confianza: ND & COMP-BARYON\par E04 / M07\\\addlinespace[3pt]
% VI-CATALOG-ROW masas 163
\textbf{163. Xi\_\allowbreak{}b(6095)\allowbreak{}0 MASS}\par Xi\_\allowbreak{}b(6095)\par {\fontsize{8.1}{9.4}\selectfont PDG2026:\allowbreak{}observable:\allowbreak{}B208:\allowbreak{}B208M:\allowbreak{}351856} & \textbf{6095.1+-0.4} MeV & \(u_+\): 0.440822\allowbreak{}16684770\allowbreak{}159\par \(u_-\): 0.440822\allowbreak{}16684770\allowbreak{}159\par Tipo: NONE\par Confianza: ND & COMP-BARYON\par E04 / M07\\\addlinespace[3pt]
% VI-CATALOG-ROW masas 164
\textbf{164. Omega\_\allowbreak{}c(3185)\allowbreak{}0 MASS}\par Omega\_\allowbreak{}c(3185)\allowbreak{}0 | Omegabar\_\allowbreak{}c(3185)\allowbreak{}0\par {\fontsize{8.1}{9.4}\selectfont PDG2026:\allowbreak{}observable:\allowbreak{}B209:\allowbreak{}B209M:\allowbreak{}351396} & \textbf{3185+7.6-1.9} MeV & \(u_+\): 7.595393\allowbreak{}34070329\allowbreak{}81\par \(u_-\): 1.933907\allowbreak{}96058137\allowbreak{}2\par Tipo: NONE\par Confianza: ND & COMP-BARYON\par E04 / M07\\\addlinespace[3pt]
% VI-CATALOG-ROW masas 165
\textbf{165. Omega\_\allowbreak{}c(3327)\allowbreak{}0 MASS}\par Omega\_\allowbreak{}c(3327)\allowbreak{}0 | Omegabar\_\allowbreak{}c(3327)\allowbreak{}0\par {\fontsize{8.1}{9.4}\selectfont PDG2026:\allowbreak{}observable:\allowbreak{}B210:\allowbreak{}B210M:\allowbreak{}351400} & \textbf{3327.1+1.2-1.8} MeV & \(u_+\): 1.220655\allowbreak{}56157337\allowbreak{}01\par \(u_-\): 1.780449\allowbreak{}38147648\allowbreak{}59\par Tipo: NONE\par Confianza: ND & COMP-BARYON\par E04 / M07\\\addlinespace[3pt]
% VI-CATALOG-ROW masas 166
\textbf{166. Omega(2109)\allowbreak{}- MASS}\par Omega(2109)\allowbreak{}- | Omegabar(2109)\allowbreak{}+\par {\fontsize{8.1}{9.4}\selectfont PDG2026:\allowbreak{}observable:\allowbreak{}B211:\allowbreak{}B211M:\allowbreak{}350514} & \textbf{2108+-5} MeV & \(u_+\): 5.277309\allowbreak{}92078350\allowbreak{}52\par \(u_-\): 5.277309\allowbreak{}92078350\allowbreak{}52\par Tipo: NONE\par Confianza: ND & COMP-BARYON\par E04 / M07\\\addlinespace[3pt]
% VI-CATALOG-ROW masas 167
\textbf{167. graviton MASS}\par graviton\par {\fontsize{8.1}{9.4}\selectfont PDG2026:\allowbreak{}observable:\allowbreak{}G033:\allowbreak{}G033M:\allowbreak{}332463} & \textbf{<1.76 E-23} eV\par Presentación: <1.76E-23 & \(u_+\): NA\par \(u_-\): NA\par Tipo: U\par Confianza: ND & GRAVIT\allowbreak{}ATIONA\allowbreak{}L\_\allowbreak{}OR\_\allowbreak{}HYPOTH\allowbreak{}ETICAL\par E04 / M10\\\addlinespace[3pt]
% VI-CATALOG-ROW masas 168
\textbf{168. omega(782) MASS}\par omega(782)\allowbreak{}0\par {\fontsize{8.1}{9.4}\selectfont PDG2026:\allowbreak{}observable:\allowbreak{}M001:\allowbreak{}M001M:\allowbreak{}333932} & \textbf{782.66+-0.13} MeV & \(u_+\): 0.129012\allowbreak{}68855075\allowbreak{}981\par \(u_-\): 0.128997\allowbreak{}04603475\allowbreak{}121\par Tipo: NONE\par Confianza: ND & COMP-MESON\par E04 / M08\\\addlinespace[3pt]
% VI-CATALOG-ROW masas 169
\textbf{169. eta'(958) MASS}\par eta\textasciicircum{}'(958)\allowbreak{}0\par {\fontsize{8.1}{9.4}\selectfont PDG2026:\allowbreak{}observable:\allowbreak{}M002:\allowbreak{}M002M:\allowbreak{}334006} & \textbf{957.78+-0.06} MeV & \(u_+\): 0.062631\allowbreak{}50681310\allowbreak{}4018\par \(u_-\): 0.062631\allowbreak{}50681310\allowbreak{}4018\par Tipo: NONE\par Confianza: ND & COMP-MESON\par E04 / M08\\\addlinespace[3pt]
% VI-CATALOG-ROW masas 170
\textbf{170. f\_\allowbreak{}0(980) MASS}\par f\_\allowbreak{}0(980)\allowbreak{}0\par {\fontsize{8.1}{9.4}\selectfont PDG2026:\allowbreak{}observable:\allowbreak{}M003:\allowbreak{}M003M1:\allowbreak{}334163} & \textbf{990+-20} MeV & \(u_+\): 20\par \(u_-\): 20\par Tipo: NONE\par Confianza: ND & COMP-MESON\par E04 / M08\\\addlinespace[3pt]
% VI-CATALOG-ROW masas 171
\textbf{171. phi(1020) MASS}\par phi(1020)\allowbreak{}0\par {\fontsize{8.1}{9.4}\selectfont PDG2026:\allowbreak{}observable:\allowbreak{}M004:\allowbreak{}M004M:\allowbreak{}334182} & \textbf{1019.460+-0.016} MeV & \(u_+\): 0.015565\allowbreak{}29663733\allowbreak{}44\par \(u_-\): 0.015565\allowbreak{}29663733\allowbreak{}44\par Tipo: NONE\par Confianza: ND & COMP-MESON\par E04 / M08\\\addlinespace[3pt]
% VI-CATALOG-ROW masas 172
\textbf{172. f\_\allowbreak{}2(1270) MASS}\par f\_\allowbreak{}2(1270)\allowbreak{}0\par {\fontsize{8.1}{9.4}\selectfont PDG2026:\allowbreak{}observable:\allowbreak{}M005:\allowbreak{}M005M:\allowbreak{}334354} & \textbf{1275.4+-0.8} MeV & \(u_+\): 0.832261\allowbreak{}60812299\allowbreak{}765\par \(u_-\): 0.832437\allowbreak{}20122632\allowbreak{}603\par Tipo: NONE\par Confianza: ND & COMP-MESON\par E04 / M08\\\addlinespace[3pt]
% VI-CATALOG-ROW masas 173
\textbf{173. f\_\allowbreak{}1(1420) MASS}\par f\_\allowbreak{}1(1420)\allowbreak{}0\par {\fontsize{8.1}{9.4}\selectfont PDG2026:\allowbreak{}observable:\allowbreak{}M006:\allowbreak{}M006M2:\allowbreak{}334554} & \textbf{1427.7+1.7-1.4} MeV & \(u_+\): 1.669286\allowbreak{}80164103\allowbreak{}2\par \(u_-\): 1.437322\allowbreak{}98578944\allowbreak{}89\par Tipo: NONE\par Confianza: ND & COMP-MESON\par E04 / M08\\\addlinespace[3pt]
% VI-CATALOG-ROW masas 174
\textbf{174. f\_\allowbreak{}1(1285) MASS}\par f\_\allowbreak{}1(1285)\allowbreak{}0\par {\fontsize{8.1}{9.4}\selectfont PDG2026:\allowbreak{}observable:\allowbreak{}M008:\allowbreak{}M008M:\allowbreak{}334403} & \textbf{1281.8+-0.5} MeV & \(u_+\): 0.457402\allowbreak{}51642395\allowbreak{}002\par \(u_-\): 0.450011\allowbreak{}00251184\allowbreak{}373\par Tipo: NONE\par Confianza: ND & COMP-MESON\par E04 / M08\\\addlinespace[3pt]
% VI-CATALOG-ROW masas 175
\textbf{175. NEUTRAL ONLY, e+ e-}\par rho(770)\allowbreak{}0 | rho(770)\allowbreak{}+ | rho(770)\allowbreak{}-\par {\fontsize{8.1}{9.4}\selectfont PDG2026:\allowbreak{}observable:\allowbreak{}M009:\allowbreak{}M009M0:\allowbreak{}333862} & \textbf{775.26+-0.23} MeV & \(u_+\): 0.234291\allowbreak{}87483587\allowbreak{}821\par \(u_-\): 0.234291\allowbreak{}87483587\allowbreak{}821\par Tipo: NONE\par Confianza: ND & COMP-MESON\par E04 / M08\\\addlinespace[3pt]
% VI-CATALOG-ROW masas 176
\textbf{176. CHARGED ONLY, tau DECAYS and e+ e-}\par rho(770)\allowbreak{}0 | rho(770)\allowbreak{}+ | rho(770)\allowbreak{}-\par {\fontsize{8.1}{9.4}\selectfont PDG2026:\allowbreak{}observable:\allowbreak{}M009:\allowbreak{}M009M5:\allowbreak{}333863} & \textbf{775.11+-0.34} MeV & \(u_+\): 0.340738\allowbreak{}77998562\allowbreak{}979\par \(u_-\): 0.340738\allowbreak{}77998562\allowbreak{}979\par Tipo: NONE\par Confianza: ND & COMP-MESON\par E04 / M08\\\addlinespace[3pt]
% VI-CATALOG-ROW masas 177
\textbf{177. MIXED CHARGES, OTHER REACTIONS}\par rho(770)\allowbreak{}0 | rho(770)\allowbreak{}+ | rho(770)\allowbreak{}-\par {\fontsize{8.1}{9.4}\selectfont PDG2026:\allowbreak{}observable:\allowbreak{}M009:\allowbreak{}M009M7:\allowbreak{}333864} & \textbf{763.0+-1.2} MeV & \(u_+\): 1.236931\allowbreak{}68768529\allowbreak{}8\par \(u_-\): 1.236931\allowbreak{}68768529\allowbreak{}8\par Tipo: NONE\par Confianza: ND & COMP-MESON\par E04 / M08\\\addlinespace[3pt]
% VI-CATALOG-ROW masas 178
\textbf{178. CHARGED ONLY, HADROPRODUCED}\par rho(770)\allowbreak{}0 | rho(770)\allowbreak{}+ | rho(770)\allowbreak{}-\par {\fontsize{8.1}{9.4}\selectfont PDG2026:\allowbreak{}observable:\allowbreak{}M009:\allowbreak{}M009M2:\allowbreak{}333865} & \textbf{766.5+-1.1} MeV & \(u_+\): 1.106853\allowbreak{}52334429\allowbreak{}6\par \(u_-\): 1.106853\allowbreak{}52334429\allowbreak{}6\par Tipo: NONE\par Confianza: ND & COMP-MESON\par E04 / M08\\\addlinespace[3pt]
% VI-CATALOG-ROW masas 179
\textbf{179. NEUTRAL ONLY, PHOTOPRODUCED}\par rho(770)\allowbreak{}0 | rho(770)\allowbreak{}+ | rho(770)\allowbreak{}-\par {\fontsize{8.1}{9.4}\selectfont PDG2026:\allowbreak{}observable:\allowbreak{}M009:\allowbreak{}M009M0P:\allowbreak{}333866} & \textbf{769.6+-0.8} MeV & \(u_+\): 0.847462\allowbreak{}60647359\allowbreak{}306\par \(u_-\): 0.828075\allowbreak{}25956483\allowbreak{}763\par Tipo: NONE\par Confianza: ND & COMP-MESON\par E04 / M08\\\addlinespace[3pt]
% VI-CATALOG-ROW masas 180
\textbf{180. NEUTRAL ONLY, OTHER REACTIONS}\par rho(770)\allowbreak{}0 | rho(770)\allowbreak{}+ | rho(770)\allowbreak{}-\par {\fontsize{8.1}{9.4}\selectfont PDG2026:\allowbreak{}observable:\allowbreak{}M009:\allowbreak{}M009M0R:\allowbreak{}333867} & \textbf{769.0+-0.9} MeV & \(u_+\): 0.850619\allowbreak{}98301285\allowbreak{}813\par \(u_-\): 0.850619\allowbreak{}98301285\allowbreak{}813\par Tipo: NONE\par Confianza: ND & COMP-MESON\par E04 / M08\\\addlinespace[3pt]
% VI-CATALOG-ROW masas 181
\textbf{181. a\_\allowbreak{}1(1260) MASS}\par a\_\allowbreak{}1(1260)\allowbreak{}0 | a\_\allowbreak{}1(1260)\allowbreak{}+ | a\_\allowbreak{}1(1260)\allowbreak{}-\par {\fontsize{8.1}{9.4}\selectfont PDG2026:\allowbreak{}observable:\allowbreak{}M010:\allowbreak{}M010M:\allowbreak{}334331} & \textbf{1230+-40} MeV & \(u_+\): 40\par \(u_-\): 40\par Tipo: NONE\par Confianza: ND & COMP-MESON\par E04 / M08\\\addlinespace[3pt]
% VI-CATALOG-ROW masas 182
\textbf{182. b\_\allowbreak{}1(1235) MASS}\par b\_\allowbreak{}1(1235)\allowbreak{}0 | b\_\allowbreak{}1(1235)\allowbreak{}+ | b\_\allowbreak{}1(1235)\allowbreak{}-\par {\fontsize{8.1}{9.4}\selectfont PDG2026:\allowbreak{}observable:\allowbreak{}M011:\allowbreak{}M011M:\allowbreak{}334309} & \textbf{1229.5+-3.2} MeV & \(u_+\): 3.169161\allowbreak{}29668438\allowbreak{}5\par \(u_-\): 3.169161\allowbreak{}29668438\allowbreak{}5\par Tipo: NONE\par Confianza: ND & COMP-MESON\par E04 / M08\\\addlinespace[3pt]
% VI-CATALOG-ROW masas 183
\textbf{183. a\_\allowbreak{}2(1320) MASS}\par a\_\allowbreak{}2(1320)\allowbreak{}0 | a\_\allowbreak{}2(1320)\allowbreak{}+ | a\_\allowbreak{}2(1320)\allowbreak{}-\par {\fontsize{8.1}{9.4}\selectfont PDG2026:\allowbreak{}observable:\allowbreak{}M012:\allowbreak{}M012M0:\allowbreak{}334469} & \textbf{1318.2+-0.6} MeV & \(u_+\): 0.609201\allowbreak{}86906481\allowbreak{}122\par \(u_-\): 0.606992\allowbreak{}92355954\allowbreak{}471\par Tipo: NONE\par Confianza: ND & COMP-MESON\par E04 / M08\\\addlinespace[3pt]
% VI-CATALOG-ROW masas 184
\textbf{184. 3pi MODE}\par a\_\allowbreak{}2(1320)\allowbreak{}0 | a\_\allowbreak{}2(1320)\allowbreak{}+ | a\_\allowbreak{}2(1320)\allowbreak{}-\par {\fontsize{8.1}{9.4}\selectfont PDG2026:\allowbreak{}observable:\allowbreak{}M012:\allowbreak{}M012M1:\allowbreak{}334470} & \textbf{1318.6+-1.3} MeV & \(u_+\): 1.325383\allowbreak{}45937344\allowbreak{}39\par \(u_-\): 1.308467\allowbreak{}95325698\allowbreak{}1\par Tipo: NONE\par Confianza: ND & COMP-MESON\par E04 / M08\\\addlinespace[3pt]
% VI-CATALOG-ROW masas 185
\textbf{185. K Kbar MODE}\par a\_\allowbreak{}2(1320)\allowbreak{}0 | a\_\allowbreak{}2(1320)\allowbreak{}+ | a\_\allowbreak{}2(1320)\allowbreak{}-\par {\fontsize{8.1}{9.4}\selectfont PDG2026:\allowbreak{}observable:\allowbreak{}M012:\allowbreak{}M012M2:\allowbreak{}334471} & \textbf{1318.1+-0.7} MeV & \(u_+\): 0.697768\allowbreak{}76905725\allowbreak{}704\par \(u_-\): 0.697768\allowbreak{}76905725\allowbreak{}704\par Tipo: NONE\par Confianza: ND & COMP-MESON\par E04 / M08\\\addlinespace[3pt]
% VI-CATALOG-ROW masas 186
\textbf{186. eta pi MODE}\par a\_\allowbreak{}2(1320)\allowbreak{}0 | a\_\allowbreak{}2(1320)\allowbreak{}+ | a\_\allowbreak{}2(1320)\allowbreak{}-\par {\fontsize{8.1}{9.4}\selectfont PDG2026:\allowbreak{}observable:\allowbreak{}M012:\allowbreak{}M012M3:\allowbreak{}334472} & \textbf{1317.7+-1.4} MeV & \(u_+\): 1.441827\allowbreak{}88627946\par \(u_-\): 1.441827\allowbreak{}88627946\par Tipo: NONE\par Confianza: ND & COMP-MESON\par E04 / M08\\\addlinespace[3pt]
% VI-CATALOG-ROW masas 187
\textbf{187. eta' pi MODE}\par a\_\allowbreak{}2(1320)\allowbreak{}0 | a\_\allowbreak{}2(1320)\allowbreak{}+ | a\_\allowbreak{}2(1320)\allowbreak{}-\par {\fontsize{8.1}{9.4}\selectfont PDG2026:\allowbreak{}observable:\allowbreak{}M012:\allowbreak{}M012M4:\allowbreak{}334473} & \textbf{1322+-7} MeV & \(u_+\): 6.676612\allowbreak{}22337908\allowbreak{}37\par \(u_-\): 7.076316\allowbreak{}43322348\allowbreak{}49\par Tipo: NONE\par Confianza: ND & COMP-MESON\par E04 / M08\\\addlinespace[3pt]
% VI-CATALOG-ROW masas 188
\textbf{188. f\_\allowbreak{}2'(1525) MASS}\par f\_\allowbreak{}2\textasciicircum{}'(1525)\allowbreak{}0\par {\fontsize{8.1}{9.4}\selectfont PDG2026:\allowbreak{}observable:\allowbreak{}M013:\allowbreak{}M013MX:\allowbreak{}334665} & \textbf{1517.3+-2.4} MeV & \(u_+\): 2.434665\allowbreak{}63247887\allowbreak{}4\par \(u_-\): 2.324152\allowbreak{}69962939\allowbreak{}91\par Tipo: NONE\par Confianza: ND & COMP-MESON\par E04 / M08\\\addlinespace[3pt]
% VI-CATALOG-ROW masas 189
\textbf{189. PRODUCED BY K+- BEAM}\par f\_\allowbreak{}2\textasciicircum{}'(1525)\allowbreak{}0\par {\fontsize{8.1}{9.4}\selectfont PDG2026:\allowbreak{}observable:\allowbreak{}M013:\allowbreak{}M013M2:\allowbreak{}334666} & \textbf{1518.0+-2.7} MeV & \(u_+\): 2.757522\allowbreak{}74044423\allowbreak{}6\par \(u_-\): 2.632632\allowbreak{}76117844\allowbreak{}79\par Tipo: NONE\par Confianza: ND & COMP-MESON\par E04 / M08\\\addlinespace[3pt]
% VI-CATALOG-ROW masas 190
\textbf{190. PRODUCED IN e+ e- ANNIHILATION AND PARTICLE DECAYS}\par f\_\allowbreak{}2\textasciicircum{}'(1525)\allowbreak{}0\par {\fontsize{8.1}{9.4}\selectfont PDG2026:\allowbreak{}observable:\allowbreak{}M013:\allowbreak{}M013M3:\allowbreak{}334667} & \textbf{1514+-4} MeV & \(u_+\): 4.424585\allowbreak{}98817588\allowbreak{}47\par \(u_-\): 4.089715\allowbreak{}45863226\allowbreak{}64\par Tipo: NONE\par Confianza: ND & COMP-MESON\par E04 / M08\\\addlinespace[3pt]
% VI-CATALOG-ROW masas 191
\textbf{191. PRODUCED IN pbar p ANNIHILATION}\par f\_\allowbreak{}2\textasciicircum{}'(1525)\allowbreak{}0\par {\fontsize{8.1}{9.4}\selectfont PDG2026:\allowbreak{}observable:\allowbreak{}M013:\allowbreak{}M013M9:\allowbreak{}334668} & \textbf{1512+-4} MeV & \(u_+\): 3.655246\allowbreak{}19448102\allowbreak{}91\par \(u_-\): 3.655246\allowbreak{}19448102\allowbreak{}91\par Tipo: NONE\par Confianza: ND & COMP-MESON\par E04 / M08\\\addlinespace[3pt]
% VI-CATALOG-ROW masas 192
\textbf{192. CENTRAL PRODUCTION}\par f\_\allowbreak{}2\textasciicircum{}'(1525)\allowbreak{}0\par {\fontsize{8.1}{9.4}\selectfont PDG2026:\allowbreak{}observable:\allowbreak{}M013:\allowbreak{}M013M4:\allowbreak{}334669} & \textbf{1515+-15} MeV & \(u_+\): 15\par \(u_-\): 15\par Tipo: NONE\par Confianza: ND & COMP-MESON\par E04 / M08\\\addlinespace[3pt]
% VI-CATALOG-ROW masas 193
\textbf{193. PRODUCED IN e p COLLISIONS}\par f\_\allowbreak{}2\textasciicircum{}'(1525)\allowbreak{}0\par {\fontsize{8.1}{9.4}\selectfont PDG2026:\allowbreak{}observable:\allowbreak{}M013:\allowbreak{}M013M10:\allowbreak{}334670} & \textbf{1512+-3.2} MeV & \(u_+\): 3.310589\allowbreak{}07144937\allowbreak{}02\par \(u_-\): 3.041381\allowbreak{}26514911\allowbreak{}01\par Tipo: NONE\par Confianza: ND & COMP-MESON\par E04 / M08\\\addlinespace[3pt]
% VI-CATALOG-ROW masas 194
\textbf{194. f\_\allowbreak{}0(500) BREIT-WIGNER MASS}\par f\_\allowbreak{}0(500)\allowbreak{}0\par {\fontsize{8.1}{9.4}\selectfont PDG2026:\allowbreak{}observable:\allowbreak{}M014:\allowbreak{}M014M:\allowbreak{}333857} & \textbf{400 to 800} MeV & \(u_+\): NA\par \(u_-\): NA\par Tipo: R\par Confianza: ND & COMP-MESON\par E04 / M08\\\addlinespace[3pt]
% VI-CATALOG-ROW masas 195
\textbf{195. rho\_\allowbreak{}3(1690) MASS}\par rho\_\allowbreak{}3(1690)\allowbreak{}0 | rho\_\allowbreak{}3(1690)\allowbreak{}+ | rho\_\allowbreak{}3(1690)\allowbreak{}-\par {\fontsize{8.1}{9.4}\selectfont PDG2026:\allowbreak{}observable:\allowbreak{}M015:\allowbreak{}M015M:\allowbreak{}334842} & \textbf{1688.8+-2.1} MeV & \(u_+\): 2.098551\allowbreak{}67236953\allowbreak{}61\par \(u_-\): 2.098551\allowbreak{}67236953\allowbreak{}61\par Tipo: NONE\par Confianza: ND & COMP-MESON\par E04 / M08\\\addlinespace[3pt]
% VI-CATALOG-ROW masas 196
\textbf{196. 2pi MODE}\par rho\_\allowbreak{}3(1690)\allowbreak{}0 | rho\_\allowbreak{}3(1690)\allowbreak{}+ | rho\_\allowbreak{}3(1690)\allowbreak{}-\par {\fontsize{8.1}{9.4}\selectfont PDG2026:\allowbreak{}observable:\allowbreak{}M015:\allowbreak{}M015M1:\allowbreak{}334843} & \textbf{1686+-4} MeV & \(u_+\): 3.975002\allowbreak{}51085426\allowbreak{}28\par \(u_-\): 3.975002\allowbreak{}51085426\allowbreak{}28\par Tipo: NONE\par Confianza: ND & COMP-MESON\par E04 / M08\\\addlinespace[3pt]
% VI-CATALOG-ROW masas 197
\textbf{197. K Kbar AND K Kbar pi MODES}\par rho\_\allowbreak{}3(1690)\allowbreak{}0 | rho\_\allowbreak{}3(1690)\allowbreak{}+ | rho\_\allowbreak{}3(1690)\allowbreak{}-\par {\fontsize{8.1}{9.4}\selectfont PDG2026:\allowbreak{}observable:\allowbreak{}M015:\allowbreak{}M015M2:\allowbreak{}334844} & \textbf{1696+-4} MeV & \(u_+\): 3.566236\allowbreak{}01326023\allowbreak{}2\par \(u_-\): 3.566236\allowbreak{}01326023\allowbreak{}2\par Tipo: NONE\par Confianza: ND & COMP-MESON\par E04 / M08\\\addlinespace[3pt]
% VI-CATALOG-ROW masas 198
\textbf{198. (4pi)\allowbreak{}+- MODE}\par rho\_\allowbreak{}3(1690)\allowbreak{}0 | rho\_\allowbreak{}3(1690)\allowbreak{}+ | rho\_\allowbreak{}3(1690)\allowbreak{}-\par {\fontsize{8.1}{9.4}\selectfont PDG2026:\allowbreak{}observable:\allowbreak{}M015:\allowbreak{}M015M3:\allowbreak{}334845} & \textbf{1686+-5} MeV & \(u_+\): 4.642848\allowbreak{}50304302\allowbreak{}04\par \(u_-\): 4.642848\allowbreak{}50304302\allowbreak{}04\par Tipo: NONE\par Confianza: ND & COMP-MESON\par E04 / M08\\\addlinespace[3pt]
% VI-CATALOG-ROW masas 199
\textbf{199. omega pi MODE}\par rho\_\allowbreak{}3(1690)\allowbreak{}0 | rho\_\allowbreak{}3(1690)\allowbreak{}+ | rho\_\allowbreak{}3(1690)\allowbreak{}-\par {\fontsize{8.1}{9.4}\selectfont PDG2026:\allowbreak{}observable:\allowbreak{}M015:\allowbreak{}M015M5:\allowbreak{}334846} & \textbf{1681+-7} MeV & \(u_+\): 6.524368\allowbreak{}54824307\allowbreak{}93\par \(u_-\): 6.524368\allowbreak{}54824307\allowbreak{}93\par Tipo: NONE\par Confianza: ND & COMP-MESON\par E04 / M08\\\addlinespace[3pt]
% VI-CATALOG-ROW masas 200
\textbf{200. eta pi+ pi- MODE}\par rho\_\allowbreak{}3(1690)\allowbreak{}0 | rho\_\allowbreak{}3(1690)\allowbreak{}+ | rho\_\allowbreak{}3(1690)\allowbreak{}-\par {\fontsize{8.1}{9.4}\selectfont PDG2026:\allowbreak{}observable:\allowbreak{}M015:\allowbreak{}M015M6:\allowbreak{}334847} & \textbf{1682+-12} MeV & \(u_+\): 12.45682\allowbreak{}19780610\allowbreak{}01\par \(u_-\): 12.45682\allowbreak{}19780610\allowbreak{}01\par Tipo: NONE\par Confianza: ND & COMP-MESON\par E04 / M08\\\addlinespace[3pt]
% VI-CATALOG-ROW masas 201
\textbf{201. f\_\allowbreak{}4(2050) MASS}\par f\_\allowbreak{}4(2050)\allowbreak{}0\par {\fontsize{8.1}{9.4}\selectfont PDG2026:\allowbreak{}observable:\allowbreak{}M016:\allowbreak{}M016M:\allowbreak{}335130} & \textbf{2018+-11} MeV & \(u_+\): 11.36985\allowbreak{}61841563\allowbreak{}4\par \(u_-\): 11.50562\allowbreak{}63269526\allowbreak{}21\par Tipo: NONE\par Confianza: ND & COMP-MESON\par E04 / M08\\\addlinespace[3pt]
% VI-CATALOG-ROW masas 202
\textbf{202. a\_\allowbreak{}4(1970) MASS}\par a\_\allowbreak{}4(1970)\allowbreak{}0 | a\_\allowbreak{}4(1970)\allowbreak{}+ | a\_\allowbreak{}4(1970)\allowbreak{}-\par {\fontsize{8.1}{9.4}\selectfont PDG2026:\allowbreak{}observable:\allowbreak{}M017:\allowbreak{}M017M:\allowbreak{}335090} & \textbf{1967+-16} MeV & \(u_+\): 16.03887\allowbreak{}62689347\allowbreak{}11\par \(u_-\): 15.89171\allowbreak{}26074592\allowbreak{}4\par Tipo: NONE\par Confianza: ND & COMP-MESON\par E04 / M08\\\addlinespace[3pt]
% VI-CATALOG-ROW masas 203
\textbf{203. K*(892) T-Matrix Pole sqrt(s)}\par K\textasciicircum{}*(892)\allowbreak{}0 | Kbar\textasciicircum{}*(892)\allowbreak{}0 | K\textasciicircum{}*(892)\allowbreak{}+ | K\textasciicircum{}*(892)\allowbreak{}-\par {\fontsize{8.1}{9.4}\selectfont PDG2026:\allowbreak{}observable:\allowbreak{}M018:\allowbreak{}M018TMP:\allowbreak{}335789} & \textbf{(890+-14) - i (26+-6)} MeV\par Presentación: 890+-14 & \(u_+\): ND\par \(u_-\): ND\par Tipo: NONE\par Confianza: ND & COMP-MESON\par E04 / M08\\\addlinespace[3pt]
% VI-CATALOG-ROW masas 204
\textbf{204. CHARGED ONLY, HADROPRODUCED}\par K\textasciicircum{}*(892)\allowbreak{}0 | Kbar\textasciicircum{}*(892)\allowbreak{}0 | K\textasciicircum{}*(892)\allowbreak{}+ | K\textasciicircum{}*(892)\allowbreak{}-\par {\fontsize{8.1}{9.4}\selectfont PDG2026:\allowbreak{}observable:\allowbreak{}M018:\allowbreak{}M018M1:\allowbreak{}335790} & \textbf{891.88+-0.23} MeV & \(u_+\): 0.231107\allowbreak{}22559041\allowbreak{}66\par \(u_-\): 0.231107\allowbreak{}22559041\allowbreak{}66\par Tipo: NONE\par Confianza: ND & COMP-MESON\par E04 / M08\\\addlinespace[3pt]
% VI-CATALOG-ROW masas 205
\textbf{205. CHARGED ONLY, PRODUCED IN tau LEPTON DECAYS}\par K\textasciicircum{}*(892)\allowbreak{}0 | Kbar\textasciicircum{}*(892)\allowbreak{}0 | K\textasciicircum{}*(892)\allowbreak{}+ | K\textasciicircum{}*(892)\allowbreak{}-\par {\fontsize{8.1}{9.4}\selectfont PDG2026:\allowbreak{}observable:\allowbreak{}M018:\allowbreak{}M018MCT:\allowbreak{}335791} & \textbf{895.5+-0.8} MeV & \(u_+\): 0.766550\allowbreak{}71586947\allowbreak{}199\par \(u_-\): 0.766550\allowbreak{}71586947\allowbreak{}199\par Tipo: NONE\par Confianza: ND & COMP-MESON\par E04 / M08\\\addlinespace[3pt]
% VI-CATALOG-ROW masas 206
\textbf{206. NEUTRAL ONLY}\par K\textasciicircum{}*(892)\allowbreak{}0 | Kbar\textasciicircum{}*(892)\allowbreak{}0 | K\textasciicircum{}*(892)\allowbreak{}+ | K\textasciicircum{}*(892)\allowbreak{}-\par {\fontsize{8.1}{9.4}\selectfont PDG2026:\allowbreak{}observable:\allowbreak{}M018:\allowbreak{}M018M2:\allowbreak{}335792} & \textbf{895.56+-0.20} MeV & \(u_+\): 0.203579\allowbreak{}38811155\allowbreak{}91\par \(u_-\): 0.205063\allowbreak{}35702397\allowbreak{}691\par Tipo: NONE\par Confianza: ND & COMP-MESON\par E04 / M08\\\addlinespace[3pt]
% VI-CATALOG-ROW masas 207
\textbf{207. K\_\allowbreak{}0*(1430) MASS}\par K\_\allowbreak{}0\textasciicircum{}*(1430)\allowbreak{}0 | Kbar\_\allowbreak{}0\textasciicircum{}*(1430)\allowbreak{}0 | K\_\allowbreak{}0\textasciicircum{}*(1430)\allowbreak{}+ | K\_\allowbreak{}0\textasciicircum{}*(1430)\allowbreak{}-\par {\fontsize{8.1}{9.4}\selectfont PDG2026:\allowbreak{}observable:\allowbreak{}M019:\allowbreak{}M019M:\allowbreak{}335878} & \textbf{1425+-50} MeV & \(u_+\): 50\par \(u_-\): 50\par Tipo: NONE\par Confianza: ND & COMP-MESON\par E04 / M08\\\addlinespace[3pt]
% VI-CATALOG-ROW masas 208
\textbf{208. pi\_\allowbreak{}2(2100) MASS}\par pi\_\allowbreak{}2(2100)\allowbreak{}0 | pi\_\allowbreak{}2(2100)\allowbreak{}+ | pi\_\allowbreak{}2(2100)\allowbreak{}-\par {\fontsize{8.1}{9.4}\selectfont PDG2026:\allowbreak{}observable:\allowbreak{}M020:\allowbreak{}M020M:\allowbreak{}335146} & \textbf{2090+-29} MeV & \(u_+\): 29.41742\allowbreak{}02707276\par \(u_-\): 29.41742\allowbreak{}02707276\par Tipo: NONE\par Confianza: ND & COMP-MESON\par E04 / M08\\\addlinespace[3pt]
% VI-CATALOG-ROW masas 209
\textbf{209. CHARGED ONLY, WITH FINAL STATE K pi}\par K\_\allowbreak{}2\textasciicircum{}*(1430)\allowbreak{}0 | Kbar\_\allowbreak{}2\textasciicircum{}*(1430)\allowbreak{}0 | K\_\allowbreak{}2\textasciicircum{}*(1430)\allowbreak{}+ | K\_\allowbreak{}2\textasciicircum{}*(1430)\allowbreak{}-\par {\fontsize{8.1}{9.4}\selectfont PDG2026:\allowbreak{}observable:\allowbreak{}M022:\allowbreak{}M022M1:\allowbreak{}335888} & \textbf{1427.3+-1.5} MeV & \(u_+\): 1.503065\allowbreak{}01170133\par \(u_-\): 1.517646\allowbreak{}48227118\allowbreak{}19\par Tipo: NONE\par Confianza: ND & COMP-MESON\par E04 / M08\\\addlinespace[3pt]
% VI-CATALOG-ROW masas 210
\textbf{210. NEUTRAL ONLY}\par K\_\allowbreak{}2\textasciicircum{}*(1430)\allowbreak{}0 | Kbar\_\allowbreak{}2\textasciicircum{}*(1430)\allowbreak{}0 | K\_\allowbreak{}2\textasciicircum{}*(1430)\allowbreak{}+ | K\_\allowbreak{}2\textasciicircum{}*(1430)\allowbreak{}-\par {\fontsize{8.1}{9.4}\selectfont PDG2026:\allowbreak{}observable:\allowbreak{}M022:\allowbreak{}M022M4:\allowbreak{}335889} & \textbf{1432.4+-1.3} MeV & \(u_+\): 1.257694\allowbreak{}08082424\allowbreak{}21\par \(u_-\): 1.257694\allowbreak{}08082424\allowbreak{}21\par Tipo: NONE\par Confianza: ND & COMP-MESON\par E04 / M08\\\addlinespace[3pt]
% VI-CATALOG-ROW masas 211
\textbf{211. K\_\allowbreak{}2(1770) MASS}\par K\_\allowbreak{}2(1770)\allowbreak{}0 | Kbar\_\allowbreak{}2(1770)\allowbreak{}0 | K\_\allowbreak{}2(1770)\allowbreak{}+ | K\_\allowbreak{}2(1770)\allowbreak{}-\par {\fontsize{8.1}{9.4}\selectfont PDG2026:\allowbreak{}observable:\allowbreak{}M023:\allowbreak{}M023M:\allowbreak{}335963} & \textbf{1773+-8} MeV & \(u_+\): 7.984155\allowbreak{}45885385\allowbreak{}04\par \(u_-\): 7.964454\allowbreak{}14241581\allowbreak{}86\par Tipo: NONE\par Confianza: ND & COMP-MESON\par E04 / M08\\\addlinespace[3pt]
% VI-CATALOG-ROW masas 212
\textbf{212. psi(4160) MASS}\par psi(4160)\par {\fontsize{8.1}{9.4}\selectfont PDG2026:\allowbreak{}observable:\allowbreak{}M025:\allowbreak{}M025M:\allowbreak{}346343} & \textbf{4191+-5} MeV & \(u_+\): 5.195998\allowbreak{}33327111\allowbreak{}33\par \(u_-\): 4.980208\allowbreak{}46264291\allowbreak{}08\par Tipo: NONE\par Confianza: ND & COMP-MESON\par E04 / M08\\\addlinespace[3pt]
% VI-CATALOG-ROW masas 213
\textbf{213. eta\_\allowbreak{}c(1S) MASS}\par eta\_\allowbreak{}c(1S)\par {\fontsize{8.1}{9.4}\selectfont PDG2026:\allowbreak{}observable:\allowbreak{}M026:\allowbreak{}M026M:\allowbreak{}343286} & \textbf{2984.09+-0.30} MeV & \(u_+\): 0.296771\allowbreak{}04771232\allowbreak{}732\par \(u_-\): 0.300210\allowbreak{}31998551\allowbreak{}229\par Tipo: NONE\par Confianza: ND & COMP-MESON\par E04 / M08\\\addlinespace[3pt]
% VI-CATALOG-ROW masas 214
\textbf{214. eta(1405) MASS}\par eta(1405)\allowbreak{}0\par {\fontsize{8.1}{9.4}\selectfont PDG2026:\allowbreak{}observable:\allowbreak{}M027:\allowbreak{}M027MX:\allowbreak{}334525} & \textbf{1408.0+1.8-1.2} MeV & \(u_+\): 1.817274\allowbreak{}60178494\allowbreak{}5\par \(u_-\): 1.231728\allowbreak{}53976773\allowbreak{}8\par Tipo: NONE\par Confianza: ND & COMP-MESON\par E04 / M08\\\addlinespace[3pt]
% VI-CATALOG-ROW masas 215
\textbf{215. eta pi pi MODE}\par eta(1405)\allowbreak{}0\par {\fontsize{8.1}{9.4}\selectfont PDG2026:\allowbreak{}observable:\allowbreak{}M027:\allowbreak{}M027M1:\allowbreak{}334526} & \textbf{1405.5+-2.3} MeV & \(u_+\): 2.261740\allowbreak{}19480541\allowbreak{}72\par \(u_-\): 2.353950\allowbreak{}70110280\allowbreak{}8\par Tipo: NONE\par Confianza: ND & COMP-MESON\par E04 / M08\\\addlinespace[3pt]
% VI-CATALOG-ROW masas 216
\textbf{216. K Kbar pi MODE (a\_\allowbreak{}0(980) pi or direct K Kbar pi)}\par eta(1405)\allowbreak{}0\par {\fontsize{8.1}{9.4}\selectfont PDG2026:\allowbreak{}observable:\allowbreak{}M027:\allowbreak{}M027M4:\allowbreak{}334527} & \textbf{1413.5+2.0-0.9} MeV & \(u_+\): 1.965228\allowbreak{}32551039\allowbreak{}79\par \(u_-\): 0.888644\allowbreak{}61966865\allowbreak{}371\par Tipo: NONE\par Confianza: ND & COMP-MESON\par E04 / M08\\\addlinespace[3pt]
% VI-CATALOG-ROW masas 217
\textbf{217. pi pi gamma MODE}\par eta(1405)\allowbreak{}0\par {\fontsize{8.1}{9.4}\selectfont PDG2026:\allowbreak{}observable:\allowbreak{}M027:\allowbreak{}M027M2:\allowbreak{}334528} & \textbf{1403+-17} MeV & \(u_+\): 16.61741\allowbreak{}43278811\allowbreak{}29\par \(u_-\): 16.61741\allowbreak{}43278811\allowbreak{}29\par Tipo: NONE\par Confianza: ND & COMP-MESON\par E04 / M08\\\addlinespace[3pt]
% VI-CATALOG-ROW masas 218
\textbf{218. phi gamma MODE}\par eta(1405)\allowbreak{}0\par {\fontsize{8.1}{9.4}\selectfont PDG2026:\allowbreak{}observable:\allowbreak{}M027:\allowbreak{}M027M7:\allowbreak{}334529} & \textbf{1422+6-8} MeV & \(u_+\): 6.262587\allowbreak{}32474047\allowbreak{}09\par \(u_-\): 8.077747\allowbreak{}21070175\allowbreak{}52\par Tipo: NONE\par Confianza: ND & COMP-MESON\par E04 / M08\\\addlinespace[3pt]
% VI-CATALOG-ROW masas 219
\textbf{219. K\_\allowbreak{}1(1270) MASS}\par K\_\allowbreak{}1(1270)\allowbreak{}0 | Kbar\_\allowbreak{}1(1270)\allowbreak{}0 | K\_\allowbreak{}1(1270)\allowbreak{}+ | K\_\allowbreak{}1(1270)\allowbreak{}-\par {\fontsize{8.1}{9.4}\selectfont PDG2026:\allowbreak{}observable:\allowbreak{}M028:\allowbreak{}M028MX:\allowbreak{}335817} & \textbf{1256+-6} MeV & \(u_+\): 6.423982\allowbreak{}47602787\allowbreak{}36\par \(u_-\): 6.423982\allowbreak{}47602787\allowbreak{}36\par Tipo: NONE\par Confianza: ND & COMP-MESON\par E04 / M08\\\addlinespace[3pt]
% VI-CATALOG-ROW masas 220
\textbf{220. PRODUCED BY K-, BACKWARD SCATTERING, HYPERON EXCHANGE}\par K\_\allowbreak{}1(1270)\allowbreak{}0 | Kbar\_\allowbreak{}1(1270)\allowbreak{}0 | K\_\allowbreak{}1(1270)\allowbreak{}+ | K\_\allowbreak{}1(1270)\allowbreak{}-\par {\fontsize{8.1}{9.4}\selectfont PDG2026:\allowbreak{}observable:\allowbreak{}M028:\allowbreak{}M028M2:\allowbreak{}335818} & \textbf{1275+-10} MeV & \(u_+\): 10\par \(u_-\): 10\par Tipo: NONE\par Confianza: ND & COMP-MESON\par E04 / M08\\\addlinespace[3pt]
% VI-CATALOG-ROW masas 221
\textbf{221. PRODUCED BY K BEAMS}\par K\_\allowbreak{}1(1270)\allowbreak{}0 | Kbar\_\allowbreak{}1(1270)\allowbreak{}0 | K\_\allowbreak{}1(1270)\allowbreak{}+ | K\_\allowbreak{}1(1270)\allowbreak{}-\par {\fontsize{8.1}{9.4}\selectfont PDG2026:\allowbreak{}observable:\allowbreak{}M028:\allowbreak{}M028M3:\allowbreak{}335819} & \textbf{1270+-10} MeV & \(u_+\): 10\par \(u_-\): 10\par Tipo: NONE\par Confianza: ND & COMP-MESON\par E04 / M08\\\addlinespace[3pt]
% VI-CATALOG-ROW masas 222
\textbf{222. PRODUCED BY BEAMS OTHER THAN K MESONS}\par K\_\allowbreak{}1(1270)\allowbreak{}0 | Kbar\_\allowbreak{}1(1270)\allowbreak{}0 | K\_\allowbreak{}1(1270)\allowbreak{}+ | K\_\allowbreak{}1(1270)\allowbreak{}-\par {\fontsize{8.1}{9.4}\selectfont PDG2026:\allowbreak{}observable:\allowbreak{}M028:\allowbreak{}M028M1:\allowbreak{}335820} & \textbf{1252+-9} MeV & \(u_+\): 8.823912\allowbreak{}33586738\allowbreak{}19\par \(u_-\): 8.823912\allowbreak{}33586738\allowbreak{}19\par Tipo: NONE\par Confianza: ND & COMP-MESON\par E04 / M08\\\addlinespace[3pt]
% VI-CATALOG-ROW masas 223
\textbf{223. PRODUCED IN tau LEPTON DECAYS}\par K\_\allowbreak{}1(1270)\allowbreak{}0 | Kbar\_\allowbreak{}1(1270)\allowbreak{}0 | K\_\allowbreak{}1(1270)\allowbreak{}+ | K\_\allowbreak{}1(1270)\allowbreak{}-\par {\fontsize{8.1}{9.4}\selectfont PDG2026:\allowbreak{}observable:\allowbreak{}M028:\allowbreak{}M028MT:\allowbreak{}335821} & \textbf{1250+-50} MeV & \(u_+\): 47.38143\allowbreak{}09619285\allowbreak{}37\par \(u_-\): 47.38143\allowbreak{}09619285\allowbreak{}37\par Tipo: NONE\par Confianza: ND & COMP-MESON\par E04 / M08\\\addlinespace[3pt]
% VI-CATALOG-ROW masas 224
\textbf{224. X(3940) MASS}\par X(3940)\par {\fontsize{8.1}{9.4}\selectfont PDG2026:\allowbreak{}observable:\allowbreak{}M029:\allowbreak{}M029M:\allowbreak{}346244} & \textbf{3942+-9} MeV & \(u_+\): 9.042768\allowbreak{}99710160\allowbreak{}19\par \(u_-\): 8.214163\allowbreak{}93085561\allowbreak{}76\par Tipo: NONE\par Confianza: ND & COMP-MESON\par E04 / M08\\\addlinespace[3pt]
% VI-CATALOG-ROW masas 225
\textbf{225. h\_\allowbreak{}1(1170) MASS}\par h\_\allowbreak{}1(1170)\allowbreak{}0\par {\fontsize{8.1}{9.4}\selectfont PDG2026:\allowbreak{}observable:\allowbreak{}M030:\allowbreak{}M030M:\allowbreak{}334306} & \textbf{1166+-6} MeV & \(u_+\): 5.830951\allowbreak{}89484530\allowbreak{}07\par \(u_-\): 5.830951\allowbreak{}89484530\allowbreak{}07\par Tipo: NONE\par Confianza: ND & COMP-MESON\par E04 / M08\\\addlinespace[3pt]
% VI-CATALOG-ROW masas 226
\textbf{226. pi- p -\kern0pt{}-> omega pi0 n}\par rho\_\allowbreak{}5(2350)\allowbreak{}0 | rho\_\allowbreak{}5(2350)\allowbreak{}+ | rho\_\allowbreak{}5(2350)\allowbreak{}-\par {\fontsize{8.1}{9.4}\selectfont PDG2026:\allowbreak{}observable:\allowbreak{}M033:\allowbreak{}M033M3:\allowbreak{}335259} & \textbf{2330+-35} MeV & \(u_+\): 35\par \(u_-\): 35\par Tipo: NONE\par Confianza: ND & COMP-MESON\par E04 / M08\\\addlinespace[3pt]
% VI-CATALOG-ROW masas 227
\textbf{227. pi\_\allowbreak{}2(1670) MASS}\par pi\_\allowbreak{}2(1670)\allowbreak{}0 | pi\_\allowbreak{}2(1670)\allowbreak{}+ | pi\_\allowbreak{}2(1670)\allowbreak{}-\par {\fontsize{8.1}{9.4}\selectfont PDG2026:\allowbreak{}observable:\allowbreak{}M034:\allowbreak{}M034M:\allowbreak{}334783} & \textbf{1670.6+2.9-1.2} MeV & \(u_+\): 2.924872\allowbreak{}85188664\par \(u_-\): 1.198769\allowbreak{}20418300\allowbreak{}89\par Tipo: NONE\par Confianza: ND & COMP-MESON\par E04 / M08\\\addlinespace[3pt]
% VI-CATALOG-ROW masas 228
\textbf{228. K\_\allowbreak{}4*(2045) MASS}\par K\_\allowbreak{}4\textasciicircum{}*(2045)\allowbreak{}0 | Kbar\_\allowbreak{}4\textasciicircum{}*(2045)\allowbreak{}0 | K\_\allowbreak{}4\textasciicircum{}*(2045)\allowbreak{}+ | K\_\allowbreak{}4\textasciicircum{}*(2045)\allowbreak{}-\par {\fontsize{8.1}{9.4}\selectfont PDG2026:\allowbreak{}observable:\allowbreak{}M035:\allowbreak{}M035M:\allowbreak{}336016} & \textbf{2048+8-9} MeV & \(u_+\): 8.064465\allowbreak{}39721915\allowbreak{}38\par \(u_-\): 8.948340\allowbreak{}05242208\allowbreak{}26\par Tipo: NONE\par Confianza: ND & COMP-MESON\par E04 / M08\\\addlinespace[3pt]
% VI-CATALOG-ROW masas 229
\textbf{229. a\_\allowbreak{}0(980) MASS}\par a\_\allowbreak{}0(980)\allowbreak{}0 | a\_\allowbreak{}0(980)\allowbreak{}+ | a\_\allowbreak{}0(980)\allowbreak{}-\par {\fontsize{8.1}{9.4}\selectfont PDG2026:\allowbreak{}observable:\allowbreak{}M036:\allowbreak{}M036MX:\allowbreak{}334171} & \textbf{980+-20} MeV & \(u_+\): 20\par \(u_-\): 20\par Tipo: NONE\par Confianza: ND & COMP-MESON\par E04 / M08\\\addlinespace[3pt]
% VI-CATALOG-ROW masas 230
\textbf{230. eta(1295) MASS}\par eta(1295)\allowbreak{}0\par {\fontsize{8.1}{9.4}\selectfont PDG2026:\allowbreak{}observable:\allowbreak{}M037:\allowbreak{}M037M:\allowbreak{}334451} & \textbf{1294+-4} MeV & \(u_+\): 3.892110\allowbreak{}11433496\allowbreak{}7\par \(u_-\): 3.892110\allowbreak{}11433496\allowbreak{}7\par Tipo: NONE\par Confianza: ND & COMP-MESON\par E04 / M08\\\addlinespace[3pt]
% VI-CATALOG-ROW masas 231
\textbf{231. f\_\allowbreak{}2(1810) MASS}\par f\_\allowbreak{}2(1810)\allowbreak{}0\par {\fontsize{8.1}{9.4}\selectfont PDG2026:\allowbreak{}observable:\allowbreak{}M038:\allowbreak{}M038M:\allowbreak{}335002} & \textbf{1815+-12} MeV & \(u_+\): 12.52864\allowbreak{}22902590\allowbreak{}8\par \(u_-\): 11.98037\allowbreak{}40701169\par Tipo: NONE\par Confianza: ND & COMP-MESON\par E04 / M08\\\addlinespace[3pt]
% VI-CATALOG-ROW masas 232
\textbf{232. K\_\allowbreak{}2(2250) MASS}\par K\_\allowbreak{}2(2250)\allowbreak{}0 | Kbar\_\allowbreak{}2(2250)\allowbreak{}0 | K\_\allowbreak{}2(2250)\allowbreak{}+ | K\_\allowbreak{}2(2250)\allowbreak{}-\par {\fontsize{8.1}{9.4}\selectfont PDG2026:\allowbreak{}observable:\allowbreak{}M040:\allowbreak{}M040M:\allowbreak{}336032} & \textbf{2247+-17} MeV & \(u_+\): 16.84303\allowbreak{}84213303\allowbreak{}8\par \(u_-\): 16.84303\allowbreak{}84213303\allowbreak{}8\par Tipo: NONE\par Confianza: ND & COMP-MESON\par E04 / M08\\\addlinespace[3pt]
% VI-CATALOG-ROW masas 233
\textbf{233. p p CENTRAL PRODUCTION}\par f\_\allowbreak{}4(2300)\allowbreak{}0\par {\fontsize{8.1}{9.4}\selectfont PDG2026:\allowbreak{}observable:\allowbreak{}M041:\allowbreak{}M041M4:\allowbreak{}335244} & \textbf{2320+-60} MeV & \(u_+\): 60\par \(u_-\): 60\par Tipo: NONE\par Confianza: ND & COMP-MESON\par E04 / M08\\\addlinespace[3pt]
% VI-CATALOG-ROW masas 234
\textbf{234. f\_\allowbreak{}2(2150) MASS, COMBINED MODES}\par f\_\allowbreak{}2(2150)\allowbreak{}0\par {\fontsize{8.1}{9.4}\selectfont PDG2026:\allowbreak{}observable:\allowbreak{}M042:\allowbreak{}M042M:\allowbreak{}335158} & \textbf{2157+-12} MeV & \(u_+\): 11.76667\allowbreak{}94864509\allowbreak{}4\par \(u_-\): 11.76667\allowbreak{}94864509\allowbreak{}4\par Tipo: NONE\par Confianza: ND & COMP-MESON\par E04 / M08\\\addlinespace[3pt]
% VI-CATALOG-ROW masas 235
\textbf{235. eta eta MODE}\par f\_\allowbreak{}2(2150)\allowbreak{}0\par {\fontsize{8.1}{9.4}\selectfont PDG2026:\allowbreak{}observable:\allowbreak{}M042:\allowbreak{}M042M3:\allowbreak{}335159} & \textbf{2157+-12} MeV & \(u_+\): 11.76667\allowbreak{}94864509\allowbreak{}4\par \(u_-\): 11.76667\allowbreak{}94864509\allowbreak{}4\par Tipo: NONE\par Confianza: ND & COMP-MESON\par E04 / M08\\\addlinespace[3pt]
% VI-CATALOG-ROW masas 236
\textbf{236. omega\_\allowbreak{}3(1670) MASS}\par omega\_\allowbreak{}3(1670)\allowbreak{}0\par {\fontsize{8.1}{9.4}\selectfont PDG2026:\allowbreak{}observable:\allowbreak{}M045:\allowbreak{}M045M:\allowbreak{}334778} & \textbf{1667+-4} MeV & \(u_+\): 3.991105\allowbreak{}60803384\allowbreak{}38\par \(u_-\): 3.991105\allowbreak{}60803384\allowbreak{}38\par Tipo: NONE\par Confianza: ND & COMP-MESON\par E04 / M08\\\addlinespace[3pt]
% VI-CATALOG-ROW masas 237
\textbf{237. Upsilon(4S) MASS}\par Upsilon(4S)\par {\fontsize{8.1}{9.4}\selectfont PDG2026:\allowbreak{}observable:\allowbreak{}M047:\allowbreak{}M047M:\allowbreak{}347744} & \textbf{10579.4+-1.2} MeV & \(u_+\): 1.189606\allowbreak{}00825384\allowbreak{}51\par \(u_-\): 1.189606\allowbreak{}00825384\allowbreak{}51\par Tipo: NONE\par Confianza: ND & COMP-MESON\par E04 / M08\\\addlinespace[3pt]
% VI-CATALOG-ROW masas 238
\textbf{238. Upsilon(3S) MASS}\par Upsilon(3S)\par {\fontsize{8.1}{9.4}\selectfont PDG2026:\allowbreak{}observable:\allowbreak{}M048:\allowbreak{}M048M:\allowbreak{}347659} & \textbf{10355.1+-0.5} MeV & \(u_+\): 0.5\par \(u_-\): 0.5\par Tipo: NONE\par Confianza: ND & COMP-MESON\par E04 / M08\\\addlinespace[3pt]
% VI-CATALOG-ROW masas 239
\textbf{239. Upsilon(1S) MASS}\par Upsilon(1S)\par {\fontsize{8.1}{9.4}\selectfont PDG2026:\allowbreak{}observable:\allowbreak{}M049:\allowbreak{}M049M:\allowbreak{}346837} & \textbf{9460.40+-0.10} MeV & \(u_+\): 0.098429\allowbreak{}82316061\allowbreak{}0927\par \(u_-\): 0.098429\allowbreak{}82316061\allowbreak{}0927\par Tipo: NONE\par Confianza: ND & COMP-MESON\par E04 / M08\\\addlinespace[3pt]
% VI-CATALOG-ROW masas 240
\textbf{240. chi\_\allowbreak{}c2(3930) MASS}\par chi\_\allowbreak{}c2(3930)\par {\fontsize{8.1}{9.4}\selectfont PDG2026:\allowbreak{}observable:\allowbreak{}M050:\allowbreak{}M050M:\allowbreak{}346228} & \textbf{3922.5+-1.0} MeV & \(u_+\): 1.019894\allowbreak{}48998559\allowbreak{}71\par \(u_-\): 1.019894\allowbreak{}48998559\allowbreak{}71\par Tipo: NONE\par Confianza: ND & COMP-MESON\par E04 / M08\\\addlinespace[3pt]
% VI-CATALOG-ROW masas 241
\textbf{241. Upsilon(2S) MASS}\par Upsilon(2S)\par {\fontsize{8.1}{9.4}\selectfont PDG2026:\allowbreak{}observable:\allowbreak{}M052:\allowbreak{}M052M:\allowbreak{}347286} & \textbf{10023.4+-0.5} MeV & \(u_+\): 0.5\par \(u_-\): 0.5\par Tipo: NONE\par Confianza: ND & COMP-MESON\par E04 / M08\\\addlinespace[3pt]
% VI-CATALOG-ROW masas 242
\textbf{242. psi(3770) MASS}\par psi(3770)\par {\fontsize{8.1}{9.4}\selectfont PDG2026:\allowbreak{}observable:\allowbreak{}M053:\allowbreak{}M053M:\allowbreak{}345866} & \textbf{3773.7+-0.7} MeV & \(u_+\): 0.695751\allowbreak{}16407149\allowbreak{}418\par \(u_-\): 0.695751\allowbreak{}16407149\allowbreak{}418\par Tipo: NONE\par Confianza: ND & COMP-MESON\par E04 / M08\\\addlinespace[3pt]
% VI-CATALOG-ROW masas 243
\textbf{243. phi\_\allowbreak{}3(1850) MASS}\par phi\_\allowbreak{}3(1850)\allowbreak{}0\par {\fontsize{8.1}{9.4}\selectfont PDG2026:\allowbreak{}observable:\allowbreak{}M054:\allowbreak{}M054M:\allowbreak{}335023} & \textbf{1854+-7} MeV & \(u_+\): 6.882472\allowbreak{}01611685\allowbreak{}27\par \(u_-\): 6.666666\allowbreak{}66666666\allowbreak{}7\par Tipo: NONE\par Confianza: ND & COMP-MESON\par E04 / M08\\\addlinespace[3pt]
% VI-CATALOG-ROW masas 244
\textbf{244. chi\_\allowbreak{}c1(1P) MASS}\par chi\_\allowbreak{}c1(1P)\par {\fontsize{8.1}{9.4}\selectfont PDG2026:\allowbreak{}observable:\allowbreak{}M055:\allowbreak{}M055M:\allowbreak{}344598} & \textbf{3510.67+-0.05} MeV & \(u_+\): 0.050739\allowbreak{}75077067\allowbreak{}1888\par \(u_-\): 0.050739\allowbreak{}75077067\allowbreak{}1888\par Tipo: NONE\par Confianza: ND & COMP-MESON\par E04 / M08\\\addlinespace[3pt]
% VI-CATALOG-ROW masas 245
\textbf{245. chi\_\allowbreak{}c0(1P) MASS}\par chi\_\allowbreak{}c0(1P)\par {\fontsize{8.1}{9.4}\selectfont PDG2026:\allowbreak{}observable:\allowbreak{}M056:\allowbreak{}M056M:\allowbreak{}344304} & \textbf{3415.50+-0.19} MeV & \(u_+\): 0.193118\allowbreak{}29210333\allowbreak{}57\par \(u_-\): 0.192405\allowbreak{}27880424\allowbreak{}149\par Tipo: NONE\par Confianza: ND & COMP-MESON\par E04 / M08\\\addlinespace[3pt]
% VI-CATALOG-ROW masas 246
\textbf{246. chi\_\allowbreak{}c2(1P) MASS}\par chi\_\allowbreak{}c2(1P)\par {\fontsize{8.1}{9.4}\selectfont PDG2026:\allowbreak{}observable:\allowbreak{}M057:\allowbreak{}M057M:\allowbreak{}344926} & \textbf{3556.17+-0.07} MeV & \(u_+\): 0.071612\allowbreak{}03246291\allowbreak{}1895\par \(u_-\): 0.071612\allowbreak{}03246291\allowbreak{}1895\par Tipo: NONE\par Confianza: ND & COMP-MESON\par E04 / M08\\\addlinespace[3pt]
% VI-CATALOG-ROW masas 247
\textbf{247. pi(1300) MASS}\par pi(1300)\allowbreak{}0 | pi(1300)\allowbreak{}+ | pi(1300)\allowbreak{}-\par {\fontsize{8.1}{9.4}\selectfont PDG2026:\allowbreak{}observable:\allowbreak{}M058:\allowbreak{}M058M:\allowbreak{}334463} & \textbf{1300+-100} MeV & \(u_+\): 100\par \(u_-\): 100\par Tipo: NONE\par Confianza: ND & COMP-MESON\par E04 / M08\\\addlinespace[3pt]
% VI-CATALOG-ROW masas 248
\textbf{248. eta\_\allowbreak{}c(2S) MASS}\par eta\_\allowbreak{}c(2S)\par {\fontsize{8.1}{9.4}\selectfont PDG2026:\allowbreak{}observable:\allowbreak{}M059:\allowbreak{}M059M:\allowbreak{}345228} & \textbf{3637.8+-0.6} MeV & \(u_+\): 0.617597\allowbreak{}63649141\allowbreak{}219\par \(u_-\): 0.620048\allowbreak{}76788434\allowbreak{}371\par Tipo: NONE\par Confianza: ND & COMP-MESON\par E04 / M08\\\addlinespace[3pt]
% VI-CATALOG-ROW masas 249
\textbf{249. K\_\allowbreak{}3*(1780) MASS}\par K\_\allowbreak{}3\textasciicircum{}*(1780)\allowbreak{}0 | Kbar\_\allowbreak{}3\textasciicircum{}*(1780)\allowbreak{}0 | K\_\allowbreak{}3\textasciicircum{}*(1780)\allowbreak{}+ | K\_\allowbreak{}3\textasciicircum{}*(1780)\allowbreak{}-\par {\fontsize{8.1}{9.4}\selectfont PDG2026:\allowbreak{}observable:\allowbreak{}M060:\allowbreak{}M060M:\allowbreak{}335975} & \textbf{1779+-8} MeV & \(u_+\): 7.821524\allowbreak{}42884524\allowbreak{}3\par \(u_-\): 7.545019\allowbreak{}38704644\allowbreak{}46\par Tipo: NONE\par Confianza: ND & COMP-MESON\par E04 / M08\\\addlinespace[3pt]
% VI-CATALOG-ROW masas 250
\textbf{250. D*(2007)\allowbreak{}0 MASS}\par D\textasciicircum{}*(2007)\allowbreak{}0 | Dbar\textasciicircum{}*(2007)\allowbreak{}0\par {\fontsize{8.1}{9.4}\selectfont PDG2026:\allowbreak{}observable:\allowbreak{}M061:\allowbreak{}M061M:\allowbreak{}337777} & \textbf{2006.86+-0.05} MeV & \(u_+\): 0.053769\allowbreak{}01475059\allowbreak{}8178\par \(u_-\): 0.053769\allowbreak{}01475059\allowbreak{}8178\par Tipo: NONE\par Confianza: ND & COMP-MESON\par E04 / M08\\\addlinespace[3pt]
% VI-CATALOG-ROW masas 251
\textbf{251. D*(2010)\allowbreak{}+- MASS}\par D\textasciicircum{}*(2010)\allowbreak{}+ | D\textasciicircum{}*(2010)\allowbreak{}-\par {\fontsize{8.1}{9.4}\selectfont PDG2026:\allowbreak{}observable:\allowbreak{}M062:\allowbreak{}M062M:\allowbreak{}337794} & \textbf{2010.27+-0.04} MeV & \(u_+\): 0.044491\allowbreak{}24559318\allowbreak{}8788\par \(u_-\): 0.044491\allowbreak{}24559318\allowbreak{}8788\par Tipo: NONE\par Confianza: ND & COMP-MESON\par E04 / M08\\\addlinespace[3pt]
% VI-CATALOG-ROW masas 252
\textbf{252. K\_\allowbreak{}1(1400) MASS}\par K\_\allowbreak{}1(1400)\allowbreak{}0 | Kbar\_\allowbreak{}1(1400)\allowbreak{}0 | K\_\allowbreak{}1(1400)\allowbreak{}+ | K\_\allowbreak{}1(1400)\allowbreak{}-\par {\fontsize{8.1}{9.4}\selectfont PDG2026:\allowbreak{}observable:\allowbreak{}M064:\allowbreak{}M064M:\allowbreak{}335845} & \textbf{1403+-7} MeV & \(u_+\): 6.873116\allowbreak{}83283743\allowbreak{}06\par \(u_-\): 6.873116\allowbreak{}83283743\allowbreak{}06\par Tipo: NONE\par Confianza: ND & COMP-MESON\par E04 / M08\\\addlinespace[3pt]
% VI-CATALOG-ROW masas 253
\textbf{253. eta rho0 AND pi+ pi- MODES}\par rho(1700)\allowbreak{}0 | rho(1700)\allowbreak{}+ | rho(1700)\allowbreak{}-\par {\fontsize{8.1}{9.4}\selectfont PDG2026:\allowbreak{}observable:\allowbreak{}M065:\allowbreak{}M065M0:\allowbreak{}334880} & \textbf{1720+-20} MeV & \(u_+\): 20\par \(u_-\): 20\par Tipo: NONE\par Confianza: ND & COMP-MESON\par E04 / M08\\\addlinespace[3pt]
% VI-CATALOG-ROW masas 254
\textbf{254. f\_\allowbreak{}2(1430) MASS}\par f\_\allowbreak{}2(1430)\allowbreak{}0\par {\fontsize{8.1}{9.4}\selectfont PDG2026:\allowbreak{}observable:\allowbreak{}M066:\allowbreak{}M066M1:\allowbreak{}334577} & \textbf{\textasciitilde{}1430} MeV & \(u_+\): ND\par \(u_-\): ND\par Tipo: NONE\par Confianza: ND & COMP-MESON\par E04 / M08\\\addlinespace[3pt]
% VI-CATALOG-ROW masas 255
\textbf{255. e+ e- PRODUCTION}\par phi(1680)\allowbreak{}0\par {\fontsize{8.1}{9.4}\selectfont PDG2026:\allowbreak{}observable:\allowbreak{}M067:\allowbreak{}M067M1:\allowbreak{}334825} & \textbf{1680+-20} MeV & \(u_+\): 20\par \(u_-\): 20\par Tipo: NONE\par Confianza: ND & COMP-MESON\par E04 / M08\\\addlinespace[3pt]
% VI-CATALOG-ROW masas 256
\textbf{256. f\_\allowbreak{}0(1710) Breit-Wigner MASS}\par f\_\allowbreak{}0(1710)\allowbreak{}0\par {\fontsize{8.1}{9.4}\selectfont PDG2026:\allowbreak{}observable:\allowbreak{}M068:\allowbreak{}M068M:\allowbreak{}334927} & \textbf{1723+7-6} MeV & \(u_+\): 6.739081\allowbreak{}38855358\allowbreak{}01\par \(u_-\): 5.561777\allowbreak{}54251369\allowbreak{}11\par Tipo: NONE\par Confianza: ND & COMP-MESON\par E04 / M08\\\addlinespace[3pt]
% VI-CATALOG-ROW masas 257
\textbf{257. J/\allowbreak{}psi(1S) MASS}\par J/\allowbreak{}psi(1S)\par {\fontsize{8.1}{9.4}\selectfont PDG2026:\allowbreak{}observable:\allowbreak{}M070:\allowbreak{}M070M:\allowbreak{}343505} & \textbf{3096.900+-0.006} MeV & \(u_+\): 0.006185\allowbreak{}88112993\allowbreak{}74487\par \(u_-\): 0.006185\allowbreak{}88112993\allowbreak{}74487\par Tipo: NONE\par Confianza: ND & COMP-MESON\par E04 / M08\\\addlinespace[3pt]
% VI-CATALOG-ROW masas 258
\textbf{258. psi(2S) MASS}\par psi(2S)\par {\fontsize{8.1}{9.4}\selectfont PDG2026:\allowbreak{}observable:\allowbreak{}M071:\allowbreak{}M071M:\allowbreak{}345299} & \textbf{3686.097+-0.011} MeV & \(u_+\): 0.010553\allowbreak{}10115798\allowbreak{}376\par \(u_-\): 0.010553\allowbreak{}10115798\allowbreak{}376\par Tipo: NONE\par Confianza: ND & COMP-MESON\par E04 / M08\\\addlinespace[3pt]
% VI-CATALOG-ROW masas 259
\textbf{259. psi(4040) MASS}\par psi(4040)\par {\fontsize{8.1}{9.4}\selectfont PDG2026:\allowbreak{}observable:\allowbreak{}M072:\allowbreak{}M072M:\allowbreak{}346249} & \textbf{4040+-4} MeV & \(u_+\): 4.299999\allowbreak{}99999999\allowbreak{}98\par \(u_-\): 4.299999\allowbreak{}99999999\allowbreak{}98\par Tipo: NONE\par Confianza: ND & COMP-MESON\par E04 / M08\\\addlinespace[3pt]
% VI-CATALOG-ROW masas 260
\textbf{260. psi(4415) MASS}\par psi(4415)\par {\fontsize{8.1}{9.4}\selectfont PDG2026:\allowbreak{}observable:\allowbreak{}M073:\allowbreak{}M073M:\allowbreak{}346684} & \textbf{4415+-5} MeV & \(u_+\): 4.527883\allowbreak{}25337319\allowbreak{}54\par \(u_-\): 4.527883\allowbreak{}25337319\allowbreak{}54\par Tipo: NONE\par Confianza: ND & COMP-MESON\par E04 / M08\\\addlinespace[3pt]
% VI-CATALOG-ROW masas 261
\textbf{261. psi(4230) MASS}\par psi(4230)\par {\fontsize{8.1}{9.4}\selectfont PDG2026:\allowbreak{}observable:\allowbreak{}M074:\allowbreak{}M074M:\allowbreak{}346460} & \textbf{4221.7+-2.5} MeV & \(u_+\): 2.456868\allowbreak{}50851105\allowbreak{}82\par \(u_-\): 2.454400\allowbreak{}62347709\allowbreak{}9\par Tipo: NONE\par Confianza: ND & COMP-MESON\par E04 / M08\\\addlinespace[3pt]
% VI-CATALOG-ROW masas 262
\textbf{262. pi(1800) MASS}\par pi(1800)\allowbreak{}0 | pi(1800)\allowbreak{}+ | pi(1800)\allowbreak{}-\par {\fontsize{8.1}{9.4}\selectfont PDG2026:\allowbreak{}observable:\allowbreak{}M075:\allowbreak{}M075M:\allowbreak{}334966} & \textbf{1810+9-11} MeV & \(u_+\): 9.404594\allowbreak{}92210335\allowbreak{}31\par \(u_-\): 11.02751\allowbreak{}39992719\allowbreak{}19\par Tipo: NONE\par Confianza: ND & COMP-MESON\par E04 / M08\\\addlinespace[3pt]
% VI-CATALOG-ROW masas 263
\textbf{263. chi\_\allowbreak{}b0(1P) MASS}\par chi\_\allowbreak{}b0(1P)\par {\fontsize{8.1}{9.4}\selectfont PDG2026:\allowbreak{}observable:\allowbreak{}M076:\allowbreak{}M076M:\allowbreak{}347129} & \textbf{9859.44+-0.42+-0.31} MeV & \(u_+\): 0.419999\allowbreak{}99999999\allowbreak{}998\par \(u_-\): 0.419999\allowbreak{}99999999\allowbreak{}998\par Tipo: NONE\par Confianza: ND & COMP-MESON\par E04 / M08\\\addlinespace[3pt]
% VI-CATALOG-ROW masas 264
\textbf{264. chi\_\allowbreak{}b1(1P) MASS}\par chi\_\allowbreak{}b1(1P)\par {\fontsize{8.1}{9.4}\selectfont PDG2026:\allowbreak{}observable:\allowbreak{}M077:\allowbreak{}M077M:\allowbreak{}347175} & \textbf{9892.78+-0.26+-0.31} MeV & \(u_+\): 0.260000\allowbreak{}00000000\allowbreak{}001\par \(u_-\): 0.260000\allowbreak{}00000000\allowbreak{}001\par Tipo: NONE\par Confianza: ND & COMP-MESON\par E04 / M08\\\addlinespace[3pt]
% VI-CATALOG-ROW masas 265
\textbf{265. chi\_\allowbreak{}b2(1P) MASS}\par chi\_\allowbreak{}b2(1P)\par {\fontsize{8.1}{9.4}\selectfont PDG2026:\allowbreak{}observable:\allowbreak{}M078:\allowbreak{}M078M:\allowbreak{}347234} & \textbf{9912.21+-0.26+-0.31} MeV & \(u_+\): 0.260000\allowbreak{}00000000\allowbreak{}001\par \(u_-\): 0.260000\allowbreak{}00000000\allowbreak{}001\par Tipo: NONE\par Confianza: ND & COMP-MESON\par E04 / M08\\\addlinespace[3pt]
% VI-CATALOG-ROW masas 266
\textbf{266. chi\_\allowbreak{}b0(2P) MASS}\par chi\_\allowbreak{}b0(2P)\par {\fontsize{8.1}{9.4}\selectfont PDG2026:\allowbreak{}observable:\allowbreak{}M079:\allowbreak{}M079M:\allowbreak{}347497} & \textbf{10232.5+-0.4+-0.5} MeV & \(u_+\): 0.400000\allowbreak{}00000000\allowbreak{}002\par \(u_-\): 0.400000\allowbreak{}00000000\allowbreak{}002\par Tipo: NONE\par Confianza: ND & COMP-MESON\par E04 / M08\\\addlinespace[3pt]
% VI-CATALOG-ROW masas 267
\textbf{267. chi\_\allowbreak{}b1(2P) MASS}\par chi\_\allowbreak{}b1(2P)\par {\fontsize{8.1}{9.4}\selectfont PDG2026:\allowbreak{}observable:\allowbreak{}M080:\allowbreak{}M080M:\allowbreak{}347543} & \textbf{10255.46+-0.22+-0.50} MeV & \(u_+\): 0.22\par \(u_-\): 0.22\par Tipo: NONE\par Confianza: ND & COMP-MESON\par E04 / M08\\\addlinespace[3pt]
% VI-CATALOG-ROW masas 268
\textbf{268. chi\_\allowbreak{}b2(2P) MASS}\par chi\_\allowbreak{}b2(2P)\par {\fontsize{8.1}{9.4}\selectfont PDG2026:\allowbreak{}observable:\allowbreak{}M081:\allowbreak{}M081M:\allowbreak{}347610} & \textbf{10268.65+-0.22+-0.50} MeV & \(u_+\): 0.22\par \(u_-\): 0.22\par Tipo: NONE\par Confianza: ND & COMP-MESON\par E04 / M08\\\addlinespace[3pt]
% VI-CATALOG-ROW masas 269
\textbf{269. f\_\allowbreak{}J(2220) MASS}\par f\_\allowbreak{}J(2220)\allowbreak{}0\par {\fontsize{8.1}{9.4}\selectfont PDG2026:\allowbreak{}observable:\allowbreak{}M082:\allowbreak{}M082M:\allowbreak{}335203} & \textbf{2231.1+-3.5} MeV & \(u_+\): 3.533524\allowbreak{}81075623\allowbreak{}18\par \(u_-\): 3.521461\allowbreak{}07791007\allowbreak{}88\par Tipo: NONE\par Confianza: ND & COMP-MESON\par E04 / M08\\\addlinespace[3pt]
% VI-CATALOG-ROW masas 270
\textbf{270. f\_\allowbreak{}1(1510) MASS}\par f\_\allowbreak{}1(1510)\allowbreak{}0\par {\fontsize{8.1}{9.4}\selectfont PDG2026:\allowbreak{}observable:\allowbreak{}M084:\allowbreak{}M084M:\allowbreak{}334658} & \textbf{1518+-5} MeV & \(u_+\): 5.216295\allowbreak{}34269634\allowbreak{}03\par \(u_-\): 5.216295\allowbreak{}34269634\allowbreak{}03\par Tipo: NONE\par Confianza: ND & COMP-MESON\par E04 / M08\\\addlinespace[3pt]
% VI-CATALOG-ROW masas 271
\textbf{271. X(1835) MASS}\par X(1835)\allowbreak{}0\par {\fontsize{8.1}{9.4}\selectfont PDG2026:\allowbreak{}observable:\allowbreak{}M085:\allowbreak{}M085M:\allowbreak{}335011} & \textbf{1834.0+4.0-3.0} MeV & \(u_+\): 4.117870\allowbreak{}74570238\allowbreak{}51\par \(u_-\): 2.984407\allowbreak{}59509600\allowbreak{}41\par Tipo: NONE\par Confianza: ND & COMP-MESON\par E04 / M08\\\addlinespace[3pt]
% VI-CATALOG-ROW masas 272
\textbf{272. f\_\allowbreak{}6(2510) MASS}\par f\_\allowbreak{}6(2510)\allowbreak{}0\par {\fontsize{8.1}{9.4}\selectfont PDG2026:\allowbreak{}observable:\allowbreak{}M089:\allowbreak{}M089M:\allowbreak{}335277} & \textbf{2470+-50} MeV & \(u_+\): 44.99999\allowbreak{}79654523\allowbreak{}51\par \(u_-\): 44.99999\allowbreak{}79654523\allowbreak{}51\par Tipo: NONE\par Confianza: ND & COMP-MESON\par E04 / M08\\\addlinespace[3pt]
% VI-CATALOG-ROW masas 273
\textbf{273. K\_\allowbreak{}3(2320) MASS}\par K\_\allowbreak{}3(2320)\allowbreak{}0 | Kbar\_\allowbreak{}3(2320)\allowbreak{}0 | K\_\allowbreak{}3(2320)\allowbreak{}+ | K\_\allowbreak{}3(2320)\allowbreak{}-\par {\fontsize{8.1}{9.4}\selectfont PDG2026:\allowbreak{}observable:\allowbreak{}M090:\allowbreak{}M090M:\allowbreak{}336034} & \textbf{2324+-24} MeV & \(u_+\): 24\par \(u_-\): 24\par Tipo: NONE\par Confianza: ND & COMP-MESON\par E04 / M08\\\addlinespace[3pt]
% VI-CATALOG-ROW masas 274
\textbf{274. K\_\allowbreak{}4(2500) MASS}\par K\_\allowbreak{}4(2500)\allowbreak{}0 | Kbar\_\allowbreak{}4(2500)\allowbreak{}0 | K\_\allowbreak{}4(2500)\allowbreak{}+ | K\_\allowbreak{}4(2500)\allowbreak{}-\par {\fontsize{8.1}{9.4}\selectfont PDG2026:\allowbreak{}observable:\allowbreak{}M091:\allowbreak{}M091M:\allowbreak{}336040} & \textbf{2490+-20} MeV & \(u_+\): 20\par \(u_-\): 20\par Tipo: NONE\par Confianza: ND & COMP-MESON\par E04 / M08\\\addlinespace[3pt]
% VI-CATALOG-ROW masas 275
\textbf{275. Upsilon(10860) MASS}\par Upsilon(10860)\par {\fontsize{8.1}{9.4}\selectfont PDG2026:\allowbreak{}observable:\allowbreak{}M092:\allowbreak{}M092M:\allowbreak{}347804} & \textbf{10885.2+2.6-1.6} MeV & \(u_+\): 2.565847\allowbreak{}11714697\allowbreak{}32\par \(u_-\): 1.568610\allowbreak{}14024743\allowbreak{}19\par Tipo: NONE\par Confianza: ND & COMP-MESON\par E04 / M08\\\addlinespace[3pt]
% VI-CATALOG-ROW masas 276
\textbf{276. Upsilon(11020) MASS}\par Upsilon(11020)\par {\fontsize{8.1}{9.4}\selectfont PDG2026:\allowbreak{}observable:\allowbreak{}M093:\allowbreak{}M093M:\allowbreak{}347905} & \textbf{11000+-4} MeV & \(u_+\): 4.023428\allowbreak{}15892643\allowbreak{}29\par \(u_-\): 4.024230\allowbreak{}73877622\allowbreak{}44\par Tipo: NONE\par Confianza: ND & COMP-MESON\par E04 / M08\\\addlinespace[3pt]
% VI-CATALOG-ROW masas 277
\textbf{277. K*(1410) MASS}\par K\textasciicircum{}*(1410)\allowbreak{}0 | Kbar\textasciicircum{}*(1410)\allowbreak{}0 | K\textasciicircum{}*(1410)\allowbreak{}+ | K\textasciicircum{}*(1410)\allowbreak{}-\par {\fontsize{8.1}{9.4}\selectfont PDG2026:\allowbreak{}observable:\allowbreak{}M094:\allowbreak{}M094M:\allowbreak{}335865} & \textbf{1414+-15} MeV & \(u_+\): 14.53984\allowbreak{}54347451\allowbreak{}6\par \(u_-\): 14.53984\allowbreak{}54347451\allowbreak{}6\par Tipo: NONE\par Confianza: ND & COMP-MESON\par E04 / M08\\\addlinespace[3pt]
% VI-CATALOG-ROW masas 278
\textbf{278. K*(1680) MASS}\par K\textasciicircum{}*(1680)\allowbreak{}0 | Kbar\textasciicircum{}*(1680)\allowbreak{}0 | K\textasciicircum{}*(1680)\allowbreak{}+ | K\textasciicircum{}*(1680)\allowbreak{}-\par {\fontsize{8.1}{9.4}\selectfont PDG2026:\allowbreak{}observable:\allowbreak{}M095:\allowbreak{}M095M:\allowbreak{}335946} & \textbf{1718+-18} MeV & \(u_+\): 17.07872\allowbreak{}69165451\allowbreak{}9\par \(u_-\): 18.69816\allowbreak{}26725943\allowbreak{}39\par Tipo: NONE\par Confianza: ND & COMP-MESON\par E04 / M08\\\addlinespace[3pt]
% VI-CATALOG-ROW masas 279
\textbf{279. K\_\allowbreak{}5*(2380) MASS}\par K\_\allowbreak{}5\textasciicircum{}*(2380)\allowbreak{}0 | Kbar\_\allowbreak{}5\textasciicircum{}*(2380)\allowbreak{}0 | K\_\allowbreak{}5\textasciicircum{}*(2380)\allowbreak{}+ | K\_\allowbreak{}5\textasciicircum{}*(2380)\allowbreak{}-\par {\fontsize{8.1}{9.4}\selectfont PDG2026:\allowbreak{}observable:\allowbreak{}M098:\allowbreak{}M098M:\allowbreak{}336036} & \textbf{2382+-24} MeV & \(u_+\): 23.60084\allowbreak{}74424118\allowbreak{}9\par \(u_-\): 23.60084\allowbreak{}74424118\allowbreak{}9\par Tipo: NONE\par Confianza: ND & COMP-MESON\par E04 / M08\\\addlinespace[3pt]
% VI-CATALOG-ROW masas 280
\textbf{280. K\_\allowbreak{}1(1650) MASS}\par K\_\allowbreak{}1(1650)\allowbreak{}0 | Kbar\_\allowbreak{}1(1650)\allowbreak{}0 | K\_\allowbreak{}1(1650)\allowbreak{}+ | K\_\allowbreak{}1(1650)\allowbreak{}-\par {\fontsize{8.1}{9.4}\selectfont PDG2026:\allowbreak{}observable:\allowbreak{}M099:\allowbreak{}M099M:\allowbreak{}335944} & \textbf{1650+-50} MeV & \(u_+\): 50\par \(u_-\): 50\par Tipo: NONE\par Confianza: ND & COMP-MESON\par E04 / M08\\\addlinespace[3pt]
% VI-CATALOG-ROW masas 281
\textbf{281. eta\_\allowbreak{}2(1870) MASS}\par eta\_\allowbreak{}2(1870)\allowbreak{}0\par {\fontsize{8.1}{9.4}\selectfont PDG2026:\allowbreak{}observable:\allowbreak{}M101:\allowbreak{}M101M:\allowbreak{}335032} & \textbf{1842+-8} MeV & \(u_+\): 8.043075\allowbreak{}30759601\allowbreak{}09\par \(u_-\): 8.043075\allowbreak{}30759601\allowbreak{}09\par Tipo: NONE\par Confianza: ND & COMP-MESON\par E04 / M08\\\addlinespace[3pt]
% VI-CATALOG-ROW masas 282
\textbf{282. phi(2170) MASS}\par phi(2170)\allowbreak{}0 | phi(2170)\allowbreak{}+ | phi(2170)\allowbreak{}-\par {\fontsize{8.1}{9.4}\selectfont PDG2026:\allowbreak{}observable:\allowbreak{}M103:\allowbreak{}M103M:\allowbreak{}335180} & \textbf{2164+-5} MeV & \(u_+\): 5.113273\allowbreak{}60016105\allowbreak{}46\par \(u_-\): 5.113273\allowbreak{}60016105\allowbreak{}46\par Tipo: NONE\par Confianza: ND & COMP-MESON\par E04 / M08\\\addlinespace[3pt]
% VI-CATALOG-ROW masas 283
\textbf{283. K\_\allowbreak{}2*(1980) MASS}\par K\_\allowbreak{}2\textasciicircum{}*(1980)\allowbreak{}0 | Kbar\_\allowbreak{}2\textasciicircum{}*(1980)\allowbreak{}0 | K\_\allowbreak{}2\textasciicircum{}*(1980)\allowbreak{}+ | K\_\allowbreak{}2\textasciicircum{}*(1980)\allowbreak{}-\par {\fontsize{8.1}{9.4}\selectfont PDG2026:\allowbreak{}observable:\allowbreak{}M104:\allowbreak{}M104M:\allowbreak{}336003} & \textbf{1990+60-50} MeV & \(u_+\): 59.10515\allowbreak{}25697060\allowbreak{}07\par \(u_-\): 45.40577\allowbreak{}22472393\allowbreak{}12\par Tipo: NONE\par Confianza: ND & COMP-MESON\par E04 / M08\\\addlinespace[3pt]
% VI-CATALOG-ROW masas 284
\textbf{284. rho(1450) MASS}\par rho(1450)\allowbreak{}0 | rho(1450)\allowbreak{}+ | rho(1450)\allowbreak{}-\par {\fontsize{8.1}{9.4}\selectfont PDG2026:\allowbreak{}observable:\allowbreak{}M105:\allowbreak{}M105M0:\allowbreak{}334592} & \textbf{1465 +-25} MeV\par Presentación: 1465+-25 & \(u_+\): 25\par \(u_-\): 25\par Tipo: NONE\par Confianza: ND & COMP-MESON\par E04 / M08\\\addlinespace[3pt]
% VI-CATALOG-ROW masas 285
\textbf{285. f\_\allowbreak{}2(2010) MASS}\par f\_\allowbreak{}2(2010)\allowbreak{}0\par {\fontsize{8.1}{9.4}\selectfont PDG2026:\allowbreak{}observable:\allowbreak{}M106:\allowbreak{}M106M:\allowbreak{}335112} & \textbf{2010+60-80} MeV & \(u_+\): 62\par \(u_-\): 76\par Tipo: NONE\par Confianza: ND & COMP-MESON\par E04 / M08\\\addlinespace[3pt]
% VI-CATALOG-ROW masas 286
\textbf{286. f\_\allowbreak{}2(2300) MASS}\par f\_\allowbreak{}2(2300)\allowbreak{}0\par {\fontsize{8.1}{9.4}\selectfont PDG2026:\allowbreak{}observable:\allowbreak{}M107:\allowbreak{}M107M:\allowbreak{}335234} & \textbf{2297+-28} MeV & \(u_+\): 28\par \(u_-\): 28\par Tipo: NONE\par Confianza: ND & COMP-MESON\par E04 / M08\\\addlinespace[3pt]
% VI-CATALOG-ROW masas 287
\textbf{287. f\_\allowbreak{}2(2340) MASS}\par f\_\allowbreak{}2(2340)\allowbreak{}0\par {\fontsize{8.1}{9.4}\selectfont PDG2026:\allowbreak{}observable:\allowbreak{}M108:\allowbreak{}M108M:\allowbreak{}335252} & \textbf{2346+21-10} MeV & \(u_+\): 21.30056\allowbreak{}24897101\allowbreak{}81\par \(u_-\): 9.742332\allowbreak{}63799634\allowbreak{}84\par Tipo: NONE\par Confianza: ND & COMP-MESON\par E04 / M08\\\addlinespace[3pt]
% VI-CATALOG-ROW masas 288
\textbf{288. h\_\allowbreak{}1(1415) MASS}\par h\_\allowbreak{}1(1415)\allowbreak{}0\par {\fontsize{8.1}{9.4}\selectfont PDG2026:\allowbreak{}observable:\allowbreak{}M109:\allowbreak{}M109M:\allowbreak{}334552} & \textbf{1409+9-8} MeV & \(u_+\): 9.419406\allowbreak{}81227334\allowbreak{}94\par \(u_-\): 7.776226\allowbreak{}72689420\allowbreak{}2\par Tipo: NONE\par Confianza: ND & COMP-MESON\par E04 / M08\\\addlinespace[3pt]
% VI-CATALOG-ROW masas 289
\textbf{289. f\_\allowbreak{}0(2200) MASS}\par f\_\allowbreak{}0(2200)\allowbreak{}0\par {\fontsize{8.1}{9.4}\selectfont PDG2026:\allowbreak{}observable:\allowbreak{}M112:\allowbreak{}M112M:\allowbreak{}335199} & \textbf{2187+-14} MeV & \(u_+\): 13.53304\allowbreak{}17595192\allowbreak{}7\par \(u_-\): 14.05775\allowbreak{}05351819\allowbreak{}79\par Tipo: NONE\par Confianza: ND & COMP-MESON\par E04 / M08\\\addlinespace[3pt]
% VI-CATALOG-ROW masas 290
\textbf{290. eta(1760) MASS}\par eta(1760)\allowbreak{}0\par {\fontsize{8.1}{9.4}\selectfont PDG2026:\allowbreak{}observable:\allowbreak{}M114:\allowbreak{}M114M:\allowbreak{}334950} & \textbf{1751+-15} MeV & \(u_+\): 14.81487\allowbreak{}49397529\allowbreak{}29\par \(u_-\): 14.98014\allowbreak{}55902391\allowbreak{}59\par Tipo: NONE\par Confianza: ND & COMP-MESON\par E04 / M08\\\addlinespace[3pt]
% VI-CATALOG-ROW masas 291
\textbf{291. eta(2225) MASS}\par eta(2225)\allowbreak{}0\par {\fontsize{8.1}{9.4}\selectfont PDG2026:\allowbreak{}observable:\allowbreak{}M115:\allowbreak{}M115M:\allowbreak{}335232} & \textbf{2221+13-10} MeV & \(u_+\): 13.27133\allowbreak{}38705345\allowbreak{}5\par \(u_-\): 9.520246\allowbreak{}07783074\allowbreak{}83\par Tipo: NONE\par Confianza: ND & COMP-MESON\par E04 / M08\\\addlinespace[3pt]
% VI-CATALOG-ROW masas 292
\textbf{292. f\_\allowbreak{}2(1640) MASS}\par f\_\allowbreak{}2(1640)\allowbreak{}0\par {\fontsize{8.1}{9.4}\selectfont PDG2026:\allowbreak{}observable:\allowbreak{}M117:\allowbreak{}M117M:\allowbreak{}334752} & \textbf{1639+-6} MeV & \(u_+\): 5.831378\allowbreak{}42111532\allowbreak{}39\par \(u_-\): 5.831378\allowbreak{}42111532\allowbreak{}39\par Tipo: NONE\par Confianza: ND & COMP-MESON\par E04 / M08\\\addlinespace[3pt]
% VI-CATALOG-ROW masas 293
\textbf{293. D\_\allowbreak{}s1(2536)\allowbreak{}+- MASS}\par D\_\allowbreak{}s1(2536)\allowbreak{}+ | D\_\allowbreak{}s1(2536)\allowbreak{}-\par {\fontsize{8.1}{9.4}\selectfont PDG2026:\allowbreak{}observable:\allowbreak{}M121:\allowbreak{}M121M:\allowbreak{}338481} & \textbf{2535.12+-0.06} MeV & \(u_+\): 0.060350\allowbreak{}97747089\allowbreak{}8563\par \(u_-\): 0.060350\allowbreak{}97747089\allowbreak{}8563\par Tipo: NONE\par Confianza: ND & COMP-MESON\par E04 / M08\\\addlinespace[3pt]
% VI-CATALOG-ROW masas 294
\textbf{294. f\_\allowbreak{}2(1565) MASS}\par f\_\allowbreak{}2(1565)\allowbreak{}0\par {\fontsize{8.1}{9.4}\selectfont PDG2026:\allowbreak{}observable:\allowbreak{}M123:\allowbreak{}M123M:\allowbreak{}334703} & \textbf{1571+-13} MeV & \(u_+\): 13.37929\allowbreak{}46324489\allowbreak{}91\par \(u_-\): 13.37929\allowbreak{}46324489\allowbreak{}91\par Tipo: NONE\par Confianza: ND & COMP-MESON\par E04 / M08\\\addlinespace[3pt]
% VI-CATALOG-ROW masas 295
\textbf{295. omega(1420) MASS}\par omega(1420)\allowbreak{}0\par {\fontsize{8.1}{9.4}\selectfont PDG2026:\allowbreak{}observable:\allowbreak{}M125:\allowbreak{}M125M:\allowbreak{}334567} & \textbf{1410+-60} MeV & \(u_+\): 60\par \(u_-\): 60\par Tipo: NONE\par Confianza: ND & COMP-MESON\par E04 / M08\\\addlinespace[3pt]
% VI-CATALOG-ROW masas 296
\textbf{296. omega(1650) MASS}\par omega(1650)\allowbreak{}0\par {\fontsize{8.1}{9.4}\selectfont PDG2026:\allowbreak{}observable:\allowbreak{}M126:\allowbreak{}M126M:\allowbreak{}334768} & \textbf{1670+-30} MeV & \(u_+\): 30\par \(u_-\): 30\par Tipo: NONE\par Confianza: ND & COMP-MESON\par E04 / M08\\\addlinespace[3pt]
% VI-CATALOG-ROW masas 297
\textbf{297. K(3100) MASS}\par K(3100)\allowbreak{}0 | Kbar(3100)\allowbreak{}0 | K(3100)\allowbreak{}+ | K(3100)\allowbreak{}-\par {\fontsize{8.1}{9.4}\selectfont PDG2026:\allowbreak{}observable:\allowbreak{}M129:\allowbreak{}M129M:\allowbreak{}336041} & \textbf{\textasciitilde{}3100} MeV & \(u_+\): ND\par \(u_-\): ND\par Tipo: NONE\par Confianza: ND & COMP-MESON\par E04 / M08\\\addlinespace[3pt]
% VI-CATALOG-ROW masas 298
\textbf{298. 3-BODY DECAYS}\par K(3100)\allowbreak{}0 | Kbar(3100)\allowbreak{}0 | K(3100)\allowbreak{}+ | K(3100)\allowbreak{}-\par {\fontsize{8.1}{9.4}\selectfont PDG2026:\allowbreak{}observable:\allowbreak{}M129:\allowbreak{}M129M1:\allowbreak{}336042} & \textbf{3054+-11} MeV & \(u_+\): 10.68148\allowbreak{}84244502\allowbreak{}6\par \(u_-\): 10.68148\allowbreak{}84244502\allowbreak{}6\par Tipo: NONE\par Confianza: ND & COMP-MESON\par E04 / M08\\\addlinespace[3pt]
% VI-CATALOG-ROW masas 299
\textbf{299. 4-BODY DECAYS}\par K(3100)\allowbreak{}0 | Kbar(3100)\allowbreak{}0 | K(3100)\allowbreak{}+ | K(3100)\allowbreak{}-\par {\fontsize{8.1}{9.4}\selectfont PDG2026:\allowbreak{}observable:\allowbreak{}M129:\allowbreak{}M129M2:\allowbreak{}336043} & \textbf{3059+-11} MeV & \(u_+\): 10.64125\allowbreak{}22648014\allowbreak{}39\par \(u_-\): 10.64125\allowbreak{}22648014\allowbreak{}39\par Tipo: NONE\par Confianza: ND & COMP-MESON\par E04 / M08\\\addlinespace[3pt]
% VI-CATALOG-ROW masas 300
\textbf{300. K\_\allowbreak{}0*(1950) MASS}\par K\_\allowbreak{}0\textasciicircum{}*(1950)\allowbreak{}0 | Kbar\_\allowbreak{}0\textasciicircum{}*(1950)\allowbreak{}0 | K\_\allowbreak{}0\textasciicircum{}*(1950)\allowbreak{}+ | K\_\allowbreak{}0\textasciicircum{}*(1950)\allowbreak{}-\par {\fontsize{8.1}{9.4}\selectfont PDG2026:\allowbreak{}observable:\allowbreak{}M134:\allowbreak{}M134M:\allowbreak{}335999} & \textbf{1957+-14} MeV & \(u_+\): 14.32022\allowbreak{}19810733\allowbreak{}29\par \(u_-\): 14.32022\allowbreak{}19810733\allowbreak{}29\par Tipo: NONE\par Confianza: ND & COMP-MESON\par E04 / M08\\\addlinespace[3pt]
% VI-CATALOG-ROW masas 301
\textbf{301. f\_\allowbreak{}2(1950) MASS}\par f\_\allowbreak{}2(1950)\allowbreak{}0\par {\fontsize{8.1}{9.4}\selectfont PDG2026:\allowbreak{}observable:\allowbreak{}M135:\allowbreak{}M135M:\allowbreak{}335075} & \textbf{1936+-12} MeV & \(u_+\): 11.73086\allowbreak{}74827569\allowbreak{}9\par \(u_-\): 11.73086\allowbreak{}74827569\allowbreak{}9\par Tipo: NONE\par Confianza: ND & COMP-MESON\par E04 / M08\\\addlinespace[3pt]
% VI-CATALOG-ROW masas 302
\textbf{302. f\_\allowbreak{}2(1910) omega omega MODE}\par f\_\allowbreak{}2(1910)\allowbreak{}0\par {\fontsize{8.1}{9.4}\selectfont PDG2026:\allowbreak{}observable:\allowbreak{}M142:\allowbreak{}M142M2:\allowbreak{}335057} & \textbf{1900+-9} MeV & \(u_+\): 9.266557\allowbreak{}49458051\allowbreak{}4\par \(u_-\): 9.266557\allowbreak{}49458051\allowbreak{}4\par Tipo: NONE\par Confianza: ND & COMP-MESON\par E04 / M08\\\addlinespace[3pt]
% VI-CATALOG-ROW masas 303
\textbf{303. h\_\allowbreak{}c(1P) MASS}\par h\_\allowbreak{}c(1P)\par {\fontsize{8.1}{9.4}\selectfont PDG2026:\allowbreak{}observable:\allowbreak{}M144:\allowbreak{}M144M:\allowbreak{}344851} & \textbf{3525.37+-0.14} MeV & \(u_+\): 0.144734\allowbreak{}63841322\allowbreak{}931\par \(u_-\): 0.144734\allowbreak{}63841322\allowbreak{}931\par Tipo: NONE\par Confianza: ND & COMP-MESON\par E04 / M08\\\addlinespace[3pt]
% VI-CATALOG-ROW masas 304
\textbf{304. K\_\allowbreak{}2(1820) MASS}\par K\_\allowbreak{}2(1820)\allowbreak{}0 | Kbar\_\allowbreak{}2(1820)\allowbreak{}0 | K\_\allowbreak{}2(1820)\allowbreak{}+ | K\_\allowbreak{}2(1820)\allowbreak{}-\par {\fontsize{8.1}{9.4}\selectfont PDG2026:\allowbreak{}observable:\allowbreak{}M146:\allowbreak{}M146M:\allowbreak{}335990} & \textbf{1819+-12} MeV & \(u_+\): 12.06762\allowbreak{}85812278\allowbreak{}79\par \(u_-\): 12.47184\allowbreak{}29929233\allowbreak{}71\par Tipo: NONE\par Confianza: ND & COMP-MESON\par E04 / M08\\\addlinespace[3pt]
% VI-CATALOG-ROW masas 305
\textbf{305. Mass (Breit-Wigner)}\par f\_\allowbreak{}0(1370)\allowbreak{}0\par {\fontsize{8.1}{9.4}\selectfont PDG2026:\allowbreak{}observable:\allowbreak{}M147:\allowbreak{}M147M:\allowbreak{}334506} & \textbf{1200 to 1500} MeV & \(u_+\): NA\par \(u_-\): NA\par Tipo: R\par Confianza: ND & COMP-MESON\par E04 / M08\\\addlinespace[3pt]
% VI-CATALOG-ROW masas 306
\textbf{306. D\_\allowbreak{}s2*(2573) MASS}\par D\_\allowbreak{}s2\textasciicircum{}*(2573)\allowbreak{}+ | D\_\allowbreak{}s2\textasciicircum{}*(2573)\allowbreak{}-\par {\fontsize{8.1}{9.4}\selectfont PDG2026:\allowbreak{}observable:\allowbreak{}M148:\allowbreak{}M148M:\allowbreak{}338509} & \textbf{2569.1+-0.8} MeV & \(u_+\): 0.827806\allowbreak{}30176758\allowbreak{}352\par \(u_-\): 0.826498\allowbreak{}31496169\allowbreak{}768\par Tipo: NONE\par Confianza: ND & COMP-MESON\par E04 / M08\\\addlinespace[3pt]
% VI-CATALOG-ROW masas 307
\textbf{307. a\_\allowbreak{}0(1450) MASS}\par a\_\allowbreak{}0(1450)\allowbreak{}0 | a\_\allowbreak{}0(1450)\allowbreak{}+ | a\_\allowbreak{}0(1450)\allowbreak{}-\par {\fontsize{8.1}{9.4}\selectfont PDG2026:\allowbreak{}observable:\allowbreak{}M149:\allowbreak{}M149M:\allowbreak{}334579} & \textbf{1439+-34} MeV & \(u_+\): 34.42623\allowbreak{}01167357\allowbreak{}38\par \(u_-\): 34.42623\allowbreak{}01167357\allowbreak{}38\par Tipo: NONE\par Confianza: ND & COMP-MESON\par E04 / M08\\\addlinespace[3pt]
% VI-CATALOG-ROW masas 308
\textbf{308. B\_\allowbreak{}J*(5732) MASS}\par B\_\allowbreak{}J\textasciicircum{}*(5732)\allowbreak{}0 | Bbar\_\allowbreak{}J\textasciicircum{}*(5732)\allowbreak{}0 | B\_\allowbreak{}J\textasciicircum{}*(5732)\allowbreak{}+ | B\_\allowbreak{}J\textasciicircum{}*(5732)\allowbreak{}-\par {\fontsize{8.1}{9.4}\selectfont PDG2026:\allowbreak{}observable:\allowbreak{}M151:\allowbreak{}M151M:\allowbreak{}342551} & \textbf{5698+-8} MeV & \(u_+\): 7.586946\allowbreak{}82434463\allowbreak{}72\par \(u_-\): 7.693025\allowbreak{}46888399\allowbreak{}34\par Tipo: NONE\par Confianza: ND & COMP-MESON\par E04 / M08\\\addlinespace[3pt]
% VI-CATALOG-ROW masas 309
\textbf{309. f\_\allowbreak{}0(1500) MASS}\par f\_\allowbreak{}0(1500)\allowbreak{}0\par {\fontsize{8.1}{9.4}\selectfont PDG2026:\allowbreak{}observable:\allowbreak{}M152:\allowbreak{}M152M:\allowbreak{}334628} & \textbf{1522+-25} MeV & \(u_+\): 25\par \(u_-\): 25\par Tipo: NONE\par Confianza: ND & COMP-MESON\par E04 / M08\\\addlinespace[3pt]
% VI-CATALOG-ROW masas 310
\textbf{310. B\_\allowbreak{}sJ*(5850) MASS}\par B\_\allowbreak{}sJ\textasciicircum{}*(5850)\allowbreak{}0 | Bbar\_\allowbreak{}sJ\textasciicircum{}*(5850)\allowbreak{}0\par {\fontsize{8.1}{9.4}\selectfont PDG2026:\allowbreak{}observable:\allowbreak{}M153:\allowbreak{}M153M:\allowbreak{}343151} & \textbf{5853+-15} MeV & \(u_+\): 15\par \(u_-\): 15\par Tipo: NONE\par Confianza: ND & COMP-MESON\par E04 / M08\\\addlinespace[3pt]
% VI-CATALOG-ROW masas 311
\textbf{311. eta\_\allowbreak{}2(1645) MASS}\par eta\_\allowbreak{}2(1645)\allowbreak{}0\par {\fontsize{8.1}{9.4}\selectfont PDG2026:\allowbreak{}observable:\allowbreak{}M154:\allowbreak{}M154M:\allowbreak{}334758} & \textbf{1617+-5} MeV & \(u_+\): 5.261314\allowbreak{}79626401\allowbreak{}97\par \(u_-\): 5.261314\allowbreak{}79626401\allowbreak{}97\par Tipo: NONE\par Confianza: ND & COMP-MESON\par E04 / M08\\\addlinespace[3pt]
% VI-CATALOG-ROW masas 312
\textbf{312. f\_\allowbreak{}0(2020) MASS}\par f\_\allowbreak{}0(2020)\allowbreak{}0\par {\fontsize{8.1}{9.4}\selectfont PDG2026:\allowbreak{}observable:\allowbreak{}M156:\allowbreak{}M156M:\allowbreak{}335120} & \textbf{1982+54.1-3.0} MeV & \(u_+\): 54.08326\allowbreak{}91319598\allowbreak{}38\par \(u_-\): 3\par Tipo: NONE\par Confianza: ND & COMP-MESON\par E04 / M08\\\addlinespace[3pt]
% VI-CATALOG-ROW masas 313
\textbf{313. D*(2640)\allowbreak{}+- MASS}\par D\textasciicircum{}*(2640)\allowbreak{}0 | Dbar\textasciicircum{}*(2640)\allowbreak{}0 | D\textasciicircum{}*(2640)\allowbreak{}+ | D\textasciicircum{}*(2640)\allowbreak{}-\par {\fontsize{8.1}{9.4}\selectfont PDG2026:\allowbreak{}observable:\allowbreak{}M158:\allowbreak{}M158M:\allowbreak{}337854} & \textbf{2637+-6} MeV & \(u_+\): 6.324555\allowbreak{}32033675\allowbreak{}9\par \(u_-\): 6.324555\allowbreak{}32033675\allowbreak{}9\par Tipo: NONE\par Confianza: ND & COMP-MESON\par E04 / M08\\\addlinespace[3pt]
% VI-CATALOG-ROW masas 314
\textbf{314. chi\_\allowbreak{}c0(3915) MASS}\par chi\_\allowbreak{}c0(3915)\par {\fontsize{8.1}{9.4}\selectfont PDG2026:\allowbreak{}observable:\allowbreak{}M159:\allowbreak{}M159M:\allowbreak{}346204} & \textbf{3922.1+-1.8} MeV & \(u_+\): 1.773070\allowbreak{}50384254\allowbreak{}29\par \(u_-\): 1.760361\allowbreak{}83430800\allowbreak{}81\par Tipo: NONE\par Confianza: ND & COMP-MESON\par E04 / M08\\\addlinespace[3pt]
% VI-CATALOG-ROW masas 315
\textbf{315. K(1630) MASS}\par K(1630)\allowbreak{}0 | Kbar(1630)\allowbreak{}0 | K(1630)\allowbreak{}+ | K(1630)\allowbreak{}-\par {\fontsize{8.1}{9.4}\selectfont PDG2026:\allowbreak{}observable:\allowbreak{}M160:\allowbreak{}M160M:\allowbreak{}335942} & \textbf{1629+-7} MeV & \(u_+\): 7\par \(u_-\): 7\par Tipo: NONE\par Confianza: ND & COMP-MESON\par E04 / M08\\\addlinespace[3pt]
% VI-CATALOG-ROW masas 316
\textbf{316. a\_\allowbreak{}1(1640) MASS}\par a\_\allowbreak{}1(1640)\allowbreak{}0 | a\_\allowbreak{}1(1640)\allowbreak{}+ | a\_\allowbreak{}1(1640)\allowbreak{}-\par {\fontsize{8.1}{9.4}\selectfont PDG2026:\allowbreak{}observable:\allowbreak{}M161:\allowbreak{}M161M:\allowbreak{}334740} & \textbf{1655+-16} MeV & \(u_+\): 15.61536\allowbreak{}24587929\allowbreak{}11\par \(u_-\): 16.74290\allowbreak{}48405724\allowbreak{}61\par Tipo: NONE\par Confianza: ND & COMP-MESON\par E04 / M08\\\addlinespace[3pt]
% VI-CATALOG-ROW masas 317
\textbf{317. a\_\allowbreak{}2(1700) MASS}\par a\_\allowbreak{}2(1700)\allowbreak{}0 | a\_\allowbreak{}2(1700)\allowbreak{}+ | a\_\allowbreak{}2(1700)\allowbreak{}-\par {\fontsize{8.1}{9.4}\selectfont PDG2026:\allowbreak{}observable:\allowbreak{}M162:\allowbreak{}M162M:\allowbreak{}334903} & \textbf{1706+-14} MeV & \(u_+\): 13.50144\allowbreak{}09315998\allowbreak{}2\par \(u_-\): 14.83907\allowbreak{}44674554\allowbreak{}29\par Tipo: NONE\par Confianza: ND & COMP-MESON\par E04 / M08\\\addlinespace[3pt]
% VI-CATALOG-ROW masas 318
\textbf{318. pi\_\allowbreak{}1(1600) T-Matrix Pole sqrt(s)}\par pi\_\allowbreak{}1(1600)\allowbreak{}0 | pi\_\allowbreak{}1(1600)\allowbreak{}+ | pi\_\allowbreak{}1(1600)\allowbreak{}-\par {\fontsize{8.1}{9.4}\selectfont PDG2026:\allowbreak{}observable:\allowbreak{}M164:\allowbreak{}M164TMP:\allowbreak{}334723} & \textbf{(1480 -\kern0pt{}- 1680) - i (150 -\kern0pt{}- 300)} MeV\par Presentación: (1480-\kern0pt{}-1680)\allowbreak{}-i(150-\kern0pt{}-300) & \(u_+\): NA\par \(u_-\): NA\par Tipo: R\par Confianza: ND & COMP-MESON\par E04 / M08\\\addlinespace[3pt]
% VI-CATALOG-ROW masas 319
\textbf{319. pi\_\allowbreak{}1(1600) MASS (eta pi mode)}\par pi\_\allowbreak{}1(1600)\allowbreak{}0 | pi\_\allowbreak{}1(1600)\allowbreak{}+ | pi\_\allowbreak{}1(1600)\allowbreak{}-\par {\fontsize{8.1}{9.4}\selectfont PDG2026:\allowbreak{}observable:\allowbreak{}M164:\allowbreak{}M164MEP:\allowbreak{}334724} & \textbf{1354+-25} MeV & \(u_+\): 25.62037\allowbreak{}14275794\allowbreak{}11\par \(u_-\): 24.37777\allowbreak{}86955451\par Tipo: NONE\par Confianza: ND & COMP-MESON\par E04 / M08\\\addlinespace[3pt]
% VI-CATALOG-ROW masas 320
\textbf{320. pi\_\allowbreak{}1(1600) MASS (non-eta pi mode)}\par pi\_\allowbreak{}1(1600)\allowbreak{}0 | pi\_\allowbreak{}1(1600)\allowbreak{}+ | pi\_\allowbreak{}1(1600)\allowbreak{}-\par {\fontsize{8.1}{9.4}\selectfont PDG2026:\allowbreak{}observable:\allowbreak{}M164:\allowbreak{}M164M:\allowbreak{}334725} & \textbf{1645+40-17} MeV & \(u_+\): 39.71975\allowbreak{}62267507\allowbreak{}18\par \(u_-\): 16.57567\allowbreak{}98763040\allowbreak{}68\par Tipo: NONE\par Confianza: ND & COMP-MESON\par E04 / M08\\\addlinespace[3pt]
% VI-CATALOG-ROW masas 321
\textbf{321. h\_\allowbreak{}1(1595) MASS}\par h\_\allowbreak{}1(1595)\allowbreak{}0\par {\fontsize{8.1}{9.4}\selectfont PDG2026:\allowbreak{}observable:\allowbreak{}M166:\allowbreak{}M166M:\allowbreak{}334720} & \textbf{1594+18-60} MeV & \(u_+\): 18.02775\allowbreak{}63773199\allowbreak{}5\par \(u_-\): 61.84658\allowbreak{}43842649\allowbreak{}08\par Tipo: NONE\par Confianza: ND & COMP-MESON\par E04 / M08\\\addlinespace[3pt]
% VI-CATALOG-ROW masas 322
\textbf{322. f\_\allowbreak{}0(2100) MASS}\par f\_\allowbreak{}0(2100)\allowbreak{}0\par {\fontsize{8.1}{9.4}\selectfont PDG2026:\allowbreak{}observable:\allowbreak{}M168:\allowbreak{}M168M:\allowbreak{}335156} & \textbf{2095+17-19} MeV & \(u_+\): 17.06049\allowbreak{}66887379\allowbreak{}81\par \(u_-\): 18.97966\allowbreak{}10074380\allowbreak{}18\par Tipo: NONE\par Confianza: ND & COMP-MESON\par E04 / M08\\\addlinespace[3pt]
% VI-CATALOG-ROW masas 323
\textbf{323. eta\_\allowbreak{}b(1S) MASS}\par eta\_\allowbreak{}b(1S)\par {\fontsize{8.1}{9.4}\selectfont PDG2026:\allowbreak{}observable:\allowbreak{}M171:\allowbreak{}M171M:\allowbreak{}346822} & \textbf{9398.7+-2.0} MeV & \(u_+\): 2.081610\allowbreak{}14572945\allowbreak{}2\par \(u_-\): 2.002370\allowbreak{}37046477\par Tipo: NONE\par Confianza: ND & COMP-MESON\par E04 / M08\\\addlinespace[3pt]
% VI-CATALOG-ROW masas 324
\textbf{324. D\_\allowbreak{}s0*(2317)\allowbreak{}+- MASS}\par D\_\allowbreak{}s0\textasciicircum{}*(2317)\allowbreak{}+ | D\_\allowbreak{}s0\textasciicircum{}*(2317)\allowbreak{}-\par {\fontsize{8.1}{9.4}\selectfont PDG2026:\allowbreak{}observable:\allowbreak{}M172:\allowbreak{}M172M:\allowbreak{}338429} & \textbf{2317.8+-0.5} MeV & \(u_+\): 0.510032\allowbreak{}96350202\allowbreak{}032\par \(u_-\): 0.510032\allowbreak{}96350202\allowbreak{}032\par Tipo: NONE\par Confianza: ND & COMP-MESON\par E04 / M08\\\addlinespace[3pt]
% VI-CATALOG-ROW masas 325
\textbf{325. D\_\allowbreak{}s1(2460)\allowbreak{}+- MASS}\par D\_\allowbreak{}s1(2460)\allowbreak{}+ | D\_\allowbreak{}s1(2460)\allowbreak{}-\par {\fontsize{8.1}{9.4}\selectfont PDG2026:\allowbreak{}observable:\allowbreak{}M173:\allowbreak{}M173M:\allowbreak{}338448} & \textbf{2459.5+-0.6} MeV & \(u_+\): 0.634430\allowbreak{}27023618\allowbreak{}292\par \(u_-\): 0.634430\allowbreak{}27023618\allowbreak{}292\par Tipo: NONE\par Confianza: ND & COMP-MESON\par E04 / M08\\\addlinespace[3pt]
% VI-CATALOG-ROW masas 326
\textbf{326. K\_\allowbreak{}0*(700) T-Matrix Pole sqrt(s)}\par K\_\allowbreak{}0\textasciicircum{}*(700)\allowbreak{}0 | Kbar\_\allowbreak{}0\textasciicircum{}*(700)\allowbreak{}0 | K\_\allowbreak{}0\textasciicircum{}*(700)\allowbreak{}+ | K\_\allowbreak{}0\textasciicircum{}*(700)\allowbreak{}-\par {\fontsize{8.1}{9.4}\selectfont PDG2026:\allowbreak{}observable:\allowbreak{}M174:\allowbreak{}M174TMP:\allowbreak{}335785} & \textbf{(630 -\kern0pt{}- 730) - i (260 -\kern0pt{}- 340)} MeV\par Presentación: (630-\kern0pt{}-730)\allowbreak{}-i(260-\kern0pt{}-340) & \(u_+\): NA\par \(u_-\): NA\par Tipo: R\par Confianza: ND & COMP-MESON\par E04 / M08\\\addlinespace[3pt]
% VI-CATALOG-ROW masas 327
\textbf{327. K\_\allowbreak{}0*(700) Breit-Wigner Mass}\par K\_\allowbreak{}0\textasciicircum{}*(700)\allowbreak{}0 | Kbar\_\allowbreak{}0\textasciicircum{}*(700)\allowbreak{}0 | K\_\allowbreak{}0\textasciicircum{}*(700)\allowbreak{}+ | K\_\allowbreak{}0\textasciicircum{}*(700)\allowbreak{}-\par {\fontsize{8.1}{9.4}\selectfont PDG2026:\allowbreak{}observable:\allowbreak{}M174:\allowbreak{}M174M:\allowbreak{}335786} & \textbf{838+-11} MeV & \(u_+\): 11.38790\allowbreak{}40790194\allowbreak{}1\par \(u_-\): 11.29733\allowbreak{}95209359\par Tipo: NONE\par Confianza: ND & COMP-MESON\par E04 / M08\\\addlinespace[3pt]
% VI-CATALOG-ROW masas 328
\textbf{328. eta(1475) MASS}\par eta(1475)\allowbreak{}0\par {\fontsize{8.1}{9.4}\selectfont PDG2026:\allowbreak{}observable:\allowbreak{}M175:\allowbreak{}M175M5:\allowbreak{}334616} & \textbf{1476+-4} MeV & \(u_+\): 3.929075\allowbreak{}02750456\allowbreak{}58\par \(u_-\): 4.010312\allowbreak{}73484770\allowbreak{}26\par Tipo: NONE\par Confianza: ND & COMP-MESON\par E04 / M08\\\addlinespace[3pt]
% VI-CATALOG-ROW masas 329
\textbf{329. chi\_\allowbreak{}c1(3872) MASS FROM J/\allowbreak{}psi X MODE}\par chi\_\allowbreak{}c1(3872)\par {\fontsize{8.1}{9.4}\selectfont PDG2026:\allowbreak{}observable:\allowbreak{}M176:\allowbreak{}M176M:\allowbreak{}346130} & \textbf{3871.64+-0.06} MeV & \(u_+\): 0.057098\allowbreak{}85902177\allowbreak{}2717\par \(u_-\): 0.057095\allowbreak{}83843793\allowbreak{}0757\par Tipo: NONE\par Confianza: ND & COMP-MESON\par E04 / M08\\\addlinespace[3pt]
% VI-CATALOG-ROW masas 330
\textbf{330. Upsilon\_\allowbreak{}2(1D) MASS}\par Upsilon\_\allowbreak{}2(1D)\par {\fontsize{8.1}{9.4}\selectfont PDG2026:\allowbreak{}observable:\allowbreak{}M177:\allowbreak{}M177M:\allowbreak{}347489} & \textbf{10163.7+-1.4} MeV & \(u_+\): 1.438601\allowbreak{}05887897\allowbreak{}99\par \(u_-\): 1.438601\allowbreak{}05887897\allowbreak{}99\par Tipo: NONE\par Confianza: ND & COMP-MESON\par E04 / M08\\\addlinespace[3pt]
% VI-CATALOG-ROW masas 331
\textbf{331. D\_\allowbreak{}1(2430)\allowbreak{}0 MASS}\par D\_\allowbreak{}1(2430)\allowbreak{}0 | Dbar\_\allowbreak{}1(2430)\allowbreak{}0\par {\fontsize{8.1}{9.4}\selectfont PDG2026:\allowbreak{}observable:\allowbreak{}M180:\allowbreak{}M180M:\allowbreak{}337827} & \textbf{2412+-9} MeV & \(u_+\): 9.174794\allowbreak{}34026708\allowbreak{}77\par \(u_-\): 9.174794\allowbreak{}34026708\allowbreak{}77\par Tipo: NONE\par Confianza: ND & COMP-MESON\par E04 / M08\\\addlinespace[3pt]
% VI-CATALOG-ROW masas 332
\textbf{332. psi(4360) MASS}\par psi(4360)\par {\fontsize{8.1}{9.4}\selectfont PDG2026:\allowbreak{}observable:\allowbreak{}M181:\allowbreak{}M181M:\allowbreak{}346635} & \textbf{4373+-7} MeV & \(u_+\): 7.026702\allowbreak{}99984180\allowbreak{}07\par \(u_-\): 7.126338\allowbreak{}98599779\allowbreak{}33\par Tipo: NONE\par Confianza: ND & COMP-MESON\par E04 / M08\\\addlinespace[3pt]
% VI-CATALOG-ROW masas 333
\textbf{333. D\_\allowbreak{}s1*(2700)\allowbreak{}+ MASS}\par D\_\allowbreak{}s1\textasciicircum{}*(2700)\allowbreak{}+ | D\_\allowbreak{}s1\textasciicircum{}*(2700)\allowbreak{}-\par {\fontsize{8.1}{9.4}\selectfont PDG2026:\allowbreak{}observable:\allowbreak{}M182:\allowbreak{}M182M:\allowbreak{}338524} & \textbf{2714+-5} MeV & \(u_+\): 5.337526\allowbreak{}18609189\allowbreak{}34\par \(u_-\): 4.578047\allowbreak{}26750868\allowbreak{}96\par Tipo: NONE\par Confianza: ND & COMP-MESON\par E04 / M08\\\addlinespace[3pt]
% VI-CATALOG-ROW masas 334
\textbf{334. pi\_\allowbreak{}2(1880) MASS}\par pi\_\allowbreak{}2(1880)\allowbreak{}0\par {\fontsize{8.1}{9.4}\selectfont PDG2026:\allowbreak{}observable:\allowbreak{}M185:\allowbreak{}M185M:\allowbreak{}335043} & \textbf{1874+26-5} MeV & \(u_+\): 25.89794\allowbreak{}60588110\allowbreak{}09\par \(u_-\): 4.798815\allowbreak{}66383704\allowbreak{}16\par Tipo: NONE\par Confianza: ND & COMP-MESON\par E04 / M08\\\addlinespace[3pt]
% VI-CATALOG-ROW masas 335
\textbf{335. B\_\allowbreak{}s2*(5840)\allowbreak{}0 MASS}\par B\_\allowbreak{}s2\textasciicircum{}*(5840)\allowbreak{}0 | Bbar\_\allowbreak{}s2\textasciicircum{}*(5840)\allowbreak{}0\par {\fontsize{8.1}{9.4}\selectfont PDG2026:\allowbreak{}observable:\allowbreak{}M186:\allowbreak{}M186M:\allowbreak{}343137} & \textbf{5839.88+-0.12} MeV & \(u_+\): 0.115984\allowbreak{}22775051\allowbreak{}18\par \(u_-\): 0.115984\allowbreak{}22775051\allowbreak{}18\par Tipo: NONE\par Confianza: ND & COMP-MESON\par E04 / M08\\\addlinespace[3pt]
% VI-CATALOG-ROW masas 336
\textbf{336. B\_\allowbreak{}s1(5830)\allowbreak{}0 MASS}\par B\_\allowbreak{}s1(5830)\allowbreak{}0 | Bbar\_\allowbreak{}s1(5830)\allowbreak{}0\par {\fontsize{8.1}{9.4}\selectfont PDG2026:\allowbreak{}observable:\allowbreak{}M187:\allowbreak{}M187M:\allowbreak{}343129} & \textbf{5828.73+-0.20} MeV & \(u_+\): 0.196515\allowbreak{}51951131\allowbreak{}191\par \(u_-\): 0.196515\allowbreak{}51951131\allowbreak{}191\par Tipo: NONE\par Confianza: ND & COMP-MESON\par E04 / M08\\\addlinespace[3pt]
% VI-CATALOG-ROW masas 337
\textbf{337. rho(1570) MASS}\par rho(1570)\allowbreak{}0\par {\fontsize{8.1}{9.4}\selectfont PDG2026:\allowbreak{}observable:\allowbreak{}M188:\allowbreak{}M188M:\allowbreak{}334714} & \textbf{1570+-70} MeV & \(u_+\): 71.69379\allowbreak{}33157396\allowbreak{}83\par \(u_-\): 71.69379\allowbreak{}33157396\allowbreak{}83\par Tipo: NONE\par Confianza: ND & COMP-MESON\par E04 / M08\\\addlinespace[3pt]
% VI-CATALOG-ROW masas 338
\textbf{338. psi(4660) MASS}\par psi(4660)\par {\fontsize{8.1}{9.4}\selectfont PDG2026:\allowbreak{}observable:\allowbreak{}M189:\allowbreak{}M189M:\allowbreak{}346767} & \textbf{4623+-10} MeV & \(u_+\): 10.20904\allowbreak{}49422513\allowbreak{}1\par \(u_-\): 10.12797\allowbreak{}27634031\allowbreak{}61\par Tipo: NONE\par Confianza: ND & COMP-MESON\par E04 / M08\\\addlinespace[3pt]
% VI-CATALOG-ROW masas 339
\textbf{339. X(4160) MASS}\par X(4160)\par {\fontsize{8.1}{9.4}\selectfont PDG2026:\allowbreak{}observable:\allowbreak{}M190:\allowbreak{}M190M:\allowbreak{}346451} & \textbf{4153+23-21} MeV & \(u_+\): 23.03765\allowbreak{}48187802\allowbreak{}9\par \(u_-\): 20.81656\allowbreak{}17804000\allowbreak{}72\par Tipo: NONE\par Confianza: ND & COMP-MESON\par E04 / M08\\\addlinespace[3pt]
% VI-CATALOG-ROW masas 340
\textbf{340. T\_\allowbreak{}c\_\allowbreak{}cbar(4050)\allowbreak{}+ MASS}\par T\_\allowbreak{}c\_\allowbreak{}cbar(4050)\allowbreak{}+\par {\fontsize{8.1}{9.4}\selectfont PDG2026:\allowbreak{}observable:\allowbreak{}M191:\allowbreak{}M191M:\allowbreak{}347990} & \textbf{4051+24-40} MeV & \(u_+\): 24.41311\allowbreak{}12314674\allowbreak{}01\par \(u_-\): 43.32435\allowbreak{}80448689\allowbreak{}37\par Tipo: NONE\par Confianza: ND & COMP-MESON\par E04 / M08\\\addlinespace[3pt]
% VI-CATALOG-ROW masas 341
\textbf{341. T\_\allowbreak{}c\_\allowbreak{}cbar(4250)\allowbreak{}+ MASS}\par T\_\allowbreak{}c\_\allowbreak{}cbar(4250)\allowbreak{}+\par {\fontsize{8.1}{9.4}\selectfont PDG2026:\allowbreak{}observable:\allowbreak{}M192:\allowbreak{}M192M:\allowbreak{}348026} & \textbf{4250+190-50} MeV & \(u_+\): 185.2997\allowbreak{}57150407\allowbreak{}49\par \(u_-\): 45.45327\allowbreak{}27094540\allowbreak{}48\par Tipo: NONE\par Confianza: ND & COMP-MESON\par E04 / M08\\\addlinespace[3pt]
% VI-CATALOG-ROW masas 342
\textbf{342. chi\_\allowbreak{}c1(4140) MASS}\par chi\_\allowbreak{}c1(4140)\par {\fontsize{8.1}{9.4}\selectfont PDG2026:\allowbreak{}observable:\allowbreak{}M193:\allowbreak{}M193M:\allowbreak{}346336} & \textbf{4146.5+-3.0} MeV & \(u_+\): 2.988542\allowbreak{}41183424\allowbreak{}7\par \(u_-\): 3.001615\allowbreak{}33279208\allowbreak{}29\par Tipo: NONE\par Confianza: ND & COMP-MESON\par E04 / M08\\\addlinespace[3pt]
% VI-CATALOG-ROW masas 343
\textbf{343. X(4350) MASS}\par X(4350)\par {\fontsize{8.1}{9.4}\selectfont PDG2026:\allowbreak{}observable:\allowbreak{}M194:\allowbreak{}M194M:\allowbreak{}346628} & \textbf{4351+-5} MeV & \(u_+\): 4.652956\allowbreak{}04965273\allowbreak{}67\par \(u_-\): 5.147815\allowbreak{}07049350\allowbreak{}04\par Tipo: NONE\par Confianza: ND & COMP-MESON\par E04 / M08\\\addlinespace[3pt]
% VI-CATALOG-ROW masas 344
\textbf{344. T\_\allowbreak{}c\_\allowbreak{}cbar\_\allowbreak{}1(4430)\allowbreak{}+ MASS}\par T\_\allowbreak{}c\_\allowbreak{}cbar\_\allowbreak{}1(4430)\allowbreak{}+\par {\fontsize{8.1}{9.4}\selectfont PDG2026:\allowbreak{}observable:\allowbreak{}M195:\allowbreak{}M195M:\allowbreak{}348030} & \textbf{4478+15-17} MeV & \(u_+\): 14.76917\allowbreak{}01609154\allowbreak{}1\par \(u_-\): 17.45344\allowbreak{}35216461\allowbreak{}89\par Tipo: NONE\par Confianza: ND & COMP-MESON\par E04 / M08\\\addlinespace[3pt]
% VI-CATALOG-ROW masas 345
\textbf{345. D\_\allowbreak{}s1*(2860)\allowbreak{}+ MASS}\par D\_\allowbreak{}s1\textasciicircum{}*(2860)\allowbreak{}+ | D\_\allowbreak{}s1\textasciicircum{}*(2860)\allowbreak{}-\par {\fontsize{8.1}{9.4}\selectfont PDG2026:\allowbreak{}observable:\allowbreak{}M196:\allowbreak{}M196M:\allowbreak{}338533} & \textbf{2859+-27} MeV & \(u_+\): 26.83281\allowbreak{}57299974\allowbreak{}81\par \(u_-\): 26.83281\allowbreak{}57299974\allowbreak{}81\par Tipo: NONE\par Confianza: ND & COMP-MESON\par E04 / M08\\\addlinespace[3pt]
% VI-CATALOG-ROW masas 346
\textbf{346. D\_\allowbreak{}sJ(3040)\allowbreak{}+ MASS}\par D\_\allowbreak{}sJ(3040)\allowbreak{}+ | D\_\allowbreak{}sJ(3040)\allowbreak{}-\par {\fontsize{8.1}{9.4}\selectfont PDG2026:\allowbreak{}observable:\allowbreak{}M197:\allowbreak{}M197M:\allowbreak{}338544} & \textbf{3044+31-9} MeV & \(u_+\): 31.04834\allowbreak{}93925200\allowbreak{}48\par \(u_-\): 9.433981\allowbreak{}13205660\allowbreak{}32\par Tipo: NONE\par Confianza: ND & COMP-MESON\par E04 / M08\\\addlinespace[3pt]
% VI-CATALOG-ROW masas 347
\textbf{347. D\_\allowbreak{}0(2550)\allowbreak{}0 MASS}\par D\_\allowbreak{}0(2550)\allowbreak{}0 | Dbar\_\allowbreak{}0(2550)\allowbreak{}0 | D\_\allowbreak{}0(2550)\allowbreak{}+ | D\_\allowbreak{}0(2550)\allowbreak{}-\par {\fontsize{8.1}{9.4}\selectfont PDG2026:\allowbreak{}observable:\allowbreak{}M198:\allowbreak{}M198M:\allowbreak{}337843} & \textbf{2549+-19} MeV & \(u_+\): 19.00400\allowbreak{}28413363\allowbreak{}41\par \(u_-\): 19.00400\allowbreak{}28413363\allowbreak{}41\par Tipo: NONE\par Confianza: ND & COMP-MESON\par E04 / M08\\\addlinespace[3pt]
% VI-CATALOG-ROW masas 348
\textbf{348. D\_\allowbreak{}1*(2600)\allowbreak{}0 MASS}\par D\_\allowbreak{}1\textasciicircum{}*(2600)\allowbreak{}0 | Dbar\_\allowbreak{}1\textasciicircum{}*(2600)\allowbreak{}0 | D\_\allowbreak{}1\textasciicircum{}*(2600)\allowbreak{}+ | D\_\allowbreak{}1\textasciicircum{}*(2600)\allowbreak{}-\par {\fontsize{8.1}{9.4}\selectfont PDG2026:\allowbreak{}observable:\allowbreak{}M199:\allowbreak{}M199M:\allowbreak{}337846} & \textbf{2627+-10} MeV & \(u_+\): 9.812824\allowbreak{}44363102\allowbreak{}79\par \(u_-\): 9.812824\allowbreak{}44363102\allowbreak{}79\par Tipo: NONE\par Confianza: ND & COMP-MESON\par E04 / M08\\\addlinespace[3pt]
% VI-CATALOG-ROW masas 349
\textbf{349. eta\_\allowbreak{}b(2S) MASS}\par eta\_\allowbreak{}b(2S)\par {\fontsize{8.1}{9.4}\selectfont PDG2026:\allowbreak{}observable:\allowbreak{}M200:\allowbreak{}M200M:\allowbreak{}347281} & \textbf{9999+-4} MeV & \(u_+\): 4.482186\allowbreak{}96620299\allowbreak{}36\par \(u_-\): 3.982461\allowbreak{}55034797\allowbreak{}5\par Tipo: NONE\par Confianza: ND & COMP-MESON\par E04 / M08\\\addlinespace[3pt]
% VI-CATALOG-ROW masas 350
\textbf{350. D\_\allowbreak{}3*(2750) MASS}\par D\_\allowbreak{}3\textasciicircum{}*(2750)\allowbreak{}0 | Dbar\_\allowbreak{}3\textasciicircum{}*(2750)\allowbreak{}0 | D\_\allowbreak{}3\textasciicircum{}*(2750)\allowbreak{}+ | D\_\allowbreak{}3\textasciicircum{}*(2750)\allowbreak{}-\par {\fontsize{8.1}{9.4}\selectfont PDG2026:\allowbreak{}observable:\allowbreak{}M203:\allowbreak{}M203M:\allowbreak{}337860} & \textbf{2763.1+-3.2} MeV & \(u_+\): 3.190853\allowbreak{}31467817\allowbreak{}88\par \(u_-\): 3.190853\allowbreak{}31467817\allowbreak{}88\par Tipo: NONE\par Confianza: ND & COMP-MESON\par E04 / M08\\\addlinespace[3pt]
% VI-CATALOG-ROW masas 351
\textbf{351. h\_\allowbreak{}b(1P) MASS}\par h\_\allowbreak{}b(1P)\par {\fontsize{8.1}{9.4}\selectfont PDG2026:\allowbreak{}observable:\allowbreak{}M204:\allowbreak{}M204M:\allowbreak{}347229} & \textbf{9899.3+-0.8} MeV & \(u_+\): 0.750768\allowbreak{}96116034\allowbreak{}797\par \(u_-\): 0.750768\allowbreak{}96116034\allowbreak{}797\par Tipo: NONE\par Confianza: ND & COMP-MESON\par E04 / M08\\\addlinespace[3pt]
% VI-CATALOG-ROW masas 352
\textbf{352. h\_\allowbreak{}b(2P) MASS}\par h\_\allowbreak{}b(2P)\par {\fontsize{8.1}{9.4}\selectfont PDG2026:\allowbreak{}observable:\allowbreak{}M205:\allowbreak{}M205M:\allowbreak{}347593} & \textbf{10259.8+-1.2} MeV & \(u_+\): 1.208304\allowbreak{}59735945\allowbreak{}71\par \(u_-\): 1.208304\allowbreak{}59735945\allowbreak{}71\par Tipo: NONE\par Confianza: ND & COMP-MESON\par E04 / M08\\\addlinespace[3pt]
% VI-CATALOG-ROW masas 353
\textbf{353. chi\_\allowbreak{}b1(3P) MASS}\par chi\_\allowbreak{}b1(3P)\par {\fontsize{8.1}{9.4}\selectfont PDG2026:\allowbreak{}observable:\allowbreak{}M206:\allowbreak{}M206M:\allowbreak{}347734} & \textbf{10513.4+-0.7} MeV & \(u_+\): 0.670074\allowbreak{}62271003\allowbreak{}815\par \(u_-\): 0.670074\allowbreak{}62271003\allowbreak{}815\par Tipo: NONE\par Confianza: ND & COMP-MESON\par E04 / M08\\\addlinespace[3pt]
% VI-CATALOG-ROW masas 354
\textbf{354. T\_\allowbreak{}b\_\allowbreak{}bbar\_\allowbreak{}1(10610)\allowbreak{}+- MASS}\par T\_\allowbreak{}b\_\allowbreak{}bbar\_\allowbreak{}1(10610)\allowbreak{}+\par {\fontsize{8.1}{9.4}\selectfont PDG2026:\allowbreak{}observable:\allowbreak{}M207:\allowbreak{}M207M:\allowbreak{}348052} & \textbf{10607.2+-2.0} MeV & \(u_+\): 2\par \(u_-\): 2\par Tipo: NONE\par Confianza: ND & COMP-MESON\par E04 / M08\\\addlinespace[3pt]
% VI-CATALOG-ROW masas 355
\textbf{355. T\_\allowbreak{}b\_\allowbreak{}bbar\_\allowbreak{}1(10610)\allowbreak{}0 MASS}\par T\_\allowbreak{}b\_\allowbreak{}bbar\_\allowbreak{}1(10610)\allowbreak{}+\par {\fontsize{8.1}{9.4}\selectfont PDG2026:\allowbreak{}observable:\allowbreak{}M207:\allowbreak{}M207M0:\allowbreak{}348053} & \textbf{10609+-6} MeV & \(u_+\): 5.656854\allowbreak{}24949238\allowbreak{}06\par \(u_-\): 5.656854\allowbreak{}24949238\allowbreak{}06\par Tipo: NONE\par Confianza: ND & COMP-MESON\par E04 / M08\\\addlinespace[3pt]
% VI-CATALOG-ROW masas 356
\textbf{356. T\_\allowbreak{}b\_\allowbreak{}bbar\_\allowbreak{}1(10650)\allowbreak{}+ MASS}\par T\_\allowbreak{}b\_\allowbreak{}bbar\_\allowbreak{}1(10650)\allowbreak{}+\par {\fontsize{8.1}{9.4}\selectfont PDG2026:\allowbreak{}observable:\allowbreak{}M208:\allowbreak{}M208M:\allowbreak{}348075} & \textbf{10652.2+-1.5} MeV & \(u_+\): 1.5\par \(u_-\): 1.5\par Tipo: NONE\par Confianza: ND & COMP-MESON\par E04 / M08\\\addlinespace[3pt]
% VI-CATALOG-ROW masas 357
\textbf{357. T\_\allowbreak{}c\_\allowbreak{}cbar\_\allowbreak{}1(3900) MASS}\par T\_\allowbreak{}c\_\allowbreak{}cbar\_\allowbreak{}1(3900)\allowbreak{}+\par {\fontsize{8.1}{9.4}\selectfont PDG2026:\allowbreak{}observable:\allowbreak{}M210:\allowbreak{}M210M:\allowbreak{}347938} & \textbf{3886.7+-2.3} MeV & \(u_+\): 2.293369\allowbreak{}19741447\allowbreak{}7\par \(u_-\): 2.273282\allowbreak{}28957628\allowbreak{}7\par Tipo: NONE\par Confianza: ND & COMP-MESON\par E04 / M08\\\addlinespace[3pt]
% VI-CATALOG-ROW masas 358
\textbf{358. psi\_\allowbreak{}2(3823) MASS}\par psi\_\allowbreak{}2(3823)\par {\fontsize{8.1}{9.4}\selectfont PDG2026:\allowbreak{}observable:\allowbreak{}M212:\allowbreak{}M212M:\allowbreak{}346100} & \textbf{3823.51+-0.34} MeV & \(u_+\): 0.339047\allowbreak{}06932929\allowbreak{}501\par \(u_-\): 0.339047\allowbreak{}06932929\allowbreak{}501\par Tipo: NONE\par Confianza: ND & COMP-MESON\par E04 / M08\\\addlinespace[3pt]
% VI-CATALOG-ROW masas 359
\textbf{359. T\_\allowbreak{}c\_\allowbreak{}cbar(4020) MASS}\par T\_\allowbreak{}c\_\allowbreak{}cbar(4020)\allowbreak{}+\par {\fontsize{8.1}{9.4}\selectfont PDG2026:\allowbreak{}observable:\allowbreak{}M213:\allowbreak{}M213M:\allowbreak{}347976} & \textbf{4024.1+-1.9} MeV & \(u_+\): 1.824690\allowbreak{}63465848\allowbreak{}7\par \(u_-\): 1.967162\allowbreak{}32389574\allowbreak{}41\par Tipo: NONE\par Confianza: ND & COMP-MESON\par E04 / M08\\\addlinespace[3pt]
% VI-CATALOG-ROW masas 360
\textbf{360. T\_\allowbreak{}c\_\allowbreak{}cbar\_\allowbreak{}0(4240)\allowbreak{}+ MASS}\par T\_\allowbreak{}c\_\allowbreak{}cbar\_\allowbreak{}0(4240)\allowbreak{}0\par {\fontsize{8.1}{9.4}\selectfont PDG2026:\allowbreak{}observable:\allowbreak{}M216:\allowbreak{}M216M:\allowbreak{}348022} & \textbf{4239+50-21} MeV & \(u_+\): 48.46648\allowbreak{}32642105\allowbreak{}38\par \(u_-\): 20.59126\allowbreak{}02819740\allowbreak{}02\par Tipo: NONE\par Confianza: ND & COMP-MESON\par E04 / M08\\\addlinespace[3pt]
% VI-CATALOG-ROW masas 361
\textbf{361. B\_\allowbreak{}c(2S)\allowbreak{}+- MASS}\par B\_\allowbreak{}c(2S)\allowbreak{}+ | B\_\allowbreak{}c(2S)\allowbreak{}-\par {\fontsize{8.1}{9.4}\selectfont PDG2026:\allowbreak{}observable:\allowbreak{}M217:\allowbreak{}M217M:\allowbreak{}343283} & \textbf{6871.2+-1.0} MeV & \(u_+\): 1.018454\allowbreak{}58192780\allowbreak{}3\par \(u_-\): 1.018454\allowbreak{}58192780\allowbreak{}3\par Tipo: NONE\par Confianza: ND & COMP-MESON\par E04 / M08\\\addlinespace[3pt]
% VI-CATALOG-ROW masas 362
\textbf{362. T\_\allowbreak{}c\_\allowbreak{}cbar(4055)\allowbreak{}+ MASS}\par T\_\allowbreak{}c\_\allowbreak{}cbar(4055)\allowbreak{}+\par {\fontsize{8.1}{9.4}\selectfont PDG2026:\allowbreak{}observable:\allowbreak{}M223:\allowbreak{}M223M:\allowbreak{}348000} & \textbf{4054+-3.2} MeV & \(u_+\): 3.162277\allowbreak{}66016838\par \(u_-\): 3.162277\allowbreak{}66016838\par Tipo: NONE\par Confianza: ND & COMP-MESON\par E04 / M08\\\addlinespace[3pt]
% VI-CATALOG-ROW masas 363
\textbf{363. D\_\allowbreak{}s3*(2860)\allowbreak{}+ MASS}\par D\_\allowbreak{}s3\textasciicircum{}*(2860)\allowbreak{}+ | D\_\allowbreak{}s3\textasciicircum{}*(2860)\allowbreak{}-\par {\fontsize{8.1}{9.4}\selectfont PDG2026:\allowbreak{}observable:\allowbreak{}M226:\allowbreak{}M226M:\allowbreak{}338535} & \textbf{2860+-7} MeV & \(u_+\): 7.000714\allowbreak{}24927485\allowbreak{}53\par \(u_-\): 7.000714\allowbreak{}24927485\allowbreak{}53\par Tipo: NONE\par Confianza: ND & COMP-MESON\par E04 / M08\\\addlinespace[3pt]
% VI-CATALOG-ROW masas 364
\textbf{364. a\_\allowbreak{}0(1950) MASS}\par a\_\allowbreak{}0(1950)\allowbreak{}0\par {\fontsize{8.1}{9.4}\selectfont PDG2026:\allowbreak{}observable:\allowbreak{}M227:\allowbreak{}M227M:\allowbreak{}335070} & \textbf{1931+-26} MeV & \(u_+\): 26.07680\allowbreak{}96208106\par \(u_-\): 26.07680\allowbreak{}96208106\par Tipo: NONE\par Confianza: ND & COMP-MESON\par E04 / M08\\\addlinespace[3pt]
% VI-CATALOG-ROW masas 365
\textbf{365. D\_\allowbreak{}2(2740)\allowbreak{}0 MASS}\par D\_\allowbreak{}2(2740)\allowbreak{}0 | Dbar\_\allowbreak{}2(2740)\allowbreak{}0 | D\_\allowbreak{}2(2740)\allowbreak{}+ | D\_\allowbreak{}2(2740)\allowbreak{}-\par {\fontsize{8.1}{9.4}\selectfont PDG2026:\allowbreak{}observable:\allowbreak{}M228:\allowbreak{}M228M:\allowbreak{}337857} & \textbf{2747+-6} MeV & \(u_+\): 6.393293\allowbreak{}07322381\allowbreak{}27\par \(u_-\): 6.393293\allowbreak{}07322381\allowbreak{}27\par Tipo: NONE\par Confianza: ND & COMP-MESON\par E04 / M08\\\addlinespace[3pt]
% VI-CATALOG-ROW masas 366
\textbf{366. D\_\allowbreak{}2*(3000)\allowbreak{}0 MASS}\par D\_\allowbreak{}2\textasciicircum{}*(3000)\allowbreak{}0 | Dbar\_\allowbreak{}2\textasciicircum{}*(3000)\allowbreak{}0 | D\_\allowbreak{}2\textasciicircum{}*(3000)\allowbreak{}+ | D\_\allowbreak{}2\textasciicircum{}*(3000)\allowbreak{}-\par {\fontsize{8.1}{9.4}\selectfont PDG2026:\allowbreak{}observable:\allowbreak{}M229:\allowbreak{}M229M:\allowbreak{}337872} & \textbf{3210+-60} MeV & \(u_+\): 56.93856\allowbreak{}33819470\allowbreak{}27\par \(u_-\): 56.93856\allowbreak{}33819470\allowbreak{}27\par Tipo: NONE\par Confianza: ND & COMP-MESON\par E04 / M08\\\addlinespace[3pt]
% VI-CATALOG-ROW masas 367
\textbf{367. T\_\allowbreak{}c\_\allowbreak{}cbar\_\allowbreak{}1(4200)\allowbreak{}+ MASS}\par T\_\allowbreak{}c\_\allowbreak{}cbar\_\allowbreak{}1(4200)\allowbreak{}+\par {\fontsize{8.1}{9.4}\selectfont PDG2026:\allowbreak{}observable:\allowbreak{}M231:\allowbreak{}M231M:\allowbreak{}348012} & \textbf{4242+-26} MeV & \(u_+\): 26.30683\allowbreak{}90095676\allowbreak{}29\par \(u_-\): 25.56726\allowbreak{}18453510\allowbreak{}91\par Tipo: NONE\par Confianza: ND & COMP-MESON\par E04 / M08\\\addlinespace[3pt]
% VI-CATALOG-ROW masas 368
\textbf{368. T\_\allowbreak{}b\_\allowbreak{}sbar(5568)\allowbreak{}+ MASS}\par T\_\allowbreak{}b\_\allowbreak{}sbar(5568)\allowbreak{}+\par {\fontsize{8.1}{9.4}\selectfont PDG2026:\allowbreak{}observable:\allowbreak{}M232:\allowbreak{}M232M:\allowbreak{}348036} & \textbf{5566.9+-3.3} MeV & \(u_+\): 3.255764\allowbreak{}11921994\allowbreak{}2\par \(u_-\): 3.324154\allowbreak{}02771893\allowbreak{}3\par Tipo: NONE\par Confianza: ND & COMP-MESON\par E04 / M08\\\addlinespace[3pt]
% VI-CATALOG-ROW masas 369
\textbf{369. chi\_\allowbreak{}c1(4274) MASS}\par chi\_\allowbreak{}c1(4274)\par {\fontsize{8.1}{9.4}\selectfont PDG2026:\allowbreak{}observable:\allowbreak{}M233:\allowbreak{}M233M:\allowbreak{}346624} & \textbf{4290+-6} MeV & \(u_+\): 5.682397\allowbreak{}90637520\allowbreak{}86\par \(u_-\): 6.461081\allowbreak{}98239139\allowbreak{}5\par Tipo: NONE\par Confianza: ND & COMP-MESON\par E04 / M08\\\addlinespace[3pt]
% VI-CATALOG-ROW masas 370
\textbf{370. chi\_\allowbreak{}c0(4500) MASS}\par chi\_\allowbreak{}c0(4500)\par {\fontsize{8.1}{9.4}\selectfont PDG2026:\allowbreak{}observable:\allowbreak{}M234:\allowbreak{}M234M:\allowbreak{}346755} & \textbf{4484+-12} MeV & \(u_+\): 12.00282\allowbreak{}00724577\allowbreak{}39\par \(u_-\): 11.98528\allowbreak{}45503755\allowbreak{}41\par Tipo: NONE\par Confianza: ND & COMP-MESON\par E04 / M08\\\addlinespace[3pt]
% VI-CATALOG-ROW masas 371
\textbf{371. chi\_\allowbreak{}c0(4700) MASS}\par chi\_\allowbreak{}c0(4700)\par {\fontsize{8.1}{9.4}\selectfont PDG2026:\allowbreak{}observable:\allowbreak{}M235:\allowbreak{}M235M:\allowbreak{}346816} & \textbf{4708+6-4} MeV & \(u_+\): 5.969036\allowbreak{}18444770\allowbreak{}42\par \(u_-\): 3.940850\allowbreak{}54655756\allowbreak{}12\par Tipo: NONE\par Confianza: ND & COMP-MESON\par E04 / M08\\\addlinespace[3pt]
% VI-CATALOG-ROW masas 372
\textbf{372. chi\_\allowbreak{}c0(3860) MASS}\par chi\_\allowbreak{}c0(3860)\par {\fontsize{8.1}{9.4}\selectfont PDG2026:\allowbreak{}observable:\allowbreak{}M237:\allowbreak{}M237M:\allowbreak{}346122} & \textbf{3862+50-35} MeV & \(u_+\): 47.70744\allowbreak{}17675062\allowbreak{}53\par \(u_-\): 34.53983\allowbreak{}20783410\allowbreak{}92\par Tipo: NONE\par Confianza: ND & COMP-MESON\par E04 / M08\\\addlinespace[3pt]
% VI-CATALOG-ROW masas 373
\textbf{373. chi\_\allowbreak{}b2(3P) MASS}\par chi\_\allowbreak{}b2(3P)\par {\fontsize{8.1}{9.4}\selectfont PDG2026:\allowbreak{}observable:\allowbreak{}M238:\allowbreak{}M238M:\allowbreak{}347741} & \textbf{10524.0+-0.8} MeV & \(u_+\): 0.778331\allowbreak{}54889160\allowbreak{}183\par \(u_-\): 0.778331\allowbreak{}54889160\allowbreak{}183\par Tipo: NONE\par Confianza: ND & COMP-MESON\par E04 / M08\\\addlinespace[3pt]
% VI-CATALOG-ROW masas 374
\textbf{374. pi\_\allowbreak{}2(2005) MASS}\par pi\_\allowbreak{}2(2005)\allowbreak{}0\par {\fontsize{8.1}{9.4}\selectfont PDG2026:\allowbreak{}observable:\allowbreak{}M239:\allowbreak{}M239M:\allowbreak{}335106} & \textbf{1963+17-27} MeV & \(u_+\): 16.66354\allowbreak{}05259114\allowbreak{}29\par \(u_-\): 27.41832\allowbreak{}31822566\allowbreak{}39\par Tipo: NONE\par Confianza: ND & COMP-MESON\par E04 / M08\\\addlinespace[3pt]
% VI-CATALOG-ROW masas 375
\textbf{375. T\_\allowbreak{}c\_\allowbreak{}cbar(4100)\allowbreak{}+ MASS}\par T\_\allowbreak{}c\_\allowbreak{}cbar(4100)\allowbreak{}+\par {\fontsize{8.1}{9.4}\selectfont PDG2026:\allowbreak{}observable:\allowbreak{}M240:\allowbreak{}M240M:\allowbreak{}348006} & \textbf{4096+-28} MeV & \(u_+\): 26.90724\allowbreak{}80941474\allowbreak{}2\par \(u_-\): 29.73213\allowbreak{}74946370\allowbreak{}09\par Tipo: NONE\par Confianza: ND & COMP-MESON\par E04 / M08\\\addlinespace[3pt]
% VI-CATALOG-ROW masas 376
\textbf{376. psi\_\allowbreak{}3(3842) MASS}\par psi\_\allowbreak{}3(3842)\par {\fontsize{8.1}{9.4}\selectfont PDG2026:\allowbreak{}observable:\allowbreak{}M241:\allowbreak{}M241M:\allowbreak{}346116} & \textbf{3842.71+-0.20} MeV & \(u_+\): 0.200000\allowbreak{}00000000\allowbreak{}001\par \(u_-\): 0.200000\allowbreak{}00000000\allowbreak{}001\par Tipo: NONE\par Confianza: ND & COMP-MESON\par E04 / M08\\\addlinespace[3pt]
% VI-CATALOG-ROW masas 377
\textbf{377. Upsilon(10753) MASS}\par Upsilon(10753)\par {\fontsize{8.1}{9.4}\selectfont PDG2026:\allowbreak{}observable:\allowbreak{}M243:\allowbreak{}M243M:\allowbreak{}347800} & \textbf{10756.6+-2.8} MeV & \(u_+\): 2.846049\allowbreak{}89415154\allowbreak{}2\par \(u_-\): 2.846049\allowbreak{}89415154\allowbreak{}2\par Tipo: NONE\par Confianza: ND & COMP-MESON\par E04 / M08\\\addlinespace[3pt]
% VI-CATALOG-ROW masas 378
\textbf{378. B\_\allowbreak{}1(5721)\allowbreak{}+ mass}\par B\_\allowbreak{}1(5721)\allowbreak{}0 | Bbar\_\allowbreak{}1(5721)\allowbreak{}0 | B\_\allowbreak{}1(5721)\allowbreak{}+ | B\_\allowbreak{}1(5721)\allowbreak{}-\par {\fontsize{8.1}{9.4}\selectfont PDG2026:\allowbreak{}observable:\allowbreak{}M244:\allowbreak{}M244M+:\allowbreak{}342539} & \textbf{5726.0+2.5-2.7} MeV & \(u_+\): 2.455967\allowbreak{}77983158\allowbreak{}4\par \(u_-\): 2.729151\allowbreak{}40828024\allowbreak{}59\par Tipo: NONE\par Confianza: ND & COMP-MESON\par E04 / M08\\\addlinespace[3pt]
% VI-CATALOG-ROW masas 379
\textbf{379. B\_\allowbreak{}1(5721)\allowbreak{}0 mass}\par B\_\allowbreak{}1(5721)\allowbreak{}0 | Bbar\_\allowbreak{}1(5721)\allowbreak{}0 | B\_\allowbreak{}1(5721)\allowbreak{}+ | B\_\allowbreak{}1(5721)\allowbreak{}-\par {\fontsize{8.1}{9.4}\selectfont PDG2026:\allowbreak{}observable:\allowbreak{}M244:\allowbreak{}M244M0:\allowbreak{}342542} & \textbf{5726.1+-1.2} MeV & \(u_+\): 1.233029\allowbreak{}59923312\allowbreak{}99\par \(u_-\): 1.262925\allowbreak{}41194705\allowbreak{}9\par Tipo: NONE\par Confianza: ND & COMP-MESON\par E04 / M08\\\addlinespace[3pt]
% VI-CATALOG-ROW masas 380
\textbf{380. B\_\allowbreak{}2*(5747)\allowbreak{}+ mass}\par B\_\allowbreak{}2\textasciicircum{}*(5747)\allowbreak{}0 | Bbar\_\allowbreak{}2\textasciicircum{}*(5747)\allowbreak{}0 | B\_\allowbreak{}2\textasciicircum{}*(5747)\allowbreak{}+ | B\_\allowbreak{}2\textasciicircum{}*(5747)\allowbreak{}-\par {\fontsize{8.1}{9.4}\selectfont PDG2026:\allowbreak{}observable:\allowbreak{}M245:\allowbreak{}M245M+:\allowbreak{}342556} & \textbf{5737.3+-0.7} MeV & \(u_+\): 0.710601\allowbreak{}65855206\allowbreak{}67\par \(u_-\): 0.732065\allowbreak{}97699796\allowbreak{}105\par Tipo: NONE\par Confianza: ND & COMP-MESON\par E04 / M08\\\addlinespace[3pt]
% VI-CATALOG-ROW masas 381
\textbf{381. B\_\allowbreak{}2*(5747)\allowbreak{}0 mass}\par B\_\allowbreak{}2\textasciicircum{}*(5747)\allowbreak{}0 | Bbar\_\allowbreak{}2\textasciicircum{}*(5747)\allowbreak{}0 | B\_\allowbreak{}2\textasciicircum{}*(5747)\allowbreak{}+ | B\_\allowbreak{}2\textasciicircum{}*(5747)\allowbreak{}-\par {\fontsize{8.1}{9.4}\selectfont PDG2026:\allowbreak{}observable:\allowbreak{}M245:\allowbreak{}M245M0:\allowbreak{}342559} & \textbf{5739.6+-0.7} MeV & \(u_+\): 0.653465\allowbreak{}68058668\allowbreak{}519\par \(u_-\): 0.653465\allowbreak{}68058668\allowbreak{}519\par Tipo: NONE\par Confianza: ND & COMP-MESON\par E04 / M08\\\addlinespace[3pt]
% VI-CATALOG-ROW masas 382
\textbf{382. B\_\allowbreak{}J(5840)\allowbreak{}+ MASS}\par B\_\allowbreak{}J(5840)\allowbreak{}0 | Bbar\_\allowbreak{}J(5840)\allowbreak{}0 | B\_\allowbreak{}J(5840)\allowbreak{}+ | B\_\allowbreak{}J(5840)\allowbreak{}-\par {\fontsize{8.1}{9.4}\selectfont PDG2026:\allowbreak{}observable:\allowbreak{}M246:\allowbreak{}M246M+:\allowbreak{}342571} & \textbf{5851+-19} MeV & \(u_+\): 19.10515\allowbreak{}94213077\allowbreak{}78\par \(u_-\): 19.10515\allowbreak{}94213077\allowbreak{}78\par Tipo: NONE\par Confianza: ND & COMP-MESON\par E04 / M08\\\addlinespace[3pt]
% VI-CATALOG-ROW masas 383
\textbf{383. B\_\allowbreak{}J(5840)\allowbreak{}0 MASS}\par B\_\allowbreak{}J(5840)\allowbreak{}0 | Bbar\_\allowbreak{}J(5840)\allowbreak{}0 | B\_\allowbreak{}J(5840)\allowbreak{}+ | B\_\allowbreak{}J(5840)\allowbreak{}-\par {\fontsize{8.1}{9.4}\selectfont PDG2026:\allowbreak{}observable:\allowbreak{}M246:\allowbreak{}M246M0:\allowbreak{}342574} & \textbf{5863+-9} MeV & \(u_+\): 8.602618\allowbreak{}99078296\allowbreak{}47\par \(u_-\): 8.602618\allowbreak{}99078296\allowbreak{}47\par Tipo: NONE\par Confianza: ND & COMP-MESON\par E04 / M08\\\addlinespace[3pt]
% VI-CATALOG-ROW masas 384
\textbf{384. eta(2370) MASS}\par eta(2370)\allowbreak{}0\par {\fontsize{8.1}{9.4}\selectfont PDG2026:\allowbreak{}observable:\allowbreak{}M247:\allowbreak{}M247M:\allowbreak{}335261} & \textbf{2377+-9} MeV & \(u_+\): 8.807211\allowbreak{}58281389\allowbreak{}9\par \(u_-\): 9.654016\allowbreak{}33399925\par Tipo: NONE\par Confianza: ND & COMP-MESON\par E04 / M08\\\addlinespace[3pt]
% VI-CATALOG-ROW masas 385
\textbf{385. B\_\allowbreak{}J(5970)\allowbreak{}+ MASS}\par B\_\allowbreak{}J(5970)\allowbreak{}0 | Bbar\_\allowbreak{}J(5970)\allowbreak{}0 | B\_\allowbreak{}J(5970)\allowbreak{}+ | B\_\allowbreak{}J(5970)\allowbreak{}-\par {\fontsize{8.1}{9.4}\selectfont PDG2026:\allowbreak{}observable:\allowbreak{}M248:\allowbreak{}M248M+:\allowbreak{}342583} & \textbf{5965+-5} MeV & \(u_+\): 4.505372\allowbreak{}31097705\allowbreak{}78\par \(u_-\): 4.505372\allowbreak{}31097705\allowbreak{}78\par Tipo: NONE\par Confianza: ND & COMP-MESON\par E04 / M08\\\addlinespace[3pt]
% VI-CATALOG-ROW masas 386
\textbf{386. B\_\allowbreak{}J(5970)\allowbreak{}0 MASS}\par B\_\allowbreak{}J(5970)\allowbreak{}0 | Bbar\_\allowbreak{}J(5970)\allowbreak{}0 | B\_\allowbreak{}J(5970)\allowbreak{}+ | B\_\allowbreak{}J(5970)\allowbreak{}-\par {\fontsize{8.1}{9.4}\selectfont PDG2026:\allowbreak{}observable:\allowbreak{}M248:\allowbreak{}M248M0:\allowbreak{}342586} & \textbf{5971+-5} MeV & \(u_+\): 5.347986\allowbreak{}85867918\allowbreak{}88\par \(u_-\): 5.347986\allowbreak{}85867918\allowbreak{}88\par Tipo: NONE\par Confianza: ND & COMP-MESON\par E04 / M08\\\addlinespace[3pt]
% VI-CATALOG-ROW masas 387
\textbf{387. D\_\allowbreak{}1*(2760)\allowbreak{}0 MASS}\par D\_\allowbreak{}1\textasciicircum{}*(2760)\allowbreak{}0 | Dbar\_\allowbreak{}1\textasciicircum{}*(2760)\allowbreak{}0 | D\_\allowbreak{}1\textasciicircum{}*(2760)\allowbreak{}+ | D\_\allowbreak{}1\textasciicircum{}*(2760)\allowbreak{}-\par {\fontsize{8.1}{9.4}\selectfont PDG2026:\allowbreak{}observable:\allowbreak{}M249:\allowbreak{}M249M:\allowbreak{}337868} & \textbf{2781+-22} MeV & \(u_+\): 22.20360\allowbreak{}33111745\allowbreak{}18\par \(u_-\): 22.20360\allowbreak{}33111745\allowbreak{}18\par Tipo: NONE\par Confianza: ND & COMP-MESON\par E04 / M08\\\addlinespace[3pt]
% VI-CATALOG-ROW masas 388
\textbf{388. T\_\allowbreak{}cbar\_\allowbreak{}sbar\_\allowbreak{}0*(2870)\allowbreak{}0 MASS}\par T\_\allowbreak{}cbar\_\allowbreak{}sbar\_\allowbreak{}0\textasciicircum{}*(2870)\allowbreak{}0\par {\fontsize{8.1}{9.4}\selectfont PDG2026:\allowbreak{}observable:\allowbreak{}M250:\allowbreak{}M250M:\allowbreak{}347916} & \textbf{2874+-11} MeV & \(u_+\): 10.73712\allowbreak{}74384021\allowbreak{}71\par \(u_-\): 10.73712\allowbreak{}74384021\allowbreak{}71\par Tipo: NONE\par Confianza: ND & COMP-MESON\par E04 / M08\\\addlinespace[3pt]
% VI-CATALOG-ROW masas 389
\textbf{389. T\_\allowbreak{}cbar\_\allowbreak{}sbar\_\allowbreak{}1*(2900)\allowbreak{}0 MASS}\par T\_\allowbreak{}cbar\_\allowbreak{}sbar\_\allowbreak{}1\textasciicircum{}*(2900)\allowbreak{}0\par {\fontsize{8.1}{9.4}\selectfont PDG2026:\allowbreak{}observable:\allowbreak{}M251:\allowbreak{}M251M:\allowbreak{}347922} & \textbf{2904+-5} MeV & \(u_+\): 5.099019\allowbreak{}51359278\allowbreak{}36\par \(u_-\): 5.099019\allowbreak{}51359278\allowbreak{}36\par Tipo: NONE\par Confianza: ND & COMP-MESON\par E04 / M08\\\addlinespace[3pt]
% VI-CATALOG-ROW masas 390
\textbf{390. D\_\allowbreak{}0*(2300) MASS}\par D\_\allowbreak{}0\textasciicircum{}*(2300)\allowbreak{}0 | Dbar\_\allowbreak{}0\textasciicircum{}*(2300)\allowbreak{}0 | D\_\allowbreak{}0\textasciicircum{}*(2300)\allowbreak{}+ | D\_\allowbreak{}0\textasciicircum{}*(2300)\allowbreak{}-\par {\fontsize{8.1}{9.4}\selectfont PDG2026:\allowbreak{}observable:\allowbreak{}M252:\allowbreak{}M252M:\allowbreak{}337813} & \textbf{2343+-10} MeV & \(u_+\): 9.654442\allowbreak{}07944618\allowbreak{}43\par \(u_-\): 9.654442\allowbreak{}07944618\allowbreak{}43\par Tipo: NONE\par Confianza: ND & COMP-MESON\par E04 / M08\\\addlinespace[3pt]
% VI-CATALOG-ROW masas 391
\textbf{391. D\_\allowbreak{}1(2420) MASS}\par D\_\allowbreak{}1(2420)\allowbreak{}0 | Dbar\_\allowbreak{}1(2420)\allowbreak{}0 | D\_\allowbreak{}1(2420)\allowbreak{}+ | D\_\allowbreak{}1(2420)\allowbreak{}-\par {\fontsize{8.1}{9.4}\selectfont PDG2026:\allowbreak{}observable:\allowbreak{}M253:\allowbreak{}M253M:\allowbreak{}337817} & \textbf{2422.1+-0.6} MeV & \(u_+\): 0.575022\allowbreak{}82665337\allowbreak{}301\par \(u_-\): 0.575022\allowbreak{}82665337\allowbreak{}301\par Tipo: NONE\par Confianza: ND & COMP-MESON\par E04 / M08\\\addlinespace[3pt]
% VI-CATALOG-ROW masas 392
\textbf{392. D\_\allowbreak{}2*(2460) MASS}\par D\_\allowbreak{}2\textasciicircum{}*(2460)\allowbreak{}0 | Dbar\_\allowbreak{}2\textasciicircum{}*(2460)\allowbreak{}0 | D\_\allowbreak{}2\textasciicircum{}*(2460)\allowbreak{}+ | D\_\allowbreak{}2\textasciicircum{}*(2460)\allowbreak{}-\par {\fontsize{8.1}{9.4}\selectfont PDG2026:\allowbreak{}observable:\allowbreak{}M254:\allowbreak{}M254M:\allowbreak{}337830} & \textbf{2461.1+-0.8} MeV & \(u_+\): 0.797518\allowbreak{}42909338\allowbreak{}577\par \(u_-\): 0.797518\allowbreak{}42909338\allowbreak{}577\par Tipo: NONE\par Confianza: ND & COMP-MESON\par E04 / M08\\\addlinespace[3pt]
% VI-CATALOG-ROW masas 393
\textbf{393. X(1750) MASS}\par X(1750)\allowbreak{}0\par {\fontsize{8.1}{9.4}\selectfont PDG2026:\allowbreak{}observable:\allowbreak{}M255:\allowbreak{}M255M:\allowbreak{}334941} & \textbf{1753.8+-2.7} MeV & \(u_+\): 2.676722\allowbreak{}05618767\allowbreak{}3\par \(u_-\): 2.734124\allowbreak{}27080383\allowbreak{}49\par Tipo: NONE\par Confianza: ND & COMP-MESON\par E04 / M08\\\addlinespace[3pt]
% VI-CATALOG-ROW masas 394
\textbf{394. D\_\allowbreak{}s0(2590)\allowbreak{}+ MASS}\par D\_\allowbreak{}s0(2590)\allowbreak{}+ | D\_\allowbreak{}s0(2590)\allowbreak{}-\par {\fontsize{8.1}{9.4}\selectfont PDG2026:\allowbreak{}observable:\allowbreak{}M256:\allowbreak{}M256M:\allowbreak{}338520} & \textbf{2591+-9} MeV & \(u_+\): 9.219544\allowbreak{}45729288\allowbreak{}71\par \(u_-\): 9.219544\allowbreak{}45729288\allowbreak{}71\par Tipo: NONE\par Confianza: ND & COMP-MESON\par E04 / M08\\\addlinespace[3pt]
% VI-CATALOG-ROW masas 395
\textbf{395. B\_\allowbreak{}sJ(6063)\allowbreak{}0 MASS}\par B\_\allowbreak{}sJ(6063)\allowbreak{}0 | Bbar\_\allowbreak{}sJ(6063)\allowbreak{}0\par {\fontsize{8.1}{9.4}\selectfont PDG2026:\allowbreak{}observable:\allowbreak{}M257:\allowbreak{}M257M:\allowbreak{}343153} & \textbf{6063.5+-1.4} MeV & \(u_+\): 1.442220\allowbreak{}51018559\allowbreak{}6\par \(u_-\): 1.442220\allowbreak{}51018559\allowbreak{}6\par Tipo: NONE\par Confianza: ND & COMP-MESON\par E04 / M08\\\addlinespace[3pt]
% VI-CATALOG-ROW masas 396
\textbf{396. B\_\allowbreak{}sJ(6114)\allowbreak{}0 MASS}\par B\_\allowbreak{}sJ(6114)\allowbreak{}0 | Bbar\_\allowbreak{}sJ(6114)\allowbreak{}0\par {\fontsize{8.1}{9.4}\selectfont PDG2026:\allowbreak{}observable:\allowbreak{}M258:\allowbreak{}M258M:\allowbreak{}343157} & \textbf{6114+-6} MeV & \(u_+\): 5.830951\allowbreak{}89484530\allowbreak{}07\par \(u_-\): 5.830951\allowbreak{}89484530\allowbreak{}07\par Tipo: NONE\par Confianza: ND & COMP-MESON\par E04 / M08\\\addlinespace[3pt]
% VI-CATALOG-ROW masas 397
\textbf{397. T\_\allowbreak{}c cbar sbar 1(4000) MASS}\par T\_\allowbreak{}c\_\allowbreak{}cbar\_\allowbreak{}sbar\_\allowbreak{}1(4000)\allowbreak{}0\par {\fontsize{8.1}{9.4}\selectfont PDG2026:\allowbreak{}observable:\allowbreak{}M259:\allowbreak{}M259M:\allowbreak{}347963} & \textbf{3980 -\kern0pt{}- 4010} MeV\par Presentación: 3980-\kern0pt{}-4010 & \(u_+\): NA\par \(u_-\): NA\par Tipo: R\par Confianza: ND & COMP-MESON\par E04 / M08\\\addlinespace[3pt]
% VI-CATALOG-ROW masas 398
\textbf{398. T\_\allowbreak{}c cbar sbar 1(4220)\allowbreak{}+ MASS}\par T\_\allowbreak{}c\_\allowbreak{}cbar\_\allowbreak{}sbar\_\allowbreak{}1(4220)\allowbreak{}+\par {\fontsize{8.1}{9.4}\selectfont PDG2026:\allowbreak{}observable:\allowbreak{}M260:\allowbreak{}M260M:\allowbreak{}348018} & \textbf{4220+50-40} MeV & \(u_+\): 49.24428\allowbreak{}90089805\allowbreak{}23\par \(u_-\): 38.41874\allowbreak{}54245970\allowbreak{}91\par Tipo: NONE\par Confianza: ND & COMP-MESON\par E04 / M08\\\addlinespace[3pt]
% VI-CATALOG-ROW masas 399
\textbf{399. chi\_\allowbreak{}c1(4685) MASS}\par chi\_\allowbreak{}c1(4685)\par {\fontsize{8.1}{9.4}\selectfont PDG2026:\allowbreak{}observable:\allowbreak{}M261:\allowbreak{}M261M:\allowbreak{}346810} & \textbf{4677+13-15} MeV & \(u_+\): 13.28238\allowbreak{}29672178\allowbreak{}6\par \(u_-\): 15.14496\allowbreak{}34763239\allowbreak{}69\par Tipo: NONE\par Confianza: ND & COMP-MESON\par E04 / M08\\\addlinespace[3pt]
% VI-CATALOG-ROW masas 400
\textbf{400. X(4630) MASS}\par X(4630)\par {\fontsize{8.1}{9.4}\selectfont PDG2026:\allowbreak{}observable:\allowbreak{}M262:\allowbreak{}M262M:\allowbreak{}346761} & \textbf{4636+29-100} MeV & \(u_+\): 29.19090\allowbreak{}74314939\allowbreak{}4\par \(u_-\): 102.4064\allowbreak{}25525287\allowbreak{}7\par Tipo: NONE\par Confianza: ND & COMP-MESON\par E04 / M08\\\addlinespace[3pt]
% VI-CATALOG-ROW masas 401
\textbf{401. a\_\allowbreak{}0(1710) MASS}\par a\_\allowbreak{}0(1710)\allowbreak{}0\par {\fontsize{8.1}{9.4}\selectfont PDG2026:\allowbreak{}observable:\allowbreak{}M263:\allowbreak{}M263M:\allowbreak{}334917} & \textbf{1713+-19} MeV & \(u_+\): 18.62757\allowbreak{}46229363\allowbreak{}29\par \(u_-\): 18.62757\allowbreak{}46229363\allowbreak{}29\par Tipo: NONE\par Confianza: ND & COMP-MESON\par E04 / M08\\\addlinespace[3pt]
% VI-CATALOG-ROW masas 402
\textbf{402. f\_\allowbreak{}0(1810) Breit-Wigner MASS}\par f\_\allowbreak{}0(1810)\allowbreak{}0\par {\fontsize{8.1}{9.4}\selectfont PDG2026:\allowbreak{}observable:\allowbreak{}M264:\allowbreak{}M264M:\allowbreak{}334992} & \textbf{1785+-13} MeV & \(u_+\): 13.47812\allowbreak{}30224536\allowbreak{}11\par \(u_-\): 13.30532\allowbreak{}37583047\allowbreak{}5\par Tipo: NONE\par Confianza: ND & COMP-MESON\par E04 / M08\\\addlinespace[3pt]
% VI-CATALOG-ROW masas 403
\textbf{403. T\_\allowbreak{}c\_\allowbreak{}c(3875)\allowbreak{}+ T-Matrix Pole sqrt(s)}\par T\_\allowbreak{}c\_\allowbreak{}c(3875)\allowbreak{}+\par {\fontsize{8.1}{9.4}\selectfont PDG2026:\allowbreak{}observable:\allowbreak{}M265:\allowbreak{}M265TMG:\allowbreak{}347934} & \textbf{(3874.75+-0.10) - i (0.024+0.001-0.007)} MeV\par Presentación: 3874.75+-0.10 & \(u_+\): ND\par \(u_-\): ND\par Tipo: NONE\par Confianza: ND & COMP-MESON\par E04 / M08\\\addlinespace[3pt]
% VI-CATALOG-ROW masas 404
\textbf{404. f\_\allowbreak{}0(2470) MASS}\par f\_\allowbreak{}0(2470)\allowbreak{}0\par {\fontsize{8.1}{9.4}\selectfont PDG2026:\allowbreak{}observable:\allowbreak{}M266:\allowbreak{}M266M:\allowbreak{}335273} & \textbf{2470+6-7} MeV & \(u_+\): 5.656854\allowbreak{}24949238\allowbreak{}06\par \(u_-\): 7.211102\allowbreak{}55092797\allowbreak{}82\par Tipo: NONE\par Confianza: ND & COMP-MESON\par E04 / M08\\\addlinespace[3pt]
% VI-CATALOG-ROW masas 405
\textbf{405. eta\_\allowbreak{}1(1855) MASS}\par eta\_\allowbreak{}1(1855)\allowbreak{}0\par {\fontsize{8.1}{9.4}\selectfont PDG2026:\allowbreak{}observable:\allowbreak{}M267:\allowbreak{}M267M:\allowbreak{}335028} & \textbf{1855+11-9} MeV & \(u_+\): 10.81665\allowbreak{}38263919\allowbreak{}7\par \(u_-\): 9.055385\allowbreak{}13813741\allowbreak{}73\par Tipo: NONE\par Confianza: ND & COMP-MESON\par E04 / M08\\\addlinespace[3pt]
% VI-CATALOG-ROW masas 406
\textbf{406. T\_\allowbreak{}c\_\allowbreak{}c\_\allowbreak{}cbar\_\allowbreak{}cbar(6900)\allowbreak{}0 MASS}\par T\_\allowbreak{}c\_\allowbreak{}c\_\allowbreak{}cbar\_\allowbreak{}cbar(6900)\allowbreak{}+\par {\fontsize{8.1}{9.4}\selectfont PDG2026:\allowbreak{}observable:\allowbreak{}M268:\allowbreak{}M268M:\allowbreak{}348044} & \textbf{6898+-12} MeV & \(u_+\): 11.88793\allowbreak{}88201560\allowbreak{}2\par \(u_-\): 11.51955\allowbreak{}40941920\allowbreak{}7\par Tipo: NONE\par Confianza: ND & COMP-MESON\par E04 / M08\\\addlinespace[3pt]
% VI-CATALOG-ROW masas 407
\textbf{407. T\_\allowbreak{}c\_\allowbreak{}sbar\_\allowbreak{}0*(2900)\allowbreak{}0 MASS}\par T\_\allowbreak{}c\_\allowbreak{}sbar\_\allowbreak{}0\textasciicircum{}*(2900)\allowbreak{}+\par {\fontsize{8.1}{9.4}\selectfont PDG2026:\allowbreak{}observable:\allowbreak{}M269:\allowbreak{}M269M:\allowbreak{}347926} & \textbf{2892+-21} MeV & \(u_+\): 20.51828\allowbreak{}45286831\allowbreak{}89\par \(u_-\): 20.51828\allowbreak{}45286831\allowbreak{}89\par Tipo: NONE\par Confianza: ND & COMP-MESON\par E04 / M08\\\addlinespace[3pt]
% VI-CATALOG-ROW masas 408
\textbf{408. omega(2220) MASS}\par omega(2220)\allowbreak{}0\par {\fontsize{8.1}{9.4}\selectfont PDG2026:\allowbreak{}observable:\allowbreak{}M270:\allowbreak{}M270M:\allowbreak{}335221} & \textbf{2188+-21} MeV & \(u_+\): 20.57966\allowbreak{}47095509\allowbreak{}4\par \(u_-\): 20.57966\allowbreak{}47095509\allowbreak{}4\par Tipo: NONE\par Confianza: ND & COMP-MESON\par E04 / M08\\\addlinespace[3pt]
% VI-CATALOG-ROW masas 409
\textbf{409. T\_\allowbreak{}c c cbar cbar(6600) MASS}\par T\_\allowbreak{}c c cbar cbar(6600)\allowbreak{}+\par {\fontsize{8.1}{9.4}\selectfont PDG2026:\allowbreak{}observable:\allowbreak{}M272:\allowbreak{}M272M:\allowbreak{}348040} & \textbf{6646+28-25} MeV & \(u_+\): 28.34909\allowbreak{}29652609\allowbreak{}9\par \(u_-\): 24.50126\allowbreak{}55281415\allowbreak{}91\par Tipo: NONE\par Confianza: ND & COMP-MESON\par E04 / M08\\\addlinespace[3pt]
% VI-CATALOG-ROW masas 410
\textbf{410. T\_\allowbreak{}c c cbar cbar(7100) MASS}\par T\_\allowbreak{}c c cbar cbar(7100)\allowbreak{}+\par {\fontsize{8.1}{9.4}\selectfont PDG2026:\allowbreak{}observable:\allowbreak{}M273:\allowbreak{}M273M:\allowbreak{}348048} & \textbf{7134+60-29} MeV & \(u_+\): 63.12685\allowbreak{}64083465\allowbreak{}19\par \(u_-\): 29.15475\allowbreak{}94742264\allowbreak{}98\par Tipo: NONE\par Confianza: ND & COMP-MESON\par E04 / M08\\\addlinespace[3pt]
% VI-CATALOG-ROW masas 411
\textbf{411. d-QUARK MASS}\par d | dbar\par {\fontsize{8.1}{9.4}\selectfont PDG2026:\allowbreak{}observable:\allowbreak{}Q001:\allowbreak{}Q001M:\allowbreak{}333589} & \textbf{4.70+-0.07} MeV & \(u_+\): 0.070000\allowbreak{}00000000\allowbreak{}0007\par \(u_-\): 0.070000\allowbreak{}00000000\allowbreak{}0007\par Tipo: NONE\par Confianza: 90 & SM-Q-D1\par E02 / M03\\\addlinespace[3pt]
% VI-CATALOG-ROW masas 412
\textbf{412. u-QUARK MASS}\par u | ubar\par {\fontsize{8.1}{9.4}\selectfont PDG2026:\allowbreak{}observable:\allowbreak{}Q002:\allowbreak{}Q002M:\allowbreak{}333588} & \textbf{2.16+-0.07} MeV & \(u_+\): 0.070000\allowbreak{}00000000\allowbreak{}0007\par \(u_-\): 0.070000\allowbreak{}00000000\allowbreak{}0007\par Tipo: NONE\par Confianza: 90 & SM-Q-U1\par E02 / M03\\\addlinespace[3pt]
% VI-CATALOG-ROW masas 413
\textbf{413. s-QUARK MASS}\par s | sbar\par {\fontsize{8.1}{9.4}\selectfont PDG2026:\allowbreak{}observable:\allowbreak{}Q003:\allowbreak{}Q003M:\allowbreak{}333590} & \textbf{92.9 +-0.7} MeV\par Presentación: 92.9+-0.7 & \(u_+\): 0.699999\allowbreak{}99999999\allowbreak{}996\par \(u_-\): 0.699999\allowbreak{}99999999\allowbreak{}996\par Tipo: NONE\par Confianza: 90 & SM-Q-D2\par E02 / M03\\\addlinespace[3pt]
% VI-CATALOG-ROW masas 414
\textbf{414. c-QUARK MASS}\par c | cbar\par {\fontsize{8.1}{9.4}\selectfont PDG2026:\allowbreak{}observable:\allowbreak{}Q004:\allowbreak{}Q004M:\allowbreak{}333604} & \textbf{1.2729+-0.0045} GeV & \(u_+\): 0.004499\allowbreak{}99999999\allowbreak{}99997\par \(u_-\): 0.004499\allowbreak{}99999999\allowbreak{}99997\par Tipo: NONE\par Confianza: 90 & SM-Q-U2\par E02 / M03\\\addlinespace[3pt]
% VI-CATALOG-ROW masas 415
\textbf{415. b-QUARK MASS}\par b | bbar\par {\fontsize{8.1}{9.4}\selectfont PDG2026:\allowbreak{}observable:\allowbreak{}Q005:\allowbreak{}Q005M:\allowbreak{}333611} & \textbf{4.186+-0.006} GeV & \(u_+\): 0.006000\allowbreak{}00000000\allowbreak{}00001\par \(u_-\): 0.006000\allowbreak{}00000000\allowbreak{}00001\par Tipo: NONE\par Confianza: 90 & SM-Q-D3\par E02 / M03\\\addlinespace[3pt]
% VI-CATALOG-ROW masas 416
\textbf{416. t-Quark Mass (Direct Measurements)}\par t | tbar\par {\fontsize{8.1}{9.4}\selectfont PDG2026:\allowbreak{}observable:\allowbreak{}Q007:\allowbreak{}Q007TP:\allowbreak{}333614} & \textbf{172.60+-0.27} GeV & \(u_+\): 0.271375\allowbreak{}00202757\allowbreak{}259\par \(u_-\): 0.273047\allowbreak{}80674185\allowbreak{}603\par Tipo: NONE\par Confianza: ND & SM-Q-U3\par E02 / M03\\\addlinespace[3pt]
% VI-CATALOG-ROW masas 417
\textbf{417. t-Quark Mass from Cross-Section Measurements}\par t | tbar\par {\fontsize{8.1}{9.4}\selectfont PDG2026:\allowbreak{}observable:\allowbreak{}Q007:\allowbreak{}Q007TP2:\allowbreak{}333615} & \textbf{162.5+2.1-1.5} GeV & \(u_+\): 2.131613\allowbreak{}70878597\allowbreak{}99\par \(u_-\): 1.532433\allowbreak{}43055240\allowbreak{}01\par Tipo: NONE\par Confianza: ND & QUARK\_\allowbreak{}EXTENDED\par E04 / M09\\\addlinespace[3pt]
% VI-CATALOG-ROW masas 418
\textbf{418. t-Quark Pole Mass from Cross-Section Measurements}\par t | tbar\par {\fontsize{8.1}{9.4}\selectfont PDG2026:\allowbreak{}observable:\allowbreak{}Q007:\allowbreak{}Q007TP4:\allowbreak{}333616} & \textbf{172.1+-0.6} GeV & \(u_+\): 0.614086\allowbreak{}44700147\allowbreak{}419\par \(u_-\): 0.589677\allowbreak{}62064252\allowbreak{}371\par Tipo: NONE\par Confianza: ND & QUARK\_\allowbreak{}EXTENDED\par E04 / M09\\\addlinespace[3pt]
% VI-CATALOG-ROW masas 419
\textbf{419. gamma MASS}\par gamma\par {\fontsize{8.1}{9.4}\selectfont PDG2026:\allowbreak{}observable:\allowbreak{}S000:\allowbreak{}S000M:\allowbreak{}332461} & \textbf{<1 E-18} eV\par Presentación: <1E-18 & \(u_+\): NA\par \(u_-\): NA\par Tipo: U\par Confianza: ND & SM-GA-EM\par E05 / M02\\\addlinespace[3pt]
% VI-CATALOG-ROW masas 420
\textbf{420. e MASS (atomic mass units u)}\par e- | e+\par {\fontsize{8.1}{9.4}\selectfont PDG2026:\allowbreak{}observable:\allowbreak{}S003:\allowbreak{}S003AMU:\allowbreak{}332837} & \textbf{(5.485799\allowbreak{}09044+-0.000000\allowbreak{}00010)\allowbreak{}E-4} u\par Presentación: 548.5799\allowbreak{}09044+-0.000000\allowbreak{}010 & \(u_+\): 9.699999\allowbreak{}99999999\allowbreak{}91e-15\par \(u_-\): 9.699999\allowbreak{}99999999\allowbreak{}91e-15\par Tipo: NONE\par Confianza: ND & ELEMEN\allowbreak{}TARY\_\allowbreak{}OR\_\allowbreak{}EXTENDED\par E04 / M05\\\addlinespace[3pt]
% VI-CATALOG-ROW masas 421
\textbf{421. e MASS}\par e- | e+\par {\fontsize{8.1}{9.4}\selectfont PDG2026:\allowbreak{}observable:\allowbreak{}S003:\allowbreak{}S003M:\allowbreak{}332838} & \textbf{0.510998\allowbreak{}95069+-0.000000\allowbreak{}00016} MeV & \(u_+\): 1.599999\allowbreak{}99999999\allowbreak{}99e-10\par \(u_-\): 1.599999\allowbreak{}99999999\allowbreak{}99e-10\par Tipo: NONE\par Confianza: ND & SM-LC-1\par E06 / M01\\\addlinespace[3pt]
% VI-CATALOG-ROW masas 422
\textbf{422. mu MASS (atomic mass units u)}\par mu- | mu+\par {\fontsize{8.1}{9.4}\selectfont PDG2026:\allowbreak{}observable:\allowbreak{}S004:\allowbreak{}S004AMU:\allowbreak{}332848} & \textbf{0.113428\allowbreak{}9257+-0.000000\allowbreak{}0025} u & \(u_+\): 2.500000\allowbreak{}00000000\allowbreak{}01e-09\par \(u_-\): 2.500000\allowbreak{}00000000\allowbreak{}01e-09\par Tipo: NONE\par Confianza: ND & ELEMEN\allowbreak{}TARY\_\allowbreak{}OR\_\allowbreak{}EXTENDED\par E04 / M05\\\addlinespace[3pt]
% VI-CATALOG-ROW masas 423
\textbf{423. mu MASS}\par mu- | mu+\par {\fontsize{8.1}{9.4}\selectfont PDG2026:\allowbreak{}observable:\allowbreak{}S004:\allowbreak{}S004M:\allowbreak{}332849} & \textbf{105.6583\allowbreak{}755+-0.000002\allowbreak{}3} MeV & \(u_+\): 2.3e-06\par \(u_-\): 2.3e-06\par Tipo: NONE\par Confianza: ND & SM-LC-2\par E03 / M04\\\addlinespace[3pt]
% VI-CATALOG-ROW masas 424
\textbf{424. pi+- MASS}\par pi+ | pi-\par {\fontsize{8.1}{9.4}\selectfont PDG2026:\allowbreak{}observable:\allowbreak{}S008:\allowbreak{}S008M:\allowbreak{}333660} & \textbf{139.5703\allowbreak{}9+-0.00018} MeV & \(u_+\): 0.000182\allowbreak{}00716040\allowbreak{}826009\par \(u_-\): 0.000182\allowbreak{}00716040\allowbreak{}826009\par Tipo: NONE\par Confianza: ND & ELEMEN\allowbreak{}TARY\_\allowbreak{}OR\_\allowbreak{}EXTENDED\par E04 / M05\\\addlinespace[3pt]
% VI-CATALOG-ROW masas 425
\textbf{425. pi0 MASS}\par pi0\par {\fontsize{8.1}{9.4}\selectfont PDG2026:\allowbreak{}observable:\allowbreak{}S009:\allowbreak{}S009M:\allowbreak{}333693} & \textbf{134.9768+-0.0005} MeV & \(u_+\): 0.000485\allowbreak{}67928228\allowbreak{}423351\par \(u_-\): 0.000485\allowbreak{}67928228\allowbreak{}423351\par Tipo: NONE\par Confianza: ND & ELEMEN\allowbreak{}TARY\_\allowbreak{}OR\_\allowbreak{}EXTENDED\par E04 / M05\\\addlinespace[3pt]
% VI-CATALOG-ROW masas 426
\textbf{426. K+- MASS}\par K+ | K-\par {\fontsize{8.1}{9.4}\selectfont PDG2026:\allowbreak{}observable:\allowbreak{}S010:\allowbreak{}S010M:\allowbreak{}335281} & \textbf{493.677+-0.015} MeV & \(u_+\): 0.015405\allowbreak{}66522607\allowbreak{}1509\par \(u_-\): 0.015405\allowbreak{}66522607\allowbreak{}1509\par Tipo: NONE\par Confianza: ND & ELEMEN\allowbreak{}TARY\_\allowbreak{}OR\_\allowbreak{}EXTENDED\par E04 / M05\\\addlinespace[3pt]
% VI-CATALOG-ROW masas 427
\textbf{427. K0 MASS}\par K0 | Kbar0\par {\fontsize{8.1}{9.4}\selectfont PDG2026:\allowbreak{}observable:\allowbreak{}S011:\allowbreak{}S011M:\allowbreak{}335501} & \textbf{497.611+-0.013} MeV & \(u_+\): 0.012990\allowbreak{}19975642\allowbreak{}42\par \(u_-\): 0.012990\allowbreak{}19975642\allowbreak{}42\par Tipo: NONE\par Confianza: ND & ELEMEN\allowbreak{}TARY\_\allowbreak{}OR\_\allowbreak{}EXTENDED\par E04 / M05\\\addlinespace[3pt]
% VI-CATALOG-ROW masas 428
\textbf{428. eta MASS}\par eta\par {\fontsize{8.1}{9.4}\selectfont PDG2026:\allowbreak{}observable:\allowbreak{}S014:\allowbreak{}S014M:\allowbreak{}333729} & \textbf{547.862+-0.017} MeV & \(u_+\): 0.017334\allowbreak{}21371191\allowbreak{}4789\par \(u_-\): 0.017334\allowbreak{}21371191\allowbreak{}4789\par Tipo: NONE\par Confianza: ND & ELEMEN\allowbreak{}TARY\_\allowbreak{}OR\_\allowbreak{}EXTENDED\par E04 / M05\\\addlinespace[3pt]
% VI-CATALOG-ROW masas 429
\textbf{429. p MASS (atomic mass units u)}\par p | pbar\par {\fontsize{8.1}{9.4}\selectfont PDG2026:\allowbreak{}observable:\allowbreak{}S016:\allowbreak{}S016AMU:\allowbreak{}348093} & \textbf{1.007276\allowbreak{}4665789+-0.000000\allowbreak{}0000083} u & \(u_+\): 8.299999\allowbreak{}99999999\allowbreak{}98e-12\par \(u_-\): 8.299999\allowbreak{}99999999\allowbreak{}98e-12\par Tipo: NONE\par Confianza: ND & ELEMEN\allowbreak{}TARY\_\allowbreak{}OR\_\allowbreak{}EXTENDED\par E04 / M05\\\addlinespace[3pt]
% VI-CATALOG-ROW masas 430
\textbf{430. p MASS (MeV)}\par p | pbar\par {\fontsize{8.1}{9.4}\selectfont PDG2026:\allowbreak{}observable:\allowbreak{}S016:\allowbreak{}S016M:\allowbreak{}348095} & \textbf{938.2720\allowbreak{}8943+-0.000000\allowbreak{}29} MeV & \(u_+\): 2.899999\allowbreak{}99999999\allowbreak{}98e-07\par \(u_-\): 2.899999\allowbreak{}99999999\allowbreak{}98e-07\par Tipo: NONE\par Confianza: ND & ELEMEN\allowbreak{}TARY\_\allowbreak{}OR\_\allowbreak{}EXTENDED\par E04 / M05\\\addlinespace[3pt]
% VI-CATALOG-ROW masas 431
\textbf{431. n MASS (atomic mass units u)}\par n | nbar\par {\fontsize{8.1}{9.4}\selectfont PDG2026:\allowbreak{}observable:\allowbreak{}S017:\allowbreak{}S017AMU:\allowbreak{}348291} & \textbf{1.008664\allowbreak{}9161+-0.000000\allowbreak{}0004} u & \(u_+\): 4.000000\allowbreak{}00000000\allowbreak{}01e-10\par \(u_-\): 4.000000\allowbreak{}00000000\allowbreak{}01e-10\par Tipo: NONE\par Confianza: ND & ELEMEN\allowbreak{}TARY\_\allowbreak{}OR\_\allowbreak{}EXTENDED\par E04 / M05\\\addlinespace[3pt]
% VI-CATALOG-ROW masas 432
\textbf{432. n MASS (MeV)}\par n | nbar\par {\fontsize{8.1}{9.4}\selectfont PDG2026:\allowbreak{}observable:\allowbreak{}S017:\allowbreak{}S017M:\allowbreak{}348292} & \textbf{939.5654\allowbreak{}219+-0.000000\allowbreak{}5} MeV & \(u_+\): 4.799999\allowbreak{}99999999\allowbreak{}96e-07\par \(u_-\): 4.799999\allowbreak{}99999999\allowbreak{}96e-07\par Tipo: NONE\par Confianza: ND & ELEMEN\allowbreak{}TARY\_\allowbreak{}OR\_\allowbreak{}EXTENDED\par E04 / M05\\\addlinespace[3pt]
% VI-CATALOG-ROW masas 433
\textbf{433. Lambda MASS}\par Lambda0 | Lambdabar0\par {\fontsize{8.1}{9.4}\selectfont PDG2026:\allowbreak{}observable:\allowbreak{}S018:\allowbreak{}S018M:\allowbreak{}349367} & \textbf{1115.683+-0.006} MeV & \(u_+\): 0.006437\allowbreak{}63777862\allowbreak{}19867\par \(u_-\): 0.006437\allowbreak{}63777862\allowbreak{}19867\par Tipo: NONE\par Confianza: ND & ELEMEN\allowbreak{}TARY\_\allowbreak{}OR\_\allowbreak{}EXTENDED\par E04 / M05\\\addlinespace[3pt]
% VI-CATALOG-ROW masas 434
\textbf{434. Sigma+ MASS}\par Sigma+ | Sigmabar-\par {\fontsize{8.1}{9.4}\selectfont PDG2026:\allowbreak{}observable:\allowbreak{}S019:\allowbreak{}S019M:\allowbreak{}349781} & \textbf{1189.37+-0.07} MeV & \(u_+\): 0.069622\allowbreak{}88561220\allowbreak{}7835\par \(u_-\): 0.069622\allowbreak{}88561220\allowbreak{}7835\par Tipo: NONE\par Confianza: ND & ELEMEN\allowbreak{}TARY\_\allowbreak{}OR\_\allowbreak{}EXTENDED\par E04 / M05\\\addlinespace[3pt]
% VI-CATALOG-ROW masas 435
\textbf{435. Sigma- MASS}\par Sigma- | Sigmabar+\par {\fontsize{8.1}{9.4}\selectfont PDG2026:\allowbreak{}observable:\allowbreak{}S020:\allowbreak{}S020M:\allowbreak{}349844} & \textbf{1197.449+-0.029} MeV & \(u_+\): 0.029158\allowbreak{}50307660\allowbreak{}851\par \(u_-\): 0.029158\allowbreak{}50307660\allowbreak{}851\par Tipo: NONE\par Confianza: ND & ELEMEN\allowbreak{}TARY\_\allowbreak{}OR\_\allowbreak{}EXTENDED\par E04 / M05\\\addlinespace[3pt]
% VI-CATALOG-ROW masas 436
\textbf{436. Sigma0 MASS}\par Sigma0 | Sigmabar0\par {\fontsize{8.1}{9.4}\selectfont PDG2026:\allowbreak{}observable:\allowbreak{}S021:\allowbreak{}S021M:\allowbreak{}349828} & \textbf{1192.642+-0.024} MeV & \(u_+\): 0.023595\allowbreak{}69142575\allowbreak{}4781\par \(u_-\): 0.023595\allowbreak{}69142575\allowbreak{}4781\par Tipo: NONE\par Confianza: ND & ELEMEN\allowbreak{}TARY\_\allowbreak{}OR\_\allowbreak{}EXTENDED\par E04 / M05\\\addlinespace[3pt]
% VI-CATALOG-ROW masas 437
\textbf{437. Xi- MASS}\par Xi- | Xibar+\par {\fontsize{8.1}{9.4}\selectfont PDG2026:\allowbreak{}observable:\allowbreak{}S022:\allowbreak{}S022M:\allowbreak{}350344} & \textbf{1321.71+-0.07} MeV & \(u_+\): 0.066535\allowbreak{}27362254\allowbreak{}3412\par \(u_-\): 0.066535\allowbreak{}27362254\allowbreak{}3412\par Tipo: NONE\par Confianza: ND & ELEMEN\allowbreak{}TARY\_\allowbreak{}OR\_\allowbreak{}EXTENDED\par E04 / M05\\\addlinespace[3pt]
% VI-CATALOG-ROW masas 438
\textbf{438. Xibar+ MASS}\par Xi- | Xibar+\par {\fontsize{8.1}{9.4}\selectfont PDG2026:\allowbreak{}observable:\allowbreak{}S022:\allowbreak{}S022M1:\allowbreak{}350346} & \textbf{1321.71+-0.07} MeV & \(u_+\): 0.066535\allowbreak{}27362254\allowbreak{}3412\par \(u_-\): 0.066535\allowbreak{}27362254\allowbreak{}3412\par Tipo: NONE\par Confianza: ND & ELEMEN\allowbreak{}TARY\_\allowbreak{}OR\_\allowbreak{}EXTENDED\par E04 / M05\\\addlinespace[3pt]
% VI-CATALOG-ROW masas 439
\textbf{439. Xi0 MASS}\par Xi0 | Xibar0\par {\fontsize{8.1}{9.4}\selectfont PDG2026:\allowbreak{}observable:\allowbreak{}S023:\allowbreak{}S023M:\allowbreak{}350293} & \textbf{1314.86+-0.20} MeV & \(u_+\): 0.203431\allowbreak{}94946851\allowbreak{}781\par \(u_-\): 0.203431\allowbreak{}94946851\allowbreak{}781\par Tipo: NONE\par Confianza: ND & ELEMEN\allowbreak{}TARY\_\allowbreak{}OR\_\allowbreak{}EXTENDED\par E04 / M05\\\addlinespace[3pt]
% VI-CATALOG-ROW masas 440
\textbf{440. Omega- MASS}\par Omega- | Omegabar+\par {\fontsize{8.1}{9.4}\selectfont PDG2026:\allowbreak{}observable:\allowbreak{}S024:\allowbreak{}S024M:\allowbreak{}350469} & \textbf{1672.45+-0.29} MeV & \(u_+\): 0.293146\allowbreak{}31976651\allowbreak{}251\par \(u_-\): 0.293146\allowbreak{}31976651\allowbreak{}251\par Tipo: NONE\par Confianza: ND & ELEMEN\allowbreak{}TARY\_\allowbreak{}OR\_\allowbreak{}EXTENDED\par E04 / M05\\\addlinespace[3pt]
% VI-CATALOG-ROW masas 441
\textbf{441. Sigma\_\allowbreak{}b+ MASS}\par Sigma\_\allowbreak{}b()\allowbreak{}- | Sigmabar\_\allowbreak{}b()\allowbreak{}+ | Sigma\_\allowbreak{}b()\allowbreak{}0 | Sigmabar\_\allowbreak{}b()\allowbreak{}0 | Sigma\_\allowbreak{}b()\allowbreak{}+ | Sigmabar\_\allowbreak{}b()\allowbreak{}-\par {\fontsize{8.1}{9.4}\selectfont PDG2026:\allowbreak{}observable:\allowbreak{}S026:\allowbreak{}S026M+:\allowbreak{}351717} & \textbf{5810.56+-0.25} MeV & \(u_+\): 0.252740\allowbreak{}62207162\allowbreak{}651\par \(u_-\): 0.252635\allowbreak{}62022485\allowbreak{}619\par Tipo: NONE\par Confianza: ND & ELEMEN\allowbreak{}TARY\_\allowbreak{}OR\_\allowbreak{}EXTENDED\par E04 / M05\\\addlinespace[3pt]
% VI-CATALOG-ROW masas 442
\textbf{442. Sigma\_\allowbreak{}b- MASS}\par Sigma\_\allowbreak{}b()\allowbreak{}- | Sigmabar\_\allowbreak{}b()\allowbreak{}+ | Sigma\_\allowbreak{}b()\allowbreak{}0 | Sigmabar\_\allowbreak{}b()\allowbreak{}0 | Sigma\_\allowbreak{}b()\allowbreak{}+ | Sigmabar\_\allowbreak{}b()\allowbreak{}-\par {\fontsize{8.1}{9.4}\selectfont PDG2026:\allowbreak{}observable:\allowbreak{}S026:\allowbreak{}S026M-:\allowbreak{}351718} & \textbf{5815.64+-0.27} MeV & \(u_+\): 0.274606\allowbreak{}54053367\allowbreak{}178\par \(u_-\): 0.274495\allowbreak{}00386905\allowbreak{}717\par Tipo: NONE\par Confianza: ND & ELEMEN\allowbreak{}TARY\_\allowbreak{}OR\_\allowbreak{}EXTENDED\par E04 / M05\\\addlinespace[3pt]
% VI-CATALOG-ROW masas 443
\textbf{443. D+- MASS}\par D+ | D-\par {\fontsize{8.1}{9.4}\selectfont PDG2026:\allowbreak{}observable:\allowbreak{}S031:\allowbreak{}S031M:\allowbreak{}336044} & \textbf{1869.66+-0.05} MeV & \(u_+\): 0.046806\allowbreak{}02241160\allowbreak{}9027\par \(u_-\): 0.046806\allowbreak{}02241160\allowbreak{}9027\par Tipo: NONE\par Confianza: ND & ELEMEN\allowbreak{}TARY\_\allowbreak{}OR\_\allowbreak{}EXTENDED\par E04 / M05\\\addlinespace[3pt]
% VI-CATALOG-ROW masas 444
\textbf{444. D0 MASS}\par D0 | Dbar0\par {\fontsize{8.1}{9.4}\selectfont PDG2026:\allowbreak{}observable:\allowbreak{}S032:\allowbreak{}S032M:\allowbreak{}336689} & \textbf{1864.84+-0.04} MeV & \(u_+\): 0.044458\allowbreak{}70561980\allowbreak{}1281\par \(u_-\): 0.044458\allowbreak{}70561980\allowbreak{}1281\par Tipo: NONE\par Confianza: ND & ELEMEN\allowbreak{}TARY\_\allowbreak{}OR\_\allowbreak{}EXTENDED\par E04 / M05\\\addlinespace[3pt]
% VI-CATALOG-ROW masas 445
\textbf{445. Lambda\_\allowbreak{}c+ MASS}\par Lambda\_\allowbreak{}c()\allowbreak{}+ | Lambdabar\_\allowbreak{}c()\allowbreak{}-\par {\fontsize{8.1}{9.4}\selectfont PDG2026:\allowbreak{}observable:\allowbreak{}S033:\allowbreak{}S033M:\allowbreak{}350527} & \textbf{2286.46+-0.14} MeV & \(u_+\): 0.139999\allowbreak{}99999999\allowbreak{}89\par \(u_-\): 0.139999\allowbreak{}99999999\allowbreak{}89\par Tipo: NONE\par Confianza: ND & ELEMEN\allowbreak{}TARY\_\allowbreak{}OR\_\allowbreak{}EXTENDED\par E04 / M05\\\addlinespace[3pt]
% VI-CATALOG-ROW masas 446
\textbf{446. D\_\allowbreak{}s+- MASS}\par D\_\allowbreak{}s()\allowbreak{}+ | D\_\allowbreak{}s()\allowbreak{}-\par {\fontsize{8.1}{9.4}\selectfont PDG2026:\allowbreak{}observable:\allowbreak{}S034:\allowbreak{}S034M:\allowbreak{}337875} & \textbf{1968.35+-0.07} MeV & \(u_+\): 0.067861\allowbreak{}20952959\allowbreak{}6483\par \(u_-\): 0.067861\allowbreak{}20952959\allowbreak{}6483\par Tipo: NONE\par Confianza: ND & ELEMEN\allowbreak{}TARY\_\allowbreak{}OR\_\allowbreak{}EXTENDED\par E04 / M05\\\addlinespace[3pt]
% VI-CATALOG-ROW masas 447
\textbf{447. tau MASS}\par tau- | tau+\par {\fontsize{8.1}{9.4}\selectfont PDG2026:\allowbreak{}observable:\allowbreak{}S035:\allowbreak{}S035M:\allowbreak{}332893} & \textbf{1776.93+-0.09} MeV & \(u_+\): 0.084914\allowbreak{}95391323\allowbreak{}0548\par \(u_-\): 0.087580\allowbreak{}19173411\allowbreak{}4063\par Tipo: NONE\par Confianza: ND & SM-LC-3\par E03 / M04\\\addlinespace[3pt]
% VI-CATALOG-ROW masas 448
\textbf{448. mass(Lambda\_\allowbreak{}b0)}\par Lambda\_\allowbreak{}b()\allowbreak{}0 | Lambdabar\_\allowbreak{}b()\allowbreak{}0\par {\fontsize{8.1}{9.4}\selectfont PDG2026:\allowbreak{}observable:\allowbreak{}S040:\allowbreak{}S040M:\allowbreak{}351420} & \textbf{5619.57+-0.16} MeV & \(u_+\): 0.157285\allowbreak{}60301875\allowbreak{}261\par \(u_-\): 0.157285\allowbreak{}60301875\allowbreak{}261\par Tipo: NONE\par Confianza: ND & ELEMEN\allowbreak{}TARY\_\allowbreak{}OR\_\allowbreak{}EXTENDED\par E04 / M05\\\addlinespace[3pt]
% VI-CATALOG-ROW masas 449
\textbf{449. B+- MASS}\par B+ | B-\par {\fontsize{8.1}{9.4}\selectfont PDG2026:\allowbreak{}observable:\allowbreak{}S041:\allowbreak{}S041M:\allowbreak{}338546} & \textbf{5279.41+-0.07} MeV & \(u_+\): 0.071087\allowbreak{}97774343\allowbreak{}7336\par \(u_-\): 0.071087\allowbreak{}97774343\allowbreak{}7336\par Tipo: NONE\par Confianza: ND & ELEMEN\allowbreak{}TARY\_\allowbreak{}OR\_\allowbreak{}EXTENDED\par E04 / M05\\\addlinespace[3pt]
% VI-CATALOG-ROW masas 450
\textbf{450. B0 MASS}\par B0 | Bbar0\par {\fontsize{8.1}{9.4}\selectfont PDG2026:\allowbreak{}observable:\allowbreak{}S042:\allowbreak{}S042M:\allowbreak{}340312} & \textbf{5279.72+-0.08} MeV & \(u_+\): 0.084359\allowbreak{}43092202\allowbreak{}2555\par \(u_-\): 0.084359\allowbreak{}43092202\allowbreak{}2555\par Tipo: NONE\par Confianza: ND & ELEMEN\allowbreak{}TARY\_\allowbreak{}OR\_\allowbreak{}EXTENDED\par E04 / M05\\\addlinespace[3pt]
% VI-CATALOG-ROW masas 451
\textbf{451. W MASS}\par W+ | W-\par {\fontsize{8.1}{9.4}\selectfont PDG2026:\allowbreak{}observable:\allowbreak{}S043:\allowbreak{}S043M:\allowbreak{}332464} & \textbf{80.3625 +-0.0077} GeV\par Presentación: 80.3625+-0.0077 & \(u_+\): 0.007700\allowbreak{}00000000\allowbreak{}00002\par \(u_-\): 0.007700\allowbreak{}00000000\allowbreak{}00002\par Tipo: NONE\par Confianza: ND & SM-GA-W\par E01 / M06\\\addlinespace[3pt]
% VI-CATALOG-ROW masas 452
\textbf{452. Z MASS}\par Z0\par {\fontsize{8.1}{9.4}\selectfont PDG2026:\allowbreak{}observable:\allowbreak{}S044:\allowbreak{}S044M:\allowbreak{}332513} & \textbf{91.1879+-0.0020} GeV & \(u_+\): 0.001966\allowbreak{}89903262\allowbreak{}91501\par \(u_-\): 0.001966\allowbreak{}89903262\allowbreak{}91501\par Tipo: NONE\par Confianza: ND & SM-GA-Z\par E01 / M06\\\addlinespace[3pt]
% VI-CATALOG-ROW masas 453
\textbf{453. Xi\_\allowbreak{}c+ MASS}\par Xi\_\allowbreak{}c()\allowbreak{}+ | Xibar\_\allowbreak{}c()\allowbreak{}-\par {\fontsize{8.1}{9.4}\selectfont PDG2026:\allowbreak{}observable:\allowbreak{}S045:\allowbreak{}S045M:\allowbreak{}351029} & \textbf{2467.79+-0.15} MeV & \(u_+\): 0.147837\allowbreak{}97163412\allowbreak{}299\par \(u_-\): 0.151521\allowbreak{}11693025\allowbreak{}43\par Tipo: NONE\par Confianza: ND & ELEMEN\allowbreak{}TARY\_\allowbreak{}OR\_\allowbreak{}EXTENDED\par E04 / M05\\\addlinespace[3pt]
% VI-CATALOG-ROW masas 454
\textbf{454. Omega\_\allowbreak{}c0 MASS}\par Omega\_\allowbreak{}c()\allowbreak{}0 | Omegabar\_\allowbreak{}c()\allowbreak{}0\par {\fontsize{8.1}{9.4}\selectfont PDG2026:\allowbreak{}observable:\allowbreak{}S047:\allowbreak{}S047M:\allowbreak{}351327} & \textbf{2695.3+-0.4} MeV & \(u_+\): 0.403933\allowbreak{}44242041\allowbreak{}429\par \(u_-\): 0.403933\allowbreak{}44242041\allowbreak{}429\par Tipo: NONE\par Confianza: ND & ELEMEN\allowbreak{}TARY\_\allowbreak{}OR\_\allowbreak{}EXTENDED\par E04 / M05\\\addlinespace[3pt]
% VI-CATALOG-ROW masas 455
\textbf{455. Xi\_\allowbreak{}c0 MASS}\par Xi\_\allowbreak{}c()\allowbreak{}0 | Xibar\_\allowbreak{}c()\allowbreak{}0\par {\fontsize{8.1}{9.4}\selectfont PDG2026:\allowbreak{}observable:\allowbreak{}S048:\allowbreak{}S048M:\allowbreak{}351116} & \textbf{2470.50+-0.25} MeV & \(u_+\): 0.246090\allowbreak{}03107560\allowbreak{}39\par \(u_-\): 0.259362\allowbreak{}93704487\allowbreak{}828\par Tipo: NONE\par Confianza: ND & ELEMEN\allowbreak{}TARY\_\allowbreak{}OR\_\allowbreak{}EXTENDED\par E04 / M05\\\addlinespace[3pt]
% VI-CATALOG-ROW masas 456
\textbf{456. Omega\_\allowbreak{}c(2770)\allowbreak{}0 MASS}\par Omega\_\allowbreak{}c(2770)\allowbreak{}0 | Omegabar\_\allowbreak{}c(2770)\allowbreak{}0\par {\fontsize{8.1}{9.4}\selectfont PDG2026:\allowbreak{}observable:\allowbreak{}S053:\allowbreak{}S053M:\allowbreak{}351372} & \textbf{2766.0+0.9-1.0} MeV & \(u_+\): 0.904428\allowbreak{}79377866\allowbreak{}114\par \(u_-\): 1.024667\allowbreak{}34765194\allowbreak{}1\par Tipo: NONE\par Confianza: ND & ELEMEN\allowbreak{}TARY\_\allowbreak{}OR\_\allowbreak{}EXTENDED\par E04 / M05\\\addlinespace[3pt]
% VI-CATALOG-ROW masas 457
\textbf{457. Xi\_\allowbreak{}c'+ MASS}\par Xi\_\allowbreak{}c\textasciicircum{}'()\allowbreak{}+ | Xibar\_\allowbreak{}c\textasciicircum{}'()\allowbreak{}-\par {\fontsize{8.1}{9.4}\selectfont PDG2026:\allowbreak{}observable:\allowbreak{}S058:\allowbreak{}S058M:\allowbreak{}351195} & \textbf{2578.3+-0.4} MeV & \(u_+\): 0.438637\allowbreak{}52697991\allowbreak{}182\par \(u_-\): 0.438637\allowbreak{}52697991\allowbreak{}182\par Tipo: NONE\par Confianza: ND & ELEMEN\allowbreak{}TARY\_\allowbreak{}OR\_\allowbreak{}EXTENDED\par E04 / M05\\\addlinespace[3pt]
% VI-CATALOG-ROW masas 458
\textbf{458. Xi\_\allowbreak{}c'0 MASS}\par Xi\_\allowbreak{}c\textasciicircum{}'()\allowbreak{}0 | Xibar\_\allowbreak{}c\textasciicircum{}'()\allowbreak{}0\par {\fontsize{8.1}{9.4}\selectfont PDG2026:\allowbreak{}observable:\allowbreak{}S059:\allowbreak{}S059M:\allowbreak{}351200} & \textbf{2578.8+-0.5} MeV & \(u_+\): 0.480228\allowbreak{}71688352\allowbreak{}092\par \(u_-\): 0.486976\allowbreak{}31770634\allowbreak{}259\par Tipo: NONE\par Confianza: ND & ELEMEN\allowbreak{}TARY\_\allowbreak{}OR\_\allowbreak{}EXTENDED\par E04 / M05\\\addlinespace[3pt]
% VI-CATALOG-ROW masas 459
\textbf{459. Sigma\_\allowbreak{}b*+ MASS}\par Sigma\_\allowbreak{}b\textasciicircum{}*()\allowbreak{}- | Sigmabar\_\allowbreak{}b\textasciicircum{}*()\allowbreak{}+ | Sigma\_\allowbreak{}b\textasciicircum{}*()\allowbreak{}0 | Sigmabar\_\allowbreak{}b\textasciicircum{}*()\allowbreak{}0 | Sigma\_\allowbreak{}b\textasciicircum{}*()\allowbreak{}+ | Sigmabar\_\allowbreak{}b\textasciicircum{}*()\allowbreak{}-\par {\fontsize{8.1}{9.4}\selectfont PDG2026:\allowbreak{}observable:\allowbreak{}S062:\allowbreak{}S062M+:\allowbreak{}351724} & \textbf{5830.32+-0.27} MeV & \(u_+\): 0.274729\allowbreak{}15359118\allowbreak{}399\par \(u_-\): 0.275017\allowbreak{}43053510\allowbreak{}241\par Tipo: NONE\par Confianza: ND & ELEMEN\allowbreak{}TARY\_\allowbreak{}OR\_\allowbreak{}EXTENDED\par E04 / M05\\\addlinespace[3pt]
% VI-CATALOG-ROW masas 460
\textbf{460. Sigma\_\allowbreak{}b*- MASS}\par Sigma\_\allowbreak{}b\textasciicircum{}*()\allowbreak{}- | Sigmabar\_\allowbreak{}b\textasciicircum{}*()\allowbreak{}+ | Sigma\_\allowbreak{}b\textasciicircum{}*()\allowbreak{}0 | Sigmabar\_\allowbreak{}b\textasciicircum{}*()\allowbreak{}0 | Sigma\_\allowbreak{}b\textasciicircum{}*()\allowbreak{}+ | Sigmabar\_\allowbreak{}b\textasciicircum{}*()\allowbreak{}-\par {\fontsize{8.1}{9.4}\selectfont PDG2026:\allowbreak{}observable:\allowbreak{}S062:\allowbreak{}S062M-:\allowbreak{}351725} & \textbf{5834.74+-0.30} MeV & \(u_+\): 0.298160\allowbreak{}79945753\allowbreak{}06\par \(u_-\): 0.298558\allowbreak{}05603668\allowbreak{}217\par Tipo: NONE\par Confianza: ND & ELEMEN\allowbreak{}TARY\_\allowbreak{}OR\_\allowbreak{}EXTENDED\par E04 / M05\\\addlinespace[3pt]
% VI-CATALOG-ROW masas 461
\textbf{461. Omega\_\allowbreak{}b- MASS}\par Omega\_\allowbreak{}b()\allowbreak{}- | Omegabar\_\allowbreak{}b()\allowbreak{}+\par {\fontsize{8.1}{9.4}\selectfont PDG2026:\allowbreak{}observable:\allowbreak{}S063:\allowbreak{}S063M:\allowbreak{}351884} & \textbf{6045.8+-0.8} MeV & \(u_+\): 0.765410\allowbreak{}57086097\allowbreak{}497\par \(u_-\): 0.765410\allowbreak{}57086097\allowbreak{}497\par Tipo: NONE\par Confianza: ND & ELEMEN\allowbreak{}TARY\_\allowbreak{}OR\_\allowbreak{}EXTENDED\par E04 / M05\\\addlinespace[3pt]
% VI-CATALOG-ROW masas 462
\textbf{462. Xi\_\allowbreak{}cc+ MASS}\par Xi\_\allowbreak{}cc()\allowbreak{}+ | Xibar\_\allowbreak{}cc()\allowbreak{}-\par {\fontsize{8.1}{9.4}\selectfont PDG2026:\allowbreak{}observable:\allowbreak{}S065:\allowbreak{}S065M:\allowbreak{}351404} & \textbf{3518.9+-0.9} MeV & \(u_+\): 0.948683\allowbreak{}29805051\allowbreak{}377\par \(u_-\): 0.948683\allowbreak{}29805051\allowbreak{}377\par Tipo: NONE\par Confianza: ND & ELEMEN\allowbreak{}TARY\_\allowbreak{}OR\_\allowbreak{}EXTENDED\par E04 / M05\\\addlinespace[3pt]
% VI-CATALOG-ROW masas 463
\textbf{463. Xi\_\allowbreak{}cc++ MASS}\par Xi\_\allowbreak{}cc()\allowbreak{}++ | Xi\_\allowbreak{}cc()\allowbreak{}-\kern0pt{}-\par {\fontsize{8.1}{9.4}\selectfont PDG2026:\allowbreak{}observable:\allowbreak{}S068:\allowbreak{}S068M:\allowbreak{}351409} & \textbf{3621.6+-0.4} MeV & \(u_+\): 0.378021\allowbreak{}16342871\allowbreak{}6\par \(u_-\): 0.378021\allowbreak{}16342871\allowbreak{}6\par Tipo: NONE\par Confianza: ND & ELEMEN\allowbreak{}TARY\_\allowbreak{}OR\_\allowbreak{}EXTENDED\par E04 / M05\\\addlinespace[3pt]
% VI-CATALOG-ROW masas 464
\textbf{464. Xi\_\allowbreak{}b- MASS}\par Xi\_\allowbreak{}b()\allowbreak{}- | Xibar\_\allowbreak{}b()\allowbreak{}+\par {\fontsize{8.1}{9.4}\selectfont PDG2026:\allowbreak{}observable:\allowbreak{}S069:\allowbreak{}S069M-:\allowbreak{}351742} & \textbf{5797.0+-0.4} MeV & \(u_+\): 0.436742\allowbreak{}65363571\allowbreak{}568\par \(u_-\): 0.436742\allowbreak{}65363571\allowbreak{}568\par Tipo: NONE\par Confianza: ND & ELEMEN\allowbreak{}TARY\_\allowbreak{}OR\_\allowbreak{}EXTENDED\par E04 / M05\\\addlinespace[3pt]
% VI-CATALOG-ROW masas 465
\textbf{465. Xi\_\allowbreak{}b0 MASS}\par Xi\_\allowbreak{}b()\allowbreak{}0 | Xibar\_\allowbreak{}b()\allowbreak{}0\par {\fontsize{8.1}{9.4}\selectfont PDG2026:\allowbreak{}observable:\allowbreak{}S070:\allowbreak{}S070M0:\allowbreak{}351793} & \textbf{5791.7+-0.4} MeV & \(u_+\): 0.413429\allowbreak{}61949753\allowbreak{}559\par \(u_-\): 0.413429\allowbreak{}61949753\allowbreak{}559\par Tipo: NONE\par Confianza: ND & ELEMEN\allowbreak{}TARY\_\allowbreak{}OR\_\allowbreak{}EXTENDED\par E04 / M05\\\addlinespace[3pt]
% VI-CATALOG-ROW masas 466
\textbf{466. D\_\allowbreak{}s*+- MASS}\par D\_\allowbreak{}s\textasciicircum{}*()\allowbreak{}+ | D\_\allowbreak{}s\textasciicircum{}*()\allowbreak{}-\par {\fontsize{8.1}{9.4}\selectfont PDG2026:\allowbreak{}observable:\allowbreak{}S074:\allowbreak{}S074M:\allowbreak{}338415} & \textbf{2112.2+-0.4} MeV & \(u_+\): 0.395186\allowbreak{}70911609\allowbreak{}538\par \(u_-\): 0.395186\allowbreak{}70911609\allowbreak{}538\par Tipo: NONE\par Confianza: ND & ELEMEN\allowbreak{}TARY\_\allowbreak{}OR\_\allowbreak{}EXTENDED\par E04 / M05\\\addlinespace[3pt]
% VI-CATALOG-ROW masas 467
\textbf{467. B* MASS}\par B\textasciicircum{}*0 | Bbar\textasciicircum{}*0 | B\textasciicircum{}*+ | B\textasciicircum{}*-\par {\fontsize{8.1}{9.4}\selectfont PDG2026:\allowbreak{}observable:\allowbreak{}S085:\allowbreak{}S085M:\allowbreak{}342531} & \textbf{5324.75+-0.20} MeV & \(u_+\): 0.203385\allowbreak{}63693989\allowbreak{}911\par \(u_-\): 0.203385\allowbreak{}63693989\allowbreak{}911\par Tipo: NONE\par Confianza: ND & ELEMEN\allowbreak{}TARY\_\allowbreak{}OR\_\allowbreak{}EXTENDED\par E04 / M05\\\addlinespace[3pt]
% VI-CATALOG-ROW masas 468
\textbf{468. B\_\allowbreak{}s0 MASS}\par B\_\allowbreak{}s()\allowbreak{}0 | Bbar\_\allowbreak{}s()\allowbreak{}0\par {\fontsize{8.1}{9.4}\selectfont PDG2026:\allowbreak{}observable:\allowbreak{}S086:\allowbreak{}S086M:\allowbreak{}342595} & \textbf{5366.93+-0.10} MeV & \(u_+\): 0.099351\allowbreak{}67389588\allowbreak{}262\par \(u_-\): 0.099351\allowbreak{}67389588\allowbreak{}262\par Tipo: NONE\par Confianza: ND & ELEMEN\allowbreak{}TARY\_\allowbreak{}OR\_\allowbreak{}EXTENDED\par E04 / M05\\\addlinespace[3pt]
% VI-CATALOG-ROW masas 469
\textbf{469. B\_\allowbreak{}s* MASS}\par B\_\allowbreak{}s\textasciicircum{}*0 | Bbar\_\allowbreak{}s\textasciicircum{}*0\par {\fontsize{8.1}{9.4}\selectfont PDG2026:\allowbreak{}observable:\allowbreak{}S087:\allowbreak{}S087M:\allowbreak{}343123} & \textbf{5415.4+-1.4} MeV & \(u_+\): 1.435389\allowbreak{}23919178\allowbreak{}2\par \(u_-\): 1.435389\allowbreak{}23919178\allowbreak{}2\par Tipo: NONE\par Confianza: ND & ELEMEN\allowbreak{}TARY\_\allowbreak{}OR\_\allowbreak{}EXTENDED\par E04 / M05\\\addlinespace[3pt]
% VI-CATALOG-ROW masas 470
\textbf{470. B\_\allowbreak{}c+ MASS}\par B\_\allowbreak{}c()\allowbreak{}+ | B\_\allowbreak{}c()\allowbreak{}-\par {\fontsize{8.1}{9.4}\selectfont PDG2026:\allowbreak{}observable:\allowbreak{}S091:\allowbreak{}S091M:\allowbreak{}343161} & \textbf{6274.47+-0.32} MeV & \(u_+\): 0.319061\allowbreak{}12267087\allowbreak{}638\par \(u_-\): 0.319061\allowbreak{}12267087\allowbreak{}638\par Tipo: NONE\par Confianza: ND & ELEMEN\allowbreak{}TARY\_\allowbreak{}OR\_\allowbreak{}EXTENDED\par E04 / M05\\\addlinespace[3pt]
% VI-CATALOG-ROW masas 471
\textbf{471. H MASS}\par H\par {\fontsize{8.1}{9.4}\selectfont PDG2026:\allowbreak{}observable:\allowbreak{}S126:\allowbreak{}S126M:\allowbreak{}332733} & \textbf{125.13+-0.11} GeV & \(u_+\): 0.111721\allowbreak{}44776548\allowbreak{}799\par \(u_-\): 0.111721\allowbreak{}44776548\allowbreak{}799\par Tipo: NONE\par Confianza: ND & SM-H-1\par E01 / M06\\\addlinespace[3pt]
\end{longtable}\endgroup
\clearpage\subsection{Anchuras: 384 registros}
\begingroup\fontsize{9.2}{11.1}\selectfont\setlength{\tabcolsep}{5pt}\setlength{\extrarowheight}{3pt}\renewcommand{\arraystretch}{1.08}
\begin{longtable}{@{}>{\raggedright\arraybackslash}p{.445\linewidth}>{\raggedright\arraybackslash}p{.205\linewidth}>{\raggedright\arraybackslash}p{.18\linewidth}>{\raggedright\arraybackslash}p{\dimexpr.17\linewidth-6\tabcolsep\relax}@{}}
\toprule
Identidad y especies & Valor y unidad & Incertidumbre & Sector y estado\\\midrule
\endfirsthead
\toprule
Identidad y especies & Valor y unidad & Incertidumbre & Sector y estado\\\midrule
\endhead
\bottomrule
\endfoot
% VI-CATALOG-ROW anchuras 1
\textbf{1. N(1895) BREIT-WIGNER WIDTH}\par N(1895)\allowbreak{}0 | Nbar(1895)\allowbreak{}0 | N(1895)\allowbreak{}+ | Nbar(1895)\allowbreak{}-\par {\fontsize{8.1}{9.4}\selectfont PDG2026:\allowbreak{}observable:\allowbreak{}B004:\allowbreak{}B004W:\allowbreak{}348706} & \textbf{80 to 120 to 200} MeV\par Presentación: 80 to 200 (\textasciitilde{}120) & \(u_+\): NA\par \(u_-\): NA\par Tipo: R\par Confianza: ND & COMP-BARYON\par E04 / M07\\\addlinespace[3pt]
% VI-CATALOG-ROW anchuras 2
\textbf{2. N(2060) BREIT-WIGNER WIDTH}\par N(2060)\allowbreak{}0 | Nbar(2060)\allowbreak{}0 | N(2060)\allowbreak{}+ | Nbar(2060)\allowbreak{}-\par {\fontsize{8.1}{9.4}\selectfont PDG2026:\allowbreak{}observable:\allowbreak{}B005:\allowbreak{}B005W:\allowbreak{}348837} & \textbf{300 to 400 to 450} MeV\par Presentación: 300 to 450 (\textasciitilde{}400) & \(u_+\): NA\par \(u_-\): NA\par Tipo: R\par Confianza: ND & COMP-BARYON\par E04 / M07\\\addlinespace[3pt]
% VI-CATALOG-ROW anchuras 3
\textbf{3. Delta(1700) BREIT-WIGNER WIDTH}\par Delta(1700)\allowbreak{}- | Deltabar(1700)\allowbreak{}+ | Delta(1700)\allowbreak{}0 | Deltabar(1700)\allowbreak{}0 | Delta(1700)\allowbreak{}+ | Deltabar(1700)\allowbreak{}- | Delta(1700)\allowbreak{}++ | Deltabar(1700)\allowbreak{}-\kern0pt{}-\par {\fontsize{8.1}{9.4}\selectfont PDG2026:\allowbreak{}observable:\allowbreak{}B010:\allowbreak{}B010W:\allowbreak{}349098} & \textbf{220 to 300 to 380} MeV\par Presentación: 220 to 380 (\textasciitilde{}300) & \(u_+\): NA\par \(u_-\): NA\par Tipo: R\par Confianza: ND & COMP-BARYON\par E04 / M07\\\addlinespace[3pt]
% VI-CATALOG-ROW anchuras 4
\textbf{4. Delta(1905) BREIT-WIGNER WIDTH}\par Delta(1905)\allowbreak{}- | Deltabar(1905)\allowbreak{}+ | Delta(1905)\allowbreak{}0 | Deltabar(1905)\allowbreak{}0 | Delta(1905)\allowbreak{}+ | Deltabar(1905)\allowbreak{}- | Delta(1905)\allowbreak{}++ | Deltabar(1905)\allowbreak{}-\kern0pt{}-\par {\fontsize{8.1}{9.4}\selectfont PDG2026:\allowbreak{}observable:\allowbreak{}B011:\allowbreak{}B011W:\allowbreak{}349152} & \textbf{270 to 330 to 400} MeV\par Presentación: 270 to 400 (\textasciitilde{}330) & \(u_+\): NA\par \(u_-\): NA\par Tipo: R\par Confianza: ND & COMP-BARYON\par E04 / M07\\\addlinespace[3pt]
% VI-CATALOG-ROW anchuras 5
\textbf{5. Delta(1910) BREIT-WIGNER WIDTH}\par Delta(1910)\allowbreak{}- | Deltabar(1910)\allowbreak{}+ | Delta(1910)\allowbreak{}0 | Deltabar(1910)\allowbreak{}0 | Delta(1910)\allowbreak{}+ | Deltabar(1910)\allowbreak{}- | Delta(1910)\allowbreak{}++ | Deltabar(1910)\allowbreak{}-\kern0pt{}-\par {\fontsize{8.1}{9.4}\selectfont PDG2026:\allowbreak{}observable:\allowbreak{}B012:\allowbreak{}B012W:\allowbreak{}349183} & \textbf{200 to 300 to 400} MeV\par Presentación: 200 to 400 (\textasciitilde{}300) & \(u_+\): NA\par \(u_-\): NA\par Tipo: R\par Confianza: ND & COMP-BARYON\par E04 / M07\\\addlinespace[3pt]
% VI-CATALOG-ROW anchuras 6
\textbf{6. Delta(1930) BREIT-WIGNER WIDTH}\par Delta(1930)\allowbreak{}- | Deltabar(1930)\allowbreak{}+ | Delta(1930)\allowbreak{}0 | Deltabar(1930)\allowbreak{}0 | Delta(1930)\allowbreak{}+ | Deltabar(1930)\allowbreak{}- | Delta(1930)\allowbreak{}++ | Deltabar(1930)\allowbreak{}-\kern0pt{}-\par {\fontsize{8.1}{9.4}\selectfont PDG2026:\allowbreak{}observable:\allowbreak{}B013:\allowbreak{}B013W:\allowbreak{}349250} & \textbf{200 to 300 to 400} MeV\par Presentación: 200 to 400 (\textasciitilde{}300) & \(u_+\): NA\par \(u_-\): NA\par Tipo: R\par Confianza: ND & COMP-BARYON\par E04 / M07\\\addlinespace[3pt]
% VI-CATALOG-ROW anchuras 7
\textbf{7. N(1710) BREIT-WIGNER WIDTH}\par N(1710)\allowbreak{}0 | Nbar(1710)\allowbreak{}0 | N(1710)\allowbreak{}+ | Nbar(1710)\allowbreak{}-\par {\fontsize{8.1}{9.4}\selectfont PDG2026:\allowbreak{}observable:\allowbreak{}B014:\allowbreak{}B014W:\allowbreak{}348536} & \textbf{80 to 140 to 200} MeV\par Presentación: 80 to 200 (\textasciitilde{}140) & \(u_+\): NA\par \(u_-\): NA\par Tipo: R\par Confianza: ND & COMP-BARYON\par E04 / M07\\\addlinespace[3pt]
% VI-CATALOG-ROW anchuras 8
\textbf{8. N(1720) BREIT-WIGNER WIDTH}\par N(1720)\allowbreak{}0 | Nbar(1720)\allowbreak{}0 | N(1720)\allowbreak{}+ | Nbar(1720)\allowbreak{}-\par {\fontsize{8.1}{9.4}\selectfont PDG2026:\allowbreak{}observable:\allowbreak{}B015:\allowbreak{}B015W:\allowbreak{}348570} & \textbf{150 to 250 to 400} MeV\par Presentación: 150 to 400 (\textasciitilde{}250) & \(u_+\): NA\par \(u_-\): NA\par Tipo: R\par Confianza: ND & COMP-BARYON\par E04 / M07\\\addlinespace[3pt]
% VI-CATALOG-ROW anchuras 9
\textbf{9. N(1875) BREIT-WIGNER WIDTH}\par N(1875)\allowbreak{}0 | Nbar(1875)\allowbreak{}0 | N(1875)\allowbreak{}+ | Nbar(1875)\allowbreak{}-\par {\fontsize{8.1}{9.4}\selectfont PDG2026:\allowbreak{}observable:\allowbreak{}B016:\allowbreak{}B016W:\allowbreak{}348634} & \textbf{120 to 200 to 250} MeV\par Presentación: 120 to 250 (\textasciitilde{}200) & \(u_+\): NA\par \(u_-\): NA\par Tipo: R\par Confianza: ND & COMP-BARYON\par E04 / M07\\\addlinespace[3pt]
% VI-CATALOG-ROW anchuras 10
\textbf{10. N(1990) BREIT-WIGNER WIDTH}\par N(1990)\allowbreak{}0 | Nbar(1990)\allowbreak{}0 | N(1990)\allowbreak{}+ | Nbar(1990)\allowbreak{}-\par {\fontsize{8.1}{9.4}\selectfont PDG2026:\allowbreak{}observable:\allowbreak{}B017:\allowbreak{}B017W:\allowbreak{}348779} & \textbf{200 to 300 to 400} MeV\par Presentación: 200 to 400 (\textasciitilde{}300) & \(u_+\): NA\par \(u_-\): NA\par Tipo: R\par Confianza: ND & COMP-BARYON\par E04 / M07\\\addlinespace[3pt]
% VI-CATALOG-ROW anchuras 11
\textbf{11. N(1700) BREIT-WIGNER WIDTH}\par N(1700)\allowbreak{}0 | Nbar(1700)\allowbreak{}0 | N(1700)\allowbreak{}+ | Nbar(1700)\allowbreak{}-\par {\fontsize{8.1}{9.4}\selectfont PDG2026:\allowbreak{}observable:\allowbreak{}B018:\allowbreak{}B018W:\allowbreak{}348500} & \textbf{100 to 200 to 300} MeV\par Presentación: 100 to 300 (\textasciitilde{}200) & \(u_+\): NA\par \(u_-\): NA\par Tipo: R\par Confianza: ND & COMP-BARYON\par E04 / M07\\\addlinespace[3pt]
% VI-CATALOG-ROW anchuras 12
\textbf{12. Delta(1600) BREIT-WIGNER WIDTH}\par Delta(1600)\allowbreak{}- | Deltabar(1600)\allowbreak{}+ | Delta(1600)\allowbreak{}0 | Deltabar(1600)\allowbreak{}0 | Delta(1600)\allowbreak{}+ | Deltabar(1600)\allowbreak{}- | Delta(1600)\allowbreak{}++ | Deltabar(1600)\allowbreak{}-\kern0pt{}-\par {\fontsize{8.1}{9.4}\selectfont PDG2026:\allowbreak{}observable:\allowbreak{}B019:\allowbreak{}B019W:\allowbreak{}349051} & \textbf{200 to 250 to 300} MeV\par Presentación: 200 to 300 (\textasciitilde{}250) & \(u_+\): NA\par \(u_-\): NA\par Tipo: R\par Confianza: ND & COMP-BARYON\par E04 / M07\\\addlinespace[3pt]
% VI-CATALOG-ROW anchuras 13
\textbf{13. Xi(1620) WIDTH}\par Xi(1620)\allowbreak{}- | Xibar(1620)\allowbreak{}+ | Xi(1620)\allowbreak{}0 | Xibar(1620)\allowbreak{}0\par {\fontsize{8.1}{9.4}\selectfont PDG2026:\allowbreak{}observable:\allowbreak{}B021:\allowbreak{}B021W:\allowbreak{}350411} & \textbf{32+8-9} MeV & \(u_+\): 8.149907\allowbreak{}01245417\allowbreak{}88\par \(u_-\): 9.411978\allowbreak{}92698348\allowbreak{}86\par Tipo: NONE\par Confianza: ND & COMP-BARYON\par E04 / M07\\\addlinespace[3pt]
% VI-CATALOG-ROW anchuras 14
\textbf{14. N(2120) BREIT-WIGNER WIDTH}\par N(2120)\allowbreak{}0 | Nbar(2120)\allowbreak{}0 | N(2120)\allowbreak{}+ | Nbar(2120)\allowbreak{}-\par {\fontsize{8.1}{9.4}\selectfont PDG2026:\allowbreak{}observable:\allowbreak{}B024:\allowbreak{}B024W:\allowbreak{}348907} & \textbf{260 to 300 to 360} MeV\par Presentación: 260 to 360 (\textasciitilde{}300) & \(u_+\): NA\par \(u_-\): NA\par Tipo: R\par Confianza: ND & COMP-BARYON\par E04 / M07\\\addlinespace[3pt]
% VI-CATALOG-ROW anchuras 15
\textbf{15. Delta(1900) BREIT-WIGNER WIDTH}\par Delta(1900)\allowbreak{}- | Deltabar(1900)\allowbreak{}+ | Delta(1900)\allowbreak{}0 | Deltabar(1900)\allowbreak{}0 | Delta(1900)\allowbreak{}+ | Deltabar(1900)\allowbreak{}- | Delta(1900)\allowbreak{}++ | Deltabar(1900)\allowbreak{}-\kern0pt{}-\par {\fontsize{8.1}{9.4}\selectfont PDG2026:\allowbreak{}observable:\allowbreak{}B030:\allowbreak{}B030W:\allowbreak{}349127} & \textbf{180 to 250 to 320} MeV\par Presentación: 180 to 320 (\textasciitilde{}250) & \(u_+\): NA\par \(u_-\): NA\par Tipo: R\par Confianza: ND & COMP-BARYON\par E04 / M07\\\addlinespace[3pt]
% VI-CATALOG-ROW anchuras 16
\textbf{16. Sigma(1620) WIDTH}\par Sigma(1620)\allowbreak{}- | Sigmabar(1620)\allowbreak{}+ | Sigma(1620)\allowbreak{}0 | Sigmabar(1620)\allowbreak{}0 | Sigma(1620)\allowbreak{}+ | Sigmabar(1620)\allowbreak{}-\par {\fontsize{8.1}{9.4}\selectfont PDG2026:\allowbreak{}observable:\allowbreak{}B032:\allowbreak{}B032W:\allowbreak{}349903} & \textbf{40 to 70 to 100} MeV\par Presentación: 40 to 100 (\textasciitilde{}70) & \(u_+\): NA\par \(u_-\): NA\par Tipo: R\par Confianza: ND & COMP-BARYON\par E04 / M07\\\addlinespace[3pt]
% VI-CATALOG-ROW anchuras 17
\textbf{17. MIXED CHARGES}\par Delta(1232)\allowbreak{}- | Deltabar(1232)\allowbreak{}+ | Delta(1232)\allowbreak{}0 | Deltabar(1232)\allowbreak{}0 | Delta(1232)\allowbreak{}+ | Deltabar(1232)\allowbreak{}- | Delta(1232)\allowbreak{}++ | Deltabar(1232)\allowbreak{}-\kern0pt{}-\par {\fontsize{8.1}{9.4}\selectfont PDG2026:\allowbreak{}observable:\allowbreak{}B033:\allowbreak{}B033W:\allowbreak{}349035} & \textbf{114 TO 117 TO 120} MeV\par Presentación: 114 to 120 (\textasciitilde{}117) & \(u_+\): NA\par \(u_-\): NA\par Tipo: R\par Confianza: ND & COMP-BARYON\par E04 / M07\\\addlinespace[3pt]
% VI-CATALOG-ROW anchuras 18
\textbf{18. Sigma(2070) WIDTH}\par Sigma(2070)\allowbreak{}- | Sigmabar(2070)\allowbreak{}+ | Sigma(2070)\allowbreak{}0 | Sigmabar(2070)\allowbreak{}0 | Sigma(2070)\allowbreak{}+ | Sigmabar(2070)\allowbreak{}-\par {\fontsize{8.1}{9.4}\selectfont PDG2026:\allowbreak{}observable:\allowbreak{}B034:\allowbreak{}B034W:\allowbreak{}350187} & \textbf{100 to 200 to 300} MeV\par Presentación: 100 to 300 (\textasciitilde{}200) & \(u_+\): NA\par \(u_-\): NA\par Tipo: R\par Confianza: ND & COMP-BARYON\par E04 / M07\\\addlinespace[3pt]
% VI-CATALOG-ROW anchuras 19
\textbf{19. Lambda(2110) WIDTH}\par Lambda(2110)\allowbreak{}0 | Lambdabar(2110)\allowbreak{}0\par {\fontsize{8.1}{9.4}\selectfont PDG2026:\allowbreak{}observable:\allowbreak{}B035:\allowbreak{}B035W:\allowbreak{}349764} & \textbf{200 to 250 to 300} MeV\par Presentación: 200 to 300 (\textasciitilde{}250) & \(u_+\): NA\par \(u_-\): NA\par Tipo: R\par Confianza: ND & COMP-BARYON\par E04 / M07\\\addlinespace[3pt]
% VI-CATALOG-ROW anchuras 20
\textbf{20. Lambda(1800) WIDTH}\par Lambda(1800)\allowbreak{}0 | Lambdabar(1800)\allowbreak{}0\par {\fontsize{8.1}{9.4}\selectfont PDG2026:\allowbreak{}observable:\allowbreak{}B036:\allowbreak{}B036W:\allowbreak{}349556} & \textbf{150 to 200 to 250} MeV\par Presentación: 150 to 250 (\textasciitilde{}200) & \(u_+\): NA\par \(u_-\): NA\par Tipo: R\par Confianza: ND & COMP-BARYON\par E04 / M07\\\addlinespace[3pt]
% VI-CATALOG-ROW anchuras 21
\textbf{21. PRODUCTION EXPERIMENTS}\par Lambda(1405)\allowbreak{}0 | Lambdabar(1405)\allowbreak{}0\par {\fontsize{8.1}{9.4}\selectfont PDG2026:\allowbreak{}observable:\allowbreak{}B037:\allowbreak{}B037W1:\allowbreak{}349428} & \textbf{50.5+-2.0} MeV & \(u_+\): 1.961161\allowbreak{}35138184\par \(u_-\): 1.961161\allowbreak{}35138184\par Tipo: NONE\par Confianza: ND & COMP-BARYON\par E04 / M07\\\addlinespace[3pt]
% VI-CATALOG-ROW anchuras 22
\textbf{22. Lambda(1520) WIDTH}\par Lambda(1520)\allowbreak{}0 | Lambdabar(1520)\allowbreak{}0\par {\fontsize{8.1}{9.4}\selectfont PDG2026:\allowbreak{}observable:\allowbreak{}B038:\allowbreak{}B038W:\allowbreak{}349439} & \textbf{15 to 16 to 17} MeV\par Presentación: 15 to 17 (\textasciitilde{}16) & \(u_+\): NA\par \(u_-\): NA\par Tipo: R\par Confianza: ND & COMP-BARYON\par E04 / M07\\\addlinespace[3pt]
% VI-CATALOG-ROW anchuras 23
\textbf{23. Lambda(1820) WIDTH}\par Lambda(1820)\allowbreak{}0 | Lambdabar(1820)\allowbreak{}0\par {\fontsize{8.1}{9.4}\selectfont PDG2026:\allowbreak{}observable:\allowbreak{}B039:\allowbreak{}B039W:\allowbreak{}349595} & \textbf{70 TO 80 TO 90} MeV\par Presentación: 70 to 90 (\textasciitilde{}80) & \(u_+\): NA\par \(u_-\): NA\par Tipo: R\par Confianza: ND & COMP-BARYON\par E04 / M07\\\addlinespace[3pt]
% VI-CATALOG-ROW anchuras 24
\textbf{24. Lambda(1670) WIDTH}\par Lambda(1670)\allowbreak{}0 | Lambdabar(1670)\allowbreak{}0\par {\fontsize{8.1}{9.4}\selectfont PDG2026:\allowbreak{}observable:\allowbreak{}B040:\allowbreak{}B040W:\allowbreak{}349495} & \textbf{25 to 30 to 35} MeV\par Presentación: 25 to 35 (\textasciitilde{}30) & \(u_+\): NA\par \(u_-\): NA\par Tipo: R\par Confianza: ND & COMP-BARYON\par E04 / M07\\\addlinespace[3pt]
% VI-CATALOG-ROW anchuras 25
\textbf{25. Lambda(2100) WIDTH}\par Lambda(2100)\allowbreak{}0 | Lambdabar(2100)\allowbreak{}0\par {\fontsize{8.1}{9.4}\selectfont PDG2026:\allowbreak{}observable:\allowbreak{}B041:\allowbreak{}B041W:\allowbreak{}349742} & \textbf{100 TO 200 TO 250} MeV\par Presentación: 100 to 250 (\textasciitilde{}200) & \(u_+\): NA\par \(u_-\): NA\par Tipo: R\par Confianza: ND & COMP-BARYON\par E04 / M07\\\addlinespace[3pt]
% VI-CATALOG-ROW anchuras 26
\textbf{26. Lambda(2350) WIDTH}\par Lambda(2350)\allowbreak{}0 | Lambdabar(2350)\allowbreak{}0\par {\fontsize{8.1}{9.4}\selectfont PDG2026:\allowbreak{}observable:\allowbreak{}B042:\allowbreak{}B042W:\allowbreak{}349776} & \textbf{100 TO 150 TO 250} MeV\par Presentación: 100 to 250 (\textasciitilde{}150) & \(u_+\): NA\par \(u_-\): NA\par Tipo: R\par Confianza: ND & COMP-BARYON\par E04 / M07\\\addlinespace[3pt]
% VI-CATALOG-ROW anchuras 27
\textbf{27. Sigma(1385)\allowbreak{}+ WIDTH}\par Sigma(1385)\allowbreak{}- | Sigmabar(1385)\allowbreak{}+ | Sigma(1385)\allowbreak{}0 | Sigmabar(1385)\allowbreak{}0 | Sigma(1385)\allowbreak{}+ | Sigmabar(1385)\allowbreak{}-\par {\fontsize{8.1}{9.4}\selectfont PDG2026:\allowbreak{}observable:\allowbreak{}B043:\allowbreak{}B043W+:\allowbreak{}349879} & \textbf{36.2+-0.7} MeV & \(u_+\): 0.691323\allowbreak{}95371450\allowbreak{}484\par \(u_-\): 0.706567\allowbreak{}46664477\allowbreak{}266\par Tipo: NONE\par Confianza: ND & COMP-BARYON\par E04 / M07\\\addlinespace[3pt]
% VI-CATALOG-ROW anchuras 28
\textbf{28. Sigma(1385)\allowbreak{}0 WIDTH}\par Sigma(1385)\allowbreak{}- | Sigmabar(1385)\allowbreak{}+ | Sigma(1385)\allowbreak{}0 | Sigmabar(1385)\allowbreak{}0 | Sigma(1385)\allowbreak{}+ | Sigmabar(1385)\allowbreak{}-\par {\fontsize{8.1}{9.4}\selectfont PDG2026:\allowbreak{}observable:\allowbreak{}B043:\allowbreak{}B043W0:\allowbreak{}349880} & \textbf{44+-8} MeV & \(u_+\): 8.168042\allowbreak{}15502619\allowbreak{}97\par \(u_-\): 8.168042\allowbreak{}15502619\allowbreak{}97\par Tipo: NONE\par Confianza: ND & COMP-BARYON\par E04 / M07\\\addlinespace[3pt]
% VI-CATALOG-ROW anchuras 29
\textbf{29. Sigma(1385)\allowbreak{}- WIDTH}\par Sigma(1385)\allowbreak{}- | Sigmabar(1385)\allowbreak{}+ | Sigma(1385)\allowbreak{}0 | Sigmabar(1385)\allowbreak{}0 | Sigma(1385)\allowbreak{}+ | Sigmabar(1385)\allowbreak{}-\par {\fontsize{8.1}{9.4}\selectfont PDG2026:\allowbreak{}observable:\allowbreak{}B043:\allowbreak{}B043W-:\allowbreak{}349881} & \textbf{39.4+-2.1} MeV & \(u_+\): 2.080864\allowbreak{}71886561\allowbreak{}2\par \(u_-\): 2.080864\allowbreak{}71886561\allowbreak{}2\par Tipo: NONE\par Confianza: ND & COMP-BARYON\par E04 / M07\\\addlinespace[3pt]
% VI-CATALOG-ROW anchuras 30
\textbf{30. Sigma(1670) WIDTH}\par Sigma(1670)\allowbreak{}- | Sigmabar(1670)\allowbreak{}+ | Sigma(1670)\allowbreak{}0 | Sigmabar(1670)\allowbreak{}0 | Sigma(1670)\allowbreak{}+ | Sigmabar(1670)\allowbreak{}-\par {\fontsize{8.1}{9.4}\selectfont PDG2026:\allowbreak{}observable:\allowbreak{}B044:\allowbreak{}B044W:\allowbreak{}349947} & \textbf{40 to 70 to 100} MeV\par Presentación: 40 to 100 (\textasciitilde{}70) & \(u_+\): NA\par \(u_-\): NA\par Tipo: R\par Confianza: ND & COMP-BARYON\par E04 / M07\\\addlinespace[3pt]
% VI-CATALOG-ROW anchuras 31
\textbf{31. Sigma(1775) WIDTH}\par Sigma(1775)\allowbreak{}- | Sigmabar(1775)\allowbreak{}+ | Sigma(1775)\allowbreak{}0 | Sigmabar(1775)\allowbreak{}0 | Sigma(1775)\allowbreak{}+ | Sigmabar(1775)\allowbreak{}-\par {\fontsize{8.1}{9.4}\selectfont PDG2026:\allowbreak{}observable:\allowbreak{}B045:\allowbreak{}B045W:\allowbreak{}349984} & \textbf{105 TO 120 TO 135} MeV\par Presentación: 105 to 135 (\textasciitilde{}120) & \(u_+\): NA\par \(u_-\): NA\par Tipo: R\par Confianza: ND & COMP-BARYON\par E04 / M07\\\addlinespace[3pt]
% VI-CATALOG-ROW anchuras 32
\textbf{32. Sigma(1915) WIDTH}\par Sigma(1915)\allowbreak{}- | Sigmabar(1915)\allowbreak{}+ | Sigma(1915)\allowbreak{}0 | Sigmabar(1915)\allowbreak{}0 | Sigma(1915)\allowbreak{}+ | Sigmabar(1915)\allowbreak{}-\par {\fontsize{8.1}{9.4}\selectfont PDG2026:\allowbreak{}observable:\allowbreak{}B046:\allowbreak{}B046W:\allowbreak{}350081} & \textbf{80 TO 120 TO 160} MeV\par Presentación: 80 to 160 (\textasciitilde{}120) & \(u_+\): NA\par \(u_-\): NA\par Tipo: R\par Confianza: ND & COMP-BARYON\par E04 / M07\\\addlinespace[3pt]
% VI-CATALOG-ROW anchuras 33
\textbf{33. Sigma(2030) WIDTH}\par Sigma(2030)\allowbreak{}- | Sigmabar(2030)\allowbreak{}+ | Sigma(2030)\allowbreak{}0 | Sigmabar(2030)\allowbreak{}0 | Sigma(2030)\allowbreak{}+ | Sigmabar(2030)\allowbreak{}-\par {\fontsize{8.1}{9.4}\selectfont PDG2026:\allowbreak{}observable:\allowbreak{}B047:\allowbreak{}B047W:\allowbreak{}350167} & \textbf{150 TO 180 TO 200} MeV\par Presentación: 150 to 200 (\textasciitilde{}180) & \(u_+\): NA\par \(u_-\): NA\par Tipo: R\par Confianza: ND & COMP-BARYON\par E04 / M07\\\addlinespace[3pt]
% VI-CATALOG-ROW anchuras 34
\textbf{34. Sigma(2250) WIDTH}\par Sigma(2250)\allowbreak{}- | Sigmabar(2250)\allowbreak{}+ | Sigma(2250)\allowbreak{}0 | Sigmabar(2250)\allowbreak{}0 | Sigma(2250)\allowbreak{}+ | Sigmabar(2250)\allowbreak{}-\par {\fontsize{8.1}{9.4}\selectfont PDG2026:\allowbreak{}observable:\allowbreak{}B048:\allowbreak{}B048W:\allowbreak{}350278} & \textbf{60 TO 100 TO 150} MeV\par Presentación: 60 to 150 (\textasciitilde{}100) & \(u_+\): NA\par \(u_-\): NA\par Tipo: R\par Confianza: ND & COMP-BARYON\par E04 / M07\\\addlinespace[3pt]
% VI-CATALOG-ROW anchuras 35
\textbf{35. Xi(1530)\allowbreak{}0 WIDTH}\par Xi(1530)\allowbreak{}- | Xibar(1530)\allowbreak{}+ | Xi(1530)\allowbreak{}0 | Xibar(1530)\allowbreak{}0\par {\fontsize{8.1}{9.4}\selectfont PDG2026:\allowbreak{}observable:\allowbreak{}B049:\allowbreak{}B049W0:\allowbreak{}350405} & \textbf{9.1+-0.5} MeV & \(u_+\): 0.480582\allowbreak{}71927942\allowbreak{}953\par \(u_-\): 0.480582\allowbreak{}71927942\allowbreak{}953\par Tipo: NONE\par Confianza: ND & COMP-BARYON\par E04 / M07\\\addlinespace[3pt]
% VI-CATALOG-ROW anchuras 36
\textbf{36. Xi(1530)\allowbreak{}- WIDTH}\par Xi(1530)\allowbreak{}- | Xibar(1530)\allowbreak{}+ | Xi(1530)\allowbreak{}0 | Xibar(1530)\allowbreak{}0\par {\fontsize{8.1}{9.4}\selectfont PDG2026:\allowbreak{}observable:\allowbreak{}B049:\allowbreak{}B049W-:\allowbreak{}350406} & \textbf{9.9+1.7-1.9} MeV & \(u_+\): 1.731350\allowbreak{}23766534\allowbreak{}4\par \(u_-\): 1.930203\allowbreak{}24724418\allowbreak{}7\par Tipo: NONE\par Confianza: ND & COMP-BARYON\par E04 / M07\\\addlinespace[3pt]
% VI-CATALOG-ROW anchuras 37
\textbf{37. Xi(1820) WIDTH}\par Xi(1820)\allowbreak{}- | Xibar(1820)\allowbreak{}+ | Xi(1820)\allowbreak{}0 | Xibar(1820)\allowbreak{}0\par {\fontsize{8.1}{9.4}\selectfont PDG2026:\allowbreak{}observable:\allowbreak{}B050:\allowbreak{}B050W:\allowbreak{}350422} & \textbf{24+15-10} MeV & \(u_+\): 15\par \(u_-\): 10\par Tipo: NONE\par Confianza: ND & COMP-BARYON\par E04 / M07\\\addlinespace[3pt]
% VI-CATALOG-ROW anchuras 38
\textbf{38. Xi(1950) WIDTH}\par Xi(1950)\allowbreak{}- | Xibar(1950)\allowbreak{}+ | Xi(1950)\allowbreak{}0 | Xibar(1950)\allowbreak{}0\par {\fontsize{8.1}{9.4}\selectfont PDG2026:\allowbreak{}observable:\allowbreak{}B052:\allowbreak{}B052W:\allowbreak{}350438} & \textbf{60+-20} MeV & \(u_+\): 20\par \(u_-\): 20\par Tipo: NONE\par Confianza: ND & COMP-BARYON\par E04 / M07\\\addlinespace[3pt]
% VI-CATALOG-ROW anchuras 39
\textbf{39. Lambda(1690) WIDTH}\par Lambda(1690)\allowbreak{}0 | Lambdabar(1690)\allowbreak{}0\par {\fontsize{8.1}{9.4}\selectfont PDG2026:\allowbreak{}observable:\allowbreak{}B055:\allowbreak{}B055W:\allowbreak{}349519} & \textbf{60 to 70 to 80} MeV\par Presentación: 60 to 80 (\textasciitilde{}70) & \(u_+\): NA\par \(u_-\): NA\par Tipo: R\par Confianza: ND & COMP-BARYON\par E04 / M07\\\addlinespace[3pt]
% VI-CATALOG-ROW anchuras 40
\textbf{40. Lambda(1830) WIDTH}\par Lambda(1830)\allowbreak{}0 | Lambdabar(1830)\allowbreak{}0\par {\fontsize{8.1}{9.4}\selectfont PDG2026:\allowbreak{}observable:\allowbreak{}B056:\allowbreak{}B056W:\allowbreak{}349615} & \textbf{60 to 90 to 120} MeV\par Presentación: 60 to 120 (\textasciitilde{}90) & \(u_+\): NA\par \(u_-\): NA\par Tipo: R\par Confianza: ND & COMP-BARYON\par E04 / M07\\\addlinespace[3pt]
% VI-CATALOG-ROW anchuras 41
\textbf{41. Sigma(1750) WIDTH}\par Sigma(1750)\allowbreak{}- | Sigmabar(1750)\allowbreak{}+ | Sigma(1750)\allowbreak{}0 | Sigmabar(1750)\allowbreak{}0 | Sigma(1750)\allowbreak{}+ | Sigmabar(1750)\allowbreak{}-\par {\fontsize{8.1}{9.4}\selectfont PDG2026:\allowbreak{}observable:\allowbreak{}B057:\allowbreak{}B057W:\allowbreak{}349967} & \textbf{100 to 150 to 200} MeV\par Presentación: 100 to 200 (\textasciitilde{}150) & \(u_+\): NA\par \(u_-\): NA\par Tipo: R\par Confianza: ND & COMP-BARYON\par E04 / M07\\\addlinespace[3pt]
% VI-CATALOG-ROW anchuras 42
\textbf{42. Lambda(1890) WIDTH}\par Lambda(1890)\allowbreak{}0 | Lambdabar(1890)\allowbreak{}0\par {\fontsize{8.1}{9.4}\selectfont PDG2026:\allowbreak{}observable:\allowbreak{}B060:\allowbreak{}B060W:\allowbreak{}349636} & \textbf{80 to 120 to 160} MeV\par Presentación: 80 to 160 (\textasciitilde{}120) & \(u_+\): NA\par \(u_-\): NA\par Tipo: R\par Confianza: ND & COMP-BARYON\par E04 / M07\\\addlinespace[3pt]
% VI-CATALOG-ROW anchuras 43
\textbf{43. N(1440) BREIT-WIGNER WIDTH}\par N(1440)\allowbreak{}0 | Nbar(1440)\allowbreak{}0 | N(1440)\allowbreak{}+ | Nbar(1440)\allowbreak{}-\par {\fontsize{8.1}{9.4}\selectfont PDG2026:\allowbreak{}observable:\allowbreak{}B061:\allowbreak{}B061W:\allowbreak{}348327} & \textbf{250 to 350 to 450} MeV\par Presentación: 250 to 450 (\textasciitilde{}350) & \(u_+\): NA\par \(u_-\): NA\par Tipo: R\par Confianza: ND & COMP-BARYON\par E04 / M07\\\addlinespace[3pt]
% VI-CATALOG-ROW anchuras 44
\textbf{44. N(1520) BREIT-WIGNER WIDTH}\par N(1520)\allowbreak{}0 | Nbar(1520)\allowbreak{}0 | N(1520)\allowbreak{}+ | Nbar(1520)\allowbreak{}-\par {\fontsize{8.1}{9.4}\selectfont PDG2026:\allowbreak{}observable:\allowbreak{}B062:\allowbreak{}B062W:\allowbreak{}348350} & \textbf{100 to 110 to 120} MeV\par Presentación: 100 to 120 (\textasciitilde{}110) & \(u_+\): NA\par \(u_-\): NA\par Tipo: R\par Confianza: ND & COMP-BARYON\par E04 / M07\\\addlinespace[3pt]
% VI-CATALOG-ROW anchuras 45
\textbf{45. N(1535) BREIT-WIGNER WIDTH}\par N(1535)\allowbreak{}0 | Nbar(1535)\allowbreak{}0 | N(1535)\allowbreak{}+ | Nbar(1535)\allowbreak{}-\par {\fontsize{8.1}{9.4}\selectfont PDG2026:\allowbreak{}observable:\allowbreak{}B063:\allowbreak{}B063W:\allowbreak{}348384} & \textbf{125 TO 150 TO 175} MeV\par Presentación: 125 to 175 (\textasciitilde{}150) & \(u_+\): NA\par \(u_-\): NA\par Tipo: R\par Confianza: ND & COMP-BARYON\par E04 / M07\\\addlinespace[3pt]
% VI-CATALOG-ROW anchuras 46
\textbf{46. N(1675) BREIT-WIGNER WIDTH}\par N(1675)\allowbreak{}0 | Nbar(1675)\allowbreak{}0 | N(1675)\allowbreak{}+ | Nbar(1675)\allowbreak{}-\par {\fontsize{8.1}{9.4}\selectfont PDG2026:\allowbreak{}observable:\allowbreak{}B064:\allowbreak{}B064W:\allowbreak{}348436} & \textbf{130 to 145 to 160} MeV\par Presentación: 130 to 160 (\textasciitilde{}145) & \(u_+\): NA\par \(u_-\): NA\par Tipo: R\par Confianza: ND & COMP-BARYON\par E04 / M07\\\addlinespace[3pt]
% VI-CATALOG-ROW anchuras 47
\textbf{47. N(1680) BREIT-WIGNER WIDTH}\par N(1680)\allowbreak{}0 | Nbar(1680)\allowbreak{}0 | N(1680)\allowbreak{}+ | Nbar(1680)\allowbreak{}-\par {\fontsize{8.1}{9.4}\selectfont PDG2026:\allowbreak{}observable:\allowbreak{}B065:\allowbreak{}B065W:\allowbreak{}348468} & \textbf{115 to 120 to 130} MeV\par Presentación: 115 to 130 (\textasciitilde{}120) & \(u_+\): NA\par \(u_-\): NA\par Tipo: R\par Confianza: ND & COMP-BARYON\par E04 / M07\\\addlinespace[3pt]
% VI-CATALOG-ROW anchuras 48
\textbf{48. N(1650) BREIT-WIGNER WIDTH}\par N(1650)\allowbreak{}0 | Nbar(1650)\allowbreak{}0 | N(1650)\allowbreak{}+ | Nbar(1650)\allowbreak{}-\par {\fontsize{8.1}{9.4}\selectfont PDG2026:\allowbreak{}observable:\allowbreak{}B066:\allowbreak{}B066W:\allowbreak{}348411} & \textbf{100 to 125 to 150} MeV\par Presentación: 100 to 150 (\textasciitilde{}125) & \(u_+\): NA\par \(u_-\): NA\par Tipo: R\par Confianza: ND & COMP-BARYON\par E04 / M07\\\addlinespace[3pt]
% VI-CATALOG-ROW anchuras 49
\textbf{49. Sigma(1880) WIDTH}\par Sigma(1880)\allowbreak{}- | Sigmabar(1880)\allowbreak{}+ | Sigma(1880)\allowbreak{}0 | Sigmabar(1880)\allowbreak{}0 | Sigma(1880)\allowbreak{}+ | Sigmabar(1880)\allowbreak{}-\par {\fontsize{8.1}{9.4}\selectfont PDG2026:\allowbreak{}observable:\allowbreak{}B067:\allowbreak{}B067W:\allowbreak{}350000} & \textbf{100 to 200 to 300} MeV\par Presentación: 100 to 300 (\textasciitilde{}200) & \(u_+\): NA\par \(u_-\): NA\par Tipo: R\par Confianza: ND & COMP-BARYON\par E04 / M07\\\addlinespace[3pt]
% VI-CATALOG-ROW anchuras 50
\textbf{50. Xi(2030) WIDTH}\par Xi(2030)\allowbreak{}- | Xibar(2030)\allowbreak{}+ | Xi(2030)\allowbreak{}0 | Xibar(2030)\allowbreak{}0\par {\fontsize{8.1}{9.4}\selectfont PDG2026:\allowbreak{}observable:\allowbreak{}B068:\allowbreak{}B068W:\allowbreak{}350445} & \textbf{20+15-5} MeV & \(u_+\): 15\par \(u_-\): 5\par Tipo: NONE\par Confianza: ND & COMP-BARYON\par E04 / M07\\\addlinespace[3pt]
% VI-CATALOG-ROW anchuras 51
\textbf{51. N(2190) BREIT-WIGNER WIDTH}\par N(2190)\allowbreak{}0 | Nbar(2190)\allowbreak{}0 | N(2190)\allowbreak{}+ | Nbar(2190)\allowbreak{}-\par {\fontsize{8.1}{9.4}\selectfont PDG2026:\allowbreak{}observable:\allowbreak{}B071:\allowbreak{}B071W:\allowbreak{}348954} & \textbf{300 to 400 to 500} MeV\par Presentación: 300 to 500 (\textasciitilde{}400) & \(u_+\): NA\par \(u_-\): NA\par Tipo: R\par Confianza: ND & COMP-BARYON\par E04 / M07\\\addlinespace[3pt]
% VI-CATALOG-ROW anchuras 52
\textbf{52. Lambda(1810) WIDTH}\par Lambda(1810)\allowbreak{}0 | Lambdabar(1810)\allowbreak{}0\par {\fontsize{8.1}{9.4}\selectfont PDG2026:\allowbreak{}observable:\allowbreak{}B077:\allowbreak{}B077W:\allowbreak{}349577} & \textbf{50 to 110 to 170} MeV\par Presentación: 50 to 170 (\textasciitilde{}110) & \(u_+\): NA\par \(u_-\): NA\par Tipo: R\par Confianza: ND & COMP-BARYON\par E04 / M07\\\addlinespace[3pt]
% VI-CATALOG-ROW anchuras 53
\textbf{53. Sigma(1660) WIDTH}\par Sigma(1660)\allowbreak{}- | Sigmabar(1660)\allowbreak{}+ | Sigma(1660)\allowbreak{}0 | Sigmabar(1660)\allowbreak{}0 | Sigma(1660)\allowbreak{}+ | Sigmabar(1660)\allowbreak{}-\par {\fontsize{8.1}{9.4}\selectfont PDG2026:\allowbreak{}observable:\allowbreak{}B079:\allowbreak{}B079W:\allowbreak{}349921} & \textbf{100 to 200 to 300} MeV\par Presentación: 100 to 300 (\textasciitilde{}200) & \(u_+\): NA\par \(u_-\): NA\par Tipo: R\par Confianza: ND & COMP-BARYON\par E04 / M07\\\addlinespace[3pt]
% VI-CATALOG-ROW anchuras 54
\textbf{54. Delta(1620) BREIT-WIGNER WIDTH}\par Delta(1620)\allowbreak{}- | Deltabar(1620)\allowbreak{}+ | Delta(1620)\allowbreak{}0 | Deltabar(1620)\allowbreak{}0 | Delta(1620)\allowbreak{}+ | Deltabar(1620)\allowbreak{}- | Delta(1620)\allowbreak{}++ | Deltabar(1620)\allowbreak{}-\kern0pt{}-\par {\fontsize{8.1}{9.4}\selectfont PDG2026:\allowbreak{}observable:\allowbreak{}B082:\allowbreak{}B082W:\allowbreak{}349079} & \textbf{110 to 130 to 150} MeV\par Presentación: 110 to 150 (\textasciitilde{}130) & \(u_+\): NA\par \(u_-\): NA\par Tipo: R\par Confianza: ND & COMP-BARYON\par E04 / M07\\\addlinespace[3pt]
% VI-CATALOG-ROW anchuras 55
\textbf{55. Delta(1950) BREIT-WIGNER WIDTH}\par Delta(1950)\allowbreak{}- | Deltabar(1950)\allowbreak{}+ | Delta(1950)\allowbreak{}0 | Deltabar(1950)\allowbreak{}0 | Delta(1950)\allowbreak{}+ | Deltabar(1950)\allowbreak{}- | Delta(1950)\allowbreak{}++ | Deltabar(1950)\allowbreak{}-\kern0pt{}-\par {\fontsize{8.1}{9.4}\selectfont PDG2026:\allowbreak{}observable:\allowbreak{}B083:\allowbreak{}B083W:\allowbreak{}349309} & \textbf{235 to 285 to 335} MeV\par Presentación: 235 to 335 (\textasciitilde{}285) & \(u_+\): NA\par \(u_-\): NA\par Tipo: R\par Confianza: ND & COMP-BARYON\par E04 / M07\\\addlinespace[3pt]
% VI-CATALOG-ROW anchuras 56
\textbf{56. Delta(2420) BREIT-WIGNER WIDTH}\par Delta(2420)\allowbreak{}- | Deltabar(2420)\allowbreak{}+ | Delta(2420)\allowbreak{}0 | Deltabar(2420)\allowbreak{}0 | Delta(2420)\allowbreak{}+ | Deltabar(2420)\allowbreak{}- | Delta(2420)\allowbreak{}++ | Deltabar(2420)\allowbreak{}-\kern0pt{}-\par {\fontsize{8.1}{9.4}\selectfont PDG2026:\allowbreak{}observable:\allowbreak{}B084:\allowbreak{}B084W:\allowbreak{}349360} & \textbf{300 to 500 to 700} MeV\par Presentación: 300 to 700 (\textasciitilde{}500) & \(u_+\): NA\par \(u_-\): NA\par Tipo: R\par Confianza: ND & COMP-BARYON\par E04 / M07\\\addlinespace[3pt]
% VI-CATALOG-ROW anchuras 57
\textbf{57. N(1880) BREIT-WIGNER WIDTH}\par N(1880)\allowbreak{}0 | Nbar(1880)\allowbreak{}0 | N(1880)\allowbreak{}+ | Nbar(1880)\allowbreak{}-\par {\fontsize{8.1}{9.4}\selectfont PDG2026:\allowbreak{}observable:\allowbreak{}B087:\allowbreak{}B087W:\allowbreak{}348672} & \textbf{200 to 300 to 400} MeV\par Presentación: 200 to 400 (\textasciitilde{}300) & \(u_+\): NA\par \(u_-\): NA\par Tipo: R\par Confianza: ND & COMP-BARYON\par E04 / M07\\\addlinespace[3pt]
% VI-CATALOG-ROW anchuras 58
\textbf{58. Sigma(2080) WIDTH}\par Sigma(2080)\allowbreak{}- | Sigmabar(2080)\allowbreak{}+ | Sigma(2080)\allowbreak{}0 | Sigmabar(2080)\allowbreak{}0 | Sigma(2080)\allowbreak{}+ | Sigmabar(2080)\allowbreak{}-\par {\fontsize{8.1}{9.4}\selectfont PDG2026:\allowbreak{}observable:\allowbreak{}B088:\allowbreak{}B088W:\allowbreak{}350189} & \textbf{100 to 170 to 240} MeV\par Presentación: 100 to 240 (\textasciitilde{}170) & \(u_+\): NA\par \(u_-\): NA\par Tipo: R\par Confianza: ND & COMP-BARYON\par E04 / M07\\\addlinespace[3pt]
% VI-CATALOG-ROW anchuras 59
\textbf{59. N(2220) BREIT-WIGNER WIDTH}\par N(2220)\allowbreak{}0 | Nbar(2220)\allowbreak{}0 | N(2220)\allowbreak{}+ | Nbar(2220)\allowbreak{}-\par {\fontsize{8.1}{9.4}\selectfont PDG2026:\allowbreak{}observable:\allowbreak{}B090:\allowbreak{}B090W:\allowbreak{}348987} & \textbf{350 TO 400 TO 500} MeV\par Presentación: 350 to 500 (\textasciitilde{}400) & \(u_+\): NA\par \(u_-\): NA\par Tipo: R\par Confianza: ND & COMP-BARYON\par E04 / M07\\\addlinespace[3pt]
% VI-CATALOG-ROW anchuras 60
\textbf{60. Sigma(1910) WIDTH}\par Sigma(1910)\allowbreak{}- | Sigmabar(1910)\allowbreak{}+ | Sigma(1910)\allowbreak{}0 | Sigmabar(1910)\allowbreak{}0 | Sigma(1910)\allowbreak{}+ | Sigmabar(1910)\allowbreak{}-\par {\fontsize{8.1}{9.4}\selectfont PDG2026:\allowbreak{}observable:\allowbreak{}B098:\allowbreak{}B098W:\allowbreak{}350051} & \textbf{150 to 220 to 300} MeV\par Presentación: 150 to 300 (\textasciitilde{}220) & \(u_+\): NA\par \(u_-\): NA\par Tipo: R\par Confianza: ND & COMP-BARYON\par E04 / M07\\\addlinespace[3pt]
% VI-CATALOG-ROW anchuras 61
\textbf{61. Lambda(1600) WIDTH}\par Lambda(1600)\allowbreak{}0 | Lambdabar(1600)\allowbreak{}0\par {\fontsize{8.1}{9.4}\selectfont PDG2026:\allowbreak{}observable:\allowbreak{}B101:\allowbreak{}B101W:\allowbreak{}349473} & \textbf{150 to 200 to 250} MeV\par Presentación: 150 to 250 (\textasciitilde{}200) & \(u_+\): NA\par \(u_-\): NA\par Tipo: R\par Confianza: ND & COMP-BARYON\par E04 / M07\\\addlinespace[3pt]
% VI-CATALOG-ROW anchuras 62
\textbf{62. Lambda\_\allowbreak{}c(2625)\allowbreak{}+ WIDTH}\par Lambda\_\allowbreak{}c(2625)\allowbreak{}+ | Lambdabar\_\allowbreak{}c(2625)\allowbreak{}-\par {\fontsize{8.1}{9.4}\selectfont PDG2026:\allowbreak{}observable:\allowbreak{}B102:\allowbreak{}B102W:\allowbreak{}350937} & \textbf{< 0.52} MeV\par Presentación: <0.52 & \(u_+\): NA\par \(u_-\): NA\par Tipo: U\par Confianza: 90 & COMP-BARYON\par E04 / M07\\\addlinespace[3pt]
% VI-CATALOG-ROW anchuras 63
\textbf{63. Sigma\_\allowbreak{}c(2455)\allowbreak{}++ WIDTH}\par Sigma\_\allowbreak{}c(2455)\allowbreak{}0 | Sigmabar\_\allowbreak{}c(2455)\allowbreak{}0 | Sigma\_\allowbreak{}c(2455)\allowbreak{}+ | Sigmabar\_\allowbreak{}c(2455)\allowbreak{}- | Sigma\_\allowbreak{}c(2455)\allowbreak{}++ | Sigmabar\_\allowbreak{}c(2455)\allowbreak{}-\kern0pt{}-\par {\fontsize{8.1}{9.4}\selectfont PDG2026:\allowbreak{}observable:\allowbreak{}B104:\allowbreak{}B104W++:\allowbreak{}350997} & \textbf{1.89+0.09-0.18} MeV & \(u_+\): 0.085968\allowbreak{}74087768\allowbreak{}8063\par \(u_-\): 0.177414\allowbreak{}42138916\allowbreak{}921\par Tipo: NONE\par Confianza: ND & COMP-BARYON\par E04 / M07\\\addlinespace[3pt]
% VI-CATALOG-ROW anchuras 64
\textbf{64. Sigma\_\allowbreak{}c(2455)\allowbreak{}+ WIDTH}\par Sigma\_\allowbreak{}c(2455)\allowbreak{}0 | Sigmabar\_\allowbreak{}c(2455)\allowbreak{}0 | Sigma\_\allowbreak{}c(2455)\allowbreak{}+ | Sigmabar\_\allowbreak{}c(2455)\allowbreak{}- | Sigma\_\allowbreak{}c(2455)\allowbreak{}++ | Sigmabar\_\allowbreak{}c(2455)\allowbreak{}-\kern0pt{}-\par {\fontsize{8.1}{9.4}\selectfont PDG2026:\allowbreak{}observable:\allowbreak{}B104:\allowbreak{}B104W+:\allowbreak{}350998} & \textbf{2.3+-0.4} MeV & \(u_+\): 0.424264\allowbreak{}06871192\allowbreak{}851\par \(u_-\): 0.424264\allowbreak{}06871192\allowbreak{}851\par Tipo: NONE\par Confianza: ND & COMP-BARYON\par E04 / M07\\\addlinespace[3pt]
% VI-CATALOG-ROW anchuras 65
\textbf{65. Sigma\_\allowbreak{}c(2455)\allowbreak{}0 WIDTH}\par Sigma\_\allowbreak{}c(2455)\allowbreak{}0 | Sigmabar\_\allowbreak{}c(2455)\allowbreak{}0 | Sigma\_\allowbreak{}c(2455)\allowbreak{}+ | Sigmabar\_\allowbreak{}c(2455)\allowbreak{}- | Sigma\_\allowbreak{}c(2455)\allowbreak{}++ | Sigmabar\_\allowbreak{}c(2455)\allowbreak{}-\kern0pt{}-\par {\fontsize{8.1}{9.4}\selectfont PDG2026:\allowbreak{}observable:\allowbreak{}B104:\allowbreak{}B104W0:\allowbreak{}350999} & \textbf{1.83+0.11-0.19} MeV & \(u_+\): 0.111341\allowbreak{}94993665\allowbreak{}61\par \(u_-\): 0.193479\allowbreak{}36378811\allowbreak{}631\par Tipo: NONE\par Confianza: ND & COMP-BARYON\par E04 / M07\\\addlinespace[3pt]
% VI-CATALOG-ROW anchuras 66
\textbf{66. MIXED CHARGES}\par Xi(1690)\allowbreak{}- | Xibar(1690)\allowbreak{}+ | Xi(1690)\allowbreak{}0 | Xibar(1690)\allowbreak{}0\par {\fontsize{8.1}{9.4}\selectfont PDG2026:\allowbreak{}observable:\allowbreak{}B105:\allowbreak{}B105W:\allowbreak{}350413} & \textbf{20+-15} MeV & \(u_+\): 15\par \(u_-\): 15\par Tipo: NONE\par Confianza: ND & COMP-BARYON\par E04 / M07\\\addlinespace[3pt]
% VI-CATALOG-ROW anchuras 67
\textbf{67. N(2250) BREIT-WIGNER WIDTH}\par N(2250)\allowbreak{}0 | Nbar(2250)\allowbreak{}0 | N(2250)\allowbreak{}+ | Nbar(2250)\allowbreak{}-\par {\fontsize{8.1}{9.4}\selectfont PDG2026:\allowbreak{}observable:\allowbreak{}B113:\allowbreak{}B113W:\allowbreak{}349005} & \textbf{300 to 500 to 600} MeV\par Presentación: 300 to 600 (\textasciitilde{}500) & \(u_+\): NA\par \(u_-\): NA\par Tipo: R\par Confianza: ND & COMP-BARYON\par E04 / M07\\\addlinespace[3pt]
% VI-CATALOG-ROW anchuras 68
\textbf{68. Sigma\_\allowbreak{}c(2520)\allowbreak{}++ WIDTH}\par Sigma\_\allowbreak{}c(2520)\allowbreak{}0 | Sigmabar\_\allowbreak{}c(2520)\allowbreak{}0 | Sigma\_\allowbreak{}c(2520)\allowbreak{}+ | Sigmabar\_\allowbreak{}c(2520)\allowbreak{}- | Sigma\_\allowbreak{}c(2520)\allowbreak{}++ | Sigmabar\_\allowbreak{}c(2520)\allowbreak{}-\kern0pt{}-\par {\fontsize{8.1}{9.4}\selectfont PDG2026:\allowbreak{}observable:\allowbreak{}B115:\allowbreak{}B115W++:\allowbreak{}351011} & \textbf{14.78+0.30-0.40} MeV & \(u_+\): 0.302210\allowbreak{}98294841\allowbreak{}742\par \(u_-\): 0.378217\allowbreak{}26734358\allowbreak{}602\par Tipo: NONE\par Confianza: ND & COMP-BARYON\par E04 / M07\\\addlinespace[3pt]
% VI-CATALOG-ROW anchuras 69
\textbf{69. Sigma\_\allowbreak{}c(2520)\allowbreak{}+ WIDTH}\par Sigma\_\allowbreak{}c(2520)\allowbreak{}0 | Sigmabar\_\allowbreak{}c(2520)\allowbreak{}0 | Sigma\_\allowbreak{}c(2520)\allowbreak{}+ | Sigmabar\_\allowbreak{}c(2520)\allowbreak{}- | Sigma\_\allowbreak{}c(2520)\allowbreak{}++ | Sigmabar\_\allowbreak{}c(2520)\allowbreak{}-\kern0pt{}-\par {\fontsize{8.1}{9.4}\selectfont PDG2026:\allowbreak{}observable:\allowbreak{}B115:\allowbreak{}B115W+:\allowbreak{}351012} & \textbf{17.2+4.0-2.2} MeV & \(u_+\): 3.860051\allowbreak{}81312375\allowbreak{}71\par \(u_-\): 2.213594\allowbreak{}36211786\allowbreak{}6\par Tipo: NONE\par Confianza: ND & COMP-BARYON\par E04 / M07\\\addlinespace[3pt]
% VI-CATALOG-ROW anchuras 70
\textbf{70. Sigma\_\allowbreak{}c(2520)\allowbreak{}0 WIDTH}\par Sigma\_\allowbreak{}c(2520)\allowbreak{}0 | Sigmabar\_\allowbreak{}c(2520)\allowbreak{}0 | Sigma\_\allowbreak{}c(2520)\allowbreak{}+ | Sigmabar\_\allowbreak{}c(2520)\allowbreak{}- | Sigma\_\allowbreak{}c(2520)\allowbreak{}++ | Sigmabar\_\allowbreak{}c(2520)\allowbreak{}-\kern0pt{}-\par {\fontsize{8.1}{9.4}\selectfont PDG2026:\allowbreak{}observable:\allowbreak{}B115:\allowbreak{}B115W0:\allowbreak{}351013} & \textbf{15.3+0.4-0.5} MeV & \(u_+\): 0.438053\allowbreak{}27113015\allowbreak{}552\par \(u_-\): 0.491725\allowbreak{}79862174\allowbreak{}309\par Tipo: NONE\par Confianza: ND & COMP-BARYON\par E04 / M07\\\addlinespace[3pt]
% VI-CATALOG-ROW anchuras 71
\textbf{71. Delta(1920) BREIT-WIGNER WIDTH}\par Delta(1920)\allowbreak{}- | Deltabar(1920)\allowbreak{}+ | Delta(1920)\allowbreak{}0 | Deltabar(1920)\allowbreak{}0 | Delta(1920)\allowbreak{}+ | Deltabar(1920)\allowbreak{}- | Delta(1920)\allowbreak{}++ | Deltabar(1920)\allowbreak{}-\kern0pt{}-\par {\fontsize{8.1}{9.4}\selectfont PDG2026:\allowbreak{}observable:\allowbreak{}B117:\allowbreak{}B117W:\allowbreak{}349209} & \textbf{240 to 300 to 360} MeV\par Presentación: 240 to 360 (\textasciitilde{}300) & \(u_+\): NA\par \(u_-\): NA\par Tipo: R\par Confianza: ND & COMP-BARYON\par E04 / M07\\\addlinespace[3pt]
% VI-CATALOG-ROW anchuras 72
\textbf{72. Lambda\_\allowbreak{}c(2595)\allowbreak{}+ WIDTH}\par Lambda\_\allowbreak{}c(2595)\allowbreak{}+ | Lambdabar\_\allowbreak{}c(2595)\allowbreak{}-\par {\fontsize{8.1}{9.4}\selectfont PDG2026:\allowbreak{}observable:\allowbreak{}B119:\allowbreak{}B119W:\allowbreak{}350920} & \textbf{2.6+-0.6} MeV & \(u_+\): 0.557584\allowbreak{}07437802\allowbreak{}593\par \(u_-\): 0.557584\allowbreak{}07437802\allowbreak{}593\par Tipo: NONE\par Confianza: ND & COMP-BARYON\par E04 / M07\\\addlinespace[3pt]
% VI-CATALOG-ROW anchuras 73
\textbf{73. N(2600) BREIT-WIGNER WIDTH}\par N(2600)\allowbreak{}0 | Nbar(2600)\allowbreak{}0 | N(2600)\allowbreak{}+ | Nbar(2600)\allowbreak{}-\par {\fontsize{8.1}{9.4}\selectfont PDG2026:\allowbreak{}observable:\allowbreak{}B120:\allowbreak{}B120W:\allowbreak{}349023} & \textbf{500 TO 650 TO 800} MeV\par Presentación: 500 to 800 (\textasciitilde{}650) & \(u_+\): NA\par \(u_-\): NA\par Tipo: R\par Confianza: ND & COMP-BARYON\par E04 / M07\\\addlinespace[3pt]
% VI-CATALOG-ROW anchuras 74
\textbf{74. Lambda\_\allowbreak{}c(2940)\allowbreak{}+ WIDTH}\par Lambda\_\allowbreak{}c(2940)\allowbreak{}0 | Lambdabar\_\allowbreak{}c(2940)\allowbreak{}0\par {\fontsize{8.1}{9.4}\selectfont PDG2026:\allowbreak{}observable:\allowbreak{}B122:\allowbreak{}B122W:\allowbreak{}350980} & \textbf{20+6-5} MeV & \(u_+\): 5.981963\allowbreak{}79520038\allowbreak{}41\par \(u_-\): 5.228538\allowbreak{}13129437\allowbreak{}66\par Tipo: NONE\par Confianza: ND & COMP-BARYON\par E04 / M07\\\addlinespace[3pt]
% VI-CATALOG-ROW anchuras 75
\textbf{75. Xi\_\allowbreak{}c(2970)\allowbreak{}+ WIDTH}\par Xi\_\allowbreak{}c(2970)\allowbreak{}0 | Xibar\_\allowbreak{}c(2970)\allowbreak{}0\par {\fontsize{8.1}{9.4}\selectfont PDG2026:\allowbreak{}observable:\allowbreak{}B130:\allowbreak{}B130W+:\allowbreak{}351281} & \textbf{27.2+2.8-3.3} MeV & \(u_+\): 2.769172\allowbreak{}69976361\allowbreak{}79\par \(u_-\): 3.314747\allowbreak{}83369252\allowbreak{}31\par Tipo: NONE\par Confianza: ND & COMP-BARYON\par E04 / M07\\\addlinespace[3pt]
% VI-CATALOG-ROW anchuras 76
\textbf{76. N(2100) BREIT-WIGNER WIDTH}\par N(2100)\allowbreak{}0 | Nbar(2100)\allowbreak{}0 | N(2100)\allowbreak{}+ | Nbar(2100)\allowbreak{}-\par {\fontsize{8.1}{9.4}\selectfont PDG2026:\allowbreak{}observable:\allowbreak{}B132:\allowbreak{}B132W:\allowbreak{}348872} & \textbf{200 to 260 to 320} MeV\par Presentación: 200 to 320 (\textasciitilde{}260) & \(u_+\): NA\par \(u_-\): NA\par Tipo: R\par Confianza: ND & COMP-BARYON\par E04 / M07\\\addlinespace[3pt]
% VI-CATALOG-ROW anchuras 77
\textbf{77. Delta(2200) BREIT-WIGNER WIDTH}\par Delta(2200)\allowbreak{}- | Deltabar(2200)\allowbreak{}+ | Delta(2200)\allowbreak{}0 | Deltabar(2200)\allowbreak{}0 | Delta(2200)\allowbreak{}+ | Deltabar(2200)\allowbreak{}- | Delta(2200)\allowbreak{}++ | Deltabar(2200)\allowbreak{}-\kern0pt{}-\par {\fontsize{8.1}{9.4}\selectfont PDG2026:\allowbreak{}observable:\allowbreak{}B135:\allowbreak{}B135W:\allowbreak{}349340} & \textbf{200 to 350 to 500} MeV\par Presentación: 200 to 500 (\textasciitilde{}350) & \(u_+\): NA\par \(u_-\): NA\par Tipo: R\par Confianza: ND & COMP-BARYON\par E04 / M07\\\addlinespace[3pt]
% VI-CATALOG-ROW anchuras 78
\textbf{78. Delta(1940) BREIT-WIGNER WIDTH}\par Delta(1940)\allowbreak{}- | Deltabar(1940)\allowbreak{}+ | Delta(1940)\allowbreak{}0 | Deltabar(1940)\allowbreak{}0 | Delta(1940)\allowbreak{}+ | Deltabar(1940)\allowbreak{}- | Delta(1940)\allowbreak{}++ | Deltabar(1940)\allowbreak{}-\kern0pt{}-\par {\fontsize{8.1}{9.4}\selectfont PDG2026:\allowbreak{}observable:\allowbreak{}B136:\allowbreak{}B136W:\allowbreak{}349273} & \textbf{300 to 400 to 500} MeV\par Presentación: 300 to 500 (\textasciitilde{}400) & \(u_+\): NA\par \(u_-\): NA\par Tipo: R\par Confianza: ND & COMP-BARYON\par E04 / M07\\\addlinespace[3pt]
% VI-CATALOG-ROW anchuras 79
\textbf{79. Omega(2250)\allowbreak{}- WIDTH}\par Omega(2250)\allowbreak{}- | Omegabar(2250)\allowbreak{}+\par {\fontsize{8.1}{9.4}\selectfont PDG2026:\allowbreak{}observable:\allowbreak{}B141:\allowbreak{}B141W:\allowbreak{}350517} & \textbf{55+-18} MeV & \(u_+\): 17.69836\allowbreak{}44476396\allowbreak{}51\par \(u_-\): 17.69836\allowbreak{}44476396\allowbreak{}51\par Tipo: NONE\par Confianza: ND & COMP-BARYON\par E04 / M07\\\addlinespace[3pt]
% VI-CATALOG-ROW anchuras 80
\textbf{80. Omega(2470)\allowbreak{}- WIDTH}\par Omega(2470)\allowbreak{}- | Omegabar(2470)\allowbreak{}+\par {\fontsize{8.1}{9.4}\selectfont PDG2026:\allowbreak{}observable:\allowbreak{}B143:\allowbreak{}B143W:\allowbreak{}350524} & \textbf{72+-33} MeV & \(u_+\): 33\par \(u_-\): 33\par Tipo: NONE\par Confianza: ND & COMP-BARYON\par E04 / M07\\\addlinespace[3pt]
% VI-CATALOG-ROW anchuras 81
\textbf{81. N(1900) BREIT-WIGNER WIDTH}\par N(1900)\allowbreak{}0 | Nbar(1900)\allowbreak{}0 | N(1900)\allowbreak{}+ | Nbar(1900)\allowbreak{}-\par {\fontsize{8.1}{9.4}\selectfont PDG2026:\allowbreak{}observable:\allowbreak{}B144:\allowbreak{}B144W:\allowbreak{}348741} & \textbf{100 to 200 to 320} MeV\par Presentación: 100 to 320 (\textasciitilde{}200) & \(u_+\): NA\par \(u_-\): NA\par Tipo: R\par Confianza: ND & COMP-BARYON\par E04 / M07\\\addlinespace[3pt]
% VI-CATALOG-ROW anchuras 82
\textbf{82. Xi\_\allowbreak{}c(2645)\allowbreak{}+ WIDTH}\par Xi\_\allowbreak{}c(2645)\allowbreak{}0 | Xibar\_\allowbreak{}c(2645)\allowbreak{}0 | Xi\_\allowbreak{}c(2645)\allowbreak{}+ | Xibar\_\allowbreak{}c(2645)\allowbreak{}-\par {\fontsize{8.1}{9.4}\selectfont PDG2026:\allowbreak{}observable:\allowbreak{}B146:\allowbreak{}B146W+:\allowbreak{}351213} & \textbf{2.14+-0.19} MeV & \(u_+\): 0.189898\allowbreak{}96013117\allowbreak{}41\par \(u_-\): 0.189898\allowbreak{}96013117\allowbreak{}41\par Tipo: NONE\par Confianza: ND & COMP-BARYON\par E04 / M07\\\addlinespace[3pt]
% VI-CATALOG-ROW anchuras 83
\textbf{83. Xi\_\allowbreak{}c(2645)\allowbreak{}0 WIDTH}\par Xi\_\allowbreak{}c(2645)\allowbreak{}0 | Xibar\_\allowbreak{}c(2645)\allowbreak{}0 | Xi\_\allowbreak{}c(2645)\allowbreak{}+ | Xibar\_\allowbreak{}c(2645)\allowbreak{}-\par {\fontsize{8.1}{9.4}\selectfont PDG2026:\allowbreak{}observable:\allowbreak{}B146:\allowbreak{}B146W0:\allowbreak{}351214} & \textbf{2.35+-0.22} MeV & \(u_+\): 0.222036\allowbreak{}03311174\allowbreak{}519\par \(u_-\): 0.222036\allowbreak{}03311174\allowbreak{}519\par Tipo: NONE\par Confianza: ND & COMP-BARYON\par E04 / M07\\\addlinespace[3pt]
% VI-CATALOG-ROW anchuras 84
\textbf{84. Xi\_\allowbreak{}c(3080)\allowbreak{}+ WIDTH}\par Xi\_\allowbreak{}c(3080)\allowbreak{}0 | Xibar\_\allowbreak{}c(3080)\allowbreak{}0\par {\fontsize{8.1}{9.4}\selectfont PDG2026:\allowbreak{}observable:\allowbreak{}B147:\allowbreak{}B147W+:\allowbreak{}351309} & \textbf{3.8+-0.9} MeV & \(u_+\): 0.939457\allowbreak{}66227499\allowbreak{}938\par \(u_-\): 0.939457\allowbreak{}66227499\allowbreak{}938\par Tipo: NONE\par Confianza: ND & COMP-BARYON\par E04 / M07\\\addlinespace[3pt]
% VI-CATALOG-ROW anchuras 85
\textbf{85. Xi\_\allowbreak{}c(3080)\allowbreak{}0 WIDTH}\par Xi\_\allowbreak{}c(3080)\allowbreak{}0 | Xibar\_\allowbreak{}c(3080)\allowbreak{}0\par {\fontsize{8.1}{9.4}\selectfont PDG2026:\allowbreak{}observable:\allowbreak{}B147:\allowbreak{}B147W0:\allowbreak{}351310} & \textbf{5.6+-2.2} MeV & \(u_+\): 2.179860\allowbreak{}10862691\allowbreak{}2\par \(u_-\): 2.179860\allowbreak{}10862691\allowbreak{}2\par Tipo: NONE\par Confianza: ND & COMP-BARYON\par E04 / M07\\\addlinespace[3pt]
% VI-CATALOG-ROW anchuras 86
\textbf{86. Xi\_\allowbreak{}c(2815)\allowbreak{}+ WIDTH}\par Xi\_\allowbreak{}c(2815)\allowbreak{}0 | Xibar\_\allowbreak{}c(2815)\allowbreak{}0 | Xi\_\allowbreak{}c(2815)\allowbreak{}+ | Xibar\_\allowbreak{}c(2815)\allowbreak{}-\par {\fontsize{8.1}{9.4}\selectfont PDG2026:\allowbreak{}observable:\allowbreak{}B148:\allowbreak{}B148W+:\allowbreak{}351240} & \textbf{2.15+-0.15} MeV & \(u_+\): 0.151932\allowbreak{}72999421\allowbreak{}351\par \(u_-\): 0.151932\allowbreak{}72999421\allowbreak{}351\par Tipo: NONE\par Confianza: ND & COMP-BARYON\par E04 / M07\\\addlinespace[3pt]
% VI-CATALOG-ROW anchuras 87
\textbf{87. Xi\_\allowbreak{}c(2815)\allowbreak{}0 WIDTH}\par Xi\_\allowbreak{}c(2815)\allowbreak{}0 | Xibar\_\allowbreak{}c(2815)\allowbreak{}0 | Xi\_\allowbreak{}c(2815)\allowbreak{}+ | Xibar\_\allowbreak{}c(2815)\allowbreak{}-\par {\fontsize{8.1}{9.4}\selectfont PDG2026:\allowbreak{}observable:\allowbreak{}B148:\allowbreak{}B148W0:\allowbreak{}351241} & \textbf{2.54+-0.25} MeV & \(u_+\): 0.247588\allowbreak{}36806279\allowbreak{}891\par \(u_-\): 0.247588\allowbreak{}36806279\allowbreak{}891\par Tipo: NONE\par Confianza: ND & COMP-BARYON\par E04 / M07\\\addlinespace[3pt]
% VI-CATALOG-ROW anchuras 88
\textbf{88. Xi\_\allowbreak{}c(2790)\allowbreak{}+ WIDTH}\par Xi\_\allowbreak{}c(2790)\allowbreak{}0 | Xibar\_\allowbreak{}c(2790)\allowbreak{}0 | Xi\_\allowbreak{}c(2790)\allowbreak{}+ | Xibar\_\allowbreak{}c(2790)\allowbreak{}-\par {\fontsize{8.1}{9.4}\selectfont PDG2026:\allowbreak{}observable:\allowbreak{}B149:\allowbreak{}B149W+:\allowbreak{}351224} & \textbf{8.9+-1.0} MeV & \(u_+\): 1\par \(u_-\): 1\par Tipo: NONE\par Confianza: ND & COMP-BARYON\par E04 / M07\\\addlinespace[3pt]
% VI-CATALOG-ROW anchuras 89
\textbf{89. Xi\_\allowbreak{}c(2790)\allowbreak{}0 WIDTH}\par Xi\_\allowbreak{}c(2790)\allowbreak{}0 | Xibar\_\allowbreak{}c(2790)\allowbreak{}0 | Xi\_\allowbreak{}c(2790)\allowbreak{}+ | Xibar\_\allowbreak{}c(2790)\allowbreak{}-\par {\fontsize{8.1}{9.4}\selectfont PDG2026:\allowbreak{}observable:\allowbreak{}B149:\allowbreak{}B149W0:\allowbreak{}351225} & \textbf{10.0+-1.1} MeV & \(u_+\): 1.063014\allowbreak{}58127346\allowbreak{}5\par \(u_-\): 1.063014\allowbreak{}58127346\allowbreak{}5\par Tipo: NONE\par Confianza: ND & COMP-BARYON\par E04 / M07\\\addlinespace[3pt]
% VI-CATALOG-ROW anchuras 90
\textbf{90. Lambda\_\allowbreak{}c(2765)\allowbreak{}+ WIDTH}\par Lambda\_\allowbreak{}c(2765)\allowbreak{}0 | Lambdabar\_\allowbreak{}c(2765)\allowbreak{}0\par {\fontsize{8.1}{9.4}\selectfont PDG2026:\allowbreak{}observable:\allowbreak{}B150:\allowbreak{}B150W:\allowbreak{}350954} & \textbf{50} MeV & \(u_+\): 0\par \(u_-\): 0\par Tipo: NONE\par Confianza: ND & COMP-BARYON\par E04 / M07\\\addlinespace[3pt]
% VI-CATALOG-ROW anchuras 91
\textbf{91. Lambda\_\allowbreak{}c(2880)\allowbreak{}+ WIDTH}\par Lambda\_\allowbreak{}c(2880)\allowbreak{}+ | Lambdabar\_\allowbreak{}c(2880)\allowbreak{}-\par {\fontsize{8.1}{9.4}\selectfont PDG2026:\allowbreak{}observable:\allowbreak{}B151:\allowbreak{}B151W:\allowbreak{}350964} & \textbf{5.6+0.8-0.6} MeV & \(u_+\): 0.772000\allowbreak{}97671302\allowbreak{}059\par \(u_-\): 0.622892\allowbreak{}22418902\allowbreak{}057\par Tipo: NONE\par Confianza: ND & COMP-BARYON\par E04 / M07\\\addlinespace[3pt]
% VI-CATALOG-ROW anchuras 92
\textbf{92. Sigma\_\allowbreak{}c(2800)\allowbreak{}++ WIDTH}\par Sigma\_\allowbreak{}c(2800)\allowbreak{}0 | Sigmabar\_\allowbreak{}c(2800)\allowbreak{}0\par {\fontsize{8.1}{9.4}\selectfont PDG2026:\allowbreak{}observable:\allowbreak{}B155:\allowbreak{}B155W++:\allowbreak{}351025} & \textbf{75+22-17} MeV & \(u_+\): 21.63330\allowbreak{}76527839\allowbreak{}41\par \(u_-\): 17.02938\allowbreak{}63659263\allowbreak{}99\par Tipo: NONE\par Confianza: ND & COMP-BARYON\par E04 / M07\\\addlinespace[3pt]
% VI-CATALOG-ROW anchuras 93
\textbf{93. Sigma\_\allowbreak{}c(2800)\allowbreak{}+ WIDTH}\par Sigma\_\allowbreak{}c(2800)\allowbreak{}0 | Sigmabar\_\allowbreak{}c(2800)\allowbreak{}0\par {\fontsize{8.1}{9.4}\selectfont PDG2026:\allowbreak{}observable:\allowbreak{}B155:\allowbreak{}B155W+:\allowbreak{}351026} & \textbf{60+60-40} MeV & \(u_+\): 63.82005\allowbreak{}95424354\allowbreak{}02\par \(u_-\): 44.41846\allowbreak{}46290256\allowbreak{}21\par Tipo: NONE\par Confianza: ND & COMP-BARYON\par E04 / M07\\\addlinespace[3pt]
% VI-CATALOG-ROW anchuras 94
\textbf{94. Sigma\_\allowbreak{}c(2800)\allowbreak{}0 WIDTH}\par Sigma\_\allowbreak{}c(2800)\allowbreak{}0 | Sigmabar\_\allowbreak{}c(2800)\allowbreak{}0\par {\fontsize{8.1}{9.4}\selectfont PDG2026:\allowbreak{}observable:\allowbreak{}B155:\allowbreak{}B155W0:\allowbreak{}351027} & \textbf{72+22-15} MeV & \(u_+\): 22.09356\allowbreak{}07974680\allowbreak{}29\par \(u_-\): 14.82346\allowbreak{}07851214\allowbreak{}41\par Tipo: NONE\par Confianza: ND & COMP-BARYON\par E04 / M07\\\addlinespace[3pt]
% VI-CATALOG-ROW anchuras 95
\textbf{95. Xi\_\allowbreak{}c(2930)\allowbreak{}+ WIDTH}\par Xi\_\allowbreak{}c(2930)\allowbreak{}0 | Xibar\_\allowbreak{}c(2930)\allowbreak{}0\par {\fontsize{8.1}{9.4}\selectfont PDG2026:\allowbreak{}observable:\allowbreak{}B156:\allowbreak{}B156A05:\allowbreak{}351265} & \textbf{15+-9} MeV & \(u_+\): 9.148223\allowbreak{}87133152\allowbreak{}76\par \(u_-\): 9.148223\allowbreak{}87133152\allowbreak{}76\par Tipo: NONE\par Confianza: ND & COMP-BARYON\par E04 / M07\\\addlinespace[3pt]
% VI-CATALOG-ROW anchuras 96
\textbf{96. Xi\_\allowbreak{}c(2930)\allowbreak{}0 WIDTH}\par Xi\_\allowbreak{}c(2930)\allowbreak{}0 | Xibar\_\allowbreak{}c(2930)\allowbreak{}0\par {\fontsize{8.1}{9.4}\selectfont PDG2026:\allowbreak{}observable:\allowbreak{}B156:\allowbreak{}B156A06:\allowbreak{}351266} & \textbf{10.2+-1.4} MeV & \(u_+\): 1.360147\allowbreak{}05087354\allowbreak{}4\par \(u_-\): 1.360147\allowbreak{}05087354\allowbreak{}4\par Tipo: NONE\par Confianza: ND & COMP-BARYON\par E04 / M07\\\addlinespace[3pt]
% VI-CATALOG-ROW anchuras 97
\textbf{97. Xi\_\allowbreak{}c(3055)\allowbreak{}+ WIDTH}\par Xi\_\allowbreak{}c(3055)\allowbreak{}0 | Xibar\_\allowbreak{}c(3055)\allowbreak{}0\par {\fontsize{8.1}{9.4}\selectfont PDG2026:\allowbreak{}observable:\allowbreak{}B157:\allowbreak{}B157W+:\allowbreak{}351299} & \textbf{8.0+-0.8} MeV & \(u_+\): 0.763918\allowbreak{}12853200\allowbreak{}158\par \(u_-\): 0.763918\allowbreak{}12853200\allowbreak{}158\par Tipo: NONE\par Confianza: ND & COMP-BARYON\par E04 / M07\\\addlinespace[3pt]
% VI-CATALOG-ROW anchuras 98
\textbf{98. Xi\_\allowbreak{}c(3055)\allowbreak{}0 WIDTH}\par Xi\_\allowbreak{}c(3055)\allowbreak{}0 | Xibar\_\allowbreak{}c(3055)\allowbreak{}0\par {\fontsize{8.1}{9.4}\selectfont PDG2026:\allowbreak{}observable:\allowbreak{}B157:\allowbreak{}B157W:\allowbreak{}351300} & \textbf{12.4+-2.3} MeV & \(u_+\): 2.282542\allowbreak{}44210266\allowbreak{}5\par \(u_-\): 2.282542\allowbreak{}44210266\allowbreak{}5\par Tipo: NONE\par Confianza: ND & COMP-BARYON\par E04 / M07\\\addlinespace[3pt]
% VI-CATALOG-ROW anchuras 99
\textbf{99. Xi\_\allowbreak{}c(3123)\allowbreak{}+ WIDTH}\par Xi\_\allowbreak{}c(3123)\allowbreak{}0 | Xibar\_\allowbreak{}c(3123)\allowbreak{}0\par {\fontsize{8.1}{9.4}\selectfont PDG2026:\allowbreak{}observable:\allowbreak{}B158:\allowbreak{}B158W+:\allowbreak{}351326} & \textbf{4+-4} MeV & \(u_+\): 3.801315\allowbreak{}56174964\allowbreak{}3\par \(u_-\): 3.801315\allowbreak{}56174964\allowbreak{}3\par Tipo: NONE\par Confianza: ND & COMP-BARYON\par E04 / M07\\\addlinespace[3pt]
% VI-CATALOG-ROW anchuras 100
\textbf{100. Xi\_\allowbreak{}b(5945)\allowbreak{}0 WIDTH}\par Xi\_\allowbreak{}b(5945)\allowbreak{}0 | Xibar\_\allowbreak{}b(5945)\allowbreak{}0\par {\fontsize{8.1}{9.4}\selectfont PDG2026:\allowbreak{}observable:\allowbreak{}B161:\allowbreak{}B161W:\allowbreak{}351845} & \textbf{0.87+-0.07} MeV & \(u_+\): 0.075069\allowbreak{}67217458\allowbreak{}3921\par \(u_-\): 0.074706\allowbreak{}30547753\allowbreak{}7885\par Tipo: NONE\par Confianza: ND & COMP-BARYON\par E04 / M07\\\addlinespace[3pt]
% VI-CATALOG-ROW anchuras 101
\textbf{101. Lambda\_\allowbreak{}b(5912)\allowbreak{}0 WIDTH}\par Lambda\_\allowbreak{}b(5912)\allowbreak{}0 | Lambdabar\_\allowbreak{}b(5912)\allowbreak{}0\par {\fontsize{8.1}{9.4}\selectfont PDG2026:\allowbreak{}observable:\allowbreak{}B162:\allowbreak{}B162W:\allowbreak{}351695} & \textbf{<0.25} MeV & \(u_+\): NA\par \(u_-\): NA\par Tipo: U\par Confianza: 90 & COMP-BARYON\par E04 / M07\\\addlinespace[3pt]
% VI-CATALOG-ROW anchuras 102
\textbf{102. Lambda\_\allowbreak{}b(5920)\allowbreak{}0 WIDTH}\par Lambda\_\allowbreak{}b(5920)\allowbreak{}0 | Lambdabar\_\allowbreak{}b(5920)\allowbreak{}0\par {\fontsize{8.1}{9.4}\selectfont PDG2026:\allowbreak{}observable:\allowbreak{}B163:\allowbreak{}B163W:\allowbreak{}351699} & \textbf{<0.19} MeV & \(u_+\): NA\par \(u_-\): NA\par Tipo: U\par Confianza: 90 & COMP-BARYON\par E04 / M07\\\addlinespace[3pt]
% VI-CATALOG-ROW anchuras 103
\textbf{103. Lambda(1710) WIDTH}\par Lambda(1710)\allowbreak{}0 | Lambdabar(1710)\allowbreak{}0\par {\fontsize{8.1}{9.4}\selectfont PDG2026:\allowbreak{}observable:\allowbreak{}B164:\allowbreak{}B164W:\allowbreak{}349533} & \textbf{180+-40} MeV & \(u_+\): 42\par \(u_-\): 42\par Tipo: NONE\par Confianza: ND & COMP-BARYON\par E04 / M07\\\addlinespace[3pt]
% VI-CATALOG-ROW anchuras 104
\textbf{104. Lambda(2050) WIDTH}\par Lambda(2050)\allowbreak{}0 | Lambdabar(2050)\allowbreak{}0\par {\fontsize{8.1}{9.4}\selectfont PDG2026:\allowbreak{}observable:\allowbreak{}B165:\allowbreak{}B165W:\allowbreak{}349652} & \textbf{490+-60} MeV & \(u_+\): 61\par \(u_-\): 61\par Tipo: NONE\par Confianza: ND & COMP-BARYON\par E04 / M07\\\addlinespace[3pt]
% VI-CATALOG-ROW anchuras 105
\textbf{105. Sigma(1780) WIDTH}\par Sigma(1780)\allowbreak{}- | Sigmabar(1780)\allowbreak{}+ | Sigma(1780)\allowbreak{}0 | Sigmabar(1780)\allowbreak{}0 | Sigma(1780)\allowbreak{}+ | Sigmabar(1780)\allowbreak{}-\par {\fontsize{8.1}{9.4}\selectfont PDG2026:\allowbreak{}observable:\allowbreak{}B166:\allowbreak{}B166W:\allowbreak{}349992} & \textbf{100 to 200 to 300} MeV\par Presentación: 100 to 300 (\textasciitilde{}200) & \(u_+\): NA\par \(u_-\): NA\par Tipo: R\par Confianza: ND & COMP-BARYON\par E04 / M07\\\addlinespace[3pt]
% VI-CATALOG-ROW anchuras 106
\textbf{106. Sigma(1900) WIDTH}\par Sigma(1900)\allowbreak{}- | Sigmabar(1900)\allowbreak{}+ | Sigma(1900)\allowbreak{}0 | Sigmabar(1900)\allowbreak{}0 | Sigma(1900)\allowbreak{}+ | Sigmabar(1900)\allowbreak{}-\par {\fontsize{8.1}{9.4}\selectfont PDG2026:\allowbreak{}observable:\allowbreak{}B167:\allowbreak{}B167W:\allowbreak{}350020} & \textbf{140 to 165 to 190} MeV\par Presentación: 140 to 190 (\textasciitilde{}165) & \(u_+\): NA\par \(u_-\): NA\par Tipo: R\par Confianza: ND & COMP-BARYON\par E04 / M07\\\addlinespace[3pt]
% VI-CATALOG-ROW anchuras 107
\textbf{107. Sigma(1940) WIDTH}\par Sigma(1940)\allowbreak{}- | Sigmabar(1940)\allowbreak{}+ | Sigma(1940)\allowbreak{}0 | Sigmabar(1940)\allowbreak{}0 | Sigma(1940)\allowbreak{}+ | Sigmabar(1940)\allowbreak{}-\par {\fontsize{8.1}{9.4}\selectfont PDG2026:\allowbreak{}observable:\allowbreak{}B168:\allowbreak{}B168W:\allowbreak{}350103} & \textbf{100 to 250 to 400} MeV\par Presentación: 100 to 400 (\textasciitilde{}250) & \(u_+\): NA\par \(u_-\): NA\par Tipo: R\par Confianza: ND & COMP-BARYON\par E04 / M07\\\addlinespace[3pt]
% VI-CATALOG-ROW anchuras 108
\textbf{108. Xi\_\allowbreak{}b'(5935)\allowbreak{}- WIDTH}\par Xi\_\allowbreak{}b\textasciicircum{}'(5935)\allowbreak{}- | Xibar\_\allowbreak{}b\textasciicircum{}'(5935)\allowbreak{}+\par {\fontsize{8.1}{9.4}\selectfont PDG2026:\allowbreak{}observable:\allowbreak{}B169:\allowbreak{}B169W:\allowbreak{}351841} & \textbf{0.03+-0.032} MeV & \(u_+\): 0.031622\allowbreak{}77660168\allowbreak{}3791\par \(u_-\): 0.031622\allowbreak{}77660168\allowbreak{}3791\par Tipo: NONE\par Confianza: ND & COMP-BARYON\par E04 / M07\\\addlinespace[3pt]
% VI-CATALOG-ROW anchuras 109
\textbf{109. Xi\_\allowbreak{}b(5955)\allowbreak{}- WIDTH}\par Xi\_\allowbreak{}b(5955)\allowbreak{}- | Xibar\_\allowbreak{}b(5955)\allowbreak{}+\par {\fontsize{8.1}{9.4}\selectfont PDG2026:\allowbreak{}observable:\allowbreak{}B170:\allowbreak{}B170W:\allowbreak{}351849} & \textbf{1.43+-0.11} MeV & \(u_+\): 0.113137\allowbreak{}08498984\allowbreak{}761\par \(u_-\): 0.113137\allowbreak{}08498984\allowbreak{}761\par Tipo: NONE\par Confianza: ND & COMP-BARYON\par E04 / M07\\\addlinespace[3pt]
% VI-CATALOG-ROW anchuras 110
\textbf{110. P\_\allowbreak{}c\_\allowbreak{}cbar(4380)\allowbreak{}+ WIDTH}\par P\_\allowbreak{}c\_\allowbreak{}cbar(4380)\allowbreak{}+\par {\fontsize{8.1}{9.4}\selectfont PDG2026:\allowbreak{}observable:\allowbreak{}B171:\allowbreak{}B171W:\allowbreak{}351952} & \textbf{210+-90} MeV & \(u_+\): 87.86353\allowbreak{}05459551\allowbreak{}86\par \(u_-\): 87.86353\allowbreak{}05459551\allowbreak{}86\par Tipo: NONE\par Confianza: ND & COMP-BARYON\par E04 / M07\\\addlinespace[3pt]
% VI-CATALOG-ROW anchuras 111
\textbf{111. P\_\allowbreak{}c\_\allowbreak{}cbar(4457)\allowbreak{}+ WIDTH}\par P\_\allowbreak{}c\_\allowbreak{}cbar(4457)\allowbreak{}+\par {\fontsize{8.1}{9.4}\selectfont PDG2026:\allowbreak{}observable:\allowbreak{}B172:\allowbreak{}B172W:\allowbreak{}351978} & \textbf{6.4+6.0-2.8} MeV & \(u_+\): 6.040695\allowbreak{}32421558\allowbreak{}28\par \(u_-\): 2.758622\allowbreak{}84482674\allowbreak{}4\par Tipo: NONE\par Confianza: ND & COMP-BARYON\par E04 / M07\\\addlinespace[3pt]
% VI-CATALOG-ROW anchuras 112
\textbf{112. Omega\_\allowbreak{}c(3000)\allowbreak{}0 WIDTH}\par Omega\_\allowbreak{}c(3000)\allowbreak{}0 | Omegabar\_\allowbreak{}c(3000)\allowbreak{}0\par {\fontsize{8.1}{9.4}\selectfont PDG2026:\allowbreak{}observable:\allowbreak{}B173:\allowbreak{}B173W:\allowbreak{}351377} & \textbf{3.8+1.6-0.4} MeV & \(u_+\): 1.606549\allowbreak{}09666651\allowbreak{}7\par \(u_-\): 0.370135\allowbreak{}11046643\allowbreak{}488\par Tipo: NONE\par Confianza: ND & COMP-BARYON\par E04 / M07\\\addlinespace[3pt]
% VI-CATALOG-ROW anchuras 113
\textbf{113. Omega\_\allowbreak{}c(3050)\allowbreak{}0 WIDTH}\par Omega\_\allowbreak{}c(3050)\allowbreak{}0 | Omegabar\_\allowbreak{}c(3050)\allowbreak{}0\par {\fontsize{8.1}{9.4}\selectfont PDG2026:\allowbreak{}observable:\allowbreak{}B174:\allowbreak{}B174W:\allowbreak{}351381} & \textbf{< 1.8} MeV\par Presentación: <1.8 & \(u_+\): NA\par \(u_-\): NA\par Tipo: U\par Confianza: 95 & COMP-BARYON\par E04 / M07\\\addlinespace[3pt]
% VI-CATALOG-ROW anchuras 114
\textbf{114. Omega\_\allowbreak{}c(3065)\allowbreak{}0 WIDTH}\par Omega\_\allowbreak{}c(3065)\allowbreak{}0 | Omegabar\_\allowbreak{}c(3065)\allowbreak{}0\par {\fontsize{8.1}{9.4}\selectfont PDG2026:\allowbreak{}observable:\allowbreak{}B175:\allowbreak{}B175W:\allowbreak{}351385} & \textbf{3.4+0.7-0.8} MeV & \(u_+\): 0.684966\allowbreak{}37416672\allowbreak{}579\par \(u_-\): 0.793858\allowbreak{}52123859\allowbreak{}457\par Tipo: NONE\par Confianza: ND & COMP-BARYON\par E04 / M07\\\addlinespace[3pt]
% VI-CATALOG-ROW anchuras 115
\textbf{115. Omega\_\allowbreak{}c(3090)\allowbreak{}0 WIDTH}\par Omega\_\allowbreak{}c(3090)\allowbreak{}0 | Omegabar\_\allowbreak{}c(3090)\allowbreak{}0\par {\fontsize{8.1}{9.4}\selectfont PDG2026:\allowbreak{}observable:\allowbreak{}B176:\allowbreak{}B176W:\allowbreak{}351389} & \textbf{8.5+0.8-1.7} MeV & \(u_+\): 0.752130\allowbreak{}30785895\allowbreak{}074\par \(u_-\): 1.678689\allowbreak{}96541946\allowbreak{}4\par Tipo: NONE\par Confianza: ND & COMP-BARYON\par E04 / M07\\\addlinespace[3pt]
% VI-CATALOG-ROW anchuras 116
\textbf{116. Omega\_\allowbreak{}c(3120)\allowbreak{}0 WIDTH}\par Omega\_\allowbreak{}c(3120)\allowbreak{}0 | Omegabar\_\allowbreak{}c(3120)\allowbreak{}0\par {\fontsize{8.1}{9.4}\selectfont PDG2026:\allowbreak{}observable:\allowbreak{}B177:\allowbreak{}B177W:\allowbreak{}351393} & \textbf{<2.5} MeV & \(u_+\): NA\par \(u_-\): NA\par Tipo: U\par Confianza: 95 & COMP-BARYON\par E04 / M07\\\addlinespace[3pt]
% VI-CATALOG-ROW anchuras 117
\textbf{117. Lambda\_\allowbreak{}c(2860)\allowbreak{}+ WIDTH}\par Lambda\_\allowbreak{}c(2860)\allowbreak{}0 | Lambdabar\_\allowbreak{}c(2860)\allowbreak{}0\par {\fontsize{8.1}{9.4}\selectfont PDG2026:\allowbreak{}observable:\allowbreak{}B178:\allowbreak{}B178W:\allowbreak{}350957} & \textbf{68+12-22} MeV & \(u_+\): 11.78049\allowbreak{}23496431\allowbreak{}2\par \(u_-\): 21.62336\allowbreak{}69903648\allowbreak{}4\par Tipo: NONE\par Confianza: ND & COMP-BARYON\par E04 / M07\\\addlinespace[3pt]
% VI-CATALOG-ROW anchuras 118
\textbf{118. Omega(2012)\allowbreak{}- WIDTH}\par Omega(2012)\allowbreak{}- | Omegabar(2012)\allowbreak{}+\par {\fontsize{8.1}{9.4}\selectfont PDG2026:\allowbreak{}observable:\allowbreak{}B179:\allowbreak{}B179W:\allowbreak{}350507} & \textbf{6.3+-2.0} MeV & \(u_+\): 2.074289\allowbreak{}54575496\allowbreak{}5\par \(u_-\): 1.919724\allowbreak{}98587065\allowbreak{}51\par Tipo: NONE\par Confianza: ND & COMP-BARYON\par E04 / M07\\\addlinespace[3pt]
% VI-CATALOG-ROW anchuras 119
\textbf{119. Xi\_\allowbreak{}b(6227)\allowbreak{}- WIDTH}\par Xi\_\allowbreak{}b(6227)\allowbreak{}- | Xi\_\allowbreak{}b(6227)\allowbreak{}+\par {\fontsize{8.1}{9.4}\selectfont PDG2026:\allowbreak{}observable:\allowbreak{}B180:\allowbreak{}B180W:\allowbreak{}351865} & \textbf{19.9+-2.6} MeV & \(u_+\): 2.580697\allowbreak{}58011278\allowbreak{}78\par \(u_-\): 2.580697\allowbreak{}58011278\allowbreak{}78\par Tipo: NONE\par Confianza: ND & COMP-BARYON\par E04 / M07\\\addlinespace[3pt]
% VI-CATALOG-ROW anchuras 120
\textbf{120. Sigma\_\allowbreak{}b(6097)\allowbreak{}+ WIDTH}\par Sigma\_\allowbreak{}b(6097)\allowbreak{}+ | Sigmabar\_\allowbreak{}b(6097)\allowbreak{}-\par {\fontsize{8.1}{9.4}\selectfont PDG2026:\allowbreak{}observable:\allowbreak{}B181:\allowbreak{}B181W:\allowbreak{}351735} & \textbf{31+-6} MeV & \(u_+\): 5.544366\allowbreak{}51025164\allowbreak{}46\par \(u_-\): 5.544366\allowbreak{}51025164\allowbreak{}46\par Tipo: NONE\par Confianza: ND & COMP-BARYON\par E04 / M07\\\addlinespace[3pt]
% VI-CATALOG-ROW anchuras 121
\textbf{121. Sigma\_\allowbreak{}b(6097)\allowbreak{}- WIDTH}\par Sigma\_\allowbreak{}b(6097)\allowbreak{}- | Sigmabar\_\allowbreak{}b(6097)\allowbreak{}+\par {\fontsize{8.1}{9.4}\selectfont PDG2026:\allowbreak{}observable:\allowbreak{}B182:\allowbreak{}B182W:\allowbreak{}351739} & \textbf{29+-4} MeV & \(u_+\): 4.295346\allowbreak{}31898290\allowbreak{}62\par \(u_-\): 4.295346\allowbreak{}31898290\allowbreak{}62\par Tipo: NONE\par Confianza: ND & COMP-BARYON\par E04 / M07\\\addlinespace[3pt]
% VI-CATALOG-ROW anchuras 122
\textbf{122. Sigma(2110) WIDTH}\par Sigma(2110)\allowbreak{}- | Sigmabar(2110)\allowbreak{}+ | Sigma(2110)\allowbreak{}0 | Sigmabar(2110)\allowbreak{}0 | Sigma(2110)\allowbreak{}+ | Sigmabar(2110)\allowbreak{}-\par {\fontsize{8.1}{9.4}\selectfont PDG2026:\allowbreak{}observable:\allowbreak{}B183:\allowbreak{}B183W:\allowbreak{}350229} & \textbf{310+120-50} MeV & \(u_+\): 122.3496\allowbreak{}30121267\allowbreak{}89\par \(u_-\): 54.76107\allowbreak{}45236418\allowbreak{}52\par Tipo: NONE\par Confianza: ND & COMP-BARYON\par E04 / M07\\\addlinespace[3pt]
% VI-CATALOG-ROW anchuras 123
\textbf{123. Sigma(2230) WIDTH}\par Sigma(2230)\allowbreak{}- | Sigmabar(2230)\allowbreak{}+ | Sigma(2230)\allowbreak{}0 | Sigmabar(2230)\allowbreak{}0 | Sigma(2230)\allowbreak{}+ | Sigmabar(2230)\allowbreak{}-\par {\fontsize{8.1}{9.4}\selectfont PDG2026:\allowbreak{}observable:\allowbreak{}B184:\allowbreak{}B184W:\allowbreak{}350254} & \textbf{350+-50} MeV & \(u_+\): 50\par \(u_-\): 50\par Tipo: NONE\par Confianza: ND & COMP-BARYON\par E04 / M07\\\addlinespace[3pt]
% VI-CATALOG-ROW anchuras 124
\textbf{124. P\_\allowbreak{}c\_\allowbreak{}cbar(4312)\allowbreak{}+ WIDTH}\par P\_\allowbreak{}c\_\allowbreak{}cbar(4312)\allowbreak{}+\par {\fontsize{8.1}{9.4}\selectfont PDG2026:\allowbreak{}observable:\allowbreak{}B185:\allowbreak{}B185W:\allowbreak{}351936} & \textbf{10+-5} MeV & \(u_+\): 4.580392\allowbreak{}99623951\allowbreak{}47\par \(u_-\): 5.247856\allowbreak{}70536077\allowbreak{}03\par Tipo: NONE\par Confianza: ND & COMP-BARYON\par E04 / M07\\\addlinespace[3pt]
% VI-CATALOG-ROW anchuras 125
\textbf{125. P\_\allowbreak{}c\_\allowbreak{}cbar(4440)\allowbreak{}+ WIDTH}\par P\_\allowbreak{}c\_\allowbreak{}cbar(4440)\allowbreak{}+\par {\fontsize{8.1}{9.4}\selectfont PDG2026:\allowbreak{}observable:\allowbreak{}B186:\allowbreak{}B186W:\allowbreak{}351956} & \textbf{21+10-11} MeV & \(u_+\): 9.984988\allowbreak{}73309329\allowbreak{}23\par \(u_-\): 11.22586\allowbreak{}29957789\allowbreak{}89\par Tipo: NONE\par Confianza: ND & COMP-BARYON\par E04 / M07\\\addlinespace[3pt]
% VI-CATALOG-ROW anchuras 126
\textbf{126. Lambda\_\allowbreak{}b(6146)\allowbreak{}0 WIDTH}\par Lambda\_\allowbreak{}b(6146)\allowbreak{}0 | Lambdabar\_\allowbreak{}b(6146)\allowbreak{}0\par {\fontsize{8.1}{9.4}\selectfont PDG2026:\allowbreak{}observable:\allowbreak{}B187:\allowbreak{}B187W:\allowbreak{}351708} & \textbf{2.9+-1.3} MeV & \(u_+\): 1.334166\allowbreak{}40641263\allowbreak{}4\par \(u_-\): 1.334166\allowbreak{}40641263\allowbreak{}4\par Tipo: NONE\par Confianza: ND & COMP-BARYON\par E04 / M07\\\addlinespace[3pt]
% VI-CATALOG-ROW anchuras 127
\textbf{127. Lambda\_\allowbreak{}b(6152)\allowbreak{}0 WIDTH}\par Lambda\_\allowbreak{}b(6152)\allowbreak{}0 | Lambdabar\_\allowbreak{}b(6152)\allowbreak{}0\par {\fontsize{8.1}{9.4}\selectfont PDG2026:\allowbreak{}observable:\allowbreak{}B188:\allowbreak{}B188W:\allowbreak{}351714} & \textbf{2.1+-0.9} MeV & \(u_+\): 0.854400\allowbreak{}37453175\allowbreak{}31\par \(u_-\): 0.854400\allowbreak{}37453175\allowbreak{}31\par Tipo: NONE\par Confianza: ND & COMP-BARYON\par E04 / M07\\\addlinespace[3pt]
% VI-CATALOG-ROW anchuras 128
\textbf{128. Lambda(2070) WIDTH}\par Lambda(2070)\allowbreak{}0 | Lambdabar(2070)\allowbreak{}0\par {\fontsize{8.1}{9.4}\selectfont PDG2026:\allowbreak{}observable:\allowbreak{}B190:\allowbreak{}B190W:\allowbreak{}349677} & \textbf{370+-50} MeV & \(u_+\): 50\par \(u_-\): 50\par Tipo: NONE\par Confianza: ND & COMP-BARYON\par E04 / M07\\\addlinespace[3pt]
% VI-CATALOG-ROW anchuras 129
\textbf{129. Lambda(2080) WIDTH}\par Lambda(2080)\allowbreak{}0 | Lambdabar(2080)\allowbreak{}0\par {\fontsize{8.1}{9.4}\selectfont PDG2026:\allowbreak{}observable:\allowbreak{}B191:\allowbreak{}B191W:\allowbreak{}349712} & \textbf{181+-29} MeV & \(u_+\): 29\par \(u_-\): 29\par Tipo: NONE\par Confianza: ND & COMP-BARYON\par E04 / M07\\\addlinespace[3pt]
% VI-CATALOG-ROW anchuras 130
\textbf{130. Xi\_\allowbreak{}c(2923)\allowbreak{}0 WIDTH}\par Xi\_\allowbreak{}c(2923)\allowbreak{}0 | Xibar\_\allowbreak{}c(2923)\allowbreak{}0\par {\fontsize{8.1}{9.4}\selectfont PDG2026:\allowbreak{}observable:\allowbreak{}B192:\allowbreak{}B192W0:\allowbreak{}351257} & \textbf{5.8+-1.3} MeV & \(u_+\): 1.307967\allowbreak{}16242523\par \(u_-\): 1.307967\allowbreak{}16242523\par Tipo: NONE\par Confianza: ND & COMP-BARYON\par E04 / M07\\\addlinespace[3pt]
% VI-CATALOG-ROW anchuras 131
\textbf{131. Xi\_\allowbreak{}c(2923)\allowbreak{}+ WIDTH}\par Xi\_\allowbreak{}c(2923)\allowbreak{}0 | Xibar\_\allowbreak{}c(2923)\allowbreak{}0\par {\fontsize{8.1}{9.4}\selectfont PDG2026:\allowbreak{}observable:\allowbreak{}B192:\allowbreak{}B192WP:\allowbreak{}351258} & \textbf{5.3+-1.7} MeV & \(u_+\): 1.664331\allowbreak{}69770932\allowbreak{}41\par \(u_-\): 1.664331\allowbreak{}69770932\allowbreak{}41\par Tipo: NONE\par Confianza: ND & COMP-BARYON\par E04 / M07\\\addlinespace[3pt]
% VI-CATALOG-ROW anchuras 132
\textbf{132. Omega\_\allowbreak{}b(6316)\allowbreak{}- WIDTH}\par Omega\_\allowbreak{}b(6316)\allowbreak{}- | Omegabar\_\allowbreak{}b(6316)\allowbreak{}+\par {\fontsize{8.1}{9.4}\selectfont PDG2026:\allowbreak{}observable:\allowbreak{}B193:\allowbreak{}B193W:\allowbreak{}351906} & \textbf{<4.2} MeV & \(u_+\): NA\par \(u_-\): NA\par Tipo: U\par Confianza: 95 & COMP-BARYON\par E04 / M07\\\addlinespace[3pt]
% VI-CATALOG-ROW anchuras 133
\textbf{133. Omega\_\allowbreak{}b(6330)\allowbreak{}- WIDTH}\par Omega\_\allowbreak{}b(6330)\allowbreak{}- | Omegabar\_\allowbreak{}b(6330)\allowbreak{}+\par {\fontsize{8.1}{9.4}\selectfont PDG2026:\allowbreak{}observable:\allowbreak{}B194:\allowbreak{}B194W:\allowbreak{}351910} & \textbf{<4.7} MeV & \(u_+\): NA\par \(u_-\): NA\par Tipo: U\par Confianza: 95 & COMP-BARYON\par E04 / M07\\\addlinespace[3pt]
% VI-CATALOG-ROW anchuras 134
\textbf{134. Omega\_\allowbreak{}b(6340)\allowbreak{}- WIDTH}\par Omega\_\allowbreak{}b(6340)\allowbreak{}- | Omegabar\_\allowbreak{}b(6340)\allowbreak{}+\par {\fontsize{8.1}{9.4}\selectfont PDG2026:\allowbreak{}observable:\allowbreak{}B195:\allowbreak{}B195W:\allowbreak{}351914} & \textbf{<1.8} MeV & \(u_+\): NA\par \(u_-\): NA\par Tipo: U\par Confianza: 95 & COMP-BARYON\par E04 / M07\\\addlinespace[3pt]
% VI-CATALOG-ROW anchuras 135
\textbf{135. Omega\_\allowbreak{}b(6350)\allowbreak{}- WIDTH}\par Omega\_\allowbreak{}b(6350)\allowbreak{}- | Omegabar\_\allowbreak{}b(6350)\allowbreak{}+\par {\fontsize{8.1}{9.4}\selectfont PDG2026:\allowbreak{}observable:\allowbreak{}B196:\allowbreak{}B196W:\allowbreak{}351918} & \textbf{<3.2} MeV & \(u_+\): NA\par \(u_-\): NA\par Tipo: U\par Confianza: 95 & COMP-BARYON\par E04 / M07\\\addlinespace[3pt]
% VI-CATALOG-ROW anchuras 136
\textbf{136. Lambda\_\allowbreak{}b(6070)\allowbreak{}0 WIDTH}\par Lambda\_\allowbreak{}b(6070)\allowbreak{}0 | Lambdabar\_\allowbreak{}b(6070)\allowbreak{}0\par {\fontsize{8.1}{9.4}\selectfont PDG2026:\allowbreak{}observable:\allowbreak{}B197:\allowbreak{}B197W:\allowbreak{}351703} & \textbf{72+-11} MeV & \(u_+\): 11.18033\allowbreak{}98874989\allowbreak{}51\par \(u_-\): 11.18033\allowbreak{}98874989\allowbreak{}51\par Tipo: NONE\par Confianza: ND & COMP-BARYON\par E04 / M07\\\addlinespace[3pt]
% VI-CATALOG-ROW anchuras 137
\textbf{137. Xi\_\allowbreak{}b(6227)\allowbreak{}0 WIDTH}\par Xi\_\allowbreak{}b(6227)\allowbreak{}0 | Xibar\_\allowbreak{}b(6227)\allowbreak{}0\par {\fontsize{8.1}{9.4}\selectfont PDG2026:\allowbreak{}observable:\allowbreak{}B199:\allowbreak{}B199W:\allowbreak{}351871} & \textbf{19+5-4} MeV & \(u_+\): 5.192301\allowbreak{}99429886\allowbreak{}77\par \(u_-\): 4.332435\allowbreak{}80448689\allowbreak{}34\par Tipo: NONE\par Confianza: ND & COMP-BARYON\par E04 / M07\\\addlinespace[3pt]
% VI-CATALOG-ROW anchuras 138
\textbf{138. Xi\_\allowbreak{}b(6100)\allowbreak{}- WIDTH}\par Xi\_\allowbreak{}b(6100)\allowbreak{}- | Xibar\_\allowbreak{}b(6100)\allowbreak{}+\par {\fontsize{8.1}{9.4}\selectfont PDG2026:\allowbreak{}observable:\allowbreak{}B200:\allowbreak{}B200W:\allowbreak{}351861} & \textbf{0.94+-0.31} MeV & \(u_+\): 0.310483\allowbreak{}49392520\allowbreak{}053\par \(u_-\): 0.310483\allowbreak{}49392520\allowbreak{}053\par Tipo: NONE\par Confianza: ND & COMP-BARYON\par E04 / M07\\\addlinespace[3pt]
% VI-CATALOG-ROW anchuras 139
\textbf{139. Xi\_\allowbreak{}b(6327)\allowbreak{}0 WIDTH}\par Xi\_\allowbreak{}b(6327)\allowbreak{}0 | Xibar\_\allowbreak{}b(6327)\allowbreak{}0\par {\fontsize{8.1}{9.4}\selectfont PDG2026:\allowbreak{}observable:\allowbreak{}B201:\allowbreak{}B201W:\allowbreak{}351875} & \textbf{< 2.56} MeV\par Presentación: <2.56 & \(u_+\): NA\par \(u_-\): NA\par Tipo: U\par Confianza: 95 & COMP-BARYON\par E04 / M07\\\addlinespace[3pt]
% VI-CATALOG-ROW anchuras 140
\textbf{140. Xi\_\allowbreak{}b(6333)\allowbreak{}0 WIDTH}\par Xi\_\allowbreak{}b(6333)\allowbreak{}0 | Xibar\_\allowbreak{}b(6333)\allowbreak{}0\par {\fontsize{8.1}{9.4}\selectfont PDG2026:\allowbreak{}observable:\allowbreak{}B202:\allowbreak{}B202W:\allowbreak{}351879} & \textbf{< 1.92} MeV\par Presentación: <1.92 & \(u_+\): NA\par \(u_-\): NA\par Tipo: U\par Confianza: 95 & COMP-BARYON\par E04 / M07\\\addlinespace[3pt]
% VI-CATALOG-ROW anchuras 141
\textbf{141. Xi\_\allowbreak{}c(2882) WIDTH}\par Xi\_\allowbreak{}c(2882)\allowbreak{}0 | Xibar\_\allowbreak{}c(2882)\allowbreak{}0\par {\fontsize{8.1}{9.4}\selectfont PDG2026:\allowbreak{}observable:\allowbreak{}B203:\allowbreak{}B203W0:\allowbreak{}351252} & \textbf{12+-8} MeV & \(u_+\): 7.856844\allowbreak{}15016614\allowbreak{}68\par \(u_-\): 7.856844\allowbreak{}15016614\allowbreak{}68\par Tipo: NONE\par Confianza: ND & COMP-BARYON\par E04 / M07\\\addlinespace[3pt]
% VI-CATALOG-ROW anchuras 142
\textbf{142. P\_\allowbreak{}c\_\allowbreak{}cbar\_\allowbreak{}s(4338)\allowbreak{}0 WIDTH}\par P\_\allowbreak{}c\_\allowbreak{}cbar\_\allowbreak{}s(4338)\allowbreak{}+\par {\fontsize{8.1}{9.4}\selectfont PDG2026:\allowbreak{}observable:\allowbreak{}B204:\allowbreak{}B204W:\allowbreak{}351946} & \textbf{7.0+-1.8} MeV & \(u_+\): 1.769180\allowbreak{}60129541\allowbreak{}3\par \(u_-\): 1.769180\allowbreak{}60129541\allowbreak{}3\par Tipo: NONE\par Confianza: ND & COMP-BARYON\par E04 / M07\\\addlinespace[3pt]
% VI-CATALOG-ROW anchuras 143
\textbf{143. P\_\allowbreak{}c\_\allowbreak{}cbar(4312)\allowbreak{}+ WIDTH}\par P\_\allowbreak{}c\_\allowbreak{}cbar(4312)\allowbreak{}+\par {\fontsize{8.1}{9.4}\selectfont PDG2026:\allowbreak{}observable:\allowbreak{}B205:\allowbreak{}B205W:\allowbreak{}352000} & \textbf{19+8-7} MeV & \(u_+\): 8.156883\allowbreak{}01652626\allowbreak{}54\par \(u_-\): 7.255202\allowbreak{}81398952\allowbreak{}7\par Tipo: NONE\par Confianza: ND & COMP-BARYON\par E04 / M07\\\addlinespace[3pt]
% VI-CATALOG-ROW anchuras 144
\textbf{144. Lambda\_\allowbreak{}c(2910)\allowbreak{}+ WIDTH}\par Lambda\_\allowbreak{}c(2910)\allowbreak{}0 | Lambdabar\_\allowbreak{}c(2910)\allowbreak{}0\par {\fontsize{8.1}{9.4}\selectfont PDG2026:\allowbreak{}observable:\allowbreak{}B206:\allowbreak{}B206W:\allowbreak{}350974} & \textbf{52+-27} MeV & \(u_+\): 27.44886\allowbreak{}15428764\allowbreak{}18\par \(u_-\): 27.44886\allowbreak{}15428764\allowbreak{}18\par Tipo: NONE\par Confianza: ND & COMP-BARYON\par E04 / M07\\\addlinespace[3pt]
% VI-CATALOG-ROW anchuras 145
\textbf{145. Xi\_\allowbreak{}b(6087)\allowbreak{}0 WIDTH}\par Xi\_\allowbreak{}b(6087)\par {\fontsize{8.1}{9.4}\selectfont PDG2026:\allowbreak{}observable:\allowbreak{}B207:\allowbreak{}B207W:\allowbreak{}351853} & \textbf{2.4+-0.5} MeV & \(u_+\): 0.519711\allowbreak{}45840745\allowbreak{}137\par \(u_-\): 0.519711\allowbreak{}45840745\allowbreak{}137\par Tipo: NONE\par Confianza: ND & COMP-BARYON\par E04 / M07\\\addlinespace[3pt]
% VI-CATALOG-ROW anchuras 146
\textbf{146. Xi\_\allowbreak{}b(6095)\allowbreak{}0 WIDTH}\par Xi\_\allowbreak{}b(6095)\par {\fontsize{8.1}{9.4}\selectfont PDG2026:\allowbreak{}observable:\allowbreak{}B208:\allowbreak{}B208W:\allowbreak{}351857} & \textbf{0.50+-0.35} MeV & \(u_+\): 0.347850\allowbreak{}54261852\allowbreak{}172\par \(u_-\): 0.347850\allowbreak{}54261852\allowbreak{}172\par Tipo: NONE\par Confianza: ND & COMP-BARYON\par E04 / M07\\\addlinespace[3pt]
% VI-CATALOG-ROW anchuras 147
\textbf{147. Omega\_\allowbreak{}c(3185)\allowbreak{}0 WIDTH}\par Omega\_\allowbreak{}c(3185)\allowbreak{}0 | Omegabar\_\allowbreak{}c(3185)\allowbreak{}0\par {\fontsize{8.1}{9.4}\selectfont PDG2026:\allowbreak{}observable:\allowbreak{}B209:\allowbreak{}B209W:\allowbreak{}351397} & \textbf{50+12-21} MeV & \(u_+\): 12.20655\allowbreak{}56157337\par \(u_-\): 21.18962\allowbreak{}01004170\allowbreak{}91\par Tipo: NONE\par Confianza: ND & COMP-BARYON\par E04 / M07\\\addlinespace[3pt]
% VI-CATALOG-ROW anchuras 148
\textbf{148. Omega\_\allowbreak{}c(3327)\allowbreak{}0 WIDTH}\par Omega\_\allowbreak{}c(3327)\allowbreak{}0 | Omegabar\_\allowbreak{}c(3327)\allowbreak{}0\par {\fontsize{8.1}{9.4}\selectfont PDG2026:\allowbreak{}observable:\allowbreak{}B210:\allowbreak{}B210W:\allowbreak{}351401} & \textbf{20+14-5} MeV & \(u_+\): 13.92838\allowbreak{}82771841\allowbreak{}19\par \(u_-\): 5.099019\allowbreak{}51359278\allowbreak{}36\par Tipo: NONE\par Confianza: ND & COMP-BARYON\par E04 / M07\\\addlinespace[3pt]
% VI-CATALOG-ROW anchuras 149
\textbf{149. Omega(2109)\allowbreak{}- WIDTH}\par Omega(2109)\allowbreak{}- | Omegabar(2109)\allowbreak{}+\par {\fontsize{8.1}{9.4}\selectfont PDG2026:\allowbreak{}observable:\allowbreak{}B211:\allowbreak{}B211W:\allowbreak{}350515} & \textbf{18+-17} MeV & \(u_+\): 17.36231\allowbreak{}55137786\allowbreak{}79\par \(u_-\): 17.36231\allowbreak{}55137786\allowbreak{}79\par Tipo: NONE\par Confianza: ND & COMP-BARYON\par E04 / M07\\\addlinespace[3pt]
% VI-CATALOG-ROW anchuras 150
\textbf{150. omega(782) WIDTH}\par omega(782)\allowbreak{}0\par {\fontsize{8.1}{9.4}\selectfont PDG2026:\allowbreak{}observable:\allowbreak{}M001:\allowbreak{}M001W:\allowbreak{}333933} & \textbf{8.68+-0.13} MeV & \(u_+\): 0.132000\allowbreak{}08707764\allowbreak{}46\par \(u_-\): 0.132000\allowbreak{}08707764\allowbreak{}46\par Tipo: NONE\par Confianza: ND & COMP-MESON\par E04 / M08\\\addlinespace[3pt]
% VI-CATALOG-ROW anchuras 151
\textbf{151. eta'(958) WIDTH}\par eta\textasciicircum{}'(958)\allowbreak{}0\par {\fontsize{8.1}{9.4}\selectfont PDG2026:\allowbreak{}observable:\allowbreak{}M002:\allowbreak{}M002W:\allowbreak{}334007} & \textbf{0.188+-0.006} MeV & \(u_+\): 0.006444\allowbreak{}22122653\allowbreak{}40033\par \(u_-\): 0.006444\allowbreak{}22122653\allowbreak{}40033\par Tipo: NONE\par Confianza: ND & COMP-MESON\par E04 / M08\\\addlinespace[3pt]
% VI-CATALOG-ROW anchuras 152
\textbf{152. f\_\allowbreak{}0(980) WIDTH}\par f\_\allowbreak{}0(980)\allowbreak{}0\par {\fontsize{8.1}{9.4}\selectfont PDG2026:\allowbreak{}observable:\allowbreak{}M003:\allowbreak{}M003W1:\allowbreak{}334164} & \textbf{10 to 100} MeV & \(u_+\): NA\par \(u_-\): NA\par Tipo: R\par Confianza: ND & COMP-MESON\par E04 / M08\\\addlinespace[3pt]
% VI-CATALOG-ROW anchuras 153
\textbf{153. phi(1020) WIDTH}\par phi(1020)\allowbreak{}0\par {\fontsize{8.1}{9.4}\selectfont PDG2026:\allowbreak{}observable:\allowbreak{}M004:\allowbreak{}M004W:\allowbreak{}334183} & \textbf{4.249+-0.013} MeV & \(u_+\): 0.013116\allowbreak{}55783805\allowbreak{}327\par \(u_-\): 0.013116\allowbreak{}55783805\allowbreak{}327\par Tipo: NONE\par Confianza: ND & COMP-MESON\par E04 / M08\\\addlinespace[3pt]
% VI-CATALOG-ROW anchuras 154
\textbf{154. f\_\allowbreak{}2(1270) WIDTH}\par f\_\allowbreak{}2(1270)\allowbreak{}0\par {\fontsize{8.1}{9.4}\selectfont PDG2026:\allowbreak{}observable:\allowbreak{}M005:\allowbreak{}M005W:\allowbreak{}334355} & \textbf{186.6+2.8-2.2} MeV & \(u_+\): 2.802312\allowbreak{}63767857\allowbreak{}5\par \(u_-\): 2.196115\allowbreak{}62862221\allowbreak{}81\par Tipo: NONE\par Confianza: ND & COMP-MESON\par E04 / M08\\\addlinespace[3pt]
% VI-CATALOG-ROW anchuras 155
\textbf{155. f\_\allowbreak{}1(1420) WIDTH}\par f\_\allowbreak{}1(1420)\allowbreak{}0\par {\fontsize{8.1}{9.4}\selectfont PDG2026:\allowbreak{}observable:\allowbreak{}M006:\allowbreak{}M006W:\allowbreak{}334555} & \textbf{55.4+-3.2} MeV & \(u_+\): 3.157214\allowbreak{}33119048\allowbreak{}68\par \(u_-\): 3.199679\allowbreak{}79091576\allowbreak{}9\par Tipo: NONE\par Confianza: ND & COMP-MESON\par E04 / M08\\\addlinespace[3pt]
% VI-CATALOG-ROW anchuras 156
\textbf{156. f\_\allowbreak{}1(1285) WIDTH}\par f\_\allowbreak{}1(1285)\allowbreak{}0\par {\fontsize{8.1}{9.4}\selectfont PDG2026:\allowbreak{}observable:\allowbreak{}M008:\allowbreak{}M008W:\allowbreak{}334404} & \textbf{23.0+-1.0} MeV & \(u_+\): 1.045581\allowbreak{}01639747\allowbreak{}5\par \(u_-\): 1.025615\allowbreak{}91020738\allowbreak{}61\par Tipo: NONE\par Confianza: ND & COMP-MESON\par E04 / M08\\\addlinespace[3pt]
% VI-CATALOG-ROW anchuras 157
\textbf{157. NEUTRAL ONLY, e+ e-}\par rho(770)\allowbreak{}0 | rho(770)\allowbreak{}+ | rho(770)\allowbreak{}-\par {\fontsize{8.1}{9.4}\selectfont PDG2026:\allowbreak{}observable:\allowbreak{}M009:\allowbreak{}M009W0:\allowbreak{}333870} & \textbf{147.4+-0.8} MeV & \(u_+\): 0.820963\allowbreak{}24211264\allowbreak{}493\par \(u_-\): 0.820963\allowbreak{}24211264\allowbreak{}493\par Tipo: NONE\par Confianza: ND & COMP-MESON\par E04 / M08\\\addlinespace[3pt]
% VI-CATALOG-ROW anchuras 158
\textbf{158. CHARGED ONLY, tau DECAYS and e+ e-}\par rho(770)\allowbreak{}0 | rho(770)\allowbreak{}+ | rho(770)\allowbreak{}-\par {\fontsize{8.1}{9.4}\selectfont PDG2026:\allowbreak{}observable:\allowbreak{}M009:\allowbreak{}M009W5:\allowbreak{}333871} & \textbf{149.1+-0.8} MeV & \(u_+\): 0.849067\allowbreak{}62313788\allowbreak{}198\par \(u_-\): 0.849067\allowbreak{}62313788\allowbreak{}198\par Tipo: NONE\par Confianza: ND & COMP-MESON\par E04 / M08\\\addlinespace[3pt]
% VI-CATALOG-ROW anchuras 159
\textbf{159. MIXED CHARGES, OTHER REACTIONS}\par rho(770)\allowbreak{}0 | rho(770)\allowbreak{}+ | rho(770)\allowbreak{}-\par {\fontsize{8.1}{9.4}\selectfont PDG2026:\allowbreak{}observable:\allowbreak{}M009:\allowbreak{}M009W7:\allowbreak{}333873} & \textbf{149.5+-1.3} MeV & \(u_+\): 1.3\par \(u_-\): 1.3\par Tipo: NONE\par Confianza: ND & COMP-MESON\par E04 / M08\\\addlinespace[3pt]
% VI-CATALOG-ROW anchuras 160
\textbf{160. CHARGED ONLY, HADROPRODUCED}\par rho(770)\allowbreak{}0 | rho(770)\allowbreak{}+ | rho(770)\allowbreak{}-\par {\fontsize{8.1}{9.4}\selectfont PDG2026:\allowbreak{}observable:\allowbreak{}M009:\allowbreak{}M009W2:\allowbreak{}333874} & \textbf{150.2+-2.4} MeV & \(u_+\): 2.392524\allowbreak{}19258698\allowbreak{}49\par \(u_-\): 2.392524\allowbreak{}19258698\allowbreak{}49\par Tipo: NONE\par Confianza: ND & COMP-MESON\par E04 / M08\\\addlinespace[3pt]
% VI-CATALOG-ROW anchuras 161
\textbf{161. NEUTRAL ONLY, PHOTOPRODUCED}\par rho(770)\allowbreak{}0 | rho(770)\allowbreak{}+ | rho(770)\allowbreak{}-\par {\fontsize{8.1}{9.4}\selectfont PDG2026:\allowbreak{}observable:\allowbreak{}M009:\allowbreak{}M009W0P:\allowbreak{}333876} & \textbf{152.3+-1.7} MeV & \(u_+\): 1.614419\allowbreak{}61384316\allowbreak{}19\par \(u_-\): 1.747655\allowbreak{}09508910\allowbreak{}81\par Tipo: NONE\par Confianza: ND & COMP-MESON\par E04 / M08\\\addlinespace[3pt]
% VI-CATALOG-ROW anchuras 162
\textbf{162. NEUTRAL ONLY, OTHER REACTIONS}\par rho(770)\allowbreak{}0 | rho(770)\allowbreak{}+ | rho(770)\allowbreak{}-\par {\fontsize{8.1}{9.4}\selectfont PDG2026:\allowbreak{}observable:\allowbreak{}M009:\allowbreak{}M009W0R:\allowbreak{}333877} & \textbf{150.9+-1.7} MeV & \(u_+\): 1.734188\allowbreak{}03219095\allowbreak{}91\par \(u_-\): 1.734188\allowbreak{}03219095\allowbreak{}91\par Tipo: NONE\par Confianza: ND & COMP-MESON\par E04 / M08\\\addlinespace[3pt]
% VI-CATALOG-ROW anchuras 163
\textbf{163. a\_\allowbreak{}1(1260) WIDTH}\par a\_\allowbreak{}1(1260)\allowbreak{}0 | a\_\allowbreak{}1(1260)\allowbreak{}+ | a\_\allowbreak{}1(1260)\allowbreak{}-\par {\fontsize{8.1}{9.4}\selectfont PDG2026:\allowbreak{}observable:\allowbreak{}M010:\allowbreak{}M010W:\allowbreak{}334333} & \textbf{250 to 600} MeV & \(u_+\): NA\par \(u_-\): NA\par Tipo: R\par Confianza: ND & COMP-MESON\par E04 / M08\\\addlinespace[3pt]
% VI-CATALOG-ROW anchuras 164
\textbf{164. b\_\allowbreak{}1(1235) WIDTH}\par b\_\allowbreak{}1(1235)\allowbreak{}0 | b\_\allowbreak{}1(1235)\allowbreak{}+ | b\_\allowbreak{}1(1235)\allowbreak{}-\par {\fontsize{8.1}{9.4}\selectfont PDG2026:\allowbreak{}observable:\allowbreak{}M011:\allowbreak{}M011W:\allowbreak{}334310} & \textbf{142+-9} MeV & \(u_+\): 8.505763\allowbreak{}01057597\allowbreak{}03\par \(u_-\): 8.505763\allowbreak{}01057597\allowbreak{}03\par Tipo: NONE\par Confianza: ND & COMP-MESON\par E04 / M08\\\addlinespace[3pt]
% VI-CATALOG-ROW anchuras 165
\textbf{165. a\_\allowbreak{}2(1320) WIDTH}\par a\_\allowbreak{}2(1320)\allowbreak{}0 | a\_\allowbreak{}2(1320)\allowbreak{}+ | a\_\allowbreak{}2(1320)\allowbreak{}-\par {\fontsize{8.1}{9.4}\selectfont PDG2026:\allowbreak{}observable:\allowbreak{}M012:\allowbreak{}M012W:\allowbreak{}334474} & \textbf{107+-5} MeV & \(u_+\): 5\par \(u_-\): 5\par Tipo: NONE\par Confianza: ND & COMP-MESON\par E04 / M08\\\addlinespace[3pt]
% VI-CATALOG-ROW anchuras 166
\textbf{166. 3pi MODE}\par a\_\allowbreak{}2(1320)\allowbreak{}0 | a\_\allowbreak{}2(1320)\allowbreak{}+ | a\_\allowbreak{}2(1320)\allowbreak{}-\par {\fontsize{8.1}{9.4}\selectfont PDG2026:\allowbreak{}observable:\allowbreak{}M012:\allowbreak{}M012W1:\allowbreak{}334476} & \textbf{105.0+1.7-1.9} MeV & \(u_+\): 1.672423\allowbreak{}94191988\allowbreak{}7\par \(u_-\): 1.852396\allowbreak{}15833847\allowbreak{}09\par Tipo: NONE\par Confianza: ND & COMP-MESON\par E04 / M08\\\addlinespace[3pt]
% VI-CATALOG-ROW anchuras 167
\textbf{167. K Kbar MODE}\par a\_\allowbreak{}2(1320)\allowbreak{}0 | a\_\allowbreak{}2(1320)\allowbreak{}+ | a\_\allowbreak{}2(1320)\allowbreak{}-\par {\fontsize{8.1}{9.4}\selectfont PDG2026:\allowbreak{}observable:\allowbreak{}M012:\allowbreak{}M012W2:\allowbreak{}334477} & \textbf{109.8+-2.4} MeV & \(u_+\): 2.374890\allowbreak{}24273604\allowbreak{}12\par \(u_-\): 2.374890\allowbreak{}24273604\allowbreak{}12\par Tipo: NONE\par Confianza: ND & COMP-MESON\par E04 / M08\\\addlinespace[3pt]
% VI-CATALOG-ROW anchuras 168
\textbf{168. eta pi MODE}\par a\_\allowbreak{}2(1320)\allowbreak{}0 | a\_\allowbreak{}2(1320)\allowbreak{}+ | a\_\allowbreak{}2(1320)\allowbreak{}-\par {\fontsize{8.1}{9.4}\selectfont PDG2026:\allowbreak{}observable:\allowbreak{}M012:\allowbreak{}M012W3:\allowbreak{}334478} & \textbf{111.1+-2.4} MeV & \(u_+\): 2.448965\allowbreak{}93502306\allowbreak{}5\par \(u_-\): 2.448965\allowbreak{}93502306\allowbreak{}5\par Tipo: NONE\par Confianza: ND & COMP-MESON\par E04 / M08\\\addlinespace[3pt]
% VI-CATALOG-ROW anchuras 169
\textbf{169. eta' pi MODE}\par a\_\allowbreak{}2(1320)\allowbreak{}0 | a\_\allowbreak{}2(1320)\allowbreak{}+ | a\_\allowbreak{}2(1320)\allowbreak{}-\par {\fontsize{8.1}{9.4}\selectfont PDG2026:\allowbreak{}observable:\allowbreak{}M012:\allowbreak{}M012W4:\allowbreak{}334479} & \textbf{119+-25} MeV & \(u_+\): 25.06315\allowbreak{}16454720\allowbreak{}01\par \(u_-\): 25.06315\allowbreak{}16454720\allowbreak{}01\par Tipo: NONE\par Confianza: ND & COMP-MESON\par E04 / M08\\\addlinespace[3pt]
% VI-CATALOG-ROW anchuras 170
\textbf{170. f\_\allowbreak{}2'(1525) WIDTH}\par f\_\allowbreak{}2\textasciicircum{}'(1525)\allowbreak{}0\par {\fontsize{8.1}{9.4}\selectfont PDG2026:\allowbreak{}observable:\allowbreak{}M013:\allowbreak{}M013WX:\allowbreak{}334671} & \textbf{72+7-6} MeV & \(u_+\): 6.882857\allowbreak{}45387930\allowbreak{}53\par \(u_-\): 5.788689\allowbreak{}04167363\allowbreak{}58\par Tipo: NONE\par Confianza: ND & COMP-MESON\par E04 / M08\\\addlinespace[3pt]
% VI-CATALOG-ROW anchuras 171
\textbf{171. PRODUCED BY K+- BEAM}\par f\_\allowbreak{}2\textasciicircum{}'(1525)\allowbreak{}0\par {\fontsize{8.1}{9.4}\selectfont PDG2026:\allowbreak{}observable:\allowbreak{}M013:\allowbreak{}M013W2:\allowbreak{}334673} & \textbf{82+-6} MeV & \(u_+\): 6.069003\allowbreak{}12634116\allowbreak{}46\par \(u_-\): 5.743699\allowbreak{}64556114\allowbreak{}8\par Tipo: NONE\par Confianza: ND & COMP-MESON\par E04 / M08\\\addlinespace[3pt]
% VI-CATALOG-ROW anchuras 172
\textbf{172. PRODUCED IN e+ e- ANNIHILATION AND PARTICLE DECAYS}\par f\_\allowbreak{}2\textasciicircum{}'(1525)\allowbreak{}0\par {\fontsize{8.1}{9.4}\selectfont PDG2026:\allowbreak{}observable:\allowbreak{}M013:\allowbreak{}M013W3:\allowbreak{}334674} & \textbf{86+-4} MeV & \(u_+\): 4.272413\allowbreak{}00467022\allowbreak{}71\par \(u_-\): 3.957898\allowbreak{}39152937\allowbreak{}39\par Tipo: NONE\par Confianza: ND & COMP-MESON\par E04 / M08\\\addlinespace[3pt]
% VI-CATALOG-ROW anchuras 173
\textbf{173. PRODUCED IN pbar p ANNIHILATION}\par f\_\allowbreak{}2\textasciicircum{}'(1525)\allowbreak{}0\par {\fontsize{8.1}{9.4}\selectfont PDG2026:\allowbreak{}observable:\allowbreak{}M013:\allowbreak{}M013W9:\allowbreak{}334675} & \textbf{77+-5} MeV & \(u_+\): 4.799999\allowbreak{}99999999\allowbreak{}98\par \(u_-\): 4.799999\allowbreak{}99999999\allowbreak{}98\par Tipo: NONE\par Confianza: ND & COMP-MESON\par E04 / M08\\\addlinespace[3pt]
% VI-CATALOG-ROW anchuras 174
\textbf{174. CENTRAL PRODUCTION}\par f\_\allowbreak{}2\textasciicircum{}'(1525)\allowbreak{}0\par {\fontsize{8.1}{9.4}\selectfont PDG2026:\allowbreak{}observable:\allowbreak{}M013:\allowbreak{}M013W4:\allowbreak{}334676} & \textbf{70+-25} MeV & \(u_+\): 25\par \(u_-\): 25\par Tipo: NONE\par Confianza: ND & COMP-MESON\par E04 / M08\\\addlinespace[3pt]
% VI-CATALOG-ROW anchuras 175
\textbf{175. PRODUCED IN e p COLLISIONS}\par f\_\allowbreak{}2\textasciicircum{}'(1525)\allowbreak{}0\par {\fontsize{8.1}{9.4}\selectfont PDG2026:\allowbreak{}observable:\allowbreak{}M013:\allowbreak{}M013W10:\allowbreak{}334677} & \textbf{83+-10} MeV & \(u_+\): 10.29563\allowbreak{}01409870\allowbreak{}01\par \(u_-\): 9.848857\allowbreak{}80179610\allowbreak{}39\par Tipo: NONE\par Confianza: ND & COMP-MESON\par E04 / M08\\\addlinespace[3pt]
% VI-CATALOG-ROW anchuras 176
\textbf{176. f\_\allowbreak{}0(500) BREIT-WIGNER WIDTH}\par f\_\allowbreak{}0(500)\allowbreak{}0\par {\fontsize{8.1}{9.4}\selectfont PDG2026:\allowbreak{}observable:\allowbreak{}M014:\allowbreak{}M014W:\allowbreak{}333858} & \textbf{100 to 800} MeV & \(u_+\): NA\par \(u_-\): NA\par Tipo: R\par Confianza: ND & COMP-MESON\par E04 / M08\\\addlinespace[3pt]
% VI-CATALOG-ROW anchuras 177
\textbf{177. 2pi, K Kbar, AND K Kbar pi MODES}\par rho\_\allowbreak{}3(1690)\allowbreak{}0 | rho\_\allowbreak{}3(1690)\allowbreak{}+ | rho\_\allowbreak{}3(1690)\allowbreak{}-\par {\fontsize{8.1}{9.4}\selectfont PDG2026:\allowbreak{}observable:\allowbreak{}M015:\allowbreak{}M015W:\allowbreak{}334848} & \textbf{161+-10} MeV & \(u_+\): 9.701480\allowbreak{}09186499\allowbreak{}04\par \(u_-\): 9.653732\allowbreak{}12649761\allowbreak{}39\par Tipo: NONE\par Confianza: ND & COMP-MESON\par E04 / M08\\\addlinespace[3pt]
% VI-CATALOG-ROW anchuras 178
\textbf{178. 2pi MODE}\par rho\_\allowbreak{}3(1690)\allowbreak{}0 | rho\_\allowbreak{}3(1690)\allowbreak{}+ | rho\_\allowbreak{}3(1690)\allowbreak{}-\par {\fontsize{8.1}{9.4}\selectfont PDG2026:\allowbreak{}observable:\allowbreak{}M015:\allowbreak{}M015W1:\allowbreak{}334849} & \textbf{186+-14} MeV & \(u_+\): 14.04672\allowbreak{}49736341\allowbreak{}89\par \(u_-\): 14.04672\allowbreak{}49736341\allowbreak{}89\par Tipo: NONE\par Confianza: ND & COMP-MESON\par E04 / M08\\\addlinespace[3pt]
% VI-CATALOG-ROW anchuras 179
\textbf{179. K Kbar AND K Kbar pi MODES}\par rho\_\allowbreak{}3(1690)\allowbreak{}0 | rho\_\allowbreak{}3(1690)\allowbreak{}+ | rho\_\allowbreak{}3(1690)\allowbreak{}-\par {\fontsize{8.1}{9.4}\selectfont PDG2026:\allowbreak{}observable:\allowbreak{}M015:\allowbreak{}M015W2:\allowbreak{}334850} & \textbf{204+-18} MeV & \(u_+\): 17.88854\allowbreak{}38199983\allowbreak{}18\par \(u_-\): 17.88854\allowbreak{}38199983\allowbreak{}18\par Tipo: NONE\par Confianza: ND & COMP-MESON\par E04 / M08\\\addlinespace[3pt]
% VI-CATALOG-ROW anchuras 180
\textbf{180. (4pi)\allowbreak{}+- MODE}\par rho\_\allowbreak{}3(1690)\allowbreak{}0 | rho\_\allowbreak{}3(1690)\allowbreak{}+ | rho\_\allowbreak{}3(1690)\allowbreak{}-\par {\fontsize{8.1}{9.4}\selectfont PDG2026:\allowbreak{}observable:\allowbreak{}M015:\allowbreak{}M015W3:\allowbreak{}334851} & \textbf{129+-10} MeV & \(u_+\): 10.28440\allowbreak{}09679589\allowbreak{}39\par \(u_-\): 10.16157\allowbreak{}27369938\allowbreak{}9\par Tipo: NONE\par Confianza: ND & COMP-MESON\par E04 / M08\\\addlinespace[3pt]
% VI-CATALOG-ROW anchuras 181
\textbf{181. omega pi MODE}\par rho\_\allowbreak{}3(1690)\allowbreak{}0 | rho\_\allowbreak{}3(1690)\allowbreak{}+ | rho\_\allowbreak{}3(1690)\allowbreak{}-\par {\fontsize{8.1}{9.4}\selectfont PDG2026:\allowbreak{}observable:\allowbreak{}M015:\allowbreak{}M015W5:\allowbreak{}334852} & \textbf{190+-40} MeV & \(u_+\): 35.52783\allowbreak{}84233964\allowbreak{}18\par \(u_-\): 35.52783\allowbreak{}84233964\allowbreak{}18\par Tipo: NONE\par Confianza: ND & COMP-MESON\par E04 / M08\\\addlinespace[3pt]
% VI-CATALOG-ROW anchuras 182
\textbf{182. eta pi+ pi- MODE}\par rho\_\allowbreak{}3(1690)\allowbreak{}0 | rho\_\allowbreak{}3(1690)\allowbreak{}+ | rho\_\allowbreak{}3(1690)\allowbreak{}-\par {\fontsize{8.1}{9.4}\selectfont PDG2026:\allowbreak{}observable:\allowbreak{}M015:\allowbreak{}M015W6:\allowbreak{}334853} & \textbf{130+-40} MeV & \(u_+\): 43.46735\allowbreak{}18760492\allowbreak{}57\par \(u_-\): 43.46735\allowbreak{}18760492\allowbreak{}57\par Tipo: NONE\par Confianza: ND & COMP-MESON\par E04 / M08\\\addlinespace[3pt]
% VI-CATALOG-ROW anchuras 183
\textbf{183. f\_\allowbreak{}4(2050) WIDTH}\par f\_\allowbreak{}4(2050)\allowbreak{}0\par {\fontsize{8.1}{9.4}\selectfont PDG2026:\allowbreak{}observable:\allowbreak{}M016:\allowbreak{}M016W:\allowbreak{}335131} & \textbf{237+-18} MeV & \(u_+\): 17.88029\allowbreak{}94325271\allowbreak{}7\par \(u_-\): 18.17953\allowbreak{}28375529\allowbreak{}09\par Tipo: NONE\par Confianza: ND & COMP-MESON\par E04 / M08\\\addlinespace[3pt]
% VI-CATALOG-ROW anchuras 184
\textbf{184. a\_\allowbreak{}4(1970) WIDTH}\par a\_\allowbreak{}4(1970)\allowbreak{}0 | a\_\allowbreak{}4(1970)\allowbreak{}+ | a\_\allowbreak{}4(1970)\allowbreak{}-\par {\fontsize{8.1}{9.4}\selectfont PDG2026:\allowbreak{}observable:\allowbreak{}M017:\allowbreak{}M017W:\allowbreak{}335091} & \textbf{324+15-18} MeV & \(u_+\): 14.53339\allowbreak{}33159854\allowbreak{}99\par \(u_-\): 18.10011\allowbreak{}67978534\allowbreak{}31\par Tipo: NONE\par Confianza: ND & COMP-MESON\par E04 / M08\\\addlinespace[3pt]
% VI-CATALOG-ROW anchuras 185
\textbf{185. CHARGED ONLY, HADROPRODUCED}\par K\textasciicircum{}*(892)\allowbreak{}0 | Kbar\textasciicircum{}*(892)\allowbreak{}0 | K\textasciicircum{}*(892)\allowbreak{}+ | K\textasciicircum{}*(892)\allowbreak{}-\par {\fontsize{8.1}{9.4}\selectfont PDG2026:\allowbreak{}observable:\allowbreak{}M018:\allowbreak{}M018W1:\allowbreak{}335795} & \textbf{48.5+-1.2} MeV & \(u_+\): 1.189202\allowbreak{}18180429\allowbreak{}3\par \(u_-\): 1.189202\allowbreak{}18180429\allowbreak{}3\par Tipo: NONE\par Confianza: ND & COMP-MESON\par E04 / M08\\\addlinespace[3pt]
% VI-CATALOG-ROW anchuras 186
\textbf{186. CHARGED ONLY, PRODUCED IN tau LEPTON DECAYS}\par K\textasciicircum{}*(892)\allowbreak{}0 | Kbar\textasciicircum{}*(892)\allowbreak{}0 | K\textasciicircum{}*(892)\allowbreak{}+ | K\textasciicircum{}*(892)\allowbreak{}-\par {\fontsize{8.1}{9.4}\selectfont PDG2026:\allowbreak{}observable:\allowbreak{}M018:\allowbreak{}M018W5:\allowbreak{}335797} & \textbf{46.2+-1.3} MeV & \(u_+\): 1.341640\allowbreak{}78649987\allowbreak{}41\par \(u_-\): 1.341640\allowbreak{}78649987\allowbreak{}41\par Tipo: NONE\par Confianza: ND & COMP-MESON\par E04 / M08\\\addlinespace[3pt]
% VI-CATALOG-ROW anchuras 187
\textbf{187. NEUTRAL ONLY}\par K\textasciicircum{}*(892)\allowbreak{}0 | Kbar\textasciicircum{}*(892)\allowbreak{}0 | K\textasciicircum{}*(892)\allowbreak{}+ | K\textasciicircum{}*(892)\allowbreak{}-\par {\fontsize{8.1}{9.4}\selectfont PDG2026:\allowbreak{}observable:\allowbreak{}M018:\allowbreak{}M018W2:\allowbreak{}335798} & \textbf{47.1+-0.5} MeV & \(u_+\): 0.493474\allowbreak{}96850837\allowbreak{}202\par \(u_-\): 0.493474\allowbreak{}96850837\allowbreak{}202\par Tipo: NONE\par Confianza: ND & COMP-MESON\par E04 / M08\\\addlinespace[3pt]
% VI-CATALOG-ROW anchuras 188
\textbf{188. K\_\allowbreak{}0*(1430) WIDTH}\par K\_\allowbreak{}0\textasciicircum{}*(1430)\allowbreak{}0 | Kbar\_\allowbreak{}0\textasciicircum{}*(1430)\allowbreak{}0 | K\_\allowbreak{}0\textasciicircum{}*(1430)\allowbreak{}+ | K\_\allowbreak{}0\textasciicircum{}*(1430)\allowbreak{}-\par {\fontsize{8.1}{9.4}\selectfont PDG2026:\allowbreak{}observable:\allowbreak{}M019:\allowbreak{}M019W:\allowbreak{}335879} & \textbf{270+-80} MeV & \(u_+\): 80\par \(u_-\): 80\par Tipo: NONE\par Confianza: ND & COMP-MESON\par E04 / M08\\\addlinespace[3pt]
% VI-CATALOG-ROW anchuras 189
\textbf{189. pi\_\allowbreak{}2(2100) WIDTH}\par pi\_\allowbreak{}2(2100)\allowbreak{}0 | pi\_\allowbreak{}2(2100)\allowbreak{}+ | pi\_\allowbreak{}2(2100)\allowbreak{}-\par {\fontsize{8.1}{9.4}\selectfont PDG2026:\allowbreak{}observable:\allowbreak{}M020:\allowbreak{}M020W:\allowbreak{}335147} & \textbf{620+-50} MeV & \(u_+\): 52.40000\allowbreak{}10252452\allowbreak{}67\par \(u_-\): 52.40000\allowbreak{}10252452\allowbreak{}67\par Tipo: NONE\par Confianza: ND & COMP-MESON\par E04 / M08\\\addlinespace[3pt]
% VI-CATALOG-ROW anchuras 190
\textbf{190. CHARGED ONLY, WITH FINAL STATE K pi}\par K\_\allowbreak{}2\textasciicircum{}*(1430)\allowbreak{}0 | Kbar\_\allowbreak{}2\textasciicircum{}*(1430)\allowbreak{}0 | K\_\allowbreak{}2\textasciicircum{}*(1430)\allowbreak{}+ | K\_\allowbreak{}2\textasciicircum{}*(1430)\allowbreak{}-\par {\fontsize{8.1}{9.4}\selectfont PDG2026:\allowbreak{}observable:\allowbreak{}M022:\allowbreak{}M022W1:\allowbreak{}335890} & \textbf{100.0+-2.2} MeV & \(u_+\): 2.236428\allowbreak{}14686197\allowbreak{}42\par \(u_-\): 2.178532\allowbreak{}63272092\allowbreak{}91\par Tipo: NONE\par Confianza: ND & COMP-MESON\par E04 / M08\\\addlinespace[3pt]
% VI-CATALOG-ROW anchuras 191
\textbf{191. NEUTRAL ONLY}\par K\_\allowbreak{}2\textasciicircum{}*(1430)\allowbreak{}0 | Kbar\_\allowbreak{}2\textasciicircum{}*(1430)\allowbreak{}0 | K\_\allowbreak{}2\textasciicircum{}*(1430)\allowbreak{}+ | K\_\allowbreak{}2\textasciicircum{}*(1430)\allowbreak{}-\par {\fontsize{8.1}{9.4}\selectfont PDG2026:\allowbreak{}observable:\allowbreak{}M022:\allowbreak{}M022W4:\allowbreak{}335892} & \textbf{109+-5} MeV & \(u_+\): 5.477817\allowbreak{}13955507\allowbreak{}94\par \(u_-\): 5.477817\allowbreak{}13955507\allowbreak{}94\par Tipo: NONE\par Confianza: ND & COMP-MESON\par E04 / M08\\\addlinespace[3pt]
% VI-CATALOG-ROW anchuras 192
\textbf{192. K\_\allowbreak{}2(1770) WIDTH}\par K\_\allowbreak{}2(1770)\allowbreak{}0 | Kbar\_\allowbreak{}2(1770)\allowbreak{}0 | K\_\allowbreak{}2(1770)\allowbreak{}+ | K\_\allowbreak{}2(1770)\allowbreak{}-\par {\fontsize{8.1}{9.4}\selectfont PDG2026:\allowbreak{}observable:\allowbreak{}M023:\allowbreak{}M023W:\allowbreak{}335964} & \textbf{186+-14} MeV & \(u_+\): 13.97802\allowbreak{}82998288\allowbreak{}69\par \(u_-\): 13.96323\allowbreak{}58217417\allowbreak{}9\par Tipo: NONE\par Confianza: ND & COMP-MESON\par E04 / M08\\\addlinespace[3pt]
% VI-CATALOG-ROW anchuras 193
\textbf{193. psi(4160) WIDTH}\par psi(4160)\par {\fontsize{8.1}{9.4}\selectfont PDG2026:\allowbreak{}observable:\allowbreak{}M025:\allowbreak{}M025W:\allowbreak{}346344} & \textbf{69+-10} MeV & \(u_+\): 10.53768\allowbreak{}52320708\allowbreak{}19\par \(u_-\): 9.602223\allowbreak{}38741453\allowbreak{}46\par Tipo: NONE\par Confianza: ND & COMP-MESON\par E04 / M08\\\addlinespace[3pt]
% VI-CATALOG-ROW anchuras 194
\textbf{194. eta\_\allowbreak{}c(1S) WIDTH}\par eta\_\allowbreak{}c(1S)\par {\fontsize{8.1}{9.4}\selectfont PDG2026:\allowbreak{}observable:\allowbreak{}M026:\allowbreak{}M026W:\allowbreak{}343287} & \textbf{30.0+-0.5} MeV & \(u_+\): 0.460621\allowbreak{}78744488\allowbreak{}053\par \(u_-\): 0.460621\allowbreak{}78744488\allowbreak{}053\par Tipo: NONE\par Confianza: ND & COMP-MESON\par E04 / M08\\\addlinespace[3pt]
% VI-CATALOG-ROW anchuras 195
\textbf{195. eta(1405) WIDTH}\par eta(1405)\allowbreak{}0\par {\fontsize{8.1}{9.4}\selectfont PDG2026:\allowbreak{}observable:\allowbreak{}M027:\allowbreak{}M027WX:\allowbreak{}334530} & \textbf{49.8+2.3-2.0} MeV & \(u_+\): 2.298064\allowbreak{}58315265\allowbreak{}21\par \(u_-\): 1.982120\allowbreak{}74261206\allowbreak{}99\par Tipo: NONE\par Confianza: ND & COMP-MESON\par E04 / M08\\\addlinespace[3pt]
% VI-CATALOG-ROW anchuras 196
\textbf{196. eta pi pi MODE}\par eta(1405)\allowbreak{}0\par {\fontsize{8.1}{9.4}\selectfont PDG2026:\allowbreak{}observable:\allowbreak{}M027:\allowbreak{}M027W1:\allowbreak{}334531} & \textbf{51.1+2.8-2.1} MeV & \(u_+\): 2.797925\allowbreak{}80324660\allowbreak{}28\par \(u_-\): 2.058132\allowbreak{}27941554\allowbreak{}99\par Tipo: NONE\par Confianza: ND & COMP-MESON\par E04 / M08\\\addlinespace[3pt]
% VI-CATALOG-ROW anchuras 197
\textbf{197. K Kbar pi MODE (a\_\allowbreak{}0(980) pi or direct K Kbar pi)}\par eta(1405)\allowbreak{}0\par {\fontsize{8.1}{9.4}\selectfont PDG2026:\allowbreak{}observable:\allowbreak{}M027:\allowbreak{}M027W4:\allowbreak{}334532} & \textbf{49+-4} MeV & \(u_+\): 3.912958\allowbreak{}95972736\allowbreak{}62\par \(u_-\): 4.108963\allowbreak{}96175021\allowbreak{}12\par Tipo: NONE\par Confianza: ND & COMP-MESON\par E04 / M08\\\addlinespace[3pt]
% VI-CATALOG-ROW anchuras 198
\textbf{198. pi pi gamma MODE}\par eta(1405)\allowbreak{}0\par {\fontsize{8.1}{9.4}\selectfont PDG2026:\allowbreak{}observable:\allowbreak{}M027:\allowbreak{}M027W2:\allowbreak{}334533} & \textbf{89+-17} MeV & \(u_+\): 17.30793\allowbreak{}25452340\allowbreak{}81\par \(u_-\): 17.30793\allowbreak{}25452340\allowbreak{}81\par Tipo: NONE\par Confianza: ND & COMP-MESON\par E04 / M08\\\addlinespace[3pt]
% VI-CATALOG-ROW anchuras 199
\textbf{199. phi gamma MODE}\par eta(1405)\allowbreak{}0\par {\fontsize{8.1}{9.4}\selectfont PDG2026:\allowbreak{}observable:\allowbreak{}M027:\allowbreak{}M027W7:\allowbreak{}334534} & \textbf{86+7-18} MeV & \(u_+\): 7.130918\allowbreak{}59440282\allowbreak{}77\par \(u_-\): 17.60823\allowbreak{}67089950\allowbreak{}3\par Tipo: NONE\par Confianza: ND & COMP-MESON\par E04 / M08\\\addlinespace[3pt]
% VI-CATALOG-ROW anchuras 200
\textbf{200. K\_\allowbreak{}1(1270) WIDTH}\par K\_\allowbreak{}1(1270)\allowbreak{}0 | Kbar\_\allowbreak{}1(1270)\allowbreak{}0 | K\_\allowbreak{}1(1270)\allowbreak{}+ | K\_\allowbreak{}1(1270)\allowbreak{}-\par {\fontsize{8.1}{9.4}\selectfont PDG2026:\allowbreak{}observable:\allowbreak{}M028:\allowbreak{}M028WX:\allowbreak{}335822} & \textbf{90+-20} MeV & \(u_+\): 20\par \(u_-\): 20\par Tipo: NONE\par Confianza: ND & COMP-MESON\par E04 / M08\\\addlinespace[3pt]
% VI-CATALOG-ROW anchuras 201
\textbf{201. PRODUCED BY K-, BACKWARD SCATTERING, HYPERON EXCHANGE}\par K\_\allowbreak{}1(1270)\allowbreak{}0 | Kbar\_\allowbreak{}1(1270)\allowbreak{}0 | K\_\allowbreak{}1(1270)\allowbreak{}+ | K\_\allowbreak{}1(1270)\allowbreak{}-\par {\fontsize{8.1}{9.4}\selectfont PDG2026:\allowbreak{}observable:\allowbreak{}M028:\allowbreak{}M028W2:\allowbreak{}335824} & \textbf{75+-15} MeV & \(u_+\): 15\par \(u_-\): 15\par Tipo: NONE\par Confianza: ND & COMP-MESON\par E04 / M08\\\addlinespace[3pt]
% VI-CATALOG-ROW anchuras 202
\textbf{202. PRODUCED BY K BEAMS}\par K\_\allowbreak{}1(1270)\allowbreak{}0 | Kbar\_\allowbreak{}1(1270)\allowbreak{}0 | K\_\allowbreak{}1(1270)\allowbreak{}+ | K\_\allowbreak{}1(1270)\allowbreak{}-\par {\fontsize{8.1}{9.4}\selectfont PDG2026:\allowbreak{}observable:\allowbreak{}M028:\allowbreak{}M028W3:\allowbreak{}335825} & \textbf{90+-8} MeV & \(u_+\): 8\par \(u_-\): 8\par Tipo: NONE\par Confianza: ND & COMP-MESON\par E04 / M08\\\addlinespace[3pt]
% VI-CATALOG-ROW anchuras 203
\textbf{203. PRODUCED BY BEAMS OTHER THAN K MESONS}\par K\_\allowbreak{}1(1270)\allowbreak{}0 | Kbar\_\allowbreak{}1(1270)\allowbreak{}0 | K\_\allowbreak{}1(1270)\allowbreak{}+ | K\_\allowbreak{}1(1270)\allowbreak{}-\par {\fontsize{8.1}{9.4}\selectfont PDG2026:\allowbreak{}observable:\allowbreak{}M028:\allowbreak{}M028W1:\allowbreak{}335826} & \textbf{126+-16} MeV & \(u_+\): 16.08195\allowbreak{}77830664\allowbreak{}29\par \(u_-\): 16.08195\allowbreak{}77830664\allowbreak{}29\par Tipo: NONE\par Confianza: ND & COMP-MESON\par E04 / M08\\\addlinespace[3pt]
% VI-CATALOG-ROW anchuras 204
\textbf{204. PRODUCED IN tau LEPTON DECAYS}\par K\_\allowbreak{}1(1270)\allowbreak{}0 | Kbar\_\allowbreak{}1(1270)\allowbreak{}0 | K\_\allowbreak{}1(1270)\allowbreak{}+ | K\_\allowbreak{}1(1270)\allowbreak{}-\par {\fontsize{8.1}{9.4}\selectfont PDG2026:\allowbreak{}observable:\allowbreak{}M028:\allowbreak{}M028WT:\allowbreak{}335827} & \textbf{260+120-110} MeV & \(u_+\): 120.4159\allowbreak{}45787922\allowbreak{}99\par \(u_-\): 106.3014\allowbreak{}58127346\allowbreak{}51\par Tipo: NONE\par Confianza: ND & COMP-MESON\par E04 / M08\\\addlinespace[3pt]
% VI-CATALOG-ROW anchuras 205
\textbf{205. X(3940) WIDTH}\par X(3940)\par {\fontsize{8.1}{9.4}\selectfont PDG2026:\allowbreak{}observable:\allowbreak{}M029:\allowbreak{}M029W:\allowbreak{}346245} & \textbf{43+28-18} MeV & \(u_+\): 28.10243\allowbreak{}81828863\allowbreak{}21\par \(u_-\): 17.56265\allowbreak{}52105673\allowbreak{}6\par Tipo: NONE\par Confianza: ND & COMP-MESON\par E04 / M08\\\addlinespace[3pt]
% VI-CATALOG-ROW anchuras 206
\textbf{206. h\_\allowbreak{}1(1170) WIDTH}\par h\_\allowbreak{}1(1170)\allowbreak{}0\par {\fontsize{8.1}{9.4}\selectfont PDG2026:\allowbreak{}observable:\allowbreak{}M030:\allowbreak{}M030W:\allowbreak{}334307} & \textbf{375+-35} MeV & \(u_+\): 34.52535\allowbreak{}30032641\allowbreak{}36\par \(u_-\): 34.52535\allowbreak{}30032641\allowbreak{}36\par Tipo: NONE\par Confianza: ND & COMP-MESON\par E04 / M08\\\addlinespace[3pt]
% VI-CATALOG-ROW anchuras 207
\textbf{207. pi- p -\kern0pt{}-> omega pi0 n}\par rho\_\allowbreak{}5(2350)\allowbreak{}0 | rho\_\allowbreak{}5(2350)\allowbreak{}+ | rho\_\allowbreak{}5(2350)\allowbreak{}-\par {\fontsize{8.1}{9.4}\selectfont PDG2026:\allowbreak{}observable:\allowbreak{}M033:\allowbreak{}M033W3:\allowbreak{}335260} & \textbf{400+-100} MeV & \(u_+\): 100\par \(u_-\): 100\par Tipo: NONE\par Confianza: ND & COMP-MESON\par E04 / M08\\\addlinespace[3pt]
% VI-CATALOG-ROW anchuras 208
\textbf{208. pi\_\allowbreak{}2(1670) WIDTH}\par pi\_\allowbreak{}2(1670)\allowbreak{}0 | pi\_\allowbreak{}2(1670)\allowbreak{}+ | pi\_\allowbreak{}2(1670)\allowbreak{}-\par {\fontsize{8.1}{9.4}\selectfont PDG2026:\allowbreak{}observable:\allowbreak{}M034:\allowbreak{}M034W:\allowbreak{}334784} & \textbf{258+8-9} MeV & \(u_+\): 7.994523\allowbreak{}19640628\allowbreak{}42\par \(u_-\): 9.032651\allowbreak{}63231475\allowbreak{}12\par Tipo: NONE\par Confianza: ND & COMP-MESON\par E04 / M08\\\addlinespace[3pt]
% VI-CATALOG-ROW anchuras 209
\textbf{209. K\_\allowbreak{}4*(2045) WIDTH}\par K\_\allowbreak{}4\textasciicircum{}*(2045)\allowbreak{}0 | Kbar\_\allowbreak{}4\textasciicircum{}*(2045)\allowbreak{}0 | K\_\allowbreak{}4\textasciicircum{}*(2045)\allowbreak{}+ | K\_\allowbreak{}4\textasciicircum{}*(2045)\allowbreak{}-\par {\fontsize{8.1}{9.4}\selectfont PDG2026:\allowbreak{}observable:\allowbreak{}M035:\allowbreak{}M035W:\allowbreak{}336017} & \textbf{199+27-19} MeV & \(u_+\): 26.50887\allowbreak{}10269003\allowbreak{}51\par \(u_-\): 19.30050\allowbreak{}75635472\allowbreak{}99\par Tipo: NONE\par Confianza: ND & COMP-MESON\par E04 / M08\\\addlinespace[3pt]
% VI-CATALOG-ROW anchuras 210
\textbf{210. a\_\allowbreak{}0(980) WIDTH}\par a\_\allowbreak{}0(980)\allowbreak{}0 | a\_\allowbreak{}0(980)\allowbreak{}+ | a\_\allowbreak{}0(980)\allowbreak{}-\par {\fontsize{8.1}{9.4}\selectfont PDG2026:\allowbreak{}observable:\allowbreak{}M036:\allowbreak{}M036W1:\allowbreak{}334172} & \textbf{50 to 100} MeV & \(u_+\): NA\par \(u_-\): NA\par Tipo: R\par Confianza: ND & COMP-MESON\par E04 / M08\\\addlinespace[3pt]
% VI-CATALOG-ROW anchuras 211
\textbf{211. eta(1295) WIDTH}\par eta(1295)\allowbreak{}0\par {\fontsize{8.1}{9.4}\selectfont PDG2026:\allowbreak{}observable:\allowbreak{}M037:\allowbreak{}M037W:\allowbreak{}334452} & \textbf{55+-5} MeV & \(u_+\): 5.366281\allowbreak{}63064819\allowbreak{}13\par \(u_-\): 5.366281\allowbreak{}63064819\allowbreak{}13\par Tipo: NONE\par Confianza: ND & COMP-MESON\par E04 / M08\\\addlinespace[3pt]
% VI-CATALOG-ROW anchuras 212
\textbf{212. f\_\allowbreak{}2(1810) WIDTH}\par f\_\allowbreak{}2(1810)\allowbreak{}0\par {\fontsize{8.1}{9.4}\selectfont PDG2026:\allowbreak{}observable:\allowbreak{}M038:\allowbreak{}M038W:\allowbreak{}335003} & \textbf{197+-22} MeV & \(u_+\): 21.97942\allowbreak{}66885127\allowbreak{}31\par \(u_-\): 22.15397\allowbreak{}03002476\allowbreak{}11\par Tipo: NONE\par Confianza: ND & COMP-MESON\par E04 / M08\\\addlinespace[3pt]
% VI-CATALOG-ROW anchuras 213
\textbf{213. K\_\allowbreak{}2(2250) WIDTH}\par K\_\allowbreak{}2(2250)\allowbreak{}0 | Kbar\_\allowbreak{}2(2250)\allowbreak{}0 | K\_\allowbreak{}2(2250)\allowbreak{}+ | K\_\allowbreak{}2(2250)\allowbreak{}-\par {\fontsize{8.1}{9.4}\selectfont PDG2026:\allowbreak{}observable:\allowbreak{}M040:\allowbreak{}M040W:\allowbreak{}336033} & \textbf{180+-30} MeV & \(u_+\): 29.99999\allowbreak{}94865718\allowbreak{}7\par \(u_-\): 29.99999\allowbreak{}94865718\allowbreak{}7\par Tipo: NONE\par Confianza: ND & COMP-MESON\par E04 / M08\\\addlinespace[3pt]
% VI-CATALOG-ROW anchuras 214
\textbf{214. p p CENTRAL PRODUCTION}\par f\_\allowbreak{}4(2300)\allowbreak{}0\par {\fontsize{8.1}{9.4}\selectfont PDG2026:\allowbreak{}observable:\allowbreak{}M041:\allowbreak{}M041W4:\allowbreak{}335245} & \textbf{250+-80} MeV & \(u_+\): 80\par \(u_-\): 80\par Tipo: NONE\par Confianza: ND & COMP-MESON\par E04 / M08\\\addlinespace[3pt]
% VI-CATALOG-ROW anchuras 215
\textbf{215. f\_\allowbreak{}2(2150) WIDTH, COMBINED MODES}\par f\_\allowbreak{}2(2150)\allowbreak{}0\par {\fontsize{8.1}{9.4}\selectfont PDG2026:\allowbreak{}observable:\allowbreak{}M042:\allowbreak{}M042W:\allowbreak{}335160} & \textbf{152+-30} MeV & \(u_+\): 30.18700\allowbreak{}23487088\allowbreak{}41\par \(u_-\): 30.18700\allowbreak{}23487088\allowbreak{}41\par Tipo: NONE\par Confianza: ND & COMP-MESON\par E04 / M08\\\addlinespace[3pt]
% VI-CATALOG-ROW anchuras 216
\textbf{216. eta eta MODE}\par f\_\allowbreak{}2(2150)\allowbreak{}0\par {\fontsize{8.1}{9.4}\selectfont PDG2026:\allowbreak{}observable:\allowbreak{}M042:\allowbreak{}M042W3:\allowbreak{}335161} & \textbf{152+-30} MeV & \(u_+\): 30.18700\allowbreak{}23487088\allowbreak{}41\par \(u_-\): 30.18700\allowbreak{}23487088\allowbreak{}41\par Tipo: NONE\par Confianza: ND & COMP-MESON\par E04 / M08\\\addlinespace[3pt]
% VI-CATALOG-ROW anchuras 217
\textbf{217. pbar p -\kern0pt{}-> pi pi}\par f\_\allowbreak{}2(2150)\allowbreak{}0\par {\fontsize{8.1}{9.4}\selectfont PDG2026:\allowbreak{}observable:\allowbreak{}M042:\allowbreak{}M042W1:\allowbreak{}335162} & \textbf{250} MeV & \(u_+\): ND\par \(u_-\): ND\par Tipo: NONE\par Confianza: ND & COMP-MESON\par E04 / M08\\\addlinespace[3pt]
% VI-CATALOG-ROW anchuras 218
\textbf{218. omega\_\allowbreak{}3(1670) WIDTH}\par omega\_\allowbreak{}3(1670)\allowbreak{}0\par {\fontsize{8.1}{9.4}\selectfont PDG2026:\allowbreak{}observable:\allowbreak{}M045:\allowbreak{}M045W:\allowbreak{}334779} & \textbf{168+-10} MeV & \(u_+\): 10.01797\allowbreak{}20379382\par \(u_-\): 10.01797\allowbreak{}20379382\par Tipo: NONE\par Confianza: ND & COMP-MESON\par E04 / M08\\\addlinespace[3pt]
% VI-CATALOG-ROW anchuras 219
\textbf{219. Upsilon(4S) WIDTH}\par Upsilon(4S)\par {\fontsize{8.1}{9.4}\selectfont PDG2026:\allowbreak{}observable:\allowbreak{}M047:\allowbreak{}M047W:\allowbreak{}347745} & \textbf{20.5+-2.5} MeV & \(u_+\): 2.473041\allowbreak{}03647041\allowbreak{}1\par \(u_-\): 2.473041\allowbreak{}03647041\allowbreak{}1\par Tipo: NONE\par Confianza: ND & COMP-MESON\par E04 / M08\\\addlinespace[3pt]
% VI-CATALOG-ROW anchuras 220
\textbf{220. Upsilon(3S) WIDTH}\par Upsilon(3S)\par {\fontsize{8.1}{9.4}\selectfont PDG2026:\allowbreak{}observable:\allowbreak{}M048:\allowbreak{}M048W:\allowbreak{}347661} & \textbf{20.32+-1.85} keV & \(u_+\): 1.850000\allowbreak{}00000000\allowbreak{}01\par \(u_-\): 1.850000\allowbreak{}00000000\allowbreak{}01\par Tipo: NONE\par Confianza: ND & COMP-MESON\par E04 / M08\\\addlinespace[3pt]
% VI-CATALOG-ROW anchuras 221
\textbf{221. Upsilon(1S) WIDTH}\par Upsilon(1S)\par {\fontsize{8.1}{9.4}\selectfont PDG2026:\allowbreak{}observable:\allowbreak{}M049:\allowbreak{}M049W:\allowbreak{}346838} & \textbf{54.02+-1.25} keV & \(u_+\): 1.25\par \(u_-\): 1.25\par Tipo: NONE\par Confianza: ND & COMP-MESON\par E04 / M08\\\addlinespace[3pt]
% VI-CATALOG-ROW anchuras 222
\textbf{222. chi\_\allowbreak{}c2(3930) WIDTH}\par chi\_\allowbreak{}c2(3930)\par {\fontsize{8.1}{9.4}\selectfont PDG2026:\allowbreak{}observable:\allowbreak{}M050:\allowbreak{}M050W:\allowbreak{}346229} & \textbf{35.2+-2.2} MeV & \(u_+\): 2.228673\allowbreak{}05911869\allowbreak{}32\par \(u_-\): 2.228673\allowbreak{}05911869\allowbreak{}32\par Tipo: NONE\par Confianza: ND & COMP-MESON\par E04 / M08\\\addlinespace[3pt]
% VI-CATALOG-ROW anchuras 223
\textbf{223. Upsilon(2S) WIDTH}\par Upsilon(2S)\par {\fontsize{8.1}{9.4}\selectfont PDG2026:\allowbreak{}observable:\allowbreak{}M052:\allowbreak{}M052W:\allowbreak{}347288} & \textbf{31.98+-2.63} keV & \(u_+\): 2.629999\allowbreak{}99999999\allowbreak{}99\par \(u_-\): 2.629999\allowbreak{}99999999\allowbreak{}99\par Tipo: NONE\par Confianza: ND & COMP-MESON\par E04 / M08\\\addlinespace[3pt]
% VI-CATALOG-ROW anchuras 224
\textbf{224. psi(3770) WIDTH}\par psi(3770)\par {\fontsize{8.1}{9.4}\selectfont PDG2026:\allowbreak{}observable:\allowbreak{}M053:\allowbreak{}M053W:\allowbreak{}345870} & \textbf{27.2+-1.0} MeV & \(u_+\): 0.959248\allowbreak{}22707886\allowbreak{}17\par \(u_-\): 0.959248\allowbreak{}22707886\allowbreak{}17\par Tipo: NONE\par Confianza: ND & COMP-MESON\par E04 / M08\\\addlinespace[3pt]
% VI-CATALOG-ROW anchuras 225
\textbf{225. phi\_\allowbreak{}3(1850) WIDTH}\par phi\_\allowbreak{}3(1850)\allowbreak{}0\par {\fontsize{8.1}{9.4}\selectfont PDG2026:\allowbreak{}observable:\allowbreak{}M054:\allowbreak{}M054W:\allowbreak{}335024} & \textbf{87+28-23} MeV & \(u_+\): 27.57879\allowbreak{}29492030\allowbreak{}59\par \(u_-\): 23.09249\allowbreak{}92747863\allowbreak{}09\par Tipo: NONE\par Confianza: ND & COMP-MESON\par E04 / M08\\\addlinespace[3pt]
% VI-CATALOG-ROW anchuras 226
\textbf{226. chi\_\allowbreak{}c1(1P) WIDTH}\par chi\_\allowbreak{}c1(1P)\par {\fontsize{8.1}{9.4}\selectfont PDG2026:\allowbreak{}observable:\allowbreak{}M055:\allowbreak{}M055W:\allowbreak{}344599} & \textbf{0.82+-0.04} MeV & \(u_+\): 0.041623\allowbreak{}41733038\allowbreak{}1482\par \(u_-\): 0.041623\allowbreak{}41733038\allowbreak{}1482\par Tipo: NONE\par Confianza: ND & COMP-MESON\par E04 / M08\\\addlinespace[3pt]
% VI-CATALOG-ROW anchuras 227
\textbf{227. chi\_\allowbreak{}c0(1P) WIDTH}\par chi\_\allowbreak{}c0(1P)\par {\fontsize{8.1}{9.4}\selectfont PDG2026:\allowbreak{}observable:\allowbreak{}M056:\allowbreak{}M056W:\allowbreak{}344305} & \textbf{12.40+-0.27} MeV & \(u_+\): 0.265077\allowbreak{}16522265\allowbreak{}817\par \(u_-\): 0.265077\allowbreak{}16522265\allowbreak{}817\par Tipo: NONE\par Confianza: ND & COMP-MESON\par E04 / M08\\\addlinespace[3pt]
% VI-CATALOG-ROW anchuras 228
\textbf{228. chi\_\allowbreak{}c2(1P) WIDTH}\par chi\_\allowbreak{}c2(1P)\par {\fontsize{8.1}{9.4}\selectfont PDG2026:\allowbreak{}observable:\allowbreak{}M057:\allowbreak{}M057W:\allowbreak{}344927} & \textbf{1.94+-0.09} MeV & \(u_+\): 0.089048\allowbreak{}41015736\allowbreak{}6226\par \(u_-\): 0.089048\allowbreak{}41015736\allowbreak{}6226\par Tipo: NONE\par Confianza: ND & COMP-MESON\par E04 / M08\\\addlinespace[3pt]
% VI-CATALOG-ROW anchuras 229
\textbf{229. pi(1300) WIDTH}\par pi(1300)\allowbreak{}0 | pi(1300)\allowbreak{}+ | pi(1300)\allowbreak{}-\par {\fontsize{8.1}{9.4}\selectfont PDG2026:\allowbreak{}observable:\allowbreak{}M058:\allowbreak{}M058W:\allowbreak{}334464} & \textbf{200 TO 600} MeV\par Presentación: 200 to 600 & \(u_+\): NA\par \(u_-\): NA\par Tipo: R\par Confianza: ND & COMP-MESON\par E04 / M08\\\addlinespace[3pt]
% VI-CATALOG-ROW anchuras 230
\textbf{230. eta\_\allowbreak{}c(2S) WIDTH}\par eta\_\allowbreak{}c(2S)\par {\fontsize{8.1}{9.4}\selectfont PDG2026:\allowbreak{}observable:\allowbreak{}M059:\allowbreak{}M059W:\allowbreak{}345229} & \textbf{11.6+-1.4} MeV & \(u_+\): 1.461405\allowbreak{}21435830\allowbreak{}21\par \(u_-\): 1.424822\allowbreak{}73720458\allowbreak{}5\par Tipo: NONE\par Confianza: ND & COMP-MESON\par E04 / M08\\\addlinespace[3pt]
% VI-CATALOG-ROW anchuras 231
\textbf{231. K\_\allowbreak{}3*(1780) WIDTH}\par K\_\allowbreak{}3\textasciicircum{}*(1780)\allowbreak{}0 | Kbar\_\allowbreak{}3\textasciicircum{}*(1780)\allowbreak{}0 | K\_\allowbreak{}3\textasciicircum{}*(1780)\allowbreak{}+ | K\_\allowbreak{}3\textasciicircum{}*(1780)\allowbreak{}-\par {\fontsize{8.1}{9.4}\selectfont PDG2026:\allowbreak{}observable:\allowbreak{}M060:\allowbreak{}M060W:\allowbreak{}335976} & \textbf{161+-17} MeV & \(u_+\): 16.77577\allowbreak{}77591419\allowbreak{}21\par \(u_-\): 17.68736\allowbreak{}93367902\allowbreak{}58\par Tipo: NONE\par Confianza: ND & COMP-MESON\par E04 / M08\\\addlinespace[3pt]
% VI-CATALOG-ROW anchuras 232
\textbf{232. D*(2007)\allowbreak{}0 WIDTH}\par D\textasciicircum{}*(2007)\allowbreak{}0 | Dbar\textasciicircum{}*(2007)\allowbreak{}0\par {\fontsize{8.1}{9.4}\selectfont PDG2026:\allowbreak{}observable:\allowbreak{}M061:\allowbreak{}M061W:\allowbreak{}337780} & \textbf{<2.1} MeV & \(u_+\): NA\par \(u_-\): NA\par Tipo: U\par Confianza: 90 & COMP-MESON\par E04 / M08\\\addlinespace[3pt]
% VI-CATALOG-ROW anchuras 233
\textbf{233. D*(2010)\allowbreak{}+- WIDTH}\par D\textasciicircum{}*(2010)\allowbreak{}+ | D\textasciicircum{}*(2010)\allowbreak{}-\par {\fontsize{8.1}{9.4}\selectfont PDG2026:\allowbreak{}observable:\allowbreak{}M062:\allowbreak{}M062W:\allowbreak{}337799} & \textbf{83.4+-1.8} keV & \(u_+\): 1.837671\allowbreak{}39449678\allowbreak{}71\par \(u_-\): 1.837671\allowbreak{}39449678\allowbreak{}71\par Tipo: NONE\par Confianza: ND & COMP-MESON\par E04 / M08\\\addlinespace[3pt]
% VI-CATALOG-ROW anchuras 234
\textbf{234. K\_\allowbreak{}1(1400) WIDTH}\par K\_\allowbreak{}1(1400)\allowbreak{}0 | Kbar\_\allowbreak{}1(1400)\allowbreak{}0 | K\_\allowbreak{}1(1400)\allowbreak{}+ | K\_\allowbreak{}1(1400)\allowbreak{}-\par {\fontsize{8.1}{9.4}\selectfont PDG2026:\allowbreak{}observable:\allowbreak{}M064:\allowbreak{}M064W:\allowbreak{}335846} & \textbf{174+-13} MeV & \(u_+\): 12.61112\allowbreak{}36551536\allowbreak{}99\par \(u_-\): 12.60116\allowbreak{}09689773\allowbreak{}61\par Tipo: NONE\par Confianza: ND & COMP-MESON\par E04 / M08\\\addlinespace[3pt]
% VI-CATALOG-ROW anchuras 235
\textbf{235. eta rho0 AND pi+ pi- MODES}\par rho(1700)\allowbreak{}0 | rho(1700)\allowbreak{}+ | rho(1700)\allowbreak{}-\par {\fontsize{8.1}{9.4}\selectfont PDG2026:\allowbreak{}observable:\allowbreak{}M065:\allowbreak{}M065W0:\allowbreak{}334881} & \textbf{250+-100} MeV & \(u_+\): 100\par \(u_-\): 100\par Tipo: NONE\par Confianza: ND & COMP-MESON\par E04 / M08\\\addlinespace[3pt]
% VI-CATALOG-ROW anchuras 236
\textbf{236. e+ e- PRODUCTION}\par phi(1680)\allowbreak{}0\par {\fontsize{8.1}{9.4}\selectfont PDG2026:\allowbreak{}observable:\allowbreak{}M067:\allowbreak{}M067W1:\allowbreak{}334826} & \textbf{150+-50} MeV & \(u_+\): 50\par \(u_-\): 50\par Tipo: NONE\par Confianza: ND & COMP-MESON\par E04 / M08\\\addlinespace[3pt]
% VI-CATALOG-ROW anchuras 237
\textbf{237. f\_\allowbreak{}0(1710) Breit-Wigner WIDTH}\par f\_\allowbreak{}0(1710)\allowbreak{}0\par {\fontsize{8.1}{9.4}\selectfont PDG2026:\allowbreak{}observable:\allowbreak{}M068:\allowbreak{}M068W:\allowbreak{}334928} & \textbf{149+-10} MeV & \(u_+\): 10.92532\allowbreak{}40627161\allowbreak{}7\par \(u_-\): 9.891723\allowbreak{}41764963\allowbreak{}65\par Tipo: NONE\par Confianza: ND & COMP-MESON\par E04 / M08\\\addlinespace[3pt]
% VI-CATALOG-ROW anchuras 238
\textbf{238. J/\allowbreak{}psi(1S) WIDTH}\par J/\allowbreak{}psi(1S)\par {\fontsize{8.1}{9.4}\selectfont PDG2026:\allowbreak{}observable:\allowbreak{}M070:\allowbreak{}M070W:\allowbreak{}343506} & \textbf{92.6+-1.7} keV & \(u_+\): 1.718251\allowbreak{}16057207\allowbreak{}09\par \(u_-\): 1.711913\allowbreak{}28371000\allowbreak{}99\par Tipo: NONE\par Confianza: ND & COMP-MESON\par E04 / M08\\\addlinespace[3pt]
% VI-CATALOG-ROW anchuras 239
\textbf{239. psi(2S) WIDTH}\par psi(2S)\par {\fontsize{8.1}{9.4}\selectfont PDG2026:\allowbreak{}observable:\allowbreak{}M071:\allowbreak{}M071W:\allowbreak{}345302} & \textbf{293+-9} keV & \(u_+\): 9.121617\allowbreak{}13271473\allowbreak{}04\par \(u_-\): 9.121617\allowbreak{}13271473\allowbreak{}04\par Tipo: NONE\par Confianza: ND & COMP-MESON\par E04 / M08\\\addlinespace[3pt]
% VI-CATALOG-ROW anchuras 240
\textbf{240. psi(4040) WIDTH}\par psi(4040)\par {\fontsize{8.1}{9.4}\selectfont PDG2026:\allowbreak{}observable:\allowbreak{}M072:\allowbreak{}M072W:\allowbreak{}346250} & \textbf{84+-12} MeV & \(u_+\): 12.30000\allowbreak{}00000000\allowbreak{}01\par \(u_-\): 12.30000\allowbreak{}00000000\allowbreak{}01\par Tipo: NONE\par Confianza: ND & COMP-MESON\par E04 / M08\\\addlinespace[3pt]
% VI-CATALOG-ROW anchuras 241
\textbf{241. psi(4415) WIDTH}\par psi(4415)\par {\fontsize{8.1}{9.4}\selectfont PDG2026:\allowbreak{}observable:\allowbreak{}M073:\allowbreak{}M073W:\allowbreak{}346685} & \textbf{110+-13} MeV & \(u_+\): 13.28782\allowbreak{}78993015\allowbreak{}7\par \(u_-\): 13.28782\allowbreak{}78993015\allowbreak{}7\par Tipo: NONE\par Confianza: ND & COMP-MESON\par E04 / M08\\\addlinespace[3pt]
% VI-CATALOG-ROW anchuras 242
\textbf{242. psi(4230) WIDTH}\par psi(4230)\par {\fontsize{8.1}{9.4}\selectfont PDG2026:\allowbreak{}observable:\allowbreak{}M074:\allowbreak{}M074W:\allowbreak{}346461} & \textbf{51+-7} MeV & \(u_+\): 6.997290\allowbreak{}38980041\allowbreak{}02\par \(u_-\): 6.882414\allowbreak{}90298323\allowbreak{}64\par Tipo: NONE\par Confianza: ND & COMP-MESON\par E04 / M08\\\addlinespace[3pt]
% VI-CATALOG-ROW anchuras 243
\textbf{243. pi(1800) WIDTH}\par pi(1800)\allowbreak{}0 | pi(1800)\allowbreak{}+ | pi(1800)\allowbreak{}-\par {\fontsize{8.1}{9.4}\selectfont PDG2026:\allowbreak{}observable:\allowbreak{}M075:\allowbreak{}M075W:\allowbreak{}334967} & \textbf{215+7-8} MeV & \(u_+\): 6.706577\allowbreak{}42096917\allowbreak{}8\par \(u_-\): 8.199899\allowbreak{}30002816\allowbreak{}08\par Tipo: NONE\par Confianza: ND & COMP-MESON\par E04 / M08\\\addlinespace[3pt]
% VI-CATALOG-ROW anchuras 244
\textbf{244. f\_\allowbreak{}J(2220) WIDTH}\par f\_\allowbreak{}J(2220)\allowbreak{}0\par {\fontsize{8.1}{9.4}\selectfont PDG2026:\allowbreak{}observable:\allowbreak{}M082:\allowbreak{}M082W:\allowbreak{}335204} & \textbf{23+8-7} MeV & \(u_+\): 8.290723\allowbreak{}39682046\allowbreak{}79\par \(u_-\): 6.995398\allowbreak{}15418370\allowbreak{}2\par Tipo: NONE\par Confianza: ND & COMP-MESON\par E04 / M08\\\addlinespace[3pt]
% VI-CATALOG-ROW anchuras 245
\textbf{245. f\_\allowbreak{}1(1510) WIDTH}\par f\_\allowbreak{}1(1510)\allowbreak{}0\par {\fontsize{8.1}{9.4}\selectfont PDG2026:\allowbreak{}observable:\allowbreak{}M084:\allowbreak{}M084W:\allowbreak{}334659} & \textbf{73+-25} MeV & \(u_+\): 25.12449\allowbreak{}04562855\allowbreak{}5\par \(u_-\): 25.12449\allowbreak{}04562855\allowbreak{}5\par Tipo: NONE\par Confianza: ND & COMP-MESON\par E04 / M08\\\addlinespace[3pt]
% VI-CATALOG-ROW anchuras 246
\textbf{246. X(1835) WIDTH}\par X(1835)\allowbreak{}0\par {\fontsize{8.1}{9.4}\selectfont PDG2026:\allowbreak{}observable:\allowbreak{}M085:\allowbreak{}M085W:\allowbreak{}335012} & \textbf{130+-50} MeV & \(u_+\): 46.52142\allowbreak{}48579634\allowbreak{}67\par \(u_-\): 48.80576\allowbreak{}18821165\allowbreak{}28\par Tipo: NONE\par Confianza: ND & COMP-MESON\par E04 / M08\\\addlinespace[3pt]
% VI-CATALOG-ROW anchuras 247
\textbf{247. f\_\allowbreak{}6(2510) WIDTH}\par f\_\allowbreak{}6(2510)\allowbreak{}0\par {\fontsize{8.1}{9.4}\selectfont PDG2026:\allowbreak{}observable:\allowbreak{}M089:\allowbreak{}M089W:\allowbreak{}335278} & \textbf{260+-40} MeV & \(u_+\): 42.42640\allowbreak{}68711928\allowbreak{}53\par \(u_-\): 42.42640\allowbreak{}68711928\allowbreak{}53\par Tipo: NONE\par Confianza: ND & COMP-MESON\par E04 / M08\\\addlinespace[3pt]
% VI-CATALOG-ROW anchuras 248
\textbf{248. K\_\allowbreak{}3(2320) WIDTH}\par K\_\allowbreak{}3(2320)\allowbreak{}0 | Kbar\_\allowbreak{}3(2320)\allowbreak{}0 | K\_\allowbreak{}3(2320)\allowbreak{}+ | K\_\allowbreak{}3(2320)\allowbreak{}-\par {\fontsize{8.1}{9.4}\selectfont PDG2026:\allowbreak{}observable:\allowbreak{}M090:\allowbreak{}M090W:\allowbreak{}336035} & \textbf{150+-30} MeV & \(u_+\): 30\par \(u_-\): 30\par Tipo: NONE\par Confianza: ND & COMP-MESON\par E04 / M08\\\addlinespace[3pt]
% VI-CATALOG-ROW anchuras 249
\textbf{249. Upsilon(10860) WIDTH}\par Upsilon(10860)\par {\fontsize{8.1}{9.4}\selectfont PDG2026:\allowbreak{}observable:\allowbreak{}M092:\allowbreak{}M092W:\allowbreak{}347805} & \textbf{37+-4} MeV & \(u_+\): 4.266562\allowbreak{}07180538\allowbreak{}96\par \(u_-\): 3.976736\allowbreak{}53123974\allowbreak{}89\par Tipo: NONE\par Confianza: ND & COMP-MESON\par E04 / M08\\\addlinespace[3pt]
% VI-CATALOG-ROW anchuras 250
\textbf{250. Upsilon(11020) WIDTH}\par Upsilon(11020)\par {\fontsize{8.1}{9.4}\selectfont PDG2026:\allowbreak{}observable:\allowbreak{}M093:\allowbreak{}M093W:\allowbreak{}347906} & \textbf{24+8-6} MeV & \(u_+\): 7.707706\allowbreak{}56457082\allowbreak{}83\par \(u_-\): 6.457155\allowbreak{}82578219\allowbreak{}66\par Tipo: NONE\par Confianza: ND & COMP-MESON\par E04 / M08\\\addlinespace[3pt]
% VI-CATALOG-ROW anchuras 251
\textbf{251. K*(1410) WIDTH}\par K\textasciicircum{}*(1410)\allowbreak{}0 | Kbar\textasciicircum{}*(1410)\allowbreak{}0 | K\textasciicircum{}*(1410)\allowbreak{}+ | K\textasciicircum{}*(1410)\allowbreak{}-\par {\fontsize{8.1}{9.4}\selectfont PDG2026:\allowbreak{}observable:\allowbreak{}M094:\allowbreak{}M094W:\allowbreak{}335866} & \textbf{232+-21} MeV & \(u_+\): 21.38237\allowbreak{}21605395\allowbreak{}6\par \(u_-\): 21.38237\allowbreak{}21605395\allowbreak{}6\par Tipo: NONE\par Confianza: ND & COMP-MESON\par E04 / M08\\\addlinespace[3pt]
% VI-CATALOG-ROW anchuras 252
\textbf{252. K*(1680) WIDTH}\par K\textasciicircum{}*(1680)\allowbreak{}0 | Kbar\textasciicircum{}*(1680)\allowbreak{}0 | K\textasciicircum{}*(1680)\allowbreak{}+ | K\textasciicircum{}*(1680)\allowbreak{}-\par {\fontsize{8.1}{9.4}\selectfont PDG2026:\allowbreak{}observable:\allowbreak{}M095:\allowbreak{}M095W:\allowbreak{}335947} & \textbf{320+-110} MeV & \(u_+\): 107.3383\allowbreak{}45980695\allowbreak{}71\par \(u_-\): 107.8075\allowbreak{}66482925\allowbreak{}91\par Tipo: NONE\par Confianza: ND & COMP-MESON\par E04 / M08\\\addlinespace[3pt]
% VI-CATALOG-ROW anchuras 253
\textbf{253. K\_\allowbreak{}5*(2380) WIDTH}\par K\_\allowbreak{}5\textasciicircum{}*(2380)\allowbreak{}0 | Kbar\_\allowbreak{}5\textasciicircum{}*(2380)\allowbreak{}0 | K\_\allowbreak{}5\textasciicircum{}*(2380)\allowbreak{}+ | K\_\allowbreak{}5\textasciicircum{}*(2380)\allowbreak{}-\par {\fontsize{8.1}{9.4}\selectfont PDG2026:\allowbreak{}observable:\allowbreak{}M098:\allowbreak{}M098W:\allowbreak{}336037} & \textbf{180+-50} MeV & \(u_+\): 48.91829\allowbreak{}92345400\allowbreak{}36\par \(u_-\): 48.91829\allowbreak{}92345400\allowbreak{}36\par Tipo: NONE\par Confianza: ND & COMP-MESON\par E04 / M08\\\addlinespace[3pt]
% VI-CATALOG-ROW anchuras 254
\textbf{254. K\_\allowbreak{}1(1650) WIDTH}\par K\_\allowbreak{}1(1650)\allowbreak{}0 | Kbar\_\allowbreak{}1(1650)\allowbreak{}0 | K\_\allowbreak{}1(1650)\allowbreak{}+ | K\_\allowbreak{}1(1650)\allowbreak{}-\par {\fontsize{8.1}{9.4}\selectfont PDG2026:\allowbreak{}observable:\allowbreak{}M099:\allowbreak{}M099W:\allowbreak{}335945} & \textbf{150+-50} MeV & \(u_+\): 50\par \(u_-\): 50\par Tipo: NONE\par Confianza: ND & COMP-MESON\par E04 / M08\\\addlinespace[3pt]
% VI-CATALOG-ROW anchuras 255
\textbf{255. eta\_\allowbreak{}2(1870) WIDTH}\par eta\_\allowbreak{}2(1870)\allowbreak{}0\par {\fontsize{8.1}{9.4}\selectfont PDG2026:\allowbreak{}observable:\allowbreak{}M101:\allowbreak{}M101W:\allowbreak{}335033} & \textbf{225+-14} MeV & \(u_+\): 14.06500\allowbreak{}75480096\allowbreak{}5\par \(u_-\): 14.06500\allowbreak{}75480096\allowbreak{}5\par Tipo: NONE\par Confianza: ND & COMP-MESON\par E04 / M08\\\addlinespace[3pt]
% VI-CATALOG-ROW anchuras 256
\textbf{256. phi(2170) WIDTH}\par phi(2170)\allowbreak{}0 | phi(2170)\allowbreak{}+ | phi(2170)\allowbreak{}-\par {\fontsize{8.1}{9.4}\selectfont PDG2026:\allowbreak{}observable:\allowbreak{}M103:\allowbreak{}M103W:\allowbreak{}335181} & \textbf{88+26-21} MeV & \(u_+\): 25.84561\allowbreak{}82632216\allowbreak{}7\par \(u_-\): 20.63586\allowbreak{}51464737\allowbreak{}9\par Tipo: NONE\par Confianza: ND & COMP-MESON\par E04 / M08\\\addlinespace[3pt]
% VI-CATALOG-ROW anchuras 257
\textbf{257. K\_\allowbreak{}2*(1980) WIDTH}\par K\_\allowbreak{}2\textasciicircum{}*(1980)\allowbreak{}0 | Kbar\_\allowbreak{}2\textasciicircum{}*(1980)\allowbreak{}0 | K\_\allowbreak{}2\textasciicircum{}*(1980)\allowbreak{}+ | K\_\allowbreak{}2\textasciicircum{}*(1980)\allowbreak{}-\par {\fontsize{8.1}{9.4}\selectfont PDG2026:\allowbreak{}observable:\allowbreak{}M104:\allowbreak{}M104W:\allowbreak{}336004} & \textbf{348+50-30} MeV & \(u_+\): 47.99501\allowbreak{}58521243\allowbreak{}11\par \(u_-\): 29.82270\allowbreak{}03927434\allowbreak{}99\par Tipo: NONE\par Confianza: ND & COMP-MESON\par E04 / M08\\\addlinespace[3pt]
% VI-CATALOG-ROW anchuras 258
\textbf{258. rho(1450) WIDTH}\par rho(1450)\allowbreak{}0 | rho(1450)\allowbreak{}+ | rho(1450)\allowbreak{}-\par {\fontsize{8.1}{9.4}\selectfont PDG2026:\allowbreak{}observable:\allowbreak{}M105:\allowbreak{}M105W0:\allowbreak{}334593} & \textbf{400 +-60} MeV\par Presentación: 400+-60 & \(u_+\): 60\par \(u_-\): 60\par Tipo: NONE\par Confianza: ND & COMP-MESON\par E04 / M08\\\addlinespace[3pt]
% VI-CATALOG-ROW anchuras 259
\textbf{259. f\_\allowbreak{}2(2010) WIDTH}\par f\_\allowbreak{}2(2010)\allowbreak{}0\par {\fontsize{8.1}{9.4}\selectfont PDG2026:\allowbreak{}observable:\allowbreak{}M106:\allowbreak{}M106W:\allowbreak{}335113} & \textbf{200+-60} MeV & \(u_+\): 67\par \(u_-\): 62\par Tipo: NONE\par Confianza: ND & COMP-MESON\par E04 / M08\\\addlinespace[3pt]
% VI-CATALOG-ROW anchuras 260
\textbf{260. f\_\allowbreak{}2(2300) WIDTH}\par f\_\allowbreak{}2(2300)\allowbreak{}0\par {\fontsize{8.1}{9.4}\selectfont PDG2026:\allowbreak{}observable:\allowbreak{}M107:\allowbreak{}M107W:\allowbreak{}335235} & \textbf{150+-40} MeV & \(u_+\): 41\par \(u_-\): 41\par Tipo: NONE\par Confianza: ND & COMP-MESON\par E04 / M08\\\addlinespace[3pt]
% VI-CATALOG-ROW anchuras 261
\textbf{261. f\_\allowbreak{}2(2340) WIDTH}\par f\_\allowbreak{}2(2340)\allowbreak{}0\par {\fontsize{8.1}{9.4}\selectfont PDG2026:\allowbreak{}observable:\allowbreak{}M108:\allowbreak{}M108W:\allowbreak{}335253} & \textbf{331+27-18} MeV & \(u_+\): 27.40609\allowbreak{}49578306\allowbreak{}5\par \(u_-\): 17.59908\allowbreak{}97768665\allowbreak{}38\par Tipo: NONE\par Confianza: ND & COMP-MESON\par E04 / M08\\\addlinespace[3pt]
% VI-CATALOG-ROW anchuras 262
\textbf{262. h\_\allowbreak{}1(1415) WIDTH}\par h\_\allowbreak{}1(1415)\allowbreak{}0\par {\fontsize{8.1}{9.4}\selectfont PDG2026:\allowbreak{}observable:\allowbreak{}M109:\allowbreak{}M109W:\allowbreak{}334553} & \textbf{78+-11} MeV & \(u_+\): 10.89476\allowbreak{}84184199\allowbreak{}31\par \(u_-\): 10.35476\allowbreak{}10268449\allowbreak{}6\par Tipo: NONE\par Confianza: ND & COMP-MESON\par E04 / M08\\\addlinespace[3pt]
% VI-CATALOG-ROW anchuras 263
\textbf{263. f\_\allowbreak{}0(2200) WIDTH}\par f\_\allowbreak{}0(2200)\allowbreak{}0\par {\fontsize{8.1}{9.4}\selectfont PDG2026:\allowbreak{}observable:\allowbreak{}M112:\allowbreak{}M112W:\allowbreak{}335200} & \textbf{210+-40} MeV & \(u_+\): 41.63865\allowbreak{}65682062\allowbreak{}37\par \(u_-\): 42.17325\allowbreak{}16055459\allowbreak{}43\par Tipo: NONE\par Confianza: ND & COMP-MESON\par E04 / M08\\\addlinespace[3pt]
% VI-CATALOG-ROW anchuras 264
\textbf{264. eta(1760) WIDTH}\par eta(1760)\allowbreak{}0\par {\fontsize{8.1}{9.4}\selectfont PDG2026:\allowbreak{}observable:\allowbreak{}M114:\allowbreak{}M114W:\allowbreak{}334951} & \textbf{240+-30} MeV & \(u_+\): 30.76700\allowbreak{}00039266\allowbreak{}5\par \(u_-\): 28.81984\allowbreak{}35999457\allowbreak{}98\par Tipo: NONE\par Confianza: ND & COMP-MESON\par E04 / M08\\\addlinespace[3pt]
% VI-CATALOG-ROW anchuras 265
\textbf{265. eta(2225) WIDTH}\par eta(2225)\allowbreak{}0\par {\fontsize{8.1}{9.4}\selectfont PDG2026:\allowbreak{}observable:\allowbreak{}M115:\allowbreak{}M115W:\allowbreak{}335233} & \textbf{185+40-20} MeV & \(u_+\): 36.89403\allowbreak{}65233648\allowbreak{}68\par \(u_-\): 19.60880\allowbreak{}60392476\par Tipo: NONE\par Confianza: ND & COMP-MESON\par E04 / M08\\\addlinespace[3pt]
% VI-CATALOG-ROW anchuras 266
\textbf{266. f\_\allowbreak{}2(1640) WIDTH}\par f\_\allowbreak{}2(1640)\allowbreak{}0\par {\fontsize{8.1}{9.4}\selectfont PDG2026:\allowbreak{}observable:\allowbreak{}M117:\allowbreak{}M117W:\allowbreak{}334753} & \textbf{100+60-40} MeV & \(u_+\): 55.00727\allowbreak{}09658098\allowbreak{}73\par \(u_-\): 40.99999\allowbreak{}90454338\allowbreak{}03\par Tipo: NONE\par Confianza: ND & COMP-MESON\par E04 / M08\\\addlinespace[3pt]
% VI-CATALOG-ROW anchuras 267
\textbf{267. D\_\allowbreak{}s1(2536)\allowbreak{}+- WIDTH}\par D\_\allowbreak{}s1(2536)\allowbreak{}+ | D\_\allowbreak{}s1(2536)\allowbreak{}-\par {\fontsize{8.1}{9.4}\selectfont PDG2026:\allowbreak{}observable:\allowbreak{}M121:\allowbreak{}M121W:\allowbreak{}338489} & \textbf{0.92+-0.05} MeV & \(u_+\): 0.049965\allowbreak{}31390494\allowbreak{}7829\par \(u_-\): 0.049965\allowbreak{}31390494\allowbreak{}7829\par Tipo: NONE\par Confianza: ND & COMP-MESON\par E04 / M08\\\addlinespace[3pt]
% VI-CATALOG-ROW anchuras 268
\textbf{268. f\_\allowbreak{}2(1565) WIDTH}\par f\_\allowbreak{}2(1565)\allowbreak{}0\par {\fontsize{8.1}{9.4}\selectfont PDG2026:\allowbreak{}observable:\allowbreak{}M123:\allowbreak{}M123W:\allowbreak{}334704} & \textbf{132+-23} MeV & \(u_+\): 22.49999\allowbreak{}92121598\allowbreak{}28\par \(u_-\): 22.49999\allowbreak{}92121598\allowbreak{}28\par Tipo: NONE\par Confianza: ND & COMP-MESON\par E04 / M08\\\addlinespace[3pt]
% VI-CATALOG-ROW anchuras 269
\textbf{269. omega(1420) WIDTH}\par omega(1420)\allowbreak{}0\par {\fontsize{8.1}{9.4}\selectfont PDG2026:\allowbreak{}observable:\allowbreak{}M125:\allowbreak{}M125W:\allowbreak{}334568} & \textbf{290+-190} MeV & \(u_+\): 190\par \(u_-\): 190\par Tipo: NONE\par Confianza: ND & COMP-MESON\par E04 / M08\\\addlinespace[3pt]
% VI-CATALOG-ROW anchuras 270
\textbf{270. omega(1650) WIDTH}\par omega(1650)\allowbreak{}0\par {\fontsize{8.1}{9.4}\selectfont PDG2026:\allowbreak{}observable:\allowbreak{}M126:\allowbreak{}M126W:\allowbreak{}334769} & \textbf{150+-50} MeV & \(u_+\): 50\par \(u_-\): 50\par Tipo: NONE\par Confianza: ND & COMP-MESON\par E04 / M08\\\addlinespace[3pt]
% VI-CATALOG-ROW anchuras 271
\textbf{271. K\_\allowbreak{}0*(1950) WIDTH}\par K\_\allowbreak{}0\textasciicircum{}*(1950)\allowbreak{}0 | Kbar\_\allowbreak{}0\textasciicircum{}*(1950)\allowbreak{}0 | K\_\allowbreak{}0\textasciicircum{}*(1950)\allowbreak{}+ | K\_\allowbreak{}0\textasciicircum{}*(1950)\allowbreak{}-\par {\fontsize{8.1}{9.4}\selectfont PDG2026:\allowbreak{}observable:\allowbreak{}M134:\allowbreak{}M134W:\allowbreak{}336000} & \textbf{170+-50} MeV & \(u_+\): 49.97763\allowbreak{}46249792\allowbreak{}67\par \(u_-\): 49.97763\allowbreak{}46249792\allowbreak{}67\par Tipo: NONE\par Confianza: ND & COMP-MESON\par E04 / M08\\\addlinespace[3pt]
% VI-CATALOG-ROW anchuras 272
\textbf{272. f\_\allowbreak{}2(1950) WIDTH}\par f\_\allowbreak{}2(1950)\allowbreak{}0\par {\fontsize{8.1}{9.4}\selectfont PDG2026:\allowbreak{}observable:\allowbreak{}M135:\allowbreak{}M135W:\allowbreak{}335076} & \textbf{464+-24} MeV & \(u_+\): 24.23360\allowbreak{}83935200\allowbreak{}51\par \(u_-\): 23.85817\allowbreak{}51115928\allowbreak{}71\par Tipo: NONE\par Confianza: ND & COMP-MESON\par E04 / M08\\\addlinespace[3pt]
% VI-CATALOG-ROW anchuras 273
\textbf{273. f\_\allowbreak{}2(1910) omega omega MODE}\par f\_\allowbreak{}2(1910)\allowbreak{}0\par {\fontsize{8.1}{9.4}\selectfont PDG2026:\allowbreak{}observable:\allowbreak{}M142:\allowbreak{}M142W2:\allowbreak{}335059} & \textbf{167+-21} MeV & \(u_+\): 20.59321\allowbreak{}07794648\allowbreak{}79\par \(u_-\): 20.59321\allowbreak{}07794648\allowbreak{}79\par Tipo: NONE\par Confianza: ND & COMP-MESON\par E04 / M08\\\addlinespace[3pt]
% VI-CATALOG-ROW anchuras 274
\textbf{274. f\_\allowbreak{}2(1910) eta eta' MODE}\par f\_\allowbreak{}2(1910)\allowbreak{}0\par {\fontsize{8.1}{9.4}\selectfont PDG2026:\allowbreak{}observable:\allowbreak{}M142:\allowbreak{}M142W3:\allowbreak{}335060} & \textbf{140+-40} MeV & \(u_+\): 41\par \(u_-\): 41\par Tipo: NONE\par Confianza: ND & COMP-MESON\par E04 / M08\\\addlinespace[3pt]
% VI-CATALOG-ROW anchuras 275
\textbf{275. h\_\allowbreak{}c(1P) WIDTH}\par h\_\allowbreak{}c(1P)\par {\fontsize{8.1}{9.4}\selectfont PDG2026:\allowbreak{}observable:\allowbreak{}M144:\allowbreak{}M144W:\allowbreak{}344852} & \textbf{0.78+-0.28} MeV & \(u_+\): 0.295465\allowbreak{}73405388\allowbreak{}321\par \(u_-\): 0.268328\allowbreak{}15729997\allowbreak{}478\par Tipo: NONE\par Confianza: ND & COMP-MESON\par E04 / M08\\\addlinespace[3pt]
% VI-CATALOG-ROW anchuras 276
\textbf{276. K\_\allowbreak{}2(1820) WIDTH}\par K\_\allowbreak{}2(1820)\allowbreak{}0 | Kbar\_\allowbreak{}2(1820)\allowbreak{}0 | K\_\allowbreak{}2(1820)\allowbreak{}+ | K\_\allowbreak{}2(1820)\allowbreak{}-\par {\fontsize{8.1}{9.4}\selectfont PDG2026:\allowbreak{}observable:\allowbreak{}M146:\allowbreak{}M146W:\allowbreak{}335991} & \textbf{264+-34} MeV & \(u_+\): 33.91116\allowbreak{}28155620\allowbreak{}37\par \(u_-\): 33.58515\allowbreak{}41639118\allowbreak{}88\par Tipo: NONE\par Confianza: ND & COMP-MESON\par E04 / M08\\\addlinespace[3pt]
% VI-CATALOG-ROW anchuras 277
\textbf{277. Full width (Breit-Wigner)}\par f\_\allowbreak{}0(1370)\allowbreak{}0\par {\fontsize{8.1}{9.4}\selectfont PDG2026:\allowbreak{}observable:\allowbreak{}M147:\allowbreak{}M147W:\allowbreak{}334507} & \textbf{200 to 500} MeV & \(u_+\): NA\par \(u_-\): NA\par Tipo: R\par Confianza: ND & COMP-MESON\par E04 / M08\\\addlinespace[3pt]
% VI-CATALOG-ROW anchuras 278
\textbf{278. D\_\allowbreak{}s2*(2573) WIDTH}\par D\_\allowbreak{}s2\textasciicircum{}*(2573)\allowbreak{}+ | D\_\allowbreak{}s2\textasciicircum{}*(2573)\allowbreak{}-\par {\fontsize{8.1}{9.4}\selectfont PDG2026:\allowbreak{}observable:\allowbreak{}M148:\allowbreak{}M148W:\allowbreak{}338511} & \textbf{16.9+-0.7} MeV & \(u_+\): 0.739720\allowbreak{}13426443\allowbreak{}928\par \(u_-\): 0.737586\allowbreak{}53054042\allowbreak{}789\par Tipo: NONE\par Confianza: ND & COMP-MESON\par E04 / M08\\\addlinespace[3pt]
% VI-CATALOG-ROW anchuras 279
\textbf{279. a\_\allowbreak{}0(1450) WIDTH}\par a\_\allowbreak{}0(1450)\allowbreak{}0 | a\_\allowbreak{}0(1450)\allowbreak{}+ | a\_\allowbreak{}0(1450)\allowbreak{}-\par {\fontsize{8.1}{9.4}\selectfont PDG2026:\allowbreak{}observable:\allowbreak{}M149:\allowbreak{}M149W:\allowbreak{}334580} & \textbf{258+-14} MeV & \(u_+\): 14.00000\allowbreak{}00989268\allowbreak{}51\par \(u_-\): 14.00000\allowbreak{}00989268\allowbreak{}51\par Tipo: NONE\par Confianza: ND & COMP-MESON\par E04 / M08\\\addlinespace[3pt]
% VI-CATALOG-ROW anchuras 280
\textbf{280. B\_\allowbreak{}J*(5732) WIDTH}\par B\_\allowbreak{}J\textasciicircum{}*(5732)\allowbreak{}0 | Bbar\_\allowbreak{}J\textasciicircum{}*(5732)\allowbreak{}0 | B\_\allowbreak{}J\textasciicircum{}*(5732)\allowbreak{}+ | B\_\allowbreak{}J\textasciicircum{}*(5732)\allowbreak{}-\par {\fontsize{8.1}{9.4}\selectfont PDG2026:\allowbreak{}observable:\allowbreak{}M151:\allowbreak{}M151W:\allowbreak{}342552} & \textbf{128+-18} MeV & \(u_+\): 18.22215\allowbreak{}84567671\allowbreak{}19\par \(u_-\): 18.22215\allowbreak{}84567671\allowbreak{}19\par Tipo: NONE\par Confianza: ND & COMP-MESON\par E04 / M08\\\addlinespace[3pt]
% VI-CATALOG-ROW anchuras 281
\textbf{281. f\_\allowbreak{}0(1500) WIDTH}\par f\_\allowbreak{}0(1500)\allowbreak{}0\par {\fontsize{8.1}{9.4}\selectfont PDG2026:\allowbreak{}observable:\allowbreak{}M152:\allowbreak{}M152W:\allowbreak{}334629} & \textbf{108+-33} MeV & \(u_+\): 33\par \(u_-\): 33\par Tipo: NONE\par Confianza: ND & COMP-MESON\par E04 / M08\\\addlinespace[3pt]
% VI-CATALOG-ROW anchuras 282
\textbf{282. B\_\allowbreak{}sJ*(5850) WIDTH}\par B\_\allowbreak{}sJ\textasciicircum{}*(5850)\allowbreak{}0 | Bbar\_\allowbreak{}sJ\textasciicircum{}*(5850)\allowbreak{}0\par {\fontsize{8.1}{9.4}\selectfont PDG2026:\allowbreak{}observable:\allowbreak{}M153:\allowbreak{}M153W:\allowbreak{}343152} & \textbf{47+-22} MeV & \(u_+\): 22\par \(u_-\): 22\par Tipo: NONE\par Confianza: ND & COMP-MESON\par E04 / M08\\\addlinespace[3pt]
% VI-CATALOG-ROW anchuras 283
\textbf{283. eta\_\allowbreak{}2(1645) WIDTH}\par eta\_\allowbreak{}2(1645)\allowbreak{}0\par {\fontsize{8.1}{9.4}\selectfont PDG2026:\allowbreak{}observable:\allowbreak{}M154:\allowbreak{}M154W:\allowbreak{}334759} & \textbf{181+-11} MeV & \(u_+\): 10.78574\allowbreak{}47992021\allowbreak{}21\par \(u_-\): 10.49166\allowbreak{}41064320\allowbreak{}9\par Tipo: NONE\par Confianza: ND & COMP-MESON\par E04 / M08\\\addlinespace[3pt]
% VI-CATALOG-ROW anchuras 284
\textbf{284. f\_\allowbreak{}0(2020) WIDTH}\par f\_\allowbreak{}0(2020)\allowbreak{}0\par {\fontsize{8.1}{9.4}\selectfont PDG2026:\allowbreak{}observable:\allowbreak{}M156:\allowbreak{}M156W:\allowbreak{}335121} & \textbf{440+-50} MeV & \(u_+\): 46.17358\allowbreak{}55224607\allowbreak{}83\par \(u_-\): 49.16299\allowbreak{}42131274\allowbreak{}19\par Tipo: NONE\par Confianza: ND & COMP-MESON\par E04 / M08\\\addlinespace[3pt]
% VI-CATALOG-ROW anchuras 285
\textbf{285. D*(2640)\allowbreak{}+- WIDTH}\par D\textasciicircum{}*(2640)\allowbreak{}0 | Dbar\textasciicircum{}*(2640)\allowbreak{}0 | D\textasciicircum{}*(2640)\allowbreak{}+ | D\textasciicircum{}*(2640)\allowbreak{}-\par {\fontsize{8.1}{9.4}\selectfont PDG2026:\allowbreak{}observable:\allowbreak{}M158:\allowbreak{}M158W:\allowbreak{}337855} & \textbf{<15} MeV & \(u_+\): NA\par \(u_-\): NA\par Tipo: U\par Confianza: 95 & COMP-MESON\par E04 / M08\\\addlinespace[3pt]
% VI-CATALOG-ROW anchuras 286
\textbf{286. chi\_\allowbreak{}c0(3915) WIDTH}\par chi\_\allowbreak{}c0(3915)\par {\fontsize{8.1}{9.4}\selectfont PDG2026:\allowbreak{}observable:\allowbreak{}M159:\allowbreak{}M159W:\allowbreak{}346205} & \textbf{20+-4} MeV & \(u_+\): 3.923709\allowbreak{}31837659\allowbreak{}9\par \(u_-\): 3.848573\allowbreak{}27141186\allowbreak{}99\par Tipo: NONE\par Confianza: ND & COMP-MESON\par E04 / M08\\\addlinespace[3pt]
% VI-CATALOG-ROW anchuras 287
\textbf{287. K(1630) WIDTH}\par K(1630)\allowbreak{}0 | Kbar(1630)\allowbreak{}0 | K(1630)\allowbreak{}+ | K(1630)\allowbreak{}-\par {\fontsize{8.1}{9.4}\selectfont PDG2026:\allowbreak{}observable:\allowbreak{}M160:\allowbreak{}M160W:\allowbreak{}335943} & \textbf{16+19-16} MeV & \(u_+\): 19\par \(u_-\): 16\par Tipo: NONE\par Confianza: ND & COMP-MESON\par E04 / M08\\\addlinespace[3pt]
% VI-CATALOG-ROW anchuras 288
\textbf{288. a\_\allowbreak{}1(1640) WIDTH}\par a\_\allowbreak{}1(1640)\allowbreak{}0 | a\_\allowbreak{}1(1640)\allowbreak{}+ | a\_\allowbreak{}1(1640)\allowbreak{}-\par {\fontsize{8.1}{9.4}\selectfont PDG2026:\allowbreak{}observable:\allowbreak{}M161:\allowbreak{}M161W:\allowbreak{}334741} & \textbf{250+-40} MeV & \(u_+\): 38.76802\allowbreak{}61069732\allowbreak{}67\par \(u_-\): 38.00446\allowbreak{}44472388\allowbreak{}31\par Tipo: NONE\par Confianza: ND & COMP-MESON\par E04 / M08\\\addlinespace[3pt]
% VI-CATALOG-ROW anchuras 289
\textbf{289. a\_\allowbreak{}2(1700) WIDTH}\par a\_\allowbreak{}2(1700)\allowbreak{}0 | a\_\allowbreak{}2(1700)\allowbreak{}+ | a\_\allowbreak{}2(1700)\allowbreak{}-\par {\fontsize{8.1}{9.4}\selectfont PDG2026:\allowbreak{}observable:\allowbreak{}M162:\allowbreak{}M162W:\allowbreak{}334904} & \textbf{380+60-50} MeV & \(u_+\): 56.80215\allowbreak{}66161437\allowbreak{}33\par \(u_-\): 49.91978\allowbreak{}68473076\allowbreak{}7\par Tipo: NONE\par Confianza: ND & COMP-MESON\par E04 / M08\\\addlinespace[3pt]
% VI-CATALOG-ROW anchuras 290
\textbf{290. pi\_\allowbreak{}1(1600) WIDTH (eta pi mode)}\par pi\_\allowbreak{}1(1600)\allowbreak{}0 | pi\_\allowbreak{}1(1600)\allowbreak{}+ | pi\_\allowbreak{}1(1600)\allowbreak{}-\par {\fontsize{8.1}{9.4}\selectfont PDG2026:\allowbreak{}observable:\allowbreak{}M164:\allowbreak{}M164WEP:\allowbreak{}334726} & \textbf{330+-35} MeV & \(u_+\): 34.72768\allowbreak{}93152504\allowbreak{}51\par \(u_-\): 34.69315\allowbreak{}16296829\allowbreak{}87\par Tipo: NONE\par Confianza: ND & COMP-MESON\par E04 / M08\\\addlinespace[3pt]
% VI-CATALOG-ROW anchuras 291
\textbf{291. pi\_\allowbreak{}1(1600) WIDTH (non-eta pi mode)}\par pi\_\allowbreak{}1(1600)\allowbreak{}0 | pi\_\allowbreak{}1(1600)\allowbreak{}+ | pi\_\allowbreak{}1(1600)\allowbreak{}-\par {\fontsize{8.1}{9.4}\selectfont PDG2026:\allowbreak{}observable:\allowbreak{}M164:\allowbreak{}M164W:\allowbreak{}334727} & \textbf{370+50-60} MeV & \(u_+\): 50.32451\allowbreak{}63918421\allowbreak{}17\par \(u_-\): 56.45471\allowbreak{}40483811\allowbreak{}28\par Tipo: NONE\par Confianza: ND & COMP-MESON\par E04 / M08\\\addlinespace[3pt]
% VI-CATALOG-ROW anchuras 292
\textbf{292. h\_\allowbreak{}1(1595) WIDTH}\par h\_\allowbreak{}1(1595)\allowbreak{}0\par {\fontsize{8.1}{9.4}\selectfont PDG2026:\allowbreak{}observable:\allowbreak{}M166:\allowbreak{}M166W:\allowbreak{}334721} & \textbf{380+90-120} MeV & \(u_+\): 92.19544\allowbreak{}45729288\allowbreak{}78\par \(u_-\): 116.6190\allowbreak{}37896905\allowbreak{}99\par Tipo: NONE\par Confianza: ND & COMP-MESON\par E04 / M08\\\addlinespace[3pt]
% VI-CATALOG-ROW anchuras 293
\textbf{293. f\_\allowbreak{}0(2100) WIDTH}\par f\_\allowbreak{}0(2100)\allowbreak{}0\par {\fontsize{8.1}{9.4}\selectfont PDG2026:\allowbreak{}observable:\allowbreak{}M168:\allowbreak{}M168W:\allowbreak{}335157} & \textbf{287+32-24} MeV & \(u_+\): 31.59491\allowbreak{}57013644\allowbreak{}21\par \(u_-\): 24.04867\allowbreak{}52700559\allowbreak{}91\par Tipo: NONE\par Confianza: ND & COMP-MESON\par E04 / M08\\\addlinespace[3pt]
% VI-CATALOG-ROW anchuras 294
\textbf{294. eta\_\allowbreak{}b(1S) WIDTH}\par eta\_\allowbreak{}b(1S)\par {\fontsize{8.1}{9.4}\selectfont PDG2026:\allowbreak{}observable:\allowbreak{}M171:\allowbreak{}M171W:\allowbreak{}346826} & \textbf{10+5-4} MeV & \(u_+\): 4.768415\allowbreak{}79595203\allowbreak{}06\par \(u_-\): 3.614819\allowbreak{}87601074\allowbreak{}91\par Tipo: NONE\par Confianza: ND & COMP-MESON\par E04 / M08\\\addlinespace[3pt]
% VI-CATALOG-ROW anchuras 295
\textbf{295. D\_\allowbreak{}s0*(2317)\allowbreak{}+- WIDTH}\par D\_\allowbreak{}s0\textasciicircum{}*(2317)\allowbreak{}+ | D\_\allowbreak{}s0\textasciicircum{}*(2317)\allowbreak{}-\par {\fontsize{8.1}{9.4}\selectfont PDG2026:\allowbreak{}observable:\allowbreak{}M172:\allowbreak{}M172W:\allowbreak{}338433} & \textbf{< 3.8} MeV\par Presentación: <3.8 & \(u_+\): NA\par \(u_-\): NA\par Tipo: U\par Confianza: 95 & COMP-MESON\par E04 / M08\\\addlinespace[3pt]
% VI-CATALOG-ROW anchuras 296
\textbf{296. D\_\allowbreak{}s1(2460)\allowbreak{}+- WIDTH}\par D\_\allowbreak{}s1(2460)\allowbreak{}+ | D\_\allowbreak{}s1(2460)\allowbreak{}-\par {\fontsize{8.1}{9.4}\selectfont PDG2026:\allowbreak{}observable:\allowbreak{}M173:\allowbreak{}M173W:\allowbreak{}338454} & \textbf{< 3.5} MeV\par Presentación: <3.5 & \(u_+\): NA\par \(u_-\): NA\par Tipo: U\par Confianza: 95 & COMP-MESON\par E04 / M08\\\addlinespace[3pt]
% VI-CATALOG-ROW anchuras 297
\textbf{297. K\_\allowbreak{}0*(700) Breit-Wigner Width}\par K\_\allowbreak{}0\textasciicircum{}*(700)\allowbreak{}0 | Kbar\_\allowbreak{}0\textasciicircum{}*(700)\allowbreak{}0 | K\_\allowbreak{}0\textasciicircum{}*(700)\allowbreak{}+ | K\_\allowbreak{}0\textasciicircum{}*(700)\allowbreak{}-\par {\fontsize{8.1}{9.4}\selectfont PDG2026:\allowbreak{}observable:\allowbreak{}M174:\allowbreak{}M174W:\allowbreak{}335787} & \textbf{463+-27} MeV & \(u_+\): 28.01802\allowbreak{}00105868\allowbreak{}49\par \(u_-\): 25.96920\allowbreak{}55984802\allowbreak{}12\par Tipo: NONE\par Confianza: ND & COMP-MESON\par E04 / M08\\\addlinespace[3pt]
% VI-CATALOG-ROW anchuras 298
\textbf{298. eta(1475) WIDTH}\par eta(1475)\allowbreak{}0\par {\fontsize{8.1}{9.4}\selectfont PDG2026:\allowbreak{}observable:\allowbreak{}M175:\allowbreak{}M175W5:\allowbreak{}334617} & \textbf{96+-9} MeV & \(u_+\): 9.164778\allowbreak{}12816133\allowbreak{}59\par \(u_-\): 8.658910\allowbreak{}56010132\allowbreak{}05\par Tipo: NONE\par Confianza: ND & COMP-MESON\par E04 / M08\\\addlinespace[3pt]
% VI-CATALOG-ROW anchuras 299
\textbf{299. chi\_\allowbreak{}c1(3872) WIDTH}\par chi\_\allowbreak{}c1(3872)\par {\fontsize{8.1}{9.4}\selectfont PDG2026:\allowbreak{}observable:\allowbreak{}M176:\allowbreak{}M176W:\allowbreak{}346132} & \textbf{1.19+-0.21} MeV & \(u_+\): 0.214461\allowbreak{}41015052\allowbreak{}739\par \(u_-\): 0.212127\allowbreak{}96092607\allowbreak{}489\par Tipo: NONE\par Confianza: ND & COMP-MESON\par E04 / M08\\\addlinespace[3pt]
% VI-CATALOG-ROW anchuras 300
\textbf{300. D\_\allowbreak{}1(2430)\allowbreak{}0 WIDTH}\par D\_\allowbreak{}1(2430)\allowbreak{}0 | Dbar\_\allowbreak{}1(2430)\allowbreak{}0\par {\fontsize{8.1}{9.4}\selectfont PDG2026:\allowbreak{}observable:\allowbreak{}M180:\allowbreak{}M180W:\allowbreak{}337828} & \textbf{314+-29} MeV & \(u_+\): 28.68695\allowbreak{}27683639\allowbreak{}7\par \(u_-\): 28.32791\allowbreak{}46860413\allowbreak{}21\par Tipo: NONE\par Confianza: ND & COMP-MESON\par E04 / M08\\\addlinespace[3pt]
% VI-CATALOG-ROW anchuras 301
\textbf{301. psi(4360) WIDTH}\par psi(4360)\par {\fontsize{8.1}{9.4}\selectfont PDG2026:\allowbreak{}observable:\allowbreak{}M181:\allowbreak{}M181W:\allowbreak{}346636} & \textbf{124+-13} MeV & \(u_+\): 13.06112\allowbreak{}47068081\allowbreak{}21\par \(u_-\): 13.28650\allowbreak{}41822744\allowbreak{}7\par Tipo: NONE\par Confianza: ND & COMP-MESON\par E04 / M08\\\addlinespace[3pt]
% VI-CATALOG-ROW anchuras 302
\textbf{302. D\_\allowbreak{}s1*(2700)\allowbreak{}+ WIDTH}\par D\_\allowbreak{}s1\textasciicircum{}*(2700)\allowbreak{}+ | D\_\allowbreak{}s1\textasciicircum{}*(2700)\allowbreak{}-\par {\fontsize{8.1}{9.4}\selectfont PDG2026:\allowbreak{}observable:\allowbreak{}M182:\allowbreak{}M182W:\allowbreak{}338525} & \textbf{122+-10} MeV & \(u_+\): 10.54341\allowbreak{}75204121\allowbreak{}29\par \(u_-\): 10.06094\allowbreak{}70530140\allowbreak{}1\par Tipo: NONE\par Confianza: ND & COMP-MESON\par E04 / M08\\\addlinespace[3pt]
% VI-CATALOG-ROW anchuras 303
\textbf{303. pi\_\allowbreak{}2(1880) WIDTH}\par pi\_\allowbreak{}2(1880)\allowbreak{}0\par {\fontsize{8.1}{9.4}\selectfont PDG2026:\allowbreak{}observable:\allowbreak{}M185:\allowbreak{}M185W:\allowbreak{}335044} & \textbf{237+33-30} MeV & \(u_+\): 33.31719\allowbreak{}74639905\allowbreak{}17\par \(u_-\): 29.50265\allowbreak{}80378861\allowbreak{}11\par Tipo: NONE\par Confianza: ND & COMP-MESON\par E04 / M08\\\addlinespace[3pt]
% VI-CATALOG-ROW anchuras 304
\textbf{304. B\_\allowbreak{}s2*(5840)\allowbreak{}0 WIDTH}\par B\_\allowbreak{}s2\textasciicircum{}*(5840)\allowbreak{}0 | Bbar\_\allowbreak{}s2\textasciicircum{}*(5840)\allowbreak{}0\par {\fontsize{8.1}{9.4}\selectfont PDG2026:\allowbreak{}observable:\allowbreak{}M186:\allowbreak{}M186W:\allowbreak{}343141} & \textbf{1.49+-0.27} MeV & \(u_+\): 0.266604\allowbreak{}79598163\allowbreak{}422\par \(u_-\): 0.266604\allowbreak{}79598163\allowbreak{}422\par Tipo: NONE\par Confianza: ND & COMP-MESON\par E04 / M08\\\addlinespace[3pt]
% VI-CATALOG-ROW anchuras 305
\textbf{305. B\_\allowbreak{}s1(5830)\allowbreak{}0 WIDTH}\par B\_\allowbreak{}s1(5830)\allowbreak{}0 | Bbar\_\allowbreak{}s1(5830)\allowbreak{}0\par {\fontsize{8.1}{9.4}\selectfont PDG2026:\allowbreak{}observable:\allowbreak{}M187:\allowbreak{}M187W:\allowbreak{}343133} & \textbf{0.5+-0.4} MeV & \(u_+\): 0.424264\allowbreak{}06871192\allowbreak{}851\par \(u_-\): 0.424264\allowbreak{}06871192\allowbreak{}851\par Tipo: NONE\par Confianza: ND & COMP-MESON\par E04 / M08\\\addlinespace[3pt]
% VI-CATALOG-ROW anchuras 306
\textbf{306. rho(1570) WIDTH}\par rho(1570)\allowbreak{}0\par {\fontsize{8.1}{9.4}\selectfont PDG2026:\allowbreak{}observable:\allowbreak{}M188:\allowbreak{}M188W:\allowbreak{}334715} & \textbf{140+-90} MeV & \(u_+\): 86.45229\allowbreak{}89862039\allowbreak{}42\par \(u_-\): 86.45229\allowbreak{}89862039\allowbreak{}42\par Tipo: NONE\par Confianza: ND & COMP-MESON\par E04 / M08\\\addlinespace[3pt]
% VI-CATALOG-ROW anchuras 307
\textbf{307. psi(4660) WIDTH}\par psi(4660)\par {\fontsize{8.1}{9.4}\selectfont PDG2026:\allowbreak{}observable:\allowbreak{}M189:\allowbreak{}M189W:\allowbreak{}346768} & \textbf{55+-9} MeV & \(u_+\): 8.889722\allowbreak{}91869921\allowbreak{}7\par \(u_-\): 8.409913\allowbreak{}96186550\allowbreak{}06\par Tipo: NONE\par Confianza: ND & COMP-MESON\par E04 / M08\\\addlinespace[3pt]
% VI-CATALOG-ROW anchuras 308
\textbf{308. X(4160) WIDTH}\par X(4160)\par {\fontsize{8.1}{9.4}\selectfont PDG2026:\allowbreak{}observable:\allowbreak{}M190:\allowbreak{}M190W:\allowbreak{}346452} & \textbf{136+60-35} MeV & \(u_+\): 56.53922\allowbreak{}09238173\allowbreak{}54\par \(u_-\): 34.62523\allowbreak{}07256621\allowbreak{}42\par Tipo: NONE\par Confianza: ND & COMP-MESON\par E04 / M08\\\addlinespace[3pt]
% VI-CATALOG-ROW anchuras 309
\textbf{309. T\_\allowbreak{}c\_\allowbreak{}cbar(4050)\allowbreak{}+ WIDTH}\par T\_\allowbreak{}c\_\allowbreak{}cbar(4050)\allowbreak{}+\par {\fontsize{8.1}{9.4}\selectfont PDG2026:\allowbreak{}observable:\allowbreak{}M191:\allowbreak{}M191W:\allowbreak{}347991} & \textbf{82+50-28} MeV & \(u_+\): 51.47815\allowbreak{}07049349\allowbreak{}97\par \(u_-\): 27.80287\allowbreak{}75489156\allowbreak{}89\par Tipo: NONE\par Confianza: ND & COMP-MESON\par E04 / M08\\\addlinespace[3pt]
% VI-CATALOG-ROW anchuras 310
\textbf{310. T\_\allowbreak{}c\_\allowbreak{}cbar(4250)\allowbreak{}+ WIDTH}\par T\_\allowbreak{}c\_\allowbreak{}cbar(4250)\allowbreak{}+\par {\fontsize{8.1}{9.4}\selectfont PDG2026:\allowbreak{}observable:\allowbreak{}M192:\allowbreak{}M192W:\allowbreak{}348027} & \textbf{180+320-70} MeV & \(u_+\): 320.5807\allowbreak{}23063630\allowbreak{}27\par \(u_-\): 72.40165\allowbreak{}74395918\allowbreak{}16\par Tipo: NONE\par Confianza: ND & COMP-MESON\par E04 / M08\\\addlinespace[3pt]
% VI-CATALOG-ROW anchuras 311
\textbf{311. chi\_\allowbreak{}c1(4140) WIDTH}\par chi\_\allowbreak{}c1(4140)\par {\fontsize{8.1}{9.4}\selectfont PDG2026:\allowbreak{}observable:\allowbreak{}M193:\allowbreak{}M193W:\allowbreak{}346337} & \textbf{19+7-5} MeV & \(u_+\): 6.558667\allowbreak{}34077530\allowbreak{}64\par \(u_-\): 5.275820\allowbreak{}25104883\allowbreak{}37\par Tipo: NONE\par Confianza: ND & COMP-MESON\par E04 / M08\\\addlinespace[3pt]
% VI-CATALOG-ROW anchuras 312
\textbf{312. X(4350) WIDTH}\par X(4350)\par {\fontsize{8.1}{9.4}\selectfont PDG2026:\allowbreak{}observable:\allowbreak{}M194:\allowbreak{}M194W:\allowbreak{}346629} & \textbf{13+18-10} MeV & \(u_+\): 18.43908\allowbreak{}89145857\allowbreak{}71\par \(u_-\): 9.848857\allowbreak{}80179610\allowbreak{}39\par Tipo: NONE\par Confianza: ND & COMP-MESON\par E04 / M08\\\addlinespace[3pt]
% VI-CATALOG-ROW anchuras 313
\textbf{313. T\_\allowbreak{}c\_\allowbreak{}cbar\_\allowbreak{}1(4430)\allowbreak{}+ WIDTH}\par T\_\allowbreak{}c\_\allowbreak{}cbar\_\allowbreak{}1(4430)\allowbreak{}+\par {\fontsize{8.1}{9.4}\selectfont PDG2026:\allowbreak{}observable:\allowbreak{}M195:\allowbreak{}M195W:\allowbreak{}348031} & \textbf{189+-29} MeV & \(u_+\): 28.96086\allowbreak{}26207506\allowbreak{}99\par \(u_-\): 29.21228\allowbreak{}16924321\allowbreak{}89\par Tipo: NONE\par Confianza: ND & COMP-MESON\par E04 / M08\\\addlinespace[3pt]
% VI-CATALOG-ROW anchuras 314
\textbf{314. D\_\allowbreak{}s1*(2860)\allowbreak{}+ WIDTH}\par D\_\allowbreak{}s1\textasciicircum{}*(2860)\allowbreak{}+ | D\_\allowbreak{}s1\textasciicircum{}*(2860)\allowbreak{}-\par {\fontsize{8.1}{9.4}\selectfont PDG2026:\allowbreak{}observable:\allowbreak{}M196:\allowbreak{}M196W:\allowbreak{}338534} & \textbf{160+-80} MeV & \(u_+\): 80.36168\allowbreak{}24114577\allowbreak{}64\par \(u_-\): 80.36168\allowbreak{}24114577\allowbreak{}64\par Tipo: NONE\par Confianza: ND & COMP-MESON\par E04 / M08\\\addlinespace[3pt]
% VI-CATALOG-ROW anchuras 315
\textbf{315. D\_\allowbreak{}sJ(3040)\allowbreak{}+ WIDTH}\par D\_\allowbreak{}sJ(3040)\allowbreak{}+ | D\_\allowbreak{}sJ(3040)\allowbreak{}-\par {\fontsize{8.1}{9.4}\selectfont PDG2026:\allowbreak{}observable:\allowbreak{}M197:\allowbreak{}M197W:\allowbreak{}338545} & \textbf{240+-60} MeV & \(u_+\): 57.80138\allowbreak{}40664737\allowbreak{}02\par \(u_-\): 54.67174\allowbreak{}77313465\allowbreak{}78\par Tipo: NONE\par Confianza: ND & COMP-MESON\par E04 / M08\\\addlinespace[3pt]
% VI-CATALOG-ROW anchuras 316
\textbf{316. D\_\allowbreak{}0(2550)\allowbreak{}0 WIDTH}\par D\_\allowbreak{}0(2550)\allowbreak{}0 | Dbar\_\allowbreak{}0(2550)\allowbreak{}0 | D\_\allowbreak{}0(2550)\allowbreak{}+ | D\_\allowbreak{}0(2550)\allowbreak{}-\par {\fontsize{8.1}{9.4}\selectfont PDG2026:\allowbreak{}observable:\allowbreak{}M198:\allowbreak{}M198W:\allowbreak{}337844} & \textbf{165+-24} MeV & \(u_+\): 23.74755\allowbreak{}33977095\allowbreak{}49\par \(u_-\): 23.74755\allowbreak{}33977095\allowbreak{}49\par Tipo: NONE\par Confianza: ND & COMP-MESON\par E04 / M08\\\addlinespace[3pt]
% VI-CATALOG-ROW anchuras 317
\textbf{317. D\_\allowbreak{}1*(2600)\allowbreak{}0 WIDTH}\par D\_\allowbreak{}1\textasciicircum{}*(2600)\allowbreak{}0 | Dbar\_\allowbreak{}1\textasciicircum{}*(2600)\allowbreak{}0 | D\_\allowbreak{}1\textasciicircum{}*(2600)\allowbreak{}+ | D\_\allowbreak{}1\textasciicircum{}*(2600)\allowbreak{}-\par {\fontsize{8.1}{9.4}\selectfont PDG2026:\allowbreak{}observable:\allowbreak{}M199:\allowbreak{}M199W:\allowbreak{}337847} & \textbf{141+-23} MeV & \(u_+\): 22.83652\allowbreak{}40441833\allowbreak{}88\par \(u_-\): 22.83652\allowbreak{}40441833\allowbreak{}88\par Tipo: NONE\par Confianza: ND & COMP-MESON\par E04 / M08\\\addlinespace[3pt]
% VI-CATALOG-ROW anchuras 318
\textbf{318. eta\_\allowbreak{}b(2S) WIDTH}\par eta\_\allowbreak{}b(2S)\par {\fontsize{8.1}{9.4}\selectfont PDG2026:\allowbreak{}observable:\allowbreak{}M200:\allowbreak{}M200W:\allowbreak{}347283} & \textbf{< 24} MeV\par Presentación: <24 & \(u_+\): NA\par \(u_-\): NA\par Tipo: U\par Confianza: 90 & COMP-MESON\par E04 / M08\\\addlinespace[3pt]
% VI-CATALOG-ROW anchuras 319
\textbf{319. D\_\allowbreak{}3*(2750) WIDTH}\par D\_\allowbreak{}3\textasciicircum{}*(2750)\allowbreak{}0 | Dbar\_\allowbreak{}3\textasciicircum{}*(2750)\allowbreak{}0 | D\_\allowbreak{}3\textasciicircum{}*(2750)\allowbreak{}+ | D\_\allowbreak{}3\textasciicircum{}*(2750)\allowbreak{}-\par {\fontsize{8.1}{9.4}\selectfont PDG2026:\allowbreak{}observable:\allowbreak{}M203:\allowbreak{}M203W:\allowbreak{}337861} & \textbf{66+-5} MeV & \(u_+\): 4.602693\allowbreak{}09173125\allowbreak{}41\par \(u_-\): 4.602693\allowbreak{}09173125\allowbreak{}41\par Tipo: NONE\par Confianza: ND & COMP-MESON\par E04 / M08\\\addlinespace[3pt]
% VI-CATALOG-ROW anchuras 320
\textbf{320. T\_\allowbreak{}b\_\allowbreak{}bbar\_\allowbreak{}1(10610)\allowbreak{}+- WIDTH}\par T\_\allowbreak{}b\_\allowbreak{}bbar\_\allowbreak{}1(10610)\allowbreak{}+\par {\fontsize{8.1}{9.4}\selectfont PDG2026:\allowbreak{}observable:\allowbreak{}M207:\allowbreak{}M207W:\allowbreak{}348054} & \textbf{18.4+-2.4} MeV & \(u_+\): 2.399999\allowbreak{}99999999\allowbreak{}99\par \(u_-\): 2.399999\allowbreak{}99999999\allowbreak{}99\par Tipo: NONE\par Confianza: ND & COMP-MESON\par E04 / M08\\\addlinespace[3pt]
% VI-CATALOG-ROW anchuras 321
\textbf{321. T\_\allowbreak{}b\_\allowbreak{}bbar\_\allowbreak{}1(10650)\allowbreak{}+ WIDTH}\par T\_\allowbreak{}b\_\allowbreak{}bbar\_\allowbreak{}1(10650)\allowbreak{}+\par {\fontsize{8.1}{9.4}\selectfont PDG2026:\allowbreak{}observable:\allowbreak{}M208:\allowbreak{}M208W:\allowbreak{}348076} & \textbf{11.5+-2.2} MeV & \(u_+\): 2.200000\allowbreak{}00000000\allowbreak{}02\par \(u_-\): 2.200000\allowbreak{}00000000\allowbreak{}02\par Tipo: NONE\par Confianza: ND & COMP-MESON\par E04 / M08\\\addlinespace[3pt]
% VI-CATALOG-ROW anchuras 322
\textbf{322. T\_\allowbreak{}c\_\allowbreak{}cbar\_\allowbreak{}1(3900) WIDTH}\par T\_\allowbreak{}c\_\allowbreak{}cbar\_\allowbreak{}1(3900)\allowbreak{}+\par {\fontsize{8.1}{9.4}\selectfont PDG2026:\allowbreak{}observable:\allowbreak{}M210:\allowbreak{}M210W:\allowbreak{}347939} & \textbf{29.5+-2.4} MeV & \(u_+\): 2.425749\allowbreak{}28658747\allowbreak{}7\par \(u_-\): 2.417401\allowbreak{}48443531\allowbreak{}49\par Tipo: NONE\par Confianza: ND & COMP-MESON\par E04 / M08\\\addlinespace[3pt]
% VI-CATALOG-ROW anchuras 323
\textbf{323. psi\_\allowbreak{}2(3823) WIDTH}\par psi\_\allowbreak{}2(3823)\par {\fontsize{8.1}{9.4}\selectfont PDG2026:\allowbreak{}observable:\allowbreak{}M212:\allowbreak{}M212W:\allowbreak{}346101} & \textbf{< 2.9} MeV\par Presentación: <2.9 & \(u_+\): NA\par \(u_-\): NA\par Tipo: U\par Confianza: 90 & COMP-MESON\par E04 / M08\\\addlinespace[3pt]
% VI-CATALOG-ROW anchuras 324
\textbf{324. T\_\allowbreak{}c\_\allowbreak{}cbar(4020) WIDTH}\par T\_\allowbreak{}c\_\allowbreak{}cbar(4020)\allowbreak{}+\par {\fontsize{8.1}{9.4}\selectfont PDG2026:\allowbreak{}observable:\allowbreak{}M213:\allowbreak{}M213W:\allowbreak{}347977} & \textbf{13+-5} MeV & \(u_+\): 5.274210\allowbreak{}46109626\allowbreak{}52\par \(u_-\): 5.274210\allowbreak{}46109626\allowbreak{}52\par Tipo: NONE\par Confianza: ND & COMP-MESON\par E04 / M08\\\addlinespace[3pt]
% VI-CATALOG-ROW anchuras 325
\textbf{325. T\_\allowbreak{}c\_\allowbreak{}cbar\_\allowbreak{}0(4240)\allowbreak{}+ WIDTH}\par T\_\allowbreak{}c\_\allowbreak{}cbar\_\allowbreak{}0(4240)\allowbreak{}0\par {\fontsize{8.1}{9.4}\selectfont PDG2026:\allowbreak{}observable:\allowbreak{}M216:\allowbreak{}M216W:\allowbreak{}348023} & \textbf{220+120-90} MeV & \(u_+\): 117.7837\allowbreak{}00060746\allowbreak{}99\par \(u_-\): 87.66413\allowbreak{}17757724\allowbreak{}61\par Tipo: NONE\par Confianza: ND & COMP-MESON\par E04 / M08\\\addlinespace[3pt]
% VI-CATALOG-ROW anchuras 326
\textbf{326. T\_\allowbreak{}c\_\allowbreak{}cbar(4055)\allowbreak{}+ WIDTH}\par T\_\allowbreak{}c\_\allowbreak{}cbar(4055)\allowbreak{}+\par {\fontsize{8.1}{9.4}\selectfont PDG2026:\allowbreak{}observable:\allowbreak{}M223:\allowbreak{}M223W:\allowbreak{}348001} & \textbf{45+-13} MeV & \(u_+\): 12.52996\allowbreak{}40861416\allowbreak{}69\par \(u_-\): 12.52996\allowbreak{}40861416\allowbreak{}69\par Tipo: NONE\par Confianza: ND & COMP-MESON\par E04 / M08\\\addlinespace[3pt]
% VI-CATALOG-ROW anchuras 327
\textbf{327. D\_\allowbreak{}s3*(2860)\allowbreak{}+ WIDTH}\par D\_\allowbreak{}s3\textasciicircum{}*(2860)\allowbreak{}+ | D\_\allowbreak{}s3\textasciicircum{}*(2860)\allowbreak{}-\par {\fontsize{8.1}{9.4}\selectfont PDG2026:\allowbreak{}observable:\allowbreak{}M226:\allowbreak{}M226W:\allowbreak{}338536} & \textbf{53+-10} MeV & \(u_+\): 9.899494\allowbreak{}93661166\allowbreak{}54\par \(u_-\): 9.899494\allowbreak{}93661166\allowbreak{}54\par Tipo: NONE\par Confianza: ND & COMP-MESON\par E04 / M08\\\addlinespace[3pt]
% VI-CATALOG-ROW anchuras 328
\textbf{328. a\_\allowbreak{}0(1950) WIDTH}\par a\_\allowbreak{}0(1950)\allowbreak{}0\par {\fontsize{8.1}{9.4}\selectfont PDG2026:\allowbreak{}observable:\allowbreak{}M227:\allowbreak{}M227W:\allowbreak{}335071} & \textbf{270+-40} MeV & \(u_+\): 36.40054\allowbreak{}94464025\allowbreak{}93\par \(u_-\): 36.40054\allowbreak{}94464025\allowbreak{}93\par Tipo: NONE\par Confianza: ND & COMP-MESON\par E04 / M08\\\addlinespace[3pt]
% VI-CATALOG-ROW anchuras 329
\textbf{329. D\_\allowbreak{}2(2740)\allowbreak{}0 WIDTH}\par D\_\allowbreak{}2(2740)\allowbreak{}0 | Dbar\_\allowbreak{}2(2740)\allowbreak{}0 | D\_\allowbreak{}2(2740)\allowbreak{}+ | D\_\allowbreak{}2(2740)\allowbreak{}-\par {\fontsize{8.1}{9.4}\selectfont PDG2026:\allowbreak{}observable:\allowbreak{}M228:\allowbreak{}M228W:\allowbreak{}337858} & \textbf{88+-19} MeV & \(u_+\): 19.43521\allowbreak{}80211727\allowbreak{}09\par \(u_-\): 19.43521\allowbreak{}80211727\allowbreak{}09\par Tipo: NONE\par Confianza: ND & COMP-MESON\par E04 / M08\\\addlinespace[3pt]
% VI-CATALOG-ROW anchuras 330
\textbf{330. D\_\allowbreak{}2*(3000)\allowbreak{}0 WIDTH}\par D\_\allowbreak{}2\textasciicircum{}*(3000)\allowbreak{}0 | Dbar\_\allowbreak{}2\textasciicircum{}*(3000)\allowbreak{}0 | D\_\allowbreak{}2\textasciicircum{}*(3000)\allowbreak{}+ | D\_\allowbreak{}2\textasciicircum{}*(3000)\allowbreak{}-\par {\fontsize{8.1}{9.4}\selectfont PDG2026:\allowbreak{}observable:\allowbreak{}M229:\allowbreak{}M229W:\allowbreak{}337873} & \textbf{190+-80} MeV & \(u_+\): 81.41252\allowbreak{}97481904\allowbreak{}79\par \(u_-\): 81.41252\allowbreak{}97481904\allowbreak{}79\par Tipo: NONE\par Confianza: ND & COMP-MESON\par E04 / M08\\\addlinespace[3pt]
% VI-CATALOG-ROW anchuras 331
\textbf{331. T\_\allowbreak{}c\_\allowbreak{}cbar\_\allowbreak{}1(4200)\allowbreak{}+ WIDTH}\par T\_\allowbreak{}c\_\allowbreak{}cbar\_\allowbreak{}1(4200)\allowbreak{}+\par {\fontsize{8.1}{9.4}\selectfont PDG2026:\allowbreak{}observable:\allowbreak{}M231:\allowbreak{}M231W:\allowbreak{}348013} & \textbf{310+-40} MeV & \(u_+\): 35.26096\allowbreak{}58006496\allowbreak{}29\par \(u_-\): 36.58705\allowbreak{}73754017\allowbreak{}27\par Tipo: NONE\par Confianza: ND & COMP-MESON\par E04 / M08\\\addlinespace[3pt]
% VI-CATALOG-ROW anchuras 332
\textbf{332. T\_\allowbreak{}b\_\allowbreak{}sbar(5568)\allowbreak{}+ WIDTH}\par T\_\allowbreak{}b\_\allowbreak{}sbar(5568)\allowbreak{}+\par {\fontsize{8.1}{9.4}\selectfont PDG2026:\allowbreak{}observable:\allowbreak{}M232:\allowbreak{}M232W:\allowbreak{}348037} & \textbf{19+9-7} MeV & \(u_+\): 8.640601\allowbreak{}83089118\allowbreak{}01\par \(u_-\): 7.186793\allowbreak{}44353237\allowbreak{}86\par Tipo: NONE\par Confianza: ND & COMP-MESON\par E04 / M08\\\addlinespace[3pt]
% VI-CATALOG-ROW anchuras 333
\textbf{333. chi\_\allowbreak{}c1(4274) WIDTH}\par chi\_\allowbreak{}c1(4274)\par {\fontsize{8.1}{9.4}\selectfont PDG2026:\allowbreak{}observable:\allowbreak{}M233:\allowbreak{}M233W:\allowbreak{}346625} & \textbf{54+-8} MeV & \(u_+\): 8.196852\allowbreak{}54949333\allowbreak{}24\par \(u_-\): 7.793383\allowbreak{}77392468\allowbreak{}21\par Tipo: NONE\par Confianza: ND & COMP-MESON\par E04 / M08\\\addlinespace[3pt]
% VI-CATALOG-ROW anchuras 334
\textbf{334. chi\_\allowbreak{}c0(4500) WIDTH}\par chi\_\allowbreak{}c0(4500)\par {\fontsize{8.1}{9.4}\selectfont PDG2026:\allowbreak{}observable:\allowbreak{}M234:\allowbreak{}M234W:\allowbreak{}346756} & \textbf{87+30-24} MeV & \(u_+\): 29.92716\allowbreak{}04050614\allowbreak{}39\par \(u_-\): 24.39965\allowbreak{}24356650\allowbreak{}01\par Tipo: NONE\par Confianza: ND & COMP-MESON\par E04 / M08\\\addlinespace[3pt]
% VI-CATALOG-ROW anchuras 335
\textbf{335. chi\_\allowbreak{}c0(4700) WIDTH}\par chi\_\allowbreak{}c0(4700)\par {\fontsize{8.1}{9.4}\selectfont PDG2026:\allowbreak{}observable:\allowbreak{}M235:\allowbreak{}M235W:\allowbreak{}346817} & \textbf{78+15-11} MeV & \(u_+\): 14.53972\allowbreak{}99640840\allowbreak{}1\par \(u_-\): 10.85293\allowbreak{}13170942\allowbreak{}49\par Tipo: NONE\par Confianza: ND & COMP-MESON\par E04 / M08\\\addlinespace[3pt]
% VI-CATALOG-ROW anchuras 336
\textbf{336. chi\_\allowbreak{}c0(3860) WIDTH}\par chi\_\allowbreak{}c0(3860)\par {\fontsize{8.1}{9.4}\selectfont PDG2026:\allowbreak{}observable:\allowbreak{}M237:\allowbreak{}M237W:\allowbreak{}346123} & \textbf{200+180-110} MeV & \(u_+\): 177.3696\allowbreak{}70462568\allowbreak{}11\par \(u_-\): 105.8914\allowbreak{}53857240\allowbreak{}09\par Tipo: NONE\par Confianza: ND & COMP-MESON\par E04 / M08\\\addlinespace[3pt]
% VI-CATALOG-ROW anchuras 337
\textbf{337. pi\_\allowbreak{}2(2005) WIDTH}\par pi\_\allowbreak{}2(2005)\allowbreak{}0\par {\fontsize{8.1}{9.4}\selectfont PDG2026:\allowbreak{}observable:\allowbreak{}M239:\allowbreak{}M239W:\allowbreak{}335107} & \textbf{370+16-90} MeV & \(u_+\): 15.91185\allowbreak{}26477216\allowbreak{}2\par \(u_-\): 94.13731\allowbreak{}62126639\allowbreak{}9\par Tipo: NONE\par Confianza: ND & COMP-MESON\par E04 / M08\\\addlinespace[3pt]
% VI-CATALOG-ROW anchuras 338
\textbf{338. T\_\allowbreak{}c\_\allowbreak{}cbar(4100)\allowbreak{}+ WIDTH}\par T\_\allowbreak{}c\_\allowbreak{}cbar(4100)\allowbreak{}+\par {\fontsize{8.1}{9.4}\selectfont PDG2026:\allowbreak{}observable:\allowbreak{}M240:\allowbreak{}M240W:\allowbreak{}348007} & \textbf{150+80-70} MeV & \(u_+\): 83.45058\allowbreak{}41801002\allowbreak{}65\par \(u_-\): 67.74215\allowbreak{}82177597\allowbreak{}94\par Tipo: NONE\par Confianza: ND & COMP-MESON\par E04 / M08\\\addlinespace[3pt]
% VI-CATALOG-ROW anchuras 339
\textbf{339. psi\_\allowbreak{}3(3842) WIDTH}\par psi\_\allowbreak{}3(3842)\par {\fontsize{8.1}{9.4}\selectfont PDG2026:\allowbreak{}observable:\allowbreak{}M241:\allowbreak{}M241W:\allowbreak{}346117} & \textbf{2.8+-0.6} MeV & \(u_+\): 0.618546\allowbreak{}68376768\allowbreak{}462\par \(u_-\): 0.618546\allowbreak{}68376768\allowbreak{}462\par Tipo: NONE\par Confianza: ND & COMP-MESON\par E04 / M08\\\addlinespace[3pt]
% VI-CATALOG-ROW anchuras 340
\textbf{340. Upsilon(10753) WIDTH}\par Upsilon(10753)\par {\fontsize{8.1}{9.4}\selectfont PDG2026:\allowbreak{}observable:\allowbreak{}M243:\allowbreak{}M243W:\allowbreak{}347801} & \textbf{29+-9} MeV & \(u_+\): 8.881441\allowbreak{}32446980\allowbreak{}81\par \(u_-\): 8.881441\allowbreak{}32446980\allowbreak{}81\par Tipo: NONE\par Confianza: ND & COMP-MESON\par E04 / M08\\\addlinespace[3pt]
% VI-CATALOG-ROW anchuras 341
\textbf{341. B\_\allowbreak{}1(5721)\allowbreak{}+ width}\par B\_\allowbreak{}1(5721)\allowbreak{}0 | Bbar\_\allowbreak{}1(5721)\allowbreak{}0 | B\_\allowbreak{}1(5721)\allowbreak{}+ | B\_\allowbreak{}1(5721)\allowbreak{}-\par {\fontsize{8.1}{9.4}\selectfont PDG2026:\allowbreak{}observable:\allowbreak{}M244:\allowbreak{}M244W+:\allowbreak{}342547} & \textbf{31+-6} MeV & \(u_+\): 5.847044\allowbreak{}69385892\par \(u_-\): 6.092104\allowbreak{}20060922\allowbreak{}43\par Tipo: NONE\par Confianza: ND & COMP-MESON\par E04 / M08\\\addlinespace[3pt]
% VI-CATALOG-ROW anchuras 342
\textbf{342. B\_\allowbreak{}1(5721)\allowbreak{}0 width}\par B\_\allowbreak{}1(5721)\allowbreak{}0 | Bbar\_\allowbreak{}1(5721)\allowbreak{}0 | B\_\allowbreak{}1(5721)\allowbreak{}+ | B\_\allowbreak{}1(5721)\allowbreak{}-\par {\fontsize{8.1}{9.4}\selectfont PDG2026:\allowbreak{}observable:\allowbreak{}M244:\allowbreak{}M244W0:\allowbreak{}342548} & \textbf{27.5+-3.4} MeV & \(u_+\): 3.422277\allowbreak{}84518650\allowbreak{}01\par \(u_-\): 3.422277\allowbreak{}84518650\allowbreak{}01\par Tipo: NONE\par Confianza: ND & COMP-MESON\par E04 / M08\\\addlinespace[3pt]
% VI-CATALOG-ROW anchuras 343
\textbf{343. B\_\allowbreak{}2*(5747)\allowbreak{}+ width}\par B\_\allowbreak{}2\textasciicircum{}*(5747)\allowbreak{}0 | Bbar\_\allowbreak{}2\textasciicircum{}*(5747)\allowbreak{}0 | B\_\allowbreak{}2\textasciicircum{}*(5747)\allowbreak{}+ | B\_\allowbreak{}2\textasciicircum{}*(5747)\allowbreak{}-\par {\fontsize{8.1}{9.4}\selectfont PDG2026:\allowbreak{}observable:\allowbreak{}M245:\allowbreak{}M245W+:\allowbreak{}342564} & \textbf{20+-5} MeV & \(u_+\): 5.468422\allowbreak{}91584328\allowbreak{}51\par \(u_-\): 5.468422\allowbreak{}91584328\allowbreak{}51\par Tipo: NONE\par Confianza: ND & COMP-MESON\par E04 / M08\\\addlinespace[3pt]
% VI-CATALOG-ROW anchuras 344
\textbf{344. B\_\allowbreak{}2*(5747)\allowbreak{}0 width}\par B\_\allowbreak{}2\textasciicircum{}*(5747)\allowbreak{}0 | Bbar\_\allowbreak{}2\textasciicircum{}*(5747)\allowbreak{}0 | B\_\allowbreak{}2\textasciicircum{}*(5747)\allowbreak{}+ | B\_\allowbreak{}2\textasciicircum{}*(5747)\allowbreak{}-\par {\fontsize{8.1}{9.4}\selectfont PDG2026:\allowbreak{}observable:\allowbreak{}M245:\allowbreak{}M245W0:\allowbreak{}342565} & \textbf{24.2+-1.7} MeV & \(u_+\): 1.695908\allowbreak{}66342803\allowbreak{}49\par \(u_-\): 1.709526\allowbreak{}43154618\allowbreak{}71\par Tipo: NONE\par Confianza: ND & COMP-MESON\par E04 / M08\\\addlinespace[3pt]
% VI-CATALOG-ROW anchuras 345
\textbf{345. B\_\allowbreak{}J(5840)\allowbreak{}+ WIDTH}\par B\_\allowbreak{}J(5840)\allowbreak{}0 | Bbar\_\allowbreak{}J(5840)\allowbreak{}0 | B\_\allowbreak{}J(5840)\allowbreak{}+ | B\_\allowbreak{}J(5840)\allowbreak{}-\par {\fontsize{8.1}{9.4}\selectfont PDG2026:\allowbreak{}observable:\allowbreak{}M246:\allowbreak{}M246W+:\allowbreak{}342577} & \textbf{220+-80} MeV & \(u_+\): 83.52245\allowbreak{}20712844\allowbreak{}05\par \(u_-\): 83.52245\allowbreak{}20712844\allowbreak{}05\par Tipo: NONE\par Confianza: ND & COMP-MESON\par E04 / M08\\\addlinespace[3pt]
% VI-CATALOG-ROW anchuras 346
\textbf{346. B\_\allowbreak{}J(5840)\allowbreak{}0 WIDTH}\par B\_\allowbreak{}J(5840)\allowbreak{}0 | Bbar\_\allowbreak{}J(5840)\allowbreak{}0 | B\_\allowbreak{}J(5840)\allowbreak{}+ | B\_\allowbreak{}J(5840)\allowbreak{}-\par {\fontsize{8.1}{9.4}\selectfont PDG2026:\allowbreak{}observable:\allowbreak{}M246:\allowbreak{}M246W0:\allowbreak{}342578} & \textbf{130+-40} MeV & \(u_+\): 38.01315\allowbreak{}56174964\allowbreak{}23\par \(u_-\): 38.01315\allowbreak{}56174964\allowbreak{}23\par Tipo: NONE\par Confianza: ND & COMP-MESON\par E04 / M08\\\addlinespace[3pt]
% VI-CATALOG-ROW anchuras 347
\textbf{347. eta(2370) WIDTH}\par eta(2370)\allowbreak{}0\par {\fontsize{8.1}{9.4}\selectfont PDG2026:\allowbreak{}observable:\allowbreak{}M247:\allowbreak{}M247W:\allowbreak{}335262} & \textbf{148+80-28} MeV & \(u_+\): 77.22224\allowbreak{}37006731\allowbreak{}61\par \(u_-\): 28.36717\allowbreak{}77206377\allowbreak{}8\par Tipo: NONE\par Confianza: ND & COMP-MESON\par E04 / M08\\\addlinespace[3pt]
% VI-CATALOG-ROW anchuras 348
\textbf{348. B\_\allowbreak{}J(5970)\allowbreak{}+ WIDTH}\par B\_\allowbreak{}J(5970)\allowbreak{}0 | Bbar\_\allowbreak{}J(5970)\allowbreak{}0 | B\_\allowbreak{}J(5970)\allowbreak{}+ | B\_\allowbreak{}J(5970)\allowbreak{}-\par {\fontsize{8.1}{9.4}\selectfont PDG2026:\allowbreak{}observable:\allowbreak{}M248:\allowbreak{}M248W+:\allowbreak{}342589} & \textbf{62+-20} MeV & \(u_+\): 20.64809\allowbreak{}26301544\allowbreak{}19\par \(u_-\): 20.22152\allowbreak{}57441381\allowbreak{}8\par Tipo: NONE\par Confianza: ND & COMP-MESON\par E04 / M08\\\addlinespace[3pt]
% VI-CATALOG-ROW anchuras 349
\textbf{349. B\_\allowbreak{}J(5970)\allowbreak{}0 WIDTH}\par B\_\allowbreak{}J(5970)\allowbreak{}0 | Bbar\_\allowbreak{}J(5970)\allowbreak{}0 | B\_\allowbreak{}J(5970)\allowbreak{}+ | B\_\allowbreak{}J(5970)\allowbreak{}-\par {\fontsize{8.1}{9.4}\selectfont PDG2026:\allowbreak{}observable:\allowbreak{}M248:\allowbreak{}M248W0:\allowbreak{}342590} & \textbf{81+-12} MeV & \(u_+\): 11.58405\allowbreak{}07320396\allowbreak{}3\par \(u_-\): 11.42146\allowbreak{}34346750\allowbreak{}6\par Tipo: NONE\par Confianza: ND & COMP-MESON\par E04 / M08\\\addlinespace[3pt]
% VI-CATALOG-ROW anchuras 350
\textbf{350. D\_\allowbreak{}1*(2760)\allowbreak{}0 WIDTH}\par D\_\allowbreak{}1\textasciicircum{}*(2760)\allowbreak{}0 | Dbar\_\allowbreak{}1\textasciicircum{}*(2760)\allowbreak{}0 | D\_\allowbreak{}1\textasciicircum{}*(2760)\allowbreak{}+ | D\_\allowbreak{}1\textasciicircum{}*(2760)\allowbreak{}-\par {\fontsize{8.1}{9.4}\selectfont PDG2026:\allowbreak{}observable:\allowbreak{}M249:\allowbreak{}M249W:\allowbreak{}337869} & \textbf{180+-40} MeV & \(u_+\): 38.27531\allowbreak{}84180092\allowbreak{}77\par \(u_-\): 38.27531\allowbreak{}84180092\allowbreak{}77\par Tipo: NONE\par Confianza: ND & COMP-MESON\par E04 / M08\\\addlinespace[3pt]
% VI-CATALOG-ROW anchuras 351
\textbf{351. T\_\allowbreak{}cbar\_\allowbreak{}sbar\_\allowbreak{}0*(2870)\allowbreak{}0 WIDTH}\par T\_\allowbreak{}cbar\_\allowbreak{}sbar\_\allowbreak{}0\textasciicircum{}*(2870)\allowbreak{}0\par {\fontsize{8.1}{9.4}\selectfont PDG2026:\allowbreak{}observable:\allowbreak{}M250:\allowbreak{}M250W:\allowbreak{}347917} & \textbf{68+-17} MeV & \(u_+\): 16.28283\allowbreak{}39322687\allowbreak{}69\par \(u_-\): 17.20879\allowbreak{}01693615\par Tipo: NONE\par Confianza: ND & COMP-MESON\par E04 / M08\\\addlinespace[3pt]
% VI-CATALOG-ROW anchuras 352
\textbf{352. T\_\allowbreak{}cbar\_\allowbreak{}sbar\_\allowbreak{}1*(2900)\allowbreak{}0 WIDTH}\par T\_\allowbreak{}cbar\_\allowbreak{}sbar\_\allowbreak{}1\textasciicircum{}*(2900)\allowbreak{}0\par {\fontsize{8.1}{9.4}\selectfont PDG2026:\allowbreak{}observable:\allowbreak{}M251:\allowbreak{}M251W:\allowbreak{}347923} & \textbf{106+-10} MeV & \(u_+\): 10.39615\allowbreak{}90631619\allowbreak{}3\par \(u_-\): 10.39615\allowbreak{}90631619\allowbreak{}3\par Tipo: NONE\par Confianza: ND & COMP-MESON\par E04 / M08\\\addlinespace[3pt]
% VI-CATALOG-ROW anchuras 353
\textbf{353. D\_\allowbreak{}0*(2300) WIDTH}\par D\_\allowbreak{}0\textasciicircum{}*(2300)\allowbreak{}0 | Dbar\_\allowbreak{}0\textasciicircum{}*(2300)\allowbreak{}0 | D\_\allowbreak{}0\textasciicircum{}*(2300)\allowbreak{}+ | D\_\allowbreak{}0\textasciicircum{}*(2300)\allowbreak{}-\par {\fontsize{8.1}{9.4}\selectfont PDG2026:\allowbreak{}observable:\allowbreak{}M252:\allowbreak{}M252W:\allowbreak{}337814} & \textbf{229+-16} MeV & \(u_+\): 16.01488\allowbreak{}68979557\allowbreak{}51\par \(u_-\): 16.01488\allowbreak{}68979557\allowbreak{}51\par Tipo: NONE\par Confianza: ND & COMP-MESON\par E04 / M08\\\addlinespace[3pt]
% VI-CATALOG-ROW anchuras 354
\textbf{354. D\_\allowbreak{}1(2420) WIDTH}\par D\_\allowbreak{}1(2420)\allowbreak{}0 | Dbar\_\allowbreak{}1(2420)\allowbreak{}0 | D\_\allowbreak{}1(2420)\allowbreak{}+ | D\_\allowbreak{}1(2420)\allowbreak{}-\par {\fontsize{8.1}{9.4}\selectfont PDG2026:\allowbreak{}observable:\allowbreak{}M253:\allowbreak{}M253W:\allowbreak{}337822} & \textbf{31.3+-1.9} MeV & \(u_+\): 1.939893\allowbreak{}76648816\allowbreak{}9\par \(u_-\): 1.931511\allowbreak{}99434054\allowbreak{}9\par Tipo: NONE\par Confianza: ND & COMP-MESON\par E04 / M08\\\addlinespace[3pt]
% VI-CATALOG-ROW anchuras 355
\textbf{355. D\_\allowbreak{}2*(2460) WIDTH}\par D\_\allowbreak{}2\textasciicircum{}*(2460)\allowbreak{}0 | Dbar\_\allowbreak{}2\textasciicircum{}*(2460)\allowbreak{}0 | D\_\allowbreak{}2\textasciicircum{}*(2460)\allowbreak{}+ | D\_\allowbreak{}2\textasciicircum{}*(2460)\allowbreak{}-\par {\fontsize{8.1}{9.4}\selectfont PDG2026:\allowbreak{}observable:\allowbreak{}M254:\allowbreak{}M254W:\allowbreak{}337837} & \textbf{47.3+-0.8} MeV & \(u_+\): 0.774313\allowbreak{}47668341\allowbreak{}055\par \(u_-\): 0.773629\allowbreak{}57674648\allowbreak{}969\par Tipo: NONE\par Confianza: ND & COMP-MESON\par E04 / M08\\\addlinespace[3pt]
% VI-CATALOG-ROW anchuras 356
\textbf{356. X(1750) WIDTH}\par X(1750)\allowbreak{}0\par {\fontsize{8.1}{9.4}\selectfont PDG2026:\allowbreak{}observable:\allowbreak{}M255:\allowbreak{}M255W:\allowbreak{}334942} & \textbf{120+-10} MeV & \(u_+\): 9.290122\allowbreak{}56224712\allowbreak{}55\par \(u_-\): 9.822200\allowbreak{}80272760\allowbreak{}93\par Tipo: NONE\par Confianza: ND & COMP-MESON\par E04 / M08\\\addlinespace[3pt]
% VI-CATALOG-ROW anchuras 357
\textbf{357. D\_\allowbreak{}s0(2590)\allowbreak{}+ WIDTH}\par D\_\allowbreak{}s0(2590)\allowbreak{}+ | D\_\allowbreak{}s0(2590)\allowbreak{}-\par {\fontsize{8.1}{9.4}\selectfont PDG2026:\allowbreak{}observable:\allowbreak{}M256:\allowbreak{}M256W:\allowbreak{}338521} & \textbf{89+-20} MeV & \(u_+\): 20\par \(u_-\): 20\par Tipo: NONE\par Confianza: ND & COMP-MESON\par E04 / M08\\\addlinespace[3pt]
% VI-CATALOG-ROW anchuras 358
\textbf{358. B\_\allowbreak{}sJ(6063)\allowbreak{}0 WIDTH}\par B\_\allowbreak{}sJ(6063)\allowbreak{}0 | Bbar\_\allowbreak{}sJ(6063)\allowbreak{}0\par {\fontsize{8.1}{9.4}\selectfont PDG2026:\allowbreak{}observable:\allowbreak{}M257:\allowbreak{}M257W:\allowbreak{}343154} & \textbf{26+-6} MeV & \(u_+\): 5.656854\allowbreak{}24949238\allowbreak{}06\par \(u_-\): 5.656854\allowbreak{}24949238\allowbreak{}06\par Tipo: NONE\par Confianza: ND & COMP-MESON\par E04 / M08\\\addlinespace[3pt]
% VI-CATALOG-ROW anchuras 359
\textbf{359. B\_\allowbreak{}sJ(6114)\allowbreak{}0 WIDTH}\par B\_\allowbreak{}sJ(6114)\allowbreak{}0 | Bbar\_\allowbreak{}sJ(6114)\allowbreak{}0\par {\fontsize{8.1}{9.4}\selectfont PDG2026:\allowbreak{}observable:\allowbreak{}M258:\allowbreak{}M258W:\allowbreak{}343158} & \textbf{66+-28} MeV & \(u_+\): 27.65863\allowbreak{}33718786\allowbreak{}61\par \(u_-\): 27.65863\allowbreak{}33718786\allowbreak{}61\par Tipo: NONE\par Confianza: ND & COMP-MESON\par E04 / M08\\\addlinespace[3pt]
% VI-CATALOG-ROW anchuras 360
\textbf{360. T\_\allowbreak{}c cbar sbar 1(4000) WIDTH}\par T\_\allowbreak{}c\_\allowbreak{}cbar\_\allowbreak{}sbar\_\allowbreak{}1(4000)\allowbreak{}0\par {\fontsize{8.1}{9.4}\selectfont PDG2026:\allowbreak{}observable:\allowbreak{}M259:\allowbreak{}M259W:\allowbreak{}347965} & \textbf{5 -\kern0pt{}- 150} MeV\par Presentación: 5-\kern0pt{}-150 & \(u_+\): NA\par \(u_-\): NA\par Tipo: R\par Confianza: ND & COMP-MESON\par E04 / M08\\\addlinespace[3pt]
% VI-CATALOG-ROW anchuras 361
\textbf{361. T\_\allowbreak{}c cbar sbar 1(4220)\allowbreak{}+ WIDTH}\par T\_\allowbreak{}c\_\allowbreak{}cbar\_\allowbreak{}sbar\_\allowbreak{}1(4220)\allowbreak{}+\par {\fontsize{8.1}{9.4}\selectfont PDG2026:\allowbreak{}observable:\allowbreak{}M260:\allowbreak{}M260W:\allowbreak{}348019} & \textbf{230+110-90} MeV & \(u_+\): 110.0590\allowbreak{}75046086\allowbreak{}09\par \(u_-\): 89.62700\allowbreak{}48590267\allowbreak{}27\par Tipo: NONE\par Confianza: ND & COMP-MESON\par E04 / M08\\\addlinespace[3pt]
% VI-CATALOG-ROW anchuras 362
\textbf{362. chi\_\allowbreak{}c1(4685) WIDTH}\par chi\_\allowbreak{}c1(4685)\par {\fontsize{8.1}{9.4}\selectfont PDG2026:\allowbreak{}observable:\allowbreak{}M261:\allowbreak{}M261W:\allowbreak{}346811} & \textbf{180+-50} MeV & \(u_+\): 49.86899\allowbreak{}22145262\allowbreak{}23\par \(u_-\): 51.68253\allowbreak{}10377275\allowbreak{}5\par Tipo: NONE\par Confianza: ND & COMP-MESON\par E04 / M08\\\addlinespace[3pt]
% VI-CATALOG-ROW anchuras 363
\textbf{363. X(4630) WIDTH}\par X(4630)\par {\fontsize{8.1}{9.4}\selectfont PDG2026:\allowbreak{}observable:\allowbreak{}M262:\allowbreak{}M262W:\allowbreak{}346762} & \textbf{350+140-100} MeV & \(u_+\): 137.9716\allowbreak{}56296596\allowbreak{}31\par \(u_-\): 102.5281\allowbreak{}69867227\allowbreak{}7\par Tipo: NONE\par Confianza: ND & COMP-MESON\par E04 / M08\\\addlinespace[3pt]
% VI-CATALOG-ROW anchuras 364
\textbf{364. a\_\allowbreak{}0(1710) WIDTH}\par a\_\allowbreak{}0(1710)\allowbreak{}0\par {\fontsize{8.1}{9.4}\selectfont PDG2026:\allowbreak{}observable:\allowbreak{}M263:\allowbreak{}M263W:\allowbreak{}334918} & \textbf{107+-15} MeV & \(u_+\): 14.82500\allowbreak{}71785796\allowbreak{}8\par \(u_-\): 14.82500\allowbreak{}71785796\allowbreak{}8\par Tipo: NONE\par Confianza: ND & COMP-MESON\par E04 / M08\\\addlinespace[3pt]
% VI-CATALOG-ROW anchuras 365
\textbf{365. f\_\allowbreak{}0(1810) Breit-Wigner WIDTH}\par f\_\allowbreak{}0(1810)\allowbreak{}0\par {\fontsize{8.1}{9.4}\selectfont PDG2026:\allowbreak{}observable:\allowbreak{}M264:\allowbreak{}M264W:\allowbreak{}334993} & \textbf{137+15-17} MeV & \(u_+\): 15.02739\allowbreak{}35810542\allowbreak{}31\par \(u_-\): 16.78598\allowbreak{}3289291\par Tipo: NONE\par Confianza: ND & COMP-MESON\par E04 / M08\\\addlinespace[3pt]
% VI-CATALOG-ROW anchuras 366
\textbf{366. T\_\allowbreak{}c\_\allowbreak{}c(3875)\allowbreak{}+ WIDTH}\par T\_\allowbreak{}c\_\allowbreak{}c(3875)\allowbreak{}+\par {\fontsize{8.1}{9.4}\selectfont PDG2026:\allowbreak{}observable:\allowbreak{}M265:\allowbreak{}M265W:\allowbreak{}347935} & \textbf{0.048+0.0020-0.0141} MeV & \(u_+\): 0.002\par \(u_-\): 0.014142\allowbreak{}13562373\allowbreak{}0951\par Tipo: NONE\par Confianza: ND & COMP-MESON\par E04 / M08\\\addlinespace[3pt]
% VI-CATALOG-ROW anchuras 367
\textbf{367. f\_\allowbreak{}0(2470) WIDTH}\par f\_\allowbreak{}0(2470)\allowbreak{}0\par {\fontsize{8.1}{9.4}\selectfont PDG2026:\allowbreak{}observable:\allowbreak{}M266:\allowbreak{}M266W:\allowbreak{}335274} & \textbf{75+14-12} MeV & \(u_+\): 14.21267\allowbreak{}04035519\par \(u_-\): 12.04159\allowbreak{}45787922\allowbreak{}99\par Tipo: NONE\par Confianza: ND & COMP-MESON\par E04 / M08\\\addlinespace[3pt]
% VI-CATALOG-ROW anchuras 368
\textbf{368. eta\_\allowbreak{}1(1855) WIDTH}\par eta\_\allowbreak{}1(1855)\allowbreak{}0\par {\fontsize{8.1}{9.4}\selectfont PDG2026:\allowbreak{}observable:\allowbreak{}M267:\allowbreak{}M267W:\allowbreak{}335029} & \textbf{188+-19} MeV & \(u_+\): 18.24828\allowbreak{}75908946\allowbreak{}59\par \(u_-\): 19.69771\allowbreak{}56035922\allowbreak{}11\par Tipo: NONE\par Confianza: ND & COMP-MESON\par E04 / M08\\\addlinespace[3pt]
% VI-CATALOG-ROW anchuras 369
\textbf{369. T\_\allowbreak{}c\_\allowbreak{}c\_\allowbreak{}cbar\_\allowbreak{}cbar(6900)\allowbreak{}0 WIDTH}\par T\_\allowbreak{}c\_\allowbreak{}c\_\allowbreak{}cbar\_\allowbreak{}cbar(6900)\allowbreak{}+\par {\fontsize{8.1}{9.4}\selectfont PDG2026:\allowbreak{}observable:\allowbreak{}M268:\allowbreak{}M268W:\allowbreak{}348045} & \textbf{161+-26} MeV & \(u_+\): 27.00038\allowbreak{}42452260\allowbreak{}8\par \(u_-\): 25.46015\allowbreak{}45852605\allowbreak{}79\par Tipo: NONE\par Confianza: ND & COMP-MESON\par E04 / M08\\\addlinespace[3pt]
% VI-CATALOG-ROW anchuras 370
\textbf{370. T\_\allowbreak{}c\_\allowbreak{}sbar\_\allowbreak{}0*(2900)\allowbreak{}0 WIDTH}\par T\_\allowbreak{}c\_\allowbreak{}sbar\_\allowbreak{}0\textasciicircum{}*(2900)\allowbreak{}+\par {\fontsize{8.1}{9.4}\selectfont PDG2026:\allowbreak{}observable:\allowbreak{}M269:\allowbreak{}M269W:\allowbreak{}347928} & \textbf{119+-29} MeV & \(u_+\): 29.06888\allowbreak{}37074972\allowbreak{}71\par \(u_-\): 29.06888\allowbreak{}37074972\allowbreak{}71\par Tipo: NONE\par Confianza: ND & COMP-MESON\par E04 / M08\\\addlinespace[3pt]
% VI-CATALOG-ROW anchuras 371
\textbf{371. T\_\allowbreak{}c\_\allowbreak{}sbar\_\allowbreak{}0*(2900)\allowbreak{}++ WIDTH}\par T\_\allowbreak{}c\_\allowbreak{}sbar\_\allowbreak{}0\textasciicircum{}*(2900)\allowbreak{}+\par {\fontsize{8.1}{9.4}\selectfont PDG2026:\allowbreak{}observable:\allowbreak{}M269:\allowbreak{}M269W++:\allowbreak{}347929} & \textbf{140+-40} MeV & \(u_+\): 36.23534\allowbreak{}18639868\allowbreak{}68\par \(u_-\): 36.23534\allowbreak{}18639868\allowbreak{}68\par Tipo: NONE\par Confianza: ND & COMP-MESON\par E04 / M08\\\addlinespace[3pt]
% VI-CATALOG-ROW anchuras 372
\textbf{372. omega(2220) WIDTH}\par omega(2220)\allowbreak{}0\par {\fontsize{8.1}{9.4}\selectfont PDG2026:\allowbreak{}observable:\allowbreak{}M270:\allowbreak{}M270W:\allowbreak{}335222} & \textbf{105+-34} MeV & \(u_+\): 33.80726\allowbreak{}76977683\allowbreak{}03\par \(u_-\): 33.80726\allowbreak{}76977683\allowbreak{}03\par Tipo: NONE\par Confianza: ND & COMP-MESON\par E04 / M08\\\addlinespace[3pt]
% VI-CATALOG-ROW anchuras 373
\textbf{373. T\_\allowbreak{}c c cbar cbar(6600) WIDTH}\par T\_\allowbreak{}c c cbar cbar(6600)\allowbreak{}+\par {\fontsize{8.1}{9.4}\selectfont PDG2026:\allowbreak{}observable:\allowbreak{}M272:\allowbreak{}M272W:\allowbreak{}348041} & \textbf{440+-70} MeV & \(u_+\): 74.67698\allowbreak{}84366600\allowbreak{}39\par \(u_-\): 68.96618\allowbreak{}94166342\allowbreak{}36\par Tipo: NONE\par Confianza: ND & COMP-MESON\par E04 / M08\\\addlinespace[3pt]
% VI-CATALOG-ROW anchuras 374
\textbf{374. T\_\allowbreak{}c c cbar cbar(7100) WIDTH}\par T\_\allowbreak{}c c cbar cbar(7100)\allowbreak{}+\par {\fontsize{8.1}{9.4}\selectfont PDG2026:\allowbreak{}observable:\allowbreak{}M273:\allowbreak{}M273W:\allowbreak{}348049} & \textbf{100+50-40} MeV & \(u_+\): 49.40647\allowbreak{}73081424\allowbreak{}99\par \(u_-\): 38.94868\allowbreak{}41883008\allowbreak{}88\par Tipo: NONE\par Confianza: ND & COMP-MESON\par E04 / M08\\\addlinespace[3pt]
% VI-CATALOG-ROW anchuras 375
\textbf{375. t-quark DECAY WIDTH}\par t | tbar\par {\fontsize{8.1}{9.4}\selectfont PDG2026:\allowbreak{}observable:\allowbreak{}Q007:\allowbreak{}Q007W:\allowbreak{}333618} & \textbf{1.42+0.19-0.15} GeV & \(u_+\): 0.190513\allowbreak{}42973542\allowbreak{}351\par \(u_-\): 0.152560\allowbreak{}41171179\allowbreak{}069\par Tipo: NONE\par Confianza: ND & QUARK\_\allowbreak{}EXTENDED\par E04 / M09\\\addlinespace[3pt]
% VI-CATALOG-ROW anchuras 376
\textbf{376. eta WIDTH}\par eta\par {\fontsize{8.1}{9.4}\selectfont PDG2026:\allowbreak{}observable:\allowbreak{}S014:\allowbreak{}S014W:\allowbreak{}333730} & \textbf{1.31+-0.05} keV & \(u_+\): 0.045198\allowbreak{}90807919\allowbreak{}2503\par \(u_-\): 0.045198\allowbreak{}90807919\allowbreak{}2503\par Tipo: NONE\par Confianza: ND & ELEMEN\allowbreak{}TARY\_\allowbreak{}OR\_\allowbreak{}EXTENDED\par E04 / M05\\\addlinespace[3pt]
% VI-CATALOG-ROW anchuras 377
\textbf{377. Sigma\_\allowbreak{}b+ WIDTH}\par Sigma\_\allowbreak{}b()\allowbreak{}- | Sigmabar\_\allowbreak{}b()\allowbreak{}+ | Sigma\_\allowbreak{}b()\allowbreak{}0 | Sigmabar\_\allowbreak{}b()\allowbreak{}0 | Sigma\_\allowbreak{}b()\allowbreak{}+ | Sigmabar\_\allowbreak{}b()\allowbreak{}-\par {\fontsize{8.1}{9.4}\selectfont PDG2026:\allowbreak{}observable:\allowbreak{}S026:\allowbreak{}S026W+:\allowbreak{}351720} & \textbf{5.0+-0.5} MeV & \(u_+\): 0.479198\allowbreak{}01016756\allowbreak{}203\par \(u_-\): 0.476604\allowbreak{}43013974\allowbreak{}058\par Tipo: NONE\par Confianza: ND & ELEMEN\allowbreak{}TARY\_\allowbreak{}OR\_\allowbreak{}EXTENDED\par E04 / M05\\\addlinespace[3pt]
% VI-CATALOG-ROW anchuras 378
\textbf{378. Sigma\_\allowbreak{}b- WIDTH}\par Sigma\_\allowbreak{}b()\allowbreak{}- | Sigmabar\_\allowbreak{}b()\allowbreak{}+ | Sigma\_\allowbreak{}b()\allowbreak{}0 | Sigmabar\_\allowbreak{}b()\allowbreak{}0 | Sigma\_\allowbreak{}b()\allowbreak{}+ | Sigmabar\_\allowbreak{}b()\allowbreak{}-\par {\fontsize{8.1}{9.4}\selectfont PDG2026:\allowbreak{}observable:\allowbreak{}S026:\allowbreak{}S026W-:\allowbreak{}351721} & \textbf{5.3+-0.5} MeV & \(u_+\): 0.551800\allowbreak{}22333948\allowbreak{}409\par \(u_-\): 0.544753\allowbreak{}71590490\allowbreak{}007\par Tipo: NONE\par Confianza: ND & ELEMEN\allowbreak{}TARY\_\allowbreak{}OR\_\allowbreak{}EXTENDED\par E04 / M05\\\addlinespace[3pt]
% VI-CATALOG-ROW anchuras 379
\textbf{379. W WIDTH}\par W+ | W-\par {\fontsize{8.1}{9.4}\selectfont PDG2026:\allowbreak{}observable:\allowbreak{}S043:\allowbreak{}S043W:\allowbreak{}332468} & \textbf{2.14+-0.05} GeV & \(u_+\): 0.053656\allowbreak{}22500314\allowbreak{}9413\par \(u_-\): 0.053656\allowbreak{}22500314\allowbreak{}9413\par Tipo: NONE\par Confianza: ND & ELEMEN\allowbreak{}TARY\_\allowbreak{}OR\_\allowbreak{}EXTENDED\par E04 / M05\\\addlinespace[3pt]
% VI-CATALOG-ROW anchuras 380
\textbf{380. Z WIDTH}\par Z0\par {\fontsize{8.1}{9.4}\selectfont PDG2026:\allowbreak{}observable:\allowbreak{}S044:\allowbreak{}S044W:\allowbreak{}332514} & \textbf{2.4955+-0.0023} GeV & \(u_+\): 0.0023\par \(u_-\): 0.0023\par Tipo: NONE\par Confianza: ND & ELEMEN\allowbreak{}TARY\_\allowbreak{}OR\_\allowbreak{}EXTENDED\par E04 / M05\\\addlinespace[3pt]
% VI-CATALOG-ROW anchuras 381
\textbf{381. Sigma\_\allowbreak{}b*+ WIDTH}\par Sigma\_\allowbreak{}b\textasciicircum{}*()\allowbreak{}- | Sigmabar\_\allowbreak{}b\textasciicircum{}*()\allowbreak{}+ | Sigma\_\allowbreak{}b\textasciicircum{}*()\allowbreak{}0 | Sigmabar\_\allowbreak{}b\textasciicircum{}*()\allowbreak{}0 | Sigma\_\allowbreak{}b\textasciicircum{}*()\allowbreak{}+ | Sigmabar\_\allowbreak{}b\textasciicircum{}*()\allowbreak{}-\par {\fontsize{8.1}{9.4}\selectfont PDG2026:\allowbreak{}observable:\allowbreak{}S062:\allowbreak{}S062W+:\allowbreak{}351729} & \textbf{9.4+-0.5} MeV & \(u_+\): 0.528012\allowbreak{}85527312\allowbreak{}969\par \(u_-\): 0.526516\allowbreak{}49982212\allowbreak{}589\par Tipo: NONE\par Confianza: ND & ELEMEN\allowbreak{}TARY\_\allowbreak{}OR\_\allowbreak{}EXTENDED\par E04 / M05\\\addlinespace[3pt]
% VI-CATALOG-ROW anchuras 382
\textbf{382. Sigma\_\allowbreak{}b*- WIDTH}\par Sigma\_\allowbreak{}b\textasciicircum{}*()\allowbreak{}- | Sigmabar\_\allowbreak{}b\textasciicircum{}*()\allowbreak{}+ | Sigma\_\allowbreak{}b\textasciicircum{}*()\allowbreak{}0 | Sigmabar\_\allowbreak{}b\textasciicircum{}*()\allowbreak{}0 | Sigma\_\allowbreak{}b\textasciicircum{}*()\allowbreak{}+ | Sigmabar\_\allowbreak{}b\textasciicircum{}*()\allowbreak{}-\par {\fontsize{8.1}{9.4}\selectfont PDG2026:\allowbreak{}observable:\allowbreak{}S062:\allowbreak{}S062W-:\allowbreak{}351730} & \textbf{10.4+-0.8} MeV & \(u_+\): 0.845898\allowbreak{}24638254\allowbreak{}64\par \(u_-\): 0.843107\allowbreak{}30525944\allowbreak{}599\par Tipo: NONE\par Confianza: ND & ELEMEN\allowbreak{}TARY\_\allowbreak{}OR\_\allowbreak{}EXTENDED\par E04 / M05\\\addlinespace[3pt]
% VI-CATALOG-ROW anchuras 383
\textbf{383. D\_\allowbreak{}s*+- WIDTH}\par D\_\allowbreak{}s\textasciicircum{}*()\allowbreak{}+ | D\_\allowbreak{}s\textasciicircum{}*()\allowbreak{}-\par {\fontsize{8.1}{9.4}\selectfont PDG2026:\allowbreak{}observable:\allowbreak{}S074:\allowbreak{}S074W:\allowbreak{}338419} & \textbf{<1.9} MeV & \(u_+\): NA\par \(u_-\): NA\par Tipo: U\par Confianza: 90 & ELEMEN\allowbreak{}TARY\_\allowbreak{}OR\_\allowbreak{}EXTENDED\par E04 / M05\\\addlinespace[3pt]
% VI-CATALOG-ROW anchuras 384
\textbf{384. H DECAY WIDTH}\par H\par {\fontsize{8.1}{9.4}\selectfont PDG2026:\allowbreak{}observable:\allowbreak{}S126:\allowbreak{}S126W:\allowbreak{}332734} & \textbf{3.0+1.5-0.7} MeV & \(u_+\): 1.452973\allowbreak{}69001768\allowbreak{}31\par \(u_-\): 0.715011\allowbreak{}91940991\allowbreak{}071\par Tipo: NONE\par Confianza: ND & ELEMEN\allowbreak{}TARY\_\allowbreak{}OR\_\allowbreak{}EXTENDED\par E04 / M05\\\addlinespace[3pt]
\end{longtable}\endgroup
\section{Procedencia y correspondencia de registros}\label{sec:trace}
Este repertorio conserva los códigos literales, las huellas y los ordinales de ambos registros de cada ruta. Los enlaces devuelven a la ficha correspondiente. Las tres claves siguientes son comunes a las 324 rutas y su significado se explica en la guía de lectura.
\par\noindent\textbf{route\_\allowbreak{}status}\par\noindent TARGET\_\allowbreak{}FREE\_\allowbreak{}INTERNAL\_\allowbreak{}ROUTE
\par\noindent\textbf{join\_\allowbreak{}status}\par\noindent EXACT\_\allowbreak{}KEY\_\allowbreak{}JOIN\_\allowbreak{}\_\allowbreak{}NO\_\allowbreak{}EXTERNAL\_\allowbreak{}MASS\_\allowbreak{}READ
\par\noindent\textbf{tower\_\allowbreak{}forward\_\allowbreak{}status}\par\noindent ROUTE\_\allowbreak{}MEMORY\_\allowbreak{}AVAILABLE\_\allowbreak{}\_\allowbreak{}ETA\_\allowbreak{}TRANSDUCER\_\allowbreak{}REQUIRES\_\allowbreak{}CERTIFICATE
\par\medskip Equivalencias de clasificación:
charged\_\allowbreak{}lepton = leptón cargado\quad 
neutrino = neutrino\quad 
quark\_\allowbreak{}down = quark de tipo abajo\quad 
quark\_\allowbreak{}up = quark de tipo arriba\quad 
lepton = leptón\quad 
quark = quark\quad 
G0\_\allowbreak{}center = G0, centro\quad 
True = sí\quad 
False = no\quad 
YES = sí\quad 
NO = no\quad 
\par\medskip
\Needspace{87pt}\phantomsection\label{trace-1}\noindent\hyperref[nav:vi-ruta-1]{\textbf{r1c1\_\allowbreak{}h-1v-1}}\quad Ordinal unificado / memoria: 1 / 4\par\nopagebreak
\begingroup\fontsize{8.3}{10.1}\selectfont\setlength{\tabcolsep}{4pt}\noindent\begin{tabularx}{\linewidth}{@{}p{.21\linewidth}X@{}}
Genealogía & effeba54\allowbreak{}64fe6134\allowbreak{}ce328276\allowbreak{}caffda45\allowbreak{}7de9d85e\allowbreak{}22f3932a\allowbreak{}a8ce593a\allowbreak{}521c1040\\
Lector & 32867437\allowbreak{}928eda79\allowbreak{}ef832344\allowbreak{}d5cde588\allowbreak{}6a4c715e\allowbreak{}968ad777\allowbreak{}46df4b89\allowbreak{}0958246e\\
Registro en 108 & \textemdash{}\\
\end{tabularx}\endgroup\par\vspace{6pt}
\Needspace{87pt}\phantomsection\label{trace-2}\noindent\hyperref[nav:vi-ruta-2]{\textbf{r1c1\_\allowbreak{}h-1v+1}}\quad Ordinal unificado / memoria: 2 / 3\par\nopagebreak
\begingroup\fontsize{8.3}{10.1}\selectfont\setlength{\tabcolsep}{4pt}\noindent\begin{tabularx}{\linewidth}{@{}p{.21\linewidth}X@{}}
Genealogía & 7cc1cc9f\allowbreak{}41e40afb\allowbreak{}8776ff84\allowbreak{}5297b2b0\allowbreak{}b8b440b8\allowbreak{}724f822c\allowbreak{}c6fca90c\allowbreak{}7639f2df\\
Lector & 85737ea3\allowbreak{}9e08488f\allowbreak{}4972aa61\allowbreak{}178a8ccb\allowbreak{}edcd6662\allowbreak{}f8b923a8\allowbreak{}009c6ffe\allowbreak{}ba7c5e7b\\
Registro en 108 & 5e13be75\allowbreak{}1c11cc8b\allowbreak{}f7ef7949\allowbreak{}82848590\allowbreak{}e3ae4f72\allowbreak{}73685430\allowbreak{}f5f6144b\allowbreak{}fc16ebd7\\
\end{tabularx}\endgroup\par\vspace{6pt}
\Needspace{87pt}\phantomsection\label{trace-3}\noindent\hyperref[nav:vi-ruta-3]{\textbf{r1c1\_\allowbreak{}h+1v-1}}\quad Ordinal unificado / memoria: 3 / 2\par\nopagebreak
\begingroup\fontsize{8.3}{10.1}\selectfont\setlength{\tabcolsep}{4pt}\noindent\begin{tabularx}{\linewidth}{@{}p{.21\linewidth}X@{}}
Genealogía & aabfcc60\allowbreak{}99652e23\allowbreak{}f82f1a50\allowbreak{}e4ad66a2\allowbreak{}8960ccdb\allowbreak{}0db84037\allowbreak{}0dfdf043\allowbreak{}e39dd2d5\\
Lector & ecd7e361\allowbreak{}6b9a9821\allowbreak{}dc08227b\allowbreak{}0714d924\allowbreak{}426ca0d1\allowbreak{}3867b735\allowbreak{}ffb5c614\allowbreak{}b0ed3390\\
Registro en 108 & \textemdash{}\\
\end{tabularx}\endgroup\par\vspace{6pt}
\Needspace{87pt}\phantomsection\label{trace-4}\noindent\hyperref[nav:vi-ruta-4]{\textbf{r1c1\_\allowbreak{}h+1v+1}}\quad Ordinal unificado / memoria: 4 / 1\par\nopagebreak
\begingroup\fontsize{8.3}{10.1}\selectfont\setlength{\tabcolsep}{4pt}\noindent\begin{tabularx}{\linewidth}{@{}p{.21\linewidth}X@{}}
Genealogía & cb1834fa\allowbreak{}28511686\allowbreak{}ade62a63\allowbreak{}27bd1ef9\allowbreak{}daabe103\allowbreak{}193e9e71\allowbreak{}583b231c\allowbreak{}288f5657\\
Lector & ce86abc7\allowbreak{}9b701a0e\allowbreak{}c9ae2123\allowbreak{}d730f8c5\allowbreak{}eb0e88da\allowbreak{}f5bfa926\allowbreak{}37854368\allowbreak{}8333fd17\\
Registro en 108 & \textemdash{}\\
\end{tabularx}\endgroup\par\vspace{6pt}
\Needspace{87pt}\phantomsection\label{trace-5}\noindent\hyperref[nav:vi-ruta-5]{\textbf{r1c2\_\allowbreak{}h-1v-1}}\quad Ordinal unificado / memoria: 5 / 8\par\nopagebreak
\begingroup\fontsize{8.3}{10.1}\selectfont\setlength{\tabcolsep}{4pt}\noindent\begin{tabularx}{\linewidth}{@{}p{.21\linewidth}X@{}}
Genealogía & 7334be27\allowbreak{}187ca745\allowbreak{}a3c7e70d\allowbreak{}9f4fea4f\allowbreak{}b6d8c31f\allowbreak{}990e2382\allowbreak{}f7f91263\allowbreak{}fb447d0c\\
Lector & 7677699b\allowbreak{}2a6f6153\allowbreak{}d139cc02\allowbreak{}ebcac4a1\allowbreak{}ff7f1894\allowbreak{}b6da0a7b\allowbreak{}1e46b74f\allowbreak{}55110f6a\\
Registro en 108 & e589bca4\allowbreak{}743ae3a8\allowbreak{}412d609f\allowbreak{}38a1cbca\allowbreak{}9ab5e815\allowbreak{}21c621f6\allowbreak{}51183be2\allowbreak{}730acc06\\
\end{tabularx}\endgroup\par\vspace{6pt}
\Needspace{87pt}\phantomsection\label{trace-6}\noindent\hyperref[nav:vi-ruta-6]{\textbf{r1c2\_\allowbreak{}h-1v+1}}\quad Ordinal unificado / memoria: 6 / 7\par\nopagebreak
\begingroup\fontsize{8.3}{10.1}\selectfont\setlength{\tabcolsep}{4pt}\noindent\begin{tabularx}{\linewidth}{@{}p{.21\linewidth}X@{}}
Genealogía & 56149b04\allowbreak{}962ed2f1\allowbreak{}a098af0d\allowbreak{}cf0fc19f\allowbreak{}f01c1fa5\allowbreak{}825b32f8\allowbreak{}99d69dea\allowbreak{}426f1ddb\\
Lector & 474a5b8c\allowbreak{}bea67b1d\allowbreak{}3d0ba8e2\allowbreak{}89cb7ef3\allowbreak{}0ed5d908\allowbreak{}887d6ba4\allowbreak{}6fa80433\allowbreak{}e81f6993\\
Registro en 108 & \textemdash{}\\
\end{tabularx}\endgroup\par\vspace{6pt}
\Needspace{87pt}\phantomsection\label{trace-7}\noindent\hyperref[nav:vi-ruta-7]{\textbf{r1c2\_\allowbreak{}h+1v-1}}\quad Ordinal unificado / memoria: 7 / 6\par\nopagebreak
\begingroup\fontsize{8.3}{10.1}\selectfont\setlength{\tabcolsep}{4pt}\noindent\begin{tabularx}{\linewidth}{@{}p{.21\linewidth}X@{}}
Genealogía & ccb61e50\allowbreak{}f720be11\allowbreak{}f86bf136\allowbreak{}d02b3658\allowbreak{}7f92d52c\allowbreak{}35df35bf\allowbreak{}4b86ef5d\allowbreak{}779f7bcf\\
Lector & e7c7e699\allowbreak{}830c4f17\allowbreak{}e88719d1\allowbreak{}51c262a7\allowbreak{}75f5cf6a\allowbreak{}fccadf3b\allowbreak{}9ce34e42\allowbreak{}6bc57b55\\
Registro en 108 & \textemdash{}\\
\end{tabularx}\endgroup\par\vspace{6pt}
\Needspace{87pt}\phantomsection\label{trace-8}\noindent\hyperref[nav:vi-ruta-8]{\textbf{r1c2\_\allowbreak{}h+1v+1}}\quad Ordinal unificado / memoria: 8 / 5\par\nopagebreak
\begingroup\fontsize{8.3}{10.1}\selectfont\setlength{\tabcolsep}{4pt}\noindent\begin{tabularx}{\linewidth}{@{}p{.21\linewidth}X@{}}
Genealogía & 72b9df13\allowbreak{}97f2987a\allowbreak{}c2c70e5f\allowbreak{}7163572e\allowbreak{}0b885482\allowbreak{}5eee3b61\allowbreak{}db397f28\allowbreak{}ba9b9ba6\\
Lector & a1c4553c\allowbreak{}b5d2563a\allowbreak{}0a9a4d4d\allowbreak{}4659a3d9\allowbreak{}c9316bb9\allowbreak{}eb975446\allowbreak{}70b4f842\allowbreak{}25f0a4f5\\
Registro en 108 & \textemdash{}\\
\end{tabularx}\endgroup\par\vspace{6pt}
\Needspace{87pt}\phantomsection\label{trace-9}\noindent\hyperref[nav:vi-ruta-9]{\textbf{r1c3\_\allowbreak{}h-1v-1}}\quad Ordinal unificado / memoria: 9 / 12\par\nopagebreak
\begingroup\fontsize{8.3}{10.1}\selectfont\setlength{\tabcolsep}{4pt}\noindent\begin{tabularx}{\linewidth}{@{}p{.21\linewidth}X@{}}
Genealogía & fae367f8\allowbreak{}fb866c76\allowbreak{}b0300f5b\allowbreak{}d2cb0dbb\allowbreak{}059ab3e1\allowbreak{}63b6df3c\allowbreak{}5e46e685\allowbreak{}e4d9fa74\\
Lector & 716a4c36\allowbreak{}600c27b5\allowbreak{}09736fc5\allowbreak{}15d19437\allowbreak{}28ac458d\allowbreak{}23500627\allowbreak{}b24d70bd\allowbreak{}413dd168\\
Registro en 108 & \textemdash{}\\
\end{tabularx}\endgroup\par\vspace{6pt}
\Needspace{87pt}\phantomsection\label{trace-10}\noindent\hyperref[nav:vi-ruta-10]{\textbf{r1c3\_\allowbreak{}h-1v+1}}\quad Ordinal unificado / memoria: 10 / 11\par\nopagebreak
\begingroup\fontsize{8.3}{10.1}\selectfont\setlength{\tabcolsep}{4pt}\noindent\begin{tabularx}{\linewidth}{@{}p{.21\linewidth}X@{}}
Genealogía & 1120c2a5\allowbreak{}d9a6516e\allowbreak{}9138153d\allowbreak{}5365283c\allowbreak{}7af167d9\allowbreak{}e5b3e46c\allowbreak{}f7ab8d55\allowbreak{}a1974cfc\\
Lector & cc519544\allowbreak{}fb45c5f0\allowbreak{}387e629c\allowbreak{}3a067329\allowbreak{}125f213d\allowbreak{}0292e993\allowbreak{}5aa46d4b\allowbreak{}cf1027b8\\
Registro en 108 & b1cc79b3\allowbreak{}54042ce1\allowbreak{}877e1c4c\allowbreak{}f5347fc4\allowbreak{}f509c235\allowbreak{}77711204\allowbreak{}cf7cddec\allowbreak{}6beea6dc\\
\end{tabularx}\endgroup\par\vspace{6pt}
\Needspace{87pt}\phantomsection\label{trace-11}\noindent\hyperref[nav:vi-ruta-11]{\textbf{r1c3\_\allowbreak{}h+1v-1}}\quad Ordinal unificado / memoria: 11 / 10\par\nopagebreak
\begingroup\fontsize{8.3}{10.1}\selectfont\setlength{\tabcolsep}{4pt}\noindent\begin{tabularx}{\linewidth}{@{}p{.21\linewidth}X@{}}
Genealogía & d8f175e5\allowbreak{}7deb40be\allowbreak{}a0135957\allowbreak{}9d16ae6e\allowbreak{}1fdd0b45\allowbreak{}8ebd8b41\allowbreak{}608e5eba\allowbreak{}1782d060\\
Lector & 29304938\allowbreak{}102d2970\allowbreak{}60297c36\allowbreak{}2f9d8138\allowbreak{}0336ecff\allowbreak{}5a81f6af\allowbreak{}fe85c413\allowbreak{}e0670568\\
Registro en 108 & \textemdash{}\\
\end{tabularx}\endgroup\par\vspace{6pt}
\Needspace{87pt}\phantomsection\label{trace-12}\noindent\hyperref[nav:vi-ruta-12]{\textbf{r1c3\_\allowbreak{}h+1v+1}}\quad Ordinal unificado / memoria: 12 / 9\par\nopagebreak
\begingroup\fontsize{8.3}{10.1}\selectfont\setlength{\tabcolsep}{4pt}\noindent\begin{tabularx}{\linewidth}{@{}p{.21\linewidth}X@{}}
Genealogía & 13e342fa\allowbreak{}83db93f3\allowbreak{}f7dc1316\allowbreak{}3653112f\allowbreak{}aecbb1dd\allowbreak{}8aa9f872\allowbreak{}d9121b96\allowbreak{}007e8839\\
Lector & 5aa200a3\allowbreak{}69599173\allowbreak{}78a47853\allowbreak{}2c886d1a\allowbreak{}d9e4958c\allowbreak{}7cd2fe47\allowbreak{}6a7797fc\allowbreak{}5306366a\\
Registro en 108 & \textemdash{}\\
\end{tabularx}\endgroup\par\vspace{6pt}
\Needspace{87pt}\phantomsection\label{trace-13}\noindent\hyperref[nav:vi-ruta-13]{\textbf{r1c4\_\allowbreak{}h-1v-1}}\quad Ordinal unificado / memoria: 13 / 16\par\nopagebreak
\begingroup\fontsize{8.3}{10.1}\selectfont\setlength{\tabcolsep}{4pt}\noindent\begin{tabularx}{\linewidth}{@{}p{.21\linewidth}X@{}}
Genealogía & db632f76\allowbreak{}3af301d2\allowbreak{}46516155\allowbreak{}fafabc14\allowbreak{}d1d401a2\allowbreak{}9c7c7b63\allowbreak{}1471a8d8\allowbreak{}0c95591b\\
Lector & 5f1e5455\allowbreak{}7f63ee41\allowbreak{}bdd234e2\allowbreak{}0df6fb36\allowbreak{}7b00a650\allowbreak{}97cb906e\allowbreak{}b634fce7\allowbreak{}66250a20\\
Registro en 108 & \textemdash{}\\
\end{tabularx}\endgroup\par\vspace{6pt}
\Needspace{87pt}\phantomsection\label{trace-14}\noindent\hyperref[nav:vi-ruta-14]{\textbf{r1c4\_\allowbreak{}h-1v+1}}\quad Ordinal unificado / memoria: 14 / 15\par\nopagebreak
\begingroup\fontsize{8.3}{10.1}\selectfont\setlength{\tabcolsep}{4pt}\noindent\begin{tabularx}{\linewidth}{@{}p{.21\linewidth}X@{}}
Genealogía & aa732765\allowbreak{}9272d308\allowbreak{}9edbc150\allowbreak{}6b9b161a\allowbreak{}e47d6711\allowbreak{}976c8a2a\allowbreak{}4a116769\allowbreak{}39d76776\\
Lector & cb535d05\allowbreak{}a86f529c\allowbreak{}6b30d46b\allowbreak{}f87f9924\allowbreak{}e83bfc06\allowbreak{}9b9f421b\allowbreak{}49a8bed1\allowbreak{}d576b346\\
Registro en 108 & \textemdash{}\\
\end{tabularx}\endgroup\par\vspace{6pt}
\Needspace{87pt}\phantomsection\label{trace-15}\noindent\hyperref[nav:vi-ruta-15]{\textbf{r1c4\_\allowbreak{}h+1v-1}}\quad Ordinal unificado / memoria: 15 / 14\par\nopagebreak
\begingroup\fontsize{8.3}{10.1}\selectfont\setlength{\tabcolsep}{4pt}\noindent\begin{tabularx}{\linewidth}{@{}p{.21\linewidth}X@{}}
Genealogía & de3386f8\allowbreak{}b28db7e8\allowbreak{}c79bad44\allowbreak{}4dd9f558\allowbreak{}fb5ce614\allowbreak{}96e054c2\allowbreak{}ab24a718\allowbreak{}8cc309d4\\
Lector & e6136c80\allowbreak{}0ce4e9ca\allowbreak{}12e80a5e\allowbreak{}90dbb682\allowbreak{}78b0c098\allowbreak{}aafeab82\allowbreak{}6ff704d0\allowbreak{}6ce3b1d1\\
Registro en 108 & \textemdash{}\\
\end{tabularx}\endgroup\par\vspace{6pt}
\Needspace{87pt}\phantomsection\label{trace-16}\noindent\hyperref[nav:vi-ruta-16]{\textbf{r1c4\_\allowbreak{}h+1v+1}}\quad Ordinal unificado / memoria: 16 / 13\par\nopagebreak
\begingroup\fontsize{8.3}{10.1}\selectfont\setlength{\tabcolsep}{4pt}\noindent\begin{tabularx}{\linewidth}{@{}p{.21\linewidth}X@{}}
Genealogía & 1c8b5cda\allowbreak{}a63332d1\allowbreak{}e8f0ae42\allowbreak{}fd515992\allowbreak{}f5d9e41e\allowbreak{}72975eaa\allowbreak{}e772af08\allowbreak{}588c657a\\
Lector & 338da10d\allowbreak{}5c8906ec\allowbreak{}03bdfade\allowbreak{}56c89b7f\allowbreak{}d502a0b0\allowbreak{}68207406\allowbreak{}3adb644b\allowbreak{}01f77606\\
Registro en 108 & ce501e32\allowbreak{}503d7df0\allowbreak{}9e2c5ae6\allowbreak{}1ff9543f\allowbreak{}cb85dba7\allowbreak{}afa7a539\allowbreak{}6fb0f369\allowbreak{}09af84dc\\
\end{tabularx}\endgroup\par\vspace{6pt}
\Needspace{87pt}\phantomsection\label{trace-17}\noindent\hyperref[nav:vi-ruta-17]{\textbf{r1c5\_\allowbreak{}h-1v-1}}\quad Ordinal unificado / memoria: 17 / 20\par\nopagebreak
\begingroup\fontsize{8.3}{10.1}\selectfont\setlength{\tabcolsep}{4pt}\noindent\begin{tabularx}{\linewidth}{@{}p{.21\linewidth}X@{}}
Genealogía & e753e7cf\allowbreak{}912988ef\allowbreak{}a9df9eda\allowbreak{}0afb8e53\allowbreak{}2a8b8f33\allowbreak{}f9217225\allowbreak{}87fb3545\allowbreak{}51c5342c\\
Lector & 9d043f03\allowbreak{}a6f51024\allowbreak{}11e0e8ec\allowbreak{}bc080c07\allowbreak{}3548ea8c\allowbreak{}e98deaf7\allowbreak{}0e4bb5ae\allowbreak{}9fbed35e\\
Registro en 108 & \textemdash{}\\
\end{tabularx}\endgroup\par\vspace{6pt}
\Needspace{87pt}\phantomsection\label{trace-18}\noindent\hyperref[nav:vi-ruta-18]{\textbf{r1c5\_\allowbreak{}h-1v+1}}\quad Ordinal unificado / memoria: 18 / 19\par\nopagebreak
\begingroup\fontsize{8.3}{10.1}\selectfont\setlength{\tabcolsep}{4pt}\noindent\begin{tabularx}{\linewidth}{@{}p{.21\linewidth}X@{}}
Genealogía & 959d954a\allowbreak{}11c429bd\allowbreak{}cfe3a0db\allowbreak{}7b21c181\allowbreak{}0828fff0\allowbreak{}868a1d0a\allowbreak{}c26ba7ef\allowbreak{}3e127805\\
Lector & b90ff070\allowbreak{}e1aae58c\allowbreak{}fca71a9a\allowbreak{}e8253522\allowbreak{}544bb269\allowbreak{}f70b257e\allowbreak{}4e6afd6c\allowbreak{}0c3c24f8\\
Registro en 108 & \textemdash{}\\
\end{tabularx}\endgroup\par\vspace{6pt}
\Needspace{87pt}\phantomsection\label{trace-19}\noindent\hyperref[nav:vi-ruta-19]{\textbf{r1c5\_\allowbreak{}h+1v-1}}\quad Ordinal unificado / memoria: 19 / 18\par\nopagebreak
\begingroup\fontsize{8.3}{10.1}\selectfont\setlength{\tabcolsep}{4pt}\noindent\begin{tabularx}{\linewidth}{@{}p{.21\linewidth}X@{}}
Genealogía & 49c3be2e\allowbreak{}b41e01bb\allowbreak{}1a1f27f1\allowbreak{}b1ec147e\allowbreak{}4db99f16\allowbreak{}7f982dc7\allowbreak{}29a2d499\allowbreak{}3527c5f8\\
Lector & 1b071d27\allowbreak{}4a26233a\allowbreak{}2a95e7db\allowbreak{}4a91c4d3\allowbreak{}e75f2fa8\allowbreak{}1ab184e1\allowbreak{}dbc739d9\allowbreak{}38ccd981\\
Registro en 108 & 8c0e7e1d\allowbreak{}d0b049c9\allowbreak{}330530f4\allowbreak{}b68eb1be\allowbreak{}6a328e08\allowbreak{}bd9c0f29\allowbreak{}3df7e20e\allowbreak{}a4956ab2\\
\end{tabularx}\endgroup\par\vspace{6pt}
\Needspace{87pt}\phantomsection\label{trace-20}\noindent\hyperref[nav:vi-ruta-20]{\textbf{r1c5\_\allowbreak{}h+1v+1}}\quad Ordinal unificado / memoria: 20 / 17\par\nopagebreak
\begingroup\fontsize{8.3}{10.1}\selectfont\setlength{\tabcolsep}{4pt}\noindent\begin{tabularx}{\linewidth}{@{}p{.21\linewidth}X@{}}
Genealogía & 1b217fd9\allowbreak{}1d203aea\allowbreak{}e19f2a5a\allowbreak{}5a90fcad\allowbreak{}b55a5bb5\allowbreak{}48ca02e5\allowbreak{}1e4e20ec\allowbreak{}78b6444c\\
Lector & 78e79216\allowbreak{}b3802312\allowbreak{}a31bc267\allowbreak{}03b19549\allowbreak{}f845de0e\allowbreak{}5c1403ce\allowbreak{}0bfc7e42\allowbreak{}78fd13d9\\
Registro en 108 & \textemdash{}\\
\end{tabularx}\endgroup\par\vspace{6pt}
\Needspace{87pt}\phantomsection\label{trace-21}\noindent\hyperref[nav:vi-ruta-21]{\textbf{r1c6\_\allowbreak{}h-1v-1}}\quad Ordinal unificado / memoria: 21 / 24\par\nopagebreak
\begingroup\fontsize{8.3}{10.1}\selectfont\setlength{\tabcolsep}{4pt}\noindent\begin{tabularx}{\linewidth}{@{}p{.21\linewidth}X@{}}
Genealogía & d52d9bc5\allowbreak{}8081da1a\allowbreak{}7b2f9e5d\allowbreak{}23e9d5a6\allowbreak{}d67401f8\allowbreak{}24eb66ed\allowbreak{}733498c2\allowbreak{}fd0dab4c\\
Lector & 560f1934\allowbreak{}27fddf83\allowbreak{}fb498c6e\allowbreak{}c7a6b2e4\allowbreak{}a683ed18\allowbreak{}b5f9d3a6\allowbreak{}32bb9743\allowbreak{}5ce13545\\
Registro en 108 & b19c7187\allowbreak{}62e5299b\allowbreak{}6bd2d56c\allowbreak{}b241cd2c\allowbreak{}def7df5b\allowbreak{}1d224f69\allowbreak{}51364c8e\allowbreak{}46e57ade\\
\end{tabularx}\endgroup\par\vspace{6pt}
\Needspace{87pt}\phantomsection\label{trace-22}\noindent\hyperref[nav:vi-ruta-22]{\textbf{r1c6\_\allowbreak{}h-1v+1}}\quad Ordinal unificado / memoria: 22 / 23\par\nopagebreak
\begingroup\fontsize{8.3}{10.1}\selectfont\setlength{\tabcolsep}{4pt}\noindent\begin{tabularx}{\linewidth}{@{}p{.21\linewidth}X@{}}
Genealogía & 0af1af6e\allowbreak{}61e27f82\allowbreak{}e9c6658e\allowbreak{}03d303a9\allowbreak{}801f5cf6\allowbreak{}609402e5\allowbreak{}810dcbc9\allowbreak{}44b8c00b\\
Lector & a28a41d6\allowbreak{}785403d9\allowbreak{}b38e510f\allowbreak{}4c5203a2\allowbreak{}344cac70\allowbreak{}366f151a\allowbreak{}6d143c80\allowbreak{}93c59464\\
Registro en 108 & \textemdash{}\\
\end{tabularx}\endgroup\par\vspace{6pt}
\Needspace{87pt}\phantomsection\label{trace-23}\noindent\hyperref[nav:vi-ruta-23]{\textbf{r1c6\_\allowbreak{}h+1v-1}}\quad Ordinal unificado / memoria: 23 / 22\par\nopagebreak
\begingroup\fontsize{8.3}{10.1}\selectfont\setlength{\tabcolsep}{4pt}\noindent\begin{tabularx}{\linewidth}{@{}p{.21\linewidth}X@{}}
Genealogía & 36295015\allowbreak{}96c2b853\allowbreak{}8e13e1e9\allowbreak{}7d2b0c69\allowbreak{}f585dbc0\allowbreak{}f13f6f9f\allowbreak{}be5d7b01\allowbreak{}2d7082ea\\
Lector & f992e4a2\allowbreak{}a1b53de9\allowbreak{}a0d7533c\allowbreak{}2a673b02\allowbreak{}d28083e7\allowbreak{}f6aaaa45\allowbreak{}587ef610\allowbreak{}4b63db2f\\
Registro en 108 & \textemdash{}\\
\end{tabularx}\endgroup\par\vspace{6pt}
\Needspace{87pt}\phantomsection\label{trace-24}\noindent\hyperref[nav:vi-ruta-24]{\textbf{r1c6\_\allowbreak{}h+1v+1}}\quad Ordinal unificado / memoria: 24 / 21\par\nopagebreak
\begingroup\fontsize{8.3}{10.1}\selectfont\setlength{\tabcolsep}{4pt}\noindent\begin{tabularx}{\linewidth}{@{}p{.21\linewidth}X@{}}
Genealogía & c5b01f4a\allowbreak{}1d0acc94\allowbreak{}193d804a\allowbreak{}45897f48\allowbreak{}9ee02521\allowbreak{}a35214a9\allowbreak{}b8608671\allowbreak{}7d71cfdc\\
Lector & 09874714\allowbreak{}51aca8db\allowbreak{}dc0c7324\allowbreak{}674e2ffa\allowbreak{}42379094\allowbreak{}12e43186\allowbreak{}2c124cac\allowbreak{}78f45504\\
Registro en 108 & \textemdash{}\\
\end{tabularx}\endgroup\par\vspace{6pt}
\Needspace{87pt}\phantomsection\label{trace-25}\noindent\hyperref[nav:vi-ruta-25]{\textbf{r1c7\_\allowbreak{}h-1v-1}}\quad Ordinal unificado / memoria: 25 / 28\par\nopagebreak
\begingroup\fontsize{8.3}{10.1}\selectfont\setlength{\tabcolsep}{4pt}\noindent\begin{tabularx}{\linewidth}{@{}p{.21\linewidth}X@{}}
Genealogía & a911ce23\allowbreak{}dfdeeccc\allowbreak{}43cb29fe\allowbreak{}5898ecd2\allowbreak{}fabdd3f9\allowbreak{}94ca7281\allowbreak{}4b4e1b69\allowbreak{}8537ef02\\
Lector & 62e6f1b1\allowbreak{}0122e376\allowbreak{}1b7f3957\allowbreak{}4d07f715\allowbreak{}c69636c5\allowbreak{}8df0f10a\allowbreak{}f2f06d2c\allowbreak{}4289e08b\\
Registro en 108 & \textemdash{}\\
\end{tabularx}\endgroup\par\vspace{6pt}
\Needspace{87pt}\phantomsection\label{trace-26}\noindent\hyperref[nav:vi-ruta-26]{\textbf{r1c7\_\allowbreak{}h-1v+1}}\quad Ordinal unificado / memoria: 26 / 27\par\nopagebreak
\begingroup\fontsize{8.3}{10.1}\selectfont\setlength{\tabcolsep}{4pt}\noindent\begin{tabularx}{\linewidth}{@{}p{.21\linewidth}X@{}}
Genealogía & f580dc82\allowbreak{}47645d38\allowbreak{}12ad5b49\allowbreak{}31a5bf89\allowbreak{}ed066bd7\allowbreak{}a6cd3321\allowbreak{}f709a411\allowbreak{}f3996c80\\
Lector & 07b8d450\allowbreak{}b7f6e399\allowbreak{}8bfcc5fd\allowbreak{}f7dc71c5\allowbreak{}0982ba49\allowbreak{}1d4c5dd5\allowbreak{}bfe6236c\allowbreak{}6a11645f\\
Registro en 108 & 111ca40e\allowbreak{}5cb5dd4e\allowbreak{}1469c8db\allowbreak{}516ceb95\allowbreak{}b20b75ec\allowbreak{}5dfc93b8\allowbreak{}40a6cce9\allowbreak{}1c92b62c\\
\end{tabularx}\endgroup\par\vspace{6pt}
\Needspace{87pt}\phantomsection\label{trace-27}\noindent\hyperref[nav:vi-ruta-27]{\textbf{r1c7\_\allowbreak{}h+1v-1}}\quad Ordinal unificado / memoria: 27 / 26\par\nopagebreak
\begingroup\fontsize{8.3}{10.1}\selectfont\setlength{\tabcolsep}{4pt}\noindent\begin{tabularx}{\linewidth}{@{}p{.21\linewidth}X@{}}
Genealogía & 342d7eac\allowbreak{}14c30ee1\allowbreak{}9c5129f2\allowbreak{}70b31c29\allowbreak{}fa591198\allowbreak{}d94f5d42\allowbreak{}f5c23629\allowbreak{}fb102c38\\
Lector & e5e14191\allowbreak{}accc501f\allowbreak{}24a061fb\allowbreak{}7067050e\allowbreak{}f0a2efd5\allowbreak{}d86cf532\allowbreak{}156226da\allowbreak{}818a11e9\\
Registro en 108 & \textemdash{}\\
\end{tabularx}\endgroup\par\vspace{6pt}
\Needspace{87pt}\phantomsection\label{trace-28}\noindent\hyperref[nav:vi-ruta-28]{\textbf{r1c7\_\allowbreak{}h+1v+1}}\quad Ordinal unificado / memoria: 28 / 25\par\nopagebreak
\begingroup\fontsize{8.3}{10.1}\selectfont\setlength{\tabcolsep}{4pt}\noindent\begin{tabularx}{\linewidth}{@{}p{.21\linewidth}X@{}}
Genealogía & e65938ec\allowbreak{}108f08f6\allowbreak{}7b38c1bf\allowbreak{}dd3f600d\allowbreak{}b58f9a17\allowbreak{}d9c59a82\allowbreak{}b75ccff0\allowbreak{}6d19ce60\\
Lector & 9905fffa\allowbreak{}baeddac8\allowbreak{}86538ee3\allowbreak{}4f7320c7\allowbreak{}891f879f\allowbreak{}e8876c47\allowbreak{}0d8a790d\allowbreak{}e73b4d13\\
Registro en 108 & \textemdash{}\\
\end{tabularx}\endgroup\par\vspace{6pt}
\Needspace{87pt}\phantomsection\label{trace-29}\noindent\hyperref[nav:vi-ruta-29]{\textbf{r1c8\_\allowbreak{}h-1v-1}}\quad Ordinal unificado / memoria: 29 / 32\par\nopagebreak
\begingroup\fontsize{8.3}{10.1}\selectfont\setlength{\tabcolsep}{4pt}\noindent\begin{tabularx}{\linewidth}{@{}p{.21\linewidth}X@{}}
Genealogía & 72789b3e\allowbreak{}fdc97b1d\allowbreak{}767a5aa7\allowbreak{}3046d084\allowbreak{}583fa5ab\allowbreak{}a373f9a0\allowbreak{}7a6fc16b\allowbreak{}1c1f04df\\
Lector & 53e27ee1\allowbreak{}34cda7bb\allowbreak{}20474d1c\allowbreak{}7af7cc89\allowbreak{}21838c09\allowbreak{}4a2dd9bc\allowbreak{}724b2d30\allowbreak{}cff5a10a\\
Registro en 108 & \textemdash{}\\
\end{tabularx}\endgroup\par\vspace{6pt}
\Needspace{87pt}\phantomsection\label{trace-30}\noindent\hyperref[nav:vi-ruta-30]{\textbf{r1c8\_\allowbreak{}h-1v+1}}\quad Ordinal unificado / memoria: 30 / 31\par\nopagebreak
\begingroup\fontsize{8.3}{10.1}\selectfont\setlength{\tabcolsep}{4pt}\noindent\begin{tabularx}{\linewidth}{@{}p{.21\linewidth}X@{}}
Genealogía & f612d6ef\allowbreak{}200a1ff5\allowbreak{}4e74c4ed\allowbreak{}7c7a5108\allowbreak{}93ca3404\allowbreak{}8d213814\allowbreak{}66db682c\allowbreak{}ca3cc661\\
Lector & 48191dd8\allowbreak{}270ba6e6\allowbreak{}8ab41300\allowbreak{}a7daa22a\allowbreak{}6daa2102\allowbreak{}e4aff594\allowbreak{}b72c5fbc\allowbreak{}01440d85\\
Registro en 108 & \textemdash{}\\
\end{tabularx}\endgroup\par\vspace{6pt}
\Needspace{87pt}\phantomsection\label{trace-31}\noindent\hyperref[nav:vi-ruta-31]{\textbf{r1c8\_\allowbreak{}h+1v-1}}\quad Ordinal unificado / memoria: 31 / 30\par\nopagebreak
\begingroup\fontsize{8.3}{10.1}\selectfont\setlength{\tabcolsep}{4pt}\noindent\begin{tabularx}{\linewidth}{@{}p{.21\linewidth}X@{}}
Genealogía & 48fc9a27\allowbreak{}e8d42d25\allowbreak{}0d4ab52c\allowbreak{}b98cf105\allowbreak{}f3250728\allowbreak{}1dc8824f\allowbreak{}92a8ce07\allowbreak{}cfd13aeb\\
Lector & ab4343ca\allowbreak{}0f3279dd\allowbreak{}31d470d4\allowbreak{}b27e6360\allowbreak{}42039822\allowbreak{}26a40d21\allowbreak{}96f33fdb\allowbreak{}10b85cc1\\
Registro en 108 & 012b2f08\allowbreak{}ceb67413\allowbreak{}89c89bac\allowbreak{}c896d7b2\allowbreak{}769a75fa\allowbreak{}04b47531\allowbreak{}35b68e45\allowbreak{}09b3ad20\\
\end{tabularx}\endgroup\par\vspace{6pt}
\Needspace{87pt}\phantomsection\label{trace-32}\noindent\hyperref[nav:vi-ruta-32]{\textbf{r1c8\_\allowbreak{}h+1v+1}}\quad Ordinal unificado / memoria: 32 / 29\par\nopagebreak
\begingroup\fontsize{8.3}{10.1}\selectfont\setlength{\tabcolsep}{4pt}\noindent\begin{tabularx}{\linewidth}{@{}p{.21\linewidth}X@{}}
Genealogía & 3aa3fba3\allowbreak{}6ac1b063\allowbreak{}64a3368d\allowbreak{}fd7a9486\allowbreak{}ddc4de29\allowbreak{}c64dc93a\allowbreak{}59368732\allowbreak{}61d07e18\\
Lector & f0c45d18\allowbreak{}09914b63\allowbreak{}308aa054\allowbreak{}a401a393\allowbreak{}3c1b9671\allowbreak{}f0dce46e\allowbreak{}834b5c6a\allowbreak{}8afb2263\\
Registro en 108 & \textemdash{}\\
\end{tabularx}\endgroup\par\vspace{6pt}
\Needspace{87pt}\phantomsection\label{trace-33}\noindent\hyperref[nav:vi-ruta-33]{\textbf{r1c9\_\allowbreak{}h-1v-1}}\quad Ordinal unificado / memoria: 33 / 36\par\nopagebreak
\begingroup\fontsize{8.3}{10.1}\selectfont\setlength{\tabcolsep}{4pt}\noindent\begin{tabularx}{\linewidth}{@{}p{.21\linewidth}X@{}}
Genealogía & ed985a2e\allowbreak{}d5407991\allowbreak{}50c7e349\allowbreak{}e2cc896d\allowbreak{}9aac716d\allowbreak{}704d2653\allowbreak{}88865e60\allowbreak{}38a3165e\\
Lector & 66eb3e8e\allowbreak{}ed039dc7\allowbreak{}4f989a1a\allowbreak{}f5a5acbe\allowbreak{}386ed984\allowbreak{}12b13ef3\allowbreak{}dcedb1c6\allowbreak{}487e9571\\
Registro en 108 & \textemdash{}\\
\end{tabularx}\endgroup\par\vspace{6pt}
\Needspace{87pt}\phantomsection\label{trace-34}\noindent\hyperref[nav:vi-ruta-34]{\textbf{r1c9\_\allowbreak{}h-1v+1}}\quad Ordinal unificado / memoria: 34 / 35\par\nopagebreak
\begingroup\fontsize{8.3}{10.1}\selectfont\setlength{\tabcolsep}{4pt}\noindent\begin{tabularx}{\linewidth}{@{}p{.21\linewidth}X@{}}
Genealogía & a02045b2\allowbreak{}e206e07f\allowbreak{}7236c381\allowbreak{}210f0d3e\allowbreak{}9dc3fcaf\allowbreak{}430a0b34\allowbreak{}9f8488ee\allowbreak{}c7605f89\\
Lector & fdda7f9e\allowbreak{}60574459\allowbreak{}88cd077d\allowbreak{}c481c836\allowbreak{}df8a633e\allowbreak{}66f5cfad\allowbreak{}3ab15b29\allowbreak{}96ff488f\\
Registro en 108 & \textemdash{}\\
\end{tabularx}\endgroup\par\vspace{6pt}
\Needspace{87pt}\phantomsection\label{trace-35}\noindent\hyperref[nav:vi-ruta-35]{\textbf{r1c9\_\allowbreak{}h+1v-1}}\quad Ordinal unificado / memoria: 35 / 34\par\nopagebreak
\begingroup\fontsize{8.3}{10.1}\selectfont\setlength{\tabcolsep}{4pt}\noindent\begin{tabularx}{\linewidth}{@{}p{.21\linewidth}X@{}}
Genealogía & a1fe11c1\allowbreak{}f44fee33\allowbreak{}9c8dfe31\allowbreak{}8ad4ce23\allowbreak{}565bd11d\allowbreak{}206e5e5b\allowbreak{}0a3f8136\allowbreak{}f94cbcff\\
Lector & 409959c2\allowbreak{}ff367348\allowbreak{}163db758\allowbreak{}c5659d70\allowbreak{}ab997d5c\allowbreak{}b215f7ff\allowbreak{}22fc64d0\allowbreak{}db140a39\\
Registro en 108 & \textemdash{}\\
\end{tabularx}\endgroup\par\vspace{6pt}
\Needspace{87pt}\phantomsection\label{trace-36}\noindent\hyperref[nav:vi-ruta-36]{\textbf{r1c9\_\allowbreak{}h+1v+1}}\quad Ordinal unificado / memoria: 36 / 33\par\nopagebreak
\begingroup\fontsize{8.3}{10.1}\selectfont\setlength{\tabcolsep}{4pt}\noindent\begin{tabularx}{\linewidth}{@{}p{.21\linewidth}X@{}}
Genealogía & 5decd17f\allowbreak{}a25c6e87\allowbreak{}374e54b0\allowbreak{}f6b6fa4c\allowbreak{}85317267\allowbreak{}086884d1\allowbreak{}e9b3cc6c\allowbreak{}2626241b\\
Lector & e8213bed\allowbreak{}e9eb8c18\allowbreak{}9e246a8a\allowbreak{}ac685589\allowbreak{}8eb38d7c\allowbreak{}085789b8\allowbreak{}527ff022\allowbreak{}f9311b9f\\
Registro en 108 & 7d1f109e\allowbreak{}1e33bf84\allowbreak{}ddf9384c\allowbreak{}f13c342f\allowbreak{}63017fe9\allowbreak{}faecd3c7\allowbreak{}3d8063ec\allowbreak{}a5fa5462\\
\end{tabularx}\endgroup\par\vspace{6pt}
\Needspace{87pt}\phantomsection\label{trace-37}\noindent\hyperref[nav:vi-ruta-37]{\textbf{r2c1\_\allowbreak{}h-1v-1}}\quad Ordinal unificado / memoria: 37 / 40\par\nopagebreak
\begingroup\fontsize{8.3}{10.1}\selectfont\setlength{\tabcolsep}{4pt}\noindent\begin{tabularx}{\linewidth}{@{}p{.21\linewidth}X@{}}
Genealogía & 42121291\allowbreak{}b4db1b25\allowbreak{}1482aaa6\allowbreak{}00938754\allowbreak{}c4e97069\allowbreak{}468ccdcb\allowbreak{}d0d654a2\allowbreak{}8f9b0873\\
Lector & 7e147ace\allowbreak{}8f35cc0e\allowbreak{}74360e29\allowbreak{}6f6f57ce\allowbreak{}18196c7f\allowbreak{}81af3828\allowbreak{}e0e513a7\allowbreak{}8d821461\\
Registro en 108 & \textemdash{}\\
\end{tabularx}\endgroup\par\vspace{6pt}
\Needspace{87pt}\phantomsection\label{trace-38}\noindent\hyperref[nav:vi-ruta-38]{\textbf{r2c1\_\allowbreak{}h-1v+1}}\quad Ordinal unificado / memoria: 38 / 39\par\nopagebreak
\begingroup\fontsize{8.3}{10.1}\selectfont\setlength{\tabcolsep}{4pt}\noindent\begin{tabularx}{\linewidth}{@{}p{.21\linewidth}X@{}}
Genealogía & e3a4eca8\allowbreak{}8fb523b4\allowbreak{}cd27f0dc\allowbreak{}bc45a95e\allowbreak{}626a0d4f\allowbreak{}ae70ca90\allowbreak{}85350d36\allowbreak{}de129e86\\
Lector & f06f9435\allowbreak{}77be9d47\allowbreak{}6fa92b2c\allowbreak{}b82203fd\allowbreak{}77be6d41\allowbreak{}f9d8efd5\allowbreak{}8303786f\allowbreak{}07a013ef\\
Registro en 108 & \textemdash{}\\
\end{tabularx}\endgroup\par\vspace{6pt}
\Needspace{87pt}\phantomsection\label{trace-39}\noindent\hyperref[nav:vi-ruta-39]{\textbf{r2c1\_\allowbreak{}h+1v-1}}\quad Ordinal unificado / memoria: 39 / 38\par\nopagebreak
\begingroup\fontsize{8.3}{10.1}\selectfont\setlength{\tabcolsep}{4pt}\noindent\begin{tabularx}{\linewidth}{@{}p{.21\linewidth}X@{}}
Genealogía & 6cd73505\allowbreak{}4b3d2851\allowbreak{}58e10171\allowbreak{}0ab447ea\allowbreak{}6f8236d0\allowbreak{}455b76b5\allowbreak{}aa102be6\allowbreak{}8ac85002\\
Lector & efe6e221\allowbreak{}b6809bf8\allowbreak{}4456087d\allowbreak{}b0f756cc\allowbreak{}315eba62\allowbreak{}05271dfe\allowbreak{}72062438\allowbreak{}b4952f78\\
Registro en 108 & \textemdash{}\\
\end{tabularx}\endgroup\par\vspace{6pt}
\Needspace{87pt}\phantomsection\label{trace-40}\noindent\hyperref[nav:vi-ruta-40]{\textbf{r2c1\_\allowbreak{}h+1v+1}}\quad Ordinal unificado / memoria: 40 / 37\par\nopagebreak
\begingroup\fontsize{8.3}{10.1}\selectfont\setlength{\tabcolsep}{4pt}\noindent\begin{tabularx}{\linewidth}{@{}p{.21\linewidth}X@{}}
Genealogía & 7654e0aa\allowbreak{}6d4ec6e5\allowbreak{}959f7706\allowbreak{}79686d0f\allowbreak{}e3c0dafd\allowbreak{}21ff7cb9\allowbreak{}4ed4ea85\allowbreak{}05f460a2\\
Lector & 39a510c4\allowbreak{}fa875ca5\allowbreak{}fc2c73e4\allowbreak{}82279f49\allowbreak{}5f57505a\allowbreak{}c7d2e998\allowbreak{}6bf08753\allowbreak{}fb317a28\\
Registro en 108 & 355e9237\allowbreak{}5547ac1e\allowbreak{}37c63290\allowbreak{}e00e0b13\allowbreak{}115df481\allowbreak{}2b567558\allowbreak{}c2432472\allowbreak{}731b13de\\
\end{tabularx}\endgroup\par\vspace{6pt}
\Needspace{87pt}\phantomsection\label{trace-41}\noindent\hyperref[nav:vi-ruta-41]{\textbf{r2c2\_\allowbreak{}h-1v-1}}\quad Ordinal unificado / memoria: 41 / 44\par\nopagebreak
\begingroup\fontsize{8.3}{10.1}\selectfont\setlength{\tabcolsep}{4pt}\noindent\begin{tabularx}{\linewidth}{@{}p{.21\linewidth}X@{}}
Genealogía & 9eb4e897\allowbreak{}dd2417fd\allowbreak{}d7ef230e\allowbreak{}161ee02a\allowbreak{}73430f2c\allowbreak{}d4c165e0\allowbreak{}b5e03f4c\allowbreak{}d9432a7c\\
Lector & 7390d3cc\allowbreak{}e6527e32\allowbreak{}50ecccb9\allowbreak{}cd6d58a5\allowbreak{}d3548fa0\allowbreak{}525ab002\allowbreak{}4b7f93b2\allowbreak{}4d292216\\
Registro en 108 & \textemdash{}\\
\end{tabularx}\endgroup\par\vspace{6pt}
\Needspace{87pt}\phantomsection\label{trace-42}\noindent\hyperref[nav:vi-ruta-42]{\textbf{r2c2\_\allowbreak{}h-1v+1}}\quad Ordinal unificado / memoria: 42 / 43\par\nopagebreak
\begingroup\fontsize{8.3}{10.1}\selectfont\setlength{\tabcolsep}{4pt}\noindent\begin{tabularx}{\linewidth}{@{}p{.21\linewidth}X@{}}
Genealogía & 1fe3a911\allowbreak{}aaa125a9\allowbreak{}905aeadc\allowbreak{}38d1dbb6\allowbreak{}3d175976\allowbreak{}4b7429b0\allowbreak{}39228794\allowbreak{}908d53b8\\
Lector & 779e5b69\allowbreak{}8afcdc09\allowbreak{}10f6d945\allowbreak{}1f0ee72b\allowbreak{}8327bc1e\allowbreak{}94ac4096\allowbreak{}d006fec2\allowbreak{}406ecf3c\\
Registro en 108 & e34326b4\allowbreak{}9595fe5f\allowbreak{}f46f40dd\allowbreak{}d0b7b0a5\allowbreak{}11a0f579\allowbreak{}53ba2953\allowbreak{}b50c9ece\allowbreak{}d0884e64\\
\end{tabularx}\endgroup\par\vspace{6pt}
\Needspace{87pt}\phantomsection\label{trace-43}\noindent\hyperref[nav:vi-ruta-43]{\textbf{r2c2\_\allowbreak{}h+1v-1}}\quad Ordinal unificado / memoria: 43 / 42\par\nopagebreak
\begingroup\fontsize{8.3}{10.1}\selectfont\setlength{\tabcolsep}{4pt}\noindent\begin{tabularx}{\linewidth}{@{}p{.21\linewidth}X@{}}
Genealogía & d8cfd857\allowbreak{}e63aa9d9\allowbreak{}c90169d1\allowbreak{}594eaa1a\allowbreak{}d2e6b70c\allowbreak{}6d138a85\allowbreak{}ea226813\allowbreak{}977a322a\\
Lector & 46e898cd\allowbreak{}af669381\allowbreak{}51ec8b48\allowbreak{}a75efafe\allowbreak{}4eba7615\allowbreak{}becf22c8\allowbreak{}9a22d7a0\allowbreak{}9b274062\\
Registro en 108 & \textemdash{}\\
\end{tabularx}\endgroup\par\vspace{6pt}
\Needspace{87pt}\phantomsection\label{trace-44}\noindent\hyperref[nav:vi-ruta-44]{\textbf{r2c2\_\allowbreak{}h+1v+1}}\quad Ordinal unificado / memoria: 44 / 41\par\nopagebreak
\begingroup\fontsize{8.3}{10.1}\selectfont\setlength{\tabcolsep}{4pt}\noindent\begin{tabularx}{\linewidth}{@{}p{.21\linewidth}X@{}}
Genealogía & bf3e75cc\allowbreak{}7e977bbf\allowbreak{}be883ce1\allowbreak{}215b4ab8\allowbreak{}50affa32\allowbreak{}42a70408\allowbreak{}9c1350ee\allowbreak{}55fde93f\\
Lector & bec293e8\allowbreak{}a9352a0c\allowbreak{}9e5b56fb\allowbreak{}6c7d30c8\allowbreak{}76b80cec\allowbreak{}47727a2f\allowbreak{}34d40d9a\allowbreak{}3e0cd3ca\\
Registro en 108 & \textemdash{}\\
\end{tabularx}\endgroup\par\vspace{6pt}
\Needspace{87pt}\phantomsection\label{trace-45}\noindent\hyperref[nav:vi-ruta-45]{\textbf{r2c3\_\allowbreak{}h-1v-1}}\quad Ordinal unificado / memoria: 45 / 48\par\nopagebreak
\begingroup\fontsize{8.3}{10.1}\selectfont\setlength{\tabcolsep}{4pt}\noindent\begin{tabularx}{\linewidth}{@{}p{.21\linewidth}X@{}}
Genealogía & 21648712\allowbreak{}dd5dcff5\allowbreak{}4d6151e8\allowbreak{}5c6297f4\allowbreak{}97049ea1\allowbreak{}8f546275\allowbreak{}a253f506\allowbreak{}f95483d2\\
Lector & 6140ac87\allowbreak{}12fe0604\allowbreak{}6683f5fd\allowbreak{}5d240eb6\allowbreak{}2685661a\allowbreak{}d9a8244e\allowbreak{}3cbe742c\allowbreak{}5051fd66\\
Registro en 108 & 6c3c21ff\allowbreak{}d2893c36\allowbreak{}62b538c0\allowbreak{}17a94d62\allowbreak{}3d66cc85\allowbreak{}b623c55c\allowbreak{}f8d08a67\allowbreak{}b4f70c93\\
\end{tabularx}\endgroup\par\vspace{6pt}
\Needspace{87pt}\phantomsection\label{trace-46}\noindent\hyperref[nav:vi-ruta-46]{\textbf{r2c3\_\allowbreak{}h-1v+1}}\quad Ordinal unificado / memoria: 46 / 47\par\nopagebreak
\begingroup\fontsize{8.3}{10.1}\selectfont\setlength{\tabcolsep}{4pt}\noindent\begin{tabularx}{\linewidth}{@{}p{.21\linewidth}X@{}}
Genealogía & 2729489d\allowbreak{}2f4f122b\allowbreak{}773697f8\allowbreak{}b3eef3fa\allowbreak{}11b69d50\allowbreak{}f3606c4a\allowbreak{}5418b665\allowbreak{}87cb1a4b\\
Lector & f6107968\allowbreak{}20399955\allowbreak{}f3b317dc\allowbreak{}077df3c8\allowbreak{}56ff4c75\allowbreak{}f80da841\allowbreak{}ce0a7221\allowbreak{}4efd06d0\\
Registro en 108 & \textemdash{}\\
\end{tabularx}\endgroup\par\vspace{6pt}
\Needspace{87pt}\phantomsection\label{trace-47}\noindent\hyperref[nav:vi-ruta-47]{\textbf{r2c3\_\allowbreak{}h+1v-1}}\quad Ordinal unificado / memoria: 47 / 46\par\nopagebreak
\begingroup\fontsize{8.3}{10.1}\selectfont\setlength{\tabcolsep}{4pt}\noindent\begin{tabularx}{\linewidth}{@{}p{.21\linewidth}X@{}}
Genealogía & c5f5a7cc\allowbreak{}8a3f497d\allowbreak{}2781fcfc\allowbreak{}abb97bed\allowbreak{}e9beb413\allowbreak{}9194e9cd\allowbreak{}554320a8\allowbreak{}d4612024\\
Lector & e8c645b2\allowbreak{}38014406\allowbreak{}bd74ee9d\allowbreak{}49df1370\allowbreak{}3368a36d\allowbreak{}854493d9\allowbreak{}9ef35505\allowbreak{}8b2a9977\\
Registro en 108 & \textemdash{}\\
\end{tabularx}\endgroup\par\vspace{6pt}
\Needspace{87pt}\phantomsection\label{trace-48}\noindent\hyperref[nav:vi-ruta-48]{\textbf{r2c3\_\allowbreak{}h+1v+1}}\quad Ordinal unificado / memoria: 48 / 45\par\nopagebreak
\begingroup\fontsize{8.3}{10.1}\selectfont\setlength{\tabcolsep}{4pt}\noindent\begin{tabularx}{\linewidth}{@{}p{.21\linewidth}X@{}}
Genealogía & 698f0a4e\allowbreak{}d1e27622\allowbreak{}225b33b8\allowbreak{}cf29ee36\allowbreak{}96219e67\allowbreak{}44484603\allowbreak{}2e3f9212\allowbreak{}d53466a3\\
Lector & 816cbe81\allowbreak{}5adfede4\allowbreak{}68d9709e\allowbreak{}c0485910\allowbreak{}767d8209\allowbreak{}7a66de13\allowbreak{}f0984f60\allowbreak{}9f5bd113\\
Registro en 108 & \textemdash{}\\
\end{tabularx}\endgroup\par\vspace{6pt}
\Needspace{87pt}\phantomsection\label{trace-49}\noindent\hyperref[nav:vi-ruta-49]{\textbf{r2c4\_\allowbreak{}h-1v-1}}\quad Ordinal unificado / memoria: 49 / 52\par\nopagebreak
\begingroup\fontsize{8.3}{10.1}\selectfont\setlength{\tabcolsep}{4pt}\noindent\begin{tabularx}{\linewidth}{@{}p{.21\linewidth}X@{}}
Genealogía & 2597b145\allowbreak{}8c7c6d27\allowbreak{}65567053\allowbreak{}3a40d232\allowbreak{}5d645709\allowbreak{}b9f9e060\allowbreak{}260b7cc2\allowbreak{}573da835\\
Lector & 1a459f0d\allowbreak{}ff0ec31f\allowbreak{}6beadf54\allowbreak{}00457dcc\allowbreak{}3102209e\allowbreak{}4b1f863b\allowbreak{}0749dcff\allowbreak{}b5f87ade\\
Registro en 108 & \textemdash{}\\
\end{tabularx}\endgroup\par\vspace{6pt}
\Needspace{87pt}\phantomsection\label{trace-50}\noindent\hyperref[nav:vi-ruta-50]{\textbf{r2c4\_\allowbreak{}h-1v+1}}\quad Ordinal unificado / memoria: 50 / 51\par\nopagebreak
\begingroup\fontsize{8.3}{10.1}\selectfont\setlength{\tabcolsep}{4pt}\noindent\begin{tabularx}{\linewidth}{@{}p{.21\linewidth}X@{}}
Genealogía & 679c8850\allowbreak{}da30d692\allowbreak{}454f0003\allowbreak{}9b4e73d5\allowbreak{}8c5eff33\allowbreak{}5d815532\allowbreak{}21a5a0b3\allowbreak{}159a5624\\
Lector & ace3d803\allowbreak{}4260254a\allowbreak{}3bb93657\allowbreak{}b2e2123e\allowbreak{}63e67d16\allowbreak{}904e7f0b\allowbreak{}632fdcb6\allowbreak{}71332da7\\
Registro en 108 & \textemdash{}\\
\end{tabularx}\endgroup\par\vspace{6pt}
\Needspace{87pt}\phantomsection\label{trace-51}\noindent\hyperref[nav:vi-ruta-51]{\textbf{r2c4\_\allowbreak{}h+1v-1}}\quad Ordinal unificado / memoria: 51 / 50\par\nopagebreak
\begingroup\fontsize{8.3}{10.1}\selectfont\setlength{\tabcolsep}{4pt}\noindent\begin{tabularx}{\linewidth}{@{}p{.21\linewidth}X@{}}
Genealogía & 6055f66a\allowbreak{}8741e114\allowbreak{}0af22cd2\allowbreak{}d82c9b62\allowbreak{}c55c4e32\allowbreak{}ae9162ca\allowbreak{}d35bad23\allowbreak{}d3cbdcc2\\
Lector & 9078ab56\allowbreak{}3bb6e210\allowbreak{}453e86e1\allowbreak{}3e18dd71\allowbreak{}d892a6b9\allowbreak{}14ff5fc6\allowbreak{}5482a490\allowbreak{}04bcd077\\
Registro en 108 & 00c21eea\allowbreak{}4789a8c0\allowbreak{}6b37c224\allowbreak{}41f75093\allowbreak{}ea0f0184\allowbreak{}ff1bc295\allowbreak{}54037724\allowbreak{}90155a81\\
\end{tabularx}\endgroup\par\vspace{6pt}
\Needspace{87pt}\phantomsection\label{trace-52}\noindent\hyperref[nav:vi-ruta-52]{\textbf{r2c4\_\allowbreak{}h+1v+1}}\quad Ordinal unificado / memoria: 52 / 49\par\nopagebreak
\begingroup\fontsize{8.3}{10.1}\selectfont\setlength{\tabcolsep}{4pt}\noindent\begin{tabularx}{\linewidth}{@{}p{.21\linewidth}X@{}}
Genealogía & 9b29c858\allowbreak{}0391aae7\allowbreak{}b3dbcfef\allowbreak{}d966550a\allowbreak{}bd79054c\allowbreak{}5e2e9520\allowbreak{}b547d06b\allowbreak{}227fe840\\
Lector & 9a3f104e\allowbreak{}10a55934\allowbreak{}eddfadbb\allowbreak{}190a4e8d\allowbreak{}0d79fb55\allowbreak{}eba7bc44\allowbreak{}03c29aac\allowbreak{}08923eec\\
Registro en 108 & \textemdash{}\\
\end{tabularx}\endgroup\par\vspace{6pt}
\Needspace{87pt}\phantomsection\label{trace-53}\noindent\hyperref[nav:vi-ruta-53]{\textbf{r2c5\_\allowbreak{}h-1v-1}}\quad Ordinal unificado / memoria: 53 / 56\par\nopagebreak
\begingroup\fontsize{8.3}{10.1}\selectfont\setlength{\tabcolsep}{4pt}\noindent\begin{tabularx}{\linewidth}{@{}p{.21\linewidth}X@{}}
Genealogía & d1fd4904\allowbreak{}0e524fea\allowbreak{}8052b4b0\allowbreak{}aaf01811\allowbreak{}92155662\allowbreak{}059c8363\allowbreak{}70670d5c\allowbreak{}2814748b\\
Lector & 83380600\allowbreak{}e5404d7e\allowbreak{}88802387\allowbreak{}96855cf6\allowbreak{}4cf95c1f\allowbreak{}b93d17d2\allowbreak{}320ce187\allowbreak{}bf98369d\\
Registro en 108 & \textemdash{}\\
\end{tabularx}\endgroup\par\vspace{6pt}
\Needspace{87pt}\phantomsection\label{trace-54}\noindent\hyperref[nav:vi-ruta-54]{\textbf{r2c5\_\allowbreak{}h-1v+1}}\quad Ordinal unificado / memoria: 54 / 55\par\nopagebreak
\begingroup\fontsize{8.3}{10.1}\selectfont\setlength{\tabcolsep}{4pt}\noindent\begin{tabularx}{\linewidth}{@{}p{.21\linewidth}X@{}}
Genealogía & 2708ff66\allowbreak{}18e713f0\allowbreak{}d4900bb9\allowbreak{}c1309696\allowbreak{}1b5c2b85\allowbreak{}4459b699\allowbreak{}c5a469f1\allowbreak{}33917fe1\\
Lector & 02b48980\allowbreak{}193d7b36\allowbreak{}66a37a97\allowbreak{}c6cc611e\allowbreak{}2af419c8\allowbreak{}f78bacae\allowbreak{}9bb362a7\allowbreak{}9e43767e\\
Registro en 108 & \textemdash{}\\
\end{tabularx}\endgroup\par\vspace{6pt}
\Needspace{87pt}\phantomsection\label{trace-55}\noindent\hyperref[nav:vi-ruta-55]{\textbf{r2c5\_\allowbreak{}h+1v-1}}\quad Ordinal unificado / memoria: 55 / 54\par\nopagebreak
\begingroup\fontsize{8.3}{10.1}\selectfont\setlength{\tabcolsep}{4pt}\noindent\begin{tabularx}{\linewidth}{@{}p{.21\linewidth}X@{}}
Genealogía & 2133fe30\allowbreak{}ab9b2c2f\allowbreak{}da2fe370\allowbreak{}cf35fabe\allowbreak{}c9776902\allowbreak{}ae7f1bf5\allowbreak{}dd7a68db\allowbreak{}23c22035\\
Lector & 7e4da1b5\allowbreak{}6102a953\allowbreak{}09b04e40\allowbreak{}b636fc43\allowbreak{}7722155b\allowbreak{}1ed9b43b\allowbreak{}c63ed9a4\allowbreak{}6c0f9c84\\
Registro en 108 & \textemdash{}\\
\end{tabularx}\endgroup\par\vspace{6pt}
\Needspace{87pt}\phantomsection\label{trace-56}\noindent\hyperref[nav:vi-ruta-56]{\textbf{r2c5\_\allowbreak{}h+1v+1}}\quad Ordinal unificado / memoria: 56 / 53\par\nopagebreak
\begingroup\fontsize{8.3}{10.1}\selectfont\setlength{\tabcolsep}{4pt}\noindent\begin{tabularx}{\linewidth}{@{}p{.21\linewidth}X@{}}
Genealogía & 7f85598f\allowbreak{}2d0761cd\allowbreak{}b1daa591\allowbreak{}c78b22b9\allowbreak{}37d0b0cf\allowbreak{}9020ac23\allowbreak{}e6f9e943\allowbreak{}fd562608\\
Lector & fb7d8d4d\allowbreak{}5037cd9c\allowbreak{}e900cca6\allowbreak{}f805cb47\allowbreak{}85c5e6dc\allowbreak{}4618e4b3\allowbreak{}fa59c593\allowbreak{}bf78eebf\\
Registro en 108 & 1523ad72\allowbreak{}32fb11fc\allowbreak{}5646852a\allowbreak{}3d2a3f11\allowbreak{}22cadd04\allowbreak{}0d970964\allowbreak{}3b2c332d\allowbreak{}82848b2a\\
\end{tabularx}\endgroup\par\vspace{6pt}
\Needspace{87pt}\phantomsection\label{trace-57}\noindent\hyperref[nav:vi-ruta-57]{\textbf{r2c6\_\allowbreak{}h-1v-1}}\quad Ordinal unificado / memoria: 57 / 60\par\nopagebreak
\begingroup\fontsize{8.3}{10.1}\selectfont\setlength{\tabcolsep}{4pt}\noindent\begin{tabularx}{\linewidth}{@{}p{.21\linewidth}X@{}}
Genealogía & 455c2b00\allowbreak{}84887303\allowbreak{}9e050631\allowbreak{}5671ab73\allowbreak{}33369975\allowbreak{}0588879e\allowbreak{}8ae6a976\allowbreak{}2ba3af54\\
Lector & 36eab0f2\allowbreak{}17dc9a30\allowbreak{}cf9ca2da\allowbreak{}7ac99f6c\allowbreak{}2e675173\allowbreak{}47899cf3\allowbreak{}bee9182c\allowbreak{}711ace32\\
Registro en 108 & \textemdash{}\\
\end{tabularx}\endgroup\par\vspace{6pt}
\Needspace{87pt}\phantomsection\label{trace-58}\noindent\hyperref[nav:vi-ruta-58]{\textbf{r2c6\_\allowbreak{}h-1v+1}}\quad Ordinal unificado / memoria: 58 / 59\par\nopagebreak
\begingroup\fontsize{8.3}{10.1}\selectfont\setlength{\tabcolsep}{4pt}\noindent\begin{tabularx}{\linewidth}{@{}p{.21\linewidth}X@{}}
Genealogía & 6f127fd9\allowbreak{}f25200c9\allowbreak{}da395c64\allowbreak{}82cb10ba\allowbreak{}06991f9a\allowbreak{}2b311706\allowbreak{}e05753d8\allowbreak{}8f848809\\
Lector & 97916a87\allowbreak{}21973f59\allowbreak{}4ba41fb3\allowbreak{}ce987a18\allowbreak{}682cae23\allowbreak{}e96d3ad8\allowbreak{}1ad92f70\allowbreak{}aa3ffda7\\
Registro en 108 & bfc92621\allowbreak{}bc7d75ba\allowbreak{}4ede7c52\allowbreak{}fd53f1a3\allowbreak{}2c441577\allowbreak{}160adc4a\allowbreak{}69490655\allowbreak{}db8b7ea3\\
\end{tabularx}\endgroup\par\vspace{6pt}
\Needspace{87pt}\phantomsection\label{trace-59}\noindent\hyperref[nav:vi-ruta-59]{\textbf{r2c6\_\allowbreak{}h+1v-1}}\quad Ordinal unificado / memoria: 59 / 58\par\nopagebreak
\begingroup\fontsize{8.3}{10.1}\selectfont\setlength{\tabcolsep}{4pt}\noindent\begin{tabularx}{\linewidth}{@{}p{.21\linewidth}X@{}}
Genealogía & a820aa46\allowbreak{}633c948e\allowbreak{}12bb41b7\allowbreak{}b41a621a\allowbreak{}d730d970\allowbreak{}e797f93a\allowbreak{}d58c7087\allowbreak{}7be24f53\\
Lector & f5bcac8c\allowbreak{}ac9dad8e\allowbreak{}e01f13e7\allowbreak{}f30b39e2\allowbreak{}243232b4\allowbreak{}d9a8da8a\allowbreak{}03297e16\allowbreak{}38b16cb4\\
Registro en 108 & \textemdash{}\\
\end{tabularx}\endgroup\par\vspace{6pt}
\Needspace{87pt}\phantomsection\label{trace-60}\noindent\hyperref[nav:vi-ruta-60]{\textbf{r2c6\_\allowbreak{}h+1v+1}}\quad Ordinal unificado / memoria: 60 / 57\par\nopagebreak
\begingroup\fontsize{8.3}{10.1}\selectfont\setlength{\tabcolsep}{4pt}\noindent\begin{tabularx}{\linewidth}{@{}p{.21\linewidth}X@{}}
Genealogía & ea920239\allowbreak{}9ce61bc1\allowbreak{}abe3e3bb\allowbreak{}f0faed08\allowbreak{}2c6e4d11\allowbreak{}f9e0ce94\allowbreak{}fb1a343e\allowbreak{}0f60b2c5\\
Lector & 9fdef0c8\allowbreak{}e1e5027d\allowbreak{}0dcdddc3\allowbreak{}848d7807\allowbreak{}3d21e8d1\allowbreak{}5c0ce90a\allowbreak{}2b53a7de\allowbreak{}4ed3563b\\
Registro en 108 & \textemdash{}\\
\end{tabularx}\endgroup\par\vspace{6pt}
\Needspace{87pt}\phantomsection\label{trace-61}\noindent\hyperref[nav:vi-ruta-61]{\textbf{r2c7\_\allowbreak{}h-1v-1}}\quad Ordinal unificado / memoria: 61 / 64\par\nopagebreak
\begingroup\fontsize{8.3}{10.1}\selectfont\setlength{\tabcolsep}{4pt}\noindent\begin{tabularx}{\linewidth}{@{}p{.21\linewidth}X@{}}
Genealogía & 93f4d12e\allowbreak{}988e9b7d\allowbreak{}84464bea\allowbreak{}475814d6\allowbreak{}748da201\allowbreak{}59cd95ca\allowbreak{}18726395\allowbreak{}02d4e22e\\
Lector & 4f6bcfea\allowbreak{}35a49b57\allowbreak{}dd08fe0a\allowbreak{}b2f39801\allowbreak{}6e0a8729\allowbreak{}31ce0832\allowbreak{}af3f2f6f\allowbreak{}123be521\\
Registro en 108 & \textemdash{}\\
\end{tabularx}\endgroup\par\vspace{6pt}
\Needspace{87pt}\phantomsection\label{trace-62}\noindent\hyperref[nav:vi-ruta-62]{\textbf{r2c7\_\allowbreak{}h-1v+1}}\quad Ordinal unificado / memoria: 62 / 63\par\nopagebreak
\begingroup\fontsize{8.3}{10.1}\selectfont\setlength{\tabcolsep}{4pt}\noindent\begin{tabularx}{\linewidth}{@{}p{.21\linewidth}X@{}}
Genealogía & fd3fc6ae\allowbreak{}2ff2d25b\allowbreak{}c3c572dc\allowbreak{}5362bf82\allowbreak{}ba08155c\allowbreak{}ab1135ce\allowbreak{}cc3e4ea8\allowbreak{}c6cabd14\\
Lector & 97b2ab07\allowbreak{}47f9cae8\allowbreak{}36c873b6\allowbreak{}9558f370\allowbreak{}f37393dc\allowbreak{}00bc1b47\allowbreak{}5ff34980\allowbreak{}99cf3f6a\\
Registro en 108 & 4ce890da\allowbreak{}3360de04\allowbreak{}fce789c2\allowbreak{}ce40701a\allowbreak{}acf357aa\allowbreak{}4d27239b\allowbreak{}5d03dfa5\allowbreak{}2c8c9f85\\
\end{tabularx}\endgroup\par\vspace{6pt}
\Needspace{87pt}\phantomsection\label{trace-63}\noindent\hyperref[nav:vi-ruta-63]{\textbf{r2c7\_\allowbreak{}h+1v-1}}\quad Ordinal unificado / memoria: 63 / 62\par\nopagebreak
\begingroup\fontsize{8.3}{10.1}\selectfont\setlength{\tabcolsep}{4pt}\noindent\begin{tabularx}{\linewidth}{@{}p{.21\linewidth}X@{}}
Genealogía & 6b038b1a\allowbreak{}3fa41286\allowbreak{}fe974ac6\allowbreak{}3c5d90cb\allowbreak{}981ae568\allowbreak{}eab928b7\allowbreak{}7b8abbe7\allowbreak{}3506531d\\
Lector & 5abbe501\allowbreak{}c8adf01b\allowbreak{}deab4cee\allowbreak{}f778d258\allowbreak{}79f7df34\allowbreak{}3fcf8dea\allowbreak{}b50624ce\allowbreak{}7852d47c\\
Registro en 108 & \textemdash{}\\
\end{tabularx}\endgroup\par\vspace{6pt}
\Needspace{87pt}\phantomsection\label{trace-64}\noindent\hyperref[nav:vi-ruta-64]{\textbf{r2c7\_\allowbreak{}h+1v+1}}\quad Ordinal unificado / memoria: 64 / 61\par\nopagebreak
\begingroup\fontsize{8.3}{10.1}\selectfont\setlength{\tabcolsep}{4pt}\noindent\begin{tabularx}{\linewidth}{@{}p{.21\linewidth}X@{}}
Genealogía & 7690711b\allowbreak{}d23b740a\allowbreak{}fdbe1c4f\allowbreak{}a7b76bb7\allowbreak{}8b49dfea\allowbreak{}24072771\allowbreak{}baed5540\allowbreak{}881ba560\\
Lector & 2e84b322\allowbreak{}980e807e\allowbreak{}05210dea\allowbreak{}2a1a280f\allowbreak{}6589f24b\allowbreak{}5d8d2b88\allowbreak{}a0a309b4\allowbreak{}0ad8d5a1\\
Registro en 108 & \textemdash{}\\
\end{tabularx}\endgroup\par\vspace{6pt}
\Needspace{87pt}\phantomsection\label{trace-65}\noindent\hyperref[nav:vi-ruta-65]{\textbf{r2c8\_\allowbreak{}h-1v-1}}\quad Ordinal unificado / memoria: 65 / 68\par\nopagebreak
\begingroup\fontsize{8.3}{10.1}\selectfont\setlength{\tabcolsep}{4pt}\noindent\begin{tabularx}{\linewidth}{@{}p{.21\linewidth}X@{}}
Genealogía & 59979b33\allowbreak{}9515cffd\allowbreak{}98571eb8\allowbreak{}6bd694e3\allowbreak{}3ea0f197\allowbreak{}3905ec0b\allowbreak{}f1c9caa7\allowbreak{}4ef63de7\\
Lector & 524796ca\allowbreak{}2925cfb1\allowbreak{}d109f28a\allowbreak{}8a8378d7\allowbreak{}78d802cf\allowbreak{}51f272b5\allowbreak{}cbd5bd45\allowbreak{}fcaa368a\\
Registro en 108 & d88dd1b6\allowbreak{}168ecd8f\allowbreak{}1b4de529\allowbreak{}649f5ba7\allowbreak{}3c37ee39\allowbreak{}b4c64650\allowbreak{}605497dd\allowbreak{}1d21e3f3\\
\end{tabularx}\endgroup\par\vspace{6pt}
\Needspace{87pt}\phantomsection\label{trace-66}\noindent\hyperref[nav:vi-ruta-66]{\textbf{r2c8\_\allowbreak{}h-1v+1}}\quad Ordinal unificado / memoria: 66 / 67\par\nopagebreak
\begingroup\fontsize{8.3}{10.1}\selectfont\setlength{\tabcolsep}{4pt}\noindent\begin{tabularx}{\linewidth}{@{}p{.21\linewidth}X@{}}
Genealogía & 354c6d9a\allowbreak{}f4e6a921\allowbreak{}677a2e1c\allowbreak{}a86c2615\allowbreak{}967aadbb\allowbreak{}a749be2f\allowbreak{}b9cbc947\allowbreak{}1af9ef30\\
Lector & 835f1e89\allowbreak{}a2323718\allowbreak{}255fdf26\allowbreak{}105c31d0\allowbreak{}c239439f\allowbreak{}72826178\allowbreak{}a5dfe560\allowbreak{}851db501\\
Registro en 108 & \textemdash{}\\
\end{tabularx}\endgroup\par\vspace{6pt}
\Needspace{87pt}\phantomsection\label{trace-67}\noindent\hyperref[nav:vi-ruta-67]{\textbf{r2c8\_\allowbreak{}h+1v-1}}\quad Ordinal unificado / memoria: 67 / 66\par\nopagebreak
\begingroup\fontsize{8.3}{10.1}\selectfont\setlength{\tabcolsep}{4pt}\noindent\begin{tabularx}{\linewidth}{@{}p{.21\linewidth}X@{}}
Genealogía & 5908d499\allowbreak{}c7cb58ca\allowbreak{}1019ee6f\allowbreak{}a049bb03\allowbreak{}5b537bb2\allowbreak{}8190db76\allowbreak{}80139634\allowbreak{}d9fbe00b\\
Lector & be902208\allowbreak{}821a34b0\allowbreak{}42d5e11f\allowbreak{}1d0fcc5e\allowbreak{}10dbf5a2\allowbreak{}f9625040\allowbreak{}dbc02e41\allowbreak{}6483f054\\
Registro en 108 & \textemdash{}\\
\end{tabularx}\endgroup\par\vspace{6pt}
\Needspace{87pt}\phantomsection\label{trace-68}\noindent\hyperref[nav:vi-ruta-68]{\textbf{r2c8\_\allowbreak{}h+1v+1}}\quad Ordinal unificado / memoria: 68 / 65\par\nopagebreak
\begingroup\fontsize{8.3}{10.1}\selectfont\setlength{\tabcolsep}{4pt}\noindent\begin{tabularx}{\linewidth}{@{}p{.21\linewidth}X@{}}
Genealogía & 84615c88\allowbreak{}02e6b1fd\allowbreak{}96b339d0\allowbreak{}d1d76101\allowbreak{}fbba5fc6\allowbreak{}4863032e\allowbreak{}5be62133\allowbreak{}8e563f29\\
Lector & feee18fc\allowbreak{}f009a483\allowbreak{}3205b507\allowbreak{}5b749373\allowbreak{}cf28f0b1\allowbreak{}13b0376d\allowbreak{}55a8f623\allowbreak{}1fe5d8b9\\
Registro en 108 & \textemdash{}\\
\end{tabularx}\endgroup\par\vspace{6pt}
\Needspace{87pt}\phantomsection\label{trace-69}\noindent\hyperref[nav:vi-ruta-69]{\textbf{r2c9\_\allowbreak{}h-1v-1}}\quad Ordinal unificado / memoria: 69 / 72\par\nopagebreak
\begingroup\fontsize{8.3}{10.1}\selectfont\setlength{\tabcolsep}{4pt}\noindent\begin{tabularx}{\linewidth}{@{}p{.21\linewidth}X@{}}
Genealogía & 10af243b\allowbreak{}dfdbc6b0\allowbreak{}4a47e1da\allowbreak{}71d5e5ca\allowbreak{}47143ea9\allowbreak{}930de740\allowbreak{}04529115\allowbreak{}4fd1ca26\\
Lector & 46188c47\allowbreak{}64cb2c18\allowbreak{}ad5e28f7\allowbreak{}d7701b20\allowbreak{}516f4776\allowbreak{}7883a606\allowbreak{}fd272ab0\allowbreak{}ce1b6e7f\\
Registro en 108 & \textemdash{}\\
\end{tabularx}\endgroup\par\vspace{6pt}
\Needspace{87pt}\phantomsection\label{trace-70}\noindent\hyperref[nav:vi-ruta-70]{\textbf{r2c9\_\allowbreak{}h-1v+1}}\quad Ordinal unificado / memoria: 70 / 71\par\nopagebreak
\begingroup\fontsize{8.3}{10.1}\selectfont\setlength{\tabcolsep}{4pt}\noindent\begin{tabularx}{\linewidth}{@{}p{.21\linewidth}X@{}}
Genealogía & 78a741ee\allowbreak{}3c7bfbe7\allowbreak{}ae864101\allowbreak{}5a71bd34\allowbreak{}f1888a67\allowbreak{}bce6ff00\allowbreak{}9113e38e\allowbreak{}c348d112\\
Lector & fc684626\allowbreak{}be7c75ff\allowbreak{}48c0e535\allowbreak{}0e38d10c\allowbreak{}764d4db5\allowbreak{}1dbaede3\allowbreak{}69429944\allowbreak{}03aa60f0\\
Registro en 108 & 4f85ae5a\allowbreak{}226daca7\allowbreak{}170f6876\allowbreak{}0b433575\allowbreak{}e50c6075\allowbreak{}5178c473\allowbreak{}034d1496\allowbreak{}819aa0a7\\
\end{tabularx}\endgroup\par\vspace{6pt}
\Needspace{87pt}\phantomsection\label{trace-71}\noindent\hyperref[nav:vi-ruta-71]{\textbf{r2c9\_\allowbreak{}h+1v-1}}\quad Ordinal unificado / memoria: 71 / 70\par\nopagebreak
\begingroup\fontsize{8.3}{10.1}\selectfont\setlength{\tabcolsep}{4pt}\noindent\begin{tabularx}{\linewidth}{@{}p{.21\linewidth}X@{}}
Genealogía & c0b14bfd\allowbreak{}df9f53d4\allowbreak{}228d9411\allowbreak{}817166fe\allowbreak{}0e3268f7\allowbreak{}dad8c998\allowbreak{}8d477618\allowbreak{}6e04fe7f\\
Lector & 78ccb54b\allowbreak{}24e9defb\allowbreak{}e6c1315c\allowbreak{}36f064a2\allowbreak{}e031e90a\allowbreak{}f5bc9cb6\allowbreak{}9113c1ba\allowbreak{}082b2a72\\
Registro en 108 & \textemdash{}\\
\end{tabularx}\endgroup\par\vspace{6pt}
\Needspace{87pt}\phantomsection\label{trace-72}\noindent\hyperref[nav:vi-ruta-72]{\textbf{r2c9\_\allowbreak{}h+1v+1}}\quad Ordinal unificado / memoria: 72 / 69\par\nopagebreak
\begingroup\fontsize{8.3}{10.1}\selectfont\setlength{\tabcolsep}{4pt}\noindent\begin{tabularx}{\linewidth}{@{}p{.21\linewidth}X@{}}
Genealogía & bb93b38f\allowbreak{}6a11bd24\allowbreak{}fe0d784b\allowbreak{}0382327e\allowbreak{}0e957737\allowbreak{}8da33db4\allowbreak{}ff0c2393\allowbreak{}36b2076c\\
Lector & 0323a244\allowbreak{}67e3bb3b\allowbreak{}19fb86ed\allowbreak{}bb720c35\allowbreak{}efac9bd2\allowbreak{}c76f856a\allowbreak{}02bbfaa6\allowbreak{}be1de3fb\\
Registro en 108 & \textemdash{}\\
\end{tabularx}\endgroup\par\vspace{6pt}
\Needspace{87pt}\phantomsection\label{trace-73}\noindent\hyperref[nav:vi-ruta-73]{\textbf{r3c1\_\allowbreak{}h-1v-1}}\quad Ordinal unificado / memoria: 73 / 76\par\nopagebreak
\begingroup\fontsize{8.3}{10.1}\selectfont\setlength{\tabcolsep}{4pt}\noindent\begin{tabularx}{\linewidth}{@{}p{.21\linewidth}X@{}}
Genealogía & b49a35c8\allowbreak{}0200681a\allowbreak{}5cd1c24a\allowbreak{}3af33bdc\allowbreak{}b3f5dfc1\allowbreak{}5693d461\allowbreak{}d3c7c6d1\allowbreak{}a094106a\\
Lector & 9cd3c1cd\allowbreak{}d9f24b93\allowbreak{}ee41db3b\allowbreak{}5db0405b\allowbreak{}d95787b2\allowbreak{}7d0b9563\allowbreak{}cb941f3b\allowbreak{}cf9ec663\\
Registro en 108 & \textemdash{}\\
\end{tabularx}\endgroup\par\vspace{6pt}
\Needspace{87pt}\phantomsection\label{trace-74}\noindent\hyperref[nav:vi-ruta-74]{\textbf{r3c1\_\allowbreak{}h-1v+1}}\quad Ordinal unificado / memoria: 74 / 75\par\nopagebreak
\begingroup\fontsize{8.3}{10.1}\selectfont\setlength{\tabcolsep}{4pt}\noindent\begin{tabularx}{\linewidth}{@{}p{.21\linewidth}X@{}}
Genealogía & aa18d948\allowbreak{}2f30ff1d\allowbreak{}1e8a4a2a\allowbreak{}55cb5650\allowbreak{}61e18e03\allowbreak{}d89d65c1\allowbreak{}83f098f5\allowbreak{}ecfa8627\\
Lector & 9f051964\allowbreak{}46d319f0\allowbreak{}bee5611c\allowbreak{}ef354231\allowbreak{}f598d0c6\allowbreak{}9a54da6b\allowbreak{}36803c84\allowbreak{}0ef9d155\\
Registro en 108 & 4f280c51\allowbreak{}aedd8c9f\allowbreak{}5e7c1710\allowbreak{}1e369b77\allowbreak{}4441e690\allowbreak{}405561db\allowbreak{}94125e58\allowbreak{}da5ec0c1\\
\end{tabularx}\endgroup\par\vspace{6pt}
\Needspace{87pt}\phantomsection\label{trace-75}\noindent\hyperref[nav:vi-ruta-75]{\textbf{r3c1\_\allowbreak{}h+1v-1}}\quad Ordinal unificado / memoria: 75 / 74\par\nopagebreak
\begingroup\fontsize{8.3}{10.1}\selectfont\setlength{\tabcolsep}{4pt}\noindent\begin{tabularx}{\linewidth}{@{}p{.21\linewidth}X@{}}
Genealogía & f93f9cc3\allowbreak{}c26fe793\allowbreak{}e3a4e633\allowbreak{}95e34edd\allowbreak{}0c959cb9\allowbreak{}bc132a58\allowbreak{}15d675d4\allowbreak{}6041d5b6\\
Lector & 8343e0ea\allowbreak{}9628938a\allowbreak{}9aa892c5\allowbreak{}93f5925d\allowbreak{}0eace614\allowbreak{}709da841\allowbreak{}94ad5609\allowbreak{}bd12b447\\
Registro en 108 & \textemdash{}\\
\end{tabularx}\endgroup\par\vspace{6pt}
\Needspace{87pt}\phantomsection\label{trace-76}\noindent\hyperref[nav:vi-ruta-76]{\textbf{r3c1\_\allowbreak{}h+1v+1}}\quad Ordinal unificado / memoria: 76 / 73\par\nopagebreak
\begingroup\fontsize{8.3}{10.1}\selectfont\setlength{\tabcolsep}{4pt}\noindent\begin{tabularx}{\linewidth}{@{}p{.21\linewidth}X@{}}
Genealogía & b5ee1bee\allowbreak{}5453ebdf\allowbreak{}2ca8a16f\allowbreak{}35b7700f\allowbreak{}e788e779\allowbreak{}16965f33\allowbreak{}ae801dd0\allowbreak{}14d7a637\\
Lector & 6aee9e8b\allowbreak{}24e2230f\allowbreak{}fb355d66\allowbreak{}6b1169ce\allowbreak{}d471e4a3\allowbreak{}b73faf2c\allowbreak{}e8a3168a\allowbreak{}2f909cb2\\
Registro en 108 & \textemdash{}\\
\end{tabularx}\endgroup\par\vspace{6pt}
\Needspace{87pt}\phantomsection\label{trace-77}\noindent\hyperref[nav:vi-ruta-77]{\textbf{r3c2\_\allowbreak{}h-1v-1}}\quad Ordinal unificado / memoria: 77 / 80\par\nopagebreak
\begingroup\fontsize{8.3}{10.1}\selectfont\setlength{\tabcolsep}{4pt}\noindent\begin{tabularx}{\linewidth}{@{}p{.21\linewidth}X@{}}
Genealogía & 5fbc72ff\allowbreak{}0a96e82d\allowbreak{}4616ba1a\allowbreak{}6ef070f7\allowbreak{}a682551f\allowbreak{}464c637a\allowbreak{}d43a97a1\allowbreak{}1f6e8a97\\
Lector & 92a383d8\allowbreak{}2ac125a3\allowbreak{}74465a81\allowbreak{}2287a783\allowbreak{}220c4993\allowbreak{}af802ee3\allowbreak{}22a89c2c\allowbreak{}20798cbd\\
Registro en 108 & \textemdash{}\\
\end{tabularx}\endgroup\par\vspace{6pt}
\Needspace{87pt}\phantomsection\label{trace-78}\noindent\hyperref[nav:vi-ruta-78]{\textbf{r3c2\_\allowbreak{}h-1v+1}}\quad Ordinal unificado / memoria: 78 / 79\par\nopagebreak
\begingroup\fontsize{8.3}{10.1}\selectfont\setlength{\tabcolsep}{4pt}\noindent\begin{tabularx}{\linewidth}{@{}p{.21\linewidth}X@{}}
Genealogía & 69e2bd91\allowbreak{}cf69da87\allowbreak{}5a43f975\allowbreak{}23596d23\allowbreak{}1b3e64e5\allowbreak{}813e5f26\allowbreak{}c0c2804e\allowbreak{}11812077\\
Lector & 3f2554e2\allowbreak{}7e8f08e1\allowbreak{}8bce3bad\allowbreak{}ea9cd8ff\allowbreak{}21745b26\allowbreak{}de177ed6\allowbreak{}08864d59\allowbreak{}1051a07e\\
Registro en 108 & \textemdash{}\\
\end{tabularx}\endgroup\par\vspace{6pt}
\Needspace{87pt}\phantomsection\label{trace-79}\noindent\hyperref[nav:vi-ruta-79]{\textbf{r3c2\_\allowbreak{}h+1v-1}}\quad Ordinal unificado / memoria: 79 / 78\par\nopagebreak
\begingroup\fontsize{8.3}{10.1}\selectfont\setlength{\tabcolsep}{4pt}\noindent\begin{tabularx}{\linewidth}{@{}p{.21\linewidth}X@{}}
Genealogía & c11bf3a5\allowbreak{}5ac78a0a\allowbreak{}7ac79598\allowbreak{}bfd55908\allowbreak{}c204e4a6\allowbreak{}0ec21403\allowbreak{}b684c336\allowbreak{}250918b2\\
Lector & 4f92d4ea\allowbreak{}7024f550\allowbreak{}1e8b0a88\allowbreak{}2bbee097\allowbreak{}36d0169b\allowbreak{}49bfb21a\allowbreak{}d0504f8c\allowbreak{}1d37c43a\\
Registro en 108 & \textemdash{}\\
\end{tabularx}\endgroup\par\vspace{6pt}
\Needspace{87pt}\phantomsection\label{trace-80}\noindent\hyperref[nav:vi-ruta-80]{\textbf{r3c2\_\allowbreak{}h+1v+1}}\quad Ordinal unificado / memoria: 80 / 77\par\nopagebreak
\begingroup\fontsize{8.3}{10.1}\selectfont\setlength{\tabcolsep}{4pt}\noindent\begin{tabularx}{\linewidth}{@{}p{.21\linewidth}X@{}}
Genealogía & 84b3ba30\allowbreak{}95f1140b\allowbreak{}146b108f\allowbreak{}d4846193\allowbreak{}32fad0f9\allowbreak{}436e0ecc\allowbreak{}1ecd0a44\allowbreak{}3acd145b\\
Lector & 6cd46053\allowbreak{}b4aac7eb\allowbreak{}750510ca\allowbreak{}7ff5aa4f\allowbreak{}f8e13be0\allowbreak{}97058396\allowbreak{}eee32adb\allowbreak{}341196d1\\
Registro en 108 & 3364070d\allowbreak{}4d0d9982\allowbreak{}5f096a5b\allowbreak{}3dd87451\allowbreak{}aa757f0f\allowbreak{}85aa6ab9\allowbreak{}74218271\allowbreak{}ed7418df\\
\end{tabularx}\endgroup\par\vspace{6pt}
\Needspace{87pt}\phantomsection\label{trace-81}\noindent\hyperref[nav:vi-ruta-81]{\textbf{r3c3\_\allowbreak{}h-1v-1}}\quad Ordinal unificado / memoria: 81 / 84\par\nopagebreak
\begingroup\fontsize{8.3}{10.1}\selectfont\setlength{\tabcolsep}{4pt}\noindent\begin{tabularx}{\linewidth}{@{}p{.21\linewidth}X@{}}
Genealogía & 90be6b64\allowbreak{}5dcb7c34\allowbreak{}479da4dd\allowbreak{}b58dc6b3\allowbreak{}31678651\allowbreak{}b48f48bf\allowbreak{}a32fa2f1\allowbreak{}cacbf2d2\\
Lector & 0a266f2d\allowbreak{}e19b3916\allowbreak{}bf94b020\allowbreak{}3007d4c4\allowbreak{}edc9fba6\allowbreak{}89170894\allowbreak{}8e1cd7c8\allowbreak{}d352827d\\
Registro en 108 & \textemdash{}\\
\end{tabularx}\endgroup\par\vspace{6pt}
\Needspace{87pt}\phantomsection\label{trace-82}\noindent\hyperref[nav:vi-ruta-82]{\textbf{r3c3\_\allowbreak{}h-1v+1}}\quad Ordinal unificado / memoria: 82 / 83\par\nopagebreak
\begingroup\fontsize{8.3}{10.1}\selectfont\setlength{\tabcolsep}{4pt}\noindent\begin{tabularx}{\linewidth}{@{}p{.21\linewidth}X@{}}
Genealogía & 970cf8e4\allowbreak{}20061a4c\allowbreak{}2e86d72e\allowbreak{}44ab8e14\allowbreak{}d201e17d\allowbreak{}5273d819\allowbreak{}f231ca3c\allowbreak{}228b99cc\\
Lector & e52e1aa4\allowbreak{}90feb02a\allowbreak{}80eacd19\allowbreak{}facbd502\allowbreak{}f5a2bead\allowbreak{}0ea1c7ba\allowbreak{}6885899e\allowbreak{}b000c242\\
Registro en 108 & \textemdash{}\\
\end{tabularx}\endgroup\par\vspace{6pt}
\Needspace{87pt}\phantomsection\label{trace-83}\noindent\hyperref[nav:vi-ruta-83]{\textbf{r3c3\_\allowbreak{}h+1v-1}}\quad Ordinal unificado / memoria: 83 / 82\par\nopagebreak
\begingroup\fontsize{8.3}{10.1}\selectfont\setlength{\tabcolsep}{4pt}\noindent\begin{tabularx}{\linewidth}{@{}p{.21\linewidth}X@{}}
Genealogía & fb908c9b\allowbreak{}7d05c874\allowbreak{}0756c16c\allowbreak{}e54c6e78\allowbreak{}90b67ab3\allowbreak{}adc021ba\allowbreak{}17bc80d1\allowbreak{}580a6a7e\\
Lector & 540ae414\allowbreak{}abfa6f64\allowbreak{}6b2b1e35\allowbreak{}635c7d7f\allowbreak{}e2834db6\allowbreak{}ead0ec11\allowbreak{}42e9fb00\allowbreak{}04b694aa\\
Registro en 108 & 8fc1d1a0\allowbreak{}17faec79\allowbreak{}2451c97d\allowbreak{}492516e4\allowbreak{}42f2ce08\allowbreak{}857f8c2a\allowbreak{}a8930ac0\allowbreak{}1f5a0821\\
\end{tabularx}\endgroup\par\vspace{6pt}
\Needspace{87pt}\phantomsection\label{trace-84}\noindent\hyperref[nav:vi-ruta-84]{\textbf{r3c3\_\allowbreak{}h+1v+1}}\quad Ordinal unificado / memoria: 84 / 81\par\nopagebreak
\begingroup\fontsize{8.3}{10.1}\selectfont\setlength{\tabcolsep}{4pt}\noindent\begin{tabularx}{\linewidth}{@{}p{.21\linewidth}X@{}}
Genealogía & f04730d7\allowbreak{}e2e26b77\allowbreak{}54d2c540\allowbreak{}697c1ccd\allowbreak{}c7a27425\allowbreak{}5c88f943\allowbreak{}077621c5\allowbreak{}e09b74ec\\
Lector & 268d8ac4\allowbreak{}9e67cc16\allowbreak{}79252262\allowbreak{}900b5b9d\allowbreak{}f37d5c42\allowbreak{}3422bb7f\allowbreak{}2b05970f\allowbreak{}b9a860da\\
Registro en 108 & \textemdash{}\\
\end{tabularx}\endgroup\par\vspace{6pt}
\Needspace{87pt}\phantomsection\label{trace-85}\noindent\hyperref[nav:vi-ruta-85]{\textbf{r3c4\_\allowbreak{}h-1v-1}}\quad Ordinal unificado / memoria: 85 / 88\par\nopagebreak
\begingroup\fontsize{8.3}{10.1}\selectfont\setlength{\tabcolsep}{4pt}\noindent\begin{tabularx}{\linewidth}{@{}p{.21\linewidth}X@{}}
Genealogía & fe19794a\allowbreak{}de36889f\allowbreak{}09cc94a7\allowbreak{}7d8a9722\allowbreak{}8b8530f2\allowbreak{}fc302f96\allowbreak{}a3cae6b7\allowbreak{}1d488126\\
Lector & 4fb57798\allowbreak{}b1f66825\allowbreak{}757136b6\allowbreak{}91435da4\allowbreak{}276f1d74\allowbreak{}75fb32d3\allowbreak{}7ada397c\allowbreak{}6bc2cc10\\
Registro en 108 & 47361123\allowbreak{}860c471c\allowbreak{}5c3e3660\allowbreak{}c7e98011\allowbreak{}f4bdf752\allowbreak{}d76ffd84\allowbreak{}b3957ffe\allowbreak{}686a0fef\\
\end{tabularx}\endgroup\par\vspace{6pt}
\Needspace{87pt}\phantomsection\label{trace-86}\noindent\hyperref[nav:vi-ruta-86]{\textbf{r3c4\_\allowbreak{}h-1v+1}}\quad Ordinal unificado / memoria: 86 / 87\par\nopagebreak
\begingroup\fontsize{8.3}{10.1}\selectfont\setlength{\tabcolsep}{4pt}\noindent\begin{tabularx}{\linewidth}{@{}p{.21\linewidth}X@{}}
Genealogía & 735a2f2a\allowbreak{}02b51b9b\allowbreak{}28d0c966\allowbreak{}1b965464\allowbreak{}ea385226\allowbreak{}29179c97\allowbreak{}92287ffe\allowbreak{}efd17733\\
Lector & 0ecdaa6d\allowbreak{}f15a6f0f\allowbreak{}268fb041\allowbreak{}bac3c594\allowbreak{}cca56daa\allowbreak{}e8f41abd\allowbreak{}c8b4d3ed\allowbreak{}d12e479b\\
Registro en 108 & \textemdash{}\\
\end{tabularx}\endgroup\par\vspace{6pt}
\Needspace{87pt}\phantomsection\label{trace-87}\noindent\hyperref[nav:vi-ruta-87]{\textbf{r3c4\_\allowbreak{}h+1v-1}}\quad Ordinal unificado / memoria: 87 / 86\par\nopagebreak
\begingroup\fontsize{8.3}{10.1}\selectfont\setlength{\tabcolsep}{4pt}\noindent\begin{tabularx}{\linewidth}{@{}p{.21\linewidth}X@{}}
Genealogía & 1931926d\allowbreak{}8929382c\allowbreak{}703b3c3e\allowbreak{}fe5ee8d5\allowbreak{}ec8768aa\allowbreak{}ba2fdc56\allowbreak{}8564d21d\allowbreak{}05a242f4\\
Lector & 2dda1905\allowbreak{}a5534096\allowbreak{}e9ad9a42\allowbreak{}7c6d8048\allowbreak{}fc9d39c0\allowbreak{}c5825306\allowbreak{}d0203b1a\allowbreak{}12572285\\
Registro en 108 & \textemdash{}\\
\end{tabularx}\endgroup\par\vspace{6pt}
\Needspace{87pt}\phantomsection\label{trace-88}\noindent\hyperref[nav:vi-ruta-88]{\textbf{r3c4\_\allowbreak{}h+1v+1}}\quad Ordinal unificado / memoria: 88 / 85\par\nopagebreak
\begingroup\fontsize{8.3}{10.1}\selectfont\setlength{\tabcolsep}{4pt}\noindent\begin{tabularx}{\linewidth}{@{}p{.21\linewidth}X@{}}
Genealogía & 44ed4403\allowbreak{}0e829f7e\allowbreak{}f3f2d253\allowbreak{}9c0ecae1\allowbreak{}af3389fc\allowbreak{}b03b6af4\allowbreak{}43c75156\allowbreak{}88564f46\\
Lector & 1ee0bda6\allowbreak{}2536a4e3\allowbreak{}64a3618e\allowbreak{}5ba602de\allowbreak{}d8f28109\allowbreak{}b25f42ce\allowbreak{}1150ff31\allowbreak{}5daeb5f9\\
Registro en 108 & \textemdash{}\\
\end{tabularx}\endgroup\par\vspace{6pt}
\Needspace{87pt}\phantomsection\label{trace-89}\noindent\hyperref[nav:vi-ruta-89]{\textbf{r3c5\_\allowbreak{}h-1v-1}}\quad Ordinal unificado / memoria: 89 / 92\par\nopagebreak
\begingroup\fontsize{8.3}{10.1}\selectfont\setlength{\tabcolsep}{4pt}\noindent\begin{tabularx}{\linewidth}{@{}p{.21\linewidth}X@{}}
Genealogía & e953714a\allowbreak{}3f1290a2\allowbreak{}21bb6aaa\allowbreak{}7ffb0b8c\allowbreak{}7ba40f3d\allowbreak{}e4825cb2\allowbreak{}31af24f8\allowbreak{}8c67487c\\
Lector & 516ac20b\allowbreak{}777e49c4\allowbreak{}9141c04f\allowbreak{}9435195b\allowbreak{}0a72bf8f\allowbreak{}d107bf11\allowbreak{}31f9797a\allowbreak{}550996f1\\
Registro en 108 & \textemdash{}\\
\end{tabularx}\endgroup\par\vspace{6pt}
\Needspace{87pt}\phantomsection\label{trace-90}\noindent\hyperref[nav:vi-ruta-90]{\textbf{r3c5\_\allowbreak{}h-1v+1}}\quad Ordinal unificado / memoria: 90 / 91\par\nopagebreak
\begingroup\fontsize{8.3}{10.1}\selectfont\setlength{\tabcolsep}{4pt}\noindent\begin{tabularx}{\linewidth}{@{}p{.21\linewidth}X@{}}
Genealogía & b34cbe3d\allowbreak{}2a71fe4f\allowbreak{}b1ebe419\allowbreak{}4d4a8771\allowbreak{}8d2bd15e\allowbreak{}d5ffc2b7\allowbreak{}22216776\allowbreak{}1d430287\\
Lector & 64735248\allowbreak{}9028a5eb\allowbreak{}55e737d1\allowbreak{}ba0d6c52\allowbreak{}67b15924\allowbreak{}d267d799\allowbreak{}5e4d3e1f\allowbreak{}03269bfd\\
Registro en 108 & 2bdf6193\allowbreak{}049e9d8f\allowbreak{}1fc69526\allowbreak{}73cd5a6d\allowbreak{}3ece96f2\allowbreak{}80eea325\allowbreak{}13f93207\allowbreak{}a69269c1\\
\end{tabularx}\endgroup\par\vspace{6pt}
\Needspace{87pt}\phantomsection\label{trace-91}\noindent\hyperref[nav:vi-ruta-91]{\textbf{r3c5\_\allowbreak{}h+1v-1}}\quad Ordinal unificado / memoria: 91 / 90\par\nopagebreak
\begingroup\fontsize{8.3}{10.1}\selectfont\setlength{\tabcolsep}{4pt}\noindent\begin{tabularx}{\linewidth}{@{}p{.21\linewidth}X@{}}
Genealogía & 257cfe72\allowbreak{}d342b833\allowbreak{}0026aca6\allowbreak{}da1825c9\allowbreak{}35e6890f\allowbreak{}3fa068db\allowbreak{}6795c936\allowbreak{}086cae39\\
Lector & fceaa589\allowbreak{}5efdf740\allowbreak{}e8a408a6\allowbreak{}5e4bd72b\allowbreak{}828faadd\allowbreak{}8731c947\allowbreak{}705e6375\allowbreak{}e4f6895d\\
Registro en 108 & \textemdash{}\\
\end{tabularx}\endgroup\par\vspace{6pt}
\Needspace{87pt}\phantomsection\label{trace-92}\noindent\hyperref[nav:vi-ruta-92]{\textbf{r3c5\_\allowbreak{}h+1v+1}}\quad Ordinal unificado / memoria: 92 / 89\par\nopagebreak
\begingroup\fontsize{8.3}{10.1}\selectfont\setlength{\tabcolsep}{4pt}\noindent\begin{tabularx}{\linewidth}{@{}p{.21\linewidth}X@{}}
Genealogía & 5144d953\allowbreak{}29eaf9c2\allowbreak{}e01d5312\allowbreak{}581dfa5a\allowbreak{}b24fac78\allowbreak{}ab7b2ea1\allowbreak{}24db8d95\allowbreak{}db7a998c\\
Lector & 0ac1f885\allowbreak{}d966f22d\allowbreak{}fb79a602\allowbreak{}b576a688\allowbreak{}b351e32d\allowbreak{}455ddca0\allowbreak{}4b3f46b9\allowbreak{}37f49b0e\\
Registro en 108 & \textemdash{}\\
\end{tabularx}\endgroup\par\vspace{6pt}
\Needspace{87pt}\phantomsection\label{trace-93}\noindent\hyperref[nav:vi-ruta-93]{\textbf{r3c6\_\allowbreak{}h-1v-1}}\quad Ordinal unificado / memoria: 93 / 96\par\nopagebreak
\begingroup\fontsize{8.3}{10.1}\selectfont\setlength{\tabcolsep}{4pt}\noindent\begin{tabularx}{\linewidth}{@{}p{.21\linewidth}X@{}}
Genealogía & 8e5c9476\allowbreak{}3ee2ea8f\allowbreak{}fb481fc1\allowbreak{}fefbf17f\allowbreak{}6c771f79\allowbreak{}14fcfcf1\allowbreak{}80fc6ba3\allowbreak{}84a48110\\
Lector & e9907b10\allowbreak{}8eaaf6db\allowbreak{}8390e712\allowbreak{}9e3c3d3f\allowbreak{}5c497165\allowbreak{}928e17a4\allowbreak{}b9f231a7\allowbreak{}ba1fa7d0\\
Registro en 108 & \textemdash{}\\
\end{tabularx}\endgroup\par\vspace{6pt}
\Needspace{87pt}\phantomsection\label{trace-94}\noindent\hyperref[nav:vi-ruta-94]{\textbf{r3c6\_\allowbreak{}h-1v+1}}\quad Ordinal unificado / memoria: 94 / 95\par\nopagebreak
\begingroup\fontsize{8.3}{10.1}\selectfont\setlength{\tabcolsep}{4pt}\noindent\begin{tabularx}{\linewidth}{@{}p{.21\linewidth}X@{}}
Genealogía & e11e9cb2\allowbreak{}d964364d\allowbreak{}36ac21f7\allowbreak{}6809e6d1\allowbreak{}f0b31214\allowbreak{}3e898c58\allowbreak{}4c45d1ef\allowbreak{}8bf497ca\\
Lector & 0b1b5faa\allowbreak{}852ad0e6\allowbreak{}071d9578\allowbreak{}6097ecab\allowbreak{}1e741a4f\allowbreak{}64d3bdd1\allowbreak{}d778f85c\allowbreak{}74963cbd\\
Registro en 108 & 406b35c8\allowbreak{}5cc7381b\allowbreak{}8d89efa8\allowbreak{}2a29dacc\allowbreak{}aa3ddcc2\allowbreak{}4246e6c0\allowbreak{}4702d902\allowbreak{}3088d3ad\\
\end{tabularx}\endgroup\par\vspace{6pt}
\Needspace{87pt}\phantomsection\label{trace-95}\noindent\hyperref[nav:vi-ruta-95]{\textbf{r3c6\_\allowbreak{}h+1v-1}}\quad Ordinal unificado / memoria: 95 / 94\par\nopagebreak
\begingroup\fontsize{8.3}{10.1}\selectfont\setlength{\tabcolsep}{4pt}\noindent\begin{tabularx}{\linewidth}{@{}p{.21\linewidth}X@{}}
Genealogía & b69e41e3\allowbreak{}adf88e39\allowbreak{}7196cba6\allowbreak{}a46da696\allowbreak{}41203067\allowbreak{}a7118c65\allowbreak{}e4b64e8a\allowbreak{}a46b34d1\\
Lector & 783dfded\allowbreak{}cf8fe84b\allowbreak{}dabe1f1b\allowbreak{}52581403\allowbreak{}e35dfe49\allowbreak{}f318adb4\allowbreak{}1101ada8\allowbreak{}b61c501c\\
Registro en 108 & \textemdash{}\\
\end{tabularx}\endgroup\par\vspace{6pt}
\Needspace{87pt}\phantomsection\label{trace-96}\noindent\hyperref[nav:vi-ruta-96]{\textbf{r3c6\_\allowbreak{}h+1v+1}}\quad Ordinal unificado / memoria: 96 / 93\par\nopagebreak
\begingroup\fontsize{8.3}{10.1}\selectfont\setlength{\tabcolsep}{4pt}\noindent\begin{tabularx}{\linewidth}{@{}p{.21\linewidth}X@{}}
Genealogía & 6288242c\allowbreak{}4687f7c0\allowbreak{}ef85b06a\allowbreak{}62c64390\allowbreak{}d926911d\allowbreak{}a1796280\allowbreak{}87bd63b7\allowbreak{}387a992d\\
Lector & 9055dce2\allowbreak{}3724fb2e\allowbreak{}e036867b\allowbreak{}505c820f\allowbreak{}79933f0a\allowbreak{}7d4ecd3d\allowbreak{}289c0b96\allowbreak{}9ca46da2\\
Registro en 108 & \textemdash{}\\
\end{tabularx}\endgroup\par\vspace{6pt}
\Needspace{87pt}\phantomsection\label{trace-97}\noindent\hyperref[nav:vi-ruta-97]{\textbf{r3c7\_\allowbreak{}h-1v-1}}\quad Ordinal unificado / memoria: 97 / 100\par\nopagebreak
\begingroup\fontsize{8.3}{10.1}\selectfont\setlength{\tabcolsep}{4pt}\noindent\begin{tabularx}{\linewidth}{@{}p{.21\linewidth}X@{}}
Genealogía & 6589a27f\allowbreak{}5db15e2f\allowbreak{}aedf8e0b\allowbreak{}f0f4b488\allowbreak{}87925aa9\allowbreak{}6b94b118\allowbreak{}940ab4b2\allowbreak{}75b4f51e\\
Lector & a00262e0\allowbreak{}cf95aa0f\allowbreak{}99b91ef0\allowbreak{}81a1ac39\allowbreak{}c8c5f9c7\allowbreak{}ba5b825e\allowbreak{}b0d5f9a8\allowbreak{}81b4f01e\\
Registro en 108 & \textemdash{}\\
\end{tabularx}\endgroup\par\vspace{6pt}
\Needspace{87pt}\phantomsection\label{trace-98}\noindent\hyperref[nav:vi-ruta-98]{\textbf{r3c7\_\allowbreak{}h-1v+1}}\quad Ordinal unificado / memoria: 98 / 99\par\nopagebreak
\begingroup\fontsize{8.3}{10.1}\selectfont\setlength{\tabcolsep}{4pt}\noindent\begin{tabularx}{\linewidth}{@{}p{.21\linewidth}X@{}}
Genealogía & 74fca7a0\allowbreak{}65152326\allowbreak{}f2680e83\allowbreak{}c4eb79ed\allowbreak{}dfceaf25\allowbreak{}5870015f\allowbreak{}6dc3f0dc\allowbreak{}d4da23d4\\
Lector & d4d5eba8\allowbreak{}2dff0efd\allowbreak{}18906a4c\allowbreak{}20b1c052\allowbreak{}b04b6346\allowbreak{}0dea63d1\allowbreak{}c3908efa\allowbreak{}b54ca465\\
Registro en 108 & \textemdash{}\\
\end{tabularx}\endgroup\par\vspace{6pt}
\Needspace{87pt}\phantomsection\label{trace-99}\noindent\hyperref[nav:vi-ruta-99]{\textbf{r3c7\_\allowbreak{}h+1v-1}}\quad Ordinal unificado / memoria: 99 / 98\par\nopagebreak
\begingroup\fontsize{8.3}{10.1}\selectfont\setlength{\tabcolsep}{4pt}\noindent\begin{tabularx}{\linewidth}{@{}p{.21\linewidth}X@{}}
Genealogía & e56e3ff2\allowbreak{}92c8e252\allowbreak{}0d3cecaa\allowbreak{}4e8d2963\allowbreak{}24ef1d8f\allowbreak{}8f682120\allowbreak{}e853d2eb\allowbreak{}302e843b\\
Lector & 809858a6\allowbreak{}f732a2d7\allowbreak{}201ea656\allowbreak{}68fee396\allowbreak{}0f55de4f\allowbreak{}c745fd45\allowbreak{}395c2dd5\allowbreak{}e8199cd5\\
Registro en 108 & \textemdash{}\\
\end{tabularx}\endgroup\par\vspace{6pt}
\Needspace{87pt}\phantomsection\label{trace-100}\noindent\hyperref[nav:vi-ruta-100]{\textbf{r3c7\_\allowbreak{}h+1v+1}}\quad Ordinal unificado / memoria: 100 / 97\par\nopagebreak
\begingroup\fontsize{8.3}{10.1}\selectfont\setlength{\tabcolsep}{4pt}\noindent\begin{tabularx}{\linewidth}{@{}p{.21\linewidth}X@{}}
Genealogía & 77168e1c\allowbreak{}5f252e65\allowbreak{}3c32d6b3\allowbreak{}899be2f0\allowbreak{}9ce151e2\allowbreak{}31c9bec5\allowbreak{}285c3e8b\allowbreak{}aa931507\\
Lector & 78e3b7b4\allowbreak{}81481612\allowbreak{}e0e698c1\allowbreak{}5bce3bb6\allowbreak{}731d92e3\allowbreak{}45512472\allowbreak{}bbfdf232\allowbreak{}57191308\\
Registro en 108 & 6124227c\allowbreak{}e899ccb2\allowbreak{}df636930\allowbreak{}314e07ff\allowbreak{}8fc23be0\allowbreak{}d2d973f8\allowbreak{}128c25ed\allowbreak{}2e55be47\\
\end{tabularx}\endgroup\par\vspace{6pt}
\Needspace{87pt}\phantomsection\label{trace-101}\noindent\hyperref[nav:vi-ruta-101]{\textbf{r3c8\_\allowbreak{}h-1v-1}}\quad Ordinal unificado / memoria: 101 / 104\par\nopagebreak
\begingroup\fontsize{8.3}{10.1}\selectfont\setlength{\tabcolsep}{4pt}\noindent\begin{tabularx}{\linewidth}{@{}p{.21\linewidth}X@{}}
Genealogía & a3c4b2e4\allowbreak{}2baf237a\allowbreak{}5e4c2f9b\allowbreak{}552ebcd2\allowbreak{}9e8bfa6d\allowbreak{}4980de49\allowbreak{}80f828b3\allowbreak{}77f1771b\\
Lector & f96905bc\allowbreak{}da720515\allowbreak{}1dd5dacf\allowbreak{}1f652576\allowbreak{}2d535f2d\allowbreak{}d174b1c3\allowbreak{}e19f351b\allowbreak{}ded335b8\\
Registro en 108 & \textemdash{}\\
\end{tabularx}\endgroup\par\vspace{6pt}
\Needspace{87pt}\phantomsection\label{trace-102}\noindent\hyperref[nav:vi-ruta-102]{\textbf{r3c8\_\allowbreak{}h-1v+1}}\quad Ordinal unificado / memoria: 102 / 103\par\nopagebreak
\begingroup\fontsize{8.3}{10.1}\selectfont\setlength{\tabcolsep}{4pt}\noindent\begin{tabularx}{\linewidth}{@{}p{.21\linewidth}X@{}}
Genealogía & 30b5ef0c\allowbreak{}47761242\allowbreak{}ea250c0b\allowbreak{}d063668a\allowbreak{}94033577\allowbreak{}35c2c59f\allowbreak{}ab1ad3ef\allowbreak{}3b33a931\\
Lector & 11b41f42\allowbreak{}21d7ba52\allowbreak{}85c94c73\allowbreak{}05abf324\allowbreak{}2a05aad8\allowbreak{}6ea35723\allowbreak{}202c943b\allowbreak{}48489ea1\\
Registro en 108 & 0dd6b250\allowbreak{}296cb9be\allowbreak{}5eb2e1d1\allowbreak{}9659528b\allowbreak{}7de17fdc\allowbreak{}8796cf0e\allowbreak{}c3cc88e3\allowbreak{}aabfa877\\
\end{tabularx}\endgroup\par\vspace{6pt}
\Needspace{87pt}\phantomsection\label{trace-103}\noindent\hyperref[nav:vi-ruta-103]{\textbf{r3c8\_\allowbreak{}h+1v-1}}\quad Ordinal unificado / memoria: 103 / 102\par\nopagebreak
\begingroup\fontsize{8.3}{10.1}\selectfont\setlength{\tabcolsep}{4pt}\noindent\begin{tabularx}{\linewidth}{@{}p{.21\linewidth}X@{}}
Genealogía & d954a2a9\allowbreak{}083ac812\allowbreak{}24a21b4c\allowbreak{}21d85c16\allowbreak{}8f7e0b0f\allowbreak{}95f77416\allowbreak{}3b3907a8\allowbreak{}c521f3ce\\
Lector & bce6affb\allowbreak{}972de833\allowbreak{}92f296ec\allowbreak{}01c79cf6\allowbreak{}5f166244\allowbreak{}1464abdb\allowbreak{}46015ef5\allowbreak{}f9986e86\\
Registro en 108 & \textemdash{}\\
\end{tabularx}\endgroup\par\vspace{6pt}
\Needspace{87pt}\phantomsection\label{trace-104}\noindent\hyperref[nav:vi-ruta-104]{\textbf{r3c8\_\allowbreak{}h+1v+1}}\quad Ordinal unificado / memoria: 104 / 101\par\nopagebreak
\begingroup\fontsize{8.3}{10.1}\selectfont\setlength{\tabcolsep}{4pt}\noindent\begin{tabularx}{\linewidth}{@{}p{.21\linewidth}X@{}}
Genealogía & 03881732\allowbreak{}293dbdb0\allowbreak{}dd5bde4d\allowbreak{}431381e4\allowbreak{}c1078b1a\allowbreak{}d88c6594\allowbreak{}eb82b5dd\allowbreak{}d3f00398\\
Lector & c585e17a\allowbreak{}9f62ed44\allowbreak{}ecbccc25\allowbreak{}3a8758fe\allowbreak{}1c652c8a\allowbreak{}544d1587\allowbreak{}54d2ae6f\allowbreak{}ddf6df97\\
Registro en 108 & \textemdash{}\\
\end{tabularx}\endgroup\par\vspace{6pt}
\Needspace{87pt}\phantomsection\label{trace-105}\noindent\hyperref[nav:vi-ruta-105]{\textbf{r3c9\_\allowbreak{}h-1v-1}}\quad Ordinal unificado / memoria: 105 / 108\par\nopagebreak
\begingroup\fontsize{8.3}{10.1}\selectfont\setlength{\tabcolsep}{4pt}\noindent\begin{tabularx}{\linewidth}{@{}p{.21\linewidth}X@{}}
Genealogía & 8baf9b09\allowbreak{}874f716a\allowbreak{}91bea123\allowbreak{}45945a5f\allowbreak{}1bd1f4f2\allowbreak{}baa14a32\allowbreak{}4556a40c\allowbreak{}e2ba26af\\
Lector & a0b0a058\allowbreak{}859733c7\allowbreak{}04c7f2fe\allowbreak{}d250076e\allowbreak{}56e37e36\allowbreak{}f21b93d7\allowbreak{}55a4c033\allowbreak{}376ccdc1\\
Registro en 108 & f669a046\allowbreak{}ed98e775\allowbreak{}5abdbfba\allowbreak{}15ed1ac5\allowbreak{}5e8a6e22\allowbreak{}8b2b4eb3\allowbreak{}0e40719f\allowbreak{}e160aaae\\
\end{tabularx}\endgroup\par\vspace{6pt}
\Needspace{87pt}\phantomsection\label{trace-106}\noindent\hyperref[nav:vi-ruta-106]{\textbf{r3c9\_\allowbreak{}h-1v+1}}\quad Ordinal unificado / memoria: 106 / 107\par\nopagebreak
\begingroup\fontsize{8.3}{10.1}\selectfont\setlength{\tabcolsep}{4pt}\noindent\begin{tabularx}{\linewidth}{@{}p{.21\linewidth}X@{}}
Genealogía & 56152a5d\allowbreak{}e4e1754d\allowbreak{}fb18df17\allowbreak{}97256c5b\allowbreak{}173ad55b\allowbreak{}e9a8457f\allowbreak{}b6046732\allowbreak{}40f1c027\\
Lector & d706d3e3\allowbreak{}faf56f94\allowbreak{}f030d2d4\allowbreak{}769dda63\allowbreak{}99824d27\allowbreak{}fd2506f0\allowbreak{}451a4918\allowbreak{}a9a901fc\\
Registro en 108 & \textemdash{}\\
\end{tabularx}\endgroup\par\vspace{6pt}
\Needspace{87pt}\phantomsection\label{trace-107}\noindent\hyperref[nav:vi-ruta-107]{\textbf{r3c9\_\allowbreak{}h+1v-1}}\quad Ordinal unificado / memoria: 107 / 106\par\nopagebreak
\begingroup\fontsize{8.3}{10.1}\selectfont\setlength{\tabcolsep}{4pt}\noindent\begin{tabularx}{\linewidth}{@{}p{.21\linewidth}X@{}}
Genealogía & 98a70374\allowbreak{}da093a19\allowbreak{}29fe532b\allowbreak{}452dc2b5\allowbreak{}94c1f88d\allowbreak{}c9efbb4f\allowbreak{}8bb82206\allowbreak{}ed8e63b8\\
Lector & 956dd120\allowbreak{}5d0eaf9c\allowbreak{}e7232de7\allowbreak{}9c5985f6\allowbreak{}2b2f20e1\allowbreak{}807aa46b\allowbreak{}76e7ce28\allowbreak{}50285b90\\
Registro en 108 & \textemdash{}\\
\end{tabularx}\endgroup\par\vspace{6pt}
\Needspace{87pt}\phantomsection\label{trace-108}\noindent\hyperref[nav:vi-ruta-108]{\textbf{r3c9\_\allowbreak{}h+1v+1}}\quad Ordinal unificado / memoria: 108 / 105\par\nopagebreak
\begingroup\fontsize{8.3}{10.1}\selectfont\setlength{\tabcolsep}{4pt}\noindent\begin{tabularx}{\linewidth}{@{}p{.21\linewidth}X@{}}
Genealogía & 4df97b7f\allowbreak{}c1cea288\allowbreak{}d4b149a4\allowbreak{}3852d6e5\allowbreak{}4a768eaa\allowbreak{}7b1e17bd\allowbreak{}091da4b9\allowbreak{}eb93eb18\\
Lector & 7b1d3368\allowbreak{}504c06f4\allowbreak{}756793ae\allowbreak{}495b0eb9\allowbreak{}f07fef95\allowbreak{}77e95bcf\allowbreak{}66b649fa\allowbreak{}d5640235\\
Registro en 108 & \textemdash{}\\
\end{tabularx}\endgroup\par\vspace{6pt}
\Needspace{87pt}\phantomsection\label{trace-109}\noindent\hyperref[nav:vi-ruta-109]{\textbf{r4c1\_\allowbreak{}h-1v-1}}\quad Ordinal unificado / memoria: 109 / 112\par\nopagebreak
\begingroup\fontsize{8.3}{10.1}\selectfont\setlength{\tabcolsep}{4pt}\noindent\begin{tabularx}{\linewidth}{@{}p{.21\linewidth}X@{}}
Genealogía & 63616877\allowbreak{}dbaa1494\allowbreak{}e1e12a93\allowbreak{}696a27ab\allowbreak{}3d0d9a09\allowbreak{}bd6c27e7\allowbreak{}a8cf49db\allowbreak{}5f58e7c8\\
Lector & 3405682b\allowbreak{}f852ce02\allowbreak{}3ebe9993\allowbreak{}31169514\allowbreak{}2ffa0a43\allowbreak{}07c6c7a9\allowbreak{}da972184\allowbreak{}6f8f6ae1\\
Registro en 108 & e58666b4\allowbreak{}f6fa6d33\allowbreak{}a51e958e\allowbreak{}f358e1ae\allowbreak{}c182e5eb\allowbreak{}5a6f84e4\allowbreak{}24583f2d\allowbreak{}27bdfb85\\
\end{tabularx}\endgroup\par\vspace{6pt}
\Needspace{87pt}\phantomsection\label{trace-110}\noindent\hyperref[nav:vi-ruta-110]{\textbf{r4c1\_\allowbreak{}h-1v+1}}\quad Ordinal unificado / memoria: 110 / 111\par\nopagebreak
\begingroup\fontsize{8.3}{10.1}\selectfont\setlength{\tabcolsep}{4pt}\noindent\begin{tabularx}{\linewidth}{@{}p{.21\linewidth}X@{}}
Genealogía & 80a08e93\allowbreak{}0c195bab\allowbreak{}65a18748\allowbreak{}d6675121\allowbreak{}6778294e\allowbreak{}ac2b155d\allowbreak{}72ae50f3\allowbreak{}9d5d7dda\\
Lector & 2037eb27\allowbreak{}ad0ed4d7\allowbreak{}ff20939e\allowbreak{}d394b026\allowbreak{}3366036e\allowbreak{}5d5ff1cb\allowbreak{}c8635bf5\allowbreak{}87eee9a2\\
Registro en 108 & \textemdash{}\\
\end{tabularx}\endgroup\par\vspace{6pt}
\Needspace{87pt}\phantomsection\label{trace-111}\noindent\hyperref[nav:vi-ruta-111]{\textbf{r4c1\_\allowbreak{}h+1v-1}}\quad Ordinal unificado / memoria: 111 / 110\par\nopagebreak
\begingroup\fontsize{8.3}{10.1}\selectfont\setlength{\tabcolsep}{4pt}\noindent\begin{tabularx}{\linewidth}{@{}p{.21\linewidth}X@{}}
Genealogía & eb9fef92\allowbreak{}470f339d\allowbreak{}df8c9ed1\allowbreak{}517e5de1\allowbreak{}b10969e9\allowbreak{}ebb16591\allowbreak{}a7f34cf5\allowbreak{}d9a76997\\
Lector & 0a8b7f1d\allowbreak{}8bfb27ed\allowbreak{}38bd2268\allowbreak{}b2ef01ad\allowbreak{}94636043\allowbreak{}402cd263\allowbreak{}c98ad3f7\allowbreak{}a6693437\\
Registro en 108 & \textemdash{}\\
\end{tabularx}\endgroup\par\vspace{6pt}
\Needspace{87pt}\phantomsection\label{trace-112}\noindent\hyperref[nav:vi-ruta-112]{\textbf{r4c1\_\allowbreak{}h+1v+1}}\quad Ordinal unificado / memoria: 112 / 109\par\nopagebreak
\begingroup\fontsize{8.3}{10.1}\selectfont\setlength{\tabcolsep}{4pt}\noindent\begin{tabularx}{\linewidth}{@{}p{.21\linewidth}X@{}}
Genealogía & 322c7aa9\allowbreak{}113fe734\allowbreak{}4258693c\allowbreak{}e39c7280\allowbreak{}2f7f0c73\allowbreak{}dc8b45c2\allowbreak{}f74ec16a\allowbreak{}33fef9e4\\
Lector & f4cec5cc\allowbreak{}7eaa2637\allowbreak{}6f615523\allowbreak{}2dee3e99\allowbreak{}8411e318\allowbreak{}92f8981e\allowbreak{}72e525ab\allowbreak{}7e0192d8\\
Registro en 108 & \textemdash{}\\
\end{tabularx}\endgroup\par\vspace{6pt}
\Needspace{87pt}\phantomsection\label{trace-113}\noindent\hyperref[nav:vi-ruta-113]{\textbf{r4c2\_\allowbreak{}h-1v-1}}\quad Ordinal unificado / memoria: 113 / 116\par\nopagebreak
\begingroup\fontsize{8.3}{10.1}\selectfont\setlength{\tabcolsep}{4pt}\noindent\begin{tabularx}{\linewidth}{@{}p{.21\linewidth}X@{}}
Genealogía & 320c90bb\allowbreak{}3283149b\allowbreak{}e781a543\allowbreak{}a1877eb0\allowbreak{}bab2603c\allowbreak{}2de7a705\allowbreak{}296b6be5\allowbreak{}be9e1da5\\
Lector & fc501338\allowbreak{}6134f896\allowbreak{}41ceb3f1\allowbreak{}0fc2600b\allowbreak{}52ba7b58\allowbreak{}3a6be481\allowbreak{}6aecffff\allowbreak{}35bfab68\\
Registro en 108 & \textemdash{}\\
\end{tabularx}\endgroup\par\vspace{6pt}
\Needspace{87pt}\phantomsection\label{trace-114}\noindent\hyperref[nav:vi-ruta-114]{\textbf{r4c2\_\allowbreak{}h-1v+1}}\quad Ordinal unificado / memoria: 114 / 115\par\nopagebreak
\begingroup\fontsize{8.3}{10.1}\selectfont\setlength{\tabcolsep}{4pt}\noindent\begin{tabularx}{\linewidth}{@{}p{.21\linewidth}X@{}}
Genealogía & 4257bc9a\allowbreak{}976f4f17\allowbreak{}5e997aba\allowbreak{}070b542a\allowbreak{}dc329417\allowbreak{}c6dae0d0\allowbreak{}8e497d6d\allowbreak{}65d51564\\
Lector & 102c954f\allowbreak{}9d43af20\allowbreak{}49a828e9\allowbreak{}72c927ba\allowbreak{}ce8ca4c9\allowbreak{}6291f65e\allowbreak{}fffa01ca\allowbreak{}e7c2fdfb\\
Registro en 108 & \textemdash{}\\
\end{tabularx}\endgroup\par\vspace{6pt}
\Needspace{87pt}\phantomsection\label{trace-115}\noindent\hyperref[nav:vi-ruta-115]{\textbf{r4c2\_\allowbreak{}h+1v-1}}\quad Ordinal unificado / memoria: 115 / 114\par\nopagebreak
\begingroup\fontsize{8.3}{10.1}\selectfont\setlength{\tabcolsep}{4pt}\noindent\begin{tabularx}{\linewidth}{@{}p{.21\linewidth}X@{}}
Genealogía & 16f46915\allowbreak{}102b6339\allowbreak{}647250e5\allowbreak{}2a7f69a3\allowbreak{}51f7ca67\allowbreak{}ea21ef1a\allowbreak{}4634567c\allowbreak{}1e82c249\\
Lector & 3e5eb22e\allowbreak{}721faa17\allowbreak{}e47903a4\allowbreak{}94789298\allowbreak{}f6831e31\allowbreak{}c48a1aca\allowbreak{}c42b56f6\allowbreak{}1a087c39\\
Registro en 108 & 976690fd\allowbreak{}5ec3fa4b\allowbreak{}f0ee2a75\allowbreak{}d8c8cf5e\allowbreak{}15e549ba\allowbreak{}4e8d2436\allowbreak{}57dff0e8\allowbreak{}cd9a4eed\\
\end{tabularx}\endgroup\par\vspace{6pt}
\Needspace{87pt}\phantomsection\label{trace-116}\noindent\hyperref[nav:vi-ruta-116]{\textbf{r4c2\_\allowbreak{}h+1v+1}}\quad Ordinal unificado / memoria: 116 / 113\par\nopagebreak
\begingroup\fontsize{8.3}{10.1}\selectfont\setlength{\tabcolsep}{4pt}\noindent\begin{tabularx}{\linewidth}{@{}p{.21\linewidth}X@{}}
Genealogía & 3108b4aa\allowbreak{}18a6e54b\allowbreak{}1b72c941\allowbreak{}27319da7\allowbreak{}c81a3c4b\allowbreak{}94e8d584\allowbreak{}2a621dea\allowbreak{}e470a1c1\\
Lector & a6f1b74f\allowbreak{}60b1279c\allowbreak{}43943235\allowbreak{}a599dd91\allowbreak{}eba83431\allowbreak{}12b3912e\allowbreak{}62c18ec7\allowbreak{}472ea819\\
Registro en 108 & \textemdash{}\\
\end{tabularx}\endgroup\par\vspace{6pt}
\Needspace{87pt}\phantomsection\label{trace-117}\noindent\hyperref[nav:vi-ruta-117]{\textbf{r4c3\_\allowbreak{}h-1v-1}}\quad Ordinal unificado / memoria: 117 / 120\par\nopagebreak
\begingroup\fontsize{8.3}{10.1}\selectfont\setlength{\tabcolsep}{4pt}\noindent\begin{tabularx}{\linewidth}{@{}p{.21\linewidth}X@{}}
Genealogía & dfba8889\allowbreak{}ea2e308d\allowbreak{}b35da1a1\allowbreak{}e3a470cb\allowbreak{}901f5b1b\allowbreak{}448eecfa\allowbreak{}8b072910\allowbreak{}f423cf30\\
Lector & afd36301\allowbreak{}52262dd7\allowbreak{}24d21cfe\allowbreak{}75259c89\allowbreak{}2dc76801\allowbreak{}5259d63a\allowbreak{}8c73a387\allowbreak{}e703340c\\
Registro en 108 & \textemdash{}\\
\end{tabularx}\endgroup\par\vspace{6pt}
\Needspace{87pt}\phantomsection\label{trace-118}\noindent\hyperref[nav:vi-ruta-118]{\textbf{r4c3\_\allowbreak{}h-1v+1}}\quad Ordinal unificado / memoria: 118 / 119\par\nopagebreak
\begingroup\fontsize{8.3}{10.1}\selectfont\setlength{\tabcolsep}{4pt}\noindent\begin{tabularx}{\linewidth}{@{}p{.21\linewidth}X@{}}
Genealogía & b9a0e830\allowbreak{}c5b79e19\allowbreak{}fac92114\allowbreak{}fb87e7d7\allowbreak{}3c89c8ab\allowbreak{}1412847c\allowbreak{}40d4194e\allowbreak{}dfafd295\\
Lector & 45d58d93\allowbreak{}81a000cc\allowbreak{}9d767cfa\allowbreak{}adb64e41\allowbreak{}2d11aefe\allowbreak{}df7a5b16\allowbreak{}fdf39ebb\allowbreak{}2427bb4f\\
Registro en 108 & \textemdash{}\\
\end{tabularx}\endgroup\par\vspace{6pt}
\Needspace{87pt}\phantomsection\label{trace-119}\noindent\hyperref[nav:vi-ruta-119]{\textbf{r4c3\_\allowbreak{}h+1v-1}}\quad Ordinal unificado / memoria: 119 / 118\par\nopagebreak
\begingroup\fontsize{8.3}{10.1}\selectfont\setlength{\tabcolsep}{4pt}\noindent\begin{tabularx}{\linewidth}{@{}p{.21\linewidth}X@{}}
Genealogía & f18a5f3b\allowbreak{}b67d88d2\allowbreak{}241efee9\allowbreak{}3e0869d0\allowbreak{}669b4af6\allowbreak{}13f1efc2\allowbreak{}a3340b82\allowbreak{}2993e9f0\\
Lector & 2fa63357\allowbreak{}47d64259\allowbreak{}eacd9cd2\allowbreak{}cfd7362a\allowbreak{}3d7b22ed\allowbreak{}15a290e4\allowbreak{}3fc0aa7d\allowbreak{}b9ebb8ee\\
Registro en 108 & \textemdash{}\\
\end{tabularx}\endgroup\par\vspace{6pt}
\Needspace{87pt}\phantomsection\label{trace-120}\noindent\hyperref[nav:vi-ruta-120]{\textbf{r4c3\_\allowbreak{}h+1v+1}}\quad Ordinal unificado / memoria: 120 / 117\par\nopagebreak
\begingroup\fontsize{8.3}{10.1}\selectfont\setlength{\tabcolsep}{4pt}\noindent\begin{tabularx}{\linewidth}{@{}p{.21\linewidth}X@{}}
Genealogía & d562cf51\allowbreak{}d4b94895\allowbreak{}d64814ea\allowbreak{}9d20c389\allowbreak{}4b082a0c\allowbreak{}3fdcbdd5\allowbreak{}29750f0f\allowbreak{}92dc6249\\
Lector & 46e7971f\allowbreak{}d3cde9a3\allowbreak{}330a90fd\allowbreak{}8d65631f\allowbreak{}09d3233d\allowbreak{}be7e2bec\allowbreak{}0108983f\allowbreak{}f0938898\\
Registro en 108 & 0445e186\allowbreak{}74117bf6\allowbreak{}22efd6fb\allowbreak{}fa1d6e4c\allowbreak{}6164b923\allowbreak{}8247bb7c\allowbreak{}762c152b\allowbreak{}b9686738\\
\end{tabularx}\endgroup\par\vspace{6pt}
\Needspace{87pt}\phantomsection\label{trace-121}\noindent\hyperref[nav:vi-ruta-121]{\textbf{r4c4\_\allowbreak{}h-1v-1}}\quad Ordinal unificado / memoria: 121 / 124\par\nopagebreak
\begingroup\fontsize{8.3}{10.1}\selectfont\setlength{\tabcolsep}{4pt}\noindent\begin{tabularx}{\linewidth}{@{}p{.21\linewidth}X@{}}
Genealogía & 77c77bed\allowbreak{}c6c70e8f\allowbreak{}673bba3f\allowbreak{}1a38ebbb\allowbreak{}ad0e9216\allowbreak{}da5a1135\allowbreak{}bf40b23f\allowbreak{}327bb557\\
Lector & 71932081\allowbreak{}a5e57d25\allowbreak{}037ce0d5\allowbreak{}b5db3847\allowbreak{}9d4de5b7\allowbreak{}9809ce8e\allowbreak{}2c0b6c8a\allowbreak{}a771cdfc\\
Registro en 108 & \textemdash{}\\
\end{tabularx}\endgroup\par\vspace{6pt}
\Needspace{87pt}\phantomsection\label{trace-122}\noindent\hyperref[nav:vi-ruta-122]{\textbf{r4c4\_\allowbreak{}h-1v+1}}\quad Ordinal unificado / memoria: 122 / 123\par\nopagebreak
\begingroup\fontsize{8.3}{10.1}\selectfont\setlength{\tabcolsep}{4pt}\noindent\begin{tabularx}{\linewidth}{@{}p{.21\linewidth}X@{}}
Genealogía & 287327ac\allowbreak{}0b7a6795\allowbreak{}1ed95a00\allowbreak{}dceb2cc4\allowbreak{}b73e9aae\allowbreak{}63f89cf3\allowbreak{}f375f2e1\allowbreak{}4608c87d\\
Lector & 015d0ecd\allowbreak{}c3d3681a\allowbreak{}13d3a6df\allowbreak{}b1164ab5\allowbreak{}5a41a27e\allowbreak{}6f103f07\allowbreak{}f68ac1e9\allowbreak{}4b8ffcd2\\
Registro en 108 & 282e5519\allowbreak{}2bd3edf9\allowbreak{}6e14e9bf\allowbreak{}0b50e347\allowbreak{}9922204d\allowbreak{}b1e0fa80\allowbreak{}756865f1\allowbreak{}31f8c5ed\\
\end{tabularx}\endgroup\par\vspace{6pt}
\Needspace{87pt}\phantomsection\label{trace-123}\noindent\hyperref[nav:vi-ruta-123]{\textbf{r4c4\_\allowbreak{}h+1v-1}}\quad Ordinal unificado / memoria: 123 / 122\par\nopagebreak
\begingroup\fontsize{8.3}{10.1}\selectfont\setlength{\tabcolsep}{4pt}\noindent\begin{tabularx}{\linewidth}{@{}p{.21\linewidth}X@{}}
Genealogía & b39d7e16\allowbreak{}04512dee\allowbreak{}f70b3918\allowbreak{}e486fa58\allowbreak{}fffe93c5\allowbreak{}ba2477a0\allowbreak{}9df55b64\allowbreak{}831ca74b\\
Lector & 8f35c9aa\allowbreak{}a9f1b198\allowbreak{}75c81199\allowbreak{}d6483238\allowbreak{}76ea87ed\allowbreak{}664f20cd\allowbreak{}b441f73a\allowbreak{}dcdccb12\\
Registro en 108 & \textemdash{}\\
\end{tabularx}\endgroup\par\vspace{6pt}
\Needspace{87pt}\phantomsection\label{trace-124}\noindent\hyperref[nav:vi-ruta-124]{\textbf{r4c4\_\allowbreak{}h+1v+1}}\quad Ordinal unificado / memoria: 124 / 121\par\nopagebreak
\begingroup\fontsize{8.3}{10.1}\selectfont\setlength{\tabcolsep}{4pt}\noindent\begin{tabularx}{\linewidth}{@{}p{.21\linewidth}X@{}}
Genealogía & a20fae0d\allowbreak{}e3205269\allowbreak{}14318b14\allowbreak{}32bf0b0e\allowbreak{}b2e61de2\allowbreak{}2cd563a0\allowbreak{}e5a17b0e\allowbreak{}a9128bc0\\
Lector & 57d09c9e\allowbreak{}bfd5b2a4\allowbreak{}289c6b10\allowbreak{}55b0fc67\allowbreak{}ad75d3af\allowbreak{}3cdcf8d8\allowbreak{}4d685d0e\allowbreak{}8828573d\\
Registro en 108 & \textemdash{}\\
\end{tabularx}\endgroup\par\vspace{6pt}
\Needspace{87pt}\phantomsection\label{trace-125}\noindent\hyperref[nav:vi-ruta-125]{\textbf{r4c5\_\allowbreak{}h-1v-1}}\quad Ordinal unificado / memoria: 125 / 128\par\nopagebreak
\begingroup\fontsize{8.3}{10.1}\selectfont\setlength{\tabcolsep}{4pt}\noindent\begin{tabularx}{\linewidth}{@{}p{.21\linewidth}X@{}}
Genealogía & 79fb9a84\allowbreak{}5241446b\allowbreak{}08dd24c6\allowbreak{}05f6b1b7\allowbreak{}bd4981c5\allowbreak{}e6fa6e00\allowbreak{}67ea36ce\allowbreak{}0e7baef1\\
Lector & fb4bcab3\allowbreak{}9bf089e5\allowbreak{}c421bf35\allowbreak{}2caad667\allowbreak{}4c9c54f0\allowbreak{}4a9d9549\allowbreak{}4401d5b5\allowbreak{}db852e28\\
Registro en 108 & \textemdash{}\\
\end{tabularx}\endgroup\par\vspace{6pt}
\Needspace{87pt}\phantomsection\label{trace-126}\noindent\hyperref[nav:vi-ruta-126]{\textbf{r4c5\_\allowbreak{}h-1v+1}}\quad Ordinal unificado / memoria: 126 / 127\par\nopagebreak
\begingroup\fontsize{8.3}{10.1}\selectfont\setlength{\tabcolsep}{4pt}\noindent\begin{tabularx}{\linewidth}{@{}p{.21\linewidth}X@{}}
Genealogía & 0c006090\allowbreak{}1278e986\allowbreak{}f81b5b52\allowbreak{}2e367331\allowbreak{}377866c4\allowbreak{}eec0a6c9\allowbreak{}a5fd33b1\allowbreak{}702579f7\\
Lector & 7359aaa7\allowbreak{}14e722cf\allowbreak{}55306e73\allowbreak{}89c0a782\allowbreak{}1640a976\allowbreak{}ddca2652\allowbreak{}c3e9de16\allowbreak{}1d84c014\\
Registro en 108 & \textemdash{}\\
\end{tabularx}\endgroup\par\vspace{6pt}
\Needspace{87pt}\phantomsection\label{trace-127}\noindent\hyperref[nav:vi-ruta-127]{\textbf{r4c5\_\allowbreak{}h+1v-1}}\quad Ordinal unificado / memoria: 127 / 126\par\nopagebreak
\begingroup\fontsize{8.3}{10.1}\selectfont\setlength{\tabcolsep}{4pt}\noindent\begin{tabularx}{\linewidth}{@{}p{.21\linewidth}X@{}}
Genealogía & 9866ca01\allowbreak{}6d61a93c\allowbreak{}8605f9d3\allowbreak{}93f5c7d1\allowbreak{}ccbf95b3\allowbreak{}9e0be07e\allowbreak{}da5205d1\allowbreak{}77bbe0a3\\
Lector & 37ed6a79\allowbreak{}26a4b489\allowbreak{}b076fae2\allowbreak{}ae21dc04\allowbreak{}dd4b9ca8\allowbreak{}e9edd354\allowbreak{}076bf251\allowbreak{}fcd73744\\
Registro en 108 & 58bd1091\allowbreak{}18dc4021\allowbreak{}43446db4\allowbreak{}3b056afb\allowbreak{}cc8d2a59\allowbreak{}91cace6f\allowbreak{}6ec3fb46\allowbreak{}70b5eff4\\
\end{tabularx}\endgroup\par\vspace{6pt}
\Needspace{87pt}\phantomsection\label{trace-128}\noindent\hyperref[nav:vi-ruta-128]{\textbf{r4c5\_\allowbreak{}h+1v+1}}\quad Ordinal unificado / memoria: 128 / 125\par\nopagebreak
\begingroup\fontsize{8.3}{10.1}\selectfont\setlength{\tabcolsep}{4pt}\noindent\begin{tabularx}{\linewidth}{@{}p{.21\linewidth}X@{}}
Genealogía & 4fcdd9bf\allowbreak{}c143fe82\allowbreak{}bee933b8\allowbreak{}25b84006\allowbreak{}0bc9cc4c\allowbreak{}b3164418\allowbreak{}cc30da5c\allowbreak{}55f863c5\\
Lector & ba031ff9\allowbreak{}02fef28f\allowbreak{}8cc55107\allowbreak{}fa551870\allowbreak{}981ae35f\allowbreak{}912431d3\allowbreak{}e4c6b06e\allowbreak{}9fd726fa\\
Registro en 108 & \textemdash{}\\
\end{tabularx}\endgroup\par\vspace{6pt}
\Needspace{87pt}\phantomsection\label{trace-129}\noindent\hyperref[nav:vi-ruta-129]{\textbf{r4c6\_\allowbreak{}h-1v-1}}\quad Ordinal unificado / memoria: 129 / 132\par\nopagebreak
\begingroup\fontsize{8.3}{10.1}\selectfont\setlength{\tabcolsep}{4pt}\noindent\begin{tabularx}{\linewidth}{@{}p{.21\linewidth}X@{}}
Genealogía & 87efef61\allowbreak{}3d4e87ec\allowbreak{}97c29419\allowbreak{}f7de6fbe\allowbreak{}a09a79bc\allowbreak{}753ae3e4\allowbreak{}04b221d8\allowbreak{}ede75b4a\\
Lector & 6e6b7dd4\allowbreak{}14d89eb1\allowbreak{}aaaf5a8e\allowbreak{}a14fac78\allowbreak{}61011753\allowbreak{}e2349a9c\allowbreak{}c217300f\allowbreak{}d1b80edf\\
Registro en 108 & \textemdash{}\\
\end{tabularx}\endgroup\par\vspace{6pt}
\Needspace{87pt}\phantomsection\label{trace-130}\noindent\hyperref[nav:vi-ruta-130]{\textbf{r4c6\_\allowbreak{}h-1v+1}}\quad Ordinal unificado / memoria: 130 / 131\par\nopagebreak
\begingroup\fontsize{8.3}{10.1}\selectfont\setlength{\tabcolsep}{4pt}\noindent\begin{tabularx}{\linewidth}{@{}p{.21\linewidth}X@{}}
Genealogía & 055c55a6\allowbreak{}583303aa\allowbreak{}c67730b6\allowbreak{}2377ba8a\allowbreak{}1d0049ac\allowbreak{}7aed4deb\allowbreak{}c98aad58\allowbreak{}81bc9bb5\\
Lector & 8fa03549\allowbreak{}920be213\allowbreak{}176b7e0c\allowbreak{}881c6482\allowbreak{}be1ed34c\allowbreak{}cd91ebfb\allowbreak{}07601982\allowbreak{}97de4e2b\\
Registro en 108 & \textemdash{}\\
\end{tabularx}\endgroup\par\vspace{6pt}
\Needspace{87pt}\phantomsection\label{trace-131}\noindent\hyperref[nav:vi-ruta-131]{\textbf{r4c6\_\allowbreak{}h+1v-1}}\quad Ordinal unificado / memoria: 131 / 130\par\nopagebreak
\begingroup\fontsize{8.3}{10.1}\selectfont\setlength{\tabcolsep}{4pt}\noindent\begin{tabularx}{\linewidth}{@{}p{.21\linewidth}X@{}}
Genealogía & 406e5327\allowbreak{}b2546eaa\allowbreak{}3992ae56\allowbreak{}c6688e01\allowbreak{}4239f25e\allowbreak{}f3712139\allowbreak{}e37360e7\allowbreak{}0307a0cf\\
Lector & c31d14fc\allowbreak{}f53e490e\allowbreak{}7b58ce90\allowbreak{}33002447\allowbreak{}67c13da2\allowbreak{}0e178665\allowbreak{}54a53b7f\allowbreak{}1e8aa2ca\\
Registro en 108 & \textemdash{}\\
\end{tabularx}\endgroup\par\vspace{6pt}
\Needspace{87pt}\phantomsection\label{trace-132}\noindent\hyperref[nav:vi-ruta-132]{\textbf{r4c6\_\allowbreak{}h+1v+1}}\quad Ordinal unificado / memoria: 132 / 129\par\nopagebreak
\begingroup\fontsize{8.3}{10.1}\selectfont\setlength{\tabcolsep}{4pt}\noindent\begin{tabularx}{\linewidth}{@{}p{.21\linewidth}X@{}}
Genealogía & edc42546\allowbreak{}30dee8ee\allowbreak{}d6f09c01\allowbreak{}001d7e5f\allowbreak{}f9ae88a5\allowbreak{}2b26f6c8\allowbreak{}4b34889e\allowbreak{}95b4b44f\\
Lector & 8e63a835\allowbreak{}c6a45d8c\allowbreak{}dfe6ec76\allowbreak{}15e7954f\allowbreak{}6d06dba1\allowbreak{}53209bd7\allowbreak{}39e0e946\allowbreak{}65079019\\
Registro en 108 & ac355fc3\allowbreak{}c7d67e23\allowbreak{}747af693\allowbreak{}446ce597\allowbreak{}3090a128\allowbreak{}8c89b61e\allowbreak{}d5e84c16\allowbreak{}4421d478\\
\end{tabularx}\endgroup\par\vspace{6pt}
\Needspace{87pt}\phantomsection\label{trace-133}\noindent\hyperref[nav:vi-ruta-133]{\textbf{r4c7\_\allowbreak{}h-1v-1}}\quad Ordinal unificado / memoria: 133 / 136\par\nopagebreak
\begingroup\fontsize{8.3}{10.1}\selectfont\setlength{\tabcolsep}{4pt}\noindent\begin{tabularx}{\linewidth}{@{}p{.21\linewidth}X@{}}
Genealogía & 4ca34566\allowbreak{}97b5beaf\allowbreak{}b2684f34\allowbreak{}6ef4fe8a\allowbreak{}825e4e9b\allowbreak{}1923a035\allowbreak{}4b9e170e\allowbreak{}0a544faa\\
Lector & e241f838\allowbreak{}e36aaa93\allowbreak{}1f974504\allowbreak{}e5963e77\allowbreak{}4300dbe4\allowbreak{}095f5408\allowbreak{}12bdb1be\allowbreak{}522d47ca\\
Registro en 108 & \textemdash{}\\
\end{tabularx}\endgroup\par\vspace{6pt}
\Needspace{87pt}\phantomsection\label{trace-134}\noindent\hyperref[nav:vi-ruta-134]{\textbf{r4c7\_\allowbreak{}h-1v+1}}\quad Ordinal unificado / memoria: 134 / 135\par\nopagebreak
\begingroup\fontsize{8.3}{10.1}\selectfont\setlength{\tabcolsep}{4pt}\noindent\begin{tabularx}{\linewidth}{@{}p{.21\linewidth}X@{}}
Genealogía & 01c95731\allowbreak{}317c519b\allowbreak{}15d30504\allowbreak{}737dbf07\allowbreak{}55697c53\allowbreak{}0400d12b\allowbreak{}6dbb06b6\allowbreak{}fddb26df\\
Lector & 0775e69a\allowbreak{}a78c62ae\allowbreak{}b37e0d80\allowbreak{}3ae05fdb\allowbreak{}216bf35b\allowbreak{}3648f5cb\allowbreak{}d17c9271\allowbreak{}9739b065\\
Registro en 108 & 4acc79e1\allowbreak{}bddf0a84\allowbreak{}46aca263\allowbreak{}fca69bd3\allowbreak{}9879caa1\allowbreak{}e6469c33\allowbreak{}d703da91\allowbreak{}72c27cd9\\
\end{tabularx}\endgroup\par\vspace{6pt}
\Needspace{87pt}\phantomsection\label{trace-135}\noindent\hyperref[nav:vi-ruta-135]{\textbf{r4c7\_\allowbreak{}h+1v-1}}\quad Ordinal unificado / memoria: 135 / 134\par\nopagebreak
\begingroup\fontsize{8.3}{10.1}\selectfont\setlength{\tabcolsep}{4pt}\noindent\begin{tabularx}{\linewidth}{@{}p{.21\linewidth}X@{}}
Genealogía & 782a318c\allowbreak{}c548e949\allowbreak{}b83af905\allowbreak{}97ea3aa0\allowbreak{}715c06a5\allowbreak{}7efc8cd4\allowbreak{}c09cc84e\allowbreak{}af42ad4c\\
Lector & f4d3eea4\allowbreak{}3826eed3\allowbreak{}747c080d\allowbreak{}0cdd48c8\allowbreak{}09877dc7\allowbreak{}fc409585\allowbreak{}fd0274bc\allowbreak{}331de472\\
Registro en 108 & \textemdash{}\\
\end{tabularx}\endgroup\par\vspace{6pt}
\Needspace{87pt}\phantomsection\label{trace-136}\noindent\hyperref[nav:vi-ruta-136]{\textbf{r4c7\_\allowbreak{}h+1v+1}}\quad Ordinal unificado / memoria: 136 / 133\par\nopagebreak
\begingroup\fontsize{8.3}{10.1}\selectfont\setlength{\tabcolsep}{4pt}\noindent\begin{tabularx}{\linewidth}{@{}p{.21\linewidth}X@{}}
Genealogía & 54436c13\allowbreak{}07ec79e7\allowbreak{}6d82397f\allowbreak{}2b17691b\allowbreak{}a5e8eca0\allowbreak{}4e1f5c11\allowbreak{}35c5a6f0\allowbreak{}7fdd51cf\\
Lector & 073fbdcd\allowbreak{}b0df31cb\allowbreak{}cc689c53\allowbreak{}9e57777a\allowbreak{}6ea89a07\allowbreak{}71eb0723\allowbreak{}8bc113fd\allowbreak{}3c52f129\\
Registro en 108 & \textemdash{}\\
\end{tabularx}\endgroup\par\vspace{6pt}
\Needspace{87pt}\phantomsection\label{trace-137}\noindent\hyperref[nav:vi-ruta-137]{\textbf{r4c8\_\allowbreak{}h-1v-1}}\quad Ordinal unificado / memoria: 137 / 140\par\nopagebreak
\begingroup\fontsize{8.3}{10.1}\selectfont\setlength{\tabcolsep}{4pt}\noindent\begin{tabularx}{\linewidth}{@{}p{.21\linewidth}X@{}}
Genealogía & 2b804e0c\allowbreak{}fd9c913a\allowbreak{}a79e7b25\allowbreak{}29b115c9\allowbreak{}1f296d3c\allowbreak{}f84b6837\allowbreak{}e5e05321\allowbreak{}bfc2a6c0\\
Lector & 98e82256\allowbreak{}254f483d\allowbreak{}b3a5842d\allowbreak{}1f7fd6a5\allowbreak{}69044ecf\allowbreak{}3c542b94\allowbreak{}17058555\allowbreak{}59737516\\
Registro en 108 & bf56e7e7\allowbreak{}a4414205\allowbreak{}3df74d41\allowbreak{}b3498a6c\allowbreak{}84055d81\allowbreak{}2f97946e\allowbreak{}d2c06740\allowbreak{}99511566\\
\end{tabularx}\endgroup\par\vspace{6pt}
\Needspace{87pt}\phantomsection\label{trace-138}\noindent\hyperref[nav:vi-ruta-138]{\textbf{r4c8\_\allowbreak{}h-1v+1}}\quad Ordinal unificado / memoria: 138 / 139\par\nopagebreak
\begingroup\fontsize{8.3}{10.1}\selectfont\setlength{\tabcolsep}{4pt}\noindent\begin{tabularx}{\linewidth}{@{}p{.21\linewidth}X@{}}
Genealogía & 03804c45\allowbreak{}4b95d8b9\allowbreak{}3e9f2789\allowbreak{}69bffd8b\allowbreak{}f62a8848\allowbreak{}88eefa4f\allowbreak{}44c9253a\allowbreak{}fb6c9aa4\\
Lector & 87681b4a\allowbreak{}0c073f38\allowbreak{}f2a84e8a\allowbreak{}65f543f2\allowbreak{}c195b938\allowbreak{}cd71606b\allowbreak{}35a72691\allowbreak{}d777d3ce\\
Registro en 108 & \textemdash{}\\
\end{tabularx}\endgroup\par\vspace{6pt}
\Needspace{87pt}\phantomsection\label{trace-139}\noindent\hyperref[nav:vi-ruta-139]{\textbf{r4c8\_\allowbreak{}h+1v-1}}\quad Ordinal unificado / memoria: 139 / 138\par\nopagebreak
\begingroup\fontsize{8.3}{10.1}\selectfont\setlength{\tabcolsep}{4pt}\noindent\begin{tabularx}{\linewidth}{@{}p{.21\linewidth}X@{}}
Genealogía & b6ebfc86\allowbreak{}18552fdd\allowbreak{}72e29c86\allowbreak{}168f9f6e\allowbreak{}2a46c3fe\allowbreak{}8d74fb9e\allowbreak{}06ff8d9e\allowbreak{}f4d8355b\\
Lector & e8e9508b\allowbreak{}84066283\allowbreak{}79069d3f\allowbreak{}617b9bea\allowbreak{}ee92634d\allowbreak{}170e0a1c\allowbreak{}f2a20c31\allowbreak{}d62a42e6\\
Registro en 108 & \textemdash{}\\
\end{tabularx}\endgroup\par\vspace{6pt}
\Needspace{87pt}\phantomsection\label{trace-140}\noindent\hyperref[nav:vi-ruta-140]{\textbf{r4c8\_\allowbreak{}h+1v+1}}\quad Ordinal unificado / memoria: 140 / 137\par\nopagebreak
\begingroup\fontsize{8.3}{10.1}\selectfont\setlength{\tabcolsep}{4pt}\noindent\begin{tabularx}{\linewidth}{@{}p{.21\linewidth}X@{}}
Genealogía & b1167825\allowbreak{}dc88e346\allowbreak{}b1033aa3\allowbreak{}68e2c68a\allowbreak{}7ce87d85\allowbreak{}4b37460a\allowbreak{}00c312e0\allowbreak{}b818fce5\\
Lector & b7f40bf3\allowbreak{}d4a01f15\allowbreak{}bdaa23dc\allowbreak{}2ebd2cdd\allowbreak{}551eca43\allowbreak{}991086c5\allowbreak{}2426a372\allowbreak{}930a411b\\
Registro en 108 & \textemdash{}\\
\end{tabularx}\endgroup\par\vspace{6pt}
\Needspace{87pt}\phantomsection\label{trace-141}\noindent\hyperref[nav:vi-ruta-141]{\textbf{r4c9\_\allowbreak{}h-1v-1}}\quad Ordinal unificado / memoria: 141 / 144\par\nopagebreak
\begingroup\fontsize{8.3}{10.1}\selectfont\setlength{\tabcolsep}{4pt}\noindent\begin{tabularx}{\linewidth}{@{}p{.21\linewidth}X@{}}
Genealogía & 1fac69da\allowbreak{}fa70b079\allowbreak{}93139f7a\allowbreak{}846f20ab\allowbreak{}b465238f\allowbreak{}4558212f\allowbreak{}54421745\allowbreak{}cb5e3918\\
Lector & f8d6c21d\allowbreak{}ad31131f\allowbreak{}c4c2fc3d\allowbreak{}70185a1f\allowbreak{}888c237d\allowbreak{}6474a021\allowbreak{}870b6519\allowbreak{}6e6d416e\\
Registro en 108 & \textemdash{}\\
\end{tabularx}\endgroup\par\vspace{6pt}
\Needspace{87pt}\phantomsection\label{trace-142}\noindent\hyperref[nav:vi-ruta-142]{\textbf{r4c9\_\allowbreak{}h-1v+1}}\quad Ordinal unificado / memoria: 142 / 143\par\nopagebreak
\begingroup\fontsize{8.3}{10.1}\selectfont\setlength{\tabcolsep}{4pt}\noindent\begin{tabularx}{\linewidth}{@{}p{.21\linewidth}X@{}}
Genealogía & 18959d06\allowbreak{}bbedaac3\allowbreak{}d5d11329\allowbreak{}81d981df\allowbreak{}9bd28a79\allowbreak{}c6da3375\allowbreak{}6e8d2149\allowbreak{}cb51db61\\
Lector & 4c038f82\allowbreak{}1f5bd14c\allowbreak{}eca5c1f4\allowbreak{}829b72b3\allowbreak{}72c1563b\allowbreak{}32e0fcc5\allowbreak{}b34cfb9e\allowbreak{}ff6992ee\\
Registro en 108 & 1e448780\allowbreak{}f5a741d6\allowbreak{}641ffc8e\allowbreak{}2ae465ac\allowbreak{}249cdecb\allowbreak{}bd7f941e\allowbreak{}44db167f\allowbreak{}2a18cb1b\\
\end{tabularx}\endgroup\par\vspace{6pt}
\Needspace{87pt}\phantomsection\label{trace-143}\noindent\hyperref[nav:vi-ruta-143]{\textbf{r4c9\_\allowbreak{}h+1v-1}}\quad Ordinal unificado / memoria: 143 / 142\par\nopagebreak
\begingroup\fontsize{8.3}{10.1}\selectfont\setlength{\tabcolsep}{4pt}\noindent\begin{tabularx}{\linewidth}{@{}p{.21\linewidth}X@{}}
Genealogía & e624c172\allowbreak{}e0c76c77\allowbreak{}cc5995ae\allowbreak{}25a89b2e\allowbreak{}951081d6\allowbreak{}31056741\allowbreak{}95861438\allowbreak{}08bab0aa\\
Lector & b10d5887\allowbreak{}020993ff\allowbreak{}069a374c\allowbreak{}70846c7d\allowbreak{}b2da5382\allowbreak{}0e650c37\allowbreak{}b95231a6\allowbreak{}b241a3cb\\
Registro en 108 & \textemdash{}\\
\end{tabularx}\endgroup\par\vspace{6pt}
\Needspace{87pt}\phantomsection\label{trace-144}\noindent\hyperref[nav:vi-ruta-144]{\textbf{r4c9\_\allowbreak{}h+1v+1}}\quad Ordinal unificado / memoria: 144 / 141\par\nopagebreak
\begingroup\fontsize{8.3}{10.1}\selectfont\setlength{\tabcolsep}{4pt}\noindent\begin{tabularx}{\linewidth}{@{}p{.21\linewidth}X@{}}
Genealogía & 555a0b57\allowbreak{}737e2212\allowbreak{}2d0fd734\allowbreak{}f5aa5516\allowbreak{}d6461c18\allowbreak{}5bde5c94\allowbreak{}fab70620\allowbreak{}44c81b31\\
Lector & 40c90b30\allowbreak{}41a8d94c\allowbreak{}a52a8953\allowbreak{}35732ae5\allowbreak{}1dc7c7f9\allowbreak{}837d5ada\allowbreak{}66d7bc12\allowbreak{}643ede6a\\
Registro en 108 & \textemdash{}\\
\end{tabularx}\endgroup\par\vspace{6pt}
\Needspace{87pt}\phantomsection\label{trace-145}\noindent\hyperref[nav:vi-ruta-145]{\textbf{r5c1\_\allowbreak{}h-1v-1}}\quad Ordinal unificado / memoria: 145 / 148\par\nopagebreak
\begingroup\fontsize{8.3}{10.1}\selectfont\setlength{\tabcolsep}{4pt}\noindent\begin{tabularx}{\linewidth}{@{}p{.21\linewidth}X@{}}
Genealogía & 44128f77\allowbreak{}694e22cd\allowbreak{}b74ce801\allowbreak{}d17ef58c\allowbreak{}b5f75058\allowbreak{}dd488325\allowbreak{}f94b8db3\allowbreak{}aef4c2f3\\
Lector & 8e414ca1\allowbreak{}b56fa4e2\allowbreak{}e54dc83e\allowbreak{}480cd2bd\allowbreak{}b5cfc205\allowbreak{}2546af04\allowbreak{}082e0f01\allowbreak{}7d24c4e5\\
Registro en 108 & \textemdash{}\\
\end{tabularx}\endgroup\par\vspace{6pt}
\Needspace{87pt}\phantomsection\label{trace-146}\noindent\hyperref[nav:vi-ruta-146]{\textbf{r5c1\_\allowbreak{}h-1v+1}}\quad Ordinal unificado / memoria: 146 / 147\par\nopagebreak
\begingroup\fontsize{8.3}{10.1}\selectfont\setlength{\tabcolsep}{4pt}\noindent\begin{tabularx}{\linewidth}{@{}p{.21\linewidth}X@{}}
Genealogía & 0f38c695\allowbreak{}8afd01ba\allowbreak{}27b6bc49\allowbreak{}a6893e48\allowbreak{}53cbeff8\allowbreak{}b05159d5\allowbreak{}12fb8b00\allowbreak{}bcb73007\\
Lector & d9b33583\allowbreak{}8b663c21\allowbreak{}1a4ec0d1\allowbreak{}0a405d91\allowbreak{}7b276d2c\allowbreak{}a4cd5695\allowbreak{}fbf07de3\allowbreak{}d36643e7\\
Registro en 108 & \textemdash{}\\
\end{tabularx}\endgroup\par\vspace{6pt}
\Needspace{87pt}\phantomsection\label{trace-147}\noindent\hyperref[nav:vi-ruta-147]{\textbf{r5c1\_\allowbreak{}h+1v-1}}\quad Ordinal unificado / memoria: 147 / 146\par\nopagebreak
\begingroup\fontsize{8.3}{10.1}\selectfont\setlength{\tabcolsep}{4pt}\noindent\begin{tabularx}{\linewidth}{@{}p{.21\linewidth}X@{}}
Genealogía & aaffc169\allowbreak{}d0598de5\allowbreak{}97a3a288\allowbreak{}2b457981\allowbreak{}816b63fb\allowbreak{}5a1853f5\allowbreak{}546db1a9\allowbreak{}85aea4ba\\
Lector & deaf6877\allowbreak{}f8a7c888\allowbreak{}0ba1417f\allowbreak{}e43eee1c\allowbreak{}66d2cde0\allowbreak{}af9e5961\allowbreak{}bbf415f2\allowbreak{}f0787f8e\\
Registro en 108 & ff9f223f\allowbreak{}aafbe5f7\allowbreak{}d396cb50\allowbreak{}240d6d52\allowbreak{}e53eaab1\allowbreak{}a745f771\allowbreak{}9f6f1503\allowbreak{}b398f6a5\\
\end{tabularx}\endgroup\par\vspace{6pt}
\Needspace{87pt}\phantomsection\label{trace-148}\noindent\hyperref[nav:vi-ruta-148]{\textbf{r5c1\_\allowbreak{}h+1v+1}}\quad Ordinal unificado / memoria: 148 / 145\par\nopagebreak
\begingroup\fontsize{8.3}{10.1}\selectfont\setlength{\tabcolsep}{4pt}\noindent\begin{tabularx}{\linewidth}{@{}p{.21\linewidth}X@{}}
Genealogía & 6bb6953e\allowbreak{}e6341f42\allowbreak{}515e8490\allowbreak{}c33965ad\allowbreak{}839bbc94\allowbreak{}6cbc777d\allowbreak{}d1d83224\allowbreak{}8ba36e66\\
Lector & 9743965c\allowbreak{}5f547f81\allowbreak{}a6eb1b5b\allowbreak{}fc13cd5e\allowbreak{}e9140eeb\allowbreak{}402f7500\allowbreak{}f7a52149\allowbreak{}283e0a47\\
Registro en 108 & \textemdash{}\\
\end{tabularx}\endgroup\par\vspace{6pt}
\Needspace{87pt}\phantomsection\label{trace-149}\noindent\hyperref[nav:vi-ruta-149]{\textbf{r5c2\_\allowbreak{}h-1v-1}}\quad Ordinal unificado / memoria: 149 / 152\par\nopagebreak
\begingroup\fontsize{8.3}{10.1}\selectfont\setlength{\tabcolsep}{4pt}\noindent\begin{tabularx}{\linewidth}{@{}p{.21\linewidth}X@{}}
Genealogía & 7cfcd177\allowbreak{}6e4330be\allowbreak{}5d0164d0\allowbreak{}74d9db7e\allowbreak{}bbc9ed42\allowbreak{}1fdd9625\allowbreak{}dbdae32b\allowbreak{}33ae0ef6\\
Lector & b369d9f2\allowbreak{}359a2a99\allowbreak{}a416d16f\allowbreak{}936b4222\allowbreak{}912fe42a\allowbreak{}a504fb2c\allowbreak{}a9cbc185\allowbreak{}be1d8c70\\
Registro en 108 & 8754d795\allowbreak{}eb4696cd\allowbreak{}b9a16e28\allowbreak{}9f13d7d0\allowbreak{}caa88f17\allowbreak{}9c1ebb00\allowbreak{}9c21bcc2\allowbreak{}36cb922c\\
\end{tabularx}\endgroup\par\vspace{6pt}
\Needspace{87pt}\phantomsection\label{trace-150}\noindent\hyperref[nav:vi-ruta-150]{\textbf{r5c2\_\allowbreak{}h-1v+1}}\quad Ordinal unificado / memoria: 150 / 151\par\nopagebreak
\begingroup\fontsize{8.3}{10.1}\selectfont\setlength{\tabcolsep}{4pt}\noindent\begin{tabularx}{\linewidth}{@{}p{.21\linewidth}X@{}}
Genealogía & ce452c82\allowbreak{}855b0221\allowbreak{}b7ba9f0c\allowbreak{}035bcb73\allowbreak{}09bb6bcd\allowbreak{}57b4db0e\allowbreak{}3abe8614\allowbreak{}8d83512e\\
Lector & b1388ecf\allowbreak{}bf1e3638\allowbreak{}e950a7e9\allowbreak{}0fed725e\allowbreak{}51d428bc\allowbreak{}70bb666d\allowbreak{}2eba9ccc\allowbreak{}3a4e8599\\
Registro en 108 & \textemdash{}\\
\end{tabularx}\endgroup\par\vspace{6pt}
\Needspace{87pt}\phantomsection\label{trace-151}\noindent\hyperref[nav:vi-ruta-151]{\textbf{r5c2\_\allowbreak{}h+1v-1}}\quad Ordinal unificado / memoria: 151 / 150\par\nopagebreak
\begingroup\fontsize{8.3}{10.1}\selectfont\setlength{\tabcolsep}{4pt}\noindent\begin{tabularx}{\linewidth}{@{}p{.21\linewidth}X@{}}
Genealogía & 2e1dadca\allowbreak{}d246fee2\allowbreak{}81f87bf1\allowbreak{}35901e69\allowbreak{}5efb2a4a\allowbreak{}05df738f\allowbreak{}6c070024\allowbreak{}07a0b1f2\\
Lector & 184e6040\allowbreak{}f83f204c\allowbreak{}85f17697\allowbreak{}1d97ea90\allowbreak{}40bf8f92\allowbreak{}63d3f91b\allowbreak{}694e0a6d\allowbreak{}3fdcf152\\
Registro en 108 & \textemdash{}\\
\end{tabularx}\endgroup\par\vspace{6pt}
\Needspace{87pt}\phantomsection\label{trace-152}\noindent\hyperref[nav:vi-ruta-152]{\textbf{r5c2\_\allowbreak{}h+1v+1}}\quad Ordinal unificado / memoria: 152 / 149\par\nopagebreak
\begingroup\fontsize{8.3}{10.1}\selectfont\setlength{\tabcolsep}{4pt}\noindent\begin{tabularx}{\linewidth}{@{}p{.21\linewidth}X@{}}
Genealogía & b277f51b\allowbreak{}0e0d2ae2\allowbreak{}24e9987f\allowbreak{}04a99ed8\allowbreak{}8a84fd08\allowbreak{}1c5d0ef6\allowbreak{}662e6eb3\allowbreak{}2b5cde7c\\
Lector & 664bc723\allowbreak{}29b1dfa4\allowbreak{}61d6b871\allowbreak{}16731273\allowbreak{}82f8d4cf\allowbreak{}08e87ae3\allowbreak{}a09df627\allowbreak{}4e4b9cb3\\
Registro en 108 & \textemdash{}\\
\end{tabularx}\endgroup\par\vspace{6pt}
\Needspace{87pt}\phantomsection\label{trace-153}\noindent\hyperref[nav:vi-ruta-153]{\textbf{r5c3\_\allowbreak{}h-1v-1}}\quad Ordinal unificado / memoria: 153 / 156\par\nopagebreak
\begingroup\fontsize{8.3}{10.1}\selectfont\setlength{\tabcolsep}{4pt}\noindent\begin{tabularx}{\linewidth}{@{}p{.21\linewidth}X@{}}
Genealogía & 628aa132\allowbreak{}9a1d669f\allowbreak{}1d699ac5\allowbreak{}9afb3d24\allowbreak{}a5b55af8\allowbreak{}129e6ebb\allowbreak{}54860dbb\allowbreak{}0a921a82\\
Lector & 8d0e0a7c\allowbreak{}9c083d59\allowbreak{}bd7752c1\allowbreak{}2edada42\allowbreak{}83a1227f\allowbreak{}99644dd9\allowbreak{}9fd52362\allowbreak{}03276d59\\
Registro en 108 & \textemdash{}\\
\end{tabularx}\endgroup\par\vspace{6pt}
\Needspace{87pt}\phantomsection\label{trace-154}\noindent\hyperref[nav:vi-ruta-154]{\textbf{r5c3\_\allowbreak{}h-1v+1}}\quad Ordinal unificado / memoria: 154 / 155\par\nopagebreak
\begingroup\fontsize{8.3}{10.1}\selectfont\setlength{\tabcolsep}{4pt}\noindent\begin{tabularx}{\linewidth}{@{}p{.21\linewidth}X@{}}
Genealogía & e1edda89\allowbreak{}7425f41c\allowbreak{}7b0d4061\allowbreak{}df3db1cc\allowbreak{}7e64b9f2\allowbreak{}8d3db5f3\allowbreak{}0063789b\allowbreak{}0434a10d\\
Lector & 3dc097d1\allowbreak{}af9e2772\allowbreak{}59c79d8c\allowbreak{}f836e831\allowbreak{}c4250729\allowbreak{}50eb258b\allowbreak{}5eadfec9\allowbreak{}0704df5e\\
Registro en 108 & f686fd28\allowbreak{}35550442\allowbreak{}1f5166a6\allowbreak{}7fe65217\allowbreak{}47f0e1dc\allowbreak{}54b72eb3\allowbreak{}9e0a7a80\allowbreak{}b67ba3f9\\
\end{tabularx}\endgroup\par\vspace{6pt}
\Needspace{87pt}\phantomsection\label{trace-155}\noindent\hyperref[nav:vi-ruta-155]{\textbf{r5c3\_\allowbreak{}h+1v-1}}\quad Ordinal unificado / memoria: 155 / 154\par\nopagebreak
\begingroup\fontsize{8.3}{10.1}\selectfont\setlength{\tabcolsep}{4pt}\noindent\begin{tabularx}{\linewidth}{@{}p{.21\linewidth}X@{}}
Genealogía & 5ebf99eb\allowbreak{}3286cd3e\allowbreak{}668d1881\allowbreak{}ac5a6d25\allowbreak{}c3fed231\allowbreak{}5f47a970\allowbreak{}c3ffd75b\allowbreak{}4e986591\\
Lector & c5b7a860\allowbreak{}701e0f33\allowbreak{}931847ef\allowbreak{}b2a992c3\allowbreak{}1bef6f0d\allowbreak{}9da926e5\allowbreak{}74e7ab3e\allowbreak{}68fb470f\\
Registro en 108 & \textemdash{}\\
\end{tabularx}\endgroup\par\vspace{6pt}
\Needspace{87pt}\phantomsection\label{trace-156}\noindent\hyperref[nav:vi-ruta-156]{\textbf{r5c3\_\allowbreak{}h+1v+1}}\quad Ordinal unificado / memoria: 156 / 153\par\nopagebreak
\begingroup\fontsize{8.3}{10.1}\selectfont\setlength{\tabcolsep}{4pt}\noindent\begin{tabularx}{\linewidth}{@{}p{.21\linewidth}X@{}}
Genealogía & 0e0d1f79\allowbreak{}3f8fb5b4\allowbreak{}55383130\allowbreak{}d2c9912e\allowbreak{}5b2700a0\allowbreak{}c2fb4ce1\allowbreak{}bd1b83c5\allowbreak{}4267422c\\
Lector & 997312cb\allowbreak{}2b9d7228\allowbreak{}95ddf8de\allowbreak{}8b30b950\allowbreak{}1c2d380d\allowbreak{}db484990\allowbreak{}6bc9cd25\allowbreak{}95449bc9\\
Registro en 108 & \textemdash{}\\
\end{tabularx}\endgroup\par\vspace{6pt}
\Needspace{87pt}\phantomsection\label{trace-157}\noindent\hyperref[nav:vi-ruta-157]{\textbf{r5c4\_\allowbreak{}h-1v-1}}\quad Ordinal unificado / memoria: 157 / 160\par\nopagebreak
\begingroup\fontsize{8.3}{10.1}\selectfont\setlength{\tabcolsep}{4pt}\noindent\begin{tabularx}{\linewidth}{@{}p{.21\linewidth}X@{}}
Genealogía & 698e05b3\allowbreak{}1e9a2622\allowbreak{}9e96d506\allowbreak{}b2415dfb\allowbreak{}f00df561\allowbreak{}46f44d71\allowbreak{}fcbea9b7\allowbreak{}8701e56b\\
Lector & 4920659c\allowbreak{}510a6cf9\allowbreak{}0e8944ec\allowbreak{}b1a49dff\allowbreak{}b1bfe16a\allowbreak{}d3249fd0\allowbreak{}b5fb8a10\allowbreak{}930a6547\\
Registro en 108 & \textemdash{}\\
\end{tabularx}\endgroup\par\vspace{6pt}
\Needspace{87pt}\phantomsection\label{trace-158}\noindent\hyperref[nav:vi-ruta-158]{\textbf{r5c4\_\allowbreak{}h-1v+1}}\quad Ordinal unificado / memoria: 158 / 159\par\nopagebreak
\begingroup\fontsize{8.3}{10.1}\selectfont\setlength{\tabcolsep}{4pt}\noindent\begin{tabularx}{\linewidth}{@{}p{.21\linewidth}X@{}}
Genealogía & a3dc2fcf\allowbreak{}597df676\allowbreak{}434ea90b\allowbreak{}519e5540\allowbreak{}b275620b\allowbreak{}60073073\allowbreak{}f9ecf517\allowbreak{}45655b0d\\
Lector & 2bc38a4f\allowbreak{}c5942519\allowbreak{}23f006d0\allowbreak{}b3ea2789\allowbreak{}df5cf2da\allowbreak{}0f7d89ba\allowbreak{}b5c40e7d\allowbreak{}fa00c9bd\\
Registro en 108 & \textemdash{}\\
\end{tabularx}\endgroup\par\vspace{6pt}
\Needspace{87pt}\phantomsection\label{trace-159}\noindent\hyperref[nav:vi-ruta-159]{\textbf{r5c4\_\allowbreak{}h+1v-1}}\quad Ordinal unificado / memoria: 159 / 158\par\nopagebreak
\begingroup\fontsize{8.3}{10.1}\selectfont\setlength{\tabcolsep}{4pt}\noindent\begin{tabularx}{\linewidth}{@{}p{.21\linewidth}X@{}}
Genealogía & 8944ea77\allowbreak{}aea3c31e\allowbreak{}b47da169\allowbreak{}2c7b6e0a\allowbreak{}3c809102\allowbreak{}f8b9002b\allowbreak{}a05b291b\allowbreak{}aa3233c1\\
Lector & 2d9f1de9\allowbreak{}bfa0686d\allowbreak{}cfd7cd7a\allowbreak{}c00fc0b9\allowbreak{}12d37ab7\allowbreak{}3f9fb467\allowbreak{}dccacfef\allowbreak{}55a9eddb\\
Registro en 108 & 23031ced\allowbreak{}1c53e342\allowbreak{}1cc8604c\allowbreak{}b7d3e603\allowbreak{}73a144d8\allowbreak{}1c02cf0c\allowbreak{}7307d379\allowbreak{}831633c3\\
\end{tabularx}\endgroup\par\vspace{6pt}
\Needspace{87pt}\phantomsection\label{trace-160}\noindent\hyperref[nav:vi-ruta-160]{\textbf{r5c4\_\allowbreak{}h+1v+1}}\quad Ordinal unificado / memoria: 160 / 157\par\nopagebreak
\begingroup\fontsize{8.3}{10.1}\selectfont\setlength{\tabcolsep}{4pt}\noindent\begin{tabularx}{\linewidth}{@{}p{.21\linewidth}X@{}}
Genealogía & 79e4be3a\allowbreak{}d5ebc675\allowbreak{}3342e428\allowbreak{}4adf56b7\allowbreak{}1d921b04\allowbreak{}cea8f219\allowbreak{}c842fa9f\allowbreak{}052e874f\\
Lector & dc5c02bc\allowbreak{}dfaf96e1\allowbreak{}2535891f\allowbreak{}9aa10bc7\allowbreak{}bd9659e0\allowbreak{}1354a2a2\allowbreak{}d5c6e9dc\allowbreak{}6ea2e751\\
Registro en 108 & \textemdash{}\\
\end{tabularx}\endgroup\par\vspace{6pt}
\Needspace{87pt}\phantomsection\label{trace-161}\noindent\hyperref[nav:vi-ruta-161]{\textbf{r5c5\_\allowbreak{}h-1v-1}}\quad Ordinal unificado / memoria: 161 / 164\par\nopagebreak
\begingroup\fontsize{8.3}{10.1}\selectfont\setlength{\tabcolsep}{4pt}\noindent\begin{tabularx}{\linewidth}{@{}p{.21\linewidth}X@{}}
Genealogía & a4e3cc2e\allowbreak{}ab44f4bb\allowbreak{}51b11970\allowbreak{}0b1b5cae\allowbreak{}bf22edbd\allowbreak{}82634847\allowbreak{}31a154bd\allowbreak{}a1c77f6b\\
Lector & 98b8a257\allowbreak{}24d3a874\allowbreak{}f2368cb2\allowbreak{}dc819bc6\allowbreak{}ce0d16ad\allowbreak{}f27017ad\allowbreak{}5f139361\allowbreak{}24662d5a\\
Registro en 108 & \textemdash{}\\
\end{tabularx}\endgroup\par\vspace{6pt}
\Needspace{87pt}\phantomsection\label{trace-162}\noindent\hyperref[nav:vi-ruta-162]{\textbf{r5c5\_\allowbreak{}h-1v+1}}\quad Ordinal unificado / memoria: 162 / 163\par\nopagebreak
\begingroup\fontsize{8.3}{10.1}\selectfont\setlength{\tabcolsep}{4pt}\noindent\begin{tabularx}{\linewidth}{@{}p{.21\linewidth}X@{}}
Genealogía & dafa9569\allowbreak{}2d1f9384\allowbreak{}c65e0afd\allowbreak{}29106c64\allowbreak{}a3e62d6f\allowbreak{}2d545d4a\allowbreak{}6b6b404e\allowbreak{}a1b2d70b\\
Lector & 2bdd7331\allowbreak{}f202139e\allowbreak{}a68552f2\allowbreak{}7c7eb9dc\allowbreak{}f445fdc5\allowbreak{}c70d95f9\allowbreak{}e3ec7c65\allowbreak{}58dbce1d\\
Registro en 108 & \textemdash{}\\
\end{tabularx}\endgroup\par\vspace{6pt}
\Needspace{87pt}\phantomsection\label{trace-163}\noindent\hyperref[nav:vi-ruta-163]{\textbf{r5c5\_\allowbreak{}h+1v-1}}\quad Ordinal unificado / memoria: 163 / 162\par\nopagebreak
\begingroup\fontsize{8.3}{10.1}\selectfont\setlength{\tabcolsep}{4pt}\noindent\begin{tabularx}{\linewidth}{@{}p{.21\linewidth}X@{}}
Genealogía & c6b14aed\allowbreak{}d3f88fa1\allowbreak{}7e0478cc\allowbreak{}431b63de\allowbreak{}ca8ee71b\allowbreak{}b8e6d79c\allowbreak{}4c811add\allowbreak{}87d03720\\
Lector & 6e152edf\allowbreak{}9b94425a\allowbreak{}404b2c13\allowbreak{}793f91e6\allowbreak{}88091f69\allowbreak{}eff6453c\allowbreak{}8efda5d6\allowbreak{}9fa8861c\\
Registro en 108 & \textemdash{}\\
\end{tabularx}\endgroup\par\vspace{6pt}
\Needspace{87pt}\phantomsection\label{trace-164}\noindent\hyperref[nav:vi-ruta-164]{\textbf{r5c5\_\allowbreak{}h+1v+1}}\quad Ordinal unificado / memoria: 164 / 161\par\nopagebreak
\begingroup\fontsize{8.3}{10.1}\selectfont\setlength{\tabcolsep}{4pt}\noindent\begin{tabularx}{\linewidth}{@{}p{.21\linewidth}X@{}}
Genealogía & 57d24c5d\allowbreak{}c0a69cd8\allowbreak{}111933d8\allowbreak{}bc40cc26\allowbreak{}0571c0ef\allowbreak{}4f8f3008\allowbreak{}f7197910\allowbreak{}ff07c7dc\\
Lector & 514b4f56\allowbreak{}201d42c1\allowbreak{}181a4a14\allowbreak{}961f9abe\allowbreak{}4a328f77\allowbreak{}51e0d100\allowbreak{}a75bfb10\allowbreak{}1ea22f75\\
Registro en 108 & 9bf67dd6\allowbreak{}633e13d0\allowbreak{}3463463e\allowbreak{}6d588097\allowbreak{}8f1459c3\allowbreak{}ba4e9352\allowbreak{}c82c30f6\allowbreak{}f5db9d13\\
\end{tabularx}\endgroup\par\vspace{6pt}
\Needspace{87pt}\phantomsection\label{trace-165}\noindent\hyperref[nav:vi-ruta-165]{\textbf{r5c6\_\allowbreak{}h-1v-1}}\quad Ordinal unificado / memoria: 165 / 168\par\nopagebreak
\begingroup\fontsize{8.3}{10.1}\selectfont\setlength{\tabcolsep}{4pt}\noindent\begin{tabularx}{\linewidth}{@{}p{.21\linewidth}X@{}}
Genealogía & a03908be\allowbreak{}0f96c387\allowbreak{}74c03f0b\allowbreak{}83c3acb9\allowbreak{}6e9e98b0\allowbreak{}83d11b55\allowbreak{}5c03fb4c\allowbreak{}6e686faf\\
Lector & a2e86299\allowbreak{}6bf61364\allowbreak{}ed0e6b6b\allowbreak{}37873be9\allowbreak{}5744ffab\allowbreak{}e1b4e869\allowbreak{}73b0917d\allowbreak{}492ddcf0\\
Registro en 108 & \textemdash{}\\
\end{tabularx}\endgroup\par\vspace{6pt}
\Needspace{87pt}\phantomsection\label{trace-166}\noindent\hyperref[nav:vi-ruta-166]{\textbf{r5c6\_\allowbreak{}h-1v+1}}\quad Ordinal unificado / memoria: 166 / 167\par\nopagebreak
\begingroup\fontsize{8.3}{10.1}\selectfont\setlength{\tabcolsep}{4pt}\noindent\begin{tabularx}{\linewidth}{@{}p{.21\linewidth}X@{}}
Genealogía & 81927e69\allowbreak{}032d8dda\allowbreak{}6c89f7e0\allowbreak{}c51591bb\allowbreak{}74a0922c\allowbreak{}58dc63b9\allowbreak{}6cc20c65\allowbreak{}08c171f8\\
Lector & b51c5e79\allowbreak{}6ab9a08f\allowbreak{}438f0908\allowbreak{}5e623965\allowbreak{}205e7f4b\allowbreak{}e2ed57ea\allowbreak{}ebc1afa1\allowbreak{}6d1504ce\\
Registro en 108 & 598cf254\allowbreak{}0b6a078e\allowbreak{}4972cb09\allowbreak{}c4a957ff\allowbreak{}723cc9c3\allowbreak{}e7d7797d\allowbreak{}e08ffcaf\allowbreak{}b6666426\\
\end{tabularx}\endgroup\par\vspace{6pt}
\Needspace{87pt}\phantomsection\label{trace-167}\noindent\hyperref[nav:vi-ruta-167]{\textbf{r5c6\_\allowbreak{}h+1v-1}}\quad Ordinal unificado / memoria: 167 / 166\par\nopagebreak
\begingroup\fontsize{8.3}{10.1}\selectfont\setlength{\tabcolsep}{4pt}\noindent\begin{tabularx}{\linewidth}{@{}p{.21\linewidth}X@{}}
Genealogía & a75845a7\allowbreak{}c9ce451a\allowbreak{}84ec8de2\allowbreak{}149e4578\allowbreak{}724b2982\allowbreak{}0b255821\allowbreak{}78837475\allowbreak{}acd97e29\\
Lector & 52f8b462\allowbreak{}7b04a1a9\allowbreak{}25e56d90\allowbreak{}21b47cdb\allowbreak{}f4e172ca\allowbreak{}6df35000\allowbreak{}72c1a8cf\allowbreak{}dd00d329\\
Registro en 108 & \textemdash{}\\
\end{tabularx}\endgroup\par\vspace{6pt}
\Needspace{87pt}\phantomsection\label{trace-168}\noindent\hyperref[nav:vi-ruta-168]{\textbf{r5c6\_\allowbreak{}h+1v+1}}\quad Ordinal unificado / memoria: 168 / 165\par\nopagebreak
\begingroup\fontsize{8.3}{10.1}\selectfont\setlength{\tabcolsep}{4pt}\noindent\begin{tabularx}{\linewidth}{@{}p{.21\linewidth}X@{}}
Genealogía & 09134b2a\allowbreak{}60678acc\allowbreak{}4e4d7fe4\allowbreak{}1328fa47\allowbreak{}5073a907\allowbreak{}14988713\allowbreak{}b330280c\allowbreak{}61906435\\
Lector & 3a010d94\allowbreak{}93ed73b6\allowbreak{}2d23d3ba\allowbreak{}f7295854\allowbreak{}40be9fb9\allowbreak{}b19ad1cf\allowbreak{}a7ca16b2\allowbreak{}970beb82\\
Registro en 108 & \textemdash{}\\
\end{tabularx}\endgroup\par\vspace{6pt}
\Needspace{87pt}\phantomsection\label{trace-169}\noindent\hyperref[nav:vi-ruta-169]{\textbf{r5c7\_\allowbreak{}h-1v-1}}\quad Ordinal unificado / memoria: 169 / 172\par\nopagebreak
\begingroup\fontsize{8.3}{10.1}\selectfont\setlength{\tabcolsep}{4pt}\noindent\begin{tabularx}{\linewidth}{@{}p{.21\linewidth}X@{}}
Genealogía & f0ba7ed7\allowbreak{}97a68175\allowbreak{}311757b8\allowbreak{}030cdecd\allowbreak{}682fb29e\allowbreak{}4eed2cda\allowbreak{}f8cfd383\allowbreak{}2be16d4c\\
Lector & 04477201\allowbreak{}ed2fa538\allowbreak{}46f5f328\allowbreak{}151d60bc\allowbreak{}cff05c49\allowbreak{}e4724a32\allowbreak{}37f24a74\allowbreak{}02a084ab\\
Registro en 108 & \textemdash{}\\
\end{tabularx}\endgroup\par\vspace{6pt}
\Needspace{87pt}\phantomsection\label{trace-170}\noindent\hyperref[nav:vi-ruta-170]{\textbf{r5c7\_\allowbreak{}h-1v+1}}\quad Ordinal unificado / memoria: 170 / 171\par\nopagebreak
\begingroup\fontsize{8.3}{10.1}\selectfont\setlength{\tabcolsep}{4pt}\noindent\begin{tabularx}{\linewidth}{@{}p{.21\linewidth}X@{}}
Genealogía & b6376735\allowbreak{}f9acad08\allowbreak{}bb8f9c39\allowbreak{}f4ca3e69\allowbreak{}be62405d\allowbreak{}36279437\allowbreak{}96f664a8\allowbreak{}a9cec6d6\\
Lector & 54983701\allowbreak{}597aea63\allowbreak{}96a38cb5\allowbreak{}9a141229\allowbreak{}50ea0616\allowbreak{}556c5847\allowbreak{}dc64e8be\allowbreak{}e3fec53f\\
Registro en 108 & \textemdash{}\\
\end{tabularx}\endgroup\par\vspace{6pt}
\Needspace{87pt}\phantomsection\label{trace-171}\noindent\hyperref[nav:vi-ruta-171]{\textbf{r5c7\_\allowbreak{}h+1v-1}}\quad Ordinal unificado / memoria: 171 / 170\par\nopagebreak
\begingroup\fontsize{8.3}{10.1}\selectfont\setlength{\tabcolsep}{4pt}\noindent\begin{tabularx}{\linewidth}{@{}p{.21\linewidth}X@{}}
Genealogía & b34a9417\allowbreak{}a09638f6\allowbreak{}bc02ffe9\allowbreak{}8471b0e2\allowbreak{}b627116e\allowbreak{}e13284a7\allowbreak{}ad74bf50\allowbreak{}74b0f4d2\\
Lector & b8df04a2\allowbreak{}3c90b576\allowbreak{}16a1bf83\allowbreak{}34a7f775\allowbreak{}7313f38e\allowbreak{}95a5c454\allowbreak{}0aa36d3e\allowbreak{}166549dd\\
Registro en 108 & \textemdash{}\\
\end{tabularx}\endgroup\par\vspace{6pt}
\Needspace{87pt}\phantomsection\label{trace-172}\noindent\hyperref[nav:vi-ruta-172]{\textbf{r5c7\_\allowbreak{}h+1v+1}}\quad Ordinal unificado / memoria: 172 / 169\par\nopagebreak
\begingroup\fontsize{8.3}{10.1}\selectfont\setlength{\tabcolsep}{4pt}\noindent\begin{tabularx}{\linewidth}{@{}p{.21\linewidth}X@{}}
Genealogía & 1a541520\allowbreak{}4e85bcd7\allowbreak{}ffabce87\allowbreak{}c95eff5f\allowbreak{}c47e613c\allowbreak{}eddd713d\allowbreak{}8f05b937\allowbreak{}18293d41\\
Lector & 3f7c2ace\allowbreak{}63014190\allowbreak{}6ef94299\allowbreak{}c92c3d34\allowbreak{}f7489791\allowbreak{}9e7ec93e\allowbreak{}685fd0b6\allowbreak{}210dd9f3\\
Registro en 108 & f764349c\allowbreak{}d2ba2f9d\allowbreak{}1ee7db9a\allowbreak{}6f0bd775\allowbreak{}d595c448\allowbreak{}0bce2698\allowbreak{}06180310\allowbreak{}0c3f18e5\\
\end{tabularx}\endgroup\par\vspace{6pt}
\Needspace{87pt}\phantomsection\label{trace-173}\noindent\hyperref[nav:vi-ruta-173]{\textbf{r5c8\_\allowbreak{}h-1v-1}}\quad Ordinal unificado / memoria: 173 / 176\par\nopagebreak
\begingroup\fontsize{8.3}{10.1}\selectfont\setlength{\tabcolsep}{4pt}\noindent\begin{tabularx}{\linewidth}{@{}p{.21\linewidth}X@{}}
Genealogía & 701966ab\allowbreak{}31f24a2e\allowbreak{}0ac479f7\allowbreak{}47422728\allowbreak{}dbafc440\allowbreak{}21012aef\allowbreak{}5bc7700e\allowbreak{}331d65b8\\
Lector & 5d22412c\allowbreak{}39797827\allowbreak{}c404556a\allowbreak{}d608c217\allowbreak{}b66ea943\allowbreak{}a89b289f\allowbreak{}458610a7\allowbreak{}0c151ad7\\
Registro en 108 & \textemdash{}\\
\end{tabularx}\endgroup\par\vspace{6pt}
\Needspace{87pt}\phantomsection\label{trace-174}\noindent\hyperref[nav:vi-ruta-174]{\textbf{r5c8\_\allowbreak{}h-1v+1}}\quad Ordinal unificado / memoria: 174 / 175\par\nopagebreak
\begingroup\fontsize{8.3}{10.1}\selectfont\setlength{\tabcolsep}{4pt}\noindent\begin{tabularx}{\linewidth}{@{}p{.21\linewidth}X@{}}
Genealogía & a03329e4\allowbreak{}0257d5ed\allowbreak{}e551debc\allowbreak{}7365527a\allowbreak{}24878f84\allowbreak{}197e9298\allowbreak{}100b0b55\allowbreak{}af57184b\\
Lector & b12e9f54\allowbreak{}798f62c1\allowbreak{}9f474b63\allowbreak{}d6834882\allowbreak{}3a829225\allowbreak{}e5a25e30\allowbreak{}8905bb4e\allowbreak{}98e7ff7d\\
Registro en 108 & e8f3e463\allowbreak{}9e854424\allowbreak{}d2658b46\allowbreak{}9d572245\allowbreak{}3265869a\allowbreak{}56f8e553\allowbreak{}445bc003\allowbreak{}7bd9882f\\
\end{tabularx}\endgroup\par\vspace{6pt}
\Needspace{87pt}\phantomsection\label{trace-175}\noindent\hyperref[nav:vi-ruta-175]{\textbf{r5c8\_\allowbreak{}h+1v-1}}\quad Ordinal unificado / memoria: 175 / 174\par\nopagebreak
\begingroup\fontsize{8.3}{10.1}\selectfont\setlength{\tabcolsep}{4pt}\noindent\begin{tabularx}{\linewidth}{@{}p{.21\linewidth}X@{}}
Genealogía & fe0e3cc0\allowbreak{}a84e1432\allowbreak{}89d0a34c\allowbreak{}2824d4fe\allowbreak{}4fddc992\allowbreak{}1c05852e\allowbreak{}b70f3627\allowbreak{}7c4ab0b0\\
Lector & 41a91226\allowbreak{}e39febce\allowbreak{}fa361982\allowbreak{}83195e11\allowbreak{}51884a54\allowbreak{}38e33053\allowbreak{}0b3217ad\allowbreak{}9b811153\\
Registro en 108 & \textemdash{}\\
\end{tabularx}\endgroup\par\vspace{6pt}
\Needspace{87pt}\phantomsection\label{trace-176}\noindent\hyperref[nav:vi-ruta-176]{\textbf{r5c8\_\allowbreak{}h+1v+1}}\quad Ordinal unificado / memoria: 176 / 173\par\nopagebreak
\begingroup\fontsize{8.3}{10.1}\selectfont\setlength{\tabcolsep}{4pt}\noindent\begin{tabularx}{\linewidth}{@{}p{.21\linewidth}X@{}}
Genealogía & 9f17a032\allowbreak{}ddfb1ca8\allowbreak{}9f4d34e1\allowbreak{}fc106e59\allowbreak{}e07a18fc\allowbreak{}1dbada65\allowbreak{}e5061f6b\allowbreak{}0974791e\\
Lector & 5cc037f3\allowbreak{}f593ef27\allowbreak{}c0b04dd6\allowbreak{}0612b27c\allowbreak{}3ee3fb21\allowbreak{}7cc66ae4\allowbreak{}95e7c6eb\allowbreak{}2c3f5800\\
Registro en 108 & \textemdash{}\\
\end{tabularx}\endgroup\par\vspace{6pt}
\Needspace{87pt}\phantomsection\label{trace-177}\noindent\hyperref[nav:vi-ruta-177]{\textbf{r5c9\_\allowbreak{}h-1v-1}}\quad Ordinal unificado / memoria: 177 / 180\par\nopagebreak
\begingroup\fontsize{8.3}{10.1}\selectfont\setlength{\tabcolsep}{4pt}\noindent\begin{tabularx}{\linewidth}{@{}p{.21\linewidth}X@{}}
Genealogía & d13b1075\allowbreak{}8f72225b\allowbreak{}c8a3bb0f\allowbreak{}29e7117f\allowbreak{}1e568076\allowbreak{}f05c67d1\allowbreak{}0e59f0ec\allowbreak{}d85a0d8d\\
Lector & 3a627fb4\allowbreak{}a5c98662\allowbreak{}723a034a\allowbreak{}9f86de71\allowbreak{}0c9cb3ab\allowbreak{}c8f8d035\allowbreak{}e035b6e5\allowbreak{}79396d47\\
Registro en 108 & 24249827\allowbreak{}8d713b13\allowbreak{}59912fec\allowbreak{}0f87443b\allowbreak{}9461fd3d\allowbreak{}d582fd98\allowbreak{}89d8a0f5\allowbreak{}668c22be\\
\end{tabularx}\endgroup\par\vspace{6pt}
\Needspace{87pt}\phantomsection\label{trace-178}\noindent\hyperref[nav:vi-ruta-178]{\textbf{r5c9\_\allowbreak{}h-1v+1}}\quad Ordinal unificado / memoria: 178 / 179\par\nopagebreak
\begingroup\fontsize{8.3}{10.1}\selectfont\setlength{\tabcolsep}{4pt}\noindent\begin{tabularx}{\linewidth}{@{}p{.21\linewidth}X@{}}
Genealogía & 05d5c4bf\allowbreak{}566052d2\allowbreak{}f9fb3f06\allowbreak{}980ca8d3\allowbreak{}107fcf46\allowbreak{}ffefa36a\allowbreak{}99432ea5\allowbreak{}1df4ffaf\\
Lector & fed58486\allowbreak{}0b86d5c5\allowbreak{}a263e009\allowbreak{}5165b1b1\allowbreak{}f3766fe8\allowbreak{}3937baad\allowbreak{}4fe826a7\allowbreak{}98377b6b\\
Registro en 108 & \textemdash{}\\
\end{tabularx}\endgroup\par\vspace{6pt}
\Needspace{87pt}\phantomsection\label{trace-179}\noindent\hyperref[nav:vi-ruta-179]{\textbf{r5c9\_\allowbreak{}h+1v-1}}\quad Ordinal unificado / memoria: 179 / 178\par\nopagebreak
\begingroup\fontsize{8.3}{10.1}\selectfont\setlength{\tabcolsep}{4pt}\noindent\begin{tabularx}{\linewidth}{@{}p{.21\linewidth}X@{}}
Genealogía & faa1ebd5\allowbreak{}78b981da\allowbreak{}a2f34f8d\allowbreak{}42e7ddcf\allowbreak{}87be1072\allowbreak{}ab756989\allowbreak{}1f0566f3\allowbreak{}aa0fcc75\\
Lector & 59f77181\allowbreak{}46a7eed5\allowbreak{}8f86cef8\allowbreak{}69bc7e00\allowbreak{}ac61585c\allowbreak{}b84db3ca\allowbreak{}040a23ab\allowbreak{}561ec874\\
Registro en 108 & \textemdash{}\\
\end{tabularx}\endgroup\par\vspace{6pt}
\Needspace{87pt}\phantomsection\label{trace-180}\noindent\hyperref[nav:vi-ruta-180]{\textbf{r5c9\_\allowbreak{}h+1v+1}}\quad Ordinal unificado / memoria: 180 / 177\par\nopagebreak
\begingroup\fontsize{8.3}{10.1}\selectfont\setlength{\tabcolsep}{4pt}\noindent\begin{tabularx}{\linewidth}{@{}p{.21\linewidth}X@{}}
Genealogía & a2caaa18\allowbreak{}85e5e1ec\allowbreak{}3a77b1b1\allowbreak{}476830d4\allowbreak{}cdea1d83\allowbreak{}ef5de104\allowbreak{}3297f20d\allowbreak{}5b040bae\\
Lector & 33b2af69\allowbreak{}d9d7267f\allowbreak{}506c5047\allowbreak{}bd415adc\allowbreak{}3eb10323\allowbreak{}7644181f\allowbreak{}f6e49fca\allowbreak{}db855252\\
Registro en 108 & \textemdash{}\\
\end{tabularx}\endgroup\par\vspace{6pt}
\Needspace{87pt}\phantomsection\label{trace-181}\noindent\hyperref[nav:vi-ruta-181]{\textbf{r6c1\_\allowbreak{}h-1v-1}}\quad Ordinal unificado / memoria: 181 / 184\par\nopagebreak
\begingroup\fontsize{8.3}{10.1}\selectfont\setlength{\tabcolsep}{4pt}\noindent\begin{tabularx}{\linewidth}{@{}p{.21\linewidth}X@{}}
Genealogía & 57238369\allowbreak{}45af498a\allowbreak{}ee1f7524\allowbreak{}fa9f3ca6\allowbreak{}25afae9a\allowbreak{}8ad27788\allowbreak{}75e0f116\allowbreak{}ddf54bd0\\
Lector & 26c99915\allowbreak{}4dc98e67\allowbreak{}ef1eb0cb\allowbreak{}e3c4af5d\allowbreak{}c49cc63f\allowbreak{}f9acd722\allowbreak{}0b474c1d\allowbreak{}35deb661\\
Registro en 108 & e29c6e53\allowbreak{}cf987fd4\allowbreak{}cd079172\allowbreak{}3fa331bd\allowbreak{}6c54a8c0\allowbreak{}e6be1441\allowbreak{}99dae0d5\allowbreak{}aa759d38\\
\end{tabularx}\endgroup\par\vspace{6pt}
\Needspace{87pt}\phantomsection\label{trace-182}\noindent\hyperref[nav:vi-ruta-182]{\textbf{r6c1\_\allowbreak{}h-1v+1}}\quad Ordinal unificado / memoria: 182 / 183\par\nopagebreak
\begingroup\fontsize{8.3}{10.1}\selectfont\setlength{\tabcolsep}{4pt}\noindent\begin{tabularx}{\linewidth}{@{}p{.21\linewidth}X@{}}
Genealogía & e39ec08b\allowbreak{}be801ae5\allowbreak{}36649422\allowbreak{}bf79b77a\allowbreak{}37f016ba\allowbreak{}9a4bf092\allowbreak{}521046f2\allowbreak{}8fab1bc6\\
Lector & 68b45d01\allowbreak{}d606e5ea\allowbreak{}f1737c7f\allowbreak{}7d9c0e21\allowbreak{}61113069\allowbreak{}76f593bf\allowbreak{}8cb7006a\allowbreak{}1f0989a7\\
Registro en 108 & \textemdash{}\\
\end{tabularx}\endgroup\par\vspace{6pt}
\Needspace{87pt}\phantomsection\label{trace-183}\noindent\hyperref[nav:vi-ruta-183]{\textbf{r6c1\_\allowbreak{}h+1v-1}}\quad Ordinal unificado / memoria: 183 / 182\par\nopagebreak
\begingroup\fontsize{8.3}{10.1}\selectfont\setlength{\tabcolsep}{4pt}\noindent\begin{tabularx}{\linewidth}{@{}p{.21\linewidth}X@{}}
Genealogía & 2262a4ea\allowbreak{}d8d34652\allowbreak{}790e70bc\allowbreak{}f39dabf4\allowbreak{}c6da0042\allowbreak{}f349313d\allowbreak{}cd38adb2\allowbreak{}41f7266b\\
Lector & e348c328\allowbreak{}36c7d0ef\allowbreak{}2ec19bed\allowbreak{}6383e925\allowbreak{}7ee19242\allowbreak{}9d72bf57\allowbreak{}548f6155\allowbreak{}29d8ae53\\
Registro en 108 & \textemdash{}\\
\end{tabularx}\endgroup\par\vspace{6pt}
\Needspace{87pt}\phantomsection\label{trace-184}\noindent\hyperref[nav:vi-ruta-184]{\textbf{r6c1\_\allowbreak{}h+1v+1}}\quad Ordinal unificado / memoria: 184 / 181\par\nopagebreak
\begingroup\fontsize{8.3}{10.1}\selectfont\setlength{\tabcolsep}{4pt}\noindent\begin{tabularx}{\linewidth}{@{}p{.21\linewidth}X@{}}
Genealogía & 7c76e8e8\allowbreak{}f5478ad5\allowbreak{}36bb1cb8\allowbreak{}f675ca9a\allowbreak{}eceb0754\allowbreak{}1dc9051b\allowbreak{}6af6607d\allowbreak{}9276aa67\\
Lector & 86da2f8b\allowbreak{}a7211912\allowbreak{}d067c900\allowbreak{}8753f8e1\allowbreak{}15b2fbcc\allowbreak{}6eca36f4\allowbreak{}2bee9c89\allowbreak{}352fd976\\
Registro en 108 & \textemdash{}\\
\end{tabularx}\endgroup\par\vspace{6pt}
\Needspace{87pt}\phantomsection\label{trace-185}\noindent\hyperref[nav:vi-ruta-185]{\textbf{r6c2\_\allowbreak{}h-1v-1}}\quad Ordinal unificado / memoria: 185 / 188\par\nopagebreak
\begingroup\fontsize{8.3}{10.1}\selectfont\setlength{\tabcolsep}{4pt}\noindent\begin{tabularx}{\linewidth}{@{}p{.21\linewidth}X@{}}
Genealogía & 7b821c47\allowbreak{}1bf4e01e\allowbreak{}7dcfca13\allowbreak{}41b32549\allowbreak{}1b18cf5a\allowbreak{}dbe433f3\allowbreak{}054bb10a\allowbreak{}d081a842\\
Lector & a37372d4\allowbreak{}b23fd893\allowbreak{}7236a0d2\allowbreak{}35736280\allowbreak{}1e9e9014\allowbreak{}0f10eb07\allowbreak{}52028f08\allowbreak{}35fa0d80\\
Registro en 108 & \textemdash{}\\
\end{tabularx}\endgroup\par\vspace{6pt}
\Needspace{87pt}\phantomsection\label{trace-186}\noindent\hyperref[nav:vi-ruta-186]{\textbf{r6c2\_\allowbreak{}h-1v+1}}\quad Ordinal unificado / memoria: 186 / 187\par\nopagebreak
\begingroup\fontsize{8.3}{10.1}\selectfont\setlength{\tabcolsep}{4pt}\noindent\begin{tabularx}{\linewidth}{@{}p{.21\linewidth}X@{}}
Genealogía & 77bf3617\allowbreak{}7edee2e3\allowbreak{}31d458ba\allowbreak{}93651e87\allowbreak{}788985b6\allowbreak{}fbdfe5b7\allowbreak{}35256fad\allowbreak{}b6b05411\\
Lector & 3be69c77\allowbreak{}c702f3b1\allowbreak{}7d2f920c\allowbreak{}2531cf2d\allowbreak{}4e4c212e\allowbreak{}395a7190\allowbreak{}5205f20c\allowbreak{}05c4892e\\
Registro en 108 & 9485907b\allowbreak{}c9b83fa7\allowbreak{}aabbe744\allowbreak{}a6a88d1d\allowbreak{}3223272e\allowbreak{}a89792d7\allowbreak{}558c0578\allowbreak{}84f8e9f0\\
\end{tabularx}\endgroup\par\vspace{6pt}
\Needspace{87pt}\phantomsection\label{trace-187}\noindent\hyperref[nav:vi-ruta-187]{\textbf{r6c2\_\allowbreak{}h+1v-1}}\quad Ordinal unificado / memoria: 187 / 186\par\nopagebreak
\begingroup\fontsize{8.3}{10.1}\selectfont\setlength{\tabcolsep}{4pt}\noindent\begin{tabularx}{\linewidth}{@{}p{.21\linewidth}X@{}}
Genealogía & a039b8fd\allowbreak{}a9943425\allowbreak{}c307ca97\allowbreak{}bc18ad16\allowbreak{}1e61a4a1\allowbreak{}2c0516fe\allowbreak{}66793b76\allowbreak{}5e700217\\
Lector & cd055221\allowbreak{}8d548e6a\allowbreak{}3dd43321\allowbreak{}97268fca\allowbreak{}2c06a5b8\allowbreak{}1c36d69a\allowbreak{}aae3d2c7\allowbreak{}17d9f0a2\\
Registro en 108 & \textemdash{}\\
\end{tabularx}\endgroup\par\vspace{6pt}
\Needspace{87pt}\phantomsection\label{trace-188}\noindent\hyperref[nav:vi-ruta-188]{\textbf{r6c2\_\allowbreak{}h+1v+1}}\quad Ordinal unificado / memoria: 188 / 185\par\nopagebreak
\begingroup\fontsize{8.3}{10.1}\selectfont\setlength{\tabcolsep}{4pt}\noindent\begin{tabularx}{\linewidth}{@{}p{.21\linewidth}X@{}}
Genealogía & 8615c240\allowbreak{}bf77cdf6\allowbreak{}ceadf5b0\allowbreak{}8213f626\allowbreak{}5f4ca614\allowbreak{}23b297ef\allowbreak{}17ac1b43\allowbreak{}dddf3c99\\
Lector & bd032ab9\allowbreak{}96bb19de\allowbreak{}9c107495\allowbreak{}3658e061\allowbreak{}c757ebca\allowbreak{}4281f4b0\allowbreak{}89db87ed\allowbreak{}939370c9\\
Registro en 108 & \textemdash{}\\
\end{tabularx}\endgroup\par\vspace{6pt}
\Needspace{87pt}\phantomsection\label{trace-189}\noindent\hyperref[nav:vi-ruta-189]{\textbf{r6c3\_\allowbreak{}h-1v-1}}\quad Ordinal unificado / memoria: 189 / 192\par\nopagebreak
\begingroup\fontsize{8.3}{10.1}\selectfont\setlength{\tabcolsep}{4pt}\noindent\begin{tabularx}{\linewidth}{@{}p{.21\linewidth}X@{}}
Genealogía & cd011cd5\allowbreak{}866d1188\allowbreak{}89051f57\allowbreak{}ac2a0d2b\allowbreak{}5ca6c7ab\allowbreak{}bc4b7d73\allowbreak{}3a0cc355\allowbreak{}15322e9b\\
Lector & 91deae21\allowbreak{}588e4bbf\allowbreak{}1be6e392\allowbreak{}88698ebd\allowbreak{}cfcb2cde\allowbreak{}1320d3f0\allowbreak{}ef090293\allowbreak{}f311c428\\
Registro en 108 & \textemdash{}\\
\end{tabularx}\endgroup\par\vspace{6pt}
\Needspace{87pt}\phantomsection\label{trace-190}\noindent\hyperref[nav:vi-ruta-190]{\textbf{r6c3\_\allowbreak{}h-1v+1}}\quad Ordinal unificado / memoria: 190 / 191\par\nopagebreak
\begingroup\fontsize{8.3}{10.1}\selectfont\setlength{\tabcolsep}{4pt}\noindent\begin{tabularx}{\linewidth}{@{}p{.21\linewidth}X@{}}
Genealogía & 7b28dd5d\allowbreak{}aca27177\allowbreak{}a2ca44f0\allowbreak{}b080b424\allowbreak{}07a6023d\allowbreak{}8a10ce5d\allowbreak{}6e034058\allowbreak{}ddeadc35\\
Lector & ccb618b1\allowbreak{}bf7c7851\allowbreak{}9fb2a0fb\allowbreak{}4e8ce96b\allowbreak{}e3fd471e\allowbreak{}b2a4f12c\allowbreak{}bfb2237f\allowbreak{}df4a50c4\\
Registro en 108 & 59349e87\allowbreak{}98fb69e2\allowbreak{}e5ae6c2c\allowbreak{}f1ca8ec3\allowbreak{}053390b2\allowbreak{}07b9dc18\allowbreak{}e92a5351\allowbreak{}bba14dc5\\
\end{tabularx}\endgroup\par\vspace{6pt}
\Needspace{87pt}\phantomsection\label{trace-191}\noindent\hyperref[nav:vi-ruta-191]{\textbf{r6c3\_\allowbreak{}h+1v-1}}\quad Ordinal unificado / memoria: 191 / 190\par\nopagebreak
\begingroup\fontsize{8.3}{10.1}\selectfont\setlength{\tabcolsep}{4pt}\noindent\begin{tabularx}{\linewidth}{@{}p{.21\linewidth}X@{}}
Genealogía & 41e436cb\allowbreak{}eb2b64cf\allowbreak{}ecd0594d\allowbreak{}db877df5\allowbreak{}834695c8\allowbreak{}c3b03374\allowbreak{}2d8bdd9e\allowbreak{}04c2774c\\
Lector & 766c9065\allowbreak{}f7080b17\allowbreak{}e85ed5f8\allowbreak{}029fb902\allowbreak{}f2ab1b56\allowbreak{}3deba91f\allowbreak{}4bc3ecd2\allowbreak{}30c44c7e\\
Registro en 108 & \textemdash{}\\
\end{tabularx}\endgroup\par\vspace{6pt}
\Needspace{87pt}\phantomsection\label{trace-192}\noindent\hyperref[nav:vi-ruta-192]{\textbf{r6c3\_\allowbreak{}h+1v+1}}\quad Ordinal unificado / memoria: 192 / 189\par\nopagebreak
\begingroup\fontsize{8.3}{10.1}\selectfont\setlength{\tabcolsep}{4pt}\noindent\begin{tabularx}{\linewidth}{@{}p{.21\linewidth}X@{}}
Genealogía & 2b62e065\allowbreak{}285f7e9d\allowbreak{}ff5a3ecd\allowbreak{}8ef42eb7\allowbreak{}3412a1ae\allowbreak{}0bd6ecfa\allowbreak{}a0f70240\allowbreak{}b91fb5e9\\
Lector & 102c4214\allowbreak{}f537ef17\allowbreak{}f9428490\allowbreak{}ac06f432\allowbreak{}b8c2bb4e\allowbreak{}74669c28\allowbreak{}4d556fbe\allowbreak{}f2f3c983\\
Registro en 108 & \textemdash{}\\
\end{tabularx}\endgroup\par\vspace{6pt}
\Needspace{87pt}\phantomsection\label{trace-193}\noindent\hyperref[nav:vi-ruta-193]{\textbf{r6c4\_\allowbreak{}h-1v-1}}\quad Ordinal unificado / memoria: 193 / 196\par\nopagebreak
\begingroup\fontsize{8.3}{10.1}\selectfont\setlength{\tabcolsep}{4pt}\noindent\begin{tabularx}{\linewidth}{@{}p{.21\linewidth}X@{}}
Genealogía & 978a72f4\allowbreak{}858eb2d5\allowbreak{}67ab4b95\allowbreak{}f6eebbde\allowbreak{}0c0eff2f\allowbreak{}c96c5e03\allowbreak{}55718f60\allowbreak{}86eb7ebd\\
Lector & 4bf0490c\allowbreak{}b60c3170\allowbreak{}3dd122be\allowbreak{}07418195\allowbreak{}8c4e70fb\allowbreak{}96f4d3f4\allowbreak{}ce35715f\allowbreak{}4e54679a\\
Registro en 108 & 09534891\allowbreak{}acc2eaf6\allowbreak{}9f6c6965\allowbreak{}a3dcbf12\allowbreak{}4972602e\allowbreak{}60fbf8fa\allowbreak{}02e0909e\allowbreak{}20d5411b\\
\end{tabularx}\endgroup\par\vspace{6pt}
\Needspace{87pt}\phantomsection\label{trace-194}\noindent\hyperref[nav:vi-ruta-194]{\textbf{r6c4\_\allowbreak{}h-1v+1}}\quad Ordinal unificado / memoria: 194 / 195\par\nopagebreak
\begingroup\fontsize{8.3}{10.1}\selectfont\setlength{\tabcolsep}{4pt}\noindent\begin{tabularx}{\linewidth}{@{}p{.21\linewidth}X@{}}
Genealogía & dcdcf9eb\allowbreak{}3b8abc3f\allowbreak{}ea552a68\allowbreak{}45f15c06\allowbreak{}b3f51f76\allowbreak{}b9bd638f\allowbreak{}95c451fd\allowbreak{}3cf68655\\
Lector & 52d4bbb4\allowbreak{}0854b103\allowbreak{}ab52b35e\allowbreak{}a76890da\allowbreak{}19bf740b\allowbreak{}b9db8f78\allowbreak{}4ae8f256\allowbreak{}fc3c4d04\\
Registro en 108 & \textemdash{}\\
\end{tabularx}\endgroup\par\vspace{6pt}
\Needspace{87pt}\phantomsection\label{trace-195}\noindent\hyperref[nav:vi-ruta-195]{\textbf{r6c4\_\allowbreak{}h+1v-1}}\quad Ordinal unificado / memoria: 195 / 194\par\nopagebreak
\begingroup\fontsize{8.3}{10.1}\selectfont\setlength{\tabcolsep}{4pt}\noindent\begin{tabularx}{\linewidth}{@{}p{.21\linewidth}X@{}}
Genealogía & 78f3ffaa\allowbreak{}921e46c6\allowbreak{}1e0a5e8f\allowbreak{}0098d060\allowbreak{}fca37d87\allowbreak{}139b6541\allowbreak{}dce6069a\allowbreak{}8bd8a27f\\
Lector & 32872bf6\allowbreak{}4adafd68\allowbreak{}5ad864ba\allowbreak{}bb2ca0d2\allowbreak{}e7551b61\allowbreak{}ccfecec1\allowbreak{}45ac607c\allowbreak{}ffe25d0d\\
Registro en 108 & \textemdash{}\\
\end{tabularx}\endgroup\par\vspace{6pt}
\Needspace{87pt}\phantomsection\label{trace-196}\noindent\hyperref[nav:vi-ruta-196]{\textbf{r6c4\_\allowbreak{}h+1v+1}}\quad Ordinal unificado / memoria: 196 / 193\par\nopagebreak
\begingroup\fontsize{8.3}{10.1}\selectfont\setlength{\tabcolsep}{4pt}\noindent\begin{tabularx}{\linewidth}{@{}p{.21\linewidth}X@{}}
Genealogía & 77cf5c93\allowbreak{}9793f9ff\allowbreak{}efd542cd\allowbreak{}fd63c675\allowbreak{}47932703\allowbreak{}378b7838\allowbreak{}2725fd39\allowbreak{}6edf7833\\
Lector & 28764df7\allowbreak{}e51b08e0\allowbreak{}cf0a4e12\allowbreak{}31218c25\allowbreak{}be1d6f46\allowbreak{}c66294bc\allowbreak{}d19eed1f\allowbreak{}6b285cce\\
Registro en 108 & \textemdash{}\\
\end{tabularx}\endgroup\par\vspace{6pt}
\Needspace{87pt}\phantomsection\label{trace-197}\noindent\hyperref[nav:vi-ruta-197]{\textbf{r6c5\_\allowbreak{}h-1v-1}}\quad Ordinal unificado / memoria: 197 / 200\par\nopagebreak
\begingroup\fontsize{8.3}{10.1}\selectfont\setlength{\tabcolsep}{4pt}\noindent\begin{tabularx}{\linewidth}{@{}p{.21\linewidth}X@{}}
Genealogía & b6badb1f\allowbreak{}3717c85f\allowbreak{}8eab3126\allowbreak{}2b0fde33\allowbreak{}8dbc4919\allowbreak{}7b8eab98\allowbreak{}20df39cf\allowbreak{}d1d22827\\
Lector & 17bd49a7\allowbreak{}de09fc63\allowbreak{}fbf85730\allowbreak{}b09b8c69\allowbreak{}da5258cd\allowbreak{}3a24e4c1\allowbreak{}1552fad8\allowbreak{}eb796267\\
Registro en 108 & \textemdash{}\\
\end{tabularx}\endgroup\par\vspace{6pt}
\Needspace{87pt}\phantomsection\label{trace-198}\noindent\hyperref[nav:vi-ruta-198]{\textbf{r6c5\_\allowbreak{}h-1v+1}}\quad Ordinal unificado / memoria: 198 / 199\par\nopagebreak
\begingroup\fontsize{8.3}{10.1}\selectfont\setlength{\tabcolsep}{4pt}\noindent\begin{tabularx}{\linewidth}{@{}p{.21\linewidth}X@{}}
Genealogía & 7493844e\allowbreak{}e9e55ec3\allowbreak{}fa446355\allowbreak{}3c671c3b\allowbreak{}fe04eb94\allowbreak{}006e369b\allowbreak{}09aa1c2f\allowbreak{}e3fd8152\\
Lector & d6f20c0a\allowbreak{}7858c075\allowbreak{}2246770f\allowbreak{}fbfd6643\allowbreak{}2cef12d3\allowbreak{}1a555564\allowbreak{}cd4c76c2\allowbreak{}195800ee\\
Registro en 108 & 247db8a5\allowbreak{}8d592a8c\allowbreak{}dd0ba517\allowbreak{}321e06c1\allowbreak{}f88eff82\allowbreak{}41629458\allowbreak{}6fe44cb5\allowbreak{}07cd1737\\
\end{tabularx}\endgroup\par\vspace{6pt}
\Needspace{87pt}\phantomsection\label{trace-199}\noindent\hyperref[nav:vi-ruta-199]{\textbf{r6c5\_\allowbreak{}h+1v-1}}\quad Ordinal unificado / memoria: 199 / 198\par\nopagebreak
\begingroup\fontsize{8.3}{10.1}\selectfont\setlength{\tabcolsep}{4pt}\noindent\begin{tabularx}{\linewidth}{@{}p{.21\linewidth}X@{}}
Genealogía & 0b8c5e26\allowbreak{}fd15e548\allowbreak{}7404c126\allowbreak{}325089f3\allowbreak{}744921e3\allowbreak{}c70229aa\allowbreak{}7c4efd95\allowbreak{}98c80084\\
Lector & bb50d7c6\allowbreak{}b121d5b9\allowbreak{}b27c2a8d\allowbreak{}f96ae10f\allowbreak{}4164b1fc\allowbreak{}3d2ae868\allowbreak{}7629f5ba\allowbreak{}15b9d059\\
Registro en 108 & \textemdash{}\\
\end{tabularx}\endgroup\par\vspace{6pt}
\Needspace{87pt}\phantomsection\label{trace-200}\noindent\hyperref[nav:vi-ruta-200]{\textbf{r6c5\_\allowbreak{}h+1v+1}}\quad Ordinal unificado / memoria: 200 / 197\par\nopagebreak
\begingroup\fontsize{8.3}{10.1}\selectfont\setlength{\tabcolsep}{4pt}\noindent\begin{tabularx}{\linewidth}{@{}p{.21\linewidth}X@{}}
Genealogía & 459225a8\allowbreak{}c62810d9\allowbreak{}2ebb16df\allowbreak{}6c284185\allowbreak{}5fdcc797\allowbreak{}3f566954\allowbreak{}9e10bdb7\allowbreak{}f48ba668\\
Lector & 3207260b\allowbreak{}e401417e\allowbreak{}ccf8d9f7\allowbreak{}56c5e76d\allowbreak{}13ba51ee\allowbreak{}efe672fe\allowbreak{}ec28d8ab\allowbreak{}04ef1ed5\\
Registro en 108 & \textemdash{}\\
\end{tabularx}\endgroup\par\vspace{6pt}
\Needspace{87pt}\phantomsection\label{trace-201}\noindent\hyperref[nav:vi-ruta-201]{\textbf{r6c6\_\allowbreak{}h-1v-1}}\quad Ordinal unificado / memoria: 201 / 204\par\nopagebreak
\begingroup\fontsize{8.3}{10.1}\selectfont\setlength{\tabcolsep}{4pt}\noindent\begin{tabularx}{\linewidth}{@{}p{.21\linewidth}X@{}}
Genealogía & 856f2c4e\allowbreak{}cede7c9f\allowbreak{}ad89891b\allowbreak{}9a493ce3\allowbreak{}36be2d4b\allowbreak{}5153ce9e\allowbreak{}cee9540f\allowbreak{}cd10311c\\
Lector & a8e2d941\allowbreak{}75cce86b\allowbreak{}4b4ebf52\allowbreak{}a6722aaa\allowbreak{}5895cc5b\allowbreak{}77acfab8\allowbreak{}bfc923d3\allowbreak{}76d84f57\\
Registro en 108 & \textemdash{}\\
\end{tabularx}\endgroup\par\vspace{6pt}
\Needspace{87pt}\phantomsection\label{trace-202}\noindent\hyperref[nav:vi-ruta-202]{\textbf{r6c6\_\allowbreak{}h-1v+1}}\quad Ordinal unificado / memoria: 202 / 203\par\nopagebreak
\begingroup\fontsize{8.3}{10.1}\selectfont\setlength{\tabcolsep}{4pt}\noindent\begin{tabularx}{\linewidth}{@{}p{.21\linewidth}X@{}}
Genealogía & 362984c5\allowbreak{}834759f4\allowbreak{}8f4db727\allowbreak{}a5e95cba\allowbreak{}6a614bb5\allowbreak{}b737be4b\allowbreak{}a5aa7e00\allowbreak{}92099305\\
Lector & 54b65075\allowbreak{}50fde869\allowbreak{}f152d240\allowbreak{}26fee643\allowbreak{}979e44fe\allowbreak{}fccb5a54\allowbreak{}fd5374ce\allowbreak{}de0c4178\\
Registro en 108 & \textemdash{}\\
\end{tabularx}\endgroup\par\vspace{6pt}
\Needspace{87pt}\phantomsection\label{trace-203}\noindent\hyperref[nav:vi-ruta-203]{\textbf{r6c6\_\allowbreak{}h+1v-1}}\quad Ordinal unificado / memoria: 203 / 202\par\nopagebreak
\begingroup\fontsize{8.3}{10.1}\selectfont\setlength{\tabcolsep}{4pt}\noindent\begin{tabularx}{\linewidth}{@{}p{.21\linewidth}X@{}}
Genealogía & ed910a02\allowbreak{}3c24aa3e\allowbreak{}30b58383\allowbreak{}d60bafc5\allowbreak{}d58eb8a1\allowbreak{}f08e0253\allowbreak{}b59d1c80\allowbreak{}4d53f35f\\
Lector & 1c9ca174\allowbreak{}be7648f8\allowbreak{}a2b7e60e\allowbreak{}181a19b3\allowbreak{}b34f7c4e\allowbreak{}0d98e91b\allowbreak{}ac5a4cba\allowbreak{}cc8c9a20\\
Registro en 108 & \textemdash{}\\
\end{tabularx}\endgroup\par\vspace{6pt}
\Needspace{87pt}\phantomsection\label{trace-204}\noindent\hyperref[nav:vi-ruta-204]{\textbf{r6c6\_\allowbreak{}h+1v+1}}\quad Ordinal unificado / memoria: 204 / 201\par\nopagebreak
\begingroup\fontsize{8.3}{10.1}\selectfont\setlength{\tabcolsep}{4pt}\noindent\begin{tabularx}{\linewidth}{@{}p{.21\linewidth}X@{}}
Genealogía & ea448f4e\allowbreak{}ea6376a5\allowbreak{}8f482d67\allowbreak{}6495495f\allowbreak{}0c7b3c6b\allowbreak{}b76ccae4\allowbreak{}0998ce25\allowbreak{}47269821\\
Lector & 7d817835\allowbreak{}fbd82223\allowbreak{}ed54da24\allowbreak{}836deff3\allowbreak{}fe6b2f2a\allowbreak{}2c237cae\allowbreak{}fadfa287\allowbreak{}194b6ac4\\
Registro en 108 & 1204da29\allowbreak{}afa516ca\allowbreak{}e33ccda6\allowbreak{}4944005b\allowbreak{}00ffe568\allowbreak{}f4826a58\allowbreak{}5823f1bb\allowbreak{}2e720ae3\\
\end{tabularx}\endgroup\par\vspace{6pt}
\Needspace{87pt}\phantomsection\label{trace-205}\noindent\hyperref[nav:vi-ruta-205]{\textbf{r6c7\_\allowbreak{}h-1v-1}}\quad Ordinal unificado / memoria: 205 / 208\par\nopagebreak
\begingroup\fontsize{8.3}{10.1}\selectfont\setlength{\tabcolsep}{4pt}\noindent\begin{tabularx}{\linewidth}{@{}p{.21\linewidth}X@{}}
Genealogía & 88f473d2\allowbreak{}cdc82fd9\allowbreak{}a950b895\allowbreak{}1b6621d7\allowbreak{}85eb978f\allowbreak{}2df99dbc\allowbreak{}0cf48b60\allowbreak{}503f0407\\
Lector & 0f39551e\allowbreak{}76068333\allowbreak{}3bc943e7\allowbreak{}a463b426\allowbreak{}ae514492\allowbreak{}913154ad\allowbreak{}8c8a50a5\allowbreak{}79ac0f4b\\
Registro en 108 & \textemdash{}\\
\end{tabularx}\endgroup\par\vspace{6pt}
\Needspace{87pt}\phantomsection\label{trace-206}\noindent\hyperref[nav:vi-ruta-206]{\textbf{r6c7\_\allowbreak{}h-1v+1}}\quad Ordinal unificado / memoria: 206 / 207\par\nopagebreak
\begingroup\fontsize{8.3}{10.1}\selectfont\setlength{\tabcolsep}{4pt}\noindent\begin{tabularx}{\linewidth}{@{}p{.21\linewidth}X@{}}
Genealogía & 616fe1c1\allowbreak{}53e8405f\allowbreak{}b9fdd203\allowbreak{}ddc30c40\allowbreak{}5976b9f5\allowbreak{}3a6b887a\allowbreak{}78eb2deb\allowbreak{}09ff3019\\
Lector & 0f506299\allowbreak{}138488a5\allowbreak{}bbb95a41\allowbreak{}ea186f87\allowbreak{}7f691de6\allowbreak{}2c43341f\allowbreak{}65029d52\allowbreak{}a295d822\\
Registro en 108 & e35a8717\allowbreak{}8dee05ba\allowbreak{}4b834fcc\allowbreak{}d04d9151\allowbreak{}72318520\allowbreak{}9b3bf306\allowbreak{}de515f71\allowbreak{}c318bc40\\
\end{tabularx}\endgroup\par\vspace{6pt}
\Needspace{87pt}\phantomsection\label{trace-207}\noindent\hyperref[nav:vi-ruta-207]{\textbf{r6c7\_\allowbreak{}h+1v-1}}\quad Ordinal unificado / memoria: 207 / 206\par\nopagebreak
\begingroup\fontsize{8.3}{10.1}\selectfont\setlength{\tabcolsep}{4pt}\noindent\begin{tabularx}{\linewidth}{@{}p{.21\linewidth}X@{}}
Genealogía & 4ff06f4f\allowbreak{}62859461\allowbreak{}a846f4d3\allowbreak{}090bec68\allowbreak{}74de8d40\allowbreak{}4dc28e4c\allowbreak{}85c00395\allowbreak{}24ebce3f\\
Lector & 82db73e1\allowbreak{}3fb33b93\allowbreak{}832d00a5\allowbreak{}d3ffae2a\allowbreak{}3ce893d4\allowbreak{}059d410f\allowbreak{}530239bd\allowbreak{}20944a08\\
Registro en 108 & \textemdash{}\\
\end{tabularx}\endgroup\par\vspace{6pt}
\Needspace{87pt}\phantomsection\label{trace-208}\noindent\hyperref[nav:vi-ruta-208]{\textbf{r6c7\_\allowbreak{}h+1v+1}}\quad Ordinal unificado / memoria: 208 / 205\par\nopagebreak
\begingroup\fontsize{8.3}{10.1}\selectfont\setlength{\tabcolsep}{4pt}\noindent\begin{tabularx}{\linewidth}{@{}p{.21\linewidth}X@{}}
Genealogía & 44eb46a7\allowbreak{}747c4aef\allowbreak{}1e300e34\allowbreak{}4b95cc35\allowbreak{}662ea5f5\allowbreak{}e7be767c\allowbreak{}3e9e1d8d\allowbreak{}753058fc\\
Lector & d37dee38\allowbreak{}8ad27b72\allowbreak{}a09d85bd\allowbreak{}19587c61\allowbreak{}3c7ccaf6\allowbreak{}048eeba3\allowbreak{}53d31d5f\allowbreak{}1b642289\\
Registro en 108 & \textemdash{}\\
\end{tabularx}\endgroup\par\vspace{6pt}
\Needspace{87pt}\phantomsection\label{trace-209}\noindent\hyperref[nav:vi-ruta-209]{\textbf{r6c8\_\allowbreak{}h-1v-1}}\quad Ordinal unificado / memoria: 209 / 212\par\nopagebreak
\begingroup\fontsize{8.3}{10.1}\selectfont\setlength{\tabcolsep}{4pt}\noindent\begin{tabularx}{\linewidth}{@{}p{.21\linewidth}X@{}}
Genealogía & 46a7097d\allowbreak{}2858df91\allowbreak{}76ba9586\allowbreak{}b269471d\allowbreak{}d295cb01\allowbreak{}1c04f97e\allowbreak{}e1f4da54\allowbreak{}cdd7a43e\\
Lector & 1a456767\allowbreak{}b18a5ec9\allowbreak{}86b3b3fb\allowbreak{}3c385626\allowbreak{}7b3045e9\allowbreak{}d69a809c\allowbreak{}24d95845\allowbreak{}977eb71f\\
Registro en 108 & \textemdash{}\\
\end{tabularx}\endgroup\par\vspace{6pt}
\Needspace{87pt}\phantomsection\label{trace-210}\noindent\hyperref[nav:vi-ruta-210]{\textbf{r6c8\_\allowbreak{}h-1v+1}}\quad Ordinal unificado / memoria: 210 / 211\par\nopagebreak
\begingroup\fontsize{8.3}{10.1}\selectfont\setlength{\tabcolsep}{4pt}\noindent\begin{tabularx}{\linewidth}{@{}p{.21\linewidth}X@{}}
Genealogía & c662417a\allowbreak{}f4f02cbb\allowbreak{}24da0df8\allowbreak{}ddc65daf\allowbreak{}21c785dc\allowbreak{}da7367f4\allowbreak{}99a615d5\allowbreak{}f1c13acc\\
Lector & 0bf138f4\allowbreak{}01402e50\allowbreak{}7cb05c52\allowbreak{}92252e90\allowbreak{}d8a25126\allowbreak{}3e471524\allowbreak{}04d52b83\allowbreak{}1e9fce30\\
Registro en 108 & \textemdash{}\\
\end{tabularx}\endgroup\par\vspace{6pt}
\Needspace{87pt}\phantomsection\label{trace-211}\noindent\hyperref[nav:vi-ruta-211]{\textbf{r6c8\_\allowbreak{}h+1v-1}}\quad Ordinal unificado / memoria: 211 / 210\par\nopagebreak
\begingroup\fontsize{8.3}{10.1}\selectfont\setlength{\tabcolsep}{4pt}\noindent\begin{tabularx}{\linewidth}{@{}p{.21\linewidth}X@{}}
Genealogía & 6244f8a5\allowbreak{}a54cd67b\allowbreak{}bef735c4\allowbreak{}354e23b0\allowbreak{}7ca1ae58\allowbreak{}e21214f4\allowbreak{}a6ff3516\allowbreak{}629e19c1\\
Lector & 65c53fc3\allowbreak{}92fa3dbf\allowbreak{}000bcb9c\allowbreak{}9539167b\allowbreak{}8ce08130\allowbreak{}0211ef00\allowbreak{}c2eee82e\allowbreak{}0d9ff65c\\
Registro en 108 & \textemdash{}\\
\end{tabularx}\endgroup\par\vspace{6pt}
\Needspace{87pt}\phantomsection\label{trace-212}\noindent\hyperref[nav:vi-ruta-212]{\textbf{r6c8\_\allowbreak{}h+1v+1}}\quad Ordinal unificado / memoria: 212 / 209\par\nopagebreak
\begingroup\fontsize{8.3}{10.1}\selectfont\setlength{\tabcolsep}{4pt}\noindent\begin{tabularx}{\linewidth}{@{}p{.21\linewidth}X@{}}
Genealogía & 4f813646\allowbreak{}1cebd5d0\allowbreak{}e1aae0b9\allowbreak{}4a131095\allowbreak{}2fec8ed3\allowbreak{}2458cf20\allowbreak{}ae000c9a\allowbreak{}d96abc39\\
Lector & 478ca35d\allowbreak{}f2cf0d45\allowbreak{}64142043\allowbreak{}734baaef\allowbreak{}374d7159\allowbreak{}a3c7775c\allowbreak{}b731a590\allowbreak{}a615a0db\\
Registro en 108 & ec0c7b15\allowbreak{}46af4e29\allowbreak{}3324a0e6\allowbreak{}fc89aa4c\allowbreak{}d160e3c7\allowbreak{}0eacc9a9\allowbreak{}b3494baf\allowbreak{}329435e3\\
\end{tabularx}\endgroup\par\vspace{6pt}
\Needspace{87pt}\phantomsection\label{trace-213}\noindent\hyperref[nav:vi-ruta-213]{\textbf{r6c9\_\allowbreak{}h-1v-1}}\quad Ordinal unificado / memoria: 213 / 216\par\nopagebreak
\begingroup\fontsize{8.3}{10.1}\selectfont\setlength{\tabcolsep}{4pt}\noindent\begin{tabularx}{\linewidth}{@{}p{.21\linewidth}X@{}}
Genealogía & d29410c3\allowbreak{}67743e16\allowbreak{}11b951c1\allowbreak{}79305418\allowbreak{}a0c07870\allowbreak{}34e64248\allowbreak{}8d40b842\allowbreak{}1a4f1b20\\
Lector & 05a266cc\allowbreak{}e4affce3\allowbreak{}09ae74ef\allowbreak{}598dad50\allowbreak{}69603262\allowbreak{}543c2e64\allowbreak{}d874382a\allowbreak{}83755978\\
Registro en 108 & \textemdash{}\\
\end{tabularx}\endgroup\par\vspace{6pt}
\Needspace{87pt}\phantomsection\label{trace-214}\noindent\hyperref[nav:vi-ruta-214]{\textbf{r6c9\_\allowbreak{}h-1v+1}}\quad Ordinal unificado / memoria: 214 / 215\par\nopagebreak
\begingroup\fontsize{8.3}{10.1}\selectfont\setlength{\tabcolsep}{4pt}\noindent\begin{tabularx}{\linewidth}{@{}p{.21\linewidth}X@{}}
Genealogía & 383fc3c8\allowbreak{}eab5e689\allowbreak{}c371ebbb\allowbreak{}9e5a07b1\allowbreak{}b4c5c0c8\allowbreak{}69c53d1e\allowbreak{}34f80c09\allowbreak{}66fe4683\\
Lector & eca848d3\allowbreak{}437105f4\allowbreak{}e3992a55\allowbreak{}c435b0eb\allowbreak{}b5e5b25c\allowbreak{}7918788f\allowbreak{}ba77fb97\allowbreak{}fa69eb3e\\
Registro en 108 & \textemdash{}\\
\end{tabularx}\endgroup\par\vspace{6pt}
\Needspace{87pt}\phantomsection\label{trace-215}\noindent\hyperref[nav:vi-ruta-215]{\textbf{r6c9\_\allowbreak{}h+1v-1}}\quad Ordinal unificado / memoria: 215 / 214\par\nopagebreak
\begingroup\fontsize{8.3}{10.1}\selectfont\setlength{\tabcolsep}{4pt}\noindent\begin{tabularx}{\linewidth}{@{}p{.21\linewidth}X@{}}
Genealogía & 5f4e2b04\allowbreak{}f8b80b45\allowbreak{}dc0a27e9\allowbreak{}68364702\allowbreak{}c50f95e4\allowbreak{}243d4268\allowbreak{}177f1092\allowbreak{}8526aee6\\
Lector & cfa35761\allowbreak{}54e8cf22\allowbreak{}0eeb6ace\allowbreak{}24ec313e\allowbreak{}c4d442e4\allowbreak{}d1da02a4\allowbreak{}bf007f30\allowbreak{}bd2c6621\\
Registro en 108 & 80250ff7\allowbreak{}621218a0\allowbreak{}3a2e2ca3\allowbreak{}9a7a6e60\allowbreak{}556d81f6\allowbreak{}2485cf81\allowbreak{}b5b4de27\allowbreak{}6a81d550\\
\end{tabularx}\endgroup\par\vspace{6pt}
\Needspace{87pt}\phantomsection\label{trace-216}\noindent\hyperref[nav:vi-ruta-216]{\textbf{r6c9\_\allowbreak{}h+1v+1}}\quad Ordinal unificado / memoria: 216 / 213\par\nopagebreak
\begingroup\fontsize{8.3}{10.1}\selectfont\setlength{\tabcolsep}{4pt}\noindent\begin{tabularx}{\linewidth}{@{}p{.21\linewidth}X@{}}
Genealogía & 79b913bc\allowbreak{}6b44d4d2\allowbreak{}00a69fd9\allowbreak{}2eecc993\allowbreak{}6a66c4d2\allowbreak{}89739af3\allowbreak{}88dbe597\allowbreak{}c8ed9751\\
Lector & 32027225\allowbreak{}b3900396\allowbreak{}db67d969\allowbreak{}92263b16\allowbreak{}a9b0766e\allowbreak{}b91a7c93\allowbreak{}9ac99ad9\allowbreak{}a97d3bf9\\
Registro en 108 & \textemdash{}\\
\end{tabularx}\endgroup\par\vspace{6pt}
\Needspace{87pt}\phantomsection\label{trace-217}\noindent\hyperref[nav:vi-ruta-217]{\textbf{r7c1\_\allowbreak{}h-1v-1}}\quad Ordinal unificado / memoria: 217 / 220\par\nopagebreak
\begingroup\fontsize{8.3}{10.1}\selectfont\setlength{\tabcolsep}{4pt}\noindent\begin{tabularx}{\linewidth}{@{}p{.21\linewidth}X@{}}
Genealogía & 213e5026\allowbreak{}6085d4cf\allowbreak{}0de7fe55\allowbreak{}83f14cd6\allowbreak{}e9fd4d82\allowbreak{}01a2fca7\allowbreak{}bbb3795d\allowbreak{}68988de6\\
Lector & bcca332d\allowbreak{}749bd269\allowbreak{}e02e92f3\allowbreak{}30430569\allowbreak{}c4525cd5\allowbreak{}e249720d\allowbreak{}b77549f9\allowbreak{}66a52809\\
Registro en 108 & \textemdash{}\\
\end{tabularx}\endgroup\par\vspace{6pt}
\Needspace{87pt}\phantomsection\label{trace-218}\noindent\hyperref[nav:vi-ruta-218]{\textbf{r7c1\_\allowbreak{}h-1v+1}}\quad Ordinal unificado / memoria: 218 / 219\par\nopagebreak
\begingroup\fontsize{8.3}{10.1}\selectfont\setlength{\tabcolsep}{4pt}\noindent\begin{tabularx}{\linewidth}{@{}p{.21\linewidth}X@{}}
Genealogía & d62df3d9\allowbreak{}6add1bcb\allowbreak{}3616655b\allowbreak{}4c5ab36d\allowbreak{}d6a959ad\allowbreak{}7680ed61\allowbreak{}f0bf5670\allowbreak{}35e46c10\\
Lector & 04f39325\allowbreak{}01d1ad67\allowbreak{}facab5d8\allowbreak{}2019498a\allowbreak{}f3313df6\allowbreak{}e346a07a\allowbreak{}dcb84e7c\allowbreak{}755faee0\\
Registro en 108 & c10b18e3\allowbreak{}f00c1efe\allowbreak{}ab677c68\allowbreak{}4d1a7bd0\allowbreak{}05d2ee4a\allowbreak{}89b23fa3\allowbreak{}bc838923\allowbreak{}718a08fa\\
\end{tabularx}\endgroup\par\vspace{6pt}
\Needspace{87pt}\phantomsection\label{trace-219}\noindent\hyperref[nav:vi-ruta-219]{\textbf{r7c1\_\allowbreak{}h+1v-1}}\quad Ordinal unificado / memoria: 219 / 218\par\nopagebreak
\begingroup\fontsize{8.3}{10.1}\selectfont\setlength{\tabcolsep}{4pt}\noindent\begin{tabularx}{\linewidth}{@{}p{.21\linewidth}X@{}}
Genealogía & ea1bc7bb\allowbreak{}3c188b9a\allowbreak{}7e9dd91a\allowbreak{}90b29f4b\allowbreak{}3fe86da7\allowbreak{}70ec0c1a\allowbreak{}40183a53\allowbreak{}264492b5\\
Lector & ad40326e\allowbreak{}62bbf34e\allowbreak{}abd69280\allowbreak{}6fb96c19\allowbreak{}a8fd8520\allowbreak{}056d702c\allowbreak{}51a3c115\allowbreak{}dc8ce1fb\\
Registro en 108 & \textemdash{}\\
\end{tabularx}\endgroup\par\vspace{6pt}
\Needspace{87pt}\phantomsection\label{trace-220}\noindent\hyperref[nav:vi-ruta-220]{\textbf{r7c1\_\allowbreak{}h+1v+1}}\quad Ordinal unificado / memoria: 220 / 217\par\nopagebreak
\begingroup\fontsize{8.3}{10.1}\selectfont\setlength{\tabcolsep}{4pt}\noindent\begin{tabularx}{\linewidth}{@{}p{.21\linewidth}X@{}}
Genealogía & 3a495db2\allowbreak{}e5a0a670\allowbreak{}7f8a97f7\allowbreak{}24743b33\allowbreak{}c5696f6e\allowbreak{}a59a145f\allowbreak{}59858765\allowbreak{}279c3c3b\\
Lector & 22a43ddf\allowbreak{}f0542b8b\allowbreak{}aedc3fef\allowbreak{}b41c9e08\allowbreak{}8a990090\allowbreak{}a8a95fc8\allowbreak{}d101d973\allowbreak{}f3a8fed9\\
Registro en 108 & \textemdash{}\\
\end{tabularx}\endgroup\par\vspace{6pt}
\Needspace{87pt}\phantomsection\label{trace-221}\noindent\hyperref[nav:vi-ruta-221]{\textbf{r7c2\_\allowbreak{}h-1v-1}}\quad Ordinal unificado / memoria: 221 / 224\par\nopagebreak
\begingroup\fontsize{8.3}{10.1}\selectfont\setlength{\tabcolsep}{4pt}\noindent\begin{tabularx}{\linewidth}{@{}p{.21\linewidth}X@{}}
Genealogía & de97a581\allowbreak{}86c1b724\allowbreak{}ae969ef3\allowbreak{}d0b46d0b\allowbreak{}58439e4a\allowbreak{}07374877\allowbreak{}72214764\allowbreak{}5eed901a\\
Lector & b5cbdb2a\allowbreak{}cc289475\allowbreak{}e1adbffd\allowbreak{}a63c9cb0\allowbreak{}8aa1f829\allowbreak{}e2f904b2\allowbreak{}bb353276\allowbreak{}a27933bf\\
Registro en 108 & \textemdash{}\\
\end{tabularx}\endgroup\par\vspace{6pt}
\Needspace{87pt}\phantomsection\label{trace-222}\noindent\hyperref[nav:vi-ruta-222]{\textbf{r7c2\_\allowbreak{}h-1v+1}}\quad Ordinal unificado / memoria: 222 / 223\par\nopagebreak
\begingroup\fontsize{8.3}{10.1}\selectfont\setlength{\tabcolsep}{4pt}\noindent\begin{tabularx}{\linewidth}{@{}p{.21\linewidth}X@{}}
Genealogía & d0d35604\allowbreak{}38d49301\allowbreak{}7163b00e\allowbreak{}45d39358\allowbreak{}96eb8554\allowbreak{}e0dddda5\allowbreak{}b8a0856c\allowbreak{}c7a13b34\\
Lector & 1de0ff8e\allowbreak{}3d24aca9\allowbreak{}0f7f4d43\allowbreak{}2eafa71c\allowbreak{}5d5aab1b\allowbreak{}c65d749f\allowbreak{}8cd1dc60\allowbreak{}2a5a7a65\\
Registro en 108 & c97a50bb\allowbreak{}057d87c6\allowbreak{}a9915463\allowbreak{}f589846d\allowbreak{}13ce5d25\allowbreak{}a64e07f2\allowbreak{}625114cb\allowbreak{}93e204f5\\
\end{tabularx}\endgroup\par\vspace{6pt}
\Needspace{87pt}\phantomsection\label{trace-223}\noindent\hyperref[nav:vi-ruta-223]{\textbf{r7c2\_\allowbreak{}h+1v-1}}\quad Ordinal unificado / memoria: 223 / 222\par\nopagebreak
\begingroup\fontsize{8.3}{10.1}\selectfont\setlength{\tabcolsep}{4pt}\noindent\begin{tabularx}{\linewidth}{@{}p{.21\linewidth}X@{}}
Genealogía & a8a820b2\allowbreak{}e3341e50\allowbreak{}8860df8f\allowbreak{}a52e8ba6\allowbreak{}be5fca90\allowbreak{}c44b0b91\allowbreak{}8cfb3e75\allowbreak{}fd28d176\\
Lector & 3c31186b\allowbreak{}3fac2aa2\allowbreak{}f596f047\allowbreak{}ea1756c7\allowbreak{}34f0ddc5\allowbreak{}9694e37a\allowbreak{}98f47d9a\allowbreak{}3bdc8475\\
Registro en 108 & \textemdash{}\\
\end{tabularx}\endgroup\par\vspace{6pt}
\Needspace{87pt}\phantomsection\label{trace-224}\noindent\hyperref[nav:vi-ruta-224]{\textbf{r7c2\_\allowbreak{}h+1v+1}}\quad Ordinal unificado / memoria: 224 / 221\par\nopagebreak
\begingroup\fontsize{8.3}{10.1}\selectfont\setlength{\tabcolsep}{4pt}\noindent\begin{tabularx}{\linewidth}{@{}p{.21\linewidth}X@{}}
Genealogía & 3bac38d1\allowbreak{}4f028f87\allowbreak{}4e98149b\allowbreak{}7274261b\allowbreak{}f05c7219\allowbreak{}f5381cfc\allowbreak{}4cbad9f6\allowbreak{}c1c850b9\\
Lector & d0f373ba\allowbreak{}b545802e\allowbreak{}ad72e058\allowbreak{}f777728e\allowbreak{}20bf2d03\allowbreak{}ce412db3\allowbreak{}853130c4\allowbreak{}b1a38f22\\
Registro en 108 & \textemdash{}\\
\end{tabularx}\endgroup\par\vspace{6pt}
\Needspace{87pt}\phantomsection\label{trace-225}\noindent\hyperref[nav:vi-ruta-225]{\textbf{r7c3\_\allowbreak{}h-1v-1}}\quad Ordinal unificado / memoria: 225 / 228\par\nopagebreak
\begingroup\fontsize{8.3}{10.1}\selectfont\setlength{\tabcolsep}{4pt}\noindent\begin{tabularx}{\linewidth}{@{}p{.21\linewidth}X@{}}
Genealogía & 50dd3ab8\allowbreak{}5f947dea\allowbreak{}000293ba\allowbreak{}18e7e5d2\allowbreak{}36ed2a24\allowbreak{}82880fce\allowbreak{}8a55b876\allowbreak{}402a96d6\\
Lector & f8cfb75a\allowbreak{}9cb25ac6\allowbreak{}b8cbb42e\allowbreak{}3fc084e2\allowbreak{}33793934\allowbreak{}4c2a9905\allowbreak{}41bae640\allowbreak{}dcee3531\\
Registro en 108 & \textemdash{}\\
\end{tabularx}\endgroup\par\vspace{6pt}
\Needspace{87pt}\phantomsection\label{trace-226}\noindent\hyperref[nav:vi-ruta-226]{\textbf{r7c3\_\allowbreak{}h-1v+1}}\quad Ordinal unificado / memoria: 226 / 227\par\nopagebreak
\begingroup\fontsize{8.3}{10.1}\selectfont\setlength{\tabcolsep}{4pt}\noindent\begin{tabularx}{\linewidth}{@{}p{.21\linewidth}X@{}}
Genealogía & 97b0c7b8\allowbreak{}352c824e\allowbreak{}a503a1ec\allowbreak{}527396cf\allowbreak{}c0d0b449\allowbreak{}00be367d\allowbreak{}d4b03945\allowbreak{}006857bf\\
Lector & b7b4d753\allowbreak{}04d95a6a\allowbreak{}60aba238\allowbreak{}61edd111\allowbreak{}2aad3aed\allowbreak{}5ee0e967\allowbreak{}3c71017c\allowbreak{}cc3594bd\\
Registro en 108 & \textemdash{}\\
\end{tabularx}\endgroup\par\vspace{6pt}
\Needspace{87pt}\phantomsection\label{trace-227}\noindent\hyperref[nav:vi-ruta-227]{\textbf{r7c3\_\allowbreak{}h+1v-1}}\quad Ordinal unificado / memoria: 227 / 226\par\nopagebreak
\begingroup\fontsize{8.3}{10.1}\selectfont\setlength{\tabcolsep}{4pt}\noindent\begin{tabularx}{\linewidth}{@{}p{.21\linewidth}X@{}}
Genealogía & 34dff485\allowbreak{}32c73a28\allowbreak{}0d1550ed\allowbreak{}80b44a2b\allowbreak{}af487d48\allowbreak{}6dd06bec\allowbreak{}76436e8e\allowbreak{}94d65f93\\
Lector & 502b383c\allowbreak{}ae291d41\allowbreak{}508f616a\allowbreak{}c645f9bc\allowbreak{}def36a9a\allowbreak{}45477b31\allowbreak{}c70d5cb7\allowbreak{}ee016951\\
Registro en 108 & \textemdash{}\\
\end{tabularx}\endgroup\par\vspace{6pt}
\Needspace{87pt}\phantomsection\label{trace-228}\noindent\hyperref[nav:vi-ruta-228]{\textbf{r7c3\_\allowbreak{}h+1v+1}}\quad Ordinal unificado / memoria: 228 / 225\par\nopagebreak
\begingroup\fontsize{8.3}{10.1}\selectfont\setlength{\tabcolsep}{4pt}\noindent\begin{tabularx}{\linewidth}{@{}p{.21\linewidth}X@{}}
Genealogía & 2c8a0ec9\allowbreak{}e3c50967\allowbreak{}fd90be0a\allowbreak{}53359711\allowbreak{}cda37fe9\allowbreak{}6004fa0d\allowbreak{}d2d2c109\allowbreak{}d3ff8203\\
Lector & e7e5c8dd\allowbreak{}a4d5d1f9\allowbreak{}b878e9cd\allowbreak{}1a1c64fb\allowbreak{}8345513e\allowbreak{}b16b32b9\allowbreak{}1933ab57\allowbreak{}1fbdbb5d\\
Registro en 108 & d232cdd0\allowbreak{}8b31d355\allowbreak{}345b6740\allowbreak{}2bf59eda\allowbreak{}9ec2f460\allowbreak{}c0b2fde0\allowbreak{}aabb2dc8\allowbreak{}b105a72c\\
\end{tabularx}\endgroup\par\vspace{6pt}
\Needspace{87pt}\phantomsection\label{trace-229}\noindent\hyperref[nav:vi-ruta-229]{\textbf{r7c4\_\allowbreak{}h-1v-1}}\quad Ordinal unificado / memoria: 229 / 232\par\nopagebreak
\begingroup\fontsize{8.3}{10.1}\selectfont\setlength{\tabcolsep}{4pt}\noindent\begin{tabularx}{\linewidth}{@{}p{.21\linewidth}X@{}}
Genealogía & 61986b61\allowbreak{}75f7b937\allowbreak{}67657333\allowbreak{}7389fee3\allowbreak{}569c3135\allowbreak{}5092cf94\allowbreak{}7843a9d9\allowbreak{}30ba1bc2\\
Lector & 5924cc1f\allowbreak{}56504b4c\allowbreak{}33f30bad\allowbreak{}d28f71c2\allowbreak{}26336afe\allowbreak{}924e0781\allowbreak{}01dceaa8\allowbreak{}e836decf\\
Registro en 108 & \textemdash{}\\
\end{tabularx}\endgroup\par\vspace{6pt}
\Needspace{87pt}\phantomsection\label{trace-230}\noindent\hyperref[nav:vi-ruta-230]{\textbf{r7c4\_\allowbreak{}h-1v+1}}\quad Ordinal unificado / memoria: 230 / 231\par\nopagebreak
\begingroup\fontsize{8.3}{10.1}\selectfont\setlength{\tabcolsep}{4pt}\noindent\begin{tabularx}{\linewidth}{@{}p{.21\linewidth}X@{}}
Genealogía & c59fd018\allowbreak{}875bb472\allowbreak{}7075a352\allowbreak{}424cec29\allowbreak{}d327972c\allowbreak{}e0436e0d\allowbreak{}6620d4c5\allowbreak{}81058224\\
Lector & e049714f\allowbreak{}fb2d3eea\allowbreak{}8b496d14\allowbreak{}6d72cdeb\allowbreak{}379c2764\allowbreak{}679bbf91\allowbreak{}37027cba\allowbreak{}dddd41f3\\
Registro en 108 & c7e0c373\allowbreak{}385cd688\allowbreak{}c8ff7aaf\allowbreak{}d12c63d3\allowbreak{}af22800c\allowbreak{}9ae593a7\allowbreak{}a3bc5952\allowbreak{}d7512202\\
\end{tabularx}\endgroup\par\vspace{6pt}
\Needspace{87pt}\phantomsection\label{trace-231}\noindent\hyperref[nav:vi-ruta-231]{\textbf{r7c4\_\allowbreak{}h+1v-1}}\quad Ordinal unificado / memoria: 231 / 230\par\nopagebreak
\begingroup\fontsize{8.3}{10.1}\selectfont\setlength{\tabcolsep}{4pt}\noindent\begin{tabularx}{\linewidth}{@{}p{.21\linewidth}X@{}}
Genealogía & 7a6c6470\allowbreak{}c763c817\allowbreak{}71cd0d88\allowbreak{}89250edb\allowbreak{}8a8675dd\allowbreak{}0391a21b\allowbreak{}31385cba\allowbreak{}4993e652\\
Lector & be05cd14\allowbreak{}c4f08902\allowbreak{}7d45446a\allowbreak{}e145d5bb\allowbreak{}d383576b\allowbreak{}0d091e16\allowbreak{}60097007\allowbreak{}ebf4fc04\\
Registro en 108 & \textemdash{}\\
\end{tabularx}\endgroup\par\vspace{6pt}
\Needspace{87pt}\phantomsection\label{trace-232}\noindent\hyperref[nav:vi-ruta-232]{\textbf{r7c4\_\allowbreak{}h+1v+1}}\quad Ordinal unificado / memoria: 232 / 229\par\nopagebreak
\begingroup\fontsize{8.3}{10.1}\selectfont\setlength{\tabcolsep}{4pt}\noindent\begin{tabularx}{\linewidth}{@{}p{.21\linewidth}X@{}}
Genealogía & 3b0c0dbf\allowbreak{}a4786567\allowbreak{}c9e927a3\allowbreak{}d9255385\allowbreak{}f88a8bf3\allowbreak{}e4ee10b3\allowbreak{}9a899004\allowbreak{}5dbb8251\\
Lector & 7e34fe3b\allowbreak{}d8853b36\allowbreak{}3e1c7195\allowbreak{}129b813c\allowbreak{}582bbfbb\allowbreak{}75ebd665\allowbreak{}1f739c23\allowbreak{}d7b1b85a\\
Registro en 108 & \textemdash{}\\
\end{tabularx}\endgroup\par\vspace{6pt}
\Needspace{87pt}\phantomsection\label{trace-233}\noindent\hyperref[nav:vi-ruta-233]{\textbf{r7c5\_\allowbreak{}h-1v-1}}\quad Ordinal unificado / memoria: 233 / 236\par\nopagebreak
\begingroup\fontsize{8.3}{10.1}\selectfont\setlength{\tabcolsep}{4pt}\noindent\begin{tabularx}{\linewidth}{@{}p{.21\linewidth}X@{}}
Genealogía & f6ed0ad9\allowbreak{}6ec6366e\allowbreak{}aed3b499\allowbreak{}14434988\allowbreak{}2d3ffcd7\allowbreak{}cd442845\allowbreak{}10c32708\allowbreak{}21d0056c\\
Lector & 25d3f4dd\allowbreak{}b064b47c\allowbreak{}9a608e17\allowbreak{}a712bb09\allowbreak{}d854df59\allowbreak{}edc3903e\allowbreak{}31a2673b\allowbreak{}50b1ae0a\\
Registro en 108 & 68cfa463\allowbreak{}5763ee62\allowbreak{}0dde045d\allowbreak{}3406cd27\allowbreak{}b8ada4eb\allowbreak{}740d3dd0\allowbreak{}2cfcde7e\allowbreak{}4775e236\\
\end{tabularx}\endgroup\par\vspace{6pt}
\Needspace{87pt}\phantomsection\label{trace-234}\noindent\hyperref[nav:vi-ruta-234]{\textbf{r7c5\_\allowbreak{}h-1v+1}}\quad Ordinal unificado / memoria: 234 / 235\par\nopagebreak
\begingroup\fontsize{8.3}{10.1}\selectfont\setlength{\tabcolsep}{4pt}\noindent\begin{tabularx}{\linewidth}{@{}p{.21\linewidth}X@{}}
Genealogía & 99ab406e\allowbreak{}cdac6329\allowbreak{}ab5ce5eb\allowbreak{}6325504e\allowbreak{}cb11f2d0\allowbreak{}4c9d8b34\allowbreak{}6db7499d\allowbreak{}7f567b1e\\
Lector & de03114b\allowbreak{}a14ea466\allowbreak{}2b556a0e\allowbreak{}8594ed8b\allowbreak{}86378685\allowbreak{}d02ee66d\allowbreak{}90537006\allowbreak{}e76e8ab8\\
Registro en 108 & \textemdash{}\\
\end{tabularx}\endgroup\par\vspace{6pt}
\Needspace{87pt}\phantomsection\label{trace-235}\noindent\hyperref[nav:vi-ruta-235]{\textbf{r7c5\_\allowbreak{}h+1v-1}}\quad Ordinal unificado / memoria: 235 / 234\par\nopagebreak
\begingroup\fontsize{8.3}{10.1}\selectfont\setlength{\tabcolsep}{4pt}\noindent\begin{tabularx}{\linewidth}{@{}p{.21\linewidth}X@{}}
Genealogía & 0386e9c2\allowbreak{}055bfd8a\allowbreak{}b8494edc\allowbreak{}4b8a16f0\allowbreak{}0671f216\allowbreak{}a4618044\allowbreak{}24aa45aa\allowbreak{}b64b22d5\\
Lector & 66b69197\allowbreak{}86b672f0\allowbreak{}7cdf87be\allowbreak{}4aefcb18\allowbreak{}b9cac24f\allowbreak{}a6eddfae\allowbreak{}27e4ddf2\allowbreak{}6caba935\\
Registro en 108 & \textemdash{}\\
\end{tabularx}\endgroup\par\vspace{6pt}
\Needspace{87pt}\phantomsection\label{trace-236}\noindent\hyperref[nav:vi-ruta-236]{\textbf{r7c5\_\allowbreak{}h+1v+1}}\quad Ordinal unificado / memoria: 236 / 233\par\nopagebreak
\begingroup\fontsize{8.3}{10.1}\selectfont\setlength{\tabcolsep}{4pt}\noindent\begin{tabularx}{\linewidth}{@{}p{.21\linewidth}X@{}}
Genealogía & b7a1b0a7\allowbreak{}f7d9328f\allowbreak{}a3ce85e1\allowbreak{}fff27ea9\allowbreak{}94c202f1\allowbreak{}387e665f\allowbreak{}13646fe3\allowbreak{}9a6f9594\\
Lector & d7423e62\allowbreak{}f53bddcb\allowbreak{}400b3cf3\allowbreak{}d5485932\allowbreak{}f1df7a7e\allowbreak{}db0186d4\allowbreak{}e92c189a\allowbreak{}60b8a179\\
Registro en 108 & \textemdash{}\\
\end{tabularx}\endgroup\par\vspace{6pt}
\Needspace{87pt}\phantomsection\label{trace-237}\noindent\hyperref[nav:vi-ruta-237]{\textbf{r7c6\_\allowbreak{}h-1v-1}}\quad Ordinal unificado / memoria: 237 / 240\par\nopagebreak
\begingroup\fontsize{8.3}{10.1}\selectfont\setlength{\tabcolsep}{4pt}\noindent\begin{tabularx}{\linewidth}{@{}p{.21\linewidth}X@{}}
Genealogía & bdd0f510\allowbreak{}981aa2bd\allowbreak{}14e11cc9\allowbreak{}e772a40c\allowbreak{}0fa9629c\allowbreak{}6c33a6c8\allowbreak{}63e12854\allowbreak{}3bb3c064\\
Lector & 094ff628\allowbreak{}32260c10\allowbreak{}76787d26\allowbreak{}e6f33ddf\allowbreak{}ede703d5\allowbreak{}dedbf9b8\allowbreak{}6f089121\allowbreak{}79c0158e\\
Registro en 108 & \textemdash{}\\
\end{tabularx}\endgroup\par\vspace{6pt}
\Needspace{87pt}\phantomsection\label{trace-238}\noindent\hyperref[nav:vi-ruta-238]{\textbf{r7c6\_\allowbreak{}h-1v+1}}\quad Ordinal unificado / memoria: 238 / 239\par\nopagebreak
\begingroup\fontsize{8.3}{10.1}\selectfont\setlength{\tabcolsep}{4pt}\noindent\begin{tabularx}{\linewidth}{@{}p{.21\linewidth}X@{}}
Genealogía & 3b43feb2\allowbreak{}9ef42944\allowbreak{}b5c6da5d\allowbreak{}1493e464\allowbreak{}1332fe29\allowbreak{}0cc733c9\allowbreak{}4b555881\allowbreak{}85f490d0\\
Lector & 77fffc5e\allowbreak{}08a1d1d2\allowbreak{}2e2d56a9\allowbreak{}caa4f750\allowbreak{}16d1c692\allowbreak{}61b9bac2\allowbreak{}8a2e8ced\allowbreak{}d44d0ddb\\
Registro en 108 & 4261f837\allowbreak{}a3e329d1\allowbreak{}99831ee7\allowbreak{}64123ea5\allowbreak{}1bb952b2\allowbreak{}5504d3f1\allowbreak{}b96372b3\allowbreak{}44872b11\\
\end{tabularx}\endgroup\par\vspace{6pt}
\Needspace{87pt}\phantomsection\label{trace-239}\noindent\hyperref[nav:vi-ruta-239]{\textbf{r7c6\_\allowbreak{}h+1v-1}}\quad Ordinal unificado / memoria: 239 / 238\par\nopagebreak
\begingroup\fontsize{8.3}{10.1}\selectfont\setlength{\tabcolsep}{4pt}\noindent\begin{tabularx}{\linewidth}{@{}p{.21\linewidth}X@{}}
Genealogía & 231bf7dd\allowbreak{}55d8a438\allowbreak{}f107de8c\allowbreak{}3e140751\allowbreak{}06170fc9\allowbreak{}68fbf990\allowbreak{}a0d8a6f5\allowbreak{}ffd71b41\\
Lector & 573249f7\allowbreak{}14bc3cb9\allowbreak{}8d64797e\allowbreak{}90987918\allowbreak{}abae9806\allowbreak{}80a392c4\allowbreak{}1e65c25c\allowbreak{}1bba6d94\\
Registro en 108 & \textemdash{}\\
\end{tabularx}\endgroup\par\vspace{6pt}
\Needspace{87pt}\phantomsection\label{trace-240}\noindent\hyperref[nav:vi-ruta-240]{\textbf{r7c6\_\allowbreak{}h+1v+1}}\quad Ordinal unificado / memoria: 240 / 237\par\nopagebreak
\begingroup\fontsize{8.3}{10.1}\selectfont\setlength{\tabcolsep}{4pt}\noindent\begin{tabularx}{\linewidth}{@{}p{.21\linewidth}X@{}}
Genealogía & 3e939392\allowbreak{}091ce7df\allowbreak{}7511a0ee\allowbreak{}da421eb0\allowbreak{}299c2ba4\allowbreak{}582a4adb\allowbreak{}f492bef6\allowbreak{}aebffe3f\\
Lector & b6443f7d\allowbreak{}3f0c4492\allowbreak{}2d917e32\allowbreak{}68c496cf\allowbreak{}f4263012\allowbreak{}96ffd17b\allowbreak{}39c8f71c\allowbreak{}592309c7\\
Registro en 108 & \textemdash{}\\
\end{tabularx}\endgroup\par\vspace{6pt}
\Needspace{87pt}\phantomsection\label{trace-241}\noindent\hyperref[nav:vi-ruta-241]{\textbf{r7c7\_\allowbreak{}h-1v-1}}\quad Ordinal unificado / memoria: 241 / 244\par\nopagebreak
\begingroup\fontsize{8.3}{10.1}\selectfont\setlength{\tabcolsep}{4pt}\noindent\begin{tabularx}{\linewidth}{@{}p{.21\linewidth}X@{}}
Genealogía & f7adee5b\allowbreak{}3133fb43\allowbreak{}299c7e7f\allowbreak{}57cfe9d2\allowbreak{}104d1295\allowbreak{}1c24558e\allowbreak{}004b0d10\allowbreak{}46825410\\
Lector & ee9f7a74\allowbreak{}820c460f\allowbreak{}fd5606aa\allowbreak{}30bdddd4\allowbreak{}0bb68754\allowbreak{}27bf381c\allowbreak{}1bb6883d\allowbreak{}6bf2cb2d\\
Registro en 108 & 2b5fcf4e\allowbreak{}e156f8e7\allowbreak{}f3d9ef4b\allowbreak{}48b12761\allowbreak{}781d5874\allowbreak{}0c7f25b3\allowbreak{}50ada993\allowbreak{}fdfb43da\\
\end{tabularx}\endgroup\par\vspace{6pt}
\Needspace{87pt}\phantomsection\label{trace-242}\noindent\hyperref[nav:vi-ruta-242]{\textbf{r7c7\_\allowbreak{}h-1v+1}}\quad Ordinal unificado / memoria: 242 / 243\par\nopagebreak
\begingroup\fontsize{8.3}{10.1}\selectfont\setlength{\tabcolsep}{4pt}\noindent\begin{tabularx}{\linewidth}{@{}p{.21\linewidth}X@{}}
Genealogía & e5a3a65c\allowbreak{}351de531\allowbreak{}521f5eb2\allowbreak{}e8a895c9\allowbreak{}38d33cc4\allowbreak{}b6092ae7\allowbreak{}a236bfde\allowbreak{}174f67e6\\
Lector & f1528715\allowbreak{}99c0f372\allowbreak{}a009df5e\allowbreak{}0bf5d2e8\allowbreak{}9a435b88\allowbreak{}231a923e\allowbreak{}31fd5fc3\allowbreak{}2cbfc7f8\\
Registro en 108 & \textemdash{}\\
\end{tabularx}\endgroup\par\vspace{6pt}
\Needspace{87pt}\phantomsection\label{trace-243}\noindent\hyperref[nav:vi-ruta-243]{\textbf{r7c7\_\allowbreak{}h+1v-1}}\quad Ordinal unificado / memoria: 243 / 242\par\nopagebreak
\begingroup\fontsize{8.3}{10.1}\selectfont\setlength{\tabcolsep}{4pt}\noindent\begin{tabularx}{\linewidth}{@{}p{.21\linewidth}X@{}}
Genealogía & b139139d\allowbreak{}502bd446\allowbreak{}d91e7d84\allowbreak{}ef637882\allowbreak{}9caef959\allowbreak{}c57c068a\allowbreak{}87604b10\allowbreak{}a6c051f8\\
Lector & d3747a41\allowbreak{}b0c5e19b\allowbreak{}f084c3f3\allowbreak{}9ce497e8\allowbreak{}25be56b5\allowbreak{}59a223f0\allowbreak{}0b148f15\allowbreak{}80acb8ed\\
Registro en 108 & \textemdash{}\\
\end{tabularx}\endgroup\par\vspace{6pt}
\Needspace{87pt}\phantomsection\label{trace-244}\noindent\hyperref[nav:vi-ruta-244]{\textbf{r7c7\_\allowbreak{}h+1v+1}}\quad Ordinal unificado / memoria: 244 / 241\par\nopagebreak
\begingroup\fontsize{8.3}{10.1}\selectfont\setlength{\tabcolsep}{4pt}\noindent\begin{tabularx}{\linewidth}{@{}p{.21\linewidth}X@{}}
Genealogía & 97a8e38d\allowbreak{}1e0bc290\allowbreak{}75a19475\allowbreak{}862f706d\allowbreak{}23c2cee0\allowbreak{}e4509a10\allowbreak{}167a3215\allowbreak{}62108b24\\
Lector & 113a5319\allowbreak{}d2bebd39\allowbreak{}24e86aff\allowbreak{}9c2f6618\allowbreak{}18e1d5da\allowbreak{}017520c4\allowbreak{}f5a24875\allowbreak{}0b426e5c\\
Registro en 108 & \textemdash{}\\
\end{tabularx}\endgroup\par\vspace{6pt}
\Needspace{87pt}\phantomsection\label{trace-245}\noindent\hyperref[nav:vi-ruta-245]{\textbf{r7c8\_\allowbreak{}h-1v-1}}\quad Ordinal unificado / memoria: 245 / 248\par\nopagebreak
\begingroup\fontsize{8.3}{10.1}\selectfont\setlength{\tabcolsep}{4pt}\noindent\begin{tabularx}{\linewidth}{@{}p{.21\linewidth}X@{}}
Genealogía & 93e2a454\allowbreak{}a1c5de5f\allowbreak{}4f4e4b0e\allowbreak{}ec1903cc\allowbreak{}6a51d31d\allowbreak{}19260e98\allowbreak{}a023f085\allowbreak{}f717f2c5\\
Lector & 5e1722d0\allowbreak{}1279d898\allowbreak{}8e0e53f3\allowbreak{}7eeed749\allowbreak{}fe741983\allowbreak{}d8db78b0\allowbreak{}9039733d\allowbreak{}0e70de18\\
Registro en 108 & \textemdash{}\\
\end{tabularx}\endgroup\par\vspace{6pt}
\Needspace{87pt}\phantomsection\label{trace-246}\noindent\hyperref[nav:vi-ruta-246]{\textbf{r7c8\_\allowbreak{}h-1v+1}}\quad Ordinal unificado / memoria: 246 / 247\par\nopagebreak
\begingroup\fontsize{8.3}{10.1}\selectfont\setlength{\tabcolsep}{4pt}\noindent\begin{tabularx}{\linewidth}{@{}p{.21\linewidth}X@{}}
Genealogía & ff6632f4\allowbreak{}59518549\allowbreak{}1ccf0bb1\allowbreak{}ee6fb096\allowbreak{}51fc7ee9\allowbreak{}d444f4a3\allowbreak{}6b243492\allowbreak{}1c63e845\\
Lector & c689f380\allowbreak{}d0c3f1f1\allowbreak{}3830f345\allowbreak{}827ac079\allowbreak{}bfef21a0\allowbreak{}290c3319\allowbreak{}a4e040b6\allowbreak{}6452df42\\
Registro en 108 & \textemdash{}\\
\end{tabularx}\endgroup\par\vspace{6pt}
\Needspace{87pt}\phantomsection\label{trace-247}\noindent\hyperref[nav:vi-ruta-247]{\textbf{r7c8\_\allowbreak{}h+1v-1}}\quad Ordinal unificado / memoria: 247 / 246\par\nopagebreak
\begingroup\fontsize{8.3}{10.1}\selectfont\setlength{\tabcolsep}{4pt}\noindent\begin{tabularx}{\linewidth}{@{}p{.21\linewidth}X@{}}
Genealogía & 572bed3e\allowbreak{}c2da43b5\allowbreak{}88e8f03c\allowbreak{}20b85deb\allowbreak{}3a323878\allowbreak{}7c652806\allowbreak{}bdaedb67\allowbreak{}e8a2e65d\\
Lector & f6d6251c\allowbreak{}2a3a0bc5\allowbreak{}1ba9e9d0\allowbreak{}62b640aa\allowbreak{}e53badbc\allowbreak{}a7baa4e8\allowbreak{}7ead4bf0\allowbreak{}ae639226\\
Registro en 108 & 2bc42341\allowbreak{}df8c6224\allowbreak{}4bc513d9\allowbreak{}a683f0a2\allowbreak{}cd6caf97\allowbreak{}ec8b8d54\allowbreak{}ada93d9a\allowbreak{}c2a0f146\\
\end{tabularx}\endgroup\par\vspace{6pt}
\Needspace{87pt}\phantomsection\label{trace-248}\noindent\hyperref[nav:vi-ruta-248]{\textbf{r7c8\_\allowbreak{}h+1v+1}}\quad Ordinal unificado / memoria: 248 / 245\par\nopagebreak
\begingroup\fontsize{8.3}{10.1}\selectfont\setlength{\tabcolsep}{4pt}\noindent\begin{tabularx}{\linewidth}{@{}p{.21\linewidth}X@{}}
Genealogía & 5468c23a\allowbreak{}aa5fe88b\allowbreak{}8dfc71dc\allowbreak{}a843b909\allowbreak{}8b3fad83\allowbreak{}d8074505\allowbreak{}2c981c39\allowbreak{}6d08d09d\\
Lector & e7b0197c\allowbreak{}a2c34747\allowbreak{}ea501fc5\allowbreak{}fc951f0e\allowbreak{}abb2154b\allowbreak{}fe7d1ca6\allowbreak{}32f93107\allowbreak{}1fa14d07\\
Registro en 108 & \textemdash{}\\
\end{tabularx}\endgroup\par\vspace{6pt}
\Needspace{87pt}\phantomsection\label{trace-249}\noindent\hyperref[nav:vi-ruta-249]{\textbf{r7c9\_\allowbreak{}h-1v-1}}\quad Ordinal unificado / memoria: 249 / 252\par\nopagebreak
\begingroup\fontsize{8.3}{10.1}\selectfont\setlength{\tabcolsep}{4pt}\noindent\begin{tabularx}{\linewidth}{@{}p{.21\linewidth}X@{}}
Genealogía & ba5642d5\allowbreak{}96d97f27\allowbreak{}564669dc\allowbreak{}98515c90\allowbreak{}5e1973cd\allowbreak{}764f435e\allowbreak{}e5c853fc\allowbreak{}d049beed\\
Lector & 268ca678\allowbreak{}5b1f8117\allowbreak{}df7ac172\allowbreak{}4f8f4b30\allowbreak{}ac60083e\allowbreak{}dec11026\allowbreak{}41a6611a\allowbreak{}dc1ffb9c\\
Registro en 108 & \textemdash{}\\
\end{tabularx}\endgroup\par\vspace{6pt}
\Needspace{87pt}\phantomsection\label{trace-250}\noindent\hyperref[nav:vi-ruta-250]{\textbf{r7c9\_\allowbreak{}h-1v+1}}\quad Ordinal unificado / memoria: 250 / 251\par\nopagebreak
\begingroup\fontsize{8.3}{10.1}\selectfont\setlength{\tabcolsep}{4pt}\noindent\begin{tabularx}{\linewidth}{@{}p{.21\linewidth}X@{}}
Genealogía & da00f1d5\allowbreak{}55eb9be2\allowbreak{}5c1ceca2\allowbreak{}ebc88d96\allowbreak{}b55a1c8c\allowbreak{}2011788d\allowbreak{}814fa585\allowbreak{}7a504b9f\\
Lector & 183882de\allowbreak{}6128614a\allowbreak{}d3eef875\allowbreak{}eb974e8c\allowbreak{}55c95155\allowbreak{}8c213f7e\allowbreak{}b721348c\allowbreak{}9354f824\\
Registro en 108 & \textemdash{}\\
\end{tabularx}\endgroup\par\vspace{6pt}
\Needspace{87pt}\phantomsection\label{trace-251}\noindent\hyperref[nav:vi-ruta-251]{\textbf{r7c9\_\allowbreak{}h+1v-1}}\quad Ordinal unificado / memoria: 251 / 250\par\nopagebreak
\begingroup\fontsize{8.3}{10.1}\selectfont\setlength{\tabcolsep}{4pt}\noindent\begin{tabularx}{\linewidth}{@{}p{.21\linewidth}X@{}}
Genealogía & ad5ee82a\allowbreak{}87bb9a61\allowbreak{}ce4360a2\allowbreak{}b14bf22c\allowbreak{}5c60f40e\allowbreak{}7e75f290\allowbreak{}ec2dde4d\allowbreak{}0e7c90fc\\
Lector & ed489aa0\allowbreak{}c8e0c09e\allowbreak{}4604ad7f\allowbreak{}7f9ebf49\allowbreak{}ec343a76\allowbreak{}ed5ea60a\allowbreak{}3369f141\allowbreak{}1b16ea84\\
Registro en 108 & \textemdash{}\\
\end{tabularx}\endgroup\par\vspace{6pt}
\Needspace{87pt}\phantomsection\label{trace-252}\noindent\hyperref[nav:vi-ruta-252]{\textbf{r7c9\_\allowbreak{}h+1v+1}}\quad Ordinal unificado / memoria: 252 / 249\par\nopagebreak
\begingroup\fontsize{8.3}{10.1}\selectfont\setlength{\tabcolsep}{4pt}\noindent\begin{tabularx}{\linewidth}{@{}p{.21\linewidth}X@{}}
Genealogía & 88eff8f2\allowbreak{}d71641e5\allowbreak{}54b89b65\allowbreak{}6e288556\allowbreak{}7c18727b\allowbreak{}ad0128f4\allowbreak{}2fa25e7b\allowbreak{}df6a031d\\
Lector & 6c785480\allowbreak{}94bcdfb0\allowbreak{}58a3bcd8\allowbreak{}68ff2cf8\allowbreak{}3e421971\allowbreak{}64de9e58\allowbreak{}aa3dab4b\allowbreak{}8edc6ad7\\
Registro en 108 & 743b8088\allowbreak{}cefbb1b3\allowbreak{}82bb7dc8\allowbreak{}6b4e63b9\allowbreak{}1a876881\allowbreak{}68e208f2\allowbreak{}080f8a04\allowbreak{}0aefcffd\\
\end{tabularx}\endgroup\par\vspace{6pt}
\Needspace{87pt}\phantomsection\label{trace-253}\noindent\hyperref[nav:vi-ruta-253]{\textbf{r8c1\_\allowbreak{}h-1v-1}}\quad Ordinal unificado / memoria: 253 / 256\par\nopagebreak
\begingroup\fontsize{8.3}{10.1}\selectfont\setlength{\tabcolsep}{4pt}\noindent\begin{tabularx}{\linewidth}{@{}p{.21\linewidth}X@{}}
Genealogía & 7f20d693\allowbreak{}558c4c77\allowbreak{}0e8c14a2\allowbreak{}edb53682\allowbreak{}418d3af8\allowbreak{}8e8147e4\allowbreak{}e46a45b5\allowbreak{}84d0af32\\
Lector & d463a304\allowbreak{}37186dc5\allowbreak{}a632dbe1\allowbreak{}fe2eb6c5\allowbreak{}c462cdb5\allowbreak{}8acc7d2d\allowbreak{}24faaa41\allowbreak{}23025187\\
Registro en 108 & \textemdash{}\\
\end{tabularx}\endgroup\par\vspace{6pt}
\Needspace{87pt}\phantomsection\label{trace-254}\noindent\hyperref[nav:vi-ruta-254]{\textbf{r8c1\_\allowbreak{}h-1v+1}}\quad Ordinal unificado / memoria: 254 / 255\par\nopagebreak
\begingroup\fontsize{8.3}{10.1}\selectfont\setlength{\tabcolsep}{4pt}\noindent\begin{tabularx}{\linewidth}{@{}p{.21\linewidth}X@{}}
Genealogía & c88b9358\allowbreak{}0f9451e0\allowbreak{}e63048f3\allowbreak{}f329ae52\allowbreak{}082f45e6\allowbreak{}5a9494ac\allowbreak{}b1942445\allowbreak{}b74dbd7d\\
Lector & 06f382a4\allowbreak{}af0319fc\allowbreak{}434d2d51\allowbreak{}09c92a77\allowbreak{}f6ff4e9e\allowbreak{}b6a30cdd\allowbreak{}9620be4b\allowbreak{}2c25677f\\
Registro en 108 & \textemdash{}\\
\end{tabularx}\endgroup\par\vspace{6pt}
\Needspace{87pt}\phantomsection\label{trace-255}\noindent\hyperref[nav:vi-ruta-255]{\textbf{r8c1\_\allowbreak{}h+1v-1}}\quad Ordinal unificado / memoria: 255 / 254\par\nopagebreak
\begingroup\fontsize{8.3}{10.1}\selectfont\setlength{\tabcolsep}{4pt}\noindent\begin{tabularx}{\linewidth}{@{}p{.21\linewidth}X@{}}
Genealogía & 210c219b\allowbreak{}baaf3bf4\allowbreak{}55eaa92f\allowbreak{}45a017dd\allowbreak{}9f816ca3\allowbreak{}b46f2b08\allowbreak{}b272a9c0\allowbreak{}6868cdc0\\
Lector & 8b66a956\allowbreak{}1a72fc2c\allowbreak{}f9d44aad\allowbreak{}0a4ab928\allowbreak{}0d6872fd\allowbreak{}24ae8ade\allowbreak{}2437e7c7\allowbreak{}1c5da763\\
Registro en 108 & 2abe5fb7\allowbreak{}70abbac7\allowbreak{}2934c668\allowbreak{}102c504b\allowbreak{}01e8b32c\allowbreak{}6bee43d0\allowbreak{}a62e45a1\allowbreak{}cbb3f5cd\\
\end{tabularx}\endgroup\par\vspace{6pt}
\Needspace{87pt}\phantomsection\label{trace-256}\noindent\hyperref[nav:vi-ruta-256]{\textbf{r8c1\_\allowbreak{}h+1v+1}}\quad Ordinal unificado / memoria: 256 / 253\par\nopagebreak
\begingroup\fontsize{8.3}{10.1}\selectfont\setlength{\tabcolsep}{4pt}\noindent\begin{tabularx}{\linewidth}{@{}p{.21\linewidth}X@{}}
Genealogía & aab792bb\allowbreak{}556e0bed\allowbreak{}999eb313\allowbreak{}b4d4b26e\allowbreak{}a63b3bad\allowbreak{}901ed813\allowbreak{}0f94fb1e\allowbreak{}c4554219\\
Lector & 72362335\allowbreak{}60f55142\allowbreak{}5afe01a2\allowbreak{}7858b2d7\allowbreak{}979db398\allowbreak{}3540ff9b\allowbreak{}4e007b9f\allowbreak{}9b5e0d04\\
Registro en 108 & \textemdash{}\\
\end{tabularx}\endgroup\par\vspace{6pt}
\Needspace{87pt}\phantomsection\label{trace-257}\noindent\hyperref[nav:vi-ruta-257]{\textbf{r8c2\_\allowbreak{}h-1v-1}}\quad Ordinal unificado / memoria: 257 / 260\par\nopagebreak
\begingroup\fontsize{8.3}{10.1}\selectfont\setlength{\tabcolsep}{4pt}\noindent\begin{tabularx}{\linewidth}{@{}p{.21\linewidth}X@{}}
Genealogía & 4d09dc43\allowbreak{}af4a5991\allowbreak{}5c1a6038\allowbreak{}f5c475cb\allowbreak{}d194cfeb\allowbreak{}4bf1d5e3\allowbreak{}57df23d9\allowbreak{}754d395f\\
Lector & 4128cbf7\allowbreak{}9b80a7af\allowbreak{}93f10ded\allowbreak{}32d0960e\allowbreak{}b5da56cb\allowbreak{}f47a110b\allowbreak{}ab03f10d\allowbreak{}446c2316\\
Registro en 108 & \textemdash{}\\
\end{tabularx}\endgroup\par\vspace{6pt}
\Needspace{87pt}\phantomsection\label{trace-258}\noindent\hyperref[nav:vi-ruta-258]{\textbf{r8c2\_\allowbreak{}h-1v+1}}\quad Ordinal unificado / memoria: 258 / 259\par\nopagebreak
\begingroup\fontsize{8.3}{10.1}\selectfont\setlength{\tabcolsep}{4pt}\noindent\begin{tabularx}{\linewidth}{@{}p{.21\linewidth}X@{}}
Genealogía & c10ac37f\allowbreak{}6c4294f8\allowbreak{}4a5c4eec\allowbreak{}b9e4d137\allowbreak{}4086b7aa\allowbreak{}e332df72\allowbreak{}101d20af\allowbreak{}f5ff88c8\\
Lector & aa9f45f6\allowbreak{}e48948c3\allowbreak{}8e93e1da\allowbreak{}669a8afd\allowbreak{}c2bfc8e2\allowbreak{}51959f2c\allowbreak{}2d2ac248\allowbreak{}ac767317\\
Registro en 108 & \textemdash{}\\
\end{tabularx}\endgroup\par\vspace{6pt}
\Needspace{87pt}\phantomsection\label{trace-259}\noindent\hyperref[nav:vi-ruta-259]{\textbf{r8c2\_\allowbreak{}h+1v-1}}\quad Ordinal unificado / memoria: 259 / 258\par\nopagebreak
\begingroup\fontsize{8.3}{10.1}\selectfont\setlength{\tabcolsep}{4pt}\noindent\begin{tabularx}{\linewidth}{@{}p{.21\linewidth}X@{}}
Genealogía & 4581338c\allowbreak{}e438211b\allowbreak{}da03b46b\allowbreak{}dd171051\allowbreak{}cbdf4527\allowbreak{}bfe09385\allowbreak{}ab84b222\allowbreak{}943678e0\\
Lector & 30e7d804\allowbreak{}18eef9d9\allowbreak{}188a1367\allowbreak{}74e1e1ba\allowbreak{}1b467b36\allowbreak{}6375105e\allowbreak{}556e2595\allowbreak{}fda118ab\\
Registro en 108 & \textemdash{}\\
\end{tabularx}\endgroup\par\vspace{6pt}
\Needspace{87pt}\phantomsection\label{trace-260}\noindent\hyperref[nav:vi-ruta-260]{\textbf{r8c2\_\allowbreak{}h+1v+1}}\quad Ordinal unificado / memoria: 260 / 257\par\nopagebreak
\begingroup\fontsize{8.3}{10.1}\selectfont\setlength{\tabcolsep}{4pt}\noindent\begin{tabularx}{\linewidth}{@{}p{.21\linewidth}X@{}}
Genealogía & 37e66f95\allowbreak{}874dba75\allowbreak{}02936c06\allowbreak{}25906398\allowbreak{}abd46242\allowbreak{}1242a711\allowbreak{}a1cebe0d\allowbreak{}a9742e20\\
Lector & a066e554\allowbreak{}a21d08ba\allowbreak{}afef2969\allowbreak{}2f0e0500\allowbreak{}bdd85887\allowbreak{}6f9bbc23\allowbreak{}4b52fb9a\allowbreak{}24bea27f\\
Registro en 108 & d7f810dc\allowbreak{}8cc7b129\allowbreak{}8f7e1c92\allowbreak{}b9f22e70\allowbreak{}5a1b79fa\allowbreak{}32aaf5a1\allowbreak{}8405a2e5\allowbreak{}fbd7422f\\
\end{tabularx}\endgroup\par\vspace{6pt}
\Needspace{87pt}\phantomsection\label{trace-261}\noindent\hyperref[nav:vi-ruta-261]{\textbf{r8c3\_\allowbreak{}h-1v-1}}\quad Ordinal unificado / memoria: 261 / 264\par\nopagebreak
\begingroup\fontsize{8.3}{10.1}\selectfont\setlength{\tabcolsep}{4pt}\noindent\begin{tabularx}{\linewidth}{@{}p{.21\linewidth}X@{}}
Genealogía & ad652dbf\allowbreak{}a6c5c813\allowbreak{}518166ce\allowbreak{}23d4443c\allowbreak{}fdf6cc6a\allowbreak{}3f22f5d7\allowbreak{}f723510f\allowbreak{}50c448b8\\
Lector & 3ed48941\allowbreak{}81a15c6c\allowbreak{}37003af6\allowbreak{}a3bffeef\allowbreak{}bd5b6c91\allowbreak{}a65783eb\allowbreak{}e3a13045\allowbreak{}7ab9a8aa\\
Registro en 108 & \textemdash{}\\
\end{tabularx}\endgroup\par\vspace{6pt}
\Needspace{87pt}\phantomsection\label{trace-262}\noindent\hyperref[nav:vi-ruta-262]{\textbf{r8c3\_\allowbreak{}h-1v+1}}\quad Ordinal unificado / memoria: 262 / 263\par\nopagebreak
\begingroup\fontsize{8.3}{10.1}\selectfont\setlength{\tabcolsep}{4pt}\noindent\begin{tabularx}{\linewidth}{@{}p{.21\linewidth}X@{}}
Genealogía & 541a3c42\allowbreak{}518d601b\allowbreak{}58431822\allowbreak{}a3572a2b\allowbreak{}e049c02a\allowbreak{}300aa4fe\allowbreak{}acd5345e\allowbreak{}f948f454\\
Lector & ad8c5049\allowbreak{}5acdfc97\allowbreak{}08f3bf9b\allowbreak{}44a935a2\allowbreak{}051df217\allowbreak{}67326d50\allowbreak{}c52b3dfe\allowbreak{}ae00d3fb\\
Registro en 108 & e9f4ee18\allowbreak{}c9a8859c\allowbreak{}65aa766d\allowbreak{}bde0bd19\allowbreak{}6a242a2f\allowbreak{}3dacd16a\allowbreak{}052d90ac\allowbreak{}cc021dd1\\
\end{tabularx}\endgroup\par\vspace{6pt}
\Needspace{87pt}\phantomsection\label{trace-263}\noindent\hyperref[nav:vi-ruta-263]{\textbf{r8c3\_\allowbreak{}h+1v-1}}\quad Ordinal unificado / memoria: 263 / 262\par\nopagebreak
\begingroup\fontsize{8.3}{10.1}\selectfont\setlength{\tabcolsep}{4pt}\noindent\begin{tabularx}{\linewidth}{@{}p{.21\linewidth}X@{}}
Genealogía & cd28c157\allowbreak{}b1840882\allowbreak{}849d580e\allowbreak{}4bfbce44\allowbreak{}8d65ae5f\allowbreak{}45fe4489\allowbreak{}71e9270d\allowbreak{}acf4fbc6\\
Lector & 6d4c84d9\allowbreak{}7a1f41dd\allowbreak{}bd2bbd27\allowbreak{}d0c2afc7\allowbreak{}7d646b03\allowbreak{}c697190d\allowbreak{}46d00c87\allowbreak{}f92617a1\\
Registro en 108 & \textemdash{}\\
\end{tabularx}\endgroup\par\vspace{6pt}
\Needspace{87pt}\phantomsection\label{trace-264}\noindent\hyperref[nav:vi-ruta-264]{\textbf{r8c3\_\allowbreak{}h+1v+1}}\quad Ordinal unificado / memoria: 264 / 261\par\nopagebreak
\begingroup\fontsize{8.3}{10.1}\selectfont\setlength{\tabcolsep}{4pt}\noindent\begin{tabularx}{\linewidth}{@{}p{.21\linewidth}X@{}}
Genealogía & e3bc25b6\allowbreak{}06504b6c\allowbreak{}91007699\allowbreak{}ff3655c1\allowbreak{}39fba11c\allowbreak{}d7689c98\allowbreak{}c705d3f9\allowbreak{}d12f38cd\\
Lector & fea2f4b6\allowbreak{}8d20b684\allowbreak{}e22997fc\allowbreak{}21a30bb3\allowbreak{}9fa2a3c3\allowbreak{}1dbf3cf0\allowbreak{}fef68224\allowbreak{}900afdb0\\
Registro en 108 & \textemdash{}\\
\end{tabularx}\endgroup\par\vspace{6pt}
\Needspace{87pt}\phantomsection\label{trace-265}\noindent\hyperref[nav:vi-ruta-265]{\textbf{r8c4\_\allowbreak{}h-1v-1}}\quad Ordinal unificado / memoria: 265 / 268\par\nopagebreak
\begingroup\fontsize{8.3}{10.1}\selectfont\setlength{\tabcolsep}{4pt}\noindent\begin{tabularx}{\linewidth}{@{}p{.21\linewidth}X@{}}
Genealogía & d773d886\allowbreak{}b771b4e9\allowbreak{}d58210f9\allowbreak{}1db044e3\allowbreak{}d5fa76dc\allowbreak{}dbba2d14\allowbreak{}4b4aa6a7\allowbreak{}816e2a74\\
Lector & 2129600f\allowbreak{}36965fdc\allowbreak{}11a79112\allowbreak{}6a3ad1e3\allowbreak{}8346cb4d\allowbreak{}0a010c35\allowbreak{}bc87bab5\allowbreak{}442f815f\\
Registro en 108 & \textemdash{}\\
\end{tabularx}\endgroup\par\vspace{6pt}
\Needspace{87pt}\phantomsection\label{trace-266}\noindent\hyperref[nav:vi-ruta-266]{\textbf{r8c4\_\allowbreak{}h-1v+1}}\quad Ordinal unificado / memoria: 266 / 267\par\nopagebreak
\begingroup\fontsize{8.3}{10.1}\selectfont\setlength{\tabcolsep}{4pt}\noindent\begin{tabularx}{\linewidth}{@{}p{.21\linewidth}X@{}}
Genealogía & aa4a5949\allowbreak{}c8430a0e\allowbreak{}c7f118f6\allowbreak{}522a38b2\allowbreak{}41e22dd0\allowbreak{}f172a9ad\allowbreak{}8b2e7d04\allowbreak{}e3373bbb\\
Lector & 00391297\allowbreak{}614e66ca\allowbreak{}2afd3f79\allowbreak{}601a5f0d\allowbreak{}ec107056\allowbreak{}08f06064\allowbreak{}4783235d\allowbreak{}bd3fa513\\
Registro en 108 & \textemdash{}\\
\end{tabularx}\endgroup\par\vspace{6pt}
\Needspace{87pt}\phantomsection\label{trace-267}\noindent\hyperref[nav:vi-ruta-267]{\textbf{r8c4\_\allowbreak{}h+1v-1}}\quad Ordinal unificado / memoria: 267 / 266\par\nopagebreak
\begingroup\fontsize{8.3}{10.1}\selectfont\setlength{\tabcolsep}{4pt}\noindent\begin{tabularx}{\linewidth}{@{}p{.21\linewidth}X@{}}
Genealogía & cbe1adfc\allowbreak{}cc203f0f\allowbreak{}6a54465b\allowbreak{}a638e151\allowbreak{}3bb7f3f9\allowbreak{}2188c76a\allowbreak{}c22a231f\allowbreak{}8fc83757\\
Lector & 8e4f93fa\allowbreak{}6fc6ecbb\allowbreak{}f4f2d29f\allowbreak{}80a71639\allowbreak{}352a11bb\allowbreak{}bd98c66c\allowbreak{}5b0772c8\allowbreak{}df05d749\\
Registro en 108 & \textemdash{}\\
\end{tabularx}\endgroup\par\vspace{6pt}
\Needspace{87pt}\phantomsection\label{trace-268}\noindent\hyperref[nav:vi-ruta-268]{\textbf{r8c4\_\allowbreak{}h+1v+1}}\quad Ordinal unificado / memoria: 268 / 265\par\nopagebreak
\begingroup\fontsize{8.3}{10.1}\selectfont\setlength{\tabcolsep}{4pt}\noindent\begin{tabularx}{\linewidth}{@{}p{.21\linewidth}X@{}}
Genealogía & ad559579\allowbreak{}4bce61df\allowbreak{}813886f1\allowbreak{}36cf73b6\allowbreak{}d5b7712e\allowbreak{}8f1009a2\allowbreak{}125b1324\allowbreak{}2b4e8566\\
Lector & cf06adfa\allowbreak{}5665f171\allowbreak{}a8b8dbe5\allowbreak{}4798e561\allowbreak{}95555e71\allowbreak{}f4986392\allowbreak{}58299d96\allowbreak{}de416911\\
Registro en 108 & 573c1bd9\allowbreak{}0e4644c8\allowbreak{}d653dbd9\allowbreak{}a8229181\allowbreak{}186b709f\allowbreak{}fd121a74\allowbreak{}08c8373b\allowbreak{}4384376c\\
\end{tabularx}\endgroup\par\vspace{6pt}
\Needspace{87pt}\phantomsection\label{trace-269}\noindent\hyperref[nav:vi-ruta-269]{\textbf{r8c5\_\allowbreak{}h-1v-1}}\quad Ordinal unificado / memoria: 269 / 272\par\nopagebreak
\begingroup\fontsize{8.3}{10.1}\selectfont\setlength{\tabcolsep}{4pt}\noindent\begin{tabularx}{\linewidth}{@{}p{.21\linewidth}X@{}}
Genealogía & 7ef6525b\allowbreak{}cd2aeb75\allowbreak{}1a6f5b43\allowbreak{}b6d6a276\allowbreak{}22407710\allowbreak{}aadc04f8\allowbreak{}f412b159\allowbreak{}57971689\\
Lector & 37d699ce\allowbreak{}b70f1c3b\allowbreak{}f1a9991c\allowbreak{}e52e7173\allowbreak{}9128ddd2\allowbreak{}096da773\allowbreak{}3ba3750a\allowbreak{}74627b60\\
Registro en 108 & \textemdash{}\\
\end{tabularx}\endgroup\par\vspace{6pt}
\Needspace{87pt}\phantomsection\label{trace-270}\noindent\hyperref[nav:vi-ruta-270]{\textbf{r8c5\_\allowbreak{}h-1v+1}}\quad Ordinal unificado / memoria: 270 / 271\par\nopagebreak
\begingroup\fontsize{8.3}{10.1}\selectfont\setlength{\tabcolsep}{4pt}\noindent\begin{tabularx}{\linewidth}{@{}p{.21\linewidth}X@{}}
Genealogía & 5168b13c\allowbreak{}9ef8c5df\allowbreak{}5479cfd8\allowbreak{}6a3b5810\allowbreak{}44463645\allowbreak{}8700009b\allowbreak{}99542fd4\allowbreak{}7675d314\\
Lector & 3b12ffa1\allowbreak{}6f7d92e6\allowbreak{}5eccdd73\allowbreak{}152a1897\allowbreak{}c1e1bb8b\allowbreak{}22b621b6\allowbreak{}6ec88ce2\allowbreak{}035f57c6\\
Registro en 108 & b139e53e\allowbreak{}bbb2939e\allowbreak{}05b99768\allowbreak{}4309cd5b\allowbreak{}b321796c\allowbreak{}24a70f00\allowbreak{}06346635\allowbreak{}0ad16680\\
\end{tabularx}\endgroup\par\vspace{6pt}
\Needspace{87pt}\phantomsection\label{trace-271}\noindent\hyperref[nav:vi-ruta-271]{\textbf{r8c5\_\allowbreak{}h+1v-1}}\quad Ordinal unificado / memoria: 271 / 270\par\nopagebreak
\begingroup\fontsize{8.3}{10.1}\selectfont\setlength{\tabcolsep}{4pt}\noindent\begin{tabularx}{\linewidth}{@{}p{.21\linewidth}X@{}}
Genealogía & eedb99f7\allowbreak{}a53750e7\allowbreak{}da1c4cfb\allowbreak{}65caff5b\allowbreak{}c7d6f7fe\allowbreak{}2d17ac33\allowbreak{}bd698f5f\allowbreak{}f5cced70\\
Lector & 3660859f\allowbreak{}8e83b529\allowbreak{}9df231cf\allowbreak{}992060bc\allowbreak{}c7cafcad\allowbreak{}6c8bb559\allowbreak{}d671dcec\allowbreak{}38ae9521\\
Registro en 108 & \textemdash{}\\
\end{tabularx}\endgroup\par\vspace{6pt}
\Needspace{87pt}\phantomsection\label{trace-272}\noindent\hyperref[nav:vi-ruta-272]{\textbf{r8c5\_\allowbreak{}h+1v+1}}\quad Ordinal unificado / memoria: 272 / 269\par\nopagebreak
\begingroup\fontsize{8.3}{10.1}\selectfont\setlength{\tabcolsep}{4pt}\noindent\begin{tabularx}{\linewidth}{@{}p{.21\linewidth}X@{}}
Genealogía & 623aeb8e\allowbreak{}00cd78e8\allowbreak{}b74b3da6\allowbreak{}59368ceb\allowbreak{}c64b59b2\allowbreak{}0668d849\allowbreak{}c58cc8e8\allowbreak{}00f3c668\\
Lector & 9f01aaf1\allowbreak{}1034da5b\allowbreak{}3d5ff388\allowbreak{}ae073672\allowbreak{}65257113\allowbreak{}f331b94b\allowbreak{}7080b8d0\allowbreak{}5e61d946\\
Registro en 108 & \textemdash{}\\
\end{tabularx}\endgroup\par\vspace{6pt}
\Needspace{87pt}\phantomsection\label{trace-273}\noindent\hyperref[nav:vi-ruta-273]{\textbf{r8c6\_\allowbreak{}h-1v-1}}\quad Ordinal unificado / memoria: 273 / 276\par\nopagebreak
\begingroup\fontsize{8.3}{10.1}\selectfont\setlength{\tabcolsep}{4pt}\noindent\begin{tabularx}{\linewidth}{@{}p{.21\linewidth}X@{}}
Genealogía & 545d27eb\allowbreak{}e80f7a1b\allowbreak{}5d6e4faf\allowbreak{}6f6614e5\allowbreak{}9a75427f\allowbreak{}dc5db81a\allowbreak{}0eff6417\allowbreak{}85391008\\
Lector & df6973d8\allowbreak{}dba7d543\allowbreak{}ad4991d2\allowbreak{}fe5c1ebf\allowbreak{}3eb9a180\allowbreak{}3d7e933f\allowbreak{}0074ba08\allowbreak{}0793ff48\\
Registro en 108 & 165e9790\allowbreak{}6c2b634f\allowbreak{}80f723d3\allowbreak{}b6db0f07\allowbreak{}d486840b\allowbreak{}dcf88079\allowbreak{}defcfc71\allowbreak{}121b3293\\
\end{tabularx}\endgroup\par\vspace{6pt}
\Needspace{87pt}\phantomsection\label{trace-274}\noindent\hyperref[nav:vi-ruta-274]{\textbf{r8c6\_\allowbreak{}h-1v+1}}\quad Ordinal unificado / memoria: 274 / 275\par\nopagebreak
\begingroup\fontsize{8.3}{10.1}\selectfont\setlength{\tabcolsep}{4pt}\noindent\begin{tabularx}{\linewidth}{@{}p{.21\linewidth}X@{}}
Genealogía & f71b59ca\allowbreak{}66689f84\allowbreak{}06ea25f2\allowbreak{}9129a8bb\allowbreak{}6aadb644\allowbreak{}7707eafc\allowbreak{}3f076a9d\allowbreak{}7ed74c3f\\
Lector & 45ad6fbe\allowbreak{}bdb7517b\allowbreak{}d4fc0760\allowbreak{}b8fc1596\allowbreak{}fbc582e6\allowbreak{}077f671d\allowbreak{}37ea2a5a\allowbreak{}9582a424\\
Registro en 108 & \textemdash{}\\
\end{tabularx}\endgroup\par\vspace{6pt}
\Needspace{87pt}\phantomsection\label{trace-275}\noindent\hyperref[nav:vi-ruta-275]{\textbf{r8c6\_\allowbreak{}h+1v-1}}\quad Ordinal unificado / memoria: 275 / 274\par\nopagebreak
\begingroup\fontsize{8.3}{10.1}\selectfont\setlength{\tabcolsep}{4pt}\noindent\begin{tabularx}{\linewidth}{@{}p{.21\linewidth}X@{}}
Genealogía & b2394c0d\allowbreak{}8b4f94d4\allowbreak{}d3db2058\allowbreak{}caf65f5c\allowbreak{}ba5037fa\allowbreak{}aba86ab8\allowbreak{}b741d761\allowbreak{}4066f5b9\\
Lector & b35fffde\allowbreak{}6d573fe8\allowbreak{}075155a8\allowbreak{}4df6a2ed\allowbreak{}2e98d9bf\allowbreak{}9fe9e09c\allowbreak{}1e0624fb\allowbreak{}83395b24\\
Registro en 108 & \textemdash{}\\
\end{tabularx}\endgroup\par\vspace{6pt}
\Needspace{87pt}\phantomsection\label{trace-276}\noindent\hyperref[nav:vi-ruta-276]{\textbf{r8c6\_\allowbreak{}h+1v+1}}\quad Ordinal unificado / memoria: 276 / 273\par\nopagebreak
\begingroup\fontsize{8.3}{10.1}\selectfont\setlength{\tabcolsep}{4pt}\noindent\begin{tabularx}{\linewidth}{@{}p{.21\linewidth}X@{}}
Genealogía & 1fa8ca36\allowbreak{}8f05cacb\allowbreak{}b63847f3\allowbreak{}d7c151b4\allowbreak{}b61b01f5\allowbreak{}663f9d4f\allowbreak{}2e416c66\allowbreak{}bb475d96\\
Lector & 7958f0da\allowbreak{}3b1c71e5\allowbreak{}6248b77b\allowbreak{}0bd939b0\allowbreak{}279b68f8\allowbreak{}ff68cb14\allowbreak{}5e28e5e0\allowbreak{}6d5a3f6e\\
Registro en 108 & \textemdash{}\\
\end{tabularx}\endgroup\par\vspace{6pt}
\Needspace{87pt}\phantomsection\label{trace-277}\noindent\hyperref[nav:vi-ruta-277]{\textbf{r8c7\_\allowbreak{}h-1v-1}}\quad Ordinal unificado / memoria: 277 / 280\par\nopagebreak
\begingroup\fontsize{8.3}{10.1}\selectfont\setlength{\tabcolsep}{4pt}\noindent\begin{tabularx}{\linewidth}{@{}p{.21\linewidth}X@{}}
Genealogía & 3e4f2795\allowbreak{}854d3604\allowbreak{}5b55509b\allowbreak{}d916303c\allowbreak{}028de5ba\allowbreak{}24d97eef\allowbreak{}b4217663\allowbreak{}1fad5420\\
Lector & 4e3bc25c\allowbreak{}d834c537\allowbreak{}9d7cd435\allowbreak{}b1067433\allowbreak{}84133c2f\allowbreak{}09dc9ced\allowbreak{}7d99880b\allowbreak{}ed745a3e\\
Registro en 108 & \textemdash{}\\
\end{tabularx}\endgroup\par\vspace{6pt}
\Needspace{87pt}\phantomsection\label{trace-278}\noindent\hyperref[nav:vi-ruta-278]{\textbf{r8c7\_\allowbreak{}h-1v+1}}\quad Ordinal unificado / memoria: 278 / 279\par\nopagebreak
\begingroup\fontsize{8.3}{10.1}\selectfont\setlength{\tabcolsep}{4pt}\noindent\begin{tabularx}{\linewidth}{@{}p{.21\linewidth}X@{}}
Genealogía & fd2d7320\allowbreak{}f5b5383c\allowbreak{}a26d8e6d\allowbreak{}224b7df3\allowbreak{}6ee61afb\allowbreak{}23e51402\allowbreak{}ebf35ae0\allowbreak{}ee13a4f9\\
Lector & 3840df2d\allowbreak{}6dfaac48\allowbreak{}e51ce48c\allowbreak{}a841735c\allowbreak{}a0b7eaab\allowbreak{}8b9a9eaa\allowbreak{}a62c6ba1\allowbreak{}a7e40f39\\
Registro en 108 & \textemdash{}\\
\end{tabularx}\endgroup\par\vspace{6pt}
\Needspace{87pt}\phantomsection\label{trace-279}\noindent\hyperref[nav:vi-ruta-279]{\textbf{r8c7\_\allowbreak{}h+1v-1}}\quad Ordinal unificado / memoria: 279 / 278\par\nopagebreak
\begingroup\fontsize{8.3}{10.1}\selectfont\setlength{\tabcolsep}{4pt}\noindent\begin{tabularx}{\linewidth}{@{}p{.21\linewidth}X@{}}
Genealogía & a15a54cb\allowbreak{}8e5aeadb\allowbreak{}f9a14d62\allowbreak{}ee241c1b\allowbreak{}33911588\allowbreak{}d46d105d\allowbreak{}0ba47dc8\allowbreak{}2b7959ba\\
Lector & 0919f3a4\allowbreak{}d99a7768\allowbreak{}67a28a6d\allowbreak{}334ec1b7\allowbreak{}17ff3b3a\allowbreak{}44917668\allowbreak{}6ee79e46\allowbreak{}03cead5a\\
Registro en 108 & f768a45c\allowbreak{}6e483031\allowbreak{}90120250\allowbreak{}88381e9a\allowbreak{}82811497\allowbreak{}d99ad157\allowbreak{}d5b1c623\allowbreak{}be89b25b\\
\end{tabularx}\endgroup\par\vspace{6pt}
\Needspace{87pt}\phantomsection\label{trace-280}\noindent\hyperref[nav:vi-ruta-280]{\textbf{r8c7\_\allowbreak{}h+1v+1}}\quad Ordinal unificado / memoria: 280 / 277\par\nopagebreak
\begingroup\fontsize{8.3}{10.1}\selectfont\setlength{\tabcolsep}{4pt}\noindent\begin{tabularx}{\linewidth}{@{}p{.21\linewidth}X@{}}
Genealogía & 86715e7e\allowbreak{}ccf63f42\allowbreak{}1cf277f3\allowbreak{}abd516b2\allowbreak{}26f7eb2d\allowbreak{}97fcf24b\allowbreak{}d37912af\allowbreak{}78d1d866\\
Lector & 622a5691\allowbreak{}a81501fa\allowbreak{}895d6474\allowbreak{}a6d9e5ae\allowbreak{}5acb0ef9\allowbreak{}6649d468\allowbreak{}75be617d\allowbreak{}bfe36c68\\
Registro en 108 & \textemdash{}\\
\end{tabularx}\endgroup\par\vspace{6pt}
\Needspace{87pt}\phantomsection\label{trace-281}\noindent\hyperref[nav:vi-ruta-281]{\textbf{r8c8\_\allowbreak{}h-1v-1}}\quad Ordinal unificado / memoria: 281 / 284\par\nopagebreak
\begingroup\fontsize{8.3}{10.1}\selectfont\setlength{\tabcolsep}{4pt}\noindent\begin{tabularx}{\linewidth}{@{}p{.21\linewidth}X@{}}
Genealogía & 0a0cf9df\allowbreak{}7293b410\allowbreak{}440f3fb0\allowbreak{}87c85c25\allowbreak{}fa5039d0\allowbreak{}fc789b9a\allowbreak{}b6bad452\allowbreak{}1fb837b5\\
Lector & 111950fe\allowbreak{}fee357c0\allowbreak{}f7efaef1\allowbreak{}ee7794ea\allowbreak{}3fc454e0\allowbreak{}1de39489\allowbreak{}ecfd6e18\allowbreak{}be9fdef9\\
Registro en 108 & \textemdash{}\\
\end{tabularx}\endgroup\par\vspace{6pt}
\Needspace{87pt}\phantomsection\label{trace-282}\noindent\hyperref[nav:vi-ruta-282]{\textbf{r8c8\_\allowbreak{}h-1v+1}}\quad Ordinal unificado / memoria: 282 / 283\par\nopagebreak
\begingroup\fontsize{8.3}{10.1}\selectfont\setlength{\tabcolsep}{4pt}\noindent\begin{tabularx}{\linewidth}{@{}p{.21\linewidth}X@{}}
Genealogía & 44c68f8b\allowbreak{}3fce4afc\allowbreak{}7439b1cd\allowbreak{}4add9d73\allowbreak{}ea2d68ae\allowbreak{}57cd4025\allowbreak{}f427e0da\allowbreak{}591dda73\\
Lector & 4453d191\allowbreak{}3b855774\allowbreak{}dce1ec81\allowbreak{}650654f1\allowbreak{}0fd4b0c3\allowbreak{}e26ed21e\allowbreak{}8dff2648\allowbreak{}54b464b5\\
Registro en 108 & \textemdash{}\\
\end{tabularx}\endgroup\par\vspace{6pt}
\Needspace{87pt}\phantomsection\label{trace-283}\noindent\hyperref[nav:vi-ruta-283]{\textbf{r8c8\_\allowbreak{}h+1v-1}}\quad Ordinal unificado / memoria: 283 / 282\par\nopagebreak
\begingroup\fontsize{8.3}{10.1}\selectfont\setlength{\tabcolsep}{4pt}\noindent\begin{tabularx}{\linewidth}{@{}p{.21\linewidth}X@{}}
Genealogía & 6ccdebca\allowbreak{}af311843\allowbreak{}d1d894e5\allowbreak{}d1c8dee1\allowbreak{}f317bb40\allowbreak{}cfb7fef1\allowbreak{}455cd39b\allowbreak{}357a6807\\
Lector & cee239d9\allowbreak{}5a64fd32\allowbreak{}c7200279\allowbreak{}9f66be6c\allowbreak{}d6167d2f\allowbreak{}7d01c473\allowbreak{}97015660\allowbreak{}8f61bd1a\\
Registro en 108 & \textemdash{}\\
\end{tabularx}\endgroup\par\vspace{6pt}
\Needspace{87pt}\phantomsection\label{trace-284}\noindent\hyperref[nav:vi-ruta-284]{\textbf{r8c8\_\allowbreak{}h+1v+1}}\quad Ordinal unificado / memoria: 284 / 281\par\nopagebreak
\begingroup\fontsize{8.3}{10.1}\selectfont\setlength{\tabcolsep}{4pt}\noindent\begin{tabularx}{\linewidth}{@{}p{.21\linewidth}X@{}}
Genealogía & 1fef105f\allowbreak{}9462f0c6\allowbreak{}29ad4505\allowbreak{}f2fa1a8a\allowbreak{}2d248638\allowbreak{}113913db\allowbreak{}d37ae667\allowbreak{}91ac19f1\\
Lector & 735623ac\allowbreak{}3cc96b60\allowbreak{}8c56fc45\allowbreak{}6cff66da\allowbreak{}f8b0193b\allowbreak{}82645553\allowbreak{}f1fa88ef\allowbreak{}b1ba6bed\\
Registro en 108 & 45951d9d\allowbreak{}048937aa\allowbreak{}e4c2cfb3\allowbreak{}a4823a1f\allowbreak{}f75bbfb0\allowbreak{}f359f391\allowbreak{}421e6990\allowbreak{}f5bba610\\
\end{tabularx}\endgroup\par\vspace{6pt}
\Needspace{87pt}\phantomsection\label{trace-285}\noindent\hyperref[nav:vi-ruta-285]{\textbf{r8c9\_\allowbreak{}h-1v-1}}\quad Ordinal unificado / memoria: 285 / 288\par\nopagebreak
\begingroup\fontsize{8.3}{10.1}\selectfont\setlength{\tabcolsep}{4pt}\noindent\begin{tabularx}{\linewidth}{@{}p{.21\linewidth}X@{}}
Genealogía & 8569908f\allowbreak{}3825cdd4\allowbreak{}6d2aeb61\allowbreak{}5280f04f\allowbreak{}11d02c6e\allowbreak{}3d9b945b\allowbreak{}706a1953\allowbreak{}b6b91832\\
Lector & 9d065f91\allowbreak{}57694f9e\allowbreak{}dc555f5e\allowbreak{}00871f36\allowbreak{}ffdfbf96\allowbreak{}d1f47fe1\allowbreak{}1deda4e3\allowbreak{}a183b594\\
Registro en 108 & \textemdash{}\\
\end{tabularx}\endgroup\par\vspace{6pt}
\Needspace{87pt}\phantomsection\label{trace-286}\noindent\hyperref[nav:vi-ruta-286]{\textbf{r8c9\_\allowbreak{}h-1v+1}}\quad Ordinal unificado / memoria: 286 / 287\par\nopagebreak
\begingroup\fontsize{8.3}{10.1}\selectfont\setlength{\tabcolsep}{4pt}\noindent\begin{tabularx}{\linewidth}{@{}p{.21\linewidth}X@{}}
Genealogía & eff03ec9\allowbreak{}f7fc2cef\allowbreak{}5cb21ba6\allowbreak{}fa45a82d\allowbreak{}a8267d30\allowbreak{}92637150\allowbreak{}b953e4aa\allowbreak{}980bbe87\\
Lector & 70dd7c0e\allowbreak{}b3931478\allowbreak{}51d6fd2d\allowbreak{}32d16fd7\allowbreak{}e393a5d3\allowbreak{}53010311\allowbreak{}ff0ef60a\allowbreak{}a076d668\\
Registro en 108 & c7dfe852\allowbreak{}a4958d07\allowbreak{}d62f9651\allowbreak{}78f30c16\allowbreak{}579632ac\allowbreak{}1d2a4764\allowbreak{}b1535ee5\allowbreak{}6b4f89ba\\
\end{tabularx}\endgroup\par\vspace{6pt}
\Needspace{87pt}\phantomsection\label{trace-287}\noindent\hyperref[nav:vi-ruta-287]{\textbf{r8c9\_\allowbreak{}h+1v-1}}\quad Ordinal unificado / memoria: 287 / 286\par\nopagebreak
\begingroup\fontsize{8.3}{10.1}\selectfont\setlength{\tabcolsep}{4pt}\noindent\begin{tabularx}{\linewidth}{@{}p{.21\linewidth}X@{}}
Genealogía & 678d6ab8\allowbreak{}d97c5077\allowbreak{}2b92a1c0\allowbreak{}c11ae3a7\allowbreak{}bc3fa033\allowbreak{}78f49a8a\allowbreak{}322f7db0\allowbreak{}e7a20776\\
Lector & 738d138b\allowbreak{}0027417c\allowbreak{}cb40bd19\allowbreak{}ddb10609\allowbreak{}eb8c8531\allowbreak{}dffbf0c1\allowbreak{}903a319e\allowbreak{}2767add0\\
Registro en 108 & \textemdash{}\\
\end{tabularx}\endgroup\par\vspace{6pt}
\Needspace{87pt}\phantomsection\label{trace-288}\noindent\hyperref[nav:vi-ruta-288]{\textbf{r8c9\_\allowbreak{}h+1v+1}}\quad Ordinal unificado / memoria: 288 / 285\par\nopagebreak
\begingroup\fontsize{8.3}{10.1}\selectfont\setlength{\tabcolsep}{4pt}\noindent\begin{tabularx}{\linewidth}{@{}p{.21\linewidth}X@{}}
Genealogía & 81affe3c\allowbreak{}95425c70\allowbreak{}af6e783d\allowbreak{}072ff892\allowbreak{}34b1446d\allowbreak{}1055fed3\allowbreak{}1ee30519\allowbreak{}2280fcb3\\
Lector & 41036aca\allowbreak{}4f9f8ad1\allowbreak{}35146507\allowbreak{}fe31266d\allowbreak{}0452840b\allowbreak{}897d87b1\allowbreak{}a5f15e66\allowbreak{}39516828\\
Registro en 108 & \textemdash{}\\
\end{tabularx}\endgroup\par\vspace{6pt}
\Needspace{87pt}\phantomsection\label{trace-289}\noindent\hyperref[nav:vi-ruta-289]{\textbf{r9c1\_\allowbreak{}h-1v-1}}\quad Ordinal unificado / memoria: 289 / 292\par\nopagebreak
\begingroup\fontsize{8.3}{10.1}\selectfont\setlength{\tabcolsep}{4pt}\noindent\begin{tabularx}{\linewidth}{@{}p{.21\linewidth}X@{}}
Genealogía & a0d3b90a\allowbreak{}93737ad6\allowbreak{}bd06c49c\allowbreak{}be1e379f\allowbreak{}6e67a39e\allowbreak{}7b60dbb4\allowbreak{}e0e0b548\allowbreak{}0d1c9730\\
Lector & dd19ac7c\allowbreak{}0ec61adb\allowbreak{}ef57b6ad\allowbreak{}840a051f\allowbreak{}b3b46e58\allowbreak{}7098f2af\allowbreak{}e5101dab\allowbreak{}a8f867b1\\
Registro en 108 & c0580ece\allowbreak{}103707e2\allowbreak{}2c2b7372\allowbreak{}0bc58713\allowbreak{}206cce80\allowbreak{}bca99c63\allowbreak{}37432129\allowbreak{}84199b55\\
\end{tabularx}\endgroup\par\vspace{6pt}
\Needspace{87pt}\phantomsection\label{trace-290}\noindent\hyperref[nav:vi-ruta-290]{\textbf{r9c1\_\allowbreak{}h-1v+1}}\quad Ordinal unificado / memoria: 290 / 291\par\nopagebreak
\begingroup\fontsize{8.3}{10.1}\selectfont\setlength{\tabcolsep}{4pt}\noindent\begin{tabularx}{\linewidth}{@{}p{.21\linewidth}X@{}}
Genealogía & 8483a8df\allowbreak{}880a8c77\allowbreak{}b15bbbdc\allowbreak{}d777ed2f\allowbreak{}7cc79e2b\allowbreak{}9d4e199d\allowbreak{}c642f2a9\allowbreak{}012041b9\\
Lector & ae1cc25e\allowbreak{}4cd5b1b8\allowbreak{}7552d23b\allowbreak{}208c33db\allowbreak{}c0ec0121\allowbreak{}a8bb290c\allowbreak{}101337e7\allowbreak{}79bb5e1d\\
Registro en 108 & \textemdash{}\\
\end{tabularx}\endgroup\par\vspace{6pt}
\Needspace{87pt}\phantomsection\label{trace-291}\noindent\hyperref[nav:vi-ruta-291]{\textbf{r9c1\_\allowbreak{}h+1v-1}}\quad Ordinal unificado / memoria: 291 / 290\par\nopagebreak
\begingroup\fontsize{8.3}{10.1}\selectfont\setlength{\tabcolsep}{4pt}\noindent\begin{tabularx}{\linewidth}{@{}p{.21\linewidth}X@{}}
Genealogía & 52b846cd\allowbreak{}7eaeff31\allowbreak{}abb00f1a\allowbreak{}e440a22d\allowbreak{}9b65cc73\allowbreak{}fa2cc099\allowbreak{}196a0ef2\allowbreak{}0d97df3c\\
Lector & 3db7a708\allowbreak{}d8ae2238\allowbreak{}a749a756\allowbreak{}23f90c80\allowbreak{}ccd6b52f\allowbreak{}a0c36042\allowbreak{}a443240c\allowbreak{}6f9ebfc0\\
Registro en 108 & \textemdash{}\\
\end{tabularx}\endgroup\par\vspace{6pt}
\Needspace{87pt}\phantomsection\label{trace-292}\noindent\hyperref[nav:vi-ruta-292]{\textbf{r9c1\_\allowbreak{}h+1v+1}}\quad Ordinal unificado / memoria: 292 / 289\par\nopagebreak
\begingroup\fontsize{8.3}{10.1}\selectfont\setlength{\tabcolsep}{4pt}\noindent\begin{tabularx}{\linewidth}{@{}p{.21\linewidth}X@{}}
Genealogía & 905ded93\allowbreak{}d9b0764d\allowbreak{}20704971\allowbreak{}2252c1ff\allowbreak{}b8506d94\allowbreak{}8be0960f\allowbreak{}3b6e398c\allowbreak{}aa4a8e25\\
Lector & ddb85623\allowbreak{}34314015\allowbreak{}284c5247\allowbreak{}ec81419d\allowbreak{}909c46c4\allowbreak{}3781c417\allowbreak{}45de130f\allowbreak{}32a0e193\\
Registro en 108 & \textemdash{}\\
\end{tabularx}\endgroup\par\vspace{6pt}
\Needspace{87pt}\phantomsection\label{trace-293}\noindent\hyperref[nav:vi-ruta-293]{\textbf{r9c2\_\allowbreak{}h-1v-1}}\quad Ordinal unificado / memoria: 293 / 296\par\nopagebreak
\begingroup\fontsize{8.3}{10.1}\selectfont\setlength{\tabcolsep}{4pt}\noindent\begin{tabularx}{\linewidth}{@{}p{.21\linewidth}X@{}}
Genealogía & da666726\allowbreak{}7b4ff85f\allowbreak{}ea2f0ff5\allowbreak{}b28c8962\allowbreak{}cba978a6\allowbreak{}c2329bf5\allowbreak{}8403d88a\allowbreak{}aab3b865\\
Lector & 47731741\allowbreak{}df561e96\allowbreak{}ed208b39\allowbreak{}1e19d518\allowbreak{}430489a7\allowbreak{}ef9c6d05\allowbreak{}607c1641\allowbreak{}14f5d9a2\\
Registro en 108 & \textemdash{}\\
\end{tabularx}\endgroup\par\vspace{6pt}
\Needspace{87pt}\phantomsection\label{trace-294}\noindent\hyperref[nav:vi-ruta-294]{\textbf{r9c2\_\allowbreak{}h-1v+1}}\quad Ordinal unificado / memoria: 294 / 295\par\nopagebreak
\begingroup\fontsize{8.3}{10.1}\selectfont\setlength{\tabcolsep}{4pt}\noindent\begin{tabularx}{\linewidth}{@{}p{.21\linewidth}X@{}}
Genealogía & 46e22be8\allowbreak{}f2160920\allowbreak{}042cedaf\allowbreak{}06032c82\allowbreak{}685a879d\allowbreak{}1f441907\allowbreak{}912b9a00\allowbreak{}df9114a3\\
Lector & a7ff5dec\allowbreak{}dcbda0dc\allowbreak{}9f2f3c09\allowbreak{}ec8cdeac\allowbreak{}01d3a3ad\allowbreak{}66cb4dfe\allowbreak{}d656bf97\allowbreak{}0900d4bb\\
Registro en 108 & 5ffebc4f\allowbreak{}dd94691d\allowbreak{}eee35c6d\allowbreak{}4dcacdee\allowbreak{}c958fddf\allowbreak{}4a3daff5\allowbreak{}98d7a1a1\allowbreak{}1d1360c0\\
\end{tabularx}\endgroup\par\vspace{6pt}
\Needspace{87pt}\phantomsection\label{trace-295}\noindent\hyperref[nav:vi-ruta-295]{\textbf{r9c2\_\allowbreak{}h+1v-1}}\quad Ordinal unificado / memoria: 295 / 294\par\nopagebreak
\begingroup\fontsize{8.3}{10.1}\selectfont\setlength{\tabcolsep}{4pt}\noindent\begin{tabularx}{\linewidth}{@{}p{.21\linewidth}X@{}}
Genealogía & c3161556\allowbreak{}f17c0edf\allowbreak{}adbc9569\allowbreak{}3b722fde\allowbreak{}3bda92e6\allowbreak{}1c398a49\allowbreak{}7d9ea6f4\allowbreak{}f45c21ff\\
Lector & 0b2ce983\allowbreak{}32ac4c69\allowbreak{}c985a795\allowbreak{}b31e585c\allowbreak{}57e037a8\allowbreak{}fbbe9508\allowbreak{}bf72fe59\allowbreak{}17ac5e28\\
Registro en 108 & \textemdash{}\\
\end{tabularx}\endgroup\par\vspace{6pt}
\Needspace{87pt}\phantomsection\label{trace-296}\noindent\hyperref[nav:vi-ruta-296]{\textbf{r9c2\_\allowbreak{}h+1v+1}}\quad Ordinal unificado / memoria: 296 / 293\par\nopagebreak
\begingroup\fontsize{8.3}{10.1}\selectfont\setlength{\tabcolsep}{4pt}\noindent\begin{tabularx}{\linewidth}{@{}p{.21\linewidth}X@{}}
Genealogía & 463b5be4\allowbreak{}23eaf1c5\allowbreak{}0aa2a7f5\allowbreak{}f851d35d\allowbreak{}af48aacc\allowbreak{}82ab2571\allowbreak{}4a83f284\allowbreak{}0170d00c\\
Lector & 3e4d2191\allowbreak{}63e182c8\allowbreak{}401f7d7d\allowbreak{}fe984bfa\allowbreak{}d7fb28bb\allowbreak{}a4b15fda\allowbreak{}36d17bdd\allowbreak{}b01bcd23\\
Registro en 108 & \textemdash{}\\
\end{tabularx}\endgroup\par\vspace{6pt}
\Needspace{87pt}\phantomsection\label{trace-297}\noindent\hyperref[nav:vi-ruta-297]{\textbf{r9c3\_\allowbreak{}h-1v-1}}\quad Ordinal unificado / memoria: 297 / 300\par\nopagebreak
\begingroup\fontsize{8.3}{10.1}\selectfont\setlength{\tabcolsep}{4pt}\noindent\begin{tabularx}{\linewidth}{@{}p{.21\linewidth}X@{}}
Genealogía & 8044e591\allowbreak{}7e13ccfb\allowbreak{}da2ceb1f\allowbreak{}21613f7c\allowbreak{}e52eaa08\allowbreak{}53ec71fd\allowbreak{}75f32203\allowbreak{}76d19cc0\\
Lector & bb47bdfe\allowbreak{}5cf32f87\allowbreak{}265f7ab2\allowbreak{}c767c464\allowbreak{}40e85eef\allowbreak{}a678090e\allowbreak{}2f952f45\allowbreak{}04fe98c1\\
Registro en 108 & \textemdash{}\\
\end{tabularx}\endgroup\par\vspace{6pt}
\Needspace{87pt}\phantomsection\label{trace-298}\noindent\hyperref[nav:vi-ruta-298]{\textbf{r9c3\_\allowbreak{}h-1v+1}}\quad Ordinal unificado / memoria: 298 / 299\par\nopagebreak
\begingroup\fontsize{8.3}{10.1}\selectfont\setlength{\tabcolsep}{4pt}\noindent\begin{tabularx}{\linewidth}{@{}p{.21\linewidth}X@{}}
Genealogía & b1f5d904\allowbreak{}7f674833\allowbreak{}5b6f26ef\allowbreak{}ffc458e4\allowbreak{}baa8c0f8\allowbreak{}a8dfa880\allowbreak{}1b0cd507\allowbreak{}1ee377e9\\
Lector & e138662c\allowbreak{}f0ef5ed6\allowbreak{}72f884e4\allowbreak{}4662fd43\allowbreak{}46e0c345\allowbreak{}559c2c58\allowbreak{}45dc6dc3\allowbreak{}0a15b0db\\
Registro en 108 & \textemdash{}\\
\end{tabularx}\endgroup\par\vspace{6pt}
\Needspace{87pt}\phantomsection\label{trace-299}\noindent\hyperref[nav:vi-ruta-299]{\textbf{r9c3\_\allowbreak{}h+1v-1}}\quad Ordinal unificado / memoria: 299 / 298\par\nopagebreak
\begingroup\fontsize{8.3}{10.1}\selectfont\setlength{\tabcolsep}{4pt}\noindent\begin{tabularx}{\linewidth}{@{}p{.21\linewidth}X@{}}
Genealogía & 3a549e80\allowbreak{}92847f10\allowbreak{}e437f1ad\allowbreak{}d78c27b8\allowbreak{}119bea4b\allowbreak{}98466f1a\allowbreak{}495c7fe8\allowbreak{}37e526ab\\
Lector & 75ff4410\allowbreak{}07b14e7d\allowbreak{}209d0846\allowbreak{}ef9fadf4\allowbreak{}9bdbc245\allowbreak{}749a8b5b\allowbreak{}cddb11c7\allowbreak{}23234236\\
Registro en 108 & \textemdash{}\\
\end{tabularx}\endgroup\par\vspace{6pt}
\Needspace{87pt}\phantomsection\label{trace-300}\noindent\hyperref[nav:vi-ruta-300]{\textbf{r9c3\_\allowbreak{}h+1v+1}}\quad Ordinal unificado / memoria: 300 / 297\par\nopagebreak
\begingroup\fontsize{8.3}{10.1}\selectfont\setlength{\tabcolsep}{4pt}\noindent\begin{tabularx}{\linewidth}{@{}p{.21\linewidth}X@{}}
Genealogía & 36d1117b\allowbreak{}7f4f744e\allowbreak{}498155bb\allowbreak{}c7dc3f33\allowbreak{}7e038bc6\allowbreak{}a59df945\allowbreak{}b16955b6\allowbreak{}5617a24f\\
Lector & 2822864e\allowbreak{}6d668c0d\allowbreak{}559901e9\allowbreak{}3691fe2f\allowbreak{}f2784b9d\allowbreak{}82824b2d\allowbreak{}5fe3c593\allowbreak{}e5cc387d\\
Registro en 108 & 55065291\allowbreak{}299c2ae2\allowbreak{}d06f2baa\allowbreak{}5791c031\allowbreak{}9edfbd41\allowbreak{}3c3e9c7a\allowbreak{}8b949c33\allowbreak{}52304579\\
\end{tabularx}\endgroup\par\vspace{6pt}
\Needspace{87pt}\phantomsection\label{trace-301}\noindent\hyperref[nav:vi-ruta-301]{\textbf{r9c4\_\allowbreak{}h-1v-1}}\quad Ordinal unificado / memoria: 301 / 304\par\nopagebreak
\begingroup\fontsize{8.3}{10.1}\selectfont\setlength{\tabcolsep}{4pt}\noindent\begin{tabularx}{\linewidth}{@{}p{.21\linewidth}X@{}}
Genealogía & ee798204\allowbreak{}ea997ed7\allowbreak{}ed2199e7\allowbreak{}b3ab7d61\allowbreak{}37f5ac35\allowbreak{}f07e503f\allowbreak{}6c0a8eda\allowbreak{}e955b848\\
Lector & 18fd69ac\allowbreak{}1d8dc23d\allowbreak{}374c9d4b\allowbreak{}781e339c\allowbreak{}278b6935\allowbreak{}a74bb76c\allowbreak{}1a2badda\allowbreak{}cb2b6fea\\
Registro en 108 & \textemdash{}\\
\end{tabularx}\endgroup\par\vspace{6pt}
\Needspace{87pt}\phantomsection\label{trace-302}\noindent\hyperref[nav:vi-ruta-302]{\textbf{r9c4\_\allowbreak{}h-1v+1}}\quad Ordinal unificado / memoria: 302 / 303\par\nopagebreak
\begingroup\fontsize{8.3}{10.1}\selectfont\setlength{\tabcolsep}{4pt}\noindent\begin{tabularx}{\linewidth}{@{}p{.21\linewidth}X@{}}
Genealogía & 642cd124\allowbreak{}82d3d647\allowbreak{}cdaec0ac\allowbreak{}dd19e9c1\allowbreak{}f2a74296\allowbreak{}75270183\allowbreak{}ed630cae\allowbreak{}a3cf535a\\
Lector & e7f63cff\allowbreak{}972e6706\allowbreak{}fe7c8ecd\allowbreak{}667eb3e0\allowbreak{}76106f38\allowbreak{}0840dadd\allowbreak{}bcaf4c59\allowbreak{}501a2e78\\
Registro en 108 & 11028c38\allowbreak{}9e35d609\allowbreak{}ddb9aa12\allowbreak{}48b79b77\allowbreak{}987b4178\allowbreak{}3a58974a\allowbreak{}3c6e6873\allowbreak{}c14e1f88\\
\end{tabularx}\endgroup\par\vspace{6pt}
\Needspace{87pt}\phantomsection\label{trace-303}\noindent\hyperref[nav:vi-ruta-303]{\textbf{r9c4\_\allowbreak{}h+1v-1}}\quad Ordinal unificado / memoria: 303 / 302\par\nopagebreak
\begingroup\fontsize{8.3}{10.1}\selectfont\setlength{\tabcolsep}{4pt}\noindent\begin{tabularx}{\linewidth}{@{}p{.21\linewidth}X@{}}
Genealogía & ff002758\allowbreak{}4bafb795\allowbreak{}fea0705b\allowbreak{}ea171be4\allowbreak{}17920f7a\allowbreak{}b7cb5f37\allowbreak{}8d17b7e4\allowbreak{}0adbb917\\
Lector & 7117e0fd\allowbreak{}d957f42a\allowbreak{}5d46ef72\allowbreak{}7105788a\allowbreak{}73ce6a86\allowbreak{}3be2fd32\allowbreak{}f11bab0a\allowbreak{}df65bbc6\\
Registro en 108 & \textemdash{}\\
\end{tabularx}\endgroup\par\vspace{6pt}
\Needspace{87pt}\phantomsection\label{trace-304}\noindent\hyperref[nav:vi-ruta-304]{\textbf{r9c4\_\allowbreak{}h+1v+1}}\quad Ordinal unificado / memoria: 304 / 301\par\nopagebreak
\begingroup\fontsize{8.3}{10.1}\selectfont\setlength{\tabcolsep}{4pt}\noindent\begin{tabularx}{\linewidth}{@{}p{.21\linewidth}X@{}}
Genealogía & b5fc9a01\allowbreak{}e3980db0\allowbreak{}70dd783c\allowbreak{}d2fcc157\allowbreak{}585ffca7\allowbreak{}f68f54ea\allowbreak{}c0c828e4\allowbreak{}aeedf880\\
Lector & 5b3951f3\allowbreak{}240c6ee9\allowbreak{}34749d6d\allowbreak{}c143cce4\allowbreak{}f650ab60\allowbreak{}7b642718\allowbreak{}14be5fd6\allowbreak{}7595c351\\
Registro en 108 & \textemdash{}\\
\end{tabularx}\endgroup\par\vspace{6pt}
\Needspace{87pt}\phantomsection\label{trace-305}\noindent\hyperref[nav:vi-ruta-305]{\textbf{r9c5\_\allowbreak{}h-1v-1}}\quad Ordinal unificado / memoria: 305 / 308\par\nopagebreak
\begingroup\fontsize{8.3}{10.1}\selectfont\setlength{\tabcolsep}{4pt}\noindent\begin{tabularx}{\linewidth}{@{}p{.21\linewidth}X@{}}
Genealogía & 77908e47\allowbreak{}2c16ed99\allowbreak{}b259dbc3\allowbreak{}97c2623e\allowbreak{}67f0871a\allowbreak{}d4764417\allowbreak{}423d1c62\allowbreak{}bd7e5dfa\\
Lector & 3d68025f\allowbreak{}fbb2f16f\allowbreak{}fc102532\allowbreak{}00612155\allowbreak{}c3edd254\allowbreak{}5ed08c9a\allowbreak{}cf932f34\allowbreak{}b8417dea\\
Registro en 108 & ee7f57b4\allowbreak{}3ff6bbe2\allowbreak{}b8b762ae\allowbreak{}c3be88c2\allowbreak{}f3eec4ca\allowbreak{}c943953b\allowbreak{}f26b06ad\allowbreak{}fc1e2088\\
\end{tabularx}\endgroup\par\vspace{6pt}
\Needspace{87pt}\phantomsection\label{trace-306}\noindent\hyperref[nav:vi-ruta-306]{\textbf{r9c5\_\allowbreak{}h-1v+1}}\quad Ordinal unificado / memoria: 306 / 307\par\nopagebreak
\begingroup\fontsize{8.3}{10.1}\selectfont\setlength{\tabcolsep}{4pt}\noindent\begin{tabularx}{\linewidth}{@{}p{.21\linewidth}X@{}}
Genealogía & 6d009acb\allowbreak{}2e4a9062\allowbreak{}d353d9d3\allowbreak{}d3b9fabb\allowbreak{}dd58df3d\allowbreak{}fffdf48f\allowbreak{}f5433029\allowbreak{}54972595\\
Lector & b47145e5\allowbreak{}c4d4c54c\allowbreak{}6e7bcbde\allowbreak{}b698559d\allowbreak{}e66f3051\allowbreak{}1a1ceb50\allowbreak{}d00e80be\allowbreak{}87055bf3\\
Registro en 108 & \textemdash{}\\
\end{tabularx}\endgroup\par\vspace{6pt}
\Needspace{87pt}\phantomsection\label{trace-307}\noindent\hyperref[nav:vi-ruta-307]{\textbf{r9c5\_\allowbreak{}h+1v-1}}\quad Ordinal unificado / memoria: 307 / 306\par\nopagebreak
\begingroup\fontsize{8.3}{10.1}\selectfont\setlength{\tabcolsep}{4pt}\noindent\begin{tabularx}{\linewidth}{@{}p{.21\linewidth}X@{}}
Genealogía & b685f141\allowbreak{}3fffbb6c\allowbreak{}c5e0b9a1\allowbreak{}5517eac7\allowbreak{}8ab70bc6\allowbreak{}6c1db71b\allowbreak{}3ab13da3\allowbreak{}1d17dc60\\
Lector & db33053c\allowbreak{}3c7df39b\allowbreak{}b195f28e\allowbreak{}adc06b92\allowbreak{}7e42eedd\allowbreak{}c3e9b8a8\allowbreak{}f84cf04b\allowbreak{}4d72ac5b\\
Registro en 108 & \textemdash{}\\
\end{tabularx}\endgroup\par\vspace{6pt}
\Needspace{87pt}\phantomsection\label{trace-308}\noindent\hyperref[nav:vi-ruta-308]{\textbf{r9c5\_\allowbreak{}h+1v+1}}\quad Ordinal unificado / memoria: 308 / 305\par\nopagebreak
\begingroup\fontsize{8.3}{10.1}\selectfont\setlength{\tabcolsep}{4pt}\noindent\begin{tabularx}{\linewidth}{@{}p{.21\linewidth}X@{}}
Genealogía & 75c7d913\allowbreak{}67c880ed\allowbreak{}973e5c53\allowbreak{}6d4f7b3d\allowbreak{}dbb4d6a6\allowbreak{}c0478a77\allowbreak{}6fb92507\allowbreak{}21a22979\\
Lector & 8a5cc992\allowbreak{}862f3e41\allowbreak{}a477ea86\allowbreak{}9ccbabd1\allowbreak{}9b444a88\allowbreak{}76e6e3b1\allowbreak{}b084264c\allowbreak{}ab2e4db1\\
Registro en 108 & \textemdash{}\\
\end{tabularx}\endgroup\par\vspace{6pt}
\Needspace{87pt}\phantomsection\label{trace-309}\noindent\hyperref[nav:vi-ruta-309]{\textbf{r9c6\_\allowbreak{}h-1v-1}}\quad Ordinal unificado / memoria: 309 / 312\par\nopagebreak
\begingroup\fontsize{8.3}{10.1}\selectfont\setlength{\tabcolsep}{4pt}\noindent\begin{tabularx}{\linewidth}{@{}p{.21\linewidth}X@{}}
Genealogía & ca419a4f\allowbreak{}b89079e2\allowbreak{}64213c9b\allowbreak{}485a320a\allowbreak{}311c24d2\allowbreak{}08068ee4\allowbreak{}1866ec2a\allowbreak{}c1b43c2c\\
Lector & a382f295\allowbreak{}1d5ce06b\allowbreak{}eded8af0\allowbreak{}111b7528\allowbreak{}cf000d77\allowbreak{}19f04317\allowbreak{}52c134cd\allowbreak{}d274a383\\
Registro en 108 & \textemdash{}\\
\end{tabularx}\endgroup\par\vspace{6pt}
\Needspace{87pt}\phantomsection\label{trace-310}\noindent\hyperref[nav:vi-ruta-310]{\textbf{r9c6\_\allowbreak{}h-1v+1}}\quad Ordinal unificado / memoria: 310 / 311\par\nopagebreak
\begingroup\fontsize{8.3}{10.1}\selectfont\setlength{\tabcolsep}{4pt}\noindent\begin{tabularx}{\linewidth}{@{}p{.21\linewidth}X@{}}
Genealogía & 106b15fe\allowbreak{}f86d47d9\allowbreak{}8a2151f4\allowbreak{}c2319629\allowbreak{}333008c7\allowbreak{}64daa083\allowbreak{}780d58ab\allowbreak{}178d9ec0\\
Lector & 97e31ef7\allowbreak{}a8d5ddaf\allowbreak{}45317072\allowbreak{}ce7a9bc9\allowbreak{}0b1be5f3\allowbreak{}2591931c\allowbreak{}12674670\allowbreak{}13eddfa2\\
Registro en 108 & \textemdash{}\\
\end{tabularx}\endgroup\par\vspace{6pt}
\Needspace{87pt}\phantomsection\label{trace-311}\noindent\hyperref[nav:vi-ruta-311]{\textbf{r9c6\_\allowbreak{}h+1v-1}}\quad Ordinal unificado / memoria: 311 / 310\par\nopagebreak
\begingroup\fontsize{8.3}{10.1}\selectfont\setlength{\tabcolsep}{4pt}\noindent\begin{tabularx}{\linewidth}{@{}p{.21\linewidth}X@{}}
Genealogía & 59a1b925\allowbreak{}f37b198a\allowbreak{}9b8d28bf\allowbreak{}47609299\allowbreak{}f18ecace\allowbreak{}ce1e78e3\allowbreak{}13294726\allowbreak{}444a07ed\\
Lector & 82e2ed9e\allowbreak{}32b01eb9\allowbreak{}363063f9\allowbreak{}5aa0e97a\allowbreak{}797114c2\allowbreak{}1c1a7b3a\allowbreak{}e4bdd8f7\allowbreak{}20fbfc4f\\
Registro en 108 & 915f2448\allowbreak{}1eb9d511\allowbreak{}15cf664e\allowbreak{}358374d4\allowbreak{}a4585b51\allowbreak{}4ac886dd\allowbreak{}a764c433\allowbreak{}b1d9a456\\
\end{tabularx}\endgroup\par\vspace{6pt}
\Needspace{87pt}\phantomsection\label{trace-312}\noindent\hyperref[nav:vi-ruta-312]{\textbf{r9c6\_\allowbreak{}h+1v+1}}\quad Ordinal unificado / memoria: 312 / 309\par\nopagebreak
\begingroup\fontsize{8.3}{10.1}\selectfont\setlength{\tabcolsep}{4pt}\noindent\begin{tabularx}{\linewidth}{@{}p{.21\linewidth}X@{}}
Genealogía & 6159b0ea\allowbreak{}3949ae42\allowbreak{}d6ac44d1\allowbreak{}67a4d71d\allowbreak{}d9d0ca4f\allowbreak{}91305ded\allowbreak{}a987869a\allowbreak{}ff7a5301\\
Lector & 2952410f\allowbreak{}bf2ee021\allowbreak{}f5736f19\allowbreak{}d28673e6\allowbreak{}bcf3f086\allowbreak{}7bc4b89a\allowbreak{}48f77673\allowbreak{}b413ea76\\
Registro en 108 & \textemdash{}\\
\end{tabularx}\endgroup\par\vspace{6pt}
\Needspace{87pt}\phantomsection\label{trace-313}\noindent\hyperref[nav:vi-ruta-313]{\textbf{r9c7\_\allowbreak{}h-1v-1}}\quad Ordinal unificado / memoria: 313 / 316\par\nopagebreak
\begingroup\fontsize{8.3}{10.1}\selectfont\setlength{\tabcolsep}{4pt}\noindent\begin{tabularx}{\linewidth}{@{}p{.21\linewidth}X@{}}
Genealogía & 8684eaf7\allowbreak{}f9ac6857\allowbreak{}4b1177dd\allowbreak{}46056f16\allowbreak{}497b29a2\allowbreak{}10b4fc30\allowbreak{}5b69f129\allowbreak{}f3516cf0\\
Lector & e45675cc\allowbreak{}b2ead11c\allowbreak{}246854fd\allowbreak{}23eab547\allowbreak{}9d5768fd\allowbreak{}fb30d0a8\allowbreak{}2cc741ae\allowbreak{}6be60f35\\
Registro en 108 & 23ef2c5c\allowbreak{}d56b7ff3\allowbreak{}5c013077\allowbreak{}bee1cb8d\allowbreak{}bd0ee658\allowbreak{}e4c9ec23\allowbreak{}93eb87b4\allowbreak{}c365f550\\
\end{tabularx}\endgroup\par\vspace{6pt}
\Needspace{87pt}\phantomsection\label{trace-314}\noindent\hyperref[nav:vi-ruta-314]{\textbf{r9c7\_\allowbreak{}h-1v+1}}\quad Ordinal unificado / memoria: 314 / 315\par\nopagebreak
\begingroup\fontsize{8.3}{10.1}\selectfont\setlength{\tabcolsep}{4pt}\noindent\begin{tabularx}{\linewidth}{@{}p{.21\linewidth}X@{}}
Genealogía & 0392a0bd\allowbreak{}824d4db3\allowbreak{}872ee7e5\allowbreak{}5a8e5e59\allowbreak{}e3b21f74\allowbreak{}9916b615\allowbreak{}dad01451\allowbreak{}654991d0\\
Lector & 93e0fc94\allowbreak{}08e34fa4\allowbreak{}8b346d2a\allowbreak{}a179d966\allowbreak{}af4513b4\allowbreak{}05ea24a2\allowbreak{}68f6ebd4\allowbreak{}8d236116\\
Registro en 108 & \textemdash{}\\
\end{tabularx}\endgroup\par\vspace{6pt}
\Needspace{87pt}\phantomsection\label{trace-315}\noindent\hyperref[nav:vi-ruta-315]{\textbf{r9c7\_\allowbreak{}h+1v-1}}\quad Ordinal unificado / memoria: 315 / 314\par\nopagebreak
\begingroup\fontsize{8.3}{10.1}\selectfont\setlength{\tabcolsep}{4pt}\noindent\begin{tabularx}{\linewidth}{@{}p{.21\linewidth}X@{}}
Genealogía & 15c25fad\allowbreak{}37a466a9\allowbreak{}092234de\allowbreak{}4e51eff2\allowbreak{}14bb4541\allowbreak{}30f89e92\allowbreak{}0ac3eabd\allowbreak{}0af86e23\\
Lector & aeafcc7a\allowbreak{}88835b5e\allowbreak{}c7f51128\allowbreak{}1b2c9269\allowbreak{}e769f3ff\allowbreak{}86a4680e\allowbreak{}2e29a986\allowbreak{}6f3cabc4\\
Registro en 108 & \textemdash{}\\
\end{tabularx}\endgroup\par\vspace{6pt}
\Needspace{87pt}\phantomsection\label{trace-316}\noindent\hyperref[nav:vi-ruta-316]{\textbf{r9c7\_\allowbreak{}h+1v+1}}\quad Ordinal unificado / memoria: 316 / 313\par\nopagebreak
\begingroup\fontsize{8.3}{10.1}\selectfont\setlength{\tabcolsep}{4pt}\noindent\begin{tabularx}{\linewidth}{@{}p{.21\linewidth}X@{}}
Genealogía & 8e4b3fc2\allowbreak{}59154927\allowbreak{}80bec7f3\allowbreak{}16245b54\allowbreak{}8a5a96f4\allowbreak{}f38c5037\allowbreak{}b6ad8f94\allowbreak{}60591d64\\
Lector & d78bc138\allowbreak{}90706959\allowbreak{}02c24558\allowbreak{}78b04823\allowbreak{}455e0d80\allowbreak{}836529f9\allowbreak{}e8807940\allowbreak{}4ea54b23\\
Registro en 108 & \textemdash{}\\
\end{tabularx}\endgroup\par\vspace{6pt}
\Needspace{87pt}\phantomsection\label{trace-317}\noindent\hyperref[nav:vi-ruta-317]{\textbf{r9c8\_\allowbreak{}h-1v-1}}\quad Ordinal unificado / memoria: 317 / 320\par\nopagebreak
\begingroup\fontsize{8.3}{10.1}\selectfont\setlength{\tabcolsep}{4pt}\noindent\begin{tabularx}{\linewidth}{@{}p{.21\linewidth}X@{}}
Genealogía & ce9297be\allowbreak{}674940a6\allowbreak{}8f160eee\allowbreak{}97674c66\allowbreak{}9a846625\allowbreak{}821f53c9\allowbreak{}ca30b05b\allowbreak{}d0b56292\\
Lector & 6ef27d74\allowbreak{}57d2c0a1\allowbreak{}51c7c2cc\allowbreak{}86034c0a\allowbreak{}b1ce15e7\allowbreak{}74f84ed5\allowbreak{}c64f617e\allowbreak{}44ec520f\\
Registro en 108 & \textemdash{}\\
\end{tabularx}\endgroup\par\vspace{6pt}
\Needspace{87pt}\phantomsection\label{trace-318}\noindent\hyperref[nav:vi-ruta-318]{\textbf{r9c8\_\allowbreak{}h-1v+1}}\quad Ordinal unificado / memoria: 318 / 319\par\nopagebreak
\begingroup\fontsize{8.3}{10.1}\selectfont\setlength{\tabcolsep}{4pt}\noindent\begin{tabularx}{\linewidth}{@{}p{.21\linewidth}X@{}}
Genealogía & 16995ba1\allowbreak{}149d7ba0\allowbreak{}bc3cfc77\allowbreak{}43494696\allowbreak{}d33ba118\allowbreak{}e19e0ecc\allowbreak{}76f644d3\allowbreak{}1711ba9f\\
Lector & bf91d1ae\allowbreak{}5ba472dd\allowbreak{}e0f327f9\allowbreak{}602e6991\allowbreak{}8e8799f0\allowbreak{}e537555d\allowbreak{}5094f8ab\allowbreak{}0aac97ca\\
Registro en 108 & d2a88d33\allowbreak{}0a7badcf\allowbreak{}0a410fef\allowbreak{}d82f9b47\allowbreak{}4dec91ba\allowbreak{}1dbe5a7b\allowbreak{}22a91f36\allowbreak{}905d51f7\\
\end{tabularx}\endgroup\par\vspace{6pt}
\Needspace{87pt}\phantomsection\label{trace-319}\noindent\hyperref[nav:vi-ruta-319]{\textbf{r9c8\_\allowbreak{}h+1v-1}}\quad Ordinal unificado / memoria: 319 / 318\par\nopagebreak
\begingroup\fontsize{8.3}{10.1}\selectfont\setlength{\tabcolsep}{4pt}\noindent\begin{tabularx}{\linewidth}{@{}p{.21\linewidth}X@{}}
Genealogía & 2231f178\allowbreak{}c6f579c0\allowbreak{}407e70bc\allowbreak{}bff9ed75\allowbreak{}0acd6efb\allowbreak{}3a4dd638\allowbreak{}1f3268e8\allowbreak{}15826311\\
Lector & 810d7ebd\allowbreak{}6b5ee7a7\allowbreak{}2a6dab27\allowbreak{}1c40aef2\allowbreak{}39894808\allowbreak{}02df5c3c\allowbreak{}6064b00f\allowbreak{}893231ca\\
Registro en 108 & \textemdash{}\\
\end{tabularx}\endgroup\par\vspace{6pt}
\Needspace{87pt}\phantomsection\label{trace-320}\noindent\hyperref[nav:vi-ruta-320]{\textbf{r9c8\_\allowbreak{}h+1v+1}}\quad Ordinal unificado / memoria: 320 / 317\par\nopagebreak
\begingroup\fontsize{8.3}{10.1}\selectfont\setlength{\tabcolsep}{4pt}\noindent\begin{tabularx}{\linewidth}{@{}p{.21\linewidth}X@{}}
Genealogía & 7c9b0456\allowbreak{}e1d3dbd9\allowbreak{}f89b9961\allowbreak{}1a50f337\allowbreak{}eefa938f\allowbreak{}02f234c4\allowbreak{}a7ad065f\allowbreak{}d62a6692\\
Lector & ec678d84\allowbreak{}ce11dc61\allowbreak{}e52a2443\allowbreak{}8e4f523b\allowbreak{}4d0ea94b\allowbreak{}94d28d2c\allowbreak{}805f94c5\allowbreak{}60a94223\\
Registro en 108 & \textemdash{}\\
\end{tabularx}\endgroup\par\vspace{6pt}
\Needspace{87pt}\phantomsection\label{trace-321}\noindent\hyperref[nav:vi-ruta-321]{\textbf{r9c9\_\allowbreak{}h-1v-1}}\quad Ordinal unificado / memoria: 321 / 324\par\nopagebreak
\begingroup\fontsize{8.3}{10.1}\selectfont\setlength{\tabcolsep}{4pt}\noindent\begin{tabularx}{\linewidth}{@{}p{.21\linewidth}X@{}}
Genealogía & 3e3bf9f3\allowbreak{}b35f1316\allowbreak{}f0d24282\allowbreak{}2cae4fe2\allowbreak{}11eb3159\allowbreak{}5576f51d\allowbreak{}063c36e2\allowbreak{}41c4cc99\\
Lector & 1f648323\allowbreak{}cbeef43a\allowbreak{}f2c169ae\allowbreak{}e5b6a999\allowbreak{}302aca3e\allowbreak{}106e2780\allowbreak{}e370b148\allowbreak{}0e3ee9e3\\
Registro en 108 & \textemdash{}\\
\end{tabularx}\endgroup\par\vspace{6pt}
\Needspace{87pt}\phantomsection\label{trace-322}\noindent\hyperref[nav:vi-ruta-322]{\textbf{r9c9\_\allowbreak{}h-1v+1}}\quad Ordinal unificado / memoria: 322 / 323\par\nopagebreak
\begingroup\fontsize{8.3}{10.1}\selectfont\setlength{\tabcolsep}{4pt}\noindent\begin{tabularx}{\linewidth}{@{}p{.21\linewidth}X@{}}
Genealogía & 5ca5179e\allowbreak{}96a7a217\allowbreak{}149fc629\allowbreak{}59891128\allowbreak{}4a679d87\allowbreak{}4c335483\allowbreak{}825169c0\allowbreak{}8628742b\\
Lector & a1fd3c28\allowbreak{}7cd8f05a\allowbreak{}77fc26d8\allowbreak{}6500b68a\allowbreak{}6f6d45a5\allowbreak{}15b1bfc3\allowbreak{}2c7201ea\allowbreak{}592f831a\\
Registro en 108 & 373f8cf1\allowbreak{}d27efd3b\allowbreak{}616a88ca\allowbreak{}8eefef42\allowbreak{}03f804fa\allowbreak{}b97be9f4\allowbreak{}d60cd5a3\allowbreak{}f094d407\\
\end{tabularx}\endgroup\par\vspace{6pt}
\Needspace{87pt}\phantomsection\label{trace-323}\noindent\hyperref[nav:vi-ruta-323]{\textbf{r9c9\_\allowbreak{}h+1v-1}}\quad Ordinal unificado / memoria: 323 / 322\par\nopagebreak
\begingroup\fontsize{8.3}{10.1}\selectfont\setlength{\tabcolsep}{4pt}\noindent\begin{tabularx}{\linewidth}{@{}p{.21\linewidth}X@{}}
Genealogía & e847061b\allowbreak{}4af737e2\allowbreak{}2cde7449\allowbreak{}fd8c7399\allowbreak{}b3568176\allowbreak{}df3655b0\allowbreak{}9f776b95\allowbreak{}b8d09a80\\
Lector & 1f6b0ca6\allowbreak{}11906a48\allowbreak{}28e0616b\allowbreak{}a7456a15\allowbreak{}244cef42\allowbreak{}f9db847f\allowbreak{}10419b48\allowbreak{}f959a5cd\\
Registro en 108 & \textemdash{}\\
\end{tabularx}\endgroup\par\vspace{6pt}
\Needspace{87pt}\phantomsection\label{trace-324}\noindent\hyperref[nav:vi-ruta-324]{\textbf{r9c9\_\allowbreak{}h+1v+1}}\quad Ordinal unificado / memoria: 324 / 321\par\nopagebreak
\begingroup\fontsize{8.3}{10.1}\selectfont\setlength{\tabcolsep}{4pt}\noindent\begin{tabularx}{\linewidth}{@{}p{.21\linewidth}X@{}}
Genealogía & 8bd3961b\allowbreak{}89885d9b\allowbreak{}c5b76d4d\allowbreak{}18cbbc1b\allowbreak{}f791e769\allowbreak{}482c1949\allowbreak{}918d5074\allowbreak{}88e02251\\
Lector & 871508bd\allowbreak{}d9589580\allowbreak{}3d299de3\allowbreak{}6cd7cd3e\allowbreak{}9eade5f2\allowbreak{}2157f159\allowbreak{}9f294895\allowbreak{}163b56ce\\
Registro en 108 & \textemdash{}\\
\end{tabularx}\endgroup\par\vspace{6pt}

\end{document}
